% SPDX-FileCopyrightText: © 2026 Idiap Research Institute <contact@idiap.ch>
%
% SPDX-FileContributor: Andre Freitas <andre.freitas@idiap.ch>
% Neuro-symbolic AI Group
%
% SPDX-FileContributor: Danilo Gusicuma <danilo.gusicuma@idiap.ch>
% Neuro-symbolic AI Group
%
% SPDX-License-Identifier: MIT
\documentclass[11pt]{article}
\usepackage[margin=1in]{geometry}
\usepackage{fontspec}
\usepackage{amsmath,amssymb,mathtools}
\usepackage{enumitem}
\usepackage[colorlinks=true,linkcolor=blue,urlcolor=blue]{hyperref}
\setlist[itemize]{leftmargin=*,itemsep=0.25em,topsep=0.25em}

\title{Timeline Report LaTeX Translation}
\author{Abducer Timeline Report LaTeX Workflow}
\date{2026-07-21T14:37:40+00:00}

\begin{document}
\maketitle
\tableofcontents
\clearpage

\section*{Report Translation Contract}
This document is the LaTeX version of the timeline HTML report. Each visible report tab is translated into a section, preserving the same tab order and internal report-card structure where applicable. The companion paper includes these same sections as an appendix.

\begin{itemize}
\item Summary $\rightarrow$ Summary
\item Complexity Strategy $\rightarrow$ Complexity Strategy
\item Timeline From cycle window $\rightarrow$ Content-Specific Full Timeline
\item Ranked Analyses $\rightarrow$ Current Ranked Next-Analysis Queue
\item Qualitative Graph $\rightarrow$ Qualitative Mechanistic Graph
\item Gedanken Experiments $\rightarrow$ Gedanken Experiments
\item PMP Probes $\rightarrow$ Parsimonious Probe Control
\item Contributions $\rightarrow$ Scientific Contribution Monitor
\item Problem \& Literature $\rightarrow$ Target Problem And Literature Evolution
\item Macro Strategies $\rightarrow$ Macro Strategies
\item Strategy Judge $\rightarrow$ Strategy Progression Judge
\item Proof Strategy Story $\rightarrow$ Accumulated Proof Strategy Story
\end{itemize}

\section*{Report Translation Validation}
These sections preserve the visible timeline report tabs as LaTeX sections. validation handoff: run \texttt{python agentic/scripts/validate\_timeline\_latex\_sections.py} before treating the translation as paper-ready.

\section{Summary}

\label{sec:report-summary}

\subsection*{Current mathematical story}

current\_obstruction: The derivation does not cover arbitrary semantic SAT branch pruning or FormulaEval rules whose influencing rows are not determined by proof syntax before use; this is exactly \(RealSATInfluenceRealizer_{N}^{arb}\).; headline: The active discovery route is now concentrated on target-preserving complexity-theoretic bridges: turning weak lower-bound, meta-complexity, proof-complexity, or derandomization interfaces into explicit threshold transfers without importing nonuniform leakage or barrier-evading slogans.; latest\_learning: Cycle 154 converts strict-prefix Qrep from a syntax convention into a replay theorem with a matching temporal falsifier. The valid contribution is a restricted \(TransparentRowSupport_{N}\) rule case and Formula Footprints supporting lemma: charged, ranked, dependency-closed dictionary compression preserves \(\Gamma\) line induction; delayed or hidden compression is exactly hidden support. This is route-level certification only, not a lower bound for \(H_{N}\) and not P vs NP progress.; next\_experiment: Run the prefix-scan verifier on one minimal proof with a valid two-row strict-prefix Qrep merge and on one identical proof with the certificate delayed by one line; compute \(\Gamma_{Qrep}\) in both cases and export the rejected line as the private-row falsifier pivot.; recent\_storyline: branch\_title: Hardness-Magnification Bridge Realization; cycle\_id: cycle-0150; learning: Cycle 150 closes a proof-delivery upgrade rather than rediscovering magnification: Qrep quotient compression is valid exactly as a batched, ranked, proof-carrying same-theory dictionary charged in \(\Gamma_{Qrep}\). The sharp falsifier is the first unranked, cyclic, advice-chosen, or hidden semantic merge. The result is a scoped parser/falsifier theorem and Formula Footprints supporting lemma; it does not prove \(H_{N}\) lower bounds or \(\mathrm{P} \ne \mathrm{NP}\).; next\_experiment: Run the Tseitin-renaming Qrep thought experiment: instantiate two batches with one certified duplicate and one unlisted duplicate, compute \(K_{Qrep}\), C\_Q(D\_Q), and \(\Gamma_{Qrep}\) explicitly, then verify that the unlisted duplicate yields the first-hidden-row private-row falsifier while the certified duplicate preserves the line induction.; branch\_title: Hardness-Magnification Bridge Realization; cycle\_id: cycle-0151; learning: Cycle 151 upgrades the batched Qrep patch with temporal prefix stability: quotient compression is valid only as append-only, ranked, proof-carrying dictionary use. The first pre-use, rebinding, cyclic, advice-chosen, or hidden semantic merge is an explicit falsifier, while full semantic pruning stays quarantined.; next\_experiment: Write the finite prefix-stable Qrep verifier: fields are batch\_id, rank, source\_row, representative\_row, certificate\_root, dependency list, first\_use\_line, and no\_rebind flag; prove the first rejected field yields the private-row falsifier.; branch\_title: Hardness-Magnification Bridge Realization: prefix-stable Qrep support closure; cycle\_id: cycle-0152; learning: Cycle 152 upgrades Qrep from an append-only convention to a prefix-verifiable rule case: quotient compression preserves the transparent parser only when certificate charge strictly precedes representative use. The exact no-go is temporal, not just acyclicity; delayed same-theory evidence is already a private-row falsifier. This closes a proof-audit loophole while preserving the nonclaim boundary: no \(H_{N}\) lower bound and no \(\mathrm{P} \ne \mathrm{NP}\) result.; next\_experiment: Write the finite verifier table for PrefixStableQrep entries: fields entry\_id, append\_time, old\_row, rep\_row, rank, certificate\_root, dependency\_list, first\_use\_time, \(charge_{to,Gamma}\), and rejection reason; then run it on one EqSAT duplicate-child example and one delayed-certificate private-row counterexample.; branch\_title: Hardness-Magnification Bridge Realization: prefix-strict Qrep support audit; cycle\_id: cycle-0153; learning: Cycle 153 upgrades Qrep from prefix-verifiable syntax to a theorem-shaped \(TransparentRowSupport_{N}\) rule case: quotient compression is valid exactly as strict-prefix, ranked, proof-carrying, dependency-closed dictionary use. The new falsifier is temporal; delayed same-theory evidence is already hidden support. This is a scoped parser/falsifier contribution, not a lower bound for \(H_{N}\) and not \(\mathrm{P} \ne \mathrm{NP}\) progress.; next\_experiment: Instantiate the strict-prefix Qrep verifier on a two-duplicate projected-dDNNF FormulaEval example and a delayed-certificate counterexample; compute \(\Gamma_{bEF}\), \(U_{N}\), \(X_{N}\), \(D_{Q}\), \(E_{pi}\), and the q/epsilon support-density consequence line by line.; branch\_title: Hardness-Magnification Bridge Realization; cycle\_id: cycle-0154; learning: Cycle 154 converts strict-prefix Qrep from a syntax convention into a replay theorem with a matching temporal falsifier. The valid contribution is a restricted \(TransparentRowSupport_{N}\) rule case and Formula Footprints supporting lemma: charged, ranked, dependency-closed dictionary compression preserves \(\Gamma\) line induction; delayed or hidden compression is exactly hidden support. This is route-level certification only, not a lower bound for \(H_{N}\) and not P vs NP progress.; next\_experiment: Run the prefix-scan verifier on one minimal proof with a valid two-row strict-prefix Qrep merge and on one identical proof with the certificate delayed by one line; compute \(\Gamma_{Qrep}\) in both cases and export the rejected line as the private-row falsifier pivot.

\subsection*{What actually changed}

\begin{itemize}
\item The marathon created durable branch-local mathematical memory.
\item Each cycle produced theorem-shaped branch outputs with equations, assumptions, falsifiers, and next experiments.
\item Skill dispatch is now explicit and validated rather than implicit in a workflow name.
\item Completed subgoal enablement is now explicit: completed or near-completed subgoals are interpreted by the gate they open for the larger target, not by completion alone.
\item Literature-gap preinvestment is now explicit: broad branch exploration needs a validated avenue warrant with closest-known-result accounting, novelty boundary, field impact, unification power, and route-decision power.
\item Field relevance and general impact are now explicit: local route work is support-only unless it names a target carrier, target gate, unification-power payoff, and demotion-ready impact test.
\item Visual conceptualization is now controlled: metaphors must preserve named theorem properties, expose distortions, and generate an audited route consequence before affecting proof search.
\item Research cycle count is 154; synthetic smoke/legacy packet count is 0.
\end{itemize}

\subsection*{How this adds to the global problem}

This move adds to the P vs NP problem because it identifies whether the transfer loses all-exponent strength through a positive target-query exponent \(\beta\); it keeps barrier status and nonuniform/advice leakage attached to the actual transfer theorem. The current branch-control decision is continue\_run\_next. Payoff criterion: For this route, state the exact barrier premise it uses or escapes: relativizing, natural, algebrizing, nonuniform, or advice-leaking; promotion requires the escape to be attached to the specific transfer theorem.

\subsection*{Payoff criteria}

\begin{itemize}
\item For this route, state the exact barrier premise it uses or escapes: relativizing, natural, algebrizing, nonuniform, or advice-leaking; promotion requires the escape to be attached to the specific transfer theorem.
\item Compute \(\beta\) and q exactly; promotion requires \(\beta=0\), subpolynomial target queries, or a genuinely nonblackbox use of the SAT circuit.
\item Record relativization, natural-proofs, algebrization, nonuniformity, and advice-leakage status for the specific transfer theorem.
\item Control gate: name the source carrier or language, target class, and circuit or resource model before claiming progress.
\item Control gate: derive the exact transfer inequality from \(\mathrm{SAT}\in \mathrm{SIZE}(n^{c})\), NP subset P/poly, or the stated target assumption to the source-carrier upper bound.
\item Control gate: compute \(\beta\) and q exactly, distinguishing all-exponent consequences from bounded target-exponent windows.
\end{itemize}

\subsection*{When to stop}

\begin{itemize}
\item Every named carrier in the branch has \(\beta>0\) and yields only a bounded target-exponent exclusion
\item The branch cannot state a source lower-bound model, target class, and transfer inequality in the same notation
\item The proposed escape is only a black-box substitution restatement with no new carrier or nonblackbox mechanism
\end{itemize}

\subsection*{Do not stop prematurely if}

\begin{itemize}
\item The branch has not yet computed \(\beta\), q, and the induced source upper-bound exponent under \(\mathrm{SAT}\in \mathrm{SIZE}(n^{c})\).
\item A subpolynomial-query or nonblackbox carrier remains available and has not been audited.
\item The latest failure sharpened the threshold-collapse lemma into a formal no-go statement.
\end{itemize}


\section{Complexity Strategy}

\label{sec:report-complexity-strategy}

This top-level section states the active proof strategy before the report moves into the full timeline. For P vs NP, strategy is controlled by source carrier, target class, reduction or transfer map, exact exponent debt, uniformity/advice discipline, and known barrier status.

\subsection*{Current macro reading}
Current macro strategy: compute the transfer exponents or turn the route into a no-go theorem.

\paragraph{Contribution pivot invariant.} If the next work only restates hardness magnification, MCSP promise, or barrier folklore, the branch is looping. If it computes a concrete transfer parameter, proves a route-specific no-go, or certifies a searchable certificate carrier with no advice leakage, it is moving toward a real scientific contribution.

\paragraph{Next decisive test.} Run the prefix-scan verifier on one minimal proof with a valid two-row strict-prefix Qrep merge and on one identical proof with the certificate delayed by one line; compute \(\Gamma_{Qrep}\) in both cases and export the rejected line as the private-row falsifier pivot.

\paragraph{When to stop.} Stop, backtrack, or quarantine when two sharpened attempts still yield only a bounded exponent window, advised hard strings, a natural/relativizing/algebrizing proof object, or a certificate-existence condition equivalent to the target lower bound.

\subsection*{Black-box transfer and exponent-window audit}

\paragraph{Route family.} Complexity lower-bound transfer

\paragraph{Macro thesis.} Treat every proposed magnification route as an explicit transfer inequality from \(\mathrm{SAT}\in \mathrm{SIZE}(n^{c})\) or NP subset P/poly to a source-carrier upper bound.

\paragraph{Conceptual proof picture.} For this complexity-theory target, read the physical picture as a proof-dynamics picture: the target circuit is a resource flowing through an interface. Positive query-length exponent \(\beta\) amplifies the target exponent c and prevents a barely-superlinear source lower bound from ruling out all polynomial-size SAT circuits.

\paragraph{Why this route matters now.} The latest cycle isolated max\{a,q+beta c\} as the active obstruction, so the next work must compute a, q, and \(\beta\) rather than restate magnification.

\paragraph{What to compute next.}
\begin{itemize}
\item For the selected carrier, compute source exponent a, interface overhead q, and target-query exponent \(\beta\).
\item Derive the induced source upper bound under \(\mathrm{SAT}\in \mathrm{SIZE}(n^{c})\).
\item Promote only if \(\beta=0\), beta=o(1), or the use of the SAT circuit is genuinely nonblackbox.
\end{itemize}

\paragraph{Success condition.}
\begin{itemize}
\item A formal transfer theorem states source object, target class, model, and exponent inequality.
\item The theorem distinguishes bounded-exponent exclusion from \(\mathrm{NP}\not\subset \mathrm{P}/\mathrm{poly}\).
\item The branch yields a computable falsifier for black-box substitution routes.
\end{itemize}

\paragraph{Failure or pivot condition.}
\begin{itemize}
\item Every tested carrier has positive \(\beta\) and only excludes c below a finite threshold.
\item The interface size q already exceeds the barely-superlinear threshold.
\item The route cannot name the model and quantifier pattern in one statement.
\end{itemize}

\paragraph{Active skills.}
\begin{itemize}
\item pvsnp-symbolic-complexity-auditor
\item pvsnp-qualitative-complexity-reasoner
\item pvsnp-conceptual-pressure-localizer
\item entailment-gap-auditor
\item progression-payoff-judge
\end{itemize}

\paragraph{Evidence from story.}
\begin{itemize}
\item branch\_title: Hardness-Magnification Bridge Realization; cycle\_id: cycle-0149; learning: Cycle 149 folds representative-stable ranked QCTSR into \(TransparentRowSupport_{N}\) as a valid rule case. The new contribution is not another magnification slogan: quotient compression is proof-carrying, rank-stratified, dependency-closed, and charged, while the first rejected merge gives a concrete private-row falsifier. Full semantic SAT pruning remains exactly at \(RealSATInfluenceRealizer_{N}^{arb}\).; next\_experiment: Write the standalone paper theorem for \(TransparentRowSupport_{N}^{Qrep}\): include the \(K_{Qrep}\) fixed-point verifier pseudocode, the \(\Gamma_{Qrep}\) ledger proof, the first-rejected-merge countermodel, and a one-page nonclaim box separating this parser boundary from \(H_{N}\) lower bounds and \(\mathrm{P} \ne \mathrm{NP}\).
\item branch\_title: Hardness-Magnification Bridge Realization; cycle\_id: cycle-0150; learning: Cycle 150 closes a proof-delivery upgrade rather than rediscovering magnification: Qrep quotient compression is valid exactly as a batched, ranked, proof-carrying same-theory dictionary charged in \(\Gamma_{Qrep}\). The sharp falsifier is the first unranked, cyclic, advice-chosen, or hidden semantic merge. The result is a scoped parser/falsifier theorem and Formula Footprints supporting lemma; it does not prove \(H_{N}\) lower bounds or \(\mathrm{P} \ne \mathrm{NP}\).; next\_experiment: Run the Tseitin-renaming Qrep thought experiment: instantiate two batches with one certified duplicate and one unlisted duplicate, compute \(K_{Qrep}\), C\_Q(D\_Q), and \(\Gamma_{Qrep}\) explicitly, then verify that the unlisted duplicate yields the first-hidden-row private-row falsifier while the certified duplicate preserves the line induction.
\item branch\_title: Hardness-Magnification Bridge Realization; cycle\_id: cycle-0151; learning: Cycle 151 upgrades the batched Qrep patch with temporal prefix stability: quotient compression is valid only as append-only, ranked, proof-carrying dictionary use. The first pre-use, rebinding, cyclic, advice-chosen, or hidden semantic merge is an explicit falsifier, while full semantic pruning stays quarantined.; next\_experiment: Write the finite prefix-stable Qrep verifier: fields are batch\_id, rank, source\_row, representative\_row, certificate\_root, dependency list, first\_use\_line, and no\_rebind flag; prove the first rejected field yields the private-row falsifier.
\item branch\_title: Hardness-Magnification Bridge Realization: prefix-stable Qrep support closure; cycle\_id: cycle-0152; learning: Cycle 152 upgrades Qrep from an append-only convention to a prefix-verifiable rule case: quotient compression preserves the transparent parser only when certificate charge strictly precedes representative use. The exact no-go is temporal, not just acyclicity; delayed same-theory evidence is already a private-row falsifier. This closes a proof-audit loophole while preserving the nonclaim boundary: no \(H_{N}\) lower bound and no \(\mathrm{P} \ne \mathrm{NP}\) result.; next\_experiment: Write the finite verifier table for PrefixStableQrep entries: fields entry\_id, append\_time, old\_row, rep\_row, rank, certificate\_root, dependency\_list, first\_use\_time, \(charge_{to,Gamma}\), and rejection reason; then run it on one EqSAT duplicate-child example and one delayed-certificate private-row counterexample.
\item branch\_title: Hardness-Magnification Bridge Realization: prefix-strict Qrep support audit; cycle\_id: cycle-0153; learning: Cycle 153 upgrades Qrep from prefix-verifiable syntax to a theorem-shaped \(TransparentRowSupport_{N}\) rule case: quotient compression is valid exactly as strict-prefix, ranked, proof-carrying, dependency-closed dictionary use. The new falsifier is temporal; delayed same-theory evidence is already hidden support. This is a scoped parser/falsifier contribution, not a lower bound for \(H_{N}\) and not \(\mathrm{P} \ne \mathrm{NP}\) progress.; next\_experiment: Instantiate the strict-prefix Qrep verifier on a two-duplicate projected-dDNNF FormulaEval example and a delayed-certificate counterexample; compute \(\Gamma_{bEF}\), \(U_{N}\), \(X_{N}\), \(D_{Q}\), \(E_{pi}\), and the q/epsilon support-density consequence line by line.
\item branch\_title: Hardness-Magnification Bridge Realization; cycle\_id: cycle-0154; learning: Cycle 154 converts strict-prefix Qrep from a syntax convention into a replay theorem with a matching temporal falsifier. The valid contribution is a restricted \(TransparentRowSupport_{N}\) rule case and Formula Footprints supporting lemma: charged, ranked, dependency-closed dictionary compression preserves \(\Gamma\) line induction; delayed or hidden compression is exactly hidden support. This is route-level certification only, not a lower bound for \(H_{N}\) and not P vs NP progress.; next\_experiment: Run the prefix-scan verifier on one minimal proof with a valid two-row strict-prefix Qrep merge and on one identical proof with the certificate delayed by one line; compute \(\Gamma_{Qrep}\) in both cases and export the rejected line as the private-row falsifier pivot.
\end{itemize}

\subsection*{Padding and projection demagnification as a no-go route}

\paragraph{Route family.} Projection carrier obstruction

\paragraph{Macro thesis.} Use the padding/projection failures as theorem material: projection carriers either give only bounded exponent windows or move the real proof burden into sparse valid-length lower bounds.

\paragraph{Conceptual proof picture.} The conceptual mechanism is proof-debt conservation. Padding changes the coordinate system, but projection back to SAT reveals that the supposedly weak source lower bound already encodes the hard target lower bound.

\paragraph{Why this route matters now.} Cycles 1-5 repeatedly converted padding optimism into exact exponent inequalities and sparse-length obligations.

\paragraph{What to compute next.}
\begin{itemize}
\item State the projection-carrier dichotomy as a standalone no-go theorem.
\item Define the length map N(n), projection overhead, and target exponent window.
\item List escape conditions: subpolynomial embedding, nonprojection carrier, or nonblackbox use.
\end{itemize}

\paragraph{Success condition.}
\begin{itemize}
\item The no-go theorem has precise hypotheses and a proof skeleton.
\item It explains why the route fails without claiming P vs NP.
\item It identifies which nonprojection alternatives remain open.
\end{itemize}

\paragraph{Failure or pivot condition.}
\begin{itemize}
\item The statement merely repeats padding folklore without a new quantified obstruction.
\item The theorem accidentally rules out nonprojection or nonblackbox routes.
\item The proof depends on a stronger source lower bound than advertised.
\end{itemize}

\paragraph{Active skills.}
\begin{itemize}
\item pvsnp-symbolic-complexity-auditor
\item pvsnp-impossibility-theorem-forger
\item contradiction-route-auditor
\item paper-contribution-readiness-judge
\item scientific-contribution-detector
\end{itemize}

\paragraph{Evidence from story.}
\begin{itemize}
\item branch\_title: Hardness-Magnification Bridge Realization; cycle\_id: cycle-0098; learning: Cycle 98 upgrades the prior packet from theorem slogan to dependency-checked inference structure: \(\Gamma_{bEF}\) is justified by separate EF-dependency, SATeq-link, bounded-SR, and sparse-row charges; \(\mathrm{Psys}_{bEF}\) gives a valid restricted falsifier theorem; full \(\mathrm{Psys}_{sym}\) is blocked exactly at \(FE_{Extract,N}\), not at generic naturalness or relativization folklore.; next\_experiment: Write the formal \(FE_{Extract,N}\) proposition in two columns: left column row-expanded FormulaEval axiom schema with \(N^{1+q+o(1)}\) charge, right column semantic-oracle version with the smallest private-row countermodel; verify which column matches each \(\mathrm{Psys}_{sym}\) rule.
\item branch\_title: Hardness-Magnification Bridge Realization; cycle\_id: cycle-0123; learning: Cycle 123 closes a proof-audit gap rather than repeating the magnification slogan: \(K_{dag}\) is an implementable least fixed point, projected-dDNNF gives a concrete transparent positive case, and the unread-enable-bit/private-row construction gives a precise necessity falsifier. The valid result is a scoped parser/falsifier theorem for beta-zero parity magnification; full semantic FormulaEval is still blocked exactly at \(RealSATInfluenceRealizer_{N}\).; next\_experiment: Write the referee-checkable trace for the projected-dDNNF example: list gate fields, parent pointers, projection maps, RowSet certificates, \(K_{0},K_{1}\),... closure additions, \(U_{N}\) and \(X_{N}\) charges, then flip one unread enable bit to exhibit the first-hidden-row countermodel and verify the \(\Gamma_{dag}\) threshold numerically for qepsilon.
\item branch\_title: Hardness-Magnification Bridge Realization; cycle\_id: cycle-0136; learning: Cycle 136 turns counterfactual declared TSR from a boundary slogan into a bitfield-level closed-support theorem: K\_CTSR(pi) is a finite fixed point, \(\Gamma_\mathrm{CTSR}\) is derived by line induction, and first-hidden-child adaptive pruning is the exact falsifier for extensional safety. The valid contribution remains a scoped parser/falsifier theorem for beta-zero parity magnification certificates; it does not prove \(H_{N}\) lower bounds or \(\mathrm{P} \ne \mathrm{NP}\).; next\_experiment: Select one concrete full semantic pruning rule beyond CTSR declarations, preferably adaptive Tseitin branch pruning with real SAT contradiction tests, and attempt either a charged \(RealSATInfluenceRealizer_{N}^{arb}\) or a minimal first-unchargeable-row countermodel.
\item branch\_title: Hardness-Magnification Bridge Realization; cycle\_id: cycle-0145; learning: Cycle 145 upgrades ranked QCTSR from a line-induction slogan to a finite verifier/minimality packet. Quotient compression is valid only when proof-carrying same-theory certificates are ranked, dependency-closed, and charged before use. The exact obstruction is the first rejected quotient merge, which yields a private-row falsifier. This strengthens the paper parser/countermodel boundary without proving \(H_{N}\) lower bounds or \(\mathrm{P} \ne \mathrm{NP}\).; next\_experiment: Instantiate the finite verifier on one toy QCTSR proof with two EqSAT-equivalent counterfactual TSR children: compute \(K_{Q}\), ranks, \(B_{Q}\), and \(\Gamma_\mathrm{QCTSR}\); then delete one certificate edge and explicitly write the two-world private-row falsifier.
\item branch\_title: Hardness-Magnification Bridge Realization; cycle\_id: cycle-0149; learning: Cycle 149 folds representative-stable ranked QCTSR into \(TransparentRowSupport_{N}\) as a valid rule case. The new contribution is not another magnification slogan: quotient compression is proof-carrying, rank-stratified, dependency-closed, and charged, while the first rejected merge gives a concrete private-row falsifier. Full semantic SAT pruning remains exactly at \(RealSATInfluenceRealizer_{N}^{arb}\).; next\_experiment: Write the standalone paper theorem for \(TransparentRowSupport_{N}^{Qrep}\): include the \(K_{Qrep}\) fixed-point verifier pseudocode, the \(\Gamma_{Qrep}\) ledger proof, the first-rejected-merge countermodel, and a one-page nonclaim box separating this parser boundary from \(H_{N}\) lower bounds and \(\mathrm{P} \ne \mathrm{NP}\).
\item branch\_title: Hardness-Magnification Bridge Realization; cycle\_id: cycle-0150; learning: Cycle 150 closes a proof-delivery upgrade rather than rediscovering magnification: Qrep quotient compression is valid exactly as a batched, ranked, proof-carrying same-theory dictionary charged in \(\Gamma_{Qrep}\). The sharp falsifier is the first unranked, cyclic, advice-chosen, or hidden semantic merge. The result is a scoped parser/falsifier theorem and Formula Footprints supporting lemma; it does not prove \(H_{N}\) lower bounds or \(\mathrm{P} \ne \mathrm{NP}\).; next\_experiment: Run the Tseitin-renaming Qrep thought experiment: instantiate two batches with one certified duplicate and one unlisted duplicate, compute \(K_{Qrep}\), C\_Q(D\_Q), and \(\Gamma_{Qrep}\) explicitly, then verify that the unlisted duplicate yields the first-hidden-row private-row falsifier while the certified duplicate preserves the line induction.
\end{itemize}

\subsection*{Meta-complexity carrier pressure}

\paragraph{Route family.} MCSP/MKtP/GapMCSP route analysis

\paragraph{Macro thesis.} Use MCSP, MKtP, or GapMCSP only after the carrier's verifier overhead, density, naturalness, advice leakage, and target transfer are stated.

\paragraph{Conceptual proof picture.} The meta-complexity carrier is a diagnostic medium: it can reveal whether a weak source lower bound really transfers, or whether constructivity, largeness, and nonuniform advice absorb the intended separation.

\paragraph{Why this route matters now.} Cycles 1, 2, 4, 5, and 6 all point to MCSP/MKtP as the next plausible nonprojection test, but cycle 6 already exposed a fixed-threshold ceiling.

\paragraph{What to compute next.}
\begin{itemize}
\item Instantiate one carrier with N, n, threshold s, verifier overhead, and circuit model.
\item Check natural-proof largeness and constructivity conditions.
\item Compute whether \(\mathrm{SAT}\in \mathrm{SIZE}(n^{c})\) yields carrier circuits below the claimed threshold.
\end{itemize}

\paragraph{Success condition.}
\begin{itemize}
\item The carrier avoids pure projection demagnification.
\item The transfer is target-preserving for \(\mathrm{NP}\not\subset \mathrm{P}/\mathrm{poly}\), not only a proxy consequence.
\item Barrier status is explicit and does not rely on hidden advice leakage.
\end{itemize}

\paragraph{Failure or pivot condition.}
\begin{itemize}
\item The induced upper bound remains N\textasciicircum{}\{q+beta c+o(1)\} with positive \(\beta\).
\item The lower-bound property is natural without an escape certificate.
\item The carrier proves only bounded-exponent or non-target consequences.
\end{itemize}

\paragraph{Active skills.}
\begin{itemize}
\item pvsnp-conceptual-pressure-localizer
\item pvsnp-symbolic-complexity-auditor
\item pvsnp-qualitative-complexity-reasoner
\item counterexample-hunter
\item kb-concept-bridge-planner
\end{itemize}

\paragraph{Evidence from story.}
\begin{itemize}
\item branch\_title: Hardness-Magnification Bridge Realization; cycle\_id: cycle-0149; learning: Cycle 149 folds representative-stable ranked QCTSR into \(TransparentRowSupport_{N}\) as a valid rule case. The new contribution is not another magnification slogan: quotient compression is proof-carrying, rank-stratified, dependency-closed, and charged, while the first rejected merge gives a concrete private-row falsifier. Full semantic SAT pruning remains exactly at \(RealSATInfluenceRealizer_{N}^{arb}\).; next\_experiment: Write the standalone paper theorem for \(TransparentRowSupport_{N}^{Qrep}\): include the \(K_{Qrep}\) fixed-point verifier pseudocode, the \(\Gamma_{Qrep}\) ledger proof, the first-rejected-merge countermodel, and a one-page nonclaim box separating this parser boundary from \(H_{N}\) lower bounds and \(\mathrm{P} \ne \mathrm{NP}\).
\item branch\_title: Hardness-Magnification Bridge Realization; cycle\_id: cycle-0150; learning: Cycle 150 closes a proof-delivery upgrade rather than rediscovering magnification: Qrep quotient compression is valid exactly as a batched, ranked, proof-carrying same-theory dictionary charged in \(\Gamma_{Qrep}\). The sharp falsifier is the first unranked, cyclic, advice-chosen, or hidden semantic merge. The result is a scoped parser/falsifier theorem and Formula Footprints supporting lemma; it does not prove \(H_{N}\) lower bounds or \(\mathrm{P} \ne \mathrm{NP}\).; next\_experiment: Run the Tseitin-renaming Qrep thought experiment: instantiate two batches with one certified duplicate and one unlisted duplicate, compute \(K_{Qrep}\), C\_Q(D\_Q), and \(\Gamma_{Qrep}\) explicitly, then verify that the unlisted duplicate yields the first-hidden-row private-row falsifier while the certified duplicate preserves the line induction.
\item branch\_title: Hardness-Magnification Bridge Realization; cycle\_id: cycle-0151; learning: Cycle 151 upgrades the batched Qrep patch with temporal prefix stability: quotient compression is valid only as append-only, ranked, proof-carrying dictionary use. The first pre-use, rebinding, cyclic, advice-chosen, or hidden semantic merge is an explicit falsifier, while full semantic pruning stays quarantined.; next\_experiment: Write the finite prefix-stable Qrep verifier: fields are batch\_id, rank, source\_row, representative\_row, certificate\_root, dependency list, first\_use\_line, and no\_rebind flag; prove the first rejected field yields the private-row falsifier.
\item branch\_title: Hardness-Magnification Bridge Realization: prefix-stable Qrep support closure; cycle\_id: cycle-0152; learning: Cycle 152 upgrades Qrep from an append-only convention to a prefix-verifiable rule case: quotient compression preserves the transparent parser only when certificate charge strictly precedes representative use. The exact no-go is temporal, not just acyclicity; delayed same-theory evidence is already a private-row falsifier. This closes a proof-audit loophole while preserving the nonclaim boundary: no \(H_{N}\) lower bound and no \(\mathrm{P} \ne \mathrm{NP}\) result.; next\_experiment: Write the finite verifier table for PrefixStableQrep entries: fields entry\_id, append\_time, old\_row, rep\_row, rank, certificate\_root, dependency\_list, first\_use\_time, \(charge_{to,Gamma}\), and rejection reason; then run it on one EqSAT duplicate-child example and one delayed-certificate private-row counterexample.
\item branch\_title: Hardness-Magnification Bridge Realization: prefix-strict Qrep support audit; cycle\_id: cycle-0153; learning: Cycle 153 upgrades Qrep from prefix-verifiable syntax to a theorem-shaped \(TransparentRowSupport_{N}\) rule case: quotient compression is valid exactly as strict-prefix, ranked, proof-carrying, dependency-closed dictionary use. The new falsifier is temporal; delayed same-theory evidence is already hidden support. This is a scoped parser/falsifier contribution, not a lower bound for \(H_{N}\) and not \(\mathrm{P} \ne \mathrm{NP}\) progress.; next\_experiment: Instantiate the strict-prefix Qrep verifier on a two-duplicate projected-dDNNF FormulaEval example and a delayed-certificate counterexample; compute \(\Gamma_{bEF}\), \(U_{N}\), \(X_{N}\), \(D_{Q}\), \(E_{pi}\), and the q/epsilon support-density consequence line by line.
\item branch\_title: Hardness-Magnification Bridge Realization; cycle\_id: cycle-0154; learning: Cycle 154 converts strict-prefix Qrep from a syntax convention into a replay theorem with a matching temporal falsifier. The valid contribution is a restricted \(TransparentRowSupport_{N}\) rule case and Formula Footprints supporting lemma: charged, ranked, dependency-closed dictionary compression preserves \(\Gamma\) line induction; delayed or hidden compression is exactly hidden support. This is route-level certification only, not a lower bound for \(H_{N}\) and not P vs NP progress.; next\_experiment: Run the prefix-scan verifier on one minimal proof with a valid two-row strict-prefix Qrep merge and on one identical proof with the certificate delayed by one line; compute \(\Gamma_{Qrep}\) in both cases and export the rejected line as the private-row falsifier pivot.
\end{itemize}

\subsection*{Nonblackbox or subpolynomial-interface escape search}

\paragraph{Route family.} Escape-condition construction

\paragraph{Macro thesis.} The search should now aim at the exact escape left open by the no-go evidence: a target-preserving carrier with subpolynomial queries or a genuinely nonblackbox use of the SAT circuit.

\paragraph{Conceptual proof picture.} The route needs a new interaction pattern with the target circuit. If the circuit is only queried as a black box at polynomially growing target length, the beta-ceiling reappears; a real escape must change that interface.

\paragraph{Why this route matters now.} Cycle 7 localized the remaining opening as \(N^{o(1)}\) target queries, \(sub-N^{1+\epsilon}\) interface size, and a valid lower-bound carrier.

\paragraph{What to compute next.}
\begin{itemize}
\item Test a succinct local-consistency or proof-complexity carrier.
\item Compute a, q, and \(\beta\) before making any magnification claim.
\item Run the projection-demagnification falsifier and barrier audit immediately.
\end{itemize}

\paragraph{Success condition.}
\begin{itemize}
\item A carrier is target-preserving and not a disguised projection route.
\item The interface has \(\beta=0\) or o(1), or the construction is nonblackbox in a formally stated sense.
\item The next theorem can be stated as either a positive transfer or a no-go lemma.
\end{itemize}

\paragraph{Failure or pivot condition.}
\begin{itemize}
\item The carrier uses polynomial target queries after all.
\item The proof obligation becomes as hard as the original P vs NP/P-poly target.
\item The route cannot avoid naturalness, advice leakage, or algebrization pressure.
\end{itemize}

\paragraph{Active skills.}
\begin{itemize}
\item pvsnp-symbolic-complexity-auditor
\item pvsnp-impossibility-theorem-forger
\item pvsnp-qualitative-complexity-reasoner
\item kb-gap-abstraction-conjecturer
\item abducer-refinement-route-ranker
\end{itemize}

\paragraph{Evidence from story.}
\begin{itemize}
\item branch\_title: Hardness-Magnification Bridge Realization: Polylog Nonblackbox Compact Consumer; cycle\_id: cycle-0013; learning: Cycle 13 turns the superlogarithmic polylog window into an exact bridge: L=(log N)\textasciicircum{}d with \(d>1\) makes SAT circuits of size \(L^{c}\) amortize to \(N^{o(1)}\), while truth-table lookup costs 2\textasciicircum{}L=N\textasciicircum{}\{omega(1)\} and no longer supplies the old falsifier. The valid threshold is \(\rho<1+\epsilon\) for a target-clean nonblackbox compact consumer. The new hard obstruction is not folklore but the absence of an explicit carrier with no advice leakage and a non-natural \(N^{1+\epsilon}\) general-circuit lower bound.; next\_experiment: Instantiate \(PolyLogAudit_{N,d,V}\) with a universal-circuit verifier, compute the exact overhead exponent \(\rho\) and auxiliary advice length, and reject the instance immediately if \(H_{N}\) depends on a hidden SAT circuit or collapses to a black-box SAT query family.
\item branch\_title: Hardness-Magnification Bridge Realization: approximate relaxed bag separators; cycle\_id: cycle-0031; learning: Cycle 31 converts relaxed bag separators into an approximate threshold lemma: E bad bags cost NE gates, so errors below \(N^\epsilon\) are nonuniformly patchable. The valid source obligation is now an approximate distribution-local separator lower bound, not merely promise SAT hardness.; next\_experiment: Try a collision-balanced generator and estimate whether majority/hash-based local predicates yield \(s=N^{o(1)}\) with \(E<N^\epsilon\); if yes, quarantine this carrier family.
\item branch\_title: Relativizing locality-extraction gate for generated-OR carriers; cycle\_id: cycle-0048; learning: Cycle 48 makes the obstruction formal: separator-denominator hardness is only an upper-bound falsifier, not a valid lower-bound certificate for \(H_{N}\). The surviving route is direct C(H\_N) lower bounds or a nonrelativizing locality-extraction theorem with explicit threshold loss.; next\_experiment: Try to formulate a real-SAT locality extractor using witness self-reducibility: from a small circuit for \(H_{N}\), isolate a generated satisfiable query and build a one-sided predicate; audit algebrization and \(N^{o(1)}\) loss.
\item branch\_title: Hardness-Magnification Bridge Realization; cycle\_id: cycle-0081; learning: Cycle 81 separates the scoped occurrence-local theorem from the full \(\mathrm{EF}+\mathrm{SAT}_{U}\) ambition. The valid contribution is now sharper: private row spoofing is a genuine falsifier for occurrence-local functional certificates, while full \(\mathrm{EF}+\mathrm{SAT}_{U}\) requires a nonrelativizing semantic-to-occurrence extraction theorem because real \(\mathrm{SAT}_{U}\) axioms are not invariant under independent label flips.; next\_experiment: Formalize the semantic-coupling-width thought experiment as a lemma: for paired formulas F and F OR z, show exactly which \(\mathrm{SAT}_{U}\) implication lines break private-label invariance and whether their violations can be charged to complete forced rows below the \(N^{\epsilon}\) patch threshold.
\item branch\_title: Hardness-Magnification Bridge Realization; cycle\_id: cycle-0094; learning: Cycle 94 closes a concrete support-transparency subcase: bit-explicit EF extension variables can be parsed without support declarations, giving a valid transcript falsifier theorem. The full formula-semantic extraction fork is now sharper: the obstruction is not arbitrary EF extension variables but succinct FormulaEval-SAT comprehension, which defeats size-only occurrence parsing unless its row-index expansion is charged.; next\_experiment: Formalize \(\mathrm{Psys}_{sym}\) on finite length N: specify whether FormulaEval-SAT comprehension is a proof line, axiom schema, or extension rule; then prove either an expansion-charge parser with cost \(N^{1+q+o(1)}\) or a minimal countermodel showing surface occurrence completion is unsound.
\item branch\_title: Hardness-Magnification Bridge Realization; cycle\_id: cycle-0107; learning: Cycle 107 adds a real validation upgrade: transparent FormulaEval selectors no longer require interval-union structure for a valid theorem, because full syntactic scanning gives a certified \(B_{FE}=F_{FE}\) \(N^{1+q+o(1)}\) charge. This closes the \(q<\epsilon\) falsifier for a broader selector fragment and isolates the only missing improvement as \(AC0PrefixNonemptiness_{N}\) for output-sensitive enumeration. Full semantic FormulaEval remains quarantined by the private-row falsifier, not by generic barrier folklore.; next\_experiment: Formalize \(AC0PrefixNonemptiness_{N}\) as a finite lemma: for \(\mathrm{AC0}\) selectors on q log N row bits, test whether descriptor structure yields \(N^{o(1)}\) prefix emptiness; if not, state the selector-scan theorem as final and add \(\mathrm{AC0}\) output-sensitive enumeration to open problems.
\end{itemize}


\section{Content-Specific Full Timeline}

\label{sec:report-timeline}

Visible timeline starts at cycle-0001, where the run shifted into content-specific mathematical work. Each cycle is rendered as a report card with learning, pressure, derivation, equations, falsifiers, and next analysis.

\subsection*{cycle-0001: Hardness-Magnification Bridge Realization}

\paragraph{Execution.} research; azure\_responses\_api; attempt primary; 900.005 seconds; 3 lemmas, 3 equations, 6 derivation steps, 5 falsifiers.

\paragraph{Learning.} The exact quantitative transfer is C(Pad\_k(SAT)\_N)<=O(N+N\textasciicircum{}\{c/k\}); hence the barely-superlinear lower bound magnifies only up to target exponent k(1+epsilon). The valid bridge is useful as an exponent-calibrated route lemma, while the fixed-k P/poly claim is blocked by a concrete threshold falsifier rather than by vague folklore.

\paragraph{Pressure.} The bridge does not reach NP notin P/poly unless one proves a cofinal family of padded lower bounds or finds a non-artificial carrier whose interface compresses every polynomial target exponent below \(N^{1+\epsilon}\).

\paragraph{Derivation steps.}
\begin{enumerate}
\item Fix \(L=\mathrm{SAT}\), constants \(k\ge 2\) and \(\epsilon>0\), and legal source length \(N=n^{k}\).
\item Assume toward contradiction that SAT has \(O(n^{c})-size\) circuits for some c<k(1+epsilon).
\item On a padded input y of length N, check that the last N-n bits are zero using O(N) gates and feed the first n bits into the SAT circuit.
\item The composed circuit has size O(N+n\textasciicircum{}c)=O(N+N\textasciicircum{}\{c/k\}). Since \(max(1,c/k)<1+\epsilon\), this is below \(N^{1+\epsilon}\) for large N.
\item This contradicts the source lower bound for Pad\_k(SAT), yielding SAT notin SIZE(n\textasciicircum{}c).
\item Auditing the P/poly target shows that fixed k leaves all exponents c>=k(1+epsilon) untouched.
\end{enumerate}

\paragraph{Equation 1: \(E1_{padding,length}\).}
\[
N=n^k,\qquad Pad_k(L)_N={x0^{N-n}:x\in L_n}.
\]
\emph{Valid when:} Valid when the source lower-bound model is explicitly restricted to legal padded lengths \(N=n^{k}\).

\paragraph{Equation 2: \(E2_{composition,upper,bound}\).}
\[
C(Pad_k(L)_N)\le O(N)+C(L_n)\le O(N+N^{c/k})\quad\text{if }C(L_n)\le O(n^c).
\]
\emph{Valid when:} Applies after the padding checker and target circuit use the same nonuniform Boolean basis.

\paragraph{Equation 3: \(E3_{exponent,transfer}\).}
\[
C(Pad_k(SAT)_N)> N^{1+\epsilon}\Longrightarrow SAT\notin SIZE(n^c)\quad\text{for every }c< k(1+\epsilon).
\]
\emph{Valid when:} Requires the target-preserving composition bound and a strict exponent gap absorbing constants asymptotically.

\paragraph{Next analysis.} Instantiate the same exponent audit for a padded MKtP or GapMCSP carrier: name the verifier, compute Cook-Levin or direct-composition length blow-up, and test whether assuming \(\mathrm{SAT}\in \mathrm{SIZE}(n^{c})\) yields source circuits below \(N^{1+\epsilon}\) without model mismatch or nonuniform leakage.

\paragraph{Falsifiers.}
\begin{itemize}
\item falsifier: threshold collapse occurs when the target exponent c is at least k(1+epsilon), because the induced source circuit may exceed \(N^{1+\epsilon}\).
\item falsifier: model mismatch occurs if the source lower bound is for \(\mathrm{AC0}\), formulas, or uniform circuits while the composed circuit is unrestricted nonuniform P/poly.
\item falsifier: nonuniform leakage occurs if the padding parser or length selector uses advice not counted in the source circuit size.
\item falsifier: an infinitely-often source lower bound is not valid for excluding eventual SIZE(n\textasciicircum{}c) target circuits without matching quantifiers.
\item falsifier: a natural-proofs-style source argument for unrestricted circuits must explain why it is non-natural or why pseudorandom-function assumptions are avoided.
\end{itemize}

\subsection*{cycle-0002: Hardness-Magnification Bridge Realization}

\paragraph{Execution.} research; azure\_responses\_api; attempt primary; 900.003 seconds; 3 lemmas, 5 equations, 7 derivation steps, 5 falsifiers.

\paragraph{Learning.} Cycle 2 turns the previous threshold collapse into an exact quantifier repair. The valid route lemma is: fixed k gives only a finite SAT exponent lower bound, while a common-epsilon lower bound for unbounded padded carriers \(H_{k}\) would imply \(\mathrm{NP}\not\subset \mathrm{P}/\mathrm{poly}\). The decisive falsifier is not folklore but the inequality \(c/k\ge 1+\epsilon\), plus the O(N) checker floor.

\paragraph{Pressure.} No known source theorem proves \(C_{H_{k}}(N)>N^{1+\epsilon}\) for unbounded k with a common \(\epsilon\) against general nonuniform circuits.

\paragraph{Derivation steps.}
\begin{enumerate}
\item Assume for contradiction that SAT has circuits of size O(n\textasciicircum{}c) for some fixed c.
\item Choose k from the source hypothesis with \(c/k<1+\epsilon\).
\item Construct a circuit for \(H_{k}\) on length \(N=n^{k}\) by checking the padding pattern in O(N) gates and running the SAT circuit on the first \(n=N^{1/k}\) bits.
\item Obtain C\_\{H\_k\}(N)=O(N\textasciicircum{}\{max\{1,c/k\}\}).
\item Since \(max{1,c/k}<1+\epsilon\), the upper bound is eventually below \(N^{1+\epsilon}\).
\item Contradict the asserted source lower bound for \(H_{k}\) at that k.
\item Because c was arbitrary only under the forall c exists k source quantifier, derive \(\mathrm{SAT}\notin \mathrm{P}/\mathrm{poly}\) conditionally.
\end{enumerate}

\paragraph{Equation 1: length\_relation.}
\[
N=n^k,\quad R_k(x)=x1 0^{n^k-n-1}
\]
\emph{Valid when:} Valid when k is fixed and padded lengths \(N=n^{k}\) are the only target lengths considered.

\paragraph{Equation 2: upper\_transfer.}
\[
C_{H_k}(N)\le aN+b\,C_{SAT}(N^{1/k})
\]
\emph{Valid when:} Applies when the total padded language includes the O(N) off-image validity checker and constants are model-fixed.

\paragraph{Equation 3: exponent\_transfer.}
\[
C_{SAT}(n)=O(n^c)\Rightarrow C_{H_k}(N)=O(N^{\max{1,c/k}})
\]
\emph{Valid when:} Requires the padding reduction to preserve nonuniform length parameters and circuit-size accounting without extra advice leakage.

\paragraph{Equation 4: strict\_margin.}
\[
\Delta=1+\epsilon-\max{1,c/k}> 0\Rightarrow C_{H_k}(N)=O(N^{1+\epsilon-\Delta})< N^{1+\epsilon}\text{ eventually}
\]
\emph{Valid when:} Valid when the exponent margin is strict, so asymptotic constants are absorbed for large padded lengths.

\paragraph{Equation 5: uniform\_k\_consequence.}
\[
(\exists\epsilon> 0)(\forall c)(\exists k> c/(1+\epsilon))\ C_{H_k}(N)> N^{1+\epsilon}\Rightarrow NP\nsubseteq P/poly
\]
\emph{Valid when:} Valid only when the source lower bound is available for unbounded padding parameters, not one fixed k.

\paragraph{Next analysis.} Test MCSP or proof-complexity carriers for producing the unbounded-k source family: formalize whether a weak lower bound for MCSP[2\textasciicircum{}\{N\textasciicircum{}\{1/k\}\}] transfers to \(C_{H_{k}}(N)>N^{1+\epsilon}\) without becoming natural or losing the strict exponent margin.

\paragraph{Falsifiers.}
\begin{itemize}
\item falsifier: fixed k only yields a finite exponent ceiling c<k(1+epsilon), so arbitrary polynomial advice sizes survive.
\item falsifier: threshold \(C_{H_{k}}(N)>N+N^{0.9}\) is absorbed by the O(N) padding checker and transfers no SAT exponent lower bound.
\item falsifier: a lower bound for a compressed n-bit version of SAT merely restates the target and is not a valid source model.
\item falsifier: if lower bounds hold for only finitely many k, every c above (1+epsilon) max k remains unaddressed.
\item falsifier: a large constructive property proving the source lower bound against general circuits would meet the natural-proofs obstruction under standard PRF assumptions.
\end{itemize}

\subsection*{cycle-0003: Hardness-Magnification Bridge Realization}

\paragraph{Execution.} research; azure\_responses\_api; attempt compact\_retry; 900.006 seconds; 1 lemmas, 3 equations, 4 derivation steps, 4 falsifiers.

\paragraph{Learning.} Cycle 3 finds a valid single-carrier repair: growth padding \(N_{m}=(m+1)^{g(m+1)}\) converts an i.o. \(N^{1+\epsilon}\) lower bound for \(H_{g}\) into exclusion of every SAT exponent c. This is route-level progress, not P vs NP; the hard obligation is now the sparse valid-length lower bound.

\paragraph{Pressure.} The required source lower bound is a barely-superlinear lower bound for a sparse explicit NP carrier on valid lengths; proving it already encodes superpolynomial SAT lower bounds in the unpadded parameter.

\paragraph{Derivation steps.}
\begin{enumerate}
\item Assume toward contradiction that SAT has P/poly circuits of size \(m^{c}\) for some fixed c.
\item Build the \(H_{g}\) circuit at valid length \(N_{m}\) by checking the delimiter/padding pattern in O(N\_m) gates and feeding the prefix x to the \(\mathrm{SAT}_{m}\) circuit.
\item Since \(g(m)\to infty\), eventually \(c/g(m+1)<1+\epsilon/2\), so the transferred circuit has size at most \(N_{m}^{1+\epsilon}\).
\item This contradicts an infinitely-often valid-length lower bound \(C(H_{g},N_{m})>N_{m}^{1+\epsilon}\); therefore no fixed c exists for SAT.
\end{enumerate}

\paragraph{Equation 1: valid\_lengths.}
\[
N_m=(m+1)^{g(m+1)},\quad H_g\cap{0,1}^{N_m}={x\#0^{N_m-m-1}:x\in SAT_m}
\]
\emph{Valid when:} Valid when the delimiter format fixes m and g is computable within polynomial time in \(N_{m}\).

\paragraph{Equation 2: upper\_transfer.}
\[
C_{H_g}(N_m)\le a\,(N_m+C_{SAT}(m))\le a\,(N_m+(m+1)^c)=a\,(N_m+N_m^{c/g(m+1)})
\]
\emph{Valid when:} Applies after assuming SAT has \(m^{c}-size\) circuits and the parser uses only linear-size circuitry.

\paragraph{Equation 3: threshold\_absorption.}
\[
g(m+1)> c/(1+\epsilon/2)\ \Longrightarrow\ C_{H_g}(N_m)\le N_m^{1+\epsilon}\text{ for all large }m
\]
\emph{Valid when:} Requires sufficiently large valid lengths so constants and the O(N\_m) checker floor are absorbed.

\paragraph{Next analysis.} Instantiate g(r)=ceil(log log(r+4)); prove the verifier and length test are poly(N), then analyze whether valid-slice lower bounds for \(H_{g}\) admit a proof-complexity or MCSP carrier formulation.

\paragraph{Falsifiers.}
\begin{itemize}
\item falsifier: If g is bounded by k, the consequence collapses to ruling out only SAT circuits below exponent about k(1+epsilon).
\item falsifier: If the source claim is all-length rather than valid-length/i.o., invalid lengths have trivial rejecting circuits and kill the lower bound.
\item falsifier: If a bundled carrier includes small k slices, those slices dominate the threshold and reintroduce the finite-exponent obstruction.
\item falsifier: If length validity or prefix recovery needs super-polynomial time in N, \(H_{g}\) may leave NP and the target-preserving transfer is invalid.
\end{itemize}

\subsection*{cycle-0004: Hardness-Magnification Bridge Realization}

\paragraph{Execution.} research; azure\_responses\_api; attempt compact\_retry; 900.012 seconds; 1 lemmas, 3 equations, 4 derivation steps, 4 falsifiers.

\paragraph{Learning.} Cycle 4 proves a valid projection-sandwich lemma and demagnifies the sparse-padding route: \(N^{1+eps}\) lower bounds for \(H_{g}\) equal \(m^{(1+eps)g}\) SAT lower bounds. The bridge is correct but not a weak-source path toward P vs NP.

\paragraph{Pressure.} The bridge is valid, but its source obligation is not barely nontrivial in the original SAT length; it already asks for \(m^{\omega(1)}\) SAT lower bounds.

\paragraph{Derivation steps.}
\begin{enumerate}
\item Fix a valid length \(N_{m}\) and define \(H_{g}\) by copying the first m input bits into the \(\mathrm{SAT}_{m}\) instance; invalid lengths are not part of the source lower-bound claim.
\item If \(\mathrm{SAT}_{m}\) has a circuit of size \(m^{c}\), wire the first m inputs into it and ignore all padding bits, giving an \(H_{g}\) circuit of size m\textasciicircum{}c+O(m).
\item If \(H_{g}\) at \(N_{m}\) has a circuit of size s, restrict the padding bits to constants; the resulting m-input circuit decides \(\mathrm{SAT}_{m}\) with size s+O(1).
\item Substitute \(N_{m}=(m+1)^{g(m+1)}\). A barely superlinear bound in N is a superpolynomial SAT lower bound when g is unbounded.
\end{enumerate}

\paragraph{Equation 1: valid\_length\_carrier.}
\[
N_m=(m+1)^{g(m+1)},\quad H_g(x)=SAT_m(x_1,\ldots,x_m)\text{ for }|x|=N_m
\]
\emph{Valid when:} Valid when the source lower-bound model is explicitly restricted to these padded valid lengths.

\paragraph{Equation 2: projection\_sandwich.}
\[
C_m(SAT)\le C_{N_m}(H_g)+O(1),\quad C_{N_m}(H_g)\le C_m(SAT)+O(m)
\]
\emph{Valid when:} Valid when the padding map is a projection and circuits may nonuniformly know the length.

\paragraph{Equation 3: threshold\_demagnification.}
\[
C_{N_m}(H_g)> N_m^{1+\epsilon}\Rightarrow C_m(SAT)> (m+1)^{(1+\epsilon)g(m+1)}-O(m)
\]
\emph{Valid when:} Applies after the reduction preserves the valid length and the lower-bound threshold exactly.

\paragraph{Next analysis.} Test an MCSP or KT-complexity carrier with source model, threshold \(N^{1+eps}\), and target \(\mathrm{NP}\not\subset \mathrm{P}/\mathrm{poly}\); compute the reduction overhead and state whether the needed proof must be non-natural or non-algebrizing.

\paragraph{Falsifiers.}
\begin{itemize}
\item falsifier: if the lower bound is asserted on all lengths, invalid lengths are constant slices and the statement becomes false.
\item falsifier: if the encoding is not a projection and needs an Omega(N) decoder, the O(N) floor collapses the \(N^{1+eps}\) threshold.
\item falsifier: if g is bounded on the lower-bound subsequence, the consequence excludes only finitely many SAT exponents.
\item falsifier: a large constructive proof property would meet the natural-proofs obstruction; this route needs a specific non-natural lower-bound argument.
\end{itemize}

\subsection*{cycle-0005: Hardness-Magnification Bridge Realization}

\paragraph{Execution.} research; azure\_responses\_api; attempt primary; 900.006 seconds; 2 lemmas, 5 equations, 7 derivation steps, 4 falsifiers.

\paragraph{Learning.} Cycle 5 converts the padding quarantine into a general projection-carrier dichotomy. The valid bridge says \(N^{1+\epsilon}\) lower bounds transfer through ell(N) only when \(ell(N)^{c}\) is below the threshold. Positive polynomial embedding length gives only a bounded exponent window; subpolynomial embedding length gives all exponents but turns the source lower bound into \(ell(N)^{\omega(1)}\). Thus pure projection magnification fails as a weak-source path, with an explicit falsifier rather than a slogan.

\paragraph{Pressure.} Projection carriers cannot simultaneously keep the lower bound weak in the semantic source size and transfer it to exclusion of all polynomial SAT exponents.

\paragraph{Derivation steps.}
\begin{enumerate}
\item Generalize the padded SAT carrier to a projection family \(H_{N}\) with embedded SAT length ell(N) and linear validity/decoder overhead.
\item Assume the target negation: SAT has nonuniform circuits of size \(m^{c}\) for a fixed exponent c at every length m.
\item Compose the \(\mathrm{SAT}_{ell(N)}\) circuit with \(pi_{N}\) and \(V_{N}\), obtaining SIZE(H\_N)<=O(N log N)+ell(N)\textasciicircum{}c.
\item Combine with the source lower bound \(\mathrm{SIZE}(H_{N})>K\) \(N^{1+\epsilon}\) on infinitely many N.
\item If \(ell(N)=N^{\alpha+o(1)}\), contradiction is forced only for \(c<(1+\epsilon)/\alpha\).
\item To rule out all fixed c, require \(ell(N)=N^{o(1)}\) on the lower-bound support.
\item But then the source threshold \(N^{1+\epsilon}\) equals \(ell(N)^{\omega(1)}\), so the lower bound is no longer barely nontrivial relative to the embedded SAT instance size.
\end{enumerate}

\paragraph{Equation 1: projection\_upper\_bound.}
\[
\mathrm{SIZE}(H_N) \le A N+B\,\mathrm{SIZE}(\mathrm{SAT}_{\ell(N)})+O(N\log N).
\]
\emph{Valid when:} Valid when the source carrier is a fixed-output projection of SAT with linear decoding overhead.

\paragraph{Equation 2: exponent\_transfer.}
\[
\mathrm{SIZE}(\mathrm{SAT}_m)\le m^c\ \forall m \quad\Rightarrow\quad \mathrm{SIZE}(H_N)\le O(N\log N)+\ell(N)^c.
\]
\emph{Valid when:} Applies after the target SAT circuit family is available at every embedded length.

\paragraph{Equation 3: finite\_window.}
\[
\ell(N)=N^{\alpha+o(1)}\ \Rightarrow\ \mathrm{SIZE}(H_N)\le N^{\alpha c+o(1)}+O(N\log N).
\]
\emph{Valid when:} Valid when the embedding length has a stable polynomial exponent along the lower-bound support.

\paragraph{Equation 4: contradiction\_condition.}
\[
K N^{1+\varepsilon}< \mathrm{SIZE}(H_N)\ \mathrm{i.o.}\ \wedge\ \alpha c< 1+\varepsilon\ \Rightarrow\ \mathrm{SAT}\notin \mathrm{SIZE}(n^c).
\]
\emph{Valid when:} Requires the lower-bound lengths to be infinite and the overhead to be below the threshold.

\paragraph{Equation 5: all\_exponents\_cost.}
\[
(\forall c\;\ell(N)^c=o(N^{1+\varepsilon}))\ \Longleftrightarrow\ \ell(N)=N^{o(1)},\quad \ell(N)=N^{o(1)}\Rightarrow N^{1+\varepsilon}=\ell(N)^{\omega(1)}.
\]
\emph{Valid when:} Valid when ell(N) is unbounded and asymptotics are taken on the same infinite support.

\paragraph{Next analysis.} Audit an MCSP-based route: set source object \(\mathrm{MCSP}_{n,s}\) on \(N=2^{n}-bit\) truth tables, choose a concrete circuit model for the lower bound, state the exact threshold \(N^{1+\epsilon}\) or smaller, and test whether the known magnification transfer is non-natural, target-preserving, and strong enough to imply \(\mathrm{SAT}\notin \mathrm{P}/\mathrm{poly}\) rather than only EXP or subexponential derandomization consequences.

\paragraph{Falsifiers.}
\begin{itemize}
\item falsifier: \(ell(N)=N^{1/2}\), \(\epsilon=0.1\), and target exponent \(c=3\) leave the composed \(H_{N}\) upper bound above \(N^{1.1}\), so the bridge cannot refute \(\mathrm{SAT}\in \mathrm{SIZE}(n^{3})\).
\item falsifier: interface overhead \(N^{1+\delta}\) with \(\delta\ge \epsilon\) masks any target-circuit saving and collapses the quantitative transfer.
\item falsifier: a decoder using N-bit nonuniform advice can make \(H_{N}\) hard for reasons unrelated to SAT, causing nonuniform leakage and invalid target preservation.
\item falsifier: proving \(H_{N}>N^{1+\epsilon}\) only at lengths where the projection image is empty or bounded gives no target-class consequence.
\end{itemize}

\subsection*{cycle-0006: Hardness-Magnification Bridge Realization, Cycle 6: MCSP verifier transfer audit}

\paragraph{Execution.} research; azure\_responses\_api; attempt primary; 900.016 seconds; 3 lemmas, 4 equations, 7 derivation steps, 4 falsifiers.

\paragraph{Learning.} Cycle 6 shows that moving from projection padding to a nonprojection MCSP carrier does not by itself create major magnification. The valid transfer is \(C(\mathrm{MCSP}_{N^\alpha},N)=0\) caps any fixed barely-superlinear bridge below all-exponent P/poly separation. The failure is a concrete threshold falsifier, not barrier folklore.

\paragraph{Pressure.} For \(\mathrm{NP}\not\subset \mathrm{P}/\mathrm{poly}\) one must contradict every constant c, but fixed \(\epsilon\) and \(\alpha\) certify only c<(1+epsilon)/(1+alpha). Even \(\alpha=0\) leaves the input-reading ceiling \(c<1+\epsilon\).

\paragraph{Derivation steps.}
\begin{enumerate}
\item Fix \(\alpha\ge 0\) and \(K_{\alpha,N}=\mathrm{MCSP}_{N^\alpha}\) on valid truth-table lengths N.
\item Assume for contradiction that SAT has nonuniform circuits of size O(m\textasciicircum{}c).
\item Given an N-bit truth table T, construct the Cook-Levin formula \(Phi_{T}\) expressing that there exists a circuit of size at most \(N^\alpha\) matching T on all N inputs.
\item The construction has length M(N)=O(N\textasciicircum{}\{1+alpha\} polylog N) because it enumerates every table position and evaluates the guessed circuit locally.
\item A circuit for \(K_{\alpha,N}\) hardwires the SAT circuit for length M(N), computes \(Phi_{T}\), and outputs the SAT circuit's answer.
\item Therefore \(C(K_{\alpha,N})\le N^{c(1+\alpha)+o(1)}\).
\item A lower bound \(C(K_{\alpha,N})>N^{1+\epsilon}\) contradicts the SAT hypothesis exactly when \(c(1+\alpha)<1+\epsilon\).
\end{enumerate}

\paragraph{Equation 1: verifier\_time.}
\[
T_\alpha(N)=O(N\,s(N)\,\mathrm{polylog}\,N)=N^{1+\alpha+o(1)}
\]
\emph{Valid when:} Valid when MCSP membership is verified by exhaustive truth-table checking of the guessed small circuit.

\paragraph{Equation 2: induced\_upper\_bound.}
\[
\mathrm{SAT}\in\mathrm{SIZE}(m^c)\Rightarrow C(K_{\alpha,N})\le N^{c(1+\alpha)+o(1)}
\]
\emph{Valid when:} Applies after the Cook-Levin reduction preserves uniformity and the SAT circuit is substituted nonuniformly.

\paragraph{Equation 3: transfer\_threshold.}
\[
C(K_{\alpha,N})> N^{1+\epsilon}\Rightarrow \mathrm{SAT}\notin\mathrm{SIZE}(n^c)\quad\text{for }c< \frac{1+\epsilon}{1+\alpha}
\]
\emph{Valid when:} Requires the source lower bound and the induced SAT-based upper bound to hold on aligned valid lengths.

\paragraph{Equation 4: major\_target\_ceiling.}
\[
\sup c\ \text{certified by fixed }(\epsilon,\alpha)\le \frac{1+\epsilon}{1+\alpha}\le 1+\epsilon
\]
\emph{Valid when:} Valid when \(\alpha\) is nonnegative, as forced by input-reading verification of N-bit instances.

\paragraph{Next analysis.} Test a self-referential meta-complexity carrier, such as Gap-MCSP or MKtP with succinctly generated truth tables, where the target SAT circuit is used nonblack-box. Compute whether the induced carrier upper bound can avoid the multiplicative c factor; if not, quarantine it by the same beta-ceiling lemma.

\paragraph{Falsifiers.}
\begin{itemize}
\item falsifier: choose SAT circuit exponent c>=(1+epsilon)/(1+alpha); then the induced MCSP upper bound is not below \(N^{1+\epsilon}\).
\item falsifier: any hidden reduction overhead \(M(N)=N^{\beta+o(1)}\) with \(\beta>1+\alpha\) further shrinks the certified target range to \(c<(1+\epsilon)/\beta\).
\item falsifier: nonuniform leakage occurs if the MCSP-to-SAT construction uses length-specific advice not already present in the hypothesized SAT circuit family.
\item falsifier: invoking natural proofs alone is invalid here because the small-threshold MCSP accepted set is not large; the route fails quantitatively.
\end{itemize}

\subsection*{cycle-0007: Hardness-Magnification Bridge Realization: black-box substitution cap}

\paragraph{Execution.} research; azure\_responses\_api; attempt primary; 900.005 seconds; 2 lemmas, 3 equations, 6 derivation steps, 4 falsifiers.

\paragraph{Learning.} Cycle 7 generalizes the MCSP and padding failures to an oracle-interface theorem. A barely-superlinear carrier lower bound transfers through a black-box SAT interface as max\{a,q+beta c\}; hence positive target-query length exponent \(\beta\) is the decisive threshold-collapse falsifier. The remaining valid route must either use subpolynomial target queries with a genuinely nonprojection carrier or abandon black-box substitution.

\paragraph{Pressure.} The only unquarantined escape is to construct a nonprojection carrier with \(N^{o(1)}\) target query length, \(sub-N^{1+\epsilon}\) interface size, and a valid \(N^{1+\epsilon}\) nonuniform lower bound.

\paragraph{Derivation steps.}
\begin{enumerate}
\item Assume SAT has nonuniform circuits of size O(m\textasciicircum{}c) for every queried formula length m.
\item Model the carrier algorithm as an oracle circuit \(A_{N}\) with explicit ordinary gates and explicit SAT-query gates; adaptivity is absorbed into the circuit DAG.
\item Replace each SAT-query gate of length at most \(N^{\beta+o(1)}\) by the corresponding SAT circuit of size N\textasciicircum{}\{beta c+o(1)\}.
\item The resulting \(L_{N}\) circuit has size at most N\textasciicircum{}\{a+o(1)\}+N\textasciicircum{}\{q+beta c+o(1)\}.
\item Contrapose against \(\mathrm{SIZE}(L_{N})>N^{1+\epsilon}\): only exponents c satisfying q+beta \(c<1+\epsilon\) and \(a<1+\epsilon\) are ruled out.
\item Therefore a positive polynomial query-length exponent \(\beta\) is a quantitative threshold-collapse falsifier for deriving \(\mathrm{NP}\not\subset \mathrm{P}/\mathrm{poly}\) from one fixed barely-superlinear source lower bound.
\end{enumerate}

\paragraph{Equation 1: substitution\_upper\_bound.}
\[
SIZE(L_N) \le N^{a+o(1)}+N^{q+o(1)}\cdot (N^{\beta+o(1)})^c = N^{\max{a,q+\beta c}+o(1)}
\]
\emph{Valid when:} Valid when the source oracle circuit explicitly counts advice, preprocessing, and every SAT oracle gate before substitution.

\paragraph{Equation 2: contrapositive\_window.}
\[
SIZE(L_N)> N^{1+\epsilon} \Rightarrow SAT \notin SIZE(m^c) \text{ for } c< \frac{1+\epsilon-q}{\beta}, \ a< 1+\epsilon, \ \beta> 0
\]
\emph{Valid when:} Applies after the reduction preserves target lengths and the lower bound holds on matching valid source lengths.

\paragraph{Equation 3: all\_exponent\_obstruction.}
\[
\forall c\;[SAT\notin SIZE(m^c)] \text{ from fixed } \epsilon \text{ requires } \beta=0 \text{ and } \max{a,q}< 1+\epsilon
\]
\emph{Valid when:} Requires fixed \(\epsilon\) and black-box substitution as the only target-preserving complexity-class transfer mechanism.

\paragraph{Next analysis.} Test a succinct local-consistency or proof-complexity carrier with \(N^{o(1)}\) SAT queries: compute a,q,beta, then run the projection-demagnification falsifier before claiming any magnification.

\paragraph{Falsifiers.}
\begin{itemize}
\item falsifier: \(\beta>0\) leaves only the finite window \(c<(1+\epsilon-q)/\beta\), so all-exponent NP/poly separation cannot follow.
\item falsifier: \(q\ge 1+\epsilon\) means the oracle interface alone already exceeds the source lower-bound threshold.
\item falsifier: \(a\ge 1+\epsilon\) means preprocessing or advice leakage consumes the entire barely-superlinear margin.
\item falsifier: projection or padding carriers with subpolynomial query length demagnify the source lower bound into a strong target lower bound, as in the earlier projection-sandwich quarantine.
\end{itemize}

\subsection*{cycle-0008: Hardness-Magnification Bridge Realization, cycle 8: subpolynomial-query threshold repair}

\paragraph{Execution.} research; azure\_responses\_api; attempt primary; 111.459 seconds; 3 lemmas, 4 equations, 7 derivation steps, 6 falsifiers.

\paragraph{Learning.} Cycle 8 repairs the remaining route exactly: subpolynomial SAT-query length is sufficient to remove the c-dependent threshold term, so \(N^{1+\epsilon}\) lower bounds magnify only when \(max(a,q)<1+\epsilon\). The valid contribution is a quantified bridge plus a model-preservation gate; the unresolved hard task is finding a concrete carrier satisfying the interface and source lower bound without natural-proof or nonuniform-leakage failure.

\paragraph{Pressure.} No concrete non-natural carrier K is known with both a near-linear subpolynomial-query SAT interface and an \(N^{1+\epsilon}\) nonuniform circuit lower bound. MCSP-like truth-table carriers fail because their SAT-verification queries have polynomial, usually linear, dependence on N.

\paragraph{Derivation steps.}
\begin{enumerate}
\item Fix an arbitrary constant c and assume for contradiction that SAT has nonuniform circuits of size \(n^{c}\) for all sufficiently large n.
\item Pad all SAT queries made by the \(K_{N}\) oracle circuit to one length s(N)+O(1); finite small exceptions are hardwired and do not affect the exponent.
\item Replace each SAT oracle gate by the corresponding SAT circuit. Adaptivity causes no exponential blowup because the oracle computation is already represented as a circuit DAG with oracle gates.
\item The resulting \(K_{N}\) circuit has size at most A(N)+Q(N)(s(N)\textasciicircum{}c+O(s(N)))+O(1).
\item With \(A=N^{a+o(1)}\), \(Q=N^{q+o(1)}\), and \(s=N^{o(1)}\), this size is \(N^{max(a,q)+o(1)}\) for the fixed c.
\item Choose \(\eta>0\) with \(max(a,q)+\eta<1+\epsilon\). For all large N the substituted \(K_{N}\) circuits have size \(<N^{1+\epsilon}\), contradicting the source lower bound.
\item Since c was arbitrary, SAT is not in SIZE(n\textasciicircum{}c) for any fixed c; therefore \(\mathrm{SAT}\notin \mathrm{P}/\mathrm{poly}\) and \(\mathrm{NP}\not\subset \mathrm{P}/\mathrm{poly}\).
\end{enumerate}

\paragraph{Equation 1: oracle-substitution-size.}
\[
C_K(N) \le A(N)+Q(N)\cdot(C_{SAT}(s(N))+O(s(N)))+O(1).
\]
\emph{Valid when:} Valid when the source oracle circuit explicitly counts all query-building wires outside the SAT oracle gates.

\paragraph{Equation 2: polynomial-query-transfer.}
\[
C_{SAT}(n)\le n^c,\ A(N)\le N^{a+o(1)},\ Q(N)\le N^{q+o(1)},\ s(N)=N^{\beta(N)} \Rightarrow C_K(N)\le N^{\max{a,q+\beta(N)\max(c,1)}+o(1)}.
\]
\emph{Valid when:} Applies after the reduction preserves nonuniformity and pads all SAT queries to one controlled length.

\paragraph{Equation 3: subpoly-query-collapse.}
\[
s(N)=N^{o(1)} \Rightarrow C_K(N)\le N^{\max{a,q}+o(1)} \quad \text{for each fixed } c.
\]
\emph{Valid when:} Valid when c is fixed before taking N large, so c times the query-length exponent is still o(1).

\paragraph{Equation 4: magnified-consequence.}
\[
C_K(N)> N^{1+\epsilon}\ \wedge\ \max{a,q}< 1+\epsilon\ \Rightarrow\ \forall c\ C_{SAT}(n)\not\le n^c.
\]
\emph{Valid when:} Requires the target-preserving complexity-class transfer and an eventual source lower bound on the same circuit basis.

\paragraph{Next analysis.} Search for non-truth-table carriers in \(\mathrm{P}^\mathrm{SAT}\) with near-linear ordinary work and subpolynomial SAT queries, e.g. succinct consistency or decompression-check problems, then compute A,Q,s and test against L8.1 before discussing P vs NP.

\paragraph{Falsifiers.}
\begin{itemize}
\item falsifier: \(s(N)=N^\beta\) for \(\beta>0\) causes threshold collapse and excludes only \(c<(1+\epsilon-q)/\beta\).
\item falsifier: \(q\ge 1+\epsilon\) or ordinary overhead \(a\ge 1+\epsilon\) leaves no room below the source lower-bound threshold.
\item falsifier: the source theorem is uniform-only or restricted-model, while \(\mathrm{SAT}\notin \mathrm{P}/\mathrm{poly}\) requires a nonuniform target-preserving contradiction.
\item falsifier: a hidden advice string used to choose queries or decompressions is not counted in A(N), creating nonuniform leakage.
\item falsifier: the proposed carrier is a large constructive useful truth-table property; under standard PRF assumptions a natural proof of its P/poly usefulness is blocked.
\item falsifier: if both the bridge and the source lower-bound proof relativize, oracle worlds with \(\mathrm{NP}^{O}\) subset \(\mathrm{P}^{O}/\mathrm{poly}\) prevent promotion to a main P vs NP proof.
\end{itemize}

\subsection*{cycle-0009: Hardness-Magnification Bridge Realization: subpolynomial-interface demagnification}

\paragraph{Execution.} research; azure\_responses\_api; attempt primary; 166.602 seconds; 2 lemmas, 4 equations, 6 derivation steps, 5 falsifiers.

\paragraph{Learning.} Cycle 9 converts the cycle-8 repair into a route failure certificate. Subpolynomial SAT queries remove the c-dependent upper-bound term, but the same transfer read backward says that proving the source \(N^{1+\epsilon}\) lower bound in general SIZE already proves \(\mathrm{SAT}_{\le L}\) has \(L^{\omega(1)}\) circuits lower bounds. The remaining valid work is not carrier search inside this black-box same-model interface; it is either a restricted-model closure theorem or a nonblackbox carrier.

\paragraph{Pressure.} The branch has no concrete carrier that is simultaneously a weak same-model source lower bound and target-preserving. Same-model sources demagnify; restricted sources need an unproved substitution-closure theorem.

\paragraph{Derivation steps.}
\begin{enumerate}
\item Start from the repaired cycle-8 interface: reduction overhead exponent a, query-count exponent q, and SAT query length \(L(N)=N^{o(1)}\) with \(max(a,q)<1+\epsilon\).
\item Inline a hypothetical \(size-L^{c}\) SAT circuit into every SAT oracle gate of the carrier reduction, obtaining C(K\_N)<=s\_R(N)+Q(N)L(N)\textasciicircum{}c+O(Q(N)L(N)).
\item Observe that this proves the valid target transfer: an eventual \(N^{1+\epsilon}\) lower bound for \(K_{N}\) excludes every fixed c because L(N)\textasciicircum{}c=N\textasciicircum{}\{o(1)\}.
\item Read the same inequality contrapositively: the source lower bound implies C(SAT\_\{(N\textasciicircum{}\{1+epsilon\}-s\_R-O(QL))/Q.
\item For \(Q\le N^{q}\) with \(q<1+\epsilon\) and \(L=N^{o(1)}\), the induced SAT lower bound is N\textasciicircum{}\{1+epsilon-q-o(1)\}=L\textasciicircum{}\{omega(1)\}, already superpolynomial at the target query length.
\item Add the unconditional lookup upper bound \(C(K_{N})=\) (1+epsilon-q-o(1))log\_2 N for any possible \(N^{1+\epsilon}\) source lower bound.
\end{enumerate}

\paragraph{Equation 1: upper\_transfer.}
\[
C_{\mathrm{SIZE}}(K_N) \le s_R(N)+Q(N)\,C_{\mathrm{SIZE}}(\mathrm{SAT}_{\le L(N)})+O(Q(N)L(N))
\]
\emph{Valid when:} Applies when every oracle answer is replaced by a target SAT circuit and the reduction overhead is counted explicitly.

\paragraph{Equation 2: contrapositive\_demagnification.}
\[
C_{\mathrm{SIZE}}(K_N)> N^{1+\epsilon} \Rightarrow C_{\mathrm{SIZE}}(\mathrm{SAT}_{\le L(N)})> \frac{N^{1+\epsilon}-s_R(N)-O(Q(N)L(N))}{Q(N)}
\]
\emph{Valid when:} Valid when the source lower bound exceeds the complete substituted reduction size threshold in the same circuit model.

\paragraph{Equation 3: subpoly\_query\_consequence.}
\[
L(N)=N^{o(1)},\ Q(N)\le N^q,\ s_R(N)\le N^a,\ a,q< 1+\epsilon \Rightarrow C_{\mathrm{SIZE}}(\mathrm{SAT}_{\le L(N)})\ge N^{1+\epsilon-q-o(1)}=L(N)^{\omega(1)}
\]
\emph{Valid when:} Requires the subpolynomial query-length regime and polynomially bounded query count below the claimed source threshold.

\paragraph{Equation 4: lookup\_falsifier\_bound.}
\[
C_{\mathrm{SIZE}}(K_N)\le s_R(N)+O(Q(N)L(N)2^{L(N)})
\]
\emph{Valid when:} Valid when the source model is unrestricted nonuniform circuits and lookup tables may encode SAT truth tables.

\paragraph{Next analysis.} Pivot to a nonblackbox object: test a proof-complexity or derandomization carrier where SAT in P/poly yields the required source upper bound through a named nonrelativizing theorem rather than through black-box SAT-query substitution.

\paragraph{Falsifiers.}
\begin{itemize}
\item same-model demagnification falsifier: an \(N^{1+\epsilon}\) lower bound for \(K_{N}\) entails \(\mathrm{SAT}_{\le L(N)}\) lower bounds of size \(L(N)^{\omega(1)}\), so the source is not barely nontrivial relative to the target.
\item lookup-table falsifier: if Q(N)L(N)2\textasciicircum{}\{L(N)\}=o(N\textasciicircum{}\{1+epsilon\}) after reduction overhead, unrestricted nonuniform circuits decide \(K_{N}\) below the claimed threshold.
\item model-mismatch falsifier: replacing the source model by \(\mathrm{AC0}\), formulas, or another restricted class invalidates the transfer unless a quantitative substitution-closure lemma \(Phi_{M}\) is proved.
\item nonuniform leakage falsifier: if \(R_{N}\) or the compressor is chosen with advice encoding a hard truth table, the carrier lower bound is not target-preserving.
\item relativizing-interface falsifier: a proof using only black-box oracle substitution remains within a relativizing interface and cannot by itself overcome the known oracle barrier for the main P vs NP conjecture.
\end{itemize}

\subsection*{cycle-0010: Hardness-Magnification Bridge Realization}

\paragraph{Execution.} research; azure\_responses\_api; attempt primary; 142.192 seconds; 2 lemmas, 4 equations, 8 derivation steps, 4 falsifiers.

\paragraph{Learning.} Cycle 10 proves a valid restricted-model closure bridge and names the precise target-preservation failure. Subpolynomial queries plus \(max(a,q)<1+\epsilon\) magnify a \(\Gamma\) lower bound to SAT notin \(\Gamma/\mathrm{poly}\) for every fixed exponent. This is real route-level progress, but it is not NP notin P/poly: the missing object is a nonuniform P/poly-to-Gamma lift with subthreshold tax, and its absence is the decisive falsifier for advertising the bridge as a P vs NP route.

\paragraph{Pressure.} The target-preserving lift from arbitrary P/poly SAT circuits to \(\Gamma-circuits\) is absent; without it, restricted-source magnification leaks the target model.

\paragraph{Derivation steps.}
\begin{enumerate}
\item Fix \(\Gamma\) and an oracle \(\Gamma-circuit\) \(A_{N}\) for \(H_{N}\) with \(N^{a}\) internal gates and at most \(N^{q}\) SAT queries of length at most L(N).
\item Assume the restricted target negation: \(\mathrm{SAT}_{\le n}\) has \(\Gamma-circuits\) of size \(n^{c}\) for all large n.
\item Replace each SAT oracle gate in \(A_{N}\) by the corresponding \(\Gamma-circuit\) for \(\mathrm{SAT}_{\le L(N)}\) and use \(\Gamma\) substitution closure.
\item Obtain \(S_{Gamma}(H_{N})\le \)\(\widetilde{O}(N^{a}+N^{q} L(N)\)\(^{c}\)).
\item For \(L=N^\beta\), compare the exponent max(a,q+beta c) with 1+epsilon to exclude only c below (1+epsilon-q)/beta.
\item For \(L=N^{o(1)}\), obtain \(S_{Gamma}(H_{N})\le N^{max(a,q)+o(1)}\) for each fixed c, contradicting the source bound if \(max(a,q)<1+\epsilon\).
\item Attempt to replace the restricted target assumption SAT in \(\Gamma/\mathrm{poly}\) by SAT in P/poly. This requires converting arbitrary SAT circuits into \(\Gamma-circuits\) on the queried lengths.
\item No such conversion follows from the interface. Therefore the valid conclusion is \(\Gamma/\mathrm{poly}\) separation, while the P/poly consequence remains unproved.
\end{enumerate}

\paragraph{Equation 1: substitution\_upper\_bound.}
\[
S_{\Gamma}(H_N) \le \widetilde{O}\bigl(N^a+N^q\,S_{\Gamma}(SAT_{\le L(N)})\bigr)
\]
\emph{Valid when:} Valid when \(\Gamma\) is explicitly closed under oracle-gate substitution and query-generation circuits are counted.

\paragraph{Equation 2: polynomial\_query\_length\_transfer.}
\[
L(N)=N^{\beta},\ S_{\Gamma}(SAT_n)\le n^c\ \Longrightarrow\ S_{\Gamma}(H_N)\le N^{\max{a,q+\beta c}+o(1)}
\]
\emph{Valid when:} Applies when query lengths have fixed polynomial exponent \(\beta\) and c is fixed.

\paragraph{Equation 3: finite\_exponent\_exclusion.}
\[
S_{\Gamma}(H_N)> N^{1+\epsilon}\ \Longrightarrow\ SAT\notin \Gamma[n^c]\quad \text{for } c< (1+\epsilon-q)/\beta \text{ and } a< 1+\epsilon
\]
\emph{Valid when:} Requires positive \(\beta\) and compares the transferred exponent with the source threshold.

\paragraph{Equation 4: subpolynomial\_query\_magnification.}
\[
L(N)=N^{o(1)},\ \max{a,q}< 1+\epsilon\ \Longrightarrow\ SAT\notin \Gamma[n^c]\text{ for every fixed }c
\]
\emph{Valid when:} Valid when the o(1) term is uniform for each fixed target exponent c.

\paragraph{Next analysis.} Test a non-natural nonblackbox carrier whose query length is \(N^{o(1)}\) and whose use of a SAT circuit is semantic rather than oracle substitution; compute whether its lift tax \(\delta\) is below the source slack \(\epsilon\).

\paragraph{Falsifiers.}
\begin{itemize}
\item falsifier: if \(\beta>0\), the transfer excludes only \(c<(1+\epsilon-q)/\beta\), so all-exponent magnification collapses.
\item falsifier: if \(max(a,q)\ge 1+\epsilon\), even subpolynomial SAT queries do not beat the source threshold.
\item falsifier: if \(\Gamma\) is weaker than P/poly and no P/poly-to-Gamma lift is valid, the target class silently changes to \(\Gamma/\mathrm{poly}\).
\item falsifier: a natural-proof style dense useful property computable in \(\Gamma/\mathrm{poly}\) would not be a credible P/poly route under standard pseudorandom-function assumptions; it must be non-natural or outside that barrier's scope.
\end{itemize}

\subsection*{cycle-0011: Hardness-Magnification Bridge Realization: lookup-lift collapse for restricted target repair}

\paragraph{Execution.} research; azure\_responses\_api; attempt primary; 128.594 seconds; 3 lemmas, 4 equations, 5 derivation steps, 5 falsifiers.

\paragraph{Learning.} Cycle 11 turns the cycle-10 target split into a sharper impossibility triangle. To lift \(\Gamma\) lower bounds to P/poly by unrolling arbitrary SAT circuits, queries must be so short that every SAT query can also be answered by unconditional lookup inside \(\Gamma\). That gives \(H_{N}\) a subthreshold \(\Gamma\) upper bound and falsifies the source lower bound. Intermediate log-query lengths either collapse the source or preserve only exponent \(c\le 1\). Thus black-box formula-style P/poly repair is quarantined; progress now requires a genuinely nonblackbox carrier or a proof object where target circuits are used without independent short-query lookup.

\paragraph{Pressure.} The missing bridge is not merely an efficient P/poly-to-Gamma simulation; it must be efficient for SAT-derived query computations while not enabling small lookup for all short query predicates. Generic unrolling cannot satisfy both requirements.

\paragraph{Derivation steps.}
\begin{enumerate}
\item Start from the cycle-10 restricted bridge: a \(\Gamma\) lower bound for \(H_{N}\) magnifies to SAT notin \(\Gamma/\mathrm{poly}\) when \(H_{N}\) has subpolynomial SAT queries and \(max(a,q)<1+\epsilon\).
\item Try to repair the missing P/poly target by setting \(\Gamma\) to De Morgan formulas and converting any hypothesized length-L SAT circuit of size \(L^{c}\) into a formula of size 2\textasciicircum{}\{O(L\textasciicircum{}c)\}.
\item To preserve every fixed P/poly exponent c, require L(N)\textasciicircum{}c=o(log N) for every fixed c, hence L(N)=(log N)\textasciicircum{}\{o(1)\}.
\item In that same regime, \(\mathrm{SAT}_{\le L(N)}\) has an unconditional DNF lookup of size 2\textasciicircum{}\{O(L(N))\}poly(L(N))=N\textasciicircum{}\{o(1)\}.
\item Substituting this lookup into the \(\Gamma-oracle\) skeleton gives \(S_{Gamma}(H_{N})\le N^{max(a,q)+o(1)}\), contradicting the desired \(N^{1+\epsilon}\) source lower bound whenever \(max(a,q)<1+\epsilon\).
\end{enumerate}

\paragraph{Equation 1: oracle-substitution-bound.}
\[
S_{\Gamma}(H_N) \le O(N^a)+N^q T_{\Gamma}(SAT_{\le L(N)})
\]
\emph{Valid when:} Valid when the source carrier is represented by a \(\Gamma\) oracle skeleton and each oracle occurrence is replaced independently.

\paragraph{Equation 2: lookup-upper-bound.}
\[
T_{\Gamma}(SAT_{\le L}) \le 2^{O(L)}\operatorname{poly}(L)
\]
\emph{Valid when:} Valid when \(\Gamma\) contains DNF or decision-tree lookup implementations for every Boolean predicate on L bits.

\paragraph{Equation 3: ultrashort-collapse.}
\[
L(N)=(\log N)^{o(1)} \Rightarrow 2^{O(L(N))}\operatorname{poly}(L(N))=N^{o(1)}
\]
\emph{Valid when:} Valid in the ultra-short query regime needed to preserve every fixed target exponent under exponential unrolling.

\paragraph{Equation 4: source-lower-bound-contradiction.}
\[
S_{\Gamma}(H_N) \le N^{\max{a,q}+o(1)} < N^{1+\epsilon}
\]
\emph{Valid when:} Applies after the threshold assumption \(max{a,q}<1+\epsilon\) is combined with the lookup substitution bound.

\paragraph{Next analysis.} Test a nonblackbox proof-complexity carrier: encode a SAT circuit as a verifier object whose correctness is checked globally rather than by short oracle leaves, and compute whether the verifier overhead stays below \(N^{1+\epsilon}\) without allowing \(\mathrm{SAT}_{\le L}\) truth-table substitution.

\paragraph{Falsifiers.}
\begin{itemize}
\item falsifier: ultra-short query lookup gives an unconditional \(\Gamma\) upper bound below the proposed barely-superlinear source lower bound.
\item falsifier: generic circuit unrolling preserves all fixed P/poly exponents only in the same ultra-short regime where lookup collapse is valid.
\item falsifier: log-length queries with \(\alpha\ge 1\) avoid \(N^{o(1)}\) lookup only by preserving at most \(c\le 1\), which is not a meaningful P/poly target for SAT.
\item falsifier: any independent oracle-leaf substitution argument remains relativizing and cannot certify the needed nonblackbox target transfer.
\item falsifier: \(\Gamma\) classes lacking DNF lookup also lack the advertised generic P/poly lift unless an explicit nonnatural replacement is proved.
\end{itemize}

\subsection*{cycle-0012: Hardness-Magnification Bridge Realization}

\paragraph{Execution.} research; azure\_responses\_api; attempt primary; 140.028 seconds; 2 lemmas, 4 equations, 6 derivation steps, 4 falsifiers.

\paragraph{Learning.} Cycle 12 identifies a real surviving window: superlogarithmic subpolynomial target lengths, especially L=(log N)\textasciicircum{}d with \(d>1\). In this window SAT circuits of size \(L^{c}\) amortize to \(N^{o(1)}\), but SAT truth tables cost \(2^{L}\) and no longer give the old subthreshold lookup falsifier. The bridge is valid only for nonblackbox compact-circuit consumers; it is not a P vs NP proof without the explicit carrier and lower bound.

\paragraph{Pressure.} No explicit carrier \(H_{N}\) and evaluator \(E_{N}\) are yet exhibited whose computation genuinely exploits compact SAT circuits rather than ordinary oracle answers, and no \(N^{1+\epsilon}\) lower bound for such an \(H_{N}\) is known.

\paragraph{Derivation steps.}
\begin{enumerate}
\item Fix \(d>1\) and L(N)=(ceil log N)\textasciicircum{}d. Assume for contradiction \(\mathrm{SAT}\in \mathrm{SIZE}(n^{c})\) for some fixed c, eventually.
\item For each large N choose a correct \(\mathrm{SAT}_{L(N)}\) circuit \(D_{L(N)}\) of size at most \(L(N)^{c}\). Nonuniformity allows this circuit to be hardwired into the \(H_{N}\) upper-bound circuit.
\item Substitute \(D_{L(N)}\) into \(E_{N}\). By the interface bound, \(C(H_{N})\le K\) \(N^{a}\) \(L(N)^{bc+r}\).
\item Since \(L(N)=N^{o(1)}\), for any \(\delta>0\) and all large N, L(N)\textasciicircum{}\{bc+r\}<=N\textasciicircum{}delta. Choose delta<(1+epsilon-a).
\item Then \(\mathrm{SAT}\in \mathrm{SIZE}(n^{c})\) implies \(C(H_{N})N^{1+\epsilon}\).
\item Check against the old black-box falsifier: the unconditional lookup implementation has size N\textasciicircum{}a2\textasciicircum{}\{L(N)\}, which is superpolynomial for L=(log N)\textasciicircum{}d, \(d>1\), so it does not itself refute the source lower bound.
\end{enumerate}

\paragraph{Equation 1: conditional\_upper.}
\[
C(H_N) \leq K N^a |D_{L(N)}|^b L(N)^r
\]
\emph{Valid when:} Valid when the evaluator has fixed nonuniform circuits and the only target-dependent advice is an exact SAT circuit \(D_{L(N)}\) substituted into it.

\paragraph{Equation 2: target\_amortization.}
\[
C(SAT_L) \leq L^c \;\Rightarrow\; C(H_N) \leq K N^a L(N)^{bc+r}=N^{a+o(1)}
\]
\emph{Valid when:} Applies when L(N) is subpolynomial in N and c,b,r are fixed constants independent of the carrier length.

\paragraph{Equation 3: magnification\_contrapositive.}
\[
(\exists^\infty N)\; C(H_N)> N^{1+\epsilon} \;\Rightarrow\; SAT\notin P/poly
\]
\emph{Valid when:} Requires the conditional upper bound for all large N under SAT in P/poly and the strict threshold gap \(a<1+\epsilon\).

\paragraph{Equation 4: lookup\_boundary.}
\[
N^a2^{L(N)}< N^{1+\epsilon} \text{ if } L(N)< (1+\epsilon-a-o(1))\log_2 N
\]
\emph{Valid when:} Valid when the interface can be replaced by a full SAT truth-table lookup rather than a compact nonblackbox target circuit.

\paragraph{Next analysis.} Instantiate \(H_{N}\) as a proof-checking or self-reduction transcript language whose verifier can call a compact \(\mathrm{SAT}_{L}\) circuit nonblackbox; compute its exact overhead a and test whether full truth-table substitution or input-encoded D leakage gives a subthreshold upper bound.

\paragraph{Falsifiers.}
\begin{itemize}
\item falsifier: logarithmic L(N) permits full SAT truth-table substitution of size below \(N^{1+\epsilon}\), collapsing the carrier lower bound unconditionally.
\item falsifier: if the target circuit D is part of the input, then input-reading or nonuniform leakage reintroduces the positive target-size exponent and destroys all-exponent transfer.
\item falsifier: if \(E_{N}\) is merely a black-box SAT-query algorithm, superlogarithmic L blocks lookup but gives no reason that compact SAT circuits can be used without duplicating queries above threshold.
\item falsifier: if the source lower bound is only for formulas, \(\mathrm{AC0}\), or another \(\Gamma\) not containing the substituted evaluator, the conclusion is only SAT notin \(\Gamma/\mathrm{poly}\), not NP notin P/poly.
\end{itemize}

\subsection*{cycle-0013: Hardness-Magnification Bridge Realization: Polylog Nonblackbox Compact Consumer}

\paragraph{Execution.} research; azure\_responses\_api; attempt primary; 155.298 seconds; 3 lemmas, 4 equations, 6 derivation steps, 6 falsifiers.

\paragraph{Learning.} Cycle 13 turns the superlogarithmic polylog window into an exact bridge: L=(log N)\textasciicircum{}d with \(d>1\) makes SAT circuits of size \(L^{c}\) amortize to \(N^{o(1)}\), while truth-table lookup costs 2\textasciicircum{}L=N\textasciicircum{}\{omega(1)\} and no longer supplies the old falsifier. The valid threshold is \(\rho<1+\epsilon\) for a target-clean nonblackbox compact consumer. The new hard obstruction is not folklore but the absence of an explicit carrier with no advice leakage and a non-natural \(N^{1+\epsilon}\) general-circuit lower bound.

\paragraph{Pressure.} No explicit carrier \(H_{N}\) is yet supplied with both a target-clean polylog nonblackbox compact consumer and an i.o. \(N^{1+\epsilon}\) general-circuit lower bound. The lower-bound proof also must be non-natural or otherwise outside the natural-proofs obstruction.

\paragraph{Derivation steps.}
\begin{enumerate}
\item Fix a target exponent c and assume, for contradiction, that \(\mathrm{SAT}_{m}\) has circuits \(D_{m}\) of size \(m^{c}\) for all sufficiently large m.
\item For each sufficiently large valid carrier length N, set m=L(N)=floor((log N)\textasciicircum{}d), so the assumed SAT family supplies \(D_{m}\).
\item Run the compact nonblackbox consumer \(T_{N}\) with \(D_{m}\) wired into its universal evaluator; no SAT truth table and no independent advice are allowed.
\item Because m\textasciicircum{}c=(log N)\textasciicircum{}\{dc\}=N\textasciicircum{}\{o(1)\}, the produced \(H_{N}\) circuit has size \(N^\rho\) poly(m\textasciicircum{}c,m,log N)=N\textasciicircum{}\{rho+o(1)\}.
\item Choose \(\delta>0\) with \(\rho+deltaN^{1+\epsilon}\) on infinitely many valid lengths.
\item Since c was arbitrary and \(\rho\) does not depend on c, the conditional bridge excludes every fixed polynomial exponent for SAT circuits.
\end{enumerate}

\paragraph{Equation 1: polylog target length.}
\[
L(N)=\lfloor (\log N)^d\rfloor,\quad d> 1
\]
\emph{Valid when:} Valid when d is fixed and N ranges over the carrier's valid input lengths.

\paragraph{Equation 2: fixed exponent amortization.}
\[
L(N)^c=(\log N)^{dc}=N^{o(1)}
\]
\emph{Valid when:} Valid when c and d are constants, not growing with N.

\paragraph{Equation 3: compact substitution bound.}
\[
\mathrm{SIZE}(H_N)\le N^{\rho}\operatorname{poly}(L(N)^c,L(N),\log N)=N^{\rho+o(1)}
\]
\emph{Valid when:} Applies after the reduction preserves the source circuit model and target circuit descriptions.

\paragraph{Equation 4: magnification threshold.}
\[
\rho< 1+\varepsilon\ \wedge\ \mathrm{SIZE}(H_N)> N^{1+\varepsilon}\text{ i.o.}\ \Rightarrow\ \mathrm{SAT}\notin\mathrm{SIZE}(n^c)
\]
\emph{Valid when:} Requires the source lower bound and compact upper bound to hold on overlapping valid lengths.

\paragraph{Next analysis.} Instantiate \(PolyLogAudit_{N,d,V}\) with a universal-circuit verifier, compute the exact overhead exponent \(\rho\) and auxiliary advice length, and reject the instance immediately if \(H_{N}\) depends on a hidden SAT circuit or collapses to a black-box SAT query family.

\paragraph{Falsifiers.}
\begin{itemize}
\item falsifier: compact consumer exponent \(\rho\) reaches or exceeds 1+epsilon, causing quantitative threshold collapse.
\item falsifier: L(N)=O(log N), because hardwired SAT truth tables give an unconditional lookup simulation.
\item falsifier: the carrier reduces to a black-box projection SAT(G(x)) in the same model, reviving the demagnification from earlier cycles.
\item falsifier: auxiliary advice \(A_{N}\) is not uniformly computable from N and the hypothesized SAT circuit \(D_{L(N)}\).
\item falsifier: the produced upper circuit is not in the same SIZE model as the claimed source lower bound.
\item falsifier: a dense constructive natural property is the only known path to the general-circuit lower bound, triggering the natural-proofs barrier under standard PRG assumptions.
\end{itemize}

\subsection*{cycle-0014: Hardness-Magnification Bridge Realization: polylog generated-OR carrier}

\paragraph{Execution.} research; azure\_responses\_api; attempt primary; 142.763 seconds; 3 lemmas, 4 equations, 6 derivation steps, 6 falsifiers.

\paragraph{Learning.} Cycle 14 identifies a concrete carrier template, \(H_{N}=OR\) of \(N^{q}\) generated SAT instances of length (log N)\textasciicircum{}d, and computes the exact transfer: \(\mathrm{SAT}\in \mathrm{SIZE}(n^{c})\) gives \(H_{N}\) in SIZE(N\textasciicircum{}\{max\{rho,q\}+o(1)\}). The valid magnification window is max\{rho,q\}1, exhaustive lookup is \(N^{q+\omega(1)}\), so the previous short-query falsifier is avoided, but the source lower bound remains the hard non-natural obligation.

\paragraph{Pressure.} The bridge is valid, but no explicit target-clean \(G_{N},A_{N}\) is known with a provable unrestricted circuit lower bound \(\mathrm{SIZE}(H_{N})>N^{1+\epsilon}\). A proof of that source bound must also avoid natural-proof density, relativizing oracle, and advice-leakage obstructions.

\paragraph{Derivation steps.}
\begin{enumerate}
\item Fix constants \(d>1\), q, \(\rho\), \(\epsilon\) with \(max{\rho,q}<1+\epsilon\), and define \(H_{N}\) by a target-clean generated OR or more general consumer over \(M=N^{q}\) length-L SAT instances.
\item Assume toward contradiction that SAT is in P/poly, so for some c every sufficiently large length l has SAT circuits \(C_{l}\) of size O(l\textasciicircum{}c).
\item For each valid N, copy \(C_{L(N)}\) into every generated target slot and wire these copies to \(A_{N}\); the generator/consumer contributes \(N^\rho\) gates and the target copies contribute \(N^{q}\) \(L^{c}\) gates.
\item Since L=(log N)\textasciicircum{}d, the copy cost is \(N^{q+o(1)}\) for fixed c,d. Hence \(\mathrm{SIZE}(H_{N})\le N^{max{\rho,q}+o(1)}\).
\item Choose \(\delta>0\) with \(max{\rho,q}+\delta<1+\epsilon\). For all large valid N the composed circuit has size below \(N^{1+\epsilon}\), contradicting the i.o. source lower bound.
\item Audit the unconditional falsifier: replacing each SAT call by exhaustive search costs N\textasciicircum{}rho+N\textasciicircum{}q2\textasciicircum{}L poly(L), which is \(N^{q+\omega(1)}\) for \(d>1\), so it does not itself refute an \(N^{1+\epsilon}\) lower bound.
\end{enumerate}

\paragraph{Equation 1: carrier\_definition.}
\[
H_N(x)=A_N\bigl(x,\operatorname{SAT}_L(G_N(x,1)),\ldots,\operatorname{SAT}_L(G_N(x,M))\bigr)
\]
\emph{Valid when:} Valid when \(G_{N}\) emits exactly M generated SAT instances and \(A_{N}\) consumes only their SAT bits.

\paragraph{Equation 2: composition\_upper\_bound.}
\[
\operatorname{SIZE}(H_N)\le O\bigl(N^\rho+M(N)(L(N)^c+L(N))\bigr)
\]
\emph{Valid when:} Valid when \(\mathrm{SAT}_{L}\) has size O(L\textasciicircum{}c) circuits and each target call is substituted as a separate subcircuit.

\paragraph{Equation 3: polylog\_amortization.}
\[
M(N)L(N)^c=N^q(\log N)^{dc}=N^{q+o(1)}
\]
\emph{Valid when:} Valid when c and d are fixed constants and N tends through the chosen valid lengths.

\paragraph{Equation 4: threshold\_contradiction.}
\[
\max{\rho,q}< 1+\epsilon\ \Longrightarrow\ \operatorname{SIZE}(H_N)< N^{1+\epsilon}\text{ for all large valid }N
\]
\emph{Valid when:} Applies after choosing a positive margin between max\{rho,q\} and the claimed source lower-bound exponent.

\paragraph{Next analysis.} Instantiate \(G_{N}\) as a near-linear Cook-Levin micro-instance generator with \(M=N\) and L=(log N)\textasciicircum{}d; compute exact \(\rho\), verify \(H_{N}\) is in NP with poly(N) uniform verification, and test whether any structural simplification gives an \(N^{1+\epsilon}\) upper-bound falsifier.

\paragraph{Falsifiers.}
\begin{itemize}
\item falsifier: threshold collapse occurs if \(max{\rho,q}\ge 1+\epsilon\), because SAT circuits then give no contradiction to the claimed source lower bound.
\item falsifier: lookup collapse occurs when N\textasciicircum{}rho+N\textasciicircum{}q2\textasciicircum{}\{L(N)\}poly(L(N))<N\textasciicircum{}\{1+epsilon\} on infinitely many valid lengths.
\item falsifier: nonuniform leakage occurs if \(G_{N}\) or \(A_{N}\) encodes \(\mathrm{SAT}_{L}\) advice not charged to \(N^\rho\).
\item falsifier: model mismatch occurs if the source lower bound is for formulas, \(\mathrm{AC0}\), or another restricted model while the target substitution requires unrestricted P/poly circuits.
\item falsifier: a large constructive natural property proving the source lower bound would collide with pseudorandom-function assumptions at the same \(N^{1+\epsilon}\) scale.
\item falsifier: a fully relativizing source lower-bound proof would be suspect because oracle worlds with \(\mathrm{NP}^{O}\) subset \(\mathrm{P}^{O}/\mathrm{poly}\) should invalidate the target consequence.
\end{itemize}

\subsection*{cycle-0015: Hardness-Magnification Bridge Realization}

\paragraph{Execution.} research; azure\_responses\_api; attempt compact\_retry; 289.508 seconds; 1 lemmas, 3 equations, 4 derivation steps, 4 falsifiers.

\paragraph{Learning.} Cycle 15 adds the exact lookup boundary. Polylog queries always preserve all SAT polynomial exponents in the substitution term, but source plausibility requires q+A(log N)\textasciicircum{}\{d-1\} not fall below 1+epsilon. Thus d1 is safely valid from lookup.

\paragraph{Pressure.} No explicit no-advice generator \(G_{N}\) is known for which \(H_{N}\) has an i.o. \(N^{1+\epsilon}\) general-circuit lower bound by a non-natural method.

\paragraph{Derivation steps.}
\begin{enumerate}
\item Fix \(G_{N}\) of size \(N^\rho\) producing \(m=N^{q}\) SAT queries of length L=A(log N)\textasciicircum{}d and define \(H_{N}\) as their OR.
\item Assume SAT has circuits of size \(n^{c}\). Substitute the length-L circuit \(S_{L}\) into every query and add the generator plus OR overhead.
\item Since L\textasciicircum{}c=N\textasciicircum{}\{o(1)\}, the composed circuit has size \(N^{max{\rho,q}+o(1)}\), contradicting any i.o. \(H_{N}\) lower bound above \(N^{1+\epsilon}\) when \(max{\rho,q}<1+\epsilon\).
\item Audit the independent lookup circuit for \(\mathrm{SAT}_{L}\). It is subthreshold exactly when q+A(log N)\textasciicircum{}\{d-1\}<1+epsilon eventually, producing a source-level falsifier.
\end{enumerate}

\paragraph{Equation 1: carrier.}
\[
H_N(x)=\bigvee_{i\le N^q}\mathrm{SAT}(G_{N,i}(x)),\qquad |G_{N,i}(x)|=L(N)=A(\log_2 N)^d
\]
\emph{Valid when:} Valid when \(G_{N}\) is a uniform generator whose outputs are well-formed SAT encodings of length L(N).

\paragraph{Equation 2: sat\_substitution\_upper.}
\[
\mathrm{SAT}\in \mathrm{SIZE}(n^c)\Rightarrow C(H_N)\le N^\rho+N^qL(N)^c+O(N^q)=N^{\max{\rho,q}+o(1)}
\]
\emph{Valid when:} Applies after the assumed SAT circuit is length-uniform nonuniform advice for exactly L(N)-bit formulas.

\paragraph{Equation 3: lookup\_audit.}
\[
C_{\mathrm{lookup}}(H_N)\le N^\rho+N^q2^{L(N)},\quad \lambda(N)=\log_N2^{L(N)}=A(\log_2N)^{d-1}
\]
\emph{Valid when:} Requires the lookup audit to be compared with the same \(N^{1+\epsilon}\) source threshold.

\paragraph{Next analysis.} Instantiate the boundary cases \(d=1\) with \(A=1+\epsilon-q\) and \(d=2\); search for a generator whose \(H_{N}\) avoids known small-circuit constructions without smuggling SAT advice.

\paragraph{Falsifiers.}
\begin{itemize}
\item Lookup falsifier: if \(d<1\), or \(d=1\) with \(q+A<1+\epsilon\), the table for \(\mathrm{SAT}_{L}\) gives \(C(H_{N})<N^{1+\epsilon}\) and the source premise is impossible.
\item Threshold falsifier: if \(\rho\ge 1+\epsilon\) or \(q\ge 1+\epsilon\), SAT-substitution no longer yields a contradiction to the \(H_{N}\) lower bound.
\item Advice-leakage falsifier: if \(G_{N}\) or the lower-bound proof embeds satisfiability bits for the generated query set, the target-preserving interface to SAT in P/poly is invalid.
\item Natural-proof falsifier: if the \(H_{N}\) lower bound is obtained via a large constructive property useful against SIZE(N\textasciicircum{}\{1+epsilon\}), standard PRF assumptions block it for general circuits.
\end{itemize}

\subsection*{cycle-0016: Hardness-Magnification Bridge Realization}

\paragraph{Execution.} research; azure\_responses\_api; attempt compact\_retry; 244.266 seconds; 1 lemmas, 3 equations, 4 derivation steps, 4 falsifiers.

\paragraph{Learning.} Cycle 16 sharpens the generated-OR route: the valid transfer to \(\mathrm{SAT}\notin \mathrm{P}/\mathrm{poly}\) needs only \(q,\rho<1+\epsilon\), but source plausibility additionally requires variable entropy q+v/log N at least 1+epsilon. The old \(2^{L}\) lookup audit was too weak; the real falsifier is low variable count.

\paragraph{Pressure.} No explicit target-clean generator is known whose OR carrier has a non-natural unrestricted \(N^{1+\epsilon}\) circuit lower bound while also passing the variable-entropy gate.

\paragraph{Derivation steps.}
\begin{enumerate}
\item Fix \(H_{N}\) with \(M=N^{q}\) generated SAT queries, length L=(log N)\textasciicircum{}d, generator cost \(N^\rho\), and \(max{q,\rho}<1+\epsilon\).
\item Assume SAT has circuits of size \(n^{c}\). Substitute the length-L SAT circuit into every query; since L\textasciicircum{}c=N\textasciicircum{}\{o(1)\}, \(H_{N}\) has size \(N^{max{q,\rho}+o(1)}\).
\item For large valid N this contradicts \(\mathrm{SIZE}(H_{N})>N^{1+\epsilon}\); hence the bridge would imply \(\mathrm{SAT}\notin \mathrm{P}/\mathrm{poly}\), not merely a weak restatement of P vs NP.
\item Audit the source: independent of SAT circuits, each query can be solved by enumerating v(N) variables, so the lower bound is impossible below the threshold \(q+v(N)/log_{2}\) N \(\ge \) 1+epsilon.
\end{enumerate}

\paragraph{Equation 1: carrier\_definition.}
\[
H_N(x)=\bigvee_{i=1}^{M} SAT(\Phi_i(x)),\quad M=N^q,\quad |\Phi_i(x)|\le L=(\log N)^d
\]
\emph{Valid when:} Valid when the compact generator computes each queried SAT instance uniformly without embedding SAT advice.

\paragraph{Equation 2: sat\_substitution\_upper\_bound.}
\[
SAT\in SIZE(n^c)\Rightarrow SIZE(H_N)\le N^{\rho}+N^q L^c\,polylog(N)=N^{\max{\rho,q}+o(1)}
\]
\emph{Valid when:} Applies after the reduction preserves nonuniform target circuits and q,rho are constants independent of c.

\paragraph{Equation 3: variable\_enumeration\_falsifier.}
\[
SIZE(H_N)\le N^{\rho}+N^q2^{v(N)}poly(L)=N^{\max{\rho,q+v(N)/\log_2N}+o(1)}
\]
\emph{Valid when:} Valid when generated formulas have at most v(N) variables and deterministic exhaustive assignment evaluation is allowed.

\paragraph{Next analysis.} Instantiate \(G_{N}\) with high-essential-variable polylog 3CNF queries and test whether any structural simplification gives \(\mathrm{SIZE}(H_{N})\le N^{1+\epsilon-o(1)}\) despite v(N) superlogarithmic.

\paragraph{Falsifiers.}
\begin{itemize}
\item Entropy falsifier: if generated formulas use only \(v(N)<(1+\epsilon-q-\eta)log_{2}\) N variables, exhaustive assignment gives subthreshold circuits for \(H_{N}\).
\item Advice-leakage falsifier: if \(G_{N}\) encodes SAT truth bits or hard functions as nonuniform advice, the transfer no longer certifies \(\mathrm{SAT}\notin \mathrm{P}/\mathrm{poly}\).
\item Sparsity falsifier: if the query family lands in an easy SAT fragment or repeats \(N^{o(1)}\) essential formulas, \(H_{N}\) may have circuits below \(N^{1+\epsilon}\).
\item Naturality falsifier: a large constructive property proving the \(H_{N}\) lower bound would collide with standard PRG-based natural-proofs barriers unless the proof is non-natural or outside their scope.
\end{itemize}

\subsection*{cycle-0017: Hardness-Magnification Bridge Realization: positioned-range entropy audit}

\paragraph{Execution.} research; azure\_responses\_api; attempt primary; 103.969 seconds; 3 lemmas, 5 equations, 6 derivation steps, 6 falsifiers.

\paragraph{Learning.} Cycle 17 adds a sharper necessary condition for the generated-OR route: not only \(q,\rho<1+\epsilon\), but the positioned effective range entropy sum\_g \(min{|R_{g}|,2^{v}}\) must reach the lower-bound threshold. This converts low-range query generation into an explicit falsifier and preserves the valid target transfer only for high-entropy, target-clean carriers.

\paragraph{Pressure.} No explicit target-clean carrier is known with q,rho=1+epsilon-o(1), and a non-natural \(N^{1+\epsilon}\) lower bound against general circuits.

\paragraph{Derivation steps.}
\begin{enumerate}
\item Fix \(K_{N}\) with ordinary gates plus SAT oracle gates \(g=1,...,t_{N}\), \(t_{N}\le N^{q}\), query length L\_N=(log N)\textasciicircum{}d, and overhead size \(N^\rho\).
\item Assume SAT has circuits \(C_{L}\) of size \(L^{c}\). Replace every SAT gate by \(C_{L_{N}}\); the resulting \(H_{N}\) circuit has size N\textasciicircum{}rho+t\_NL\_N\textasciicircum{}c=N\textasciicircum{}\{max\{rho,q\}+o(1)\}.
\item Thus an eventual lower bound \(H_{N}\) notin SIZE(N\textasciicircum{}\{1+epsilon\}) contradicts \(\mathrm{SAT}\in \mathrm{SIZE}(n^{c})\) for each fixed c if \(max{\rho,q}<1+\epsilon\).
\item Audit the source claim without assuming SAT circuits. For each oracle position g, replace SAT on the finite image \(R_{g,N}\) by a DNF table for satisfiable members of \(R_{g,N}\), costing O(|R\_\{g,N\}|L\_N), or by exhaustive assignment enumeration, costing 2\textasciicircum{}\{v\_N\}poly(L\_N).
\item Summing over positions gives \(H_{N}\) in size O(N\textasciicircum{}rho)+poly(L\_N) sum\_g \(min{|R_{g,N}|,2^{v_{N}}}\).
\item Therefore any proposed source lower bound with \(max{\rho,e_{eff}}<1+\epsilon\) by a margin is internally inconsistent; high positioned effective entropy becomes a necessary certificate.
\end{enumerate}

\paragraph{Equation 1: target-substitution-size.}
\[
S_{\mathrm{SAT}}(N,c) \le N^{\rho}+t_N L_N^c\,\mathrm{poly}(L_N)=N^{\max{\rho,q}+o(1)}
\]
\emph{Valid when:} Valid when SAT has nonuniform circuits of size \(n^{c}\) and every oracle query length is L\_N=(log N)\textasciicircum{}d.

\paragraph{Equation 2: positioned-effective-entropy.}
\[
e_{\mathrm{eff}}(N)=\log_N\left(\sum_{g\in K_N}\min{|R_{g,N}|,2^{v_N}}\right)
\]
\emph{Valid when:} Valid when each oracle position has a well-defined image over all N-bit inputs of \(H_{N}\).

\paragraph{Equation 3: unconditional-range-upper-bound.}
\[
S_{\mathrm{uncond}}(N) \le O(N^{\rho})+\mathrm{poly}(L_N)\sum_g\min{|R_{g,N}|,2^{v_N}}=N^{\max{\rho,e_{\mathrm{eff}}}+o(1)}
\]
\emph{Valid when:} Applies when oracle gates may be replaced by per-position range DNFs or exhaustive assignment checks.

\paragraph{Equation 4: magnification-implication.}
\[
H_N\notin SIZE(N^{1+\epsilon})\ \wedge\ \max{\rho,q}< 1+\epsilon\ \Longrightarrow\ SAT\notin P/poly
\]
\emph{Valid when:} Requires the source lower bound to be valid in the same nonuniform circuit model used by the target substitution.

\paragraph{Equation 5: source-falsifier-threshold.}
\[
\max{\rho,e_{\mathrm{eff}}}< 1+\epsilon-\delta\ \Longrightarrow\ H_N\in SIZE(N^{1+\epsilon-\delta/2})
\]
\emph{Valid when:} Valid when the entropy bound holds eventually and polylogarithmic replacement factors are absorbed by the \(\delta\) margin.

\paragraph{Next analysis.} Instantiate \(K_{N}\) from a specific extractor-like generator of polylog-length CNFs and compute \(|R_{g}|\), \(v_{N}\), q, and \(\rho\) exactly to see whether the positioned entropy threshold can be met without advice leakage.

\paragraph{Falsifiers.}
\begin{itemize}
\item threshold collapse: \(max{\rho,q}\ge 1+\epsilon\) makes SAT substitution too large to contradict the source lower bound.
\item positioned range entropy falsifier: sum\_g min\{|R\_g|,2\textasciicircum{}v\}<N\textasciicircum{}\{1+epsilon-delta\} yields an unconditional subthreshold circuit for \(H_{N}\).
\item low variable entropy: \(q+v_{N}/log\) \(N<1+\epsilon\) by a margin lets exhaustive assignment simulation kill the source claim.
\item nonuniform leakage: \(K_{N}\) may hide SAT answers in advice-like query selection or gate wiring, invalidating target preservation.
\item model mismatch: a lower bound only for a restricted circuit class \(\Gamma\) does not imply \(\mathrm{NP}\not\subset \mathrm{P}/\mathrm{poly}\) without a separate lift.
\item natural-proof obstruction: a large constructive property proving the source lower bound against P/poly would collide with standard PRG-based natural proofs barriers.
\end{itemize}

\subsection*{cycle-0018: Hardness-Magnification Bridge Realization: positioned entropy cover gate}

\paragraph{Execution.} research; azure\_responses\_api; attempt primary; 184.781 seconds; 3 lemmas, 4 equations, 5 derivation steps, 5 falsifiers.

\paragraph{Learning.} Cycle 18 proves the exact positioned entropy cover upper bound. The route now has two independent gates: target transfer needs \(max{\rho,q}=N^{1+\epsilon-o(1)}\) after charging the generator. Low variable count or low positioned range is a concrete falsifier, not a folklore barrier.

\paragraph{Pressure.} The derivation proves only an upper-bound/falsifier lemma. It does not prove \(C(H_{N})>N^{1+\epsilon}\) for any explicit high-entropy carrier.

\paragraph{Derivation steps.}
\begin{enumerate}
\item Fix N and position i; let \(R_{i}\) be the set of formulas reachable as x varies.
\item For this position, build a nonuniform evaluator either by hardwiring SAT(y) for every y in \(R_{i}\) and comparing Gamma\_\{N,i\}(x) to the dictionary, or by brute-force checking all \(2^{v_{i}}\) assignments of the produced formula.
\item The cheaper evaluator costs \(min{|R_{i}|,2^{v_{i}}}\) \(L_{N}^{O(1)}\) gates, after \(\Gamma_{N}\) has produced the query bits.
\item Run the batched generator once, evaluate all \(M_{N}\) positioned queries, and OR the \(M_{N}\) answers.
\item This gives C(H\_N)<=N\textasciicircum{}\{rho+o(1)\}+N\textasciicircum{}\{q+o(1)\}+B\_N \(L_{N}^{O(1)}\); taking the contrapositive yields the entropy threshold gate.
\end{enumerate}

\paragraph{Equation 1: carrier\_definition.}
\[
H_N(x)=\bigvee_{i=1}^{M_N} SAT(\Gamma_{N,i}(x)),\quad M_N=N^{q+o(1)},\quad L_N=(\log N)^d
\]
\emph{Valid when:} Valid when the batched generator exposes all query positions and the SAT predicate is evaluated only on formulas of length \(L_{N}\).

\paragraph{Equation 2: positioned\_effective\_entropy.}
\[
B_N=\sum_{i=1}^{M_N}\min{|R_i|,2^{v_i}},\quad \theta_N=\log_N B_N
\]
\emph{Valid when:} Valid when each position has its own range dictionary and variable bound, so ranges from different positions are not merged for free.

\paragraph{Equation 3: cover\_upper\_bound.}
\[
C(H_N)\le N^{\rho+o(1)}+N^{q+o(1)}+B_N\cdot L_N^{O(1)}
\]
\emph{Valid when:} Applies when nonuniform circuits may hardwire per-position dictionaries or instead brute-force all variables of each produced formula.

\paragraph{Equation 4: entropy\_gate\_contrapositive.}
\[
C(H_N)> N^{1+\epsilon}\ \wedge\ \max{\rho,q}< 1+\epsilon\ \Rightarrow\ \theta_N\ge 1+\epsilon-o(1)
\]
\emph{Valid when:} Requires the source lower bound, generator-size accounting, and positioned cover upper bound to hold on the same valid lengths.

\paragraph{Next analysis.} Instantiate the affine-disperser generated-OR family and compute \(B_{N}\) exactly as a function of \(\alpha\), q, and clause density; then audit whether any known lower-bound technique for this restricted \(H_{N}\) avoids the natural-proofs density condition.

\paragraph{Falsifiers.}
\begin{itemize}
\item positioned entropy falsifier: \(B_{N}\le N^{1+\epsilon-\delta}\) on infinitely many valid lengths gives subthreshold circuits for \(H_{N}\).
\item generator-cost falsifier: charged batched generation size \(S_{charged}(\Gamma_{N})\ge N^{1+\epsilon-o(1)}\) destroys the barely-superlinear magnification window.
\item target-transfer falsifier: \(max{\rho,q}\ge 1+\epsilon\) means SAT circuit substitution no longer contradicts the source lower bound.
\item variable-count falsifier: if \(v_{i}\le \alpha\) log N uniformly and \(q+\alpha<1+\epsilon\), then \(B_{N}\le N^{q+\alpha+o(1)}\) and the source lower bound is impossible.
\item natural-proof falsifier: a dense P/poly-constructive property proving the \(H_{N}\) lower bound is blocked under standard PRF assumptions, unless the certificate is non-natural or sufficiently sparse.
\end{itemize}

\subsection*{cycle-0019: Hardness-Magnification Bridge Realization: semantic range leakage gate}

\paragraph{Execution.} research; azure\_responses\_api; attempt primary; 174.05 seconds; 3 lemmas, 4 equations, 7 derivation steps, 7 falsifiers.

\paragraph{Learning.} Cycle 19 sharpens the generated-OR route: positioned entropy is not the decisive source gate. The valid bridge uses the semantic cover \(S_{N}=sum_{j}\) C\_\{R\_\{N,j\}\}(SAT). Subthreshold \(S_{N}\) is a concrete falsifier for \(H_{N}\) lower bounds, while \(q,\rho<1+\epsilon\) still gives the target-preserving implication from a valid \(H_{N}\) lower bound to \(\mathrm{SAT}\notin \mathrm{P}/\mathrm{poly}\).

\paragraph{Pressure.} The derivation does not supply an explicit \(G_{N}\) for which both \(H_{N}\) has a general-circuit lower bound \(N^{1+\epsilon}\) and \(S_{N}=sum_{j}\) C\_\{R\_\{N,j\}\}(SAT) is provably above \(N^{1+\epsilon-o(1)}\).

\paragraph{Derivation steps.}
\begin{enumerate}
\item Let \(R_{N,j}\) be the j-th positioned range of the bulk generator \(G_{N}\).
\item For each j, choose a minimum L-input circuit \(D_{N,j}\) agreeing with SAT on \(R_{N,j}\).
\item Build a circuit for \(H_{N}\) by computing the full list G\_N(x), feeding the j-th formula to \(D_{N,j}\), and OR-ing all M outputs.
\item This gives C(H\_N)<=O(C(G\_N)+sum\_j \(C_{R_{N,j}}(\mathrm{SAT})+M\) L).
\item If SAT has global circuits of size \(n^{c}\), take every \(D_{N,j}\) to be the restriction of the L-bit SAT circuit; since L=(log N)\textasciicircum{}d, the substitution term is \(N^{q+o(1)}\).
\item Thus a valid \(N^{1+\epsilon}\) lower bound for \(H_{N}\) with \(q,\rho<1+\epsilon\) refutes every fixed polynomial-size SAT circuit family.
\item Conversely, the same upper bound proves that any subthreshold aggregate positioned range decider is a falsifier for the claimed source lower bound.
\end{enumerate}

\paragraph{Equation 1: semantic\_upper\_bound.}
\[
C(H_N) \le O\left(C(G_N)+\sum_{j\le M} C_{R_{N,j}}(\mathrm{SAT})+M L\right)
\]
\emph{Valid when:} Valid when \(G_{N}\) is a bulk circuit outputting all positioned queries and circuit size counts shared generator gates once.

\paragraph{Equation 2: global\_sat\_substitution.}
\[
\mathrm{SAT}\in \mathrm{SIZE}(n^c) \Rightarrow C(H_N) \le N^{\rho}+N^q (\log N)^{dc}+N^{q+o(1)}=N^{\max{\rho,q}+o(1)}
\]
\emph{Valid when:} Applies after the generated formulas have length L=(log N)\textasciicircum{}d and every fixed SAT circuit exponent contributes only \(N^{o(1)}\) overhead.

\paragraph{Equation 3: magnification\_consequence.}
\[
C(H_N)> N^{1+\epsilon}\ \wedge\ \max{\rho,q}< 1+\epsilon \Rightarrow \mathrm{SAT}\notin \mathrm{P/poly}
\]
\emph{Valid when:} Requires the lower bound for \(H_{N}\) to hold eventually and the target transfer to preserve nonuniform circuit size.

\paragraph{Equation 4: semantic\_gate.}
\[
C(H_N)> N^{1+\epsilon}\ \wedge\ q,\rho< 1+\epsilon \Rightarrow \sum_{j\le M} C_{R_{N,j}}(\mathrm{SAT})\ge N^{1+\epsilon-o(1)}
\]
\emph{Valid when:} Valid when q and \(\rho\) are both strictly below the claimed source lower-bound exponent 1+epsilon.

\paragraph{Next analysis.} Instantiate \(G_{N}\) with a concrete expander or extractor-based 3-CNF generator, compute \(R_{N,j}\) structure, and test whether SAT on those positioned ranges collapses to a known easy subclass or requires a new semantic lower-bound certificate.

\paragraph{Falsifiers.}
\begin{itemize}
\item semantic leakage falsifier: SAT restricted to the positioned ranges has aggregate circuit complexity below \(N^{1+\epsilon-o(1)}\).
\item threshold collapse: max\{q,rho\} reaches or exceeds 1+epsilon, so SAT substitution no longer contradicts the source lower bound.
\item bulk-generator mismatch: \(G_{N}\) is only random-access with per-query cost \(N^\rho\), changing the positive upper bound to about \(N^{q+\rho}\).
\item nonuniform leakage: the construction or proof accidentally supplies the positioned deciders \(D_{N,j}\) as advice for \(H_{N}\).
\item low-label-complexity range: high \(B_{N}\) is achieved using formulas from classes where satisfiability is easy, such as Horn, 2-CNF, bounded treewidth, or planted trivially satisfiable instances.
\item natural-proof obstruction: a constructive dense certificate that many SAT range labels evade small circuits would be natural; the surviving route must be sparse, non-natural, or use a nonrelativizing proof object.
\item oracle falsifier probe: if an oracle makes all generated range labels simple while preserving the black-box interface, then a relativizing proof of the semantic gate cannot be valid.
\end{itemize}

\subsection*{cycle-0020: Hardness-Magnification Bridge Realization: Semantic Block-Cover Gate}

\paragraph{Execution.} research; azure\_responses\_api; attempt primary; 331.928 seconds; 2 lemmas, 4 equations, 7 derivation steps, 6 falsifiers.

\paragraph{Learning.} Cycle 20 sharpens the generated-OR route: the valid target bridge is unchanged, but the source plausibility gate must minimize over semantic block covers. Singleton range complexity is too weak because OR compression and nonuniform block advice can falsify an \(N^{1+\epsilon}\) source lower bound. This is route-level progress, not a proof of \(\mathrm{P} \ne \mathrm{NP}\).

\paragraph{Pressure.} No direct-sum theorem converts \(M_{N}\ge N^{1+\epsilon-o(1)}\) into \(C(H_{N})>N^{1+\epsilon}\). A circuit for \(H_{N}\) may compute the union of positive preimages without decomposing along any chosen partition.

\paragraph{Derivation steps.}
\begin{enumerate}
\item Fix N, \(t=N^{q}\), L=(log N)\textasciicircum{}d, and generator \(G_{N}\) producing all query strings G\_\{N,1\}(x),...,G\_\{N,t\}(x).
\item Assume \(\mathrm{SAT}_{L}\) has circuits of size \(L^{c}\). Compute \(G_{N}\) once, feed the same \(\mathrm{SAT}_{L}\) circuit to each of the t positions, and OR the answers; this gives size N\textasciicircum{}\{rho+o(1)\}+N\textasciicircum{}q \(L^{c+O(1)}\).
\item Since L is polylogarithmic, any fixed SAT P/poly exponent contributes only \(N^{o(1)}\) per query, so \(max{q,\rho}<1+\epsilon\) gives the valid contrapositive target implication.
\item For the source audit, take any nonuniform partition Pi of positions. For each block B, use an optimal range-restricted circuit \(D_{B}\) of size \(K_{B}\) deciding whether any generated formula in that block is satisfiable.
\item Compute \(G_{N}\) once, route each block tuple to \(D_{B}\), then OR the block outputs. This yields C(H\_N)<=N\textasciicircum{}\{rho+o(1)\}+N\textasciicircum{}\{q+o(1)\}+sum\_B \(K_{B}\).
\item Therefore a claimed lower bound \(C(H_{N})>N^{1+\epsilon}\) forces \(M_{N}\ge N^{1+\epsilon-o(1)}\) whenever generator and routing costs have fixed slack below the threshold.
\item This repairs the previous semantic-cover gate: singleton \(S_{N}\) is only the partition into singletons and can miss block-level OR compression.
\end{enumerate}

\paragraph{Equation 1: target\_substitution.}
\[
SAT_L\in SIZE(L^c)\ \Rightarrow\ C(H_N)\le N^{\rho+o(1)}+N^q L^{c+O(1)}=N^{\max{\rho,q}+o(1)}.
\]
\emph{Valid when:} Valid when the generator is shareable, all queries have length L=(log N)\textasciicircum{}d, and one nonuniform \(\mathrm{SAT}_{L}\) circuit is reused for every generated query position.

\paragraph{Equation 2: magnification\_contrapositive.}
\[
[\ C(H_N)> N^{1+\epsilon}\ \text{eventually}\ \wedge\ \max{q,\rho}< 1+\epsilon\ ]\ \Rightarrow\ SAT\notin P/poly.
\]
\emph{Valid when:} Applies when SAT in P/poly would give some fixed exponent c and polylogarithmic query length makes L\textasciicircum{}c=N\textasciicircum{}\{o(1)\} under the target-preserving interface.

\paragraph{Equation 3: block\_cover\_upper\_bound.}
\[
C(H_N)\le N^{\rho+o(1)}+N^{q+o(1)}+\sum_{B\in\Pi} K_B.
\]
\emph{Valid when:} Valid when Pi is any nonuniform partition of query positions and each \(K_{B}\) circuit decides the block OR predicate only on its generated tuple range \(R_{B}\).

\paragraph{Equation 4: necessary\_block\_gate.}
\[
C(H_N)> N^{1+\epsilon}\ \Rightarrow\ M_N:=\min_{\Pi}\sum_{B\in\Pi}K_B\ge N^{1+\epsilon-o(1)}.
\]
\emph{Valid when:} Requires the generator and routing overheads to be below \(N^{1+\epsilon}\) by fixed polynomial slack so that a subthreshold block cover would contradict the source lower bound.

\paragraph{Next analysis.} Instantiate \(M_{N}\) for a nontrivial generator without dominated coordinates, e.g. a universal verifier family whose queries are Cook-Levin encodings of polylog-size computations, and compute whether block predicates admit compression below \(N^{1+\epsilon}\).

\paragraph{Falsifiers.}
\begin{itemize}
\item block semantic collapse: \(M_{N}<N^{1+\epsilon-o(1)}\) for a nonuniform partition even though singleton \(S_{N}\) appears large
\item dominating coordinate: each block contains an always-satisfiable generated formula, making \(H_{N}\) or block predicates easy
\item threshold collapse: \(max{q,\rho}\ge 1+\epsilon\) invalidates the target-preserving magnification inequality
\item nonuniform leakage: a source lower-bound argument ignores range-restricted per-block circuits that are legal in the general-circuit source model
\item model mismatch: proving lower bounds only for uniform algorithms or restricted circuits does not imply the required nonuniform SIZE lower bound for \(H_{N}\)
\item natural-proof pressure: a constructive large property proving the \(H_{N}\) lower bound against P/poly would need to be non-natural or outside the standard PRG-based natural proofs obstruction
\end{itemize}

\subsection*{cycle-0021: Hardness-Magnification Bridge Realization: Shared-Generator Semantic Block-Cover Gate}

\paragraph{Execution.} research; azure\_responses\_api; attempt primary; 135.544 seconds; 2 lemmas, 4 equations, 6 derivation steps, 6 falsifiers.

\paragraph{Learning.} Cycle 21 replaces the previous block-cover warning by an exact shared-generator semantic-cover lemma. The valid target bridge is unchanged, but the source lower bound must beat \(C(G_{N})+min_{Pi}\) \(sum_{B}\) C\_\{R\_B\}(OR-SAT\_B). This creates a concrete falsifier: any subthreshold semantic block partition kills the generated-OR carrier before it can contribute to \(\mathrm{SAT}\notin \mathrm{P}/\mathrm{poly}\).

\paragraph{Pressure.} No candidate generator is known for which \(min_{Pi}\) \(sum_{B}\) C\_\{R\_B\}(OR-SAT\_B) can be proved at \(N^{1+\epsilon-o(1)}\) in general circuits.

\paragraph{Derivation steps.}
\begin{enumerate}
\item Let \(M=N^{q}\) and let G\_N(x) output the M formulas y\_1(x),...,y\_M(x) of length L=(log N)\textasciicircum{}d using \(N^{\rho+o(1)}\) gates.
\item Fix a partition Pi of [M]. For each block B, define \(R_{B}\) as the set of generated tuples y\_B(x) and choose a smallest circuit \(Q_{B}\) agreeing with OR\_\{j in B\} SAT(y\_j) on \(R_{B}\).
\item Compose one copy of \(G_{N}\) with all \(Q_{B}\) circuits, wire each \(Q_{B}\) to its block tuple, and OR the block answers. This computes \(H_{N}\) exactly.
\item The resulting size is \(N^{\rho+o(1)}+sum_{B}\) C\_\{R\_B\}(OR-SAT\_B)+O(|Pi|), proving the shared-generator cover upper bound.
\item For singleton blocks, if SAT has \(L^{c}-size\) circuits, then each block cost is at most \(L^{c}\), giving C(H\_N)<=N\textasciicircum{}\{rho+o(1)\}+N\textasciicircum{}qL\textasciicircum{}c and hence the target-preserving contradiction when \(max{\rho,q}<1+\epsilon\).
\item Therefore a valid source lower bound must simultaneously rule out all nonuniform semantic block covers; otherwise the cover is an explicit falsifier.
\end{enumerate}

\paragraph{Equation 1: \(E1_{target,substitution,upper,bound}\).}
\[
C(H_N) \le N^{\rho+o(1)}+N^q\,C(SAT_L)+O(N^q) \le N^{\max{\rho,q}+o(1)}
\]
\emph{Valid when:} Valid when SAT has circuits of size \(L^{c}\) for a fixed c and L=(log N)\textasciicircum{}d makes L\textasciicircum{}c=N\textasciicircum{}\{o(1)\}.

\paragraph{Equation 2: \(E2_{target,preserving,magnification}\).}
\[
\big(C(H_N)> N^{1+\epsilon}\ \text{and}\ \max{\rho,q}< 1+\epsilon\big)\Rightarrow SAT\notin P/poly
\]
\emph{Valid when:} Applies after the generated-OR carrier is target-clean and the lower bound holds against nonuniform circuits for all sufficiently large N.

\paragraph{Equation 3: \(E3_{shared,block,cover,upper,bound}\).}
\[
C(H_N)\le C(G_N)+\min_{\Pi}\left(\sum_{B\in\Pi} C_{R_B}(OR\text{-}SAT_B)+O(|\Pi|)\right)
\]
\emph{Valid when:} Valid when each block circuit only needs to agree on the realized tuple range emitted by the shared generator.

\paragraph{Equation 4: \(E4_{source,gate,consequence}\).}
\[
C(H_N)> N^{1+\epsilon}\Rightarrow \forall\Pi:\ C(G_N)+\sum_{B\in\Pi} C_{R_B}(OR\text{-}SAT_B)\ge N^{1+\epsilon-o(1)}
\]
\emph{Valid when:} Requires the claimed source lower bound to be stronger than every semantic block-cover upper bound, including nonuniform block advice.

\paragraph{Next analysis.} Run the block-size thought experiment for \(b=N^\theta\): derive upper bounds on U(theta) from range cardinality, feature compression, and OR compression, then identify whether any \(\theta\) gives N\textasciicircum{}\{rho\}+N\textasciicircum{}\{q-theta\}U(theta)<N\textasciicircum{}\{1+epsilon\}.

\paragraph{Falsifiers.}
\begin{itemize}
\item semantic block-cover collapse: some partition Pi has \(C(G_{N})+sum_{B}\) \(C_{R_{B}}(OR-\mathrm{SAT}_{B})<N^{1+\epsilon-\delta}\)
\item range-compression falsifier: block OR-SAT values factor through short features \(A_{B}\) with total cost below the claimed threshold
\item nonuniform block advice leakage: \(Q_{B}\) encodes the satisfiability pattern on the realized range and bypasses the intended SAT hardness
\item generator output-size tax: producing all query bits already forces \(\rho\) or q to violate \(max{\rho,q}<1+\epsilon\)
\item lookup boundary failure: \(d\le 1\) with constants allowing truth-table lookup below \(N^{1+\epsilon}\)
\item natural-proof obstruction: a large constructive method proving the source lower bound for many generated carriers would distinguish pseudorandom functions, so a valid proof must be range-specific or non-natural
\end{itemize}

\subsection*{cycle-0022: Hardness-Magnification Bridge Realization}

\paragraph{Execution.} research; azure\_responses\_api; attempt primary; 128.185 seconds; 3 lemmas, 4 equations, 6 derivation steps, 5 falsifiers.

\paragraph{Learning.} Cycle 22 isolates the next exact gate: the generated-OR bridge is valid only if the source lower bound defeats nonuniform promise solvers on the actual SAT range. The new quantitative obstruction is \(T_{N}>N^{1+\epsilon-q-o(1)}\) for SAT restricted to \(U_{N}\), plus the earlier semantic block-cover threshold. This is route-level progress, not a proof of \(\mathrm{P} \ne \mathrm{NP}\).

\paragraph{Pressure.} No current argument proves the required promise-SAT or block-cover lower bound for an explicit high-entropy generated range \(U_{N}\); entropy and query length do not imply semantic circuit hardness.

\paragraph{Derivation steps.}
\begin{enumerate}
\item Fix c and assume SAT has length-n circuits of size \(n^{c}\).
\item For each N, copy the length-L SAT circuit into each of the \(M=N^{q}\) query positions produced by \(G_{N}\).
\item The resulting circuit for \(H_{N}\) has size N\textasciicircum{}rho+N\textasciicircum{}q(log N)\textasciicircum{}\{dc\}+O(N\textasciicircum{}qL)=N\textasciicircum{}\{max\{rho,q\}+o(1)\}.
\item If \(max{\rho,q}N^{1+\epsilon}\) contradicts the assumed SAT circuits, giving the route-level target consequence \(\mathrm{SAT}\notin \mathrm{P}/\mathrm{poly}\).
\item Now audit the source claim: if SAT restricted to \(U_{N}\) has promise circuits of size \(T_{N}\), the same substitution gives C(H\_N)<=N\textasciicircum{}rho+N\textasciicircum{}qT\_N+O(N\textasciicircum{}qL).
\item Therefore any source lower bound must also prove \(C_{U_{N}}(\mathrm{SAT})>N^{1+\epsilon-q-o(1)}\) and must beat all semantic block covers.
\end{enumerate}

\paragraph{Equation 1: carrier\_definition.}
\[
H_N(x)=\bigvee_{j=1}^{M} \mathrm{SAT}(G_{N,j}(x)),\quad M=N^q,\quad L=(\log N)^d
\]
\emph{Valid when:} Valid when \(G_{N}\) is a fixed nonuniform generator family and x ranges over N-bit seeds.

\paragraph{Equation 2: target\_substitution\_transfer.}
\[
\mathrm{SAT}\in \mathrm{SIZE}(n^c)\Rightarrow C(H_N)\le N^\rho+N^q(\log N)^{dc}+O(N^qL) = N^{\max{\rho,q}+o(1)}
\]
\emph{Valid when:} Applies when each generated query has length L and the SAT circuit is copied for every query position.

\paragraph{Equation 3: restricted\_solver\_falsifier\_bound.}
\[
C_{U_N}(\mathrm{SAT})\le T_N\Rightarrow C(H_N)\le N^\rho+N^qT_N+O(N^qL)
\]
\emph{Valid when:} Valid when one length-L circuit decides SAT correctly on the generated union range \(U_{N}\).

\paragraph{Equation 4: necessary\_range\_hardness.}
\[
C(H_N)> N^{1+\epsilon}\ \wedge\ \rho< 1+\epsilon\Rightarrow C_{U_N}(\mathrm{SAT})> N^{1+\epsilon-q-o(1)}
\]
\emph{Valid when:} Requires the generated-OR lower bound and the copied restricted solver upper bound to be compared at the same N.

\paragraph{Next analysis.} Instantiate \(G_{N}\) with a high-entropy family of polylog-size 3-CNF formulas and compute the best promise solver or block-cover upper bound; vary \(\tau\) to locate the first threshold-respecting candidate.

\paragraph{Falsifiers.}
\begin{itemize}
\item falsifier: \(q\ge 1+\epsilon\) or \(\rho\ge 1+\epsilon\) collapses the target threshold before SAT hardness transfers.
\item falsifier: a promise solver for SAT on \(U_{N}\) of size at most \(N^{1+\epsilon-q-o(1)}\) gives a direct small circuit for \(H_{N}\).
\item falsifier: a semantic block partition with total cost below \(N^{1+\epsilon-o(1)}\) kills the generated-OR lower bound.
\item falsifier: a generator whose formulas lie in a syntactically rich but semantically easy SAT family defeats entropy-only lower-bound arguments.
\item falsifier: advice depending on N but not on x is allowed in P/poly and must be charged in every valid source-model audit.
\end{itemize}

\subsection*{cycle-0023: Hardness-magnification bridge realization: generated-OR range-balance gate}

\paragraph{Execution.} research; azure\_responses\_api; attempt primary; 122.213 seconds; 3 lemmas, 4 equations, 6 derivation steps, 5 falsifiers.

\paragraph{Learning.} Cycle 23 adds an exact range-label balance gate. Even with q,rho in the valid target window, a generated-OR carrier is source-plausible only if both satisfiable and unsatisfiable parts of its actual generated range exceed \(N^{1+\epsilon-q-o(1)}\). This is a concrete falsifier derived from equality-DNF membership circuits, independent of vague barrier folklore.

\paragraph{Pressure.} No explicit generator is known whose \(H_{N}\) beats unrestricted \(N^{1+\epsilon}\) circuits while also defeating range-label imbalance, promise solvers on \(U_{N}\), and semantic block-cover compression.

\paragraph{Derivation steps.}
\begin{enumerate}
\item Fix N and let K\textasciicircum{}+=|U\_N\textasciicircum{}+|, K\textasciicircum{}-=|U\_N\textasciicircum{}-|.
\item If K\textasciicircum{}+<=K\textasciicircum{}-, build a membership circuit Pos\_N(y)=OR\_\{z in U\_N\textasciicircum{}+\}[y=z] of size O(K\textasciicircum{}+ L). Then \(H_{N}(x)=OR_{j}\) Pos\_N(G\_\{N,j\}(x)) on every x.
\item If K\textasciicircum{}-<K\textasciicircum{}+, build Neg\_N(y)=OR\_\{z in U\_N\textasciicircum{}-\}[y=z] of size O(K\textasciicircum{}- L). Then \(H_{N}(x)=NOT\) \(AND_{j}\) Neg\_N(G\_\{N,j\}(x)), since every generated query lies in \(U_{N}\).
\item Composing either membership circuit with the charged all-query generator gives C(H\_N)<=C(G\_N)+O(M L min\{K\textasciicircum{}+,K\textasciicircum{}-\})+O(M).
\item With \(M=N^{q+o(1)}\), \(L=N^{o(1)}\), and \(C(G_{N})N^{1+\epsilon}\) lower bound with \(rhoN^{1+\epsilon-q-o(1)}\).
\item The target substitution calculation is unchanged: if SAT has size \(n^{c}\) circuits, composing those circuits with \(G_{N}\) gives \(C(H_{N})\le N^{max{\rho,q}+o(1)}\), contradicting the source lower bound when \(max{\rho,q}<1+\epsilon\).
\end{enumerate}

\paragraph{Equation 1: carrier\_definition.}
\[
H_N(x)=\bigvee_{j=1}^{M} SAT(G_{N,j}(x)),\qquad M=N^{q+o(1)},\quad L=(\log N)^d.
\]
\emph{Valid when:} Valid when the generated queries are fixed nonuniformly for length N and their encoding length is polylogarithmic in the carrier input length.

\paragraph{Equation 2: range\_balance\_upper\_bound.}
\[
C(H_N)\le C(G_N)+O\!\left(M L\min{|U_N^+|,|U_N^-|}\right)+O(M).
\]
\emph{Valid when:} Valid when membership in the smaller labeled generated range is implemented by equality-DNF circuits and all generator cost is charged once.

\paragraph{Equation 3: necessary\_balance\_threshold.}
\[
C(H_N)> N^{1+\epsilon}\ \wedge\ C(G_N)\le N^{\rho+o(1)}\ \wedge\ \rho< 1+\epsilon\Rightarrow \min{|U_N^+|,|U_N^-|}> N^{1+\epsilon-q-o(1)}.
\]
\emph{Valid when:} Applies after the range-balance upper bound is compared with the claimed barely-superlinear source lower-bound threshold.

\paragraph{Equation 4: target\_substitution\_bound.}
\[
SAT\in SIZE(n^c)\Rightarrow C(H_N)\le N^{\rho+o(1)}+N^{q+o(1)}L^c=N^{\max{\rho,q}+o(1)}.
\]
\emph{Valid when:} Requires the target-preserving transfer where each generated SAT query is answered by the same nonuniform SAT circuit without extra advice leakage.

\paragraph{Next analysis.} Instantiate a template generator for polylog-length 3CNFs and estimate \(\theta=min\) label-support exponent; then compare it with 1+epsilon-q and with the best promise-SAT solver size on the generated template range.

\paragraph{Falsifiers.}
\begin{itemize}
\item range-label imbalance falsifier: either \(|U_{N}^{+}|\) or \(|U_{N}^{-}|\) is at most \(N^{1+\epsilon-q-o(1)}\), yielding a valid subthreshold circuit for \(H_{N}\).
\item promise-solver falsifier: SAT restricted to \(U_{N}\) has nonuniform circuits of size at most \(N^{1+\epsilon-q-o(1)}\).
\item semantic block-cover falsifier: a partition of query positions has total charged OR-SAT block complexity below \(N^{1+\epsilon-o(1)}\).
\item threshold-collapse falsifier: \(q\ge 1+\epsilon\) or \(\rho\ge 1+\epsilon\) destroys the barely-superlinear magnification window.
\item nonuniform leakage falsifier: uncharged advice encoding SAT labels on the generated range makes the source lower bound ill-posed.
\end{itemize}

\subsection*{cycle-0024: Hardness-Magnification Bridge Realization: hitting-cover gate for generated-OR carriers}

\paragraph{Execution.} research; azure\_responses\_api; attempt primary; 156.244 seconds; 2 lemmas, 4 equations, 6 derivation steps, 7 falsifiers.

\paragraph{Learning.} Cycle 24 sharpens the generated-OR route from label balance to exact hitting-cover hardness. The valid target bridge is unchanged, but the source lower bound must now defeat nonuniform covers of positive fibers and negative product fibers of total size \(N^{1+\epsilon-o(1)}\). This gives a new concrete falsifier and shows that high semantic range entropy alone cannot justify the carrier.

\paragraph{Pressure.} No current argument lower-bounds the exact hitting-cover parameter \(\kappa\) for an explicit target-clean generator; raw semantic range size, variable entropy, and label balance all allow small-cover falsifiers.

\paragraph{Derivation steps.}
\begin{enumerate}
\item Defined positive hitting covers and negative product covers on the actual generated SAT ranges \(R_{j}\).
\item Built a circuit by first computing the shared generator \(G_{N}\), then testing each output against the selected nonuniform label sets with equality ORs of cost O(L|K\_j|).
\item For a positive cover, correctness follows because negative inputs generate only unsatisfiable formulas, while every positive input is hit by some selected satisfiable formula.
\item For a negative cover, the AND of unsatisfiable-membership tests computes the complement of \(H_{N}^{G}\) exactly, because every positive input has at least one satisfiable query outside every \(K_{j}^{-}\).
\item Combined the cover circuit with the existing SAT-substitution upper bound to isolate two independent gates: target transfer \(max{\rho,q}N^{1+\epsilon-o(1)}\).
\item Tried to derive \(\kappa\) lower bounds from cycle-23 label balance, but a planted-template plus decoy construction separates label balance from hitting-cover size.
\end{enumerate}

\paragraph{Equation 1: generated\_or\_definition.}
\[
H_N^G(x)=\bigvee_{j=1}^{M} SAT(G_{N,j}(x)),\quad M=N^q,\quad L=(\log N)^d
\]
\emph{Valid when:} Valid when \(G_{N}\) is a shared multi-output generator and each output string encodes a SAT instance of length L.

\paragraph{Equation 2: target\_substitution\_upper\_bound.}
\[
SAT\in SIZE(n^c)\Rightarrow C(H_N^G)\le O(N^\rho+M L^c+M)=N^{\max{\rho,q}+o(1)}
\]
\emph{Valid when:} Applies after the reduction preserves nonuniform size accounting and the query length is polylogarithmic in N.

\paragraph{Equation 3: hitting\_cover\_upper\_bound.}
\[
C(H_N^G)\le O(N^\rho+M+L\cdot \min{\kappa^+_N,\kappa^-_N})
\]
\emph{Valid when:} Valid when the selected covers are exact on positive or negative fibers and membership is tested by equality circuits.

\paragraph{Equation 4: necessary\_cover\_threshold.}
\[
C(H_N^G)> N^{1+\epsilon}\Rightarrow \min{\kappa^+_N,\kappa^-_N}> \Omega((N^{1+\epsilon}-N^\rho-M)/L)=N^{1+\epsilon-o(1)}
\]
\emph{Valid when:} Requires \(\rho\) and q to stay below 1+epsilon so generator and OR costs do not already meet the threshold.

\paragraph{Next analysis.} Formalize the minimum cover parameter as a bipartite set-cover problem between input fibers and generated SAT labels; test whether known disperser/extractor-style generators can force \(kappa_{N}>N^{1+\epsilon-o(1)}\) without reintroducing nonuniform leakage.

\paragraph{Falsifiers.}
\begin{itemize}
\item hitting-cover collapse: \(kappa_{N}\) \(\le \) \(N^{1+\epsilon-o(1)}/L\) gives a direct subthreshold circuit for \(H_{N}^{G}\)
\item planted satisfiable template: every positive input emits one of few canonical SAT formulas despite large decoy ranges
\item negative product cover: all negative inputs use a small per-position family of unsatisfiable formulas
\item semantic block-cover compression from earlier cycles can dominate singleton cover counts
\item nonuniform leakage through length-dependent equality tables for selected generated labels
\item model mismatch if \(G_{N}\) must be replicated M times rather than charged once as a shared multi-output circuit
\item target threshold collapse if max\{rho,q\} \(\ge \) 1+epsilon, independent of cover size
\end{itemize}

\subsection*{cycle-0025: Generated-OR magnification: selector-range leakage gate}

\paragraph{Execution.} research; azure\_responses\_api; attempt primary; 108.203 seconds; 3 lemmas, 3 equations, 5 derivation steps, 5 falsifiers.

\paragraph{Learning.} Cycle 25 adds a shared-selector leakage gate. A generated-OR carrier can be easy via one compact selector and a promise SAT solver on only the selected image, even when additive semantic or hitting covers look expensive. Thus the valid source obligation is not merely high range entropy or label balance; it is lower bounding every selector-plus-selected-range-solver pair above \(N^{1+\epsilon-o(1)}\) while preserving \(q,\rho<1+\epsilon\) for the target bridge.

\paragraph{Pressure.} Need a direct-sum or shared-computation lower bound converting any small selector/promise-solver circuit into a subthreshold cycle-24 cover. No such theorem is available, and it is plausibly false for structured selectors.

\paragraph{Derivation steps.}
\begin{enumerate}
\item Start with the cycle-24 condition that positive fibers and negative product fibers have no subthreshold nonuniform hitting cover.
\item Try to conclude \(C(H_{N})>N^{1+\epsilon}\) by arguing every \(H_{N}\) circuit induces such a cover.
\item Observe that a circuit can first compute a shared selector Sel\_N(x) and query only G\_N(x,Sel\_N(x)).
\item The induced cover over decoded selector values may have large additive cost because shared gates are replicated when charged per fiber.
\item Therefore additive hitting-cover hardness does not exclude selector-plus-promise-solver circuits.
\end{enumerate}

\paragraph{Equation 1: selector\_upper\_bound.}
\[
C(H_N) \le C(G_N)+C(Sel_N)+C(P_N)+N^{o(1)}
\]
\emph{Valid when:} Valid when \(Sel_{N}\) is a genuine positive selector and \(\mathrm{P}_{N}\) is correct on the complete selected image \(V_{S}\).

\paragraph{Equation 2: source\_obstruction.}
\[
C(H_N)> N^{1+\epsilon} \Rightarrow C(Sel_N)+C(P_N) \ge N^{1+\epsilon-o(1)}
\]
\emph{Valid when:} Applies after generator overhead is below threshold and the asserted lower bound is in the same nonuniform circuit model.

\paragraph{Equation 3: target\_transfer\_unchanged.}
\[
SAT\in SIZE(n^c) \Rightarrow C(H_N) \le N^{\rho+o(1)}+N^q(\log N)^{dc}=N^{\max{\rho,q}+o(1)}
\]
\emph{Valid when:} Requires fixed target exponent c, polylogarithmic query length, and generator substitution preserving the nonuniform size accounting.

\paragraph{Next analysis.} Instantiate \(Sel_{N}\) for simple generator templates, including lexicographic first-satisfying-query generators and hash-bucketed query generators, and compute whether selected-image lookup already drops below \(N^{1+\epsilon}\).

\paragraph{Falsifiers.}
\begin{itemize}
\item Small valid selector plus small promise solver on \(V_{S}\) computes \(H_{N}\) below threshold.
\item Selected-image equality-DNF lookup with \(|V_{S}|<N^{1+\epsilon-o(1)}\) falsifies the claimed source lower bound.
\item Shared computation collapse: additive hitting-cover lower bounds may miss one compact selector circuit.
\item Nonuniform leakage through hardwired SAT labels on the selected image remains a route-level falsifier.
\item If \(q\ge 1+\epsilon\) or \(\rho\ge 1+\epsilon\), the target transfer threshold collapses even before selector leakage is considered.
\end{itemize}

\subsection*{cycle-0026: Hardness-Magnification Bridge Realization}

\paragraph{Execution.} research; azure\_responses\_api; attempt primary; 136.529 seconds; 2 lemmas, 5 equations, 6 derivation steps, 6 falsifiers.

\paragraph{Learning.} Cycle 26 adds an approximate-selector gate. Exact witness selection was too weak as a source audit: even a selector wrong on \(N^{\delta}\) inputs collapses \(H_{N}\) below \(N^{1+\epsilon}\) when \(deltaN^{1+\epsilon-o(1)}\), while the clean target bridge \(q,\rho<1+\epsilon\) remains intact.

\paragraph{Pressure.} No concrete \(E_{N}\) is known for which \(Lambda_{N}>N^{1+\epsilon-o(1)}\) can be proved in unrestricted circuits; proving it appears to require non-natural, selector-resistant lower-bound machinery rather than entropy counting.

\paragraph{Derivation steps.}
\begin{enumerate}
\item Fix \(sigma_{N}\) and \(\mathrm{P}_{N}\), and compute A(x)=P\_N(E\_N(x,sigma\_N(x))).
\item On inputs outside \(B_{sigma,\mathrm{P}}\), A(x)=H\_N(x) by definition of the bad set.
\item For each bad input b, add an equality gate EQ\_b(x); XOR A with OR\_\{b in B\} EQ\_b(x) to flip exactly the wrong outputs.
\item The correction costs O(N|B|), while selector, extractor, and promise solver costs add linearly up to standard fan-in overhead.
\item Thus any sparse-error approximate selector with small selected-range SAT solver gives a subthreshold circuit for \(H_{N}\) and refutes the proposed source lower bound.
\item Combining this with the clean target bridge shows the remaining route obligation is a minimax lower bound against all approximate selector-plus-promise-solver pairs.
\end{enumerate}

\paragraph{Equation 1: generated\_or\_carrier.}
\[
H_N(x)=\bigvee_{j< t(N)} SAT(E_N(x,j)),\qquad t(N)=N^q,\quad L(N)=(\log N)^d.
\]
\emph{Valid when:} Valid when the source carrier is the explicit generated-OR family with named query count, query length, and coordinate extractor.

\paragraph{Equation 2: clean\_target\_substitution.}
\[
SAT_L\in SIZE(L^c)\ \Longrightarrow\ C(H_N)\le N^{\rho+o(1)}+N^q L^c=N^{\max{\rho,q}+o(1)}.
\]
\emph{Valid when:} Applies after the reduction preserves nonuniform size accounting and \(L(N)^{c}\) is subpolynomial for every fixed target exponent c.

\paragraph{Equation 3: approx\_selector\_upper\_bound.}
\[
C(H_N)\le C(\sigma_N)+C(E_N)+C(P_N)+O(L+\log t)+O(N\,|B_{\sigma,P}|).
\]
\emph{Valid when:} Valid when \(\mathrm{P}_{N}\) agrees with SAT on the selected image and the remaining disagreements are corrected by nonuniform equality tests.

\paragraph{Equation 4: bad\_set\_threshold.}
\[
|B_{\sigma,P}|\le N^{\delta}\ \Longrightarrow\ O(N|B_{\sigma,P}|)\le N^{1+\delta+o(1)}.
\]
\emph{Valid when:} Requires the correction to be by explicit exceptional-input advice in unrestricted circuits over N-bit carrier inputs.

\paragraph{Equation 5: source\_obligation\_minimax.}
\[
\Lambda_N:=\min_{\sigma,P}{C(\sigma_N)+C(E_N)+C(P_N)+N|B_{\sigma,P}|}> N^{1+\epsilon-o(1)}.
\]
\emph{Valid when:} Necessary when an \(N^{1+\epsilon}\) lower bound for \(H_{N}\) is asserted in the same unrestricted circuit model.

\paragraph{Next analysis.} Instantiate the thought experiment with a block generator \(E_{N}\) whose coordinates are grouped into planted easy and hard ranges; compute the smallest \(\delta\) forced by coverage of positive inputs and test whether selected-range promise SAT stays subthreshold.

\paragraph{Falsifiers.}
\begin{itemize}
\item approximate-selector falsifier: a selector correct except on \(N^{\epsilon-\eta}\) carrier inputs plus a subthreshold selected-range promise solver gives \(H_{N}\) circuits below \(N^{1+\epsilon}\)
\item exact-selector leakage from cycle 25 remains a special case with empty bad set
\item threshold collapse if \(max{q,\rho}\ge 1+\epsilon\) in the clean target substitution
\item lookup falsifier when query length and variable entropy make SAT on the generated range tabulatable below the source threshold
\item nonuniform leakage through a compact selected image \(U_{sigma}\) with easy promise SAT despite high full-range entropy
\item natural-proof risk: a large constructive property proving selector-cover hardness for many generators would conflict with standard P/poly pseudorandom-function assumptions
\end{itemize}

\subsection*{cycle-0027: Hardness-Magnification Bridge Realization: approximate list-selector leakage gate}

\paragraph{Execution.} research; azure\_responses\_api; attempt primary; 140.568 seconds; 2 lemmas, 4 equations, 4 derivation steps, 5 falsifiers.

\paragraph{Learning.} Cycle 27 upgrades the source audit from approximate single selectors to approximate list selectors. The valid target-preserving bridge is unchanged, but a short list plus a promise solver can decide \(H_{N}\) without computing a canonical witness. Therefore any contribution-making lower-bound attempt must lower-bound \(Lambda_{N}^{list}\); otherwise the route has a concrete nonuniform-leakage falsifier.

\paragraph{Pressure.} Witness canonicalization is semantic; it is not available from the selector circuit alone and can be bypassed by list-wise OR evaluation.

\paragraph{Derivation steps.}
\begin{enumerate}
\item Start with an approximate t-list selector \(A_{N}\) and attempt to pick the lexicographically first selected satisfiable query.
\item To choose that index the circuit must know SAT labels on all selected formulas, so the reduction imports a promise solver \(S_{N}\) on \(U_{A}\).
\item Once \(S_{N}\) is imported, the decision circuit ORs the t answers and no longer needs a single witness index.
\item The cost is \(C(A_{N})+t_{N}\) \(C(S_{N})+exception\) correction, which can be subthreshold even if every single selector is hard.
\end{enumerate}

\paragraph{Equation 1: target\_substitution\_bridge.}
\[
SAT_{L_N}\in SIZE(L_N^c) \Rightarrow C(H_N)\le N^{\max{\rho,q}+o(1)}
\]
\emph{Valid when:} Valid when the source carrier uses the named batched generator interface and c,d are fixed constants.

\paragraph{Equation 2: list\_selector\_leakage\_bound.}
\[
C(H_N)\le O(C(A_N)+C(G_N^{batch})+t_N(C(S_N)+L_N^{O(1)})+N|E_A|)
\]
\emph{Valid when:} Applies when the selected formulas are generated inside the same nonuniform circuit and all selector errors are explicitly corrected.

\paragraph{Equation 3: exponent\_form\_batched.}
\[
C(H_N)\le N^{\max{a,\rho,\tau+b,1+\delta}+o(1)}
\]
\emph{Valid when:} Valid when \(G_{N}\) has batched cost \(N^\rho\) independent of the chosen list length exponent \(\tau\).

\paragraph{Equation 4: source\_obligation\_minimax.}
\[
\Lambda_N^{list}(\tau,\delta):=\min_{A,S,E}\big(C(A)+C(G_N^{batch})+t_NC(S)+N|E|\big)> N^{1+\epsilon-o(1)}
\]
\emph{Valid when:} Requires the lower-bound proof to quantify over all nonuniform selectors, solvers, and exception sets in the source model.

\paragraph{Next analysis.} Instantiate \(G_{N}\) with a high-entropy expander/hash-generated SAT family and compute whether any tau=O(epsilon/2) list selector has a selected image admitting a subthreshold promise solver.

\paragraph{Falsifiers.}
\begin{itemize}
\item Approximate list selector falsifier: \(max{a,\rho,\tau+b,1+\delta}<1+\epsilon\) yields a small circuit for \(H_{N}\).
\item Target threshold falsifier: \(q\ge 1+\epsilon\) or \(\rho\ge 1+\epsilon\) destroys the magnification implication even if the source lower bound is true.
\item Generator-interface falsifier: if selected generation costs \(N^{\rho+\tau}\) rather than \(N^\rho\), all exponent inequalities must be recomputed before calling the bridge valid.
\item Promise-range falsifier: SAT may be easy on the selected image \(U_{A}\) although hard on the full generated range.
\item Exception-leakage falsifier: \(|E_{A}|\le N^\delta\) with \(\delta<\epsilon\) is removable by equality advice at cost below the claimed lower-bound threshold.
\end{itemize}

\subsection*{cycle-0028: Hardness-Magnification Bridge Realization}

\paragraph{Execution.} research; azure\_responses\_api; attempt primary; 129.075 seconds; 3 lemmas, 4 equations, 7 derivation steps, 5 falsifiers.

\paragraph{Learning.} Cycle 28 identifies an additional leakage mode: many imperfect list selectors can collectively cover the positive fiber. The valid source audit is the ensemble cover complexity \(Xi_{N}^{ens}\); single-selector minimax lower bounds are insufficient for this generated-OR magnification route.

\paragraph{Pressure.} No lower-bound method is known for \(Xi_{N}^{ens}\) against fully nonuniform selected-image promise solvers; ensemble covers can exploit semantic clusters invisible to entropy counts.

\paragraph{Derivation steps.}
\begin{enumerate}
\item Fix \(H_{N}(x)=OR_{j}\) SAT(G\_N(x,j)) with \(M=N^{q}\) and L=(log N)\textasciicircum{}d.
\item Given an ensemble, build C\_i(x)=OR\_\{j in A\_i(x)\} B\_i(G\_N(x,j)) for each selector-solver pair.
\item If \(H_{N}(x)=0\), every generated formula is unsatisfiable, so all \(C_{i}(x)=0\) by selected-image correctness.
\item If \(H_{N}(x)=1\) and x is covered, some \(C_{i}(x)=1\) because its list contains a satisfiable generated formula.
\item Patch the E uncovered positive strings by an equality-DNF of size O(NE).
\item Share the generator access circuitry once and sum selector, solver-copy, and patch costs to obtain the ensemble upper bound.
\item Under \(\mathrm{SAT}\in \mathrm{SIZE}(n^{c})\), choose one selector listing all \(N^{q}\) indices and recover \(C(H_{N})\le N^{max{\rho,q}+o(1)}\).
\end{enumerate}

\paragraph{Equation 1: generated\_or\_definition.}
\[
H_N(x)=\bigvee_{j\in[M]} SAT(G_N(x,j)),\quad M=N^q,\quad L=(\log N)^d.
\]
\emph{Valid when:} Valid when the source generator is explicitly named and the \(H_{N}\) input length is N.

\paragraph{Equation 2: ensemble\_cover\_upper\_bound.}
\[
C(H_N)\le O\left(N^\rho+\sum_{i=1}^R(S_i+K_i(T_i+N^{o(1)}))+NE\right).
\]
\emph{Valid when:} Applies when each selected-image solver is correct and residual positives are patched by equality circuits.

\paragraph{Equation 3: exponent\_form.}
\[
R\le N^\theta,\ S_i\le N^\alpha,\ K_i\le N^\lambda,\ T_i\le N^\tau,\ E\le N^\delta\Rightarrow C(H_N)\le N^{\max{\rho,\theta+\alpha,\theta+\lambda+\tau,1+\delta}+o(1)}.
\]
\emph{Valid when:} Requires uniform exponent bounds for every ensemble member and polylogarithmic generator overhead.

\paragraph{Equation 4: target\_all\_query\_special\_case.}
\[
SAT\in SIZE(n^c)\Rightarrow C(H_N)\le N^{\max{\rho,q}+o(1)}.
\]
\emph{Valid when:} Valid when the selector lists all \(N^{q}\) queries and L\textasciicircum{}c=N\textasciicircum{}\{o(1)\} for fixed c.

\paragraph{Next analysis.} Estimate \(Xi_{N}^{ens}\) for a random high-entropy generator by a union-bound model over ensembles, then isolate which part of the argument can survive explicitness and avoid natural-proof leakage.

\paragraph{Falsifiers.}
\begin{itemize}
\item falsifier: an ensemble of \(N^\theta\) list selectors with \(\theta+\lambda+\tau<1+\epsilon\) and residual \(E<N^{\epsilon-o(1)}\) collapses the source lower bound.
\item falsifier: selected-image SAT is promise-easy even though unrestricted SAT on L-bit formulas is hard.
\item falsifier: \(q\ge 1+\epsilon\) or \(\rho\ge 1+\epsilon\) causes quantitative threshold collapse in the target bridge.
\item falsifier: a natural dense property proving \(Xi_{N}^{ens}\) large would collide with standard natural-proofs concerns unless it is generator-specific and non-natural.
\item falsifier: hidden advice in selectors depending on x beyond the input circuit model invalidates target preservation.
\end{itemize}

\subsection*{cycle-0029: Hardness-Magnification Bridge Realization: one-sided separating-ensemble leakage gate}

\paragraph{Execution.} research; azure\_responses\_api; attempt primary; 295.993 seconds; 3 lemmas, 4 equations, 6 derivation steps, 4 falsifiers.

\paragraph{Learning.} Cycle 29 sharpens the generated-OR source audit from ensemble list selectors to one-sided selected-image separators. The valid target bridge still needs only \(q,\rho<1+\epsilon\), but the source lower bound must now beat \(Psi_{N}^{sep}\); small local labelers are a concrete falsifier even when they do not solve SAT globally.

\paragraph{Pressure.} No available lower-bound method rules out arbitrary selected-image separators \(B_{t}\). These labelers need not be SAT algorithms, so the target SAT hardness hypothesis does not exclude them.

\paragraph{Derivation steps.}
\begin{enumerate}
\item Fix an ensemble of R selectors \(A_{t}\), each producing s indices into the \(N^{q}\) generated SAT queries.
\item For each pair (t,u), build a circuit that computes j=A\_t(x,u), then y=G\_N(x,j), then the selected-image label B\_t(y).
\item OR all R*s outputs. By the separation promise, this computes \(H_{N}\) on all x outside \(E_{N}\).
\item Hardwire the exceptional inputs by equality gates and override the answer on \(E_{N}\), with cost \(N|E_{N}|\) up to subpolynomial factors.
\item The resulting circuit has size at most R*s*(N\textasciicircum{}a+N\textasciicircum{}rho+N\textasciicircum{}b) plus exception cost, giving N\textasciicircum{}\{r+ell+max\{a,rho,b\}+o(1)\}+N\textasciicircum{}\{1+delta+o(1)\}.
\item Therefore any claimed \(N^{1+\epsilon}\) lower bound for \(H_{N}\) is valid only if every such separating ensemble exceeds the same exponent threshold.
\end{enumerate}

\paragraph{Equation 1: carrier-definition.}
\[
H_N(x)=\bigvee_{j< N^q} SAT(G_N(x,j)),\qquad L=(\log N)^d.
\]
\emph{Valid when:} Valid when the generated carrier and SAT-query length map are fixed before proving any lower bound.

\paragraph{Equation 2: separator-circuit.}
\[
D_N(x)=\bigvee_{t< R}\bigvee_{u< s} B_t(G_N(x,A_t(x,u))).
\]
\emph{Valid when:} Valid when the selector ensemble separates \(H_{N}\) correctly outside the explicitly hardwired exception set.

\paragraph{Equation 3: separator-upper-bound.}
\[
C(H_N)\le N^{r+\ell+\max{a,\rho,b}+o(1)}+N^{1+\delta+o(1)}.
\]
\emph{Valid when:} Applies after all selector, generator, labeler, OR fanin, and exception-correction costs are charged nonuniformly.

\paragraph{Equation 4: source-audit-threshold.}
\[
\Psi_N^{sep}> N^{1+\epsilon-o(1)}.
\]
\emph{Valid when:} Required when the advertised source lower bound is meant to survive all one-sided separating-ensemble leakage.

\paragraph{Next analysis.} Try to lower-bound or construct \(Psi_{N}^{sep}\) for a random-like generator: estimate whether selected positive and negative images admit \(N^{b}-size\) separators when \(R=N^{r}\) and \(s=N^{ell}\) are below the magnification threshold.

\paragraph{Falsifiers.}
\begin{itemize}
\item A one-sided separating ensemble with \(r+ell+max{a,\rho,b}<1+\epsilon\) and exception exponent \(\delta<\epsilon\) falsifies the source lower-bound claim.
\item A small selected-image labeler \(B_{t}\) that is not a valid SAT solver still leaks enough nonuniform information to collapse \(H_{N}\).
\item If \(q\ge 1+\epsilon\) or \(\rho\ge 1+\epsilon\), the target bridge has threshold collapse even before source leakage is considered.
\item If the length map L=(log N)\textasciicircum{}d does not cover SAT lengths with controlled rounding, the target-preserving consequence is invalid.
\end{itemize}

\subsection*{cycle-0030: Hardness-Magnification Bridge Realization}

\paragraph{Execution.} research; azure\_responses\_api; attempt primary; 154.031 seconds; 3 lemmas, 3 equations, 5 derivation steps, 5 falsifiers.

\paragraph{Learning.} Cycle 30 replaces the one-sided selected-image audit by the weaker and more dangerous relaxed bag-separator audit. The valid target bridge is unchanged, but the source lower bound must defeat predicates that merely hit one positive witness per positive bag and avoid all negative-bag formulas. This exposes a concrete nonuniform-leakage falsifier not captured by promise SAT hardness.

\paragraph{Pressure.} The converse is false or at least unproved: absence of small relaxed formula separators does not imply \(H_{N}\) is hard, because an arbitrary \(H_{N}\) circuit may exploit input-level structure that does not factor through a formula predicate.

\paragraph{Derivation steps.}
\begin{enumerate}
\item Define \(A_{N}\) on \(length-L_{N}\) formulas with negative avoidance and positive hitting over generated bags.
\item Build a circuit for \(H_{N}\) by computing G\_N(x), applying \(A_{N}\) to each of the \(t_{N}\) formulas, and ORing the answers.
\item Size accounting gives C(H\_N)<=N\textasciicircum{}\{rho+o(1)\}+N\textasciicircum{}q(size(A\_N)+poly(L\_N)).
\item Therefore a claimed lower bound \(C(H_{N})>N^{1+\epsilon-o(1)}\) forces \(size(A_{N})>N^{1+\epsilon-q-o(1)}\) whenever \(\rho<1+\epsilon\).
\item Taking \(A_{N}\) to be a genuine SAT circuit recovers the target bridge: \(\mathrm{SAT}\in \mathrm{SIZE}(L^{c})\) would imply \(C(H_{N})\le N^{max{\rho,q}+o(1)}\).
\end{enumerate}

\paragraph{Equation 1: relaxed\_separator\_conditions.}
\[
H_N(x)=0 \Rightarrow \sum_{j\le t_N} A_N(\phi_{N,j}(x))=0,\qquad H_N(x)=1 \Rightarrow \sum_{j\le t_N} A_N(\phi_{N,j}(x))\ge 1.
\]
\emph{Valid when:} Valid when \(H_{N}\) is exactly the OR of SAT labels on the generated bag and \(A_{N}\) is only required to hit each positive bag once.

\paragraph{Equation 2: separator\_upper\_bound.}
\[
C(H_N)\le C(G_N)+O(t_N(\operatorname{size}(A_N)+L_N))=N^{\max{\rho,q+\tau}+o(1)}\quad\text{if }\operatorname{size}(A_N)=N^\tau.
\]
\emph{Valid when:} Applies after the generator outputs explicit formulas and the same separator circuit is reused on all \(t_{N}\) query positions.

\paragraph{Equation 3: necessary\_source\_audit.}
\[
C(H_N)> N^{1+\epsilon-o(1)}\ \Longrightarrow\ \operatorname{SepHit}_N> N^{1+\epsilon-q-o(1)}\quad\text{provided }\rho< 1+\epsilon.
\]
\emph{Valid when:} Requires the generated-OR source lower bound to be measured against unrestricted nonuniform circuits with subthreshold generator overhead.

\paragraph{Next analysis.} Formalize \(SepHit_{N}\) as a bipartite hitting-avoidance problem between positive bags and negative formula occurrences, then test whether random-looking generators satisfy the required lower bound by a counting argument and where explicitness breaks the argument.

\paragraph{Falsifiers.}
\begin{itemize}
\item A planted syntactic marker present in one satisfiable formula of every positive bag and absent from all negative bags gives a small separator falsifier.
\item A statistical artifact distinguishing positive witness formulas from negative-bag formulas yields \(C(H_{N})\le N^{q}\) poly(L\_N) despite possible SAT hardness on \(U_{N}\).
\item If \(q\ge 1+\epsilon\), even superconstant separator lower bounds cannot preserve the barely-superlinear threshold.
\item If \(\rho\ge 1+\epsilon\), the generator alone violates the valid magnification window.
\item If the source proof only lower-bounds promise SAT on \(U_{N}\), it misses relaxed separators and leaves nonuniform leakage open.
\end{itemize}

\subsection*{cycle-0031: Hardness-Magnification Bridge Realization: approximate relaxed bag separators}

\paragraph{Execution.} research; azure\_responses\_api; attempt compact\_retry; 989.176 seconds; 1 lemmas, 3 equations, 4 derivation steps, 4 falsifiers.

\paragraph{Learning.} Cycle 31 converts relaxed bag separators into an approximate threshold lemma: E bad bags cost NE gates, so errors below \(N^\epsilon\) are nonuniformly patchable. The valid source obligation is now an approximate distribution-local separator lower bound, not merely promise SAT hardness.

\paragraph{Pressure.} We still lack a lower bound against all distribution-local approximate bag separators; global promise-SAT hardness does not rule out local predicates tuned to the generated bags.

\paragraph{Derivation steps.}
\begin{enumerate}
\item Assume \(H_{N}\) has \(C(H_{N})>N^{1+\epsilon-o(1)}\) and \(q,\rho<1+\epsilon\) on the same valid lengths.
\item Given \(A_{N}\), compute all generated formulas with the shared \(N^\rho\) generator, apply \(A_{N}\) to each of \(N^{q}\) queries, and OR the answers.
\item The resulting circuit agrees with \(H_{N}\) except on E bags; add equality tests for those inputs and hardwire the correct bit, costing O(NE).
\item If \(s=N^{sigma}\) and \(E=N^\eta\) with \(max{\rho,q+sigma,1+\eta}<1+\epsilon\), the patched circuit contradicts the source lower bound.
\end{enumerate}

\paragraph{Equation 1: carrier.}
\[
H_N(x)=\bigvee_{j\le N^q} SAT(G_{N,j}(x)),\quad |G_{N,j}(x)|=(\log N)^d
\]
\emph{Valid when:} Valid when the carrier uses fixed generated query bags and shared generator parameters.

\paragraph{Equation 2: separator\_collapse.}
\[
C(H_N)\le O(N^\rho+N^q s+NE)
\]
\emph{Valid when:} Applies when a size-s bag separator fails on at most E input bags and exceptions are patched nonuniformly.

\paragraph{Equation 3: audit\_threshold.}
\[
s=N^\sigma,E=N^\eta\Rightarrow \max{\rho,q+\sigma,1+\eta}\ge 1+\epsilon-o(1)
\]
\emph{Valid when:} Requires the source lower bound and collapse inequality in the same circuit model and valid-length family.

\paragraph{Next analysis.} Try a collision-balanced generator and estimate whether majority/hash-based local predicates yield \(s=N^{o(1)}\) with \(E<N^\epsilon\); if yes, quarantine this carrier family.

\paragraph{Falsifiers.}
\begin{itemize}
\item Approximate separator falsifier: a local \(A_{N}\) with \(s<N^{1+\epsilon-q-o(1)}\) and \(E<N^{\epsilon-o(1)}\) collapses \(H_{N}\) below threshold.
\item Exception-patching falsifier: even rare bag errors are harmless to circuits when \(NE<N^{1+\epsilon-o(1)}\).
\item Distribution leakage falsifier: \(A_{N}\) only needs to be valid on generated bags, not on all SAT instances.
\item Natural-proof barrier probe: a dense constructive proof excluding all small \(A_{N}\) would look natural unless it is explicitly range-specific or nonconstructive.
\end{itemize}

\subsection*{cycle-0032: Generated-OR magnification: correction-residual separator gate}

\paragraph{Execution.} research; azure\_responses\_api; attempt primary; 124.813 seconds; 3 lemmas, 5 equations, 5 derivation steps, 5 falsifiers.

\paragraph{Learning.} Cycle 32 converts the approximate bag-separator gate into an exact correction-residual gate. The valid target bridge is unchanged, but the source lower bound must prove Phi\_N>1+epsilon-o(1): every small local formula predicate must leave a residual \(H_{N}\) xor \(S_{A}\) of high circuit complexity. Sparse error patching is only one falsifier; structured residual leakage is the stronger obstruction.

\paragraph{Pressure.} No known reduction converts a small residual correction \(R_{A}\) into a small SAT solver or separator on the generated formula range. This is the current formal gap.

\paragraph{Derivation steps.}
\begin{enumerate}
\item Start with the previous approximate relaxed bag-separator lemma: if A is correct on all but E source bags, then sparse equality patching gives C(H\_N)<=N\textasciicircum{}\{rho+o(1)\}+N\textasciicircum{}q C(A)+O(NE).
\item Replace the support-size patch by the exact identity \(H_{N}=S_{A}\) xor \(R_{A}\), where \(R_{A}=H_{N}\) xor \(S_{A}\). This yields the sharper upper bound in L32.1 without assuming the error set is sparse.
\item Try to derive \(C(R_{A})>N^{1+\epsilon}\) from hardness of SAT on the generated range \(U_{N}={G_{N}(x,j)}\).
\item The derivation fails because \(R_{A}\) is a function of whole source bags x, not a predicate on individual formulas. It can encode global correlations among generated queries and may be small even if every range-restricted SAT predicate is hard.
\item Thus promise SAT hardness, semantic range entropy, and label balance do not imply the needed source lower bound unless they also rule out low-complexity residual corrections.
\end{enumerate}

\paragraph{Equation 1: carrier.}
\[
H_N(x)=\bigvee_{j=1}^{N^q} SAT(G_N(x,j))
\]
\emph{Valid when:} Valid when the generated bag interface and SAT encoding are fixed for input length N.

\paragraph{Equation 2: separator\_residual.}
\[
S_A(x)=\bigvee_{j=1}^{N^q} A(G_N(x,j)),\qquad R_A=H_N\oplus S_A
\]
\emph{Valid when:} Valid when A is any nonuniform predicate on one generated L-bit formula.

\paragraph{Equation 3: correction\_upper\_bound.}
\[
C(H_N)\le N^{\rho+o(1)}+N^q C(A)+C(R_A)+N^{q+o(1)}
\]
\emph{Valid when:} Applies after composing A with every generated query and xoring the residual correction circuit.

\paragraph{Equation 4: necessary\_gate.}
\[
C(H_N)> N^{1+\epsilon-o(1)}\Rightarrow \forall A:\max{\rho,q+\log_N C(A),\log_N C(R_A)}\ge 1+\epsilon-o(1)
\]
\emph{Valid when:} Requires the source lower bound to be against the same nonuniform circuit model used by the correction construction.

\paragraph{Equation 5: sparse\_error\_patch.}
\[
|\operatorname{supp}(R_A)|\le N^\delta\Rightarrow C(R_A)\le N^{1+\delta+o(1)}
\]
\emph{Valid when:} Valid when erroneous bags are patched by equality tests on N-bit source inputs.

\paragraph{Next analysis.} Instantiate \(Phi_{N}\) for a tagged generated-OR carrier and compute whether simple tag predicates produce subthreshold residual circuits; then test an untagged high-entropy generator for the same residual leakage.

\paragraph{Falsifiers.}
\begin{itemize}
\item correction-residual falsifier: exists A with \(q+log_{N}\) \(C(A)<1+\epsilon\) and C(H\_N xor S\_A)<N\textasciicircum{}\{1+epsilon-o(1)\}.
\item structured-error falsifier: the bad-bag set is large but has a simple residual circuit, so counting errors would falsely suggest hardness.
\item threshold-collapse falsifier: \(q\ge 1+\epsilon\) or \(\rho\ge 1+\epsilon\) makes the target transfer too expensive.
\item nonuniform-leakage falsifier: residual advice about source bags repairs an imperfect local separator without yielding any SAT solver.
\item natural-proof obstruction: a large constructive property separating \(H_{N}\) from \(N^{1+\epsilon}\) circuits would conflict with standard pseudorandom-function barriers unless the proof is carrier-specific and non-large.
\end{itemize}

\subsection*{cycle-0033: Hardness-Magnification Bridge Realization}

\paragraph{Execution.} research; azure\_responses\_api; attempt primary; 143.841 seconds; 3 lemmas, 6 equations, 6 derivation steps, 5 falsifiers.

\paragraph{Learning.} Cycle 33 strengthens the generated-OR source obligation from single residual correction to feature-aggregator residual hardness. This is a contribution-making repair: it preserves the exact target bridge but exposes a broader nonuniform-leakage falsifier class that prior selector and separator audits did not cover.

\paragraph{Pressure.} The missing proof obligation is a lower bound on C(H\_N xor \(S_{A,B}\)) for all subthreshold feature families and aggregators. Promise SAT hardness on the generated range does not imply this residual lower bound, because B may exploit many weak local features jointly.

\paragraph{Derivation steps.}
\begin{enumerate}
\item Start with any feature family A on L-bit formulas and any aggregator B over the \(M_{N}\) by t feature matrix.
\item Compose \(G_{N}\) with all \(A_{i}\) on all \(M_{N}\) positions; by circuit composition this costs N\textasciicircum{}\{rho+o(1)\}+N\textasciicircum{}\{q+alpha+o(1)\}.
\item Feed the resulting feature matrix to B, obtaining \(S_{A,B}\) with total size at most N\textasciicircum{}\{rho+o(1)\}+N\textasciicircum{}\{q+alpha+o(1)\}+N\textasciicircum{}\{beta+o(1)\}.
\item If a residual circuit \(R_{A,B}=H_{N}\) xor \(S_{A,B}\) has size \(N^{\gamma+o(1)}\), combine \(S_{A,B}\), \(R_{A,B}\), and one xor gate to compute \(H_{N}\).
\item Thus any feature-aggregator-residual certificate with all exponents below 1+epsilon contradicts the claimed \(N^{1+\epsilon-o(1)}\) source lower bound.
\item For target transfer, choose \(t=1\), \(A_{1}\) equal to the hypothesized \(\mathrm{SAT}_{L}\) circuit of size L\textasciicircum{}c=N\textasciicircum{}\{o(1)\}, and \(B=OR_{M_{N}}\); the residual is zero and the cost is \(N^{max{\rho,q}+o(1)}\).
\end{enumerate}

\paragraph{Equation 1: carrier\_parameters.}
\[
M_N=N^q, L=(log N)^d, C(G_N)=N^{rho+o(1)}
\]
\emph{Valid when:} Valid when the generated-bag carrier fixes polynomially many polylogarithmic SAT queries.

\paragraph{Equation 2: feature\_surrogate.}
\[
S_{A,B}(x)=B((A_i(G_{N,j}(x)))_{j\le M_N,i\le t})
\]
\emph{Valid when:} Valid when all local features are evaluated on generated formulas before aggregation.

\paragraph{Equation 3: surrogate\_upper\_bound.}
\[
C(S_{A,B}) \le N^{rho+o(1)} + N^{q+alpha+o(1)} + N^{beta+o(1)}
\]
\emph{Valid when:} Valid when \(G_{N}\), all local features, and B are circuits in the same SIZE model.

\paragraph{Equation 4: residual\_reconstruction.}
\[
C(H_N) \le O(C(S_{A,B}) + C(R_{A,B}) + N), where R_{A,B}=H_N xor S_{A,B}
\]
\emph{Valid when:} Valid when the residual circuit is allowed to correct the surrogate on all inputs.

\paragraph{Equation 5: threshold\_contradiction.}
\[
max{rho,q+alpha,beta,gamma}< 1+epsilon and C(R_{A,B})\le N^{gamma+o(1)} imply C(H_N)< N^{1+epsilon-o(1)}
\]
\emph{Valid when:} Applies when every charged component stays below the claimed source lower-bound exponent.

\paragraph{Equation 6: feature\_residual\_exponent.}
\[
Phi_N^{feat}=inf_{A,B} max{rho,q+alpha(A),beta(B),log_N C(H_N xor S_{A,B})}
\]
\emph{Valid when:} Valid when the infimum ranges over nonuniform feature families and bag aggregators.

\paragraph{Next analysis.} Formalize the restricted case where B is a depth-2 threshold or DNF over \(t=N^\tau\) local features, compute the exact q+alpha+beta leakage threshold, and search for an explicit low-residual construction on expander-walk generated SAT bags.

\paragraph{Falsifiers.}
\begin{itemize}
\item feature-compression falsifier: \(t=N^\tau\) local predicates with \(q+\alpha<1+\epsilon\) and a small aggregator predict \(H_{N}\) up to subthreshold residual complexity.
\item advice-coded aggregator falsifier: B stores labels for large generated-range blocks with \(\beta<1+\epsilon\), collapsing the source lower bound without solving SAT globally.
\item low-dimensional generator falsifier: \(G_{N}\) maps bags into a feature space where OR-SAT is decided by a small lookup or threshold over features.
\item structured residual falsifier: \(H_{N}\) xor \(S_{A,B}\) is easy even when its error support is dense, so sparse-error audits miss the collapse.
\item barrier audit: a dense constructive proof that all feature-aggregator residuals are hard would look natural unless it exploits non-natural SAT semantics or another explicit non-natural certificate.
\end{itemize}

\subsection*{cycle-0034: Hardness-Magnification Bridge Realization: tuple-feature residual gate}

\paragraph{Execution.} research; azure\_responses\_api; attempt primary; 167.842 seconds; 3 lemmas, 4 equations, 7 derivation steps, 7 falsifiers.

\paragraph{Learning.} Cycle 34 adds a quantitative tuple-arity term to the generated-OR audit. The valid target bridge still needs only \(max{\rho,q}<1+\epsilon\), but the source lower bound must rule out every subthreshold tuple-feature transcript with exponent s+tq+alpha and residual exponent \(\lambda\). This yields a finite dense-arity falsifier search and a sharper nonuniform-leakage obstruction, without claiming \(\mathrm{P} \ne \mathrm{NP}\).

\paragraph{Pressure.} The derivation gives an exact upper-bound falsifier class, but it does not prove that arbitrary small circuits for \(H_{N}\) factor through low-arity tuple transcripts. The remaining source obligation is a genuine unrestricted-circuit lower bound or a non-natural theorem ruling out all residual systems and all non-factored circuits.

\paragraph{Derivation steps.}
\begin{enumerate}
\item Fix the generated-OR carrier \(H_{N}\) from the previous cycles with \(M_{N}=N^{q}\) generated SAT instances of length L=(log N)\textasciicircum{}d and batch generator size \(N^{\rho+o(1)}\).
\item Choose \(r=N^{s}\) arity-t predicates \(\mathrm{P}_{i}\). After computing the batch B\_N(x), enumerate all \(M_{N}^{t}\) ordered tuples and OR the predicate outputs to obtain transcript bits z\_i(x).
\item The transcript circuit size is at most \(N^{\rho+o(1)}\) plus r \(M_{N}^{t}\) \(N^{\alpha+o(1)}\), giving exponent max\{rho,s+tq+alpha\}.
\item Apply an aggregator \(A_{N}\) of size \(N^\gamma\) to obtain S\_N(x). If a residual circuit \(E_{N}\) of size \(N^\lambda\) computes H\_N(x) xor S\_N(x), then \(H_{N}\) is computed by \(S_{N}\) xor \(E_{N}\).
\item Summing these components gives C(H\_N)<=N\textasciicircum{}\{rho+o(1)\}+N\textasciicircum{}\{s+tq+alpha+o(1)\}+N\textasciicircum{}gamma+N\textasciicircum{}lambda.
\item Thus any claimed source lower bound \(C(H_{N})>N^{1+\epsilon}\) must imply theta\_t>=1+epsilon-o(1) for every such tuple-feature residual system.
\item Independently, if SAT has \(n^{c}-size\) circuits, substituting them into the \(M_{N}\) polylogarithmic generated queries gives \(C(H_{N})\le N^{max{\rho,q}+o(1)}\), so the target bridge remains valid when \(max{\rho,q}<1+\epsilon\).
\end{enumerate}

\paragraph{Equation 1: tuple\_feature\_upper\_bound.}
\[
C(H_N) \le N^{\rho+o(1)}+N^{s+tq+\alpha+o(1)}+N^\gamma+N^\lambda.
\]
\emph{Valid when:} Valid when the batch generator is shared once, tuple features are densely enumerated over all \(M_{N}^{t}\) tuples, and \(E_{N}\) exactly computes \(H_{N}\) xor \(S_{N}\).

\paragraph{Equation 2: subthreshold\_falsifier\_exponent.}
\[
\theta_t:=\max{\rho,\,s+tq+\alpha,\,\gamma,\,\lambda}< 1+\epsilon \;\Rightarrow\; C(H_N)\le N^{1+\epsilon-\Omega(1)}.
\]
\emph{Valid when:} Valid when all exponent inequalities have a fixed margin and polynomially many additive circuit components are absorbed into the maximum exponent.

\paragraph{Equation 3: target\_substitution\_bridge.}
\[
SAT\in SIZE(n^c) \;\Rightarrow\; C(H_N)\le N^{\max{\rho,q}+o(1)}.
\]
\emph{Valid when:} Applies after SAT circuits are substituted only into length L=(log N)\textasciicircum{}d generated queries, so L\textasciicircum{}c=N\textasciicircum{}\{o(1)\} for each fixed c.

\paragraph{Equation 4: route\_magnification\_consequence.}
\[
C(H_N)> N^{1+\epsilon}\ \text{eventually and}\ \max{\rho,q}< 1+\epsilon \;\Rightarrow\; SAT\notin P/poly.
\]
\emph{Valid when:} Requires an eventually valid unrestricted-circuit lower bound for the named carrier and fixed target-preserving parameters q,rho,epsilon with positive slack.

\paragraph{Next analysis.} Instantiate the pair-pattern thought experiment for \(t=2\): choose q with \(q<1+\epsilon\) and test whether \(2q+\alpha<1+\epsilon\) permits a subthreshold pair separator \(S_{N}\); compute the residual complexity C(H\_N xor \(S_{N}\)) and check for hidden SAT lookup.

\paragraph{Falsifiers.}
\begin{itemize}
\item tuple-feature falsifier: some arity t system has \(max{\rho,s+tq+\alpha,\gamma,\lambda}<1+\epsilon\) and an exact or low-cost residual for \(H_{N}\).
\item sparse-tuple falsifier: the carrier admits a selector or algebraic shortcut computing the same tuple transcript below the dense \(M_{N}^{t}\) exponent.
\item target-window falsifier: \(max{\rho,q}\ge 1+\epsilon\), so the SAT-circuit substitution no longer contradicts the claimed source lower bound.
\item lookup falsifier: L or the generated range is small enough that SAT labels on the relevant range are nonuniformly tabulated below \(N^{1+\epsilon}\).
\item nonuniform-leakage falsifier: advice encodes a hitting family for positive bags or a separator for negative bags with subthreshold residual correction.
\item natural-proof falsifier: the only available source proof is a large constructive property of truth tables, leaving the route blocked under standard PRF-based natural-proofs barriers.
\item model-mismatch falsifier: the tuple lower bound is proved only for restricted circuits while the target bridge requires unrestricted P/poly substitution.
\end{itemize}

\subsection*{cycle-0035: Hardness-Magnification Bridge Realization: high-q arity cutoff for generated-OR carriers}

\paragraph{Execution.} research; azure\_responses\_api; attempt primary; 149.906 seconds; 2 lemmas, 6 equations, 5 derivation steps, 6 falsifiers.

\paragraph{Learning.} Cycle 35 finds a parameter repair: set q in ((1+epsilon)/2,1+epsilon) to preserve the generated-OR target bridge while forcing every dense tuple-feature transcript of arity at least two above the \(N^{1+\epsilon}\) threshold. The route-level lower-bound obligation is now the sharper unary residual minimax gate, not an infinite tuple-arity audit.

\paragraph{Pressure.} No proof is known that an explicit high-q generator defeats all unary range-local transcripts and subthreshold residual corrections. High q may even help unary predicates by presenting many chances to hit a satisfiable generated formula.

\paragraph{Derivation steps.}
\begin{enumerate}
\item Instantiate \(H_{N}\) with \(M_{N}=N^{q}\) and L\_N=(log N)\textasciicircum{}d, \(d>1\), so SAT circuits of size \(L_{N}^{c}\) are \(N^{o(1)}\) after substitution.
\item Assume \(\mathrm{SAT}\in \mathrm{SIZE}(n^{c})\). Composing the SAT circuit with the bulk generator and the OR over \(M_{N}\) queries gives \(C(H_{N})\le N^{max{\rho,q}+o(1)}\); hence a valid \(H_{N}\) lower bound above \(N^{1+\epsilon}\) contradicts every fixed c when \(max{\rho,q}<1+\epsilon\).
\item Model nonuniform leakage by \(N^{s}\) predicates on t-tuples of generated formulas, each composed predicate costing \(N^{\alpha+o(1)}\). Dense enumeration gives \(N^{s+tq+\alpha+o(1)}\) transcript cost.
\item If a size \(N^{a}\) aggregator plus residual correction of size \(N^\lambda\) recovers \(H_{N}\), then \(H_{N}\) has size \(N^{max{\rho,s+tq+\alpha,a,\lambda}+o(1)}\). Therefore any source lower bound must force this maximum to be at least 1+epsilon-o(1).
\item Choose \(q>(1+\epsilon)/2\). Since \(s,\alpha\ge 0\), every \(t\ge 2\) dense tuple transcript has exponent \(s+tq+\alpha>1+\epsilon\), so it cannot be a subthreshold falsifier. The audit collapses to \(t=1\) transcripts and residual corrections.
\end{enumerate}

\paragraph{Equation 1: carrier\_parameters.}
\[
M_N=N^q,\quad L_N=(\log N)^d,\quad C(G_N)\le N^{\rho+o(1)}
\]
\emph{Valid when:} Valid when the carrier family has fixed exponents q,rho,d and the generator cost is charged in the N-input circuit model.

\paragraph{Equation 2: target\_substitution.}
\[
SAT\in SIZE(n^c)\ \Longrightarrow\ C(H_N)\le N^{\max{\rho,q}+o(1)}
\]
\emph{Valid when:} Applies when \(d>1\) is fixed, c is constant, and each \(length-L_{N}\) SAT circuit contributes only L\_N\textasciicircum{}c=N\textasciicircum{}\{o(1)\} overhead.

\paragraph{Equation 3: dense\_tuple\_transcript\_cost.}
\[
C(T_{s,t,\alpha})\le N^{s+tq+\alpha+o(1)}
\]
\emph{Valid when:} Valid when \(N^{s}\) feature predicates are evaluated on all \(N^{tq}\) generated t-tuples and each composed predicate costs \(N^{\alpha+o(1)}\).

\paragraph{Equation 4: compression\_falsifier\_bound.}
\[
C(H_N)\le N^{\max{\rho,s+tq+\alpha,a,\lambda}+o(1)}
\]
\emph{Valid when:} Valid when an aggregator of size \(N^{a}\) on the transcript and a residual correction of size \(N^\lambda\) exactly recover \(H_{N}\).

\paragraph{Equation 5: source\_obligation.}
\[
C(H_N)> N^{1+\epsilon-o(1)}\Rightarrow \max{\rho,s+tq+\alpha,a,\lambda}\ge 1+\epsilon-o(1)
\]
\emph{Valid when:} Requires the claimed source lower bound to hold against every nonuniform transcript, aggregator, and correction family in the named source model.

\paragraph{Equation 6: high\_q\_cutoff.}
\[
q> (1+\epsilon)/2\ \wedge\ t\ge2\ \Longrightarrow\ s+tq+\alpha> 1+\epsilon
\]
\emph{Valid when:} Valid when feature-count and feature-composition exponents are nonnegative and dense enumeration of all t-tuples is charged.

\paragraph{Next analysis.} Formalize unary transcripts as monotone OR-of-range predicates, then test whether random-looking high-q generators make every subthreshold unary transcript leave residual circuit complexity above \(N^{1+\epsilon-o(1)}\) or whether a selected-image separator falsifier appears.

\paragraph{Falsifiers.}
\begin{itemize}
\item Unary feature residual compression with exponent \(max{\rho,s+q+\alpha,a,\lambda}<1+\epsilon\) is the decisive remaining falsifier.
\item If \(q\le (1+\epsilon)/2\), binary tuple-feature transcripts can still be subthreshold and the high-q repair is invalid.
\item If \(\rho\ge 1+\epsilon\), the generator alone destroys the target-preserving bridge.
\item If the generated range has a small promise SAT separator or low positive-fiber hitting cover, \(H_{N}\) may be easy despite global SAT hardness.
\item If a proposed lower-bound proof yields a dense constructive truth-table property, the natural-proofs obstruction applies; the route must be range-dependent or otherwise non-natural.
\item If the generator leaks SAT labels through advice or indexing structure, the source lower bound collapses by nonuniform leakage.
\end{itemize}

\subsection*{cycle-0036: Hardness-Magnification Bridge Realization: position-indexed unary residual gate}

\paragraph{Execution.} research; azure\_responses\_api; attempt primary; 152.758 seconds; 2 lemmas, 4 equations, 6 derivation steps, 5 falsifiers.

\paragraph{Learning.} Cycle 36 repairs the high-q unary reduction: binary and higher tuple features are threshold-blocked when \(q>(1+\epsilon)/2\), but shared-unary audits are not sound for general circuits. The valid remaining source obligation is a position-indexed unary transcript plus residual lower bound with total local budget \(U_{N}\). This is a sharper route-level gate and supplies an explicit falsifier class, without claiming \(\mathrm{P} \ne \mathrm{NP}\).

\paragraph{Pressure.} No lower-bound technique is supplied for the position-indexed unary residual minimax quantity. Proving it for an explicit \(H_{N}\) is essentially a general-circuit lower bound for a highly structured generated OR, not a consequence of the magnification bridge.

\paragraph{Derivation steps.}
\begin{enumerate}
\item Start with H\_N(x)=OR\_\{jN\textasciicircum{}\{1+epsilon\}.
\item The target bridge remains valid: \(\mathrm{SAT}\in \mathrm{SIZE}(n^{c})\) gives \(H_{N}\) circuits of size \(N^{max{\rho,q}+o(1)}\), so a valid \(H_{N}\) lower bound with \(max{\rho,q}<1+\epsilon\) implies SAT notin SIZE(poly).
\item Use the cycle-35 parameter \(q>(1+\epsilon)/2\). Any explicit tuple transcript that enumerates all pairs or higher arity tuples costs at least \(N^{2q-o(1)}\), already above \(N^{1+\epsilon}\).
\item Try to conclude that only a shared unary predicate A(F) must be ruled out. This fails because the nonuniform circuit for \(H_{N}\) can attach different small predicates \(A_{j}\) to different generated positions without paying the cost of a uniform rule for j.
\item Repair the audit by defining total-budget position-indexed unary transcripts and allowing an arbitrary aggregator plus residual corrector.
\item Derive the upper bound \(C(H_{N})\le N^{\rho+o(1)}+sum_{j}\) C(A\_j)+C(B)+C(R). Therefore any subthreshold position-indexed decomposition is a concrete falsifier for the claimed source lower bound.
\end{enumerate}

\paragraph{Equation 1: target\_substitution.}
\[
C(H_N) \le N^{\rho+o(1)}+N^q\,L^c = N^{\max{\rho,q}+o(1)}
\]
\emph{Valid when:} Valid when SAT has nonuniform circuits of size \(n^{c}\) on L-bit inputs and the generator emits all queries with size \(N^\rho\).

\paragraph{Equation 2: high\_q\_tuple\_floor.}
\[
tq> 1+\epsilon \quad \text{for every } t\ge 2 \text{ if } q> (1+\epsilon)/2
\]
\emph{Valid when:} Applies in the high-q parameter regime before charging any predicate or aggregator overhead to tuple features.

\paragraph{Equation 3: position\_unary\_upper\_bound.}
\[
C(H_N) \le N^{\rho+o(1)}+U_N+C(B_N)+C(R_N)
\]
\emph{Valid when:} Valid when \(H_{N}\) equals B\_N((A\_\{N,j\}(G\_\{N,j\}(x)))\_\{j\textbackslash\{\}le M\}) \(\oplus\) R\_N(x) as a Boolean function.

\paragraph{Equation 4: source\_obligation.}
\[
\inf_{A_{j},B,R}\max{\rho,\log_N U_N,\log_N C(B),\log_N C(R)}> 1+\epsilon-o(1)
\]
\emph{Valid when:} Requires the target-preserving lower-bound claim to survive all nonuniform position-indexed unary transcript decompositions and residual corrections.

\paragraph{Next analysis.} Formalize \(Iota_{N}^{pos}\) as a minimax transcript-residual measure and test whether expander-style slot mixing can rule out all low-budget position-indexed unary decompositions without becoming a natural-proof argument.

\paragraph{Falsifiers.}
\begin{itemize}
\item position-indexed unary predicates with total budget below \(N^{1+\epsilon}\)
\item an aggregator \(B_{N}\) of size below \(N^{1+\epsilon}\) whose transcript classes are nearly monochromatic for \(H_{N}\)
\item a residual corrector \(R_{N}\) of size below \(N^{1+\epsilon}\) covering mixed transcript classes
\item generator leakage where \(G_{N,j}\) makes SAT labels simple for most fixed positions
\item threshold collapse if q<=((1+epsilon)/2), since binary tuple features may become subthreshold
\end{itemize}

\subsection*{cycle-0037: Hardness-Magnification Bridge Realization: approximate witness-cover gate}

\paragraph{Execution.} research; azure\_responses\_api; attempt primary; 158.44 seconds; 3 lemmas, 4 equations, 6 derivation steps, 5 falsifiers.

\paragraph{Learning.} Cycle 37 sharpens the position-indexed unary residual obligation into an approximate witness-cover threshold. In the high-q generated-OR route, dense tuple leakage is threshold-blocked, but unary assignment covers of total size below \(N^{1+\epsilon}\) plus fewer than \(N^\epsilon\) missed positives still give exact circuits for \(H_{N}\) after equality patching. This is route-level progress: a valid candidate source lower bound must first prove high approximate witness-cover complexity, while the main \(\mathrm{P} \ne \mathrm{NP}\) conjecture remains untouched.

\paragraph{Pressure.} No explicit generated-OR family is known to have approximate position-indexed witness-cover complexity above this threshold, and defeating witness covers is still weaker than defeating arbitrary position-indexed unary predicates plus residual circuits.

\paragraph{Derivation steps.}
\begin{enumerate}
\item Fix nonuniform assignment sets \(A_{N,j}\) and build B\_A(x) by evaluating each generated formula F\_\{N,j\}(x) on all assignments in \(A_{N,j}\).
\item Use the multioutput generator once, costing \(N^{\rho+o(1)}\), and share its formula bits across all assignment evaluations.
\item Each assignment evaluation costs \(L_{N}^{O(1)}\), so the total certificate-test cost is \(L_{N}^{O(1)}W_{N}\) plus the final OR over \(M_{N}\) positions.
\item If \(B_{A}\) misses \(E_{N}\) positive inputs of \(H_{N}\), add equality tests for those N-bit inputs and OR them into \(B_{A}\), costing O(N \(E_{N}\)).
\item The resulting circuit computes \(H_{N}\) exactly with size N\textasciicircum{}\{rho+o(1)\}+N\textasciicircum{}\{q+o(1)\}+L\_N\textasciicircum{}\{O(1)\}W\_N+O(N \(E_{N}\)).
\item Therefore any claimed \(N^{1+\epsilon-o(1)}\) lower bound for \(H_{N}\) implies the necessary exponent inequality max\{rho,q,omega,1+eta\}>=1+epsilon-o(1).
\end{enumerate}

\paragraph{Equation 1: carrier.}
\[
H_N(x)=\bigvee_{j\le M_N} SAT(F_{N,j}(x)),\quad M_N=N^q,\quad L_N=(\log N)^d.
\]
\emph{Valid when:} Valid when the source carrier is the generated-OR family with explicit multioutput generation and fixed polynomially many indexed queries.

\paragraph{Equation 2: cover\_predictor.}
\[
B_A(x)=\bigvee_{j\le M_N}\bigvee_{a\in A_{N,j}} Eval(F_{N,j}(x),a).
\]
\emph{Valid when:} Valid when \(A_{N,j}\) is position indexed and Eval tests satisfaction for the generated formula encoding.

\paragraph{Equation 3: upper\_bound.}
\[
SIZE(H_N)\le N^{\rho+o(1)}+N^{q+o(1)}+L_N^{O(1)}W_N+O(NE_N).
\]
\emph{Valid when:} Applies when the cover predictor is corrected by hardwired equality tests for all uncovered positive inputs.

\paragraph{Equation 4: necessary\_threshold.}
\[
\max{\rho,q,\omega,1+\eta}\ge 1+\epsilon-o(1),\quad W_N=N^{\omega+o(1)},\ E_N=N^{\eta+o(1)}.
\]
\emph{Valid when:} Requires comparing the explicit upper bound against a claimed \(N^{1+\epsilon-o(1)}\) lower bound in general circuits.

\paragraph{Next analysis.} Instantiate the thought experiment with a concrete planted/anti-planted CNF generator and compute the minimum position-indexed assignment hitting exponent h for its satisfiable range; compare h to 1+epsilon-q.

\paragraph{Falsifiers.}
\begin{itemize}
\item falsifier: exact per-position assignment hitting sets with total \(W_{N}<N^{1+\epsilon-\delta}\).
\item falsifier: approximate witness covers with \(W_{N}<N^{1+\epsilon-\delta}\) and fewer than \(N^{\epsilon-\delta}\) missed positive inputs.
\item falsifier: planted-solution or low-witness-entropy generated ranges where one small assignment list hits every satisfiable formula in each position range.
\item falsifier: nonuniform advice choosing \(A_{N,j}\) separately for every N and j; this leakage is valid in the source SIZE model and must be charged.
\item falsifier: generator implementation whose full multioutput cost already reaches \(N^{1+\epsilon}\), destroying the target-preserving threshold.
\end{itemize}

\subsection*{cycle-0038: Hardness-Magnification Bridge Realization: realized-witness projection gate}

\paragraph{Execution.} research; azure\_responses\_api; attempt primary; 127.795 seconds; 2 lemmas, 4 equations, 6 derivation steps, 5 falsifiers.

\paragraph{Learning.} Cycle 38 isolates a canonical unary leakage mechanism: all realized satisfying witnesses per generated position form a zero-error cover. Thus a valid high-q generated-OR route needs W\_N\textasciicircum{}*>N\textasciicircum{}\{1+epsilon-o(1)\} or a stronger approximate-cover lower bound. The target substitution bridge survives, but source plausibility is now governed by witness-projection entropy rather than formula-range entropy alone.

\paragraph{Pressure.} This proves only a necessary gate for the carrier lower bound. It does not supply the hard lower bound against all possible approximate witness covers.

\paragraph{Derivation steps.}
\begin{enumerate}
\item Fix nonuniform lists \(Y_{j}\) and build one shared circuit for \(G_{N}\) that outputs all formulas F\_\{N,j\}(x).
\item For every pair (j,a) with a in \(Y_{j}\), attach a small evaluator computing whether a satisfies F\_\{N,j\}(x).
\item OR all evaluator outputs; this never accepts a negative input because every generated formula in a negative bag is unsatisfiable.
\item If a positive input is missed by all lists, add an equality patch for that x; \(E_{N}\) missed positives cost \(E_{N}*N^{1+o(1)}\) gates.
\item The resulting size is N\textasciicircum{}\{rho+o(1)\}+(sum\_j |Y\_j|)N\textasciicircum{}\{o(1)\}+E\_N \(N^{1+o(1)}\).
\item Therefore an \(N^{1+\epsilon-o(1)}\) source lower bound forces either total cover size at least \(N^{1+\epsilon-o(1)}\) or at least \(N^{\epsilon-o(1)}\) missed positives.
\end{enumerate}

\paragraph{Equation 1: target\_substitution.}
\[
SAT\in SIZE(n^c)\Rightarrow C(H_N)\le N^{\max{\rho,q}+o(1)}
\]
\emph{Valid when:} Valid when the batch generator outputs all queried formulas with shared cost \(N^{\rho+o(1)}\) and each SAT replacement circuit is used on L=(log N)\textasciicircum{}d bit instances.

\paragraph{Equation 2: witness\_cover\_upper\_bound.}
\[
C(H_N)\le N^{\rho+o(1)}+\left(\sum_{j\le N^q}|Y_j|\right)N^{o(1)}+E_N\,N^{1+o(1)}
\]
\emph{Valid when:} Applies after the witness lists hit every positive bag except \(E_{N}\) inputs and negative inputs remain rejected by direct formula evaluation.

\paragraph{Equation 3: projection\_entropy\_threshold.}
\[
\sum_{j\le N^q}|Y_j|\ge N^{1+\epsilon-o(1)}\quad\text{or}\quad E_N\ge N^{\epsilon-o(1)}
\]
\emph{Valid when:} Requires a claimed source lower bound \(C(H_{N})>N^{1+\epsilon-o(1)}\) with \(\rho<1+\epsilon\) and evaluator overhead only \(N^{o(1)}\).

\paragraph{Equation 4: average\_support\_gate.}
\[
\frac{1}{N^q}\sum_{j\le N^q}|Y_j|\ge N^{1+\epsilon-q-o(1)}
\]
\emph{Valid when:} Valid when the misses are subpatchable and the total cover lower bound is converted into a per-position average witness-support obligation.

\paragraph{Next analysis.} Instantiate the planted-pool thought experiment and compute \(W_{N}^*\) for affine, low-degree, and pseudorandom assignment-pool generators; check whether any preserves \(q+a>1+\epsilon\) without introducing natural-proof or advice-leakage collapse.

\paragraph{Falsifiers.}
\begin{itemize}
\item realized witness projection total \(W_{N}^*\) below \(N^{1+\epsilon-o(1)}\)
\item approximate witness lists with fewer than \(N^{\epsilon-o(1)}\) missed positives and subthreshold total size
\item positive fiber size below \(N^{\epsilon-o(1)}\), since equality patching alone makes \(H_{N}\) small
\item batch generator not shareable, causing the source and target interfaces to mismatch
\item \(\rho\) reaching 1+epsilon, which hides witness-cover leakage behind generator cost
\end{itemize}

\subsection*{cycle-0039: High-q generated-OR carrier: partial-certificate cover gate}

\paragraph{Execution.} research; azure\_responses\_api; attempt primary; 970.595 seconds; 3 lemmas, 5 equations, 5 derivation steps, 5 falsifiers.

\paragraph{Learning.} Cycle 39 sharpens the generated-OR source audit: full realized witness entropy is not the right gate because partial satisfying terms can compress many witnesses. The valid target bridge remains \(max{\rho,q}<1+\epsilon\), but the route now needs a superthreshold lower bound against position-indexed partial-term covers, and ultimately against sound local certificate covers.

\paragraph{Pressure.} No valid extraction theorem converts an arbitrary small circuit for \(H_{N}\) into a small full-witness set, and partial certificates give a concrete counter-mechanism.

\paragraph{Derivation steps.}
\begin{enumerate}
\item Start with the previous necessary condition W\_N\textasciicircum{}*>N\textasciicircum{}\{1+epsilon-o(1)\} for realized full satisfying witnesses.
\item Try to infer \(\mathrm{SIZE}(H_{N})>N^{1+\epsilon}\) from high \(W_{N}^*\) plus the valid target bridge \(max{\rho,q}<1+\epsilon\).
\item Observe that a circuit for \(H_{N}\) need not enumerate full witnesses; it may hardwire a partial term or another sound local certificate per position and verify it against the generated CNF.
\item Construct the circuit OR\_\{j,p in P\_j\} TermSat(F\_\{N,j\}(x),p) after computing all F\_\{N,j\}(x) once; this decides \(H_{N}\) whenever P is a complete satisfying-term cover.
\item Conclude that full-witness entropy is only a weak necessary condition and must be replaced by \(tau_{N}^{term}\), or more robustly \(K_{N}^{cert}\).
\end{enumerate}

\paragraph{Equation 1: generated\_or\_carrier.}
\[
H_N(x)=\bigvee_{j=1}^{M}\operatorname{SAT}(F_{N,j}(x)),\qquad M=N^q,\qquad |F_{N,j}|=L=(\log N)^d.
\]
\emph{Valid when:} Valid when the source object is the generated-OR CNF carrier with shared generator outputs and fixed q,d.

\paragraph{Equation 2: partial\_term\_cover\_complexity.}
\[
\tau_N^{\mathrm{term}}=\min_{\mathcal{P}}\sum_{j\le M}\sum_{p\in\mathcal{P}_j}(1+|p|).
\]
\emph{Valid when:} Valid when each hardwired partial assignment p must satisfy every clause it certifies in the actual generated CNF range.

\paragraph{Equation 3: cover\_upper\_bound.}
\[
\operatorname{SIZE}(H_N)\le C(G_N)+\tau_N^{\mathrm{term}}L^{O(1)}.
\]
\emph{Valid when:} Applies when the generator computes all query blocks once and the circuit verifies each hardwired term locally.

\paragraph{Equation 4: subthreshold\_falsifier.}
\[
\tau_N^{\mathrm{term}}\le N^{1+\epsilon-\delta}\ \wedge\ C(G_N)\le N^{1+\epsilon-\delta}\Rightarrow \operatorname{SIZE}(H_N)\le N^{1+\epsilon-\delta/2}.
\]
\emph{Valid when:} Requires polylogarithmic query length so the verifier overhead is swallowed by the exponent slack.

\paragraph{Equation 5: target\_substitution\_bridge.}
\[
\operatorname{SAT}\in\operatorname{SIZE}(n^c)\Rightarrow \operatorname{SIZE}(H_N)\le N^{\max{\rho,q}+o(1)}.
\]
\emph{Valid when:} Applies after the reduction preserves nonuniform advice accounting, shared generation, and the q,rho threshold parameters.

\paragraph{Next analysis.} Formalize \(tau_{N}^{term}\) as a set-cover problem over generated CNF ranges and test whether random-looking high-q generators can force tau\_N\textasciicircum{}\{term\}>N\textasciicircum{}\{1+epsilon\} without making \(G_{N}\) nonexplicit or violating \(\rho<1+\epsilon\).

\paragraph{Falsifiers.}
\begin{itemize}
\item Small position-indexed satisfying-term cover tau\_N\textasciicircum{}\{term\}<=N\textasciicircum{}\{1+epsilon-delta\} is a direct falsifier for the claimed source lower bound.
\item Large full-witness projection entropy \(W_{N}^*\) with tiny planted partial terms falsifies the cycle-38 audit as sufficient evidence.
\item A compact sound local certificate cover \(K_{N}^{cert}\) below threshold is a stronger nonuniform-leakage falsifier than promise-SAT easiness.
\item If the query encoding is not CNF or lacks a cheap sound partial-certificate verifier, the term-cover lemma is not valid and must be replaced by the certificate-cover version.
\item If \(\rho\) or q reaches 1+epsilon, the target-preserving bridge collapses independently of the source audit.
\end{itemize}

\subsection*{cycle-0040: Hardness-Magnification Bridge Realization: sound local certificate cover gate}

\paragraph{Execution.} research; azure\_responses\_api; attempt primary; 120.336 seconds; 3 lemmas, 4 equations, 7 derivation steps, 6 falsifiers.

\paragraph{Learning.} Cycle 40 sharpens the source audit from partial certificate covers to approximate sound local certificate covers. The exact upper bound C(H\_N)<=C(G\_N)+O(K(s+L\textasciicircum{}\{O(1)\})+EN) makes any subthreshold SLC cover a concrete falsifier. The target bridge remains valid, but a contribution now requires a superthreshold lower bound against one-sided local predicates on the generated range, not just high witness or formula entropy.

\paragraph{Pressure.} No current argument proves that every one-sided local certificate predicate family of total size below \(N^{1+\epsilon-o(1)}\) misses at least \(N^{\epsilon-o(1)}\) positive bags on the actual generated range.

\paragraph{Derivation steps.}
\begin{enumerate}
\item Fix \(\epsilon\), q in ((1+epsilon)/2,1+epsilon), rho1.
\item Define \(H_{N}\) as the OR of \(N^{q}\) generated 3CNF SAT instances and keep \(C(G_{N})\le N^{max{\rho,q}+o(1)}\) as the shared interface cost.
\item Assume a (K,E,s)-SLC cover. Build a circuit that computes G\_N(x), evaluates every \(A_{i}\) on its indexed formula \(G_{N}(x)_{j_{i}}\), and ORs the K one-sided outputs.
\item Soundness of \(A_{i}\) prevents false positives on H\_N\textasciicircum{}\{-1\}(0); completeness except E positives leaves only E missed x values.
\item Patch the E missed positive inputs by hardwired equality tests of size O(EN), giving C(H\_N)<=C(G\_N)+O(K(s+L\textasciicircum{}\{O(1)\})+EN).
\item Therefore a proposed source lower bound \(C(H_{N})>N^{1+\epsilon-o(1)}\) must imply a superthreshold lower bound against SLC covers, not merely against explicit witnesses or partial terms.
\item The old target bridge still works: if SAT has size \(n^{c}\) circuits, each length-L query costs L\textasciicircum{}c=N\textasciicircum{}\{o(1)\}, so \(H_{N}\) has size \(N^{max{\rho,q}+o(1)}\).
\end{enumerate}

\paragraph{Equation 1: carrier.}
\[
H_N(x)=\bigvee_{j=1}^{\lfloor N^q\rfloor} SAT(G_N(x)_j),\qquad L=(\log N)^d.
\]
\emph{Valid when:} Valid when the source carrier is the generated-OR family with shared generator outputs and length-L 3CNF instances.

\paragraph{Equation 2: slc-upper-bound.}
\[
C(H_N)\le C(G_N)+O(K(s+L^{O(1)})+EN).
\]
\emph{Valid when:} Valid when the SLC predicates are one-sided sound and miss only E positive carrier inputs.

\paragraph{Equation 3: threshold-collapse-test.}
\[
K=N^\sigma,\ E=N^\delta,\ s=N^{o(1)},\ \max{\rho,q,\sigma,1+\delta}< 1+\epsilon\ \Rightarrow\ C(H_N)< N^{1+\epsilon-o(1)}.
\]
\emph{Valid when:} Valid when all four exponent contributions are bounded below the desired source lower-bound threshold.

\paragraph{Equation 4: necessary-source-gate.}
\[
C(H_N)> N^{1+\epsilon-o(1)}\ \Rightarrow\ \text{every }N^{o(1)}\text{-local SLC cover has }K> N^{1+\epsilon-o(1)}\text{ or }E> N^{\epsilon-o(1)}.
\]
\emph{Valid when:} Applies after q and \(\rho\) are already fixed below the magnification threshold.

\paragraph{Next analysis.} Compute the SLC cover exponent for generated formulas with bounded backdoor width \(b=\alpha\) log N and for small semantic classes such as Horn, 2-SAT, and bounded-width CNF; identify whether the threshold \(\alpha\ge 1+\epsilon-q\) is achievable without destroying the SAT-substitution bridge.

\paragraph{Falsifiers.}
\begin{itemize}
\item falsifier: a subthreshold (\(K,E,N^{o(1)}\)) SLC cover with \(K<N^{1+\epsilon-o(1)}\) and \(E<N^{\epsilon-o(1)}\) gives an explicit small circuit for \(H_{N}\).
\item Low-width backdoor falsifier: if every positive emitted formula has a satisfying term of width \(b=\alpha\) log N and \(q+\alpha<1+\epsilon\), enumerating all such terms collapses the source lower bound.
\item Compressed-certificate falsifier: a size-poly(L) one-sided predicate can represent many partial terms, so large term count alone is not a valid source certificate.
\item Lookup-boundary falsifier: if L=O(log N) with \(q+2^{L}\) exponent below 1+epsilon, local truth-table predicates replace genuine SAT hardness.
\item Generator-bias falsifier: if positives lie in a small easy satisfiable subclass such as Horn, 2-SAT, bounded-width, or bounded-backdoor CNFs, SLC cover size may be subthreshold.
\item Nonuniform leakage falsifier: predicates \(A_{i}\) may contain advice tailored to the generated range; the source lower bound must defeat this advice rather than only uniform algorithms.
\end{itemize}

\subsection*{cycle-0041: One-sided promise-separator gate for high-q generated-OR magnification}

\paragraph{Execution.} research; azure\_responses\_api; attempt primary; 170.751 seconds; 3 lemmas, 5 equations, 7 derivation steps, 6 falsifiers.

\paragraph{Learning.} Cycle 41 isolates the exact semantic object behind the sound-local-certificate audit: Sep\_N\textasciicircum{}\{1s\}(E). Random term-cover calculations support the high-q polylog design, but they do not imply the needed separator lower bound. The main new obstruction is the term-cover versus one-sided semantic-separator gap.

\paragraph{Pressure.} No valid extractor converts arbitrary one-sided semantic separators into partial witnesses; range closure under restrictions is absent, so term-cover calibration cannot certify \(Sep_{N}^{1s}\) lower bounds.

\paragraph{Derivation steps.}
\begin{enumerate}
\item Start with \(H_{N}(x)=OR_{j}\) SAT(G\_N(x,j)) and retain \(q,\rho<1+\epsilon\), L\_N=(log N)\textasciicircum{}d.
\item Substitute an assumed SAT circuit of size \(n^{c}\) into every emitted formula; since L\_N\textasciicircum{}c=N\textasciicircum{}\{o(1)\}, obtain \(C(H_{N})\le N^{max{\rho,q}+o(1)}\).
\item Define \(D_{j}\) as one-sided promise separators for satisfiable formulas in the j-th generated range; \(OR_{j}\) D\_j(G\_N(x,j)) has no false positives for \(H_{N}\).
\item Patch E missed positive inputs by equality tests on x, giving C(H\_N)<=C(G\_N)+O(sum\_j \(|D_{j}|+E\) N).
\item Therefore any valid \(N^{1+\epsilon}\) lower bound for \(H_{N}\) must lower-bound Sep\_N\textasciicircum{}\{1s\}(N\textasciicircum{}\{epsilon-eta\}) above \(N^{1+\epsilon-\eta-o(1)}\) for every fixed \(\eta>0\).
\item A random k-CNF calculation shows partial-assignment term covers need 2\textasciicircum{}\{Omega((log N)\textasciicircum{}d)\} terms per position, far above the allowed \(N^{1+\epsilon-q}\) budget.
\item Attempt to promote term-cover hardness to separator hardness by SAT self-reduction fails because \(D_{j}\) is promised only on generated formulas, and restrictions of generated formulas need not remain in the generated range.
\end{enumerate}

\paragraph{Equation 1: carrier.}
\[
H_N(x)=OR_{j=1}^{M_N} SAT(G_N(x,j)), M_N=N^q, L_N=(log N)^d
\]
\emph{Valid when:} Valid when the source object is the generated-OR carrier with fixed position count and polylogarithmic emitted SAT length.

\paragraph{Equation 2: target-substitution.}
\[
SAT in SIZE(n^c) => C(H_N) \le C(G_N)+O(M_N L_N^c+M_N)=N^{max{rho,q}+o(1)}
\]
\emph{Valid when:} Applies after the reduction preserves nonuniformity and each emitted SAT query has length L\_N=(log N)\textasciicircum{}d.

\paragraph{Equation 3: separator-composition.}
\[
C(H_N) \le C(G_N)+O(Sep_N^{1s}(E)+E N)
\]
\emph{Valid when:} Valid when the predicates are one-sided sound and the E missed positive inputs are patched by equality gates.

\paragraph{Equation 4: necessary-source-gate.}
\[
C(H_N)> N^{1+epsilon} => Sep_N^{1s}(N^{epsilon-eta})> N^{1+epsilon-eta-o(1)} for every fixed eta> 0
\]
\emph{Valid when:} Requires \(C(G_{N})\le N^{\rho+o(1)}\) with \(\rho<1+\epsilon\) and uses the explicit equality-patching cost.

\paragraph{Equation 5: magnification-consequence.}
\[
[C(H_N)> N^{1+epsilon} and max{rho,q}< 1+epsilon] => SAT notin SIZE(n^c) for every fixed c
\]
\emph{Valid when:} Valid when the lower bound is for the same \(H_{N}\) family and target substitution incurs only \(N^{o(1)}\) SAT-circuit overhead.

\paragraph{Next analysis.} Test a range-closure repair: augment \(G_{N}\) so restrictions of emitted formulas also appear in controlled positions, compute the q/rho overhead, and see whether this lets a one-sided separator self-reduce to a term cover without breaking the \(N^{1+\epsilon}\) threshold.

\paragraph{Falsifiers.}
\begin{itemize}
\item threshold collapse if \(q\ge 1+\epsilon\) or \(\rho\ge 1+\epsilon\)
\item nonuniform leakage through small one-sided promise separators \(D_{j}\)
\item model mismatch between syntactic term covers and semantic separators
\item range-closure failure blocking witness extraction from \(D_{j}\)
\item random-label or random-formula evidence that is not target-preserving for SAT
\item natural-proof-style largeness if the separator lower bound is proved only by random counting
\end{itemize}

\subsection*{cycle-0042: Cycle 42: distributional cap certificate for one-sided semantic separators}

\paragraph{Execution.} research; azure\_responses\_api; attempt primary; 149.719 seconds; 3 lemmas, 6 equations, 5 derivation steps, 5 falsifiers.

\paragraph{Learning.} Cycle 42 turns the one-sided semantic-separator obstruction into an exact distributional cap profile A\_N(s). The valid source route now needs a superthreshold cap-to-separator lower bound, not merely random witness or term-cover entropy. The target bridge is preserved, but any large one-sided cap is a concrete falsifier for the generated-OR lower bound.

\paragraph{Pressure.} Term-cover randomness does not imply a cap bound for arbitrary one-sided semantic predicates. The missing theorem is a one-sided conditional pseudorandomness statement for the distribution of generated formulas inside positive bags.

\paragraph{Derivation steps.}
\begin{enumerate}
\item Instantiate F\_\{N,j\}(x) as polylog-length k-CNF formulas with \(v=\beta\) log N variables and density near the satisfiability threshold; set \(Q=N^{q}\) with q in ((1+epsilon)/2,1+epsilon).
\item Try to upper-bound A\_N(s) by first bounding hit probability of assignment or partial-term predicates under \(mu_{N}^{+}\).
\item Attempt a union bound over all size-s L-input circuits P, using the hope that any large accepted set contains generated unsatisfiable formulas and hence violates one-sided soundness.
\item The count of circuits is 2\textasciicircum{}\{O(s log s)\}, and for \(s=N^\alpha\) this is far too large for term-cover estimates over individual witnesses.
\item More seriously, one-sided soundness is semantic on the generated range: a circuit may recognize a correlated satisfiable slice with no small witness-term description.
\end{enumerate}

\paragraph{Equation 1: carrier\_definition.}
\[
H_N(x)=OR_{j\le N^q} SAT(F_{N,j}(x)), |F_{N,j}(x)|=L=(log N)^d
\]
\emph{Valid when:} Valid when the carrier is the generated-OR source problem with fixed high-q parameters and a named SAT encoding.

\paragraph{Equation 2: target\_substitution\_bridge.}
\[
SAT in SIZE(n^c) => C(H_N) \le N^{rho+o(1)} + N^q (log N)^{dc+O(1)} = N^{max{rho,q}+o(1)}
\]
\emph{Valid when:} Applies after the reduction preserves the generator size, query length, and nonuniform circuit-size accounting.

\paragraph{Equation 3: cap\_profile.}
\[
A_N(s)=max_{j,P one-sided, size(P)\le s} Pr_{x~mu_N^+}[P(F_{N,j}(x))=1]
\]
\emph{Valid when:} Valid when one-sided soundness is imposed on the actual generated range, not on an informal global SAT promise.

\paragraph{Equation 4: union\_cover\_lower\_bound.}
\[
K A_N(s) \ge 1 - beta_N(E)
\]
\emph{Valid when:} Valid when K local predicates miss at most E positive inputs and \(mu_{N}^{+}\) gives those exceptions total mass beta\_N(E).

\paragraph{Equation 5: separator\_lower\_bound.}
\[
Sep_N^{1s}(E;s) \ge (1-beta_N(E))(s+L^{O(1)})/A_N(s) + EN
\]
\emph{Valid when:} Requires the cap bound for every allowed local predicate size and the standard equality patching charge for missed positives.

\paragraph{Equation 6: route\_threshold.}
\[
A_N(s) \le (1-beta_N(E))(s+L^{O(1)})/N^{1+epsilon-o(1)} => Sep_N^{1s}(E;s) \ge N^{1+epsilon-o(1)}
\]
\emph{Valid when:} Applies in the source lower-bound regime where separator cost is the only remaining known nonuniform leakage channel.

\paragraph{Next analysis.} Formalize the cap profile for the unmarked random k-CNF generated-OR thought experiment: compute assignment-test hit exponents as a function of \(\beta\), then identify the first semantic predicate class beyond terms that can violate or support \(A_{N}(s)\le s/N^{1+\epsilon-o(1)}\).

\paragraph{Falsifiers.}
\begin{itemize}
\item semantic planted-slice falsifier: a small predicate recognizes a satisfiable subset hit by most positive bags while avoiding all negative bags.
\item large cap-profile falsifier: for some s, A\_N(s) is much larger than \(s/N^{1+\epsilon}\), yielding a subthreshold separator.
\item exception concentration falsifier: the chosen positive distribution assigns large mass to E patchable positives, making beta\_N(E) too large for the cover lower bound.
\item nonuniform leakage falsifier: the generator range at one position admits a small promise SAT labeler even though global SAT on length L remains hard by exhaustive lookup.
\item threshold-collapse falsifier: \(q\ge 1+\epsilon\) or \(\rho\ge 1+\epsilon\) destroys the target-preserving substitution bridge.
\end{itemize}

\subsection*{cycle-0043: Generated-OR residual cap gate for hardness magnification}

\paragraph{Execution.} research; azure\_responses\_api; attempt compact\_retry; 268.905 seconds; 1 lemmas, 3 equations, 4 derivation steps, 4 falsifiers.

\paragraph{Learning.} Cycle 43 converts the cap profile into a residual fractional covering threshold. A valid source lower bound must rule out greedy one-sided cap covers with \(B_{N}(s,E)<N^{\epsilon-sigma-o(1)}\); absence of a large first cap is not enough.

\paragraph{Pressure.} No known method proves worst-case residual cap denominator \(B_{N}(s,E)>N^{\epsilon-sigma-o(1)}\) for semantic one-sided caps with nonuniform advice.

\paragraph{Derivation steps.}
\begin{enumerate}
\item Assume the generated-OR carrier and the target bridge \(max{\rho,q}<1+\epsilon\) are fixed.
\item Define B\_N(s,E) by adversarial residual positives, not by the first cap or an average distribution.
\item Apply greedy covering: residual size shrinks by a factor at most \(1-1/B_{N}\) until at most E positives remain.
\item OR the K caps and patch the E misses; comparing exponents gives the threshold \(b+1+sigma<1+\epsilon\) and \(\delta<\epsilon\).
\end{enumerate}

\paragraph{Equation 1: target\_substitution.}
\[
C(H_N) \le N^{\max{\rho,q}+o(1)}
\]
\emph{Valid when:} Valid when \(\mathrm{SAT}_{L}\) has \(L^{c}-size\) nonuniform circuits for fixed c, L=(log N)\textasciicircum{}d with \(d>1\), and \(G_{N}\) is shared once.

\paragraph{Equation 2: residual\_greedy\_cover.}
\[
K \le \lceil B_N(s,E)(N\ln 2+1)\rceil
\]
\emph{Valid when:} Valid when every residual positive set larger than E has a one-sided cap hitting at least a 1/B\_N(s,E) fraction.

\paragraph{Equation 3: cap\_upper\_bound.}
\[
C(H_N) \le N^\rho+O(K(s+L^{O(1)})+EN)
\]
\emph{Valid when:} Applies after caps are one-sided on the generated range, share \(G_{N}\), and the remaining E positive inputs are patched by equality tests.

\paragraph{Next analysis.} Try to construct residual caps for the planted/random 3CNF generator using partial-assignment certificates, and compute the resulting exponent b.

\paragraph{Falsifiers.}
\begin{itemize}
\item Residual-cap falsifier: \(s=N^{sigma}\), \(E=N^\delta\), \(B_{N}(s,E)=N^{b}\) with \(b+1+sigma<1+\epsilon\) and \(\delta<\epsilon\) gives small circuits for \(H_{N}\).
\item Target-threshold falsifier: \(q\ge 1+\epsilon\) or \(\rho\ge 1+\epsilon\) prevents SAT substitution from contradicting an \(N^{1+\epsilon}\) source lower bound.
\item Nonuniform-leakage falsifier: caps may be chosen after seeing \(G_{N}\) and the generated range; random term entropy is invalid unless it holds for adversarial residual sets.
\item Naturalness probe: a dense constructive property excluding all small caps risks the natural-proofs obstruction, so the certificate must be non-natural or distribution-local.
\end{itemize}

\subsection*{cycle-0044: Hardness-Magnification Bridge Realization: residual cap-denominator gate}

\paragraph{Execution.} research; azure\_responses\_api; attempt primary; 305.826 seconds; 3 lemmas, 5 equations, 6 derivation steps, 6 falsifiers.

\paragraph{Learning.} Cycle 44 converts the residual cap profile into a greedy cap-denominator theorem. A generated-OR source lower bound must now produce hard residual positive distributions against all one-sided local predicates; otherwise \(D_{N}\) \(N^{sigma}\) and E N give an explicit subthreshold upper-bound falsifier. This is route-level progress, not a main \(\mathrm{P} \ne \mathrm{NP}\) proof.

\paragraph{Pressure.} No lower bound is known showing that every residual positive distribution fools all one-sided generated-range predicates of size \(N^{sigma}\) at hit probability \(N^{-(1+\epsilon-sigma)}\).

\paragraph{Derivation steps.}
\begin{enumerate}
\item Fix sigma and the family \(\mathrm{P}_{sigma}^{1s}\) of \(size-N^{sigma}\) local predicates sound on the generated range.
\item Start with R\_0=H\_N\textasciicircum{}\{-1\}(1). While \(|R_{t}|>E\), choose \(A_{t}\) with \(cap_{R_{t}}(A_{t})\ge |R_{t}|/D_{N}\) and remove the hit positives.
\item After K steps, \(|R_{K}|=D_{N}\) log(|H\_N\textasciicircum{}\{-1\}(1)|/E), so K=O(D\_N N) is enough without needing to enumerate the universe uniformly.
\item Build the circuit: compute the generator, evaluate all chosen predicates, OR their hits, and patch the E missed positives by equality tests.
\item The resulting upper bound is below \(N^{1+\epsilon}\) when \(\rho<1+\epsilon\), \(D_{N}\) \(N^{sigma}\) is subthreshold up to logs, and E N is subthreshold.
\item Contraposition yields the hard-residual distribution obligation for any valid source lower bound.
\end{enumerate}

\paragraph{Equation 1: carrier.}
\[
H_N(x)=\bigvee_{j=1}^{N^q} SAT(G_N(x,j)),\qquad |G_N(x,j)|=(\log N)^d.
\]
\emph{Valid when:} Valid when the source carrier is the explicit generated-OR family with polylogarithmic SAT queries and fixed q,d parameters.

\paragraph{Equation 2: target\_substitution.}
\[
SAT\in SIZE(n^c)\ \Longrightarrow\ C(H_N)\le N^{\max{\rho,q}+o(1)}.
\]
\emph{Valid when:} Applies after the generator size is \(N^{\rho+o(1)}\) and each polylogarithmic SAT query receives a size \(L^{c}\) circuit.

\paragraph{Equation 3: greedy\_cover\_size.}
\[
K\le D_N\log |H_N^{-1}(1)|\le O(D_N N).
\]
\emph{Valid when:} Valid when every residual positive set larger than E has a sound predicate covering at least a \(1/D_{N}\) fraction.

\paragraph{Equation 4: source\_upper\_bound.}
\[
C(H_N)\le N^{\rho+o(1)}+N^{\sigma+o(1)}D_N\,polylog(2^N)+E N^{1+o(1)}.
\]
\emph{Valid when:} Requires that the greedy predicates are one-sided sound on the generated range and missed positives are patched by equality tests.

\paragraph{Equation 5: falsifier\_threshold.}
\[
D_N\le N^{1+\epsilon-\sigma-o(1)}/polylog(N),\quad E\le N^{\epsilon-o(1)}\ \Longrightarrow\ C(H_N)\le N^{1+\epsilon-o(1)}.
\]
\emph{Valid when:} Valid in the threshold regime where \(\rho\) is below 1+epsilon and the cover and patching terms dominate the audit.

\paragraph{Next analysis.} Instantiate the cap audit for partial-assignment predicates of width w on L=(log N)\textasciicircum{}2 formulas: compute theta(w,sigma) for random-looking generated ranges, then check whether sigma+theta can stay above 1+epsilon while q,rho remain in the target-preserving window.

\paragraph{Falsifiers.}
\begin{itemize}
\item residual cap denominator \(D_{N}\) below \(N^{1+\epsilon-sigma-o(1)}\) gives a concrete falsifier for the source lower bound
\item exception set E below \(N^{\epsilon-o(1)}\) is patchable by equality gates
\item nonuniform leakage through one-sided predicates sound only on the generated range
\item threshold collapse if q or \(\rho\) reaches 1+epsilon before the target substitution is applied
\item model mismatch if the cap lower bound is proved only for uniform predicates while the source model is nonuniform circuits
\item natural-proof leakage if the residual certificate becomes a large constructive global property
\end{itemize}

\subsection*{cycle-0045: Hardness-Magnification Bridge Realization: Greedy Cap Dual Gate}

\paragraph{Execution.} research; azure\_responses\_api; attempt primary; 184.453 seconds; 3 lemmas, 4 equations, 7 derivation steps, 6 falsifiers.

\paragraph{Learning.} Cycle 45 adds the missing logarithmic greedy-cover charge and converts the cap profile into an exact threshold: \(D_{sigma,E}<N^{\epsilon-sigma-o(1)}\) is a concrete upper-bound falsifier for the generated-OR source lower bound. A valid route must now prove hard residual positive distributions against all one-sided separators and also close the normal-form gap from arbitrary circuits to separator covers.

\paragraph{Pressure.} The derivation only rules out separator-cover upper bounds. It does not prove that every subthreshold exact circuit for \(H_{N}\) has a valid one-sided separator normal form.

\paragraph{Derivation steps.}
\begin{enumerate}
\item Let P\_N=H\_N\textasciicircum{}\{-1\}(1) and fix the valid separator family \(S_{sigma}\) of size-cost at most \(N^{sigma}\).
\item Assume every residual R subseteq \(\mathrm{P}_{N}\) with \(|R|>E\) has a separator A hitting at least |R|/D of R.
\item Greedy selection removes a 1/D fraction of the current residual, so after at most D log(|P\_N|/E) \(\le \) \(N^{1+o(1)}D\) choices it leaves at most E positives uncovered.
\item OR the chosen one-sided separators. Soundness gives no false positives on H\_N\textasciicircum{}\{-1\}(0); patch the remaining E positives by equality circuits of size O(EN).
\item The resulting upper bound is C(H\_N)<=N\textasciicircum{}\{rho+o(1)\}+N\textasciicircum{}\{1+sigma+o(1)\}D+O(EN).
\item If \(D<N^{\epsilon-sigma-o(1)}\) and \(E<N^{\epsilon-o(1)}\), this is below \(N^{1+\epsilon-o(1)}\), contradicting the proposed source lower bound.
\item Therefore a valid source lower-bound proof must either prove \(D\ge N^{\epsilon-sigma-o(1)}\) for all \(sigma<\epsilon\) or provide hard residual distributions with cap at most \(N^{-(\epsilon-sigma)+o(1)}\).
\end{enumerate}

\paragraph{Equation 1: target\_substitution\_upper\_bound.}
\[
C(H_N)\le N^{\rho+o(1)}+N^q L^c=N^{\max{\rho,q}+o(1)}
\]
\emph{Valid when:} Valid when L=(log N)\textasciicircum{}d with \(d>1\), \(\mathrm{SAT}_{L}\) has nonuniform circuits of size \(L^{c}\), and the generator for all \(N^{q}\) queries has size \(N^{\rho+o(1)}\).

\paragraph{Equation 2: separator\_cover\_upper\_bound.}
\[
C(H_N)\le C(G_N)+O(KN^{\sigma}+EN)
\]
\emph{Valid when:} Valid when the separators are semantically one-sided on the generated range, cover all but E positive bags, and equality patches for missed inputs cost O(N) each.

\paragraph{Equation 3: greedy\_cover\_size.}
\[
K\le \left\lceil D_{\sigma,E}\log\frac{|H_N^{-1}(1)|}{\max{1,E}}\right\rceil\le N^{1+o(1)}D_{\sigma,E}
\]
\emph{Valid when:} Valid when every residual positive set larger than E has a separator cap at least its reciprocal denominator, so greedy removal is legitimate.

\paragraph{Equation 4: cap\_denominator\_falsifier\_threshold.}
\[
D_{\sigma,E}< N^{\epsilon-\sigma-o(1)}\ \wedge\ E< N^{\epsilon-o(1)}\ \Longrightarrow\ C(H_N)< N^{1+\epsilon-o(1)}
\]
\emph{Valid when:} Valid contrapositively after the log \(|\mathrm{P}_{N}|\) \(\le \) N factor, generator cost, and equality-error corrections are all charged below the source threshold.

\paragraph{Next analysis.} Work in a restricted circuit model such as De Morgan formulas or \(\mathrm{AC0}[p]\) and test whether exact circuits for \(H_{N}\) admit an approximate one-sided separator normal form; compute the resulting \(D_{sigma,E}\) threshold for the same high-q polylog generated-OR carrier.

\paragraph{Falsifiers.}
\begin{itemize}
\item cap-denominator falsifier: some \(sigma<\epsilon\) and \(E<N^{\epsilon-o(1)}\) have \(D_{sigma,E}<N^{\epsilon-sigma-o(1)}\), yielding a subthreshold circuit for \(H_{N}\).
\item normal-form falsifier: an exact small circuit computes \(H_{N}\) using two-sided global information that cannot be represented by valid one-sided separator covers.
\item one-sidedness falsifier: a proposed local predicate accepts even one generated negative bag, invalidating the separator-cover upper bound and the residual distribution certificate.
\item nonuniform leakage falsifier: the generator or separator list hides advice selecting easy positive fibers, reducing the cover cost below \(N^{1+\epsilon}\).
\item barrier falsifier: a proof of the residual cap lower bound via a large constructive property of SAT truth tables would run into the natural-proofs obstruction under pseudorandom functions; the current semantic promise object is only outside that barrier's direct scope.
\item lookup-threshold falsifier: if \(d\le 1\) or variable entropy drops, exhaustive SAT lookup or partial certificate covers can re-enter below the source threshold.
\end{itemize}

\subsection*{cycle-0046: Hardness-Magnification Bridge Realization: normal-form extraction gate}

\paragraph{Execution.} research; azure\_responses\_api; attempt primary; 178.415 seconds; 3 lemmas, 5 equations, 6 derivation steps, 5 falsifiers.

\paragraph{Learning.} Cycle 46 converts the cycle-45 normal-form warning into a precise obstruction. Greedy cap and separator denominators give valid upper-bound falsifiers, but separator hardness is not by itself a source lower bound. A planted-control generated-OR carrier can have \(H_{N}\) easy while local formula separators remain hard, so the route now needs a nonrelativizing locality-extraction theorem or a direct arbitrary-circuit lower bound for \(H_{N}\).

\paragraph{Pressure.} No valid locality-extraction theorem maps arbitrary small circuits for \(H_{N}\) into one-sided formula-local separator covers with subthreshold overhead.

\paragraph{Derivation steps.}
\begin{enumerate}
\item Assume a circuit C of size \(S<N^{1+\epsilon}\) computes \(H_{N}\) exactly.
\item Try to cover positive inputs x by choosing a witness position j with SAT(G\_N(x,j)).
\item Attempt to define a local predicate \(\mathrm{P}_{j,C}\) on y=G\_N(x,j) that is one-sided sound and accepts enough chosen witnesses.
\item The only available predicate from C is global: it depends on x and on correlations among all generated positions, not on y alone.
\item DNF expansion of C gives global accepting regions on x, but conversion to formula-local predicates can cost \(2^{S}\) or require inverting \(G_{N}\).
\item A planted-control carrier makes the obstruction exact: \(H_{N}=b_{N}\) is easy, while local formula labels may remain separator-hard.
\end{enumerate}

\paragraph{Equation 1: generated\_or\_carrier.}
\[
H_N(x)=vee_{j=1}^{N^q} SAT(G_N(x,j)), |G_N(x,j)|=L=(log N)^d
\]
\emph{Valid when:} Valid when the generated query family has a shared circuit \(G_{N}\) and all positions have the same polylogarithmic target length.

\paragraph{Equation 2: target\_substitution\_bridge.}
\[
SAT in SIZE(n^c) => C(H_N) \le N^rho + N^q L^c + O(N^q) = N^{max(rho,q)+o(1)}
\]
\emph{Valid when:} Applies when SAT on length L has nonuniform circuits of size \(L^{c}\) and \(G_{N}\) is reused without position-specific advice.

\paragraph{Equation 3: separator\_cover\_upper\_bound.}
\[
Sep_N^{1s}(K,s,E) => C(H_N) \le C(G_N)+O(K(s+L^{O(1)})+E N)
\]
\emph{Valid when:} Requires K position-indexed one-sided local predicates that hit every positive bag except E explicitly patchable inputs.

\paragraph{Equation 4: planted\_control\_collapse.}
\[
[(b_N(x)=1 => H_N(x)=1) and (b_N(x)=0 => H_N(x)=0)] => H_N=b_N and C(H_N)\le C(b_N)+O(1)
\]
\emph{Valid when:} Valid when the control bit semantically determines the OR label of every generated bag, even if formula labels look locally random.

\paragraph{Equation 5: subthreshold\_separator\_falsifier.}
\[
K(s+L^{O(1)})+E N = o(N^{1+epsilon}) => the claimed C(H_N)> N^{1+epsilon-o(1)} source lower bound is false
\]
\emph{Valid when:} Applies only after the separator family is sound on the generated universe and the residual errors are patchable by equality gates.

\paragraph{Next analysis.} Formalize a restricted no-control generator class, such as block-product independent seeds with no shared label bit, and test whether an S-size circuit for \(OR_{j}\) SAT(G(x,j)) can be localized to a one-sided separator cover; compute the exact overhead in q,rho,sigma,E.

\paragraph{Falsifiers.}
\begin{itemize}
\item planted-control falsifier: a simple global bit determines the OR label while individual generated formulas remain locally hard to separate.
\item global-correlation circuit falsifier: C computes a relation among generated positions without identifying any satisfiable formula witness.
\item relativization falsifier: random oracle labels can make separator complexity high while \(H_{N}\) remains equal to an easy control bit.
\item threshold falsifier: any Sep\_N\textasciicircum{}\{1s\}(K,s,E) with K(s+L\textasciicircum{}\{O(1)\})+EN=o(N\textasciicircum{}\{1+epsilon\}) invalidates the source lower bound.
\item natural-proof leakage falsifier: a constructive dense separator-hardness certificate for SAT truth tables may fall inside the Razborov-Rudich obstruction.
\end{itemize}

\subsection*{cycle-0047: Hardness-Magnification Bridge Realization: locality-extraction obstruction}

\paragraph{Execution.} research; azure\_responses\_api; attempt primary; 197.268 seconds; 3 lemmas, 4 equations, 5 derivation steps, 4 falsifiers.

\paragraph{Learning.} Cycle 47 converts the locality-extraction warning into a formal relativizing obstruction. A planted-control oracle gives hard one-sided local separators but an O(1)-size generated-OR carrier, so separator hardness is not a valid source lower-bound certificate by itself. The target substitution bridge survives unchanged; the route must now prove a nonrelativizing extraction theorem for real SAT or attack C(H\_N) directly.

\paragraph{Pressure.} There is no valid black-box operation that extracts a formula-local one-sided separator from an arbitrary global circuit for \(H_{N}\). The missing step must exploit nonrelativizing real-SAT structure or be replaced by a direct arbitrary-circuit lower bound for \(H_{N}\).

\paragraph{Derivation steps.}
\begin{enumerate}
\item Assume a circuit C of size T computes \(H_{N}(x)=OR_{j}\) SAT(G\_\{N,j\}(x)).
\item For each positive input x, choose a satisfiable generated query \(F_{x,j}\); try to define a local predicate p(F) that accepts one such F and rejects all formulas from negative bags.
\item The naive definition \(p(F)=1\) iff there exists x,j with \(G_{N,j}(x)=F\) and \(C(x)=1\) is not one-sided, because C may be using nonlocal control information rather than any property of F.
\item Adding the condition that F never appears in a negative bag requires hard semantic range membership and may cost as much as the original \(H_{N}\) circuit.
\item The planted-control oracle realizes this failure sharply: C reads the control bit c, while a local separator must solve an independent random classification problem on the codewords.
\end{enumerate}

\paragraph{Equation 1: generator-size-threshold.}
\[
R=N^{\theta},\quad |G_N|=O(RL)=N^{\theta+o(1)}< N^{1+\epsilon}
\]
\emph{Valid when:} Valid when d is fixed, \(\theta\) is below 1+epsilon, and the codebook table is allowed as nonuniform generator advice.

\paragraph{Equation 2: planted-control-easy-carrier.}
\[
H_N^U(c,y,w)=U(G_N(c,y,w))=c,\quad C(H_N^U)=O(1)
\]
\emph{Valid when:} Valid when the oracle labels assign one to every positive codeword and zero to every negative codeword.

\paragraph{Equation 3: separator-counting-bound.}
\[
\Pr[\exists p: |p|\le S,\ p(a_y)=1,\ p(b_y)=0\ \forall y\le R]\le 2^{O(S\log S)-2R}
\]
\emph{Valid when:} Valid when the oracle codewords are sampled independently and S log S=o(R).

\paragraph{Equation 4: separator-hardness-window.}
\[
S=N^{\sigma},\quad \sigma< \theta\ \Longrightarrow\ \operatorname{Sep}^{1s}_N(0)> N^{\sigma}
\]
\emph{Valid when:} Valid when \(2^{L}\) is much larger than R and the union bound over size-S circuits succeeds.

\paragraph{Next analysis.} Try a real-SAT extraction using self-reducibility or witness isolation: from a small circuit for \(H_{N}\), define a local predicate that accepts formulas whose forced witness path changes the global OR value, compute the overhead, and test it against the planted-control oracle falsifier.

\paragraph{Falsifiers.}
\begin{itemize}
\item falsifier: planted-control oracle with \(\theta\) in (sigma,1+epsilon) makes \(H_{N}\) easy while zero-error local separators require size greater than \(N^{sigma}\).
\item falsifier: any relativizing separator-promotion theorem would prove a false implication in the oracle codebook model.
\item falsifier: nonuniform generator advice can encode hard local classification while the global carrier depends only on a cheap control bit.
\item falsifier: if \(max{\rho,q}\ge 1+\epsilon\), even the valid target substitution bridge loses its contradiction threshold.
\end{itemize}

\subsection*{cycle-0048: Relativizing locality-extraction gate for generated-OR carriers}

\paragraph{Execution.} research; azure\_responses\_api; attempt compact\_retry; 243.551 seconds; 1 lemmas, 3 equations, 4 derivation steps, 4 falsifiers.

\paragraph{Learning.} Cycle 48 makes the obstruction formal: separator-denominator hardness is only an upper-bound falsifier, not a valid lower-bound certificate for \(H_{N}\). The surviving route is direct C(H\_N) lower bounds or a nonrelativizing locality-extraction theorem with explicit threshold loss.

\paragraph{Pressure.} No known nonrelativizing locality-extraction theorem converts an arbitrary small \(H_{N}\) circuit into local SAT-range separators with only \(N^{o(1)}\) loss.

\paragraph{Derivation steps.}
\begin{enumerate}
\item Fix \(d>1\) and \(q,rhoN^{1+\epsilon}\) lower bound.
\item Try to replace the direct lower bound by proving every one-sided local separator cover has D \(N^{sigma}\) log \(N+EN\ge N^{1+\epsilon-o(1)}\).
\item For this replacement to be valid, the contrapositive must extract a small separator cover from any arbitrary circuit A computing \(H_{N}\) with size below \(N^{1+\epsilon}\).
\item A planted-control oracle sets one generated formula to a trivial satisfiable/unsatisfiable control depending on b(x), so H\_N=b(x) is easy while unrelated generated SAT ranges keep separator covers hard.
\end{enumerate}

\paragraph{Equation 1: carrier.}
\[
H_N(x)=\bigvee_{j\le N^q} SAT(G_N(x,j)),\quad |G_N(x,j)|=L=(\log N)^d.
\]
\emph{Valid when:} Valid when the carrier generator is uniform and emits length-L CNF instances with no nonuniform advice.

\paragraph{Equation 2: target-substitution.}
\[
SAT\in SIZE(n^c)\Rightarrow C(H_N)\le N^{\rho+o(1)}+N^qL^{O(c)}=N^{\max{\rho,q}+o(1)}.
\]
\emph{Valid when:} Applies after each generated query is substituted by the assumed SAT circuit and L is polylogarithmic.

\paragraph{Equation 3: separator-upper-bound.}
\[
C(H_N)\le C(G_N)+O(DN^\sigma\log N+EN).
\]
\emph{Valid when:} Requires a one-sided separator cover with D local predicates of size \(N^{sigma}\) and E patchable bad bags.

\paragraph{Next analysis.} Try to formulate a real-SAT locality extractor using witness self-reducibility: from a small circuit for \(H_{N}\), isolate a generated satisfiable query and build a one-sided predicate; audit algebrization and \(N^{o(1)}\) loss.

\paragraph{Falsifiers.}
\begin{itemize}
\item Planted-control oracle falsifier: \(H_{N}\) has O(1) oracle-size circuits although one-sided separators on noncontrol ranges remain hard.
\item Normal-form falsifier: an arbitrary \(H_{N}\) circuit may compute the OR globally without exposing any local formula predicate usable as a separator.
\item Nonuniform leakage falsifier: advice can name the control coordinate or residual exceptions unless its cost is charged inside \(\rho\) or EN.
\item Threshold falsifier: if \(max{q,\rho}\ge 1+\epsilon\), even a valid \(H_{N}\) lower bound no longer magnifies to \(\mathrm{SAT}\notin \mathrm{P}/\mathrm{poly}\).
\end{itemize}

\subsection*{cycle-0049: Hardness-Magnification Bridge Realization: locality extraction threshold}

\paragraph{Execution.} research; azure\_responses\_api; attempt primary; 140.845 seconds; 3 lemmas, 4 equations, 5 derivation steps, 5 falsifiers.

\paragraph{Learning.} Cycle 49 converts the surviving generated-OR route into an exact extraction-threshold problem. Separator-denominator hardness is a valid source certificate only if every small H circuit yields covers with denominator S/N\textasciicircum{}\{1+sigma\}+N\textasciicircum{}kappa and \(\kappa<\epsilon-sigma\). Neutralization gives a concrete extraction but is a demagnifying falsifier, because it asks for restricted-SAT hardness rather than a barely-superlinear H lower bound.

\paragraph{Pressure.} There is no valid operation converting an arbitrary global H circuit into local SAT-range separators without either neutralizing the generator or using a new real-SAT, nonrelativizing locality principle.

\paragraph{Derivation steps.}
\begin{enumerate}
\item Assume a circuit A of size \(S<N^{1+\epsilon}\) decides \(H_{N}^{G}\).
\item Use A to define a residual positive distribution over generated bags and try to build one-sided local predicates accepting at least one satisfiable query per positive bag.
\item Attempt SAT self-reduction on a generated formula y=G\_N(x,j); this manipulates assignments inside y but does not produce nearby inputs x' that alter only y.
\item Add a hypothetical neutralizing map E\_j(y); then A(E\_j(y)) exactly decides SAT on \(R_{j}\), giving extraction but also demagnifying the source obligation to restricted SAT hardness.
\item Therefore the useful bridge is not separator hardness alone but the quantitative nonrelativizing extraction bound D<=S/N\textasciicircum{}\{1+sigma\}+N\textasciicircum{}kappa with \(\kappa<\epsilon-sigma\).
\end{enumerate}

\paragraph{Equation 1: generated\_or\_carrier.}
\[
H_N^G(x)=\bigvee_{j=1}^{N^q} SAT(G_N(x,j)),\quad |G_N(x,j)|=(\log N)^d
\]
\emph{Valid when:} Valid when the source object is the high-q generated-OR carrier with explicitly bounded generator complexity.

\paragraph{Equation 2: target\_substitution\_bound.}
\[
SAT\in SIZE(n^c)\Rightarrow C(H_N^G)\le N^\rho+N^q(\log N)^{dc}+O(N^q)=N^{\max{\rho,q}+o(1)}
\]
\emph{Valid when:} Valid when SAT circuits are nonuniform, c is fixed, and generated queries have polylogarithmic length.

\paragraph{Equation 3: needed\_extraction\_bound.}
\[
A\in SIZE(S),\ A=H_N^G\Rightarrow D_{\sigma,E}(G;A)\le S/N^{1+\sigma}\cdot polylog(N)+N^\kappa
\]
\emph{Valid when:} Valid when a nonrelativizing locality-extraction theorem converts arbitrary global H circuits into one-sided local covers.

\paragraph{Equation 4: separator\_to\_source\_threshold.}
\[
D_{\sigma,E}(G)> N^{\epsilon-\sigma+\omega(1)},\ E< N^{\epsilon-o(1)},\ \kappa< \epsilon-\sigma\Rightarrow C(H_N^G)> N^{1+\epsilon-o(1)}
\]
\emph{Valid when:} Valid when the denominator lower bound holds for all residual covers produced by the extraction theorem.

\paragraph{Next analysis.} Formalize one candidate extraction using SAT self-reducibility plus generator local decodability, then compute whether the required decodability already implies neutralizing embeddings and hence demagnification.

\paragraph{Falsifiers.}
\begin{itemize}
\item falsifier: extraction tax \(\kappa\ge \epsilon-sigma\) makes the separator lower bound too weak to force \(C(H_{N})>N^{1+\epsilon}\).
\item falsifier: neutralizing embeddings provide a valid extraction but turn the source lower bound into a direct restricted-SAT lower bound at threshold \(N^{1+\epsilon}\).
\item falsifier: a planted-control oracle can keep one-sided local separators hard while making \(H_{N}\) easy, so any relativizing extraction theorem is false.
\item falsifier: a subthreshold promise solver on the actual generated range gives C(H\_N)<=N\textasciicircum{}rho+N\textasciicircum{}q \(T_{N}+N^{q+o(1)}\).
\item falsifier: a dense constructive direct-H lower-bound property triggers the natural-proofs obstruction under standard PRF hardness assumptions.
\end{itemize}

\subsection*{cycle-0050: Hardness-Magnification Bridge Realization: Relativizing Extraction Falsified}

\paragraph{Execution.} research; azure\_responses\_api; attempt primary; 158.599 seconds; 2 lemmas, 4 equations, 8 derivation steps, 6 falsifiers.

\paragraph{Learning.} Cycle 50 gives a precise obstruction: small \(H_{N}\) circuits can coexist with superthreshold one-sided separator complexity in a structured singleton oracle. Thus the extraction bridge must be explicitly nonrelativizing and real-SAT-specific; separator hardness alone is not a valid source lower-bound certificate.

\paragraph{Pressure.} The argument is an oracle obstruction, not a real-SAT lower bound. Real CNF satisfiability may have nonrelativizing syntactic structure unavailable to singleton oracle labels.

\paragraph{Derivation steps.}
\begin{enumerate}
\item Replace SAT labels on the generated range by an oracle A while preserving the generated-OR interface \(H_{N}^{A}(x)=OR_{j}\) A(G\_N(x,j)).
\item Choose a structured slab \(X_{N}\) of \(R=N^{a+o(1)}\) inputs with O(N)-size membership circuit.
\item Use \(d>1\) to inject \(X_{N}\) x [N\textasciicircum{}q] into the polylogarithmic label universe.
\item Plant exactly one positive label per slab bag according to an arbitrary selector \(f:X_{N}\to [N^{q}]\), and plant no positives outside the slab.
\item Then \(H_{N}^{A}\) is exactly the slab membership predicate, so C(H\_N\textasciicircum{}A)=O(N) despite arbitrary hidden choices of f.
\item A one-sided separator cover for all but E slab positives determines f on at least R-E inputs; otherwise it either misses the bag or accepts an oracle-negative label.
\item Counting covers of total size T gives T log T \(\ge \) Omega((R-E)log N\textasciicircum{}q)-O(Rh(E/R)).
\item For E1+epsilon, every valid cover has size \(N^{a-o(1)}\), contradicting any promised \(sub-N^{1+\epsilon}\) relativizing extraction.
\end{enumerate}

\paragraph{Equation 1: slab\_capacity.}
\[
R=|X_N|=N^{a+o(1)},\quad M=N^q,\quad RM=N^{a+q+o(1)}\le 2^L
\]
\emph{Valid when:} Valid when \(d>1\) and a,q are fixed constants, so the polylogarithmic formula universe contains the polynomial slab range.

\paragraph{Equation 2: easy\_carrier.}
\[
H_N^A(x)=\bigvee_{j\le M}A(G_N(x,j))=\mathbf{1}_{x\in X_N},\quad C(H_N^A)=O(N)
\]
\emph{Valid when:} Valid when the oracle plants one positive label in each slab bag and no positive labels outside the slab.

\paragraph{Equation 3: selector\_cover\_counting.}
\[
T\log T \ge \Omega((R-E)\log M)-O(Rh(E/R))
\]
\emph{Valid when:} Valid when a one-sided cover of size T must encode a partial selector agreeing with f on at least R-E slab inputs.

\paragraph{Equation 4: threshold\_contradiction.}
\[
E< N^\epsilon,\ a> 1+\epsilon\ \Longrightarrow\ T\ge N^{a-o(1)}> N^{1+\epsilon}
\]
\emph{Valid when:} Applies when the extraction promises subthreshold separator cost and the slab exponent is chosen above the lower-bound threshold.

\paragraph{Next analysis.} Audit real-SAT self-reduction or proof-certificate structure as a possible nonrelativizing extractor, and compute whether its verifier/certificate overhead stays below \(N^{1+\epsilon}\) rather than collapsing to restricted-SAT hardness.

\paragraph{Falsifiers.}
\begin{itemize}
\item structured-slab singleton oracle falsifier for subthreshold relativizing locality extraction
\item threshold collapse if the slab exponent a is not above 1+epsilon
\item model mismatch if the extraction treats formula labels as pair-oracle cells rather than real CNF strings
\item nonuniform leakage through an explicit selector table when the covered positive slab has size at most \(N^{1+\epsilon-o(1)}\)
\item generator collisions or low semantic range that allow a compact one-sided cover
\item neutralization that proves restricted-SAT hardness instead of a barely-superlinear \(H_{N}\) lower bound
\end{itemize}

\subsection*{cycle-0051: Hardness-Magnification Bridge Realization - cycle 51 direct-sum demagnification audit}

\paragraph{Execution.} research; azure\_responses\_api; attempt primary; 172.807 seconds; 3 lemmas, 5 equations, 6 derivation steps, 5 falsifiers.

\paragraph{Learning.} Cycle 51 identifies a new quantitative conservation law: the same gap \(g_{N}\) that protects the target transfer against all SAT polynomial exponents becomes the per-slot lower-bound burden for any direct-sum source proof. Hence slot-wise SAT hardness cannot be the weak source; a valid route must prove aggregate \(H_{N}\) hardness or supply a nonrelativizing extraction theorem.

\paragraph{Pressure.} This proves only a quarantine for factorized source certificates. It does not rule out a genuinely aggregate arbitrary-circuit lower bound for \(H_{N}\) or a nonrelativizing extraction theorem using real SAT structure.

\paragraph{Derivation steps.}
\begin{enumerate}
\item Let \(M_{N}=N^{q_{N}}\) and L=(log N)\textasciicircum{}d. Under \(\mathrm{SAT}_{L}\) in SIZE(L\textasciicircum{}c), substitute one SAT circuit into every generated slot; this gives \(C(H_{N})\le N^{\rho+o(1)}+M_{N}\) \(L^{c+O(1)}\).
\item To contradict a source lower bound \(C(H_{N})>N^{1+\epsilon}\), the carrier must satisfy \(\rho<1+\epsilon\) and \(M_{N}\) L\textasciicircum{}c=o(N\textasciicircum{}\{1+epsilon\}) for each fixed c. Equivalently \(g_{N}=N^{1+\epsilon}/M_{N}\) is \(L^{\omega(1)}\).
\item Assume the source proof factors through a per-slot SAT lower bound b(L) with at most linear aggregation: \(C(H_{N})\ge M_{N}\) \(b(L)/N^{o(1)}\).
\item For this to exceed \(N^{1+\epsilon}\), it must have \(b(L)>g_{N}/N^{o(1)}\).
\item Combining the target-preserving condition with the source requirement forces \(b(L)=L^{\omega(1)}\). For fixed q\_N=2\textasciicircum{}\{Omega(L\textasciicircum{}\{1/d\})\}.
\item Therefore projection-slice proofs and direct-sum-from-SAT proofs do not realize weak hardness magnification for this carrier.
\end{enumerate}

\paragraph{Equation 1: target\_substitution\_upper\_bound.}
\[
C(H_N) \le N^{\rho+o(1)} + M_N L^{c+O(1)}
\]
\emph{Valid when:} Valid when the generator circuit is shared once and each generated SAT query is answered by the same length-L circuit family.

\paragraph{Equation 2: target\_gap.}
\[
g_N := \frac{N^{1+\epsilon}}{M_N}=N^{1+\epsilon-q_N}
\]
\emph{Valid when:} Applies when the carrier threshold is \(N^{1+\epsilon}\) and the OR fan-in is \(M_{N}=N^{q_{N}}\).

\paragraph{Equation 3: p\_poly\_preservation.}
\[
\forall c\; M_N L^c=o(N^{1+\epsilon}) \quad\Longleftrightarrow\quad \forall c\; L^c=o(g_N)
\]
\emph{Valid when:} Valid when the target consequence is \(\mathrm{SAT}\notin \mathrm{P}/\mathrm{poly}\) rather than exclusion of only one fixed exponent.

\paragraph{Equation 4: direct\_sum\_source\_requirement.}
\[
M_N b(L)/N^{o(1)} > N^{1+\epsilon} \quad\Rightarrow\quad b(L) > g_N/N^{o(1)}
\]
\emph{Valid when:} Requires the proposed source lower bound to factor through \(M_{N}\) local SAT costs with only subpolynomial losses.

\paragraph{Equation 5: fixed\_gap\_demagnification.}
\[
q_N\le 1+\epsilon-\delta \Rightarrow g_N\ge 2^{\delta L^{1/d}-o(L^{1/d})}
\]
\emph{Valid when:} Valid when \(d>1\) is fixed and logarithms define L=(log N)\textasciicircum{}d up to constant factors.

\paragraph{Next analysis.} Try to formalize an aggregate lower-bound normal form for \(H_{N}\) that distinguishes global OR computation from \(M_{N}\) independent \(\mathrm{SAT}_{L}\) copies, and test it against the planted-control oracle before using separator hardness.

\paragraph{Falsifiers.}
\begin{itemize}
\item direct-sum demagnification falsifier: a source proof of the form \(M_{N}\) b(L) requires \(b(L)>g_{N}\), while target preservation requires \(g_{N}=L^{\omega(1)}\).
\item projection-slice falsifier: the special case \(M_{N}=1\) demands a length-L SAT lower bound of size N\textasciicircum{}\{1+epsilon\}=2\textasciicircum{}\{Omega(L\textasciicircum{}\{1/d\})\}.
\item near-boundary finite-exponent falsifier: if \(g_{N}=L^{B}\), the bridge is valid only against SAT circuits of exponent below B-O(1), not against P/poly.
\item relativizing separator falsifier: hard one-sided local separators can coexist with tiny \(H_{N}\) circuits in the planted-control oracle.
\item nonuniform sharing falsifier: multiple editable slices do not by themselves yield an \(M_{N}-fold\) lower bound for arbitrary circuits without an explicit direct-sum theorem.
\end{itemize}

\subsection*{cycle-0052: Hardness-Magnification Bridge Realization: XOR Carrier Repair}

\paragraph{Execution.} research; azure\_responses\_api; attempt primary; 151.347 seconds; 3 lemmas, 4 equations, 5 derivation steps, 6 falsifiers.

\paragraph{Learning.} Cycle 52 contributes a concrete repair: XOR aggregation preserves the valid polylog-query magnification threshold while invalidating the OR-specific one-sided separator obstruction. The new source burden is sharper: prove aggregate parity-carrier hardness or rule out subthreshold parity sketches and sparse corrections. This is route-level progress only, with explicit falsifier tests.

\paragraph{Pressure.} There is no nonrelativizing theorem extracting a parity sketch, low-rank transcript, or other analyzable normal form from an arbitrary small circuit for \(H_{N}^{oplus}\); direct C(H\_N\textasciicircum{}oplus)>N\textasciicircum{}\{1+epsilon\} remains unproved.

\paragraph{Derivation steps.}
\begin{enumerate}
\item Kept the polylog query window L=(log N)\textasciicircum{}d with \(d>1\) so SAT polynomial-size circuits substitute as \(N^{o(1)}\) per query and exhaustive lookup is not automatically subthreshold.
\item Replaced the generated-OR carrier by \(H_{N}^{oplus}\), the parity of \(M=N^{q}\) generated SAT labels.
\item Computed the target transfer C(H\_N\textasciicircum{}oplus)<=C(G\_N)+O(M(s(L)+L)+M), yielding \(N^{max{\rho,q}+o(1)}\) under \(\mathrm{SAT}\in \mathrm{SIZE}(L^{c})\).
\item Checked the previous one-sided separator machinery: it decides existential positivity of OR-bags but loses the parity of the number of satisfiable slots, so it cannot certify or refute XOR hardness without extra label information.
\item Isolated the new leakage normal form: any small feature or affine sketch plus sparse corrections gives an exact subthreshold circuit for \(H_{N}^{oplus}\).
\end{enumerate}

\paragraph{Equation 1: carrier\_definition.}
\[
H^{\oplus}_{N,G}(x)=\bigoplus_{j=1}^{M} SAT(G_{N,j}(x)),\quad M=N^q,\quad L=(\log N)^d.
\]
\emph{Valid when:} Valid when \(G_{N,j}\) output length-L encodings of SAT instances and the XOR fan-in M is explicitly charged.

\paragraph{Equation 2: substitution\_upper\_bound.}
\[
C(H^{\oplus}_{N,G})\le C(G_N)+O(M\cdot(s(L)+L)+M).
\]
\emph{Valid when:} Applies after one size-s(L) SAT circuit is substituted into every generated query without advice depending on x.

\paragraph{Equation 3: polylog\_transfer.}
\[
s(L)\le L^c,\ C(G_N)\le N^{\rho}\ \Rightarrow\ C(H^{\oplus}_{N,G})\le N^{\max{\rho,q}+o(1)}.
\]
\emph{Valid when:} Valid when d is fixed, c is fixed, and L=(log N)\textasciicircum{}d remains \(N^{o(1)}\).

\paragraph{Equation 4: target\_consequence.}
\[
\big(C(H^{\oplus}_{N,G})> N^{1+\epsilon}\ \mathrm{i.o.}\ \wedge\ \max{\rho,q}< 1+\epsilon\big)\Rightarrow SAT\notin P/poly.
\]
\emph{Valid when:} Requires the lower bound on the source carrier for infinitely many N and the target-preserving threshold gap to be fixed.

\paragraph{Next analysis.} In the ideal random-label model for \(G_{N}\), count subthreshold affine sketches of size \(N^\alpha\) and estimate the tau,alpha,error threshold at which a union bound could certify parity-sketch resistance without pretending to prove the actual explicit carrier lower bound.

\paragraph{Falsifiers.}
\begin{itemize}
\item threshold collapse: \(max{q,\rho}\ge 1+\epsilon\) destroys the target-preserving gap
\item polylog lookup collapse when \(d\le 1\) and q plus the lookup exponent falls below 1+epsilon
\item nonuniform leakage from \(G_{N}\) that encodes SAT labels instead of generating target-clean queries
\item subthreshold affine-sketch or feature-aggregator predictor for \(H_{N}^{oplus}\) with fewer than \(N^\epsilon\) patchable errors
\item complete per-slot promise solver on the generated range of total size below \(N^{1+\epsilon}\)
\item natural-proof style source argument that is dense and constructive against \(N^{1+\epsilon}\) circuits under standard pseudorandomness assumptions
\end{itemize}

\subsection*{cycle-0053: Hardness-Magnification Bridge Realization: XOR carrier affine-incidence audit}

\paragraph{Execution.} research; azure\_responses\_api; attempt primary; 128.573 seconds; 3 lemmas, 6 equations, 4 derivation steps, 7 falsifiers.

\paragraph{Learning.} Cycle 53 converts the XOR repair into an exact affine-incidence conservation law. The carrier equals the \(F_{2}\) inner product of the SAT label vector with a generated parity-incidence vector; low-dimensional computable incidence structure makes the carrier easy by hardwiring basis parities. This is a new concrete falsifier for the source lower bound, while the polylog-query target bridge remains valid.

\paragraph{Pressure.} A small circuit for \(H_{N}^{oplus}\) computes only the scalar inner product ; it need not reveal p\_x, a low-rank coordinate map, or any local SAT separator. The extraction is false in planted-label oracle worlds and would need real-SAT, nonrelativizing content.

\paragraph{Derivation steps.}
\begin{enumerate}
\item Rewrite the XOR carrier as \(H_{N}^{oplus}(x)=\) over \(F_{2}\).
\item Observe that a low-cost affine-incidence sketch gives a small circuit by hardwiring basis parities gamma\_i.
\item Try to prove the contrapositive: if all such sketches have \(A_{N}+K+E_{NN}>N^{1+\epsilon}\), then C(H\_N\textasciicircum{}oplus)>N\textasciicircum{}\{1+epsilon\}.
\item Audit the extraction direction from an arbitrary small H circuit to an affine sketch of p\_x.
\end{enumerate}

\paragraph{Equation 1: xor\_carrier.}
\[
H_N^{\oplus}(x)=\bigoplus_{j=1}^{M} SAT(G_N(x,j)),\quad M=N^q,\quad L=(\log N)^d.
\]
\emph{Valid when:} Valid when \(G_{N}\) outputs length-L SAT instances and the XOR is taken over \(F_{2}\) with nonadaptive slots.

\paragraph{Equation 2: incidence\_inner\_product.}
\[
H_N^{\oplus}(x)=\langle p_x,s\rangle_{\mathbb{F}_2},\quad p_x(y)=|{j:G_N(x,j)=y}|\bmod 2.
\]
\emph{Valid when:} Valid when \(U_{N}\) is the generated formula universe and s(y) records the real SAT label.

\paragraph{Equation 3: sketch\_representation.}
\[
p_x=\sum_{i=1}^{K} a_i(x)b_i \text{ for all } x\notin B_N,\quad |B_N|=E_N.
\]
\emph{Valid when:} Valid when coefficient circuits compute a\_i and exceptions are measured over N-bit inputs.

\paragraph{Equation 4: hardwired\_basis\_parities.}
\[
\gamma_i=\langle b_i,s\rangle_{\mathbb{F}_2},\quad H_N^{\oplus}(x)=\bigoplus_{i:\gamma_i=1} a_i(x) \text{ for } x\notin B_N.
\]
\emph{Valid when:} Valid when nonuniform circuits may hardwire the K basis parities as advice constants.

\paragraph{Equation 5: upper\_bound\_falsifier.}
\[
C(H_N^{\oplus})\le O(A_N+K+E_NN).
\]
\emph{Valid when:} Valid when bad inputs are patched by equality tests and fan-in two XOR gates are allowed.

\paragraph{Equation 6: source\_gate.}
\[
\forall \text{ affine-incidence sketches }\quad A_N+K+E_NN> N^{1+\epsilon-o(1)}.
\]
\emph{Valid when:} Valid when the claimed source lower bound is C(H\_N\textasciicircum{}oplus)>N\textasciicircum{}\{1+epsilon-o(1)\} in general nonuniform circuits.

\paragraph{Next analysis.} Instantiate \(G_{N,\theta}\) with linear, pairwise-independent, and extractor-style query maps; compute K, coefficient-circuit cost \(A_{N}\), duplicate parity cancellations, and whether any case survives the \(N^{1+\epsilon}\) affine-sketch falsifier threshold.

\paragraph{Falsifiers.}
\begin{itemize}
\item affine-incidence sketch with \(K+A_{N}+E_{NN}\) below \(N^{1+\epsilon-o(1)}\)
\item duplicate or even-multiplicity query generation causing \(p_{x}=0\) on many inputs
\item small generated universe with easy standard-basis coefficients, allowing all SAT labels on \(U_{N}\) to be hardwired
\item sparse correction set \(E_{N}<N^{\epsilon-o(1)}\) patched by equality tests
\item target threshold collapse when \(max{q,\rho}\ge 1+\epsilon\)
\item rank-only or entropy-only argument that ignores the actual SAT label vector
\item natural-proof risk if the affine-sketch exclusion is large and efficiently testable over random label tables rather than range-specific to SAT
\end{itemize}

\subsection*{cycle-0054: Hardness-Magnification Bridge Realization: XOR Carrier Approximate Affine-Sketch Gate}

\paragraph{Execution.} research; azure\_responses\_api; attempt primary; 150.117 seconds; 3 lemmas, 4 equations, 5 derivation steps, 6 falsifiers.

\paragraph{Learning.} Cycle 54 strengthens the XOR route with an approximate piecewise-affine incidence-sketch conservation law. Bad seeds cost EN gates to patch, and a sketch of total cost s+Kr+EN below \(N^{1+\epsilon}\) is a concrete falsifier for the source lower bound. High incidence rank alone is not valid lower-bound evidence because the SAT label functional may be simple on the generated image.

\paragraph{Pressure.} Need a direct lower bound for the SAT-label pullback x \(\to \) , or a non-natural real-SAT theorem showing that every subthreshold \(H_{N}\) circuit induces a subthreshold approximate affine-incidence sketch. Current rank and sketch audits are only falsifiers.

\paragraph{Derivation steps.}
\begin{enumerate}
\item Represent the XOR carrier as \(H_{N}(x)=\) over \(F_{2}\).
\item Try to prove \(C(H_{N})>N^{1+\epsilon}\) from the claim that A\_N(\{0,1\}\textasciicircum{}N) has no low-dimensional affine sketch.
\item Observe that any exact or approximate piecewise-affine sketch yields a small \(H_{N}\) circuit by hardwiring the SAT inner products of the affine bases.
\item Strengthen the necessary audit to piecewise-affine sketches with exceptional-set correction cost EN.
\item Test whether large rank alone suffices; it fails because the fixed SAT label vector may be simple on the generated image.
\end{enumerate}

\paragraph{Equation 1: incidence\_pullback.}
\[
H_N(x)=\langle A_N(x),Y_{SAT}\rangle_{\mathbb{F}_2}
\]
\emph{Valid when:} Valid when \(A_{N}\) records generated formulas with parity multiplicity and \(Y_{\mathrm{SAT}}\) is the fixed truth-label vector on L-bit formulas.

\paragraph{Equation 2: piecewise\_affine\_evaluation.}
\[
H_N(x)=\lambda_t\oplus\bigoplus_{i\le r} z_i\mu_{t,i},\quad \lambda_t=\langle a_t,Y_{SAT}\rangle,\ \mu_{t,i}=\langle b_{t,i},Y_{SAT}\rangle
\]
\emph{Valid when:} Valid when x is nonexceptional and the sketch representation of A\_N(x) holds inside the selected affine piece.

\paragraph{Equation 3: sketch\_upper\_bound.}
\[
C(H_N)\le O(s+Kr+K\log K+EN)
\]
\emph{Valid when:} Applies when the sketch circuit computes piece indices and coordinates and exceptional seeds are corrected by explicit equality tests.

\paragraph{Equation 4: source\_falsifier\_threshold.}
\[
s+Kr+EN\le N^{1+\epsilon-\Omega(1)}\ \Longrightarrow\ C(H_N)\not> N^{1+\epsilon-o(1)}
\]
\emph{Valid when:} Valid in the barely-superlinear source regime where lower bounds are measured against unrestricted nonuniform circuits of input length N.

\paragraph{Next analysis.} Instantiate the hash-based XOR generator and compute whether its incidence map has an explicit low-cost piecewise-affine cover; if none is found, formulate the exact extraction theorem needed to turn small \(H_{N}\) circuits into such covers.

\paragraph{Falsifiers.}
\begin{itemize}
\item falsifier: an exact affine incidence representation of rank \(r<N^{1+\epsilon-o(1)}\) makes \(H_{N}\) easy by hardwiring basis parities.
\item falsifier: an approximate piecewise-affine sketch with \(s+Kr+EN<N^{1+\epsilon-o(1)}\) kills the source lower bound after equality patching.
\item falsifier: a linear or low-degree algebraic generator whose parity-incidence map has computable rank O(N polylog N) collapses the XOR carrier to near-linear size.
\item falsifier: a SAT-monochromatic generated range makes \(H_{N}\) constant even if the incidence image has high rank.
\item falsifier: \(q\ge 1+\epsilon\) or \(\rho\ge 1+\epsilon\) causes quantitative target-threshold collapse.
\item falsifier: \(d\le 1\) with small lookup exponent may reintroduce exhaustive SAT-label leakage below the source threshold.
\end{itemize}

\subsection*{cycle-0055: Hardness-Magnification Bridge Realization}

\paragraph{Execution.} research; azure\_responses\_api; attempt primary; 288.009 seconds; 3 lemmas, 5 equations, 5 derivation steps, 5 falsifiers.

\paragraph{Learning.} Cycle 55 adds the SAT-preserving quotient normal form for XOR carriers. The valid target bridge is still \(max{\rho,q}<1+\epsilon\), but the source lower bound must beat quotient piecewise-affine incidence sketches, not raw range or rank. Symmetry-pair cancellation is a concrete falsifier that can make a high-entropy generator compute a constant carrier; this is route-level progress only, not a proof of \(\mathrm{P} \ne \mathrm{NP}\).

\paragraph{Pressure.} Label-sensitive extraction gap: a small arbitrary circuit for the single function x \(\to \)  need not reveal a small quotient-incidence sketch, and high quotient rank alone does not preclude the SAT-labelled linear functional from being simple.

\paragraph{Derivation steps.}
\begin{enumerate}
\item Write the XOR carrier as the \(F_{2}\) inner product of the SAT truth vector \(tau_{L}\) with the generated incidence vector a\_N(x).
\item Prove the polylog substitution upper bound: SAT circuits of size \(n^{c}\) yield \(H_{N}\) circuits of size \(N^{max{\rho,q}+o(1)}\).
\item Try to certify \(C(H_{N})>N^{1+\epsilon}\) from high rank or expansion of the incidence map x \(\to \) a\_N(x).
\item Normalize by quotienting formula space under SAT-preserving symmetries and by cancelling even multiplicities.
\item Observe that a piecewise affine quotient sketch of cost below \(N^{1+\epsilon}\) directly computes \(H_{N}\) after hardwiring basis parities and patching sparse exceptions.
\end{enumerate}

\paragraph{Equation 1: xor\_incidence\_identity.}
\[
H_N(x)=\bigoplus_{j=1}^{m}\mathrm{SAT}(\phi_{N,j}(x))=\langle \tau_L,a_N(x)\rangle_{\mathbb{F}_2}
\]
\emph{Valid when:} Valid when \(G_{N}\) outputs syntactic formula encodings and repeated formulas are counted with parity over \(F_{2}\).

\paragraph{Equation 2: polylog\_substitution\_bound.}
\[
\mathrm{SAT}\in \mathrm{SIZE}(n^c)\ \Rightarrow\ C(H_N)\le O(C(G_N)+m\,L^c+mL)=N^{\max{\rho,q}+o(1)}
\]
\emph{Valid when:} Applies after SAT has \(n^{c}-size\) circuits and the generator has \(N^\rho-size\) circuits with \(m=N^{q}\) polylog-length queries.

\paragraph{Equation 3: target\_magnification\_threshold.}
\[
\max{\rho,q}< 1+\epsilon\ \wedge\ C(H_N)> N^{1+\epsilon-o(1)}\ \Rightarrow\ \mathrm{SAT}\notin\mathrm{P/poly}
\]
\emph{Valid when:} Valid for infinitely many N against every fixed target exponent c because L=(log N)\textasciicircum{}d makes L\textasciicircum{}c=N\textasciicircum{}\{o(1)\}.

\paragraph{Equation 4: sat\_preserving\_quotient\_identity.}
\[
H_N(x)=\langle \tau_L^{\sim},a_N^{\sim}(x)\rangle_{\mathbb{F}_2}
\]
\emph{Valid when:} Requires the equivalence relation to preserve SAT labels on every quotient class used by the generated formulas.

\paragraph{Equation 5: quotient\_sketch\_upper\_bound.}
\[
C(H_N)\le O(\mathsf{QSk}_{\sim,N}(E)+E N)
\]
\emph{Valid when:} Requires a piecewise affine sketch of the quotient incidence vector correct outside E named seeds and allows nonuniform hardwiring of the resulting SAT parities.

\paragraph{Next analysis.} Compute QSk\_\{sim,N\}(E) for quotient relations generated by equality, variable renaming, clause permutation, tautological padding, and duplicate-pair involutions on the current high-q XOR carrier template; record the smallest residual exponent u and pairing-certificate cost.

\paragraph{Falsifiers.}
\begin{itemize}
\item falsifier: an easy SAT-preserving quotient collapses the odd incidence vectors to subthreshold piecewise affine dimension while raw formula range remains high.
\item falsifier: fixed-point-free duplicate, variable-renaming, or clause-permutation pairings cancel most XOR slots before any SAT hardness is used.
\item falsifier: \(QSk_{sim,N}(E)+EN<N^{1+\epsilon-o(1)}\) for some \(E<N^{\epsilon-o(1)}\) gives an explicit nonuniform circuit for \(H_{N}\).
\item falsifier: high quotient rank is aligned with a low-cost SAT-labelled linear functional, so rank-only evidence is not valid source lower-bound evidence.
\item falsifier: a proposed quotient relation or sketch uses SAT truth-table advice inside \(G_{N}\), causing nonuniform leakage and breaking the target-clean interface.
\end{itemize}

\subsection*{cycle-0056: Hardness-Magnification Bridge Realization: XOR quotient codebook gate}

\paragraph{Execution.} research; azure\_responses\_api; attempt primary; 180.382 seconds; 3 lemmas, 4 equations, 5 derivation steps, 7 falsifiers.

\paragraph{Learning.} Cycle 56 adds a sharper conservation law for XOR carriers: after SAT-preserving parity quotienting, \(H_{N}\) is easy whenever the incidence map has either low affine rank or a low-cardinality computable transcript, with patchable errors charged by EN. Thus high quotient entropy or high rank alone is not valid source evidence; the source obligation is a combined codebook-or-rank minimax lower bound, while the target bridge remains unchanged.

\paragraph{Pressure.} The derivation does not show that every small arbitrary circuit for \(H_{N}\) yields such a quotient sketch. It is an upper-bound falsifier, not a lower-bound certificate.

\paragraph{Derivation steps.}
\begin{enumerate}
\item Normalize each generated multiset by deleting pairs with equal SAT-preserving quotient class; the remaining parity-incidence vector is a\_N(x).
\item For an affine cell, fix a basis \(b_{i,1},...,b_{i,r_{i}}\) and hardwire the bits  plus the affine offset bit; compute the coordinate decoder and XOR the selected constants.
\item For a codebook cell, hardwire \(p_{i,t}=\) for every transcript value t in [R\_i]; compute t\_i(x) and evaluate a table circuit of size O(R\_i log \(R_{i}\)).
\item Use the cheaper implementation on each cell, multiplex by the cell locator, and patch each of the E uncovered seeds by an equality test on the N-bit input.
\item The resulting circuit size is O(S+sum\_i min\{r\_i,R\_i log R\_i\}+EN), so any sketch below \(N^{1+\epsilon}\) contradicts the desired source lower bound.
\end{enumerate}

\paragraph{Equation 1: xor-quotient-normal-form.}
\[
H_N(x)=\langle \lambda_N,a_N(x)\rangle_{\mathbb{F}_2},\qquad a_N(x)\in\mathbb{F}_2^{Q_N}.
\]
\emph{Valid when:} Valid when generated formulas are reduced only by odd parity cancellation and by a quotient preserving SAT labels.

\paragraph{Equation 2: codebook-rank-upper-bound.}
\[
C(H_N)\le O\!\left(S+\sum_i \min{r_i, R_i\log R_i}+EN\right).
\]
\emph{Valid when:} Applies when the sketch covers all but E seeds and each bad seed is patchable by an equality circuit of size O(N).

\paragraph{Equation 3: target-substitution-bound.}
\[
SAT\in SIZE(n^c)\Rightarrow C(H_N)\le N^\rho+N^q L^c+N^q\operatorname{polylog}L=N^{\max{\rho,q}+o(1)}.
\]
\emph{Valid when:} Requires polylogarithmic query length, fixed target exponent c, and a generator circuit of size \(N^\rho\).

\paragraph{Equation 4: magnification-threshold.}
\[
C(H_N)> N^{1+\epsilon}\ \wedge\ \max{\rho,q}< 1+\epsilon\Rightarrow \forall c\; SAT\notin SIZE(n^c).
\]
\emph{Valid when:} Valid when the source lower bound holds for infinitely many large N and the query lengths L tend to infinity.

\paragraph{Next analysis.} Instantiate a candidate high-q generator and compute the exponent of its SAT-preserving odd-support transcript range versus affine span rank; reject it if min\{theta,rank-exponent\}+locator-cost falls below 1+epsilon.

\paragraph{Falsifiers.}
\begin{itemize}
\item Subthreshold quotient codebook sketch: \(S+sum_{i}\) \(R_{i}\) log \(R_{i}+EN<N^{1+\epsilon-o(1)}\) makes \(H_{N}\) easy by hardwired parity values.
\item Low affine rank after SAT-preserving quotient: \(S+sum_{i}\) \(r_{i}+EN<N^{1+\epsilon-o(1)}\) makes \(H_{N}\) easy by hardwired basis parities.
\item Symmetry-pair cancellation: most generated SAT instances cancel modulo 2 inside \(Q_{N}\), causing constant or low-transcript carriers despite high raw formula entropy.
\item Target threshold collapse: \(max{\rho,q}\ge 1+\epsilon\) invalidates the magnification implication even if the source lower bound is true.
\item Model mismatch falsifier: quotienting formulas by a relation that is not SAT-preserving breaks the identity \(H_{N}=\).
\item Nonuniform leakage through a short transcript table may decide \(H_{N}\) without computing SAT on any new formula.
\item Natural-proof falsifier: a proof using only random-label hardness of \(lambda_{N}\) would not identify real SAT and is not valid route evidence under standard pseudorandomness barriers.
\end{itemize}

\subsection*{cycle-0057: Hardness-Magnification Bridge Realization: XOR quotient carrier, low-degree feature-sketch gate}

\paragraph{Execution.} research; azure\_responses\_api; attempt primary; 171.561 seconds; 3 lemmas, 5 equations, 5 derivation steps, 5 falsifiers.

\paragraph{Learning.} Cycle 57 adds a nonlinear conservation law for the XOR carrier. The source lower bound must beat low-degree quotient-feature sketches with threshold \(S_{T}+binom(r,\le t)+EN\), so high rank or high codebook entropy is not valid evidence. The target-preserving bridge is unchanged, but the source obligation is narrower and more nonrelativizing.

\paragraph{Pressure.} No method is known to prove that the real SAT-labelled quotient functional avoids all subthreshold low-degree feature sketches. Random-incidence or random-label counting does not transfer to SAT.

\paragraph{Derivation steps.}
\begin{enumerate}
\item Start with the quotient incidence representation \(H_{N}(z)=\) over \(F_{2}\), where \(y_{N}\) is the SAT label vector on generated quotient classes.
\item Cycle 56 handled the case A\_N(z) lies in a low-rank affine subspace or in a small computable codebook; then hardwired basis parities or table entries compute \(H_{N}\).
\item Generalize: suppose a small circuit \(T_{N}\) computes \(r_{N}\) features and the composed label functional  equals a degree-t polynomial P\_N(T\_N(z)) except on \(E_{N}\) inputs.
\item Implement \(T_{N}\) once, hardwire all coefficients of \(\mathrm{P}_{N}\), compute all degree-at-most-t monomials, XOR selected monomials, and patch the \(E_{N}\) exceptional inputs by equality tests.
\item This yields \(C(H_{N})N^{1+\epsilon}\) is falsified whenever max\{s,alpha t,1+eta\}<1+epsilon.
\end{enumerate}

\paragraph{Equation 1: xor\_carrier.}
\[
H_N^{xor}(z)=\bigoplus_{j=1}^{\lfloor N^q\rfloor} SAT(F_{N,j}(z))
\]
\emph{Valid when:} Valid when the carrier is the repaired XOR aggregation route and every generated query has length L=(log N)\textasciicircum{}d.

\paragraph{Equation 2: target\_substitution.}
\[
SAT\in SIZE(n^c)\Rightarrow C(H_N^{xor})\le N^\rho+N^q\,L^{c+O(1)}=N^{\max{\rho,q}+o(1)}
\]
\emph{Valid when:} Applies after the generator size is \(N^\rho\), q is fixed below the threshold, and L is polylogarithmic in N.

\paragraph{Equation 3: magnification\_implication.}
\[
C(H_N^{xor})> N^{1+\epsilon}\ \wedge\ \max{\rho,q}< 1+\epsilon\Rightarrow SAT\notin P/poly
\]
\emph{Valid when:} Requires the target-preserving complexity-class transfer for every fixed SAT circuit exponent c with no advice leakage.

\paragraph{Equation 4: feature\_sketch\_upper\_bound.}
\[
C(H_N^{xor})\le S_T+O\!\left(\sum_{i=0}^{t}\binom{r_N}{i}\right)+O(E_NN)
\]
\emph{Valid when:} Valid when \(H_{N}\) agrees off \(E_{N}\) inputs with a degree-t polynomial in the computed feature vector T\_N(z).

\paragraph{Equation 5: constant\_degree\_threshold.}
\[
r_N=N^\alpha,\ S_T=N^s,\ E_N=N^\eta\Rightarrow C(H_N^{xor})\le N^{\max{s,\alpha t,1+\eta}+o(1)}
\]
\emph{Valid when:} Applies for constant t and polynomial-size parameters so the monomial count has exponent \(\alpha\) times t.

\paragraph{Next analysis.} Formalize the degree-2 case: characterize whether Tseitin-style parity, formula-size, variable-incidence, or clause-type features give \(r=N^\alpha\) with \(2alpha<1+\epsilon\) for the generated length-(log N)\textasciicircum{}d SAT queries.

\paragraph{Falsifiers.}
\begin{itemize}
\item A degree-t quotient-feature sketch with \(r=N^\alpha\) and \(\alpha\) \(t<1+\epsilon\) is a concrete falsifier for the source lower bound.
\item An error set of size \(E<N^\epsilon\) is patchable by O(EN) equality gates and cannot be ignored.
\item A SAT-label annihilator making  depend on few nonlinear features defeats high raw quotient rank.
\item A proof that only shows random labels avoid the sketch class is not valid for real SAT without a nonrelativizing transfer certificate.
\item Any hidden hardwiring of the SAT truth table into \(T_{N}\) is nonuniform leakage and invalidates the claimed source model.
\end{itemize}

\subsection*{cycle-0058: Hardness-magnification bridge via XOR carriers and sparse quotient-polynomial falsifiers}

\paragraph{Execution.} research; azure\_responses\_api; attempt primary; 140.312 seconds; 3 lemmas, 4 equations, 5 derivation steps, 6 falsifiers.

\paragraph{Learning.} Cycle 58 repairs the XOR-carrier audit: low degree was the wrong nonlinear measure. The valid upper-bound falsifier is sparse quotient-polynomial support with threshold \(S_{T}+A_{N}\) \(r_{N}+E_{N}\) \(N<N^{1+\epsilon}\). This preserves the polylog-query target bridge but narrows the source obligation and exposes a new extraction gap: sparse-sketch hardness is not yet general-circuit hardness.

\paragraph{Pressure.} No valid derivation converts an arbitrary small \(H_{N}\) circuit into a charged sparse quotient-polynomial sketch. Without this extraction, sketch lower bounds are falsifier tests, not source lower bounds.

\paragraph{Derivation steps.}
\begin{enumerate}
\item Start from the quotient normal form \(H_{N}=\) over \(F_{2}\) for the XOR carrier.
\item Observe that cycle 57's low-degree feature bound gives \(C(H_{N})\le S_{T}+binom(r,\le t)+E_{N}\) N when the final combiner has degree t.
\item Test the converse direction: high degree does not imply high circuit cost, because a degree-r monomial needs only O(r) gates after feature computation.
\item Replace the degree parameter by monomial support \(A_{N}\) and derive C(H\_N)<=S\_T+O(A\_N \(r_{N}+E_{N}\) N).
\item Check target preservation separately: \(\mathrm{SAT}\in \mathrm{SIZE}(n^{c})\) still gives \(C(H_{N})\le N^{max{\rho,q}+o(1)}\) for polylogarithmic queries.
\end{enumerate}

\paragraph{Equation 1: quotient-inner-product-normal-form.}
\[
H_N(x)=\bigoplus_{u\in U_N} a_u I_{N,u}(x)=\langle a_N,I_N(x)\rangle_{\mathbb{F}_2}.
\]
\emph{Valid when:} Valid when \(U_{N}\) is a SAT-preserving quotient of the generated formulas and \(I_{N}\) records parity incidence per quotient class.

\paragraph{Equation 2: target-substitution-upper-bound.}
\[
SAT\in SIZE(n^c)\ \Rightarrow\ C(H_N)\le N^{\rho+o(1)}+N^q L_N^{c+O(1)}=N^{\max{\rho,q}+o(1)}.
\]
\emph{Valid when:} Applies when SAT on length \(L_{N}\) formulas has nonuniform size \(L_{N}^{c}\) and the generator is target-clean.

\paragraph{Equation 3: sparse-sketch-upper-bound.}
\[
H_N=P_N(T_N)\text{ off }B_N,\ |B_N|=E_N\ \Rightarrow\ C(H_N)\le S_T+O(A_N r_N+E_N N).
\]
\emph{Valid when:} Requires that all feature circuits are charged once, monomial gates share them, and exceptions are patched by equality tests.

\paragraph{Equation 4: exponent-falsifier-threshold.}
\[
S_T\le N^\sigma,\ A_N\le N^\alpha,\ r_N\le N^\gamma,\ E_N\le N^\eta\ \Rightarrow\ C(H_N)\le N^{\max{\sigma,\alpha+\gamma,1+\eta}+o(1)}.
\]
\emph{Valid when:} Valid when the parameter exponents are asymptotic and \(r_{N}\) is at most \(N^{\gamma+o(1)}\).

\paragraph{Next analysis.} Formalize the weakest nontrivial extraction candidate: restrict arbitrary \(H_{N}\) circuits to circuits whose wires are invariant under SAT-preserving quotient symmetries, and test whether such invariant circuits collapse to sparse quotient-polynomial sketches with subthreshold loss.

\paragraph{Falsifiers.}
\begin{itemize}
\item Sparse quotient-polynomial sketch with \(S_{T}+A_{N}\) \(r_{N}+E_{N}\) \(N<N^{1+\epsilon-o(1)}\).
\item High-degree but one-monomial combiner, showing that degree-only hardness evidence is not valid.
\item Uncharged feature map that directly computes \(H_{N}\), collapsing the audit into a tautology rather than a route-specific falsifier.
\item Target threshold collapse when \(max{\rho,q}\ge 1+\epsilon\) or when \(L_{N}\) is not polylogarithmic.
\item Nonuniform leakage of the SAT label vector into the generator or feature circuits without being charged.
\item Relativizing planted-control oracle where \(H_{N}\) is easy while quotient-sketch or separator measures remain hard.
\end{itemize}

\subsection*{cycle-0059: Hardness-Magnification Bridge Realization: XOR carrier sparse-sketch audit}

\paragraph{Execution.} research; azure\_responses\_api; attempt primary; 296.472 seconds; 3 lemmas, 3 equations, 6 derivation steps, 6 falsifiers.

\paragraph{Learning.} Cycle 59 identifies a conservation law for the XOR-carrier audit: unrestricted sparse quotient-polynomial hardness is not a simplification of the source lower bound, since \(T=H_{N}\) is an admissible one-bit quotient transcript. The valid contribution is a sharper quarantine: sparse sketches remain concrete upper-bound falsifiers, while any proof-relevant sketch lower bound now needs a named label-free transcript class plus a nonrelativizing extraction theorem. The main \(\mathrm{P} \ne \mathrm{NP}\) conjecture is untouched.

\paragraph{Pressure.} No nonrelativizing locality-extraction theorem is known that converts an arbitrary small circuit for the real SAT-XOR carrier into a subthreshold label-free incidence transcript; relativizing planted-control examples block the naive argument.

\paragraph{Derivation steps.}
\begin{enumerate}
\item Start with the XOR carrier \(H_{N}(x)=\) over \(F_{2}\) and keep the stable target bridge \(max{\rho,q}<1+\epsilon\).
\item Try to replace the direct general-circuit lower bound \(C(H_{N})>N^{1+\epsilon}\) by hardness against sparse quotient-polynomial sketches.
\item Any sketch with transcript cost \(S_{T}\), support A, dimension r, and E errors gives C(H\_N)<=S\_T+O(A r+E N). This is a valid upper-bound falsifier.
\item For the reverse direction, if unrestricted quotient transcripts are allowed, choose T(x)=H\_N(x), \(r=1\), and \(\mathrm{P}(z)=z\). Then SQP\_all\_N(H\_N)<=C(H\_N)+O(1).
\item Thus unrestricted sparse-sketch hardness is equivalent to the original source lower bound and does not reduce the proof burden.
\item Restricting transcripts to label-free incidence features makes the falsifier meaningful, but the derivation then needs an extraction theorem from arbitrary small \(H_{N}\) circuits to those features.
\end{enumerate}

\paragraph{Equation 1: xor-carrier-incidence-form.}
\[
H_N(x)=\bigoplus_{j\le m_N} SAT(G_N(x,j))=\langle a_N(x),y_N\rangle_{\mathbb{F}_2}
\]
\emph{Valid when:} Valid when the generator outputs only length \(L_{N}\) SAT formulas and duplicate formulas are counted modulo two.

\paragraph{Equation 2: sparse-sketch-upper-bound.}
\[
C(H_N)\le S_T+O(A r+E N)
\]
\emph{Valid when:} Valid when P has A monomials over r transcript bits and differs from \(H_{N}\) on at most E inputs.

\paragraph{Equation 3: unrestricted-sqp-equivalence.}
\[
SQP^{all}_N(H_N)\le C(H_N)+O(1)\quad\text{and}\quad C(H_N)\le O(SQP^{all}_N(H_N))
\]
\emph{Valid when:} Valid when the audit permits arbitrary quotient-respecting transcripts, including the one-bit transcript T(x)=H\_N(x).

\paragraph{Next analysis.} Formalize the local-linear incidence transcript class for the XOR SAT carrier and compute the minimum sparse \(F_{2}\) support needed under a random-generator model, while auditing whether any extraction from arbitrary circuits to that class can avoid the planted-control relativization falsifier.

\paragraph{Falsifiers.}
\begin{itemize}
\item tautological transcript falsifier: \(T=H_{N}\), \(r=1\), \(\mathrm{P}(z)=z\) collapses any unrestricted sparse-sketch audit
\item structured sparse-polynomial falsifier: \(S_{T}+A\) r+E \(N<N^{1+\epsilon-o(1)}\) gives an explicit upper bound for \(H_{N}\)
\item nonuniform leakage through SAT-label-dependent transcript features
\item relativizing planted-control oracle blocking locality extraction
\item threshold collapse if max\{rho,q\} reaches 1+epsilon
\item model mismatch if the source lower bound is proved only for structured sketches but advertised as a general-circuit lower bound
\end{itemize}

\subsection*{cycle-0060: Constructive Label-Free Sketch Gate for XOR Magnification}

\paragraph{Execution.} research; azure\_responses\_api; attempt primary; 161.208 seconds; 3 lemmas, 4 equations, 5 derivation steps, 7 falsifiers.

\paragraph{Learning.} Cycle 60 separates two notions previously conflated. Sparse quotient-polynomial sketches are valid upper-bound falsifiers only when the transcript is constructively label-free; otherwise \(T_{N}=H_{N}\) collapses the audit. The target bridge \(max{\rho,q}<1+\epsilon\) survives, but sketch-hardness is not a source lower bound without a nonrelativizing extraction theorem or a direct C(H\_N) lower bound.

\paragraph{Pressure.} No theorem converts an arbitrary small circuit for the parity of SAT labels on generated queries into a constructively label-free low-support incidence sketch with subthreshold residual cost.

\paragraph{Derivation steps.}
\begin{enumerate}
\item Write \(H_{N}\) as the \(F_{2}\) inner product of the generated incidence vector a\_N(x) with the SAT label vector \(y_{N}\).
\item Assume a subthreshold circuit \(D_{N}\) of size \(N^{1+\epsilon-\delta}\) computes \(H_{N}\) and try to use its gate transcript as \(T_{N}\) with \(\mathrm{P}_{N}\) equal to the final gate.
\item This immediately violates label-freeness: \(D_{N}\) may encode information about \(y_{N}\) or about the carrier truth table, even if it never explicitly queries SAT.
\item Try expanding \(D_{N}\) as an algebraic normal form over seed bits or incidence coordinates; the support can blow up to \(2^{S}\) and gives no \(N^{1+\epsilon}\) threshold preservation.
\item Try quotienting incidences by syntactic SAT-preserving symmetries; this gives useful upper-bound falsifiers but not an extraction from arbitrary circuits.
\end{enumerate}

\paragraph{Equation 1: xor-incidence-form.}
\[
H_N(x)=\bigoplus_{j=1}^{m_N}\mathrm{SAT}(F_{N,j}(x))=\langle a_N(x),y_N\rangle_{\mathbb{F}_2},\quad m_N=N^q,\quad L_N=(\log N)^d.
\]
\emph{Valid when:} Valid when \(G_{N}\) outputs \(m_{N}\) \(length-L_{N}\) formulas and \(a_{N}\) records their mod-two incidences in the generated universe.

\paragraph{Equation 2: label-free-sketch-upper-bound.}
\[
C(H_N)\le S_T+O(A_N r_N+E_N N).
\]
\emph{Valid when:} Valid when \(T_{N}\) is constructively label-free, \(\mathrm{P}_{N}\) has \(A_{N}\) monomials, and equality gates patch exactly \(E_{N}\) exceptional seeds.

\paragraph{Equation 3: target-substitution-upper-bound.}
\[
\mathrm{SAT}\in \mathrm{SIZE}(n^c)\ \Longrightarrow\ C(H_N)\le N^{\rho+o(1)}+N^q L_N^c=N^{\max{\rho,q}+o(1)}.
\]
\emph{Valid when:} Applies after the generator has size \(N^{\rho+o(1)}\) and every generated query has polylogarithmic length \(L_{N}\).

\paragraph{Equation 4: magnification-threshold.}
\[
C(H_N)> N^{1+\epsilon}\ \text{i.o.}\ \land\ \max{\rho,q}< 1+\epsilon\ \Longrightarrow\ \mathrm{SAT}\notin \mathrm{P/poly}.
\]
\emph{Valid when:} Requires the lower bound to hold infinitely often against nonuniform circuits and fixed q,rho below the threshold.

\paragraph{Next analysis.} Define the first concrete compiler class \(T^{inc-lin}\): \(T_{N}(x)=B_{N}\) a\_N(x) for explicit SAT-preserving quotient matrices \(B_{N}\) of row sparsity polylog N, then compute the minimum \(S_{T}+A_{N}\) \(r_{N}\) for representing \(H_{N}\) by sparse \(\mathrm{P}_{N}\) in the high-q, \(d=2\) regime.

\paragraph{Falsifiers.}
\begin{itemize}
\item Output-dependent transcript leakage: \(T_{N}=H_{N}\) makes sparse-sketch hardness vacuous.
\item SAT-label advice in the transcript compiler below the tested threshold invalidates label-freeness.
\item Constructively label-free \(T_{N}\) with \(S_{T}+A_{N}\) \(r_{N}+E_{N}\) \(N<N^{1+\epsilon-o(1)}\) is a direct source lower-bound falsifier.
\item Symmetry-pair cancellation after SAT-preserving quotienting can make a high-entropy XOR carrier constant or sparse.
\item Threshold collapse occurs if \(max{\rho,q}\ge 1+\epsilon\), since SAT substitution no longer contradicts the source lower bound.
\item Lookup collapse reappears for \(d\le 1\) when the generated truth-table cost can fall below the claimed source threshold.
\item A relativizing extraction attempt is blocked if small \(H_{N}\) circuits use oracle/advice correlations not recoverable from incidence structure.
\end{itemize}

\subsection*{cycle-0061: Hardness-Magnification Bridge Realization}

\paragraph{Execution.} research; azure\_responses\_api; attempt primary; 94.538 seconds; 2 lemmas, 5 equations, 6 derivation steps, 5 falsifiers.

\paragraph{Learning.} Cycle 61 upgrades the XOR-carrier audit from vague label-free transcripts to a costed global-codebook normal form. The valid new contribution is an explicit upper-bound falsifier \(C(H_{N})\le S_{T}+S_{D}+\)\(\widetilde{O}(U_{N}+A_{Nr,N})\)\(+E_{NN}\). This closes a leakage loophole in sparse quotient sketches, while exposing the exact remaining gap: absence of label-free sketches is not a general-circuit lower bound without a nonrelativizing extraction theorem.

\paragraph{Pressure.} No extraction theorem converts arbitrary \(H_{N}\) circuits into constructively label-free codebook sketches. The trivial transcript \(T_{N}=H_{N}\) is small whenever \(H_{N}\) is small but is not label-free and may encode SAT-label information semantically.

\paragraph{Derivation steps.}
\begin{enumerate}
\item Fix the surviving escape regime L\_N=(log N)\textasciicircum{}d with \(d>1\), \(M_{N}=N^{q}\), and \(max{\rho,q}<1+\epsilon\).
\item Substitute a hypothetical size \(n^{c}\) SAT circuit into all \(N^{q}\) query slots; since L\_N\textasciicircum{}c=N\textasciicircum{}\{o(1)\}, obtain \(C(H_{N})\le N^{max{\rho,q}+o(1)}\).
\item Define \(LFSketch_{N}\) by an explicit codebook of \(U_{N}\) \(length-L_{N}\) formulas, a label-free transcript \(T_{N}\), and a decoder producing a sparse polynomial in hardwired SAT labels of codebook entries.
\item Build a nonuniform circuit for \(H_{N}\) by computing \(T_{N}\), evaluating the decoder, looking up hardwired codebook labels, evaluating at most \(A_{N}\) monomials, and patching \(E_{N}\) exceptional source inputs.
\item Conclude that any \(LFSketch_{N}\) with \(S_{T}+S_{D}+\)\(\widetilde{O}(U_{N}+A_{Nr,N})\)\(+E_{NN}<N^{1+\epsilon}\) is a concrete upper-bound falsifier for the desired source lower bound.
\item Attempt to prove the converse: small arbitrary circuits for \(H_{N}\) should yield small label-free sketches, making LFSketch hardness a source lower-bound certificate.
\end{enumerate}

\paragraph{Equation 1: xor\_carrier\_definition.}
\[
H_N(x)=\bigoplus_{j=1}^{N^q} SAT(\Phi_{N,j}(x)),\qquad |\Phi_{N,j}(x)|=L_N=(\log N)^d.
\]
\emph{Valid when:} Valid when the source object is the generated XOR-SAT carrier with fixed q,d and generator outputs well-formed SAT instances for every valid N.

\paragraph{Equation 2: polylog\_target\_substitution.}
\[
SAT\in SIZE(n^c)\ \Longrightarrow\ C(H_N)\le N^\rho+N^q L_N^c\,polylog(N)=N^{\max{\rho,q}+o(1)}.
\]
\emph{Valid when:} Applies after the controlled interface substitutes one nonuniform SAT circuit of length \(L_{N}\) into every generated query slot with \(L_{N}\) subpolynomial.

\paragraph{Equation 3: magnification\_window.}
\[
C(H_N)> N^{1+\epsilon}\ i.o.\ \wedge\ \max{\rho,q}< 1+\epsilon\ \Longrightarrow\ \forall c\ C(SAT_n)\not\le n^c.
\]
\emph{Valid when:} Requires the target-preserving condition, cofinal valid source lengths, and a source lower bound exceeding the substitution upper bound for each fixed c.

\paragraph{Equation 4: label\_free\_codebook\_upper\_bound.}
\[
C(H_N)\le S_T+S_D+\widetilde{O}(U_N+A_N r_N)+O(E_NN).
\]
\emph{Valid when:} Valid when a constructively label-free transcript and decoder express \(H_{N}\) outside \(E_{N}\) bad inputs as a sparse polynomial in hardwired SAT labels of the enumerated codebook.

\paragraph{Equation 5: source\_falsifier\_threshold.}
\[
S_T+S_D+\widetilde{O}(U_N+A_Nr_N)+E_NN< N^{1+\epsilon-o(1)}\ \Longrightarrow\ C(H_N)< N^{1+\epsilon}.
\]
\emph{Valid when:} Valid in the nonuniform source model where codebook labels may be hardwired but the transcript construction itself contains no SAT-label advice.

\paragraph{Next analysis.} Instantiate a candidate generator and compute \(U_{N},A_{N},r_{N},E_{N}\) for its SAT-preserving quotient codebook; if the cost is subthreshold, quarantine that generator, otherwise formulate the exact extraction theorem it would require.

\paragraph{Falsifiers.}
\begin{itemize}
\item label-dependent one-bit transcript falsifier: omitting label-freeness admits T\_N(x)=H\_N(x) and destroys the sketch audit.
\item small global codebook falsifier: \(U_{N}=N^{1+\epsilon-\omega(1)}\) with cheap decoder lets the circuit hardwire all relevant SAT labels.
\item patchable-error falsifier: \(E_{N}<N^{\epsilon-o(1)}\) bad source inputs cost only \(E_{NN}<N^{1+\epsilon-o(1)}\) gates.
\item positive-beta rollback falsifier: if \(L_{N}=N^\beta\) with \(\beta>0\) then \(\mathrm{SAT}\in \mathrm{SIZE}(n^{c})\) gives C(H\_N)<=N\textasciicircum{}\{q+beta c+o(1)\}, so only bounded c are excluded.
\item nonuniform leakage falsifier: generator or transcript advice that depends on SAT labels can compress the carrier while preserving apparent range entropy.
\end{itemize}

\subsection*{cycle-0062: Hardness-Magnification Bridge Realization}

\paragraph{Execution.} research; azure\_responses\_api; attempt primary; 600.009 seconds; 2 lemmas, 5 equations, 6 derivation steps, 5 falsifiers.

\paragraph{Learning.} Cycle 62 adds a new fork: high incidence rank can make the XOR carrier lower-bound property large, so an efficient semantic certificate would be natural, while low rank is an upper-bound falsifier. The polylog \(\beta=0\) bridge remains valid, but the surviving source proof must be non-natural, non-large, proof-mediated, or direct.

\paragraph{Pressure.} The derivation does not produce the actual source lower bound. It only says that high-rank, efficiently checkable, large certificates are barrier-exposed, while absence of label-free sketches is still not equivalent to arbitrary-circuit hardness without a nonrelativizing extraction theorem.

\paragraph{Derivation steps.}
\begin{enumerate}
\item Fix \(U_{N}\) as the L-bit formula universe touched by \(G_{N}\) and write \(H_{N,G,y}=A_{N}\) y over \(F_{2}\) on the truth-table rows indexed by x in \({0,1}^{N}\).
\item For any Boolean function h on \({0,1}^{N}\), the set of labels y satisfying \(A_{N}\) \(y=h\) is empty or an affine subspace of codimension \(r_{N}\), hence has uniform measure \(2^{-r_{N}}\).
\item Union-bound over all size-S source circuits h. Using \#SIZE(S)<=2\textasciicircum{}\{K S log S\} gives Pr\_y[C(H\_\{N,G,y\})<=S]<=2\textasciicircum{}\{K S log S-r\_N\}.
\item If \(r_{N}\) exceeds K S log S by \(L^{k}\), the lower-bound property \(C(H)>S\) is overwhelmingly large over L-bit truth tables.
\item If \(\mathrm{SAT}_{L}\) had \(L^{c}-size\) circuits, the valid polylog substitution gives \(C(H_{N,G,\mathrm{SAT}})N^{1+\epsilon-o(1)}\) rejects every \(L^{c}-size\) truth table when \(q,\rho<1+\epsilon\).
\item Thus a \(\mathrm{poly}(2^{L})-constructive\) verifier accepting the actual SAT truth table and certifying \(R_{N}\) would be a natural property against SIZE(L\textasciicircum{}c), conditional on the usual PRF barrier assumptions.
\end{enumerate}

\paragraph{Equation 1: xor\_carrier\_linearization.}
\[
H_{N,G,y}(x)=\langle a_N(x),y\rangle_{\mathbb{F}_2}=\bigoplus_{j=1}^{N^q} y_{G_N(x,j)}
\]
\emph{Valid when:} Valid when \(G_{N}\) outputs only L-bit formula encodings and \(a_{N}\) records queried positions with parity multiplicity over \(F_{2}\).

\paragraph{Equation 2: polylog\_target\_substitution.}
\[
\mathrm{SAT}_L\in \mathrm{SIZE}(L^c)\Rightarrow C(H_{N,G,\mathrm{SAT}})\le O(N^\rho+N^q L^c\operatorname{polylog}N)=N^{\max{\rho,q}+o(1)}
\]
\emph{Valid when:} Applies when L=(log N)\textasciicircum{}d, c is fixed, target circuits are nonuniform, and the generator has size \(N^\rho\) without SAT-label advice.

\paragraph{Equation 3: direct\_circuit\_density\_bound.}
\[
\Pr_y[C(H_{N,G,y})\le S]\le 2^{K S\log S-r_N}
\]
\emph{Valid when:} Valid when y is uniform over L-bit truth tables, \(r_{N}\) is the \(F_{2}\) rank of the incidence rows, and size-S circuits are counted by 2\textasciicircum{}\{K S log S\}.

\paragraph{Equation 4: largeness\_threshold.}
\[
r_N\ge K S\log S+L^k\Rightarrow \Pr_y[C(H_{N,G,y})> S]\ge 1-2^{-L^k}
\]
\emph{Valid when:} Requires the incidence-rank gap to dominate the circuit-counting entropy and chooses k within the remaining explicit rank margin.

\paragraph{Equation 5: naturalness\_fork.}
\[
\mathrm{Constructive}(R_N)\wedge \mathrm{Large}(R_N)\wedge \mathrm{Useful}_{L^c}(R_N)\Rightarrow R_N\text{ is a natural property against }\mathrm{SIZE}(L^c)
\]
\emph{Valid when:} Applies when \(R_{N}\) depends only on the L-bit truth table, is decidable in poly(2\textasciicircum{}L) time, and rejects all \(L^{c}-size\) truth tables.

\paragraph{Next analysis.} Instantiate the random-affine generator with L=(log N)\textasciicircum{}2 and q in ((1+epsilon)/2,1+epsilon); compute \(r_{N}\) for the induced sparse incidence matrix and test whether the rank gap forces naturalness or low-rank falsification.

\paragraph{Falsifiers.}
\begin{itemize}
\item falsifier: r\_N<=O(N\textasciicircum{}\{1+epsilon\} log N) gives too little rank for the density-counting route and revives low-rank affine sketch upper bounds.
\item falsifier: any transcript or codebook depending on SAT labels, including \(T_{N}=H_{N}\), violates label-freeness and collapses the audit into nonuniform leakage.
\item falsifier: a \(\mathrm{poly}(2^{L})-constructive\) large certificate for the actual SAT label vector becomes a natural-property route under PRF assumptions, so it is not a valid barrier escape.
\item falsifier: if a repaired carrier has target-query length \(N^\beta\) with \(\beta>0\), the q+beta c transfer cap blocks all-exponent SAT lower bounds.
\item falsifier: hidden advice in \(G_{N}\), the codebook, or the verifier above the \(N^{1+\epsilon}\) budget invalidates the source lower-bound accounting.
\end{itemize}

\subsection*{cycle-0063: Hardness-Magnification Bridge Realization}

\paragraph{Execution.} research; azure\_responses\_api; attempt primary; 600.01 seconds; 2 lemmas, 5 equations, 7 derivation steps, 5 falsifiers.

\paragraph{Learning.} Cycle 63 upgrades the rank-density fork into a quantitative gate: low constructive rank is a direct upper-bound falsifier, while high rank makes any efficient useful certificate large in label space. The \(\beta=0\) magnification bridge is still valid, but rank-only evidence is quarantined unless the source proof is direct or explicitly non-natural and nonrelativizing.

\paragraph{Pressure.} The branch still lacks a direct non-natural lower bound for \(H_{N}^{\mathrm{SAT}}\) or a nonrelativizing theorem extracting a label-free sketch from every small \(H_{N}\) circuit.

\paragraph{Derivation steps.}
\begin{enumerate}
\item Fix L=(log N)\textasciicircum{}d with \(d>1\) and define \(A_{N}\) so that \(H_{N}^{ell}=A_{N}\) ell is the XOR of \(N^{q+o(1)}\) generated length-L labels per source input.
\item Assume \(\mathrm{SAT}_{L}\) has circuits of size \(L^{c}\). Substitute the circuit into each query and add the generator cost to get C(H\_N\textasciicircum{}\{SAT\})<=N\textasciicircum{}\{max\{rho,q\}+o(1)\}.
\item If \(max{\rho,q}N^{1+\epsilon}\) is available, the substitution contradicts \(\mathrm{SAT}_{L}\) in SIZE(L\textasciicircum{}c) for that fixed c; since c was arbitrary, this is the valid \(\beta=0\) bridge to \(\mathrm{SAT}\notin \mathrm{P}/\mathrm{poly}\).
\item Audit the source proof. If rows of \(A_{N}\) have a basis-coordinate decoder of size \(D_{N}\) and rank \(r_{N}\) with \(D_{N}+O(r_{N})<N^{1+\epsilon}\), hardwire the \(r_{N}\) SAT basis parities and compute \(H_{N}^{\mathrm{SAT}}\) below threshold; this is a concrete falsifier.
\item If \(r_{N}\) is much larger than \(s_{N}\) log(s\_N+N), random label vectors produce high-complexity source functions with probability at least 1-2\textasciicircum{}\{O(s\_N log(s\_N+N))-r\_N\}. Thus the semantic property C(A\_N ell)>s\_N is large in label space.
\item Any certificate \(Q_{N}\) contained in this property, decidable from the label truth table in poly(2\textasciicircum{}L), accepting SAT, and rejecting every label function with SIZE(L\textasciicircum{}c) circuits is a natural-proof-style object against SIZE(L\textasciicircum{}c). The route must then state a non-natural escape or remain only outside the known barrier conditionally.
\item Finally, rank-only hypotheses are relativizing-invalid: setting all labels to zero preserves \(A_{N}\) and its rank but makes \(H_{N}\) constant.
\end{enumerate}

\paragraph{Equation 1: xor\_carrier\_definition.}
\[
H_N^{\ell}(x)=\bigoplus_{j\le M_N}\ell(G_N(x,j))=(A_N\ell)(x)
\]
\emph{Valid when:} Valid when \(G_{N}\) is label-free, repeated generated formulas are counted modulo two, and ell is any Boolean label vector on the generated formula universe.

\paragraph{Equation 2: beta\_zero\_target\_substitution.}
\[
C(H_N^{SAT})\le N^{\rho+o(1)}+N^{q+o(1)}L^c=N^{\max{\rho,q}+o(1)}
\]
\emph{Valid when:} Applies after SAT has SIZE(L\textasciicircum{}c) circuits, L=(log N)\textasciicircum{}d is polylogarithmic, and generator/advice overhead is charged in \(\rho\).

\paragraph{Equation 3: magnification\_implication.}
\[
C(H_N^{SAT})> N^{1+\epsilon}\ \wedge\ \max{\rho,q}< 1+\epsilon\ \Rightarrow\ SAT\notin SIZE(n^c)\text{ for every fixed }c
\]
\emph{Valid when:} Requires the same carrier family, \(\beta\) zero target-query length, and a valid source lower bound for the actual SAT labels.

\paragraph{Equation 4: rank\_density\_count.}
\[
\Pr_{\ell\leftarrow{0,1}^{U_N}}[C(A_N\ell)\le s_N]\le 2^{O(s_N\log(s_N+N))-r_N}
\]
\emph{Valid when:} Valid when \(A_{N}\) has \(F_{2}\) rank \(r_{N}\) and low-complexity N-input source functions are counted by circuits of size at most \(s_{N}\).

\paragraph{Equation 5: constructive\_low\_rank\_falsifier.}
\[
C(H_N^{SAT})\le D_N+O(r_N)
\]
\emph{Valid when:} Requires an explicit row-coordinate decoder of size \(D_{N}\) and hardwiring the \(r_{N}\) basis parities of the actual SAT label vector as advice bits.

\paragraph{Next analysis.} Perform the finite prior-art boundary check for the label-space naturalness claim, then instantiate the thought experiment with a concrete pairwise-independent generator and compute \(r_{N}\), quotient rank, and any row-coordinate decoder cost \(D_{N}\).

\paragraph{Falsifiers.}
\begin{itemize}
\item positive \(\beta\) black-box transfer cap q+beta*c+o(1), which blocks all-exponent target preservation unless \(\beta=0\) or subpolynomial
\item constructive low-rank row-coordinate decoder with \(D_{N}+O(r_{N})<N^{1+\epsilon}\)
\item label-free global codebook or sparse quotient-polynomial sketch below the source threshold
\item rank-only proof surviving replacement of SAT labels by the zero oracle
\item large \(\mathrm{poly}(2^{L})-constructive\) certificate useful against SIZE(L\textasciicircum{}c), triggering natural-proof exposure rather than a barrier escape
\end{itemize}

\subsection*{cycle-0064: Hardness-Magnification Bridge Realization}

\paragraph{Execution.} research; azure\_responses\_api; attempt primary; 600.009 seconds; 3 lemmas, 4 equations, 5 derivation steps, 5 falsifiers.

\paragraph{Learning.} Cycle 64 upgrades the XOR route from rank-density intuition to a fork theorem: low constructive incidence rank is an explicit upper-bound falsifier, while any large truth-table constructive anti-alignment proof for the SAT label pulls back through the valid \(\beta=0\) bridge to a natural property. The remaining source route must be direct, nonlarge, nonconstructive, proof-mediated, or explicitly nonrelativizing; this is route-level progress, not \(\mathrm{P} \ne \mathrm{NP}\).

\paragraph{Pressure.} The route still lacks a nonlarge, nonconstructive, proof-mediated, or direct argument that the actual SAT label avoids the easy induced-XOR label set.

\paragraph{Derivation steps.}
\begin{enumerate}
\item Model the XOR carrier as a linear map from formula-label vectors y over \(U_{N}\) to N-input truth tables \(H_{N}^{y}\).
\item Assume a certificate property \(\mathrm{P}_{N}\) contains \(y_{N}^{\mathrm{SAT}}\) and excludes every y with \(C(H_{N}^{y})\le S_{N}\).
\item Pull \(\mathrm{P}_{N}\) back to a property of all \(L_{N}-bit\) Boolean functions by testing only the coordinates in \(U_{N}\); density is preserved exactly under random restriction.
\item If an \(L_{N}-bit\) function f has circuits of size \(L_{N}^{c}\), substituting that circuit into the generator computes \(H_{N}^{f}\) in size \(N^{max{q,\rho}+o(1)}<S_{N}\).
\item Therefore the pullback property rejects every SIZE(L\_N\textasciicircum{}c) function but accepts SAT, and if it is constructive and large it is a natural property useful against polynomial-size circuits at length \(L_{N}\).
\end{enumerate}

\paragraph{Equation 1: xor\_label\_function.}
\[
H_N^y(x)=\bigoplus_{j\le Q_N} y(G_N(x,j))=\langle a_x,y\rangle_{\mathbb{F}_2}
\]
\emph{Valid when:} Valid when the carrier is the generated XOR construction and labels are assigned only to generated \(L_{N}-bit\) formulas.

\paragraph{Equation 2: target\_substitution\_upper\_bound.}
\[
SAT_{L_N}\in SIZE(L_N^c)\Rightarrow C(H_N^{SAT})\le N^{\rho}+Q_N\,L_N^{O(c)}=N^{\max{\rho,q}+o(1)}< S_N
\]
\emph{Valid when:} Applies after the generator overhead is \(N^\rho\), query count is \(N^{q}\), and \(\beta\) equals zero.

\paragraph{Equation 3: pullback\_density\_preservation.}
\[
\Pr_F[\widehat{P_{N}}(F)]=|P_N|/2^{|U_N|}
\]
\emph{Valid when:} Valid when F is uniformly random over all \(L_{N}-bit\) truth tables and restriction to \(U_{N}\) is uniform.

\paragraph{Equation 4: natural\_usefulness.}
\[
P_N\cap{y:C(H_N^y)\le S_N}=\varnothing\ \wedge\ y_N^{SAT}\in P_N\Rightarrow \widehat{P_{N}}(SAT)=1\ \wedge\ \widehat{P_{N}}(f)=0\ \forall f\in SIZE(L_N^c)
\]
\emph{Valid when:} Requires the target-preserving substitution upper bound to put every small f inside the easy induced-XOR set.

\paragraph{Next analysis.} Instantiate L\_N=(log N)\textasciicircum{}2 and q in ((1+epsilon)/2,1+epsilon); compute achievable \(M_{N}\) and \(r_{N}\) for an explicit generator family, then test for a label-free coordinate transcript or prove any proposed SAT anti-alignment certificate has density below \(2^{-\omega(L_{N})}\).

\paragraph{Falsifiers.}
\begin{itemize}
\item falsifier: positive \(\beta\) gives the transfer cap q+beta*c+o(1) and cannot exclude all SAT exponents.
\item falsifier: an exact label-free coordinate transcript with \(S_{T}+r_{N}<N^{1+\epsilon-o(1)}\) makes every XOR label easy.
\item falsifier: a large truth-table constructive anti-alignment certificate pulls back to a natural property under the \(\beta=0\) bridge.
\item falsifier: rank below O(S\_N log \(S_{N}\)) gives no distributional hard-label mass.
\item falsifier: any SAT-label advice embedded in \(G_{N}\) or \(\mathrm{P}_{N}\) invalidates the target-preserving condition.
\end{itemize}

\subsection*{cycle-0065: Hardness-Magnification Bridge Realization}

\paragraph{Execution.} research; azure\_responses\_api; attempt primary; 600.005 seconds; 3 lemmas, 4 equations, 7 derivation steps, 4 falsifiers.

\paragraph{Learning.} Cycle 65 closes a finite validation gap: short searchable proof certificates do not automatically escape natural proofs in the \(\beta=0\) XOR carrier. If they are large and sound, they form a valid natural property; if rank is low, there is a separate upper-bound falsifier. The surviving route is narrower: direct \(H_{N}\) lower bounds or nonlarge, nonconstructive, explicitly nonrelativizing proof objects.

\paragraph{Pressure.} The argument is conditional as a barrier diagnosis: it uses the standard PRF-based natural-proofs obstruction and does not rule out nonlarge SAT-specific proofs, \(\omega(L_{N})-length\) proofs, nonconstructive verifiers, or nonrelativizing proof-system semantics.

\paragraph{Derivation steps.}
\begin{enumerate}
\item Fix the XOR carrier \(H_{N,y}=A_{N}\) y over \(F_{2}\) and let r\_N=rank(A\_N). The map y \(\to \) \(H_{N,y}\) has kernel size \(2^{M_{N}-r_{N}}\).
\item Count N-input Boolean functions with circuits of size at most \(S_{N}\) as at most 2\textasciicircum{}\{O(S\_N log S\_N)\}. Therefore the easy-label density is at most 2\textasciicircum{}\{O(S\_N log S\_N)-r\_N\}.
\item If \(r_{N}\) exceeds O(S\_N log S\_N)+Omega(log \(M_{N}\)), then the hardness property \({y:C(H_{N,y})>S_{N}}\) is large. This is evidence of naturalness, not evidence for SAT.
\item For a proof relation \(R_{N}\) with \(|pi|\le aL_{N}\) and verifier time \(M_{N}^{b}\), the existential property \(Q_{N}(y)=exists\) pi R\_N(y,pi) is decidable in \(M_{N}^{O(1)}\) time by exhaustive proof search.
\item Soundness \(R_{N}\Rightarrow C(H_{N,y})>S_{N}\) plus \(max{\rho,q}<1+\epsilon\) implies \(Q_{N}\) contains no label function y with SIZE(L\_N\textasciicircum{}c) circuits for any fixed c, because substituting such a circuit into \(G_{N}\) yields \(H_{N,y}\) of size below \(S_{N}\).
\item Thus any large accepted set \(Q_{N}\) is constructive, large, and useful: exactly a natural property after encoding formula-code labels as \(length-L_{N}\) truth tables.
\item Validity regime check: the derivation uses the definitions and assumptions listed in the candidate lemmas, preserves the same source object and target object, and is valid only in the stated P vs NP branch model with its declared circuit class, length regime, and advice/uniformity conditions.
\end{enumerate}

\paragraph{Equation 1: beta\_zero\_substitution.}
\[
C(H_{N,y}) \le N^{\rho+o(1)}+N^{q+o(1)}L_N^c=N^{\max{\rho,q}+o(1)}
\]
\emph{Valid when:} Valid when \(G_{N}\) is explicit, target queries have length \(L_{N}\), and the substituted target circuit carries no additional advice.

\paragraph{Equation 2: rank\_counting\_large\_hardness\_property.}
\[
\Pr_{y\leftarrow{0,1}^{M_N}}[C(H_{N,y})\le S_N]\le 2^{O(S_N\log S_N)-r_N}
\]
\emph{Valid when:} Valid when \(A_{N}\) is viewed as an \(F_{2}\) linear map from label vectors to carrier truth tables and small N-input circuits are counted nonuniformly.

\paragraph{Equation 3: search\_time\_for\_proofs.}
\[
T_{decide}(\mathcal{Q}_N)\le 2^{aL_N}M_N^b=M_N^{a+b}
\]
\emph{Valid when:} Valid when certificates have at most \(aL_{N}\) bits, the verifier runs in \(M_{N}^{b}\) time, and \(M_{N}\) equals \(2^{L_{N}}\).

\paragraph{Equation 4: usefulness\_against\_target\_circuits.}
\[
\mathcal{Q}_N\cap SIZE(L_N^c)=\emptyset \quad \text{for every fixed }c
\]
\emph{Valid when:} Applies after the beta-zero substitution bridge preserves the same circuit model, the same advice convention, and the threshold \(S_{N}=N^{1+\epsilon}\).

\paragraph{Next analysis.} Instantiate \(R_{N}\) with a named proof system, e.g. Extended Frege or IPS, and test whether an LB certificate for \(C(H_{N,y})>N^{1+\epsilon}\) has \(O(L_{N})-bit\) searchable descriptions, inverse-polynomial label density, and truth-table soundness; otherwise record the precise escape parameter.

\paragraph{Falsifiers.}
\begin{itemize}
\item falsifier: positive \(\beta\) gives C(H\_N)<=N\textasciicircum{}\{q+beta c+o(1)\} and caps the claimed all-exponent transfer.
\item falsifier: low constructive incidence rank gives C(H\_\{N,SAT\})<=S\_T+O(r\_N), killing the \(N^{1+\epsilon}\) source lower bound.
\item falsifier: an \(O(L_{N})-searchable\) inverse-polynomially large proof certificate becomes a natural property, so it is not a valid barrier-escaping proof route under PRF assumptions.
\item falsifier: label-dependent transcripts such as \(T_{N}=H_{N}\) collapse the audit and are nonuniform leakage, not source evidence.
\end{itemize}

\subsection*{cycle-0066: Hardness-Magnification Bridge Realization}

\paragraph{Execution.} research; azure\_responses\_api; attempt primary; 600.006 seconds; 3 lemmas, 5 equations, 8 derivation steps, 5 falsifiers.

\paragraph{Learning.} Cycle 66 sharpens the proof-mediated XOR route: searchable certificates do not automatically trigger the standard natural-proofs barrier. The valid obstruction needs parametric soundness, constructivity, large pulled-back density, and a density-preserving label-realization map. Without these, the barrier claim is invalid, but the source proof must be direct, nonlarge, nonconstructive, or explicitly nonrelativizing.

\paragraph{Pressure.} No explicit density-preserving label-realization map is known for the generated SAT universe \(U_{N}\), and soundness of plausible proof predicates is usually shown only for the actual SAT label rather than for every artificial label y.

\paragraph{Derivation steps.}
\begin{enumerate}
\item Parameterize the XOR carrier by an arbitrary label vector y so that the real SAT carrier is only the specialization \(y=\mathrm{SAT}|_{U_{N}}\).
\item Assume a searchable Psys predicate \(A_{N}(y)=exists\) pi V\_N(y,pi) with same-theory parametric soundness \(A_{N}(y)\Rightarrow C(H_{N,y})>N^{1+\epsilon}\).
\item If \(\mathrm{SAT}_{L}\) has size \(L^{c}\) circuits, substitute those circuits into the real carrier and obtain C(H\_N\textasciicircum{}\{SAT\})<=N\textasciicircum{}\{max\{rho,q\}+o(1)\}, preserving the \(\beta=0\) target bridge when \(max{\rho,q}<1+\epsilon\).
\item To convert \(A_{N}\) into a standard natural property, introduce a realization map \(R_{n}\) from n-bit function truth tables to label vectors and require \(H_{N,R_{n}(f)}=f\circ tau_{n}\) up to overhead \(chi_{N}\).
\item Constructivity follows from computing \(R_{n}\) and searching pi within the declared truth-table time bound; usefulness follows because a small circuit for f would give a small circuit for \(H_{N,R_{n}(f)}\).
\item Largeness follows only if the pullback density nu\_n is at least \(2^{-O(n)}\), or if a distortion bound \(D_{N}\) transfers label-space density \(mu_{N}\) to nu\_n.
\item Therefore a large searchable sound proof predicate plus label realization yields a natural property; absence of realization or largeness is an escape condition but also a new source proof burden.
\item Validity regime check: the derivation uses the definitions and assumptions listed in the candidate lemmas, preserves the same source object and target object, and is valid only in the stated P vs NP branch model with its declared circuit class, length regime, and advice/uniformity conditions.
\end{enumerate}

\paragraph{Equation 1: beta\_zero\_substitution.}
\[
C(H_N^{SAT}) \le N^{\rho+o(1)}+N^q L^c=N^{\max{\rho,q}+o(1)}
\]
\emph{Valid when:} Valid when L=(log N)\textasciicircum{}d, \(\mathrm{SAT}_{L}\) has size \(L^{c}\) circuits, and the generator plus incidence computation is label-free with cost \(N^{\rho+o(1)}\).

\paragraph{Equation 2: parametric\_carrier.}
\[
H_{N,y}(z)=\bigoplus_{u\in U_N} I_N(z,u)y_u
\]
\emph{Valid when:} Valid when \(U_{N}\) is the fixed generated formula universe and \(I_{N}\) is computed independently of the SAT label vector y.

\paragraph{Equation 3: accepted\_label\_density.}
\[
\mu_N=2^{-|U_N|}\left|{y:\exists \pi, |\pi|\le \ell_N, V_N(y,\pi)=1}\right|
\]
\emph{Valid when:} Valid when the verifier is evaluated over artificial label vectors y and proof length cutoff \(ell_{N}\) is fixed before seeing y.

\paragraph{Equation 4: pullback\_density.}
\[
\nu_n=\Pr_f[A_N(R_n(f))=1]\ge D_N^{-1}\mu_N
\]
\emph{Valid when:} Applies after a density-distortion certificate \(R_{n},D_{N}\) proves that accepted label sets pull back to truth-table functions with bounded loss.

\paragraph{Equation 5: usefulness\_transfer.}
\[
A_N(R_n(f))=1\Rightarrow C(f)> N^{1+\epsilon}/\chi_N
\]
\emph{Valid when:} Requires parametric soundness for \(H_{N,y}\) and a realization certificate showing circuits for f induce circuits for \(H_{N,R_{n}(f)}\) with overhead \(chi_{N}\).

\paragraph{Next analysis.} Instantiate \(R_{n}\) for one high-q polylog XOR carrier or prove a finite impossibility: no label-free density-preserving realization can coexist with \(q,\rho<1+\epsilon\) and generator universe \(U_{N}\) of SAT formulas length (log N)\textasciicircum{}2.

\paragraph{Falsifiers.}
\begin{itemize}
\item falsifier: positive target-query exponent \(\beta\) gives the upper bound N\textasciicircum{}\{q+beta c+o(1)\} and caps the target exponents before reaching \(\mathrm{SAT}\notin \mathrm{P}/\mathrm{poly}\).
\item falsifier: low constructive incidence rank or a global label-free codebook yields an explicit small circuit for the XOR carrier.
\item falsifier: a large searchable sound proof predicate with density-preserving label realization becomes a natural property and is blocked under standard PRF-style assumptions.
\item falsifier: no density-preserving realization map means label-space largeness is not a valid standard natural-proofs argument.
\item falsifier: verifier advice or proof advice encoding \(\mathrm{SAT}|_{U_{N}}\) destroys the target-preserving condition by nonuniform leakage.
\end{itemize}

\subsection*{cycle-0067: Hardness-Magnification Bridge Realization}

\paragraph{Execution.} research; azure\_responses\_api; attempt primary; 600.01 seconds; 2 lemmas, 5 equations, 8 derivation steps, 6 falsifiers.

\paragraph{Learning.} Cycle 67 upgrades the previous label-realization warning into an exact restriction-extension theorem. For \(\beta=0\) XOR carriers, dense constructive label-mediated source lower bounds naturalize automatically when \(U_{N}\) is uniformly enumerable and soundness is in the same circuit model. The valid bridge survives, but any proof-mediated route must now be direct, nonlarge, nonconstructive, or explicitly nonrelativizing; ambient source-property largeness is not enough because density must be image-relative.

\paragraph{Pressure.} No lower bound for C(H\_N) is proved. The derivation only certifies that one dense constructive label-mediated proof shape would collide with natural-proofs methodology, while image-relative and nonlarge certificates remain open.

\paragraph{Derivation steps.}
\begin{enumerate}
\item Fix the \(\beta=0\) XOR carrier with L=(log N)\textasciicircum{}d, \(d>1\), \(N^{q}\) generated queries, and overhead \(N^{\rho+o(1)}\).
\item For any target circuit D of size \(L^{c}\), substitute D into the generator and XOR combiner to obtain C(A\_N y\_D)<=N\textasciicircum{}\{max\{rho,q\}+o(1)\}.
\item If \(C(H_{N})>N^{1+\epsilon}\) and \(max{\rho,q}<1+\epsilon\), then no such D computes \(\mathrm{SAT}_{L}\) on the valid lengths; this is the route-level target bridge, not \(\mathrm{P} \ne \mathrm{NP}\).
\item Assume the proposed source proof is mediated by \(Q_{N}\) over labels on \(U_{N}\), with \(Q_{N}(y^{\mathrm{SAT}})=1\) and \(Q_{N}(y)\Rightarrow C(A_{Ny})>N^{1+\epsilon}\).
\item Define the full-target property widehat \(Q_{N}(F)=1\) iff F restricted to \(U_{N}\) lies in \(Q_{N}\).
\item Uniform random full truth tables restrict to uniform labels on \(U_{N}\), so density(widehat Q\_N)=density(Q\_N) with no loss.
\item Source soundness plus the \(\beta=0\) transfer makes widehat \(Q_{N}\) reject every SIZE(L\textasciicircum{}c) target circuit, while accepting \(\mathrm{SAT}_{L}\).
\item Hence a large constructive \(Q_{N}\) naturalizes; the surviving proof-mediated route must be nonlarge, nonconstructive, not label-factorized, or explicitly nonrelativizing and direct.
\end{enumerate}

\paragraph{Equation 1: carrier-linearization.}
\[
H_y(x)=\bigoplus_{j\le N^q} y(\phi_{N,j}(x))=(A_N y)(x),\quad H_N=A_N y^{SAT}_{U_N}
\]
\emph{Valid when:} Valid when \(U_{N}\) contains every generated query string and duplicate queries are quotient-counted modulo parity.

\paragraph{Equation 2: beta-zero-transfer.}
\[
C(A_N y_D) \le N^{\rho+o(1)}+N^q L^c\,poly(L)=N^{\max{\rho,q}+o(1)}
\]
\emph{Valid when:} Valid when D is a SIZE(L\textasciicircum{}c) target circuit, L=(log N)\textasciicircum{}d, and substitution uses no advice except the fixed generator.

\paragraph{Equation 3: restriction-density.}
\[
\Pr_{F:{0,1}^L\to{0,1}}[F|_{U_N}\in Q_N]=|Q_N|/2^{|U_N|}=\delta_N
\]
\emph{Valid when:} Valid when full target truth tables are sampled uniformly and \(U_{N}\) is fixed independently of the sampled labels.

\paragraph{Equation 4: naturalization-gate.}
\[
\delta_N\ge 2^{-O(L)},\ T_Q\le 2^{O(L)},\ Q_N\Rightarrow C(A_N y)> N^{1+\epsilon}\ \Longrightarrow\ \widehat{Q_{N}}\text{ is constructive and }c\text{-useful}
\]
\emph{Valid when:} Valid when source soundness is in the same nonuniform circuit model as the beta-zero transfer.

\paragraph{Equation 5: magnification-consequence.}
\[
C(H_N)> N^{1+\epsilon}\ \wedge\ \max{\rho,q}< 1+\epsilon\ \Longrightarrow\ SAT_L\notin SIZE(L^c)
\]
\emph{Valid when:} Applies after the carrier lower bound holds on the chosen valid lengths and L\textasciicircum{}c=N\textasciicircum{}\{o(1)\} for fixed c.

\paragraph{Next analysis.} Instantiate \(U_{N}\) for one explicit polylog-query XOR generator, compute whether \(U_{N}\) is enumerable in poly(2\textasciicircum{}L), and classify a candidate certificate as dense-label, image-relative, or nonlarge.

\paragraph{Falsifiers.}
\begin{itemize}
\item positive \(\beta\) falsifier: if target query length is \(N^\beta\) with \(\beta>0\), the transfer N\textasciicircum{}\{q+beta c+o(1)\} cannot exclude all polynomial SAT exponents.
\item low constructive incidence rank falsifier: a small row-coordinate transcript lets the XOR carrier hardwire basis parities and stay below \(N^{1+\epsilon}\).
\item label-free codebook falsifier: a compact quotient transcript plus sparse corrections gives \(C(H_{N})\le S_{T}+S_{D}+\)\(\widetilde{O}(U_{N}+A_{N} r_{N})\)\(+E_{NN}\).
\item dense constructive label predicate falsifier for barrier-escape claims: if \(Q_{N}\) is large, constructive, and source-sound, restriction extension makes it a natural property.
\item image-only density falsifier: largeness of \(\mathrm{P}_{N}\) in the full source truth-table cube is irrelevant unless \(\mathrm{P}_{N}\) is large on im(A\_N).
\item nonuniform leakage falsifier: if \(U_{N}\) or \(Q_{N}\) needs hidden advice not charged in poly(2\textasciicircum{}L), constructivity and target preservation fail.
\end{itemize}

\subsection*{cycle-0068: Hardness-Magnification Bridge Realization}

\paragraph{Execution.} research; azure\_responses\_api; attempt primary; 600.005 seconds; 3 lemmas, 4 equations, 7 derivation steps, 6 falsifiers.

\paragraph{Learning.} Cycle 68 separates two notions previously conflated: label-mediated certificates for the \(\beta=0\) XOR carrier are large inside the realization image, but the image has ambient density at most 2\textasciicircum{}\{-2\textasciicircum{}N+2\textasciicircum{}\{o(N)\}\}. Thus the valid bridge to \(\mathrm{SAT}\notin \mathrm{P}/\mathrm{poly}\) survives, while the standard natural-proofs barrier requires a separate constructive ambient extension; this is a sharper route-status theorem, not progress on the main conjecture itself.

\paragraph{Pressure.} No direct lower bound for the real SAT-label carrier \(H_{N}\) is obtained, and no polynomial-time density-preserving extension from Im(R\_N) to ambient N-bit truth tables is known.

\paragraph{Derivation steps.}
\begin{enumerate}
\item Fix \(M=N^{q}\), L=(log N)\textasciicircum{}d with \(d>1\), and let \(G_{N}\) have size \(N^{\rho+o(1)}\).
\item Assume SAT on L-bit formulas has circuits of size \(L^{c}\). Substitute that circuit into all M XOR positions and compute \(G_{N}\) uniformly, giving \(C(H_{N})\le N^{max{\rho,q}+o(1)}\).
\item Thus a valid source lower bound \(C(H_{N})>N^{1+\epsilon}\) contradicts \(\mathrm{SAT}\in \mathrm{SIZE}(n^{c})\) for every fixed c when \(max{\rho,q}<1+\epsilon\).
\item Define \(U_{N}\) as the generated formula universe and \(R_{N}(y)(x)=xor_{j}\) y(G\_N(x,j)). Since formula strings have length L, |U\_N|<=2\textasciicircum{}\{O(L)\}=2\textasciicircum{}\{o(N)\}.
\item The image Im(R\_N) has at most \(2^{|U_{N}|}\) carrier functions, while the ambient N-bit function space has size 2\textasciicircum{}\{2\textasciicircum{}N\}; hence any property supported only on Im(R\_N) has density at most 2\textasciicircum{}\{|U\_N|-2\textasciicircum{}N\}.
\item Therefore a label-mediated proof can be image-relative large and still fail standard natural-proofs largeness by an exponential margin.
\item For XOR maps, equal fiber sizes show label density transfers to image-relative density, so the exact missing certificate is not density inside the image but a constructive ambient extension preserving soundness.
\end{enumerate}

\paragraph{Equation 1: beta0\_target\_transfer.}
\[
C(H_N) \le N^{\rho+o(1)}+N^q L^c = N^{\max{\rho,q}+o(1)}
\]
\emph{Valid when:} Valid when \(\beta=0\) means L=(log N)\textasciicircum{}d and SAT circuits of size \(L^{c}\) are substituted independently into the \(M=N^{q}\) generated queries.

\paragraph{Equation 2: all\_exponent\_consequence.}
\[
C(H_N)> N^{1+\epsilon} \;\wedge\; \max{\rho,q}< 1+\epsilon \Rightarrow \forall c\; SAT\notin SIZE(n^c)
\]
\emph{Valid when:} Applies after the reduction preserves nonuniform circuit model identity, generator overhead \(\rho\), query multiplicity q, and the fixed polylogarithmic query length.

\paragraph{Equation 3: realization\_image\_density.}
\[
\mu_{amb}(Im(R_N)) \le 2^{|U_N|-2^N} \le 2^{-2^N+2^{O((\log N)^d)}}
\]
\emph{Valid when:} Requires source functions on N input bits and carriers determined solely by labels on the generated formula universe \(U_{N}\).

\paragraph{Equation 4: standard\_largeness\_failure.}
\[
|U_N|=2^{o(N)} \Rightarrow \mu_{amb}(Im(R_N))< 2^{-N^a}\quad\text{for every fixed }a
\]
\emph{Valid when:} Valid when d is fixed, L=(log N)\textasciicircum{}d, and natural-proofs largeness is the standard ambient \(2^{-\mathrm{poly}(N)}\) threshold.

\paragraph{Next analysis.} Formalize the extension problem: given \(R_{N}\) and a sound searchable label certificate, prove either that no poly(N)-time ambient extension of density \(2^{-\mathrm{poly}(N)}\) exists for this certificate class, or that such an extension exists and therefore naturalizes the route.

\paragraph{Falsifiers.}
\begin{itemize}
\item Positive \(\beta\) target-query length gives the q+beta c+o(1) substitution cap and blocks all-exponent transfer.
\item \(q\ge 1+\epsilon\) or \(\rho\ge 1+\epsilon\) makes the substituted SAT upper bound compatible with the proposed source lower bound.
\item A low-cost label-free quotient transcript or affine incidence sketch still gives a direct upper-bound falsifier for the XOR carrier.
\item A constructive density-preserving ambient extension turns an image-relative certificate into a standard natural property and triggers the natural-proofs obstruction in the same circuit model.
\item Nonuniform leakage through precomputed SAT labels for \(U_{N}\) invalidates the controlled interface.
\item If d is not fixed polylogarithmic, the image-density calculation and \(\beta=0\) classification must be recomputed before it is valid.
\end{itemize}

\subsection*{cycle-0069: Hardness-Magnification Bridge Realization}

\paragraph{Execution.} research; azure\_responses\_api; attempt primary; 600.009 seconds; 3 lemmas, 4 equations, 6 derivation steps, 5 falsifiers.

\paragraph{Learning.} Cycle 69 turns the image-density ambiguity into a precise threshold theorem. A label-mediated \(\beta=0\) XOR certificate naturalizes through restriction to \(U_{N}\) exactly when its label density is at least \(2^{-\mathrm{poly}(L)}\) and its verifier is \(2^{O(L)}-constructive\). The earlier ambient-tiny source image is not by itself a valid escape; the real escapes are nonlarge density, nonconstructivity, advice leakage, or explicitly nonrelativizing direct lower bounds.

\paragraph{Pressure.} No construction is given of a \(\mathrm{P}_{N}\) that accepts \(\mathrm{SAT}|U_{N}\) while remaining sound; the derivation only classifies when such a proof-mediated route would naturalize.

\paragraph{Derivation steps.}
\begin{enumerate}
\item Fix L=(log N)\textasciicircum{}d, \(d>1\), and a uniform distinct query set \(U_{N}\) of size \(M_{N}=N^{q+o(1)}\).
\item Given a label property \(\mathrm{P}_{N}\), define \(Q_{N}\) on full L-bit truth tables by reading only coordinates in \(U_{N}\).
\item A uniformly random full truth table induces a uniformly random label vector on \(U_{N}\); hence ambient target density equals label density \(delta_{N}\) exactly.
\item If f has \(L^{c}-size\) circuits, \(\beta=0\) substitution gives \(H_{N,f|U_{N}}\) size \(N^{max{\rho,q}+o(1)}\); with \(max{\rho,q}<1+\epsilon\), soundness forces \(Q_{N}(f)=0\).
\item If \(\mathrm{P}_{N}\) is decidable in \(2^{O(L)}\) and \(delta_{N}\ge 2^{-\mathrm{poly}(L)}\), then \(Q_{N}\) is a constructive large property useful against \(L^{c}-size\) circuits.
\item Validity regime check: the derivation uses the definitions and assumptions listed in the candidate lemmas, preserves the same source object and target object, and is valid only in the stated P vs NP branch model with its declared circuit class, length regime, and advice/uniformity conditions.
\end{enumerate}

\paragraph{Equation 1: label\_space\_size.}
\[
M_N=|U_N|=N^{q+o(1)}=2^{qL^{1/d}+o(L^{1/d})}
\]
\emph{Valid when:} Valid when L=(log N)\textasciicircum{}d with fixed \(d>1\) and the generator emits \(N^{q+o(1)}\) distinct target formulas.

\paragraph{Equation 2: beta\_zero\_substitution.}
\[
C(H_{N,f|U_N})\le N^{\max{\rho,q}+o(1)}\cdot L^{O(c)}=N^{\max{\rho,q}+o(1)}
\]
\emph{Valid when:} Applies when target queries have polylogarithmic length and f has \(L^{c}-size\) circuits for fixed c.

\paragraph{Equation 3: restriction\_density\_preservation.}
\[
\Pr_{f\leftarrow{0,1}^{2^L}}[Q_N(f)=1]=|P_N|/2^{M_N}=\delta_N
\]
\emph{Valid when:} Valid when the coordinates in \(U_{N}\) are distinct so a random full truth table induces a uniform label vector.

\paragraph{Equation 4: natural\_largeness\_threshold.}
\[
\delta_N\ge 2^{-L^a}\quad\Longrightarrow\quad Q_N\text{ is }2^{-\mathrm{poly}(L)}\text{-large}
\]
\emph{Valid when:} Requires a fixed constant a and the standard target-length naturalness convention for L-bit Boolean functions.

\paragraph{Next analysis.} Instantiate \(U_{N}\) with L=(log N)\textasciicircum{}2 and q in ((1+epsilon)/2,1+epsilon); compute whether candidate proof systems give certificate density above \(2^{-\mathrm{poly}(L)}\) or force a nonlarge escape, then run the prior-art boundary check.

\paragraph{Falsifiers.}
\begin{itemize}
\item falsifier: \(\beta>0\) with transfer exponent q+beta c reaches the source threshold before all SAT exponents are excluded.
\item falsifier: low constructive rank or label-free codebook for the XOR incidence map gives a direct upper bound for \(H_{N}\).
\item falsifier: \(\mathrm{P}_{N}\) has density at least \(2^{-\mathrm{poly}(L)}\) and verifier time \(2^{O(L)}\), yielding a valid target-restriction natural property rather than a barrier escape.
\item falsifier: \(U_{N}\) or the certificate verifier uses nonuniform advice not reproducible by the target-side natural property.
\item falsifier: soundness is proved only for realized SAT labels and not for all small target functions f, breaking useful-property naturalization.
\end{itemize}

\subsection*{cycle-0070: Hardness-Magnification Bridge Realization}

\paragraph{Execution.} research; azure\_responses\_api; attempt primary; 600.002 seconds; 2 lemmas, 4 equations, 6 derivation steps, 5 falsifiers.

\paragraph{Learning.} Cycle 70 corrects the naturalness gate for the \(\beta=0\) XOR route. A dense label certificate cylinder-extends to a valid classical natural property exactly at density \(2^{-O(L)}\) with \(2^{O(L)}\) constructivity; source-image ambient sparsity is irrelevant after restriction. The real escapes are nonlarge density below the classical cutoff, nonconstructivity, advice leakage, or a direct nonrelativizing \(H_{N}\) lower bound.

\paragraph{Pressure.} This does not prove that such \(Q_{N}\) exists, and it does not turn sparse or nonconstructive certificates into barriers. It only gives a valid no-go for dense, constructive, same-model label certificates.

\paragraph{Derivation steps.}
\begin{enumerate}
\item Fix N and L=(log N)\textasciicircum{}d, \(d\ge 1\); enumerate \(U_{N}\) label-free and define R\_L(f)=Q\_N(f|U\_N).
\item For uniformly random truth table f on \({0,1}^{L}\), the restriction \(f|U_{N}\) is uniform over \({0,1}^{U_{N}}\); hence density is exactly preserved as \(delta_{N}\).
\item Given a truth table for f, enumerate \(U_{N}\), read the indexed bits, and run \(Q_{N}\); the total time is \(2^{O(L)}\) under \(|U_{N}|\le N^{q}\).
\item If \(R_{L}(f)=1\) and f had size \(L^{c}\), substitute that circuit into the XOR carrier to get C\_N(H\_\{f|U\_N\})<=N\textasciicircum{}\{max\{q,rho\}+o(1)\}<N\textasciicircum{}\{1+epsilon\}, contradicting \(Q_{N}\) soundness.
\item Since \(Q_{N}\) accepts \(\mathrm{SAT}|U_{N}\), \(R_{L}\) accepts \(\mathrm{SAT}_{L}\); thus any classically large constructive \(Q_{N}\) would naturalize the proof-mediated source lower bound.
\item Validity regime check: the derivation uses the definitions and assumptions listed in the candidate lemmas, preserves the same source object and target object, and is valid only in the stated P vs NP branch model with its declared circuit class, length regime, and advice/uniformity conditions.
\end{enumerate}

\paragraph{Equation 1: substitution\_upper\_bound.}
\[
C_N(H_{f|U_N}) \le N^\rho + N^q\,C_L(f)\operatorname{poly}(L)+N^{o(1)} = N^{\max{\rho,q}+o(1)} \quad \text{if } C_L(f)\le L^c.
\]
\emph{Valid when:} Valid when L=(log N)\textasciicircum{}d with fixed \(d\ge 1\), c is fixed, and the carrier accesses target labels only through the generated slots.

\paragraph{Equation 2: cylinder\_density\_preservation.}
\[
\Pr_{f:{0,1}^L\to{0,1}}[R_L(f)=1]=\frac{|Q_N|}{2^{|U_N|}}=\delta_N.
\]
\emph{Valid when:} Valid when \(R_{L}\) depends only on the restriction of the full truth table to the enumerated range \(U_{N}\).

\paragraph{Equation 3: constructivity\_transfer.}
\[
T_R(L) \le T_Q(N)+O(|U_N|\operatorname{poly}(L)) \le 2^{O(L)}.
\]
\emph{Valid when:} Valid when |U\_N|<=N\textasciicircum{}q=2\textasciicircum{}\{O(L)\} and \(Q_{N}\) is decidable from the restricted label table within \(2^{O(L)}\) time.

\paragraph{Equation 4: classical\_largeness\_cutoff.}
\[
\delta_N\ge 2^{-O(L)} \Rightarrow R_L \text{ is RR-large};\quad \delta_N=2^{-L^\theta},\ \theta> 1 \Rightarrow R_L \text{ is not classically large}.
\]
\emph{Valid when:} Applies when naturalness is measured at target input length L with the standard Razborov-Rudich density threshold \(2^{-O(L)}\).

\paragraph{Next analysis.} Instantiate the density-tier audit on a toy proof-certificate family: certificates describing affine anti-alignment witnesses over \(U_{N}\). Compute \(delta_{N}\), verifier time, and whether soundness reaches arbitrary \(C_{N}(H_{y})>N^{1+\epsilon}\) or only sketch hardness.

\paragraph{Falsifiers.}
\begin{itemize}
\item falsifier: \(delta_{N}\ge 2^{-O(L)}\) with \(T_{Q}\le 2^{O(L)}\) means the certificate route is a classical natural property, not a valid naturalness escape.
\item falsifier: \(max{q,\rho}\ge 1+\epsilon\) collapses the source threshold before any target lower bound is obtained.
\item falsifier: \(U_{N}\) enumeration or \(Q_{N}\) membership uses nonuniform advice encoding SAT labels, breaking the label-free restriction extension.
\item falsifier: source soundness proves only local separator hardness or rank hardness, not arbitrary circuit lower bounds for \(H_{y}\).
\item falsifier: \(d<1\) with \(|U_{N}|=2^{\omega(L)}\) can destroy \(2^{O(L)}\) constructivity of the cylinder property even when \(\beta\) remains subpolynomial.
\end{itemize}

\subsection*{cycle-0071: Hardness-Magnification Bridge Realization}

\paragraph{Execution.} research; azure\_responses\_api; attempt primary; 600.015 seconds; 3 lemmas, 5 equations, 7 derivation steps, 5 falsifiers.

\paragraph{Learning.} Cycle 71 upgrades the \(\beta=0\) XOR naturalness audit: density alone is not the decisive condition. Dense searchable certificates naturalize only when their soundness is functional over arbitrary label vectors. Therefore the surviving proof-mediated route must be sparse, nonconstructive, SAT-semantic/nonrelativizing, advice-free in a stronger sense, or a direct \(H_{N}\) lower bound.

\paragraph{Pressure.} This derivation needs functional soundness for all label vectors and a real accepted-label density estimate. SAT-only certificates and nonlarge accepted sets are outside the no-go.

\paragraph{Derivation steps.}
\begin{enumerate}
\item Fix f with \(C(f)\le L^{c}\) and set \(y=f|U_{N}\).
\item Substitute the f-circuit into every generated query of the XOR carrier; since L=polylog(N), the substituted target-circuit cost is \(N^{q+o(1)}\) and the generator cost is \(N^{\rho+o(1)}\).
\item If y is accepted by a functionally sound certificate, then C(H\_N\textasciicircum{}y)>N\textasciicircum{}\{1+epsilon\}.
\item When \(max{\rho,q}<1+\epsilon\), the previous two inequalities contradict each other, so no SIZE(L\textasciicircum{}c) function lies in the pullback property \(\mathrm{P}_{N}\).
\item For uniformly random f, \(f|U_{N}\) is uniform on \({0,1}^{U_{N}}\); hence the pullback density equals \(delta_{N}\).
\item If \(delta_{N}\ge 2^{-O(L)}\) and the verifier is \(2^{O(L)}-constructive\), \(\mathrm{P}_{N}\) is a classical natural property useful against SIZE(L\textasciicircum{}c).
\item Validity regime check: the derivation uses the definitions and assumptions listed in the candidate lemmas, preserves the same source object and target object, and is valid only in the stated P vs NP branch model with its declared circuit class, length regime, and advice/uniformity conditions.
\end{enumerate}

\paragraph{Equation 1: xor\_carrier.}
\[
H_N^y(x)=\bigoplus_{u\in U_N} a_N(x,u)y_u
\]
\emph{Valid when:} Valid when \(U_{N}\) is the distinct generated query set and \(a_{N}\) is the advice-free quotient incidence predicate.

\paragraph{Equation 2: beta\_zero\_substitution.}
\[
C(H_N^{f|U_N}) \le N^{\max{\rho,q}+o(1)} \quad \text{if } C(f)\le L^c
\]
\emph{Valid when:} Applies when target queries have length L=polylog(N) and \(\beta=0\) removes the c-dependent exponent term.

\paragraph{Equation 3: density\_pullback.}
\[
\Pr_F[P_N(F)=1]=\delta_N=|A_N|/2^{|U_N|}
\]
\emph{Valid when:} Valid when F is a uniformly random L-bit Boolean function and \(U_{N}\) contains distinct query strings.

\paragraph{Equation 4: naturalness\_cutoff.}
\[
\delta_N\ge 2^{-K L},\; T_V\le 2^{K L} \Rightarrow P_N \text{ is classically large and constructive}
\]
\emph{Valid when:} Requires the usual truth-table constructivity scale and the density cutoff measured over all L-bit functions.

\paragraph{Equation 5: usefulness\_contradiction.}
\[
P_N(f)=1 \wedge C(f)\le L^c \Rightarrow N^{1+\epsilon}< C(H_N^{f|U_N})\le N^{\max{\rho,q}+o(1)}
\]
\emph{Valid when:} Applies after functional soundness and the beta-zero transfer hold with \(max{\rho,q}<1+\epsilon\).

\paragraph{Next analysis.} Instantiate \(V_{N}\) with one concrete proof-system or MCSP-style certificate, compute \(kappa_{N}\) for accepted labels, and decide whether soundness is functional over arbitrary y or essentially SAT-semantic.

\paragraph{Falsifiers.}
\begin{itemize}
\item falsifier: positive \(\beta\) in a black-box target substitution revives the q+beta*c threshold cap and prevents all-exponent transfer.
\item falsifier: low constructive quotient rank or a small label-free codebook gives an explicit upper bound for the XOR carrier.
\item falsifier: functional soundness plus accepted-label density at least \(2^{-O(L)}\) and \(2^{O(L)}\) constructivity naturalizes the proof certificate.
\item falsifier: nonuniform row selection or advice-dependent \(U_{N}\) invalidates the density pullback and the target-preserving bridge.
\item falsifier: SAT-only soundness without a nonrelativizing invariant gives no natural property and no certified source lower bound.
\end{itemize}

\subsection*{cycle-0072: Hardness-Magnification Bridge Realization}

\paragraph{Execution.} research; azure\_responses\_api; attempt primary; 600.005 seconds; 2 lemmas, 3 equations, 6 derivation steps, 6 falsifiers.

\paragraph{Learning.} Cycle 72 upgrades the \(\beta=0\) XOR route from a vague naturalness warning to a dichotomy: dense functional searchable certificates become valid natural properties, while nonfunctional proof objects must pass an explicit row-selection leakage audit. The target bridge survives, but the source proof space is narrowed to sparse, nonconstructive, SAT-semantic nonrelativizing, or direct \(H_{N}\) lower bounds.

\paragraph{Pressure.} No theorem yet converts SAT-semantic, nonrelativizing certificates into functional certificates or into charged row-selection transcripts. That is the exact remaining gap, not a proof of \(\mathrm{P} \ne \mathrm{NP}\).

\paragraph{Derivation steps.}
\begin{enumerate}
\item Instantiate the surviving carrier with L=(log N)\textasciicircum{}d, \(d>1\), \(T_{N}=N^{q}\), \(q<1+\epsilon\), and generator size \(N^\rho\) with \(\rho<1+\epsilon\).
\item Assume a certificate predicate \(\mathrm{P}_{N}\) is dense, \(2^{O(L)}-constructive\), and functionally sound for every label vector ell. Define \(R_{L}(f)=1\) iff exists pi P\_N(ell\_f,pi).
\item Constructivity of \(R_{L}\) is polynomial in the truth-table length \(2^{L}\), and density is at least \(2^{-O(L)}\), so \(R_{L}\) is large in the Razborov-Rudich sense.
\item If f has circuits of size \(L^{c}\), substituting f into the XOR generator gives C(H\_\{N,f\})<=N\textasciicircum{}\{max\{q,rho\}+o(1)\}<N\textasciicircum{}\{1+epsilon\}; functional soundness therefore forces \(R_{L}(f)=0\). This proves usefulness against \(L^{c}\) circuits for every fixed c.
\item If \(\mathrm{P}_{N}\) is not functional, extract the charged row-selection transcript it uses. A transcript with \(A_{N}+S_{N}+min{R_{N},K_{NC}(\mathrm{SAT}_{L})}+E_{NN}\) below \(N^{1+\epsilon}\) is a concrete falsifier of the claimed source lower-bound route.
\item Validity regime check: the derivation uses the definitions and assumptions listed in the candidate lemmas, preserves the same source object and target object, and is valid only in the stated P vs NP branch model with its declared circuit class, length regime, and advice/uniformity conditions.
\end{enumerate}

\paragraph{Equation 1: substitution\_upper\_bound\_beta0.}
\[
C(H_{N,f}) \le C(G_N)+O(T_N\,C(f)+T_N L^{O(1)}) = N^{\max{\rho,q}+o(1)} \quad \text{if } C(f)\le L^c.
\]
\emph{Valid when:} Applies when \(G_{N}\) has size \(N^\rho\), \(T_{N}=N^{q}\), and the target truth-table function f has circuits of size \(L^{c}\) with L=(log N)\textasciicircum{}d.

\paragraph{Equation 2: threshold\_separation\_window.}
\[
\max{\rho,q}< 1+\epsilon \implies N^{\max{\rho,q}+o(1)}< N^{1+\epsilon}.
\]
\emph{Valid when:} Requires fixed \(\epsilon>0\), fixed q and \(\rho\) below threshold, and N large enough that the o(1) term is smaller than the threshold margin.

\paragraph{Equation 3: natural\_largeness\_cutoff.}
\[
\delta_N=|A_N|/2^{2^L}\ge 2^{-aL}=(2^L)^{-a}.
\]
\emph{Valid when:} Valid when certificates are counted over complete L-bit truth tables and constructivity is polynomial in the truth-table length \(2^{L}\).

\paragraph{Next analysis.} Instantiate the selected-probe audit for one named proof-system candidate: require \(q=1+\epsilon/2\), \(\rho=1+\epsilon/3\), L=(log N)\textasciicircum{}2, then compute whether its certificates force \(A_{N}+S_{N}+R_{N}+E_{NN}\) or \(A_{N}+S_{N}+K_{NL}^{c}+E_{NN}\) below \(N^{1+\epsilon}\).

\paragraph{Falsifiers.}
\begin{itemize}
\item falsifier: positive target-query exponent \(\beta\) gives the old transfer cap q+beta c and cannot exclude all SAT exponents.
\item falsifier: a dense searchable certificate with functional soundness is a valid natural property, so it is barred under standard PRF assumptions rather than an escape route.
\item falsifier: label-dependent row selection with \(A_{N}+S_{N}+R_{N}+E_{NN}<N^{1+\epsilon}\) gives an unconditional small circuit for the source carrier.
\item falsifier: under \(\mathrm{SAT}_{L}\) in SIZE(L\textasciicircum{}c), a selected-probe transcript with A\_N+S\_N+K\_NL\textasciicircum{}c+E\_NN<N\textasciicircum{}\{1+epsilon\} collapses the attempted contradiction before reaching all-exponent separation.
\item falsifier: density below \(2^{-\omega(L)}\) avoids classical naturalization but also breaks the promised searchable-large certificate route unless a separate nonlarge proof system is specified.
\item falsifier: advice leakage in \(\mathrm{P}_{N}\) invalidates both the natural-property derivation and the claimed target-preserving bridge unless the advice is explicitly charged.
\end{itemize}

\subsection*{cycle-0073: Hardness-Magnification Bridge Realization}

\paragraph{Execution.} research; azure\_responses\_api; attempt primary; 600.005 seconds; 3 lemmas, 4 equations, 7 derivation steps, 5 falsifiers.

\paragraph{Learning.} Cycle 73 adds a concrete row-deletion residual threshold to the surviving \(\beta=0\) XOR magnification route. A claimed \(N^{1+\epsilon}\) source lower bound is invalid if the carrier has a subthreshold circuit wrong on fewer than \(N^{\epsilon-o(1)}\) source rows, because equality patching costs O(E\_N N). Thus a proof-mediated nonfunctional escape from natural proofs must be robust against sparse-row leakage; this is route-level progress, not a P vs NP proof.

\paragraph{Pressure.} No theorem yet converts an arbitrary SAT-semantic nonfunctional certificate into either a functional natural property or a sparse-row residual circuit. The row gate is a necessary falsifier audit, not a complete proof-system lower-bound framework.

\paragraph{Derivation steps.}
\begin{enumerate}
\item Fix \(H_{N}^{\mathrm{SAT}}(x)=xor_{j\le M_{N}}\) SAT(G\_N(x,j)) with \(M_{N}=N^{q+o(1)}\), L\_N=(log N)\textasciicircum{}d, and \(C(G_{N})\le N^{\rho+o(1)}\).
\item Assume SAT has circuits of size \(n^{c}\) at length \(L_{N}\). Substituting the same SAT circuit into every generated slot gives C(H\_N\textasciicircum{}\{SAT\})<=N\textasciicircum{}\{rho+o(1)\}+N\textasciicircum{}\{q+o(1)\}L\_N\textasciicircum{}c.
\item Since L\_N\textasciicircum{}c=N\textasciicircum{}\{o(1)\}, the transfer is \(N^{max{\rho,q}+o(1)}\). Therefore C(H\_N\textasciicircum{}\{SAT\})>N\textasciicircum{}\{1+epsilon\} with \(max{\rho,q}<1+\epsilon\) implies SAT notin SIZE(n\textasciicircum{}c) for every fixed c.
\item Now suppose a residual circuit \(D_{N}\) of size \(S_{D}\) agrees with \(H_{N}^{\mathrm{SAT}}\) outside \(R_{N}\), with \(|R_{N}|=E_{N}\). For each exceptional x, add an equality test to x and xor the correction bit. This costs O(E\_N N) gates.
\item Thus C(H\_N\textasciicircum{}\{SAT\})N\textasciicircum{}\{1+epsilon-o(1)\} is false.
\item Hence a nonfunctional proof object that escapes label-functional naturalization must still prove residual hardness after deleting \(N^{\epsilon-o(1)}\) rows; otherwise it supplies its own falsifier.
\item Validity regime check: the derivation uses the definitions and assumptions listed in the candidate lemmas, preserves the same source object and target object, and is valid only in the stated P vs NP branch model with its declared circuit class, length regime, and advice/uniformity conditions.
\end{enumerate}

\paragraph{Equation 1: xor\_carrier.}
\[
H_N^{SAT}(x)=\bigoplus_{j=1}^{M_N} SAT(G_N(x,j)),\quad M_N\le N^{q+o(1)},\ L_N=(\log N)^d
\]
\emph{Valid when:} Valid when the source object is the explicitly generated beta-zero XOR carrier and every generated query has polylogarithmic target length.

\paragraph{Equation 2: target\_substitution\_upper\_bound.}
\[
SAT\in SIZE(n^c)\ \Rightarrow\ C(H_N^{SAT})\le N^{\rho+o(1)}+N^{q+o(1)}L_N^c=N^{\max{\rho,q}+o(1)}
\]
\emph{Valid when:} Applies after the reduction preserves nonuniform circuit model identity and the target query length satisfies \(L_{N}=N^{o(1)}\).

\paragraph{Equation 3: row\_patch\_upper\_bound.}
\[
D_N=H_N^{SAT}\text{ off }R_N\ \Rightarrow\ C(H_N^{SAT})\le S_D+O(E_NN)+O(1),\quad E_N=|R_N|
\]
\emph{Valid when:} Valid when the residual circuit is fixed for length N and exceptional source rows are patched by equality tests.

\paragraph{Equation 4: source\_survival\_threshold.}
\[
C(H_N^{SAT})> N^{1+\epsilon}\ \Rightarrow\ S_D\ge N^{1+\epsilon-o(1)}\text{ or }E_N\ge N^{\epsilon-o(1)}
\]
\emph{Valid when:} Requires the row-patch upper bound and the claimed barely-superlinear source lower bound in the same circuit model.

\paragraph{Next analysis.} Instantiate a named proof system Psys for \(\beta=0\) XOR certificates and test whether any accepted SAT-semantic certificate can be transformed into either a functional label property or a residual circuit wrong on fewer than \(N^{\epsilon-o(1)}\) rows.

\paragraph{Falsifiers.}
\begin{itemize}
\item sparse row residual compression with \(E_{N}<N^{\epsilon-o(1)}\) and \(S_{D}<N^{1+\epsilon-o(1)}\)
\item positive-beta black-box substitution cap q+beta c preventing all-exponent transfer
\item functional dense certificate naturalization when soundness holds for arbitrary label vectors
\item row-selection advice leakage hidden inside a SAT-semantic proof object
\item low constructive incidence rank or label-free quotient sketch making the XOR carrier easy
\end{itemize}

\subsection*{cycle-0074: Hardness-Magnification Bridge Realization: beta-zero XOR row-residual audit}

\paragraph{Execution.} research; azure\_responses\_api; attempt primary; 600.001 seconds; 3 lemmas, 5 equations, 6 derivation steps, 5 falsifiers.

\paragraph{Learning.} Cycle 74 upgrades the focus packet: the beta-zero XOR bridge and row-residual patching bound are theorem-shaped and valid under explicit batched-generator assumptions. The live obstruction is now narrower: prove an \(\mathrm{EF}+\mathrm{SAT}_{U}\) semantic extraction/no-extraction theorem, or quarantine proof-mediated certificates that cannot certify residual robustness.

\paragraph{Pressure.} No proof-analysis theorem currently forces a semantic \(\mathrm{EF}+\mathrm{SAT}_{U}\) certificate either to functionalize over all label tables or to localize its nonfunctionality to fewer than \(N^{\epsilon-o(1)}\) source rows. EF may use global SAT structure shared by many rows through the same formula labels.

\paragraph{Derivation steps.}
\begin{enumerate}
\item Fix a batched generator \(G_{N}\) with \(M_{N}\le N^{q+o(1)}\), L\_N=(log N)\textasciicircum{}d, and circuit size \(N^{\rho+o(1)}\); define \(H_{N}^\mathrm{SAT}\) by XOR aggregation.
\item Assume SAT has circuits of size \(n^{c}\). Substitute the \(length-L_{N}\) SAT circuit into every generated query and obtain C(H\_N\textasciicircum{}SAT)<=N\textasciicircum{}\{max\{rho,q\}+o(1)\} because L\_N\textasciicircum{}c=N\textasciicircum{}\{o(1)\}.
\item If \(max{\rho,q}N^{1+\epsilon}\) contradicts SAT in P/poly; this is a valid route-level bridge, not a proof of the source lower bound.
\item Let \(D_{N}\) agree with \(H_{N}^\mathrm{SAT}\) outside \(R_{N}\). Hardwire equality tests for all rows in \(R_{N}\) where \(D_{N}\) is wrong, XOR the correction, and get C(H\_N\textasciicircum{}SAT)<=C(D\_N)+O(|R\_N|N).
\item Instantiate Psys as \(\mathrm{EF}+\mathrm{SAT}_{U}\). If its accepted certificates are functional over arbitrary label tables and dense at \(2^{-O(L_{N})}\), the certificate property naturalizes. If soundness is only for the real SAT label table, the route escapes that barrier but must prove row-residual robustness.
\item Tried to derive a sparse \(D_{N},R_{N}\) from an arbitrary nonfunctional \(\mathrm{EF}+\mathrm{SAT}_{U}\) proof by collecting rows whose SAT-semantic axioms are used essentially.
\end{enumerate}

\paragraph{Equation 1: carrier\_definition.}
\[
H_N^{SAT}(x)=\bigoplus_{j=1}^{M_N} SAT(G_N(x,j)),\quad M_N\le N^{q+o(1)},\quad L_N=(\log N)^d.
\]
\emph{Valid when:} Valid when each G\_N(x,j) is a well-formed \(length-L_{N}\) SAT instance and XOR is implemented by ordinary Boolean circuits.

\paragraph{Equation 2: beta\_zero\_substitution.}
\[
SAT\in SIZE(n^c)\ \Longrightarrow\ C(H_N^{SAT})\le N^{\rho+o(1)}+M_N L_N^c+O(M_N)=N^{\max{\rho,q}+o(1)}.
\]
\emph{Valid when:} Applies after assuming a length-n SAT circuit of size \(n^{c}\) and substituting it uniformly into every generated query slot.

\paragraph{Equation 3: target\_preserving\_implication.}
\[
C(H_N^{SAT})> N^{1+\epsilon}\ \text{i.o.}\ \wedge\ \max{\rho,q}< 1+\epsilon\ \Longrightarrow\ SAT\notin P/poly.
\]
\emph{Valid when:} Applies when the source lower bound occurs on infinitely many valid source lengths and the substitution upper bound would hold on every such length.

\paragraph{Equation 4: row\_patching\_bound.}
\[
D_N=H_N^{SAT}\ \text{on}\ {0,1}^N\setminus R_N\ \Longrightarrow\ C(H_N^{SAT})\le C(D_N)+O(|R_N|N).
\]
\emph{Valid when:} Valid when \(R_{N}\) is explicitly hardwired and equality to an N-bit source row costs O(N) gates in the same circuit model.

\paragraph{Equation 5: residual\_falsifier\_threshold.}
\[
C(D_N)+O(|R_N|N)< N^{1+\epsilon}\ \Longrightarrow\ (D_N,R_N)\ \text{falsifies}\ C(H_N^{SAT})> N^{1+\epsilon}.
\]
\emph{Valid when:} Requires C(D\_N) and \(|R_{N}|N\) to be measured at the same source length and against the same exponent threshold.

\paragraph{Next analysis.} Formalize \(\mathrm{EF}+\mathrm{SAT}_{U}\) syntax with label table \(U_{N}\), prove the functional dense naturalization lemma cleanly, then attempt a finite extraction theorem: every nonfunctional accepted proof either identifies a row set \(R_{N}\) and approximator \(D_{N}\) or uses a nonrelativizing SAT principle outside the row-residual audit.

\paragraph{Falsifiers.}
\begin{itemize}
\item falsifier: a circuit \(D_{N}\) of size \(N^{1+\epsilon-\gamma}\) agreeing with \(H_{N}^\mathrm{SAT}\) except on at most \(N^{\epsilon-\gamma}\) source rows.
\item falsifier: any positive-beta black-box interface claiming all-exponent SAT lower bounds from a fixed \(N^{1+\epsilon}\) source lower bound.
\item falsifier: hidden advice in \(G_{N}\) or in the certificate search procedure that makes the source lower bound target-equivalent.
\item falsifier: a dense functional \(\mathrm{EF}+\mathrm{SAT}_{U}\) certificate family, which is valid as a natural property but invalid as a claimed natural-proofs escape.
\item falsifier: a batched-generator accounting failure where producing all queries costs \(M_{N}\) \(N^{\rho}\) rather than N\textasciicircum{}\{rho\}+N\textasciicircum{}\{q+o(1)\}.
\end{itemize}

\subsection*{cycle-0075: Hardness-Magnification Bridge Realization}

\paragraph{Execution.} research; azure\_responses\_api; attempt compact\_retry; 600.01 seconds; 1 lemmas, 3 equations, 5 derivation steps, 4 falsifiers.

\paragraph{Learning.} Cycle 75 turns the Psys gap into a precise \(\mathrm{EF}+\mathrm{SAT}_{U}\) audit. Ordinary semantic soundness is insufficient: a valid proof-mediated beta-zero XOR route must either be nonlarge/nonconstructive/direct or certify row-residual robustness against subthreshold patching.

\paragraph{Pressure.} No theorem yet extracts row-residual robustness from short \(\mathrm{EF}+\mathrm{SAT}_{U}\) semantic proofs, and no construction yet proves such extraction impossible.

\paragraph{Derivation steps.}
\begin{enumerate}
\item Assume SAT has size \(n^{c}\) circuits. Substitute the circuit into each \(L_{N}-bit\) generated query of \(H_{N}\) and charge the generator once, giving the beta-zero transfer bound.
\item If \(q,rhoN^{1+\epsilon-o(1)}\), the transfer contradicts \(\mathrm{SAT}\in \mathrm{SIZE}(n^{c})\) for every fixed c; this is route-level \(\mathrm{SAT}\notin \mathrm{P}/\mathrm{poly}\) only from the source bound.
\item Given \(D_{N}\) agreeing with \(H_{N}\) off \(R_{N}\), patch \(D_{N}\) by equality tests for rows in \(R_{N}\). If \(S_{D}+O(|R_{N}|N)<N^{1+\epsilon}\), the source lower bound is false.
\item For \(\mathrm{EF}+\mathrm{SAT}_{U}\) certificates, dense functional soundness yields a natural label property; nonfunctional SAT-semantic soundness must therefore certify the row-residual gate or remain only a heuristic witness.
\item Validity regime check: the derivation uses the definitions and assumptions listed in the candidate lemmas, preserves the same source object and target object, and is valid only in the stated P vs NP branch model with its declared circuit class, length regime, and advice/uniformity conditions.
\end{enumerate}

\paragraph{Equation 1: beta0\_substitution.}
\[
C(H_N^{SAT}) \le C(G_N)+M_N\cdot L_N^c+O(M_N\log M_N)=N^{\max{\rho,q}+o(1)}
\]
\emph{Valid when:} Valid when L\_N=(log N)\textasciicircum{}d, \(M_{N}\le N^{q+o(1)}\), \(C(G_{N})\le N^{\rho+o(1)}\), and a size \(n^{c}\) SAT circuit is substituted nonuniformly into every generated query.

\paragraph{Equation 2: row\_residual\_patch.}
\[
D_N=H_N^{SAT}\text{ off }R_N \Rightarrow C(H_N^{SAT})\le S_D+O(|R_N|N)
\]
\emph{Valid when:} Valid when unrestricted circuits may hardwire the exceptional N-bit rows and compare equality to each row before output correction.

\paragraph{Equation 3: robustness\_gate.}
\[
\forall(D_N,R_N)\;[D_N=H_N^{SAT}\text{ off }R_N \Rightarrow S_D+O(|R_N|N)\ge N^{1+\epsilon-o(1)}]
\]
\emph{Valid when:} Requires a claimed source lower bound C(H\_N\textasciicircum{}SAT)>N\textasciicircum{}\{1+epsilon-o(1)\} in the same nonuniform circuit model.

\paragraph{Next analysis.} Instantiate \(\mathrm{EF}+\mathrm{SAT}_{U}\) proofs as circuits over \(\mathrm{SAT}_{U}\) label atoms, define functionalization formally, and try to construct a \(\theta<\epsilon\) residual-leaky transcript or prove extraction to the robustness gate.

\paragraph{Falsifiers.}
\begin{itemize}
\item Subthreshold residual falsifier: find \(D_{N},R_{N}\) with \(D=H\) off R and \(S_{D}+O(|R|N)<N^{1+\epsilon-o(1)}\).
\item Naturalization falsifier: \(\mathrm{EF}+\mathrm{SAT}_{U}\) certificates are dense, \(2^{O(L)}\) searchable, and functionally sound over all label vectors on \(U_{N}\).
\item Advice leakage falsifier: the certificate or generator encodes an exceptional row set of size \(N^{\theta}\) with \(\theta<\epsilon\).
\item Beta/q falsifier: any hidden query-length exponent \(\beta>0\) or overhead \(max{\rho,q}\ge 1+\epsilon\) collapses the all-exponent target transfer.
\end{itemize}

\subsection*{cycle-0076: Hardness-Magnification Bridge Realization}

\paragraph{Execution.} research; azure\_responses\_api; attempt compact\_retry; 600.005 seconds; 1 lemmas, 3 equations, 4 derivation steps, 4 falsifiers.

\paragraph{Learning.} Cycle 76 closes the Psys instantiation to \(\mathrm{EF}+\mathrm{SAT}_{U}\) and isolates the exact missing certificate strength: ordinary SAT-semantic soundness is not valid evidence unless it proves row-residual robustness above the patching threshold.

\paragraph{Pressure.} No extraction/no-extraction theorem yet maps short \(\mathrm{EF}+\mathrm{SAT}_{U}\) proofs to functional label properties or to row-residual robustness against all subthreshold \(D_{N}\).

\paragraph{Derivation steps.}
\begin{enumerate}
\item Substitute any SIZE(L\textasciicircum{}c) SAT circuit into every generated query; \(\beta=0\) removes the c-dependent exponent and gives the valid target bridge when \(max{\rho,q}<1+\epsilon\).
\item Split \(\mathrm{EF}+\mathrm{SAT}_{U}\) certificates: dense functional soundness over arbitrary label vectors triggers the prior naturalization audit; SAT-semantic soundness alone avoids that barrier but loses functional control.
\item For nonfunctional certificates, impose the row-residual test: any \(D_{N}\) correct off \(R_{N}\) yields an exact \(H_{N}\) circuit by equality patching, with exponent max\{s,e+1\}.
\item Tried to extract from a short \(\mathrm{EF}+\mathrm{SAT}_{U}\) derivation either a functional label property or an explicit \(Robust_{N}\) certificate; ordinary semantic \(\mathrm{SAT}_{U}\) axioms do not force either outcome.
\end{enumerate}

\paragraph{Equation 1: carrier\_definition.}
\[
H_N^{\mathrm{SAT}}(x)=\bigoplus_{j=1}^{M_N} \mathrm{SAT}(G_N(x,j)),\quad M_N\le N^{q+o(1)},\ L_N=(\log N)^d
\]
\emph{Valid when:} Valid when the source object is the beta-zero XOR carrier with a uniformly generated polylogarithmic SAT-query family.

\paragraph{Equation 2: beta0\_transfer.}
\[
\mathrm{SAT}_{L_N}\in \mathrm{SIZE}(L_N^c)\Rightarrow C(H_N)\le N^{\rho+o(1)}+N^{q+o(1)}L_N^c=N^{\max{\rho,q}+o(1)}
\]
\emph{Valid when:} Applies after the reduction preserves advice-freeness, fixed target exponent c, and polylogarithmic query length \(L_{N}\).

\paragraph{Equation 3: row\_residual\_patch.}
\[
D_N=H_N\text{ off }R_N,\ C(D_N)=N^{s+o(1)},\ |R_N|=N^{e+o(1)}\Rightarrow C(H_N)\le N^{\max{s,e+1}+o(1)}
\]
\emph{Valid when:} Requires source rows to be N-bit strings so each exceptional row is patchable by an O(N)-gate equality test.

\paragraph{Next analysis.} Formalize Robust\_N(s,e) for \(\mathrm{EF}+\mathrm{SAT}_{U}\) and test whether a short proof can remain sound after replacing \(H_{N}\) by \(D_{N}\) off a sparse row set.

\paragraph{Falsifiers.}
\begin{itemize}
\item falsifier: a subthreshold \(D_{N}\) with \(C(D_{N})<N^{1+\epsilon-\eta}\) and fewer than \(N^{\epsilon-\eta}\) wrong source rows gives \(C(H_{N})<N^{1+\epsilon-o(1)}\).
\item falsifier: a dense searchable \(\mathrm{EF}+\mathrm{SAT}_{U}\) certificate that is functional over label vectors becomes a valid natural-property object for the label model.
\item falsifier: generator or verifier advice that selects exceptional rows nonuniformly breaks the target-preserving condition and invalidates the bridge.
\item falsifier: any hidden positive-beta query growth restores the transfer cap q+beta c and cannot exclude all SAT polynomial exponents.
\end{itemize}

\subsection*{cycle-0077: Hardness-Magnification Bridge Realization}

\paragraph{Execution.} research; azure\_responses\_api; attempt primary; 600.003 seconds; 3 lemmas, 5 equations, 6 derivation steps, 6 falsifiers.

\paragraph{Learning.} Cycle 77 sharpens the \(\mathrm{EF}+\mathrm{SAT}_{U}\) packet: direct lower-bound certificates automatically pass the explicit row-residual audit, so the real gap is not patching but semantic extraction. A new occurrence-oracle row-touch no-go gives a concrete falsifier for relativizing local certificate methods and leaves only SAT-semantic nonrelativizing, nonlarge/nonconstructive, or direct routes alive.

\paragraph{Pressure.} No theorem yet shows that actual \(\mathrm{EF}+\mathrm{SAT}_{U}\) certificates for the intended \(\mathrm{SAT}_{U}\) carrier either reduce to the row-touch occurrence-local model or yield a nonrelativizing extraction from certificate acceptance to LB\_N(T\_N).

\paragraph{Derivation steps.}
\begin{enumerate}
\item Formalized LB\_N(T\_N) as the direct statement \(C(H_{N}^\mathrm{SAT})>T_{N}\) in the same nonuniform circuit model used by the beta-zero transfer.
\item Constructed the standard equality-patch circuit: compute \(D_{N}\), test x against each exceptional row r in \(R_{N}\), and xor the needed correction bit on matches.
\item Derived that any direct proof of LB\_N(T\_N) automatically excludes all explicit residual pairs (\(D_{N},R_{N}\)) with \(|D_{N}|+aN|R_{N}|<T_{N}\).
\item Tried to lift ordinary \(\mathrm{EF}+\mathrm{SAT}_{U}\) semantic certificates to LB\_N(T\_N) by treating accepted certificates as evidence for row-residual robustness.
\item Found the obstruction: if the certificate is only row-local and occurrence-relativizing, unseen occurrence labels can be assigned to make the carrier trivial outside the touched rows while preserving transcript acceptance.
\item Separated the remaining live case as a nonrelativizing SAT-realization extraction problem, not as a consequence of ordinary semantic soundness.
\end{enumerate}

\paragraph{Equation 1: carrier\_definition.}
\[
H_N^{SAT}(x)=\bigoplus_{j\le M_N} SAT(G_N(x,j)),\quad |G_N(x,j)|=L_N=(\log N)^d.
\]
\emph{Valid when:} Valid when \(G_{N}\) is a fixed advice-free batched generator and rows are N-bit source inputs.

\paragraph{Equation 2: beta\_zero\_transfer.}
\[
C(H_N^{SAT})\le C(G_N)+O(M_N\,C_{SAT}(L_N)+M_NL_N)\le N^{\max{\rho,q}+o(1)}.
\]
\emph{Valid when:} Applies when SAT has \(L^{c}-size\) circuits for fixed c and L\_N=(log N)\textasciicircum{}d makes L\_N\textasciicircum{}c=N\textasciicircum{}\{o(1)\}.

\paragraph{Equation 3: explicit\_row\_patching.}
\[
D_N=H_N^{SAT}\text{ off }R_N\quad\Rightarrow\quad C(H_N^{SAT})\le |D_N|+aN|R_N|+O(|R_N|).
\]
\emph{Valid when:} Requires \(D_{N}\) to agree with \(H_{N}^{\mathrm{SAT}}\) outside the explicitly listed row set \(R_{N}\) and constants to be charged nonuniformly.

\paragraph{Equation 4: residual\_threshold.}
\[
|D_N|+aNE_N< T_N\quad\Rightarrow\quad \neg LB_N(T_N)\text{ for }E_N=|R_N|.
\]
\emph{Valid when:} Valid when the source lower-bound claim is an exact circuit lower bound for \(H_{N}^{\mathrm{SAT}}\) at threshold \(T_{N}\).

\paragraph{Equation 5: conditional\_target\_consequence.}
\[
LB_N(N^{1+\epsilon-o(1)})\wedge \max{\rho,q}< 1+\epsilon\quad\Rightarrow\quad \forall c\; SAT\notin SIZE(n^c)\text{ on matched lengths}.
\]
\emph{Valid when:} Applies after the beta-zero transfer preserves advice accounting and the lower bound holds on the same infinite or eventual length set.

\paragraph{Next analysis.} Instantiate Touch(Pi) for a concrete \(\mathrm{EF}+\mathrm{SAT}_{U}\) encoding of circuit lower-bound certificates: count whether proofs of size \(N^{k}\) can mention \(sub-N^\epsilon\) rows, and test if formula-equality/SAT-realization axioms break the occurrence-oracle extension.

\paragraph{Falsifiers.}
\begin{itemize}
\item explicit residual falsifier: a circuit \(D_{N}\) and explicit bad row list \(R_{N}\) with \(|D_{N}|+aN|R_{N}|<N^{1+\epsilon-o(1)}\) invalidates any claimed direct source lower bound.
\item positive-beta falsifier: any repaired carrier with target query length \(N^\beta\) for \(\beta>0\) reintroduces the q+beta*c transfer cap.
\item occurrence-oracle falsifier: a row-local certificate touching \(B_{N}<N^{\epsilon-o(1)}\) rows is compatible with an easy occurrence-label carrier.
\item functional-density falsifier: a dense constructive certificate property over label tables becomes a natural-property route unless it is nonlarge, nonconstructive, or nonfunctional.
\item advice-leakage falsifier: uncharged generator, proof, or row-selection advice can hide the residual repair below the recorded threshold.
\item target-equivalence falsifier: a direct EF proof of LB\_N(T\_N) may already encode the desired SAT lower bound, so it is not a weak-source magnification achievement.
\end{itemize}

\subsection*{cycle-0078: Hardness-Magnification Bridge Realization}

\paragraph{Execution.} research; azure\_responses\_api; attempt primary; 600.008 seconds; 3 lemmas, 4 equations, 6 derivation steps, 5 falsifiers.

\paragraph{Learning.} Cycle 78 converts the row-residual audit into a theorem-shaped no-go: functional row-local certificates for the beta-zero XOR carrier require \(N^{\epsilon-o(1)}\) touched rows at source threshold \(N^{1+\epsilon-o(1)}\). This sharpens the \(\mathrm{EF}+\mathrm{SAT}_{U}\) obstruction from ordinary soundness to a concrete extraction fork: either the proof is direct or nonlarge/nonconstructive, or it must use real SAT semantics in a nonrelativizing way that defeats the private-row falsifier.

\paragraph{Pressure.} The derivation is complete only for functional row-value certificates. \(\mathrm{EF}+\mathrm{SAT}_{U}\) proofs are SAT-semantic over formula labels, so a nonrelativizing extraction theorem or a private-row no-go formalization is still needed.

\paragraph{Derivation steps.}
\begin{enumerate}
\item Assume a row-value certificate A\_N(pi,h) is functional and \(s_{N}-sound\): \(A_{N}(pi,h)=1\) implies \(C(h)>s_{N}\) for all Boolean h on N-bit rows.
\item Let \(R_{pi}\) be the set of rows whose h-values can affect acceptance, with \(|R_{pi}|=T_{N}\), and suppose \(A_{N}(pi,H_{N}^\mathrm{SAT})=1\).
\item Define \(h_{R}\) to agree with \(H_{N}^\mathrm{SAT}\) on \(R_{pi}\) and to be zero outside \(R_{pi}\). Acceptance is unchanged because \(A_{N}\) only sees rows in \(R_{pi}\).
\item Build \(h_{R}\) by OR-ing equality tests for the rows in \(R_{pi}\) and hardwiring the corresponding bits, giving C(h\_R)<=O(T\_N N).
\item If T\_N=o(s\_N/N), then \(C(h_{R})<s_{N}\), contradicting functional soundness. Hence functional certificates proving the source threshold must touch \(N^{\epsilon-o(1)}\) rows.
\item Combine with the beta-zero bridge: a valid source lower bound still transfers to \(\mathrm{SAT}\notin \mathrm{P}/\mathrm{poly}\) when \(max{\rho,q}<1+\epsilon\), but the proof-mediated route must now pass the row-touch or SAT-semantic extraction fork.
\end{enumerate}

\paragraph{Equation 1: beta\_zero\_xor\_carrier.}
\[
H_N^SAT(x)=bigoplus_{j\le M_N} SAT(G_N(x,j))
\]
\emph{Valid when:} Valid when \(G_{N}\) is the named batched generator and every query length is L\_N=(log N)\textasciicircum{}d.

\paragraph{Equation 2: target\_substitution\_upper\_bound.}
\[
C(H_N^SAT) \le N^{rho+o(1)} + N^{q+o(1)} L_N^c = N^{max{rho,q}+o(1)}
\]
\emph{Valid when:} Valid when SAT on \(L_{N}-bit\) inputs has circuits of size \(L_{N}^{c}\) and \(\beta\) equals zero.

\paragraph{Equation 3: row\_residual\_patching.}
\[
C(H_N^SAT) \le S_D + O(|R_N| N)
\]
\emph{Valid when:} Applies when \(D_{N}\) agrees with \(H_{N}^\mathrm{SAT}\) outside \(R_{N}\) and row equality tests are available nonuniformly.

\paragraph{Equation 4: functional\_touch\_lower\_bound.}
\[
T_N \ge \Omega(s_N/N) = N^{epsilon-o(1)} for s_N=N^{1+epsilon-o(1)}
\]
\emph{Valid when:} Requires functional soundness over arbitrary row functions and acceptance depending on only \(T_{N}\) row values.

\paragraph{Next analysis.} Instantiate \(\mathrm{EF}+\mathrm{SAT}_{U}\) row occurrence access formally: define the touched-row set of a proof verifier, test whether standard EF reflection over \(\mathrm{SAT}_{U}\) extracts to row-value functionality, and if not construct the private-row oracle countermodel as a finite theorem.

\paragraph{Falsifiers.}
\begin{itemize}
\item A subthreshold residual circuit \(D_{N}\) with fewer than \(N^{\epsilon-o(1)}\) wrong rows is a falsifier for the claimed source lower bound.
\item An accepted functional certificate touching o(N\textasciicircum{}\{epsilon\}) rows is a falsifier for proof-mediated soundness at threshold \(N^{1+\epsilon-o(1)}\).
\item A relativizing occurrence-local \(\mathrm{EF}+\mathrm{SAT}_{U}\) argument that survives the private-row oracle is a falsifier for the claimed SAT-semantic escape.
\item A positive target-query exponent \(\beta\) reintroduces the transfer cap q+beta c and falsifies all-exponent target exclusion.
\item Any hidden nonuniform list of exceptional rows larger than the accounted proof size is advice leakage and invalidates the route certificate.
\end{itemize}

\subsection*{cycle-0079: Hardness-Magnification Bridge Realization: beta-zero XOR row-forcing audit}

\paragraph{Execution.} research; azure\_responses\_api; attempt primary; 600.004 seconds; 3 lemmas, 4 equations, 6 derivation steps, 5 falsifiers.

\paragraph{Learning.} Cycle 79 replaces vague SAT-semantic escape language by the full-row-forcing parameter ForcedRows(Q). A proof-mediated beta-zero XOR certificate that treats generated SAT instances as occurrence labels must force at least \(N^{\epsilon-o(1)}\) complete parity rows; otherwise a private occurrence-label spoof gives a subthreshold sparse-row circuit. The valid contribution is a scoped no-go and an exact extraction obligation for full \(\mathrm{EF}+\mathrm{SAT}_{U}\), not progress on the main conjecture itself.

\paragraph{Pressure.} The derivation is a proof only for \(\mathrm{EF}+\mathrm{SAT}_{U}^{occ}\) or functional label-local certificates. Full \(\mathrm{EF}+\mathrm{SAT}_{U}\) may use real formula syntax to constrain unmentioned generated instances, so an extraction theorem for \(Q_{pi}\) or a nonrelativizing semantic certificate is still missing.

\paragraph{Derivation steps.}
\begin{enumerate}
\item Fix an accepted occurrence-relativizing certificate pi and let \(Q_{pi}\) be the generated SAT occurrences whose labels affect its transcript.
\item Define \(A_{Q}\) to agree with real SAT on \(Q_{pi}\). For each source row x not in ForcedRows(Q\_pi), choose one unmentioned coordinate j\_x and set its label so that xor\_j \(A_{Q}(G_{N}(x,j))=0\).
\item Then \(H_{N}^{A_{Q}}\) is supported only on ForcedRows(Q\_pi). A circuit computes it by checking whether x is one of these rows and hardwiring the residual parity values; the size is N\textasciicircum{}\{rho+o(1)\}+O(|ForcedRows(Q\_pi)| log N).
\item If \(|ForcedRows(Q_{pi})|<N^{\epsilon-o(1)}\) and \(\rho<1+\epsilon\), this circuit has size below \(N^{1+\epsilon-o(1)}\).
\item Occurrence-relativizing transcript invariance makes pi accept equally under SAT and under \(A_{Q}\), so functional soundness would assert the same lower bound for the easy spoofed carrier. Contradiction.
\item Validity regime check: the derivation uses the definitions and assumptions listed in the candidate lemmas, preserves the same source object and target object, and is valid only in the stated P vs NP branch model with its declared circuit class, length regime, and advice/uniformity conditions.
\end{enumerate}

\paragraph{Equation 1: beta0\_substitution\_upper\_bound.}
\[
C(H_N^{SAT}) \le N^{\rho+o(1)} + M_N\,L_N^c = N^{\max{\rho,q}+o(1)}
\]
\emph{Valid when:} Valid when SAT on length \(L_{N}\) formulas has circuits of size \(L_{N}^{c}\) and the beta-zero generator makes only \(M_{N}\le N^{q+o(1)}\) polylog-length queries.

\paragraph{Equation 2: sparse\_row\_patch\_upper\_bound.}
\[
C(F) \le S_D + O(E_N n_N) + N^{o(1)} \le S_D + N^{1+o(1)}E_N/N
\]
\emph{Valid when:} Applies when a circuit D of size \(S_{D}\) agrees with the row function F except on \(E_{N}\) explicitly hardwired rows of \(n_{N}=log\) N bits.

\paragraph{Equation 3: forced\_row\_spoof\_size.}
\[
C(H_N^{A_Q}) \le N^{\rho+o(1)} + O(|ForcedRows(Q)|\,n_N)
\]
\emph{Valid when:} Requires the occurrence-label spoof to zero every row outside ForcedRows(Q) and allows the remaining forced rows to be patched by equality gates.

\paragraph{Equation 4: row\_forcing\_threshold.}
\[
|ForcedRows(Q_\pi)| \ge N^{\epsilon-o(1)}
\]
\emph{Valid when:} Necessary when an occurrence-relativizing certificate transcript soundly asserts \(C(H_{N})>N^{1+\epsilon-o(1)}\) against all label assignments consistent with its queried occurrences.

\paragraph{Next analysis.} Instantiate \(\mathrm{Psys}=\mathrm{EF}+\mathrm{SAT}_{U}^{occ}\): define the proof parser that extracts \(Q_{pi}\), prove acceptance transcript invariance under \(A_{Q}\), and compute whether known EF encodings force full rows or merely individual generated formula occurrences.

\paragraph{Falsifiers.}
\begin{itemize}
\item falsifier: a subthreshold residual circuit \(D_{N}\) with fewer than \(N^{\epsilon-o(1)}\) bad rows gives C(H\_N\textasciicircum{}SAT)<N\textasciicircum{}\{1+epsilon-o(1)\} by equality patching.
\item falsifier: an accepted occurrence-relativizing certificate pi with \(|ForcedRows(Q_{pi})|<N^{\epsilon-o(1)}\) is unsound for the claimed source lower bound.
\item falsifier: positive \(\beta\) target-query length reintroduces the transfer cap q+beta c and destroys all-exponent SAT lower-bound transfer.
\item falsifier: if the certificate proof string or verifier advice encodes more than \(N^{1+\epsilon-o(1)}\) row information, the source lower bound is being assumed rather than certified.
\item falsifier: ordinary \(\mathrm{EF}+\mathrm{SAT}_{U}\) soundness without row-residual robustness does not imply a valid source lower bound.
\end{itemize}

\subsection*{cycle-0080: Hardness-Magnification Bridge Realization}

\paragraph{Execution.} research; azure\_responses\_api; attempt primary; 600.005 seconds; 4 lemmas, 4 equations, 7 derivation steps, 5 falsifiers.

\paragraph{Learning.} Cycle 80 discharges the scoped \(\mathrm{EF}+\mathrm{SAT}_{U}^{occ}\) proof packet: occurrence-local functional certificates for the beta-zero XOR carrier must force at least \(N^{\epsilon-o(1)}\) complete parity rows, or a private occurrence-label spoof plus sparse-row equality patch gives a valid subthreshold falsifier. The live gap is no longer ordinary soundness; it is the full \(\mathrm{EF}+\mathrm{SAT}_{U}\) extraction of real SAT-semantic reasoning into syntactic \(ForcedRows\), or a deliberate restriction of the paper theorem to \(\mathrm{EF}+\mathrm{SAT}_{U}^{occ}\).

\paragraph{Pressure.} The derivation is valid for \(\mathrm{EF}+\mathrm{SAT}_{U}^{occ}\). It is blocked for full \(\mathrm{EF}+\mathrm{SAT}_{U}\) by the missing extraction lemma from real SAT semantic derivations to syntactic row forcing; full EF may use formula structure in a nonrelativizing way that private occurrence-label spoofing does not preserve.

\paragraph{Derivation steps.}
\begin{enumerate}
\item Define U\_N=\{(x,j):x in \({0,1}^{N}\), j<=M\_N\} and replace each SAT(G\_N(x,j)) by an independent occurrence label \(a_{x,j}\).
\item Parse an accepted occurrence-local transcript pi into \(Q_{pi}\) by collecting every queried occurrence, every locally derived equation, and every explicit row-parity tag; define ForcedRows(Q\_pi) as rows whose parity is fixed under all completions of \(Q_{pi}\).
\item Prove parser invariance by induction on the \(\mathrm{EF}+\mathrm{SAT}_{U}^{occ}\) derivation: each rule depends only on premises already in \(Q_{pi}\), so preserving \(Q_{pi}\) preserves acceptance.
\item For every unforced row choose a private completion making xor\_j \(a_{x,j}=0\) while preserving \(Q_{pi}\); forced rows keep their \(Q_{pi}-determined\) parity.
\item The spoofed carrier \(H_{A^{0}}\) is supported on R=ForcedRows(Q\_pi), so a nonuniform circuit hardwires R and uses O(N)-size equality tests to patch those rows, with total size S\_Q+O(|R|N).
\item If pi is functionally sound for the claimed lower bound \(s_{N}\), then \(H_{A^{0}}\) must have circuit complexity greater than \(s_{N}\); the patch circuit contradicts this whenever \(S_{Q}+O(|R|N)<s_{N}\).
\item Conclude that any valid occurrence-local functional certificate at the \(N^{1+\epsilon-o(1)}\) threshold must force \(N^{\epsilon-o(1)}\) complete rows, unless it is nonfunctional, nonconstructive, too large, or uses real SAT semantics outside \(\mathrm{EF}+\mathrm{SAT}_{U}^{occ}\).
\end{enumerate}

\paragraph{Equation 1: occurrence carrier.}
\[
H_A(x)=\bigoplus_{j\le M_N} a_{x,j}
\]
\emph{Valid when:} Valid when A is an arbitrary occurrence-label assignment over \(U_{N}\) and the row parity carrier is defined syntactically.

\paragraph{Equation 2: sparse-row spoof upper bound.}
\[
C(H_{A^0})\le S_Q+O(|ForcedRows(Q_\pi)|\,N)
\]
\emph{Valid when:} Valid when \(Q_{pi}\) is occurrence-local, \(S_{Q}\) bounds its evaluator, and only forced rows remain nonzero after spoofing.

\paragraph{Equation 3: forced-row threshold.}
\[
|ForcedRows(Q_\pi)|\ge \Omega((s_N-S_Q)/N)
\]
\emph{Valid when:} Valid when pi is functionally sound for all completions consistent with \(Q_{pi}\) and the claimed lower bound threshold is \(s_{N}\).

\paragraph{Equation 4: beta-zero target substitution.}
\[
C(H_N^{SAT})\le N^{\rho+o(1)}+M_N L_N^c=N^{\max{\rho,q}+o(1)}
\]
\emph{Valid when:} Valid when SAT has circuits of size \(n^{c}\), L\_N=(log N)\textasciicircum{}d, and the generator circuit has size \(N^{\rho+o(1)}\).

\paragraph{Next analysis.} Attempt the full \(\mathrm{EF}+\mathrm{SAT}_{U}\) boundary lemma: inspect a standard encoding of \(\mathrm{SAT}_{U}\) extension axioms and decide whether each semantic rule is occurrence-local, charges enough advice into \(S_{Q}\), or yields a nonrelativizing complete-row extraction.

\paragraph{Falsifiers.}
\begin{itemize}
\item Subthreshold forced rows: \(|ForcedRows(Q_{pi})|<N^{\epsilon-o(1)}\) with S\_Q=o(N\textasciicircum{}\{1+epsilon\}) yields a sparse-row patch circuit and is a direct falsifier for the claimed occurrence-local certificate.
\item Parser leakage: if \(Q_{pi}\) contains a truth table, a SAT solver, or \(H_{N}\) advice, then \(S_{Q}\) must be charged; if charged size is subthreshold it falsifies the source lower bound, otherwise the certificate is not barely-superlinear.
\item Nonfunctional soundness: a certificate valid only for the real SAT label vector escapes the occurrence-label theorem but loses the functional naturalization/no-go claim until a nonrelativizing extraction is proved.
\item Global parity equations: low-cost constraints fixing many row parities without touching complete rows falsify the atomic \(ForcedRows\) definition and require a parity-rank \(ForcedRows\) repair.
\item Positive-beta transfer remains an independent falsifier: if any replacement carrier reintroduces query length \(n^\beta\) with \(\beta>0\), the window q+beta c blocks all-exponent target separation.
\end{itemize}

\subsection*{cycle-0081: Hardness-Magnification Bridge Realization}

\paragraph{Execution.} research; azure\_responses\_api; attempt primary; 600.005 seconds; 3 lemmas, 4 equations, 5 derivation steps, 5 falsifiers.

\paragraph{Learning.} Cycle 81 separates the scoped occurrence-local theorem from the full \(\mathrm{EF}+\mathrm{SAT}_{U}\) ambition. The valid contribution is now sharper: private row spoofing is a genuine falsifier for occurrence-local functional certificates, while full \(\mathrm{EF}+\mathrm{SAT}_{U}\) requires a nonrelativizing semantic-to-occurrence extraction theorem because real \(\mathrm{SAT}_{U}\) axioms are not invariant under independent label flips.

\paragraph{Pressure.} Semantic-to-occurrence shadow extraction: no proof currently shows that full \(\mathrm{EF}+\mathrm{SAT}_{U}\) reasoning about actual SAT formulas can be converted into a functional occurrence-local \(Q_{pi}\) with subthreshold row residual and \(N^{o(1)}\) overhead.

\paragraph{Derivation steps.}
\begin{enumerate}
\item Start with the valid \(\mathrm{EF}+\mathrm{SAT}_{U}^{occ}\) theorem: if a functional occurrence-local transcript pi has fewer than \(N^{\epsilon-o(1)}\) complete forced rows, construct a private-label spoof \(A_{Q}\) and a parity circuit \(D_{N}\) differing from \(H_{N}^{\mathrm{SAT}}\) only on those rows.
\item Try to lift this to full \(\mathrm{EF}+\mathrm{SAT}_{U}\) by projecting each full SAT atom SAT(G\_N(x,j)) to a private occurrence label \(y_{x,j}\) and parsing the proof into \(Q_{pi}\).
\item Check transcript invariance: in \(\mathrm{EF}+\mathrm{SAT}_{U}^{occ}\), \(A_{Q}\) preserves all parsed constraints outside ForcedRows(Q\_pi); in full \(\mathrm{EF}+\mathrm{SAT}_{U}\), a line asserting SAT\_U(F) or a witness equivalence is tied to the real syntax of F.
\item Construct the countermodel to naive invariance: choose an occurrence F=G\_N(x,j) and set \(y_{F}\) to 1-SAT(F). The occurrence shadow can still satisfy uninterpreted constraints, but the full \(\mathrm{SAT}_{U}\) axiom for F is false.
\item Conclude that sparse-row patching remains a valid falsifier for occurrence-local certificates, but it cannot refute full \(\mathrm{EF}+\mathrm{SAT}_{U}\) certificates unless a shadow-extraction normal form covers every SAT-semantic violation by forced rows plus a subthreshold residual.
\end{enumerate}

\paragraph{Equation 1: beta\_zero\_xor\_carrier.}
\[
H_N^{\mathrm{SAT}}(x)=\bigoplus_{j\le M_N}\mathrm{SAT}(G_N(x,j)),\quad M_N\le N^{q+o(1)},\quad L_N=(\log N)^d
\]
\emph{Valid when:} Valid when the source object is the beta-zero generated XOR carrier with polylogarithmic target instances and an explicit row index j.

\paragraph{Equation 2: shadow\_extraction\_requirement.}
\[
\mathrm{Ext}(\pi)=(Q_\pi,\mathrm{ForcedRows}(Q_\pi),R_\pi),\quad |Q_\pi|\le |\pi|N^{o(1)},\quad |R_\pi|< N^{\epsilon-o(1)}
\]
\emph{Valid when:} Valid when a proposed full-EF transfer claims to reduce real SAT-semantic reasoning to occurrence-local constraints without exceeding the sparse-row patch threshold.

\paragraph{Equation 3: semantic\_invariance\_obligation.}
\[
A_Q\models Q_\pi\ \Longrightarrow\ \pi\ \text{remains sound except on}\ \mathrm{ForcedRows}(Q_\pi)\cup R_\pi
\]
\emph{Valid when:} Valid when private occurrence-label perturbations are required to simulate full \(\mathrm{SAT}_{U}\) lines and all violations are covered by syntactically forced complete rows plus a residual set.

\paragraph{Equation 4: \(full_{\mathrm{EF},blocking,fact}\).}
\[
\exists F,b:\ b\ne \mathrm{SAT}(F)\ \Rightarrow\ (y_F=b)\ \text{may satisfy an occurrence shadow but violates SAT\_U(F)}
\]
\emph{Valid when:} Valid when full \(\mathrm{EF}+\mathrm{SAT}_{U}\) contains semantic axioms tying SAT labels to actual formula syntax rather than to independent occurrence variables.

\paragraph{Next analysis.} Formalize the semantic-coupling-width thought experiment as a lemma: for paired formulas F and F OR z, show exactly which \(\mathrm{SAT}_{U}\) implication lines break private-label invariance and whether their violations can be charged to complete forced rows below the \(N^{\epsilon}\) patch threshold.

\paragraph{Falsifiers.}
\begin{itemize}
\item Full \(\mathrm{EF}+\mathrm{SAT}_{U}\) promotion falsifier: a \(\mathrm{SAT}_{U}\) axiom line whose truth depends on real formula syntax is not preserved by private occurrence-label perturbations.
\item Source lower-bound falsifier: any \(D_{N}\) agreeing with \(H_{N}^{\mathrm{SAT}}\) off fewer than \(N^{\epsilon-o(1)}\) rows and having \(C(D_{N})<N^{1+\epsilon-o(1)}\) gives a valid equality-patch upper bound.
\item Parser falsifier: Ext hardwires the real SAT label vector or \(H_{N}^{\mathrm{SAT}}\), making the certificate target-equivalent rather than proof-extractable.
\item Threshold falsifier: if \(max{\rho,q}\ge 1+\epsilon\), the beta-zero substitution bridge no longer yields the claimed all-exponent SAT consequence.
\item Naturalization falsifier: a dense constructive functional label certificate over \(U_{N}\) becomes a natural-property-style object unless it is nonlarge, nonconstructive, or genuinely SAT-semantic.
\end{itemize}

\subsection*{cycle-0082: Hardness-Magnification Bridge Realization}

\paragraph{Execution.} research; azure\_responses\_api; attempt compact\_retry; 600.007 seconds; 1 lemmas, 3 equations, 4 derivation steps, 4 falsifiers.

\paragraph{Learning.} Cycle 82 upgrades the proof-delivery packet by isolating a row-normal parser theorem: occurrence-local functional certificates are valid only if they force \(N^{\epsilon-o(1)}\) complete rows. The remaining gap is exactly full \(\mathrm{EF}+\mathrm{SAT}_{U}\) semantic-to-occurrence extraction.

\paragraph{Pressure.} The derivation is valid for row-normal occurrence-local transcripts; full \(\mathrm{EF}+\mathrm{SAT}_{U}\) may use real SAT syntax to create nonrelativizing dependencies not represented by \(Q_{pi}\).

\paragraph{Derivation steps.}
\begin{enumerate}
\item Parse pi into \(Q_{pi}\) by listing occurrence labels, extension equations, rule justifications, and explicit complete-row parity nodes; this is independent of the desired lower bound.
\item For x not in ForcedRows(pi), flip or set \(y_{x,j}\) privately so xor\_j \(y_{x,j}=0\) while preserving all row-normal dependencies appearing in \(Q_{pi}\).
\item Define spoofed carrier \(H^{A}\) to be zero off ForcedRows(pi); equality gates patch the forced rows, giving size \(O(|ForcedRows(pi)|N)+N^{o(1)}\).
\item Functional soundness would accept pi for \(H^{A}\), so if \(|ForcedRows(pi)|<N^{\epsilon-o(1)}\) it certifies hardness for an easy function, contradiction.
\end{enumerate}

\paragraph{Equation 1: beta0\_transfer.}
\[
C(H_N^{SAT}) \le N^{rho+o(1)} + M_N L_N^c = N^{max{rho,q}+o(1)}
\]
\emph{Valid when:} Valid when SAT has size \(L_{N}^{c}\) circuits and the beta-zero generator uses \(M_{N}=N^{q+o(1)}\) queries.

\paragraph{Equation 2: sparse\_row\_patch.}
\[
C(H_N) \le S_D + O(E_N N)
\]
\emph{Valid when:} Valid when \(D_{N}\) is an explicit circuit and errors are at most \(E_{N}\) source inputs.

\paragraph{Equation 3: forced\_row\_threshold.}
\[
|ForcedRows(pi)| \ge N^{epsilon-o(1)}
\]
\emph{Valid when:} Valid when the row-normal parser and private-label invariance assumptions are established.

\paragraph{Next analysis.} Construct the smallest EF extension-rule countermodel that hides row parity without a row-normal node, or prove a normalization that exposes such parity with subthreshold overhead.

\paragraph{Falsifiers.}
\begin{itemize}
\item falsifier: a transcript with fewer than \(N^{\epsilon-o(1)}\) forced rows and row-normal invariance gives a valid subthreshold spoof circuit.
\item falsifier: hidden extension gates computing row parity without syntactic exposure break the parser lemma and force a restriction of the theorem.
\item falsifier: SAT-semantic rules using formula syntax rather than occurrence labels evade private-label flips and require nonrelativizing extraction.
\item falsifier: q or \(\rho\) reaching 1+epsilon collapses the target-preserving beta-zero transfer window.
\end{itemize}

\subsection*{cycle-0083: Hardness-Magnification Bridge Realization}

\paragraph{Execution.} research; azure\_responses\_api; attempt primary; 600.01 seconds; 3 lemmas, 3 equations, 7 derivation steps, 5 falsifiers.

\paragraph{Learning.} Cycle 83 upgrades the proof-delivery packet from a parser slogan to a conditional theorem: row-separable row-normal functional occurrence certificates must force \(N^{\epsilon-o(1)}\) complete XOR rows, or private occurrence-label spoofing yields a valid subthreshold circuit for the induced carrier. The remaining open obligation is no longer ordinary soundness but the exact row-normalization/extraction step for broader \(\mathrm{EF}+\mathrm{SAT}_{U}\).

\paragraph{Pressure.} The proof is conditional on a row-separable row-normal parser. Arbitrary full \(\mathrm{EF}+\mathrm{SAT}_{U}\) can use real formula semantics, and arbitrary occurrence-local EF may contain short cross-row parity constraints unless the calculus forbids or charges them.

\paragraph{Derivation steps.}
\begin{enumerate}
\item Fix N, \(\epsilon\), q, \(\rho\) with \(max{q,\rho}<1+\epsilon\) and define \(H_{N}^\mathrm{SAT}(x)=xor_{j\le M_{N}}\) SAT(G\_N(x,j)).
\item Parse an accepted row-normal transcript pi into \(Q_{pi}\); define ForcedRows(Q\_pi) as rows whose parity is fixed over all assignments satisfying \(Q_{pi}\).
\item Use parser invariance: acceptance of pi depends only on \(Q_{pi}\), so any private label assignment A satisfying \(Q_{pi}\) preserves the accepted transcript.
\item If \(Q_{pi}\) is row-separable, choose \(A_{Q}\) so every unforced row has parity 0 and every forced row has its \(Q_{pi}-forced\) parity.
\item Build a circuit for \(H_{A_{Q}}\): output 0 by default and patch the forced rows with equality tests; size is \(O(C(Q_{pi}))+O(|ForcedRows|N)+N^{o(1)}\).
\item Functional certificate soundness says the same accepted pi certifies \(C(H_{A_{Q}})>N^{1+\epsilon-o(1)}\). If \(|ForcedRows|<N^{\epsilon-o(1)}\) and \(C(Q_{pi})=N^{o(1)}\), this contradicts the patch circuit.
\item Therefore any sound row-normal functional occurrence-local certificate must force \(N^{\epsilon-o(1)}\) complete rows, or leave the method class.
\end{enumerate}

\paragraph{Equation 1: beta\_zero\_transfer.}
\[
SAT in SIZE(n^c) => C(H_N^SAT) \le N^{rho+o(1)} + M_N L_N^c = N^{max{rho,q}+o(1)}
\]
\emph{Valid when:} Valid when generated formulas have length L\_N=(log N)\textasciicircum{}d, \(M_{N}\le N^{q+o(1)}\), \(\beta=0\), and one SAT circuit of size \(L_{N}^{c}\) is reused uniformly across all generated instances.

\paragraph{Equation 2: private\_spoof\_upper\_bound.}
\[
C(H_{A_Q}) \le O(C(Q_pi)) + O(|ForcedRows(Q_pi)| N) + N^{o(1)}
\]
\emph{Valid when:} Valid when \(Q_{pi}\) is row-separable, unforced row parities can be set to zero independently, and forced rows are patched by N-bit equality tests.

\paragraph{Equation 3: forced\_row\_threshold.}
\[
C(H_{A_Q}) > N^{1+epsilon-o(1)} and C(Q_pi)=N^{o(1)} imply |ForcedRows(Q_pi)| \ge N^{epsilon-o(1)}
\]
\emph{Valid when:} Applies after functional certificate soundness transfers the same lower-bound conclusion to every private occurrence-label assignment satisfying the parsed row-normal constraints.

\paragraph{Next analysis.} Instantiate the \(\mathrm{EF}+\mathrm{SAT}_{U}^{occ,row}\) rule ledger: for each extension, weakening, substitution, \(\mathrm{SAT}_{U}\) occurrence axiom, and derived parity line, mark whether it is \(Q_{pi}-local\), row-separable, or a charged cross-row constraint; the first uncharged cross-row rule becomes the countermodel or repair target.

\paragraph{Falsifiers.}
\begin{itemize}
\item Row-separability falsifier: a short accepted occurrence rule imposes a global cross-row parity equation that zeroes many spoof choices while \(ForcedRows\) remains small.
\item Functional-soundness falsifier: the certificate conclusion applies only to real SAT labels, so private occurrence-label spoofing no longer gives a contradiction.
\item Sparse-row residual falsifier: a circuit \(D_{N}\) agrees with \(H_{N}^\mathrm{SAT}\) off \(E<N^{\epsilon-o(1)}\) source rows, giving C(H\_N\textasciicircum{}SAT)<=C(D\_N)+O(EN) below the claimed threshold.
\item Nonuniform-advice falsifier: \(Q_{pi}\) or the verifier hides a selected set of hard rows or the SAT label vector, invalidating the advice-free parser.
\item Positive-beta falsifier: if the carrier secretly uses target queries of length \(N^\beta\) with \(\beta>0\), the substitution window becomes q+beta c and cannot exclude all SAT exponents.
\end{itemize}

\subsection*{cycle-0084: Hardness-Magnification Bridge Realization}

\paragraph{Execution.} research; azure\_responses\_api; attempt primary; 600.01 seconds; 3 lemmas, 4 equations, 8 derivation steps, 4 falsifiers.

\paragraph{Learning.} Cycle 84 turns the \(ForcedRows\) packet into a proof-shaped conditional theorem: row-normal occurrence-local functional certificates must force \(N^{\epsilon-o(1)}\) complete XOR rows, because otherwise private label spoofing plus sparse equality patching is a valid falsifier. The remaining gap is exactly semantic-to-occurrence extraction for full \(\mathrm{EF}+\mathrm{SAT}_{U}\), not ordinary soundness or generic natural-proofs folklore.

\paragraph{Pressure.} The derivation is complete only for row-normal occurrence-local functional certificates. Full \(\mathrm{EF}+\mathrm{SAT}_{U}\) remains blocked at semantic-to-occurrence extraction because \(\mathrm{SAT}_{U}\) axioms can couple labels through real formula syntax and witnesses, so private occurrence-label spoofs may not preserve acceptance.

\paragraph{Derivation steps.}
\begin{enumerate}
\item Fix N, \(\epsilon\), \(q,rhoT_{N}=N^{1+\epsilon-o(1)}\).
\item Parse pi syntactically into \(Q_{pi}\); for each source row x, compute the row view V\_x(Q\_pi) and put x in F=ForcedRows(Q\_pi) iff \(V_{x}\) has no two compatible completions of opposite parity.
\item Assume \(|F|=E<N^{\epsilon-o(1)}\). For every x notin F, choose a \(Q_{pi}-compatible\) completion with row parity 0; for x in F, leave the real SAT occurrence labels unchanged.
\item By row-separability, the rowwise completions combine into one global occurrence-label assignment a' with the same \(Q_{pi}-view\) as the real SAT labels.
\item By transcript invariance, pi accepts a'. By functional label soundness, \(C(H_{N}^{a'})>T_{N}\).
\item But \(H_{N}^{a'}\) is supported only on F, so an OR of equality tests for rows in F, with the correct hardwired output bit for each forced row, computes \(H_{N}^{a'}\) using O(EN) gates.
\item Since \(E<N^{\epsilon-o(1)}\), \(O(EN)=N^{1+\epsilon-o(1)}\) below the claimed threshold after choosing the margin in \(T_{N}\), contradiction.
\item Therefore every valid row-normal occurrence-local functional certificate must force at least \(N^{\epsilon-o(1)}\) complete rows.
\end{enumerate}

\paragraph{Equation 1: beta0\_xor\_carrier.}
\[
H_N^a(x)=\bigoplus_{j=1}^{M_N} a_{x,j},\quad a^{SAT}_{x,j}=SAT(G_N(x,j))
\]
\emph{Valid when:} Valid when \(G_{N}\) is the fixed beta-zero generator and a is any occurrence-label assignment.

\paragraph{Equation 2: beta0\_target\_transfer.}
\[
SAT\in SIZE(n^c)\Rightarrow C(H_N^{SAT})\le N^{\rho+o(1)}+M_N L_N^c=N^{\max{\rho,q}+o(1)}
\]
\emph{Valid when:} Applies after SAT circuits are substituted only into polylogarithmic queries and generator overhead is counted.

\paragraph{Equation 3: sparse\_row\_patch.}
\[
D=H_N \text{ off } R_N\Rightarrow C(H_N)\le C(D)+O(|R_N|N)
\]
\emph{Valid when:} Requires D and H to differ on explicitly listed N-bit source rows only.

\paragraph{Equation 4: forced\_row\_threshold.}
\[
|ForcedRows(Q_\pi)|\ge T_N/O(N)=N^{\epsilon-o(1)}\quad\text{for }T_N=N^{1+\epsilon-o(1)}
\]
\emph{Valid when:} Valid under functional label soundness and transcript invariance for row-normal occurrence certificates.

\paragraph{Next analysis.} Audit one standard full \(\mathrm{EF}+\mathrm{SAT}_{U}\) encoding rule: attempt to extract its occurrence-label dependency into \(Q_{pi}\) preserving \(A_{Q}-invariance\); if a witness or clause-soundness axiom breaks invariance, record the minimal rule-level countermodel and restrict the theorem statement to \(\mathrm{EF}+\mathrm{SAT}_{U}^{occ,rn}\).

\paragraph{Falsifiers.}
\begin{itemize}
\item falsifier: if \(|ForcedRows(Q_{pi})|=N^\eta\) with \(\eta<\epsilon\) and row-normal invariance holds, the spoofed carrier has an O(N\textasciicircum{}\{1+eta\}) patch circuit contradicting the source lower-bound certificate.
\item falsifier: a parser atom that hides real \(\mathrm{SAT}_{U}\) semantics rather than occurrence labels invalidates the occurrence-local theorem for full \(\mathrm{EF}+\mathrm{SAT}_{U}\) and triggers the extraction-boundary quarantine.
\item falsifier: any reintroduced positive query exponent \(\beta\) gives the older transfer cap q+beta*c+o(1), so beta-zero preservation must be rechecked for every proposed carrier variant.
\item falsifier: nonuniform advice selecting the exceptional rows without being charged can fake \(ForcedRows\) robustness; advice must be included in \(Q_{pi}\) or in the circuit-size accounting.
\end{itemize}

\subsection*{cycle-0085: Hardness-Magnification Bridge Realization}

\paragraph{Execution.} research; azure\_responses\_api; attempt primary; 1800.011 seconds; 3 lemmas, 4 equations, 6 derivation steps, 5 falsifiers.

\paragraph{Learning.} Cycle 85 sharpens the formula-semantic extraction fork: full occurrence parsing is false already for \(\mathrm{EF}+SATeq_{N}\) because EqSAT padding identities are real SAT-semantic and nonrelativizing. The valid repair is a semantic-link trichotomy: nonfunctional certificates must pay either occurrence full rows or EqSAT-linked rows at the \(N^{\epsilon-o(1)}\) threshold, otherwise sparse-row residual patching falsifies the source lower-bound certificate.

\paragraph{Pressure.} The derivation is complete only for \(\mathrm{EF}+SATeq_{N}\) where all SAT semantics factors through EqSAT links. It does not yet cover full \(\mathrm{EF}+\mathrm{SAT}_{U}\) with witness, self-reduction, or minimization rules unless those rules are decomposed into label-functional inferences plus row-chargeable semantic links.

\paragraph{Derivation steps.}
\begin{enumerate}
\item Fix \(\mathrm{EF}+SATeq_{N}\): ordinary EF reasoning is syntactic and label-functional; EqSAT is the sole nonrelativizing rule, derived from EF proofs of Boolean formula equivalence.
\item Scan a proof pi of LB\_N(H\_N\textasciicircum{}SAT,N\textasciicircum{}\{1+epsilon\}) and mark every SAT atom occurrence used by label-functional rules; these form \(Q_{pi}\).
\item For every EqSAT application SAT(F)SAT(F'), add a semantic link between the source rows whose generated formulas include F or F', when such rows are in the carrier range; off-range links are charged to proof size but not to parity-row forcing.
\item Run the previous private-row spoof only on rows outside ForcedRows(Q\_pi) union SemRows(E\_pi). On those rows, assign occurrence labels to agree with a low-cost completion while preserving all recorded EqSAT constraints.
\item If \(W=|ForcedRows\) union \(SemRows|<N^{\epsilon-o(1)}\), patch the W exceptional rows by equality tests to obtain C(H\_N\textasciicircum{}SAT)<=S\_D+O(WN), contradicting a claimed \(N^{1+\epsilon-o(1)}\) source lower bound.
\item The EqSAT padding identity shows why the older occurrence-only parser is invalid for full \(\mathrm{EF}+SATeq_{N}\): semantic equality is not invariant under independent occurrence-label flips.
\end{enumerate}

\paragraph{Equation 1: beta\_zero\_xor\_carrier.}
\[
H_N^{SAT}(x)=\bigoplus_{j=1}^{M_N} SAT(G_N(x,j)),\quad M_N=N^{q+o(1)},\quad L_N=(\log N)^d
\]
\emph{Valid when:} Valid when \(G_{N}\) is a fixed batched generator outputting \(length-L_{N}\) SAT instances and the carrier input length is N.

\paragraph{Equation 2: beta\_zero\_substitution\_transfer.}
\[
SAT\in SIZE(n^c)\Rightarrow C(H_N^{SAT})\le N^{\rho+o(1)}+N^{q+o(1)}L_N^c=N^{\max{\rho,q}+o(1)}
\]
\emph{Valid when:} Applies after the generator overhead is \(N^{\rho+o(1)}\), query length is polylogarithmic, and SAT circuits are substituted nonadaptively.

\paragraph{Equation 3: row\_residual\_patch.}
\[
D_N=H_N^{SAT}\text{ off }R_N\Rightarrow C(H_N^{SAT})\le C(D_N)+O(|R_N|N)
\]
\emph{Valid when:} Valid when \(R_{N}\) is a hardwired set of exceptional N-bit source rows and equality tests cost O(N) gates each.

\paragraph{Equation 4: semantic\_link\_threshold.}
\[
|ForcedRows(Q_\pi)\cup SemRows(E_\pi)|< N^{\epsilon-o(1)}\Rightarrow \text{certificate is row-residual-leaky}
\]
\emph{Valid when:} Requires the parser to expose all nonfunctional SAT uses as occurrence queries or EqSAT semantic links, with no hidden advice leakage.

\paragraph{Next analysis.} Attempt the finite \(\mathrm{EF}+SATeq_{N}\) parser proof: formalize proof lines as Boolean circuits over SAT-atoms, prove label-functionality for EF rules, charge every EqSAT endpoint to SemRows, and test whether SAT self-reduction axioms in full \(\mathrm{EF}+\mathrm{SAT}_{U}\) admit the same charge or force a new semantic-escape lemma.

\paragraph{Falsifiers.}
\begin{itemize}
\item falsifier: positive \(\beta\) or q+beta c transfer debt still caps black-box magnification before all SAT exponents are excluded.
\item falsifier: an occurrence-only parser for \(\mathrm{EF}+SATeq_{N}\) is invalid because EqSAT padding equivalence is not preserved by private occurrence-label completions.
\item falsifier: \(W=|ForcedRows\) union \(SemRows|<N^{\epsilon-o(1)}\) plus a residual circuit \(D_{N}\) gives a valid subthreshold equality-patching upper bound.
\item falsifier: hidden nonuniform advice in the proof pi that selects exceptional rows without being charged breaks the target-preserving certificate audit.
\item falsifier: a full \(\mathrm{EF}+\mathrm{SAT}_{U}\) rule using SAT self-reducibility may evade both functional naturalization and EqSAT-link accounting, forcing quarantine until a named semantic-link factorization is proved.
\end{itemize}

\subsection*{cycle-0086: Hardness-Magnification Bridge Realization}

\paragraph{Execution.} research; azure\_responses\_api; attempt primary; 1800.01 seconds; 3 lemmas, 5 equations, 6 derivation steps, 5 falsifiers.

\paragraph{Learning.} Cycle 86 resolves one part of the formula-semantic fork: occurrence-only extraction is false for \(\mathrm{EF}+SATeq_{N}\) because real EqSAT identities are nonrelativizing semantic links. The valid replacement object is a semantic-link residual transcript; subthreshold link coverage is a falsifier for source lower-bound certificates, while the open proof task is a normal form for arbitrary \(\mathrm{EF}+SATeq_{N}\) reasoning.

\paragraph{Pressure.} The normal form is not proved for arbitrary EF reasoning that invokes SAT self-reducibility, extension variables encoding witnesses, or nonlocal formula transformations not decomposed into explicit EqSAT links.

\paragraph{Derivation steps.}
\begin{enumerate}
\item Fix \(\mathrm{Psys}=\mathrm{EF}+SATeq_{N}\): EF plus explicitly checkable equisatisfiable formula identities such as padding, renaming, and clause duplication.
\item Test the previous parser theorem against the paired-padding generator G\_N(x,2i)=Pad(G\_N(x,2i-1)). \(\mathrm{EF}+SATeq_{N}\) proves equality of paired SAT labels and hence XOR cancellation.
\item Observe that a private occurrence-label spoof that flips only one member of a pair is not a model of \(\mathrm{EF}+SATeq_{N}\). Therefore the old occurrence-relativizing falsifier is not valid against this proof system.
\item Repair the observable from \(ForcedRows_{occ}\) to a semantic-link transcript \(T_{N}\) containing occurrence values, constants, and EqSAT equality edges. A row is covered when its parity vector is determined modulo these links.
\item Derive the residual upper bound C(H\_N\textasciicircum{}SAT)<=S\_T+S\_D+O(E\_N N). Thus any accepted certificate with subthreshold transcript and fewer than \(N^{\epsilon-o(1)}\) uncovered rows is a concrete falsifier for a claimed \(N^{1+\epsilon-o(1)}\) source lower bound.
\item Attempt to extend the repair to all \(\mathrm{EF}+SATeq_{N}\) proofs by translating each proof inference into the link transcript. The translation works for padding, renaming, duplicate clauses, and local formula isomorphisms.
\end{enumerate}

\paragraph{Equation 1: beta0\_transfer.}
\[
\mathrm{SIZE}(H_N^{SAT})\le N^{\rho+o(1)}+M_N L_N^c=N^{\max{\rho,q}+o(1)}
\]
\emph{Valid when:} Valid when SAT on \(L_{N}-bit\) formulas has circuits of size \(L_{N}^{c}\) and the generator is substituted without advice leakage or positive target-query exponent.

\paragraph{Equation 2: row\_residual\_patching.}
\[
C(H_N^{SAT})\le S_D+O(E_N N)
\]
\emph{Valid when:} Applies when a circuit \(D_{N}\) of size \(S_{D}\) agrees with \(H_{N}^\mathrm{SAT}\) outside \(E_{N}\) explicitly hardwired source rows and equality tests cost O(N).

\paragraph{Equation 3: semantic\_link\_residual.}
\[
C(H_N^{SAT})\le S_T+S_D+O(E_N N)
\]
\emph{Valid when:} Valid when the parsed \(\mathrm{EF}+SATeq_{N}\) transcript \(T_{N}\) is label-free, has circuit cost \(S_{T}\), and determines all but \(E_{N}\) rows through occurrence values or EqSAT links.

\paragraph{Equation 4: source\_threshold\_consequence.}
\[
C(H_N^{SAT})> N^{1+\epsilon-o(1)}\ \wedge\ S_T+S_D< N^{1+\epsilon-o(1)}\Rightarrow E_N\ge N^{\epsilon-o(1)}
\]
\emph{Valid when:} Requires the row-residual patching bound and a claimed barely-superlinear source lower bound at the same valid generator lengths.

\paragraph{Equation 5: pure\_occurrence\_failure\_pattern.}
\[
\mathrm{ForcedRows}_{occ}(\pi_N)=0,\quad \mathrm{LinkCoveredRows}_{EqSAT}(\pi_N)=N
\]
\emph{Valid when:} Applies to paired-padding carriers where \(\mathrm{EF}+SATeq_{N}\) proves parity cancellation through semantic equality links rather than complete occurrence-row queries.

\paragraph{Next analysis.} Work the \(\mathrm{EF}+SATeq_{N}\) normal-form proof rule by rule: padding, renaming, duplicate clauses, extension-variable abbreviations, and SAT self-reduction; either translate each into a bounded EqSAT-link transcript or isolate the first noncompressible SAT-semantic principle as the next theorem-shaped counterexample.

\paragraph{Falsifiers.}
\begin{itemize}
\item Semantic-link falsifier: a subthreshold EqSAT-link transcript covering all but fewer than \(N^{\epsilon-o(1)}\) rows gives an actual circuit upper bound by sparse-row patching.
\item Paired-padding falsifier: \(\mathrm{EF}+SATeq_{N}\) can prove XOR cancellations with zero occurrence-forced rows, so a pure occurrence parser is not valid for this named proof system.
\item Functional-naturalization falsifier: if accepted certificates define dense constructive label properties over arbitrary truth tables, the route falls back into the naturalization audit rather than a SAT-semantic escape.
\item Positive-beta falsifier remains quarantined: any repair that reintroduces \(L_{N}=N^\beta\) target queries obeys the old q+beta c transfer cap and loses all-exponent target preservation.
\item Advice-leakage falsifier: a transcript that hardwires SAT labels on the generated range is not valid source evidence, because it hides the target truth table in the certificate.
\end{itemize}

\subsection*{cycle-0087: Hardness-Magnification Bridge Realization}

\paragraph{Execution.} research; azure\_responses\_api; attempt primary; 600.005 seconds; 3 lemmas, 4 equations, 7 derivation steps, 5 falsifiers.

\paragraph{Learning.} Cycle 87 replaces the vague semantic-link normal-form obligation by an exact residual-code invariant: after EqSAT quotienting, a proof-mediated certificate that is functional over quotient labels is valid only if its touched-label coset has row-parity distance at least \(N^{\epsilon-o(1)}\). Low quotient-code distance is a concrete falsifier; arbitrary \(\mathrm{EF}+SATeq_{N}\) remains open because it may use genuinely SAT-semantic nonrelativizing principles outside the quotient parser.

\paragraph{Pressure.} The derivation does not produce a normal form for arbitrary \(\mathrm{EF}+SATeq_{N}\) proofs. Full EF can use nonrelativizing SAT-semantic principles, formula self-reductions, and derived equality chains whose acceptance is not a predicate over a bounded W,E quotient transcript.

\paragraph{Derivation steps.}
\begin{enumerate}
\item Fixed the lower-bound sentence LB\_N(S\_N): every circuit of size at most \(S_{N}\) fails to compute \(H_{N}^\mathrm{SAT}\) on some source row.
\item Parsed a hypothetical short \(\mathrm{EF}+SATeq_{N}\) proof into three observable components when possible: an EqSAT quotient graph E, a touched set W of quotient classes, and a predicate P depending on the real SAT labels only through W.
\item For the quotient-functional fragment, assumed \(\mathrm{P}(a_{W})=1\) and \(\mathrm{P}(z_{W})=1\) for every quotient-respecting z matching \(a_{W}\).
\item Defined the residual parity code generated by the row-incidence map \(A_{E}\) from quotient classes to source rows and the coset constraint \(z|_{W}=a_{W}\).
\item If the coset has a representative z with row support delta<S\_N/O(N), built an equality-patch circuit of size O(delta N)+cost(Q)<S\_N for \(H_{z}\).
\item Because P cannot distinguish z from the real labels on W, quotient-functional soundness would assert \(C(H_{z})>S_{N}\), contradiction.
\item Concluded that any valid quotient-functional certificate must prove residual distance at least \(N^{\epsilon-o(1)}\), not merely touch many raw occurrences or cite EqSAT soundness.
\end{enumerate}

\paragraph{Equation 1: semantic\_link\_carrier.}
\[
H_y(x)=\bigoplus_{j\le M_N} y_{[(x,j)]_E}
\]
\emph{Valid when:} Valid when E is the checkable EqSAT quotient graph and y assigns one Boolean label to each quotient class.

\paragraph{Equation 2: residual\_distance.}
\[
\delta_E(W,a_W)=\min_{z:\ z|_W=a_W}\left|\left\{x:\ \bigoplus_{j\le M_N} z_{[(x,j)]_E}=1\right\}\right|
\]
\emph{Valid when:} Valid when the parser reads only W and soundness is tested over all quotient-respecting label assignments.

\paragraph{Equation 3: sparse\_row\_upper\_bound.}
\[
C(H_z)\le O(\delta_E(W,a_W)N)+\operatorname{cost}(Q)
\]
\emph{Valid when:} Valid when row support is patched by equality tests of size O(N) and the transcript cost is subthreshold.

\paragraph{Equation 4: necessary\_distance\_threshold.}
\[
\delta_E(W,a_W)\ge S_N/O(N)=N^{\epsilon-o(1)}
\]
\emph{Valid when:} Valid when \(S_{N}=N^{1+\epsilon-o(1)}\) is the claimed source lower-bound threshold and parser overhead is o(S\_N).

\paragraph{Next analysis.} Instantiate the padded-copy EqSAT quotient family \(E_{k}\), compute the row-incidence matrix \(A_{E}\), and decide whether its constrained cosets have low-weight representatives below \(N^\epsilon\) or certify a valid semantic-link residual lower bound.

\paragraph{Falsifiers.}
\begin{itemize}
\item semantic-link residual distance \(delta_{E}(W,a_{W})<N^{\epsilon-o(1)}\), giving a valid sparse-row patching falsifier
\item positive \(\beta\) or q/rho threshold collapse in the target substitution bridge
\item EqSAT quotient classes reused across rows so heavily that the row-incidence coset has low-weight representatives
\item parser uses nonuniform SAT-label advice hidden in W or in the link witnesses
\item ordinary \(\mathrm{EF}+SATeq_{N}\) proof uses a named SAT-semantic principle outside quotient-functional parsing, blocking this certificate classification
\end{itemize}

\subsection*{cycle-0088: Hardness-Magnification Bridge Realization}

\paragraph{Execution.} research; azure\_responses\_api; attempt primary; 600.005 seconds; 3 lemmas, 5 equations, 7 derivation steps, 5 falsifiers.

\paragraph{Learning.} Cycle 88 identifies the first concrete SAT-semantic escape from the EqSAT quotient parser: self-reduction OR constraints are not invariant under quotient-label spoofing. The valid paper move is not another magnification slogan but a trichotomy repair: equality-only proofs use quotient residual codes, occurrence expansions pay row-forcing cost, and SR proofs require a constrained residual-code audit. Low constrained row distance is a falsifier for source certificates; no main \(\mathrm{P} \ne \mathrm{NP}\) claim is made.

\paragraph{Pressure.} The derivation does not normalize arbitrary \(\mathrm{EF}+\mathrm{SAT}_{U}\) or arbitrary SAT-semantic reasoning to EqSAT+SR constraints. A stronger proof may use search-to-decision, Cook-Levin witness encodings, or formula minimization principles not represented by the finite SR hypergraph \(K_{pi}\).

\paragraph{Derivation steps.}
\begin{enumerate}
\item Fix \(\mathrm{Psys}=\mathrm{EF}+SATeq+SR\) with SAT atoms \(Y_{u}\) and the source sentence LB(H\_N\textasciicircum{}SAT,N\textasciicircum{}\{1+epsilon-o(1)\}).
\item Reuse the stable beta-zero transfer: if SAT has \(L_{N}^{c}-size\) circuits, then \(H_{N}^\mathrm{SAT}\) has size \(N^{max{\rho,q}+o(1)}\), so \(q,\rho<1+\epsilon\) is the valid target-preserving window.
\item For proofs using only SATeq links, contract touched formulas to EqSAT components; the proof verifier is functional over quotient labels, so the cycle-87 quotient residual-code audit applies.
\item Insert one essential SR axiom \(Y_{u}\)  \(Y_{u0}\) OR \(Y_{u1}\). The assignment \(Y_{u}=0,Y_{u0}=1,Y_{u1}=0\) respects every equality quotient but violates SR, so an EqSAT-only parser accepts a completion that is not SAT-semantic.
\item Therefore full EF+SATeq+SR cannot be classified as occurrence/EqSAT-functional unless the parser also records the SR hyperedge or expands it into occurrence children.
\item Once SR hyperedges are recorded, the residual object is no longer a linear quotient coset; it is a finite Boolean constraint code \(K_{pi}\). The correct falsifier is low \(K_{pi}-constrained\) row-parity distance plus the row-patching bound.
\item For compressed block-SR, any completion-sound occurrence parser for negative parent labels must expose all \(2^{t}\) leaves or provide an equivalent semantic OR certificate, giving a tunable threshold 2\textasciicircum{}t=N\textasciicircum{}alpha.
\end{enumerate}

\paragraph{Equation 1: beta0\_xor\_carrier.}
\[
H_N^{SAT}(x)=\oplus_{j=1}^{M_N} SAT(G_N(x,j)), M_N=N^{q+o(1)}, L_N=(log N)^d
\]
\emph{Valid when:} Valid when the source object is the beta-zero batched SAT parity carrier with polylogarithmic target formulas and no positive target-query exponent.

\paragraph{Equation 2: target\_substitution\_upper\_bound.}
\[
C(H_N^{SAT}) \le N^{rho+o(1)}+M_N L_N^c = N^{max{rho,q}+o(1)}
\]
\emph{Valid when:} Valid when SAT on \(L_{N}-bit\) formulas has nonuniform circuits of size \(L_{N}^{c}\) and the generator is advice-free with size \(N^{\rho+o(1)}\).

\paragraph{Equation 3: \(SR_{semantic,constraint}\).}
\[
Y_u = Y_{u|z=0} \vee Y_{u|z=1}
\]
\emph{Valid when:} Valid when u is a Boolean formula containing variable z and \(Y_{v}\) denotes the satisfiability truth value of formula v under the real SAT interpretation.

\paragraph{Equation 4: \(\mathrm{EqSAT}_{spoof,obstruction}\).}
\[
(Y_u,Y_{u|z=0},Y_{u|z=1})=(0,1,0) satisfies all equality links but violates SR(u,z)
\]
\emph{Valid when:} Valid when the parser records only EqSAT quotient information and no additional OR-type self-reduction constraint for the touched formulas.

\paragraph{Equation 5: constraint\_row\_distance.}
\[
\Delta_K(pi)=min_{lambda models K_pi} |{x : \oplus_j lambda(G_N(x,j)) \ne H_N^{SAT}(x)}|
\]
\emph{Valid when:} Valid when \(K_{pi}\) is the finite system of EqSAT and SR constraints touched by the proof and \(\lambda\) ranges over Boolean assignments to its SAT atoms.

\paragraph{Next analysis.} Instantiate \(K_{pi}\) for local one-step SR only: bound \(|K_{pi}|\) by proof size, define its row-parity constraint-distance \(Delta_{K}\), and test whether \(Delta_{K}\) can be reduced to a known SAT promise-hardness statement or whether it is target-equivalent to the desired \(H_{N}\) lower bound.

\paragraph{Falsifiers.}
\begin{itemize}
\item falsifier: an EqSAT-only quotient parser for EF+SATeq+SR is invalid because the spoof labels (0,1,0) satisfy all equality links while violating a used self-reduction axiom.
\item falsifier: any constrained predictor \(D_{N}\) with \(S_{D}+E_{N}\) \(NN^{1+\epsilon-o(1)}\).
\item falsifier: compressed block-SR with 2\textasciicircum{}t<N\textasciicircum{}\{epsilon-o(1)\} can leave omitted-leaf or omitted-row errors patchable below the source threshold.
\item falsifier: reintroducing positive target-query exponent \(\beta\) gives the older transfer cap q+beta c and cannot imply all-exponent SAT lower bounds.
\item falsifier: nonuniform advice listing the residual bad rows or semantic children must be charged; otherwise the certificate leaks the target lower bound into the source object.
\end{itemize}

\subsection*{cycle-0089: Hardness-Magnification Bridge Realization}

\paragraph{Execution.} research; azure\_responses\_api; attempt compact\_retry; 1094.129 seconds; 1 lemmas, 3 equations, 4 derivation steps, 4 falsifiers.

\paragraph{Learning.} The cycle upgrades the SR semantic escape into a quantitative constrained-distance invariant. Occurrence spoofing is invalid after SR constraints, but any bounded-SR parser still faces the same row-patching threshold. Progress is route-level certification, not a P vs NP proof.

\paragraph{Pressure.} The missing proof is semantic-to-residual parsing for full \(\mathrm{EF}+SATeq+SR_{h}\); SR OR clauses couple labels across restrictions, so independent occurrence-label spoofing is not valid.

\paragraph{Derivation steps.}
\begin{enumerate}
\item Quotient formulas by EqSAT links, then add \(SR_{h}\) constraints for restrictions in the generated polylog formula closure; real SAT labels satisfy all constraints.
\item Assume a parser maps any short accepted proof pi to a functional transcript and decoder \(D_{pi}\) over constrained quotient labels, with cost \(S_{pi}\) and residual row set \(R_{pi}\).
\item Apply row patching: hardwire \(R_{pi}\) and correct \(D_{pi}\) by equality gates, giving C(H\_N\textasciicircum{}SAT)<=S\_pi+O(|R\_pi|N).
\item Therefore a proof certificate at threshold \(N^{1+\epsilon-o(1)}\) must force \(|R_{pi}|\ge N^{\epsilon-o(1)}\) unless \(S_{pi}\) already exceeds the source threshold.
\end{enumerate}

\paragraph{Equation 1: beta\_zero\_transfer.}
\[
C(H_N^{SAT}) \le N^{\rho+o(1)}+M_N L_N^c=N^{\max{\rho,q}+o(1)}
\]
\emph{Valid when:} Applies when SAT has size \(L_{N}^{c}\) circuits and the carrier uses only polylogarithmic beta-zero target queries.

\paragraph{Equation 2: row\_patch\_bound.}
\[
C(H_N^{SAT}) \le C(D_N)+O(|R_N|N)
\]
\emph{Valid when:} Valid when \(D_{N}\) agrees with \(H_{N}^{\mathrm{SAT}}\) on every source row outside the explicitly listed residual set \(R_{N}\).

\paragraph{Equation 3: sr\_certificate\_threshold.}
\[
\mathrm{ValidCert}_{SR_h}(\pi) \Rightarrow S_\pi+\Delta_{SR_h}(\pi)N \ge N^{1+\epsilon-o(1)}
\]
\emph{Valid when:} Requires the scoped parser to turn pi into an advice-free SR-constrained residual decoder for the same carrier.

\paragraph{Next analysis.} Fix h with 2\textasciicircum{}h=N\textasciicircum{}\{eta\}; formalize \(Delta_{SR,h}\) for the generated formula closure and try to prove the parser for \(\mathrm{EF}+SATeq+SR_{h}\) restricted to extension variables naming only quotient labels and SR clauses.

\paragraph{Falsifiers.}
\begin{itemize}
\item Low constrained row distance \(Delta_{SR,h}(pi)<N^{\epsilon-o(1)}\) with \(S_{pi}<N^{1+\epsilon-o(1)}\) is a direct row-patching falsifier.
\item SR closure overhead M\_N2\textasciicircum{}h=N\textasciicircum{}\{q\_SR+o(1)\} with \(q_{SR}\ge 1+\epsilon\) collapses the beta-zero target-preserving window.
\item A parser using nonuniform advice to list residual rows or semantic links invalidates the claimed advice-free certificate transfer.
\item A genuinely SAT-semantic EF inference outside EqSAT+bounded-SR blocks the constrained parser and quarantines this scoped no-go.
\end{itemize}

\subsection*{cycle-0090: Hardness-Magnification Bridge Realization}

\paragraph{Execution.} research; azure\_responses\_api; attempt primary; 600.001 seconds; 3 lemmas, 3 equations, 9 derivation steps, 6 falsifiers.

\paragraph{Learning.} Cycle 90 converts the SR semantic escape into a costed leaf-solver falsifier: bounded self-reduction depth is not free, but if \(q+\eta+\tau<1+\epsilon\) and the induced row error is below \(N^\epsilon\) then the source lower bound collapses by equality patching. The live inference-quality packet is now theorem-shaped for \(\mathrm{Psys}_{h}\); the only remaining full-system gap is the advice-free parser normal form for arbitrary EF extension variables.

\paragraph{Pressure.} ParserNormalForm(EF+EqSAT+SR\_h): prove that every short proof of LB(H\_N\textasciicircum{}\{SAT\}) either gives a dense functional label property, an occurrence/full-row transcript, or a bounded-SR residual code with advice-free size bound poly(|pi|,L\_N).

\paragraph{Derivation steps.}
\begin{enumerate}
\item Fix \(\mathrm{Psys}_{h}=\mathrm{EF}+\mathrm{EqSAT}+SR_{h}\) and the beta-zero carrier \(H_{N}^{\mathrm{SAT}}\).
\item Recheck the target bridge: substituting any SIZE(L\textasciicircum{}c) SAT circuit into all generated queries gives C(H\_N\textasciicircum{}\{SAT\})<=N\textasciicircum{}\{max\{rho,q\}+o(1)\}, so \(q,\rho<1+\epsilon\) is the target-preserving exponent condition.
\item Expose the SR semantic escape by allowing h rounds of SAT self-reduction before labels are evaluated.
\item Given a leaf solver \(A_{N}\) of size \(N^{\tau+o(1)}\), compute each top query by enumerating 2\textasciicircum{}\{bh\}=N\textasciicircum{}\{eta\_N+o(1)\} leaves, OR their \(A_{N}\) labels, and XOR over \(M_{N}=N^{q+o(1)}\) queries.
\item This yields C(H\_\{A,h\})<=N\textasciicircum{}\{rho+o(1)\}+N\textasciicircum{}\{q+eta\_N+tau+o(1)\}.
\item If \(H_{A,h}\) differs from \(H_{N}^{\mathrm{SAT}}\) on \(E_{N}\) rows, add equality patches for those rows and get C(H\_N\textasciicircum{}\{SAT\})<=N\textasciicircum{}\{max\{rho,q+eta\_N+tau,1+e\}+o(1)\}.
\item Therefore any claimed \(N^{1+\epsilon-o(1)}\) source lower bound must rule out all bounded-SR easy carriers with \(q+eta_{N}+\tau<1+\epsilon\) and \(e<\epsilon\).
\item Attempted extension to full EF+SAT semantics: EqSAT quotienting is parsed by identifying label variables, and explicit SR rules add OR hyperedges.
\item The extension fails for arbitrary EF extension variables because they may encode global SAT-semantic predicates not representable as bounded SR leaf evaluations or occurrence-row queries.
\end{enumerate}

\paragraph{Equation 1: beta\_zero\_target\_substitution.}
\[
C(H_N^{SAT}) \le N^{\rho+o(1)}+M_N\,L_N^c=N^{\max{\rho,q}+o(1)}
\]
\emph{Valid when:} Valid when SAT has size \(L^{c}\) circuits on all generated \(length-L_{N}\) formulas and L\_N=(log N)\textasciicircum{}d makes every fixed c contribute only \(N^{o(1)}\) overhead.

\paragraph{Equation 2: \(bounded_{SR,leaf,solver,upper,bound}\).}
\[
C(H_{A,h}) \le N^{\rho+o(1)}+N^{q+\eta_N+\tau+o(1)}
\]
\emph{Valid when:} Applies when the parsed self-reduction tree has 2\textasciicircum{}\{bh\}=N\textasciicircum{}\{eta\_N+o(1)\} leaves per generated query and \(A_{N}\) has size \(N^{\tau+o(1)}\) with no hidden advice.

\paragraph{Equation 3: row\_residual\_patch\_bound.}
\[
C(H_N^{SAT}) \le C(H_{A,h})+O(E_N N) \le N^{\max{\rho,q+\eta_N+\tau,1+e}+o(1)}
\]
\emph{Valid when:} Requires that \(H_{A,h}\) and \(H_{N}^{\mathrm{SAT}}\) disagree on exactly \(E_{N}=N^{e+o(1)}\) source rows and that equality tests for source rows cost O(N) gates each.

\paragraph{Next analysis.} Build the minimal \(\mathrm{EF}+\mathrm{EqSAT}+SR_{1}\) parser countermodel: allow one extension variable Z(F)=SAT(F\_0) OR SAT(F\_1) and test whether it is representable by the SR DAG without advice; if yes, increase to two nested extensions, if no, use it as the semantic escape separating full EF from \(\mathrm{Psys}_{h}\).

\paragraph{Falsifiers.}
\begin{itemize}
\item bounded-SR leaf solver falsifier: \(q+\eta+\tau<1+\epsilon\) and row error exponent \(e<\epsilon\) gives a subthreshold patched circuit for \(H_{N}^{\mathrm{SAT}}\).
\item sparse-row residual falsifier: any \(D_{N}\) agreeing with \(H_{N}^{\mathrm{SAT}}\) off fewer than \(N^{\epsilon-o(1)}\) rows and having subthreshold size invalidates the source lower-bound certificate.
\item functional-density falsifier: if accepted certificates are dense and truth-table constructive over quotient labels, they become natural-property evidence rather than a non-natural lower-bound proof.
\item parser-advice falsifier: a proof parser depending nonuniformly on the SAT label vector can fake extraction and is not valid for the route.
\item full-EF semantic-escape falsifier: extension variables for global SAT-semantic predicates may defeat the bounded-SR parser, forcing demotion to the restricted \(\mathrm{Psys}_{h}\) theorem.
\item positive-beta transfer falsifier remains quarantined: any interface with target query length \(N^\beta\), \(\beta>0\), has exponent q+beta c and cannot give all-exponent SAT lower bounds.
\end{itemize}

\subsection*{cycle-0091: Hardness-Magnification Bridge Realization: beta-zero XOR formula-semantic parser audit}

\paragraph{Execution.} research; azure\_responses\_api; attempt primary; 600.002 seconds; 3 lemmas, 4 equations, 8 derivation steps, 5 falsifiers.

\paragraph{Learning.} Cycle 91 upgrades the bounded-SR leaf-solver audit by adding an explicit EF extension-support charge and proving a no-go for extension-blind parsers. The valid theorem is restricted but contribution-shaped: subthreshold charged transcripts are upper-bound falsifiers, while arbitrary EF+SAT remains blocked at a precise advice-free normal-form lemma.

\paragraph{Pressure.} Prove or refute the advice-free extension-support normal form for arbitrary EF extension variables in EF+EqSAT+bounded-SR; the known derivation only handles proofs whose extension dependencies are syntactically charged.

\paragraph{Derivation steps.}
\begin{enumerate}
\item Fix \(H_{N}^\mathrm{SAT}\) with \(M_{N}=N^{q+o(1)}\), L\_N=(log N)\textasciicircum{}d, \(d>1\), and generator size \(N^{\rho+o(1)}\); this keeps the target substitution \(\beta\) at zero.
\item Assume SAT has \(L^{c}-size\) circuits. Substitute the SAT circuit into every generated formula and obtain C(H\_N\textasciicircum{}SAT)<=N\textasciicircum{}\{max\{rho,q\}+o(1)\}, so \(q,\rho<1+\epsilon\) is the valid target-preserving window.
\item Choose \(\mathrm{Psys}_{h,Ext}=\mathrm{EF}+\mathrm{EqSAT}+depth-h\) self-reduction. Parse pi by collecting visible SAT atoms, EqSAT quotient links, SR leaves, and all SAT-atom dependencies of extension definitions.
\item If the parsed transcript supplies a leaf solver of exponent \(\tau\) after leaf blowup \(\eta\) and is wrong on \(E_{N}\) rows, build \(D_{N}\) by running \(G_{N}\), evaluating the charged transcript, and xoring the row outputs.
\item Patch the residual rows \(R_{pi}\) using equality tests for the N-bit source row; the patch cost is O(E\_N N).
\item Derive C(H\_N\textasciicircum{}SAT)<=N\textasciicircum{}\{rho+o(1)\}+S\_ext+N\textasciicircum{}\{q+eta+tau+o(1)\}+O(E\_NN). If all exponents are below 1+epsilon, the claimed source lower bound is falsified.
\item Test parser robustness by defining an EF extension \(y_{r}=xor_{j}\) \(z_{r,j}\). A parser that ignores extension definitions sees one atom and no complete row, but expansion touches the full parity row; therefore extension-blind occurrence extraction is invalid.
\item Conclusion: the restricted charged-extension theorem is valid, while the full EF+SAT claim is blocked exactly at the advice-free extension-support normal-form lemma.
\end{enumerate}

\paragraph{Equation 1: carrier.}
\[
H_N^{SAT}(x)=\bigoplus_{j\le M_N} SAT(G_N(x,j)),\quad M_N\le N^{q+o(1)},\quad L_N=(\log N)^d.
\]
\emph{Valid when:} Valid when \(G_{N}\) is an explicit batched generator whose outputs are SAT instances of length \(L_{N}\) and whose source input has length N.

\paragraph{Equation 2: beta\_zero\_transfer.}
\[
SAT\in SIZE(n^c)\Rightarrow C(H_N^{SAT})\le N^{\rho+o(1)}+M_NL_N^c=N^{\max{\rho,q}+o(1)}.
\]
\emph{Valid when:} Applies when SAT on \(L_{N}-bit\) formulas has nonuniform circuits of size \(L_{N}^{c}\) and \(L_{N}\) is polylogarithmic, so the target-query exponent \(\beta\) is zero.

\paragraph{Equation 3: \(charged_{SR,residual,upper,bound}\).}
\[
C(H_N^{SAT})\le N^{\rho+o(1)}+S_{ext}(\pi)+N^{q+\eta+\tau+o(1)}+O(E_NN).
\]
\emph{Valid when:} Valid when the charged parser extracts a bounded self-reduction leaf solver of exponent \(\tau\), leaf blowup \(\eta\), extension support \(S_{ext}\), and \(E_{N}\) bad source rows.

\paragraph{Equation 4: subthreshold\_collapse.}
\[
S_{ext}\le N^{\gamma+o(1)},\ E_N\le N^{\delta+o(1)},\ \max{\rho,\gamma,q+\eta+\tau,1+\delta}< 1+\epsilon\Rightarrow C(H_N^{SAT})< N^{1+\epsilon-o(1)}.
\]
\emph{Valid when:} Requires all transcript costs to be counted in the same nonuniform circuit model and equality patches to cover exactly the residual rows.

\paragraph{Next analysis.} Instantiate \(\mathrm{Psys}_{h,Ext}^{schema}\): EF+EqSAT+depth-h SR with extension definitions generated by row-oblivious schemas. Try to prove ExtSupp(NF(pi))<=poly(|pi|), or construct a minimal nonrelativizing SAT-semantic extension that defeats this normal form.

\paragraph{Falsifiers.}
\begin{itemize}
\item Subthreshold charged transcript falsifier: S\_ext+N\textasciicircum{}\{q+eta+tau\}+E\_NN<N\textasciicircum{}\{1+epsilon-o(1)\} gives an explicit circuit for \(H_{N}^\mathrm{SAT}\).
\item Extension-blind parser falsifier: \(y_{r}=xor_{j}\) \(z_{r,j}\) hides a full row behind one EF atom unless definitions are expanded or charged.
\item Nonuniform leakage falsifier: a row-indexed advice table inside extension definitions invalidates any claimed advice-free parser.
\item Bounded-SR leaf-solver falsifier: if \(q+\eta+\tau<1+\epsilon\) and residual errors are below \(N^{\epsilon-o(1)}\), equality patching collapses the source lower bound.
\item Positive-beta falsifier remains quarantined: any return to polynomial-length SAT queries reintroduces the q+beta c transfer cap.
\end{itemize}

\subsection*{cycle-0092: Hardness-Magnification Bridge Realization}

\paragraph{Execution.} research; azure\_responses\_api; attempt primary; 600.009 seconds; 3 lemmas, 5 equations, 7 derivation steps, 6 falsifiers.

\paragraph{Learning.} Cycle 92 turns the formula-semantic fork into a sharper proof packet. For support-declared EF+SATeq+bounded-SR, subthreshold parsed support plus bounded leaf solving and sparse residual rows gives a valid small-circuit falsifier for \(\mathrm{LB}_{N}\). The full EF+SAT claim is not proved; the precise remaining obstruction is \(NSP_{N}\), the advice-free extraction of arbitrary extension-variable reasoning into charged row supports and SR leaves. This is route-level certification, not progress on the main \(\mathrm{P} \ne \mathrm{NP}\) conjecture itself.

\paragraph{Pressure.} \(NSP_{N}\): every short arbitrary EF+SATeq+SR proof of \(\mathrm{LB}_{N}\) must admit a support-complete, advice-free transcript with row distance and leaf-solver costs small enough to preserve the \(N^{1+\epsilon}\) threshold. This is unproved and may fail for extension-compressed SAT-semantic reasoning.

\paragraph{Derivation steps.}
\begin{enumerate}
\item Fixed the carrier \(H_{N}^\mathrm{SAT}\) with \(\beta=0\), \(M_{NN}^{1+\epsilon-o(1)}\).
\item Named the concrete proof model \(\mathrm{Psys}_{h}^{supp}\): EF+SATeq+bounded-SR with support-declared extension variables, no hidden advice, and charged proof/transcript/leaf-solver costs.
\item Parsed any accepted proof pi into \(T_{pi}\) containing declared row supports, EqSAT quotient links, bounded-SR leaf families, a charged leaf solver, and a residual row set \(R_{pi}\).
\item Under support adequacy, built \(D_{T}\) by computing \(G_{N}\), following \(T_{pi}\), evaluating charged SR leaves with the leaf solver, and xoring the resulting labels; this gives C(D\_T)<=S\_T+N\textasciicircum{}\{rho+o(1)\}+N\textasciicircum{}\{q+eta+tau+o(1)\}.
\item Patched the residual set \(R_{pi}\) by equality tests on N-bit source rows, yielding C(H\_N\textasciicircum{}SAT)<=C(D\_T)+O(|R\_pi|N).
\item Concluded that any accepted proof with total charge below \(N^{1+\epsilon-o(1)}\) is a falsifier for the claimed source lower-bound certificate, not evidence for \(\mathrm{P} \ne \mathrm{NP}\).
\item Separated the restricted theorem from the full EF+SAT ambition: arbitrary extension variables require \(NSP_{N}\), an advice-free normal-support parser lemma not supplied by ordinary soundness.
\end{enumerate}

\paragraph{Equation 1: carrier.}
\[
H_N^{SAT}(x)=\bigoplus_{j\le M_N} SAT(G_N(x,j)),\quad M_N\le N^{q+o(1)},\quad |G_N(x,j)|=L_N=(\log N)^d
\]
\emph{Valid when:} Valid when the carrier generator and SAT-label universe are fixed before any proof-system parsing or target-circuit substitution.

\paragraph{Equation 2: beta\_zero\_transfer.}
\[
SAT\in SIZE(n^c)\Rightarrow C(H_N^{SAT})\le N^{\rho+o(1)}+M_N\,L_N^c=N^{\max{\rho,q}+o(1)}
\]
\emph{Valid when:} Applies after the SAT circuit is nonuniform but fixed-exponent and the target query length is polylogarithmic beta-zero.

\paragraph{Equation 3: transcript\_approximator.}
\[
C(D_{T_\pi})\le S_T+N^{\rho+o(1)}+B_N\,N^{\tau+o(1)}\le S_T+N^{\rho+o(1)}+N^{q+\eta+\tau+o(1)}
\]
\emph{Valid when:} Requires support adequacy for \(\mathrm{Psys}_{h}^{supp}\) and a charged bounded-SR leaf solver on the generated leaf range.

\paragraph{Equation 4: row\_patching.}
\[
C(H_N^{SAT})\le C(D_{T_\pi})+O(|R_\pi|N)
\]
\emph{Valid when:} Requires the approximator be defined on all source inputs and disagree only on the explicitly listed residual rows.

\paragraph{Equation 5: certificate\_falsifier\_threshold.}
\[
S_T+N^{\rho+o(1)}+N^{q+\eta+\tau+o(1)}+|R_\pi|N< N^{1+\epsilon-o(1)}\Rightarrow LB_N(\epsilon)\text{ is not soundly certified by }\pi
\]
\emph{Valid when:} Valid when all proof, parser, leaf, generator, and residual-patching costs are charged in the same nonuniform circuit model.

\paragraph{Next analysis.} Try to prove \(NSP_{N}\) for the fragment EF+SATeq+SR whose extension definitions are row-affine and support-monotone; if the proof fails, extract the minimal extension rule that creates uncharged SAT-semantic support and state the final theorem scope around \(\mathrm{Psys}_{h}^{supp}\).

\paragraph{Falsifiers.}
\begin{itemize}
\item Subthreshold support transcript falsifier: S\_T+N\textasciicircum{}\{q+eta+tau\}+|R\_pi|N<N\textasciicircum{}\{1+epsilon-o(1)\} gives an explicit small circuit for \(H_{N}^\mathrm{SAT}\).
\item Bounded-SR leaf-solver falsifier: if \(q+\eta+\tau<1+\epsilon\) and residual rows are fewer than \(N^{\epsilon-o(1)}\), the proof-mediated source lower bound collapses.
\item Low constrained row-distance falsifier after EqSAT/SR quotienting: row-parity distance below \(N^{\epsilon-o(1)}\) is patchable.
\item Hidden-advice falsifier: any parser that uses uncharged SAT truth-table information invalidates the claimed advice-free extraction.
\item Functional-density falsifier: dense constructive label-functional certificates naturalize rather than escape barriers.
\item Positive-beta falsifier remains quarantined: any return to target-query length \(N^\beta\) with \(\beta>0\) caps the transfer by q+beta c.
\end{itemize}

\subsection*{cycle-0093: Hardness-Magnification Bridge Realization}

\paragraph{Execution.} research; azure\_responses\_api; attempt primary; 600.005 seconds; 3 lemmas, 4 equations, 8 derivation steps, 5 falsifiers.

\paragraph{Learning.} Cycle 93 upgrades the formula-semantic extraction fork by naming \(NSP_{N}\) as a support-transparency lemma for unrestricted EF extension variables and proving the conditional transcript falsifier around it. The valid restricted contribution is support-declared \(\mathrm{EF}+SATeq+bounded-SR_{h}\): any subthreshold parsed transcript collapses \(\mathrm{LB}_{N}\) by a small circuit plus sparse-row patching. The full EF+SAT claim remains blocked exactly at advice-free extraction of extension-variable support, not at generic soundness or barrier folklore.

\paragraph{Pressure.} \(NSP_{N}\): an advice-free normal-support parser for unrestricted \(\mathrm{EF}+SATeq+SR_{h}\) extension variables. The missing lemma must extract charged row supports and bounded-SR leaves from arbitrary extension definitions without SAT-label advice and without solving a target-equivalent semantic support problem.

\paragraph{Derivation steps.}
\begin{enumerate}
\item Fixed proof system \(\mathrm{Psys}_{h}^{un}\): EF plus extension axioms, \(\mathrm{EqSAT}_{N}\) formula-equality axioms, and bounded self-reduction OR axioms to depth h for generated SAT instances of length L\_N=(log N)\textasciicircum{}d.
\item Encoded the source sentence LB\_N(s\_N) as the propositional assertion that every \(size-s_{N}\) circuit differs from \(H_{N}^\mathrm{SAT}\) on some source row; \(s_{N}=N^{1+\epsilon-o(1)}\).
\item Separated accepted proofs into a functional branch and a nonfunctional branch. The functional branch induces a constructive label property and is routed to the beta-zero naturalization audit.
\item For the nonfunctional branch, assumed \(NSP_{N}\) produces a support transcript \(Q_{pi}\), \(K_{N}\) bounded-SR leaves, leaf solvers of size \(T_{leaf}\), and residual row set \(R_{pi}\), all advice-free from pi and \(G_{N}\).
\item Built a circuit \(D_{N}\) that evaluates \(G_{N}\), follows the parsed support transcript, invokes leaf solvers on charged SR leaves, and agrees with \(H_{N}^\mathrm{SAT}\) outside \(R_{pi}\).
\item Applied equality patching to \(R_{pi}\), yielding C(H\_N\textasciicircum{}SAT)<=N\textasciicircum{}\{rho+o(1)\}+S\_supp+K\_NT\_leaf+O(|R\_pi|N).
\item Concluded that any transcript below the exponent window \(max{\rho,sigma,q+\eta+\tau,1+\delta}<1+\epsilon\) contradicts the source lower-bound claim, so a valid proof-mediated certificate must force a superthreshold transcript or escape \(NSP_{N}\).
\item Tried to prove \(NSP_{N}\) for arbitrary EF extension variables by inducting on extension definitions. The induction fails when an extension variable aggregates row-dependent SAT-semantic information through phi\_i without a declared support summary.
\end{enumerate}

\paragraph{Equation 1: beta\_zero\_substitution.}
\[
C(H_N^{SAT}) \le N^{\rho+o(1)}+N^{q+o(1)}L_N^c=N^{\max{\rho,q}+o(1)}
\]
\emph{Valid when:} Valid when SAT has nonuniform circuits of size \(n^{c}\), the generator has size \(N^{\rho+o(1)}\), \(M_{N}\le N^{q+o(1)}\), and all queries have length L\_N=(log N)\textasciicircum{}d.

\paragraph{Equation 2: row\_residual\_patching.}
\[
C(H_N^{SAT}) \le S_D+O(|R_N|\,N)
\]
\emph{Valid when:} Valid when a circuit \(D_{N}\) agrees with \(H_{N}^{\mathrm{SAT}}\) outside an explicit residual row set \(R_{N}\) and equality tests on N-bit source rows are allowed nonuniformly.

\paragraph{Equation 3: parsed\_transcript\_upper\_bound.}
\[
C(H_N^{SAT}) \le N^{\rho+o(1)}+S_{supp}+K_NT_{leaf}+O(E_NN)
\]
\emph{Valid when:} Valid when \(NSP_{N}\) outputs charged support of size \(S_{supp}\), \(K_{N}\) bounded-SR leaves each solved with circuits of size \(T_{leaf}\), and only \(E_{N}\) rows remain unparsed.

\paragraph{Equation 4: subthreshold\_falsifier\_window.}
\[
\max{\rho,\sigma,q+\eta+\tau,1+\delta}< 1+\epsilon \Rightarrow C(H_N^{SAT})< N^{1+\epsilon-o(1)}
\]
\emph{Valid when:} Valid when \(S_{supp}=N^{sigma+o(1)}\), \(K_{N}=N^{q+\eta+o(1)}\), \(T_{leaf}=N^{\tau+o(1)}\), \(E_{N}=N^{\delta+o(1)}\), and polylogarithmic overheads are absorbed into o(1).

\paragraph{Next analysis.} Attempt an explicit \(NSP_{N}\) countermodel for unrestricted EF extension variables: construct a small row-mask extension family \(e_{R}\), prove that occurrence-only parsing misses its support, and decide whether support transparency can be repaired by charging C(Supp(e\_R)) without making the parser target-equivalent.

\paragraph{Falsifiers.}
\begin{itemize}
\item Subthreshold support-declared transcript with \(max{\rho,sigma,q+\eta+\tau,1+\delta}<1+\epsilon\) is a direct falsifier for the claimed \(N^{1+\epsilon}\) source lower bound.
\item A dense \(2^{O(L_{N})}-constructive\) functional label property is a naturalization falsifier for proof-mediated beta-zero certificates unless nonlargeness or nonconstructivity is proved.
\item An extension variable whose affected rows are not advice-free support-transparent is a falsifier for occurrence-only parser claims about unrestricted \(\mathrm{EF}+SATeq+SR_{h}\).
\item Positive target-query exponent \(\beta>0\) remains a separate transfer falsifier: q+beta c caps all-exponent SAT lower-bound consequences.
\item Low constrained row-parity distance after EqSAT/SR quotienting is a falsifier because sparse residual rows can be equality patched below threshold.
\end{itemize}

\subsection*{cycle-0094: Hardness-Magnification Bridge Realization}

\paragraph{Execution.} research; azure\_responses\_api; attempt primary; 600.004 seconds; 3 lemmas, 5 equations, 7 derivation steps, 5 falsifiers.

\paragraph{Learning.} Cycle 94 closes a concrete support-transparency subcase: bit-explicit EF extension variables can be parsed without support declarations, giving a valid transcript falsifier theorem. The full formula-semantic extraction fork is now sharper: the obstruction is not arbitrary EF extension variables but succinct FormulaEval-SAT comprehension, which defeats size-only occurrence parsing unless its row-index expansion is charged.

\paragraph{Pressure.} \(NSP_{sym}\): an advice-free parser for succinct formula-indexed SAT atoms must either expand \(N^{1+q+o(1)}\) concrete row-index labels or prove a new semantic-support theorem that remains completion-sound without occurrence expansion.

\paragraph{Derivation steps.}
\begin{enumerate}
\item Fix \(\mathrm{Psys}_{bEF}\): EF extension variables are introduced only by explicit Boolean definitions over earlier variables and concrete SAT atoms \(A_{x,j}\); SATeq and bounded-SR links are listed as concrete proof edges.
\item Define Dep\textasciicircum{}*(pi) by reverse graph traversal from the lower-bound certificate conclusion through extension definitions, EqSAT links, SR leaves, and residual row declarations. This is advice-free and polynomial in the proof DAG size.
\item From Dep\textasciicircum{}*(pi), build a circuit \(D_{pi}\) that evaluates the generator, all reached label predicates, bounded-SR leaf solvers, and the extension DAG. On rows outside \(R_{N}\) it agrees with \(H_{N}^\mathrm{SAT}\) by the transcript soundness premise.
\item Patch the \(E_{N}\) residual rows with equality tests, obtaining C(H\_N\textasciicircum{}SAT)<=N\textasciicircum{}\{rho+o(1)\}+Gamma\_bEF(pi). Therefore any subthreshold transcript is a concrete falsifier for LB\_N(1+epsilon).
\item Apply the independent beta-zero substitution audit: if SAT has \(L_{N}^{c}-size\) circuits for every fixed c, then \(H_{N}^\mathrm{SAT}\) has \(N^{max{\rho,q}+o(1)}\) circuits, so the target-preserving window remains \(q,\rho<1+\epsilon\).
\item Attempt to extend the parser to \(\mathrm{Psys}_{sym}\) with FormulaEval-SAT comprehension. A single symbolic line can range over N \(M_{N}\) generated labels; surface occurrence support is no longer completion-sound under independent occurrence-label spoofing.
\item Conclude that \(NSP_{N}\) is proved for bit-explicit \(\mathrm{EF}+SATeq+bounded-SR_{h}\), but full symbolic EF+SAT is blocked at a named semantic principle: succinct FormulaEval-SAT comprehension must be expanded, charged, or handled by a new nonrelativizing invariant.
\end{enumerate}

\paragraph{Equation 1: carrier\_definition.}
\[
H_N^{SAT}(x)=\bigoplus_{j\le M_N} SAT(G_N(x,j)),\quad M_N=N^{q+o(1)},\quad L_N=(\log N)^d
\]
\emph{Valid when:} Valid when the carrier is evaluated on exactly N source rows and every generated SAT instance has length L\_N=(log N)\textasciicircum{}d.

\paragraph{Equation 2: beta\_zero\_substitution.}
\[
SAT\in SIZE(n^c)\Rightarrow C(H_N^{SAT})\le N^{\rho+o(1)}+M_N L_N^c=N^{\max{\rho,q}+o(1)}
\]
\emph{Valid when:} Applies after the target SAT circuit is substituted only into polylogarithmic generated queries and the generator overhead is \(N^{\rho+o(1)}\).

\paragraph{Equation 3: parsed\_charge.}
\[
\Gamma_{bEF}(\pi)=S_{dep}(\pi)+S_{link}(\pi)+N^{q+\eta+\tau+o(1)}+E_NN
\]
\emph{Valid when:} Valid when the proof is bit-explicit, extension definitions are concrete DAG clauses, and bounded self-reduction leaves are declared by the parser.

\paragraph{Equation 4: transcript\_upper\_bound.}
\[
C(H_N^{SAT})\le N^{\rho+o(1)}+\Gamma_{bEF}(\pi)
\]
\emph{Valid when:} Requires the transcript to compute \(H_{N}^\mathrm{SAT}\) outside \(E_{N}\) rows and uses equality patches for the residual source rows.

\paragraph{Equation 5: subthreshold\_falsifier.}
\[
N^{\rho+o(1)}+\Gamma_{bEF}(\pi)< N^{1+\epsilon-o(1)}\Rightarrow \neg LB_N(1+\epsilon)
\]
\emph{Valid when:} Requires \(\rho<1+\epsilon\) and the transcript upper bound to be an actual nonuniform circuit for the source carrier.

\paragraph{Next analysis.} Formalize \(\mathrm{Psys}_{sym}\) on finite length N: specify whether FormulaEval-SAT comprehension is a proof line, axiom schema, or extension rule; then prove either an expansion-charge parser with cost \(N^{1+q+o(1)}\) or a minimal countermodel showing surface occurrence completion is unsound.

\paragraph{Falsifiers.}
\begin{itemize}
\item falsifier: a bit-explicit transcript with N\textasciicircum{}\{rho+o(1)\}+Gamma\_bEF(pi)<N\textasciicircum{}\{1+epsilon-o(1)\} yields an actual small circuit for \(H_{N}^\mathrm{SAT}\) and invalidates \(\mathrm{LB}_{N}\).
\item falsifier: a sparse residual row set \(E_{N}<N^{\epsilon-o(1)}\) is patchable by equality circuits of cost O(E\_N N).
\item falsifier: positive target-query exponent \(\beta\) would reintroduce the transfer cap q+beta c and cannot support all-exponent P/poly exclusion.
\item falsifier: a symbolic FormulaEval-SAT proof with only surface occurrence support refutes size-only parser soundness unless the parser charges the full row-index expansion.
\item falsifier: a dense functional certificate triggers the naturalization side gate rather than providing a nonfunctional SAT-semantic escape.
\end{itemize}

\subsection*{cycle-0095: Hardness-Magnification Bridge Realization: bit-explicit parser and FormulaEval boundary}

\paragraph{Execution.} research; azure\_responses\_api; attempt primary; 600.004 seconds; 3 lemmas, 4 equations, 7 derivation steps, 4 falsifiers.

\paragraph{Learning.} Cycle 95 upgrades cycle 94 from a parser slogan to a restricted theorem packet: \(\mathrm{Psys}_{bEF}\) proofs with bit-explicit EF support have a valid charged transcript and a concrete \(\mathrm{LB}_{N}\) falsifier inequality. The surviving strong-system gap is exactly symbolic FormulaEval-SAT: row expansion costs \(N^{1+q-o(1)}\), while free symbolic comprehension has a private-row countermodel to surface occurrence completion.

\paragraph{Pressure.} The remaining named gap is \(FE_{expand,N}\): prove that any sound use of symbolic FormulaEval-SAT either expands to Omega(N\textasciicircum{}\{1+q-o(1)\}) row charge or is explicitly outside the parser model as a semantic oracle. \(\mathrm{Psys}_{bEF}\) itself is not blocked under the stated bit-explicit assumptions.

\paragraph{Derivation steps.}
\begin{enumerate}
\item Define \(\mathrm{Psys}_{bEF}\) lines as bit strings with typed tags: EF extension definition, SATeq link, \(bounded-SR_{h}\) invocation, Boolean inference, and final lower-bound certificate line.
\item Traverse the EF extension DAG in topological order. For each extension variable, record whether its definition depends on x bits, row-index bits, previous charged variables, or constants; add the transitive charged definitions to S\_dep(pi).
\item For each SATeq link, normalize the two explicitly encoded formulas, charge the row/formula selectors and link witnesses to S\_link(pi), and merge only the listed quotient classes.
\item For each \(bounded-SR_{h}\) line, unfold the allowed self-reduction tree to its declared frontier, charge \(N^{\tau+o(1)}\) frontier leaves and \(N^{\eta+o(1)}\) leaf-solver overhead per row block, giving the term \(N^{q+\eta+\tau+o(1)}\).
\item Build a nonuniform circuit for the transcript: compute the batched generator with size \(N^{\rho+o(1)}\), evaluate charged extension/link/SR components, XOR the nonexceptional row labels, and patch \(E_{N}\) exceptional rows by equality tests costing O(E\_N N).
\item Conclude C(H\_N\textasciicircum{}SAT)<=N\textasciicircum{}\{rho+o(1)\}+Gamma\_bEF(pi). Comparing with LB\_N(1+epsilon) gives the falsifier theorem whenever the right side is subthreshold.
\item Test \(\mathrm{Psys}_{sym}\) by adding one symbolic FormulaEval-SAT rule. The previous traversal cannot assign row support to the rule unless it is expanded into fixed-row instances; otherwise a private unexpanded-row labeling preserves all surface occurrences while violating the semantic rule.
\end{enumerate}

\paragraph{Equation 1: carrier.}
\[
M_N = N^{q+o(1)}, L_N = (log N)^d, H_N^SAT(x) = XOR_{j\le M_N} SAT(G_N(x,j))
\]
\emph{Valid when:} Valid when the beta-zero carrier has polylogarithmic formula length and the batched generator parameters are fixed independently of target SAT circuit exponent c.

\paragraph{Equation 2: parser\_charge.}
\[
Gamma_bEF(pi) = S_dep(pi) + S_link(pi) + N^{q+eta+tau+o(1)} + E_N N
\]
\emph{Valid when:} Valid when extension dependencies, SATeq links, bounded self-reduction frontiers, and exceptional residual rows are all parsed from bit-explicit proof lines.

\paragraph{Equation 3: transcript\_upper\_bound.}
\[
C(H_N^SAT) \le N^{rho+o(1)} + Gamma_bEF(pi)
\]
\emph{Valid when:} Applies after the parsed transcript computes the nonexceptional row parities and equality patching covers exactly the declared residual row set.

\paragraph{Equation 4: source\_falsifier.}
\[
N^{rho+o(1)} + Gamma_bEF(pi) < N^{1+epsilon-o(1)} => not LB_N(1+epsilon)
\]
\emph{Valid when:} Requires the source lower-bound model to be nonuniform circuit size for \(H_{N}^\mathrm{SAT}\) and the constructed transcript circuit to be valid for all N-bit inputs.

\paragraph{Next analysis.} Instantiate \(FE_{expand,N}\) on the smallest symbolic FormulaEval-SAT rule: one generator circuit G\_N(x,j), one asserted vector y\_j=SAT(G\_N(x,j)), and \(\alpha\) row instantiations. Prove the private-row countermodel for \(\alpha<q\) and the \(N^{1+q-o(1)}\) expansion charge for \(\alpha=q\).

\paragraph{Falsifiers.}
\begin{itemize}
\item falsifier: a parsed \(\mathrm{Psys}_{bEF}\) transcript with N\textasciicircum{}\{rho+o(1)\}+Gamma\_bEF(pi)<N\textasciicircum{}\{1+epsilon-o(1)\} gives an explicit small circuit for \(H_{N}^\mathrm{SAT}\) and invalidates LB\_N(1+epsilon).
\item falsifier: \(q<\epsilon\) makes full row enumeration cost \(N^{1+q+o(1)}\), below the claimed source threshold.
\item falsifier: primitive symbolic FormulaEval-SAT with \(sub-N^{1+q}\) charge defeats surface occurrence completion by private unexpanded-row label changes.
\item falsifier: nonuniform leakage through undeclared EF extension support invalidates any parser that is not bit-explicit or support-declared.
\end{itemize}

\subsection*{cycle-0096: Hardness-Magnification Bridge Realization}

\paragraph{Execution.} research; azure\_responses\_api; attempt primary; 600.005 seconds; 3 lemmas, 5 equations, 9 derivation steps, 6 falsifiers.

\paragraph{Learning.} Cycle 96 closes a sharper proof-delivery subgoal: \(\mathrm{Psys}_{bEF}\) parsing survives bit-level audit, and \(\mathrm{Psys}_{sym}\) FormulaEval-SAT has a precise expansion-or-countermodel boundary. The new contribution is not a \(\mathrm{P} \ne \mathrm{NP}\) advance; it is a valid route-level theorem/countermodel packet showing that symbolic formula evaluation either pays \(N^{1+q+o(1)}\) row expansion or must be treated as a semantic-oracle escape outside surface occurrence parsing.

\paragraph{Pressure.} \(SemanticForcing_{FE}\) remains open: a real-SAT FormulaEval semantic oracle might force private row labels by nonrelativizing reasoning not recoverable from occurrence syntax. The branch now treats that as a named semantic-oracle escape, not as a valid \(\mathrm{Psys}_{bEF}\) parser step.

\paragraph{Derivation steps.}
\begin{enumerate}
\item \(Source_{statement}\): LB\_N(1+epsilon) claims C(H\_N\textasciicircum{}SAT)>N\textasciicircum{}\{1+epsilon-o(1)\} for the beta-zero parity carrier with \(M_{N}=N^{q+o(1)}\), generator size \(N^{\rho+o(1)}\), and polylog SAT instance length.
\item \(Target_{statement}\): a parsed \(\mathrm{Psys}_{bEF}\) or expansion-checked \(\mathrm{Psys}_{sym}\) proof pi with N\textasciicircum{}\{rho+o(1)\}+Gamma(pi)<N\textasciicircum{}\{1+epsilon-o(1)\} is a falsifier for LB\_N(1+epsilon), not evidence for \(\mathrm{P} \ne \mathrm{NP}\).
\item \(Model_{context}\): source circuits are nonuniform fan-in-two Boolean circuits; proof systems are bit-level \(\mathrm{EF}+SATeq+bounded-SR_{h}\), optionally extended by classified FormulaEval-SAT; the simulation type is proof-to-transcript plus generator composition.
\item \(Quantifier_{context}\): N tends to infinity over valid carrier lengths, q,rho,epsilon,eta,tau are fixed exponents with o(1) slack, \(\beta=0\) because \(L_{N}\) is polylogarithmic, and no advice is allowed outside the proof string and charged transcript.
\item Parser step: traverse each bit-explicit EF extension definition in topological order, add its row/formula dependencies to \(S_{dep}\), quotient only explicitly linked SATeq encodings into \(S_{link}\), and expand \(bounded-SR_{h}\) leaves with cost \(N^{q+\eta+\tau+o(1)}\).
\item Expansion step: for every FormulaEval-SAT line declared with finite domain \(D_{u}\), instantiate the generated formula encodings for all (i,j) in \(D_{u}\) and add \(|D_{u}|L_{N}^{O(1)}\) to \(\Gamma_{sym}\); a universal domain \(R_{N}\) x [M\_N] costs \(N^{1+q+o(1)}\).
\item Countermodel step: if FormulaEval-SAT is free and row-universal, select an unexpanded pair (i*,j*) outside the surfaced occurrence set; flip its private SAT label while preserving all surfaced labels, SATeq links, bounded-SR leaves, and EF bit dependencies, causing the parity row for i* to change while the surface parser transcript is unchanged.
\item Support step: any certificate that avoids the subthreshold transcript falsifier must force at least \(N^{\epsilon-o(1)}\) contacted SAT positions; dividing by \(M_{N}=N^{q+o(1)}\) yields the q-trichotomy.
\item \(Reviewer_{findings}\): reduction-consistency accepts only beta-zero generator substitution; quantifier-uniformity flags free semantic oracles as uncharged advice unless classified; barrier-specificity marks the countermodel as relativizing for occurrence labels and outside real-SAT nonrelativizing forcing; proof-derivation status is theorem-shaped for expansion-checked \(\mathrm{Psys}_{sym}\) and blocked for free semantic \(\mathrm{Psys}_{sym}\).
\end{enumerate}

\paragraph{Equation 1: carrier\_parameters.}
\[
M_N=N^{q+o(1)},\quad L_N=(\log N)^d,\quad |R_N|=N
\]
\emph{Valid when:} Valid when N is the source length parameter and the beta-zero carrier exposes \(M_{N}\) indexed polylogarithmic SAT instances.

\paragraph{Equation 2: charged\_symmetric\_transcript.}
\[
\Gamma_{sym}(\pi)=S_{dep}(\pi)+S_{link}(\pi)+N^{q+\eta+\tau+o(1)}+E_NN+\sum_{u\in FE(\pi)} |D_u|L_N^{O(1)}
\]
\emph{Valid when:} Applies after each \(\mathrm{Psys}_{bEF}\) resource and each FormulaEval domain has been bit-explicitly parsed and charged.

\paragraph{Equation 3: transcript\_upper\_bound.}
\[
C(H_N^{SAT})\le N^{\rho+o(1)}+\Gamma_{sym}(\pi)
\]
\emph{Valid when:} Requires the parsed transcript to decide all but \(E_{N}\) exceptional rows and the generator circuit to be target-clean.

\paragraph{Equation 4: full\_formulaeval\_expansion\_charge.}
\[
\sum_{(i,j)\in R_N\times[M_N]} L_N^{O(1)}=N^{1+q+o(1)}
\]
\emph{Valid when:} Valid when a row-universal FormulaEval-SAT claim is made explicit over every certified row-index pair.

\paragraph{Equation 5: lb\_falsifier\_inequality.}
\[
N^{\rho+o(1)}+\Gamma_{sym}(\pi)< N^{1+\epsilon-o(1)}\;\Rightarrow\; \neg LB_N(1+\epsilon)
\]
\emph{Valid when:} Applies only in the same nonuniform circuit model where \(\mathrm{LB}_{N}\) is stated and all transcript charges are included.

\paragraph{Next analysis.} Fix one concrete \(\mathrm{Psys}_{sym}\) syntax: FormulaEval-SAT as either an axiom schema over A(i,j,u), an EF extension rule producing formula bits, a proof line with explicit \(D_{u}\), or a semantic oracle. For that syntax, prove the expansion parser or instantiate the minimal private-row countermodel with the smallest pi, one unexpanded pair (i*,j*), and a verified \(\Gamma_{sym}\) inequality.

\paragraph{Falsifiers.}
\begin{itemize}
\item \(\Gamma_{sym}\) subthreshold: N\textasciicircum{}\{rho+o(1)\}+Gamma\_sym(pi)<N\textasciicircum{}\{1+epsilon-o(1)\} gives an explicit circuit upper-bound falsifier for LB\_N(1+epsilon).
\item \(q<\epsilon\) row enumeration: \(M_{N}=N^{q+o(1)}\) positions are too few to support a barely-superlinear lower bound after equality patching.
\item free symbolic FormulaEval-SAT private-row countermodel: surface occurrence completion cannot recover unexpanded row labels.
\item model mismatch: a proof in semantic-oracle \(\mathrm{Psys}_{sym}\) is not automatically a \(\mathrm{Psys}_{bEF}\) charged transcript.
\item nonuniform leakage: hidden row domains or advice-selected supports must be charged or the parser theorem is invalid.
\item threshold collapse: full FormulaEval expansion costs \(N^{1+q+o(1)}\), exceeding the intended barely-superlinear window when q is positive and \(\epsilon\) is small.
\end{itemize}

\subsection*{cycle-0097: Hardness-Magnification Bridge Realization}

\paragraph{Execution.} research; azure\_responses\_api; attempt primary; 600.008 seconds; 3 lemmas, 5 equations, 7 derivation steps, 5 falsifiers.

\paragraph{Learning.} Cycle 97 closes the support-density subgoal for the \(\mathrm{Psys}_{bEF}\) packet: FormulaEval support must have absolute size \(N^{\epsilon-o(1)}\), hence density \(N^{\epsilon-q-o(1)}\) inside \(M_{N}=N^{q+o(1)}\) slots. This yields a sharp q-trichotomy and upgrades the FormulaEval boundary to an expansion-or-private-row-countermodel test. The result is a valid route-level falsifier theorem, not evidence for \(\mathrm{P} \ne \mathrm{NP}\).

\paragraph{Pressure.} The derivation does not prove a parser theorem for \(\mathrm{Psys}_{sym}\) when FormulaEval-SAT is a genuine semantic oracle; the named missing lemma is FormulaEval-forcing: real SAT semantics must force every unexpanded row/slot value from the charged transcript with subthreshold overhead.

\paragraph{Derivation steps.}
\begin{enumerate}
\item Fix \(H_{N}^\mathrm{SAT}(x)=xor_{j<M_{N}}\) SAT(G\_N(x,j)) with \(M_{N}=N^{q+o(1)}\), L\_N=polylog(N), and generator size \(N^{\rho+o(1)}\).
\item Run the \(\mathrm{Psys}_{bEF}\) parser on pi: traverse the EF extension-DAG, charge transitive bit dependencies to \(S_{dep}\), charge SATeq quotient links to \(S_{link}\), expand \(bounded-SR_{h}\) leaves into \(N^{q+\eta+\tau+o(1)}\), and record \(E_{N}\) exceptional source rows.
\item For any FormulaEval-SAT component that is not a semantic oracle, expand the selected slot set K\_idx(pi) over all N source rows; the bit-level row-index expansion costs \(X_{FE}(K)=N|K|L_{N}^{O(1)}\).
\item Combine the generator, transcript subcircuits, and equality patches for exceptional rows to obtain C(H\_N\textasciicircum{}SAT)<=N\textasciicircum{}\{rho+o(1)\}+Gamma\_bEF(pi).
\item Compare the FormulaEval expansion term with the source threshold: if \(|K|=N^{\kappa+o(1)}\), subthreshold parsing requires \(\kappa=N^{\epsilon-o(1)}\).
\item Substitute \(M_{N}=N^{q+o(1)}\) to get the density threshold \(|K|/M_{N}\ge N^{\epsilon-q-o(1)}\) and the q-trichotomy: qepsilon allows sparse density but not sparse absolute support.
\item For \(\mathrm{Psys}_{sym}\) without expansion, build the private-row countermodel: keep all charged SATeq/SR/surface labels fixed, flip one unexpanded parity slot at x*, and observe that the transcript is unchanged while H\_N(x*) flips.
\end{enumerate}

\paragraph{Equation 1: slot\_expansion\_charge.}
\[
X_FE(K)=N |K| L_N^{O(1)}=N^{1+kappa+o(1)} for |K|=N^{kappa+o(1)}
\]
\emph{Valid when:} Valid when FormulaEval-SAT is expanded over all N source rows for a bit-explicit slot set K of polylogarithmic formula length.

\paragraph{Equation 2: absolute\_support\_threshold.}
\[
|K_idx(pi)| \ge N^{epsilon-o(1)}
\]
\emph{Valid when:} Valid when the parsed certificate is intended to avoid the sparse-row equality-patching falsifier below the source lower-bound threshold.

\paragraph{Equation 3: density\_threshold.}
\[
|K_idx(pi)|/M_N \ge N^{epsilon-q-o(1)}
\]
\emph{Valid when:} Applies after \(M_{N}=N^{q+o(1)}\) is fixed and K\_idx(pi) counts parity slots in the same beta-zero carrier.

\paragraph{Equation 4: full\_slot\_expansion.}
\[
X_FE([M_N])=N M_N L_N^{O(1)}=N^{1+q+o(1)}
\]
\emph{Valid when:} Valid when every parity slot is surface-expanded over every source row and the generated SAT formulas have polylogarithmic length.

\paragraph{Equation 5: trichotomy\_bound.}
\[
q< epsilon => X_FE([M_N])< N^{1+epsilon-o(1)}; q=epsilon+o(1) => X_FE([M_N])=N^{1+epsilon+o(1)}; q> epsilon => |K|/M_N may be N^{epsilon-q-o(1)}
\]
\emph{Valid when:} Requires subthreshold dependency, link, generator, and SR overheads so the FormulaEval support term is the decisive threshold variable.

\paragraph{Next analysis.} Instantiate \(\mathrm{Psys}_{sym}\) in two concrete variants: FormulaEval as an EF extension rule with row-index bits and FormulaEval as a semantic oracle. For each, compute K\_idx(pi), X\_FE(K), and whether the private-row countermodel survives SATeq and bounded-SR constraints.

\paragraph{Falsifiers.}
\begin{itemize}
\item Subthreshold Gamma\_bEF(pi) is a valid falsifier for LB\_N(1+epsilon) in the same nonuniform circuit model.
\item \(q<\epsilon\) makes full FormulaEval slot expansion cost \(N^{1+q+o(1)}\), so row/slot enumeration falsifies any claimed need for superthreshold support when other costs are subthreshold.
\item A private-row FormulaEval countermodel falsifies surface occurrence completion for \(\mathrm{Psys}_{sym}\) semantic-oracle proofs.
\item Nonuniform leakage falsifier: advice that encodes SAT labels outside K\_idx(pi) invalidates the target-preserving parser interface.
\item Threshold-collapse falsifier: if \(\rho\), \(S_{dep}\), \(S_{link}\), SR overhead, or \(E_{N}\) N already exceed \(N^{1+\epsilon-o(1)}\), the route cannot certify a barely-superlinear source contradiction.
\end{itemize}

\subsection*{cycle-0098: Hardness-Magnification Bridge Realization}

\paragraph{Execution.} research; azure\_responses\_api; attempt primary; 600.009 seconds; 3 lemmas, 5 equations, 10 derivation steps, 5 falsifiers.

\paragraph{Learning.} Cycle 98 upgrades the prior packet from theorem slogan to dependency-checked inference structure: \(\Gamma_{bEF}\) is justified by separate EF-dependency, SATeq-link, bounded-SR, and sparse-row charges; \(\mathrm{Psys}_{bEF}\) gives a valid restricted falsifier theorem; full \(\mathrm{Psys}_{sym}\) is blocked exactly at \(FE_{Extract,N}\), not at generic naturalness or relativization folklore.

\paragraph{Pressure.} The restricted \(\mathrm{Psys}_{bEF}\) theorem is derivation-complete at proof-packet level, but the full \(\mathrm{Psys}_{sym}\) extension remains blocked at \(FE_{Extract,N}\): an advice-free theorem converting semantic FormulaEval-SAT comprehension into charged occurrence rows without paying \(N^{1+q+o(1)}\) or admitting the private-row countermodel.

\paragraph{Derivation steps.}
\begin{enumerate}
\item Define the source statement as LB\_N(1+epsilon): every nonuniform Boolean circuit for \(H_{N}^\mathrm{SAT}\) has size greater than \(N^{1+\epsilon-o(1)}\).
\item Define \(\mathrm{Psys}_{bEF}\) lines over bit encodings of G\_N(x,j), EF extension variables, explicit SATeq links, and \(bounded-SR_{h}\) invocations; reject any line whose row or formula dependency is only symbolic.
\item Topologically traverse the EF extension-DAG. For each extension variable v, add its declared dependency circuit/list to S\_dep(pi), mark it as charged, and recursively add predecessor dependencies once.
\item For each SATeq line, record both explicit formula encodings and add the bit-comparison or semantic-link charge to S\_link(pi); this prevents equality reasoning from becoming hidden advice.
\item For each \(bounded-SR_{h}\) invocation on a supported row, unfold at most \(N^{\eta+o(1)}\) leaves and attach a leaf solver or verifier of cost \(N^{\tau+o(1)}\); over \(M_{N}=N^{q+o(1)}\) slots this gives \(N^{q+\eta+\tau+o(1)}\).
\item Let \(E_{N}\) be the residual number of rows on which the parsed transcript does not certify the required parity behavior. Patch those rows by equality tests at cost \(E_{N}\) N.
\item Compose the generator circuit of size \(N^{\rho+o(1)}\) with the parsed transcript and patch circuits. The resulting nonuniform circuit has size at most N\textasciicircum{}\{rho+o(1)\}+Gamma\_bEF(pi).
\item Conclude the falsifier theorem: if this upper bound is below \(N^{1+\epsilon-o(1)}\), the asserted source lower bound LB\_N(1+epsilon) is false in the same circuit model.
\item For \(\mathrm{Psys}_{sym}\), test a FormulaEval-SAT line. If it is row-expanded, charge \(N*M_{N}=N^{1+q+o(1)}\) occurrences; if it is left symbolic, construct a private-row assignment that agrees on all surface occurrences but flips an unexpanded row, showing surface occurrence completion is unsound.
\item Run the live inference-quality handoff: source carrier \(H_{N}^\mathrm{SAT}\), target falsification of \(\mathrm{LB}_{N}\), proof model \(\mathrm{Psys}_{bEF}/\mathrm{Psys}_{sym}\), reduction type parser-to-circuit simulation, \(\beta=0\) length regime, nonuniform advice audit, and blocked lemma FormulaEval semantic-to-occurrence extraction are all explicitly named.
\end{enumerate}

\paragraph{Equation 1: carrier.}
\[
H_N^{SAT}(x)=\bigoplus_{j\le M_N} SAT(G_N(x,j)),\quad M_N=N^{q+o(1)}
\]
\emph{Valid when:} Valid when \(H_{N}^\mathrm{SAT}\) is the beta-zero parity carrier generated by \(G_{N}\) with \(M_{N}=N^{q+o(1)}\) and polylogarithmic formula length.

\paragraph{Equation 2: charged\_transcript.}
\[
\Gamma_{bEF}(\pi)=S_{dep}(\pi)+S_{link}(\pi)+N^{q+\eta+\tau+o(1)}+E_N N
\]
\emph{Valid when:} Valid when pi is a \(\mathrm{Psys}_{bEF}\) proof whose extension-variable dependencies are bit-explicit and whose SATeq and \(bounded-SR_{h}\) uses are parser-visible.

\paragraph{Equation 3: simulation\_upper\_bound.}
\[
C_N(H_N^{SAT})\le N^{\rho+o(1)}+\Gamma_{bEF}(\pi)
\]
\emph{Valid when:} Applies after the parser transcript is built and all generator, dependency, semantic-link, bounded-SR, and sparse-row patch costs are nonuniformly hardwired.

\paragraph{Equation 4: \(\mathrm{LB}_{falsifier}\).}
\[
N^{\rho+o(1)}+\Gamma_{bEF}(\pi)< N^{1+\epsilon-o(1)}\Rightarrow \neg LB_N(1+\epsilon)
\]
\emph{Valid when:} Requires the source lower-bound statement to mean general nonuniform circuit size greater than \(N^{1+\epsilon-o(1)}\) at length N.

\paragraph{Equation 5: support\_pressure.}
\[
|K_{idx}(\pi)|\ge N^{\epsilon-o(1)},\quad |K_{idx}(\pi)|/M_N\ge N^{\epsilon-q-o(1)}
\]
\emph{Valid when:} Valid when a claimed certificate avoids the transcript falsifier and must therefore touch enough complete parity rows to prevent equality patching.

\paragraph{Next analysis.} Write the formal \(FE_{Extract,N}\) proposition in two columns: left column row-expanded FormulaEval axiom schema with \(N^{1+q+o(1)}\) charge, right column semantic-oracle version with the smallest private-row countermodel; verify which column matches each \(\mathrm{Psys}_{sym}\) rule.

\paragraph{Falsifiers.}
\begin{itemize}
\item Uncharged symbolic FormulaEval-SAT is a falsifier for surface occurrence parsing via the private-row countermodel.
\item If N\textasciicircum{}\{rho+o(1)\}+Gamma\_bEF(pi)<N\textasciicircum{}\{1+epsilon-o(1)\}, then the claimed LB\_N(1+epsilon) source lower bound is false.
\item If \(q<\epsilon\) and bounded leaf-solving plus generator costs stay subthreshold, full row enumeration falsifies the intended source hardness.
\item If an EF extension variable carries a hidden truth table or nonuniform advice outside S\_dep(pi), the parser theorem is invalid for the target-preserving route.
\item Positive target-query exponent \(\beta\) in any black-box substitution repair reintroduces the archived q+beta*c threshold collapse.
\end{itemize}

\subsection*{cycle-0099: Hardness-Magnification Bridge Realization}

\paragraph{Execution.} research; azure\_responses\_api; attempt primary; 600.01 seconds; 3 lemmas, 6 equations, 7 derivation steps, 5 falsifiers.

\paragraph{Learning.} Cycle 99 converts the remaining \(\mathrm{Psys}_{sym}\) ambiguity into a sharper theorem/countermodel fork. FormulaEval-SAT cannot be treated as a free occurrence-completion rule: an expanded schema pays \(N^{1+q+o(1)}\) row-slot charge, while an unexpanded symbolic line has a private-row falsifier unless a nonrelativizing \(FE_{Extract,N}\) theorem is supplied. The valid contribution is a scoped parser-boundary packet, not evidence for \(\mathrm{P} \ne \mathrm{NP}\).

\paragraph{Pressure.} The derivation does not rule out a real-SAT-specific, nonrelativizing \(FE_{Extract,N}\) theorem showing that true formula evaluation semantics force the missing labels without row-slot expansion. It only proves the surface-occurrence parser no-go and the expansion charge for explicit schemas.

\paragraph{Derivation steps.}
\begin{enumerate}
\item Fix the quotient-label abstraction y\_\{i,j\}=SAT(G\_N(i,j)) after closing all explicitly charged SATeq and \(bounded-SR_{h}\) constraints from a \(\mathrm{Psys}_{sym}\) proof pi.
\item Run the \(\mathrm{Psys}_{bEF}\) parser on the bit-explicit part of pi, obtaining \(S_{dep}\), \(S_{link}\), bounded-SR leaf cost, and sparse-row residual \(E_{N}\).
\item Classify each FormulaEval-SAT occurrence. If it is an expanded axiom schema over rows \(R_{N}\) and slots [M\_N], add \(Ch_{FE}(R_{N})\ge |R_{N}|M_{NL,N}^{O(1)}\) to the transcript.
\item If a FormulaEval-SAT line is symbolic and the parser charges o(|R\_N|M\_N) row-slot atoms, choose a row-slot (i*,j*) outside the transitive \(SATeq/SR_{h}\) closure of the transcript.
\item Build two label assignments Y and Y' that agree on every charged atom and every explicit link, but differ at (i*,j*) and, if needed, one parity-completing uncharged atom in the same row.
\item The surface transcript and proof-local occurrence constraints evaluate identically under Y and Y', while H\_N\textasciicircum{}Y(i*) differs from H\_N\textasciicircum{}\{Y'\}(i*). Thus occurrence completion cannot soundly certify the row parity or the lower-bound statement.
\item Conclude the dichotomy: expanded FormulaEval pays \(N^{1+q+o(1)}\) for global row comprehension, while unexpanded FormulaEval must be declared a semantic oracle requiring a separate nonrelativizing \(FE_{Extract,N}\) theorem.
\end{enumerate}

\paragraph{Equation 1: slot\_count.}
\[
M_N=N^{q+o(1)}
\]
\emph{Valid when:} Valid when the beta-zero carrier family fixes q and all hidden factors are subpolynomial in N.

\paragraph{Equation 2: carrier\_row.}
\[
H_N^{SAT}(i)=\oplus_{j\le M_N} SAT(G_N(i,j))
\]
\emph{Valid when:} Applies when row i is a public source row and every generated instance has polylogarithmic length.

\paragraph{Equation 3: formulaeval\_expansion\_charge.}
\[
Ch_{FE}(R_N)\ge |R_N|\,M_N\,L_N^{O(1)}=|R_N|N^{q+o(1)}
\]
\emph{Valid when:} Requires the FormulaEval-SAT mechanism to range over all slots for each charged row without hidden advice.

\paragraph{Equation 4: full\_row\_expansion.}
\[
|R_N|=N^{1-o(1)}\Rightarrow Ch_{FE}(R_N)=N^{1+q+o(1)}
\]
\emph{Valid when:} Valid when FormulaEval-SAT is a global row comprehension rather than a finite listed subclaim.

\paragraph{Equation 5: \(parsed_{upper,bound,with,FE}\).}
\[
C(H_N^{SAT})\le N^{\rho+o(1)}+\Gamma_{bEF}(\pi)+Ch_{FE}(R_N)+E_NN
\]
\emph{Valid when:} Applies after the parser preserves EF dependencies, SATeq links, bounded-SR leaves, and sparse-row patching costs.

\paragraph{Equation 6: \(source_{\mathrm{LB},falsifier}\).}
\[
N^{\rho+o(1)}+\Gamma_{bEF}(\pi)+Ch_{FE}(R_N)+E_NN< N^{1+\epsilon-o(1)}\Rightarrow LB_N(1+\epsilon)\text{ is false}
\]
\emph{Valid when:} Requires the parsed transcript to compute the same carrier \(H_{N}\) on all but \(E_{N}\) explicitly patched rows.

\paragraph{Next analysis.} Write the formal \(FE_{Extract,N}\) blocked-lemma proof: define the quotient-label model after \(SATeq/SR_{h}\) closure, prove the unmentioned row-slot existence bound, construct the two assignments Y,Y', and verify that the carrier row parity changes while the parsed transcript remains identical.

\paragraph{Falsifiers.}
\begin{itemize}
\item Private-row FormulaEval countermodel: an unexpanded symbolic FormulaEval line with sub-NM charged occurrences cannot be a valid occurrence parser because one unmentioned row-slot label can flip a carrier row.
\item Subthreshold transcript falsifier: if N\textasciicircum{}\{rho+o(1)\}+Gamma\_bEF(pi)+Ch\_FE+E\_NN<N\textasciicircum{}\{1+epsilon-o(1)\}, the claimed LB\_N(1+epsilon) lower bound is false in the same circuit model.
\item \(q<\epsilon\) support falsifier: the required absolute FormulaEval support \(N^{\epsilon-o(1)}\) exceeds the available \(M_{N}=N^{q+o(1)}\) slots, so row enumeration and patching collapse the source lower bound.
\item Nonuniform leakage falsifier: any hidden advice list of row supports or SAT labels invalidates the target-preserving parser because it is not charged in \(\Gamma_{bEF}\) or \(Ch_{FE}\).
\item Semantic-oracle escape falsifier for promotion: if FormulaEval-SAT is left unclassified, the proof packet is invalid for paper promotion even if the \(\mathrm{Psys}_{bEF}\) subtheorem is correct.
\end{itemize}

\subsection*{cycle-0100: Hardness-Magnification Bridge Realization}

\paragraph{Execution.} research; azure\_responses\_api; attempt compact\_retry; 1140.976 seconds; 1 lemmas, 3 equations, 5 derivation steps, 4 falsifiers.

\paragraph{Learning.} Cycle 100 isolates the \(bounded-SR_{h}\) normal form at bit level: SR contributes exactly B\_SR(pi) to \(\Gamma_{bEF}\) for \(\mathrm{Psys}_{bEF}^{-}\). The remaining full-system obstruction is not generic soundness but \(FE_{Extract,N}\) for symbolic FormulaEval rows.

\paragraph{Pressure.} The derivation is valid only for \(\mathrm{Psys}_{bEF}^{-}\). Full \(\mathrm{Psys}_{sym}\) remains blocked at \(FE_{Extract,N}\): a symbolic FormulaEval-SAT line may select rows without exposing K\_idx(pi) or Nbr(j,u).

\paragraph{Derivation steps.}
\begin{enumerate}
\item Definition/assumption context: this derivation is interpreted inside Hardness-Magnification Bridge Realization; source object: the branch-local proof object or transfer carrier named by this cycle; target object: the target theorem, obstruction, or upper-bound/countermodel handoff named by this cycle; validity regime/model condition: Source model: \(\mathrm{Psys}_{bEF}^{-}\) verifier is uniform and accepts only binary row indices, explicit formula-bit addresses, declared EF dependency pointers, and SATeq links between named encodings.; SR regime: per row, the SR forest has at most \(N^{\eta+o(1)}\) nodes and each node exposes at most \(N^{\tau+o(1)}\) formula, pointer, and leaf-solver bits; fixed h or \(h=N^{o(1)}\) is absorbed in \(\tau\).; Advice status: pi and \(Q_{pi}\) may be hardwired into the resulting circuit, but hidden SAT labels, private row answers, and semantic FormulaEval witnesses are not allowed.; FormulaEval status: symbolic FormulaEval-SAT is excluded from \(\mathrm{Psys}_{bEF}^{-}\) unless its RowSet and touched neighborhoods are expanded and charged outside B\_SR(pi)..
\item Fix the \(SR_{h}^{bit}\) line encoding: j uses ceil(log \(M_{N}\)) bits, node u uses \(\eta\) log N+o(log N) bits, child pointers are proof-line indices, and every interval in I names generator-output bit positions.
\item The parser traverses the transitive EF DAG, recomputes each local restriction map from the public generator bits, and appends (j,u,Nbr,ptr,leaf) to \(Q_{pi}\) only after syntactic verification.
\item Line induction: EF lines add \(S_{dep}\), SATeq lines add \(S_{link}\), \(SR_{h}^{bit}\) lines add exactly their declared \(B_{SR}\) tuples, and sparse row failures are delayed to the \(E_{NN}\) equality-patching term.
\item The generator circuit of size \(N^{\rho+o(1)}\) plus hardwired \(Q_{pi}\) evaluates all charged rows and patches \(E_{N}\) residual rows, giving the displayed \(\mathrm{LB}_{N}\) falsifier inequality.
\end{enumerate}

\paragraph{Equation 1: SR-neighborhood-charge.}
\[
B_{\mathrm{SR}}(\pi):=\sum_{(j,u)\in SR(\pi)}(|Nbr(j,u)|+|ptr(j,u)|+|leaf(j,u)|)\le |K_{idx}(\pi)|N^{\eta+\tau+o(1)}\le N^{q+\eta+\tau+o(1)}
\]
\emph{Valid when:} Valid when \(\mathrm{Psys}_{bEF}^{-}\) uses only bit-explicit \(SR_{h}\) lines, K\_idx(pi) is explicitly enumerable, and each row has at most \(N^{\eta+o(1)}\) SR nodes with \(N^{\tau+o(1)}\) exposed bits.

\paragraph{Equation 2: \(\Gamma-bEF-split\).}
\[
\Gamma_{\mathrm{bEF}}(\pi)=S_{dep}(\pi)+S_{link}(\pi)+B_{\mathrm{SR}}(\pi)+E_NN
\]
\emph{Valid when:} Valid when EF dependencies, SATeq links, bounded-SR nodes, and sparse residual rows are disjoint charged transcript classes in the parsed proof.

\paragraph{Equation 3: LB-falsifier-inequality.}
\[
N^{\rho+o(1)}+\Gamma_{\mathrm{bEF}}(\pi)< N^{1+\epsilon-o(1)}\Rightarrow C(H_N^{SAT})< N^{1+\epsilon-o(1)}
\]
\emph{Valid when:} Valid when the parsed transcript plus generator computes \(H_{N}^\mathrm{SAT}\) except on \(E_{N}\) explicitly patched rows and no hidden semantic advice is used.

\paragraph{Next analysis.} Prove \(RowSetEnumerator_{N}\) for an interval-union or \(\mathrm{AC0}\) FormulaEval selector and compute alpha-charge; if enumeration fails, emit the minimal private-row countermodel.

\paragraph{Falsifiers.}
\begin{itemize}
\item falsifier: a row selector stored as a circuit for RowSet rather than enumerated j values invalidates the valid SR charge and triggers \(FE_{Extract,N}\).
\item falsifier: a child formula produced by semantic FormulaEval-SAT rather than local restriction has the private-row countermodel and is outside \(\mathrm{Psys}_{bEF}^{-}\).
\item falsifier: \(q+\eta+\tau\ge 1+\epsilon\) causes threshold collapse, so the charged transcript no longer refutes LB\_N(1+epsilon).
\item falsifier: an EF extension variable with undeclared transitive support leaks advice; the parser must quarantine the proof instead of charging \(\Gamma_{bEF}\).
\end{itemize}

\subsection*{cycle-0101: Hardness-Magnification Bridge Realization}

\paragraph{Execution.} research; azure\_responses\_api; attempt compact\_retry; 600.003 seconds; 1 lemmas, 3 equations, 4 derivation steps, 4 falsifiers.

\paragraph{Learning.} Cycle 101 upgrades \(bounded-SR_{h}\) from a named cost to a bit-normal restricted parser theorem: SR contributes \(B_{SR}(pi)\le N^{q+\eta+\tau+o(1)}\) to \(\Gamma_{bEF}\). The remaining live obstruction is exactly \(FE_{Extract,N}\) for symbolic FormulaEval rows.

\paragraph{Pressure.} The induction excludes symbolic FormulaEval-SAT; \(FE_{Extract,N}\) is still needed to enumerate RowSet(a) or force the \(N^{1+q+o(1)}\) expansion charge.

\paragraph{Derivation steps.}
\begin{enumerate}
\item Encode each SR line as (\(tag=SR,row\) j,node u,depth t,branch b,parent pointer,occurrence offsets). Row j has (q+o(1))log N bits; node/depth fields contribute the eta,tau overhead; local neighborhoods are polylog N.
\item Verifier checks parent-child formula transformations by rereading only listed neighborhoods and SATeq-normalized bit strings. No semantic SAT truth is queried; the rule contributes its declared row to \(K_{SR}\) and listed leaves to \(B_{SR}\).
\item Induct on proof lines: EF lines add \(S_{dep}\), SATeq lines add \(S_{link}\), \(SR_{h}\) lines add the bounded-SR term, and pointers compose parsed transcripts. Thus \(Q_{i}\) has exactly the partial \(\Gamma_{bEF}\) cost up to uniform o(1) encoding slack.
\item Given \(Q_{pi}\) and \(G_{N}\), build a circuit generating referenced formulas, implementing transcript leaf decisions, patching \(E_{N}\) residual rows, and XORing outputs. A subthreshold cost is a valid falsifier for LB\_N(1+epsilon).
\end{enumerate}

\paragraph{Equation 1: SR-charge.}
\[
B_{SR}(\pi)=\sum_{r\in K_{SR}(\pi)} O(N^{\eta+\tau+o(1)}) \le N^{q+\eta+\tau+o(1)}.
\]
\emph{Valid when:} Valid when RowSet(pi) is syntactically declared and each \(bounded-SR_{h}\) row has at most \(N^{\eta+\tau+o(1)}\) exposed leaf and neighborhood checks.

\paragraph{Equation 2: \(\Gamma-bEF\).}
\[
\Gamma_{bEF}(\pi)=S_{dep}(\pi)+S_{link}(\pi)+B_{SR}(\pi)+E_N N.
\]
\emph{Valid when:} Applies after EF dependencies, SATeq links, bounded-SR fields, and sparse row residuals are separately parsed with no semantic advice leakage.

\paragraph{Equation 3: Transcript-upper-bound.}
\[
C(H_N^{SAT}) \le N^{\rho+o(1)}+\Gamma_{bEF}(\pi);\quad N^{\rho+o(1)}+\Gamma_{bEF}(\pi)< N^{1+\epsilon-o(1)} \Rightarrow \neg LB_N(1+\epsilon).
\]
\emph{Valid when:} Requires the beta-zero generator substitution and a valid transcript that computes or patches every row touched by pi in the same circuit model.

\paragraph{Next analysis.} Write the finite verifier table for one \(SR_{h}\) rule: tag layout, row j encoding, node/depth fields, parent pointers, occurrence-neighborhood list, and the exact charge contributing to B\_SR(pi).

\paragraph{Falsifiers.}
\begin{itemize}
\item A semantic FormulaEval line with a private row whose truth depends on an unread enable bit: the surface parser sees no row, but the proof output changes; this falsifier blocks \(\mathrm{Psys}_{sym}\) without \(FE_{Extract,N}\).
\item If \(q+\eta+\tau\ge 1+\epsilon\), the bounded-SR charge alone reaches the lower-bound threshold, so the transcript no longer contradicts \(\mathrm{LB}_{N}\).
\item If RowSet is supplied as nonuniform advice of size o(N\textasciicircum{}epsilon), touched rows may be hidden and the advice-free normal form is invalid.
\item If SATeq uses a different formula encoding than \(G_{N}\), local neighborhoods do not certify target-preserving row identity.
\end{itemize}

\subsection*{cycle-0102: Hardness-Magnification Bridge Realization - cycle 102 \(SR_{h}\) normal form audit}

\paragraph{Execution.} research; azure\_responses\_api; attempt primary; 600.007 seconds; 3 lemmas, 4 equations, 8 derivation steps, 6 falsifiers.

\paragraph{Learning.} Cycle 102 upgrades the \(bounded-SR_{h}\) subgoal from an exponent slogan to a parser theorem for the explicit rule schema: \(SR_{h}\) contributes exactly B\_SR(pi) to \(\Gamma_{bEF}\), and the line induction gives a valid restricted \(\mathrm{LB}_{N}\) falsifier. The remaining full-system issue is sharply isolated as \(FE_{Extract,N}\) for symbolic FormulaEval-SAT, not generic naturalness, relativization, or P vs NP itself.

\paragraph{Pressure.} No blocking gap remains for the restricted \(SR_{h}^{bit}\) normal form and line-induction packet. The full \(\mathrm{Psys}_{sym}\) theorem is blocked at \(FE_{Extract,N}\) or \(RowSetEnumerator_{N}\) for symbolic FormulaEval-SAT.

\paragraph{Derivation steps.}
\begin{enumerate}
\item Define \(\mathrm{Psys}_{bEF}^{-}\) syntax with bit-explicit EF extension definitions, SATeq quotient links, and \(SR_{h}^{bit}\) records; exclude FormulaEval-SAT from this restricted theorem.
\item For an \(SR_{h}^{bit}\) record, parse fields tag, line id, parent ids, root list \(R_{ell}\), row-slot index r=(a,j), node address u, branch word z, child pointers, sign bit, and occurrence-neighborhood list O\_ell(r,u,z).
\item Use \(b_{idx}=(1+q+o(1))log\) N to encode row-slot roots and O(log \(L_{N}\)) bits for local formula occurrences; verify each SR child by local substitution on at most h variables and clause simplification.
\item Charge each verified SR local check to \(|R_{ell}||U_{ell}||V_{ell}|L_{N}^{O_{h}(1)}\); summing gives B\_SR(pi) and, under the exponent assumptions, \(B_{SR}(pi)\le N^{q+\eta+\tau+o(1)}\).
\item Run line induction: EF definitions add transitive support \(S_{dep}\), SATeq links add quotient support \(S_{link}\), \(SR_{h}\) lines add \(B_{SR}\), and residual bad rows add \(E_{N}\) N by equality patching.
\item Construct the transcript-driven circuit: compute \(G_{N}\) with size \(N^{\rho+o(1)}\), consult parsed SAT/SR transcript entries, xor the \(M_{N}\) slot values, and patch \(E_{N}\) rows.
\item Derive C(H\_N\textasciicircum{}SAT)<=N\textasciicircum{}\{rho+o(1)\}+O(Gamma\_bEF(pi)); if this is below \(N^{1+\epsilon-o(1)}\), the lower-bound sentence LB\_N(1+epsilon) is falsified in the same nonuniform model.
\item Attempt to extend the proof to \(\mathrm{Psys}_{sym}\) FormulaEval-SAT; classify the mechanism. Expansion is charged \(N^{1+q+o(1)}\), while an unexpanded semantic line admits the private-row flip unless \(FE_{Extract,N}\) is supplied.
\end{enumerate}

\paragraph{Equation 1: row\_slot\_index\_bits.}
\[
b_{idx}(N)=\lceil \log N\rceil+\lceil \log M_N\rceil=(1+q+o(1))\log N
\]
\emph{Valid when:} Valid when row-slot roots are pairs from [N] times [M\_N] and \(M_{N}=N^{q+o(1)}\) is part of the source carrier definition.

\paragraph{Equation 2: local\_occurrence\_neighborhood.}
\[
|O_\ell(r,u,z)|\le L_N^{c_h}=N^{o(1)}
\]
\emph{Valid when:} Valid when h is fixed and the SR verifier checks only bounded-depth substitution and clause simplification in polylogarithmic-size generated formulas.

\paragraph{Equation 3: \(bounded_{SR,charge,exact}\).}
\[
B_{SR}(\pi)=\sum_{\ell\in SR(\pi)} |R_\ell|\,|U_\ell|\,|V_\ell|\,L_N^{O_h(1)}
\]
\emph{Valid when:} Valid when every \(SR_{h}\) proof record explicitly lists roots, node addresses, branch words, and occurrence neighborhoods charged by the parser.

\paragraph{Equation 4: \(bounded_{SR,exponent,bound}\).}
\[
B_{SR}(\pi)\le N^{q+\eta+\tau+o(1)}
\]
\emph{Valid when:} Applies after the proof family has \(N^{q+o(1)}\) total charged roots, \(N^{\eta+o(1)}\) node overhead, \(N^{\tau+o(1)}\) leaf overhead, and only polylogarithmic local verification factors.

\paragraph{Next analysis.} Instantiate \(RowSetEnumerator_{N}\) for one transparent \(\mathrm{Psys}_{sym}^{sel}\) fragment, preferably interval-union or \(\mathrm{AC0}\) row selectors, and decide whether FormulaEval-SAT is an expanded axiom schema with \(N^{1+q+o(1)}\) charge or a semantic oracle exposed by the private-row falsifier.

\paragraph{Falsifiers.}
\begin{itemize}
\item Hidden extension support: an EF variable whose dependency set is not revealed by the transitive bit-DAG invalidates \(S_{dep}\) and defeats the parser.
\item Semantic FormulaEval-SAT: an unexpanded symbolic line can hide a private row-slot label; the private-row falsifier changes \(H_{N}^\mathrm{SAT}\) while preserving surfaced occurrences.
\item Threshold collapse: if \(\rho\) or q+eta+tau reaches 1+epsilon, \(\Gamma_{bEF}\) no longer gives a subthreshold contradiction to LB\_N(1+epsilon).
\item Nonuniform leakage: a row selector supplied as advice rather than by \(RowSetEnumerator_{N}\) can fake occurrence support without paying \(N^{1+q+o(1)}\) expansion.
\item SR leaf overreach: if the bounded-SR rule imports an uncharged SAT solver for leaves, the \(B_{SR}(pi)\le N^{q+\eta+\tau+o(1)}\) accounting is not valid.
\item Model mismatch: a claim about full EF+SAT or \(\mathrm{Psys}_{sym}\) cannot inherit the \(\mathrm{Psys}_{bEF}^{-}\) theorem unless every semantic rule is classified and charged.
\end{itemize}

\subsection*{cycle-0103: Hardness-Magnification Bridge Realization}

\paragraph{Execution.} research; azure\_responses\_api; attempt primary; 600.005 seconds; 3 lemmas, 7 equations, 8 derivation steps, 4 falsifiers.

\paragraph{Learning.} Cycle 103 discharges a concrete FormulaEval subcase: transparent interval-union selectors admit a valid \(FE_{Extract,N}\) parser with N|RowSet| expansion charge, so global FormulaEval costs \(N^{1+q+o(1)}\) and yields the q/epsilon support trichotomy. The full \(\mathrm{Psys}_{sym}\) claim remains quarantined exactly at semantic FormulaEval extraction; unexpanded symbolic FormulaEval is refuted by the private-row falsifier, not by generic barrier folklore.

\paragraph{Pressure.} The derivation does not prove \(FE_{Extract,N}\) for arbitrary \(\mathrm{Psys}_{sym}\) selectors or for real FormulaEval-SAT comprehension whose row set is defined by semantic SAT facts rather than displayed interval endpoints.

\paragraph{Derivation steps.}
\begin{enumerate}
\item Fix N, \(M_{N}=N^{q+o(1)}\), b\_N=ceil(log \(M_{N}\)), and encode each row index j by \(b_{N}\) bits; define \(\mathrm{Psys}_{sym}^{IU}\) by adding interval-union FormulaEval-SAT lines to the already parsed \(\mathrm{Psys}_{bEF}^{-}\) rules.
\item For a FormulaEval line a, read all interval endpoints from pi, merge overlaps, enumerate RowSet(a), and append tuples (line a,row j,formula-bit neighborhood from \(G_{N}\)) to \(K_{idx}^{FE}\).
\item Extend the cycle-102 line invariant: after line t, every semantic SAT-label dependency used by the first t lines is either in \(K_{idx}\), in a SATeq link charge, in a \(bounded-SR_{h}\) leaf charge \(B_{SR}\), or in the residual error set \(E_{N}\).
\item For EF extension lines, traverse the transitive definition DAG and add \(S_{dep}\); for SATeq lines, add \(S_{link}\); for \(bounded-SR_{h}\) lines, use \(B_{SR}(pi)\le N^{q+\eta+\tau+o(1)}\) from the \(SR_{h}\) bit-normal form.
\item For interval-union FormulaEval lines, the only new dependency is membership in RowSet(a); enumeration is advice-free, so the added circuit transcript cost is N|RowSet(a)| plus negligible row-encoding and occurrence-neighborhood overhead.
\item Induct over proof lines to obtain Gamma\_sym\textasciicircum{}IU(pi)=S\_dep+S\_link+B\_SR+N|K\_idx|+E\_NN+N\textasciicircum{}\{o(1)\} and no hidden semantic row dependency inside \(\mathrm{Psys}_{sym}^{IU}\).
\item Build the falsifier circuit: compute \(G_{N}\) with size \(N^{\rho+o(1)}\), use the parsed transcript for all charged rows, and patch \(E_{N}\) residual source rows by equality gates, giving C(H\_N)<=N\textasciicircum{}\{rho+o(1)\}+Gamma\_sym\textasciicircum{}IU(pi).
\item If a FormulaEval line is instead declared semantic and unexpanded, remove it from \(\mathrm{Psys}_{sym}^{IU}\); if kept without RowSet charge, choose j* outside \(K_{idx}\) and apply the private-row spoof, showing surface occurrence completion is unsound.
\end{enumerate}

\paragraph{Equation 1: carrier.}
\[
H_N^{SAT}(x)=\bigoplus_{j=1}^{M_N} SAT(G_N(x,j)),\quad M_N=N^{q+o(1)}
\]
\emph{Valid when:} Valid when the source carrier is the beta-zero generated parity problem with explicitly indexed SAT instances.

\paragraph{Equation 2: \(IU_{row,expansion}\).}
\[
\Gamma_{FE}^{IU}(\pi)=N\cdot |K_{idx}^{FE}(\pi)|+N^{o(1)}\sum_a R_a
\]
\emph{Valid when:} Valid when every FormulaEval selector is an explicit interval union whose endpoints are proof bits.

\paragraph{Equation 3: \(global_{FE,cost}\).}
\[
\Gamma_{FE}^{all}=N\cdot M_N=N^{1+q+o(1)}
\]
\emph{Valid when:} Applies when a FormulaEval line ranges over all parity slots of the generated carrier.

\paragraph{Equation 4: total\_transcript\_cost.}
\[
\Gamma_{sym}^{IU}(\pi)=S_{dep}(\pi)+S_{link}(\pi)+B_{SR}(\pi)+N|K_{idx}(\pi)|+E_NN+N^{o(1)}
\]
\emph{Valid when:} Requires the bit-explicit EF, SATeq, \(bounded-SR_{h}\), and interval-union FormulaEval parsers to share one row-index convention.

\paragraph{Equation 5: restricted\_falsifier\_bound.}
\[
C(H_N^{SAT})\le N^{\rho+o(1)}+\Gamma_{sym}^{IU}(\pi)
\]
\emph{Valid when:} Applies after the parsed transcript plus generator uniformly simulates all non-residual carrier rows.

\paragraph{Equation 6: \(\mathrm{LB}_{falsifier,threshold}\).}
\[
N^{\rho+o(1)}+\Gamma_{sym}^{IU}(\pi)< N^{1+\epsilon-o(1)}\Rightarrow LB_N(1+\epsilon)\text{ is false}
\]
\emph{Valid when:} Valid when \(\mathrm{LB}_{N}\) asserts circuit size at least \(N^{1+\epsilon-o(1)}\) in the same nonuniform circuit model.

\paragraph{Equation 7: support\_density\_pressure.}
\[
N|K_{idx}(\pi)|\ge N^{1+\epsilon-o(1)}\Rightarrow |K_{idx}(\pi)|/M_N\ge N^{\epsilon-q-o(1)}
\]
\emph{Valid when:} Requires the transcript falsifier inequality to be the only obstruction to the claimed lower-bound certificate.

\paragraph{Next analysis.} Extend \(RowSetEnumerator_{N}\) from interval-union selectors to a named \(\mathrm{AC0}\) selector fragment by full scanning over [M\_N]; compute whether scan overhead \(N^{q+o(1)}\) is dominated by N|RowSet| and identify the smallest \(\mathrm{AC0}-private-row\) countermodel if enumeration depends on semantic advice.

\paragraph{Falsifiers.}
\begin{itemize}
\item falsifier: unexpanded symbolic FormulaEval-SAT with an uncharged private row j* makes two label worlds transcript-identical but parity-opposite, so surface occurrence completion is invalid.
\item falsifier: if \(q<\epsilon\) and global FormulaEval expands to all \(M_{N}\) rows, then the \(N^{1+q+o(1)}\) row charge is still subthreshold and can collapse the claimed \(\mathrm{LB}_{N}\) lower bound when \(\rho<1+\epsilon\).
\item falsifier: if interval endpoints or selector descriptions are supplied as nonuniform advice rather than proof bits, \(RowSetEnumerator_{N}\) is not advice-free and the \(\mathrm{Psys}_{sym}^{IU}\) theorem is not valid.
\item falsifier: positive target-query exponent \(\beta\) in any attempted black-box substitution reintroduces the transfer cap q+beta c and blocks all-exponent \(\mathrm{NP}\not\subset \mathrm{P}/\mathrm{poly}\) consequences.
\end{itemize}

\subsection*{cycle-0104: Hardness-Magnification Bridge Realization}

\paragraph{Execution.} research; azure\_responses\_api; attempt primary; 600.01 seconds; 3 lemmas, 6 equations, 8 derivation steps, 5 falsifiers.

\paragraph{Learning.} Cycle 104 pins the \(bounded-SR_{h}\) normal form to concrete verifier fields rather than an exponent slogan. The restricted \(\mathrm{Psys}_{bEF}^{-}\) plus interval-union FormulaEval theorem packet is now valid as a paper-level parser/falsifier result, while full semantic FormulaEval remains quarantined at \(FE_{Extract,N}\) by a private-row countermodel.

\paragraph{Pressure.} Full \(\mathrm{Psys}_{sym}\) remains blocked exactly at \(FE_{Extract,N}\) for semantic FormulaEval selectors and at arbitrary extension variables whose row supports are not bit-explicit. The restricted \(\mathrm{Psys}_{bEF}^{-}\) plus interval-union FormulaEval packet is valid; the full semantic proof-system parser is not proved.

\paragraph{Derivation steps.}
\begin{enumerate}
\item Fix valid lengths N and the carrier \(H_{N}^\mathrm{SAT}\) with \(\beta=0\) polylogarithmic SAT instances. Define \(\mathrm{Psys}_{bEF}^{-}\) as bit-explicit \(\mathrm{EF}+SATeq+bounded-SR_{h}\), excluding symbolic FormulaEval except charged \(IU_{FEExtract}\) lines.
\item For each EF extension line, traverse the acyclic definition DAG and append its declared input-bit support to S\_dep(pi). Parent pointers are O(log |pi|); support entries are row-indexed occurrence coordinates.
\item For each SATeq link, record the pair of formula ids, row index, occurrence bijection, and local equality proof. Charge S\_link(pi) by the number of linked occurrence coordinates, not by semantic formula identity.
\item For each \(SR_{h}\) line, write records (\(tag,i,parent0,parent1,j,u,h,\alpha,p_{1},...,p_{h},leaf_{id},Delta_{j,\alpha},y_{alpha}\)). The local verifier checks that \(Delta_{j,\alpha}\) is exactly the touched occurrence-neighborhood induced by restricting formula u along branch \(\alpha\).
\item Count one SR leaf record per row and branch. With \(|K_{SR}|\) rows, \(A_{N}=2^{h_{N}}\), and \(W_{N}\) maximal listed neighborhood size, the charged bit work is O(|K\_SR|A\_NW\_N polylog N), giving \(N^{q+\eta+\tau+o(1)}\) in the global row regime.
\item Run line induction: every proof line has a transcript \(T_{i}\) whose exposed-row semantics matches the line under all row-label assignments agreeing on \(T_{i}\). EF lines union transcripts, SATeq lines add links, \(SR_{h}\) lines add exactly their leaf neighborhoods, and residual disagreement rows are isolated into \(E_{N}\).
\item Combine \(S_{dep}\), \(S_{link}\), \(B_{SR}\), and \(E_{NN}\) with the generator circuit \(N^{\rho+o(1)}\) to obtain C(H\_N)<=N\textasciicircum{}\{rho+o(1)\}+Gamma\_bEF(pi). If this is below \(N^{1+\epsilon-o(1)}\), the source lower-bound statement LB\_N(epsilon) is false in the same circuit model.
\item Integrate the already validated interval-union \(FE_{Extract,N}\) fragment by adding \(B_{FE}^{IU}=N|RowSet|+lower-order\) descriptor charge to \(\Gamma\). Do not integrate semantic FormulaEval without \(FE_{Extract,N}\); the hidden-row flip countermodel blocks occurrence completion.
\end{enumerate}

\paragraph{Equation 1: carrier.}
\[
H_N^{\mathrm{SAT}}(x)=\bigoplus_{j=1}^{M_N} \mathrm{SAT}(G_N(x,j)),\qquad M_N=N^{q+o(1)}
\]
\emph{Valid when:} Valid when \(G_{N}\) uniformly generates the row formulas at the fixed valid length N and the target-query length is polylogarithmic.

\paragraph{Equation 2: \(SR_{branch,and,neighborhood,parameters}\).}
\[
A_N=2^{h_N}=N^{\eta+o(1)},\qquad W_N=\max_{j,\alpha}|\Delta_{j,\alpha}|=N^{\tau+o(1)}
\]
\emph{Valid when:} Applies when the bounded self-reduction depth and every touched occurrence-neighborhood list are explicitly encoded in the proof line.

\paragraph{Equation 3: \(SR_{charge}\).}
\[
B_{\mathrm{SR}}(\pi)\le c\,|K_{\mathrm{SR}}(\pi)|A_NW_N\operatorname{polylog}N\le N^{q+\eta+\tau+o(1)}
\]
\emph{Valid when:} Requires the SR parser to enumerate all charged rows and all branch leaves without semantic advice or hidden formula evaluation.

\paragraph{Equation 4: transcript\_budget.}
\[
\Gamma_{\mathrm{bEF}}(\pi)=S_{\mathrm{dep}}(\pi)+S_{\mathrm{link}}(\pi)+B_{\mathrm{SR}}(\pi)+E_NN
\]
\emph{Valid when:} Valid when EF dependencies, SATeq links, bounded-SR records, and sparse residual patches are the only charged proof mechanisms.

\paragraph{Equation 5: restricted\_falsifier\_bound.}
\[
C(H_N^{\mathrm{SAT}})\le N^{\rho+o(1)}+\Gamma_{\mathrm{bEF}}(\pi)
\]
\emph{Valid when:} Applies after the parsed transcript supplies charged row computations and residual rows are corrected by equality patches.

\paragraph{Equation 6: \(\mathrm{LB}_{refutation,threshold}\).}
\[
N^{\rho+o(1)}+\Gamma_{\mathrm{bEF}}(\pi)< N^{1+\epsilon-o(1)}\ \Longrightarrow\ \neg LB_N(\epsilon)
\]
\emph{Valid when:} Valid when LB\_N(epsilon) is the exact source lower-bound claim \(C(H_{N})>N^{1+\epsilon-o(1)}\) in the same circuit model.

\paragraph{Next analysis.} Instantiate \(RowSetEnumerator_{N}\) for one richer transparent selector class beyond interval-unions, preferably sorted \(\mathrm{AC0}\) threshold-free selectors with explicit scan charge, and decide whether it expands at N|RowSet|, collapses to an \(N^{q}\) scan bound, or admits a private-row countermodel.

\paragraph{Falsifiers.}
\begin{itemize}
\item Subthreshold transcript falsifier: N\textasciicircum{}\{rho+o(1)\}+S\_dep+S\_link+B\_SR+B\_FE+E\_NN<N\textasciicircum{}\{1+epsilon-o(1)\} gives an explicit circuit upper bound for \(H_{N}\) and refutes LB\_N(epsilon).
\item Private-row falsifier: a semantic FormulaEval line with an unenumerated selected row permits two row-label worlds with identical parsed transcript but opposite carrier parity.
\item Threshold-collapse falsifier: if \(q+\eta+\tau\ge 1+\epsilon\), \(bounded-SR_{h}\) expansion alone reaches the source lower-bound scale and the proof-system audit no longer magnifies a barely-superlinear claim.
\item Nonuniform leakage falsifier: any EF extension variable whose support is not bit-explicit can encode row advice outside \(\Gamma_{bEF}\) and invalidates the parser theorem.
\item Model-mismatch falsifier: an occurrence-local parser theorem cannot be advertised for real SAT-semantic \(\mathrm{Psys}_{sym}\) unless the proof object is shown nonrelativizing and \(FE_{Extract,N}-valid\).
\end{itemize}

\subsection*{cycle-0105: Hardness-Magnification Bridge Realization}

\paragraph{Execution.} research; azure\_responses\_api; attempt primary; 600.005 seconds; 3 lemmas, 5 equations, 10 derivation steps, 5 falsifiers.

\paragraph{Learning.} Cycle 105 converts the promoted restricted packet into an auditable inference-quality theorem: \(\mathrm{Psys}_{bEF}^{-}\) plus transparent interval-union FormulaEval has a valid parser-to-circuit falsifier with \(\Gamma_{bEF}^{IU}\), while full semantic FormulaEval is not merely hard or barrier-limited but blocked by the exact \(FE_{Extract,N}^{sem}\) private-row countermodel. The contribution is a paper-level parser/falsifier boundary, not evidence for \(\mathrm{P} \ne \mathrm{NP}\).

\paragraph{Pressure.} No remaining gap for \(\mathrm{Psys}_{bEF}^{-}\) with transparent interval-union FormulaEval. Full \(\mathrm{Psys}_{sym}\) remains blocked at \(FE_{Extract,N}^{sem}\): a semantic FormulaEval rule may compute SAT facts over hidden selected formulas without outputting RowSet indices or occurrence neighborhoods.

\paragraph{Derivation steps.}
\begin{enumerate}
\item source\_statement: LB\_N(1+epsilon) is C(H\_N\textasciicircum{}SAT)>N\textasciicircum{}\{1+epsilon-o(1)\} for the beta-zero parity carrier over valid N, with nonuniform circuits and no advice beyond circuit descriptions.
\item target\_statement: accepted \(\mathrm{Psys}_{bEF}^{-}+IU\) proofs with subthreshold Gamma\_bEF\textasciicircum{}IU(pi) imply an explicit circuit for \(H_{N}^\mathrm{SAT}\) below the claimed lower bound, so they are falsified as certificates of \(\mathrm{LB}_{N}\).
\item For each EF extension line \(e:=a\) op b, store D(e)=D(a) union D(b) plus the defining gate bit fields (line id, opcode, parent pointers, output bit); transitive DAG traversal charges S\_dep(pi) once per reached occurrence dependency.
\item For each SATeq link, store the two formula pointers, the equality-witness or syntactic rewrite code, and the affected row indices; quotienting charges S\_link(pi) but does not create free hidden rows.
\item For each \(bounded-SR_{h}\) rule, use verifier fields (rule tag, parent line, row index j in [M\_N], node path u of length \(\le h\), branch variable id, child formula pointers, leaf solver id, touched occurrence-neighborhood list, local accept bit). The local verifier recomputes the child clauses and checks that every touched bit is in the listed neighborhood.
\item The \(SR_{h}\) line induction shows that replacing every SR rule by its leaf-neighborhood transcript preserves the line's Boolean value on all rows already in \(K_{idx}\); the total charge is \(B_{SR}(pi)\le N^{q+\eta+\tau+o(1)}\) because there are \(N^{q+o(1)}\) rows, \(N^{\eta+o(1)}\) leaves per row, and \(N^{\tau+o(1)}\) neighborhood cost per leaf.
\item For each interval-union FormulaEval line, \(RowSetEnumerator_{N}\) expands endpoints into explicit row indices; the parser charges N per enumerated row, giving \(N|RowSet_{FE}(pi)|\) and preventing free symbolic occurrence completion.
\item Residual disagreement on \(E_{N}\) rows is corrected by equality gates comparing x to each bad source row, costing \(E_{NN}\); this is valid only for \(E_{N}<N^{\epsilon-o(1)}\) in the falsifier regime.
\item Composing the generator circuit of size \(N^{\rho+o(1)}\) with the parsed transcript yields C(H\_N\textasciicircum{}SAT)<=N\textasciicircum{}\{rho+o(1)\}+Gamma\_bEF\textasciicircum{}IU(pi)+N\textasciicircum{}\{o(1)\}; if this is below \(N^{1+\epsilon-o(1)}\), the source lower-bound sentence is false in the same circuit model.
\item Inference-quality status: source carrier, target proof model, reduction type (parser-to-circuit simulation), length regime, exponent regime, advice status, and barrier scope are explicit; pvsnp-symbolic-complexity-auditor should only recheck endpoint enumeration and eta/tau accounting.
\end{enumerate}

\paragraph{Equation 1: carrier.}
\[
H_N^{\mathrm{SAT}}(x)=\bigoplus_{j< M_N}\mathrm{SAT}(G_N(x,j)),\quad M_N=N^{q+o(1)},\quad |G_N(x,j)|=(\log N)^d
\]
\emph{Valid when:} Valid when the source carrier is the beta-zero generated parity problem with polylogarithmic SAT instances and explicit generator bits.

\paragraph{Equation 2: transcript\_cost.}
\[
\Gamma_{bEF}^{IU}(\pi)=S_{dep}(\pi)+S_{link}(\pi)+B_{SR}(\pi)+N\,|RowSet_{FE}(\pi)|+E_NN
\]
\emph{Valid when:} Applies after every EF, SATeq, \(bounded-SR_{h}\), interval-union FormulaEval, and residual-error component has been parsed explicitly.

\paragraph{Equation 3: \(bounded_{SR,charge}\).}
\[
B_{SR}(\pi)\le N^{q+\eta+\tau+o(1)}
\]
\emph{Valid when:} Requires the bounded self-reduction verifier to expose row indices, node paths, and touched occurrence neighborhoods with \(\eta\) and \(\tau\) overheads.

\paragraph{Equation 4: parsed\_upper\_bound.}
\[
C(H_N^{\mathrm{SAT}})\le N^{\rho+o(1)}+\Gamma_{bEF}^{IU}(\pi)+N^{o(1)}
\]
\emph{Valid when:} Applies after the parsed transcript and the generator circuit are composed without nonuniform leakage or semantic FormulaEval advice.

\paragraph{Equation 5: \(\mathrm{LB}_{falsifier}\).}
\[
N^{\rho+o(1)}+\Gamma_{bEF}^{IU}(\pi)< N^{1+\epsilon-o(1)}\ \Longrightarrow\ \neg LB_N(1+\epsilon)
\]
\emph{Valid when:} Valid when \(\mathrm{LB}_{N}\) means the stated nonuniform circuit lower bound for the same carrier, lengths, and gate model.

\paragraph{Next analysis.} Run the delta-varied interval-union selector audit: instantiate \(|RowSet|=N^{q-\delta}\), verify \(RowSetEnumerator_{N}\) gate cost, compute \(\Gamma_{bEF}^{IU}\) at \(\delta=q-\epsilon\) \(\pm\) o(1), and test whether any proposed semantic selector supplies a valid \(FE_{Extract,N}\) certificate rather than triggering the private-row falsifier.

\paragraph{Falsifiers.}
\begin{itemize}
\item Subthreshold transcript falsifier: N\textasciicircum{}\{rho+o(1)\}+Gamma\_bEF\textasciicircum{}IU(pi)<N\textasciicircum{}\{1+epsilon-o(1)\} gives an explicit circuit contradicting LB\_N(1+epsilon).
\item Private-row falsifier: semantic FormulaEval with an unenumerated selected row admits two oracle labelings agreeing on the transcript but flipping the parity output.
\item \(q<\epsilon\) enumeration falsifier: global row-expanded FormulaEval costs \(N^{1+q+o(1)}\), below the lower-bound threshold when other terms are subthreshold.
\item Nonuniform leakage falsifier: if \(\eta\) or \(\tau\) is supplied by advice rather than verifier-readable bit fields, \(\Gamma_{bEF}\) undercounts the transcript and the theorem is invalid.
\item Positive-beta transfer falsifier: any black-box target circuit substitution with query length \(n^\beta\) yields exponent q+beta c+o(1), so all-exponent NP versus P/poly consequences are capped unless \(\beta=0\) or subpolynomial.
\end{itemize}

\subsection*{cycle-0106: Hardness-Magnification Bridge Realization}

\paragraph{Execution.} research; azure\_responses\_api; attempt primary; 600.004 seconds; 3 lemmas, 6 equations, 9 derivation steps, 6 falsifiers.

\paragraph{Learning.} Cycle 106 upgrades the promoted restricted result into a live inference-quality theorem packet: \(\Gamma_{bEF}^{IU}\) is derived by line induction with explicit \(SR_{h}\) and interval-union FormulaEval charges, the q/epsilon support trichotomy is a contraposed corollary, and full semantic FormulaEval remains quarantined by the private-row falsifier rather than by generic barrier folklore.

\paragraph{Pressure.} No blocking gap remains for \(\mathrm{Psys}_{bEF}^{-}\) plus transparent interval-union FormulaEval. The named unresolved gap is \(FE_{Extract,N}^{sem}\): extracting complete row supports from unexpanded semantic FormulaEval-SAT without advice or nonrelativizing SAT-specific reasoning.

\paragraph{Derivation steps.}
\begin{enumerate}
\item Define proof-line fields: opcode; line number; parent pointers; EF extension gate id; dependency bit interval; SATeq link id; \(SR_{h}\) leaf tuple (slot j, row r, variable block b, branch word z, occurrence offset u); FormulaEval selector id with interval-union endpoints; residual patch list.
\item Run transitive EF-DAG traversal: for every extension variable v, accumulate the union of source row and occurrence dependencies of its defining subformula; charge each newly exposed bit interval to S\_dep(pi) once by topological order.
\item Process SATeq lines by quotienting explicitly named formula encodings; charge every equality link and copied occurrence neighborhood to S\_link(pi); do not infer unlisted rows from semantic equality.
\item Process \(bounded-SR_{h}\) lines locally: each touched slot contributes at most \(N^{\eta+\tau+o(1)}\) verifier leaves after shared generator decoding, so the total SR charge over \(M_{N}\) slots is \(B_{SR,h}(pi)\le N^{q+\eta+\tau+o(1)}\).
\item Process interval-union FormulaEval by expanding each declared interval endpoint list into RowSet\_IU(pi); each selected row costs N for equality-test patching or row circuit hardwiring, giving the \(N|RowSet_{IU}(pi)|\) term.
\item Induct on proof lines that the semantic value used by each accepted restricted rule is determined by the accumulated charged transcript plus at most \(E_{N}\) residual rows; the induction fails deliberately for unexpanded semantic FormulaEval.
\item Compose the generator circuit of size \(N^{\rho+o(1)}\) with the charged transcript circuit and equality patches to obtain C(H\_N\textasciicircum{}SAT)<=N\textasciicircum{}\{rho+o(1)\}+O(Gamma\_bEF\textasciicircum{}IU(pi)).
\item Contrapose against LB\_N(1+epsilon): any accepted restricted certificate of that lower bound must have N\textasciicircum{}\{rho+o(1)\}+Gamma\_bEF\textasciicircum{}IU(pi)>=N\textasciicircum{}\{1+epsilon-o(1)\}; otherwise the certificate is invalid because the constructed circuit is a falsifier.
\item Run inference-quality checks: source and target statements are identical carrier/model statements; quantifiers are fixed-exponent and valid-length; advice is charged or confined to final nonuniform circuits; full semantic FormulaEval is blocked at \(FE_{Extract,N}^{sem}\).
\end{enumerate}

\paragraph{Equation 1: carrier.}
\[
H_N^{SAT}(x)=\bigoplus_{j=1}^{M_N} SAT(G_N(x,j)),\quad M_N=N^{q+o(1)},\quad |G_N|=N^{\rho+o(1)},\quad L_N=(\log N)^d
\]
\emph{Valid when:} Valid when the beta-zero carrier generator and all valid lengths are fixed before the proof family is parsed.

\paragraph{Equation 2: \(bounded_{SR,charge}\).}
\[
B_{SR_h}(\pi)\le N^{q+\eta+\tau+o(1)}
\]
\emph{Valid when:} Applies when \(bounded-SR_{h}\) touches \(M_{N}\) slots, has leaf exponent \(\eta\), and uses verifier overhead exponent \(\tau\).

\paragraph{Equation 3: charged\_transcript.}
\[
\Gamma_{bEF}^{IU}(\pi)=S_{dep}(\pi)+S_{link}(\pi)+B_{SR_h}(\pi)+N\cdot |RowSet_{IU}(\pi)|+E_NN
\]
\emph{Valid when:} Valid when EF dependencies, SATeq links, SR leaves, interval-union rows, and residual patches are separately enumerated and charged.

\paragraph{Equation 4: parser\_to\_circuit.}
\[
C(H_N^{SAT})\le N^{\rho+o(1)}+O(\Gamma_{bEF}^{IU}(\pi))
\]
\emph{Valid when:} Requires the parsed transcript to determine every generated parity row except the explicitly patched \(E_{N}\) residual rows.

\paragraph{Equation 5: \(\mathrm{LB}_{falsifier}\).}
\[
N^{\rho+o(1)}+\Gamma_{bEF}^{IU}(\pi)< N^{1+\epsilon-o(1)}\;\Longrightarrow\; \neg LB_N(1+\epsilon)
\]
\emph{Valid when:} Valid when \(\mathrm{LB}_{N}\) means a nonuniform SIZE lower bound for exactly the same carrier \(H_{N}^{\mathrm{SAT}}\).

\paragraph{Equation 6: support\_density.}
\[
|K_{idx}(\pi)|\ge N^{\epsilon-o(1)},\quad |K_{idx}(\pi)|/M_N\ge N^{\epsilon-q-o(1)}
\]
\emph{Valid when:} Applies to sound restricted certificates after the parser-to-circuit bound is contraposed below the lower-bound threshold.

\paragraph{Next analysis.} Run the finite validation audit on one concrete \(\mathrm{Psys}_{sym}\) rule: encode FormulaEval as (a) interval-union axiom schema, (b) \(\mathrm{AC0}\) selector with enumerator, and (c) opaque semantic oracle; compute \(\alpha\) and either attach \(N^{1+\alpha+o(1)}\) charge or exhibit the private-row countermodel.

\paragraph{Falsifiers.}
\begin{itemize}
\item Subthreshold Gamma\_bEF\textasciicircum{}IU(pi) gives an explicit circuit upper-bound falsifier for LB\_N(1+epsilon).
\item A hidden semantic FormulaEval row supplies the private-row falsifier to surface occurrence completion.
\item \(q<\epsilon\) with fewer than \(N^{\epsilon-o(1)}\) complete rows falsifies any claimed restricted source certificate by row enumeration and equality patching.
\item Nonuniform leakage through uncharged EF extension variables falsifies the parser unless the transitive dependency DAG is bit-explicit.
\item Positive \(\beta\) in any black-box target-circuit substitution revives the transfer cap q+beta c+o(1), so all-exponent NP/Ppoly claims are invalid outside the beta-zero window.
\item Treating EqSAT or \(SR_{h}\) as free semantic power is invalid unless the quotient links or SR leaves are charged in \(\Gamma_{bEF}^{IU}\).
\end{itemize}

\subsection*{cycle-0107: Hardness-Magnification Bridge Realization}

\paragraph{Execution.} research; azure\_responses\_api; attempt primary; 600.009 seconds; 3 lemmas, 6 equations, 6 derivation steps, 5 falsifiers.

\paragraph{Learning.} Cycle 107 adds a real validation upgrade: transparent FormulaEval selectors no longer require interval-union structure for a valid theorem, because full syntactic scanning gives a certified \(B_{FE}=F_{FE}\) \(N^{1+q+o(1)}\) charge. This closes the \(q<\epsilon\) falsifier for a broader selector fragment and isolates the only missing improvement as \(AC0PrefixNonemptiness_{N}\) for output-sensitive enumeration. Full semantic FormulaEval remains quarantined by the private-row falsifier, not by generic barrier folklore.

\paragraph{Pressure.} \(AC0PrefixNonemptiness_{N}\): no advice-free, output-sensitive enumeration theorem is proved for arbitrary \(\mathrm{AC0}\) row selectors over q log N row bits; brute force gives a valid but sometimes over-threshold scan charge.

\paragraph{Derivation steps.}
\begin{enumerate}
\item Defined \(\mathrm{Psys}_{sym}^{sel-scan}\): each FormulaEval-SAT line contains a selector descriptor a, a syntactic evaluator Sel\_a(j), and either an enumerator certificate or permission to full-scan \(j<M_{N}\).
\item For an output-sensitive enumerator, replaced the FormulaEval line by one occurrence expansion per j in RowSet(a); the charge is \(N^{1+o(1)}\) per complete row because the row-neighborhood transcript must be usable in the \(H_{N}^\mathrm{SAT}\) circuit simulation.
\item For selectors without an enumerator, scanned all \(M_{N}=N^{q+o(1)}\) row indices, evaluated Sel\_a(j) syntactically, and expanded every accepted row. This proves \(B_{FE}^{scan}=F_{FE}\) \(N^{1+q+o(1)}\) for \(F_{FE}=N^{o(1)}\).
\item Inserted \(B_{FE}^{sel}\) into the previously promoted line-induction proof for \(\mathrm{Psys}_{bEF}^{-}\): EF dependency, SATeq link, \(bounded-SR_{h}\), FormulaEval, and sparse-patch charges are additive in \(\Gamma_{bEF}^{sel}\).
\item Derived the live inference-quality source statement: a \(\mathrm{Psys}_{sym}^{sel-scan}\) proof of LB\_N(1+epsilon) with subthreshold \(\Gamma_{bEF}^{sel}\) yields a circuit for \(H_{N}^\mathrm{SAT}\) below \(N^{1+\epsilon-o(1)}\), contradicting the source statement in the same circuit model.
\item Tried to upgrade arbitrary \(\mathrm{AC0}\) selectors to output-sensitive \(RowSetEnumerator_{N}\) by recursive prefix search; the missing step is uniform \(AC0PrefixNonemptiness_{N}\) with \(N^{o(1)}\) delay, so only the full-scan fallback is currently proved.
\end{enumerate}

\paragraph{Equation 1: carrier.}
\[
H_N^{SAT}(x)=\bigoplus_{j< M_N} SAT(G_N(x,j)),\quad M_N=N^{q+o(1)},\quad |G_N(x,j)|=(\log N)^d.
\]
\emph{Valid when:} Valid when the beta-zero carrier has polylogarithmic SAT queries, explicit row index encoding, and generator size \(N^{\rho+o(1)}\).

\paragraph{Equation 2: \(output_{sensitive,FE,charge}\).}
\[
B_{FE}^{out}(\pi)=N^{1+o(1)}\sum_{\ell\in FE(\pi)} |RowSet(a_\ell)|.
\]
\emph{Valid when:} Valid when \(RowSetEnumerator_{N}\) is uniform, advice-free, and lists exactly the selected rows with subpolynomial per-output overhead.

\paragraph{Equation 3: \(full_{scan,FE,charge}\).}
\[
B_{FE}^{scan}(\pi)=F_{FE}(\pi)\,N\,M_N\,N^{o(1)}=F_{FE}(\pi)N^{1+q+o(1)}.
\]
\emph{Valid when:} Valid when selectors are decidable from proof bits and row indices but no output-sensitive enumeration certificate is available.

\paragraph{Equation 4: selector\_gamma.}
\[
\Gamma_{bEF}^{sel}(\pi)=S_{dep}(\pi)+S_{link}(\pi)+N^{q+\eta+\tau+o(1)}+B_{FE}^{sel}(\pi)+E_NN.
\]
\emph{Valid when:} Applies after EF dependencies, SATeq links, \(bounded-SR_{h}\) neighborhoods, FormulaEval expansions, and sparse row patches are all charged.

\paragraph{Equation 5: parser\_falsifier.}
\[
N^{\rho+o(1)}+\Gamma_{bEF}^{sel}(\pi)< N^{1+\epsilon-o(1)}\;\Longrightarrow\;\neg LB_N(1+\epsilon).
\]
\emph{Valid when:} Requires the target-preserving parser-to-circuit simulation in the same nonuniform circuit model as the asserted source lower bound.

\paragraph{Equation 6: support\_trichotomy.}
\[
|K_{idx}(\pi)|\ge N^{\epsilon-o(1)},\quad |K_{idx}(\pi)|/M_N\ge N^{\epsilon-q-o(1)}.
\]
\emph{Valid when:} Valid for any parsed proof that avoids the subthreshold falsifier and charges each complete contacted row by \(N^{1-o(1)}\).

\paragraph{Next analysis.} Formalize \(AC0PrefixNonemptiness_{N}\) as a finite lemma: for \(\mathrm{AC0}\) selectors on q log N row bits, test whether descriptor structure yields \(N^{o(1)}\) prefix emptiness; if not, state the selector-scan theorem as final and add \(\mathrm{AC0}\) output-sensitive enumeration to open problems.

\paragraph{Falsifiers.}
\begin{itemize}
\item semantic FormulaEval enable bit depending on a private SAT row, which breaks occurrence completion unless charged by \(FE_{Extract,N}^{sem}\)
\item \(q<\epsilon\) with \(F_{FE}=N^{o(1)}\) under full-scan extraction, since B\_FE\textasciicircum{}\{scan\}=N\textasciicircum{}\{1+q+o(1)\} is subthreshold and falsifies LB\_N(1+epsilon)
\item nonuniform selector advice that names RowSet(a) without being encoded in pi, invalidating the advice-free parser model
\item \(\mathrm{AC0}\) sparse selector with no certified prefix-nonemptiness procedure, blocking output-sensitive enumeration and forcing the weaker scan charge
\item positive \(\beta\) black-box substitution interface, where q+beta c caps any claimed all-exponent \(\mathrm{NP}\not\subset \mathrm{P}/\mathrm{poly}\) transfer
\end{itemize}

\subsection*{cycle-0108: Hardness-Magnification Bridge Realization}

\paragraph{Execution.} research; azure\_responses\_api; attempt primary; 146.471 seconds; 3 lemmas, 6 equations, 10 derivation steps, 6 falsifiers.

\paragraph{Learning.} Cycle 108 replaces selector-case accumulation by a single \(TransparentRowSupport_{N}\) theorem pair. Transparent FormulaEval is valid only after \(U_{N}+X_{N}\) charging or full scan; hidden semantic FormulaEval is refuted by a private-row falsifier unless \(RealSATInfluenceRealizer_{N}\) supplies a nonrelativizing charged-support certificate. This is a paper-level parser/falsifier boundary, not evidence for \(\mathrm{P} \ne \mathrm{NP}\).

\paragraph{Pressure.} The positive theorem is valid for \(TransparentRowSupport_{N}\) and full-scan fallback. The full \(\mathrm{Psys}_{sym}\) semantic theorem remains blocked exactly at \(RealSATInfluenceRealizer_{N}\); \(AC0PrefixNonemptiness_{N}\) is only an optimization gap for output-sensitive enumeration.

\paragraph{Derivation steps.}
\begin{enumerate}
\item Fixed the beta-zero carrier \(H_{N}^\mathrm{SAT}(x)=xor_{j\le M_{N}}\) SAT(G\_N(x,j)) with \(M_{N}=N^{q+o(1)}\), ell\_N=(log N)\textasciicircum{}d, and generator size \(N^{\rho+o(1)}\).
\item Defined \(TransparentRowSupport_{N}\) as five clauses: visible selector code, advice-free enumeration or full scan, influence closure, same-theory occurrence expansion, and no hidden support advice.
\item Inherited the cycle-102 \(\mathrm{Psys}_{bEF}^{-}\) line induction: EF dependencies, SATeq links, and bounded \(SR_{h}\) contribute Gamma\_bEF(pi), including eta,tau and touched occurrence neighborhoods.
\item For a transparent FormulaEval line a, enumerate K\_idx(a) or scan all \(M_{N}\) slots; expand every selected row into occurrence-level verifier fields; charge \(U_{a}\) plus \(N|K_{idx}(a)|ell_{N}^{O(1)}\).
\item Summed FormulaEval charges to obtain Gamma\_trans(pi)=Gamma\_bEF(pi)+U\_N(pi)+X\_N(pi), with \(X_{N}\le F_{FE}\) \(N^{1+q+o(1)}\) under full scan.
\item Extended the line induction: every later proof line depends only on generated formulas, charged row labels, and parsed transcript gates; influence closure forbids uncharged semantic dependence.
\item Constructed the falsifying circuit \(D_{pi}\) for \(H_{N}^\mathrm{SAT}\) by combining \(G_{N}\), transcript gates of size Gamma\_trans(pi), and equality patches for \(E_{N}\) residual source rows, giving \(C(H_{N})\le N^{\rho+o(1)}+\Gamma_{trans}(pi)+E_{NN}\) polylog(N).
\item Contraposed LB\_N(1+epsilon): if the displayed upper bound is below \(N^{1+\epsilon-o(1)}\), the purported proof certificate is invalid for the claimed source lower bound.
\item Proved the hidden-support no-go at row-label level: if r* affects FormulaEval but is not in \(K_{idx}\) and not influence-forced, two label assignments agree on the whole visible transcript while differing on r*; an unread enable bit can make proof acceptance fixed while carrier parity flips.
\item Separated the open semantic escape: real SAT may invalidate arbitrary label flips, so full semantic FormulaEval needs \(RealSATInfluenceRealizer_{N}\) rather than another occurrence-parser assertion.
\end{enumerate}

\paragraph{Equation 1: beta\_zero\_carrier\_parameters.}
\[
\ell_N=(\log N)^d,\quad M_N=N^{q+o(1)},\quad |G_N|=N^{\rho+o(1)}
\]
\emph{Valid when:} Valid when the carrier generator outputs polylogarithmic SAT instances and q,rho are fixed exponent parameters.

\paragraph{Equation 2: transparent\_gamma.}
\[
\Gamma_{trans}(\pi)=\Gamma_{bEF}(\pi)+U_N(\pi)+X_N(\pi)
\]
\emph{Valid when:} Applies after every FormulaEval line satisfies \(TransparentRowSupport_{N}\) or is charged by full syntactic scanning.

\paragraph{Equation 3: formulaeval\_expansion\_charge.}
\[
X_N(\pi)=\sum_{a\in FE(\pi)} N\,|K_{idx}(a)|\,\ell_N^{O(1)}\le F_{FE}(\pi)N^{1+q+o(1)}
\]
\emph{Valid when:} Valid when K\_idx(a) is advice-free enumerable and the fallback scan ranges over all \(M_{N}\) slots.

\paragraph{Equation 4: parsed\_upper\_bound.}
\[
C_{SIZE}(H_N^{SAT})\le N^{\rho+o(1)}+\Gamma_{trans}(\pi)+E_NN\ell_N^{O(1)}
\]
\emph{Valid when:} Requires the same nonuniform SIZE model, a parsed proof transcript, and only \(E_{N}\) unparsed source rows.

\paragraph{Equation 5: subthreshold\_falsifier.}
\[
N^{\rho+o(1)}+\Gamma_{trans}(\pi)+E_NN\ell_N^{O(1)}< N^{1+\epsilon-o(1)}\Rightarrow \neg LB_N(1+\epsilon)
\]
\emph{Valid when:} Applies when \(\mathrm{LB}_{N}\) asserts an \(N^{1+\epsilon-o(1)}\) lower bound for this exact carrier.

\paragraph{Equation 6: support\_density\_corollary.}
\[
|K_{idx}^{tot}|\ge N^{\epsilon-o(1)},\quad |K_{idx}^{tot}|/M_N\ge N^{\epsilon-q-o(1)}
\]
\emph{Valid when:} Valid when FormulaEval support is the only unabsorbed term preventing the parser-circuit falsifier inequality.

\paragraph{Next analysis.} Write the referee-style \(TransparentRowSupport_{N}\) definition and audit six examples as corollaries: interval-union, full scan, ROBDD, guarded-dDNNF, projected-dDNNF, and influence-complete extension rules. Do not add new selector families unless they prove a clause, refute a clause, or improve \(AC0PrefixNonemptiness_{N}\).

\paragraph{Falsifiers.}
\begin{itemize}
\item Subthreshold \(\Gamma_{trans}\) falsifier: \(N^{\rho+o(1)}+\Gamma_{trans}(pi)+E_{NN}\) \(polylog(N)<N^{1+\epsilon-o(1)}\) contradicts LB\_N(1+epsilon) for this carrier.
\item Support-density falsifier: if total FormulaEval support |K\_idx\textasciicircum{}\{tot\}|<N\textasciicircum{}\{epsilon-o(1)\}, row expansion cannot support a valid \(\mathrm{LB}_{N}\) certificate at exponent 1+epsilon.
\item \(q=N^{\epsilon-q-o(1)}\) exceeds density one, so subthreshold transparent FormulaEval certificates cannot exist.
\item Hidden-support falsifier: a private row r* or unread enable bit makes two visible transcripts agree while the carrier parity differs.
\item Nonuniform leakage falsifier: if row names or selector witnesses are supplied as semantic advice instead of parsed code, \(TransparentRowSupport_{N}\) is violated.
\item Positive-beta transfer falsifier remains active for any broadened carrier: q+beta c caps all-exponent transfer unless \(\beta=0\) or the use is genuinely nonblackbox.
\end{itemize}

\subsection*{cycle-0109: Hardness-Magnification Bridge Realization: \(TransparentRowSupport\) theorem pair}

\paragraph{Execution.} research; azure\_responses\_api; attempt primary; 138.79 seconds; 4 lemmas, 5 equations, 9 derivation steps, 6 falsifiers.

\paragraph{Learning.} Cycle 109 upgrades the theorem pair from selector consolidation to a proof-audited interface: \(TransparentRowSupport_{N}\) is a finite checkable condition, \(\Gamma_{trans}\) is derived by line induction, hidden semantic support has a concrete private-row/unread-enable-bit falsifier, and the only full \(\mathrm{Psys}_{sym}\) escape is the nonrelativizing \(RealSATInfluenceRealizer_{N}\). This is paper-level route certification, not evidence for \(\mathrm{P} \ne \mathrm{NP}\).

\paragraph{Pressure.} \(RealSATInfluenceRealizer_{N}\) remains open for genuine real-SAT semantic FormulaEval rules; the private-row countermodel proves that a relativizing or surface-occurrence parser cannot fill this gap.

\paragraph{Derivation steps.}
\begin{enumerate}
\item Fix \(H_{N}^\mathrm{SAT}\), \(M_{N}=N^{q+o(1)}\), L\_N=(log N)\textasciicircum{}d, generator size \(N^{\rho+o(1)}\), and the source lower-bound statement LB\_N(1+epsilon) in nonuniform SIZE.
\item Define \(TransparentRowSupport_{N}\) as five checkable clauses: visible selector code, advice-free row enumeration or certified full scan, influence closure, same-theory occurrence expansion, and no hidden semantic advice.
\item Run the existing \(\mathrm{Psys}_{bEF}^{-}\) line induction for bit-explicit EF, SATeq links, \(bounded-SR_{h}\) verifier fields, and sparse-row equality patches; this yields Gamma\_bEF(pi).
\item For each transparent FormulaEval line, append either \(U_{ell}\) plus \(N^{1+o(1)}|K_{ell}|\) for enumerated rows or the full-scan bound \(N^{1+q+o(1)}\); influence closure justifies replacing semantic use by explicit occurrence rows.
\item Union row-index charges across FormulaEval lines and obtain \(\Gamma_{trans}(pi)=\Gamma_{bEF}(pi)+U_{N}+X_{N}\) without adding a new selector-specific theorem.
\item Compose the generator circuit, charged transcript circuits, expanded FormulaEval rows, and \(E_{N}\) equality patches to build a SIZE circuit for \(H_{N}^\mathrm{SAT}\) of size N\textasciicircum{}\{rho+o(1)\}+Gamma\_trans(pi)+E\_NN\textasciicircum{}\{1+o(1)\}.
\item Contrapose the construction: if that size is below \(N^{1+\epsilon-o(1)}\), any accepted proof/certificate of LB\_N(1+epsilon) has a valid small-circuit falsifier.
\item For hidden FormulaEval, construct two oracle-label states agreeing on all charged rows but differing on a private row or unread enable bit; the same surface transcript cannot be a valid occurrence parser in both states.
\item Therefore transparent FormulaEval is promoted as a theorem pair, while full semantic FormulaEval is blocked exactly at \(RealSATInfluenceRealizer_{N}\).
\end{enumerate}

\paragraph{Equation 1: carrier.}
\[
H_N^{SAT}(x)=\bigoplus_{j< M_N} SAT(G_N(x,j)),\quad M_N=N^{q+o(1)},\quad |G_N(x,j)|=(\log N)^d
\]
\emph{Valid when:} Valid when the source carrier is the fixed beta-zero batched SAT parity family with uniformly generated polylogarithmic formulas.

\paragraph{Equation 2: transparent-charge.}
\[
\Gamma_{trans}(\pi)=\Gamma_{bEF}(\pi)+U_N(\pi)+X_N(\pi),\quad X_N=N^{1+o(1)}\sum_{\ell\in FE(\pi)} |K_{idx}(\ell)|
\]
\emph{Valid when:} Valid when every FormulaEval line satisfies \(TransparentRowSupport_{N}\) and each charged row expansion is same-theory occurrence explicit.

\paragraph{Equation 3: full-scan-bound.}
\[
X_N(\pi)\le F_{FE}(\pi)\,N^{1+q+o(1)}\quad\text{for certified full-scan FormulaEval lines}
\]
\emph{Valid when:} Applies when FormulaEval support is transparent but no output-sensitive enumerator is used, so all \(M_{N}\) rows are syntactically scanned.

\paragraph{Equation 4: parser-falsifier.}
\[
C_{SIZE}(H_N^{SAT})\le N^{\rho+o(1)}+\Gamma_{trans}(\pi)+E_NN^{1+o(1)}
\]
\emph{Valid when:} Requires the target-preserving parser simulation, fixed generator circuits, charged FormulaEval support, and explicit equality patches for residual rows.

\paragraph{Equation 5: subthreshold-contradiction.}
\[
N^{\rho+o(1)}+\Gamma_{trans}(\pi)+E_NN^{1+o(1)}< N^{1+\epsilon-o(1)}\Rightarrow \neg LB_N(1+\epsilon)
\]
\emph{Valid when:} Valid when \(\mathrm{LB}_{N}\) is stated in the same nonuniform SIZE model and the constructed circuit exactly computes the source carrier.

\paragraph{Next analysis.} Referee-check one projected-dDNNF FormulaEval rule against the five \(TransparentRowSupport_{N}\) clauses, compute its \(U_{N}\) and \(X_{N}\) charges, and verify that any projected hidden variable is either influence-closed or triggers the unread-enable-bit countermodel.

\paragraph{Falsifiers.}
\begin{itemize}
\item Subthreshold Gamma\_trans(pi) is a valid falsifier for LB\_N(1+epsilon) in the same SIZE model.
\item Private-row hidden support is a falsifier for uncharged FormulaEval occurrence completion.
\item Unread-enable-bit semantic access is a falsifier for selector extensionality and influence closure.
\item For \(q<\epsilon\), certified full scan costs only \(N^{1+q+o(1)}\) and becomes a quantitative falsifier when other charges are subthreshold.
\item Any \(\tau/\eta/SR_{h}\) overhead hidden in semantic advice invalidates the charged transcript theorem.
\item Positive \(\beta\) black-box substitution remains quarantined by the old q+beta c threshold-collapse falsifier.
\end{itemize}

\subsection*{cycle-0110: Hardness-Magnification Bridge Realization}

\paragraph{Execution.} research; azure\_responses\_api; attempt primary; 126.879 seconds; 4 lemmas, 4 equations, 7 derivation steps, 6 falsifiers.

\paragraph{Learning.} Cycle 110 consolidates the selector zoo into a single theorem pair. Transparent FormulaEval is valid only when every influence is surfaced and charged in \(\Gamma_{trans}\); hidden semantic support is blocked by a private-row/unread-enable-bit falsifier. The exact full-system escape is \(RealSATInfluenceRealizer_{N}\), not a generic naturalness or relativization slogan.

\paragraph{Pressure.} No remaining gap for the restricted \(TransparentRowSupport_{N}\) theorem pair. The named open gap for full \(\mathrm{Psys}_{sym}\) is \(RealSATInfluenceRealizer_{N}\): convert real SAT-semantic FormulaEval influence into explicit charged row support without advice leakage or threshold collapse.

\paragraph{Derivation steps.}
\begin{enumerate}
\item Fixed the source carrier \(H_{N}^\mathrm{SAT}\) with \(\beta=0\) polylog SAT rows, \(M_{N}=N^{q+o(1)}\), generator size \(N^{\rho+o(1)}\), and nonuniform SIZE as both source and target circuit model.
\item Defined \(TransparentRowSupport_{N}\) as a finite interface: visible selector code, advice-free RowSet enumeration or full scan, same-theory occurrence expansion, influence closure, row-index encodings, and no hidden support advice.
\item Extended the existing \(\mathrm{Psys}_{bEF}^{-}\) line induction: bEF, SATeq, and bounded \(SR_{h}\) contribute Gamma\_bEF(pi); each transparent FormulaEval line contributes \(EnumCost(K_{l})+N|K_{l}|\) or the full-scan charge; residual bad rows cost O(E\_N N).
\item Constructed the target circuit by composing \(G_{N}\), charged transcript gates, parsed \(SR_{h}\) leaf solvers, expanded FormulaEval rows, and equality patches; this yields C\_SIZE(H\_N\textasciicircum{}SAT)<=N\textasciicircum{}\{rho+o(1)\}+Gamma\_trans(pi).
\item Contraposed the bound: if the right side is below \(N^{1+\epsilon-o(1)}\), any proof pi claiming LB\_N(1+epsilon) is invalid in the same circuit model; this supplies the route-level falsifier.
\item Proved the hidden-support no-go in the occurrence abstraction: choose a private row or unread enable bit outside \(K_{pi}\), flip its label while keeping all charged rows and visible code fixed, and obtain identical parser transcript but changed FormulaEval value or carrier parity.
\item Separated the full semantic case: real SAT may forbid independent flips, so full \(\mathrm{Psys}_{sym}\) is not refuted; it is blocked exactly at \(RealSATInfluenceRealizer_{N}\), which must produce charged real-influence supports.
\end{enumerate}

\paragraph{Equation 1: carrier\_parameters.}
\[
H_N^{SAT}(x)=xor_{j\le M_N} SAT(G_N(x,j)), M_N=N^{q+o(1)}, |G_N|=N^{rho+o(1)}, |G_N(x,j)|=(log N)^d
\]
\emph{Valid when:} Valid when N ranges over accepted carrier lengths and the beta-zero generator emits polylogarithmic SAT rows uniformly.

\paragraph{Equation 2: transparent\_charge.}
\[
Gamma_trans(pi)=Gamma_bEF(pi)+U_FE(pi)+X_FE(pi)+O(E_N N)
\]
\emph{Valid when:} Valid when every FormulaEval influence is visibly enumerated or fully scanned and sparse residual rows are equality-patched.

\paragraph{Equation 3: parser\_upper\_bound.}
\[
C_SIZE(H_N^{SAT}) \le N^{rho+o(1)}+Gamma_trans(pi)
\]
\emph{Valid when:} Applies after the line induction converts all proof dependencies into charged circuit gates in the same SIZE model.

\paragraph{Equation 4: subthreshold\_falsifier.}
\[
N^{rho+o(1)}+Gamma_bEF(pi)+U_FE(pi)+X_FE(pi)+O(E_N N) < N^{1+epsilon-o(1)} => not LB_N(1+epsilon)
\]
\emph{Valid when:} Requires the source lower-bound claim to be about \(H_{N}^{\mathrm{SAT}}\) in nonuniform SIZE with identical exponent cutoff.

\paragraph{Next analysis.} Attempt \(RealSATInfluenceRealizer_{N}\) for one influence-complete extension rule: define \(RealInfl_{l}\), produce row-index witnesses and \(EnumCost+N|K_{l}|\) charge, or exhibit a semantic rule where a real SAT-sensitive row remains unchargeable, thereby demoting full \(\mathrm{Psys}_{sym}\) beyond the transparent theorem pair.

\paragraph{Falsifiers.}
\begin{itemize}
\item Subthreshold transparent transcript: N\textasciicircum{}\{rho+o(1)\}+Gamma\_trans(pi)<N\textasciicircum{}\{1+epsilon-o(1)\} gives a valid same-model circuit falsifier for LB\_N(1+epsilon).
\item Private-row falsifier: an uncharged row outside \(K_{pi}\) changes FormulaEval or parity while all visible parsed data stay fixed.
\item Unread-enable-bit falsifier: selector extensionality depends on a hidden semantic bit, so RowSet enumeration is not advice-free.
\item \(q<\epsilon\) support falsifier: fewer than \(N^{\epsilon-o(1)}\) touched rows yield an O(N|K\_idx|) patch below the claimed threshold.
\item Positive-beta transfer falsifier persists for other carriers: q+beta*c caps all-exponent SAT exclusion unless \(\beta=0\) or the use is genuinely nonblackbox.
\item Nonuniform leakage falsifier: a proof may encode \(K_{idx}\) as hidden advice unless its description is charged in \(U_{FE}\).
\end{itemize}

\subsection*{cycle-0111: Hardness-Magnification Bridge Realization}

\paragraph{Execution.} research; azure\_responses\_api; attempt primary; 117.315 seconds; 3 lemmas, 4 equations, 11 derivation steps, 6 falsifiers.

\paragraph{Learning.} Cycle 111 converts the transparent FormulaEval consolidation into a proof-delivery packet: \(\Gamma_{trans}\) is derived by line induction, hidden support has a precise private-row falsifier, and the only full semantic escape is the named nonrelativizing \(RealSATInfluenceRealizer_{N}\). This is a valid scoped parser/falsifier contribution, not progress on the main P vs NP conjecture itself.

\paragraph{Pressure.} No gap remains for the conditional \(TransparentRowSupport_{N}\) theorem pair. The full \(\mathrm{Psys}_{sym}\) extension is blocked exactly at \(RealSATInfluenceRealizer_{N}\) for hidden semantic FormulaEval lines.

\paragraph{Derivation steps.}
\begin{enumerate}
\item Source statement: For beta-zero parity carrier \(H_{N}^\mathrm{SAT}\) with \(M_{N}=N^{q+o(1)}\), polylog SAT instances, generator size \(N^{\rho+o(1)}\), and source lower-bound claim LB\_N(1+epsilon), audit proof certificates in \(\mathrm{Psys}_{sym}\).
\item Target statement: Prove a restricted parser/falsifier theorem for transparent FormulaEval and a countermodel theorem for hidden FormulaEval; do not infer \(\mathrm{P} \ne \mathrm{NP}\) or \(\mathrm{NP}\not\subset \mathrm{P}/\mathrm{poly}\).
\item Model context: nonuniform SIZE circuits for \(H_{N}\), bit-explicit \(\mathrm{Psys}_{bEF}^{-}\) proof lines, FormulaEval extension rules, \(bounded-SR_{h}\) leaves, SATeq links, and same-theory occurrence expansions.
\item Quantifier context: fixed \(\epsilon>0\); \(q,\rho<1+\epsilon\); \(\beta=0\) because query length is polylogarithmic; all o(1) losses uniform on valid lengths; no uncharged advice or semantic row leakage.
\item Line induction base: axioms, bit-explicit EF extensions, SATeq links, and \(bounded-SR_{h}\) rules already have support fields and contribute exactly their \(\Gamma_{bEF}^{-}\) charges from the prior \(SR_{h}\) normal form packet.
\item FormulaEval induction step: \(TransparentRowSupport_{N}\) supplies either an advice-free RowSet enumerator with cost \(U_{a}\) or a full scan over \(N*M_{N}\) rows, then same-theory occurrence expansion costs \(X_{a}\); influence closure ensures no omitted row can affect the transcript.
\item Circuit extraction: combine the generator circuit, all parsed local transcript circuits, and equality patches for \(E_{N}\) residual source rows to obtain size at most \(N^{\rho+o(1)}+\Gamma_{trans}(pi)+E_{N}\) \(N^{1+o(1)}\).
\item Falsifier consequence: if this upper bound is below \(N^{1+\epsilon-o(1)}\), the certificate cannot be valid evidence for LB\_N(1+epsilon) in the same SIZE model.
\item Hidden-support proof: if visibility, selector extensionality, or influence closure fails, choose a private row or unread enable bit; two completions agree on all charged transcript fields but differ on the FormulaEval value or carrier parity, so occurrence completion is invalid.
\item Barrier status: the positive transparent theorem is syntactic and relativizing; the hidden-support escape requires a nonrelativizing real-SAT influence realizer. Natural-proofs and algebrization are not the operative obstruction in this theorem pair.
\item Reviewer handoff: pvsnp-inference-quality-controller checks source/target statements; reduction-consistency checks beta-zero substitution; quantifier-uniformity checks q,rho,epsilon and advice; barrier-specificity checks private-row scope; proof-derivation-builder checks the line induction.
\end{enumerate}

\paragraph{Equation 1: transparent\_charge.}
\[
\Gamma_{\mathrm{trans}}(\pi)=\Gamma_{\mathrm{bEF}^{-}}(\pi)+U_N(\pi)+X_N(\pi)
\]
\emph{Valid when:} Valid when every FormulaEval influence is enumerated or full-scanned and occurrence expansion is performed inside the same parsed proof theory.

\paragraph{Equation 2: formulaeval\_line\_charge.}
\[
B_{FE}(a)\le \min{U_a+X_a,\ F_a N^{1+q+o(1)}}
\]
\emph{Valid when:} Applies when the selector code for line a is visible and the fallback full scan ranges over all N source rows and \(M_{N}=N^{q+o(1)}\) slots.

\paragraph{Equation 3: parser\_upper\_bound.}
\[
C_{SIZE}(H_N^{SAT})\le N^{\rho+o(1)}+\Gamma_{\mathrm{trans}}(\pi)+E_N N^{1+o(1)}
\]
\emph{Valid when:} Valid when the line-induction parser outputs circuits for every unpatched source row and residual rows are equality patched in the same nonuniform SIZE model.

\paragraph{Equation 4: support\_density\_trichotomy.}
\[
|K_{idx}(\pi)|\ge N^{\epsilon-o(1)},\quad |K_{idx}(\pi)|/M_N\ge N^{\epsilon-q-o(1)}
\]
\emph{Valid when:} Requires LB\_N(1+epsilon), beta-zero transfer, and the convention that each surfaced complete parity-slot support costs \(N^{1+o(1)}\) in \(\Gamma_{trans}\).

\paragraph{Next analysis.} Write the final paper theorem text for \(TransparentRowSupport_{N}\) with clauses T1-T6, prove the line induction lemma in referee style, then test one genuinely semantic FormulaEval rule against \(RealSATInfluenceRealizer_{N}\) or instantiate the private-row countermodel.

\paragraph{Falsifiers.}
\begin{itemize}
\item \(q<\epsilon\) support-density falsifier: \(|K_{idx}|/M_{N}\) would exceed 1, so a subthreshold transparent proof cannot certify LB\_N(1+epsilon).
\item Private-row falsifier: a hidden FormulaEval influence absent from RowSet(a) yields two charged-identical transcripts with different induced parity.
\item Unread-enable-bit falsifier: selector code may look empty while an uncharged semantic enable bit activates a row.
\item Threshold-collapse falsifier: N\textasciicircum{}\{rho+o(1)\}+Gamma\_trans+E\_NN\textasciicircum{}\{1+o(1)\} below \(N^{1+\epsilon-o(1)}\) gives a same-model upper bound against the claimed \(\mathrm{LB}_{N}\).
\item Nonuniform-leakage falsifier: \(U_{N}\) or \(X_{N}\) supplied as semantic advice rather than finite visible code invalidates the parser theorem.
\item Model-mismatch falsifier: an occurrence-local proof step is advertised as real-SAT semantic without same-theory expansion or \(RealSATInfluenceRealizer_{N}\).
\end{itemize}

\subsection*{cycle-0112: Hardness-Magnification Bridge Realization}

\paragraph{Execution.} research; azure\_responses\_api; attempt primary; 125.723 seconds; 3 lemmas, 5 equations, 8 derivation steps, 6 falsifiers.

\paragraph{Learning.} Cycle 112 turns the transparent FormulaEval packet into a finite theorem pair: transparency is exactly a charged support-realization condition, and hidden semantic support is not merely difficult but has a two-world private-row falsifier. The contribution remains a route-level parser/falsifier boundary for beta-zero magnification; the main P vs NP conjecture is untouched.

\paragraph{Pressure.} No blocking gap remains for \(\mathrm{Psys}_{sym}^{trans}\) as defined. The extension to full semantic \(\mathrm{Psys}_{sym}\) is blocked exactly at \(RealSATInfluenceRealizer_{N}\), which must convert real-SAT semantic influence into charged same-theory row support.

\paragraph{Derivation steps.}
\begin{enumerate}
\item Define \(\mathrm{Psys}_{sym}^{trans}\) as \(\mathrm{Psys}_{bEF}^{-}\) plus FormulaEval rules carrying visible selector code, RowSet enumerator or full-scan declaration, influence-closure proof, and same-theory occurrence expansion template.
\item Induct on proof lines. Bit-explicit EF lines extend the transcript by their declared support and are charged inside Gamma\_bEF(pi).
\item For bounded \(SR_{h}\) lines, use the existing \(SR_{h}\) bit-normal form: parent pointers, row-index fields, touched occurrence-neighborhood list, and eta/tau overhead contribute B\_SR(pi) inside Gamma\_bEF(pi).
\item For a transparent FormulaEval line a, run the certified RowSet enumerator. If output-sensitive enumeration is not certified, charge the full syntactic scan. Add this to U\_N(pi).
\item Expand each row in the influence closure of a into same-theory occurrence clauses and charge the total row expansion to X\_N(pi). This prevents semantic FormulaEval from acting as a free SAT oracle.
\item The final transcript gives a circuit for \(H_{N}^\mathrm{SAT}\) by wiring the generator \(G_{N}\) of size \(N^{\rho+o(1)}\), the parsed local decisions from Gamma\_trans(pi), and equality patches for \(E_{N}\) residual rows.
\item If the resulting size is below \(N^{1+\epsilon-o(1)}\), it contradicts the source lower-bound statement LB\_N(1+epsilon); hence such a proof pi is a valid falsifier for the claimed lower bound, not evidence for \(\mathrm{P} \ne \mathrm{NP}\).
\item For hidden FormulaEval, construct two label structures that agree on every charged visible row but differ on an uncharged influencing row or unread enable bit. The same transcript then supports different carrier values, so occurrence completion is invalid.
\end{enumerate}

\paragraph{Equation 1: carrier.}
\[
H_N^{SAT}(x)=\bigoplus_{j=1}^{M_N} SAT(G_N(x,j)),\quad M_N=N^{q+o(1)},\quad |G_N(x,j)|=(\log N)^{d+O(1)}.
\]
\emph{Valid when:} Valid when the source carrier is the beta-zero batched SAT parity object with polylogarithmic generated formulas.

\paragraph{Equation 2: transparent\_charge.}
\[
\Gamma_{trans}(\pi)=\Gamma_{bEF}(\pi)+U_N(\pi)+X_N(\pi).
\]
\emph{Valid when:} Valid when every FormulaEval line satisfies \(TransparentRowSupport_{N}\) and all influence rows are enumerated or full-scan charged.

\paragraph{Equation 3: parser\_upper\_bound.}
\[
C(H_N^{SAT})\le N^{\rho+o(1)}+\Gamma_{trans}(\pi)+O(E_NN\operatorname{polylog}N).
\]
\emph{Valid when:} Applies after the parser line induction materializes a same-theory transcript and residual row errors are equality-patched.

\paragraph{Equation 4: falsifier\_threshold.}
\[
N^{\rho+o(1)}+\Gamma_{trans}(\pi)+E_NN< N^{1+\epsilon-o(1)}\Rightarrow \neg LB_N(1+\epsilon).
\]
\emph{Valid when:} Requires the source lower-bound claim and the extracted circuit to be measured in the same nonuniform SIZE model.

\paragraph{Equation 5: support\_density\_corollary.}
\[
|K_{idx}(\pi)|\ge N^{\epsilon-o(1)},\qquad |K_{idx}(\pi)|/M_N\ge N^{\epsilon-q-o(1)}.
\]
\emph{Valid when:} Valid when each unforced row can be patched at cost O(N polylog N) and q is fixed in \(M_{N}=N^{q+o(1)}\).

\paragraph{Next analysis.} Write the corollary table as examples only: interval-union, full scan, ROBDD, guarded-dDNNF, projected-dDNNF, and influence-complete extension rules must each be checked against the five \(TransparentRowSupport_{N}\) clauses or demoted.

\paragraph{Falsifiers.}
\begin{itemize}
\item private-row falsifier: an uncharged generated row changes the carrier while the visible transcript is unchanged
\item unread-enable-bit falsifier: a semantic FormulaEval gate has an unparsed enable bit that toggles influence without appearing in RowSet
\item threshold falsifier: N\textasciicircum{}\{rho+o(1)\}+Gamma\_trans(pi)+E\_NN<N\textasciicircum{}\{1+epsilon-o(1)\} gives a same-model circuit refuting LB\_N(1+epsilon)
\item nonuniform leakage falsifier: RowSet enumeration depends on hidden advice not counted in \(U_{N}\) or \(X_{N}\)
\item same-theory failure falsifier: FormulaEval expansion uses an external SAT oracle rather than proof-system lines
\item positive-beta background falsifier: any black-box target substitution with exponent \(\beta>0\) still obeys q+beta c and cannot exclude all SAT exponents
\end{itemize}

\subsection*{cycle-0113: Hardness-Magnification Bridge Realization}

\paragraph{Execution.} research; azure\_responses\_api; attempt primary; 152.132 seconds; 4 lemmas, 5 equations, 7 derivation steps, 6 falsifiers.

\paragraph{Learning.} Cycle 113 converts the promoted FormulaEval boundary into a paper-ready theorem pair: transparent support is exactly a charged support-realization interface yielding \(\Gamma_{trans}\) and the \(\mathrm{LB}_{N}\) falsifier window, while hidden semantic support is blocked by a concrete private-row/unread-enable-bit falsifier. The full-system escape is precisely \(RealSATInfluenceRealizer_{N}\), not generic barrier folklore.

\paragraph{Pressure.} No blocking gap remains for the restricted \(TransparentRowSupport_{N}\) theorem pair. The full semantic \(\mathrm{Psys}_{sym}\) extension is blocked exactly at \(RealSATInfluenceRealizer_{N}\), an advice-free nonrelativizing conversion from real SAT influence to charged occurrence support.

\paragraph{Derivation steps.}
\begin{enumerate}
\item Fix N on the valid length sequence, with \(H_{N}^{\mathrm{SAT}}\), \(M_{N}=N^{q+o(1)}\), L\_N=(log N)\textasciicircum{}d, \(q,\rho<1+\epsilon\), and \(\beta=0\); use the same nonuniform SIZE model for source upper bounds and \(\mathrm{LB}_{N}\).
\item Reuse the audited \(\mathrm{Psys}_{bEF}^{-}\) line induction: bit-explicit EF dependencies, SATeq links, and \(bounded-SR_{h}\) leaves contribute exactly Gamma\_bEF(pi).
\item For each transparent FormulaEval line ell, run its visible selector enumerator at cost \(U_{ell}\) or charge a full scan \(F_{ell}\) \(N^{1+q+o(1)}\); expand every selected generated formula occurrence at cost \(X_{ell}=|K_{ell}|L_{N}^{O(1)}\) plus parser overhead.
\item Invoke influence closure: any SAT-label change outside \(K_{ell}\) leaves ell invariant in the same proof theory, so no uncharged row can affect later functional lines.
\item Induct on proof lines, adding \(U_{ell}+X_{ell}\) for transparent FormulaEval lines and prior \(\Gamma_{bEF}\) charges for bit-explicit lines, yielding a circuit/transcript for the final functional certificate of size Gamma\_trans(pi) plus generator overhead \(N^{\rho+o(1)}\).
\item Patch the \(E_{N}\) residual source rows by equality tests at cost O(E\_N N); if the resulting size is below \(N^{1+\epsilon-o(1)}\), the asserted LB\_N(1+epsilon) is falsified in the same model.
\item For a hidden FormulaEval influence, construct two label worlds A and B agreeing on all visible transcript fields but differing at one private row-slot with odd parity contribution; the checker cannot distinguish them, so occurrence parsing is invalid without \(RealSATInfluenceRealizer_{N}\).
\end{enumerate}

\paragraph{Equation 1: carrier.}
\[
H_N^{SAT}(x)=\bigoplus_{j\le M_N} SAT(G_N(x,j)),\quad M_N=N^{q+o(1)},\ L_N=(\log N)^d
\]
\emph{Valid when:} Valid when the source carrier is the beta-zero batched SAT parity family with uniformly generated polylog-length formulas.

\paragraph{Equation 2: transparent\_charge.}
\[
\Gamma_{trans}(\pi)=\Gamma_{bEF}(\pi)+U_N(\pi)+X_N(\pi)
\]
\emph{Valid when:} Applies after every FormulaEval line satisfies \(TransparentRowSupport_{N}\) or is replaced by a charged full syntactic scan.

\paragraph{Equation 3: parser\_upper\_bound.}
\[
C_{SIZE}(H_N^{SAT})\le N^{\rho+o(1)}+\Gamma_{trans}(\pi)+O(E_NN)
\]
\emph{Valid when:} Requires the final certificate to be functional over surfaced rows and all residual bad rows to be explicitly patchable.

\paragraph{Equation 4: lb\_falsifier\_window.}
\[
N^{\rho+o(1)}+\Gamma_{trans}(\pi)+O(E_NN)< N^{1+\epsilon-o(1)}\Rightarrow \neg LB_N(1+\epsilon)
\]
\emph{Valid when:} Valid when \(\mathrm{LB}_{N}\) is interpreted in the same nonuniform SIZE model and no hidden advice enters the parser.

\paragraph{Equation 5: support\_density\_pressure.}
\[
|K_{idx}(\pi)|\ge N^{\epsilon-o(1)},\quad |K_{idx}(\pi)|/M_N\ge N^{\epsilon-q-o(1)}
\]
\emph{Valid when:} Requires the private-row falsifier regime where every influential FormulaEval slot must be surfaced or sparsely patched.

\paragraph{Next analysis.} Run the finite proof-quality checklist on the paper statement: verify each \(TransparentRowSupport_{N}\) clause for interval-union, full scan, ROBDD, guarded-dDNNF, projected-dDNNF, and influence-complete examples; then attempt or refute \(RealSATInfluenceRealizer_{N}\) for one concrete semantic FormulaEval axiom.

\paragraph{Falsifiers.}
\begin{itemize}
\item private-row falsifier: an unsurfaced SAT label changes \(H_{N}\) while the visible transcript and \(\Gamma_{trans}\) remain fixed
\item unread-enable-bit falsifier: selector activation is hidden in semantic advice, so enumeration and influence closure are not valid
\item \(q<\epsilon\) support-density falsifier: the visible slot universe is too small to force \(N^{\epsilon-o(1)}\) complete support
\item subthreshold residual falsifier: \(E_{N}<N^{\epsilon-o(1)}\) bad rows are patched by O(E\_N N) equality gates
\item positive-beta transfer falsifier: any black-box target-circuit substitution with query exponent \(\beta>0\) is capped by q+beta c+o(1) and cannot exclude all SAT circuit exponents
\item hidden-advice falsifier: selector tables or FormulaEval witnesses not charged in \(U_{N}\) or \(X_{N}\) invalidate the same-model circuit bound
\end{itemize}

\subsection*{cycle-0114: Hardness-Magnification Bridge Realization}

\paragraph{Execution.} research; azure\_responses\_api; attempt primary; 152.053 seconds; 3 lemmas, 5 equations, 7 derivation steps, 5 falsifiers.

\paragraph{Learning.} Cycle 114 upgrades the promoted theorem pair into an auditable live inference packet: transparent FormulaEval is exactly charged support realization, hidden semantic FormulaEval has a two-world falsifier, and the full-system question is the named nonrelativizing \(RealSATInfluenceRealizer_{N}\) rather than generic barrier folklore.

\paragraph{Pressure.} The positive theorem is complete for \(TransparentRowSupport_{N}\), but full semantic FormulaEval remains blocked at \(RealSATInfluenceRealizer_{N}\): no advice-free nonrelativizing procedure is known that surfaces all real-SAT influences with subthreshold cost.

\paragraph{Derivation steps.}
\begin{enumerate}
\item Fix a valid length N, beta-zero carrier \(H_{N}^\mathrm{SAT}\), \(M_{N}=N^{q+o(1)}\), generator \(G_{N}\) of size \(N^{\rho+o(1)}\), and a \(\mathrm{Psys}_{bEF}^{-}\) plus FormulaEval proof pi claiming LB\_N(1+epsilon).
\item Maintain the line-induction invariant \(I_{t}\): after t lines, every row label that can affect the parsed proof state is either in the charged support \(K_{t}\) or counted as a residual row; the accumulated cost is \(B_\mathrm{EF}+B_{eq}+B_{SR}+U+X\).
\item For EF, SATeq, and \(bounded-SR_{h}\) lines, use the existing bit-explicit parser: parent pointers, row-index encodings, touched occurrence neighborhoods, and eta/tau leaf costs update \(\Gamma_{bEF}\) without hidden semantic advice.
\item For a transparent FormulaEval line ell, run the visible selector or full scan, charge \(U_{ell}\), expand each row in \(RowSet_{ell}\) into same-theory occurrence support, charge \(X_{ell}=N^{1+o(1)}|RowSet_{ell}|\), and update \(K_{t}\) by influence closure.
\item At the final line, hardwire pi and the charged transcript into a nonuniform circuit: compute \(G_{N}\), evaluate all charged local transcript components, and equality-patch \(E_{pi}\) residual source rows at cost O(E\_pi \(N^{1+o(1)}\)).
\item This gives C(H\_N\textasciicircum{}SAT)<=N\textasciicircum{}\{rho+o(1)\}+Gamma\_trans(pi)+O(E\_pi \(N^{1+o(1)}\)); contraposition yields the support-pressure and q-density corollaries for any sound \(\mathrm{LB}_{N}\) certificate.
\item For hidden semantic FormulaEval, choose an influential uncharged row r*, construct two row-label worlds or an unread-enable-bit pair agreeing on K\_vis(pi) and all visible proof checks, but differing on the parity of \(H_{N}\) at one source input. Thus the occurrence parser is falsified unless real SAT semantics is converted into charged support by \(RealSATInfluenceRealizer_{N}\).
\end{enumerate}

\paragraph{Equation 1: \(GammaTrans\).}
\[
\Gamma_{\mathrm{trans}}(\pi)=B_{\mathrm{EF}}(\pi)+B_{\mathrm{eq}}(\pi)+B_{\mathrm{SR}}(\pi)+U_{\pi}+X_{\pi}
\]
\emph{Valid when:} Valid when the proof system line types are bit-explicit or satisfy \(TransparentRowSupport_{N}\) with all FormulaEval support charges surfaced.

\paragraph{Equation 2: FormulaEvalExpansion.}
\[
X_{\pi}=\sum_{\ell\in FE(\pi)} N^{1+o(1)}\,|RowSet_{\ell}|
\]
\emph{Valid when:} Applies after every FormulaEval selector has an advice-free RowSet enumeration or the full-scan alternative is charged.

\paragraph{Equation 3: ParsedCircuitUpperBound.}
\[
C(H_N^{SAT})\le N^{\rho+o(1)}+\Gamma_{\mathrm{trans}}(\pi)+O(E_{\pi}N^{1+o(1)})
\]
\emph{Valid when:} Requires the target-preserving beta-zero generator, same-theory expansion, and equality patching of only \(E_{pi}\) residual source rows.

\paragraph{Equation 4: SoundCertificatePressure.}
\[
\Gamma_{\mathrm{trans}}(\pi)+E_{\pi}N^{1+o(1)}\ge N^{1+\epsilon-o(1)}
\]
\emph{Valid when:} Valid for sound certificates of LB\_N(1+epsilon) once the previous circuit upper bound is accepted in the same SIZE model.

\paragraph{Equation 5: SupportDensityCorollary.}
\[
|K_{idx}(\pi)|\ge N^{\epsilon-o(1)},\quad |K_{idx}(\pi)|/M_N\ge N^{\epsilon-q-o(1)}
\]
\emph{Valid when:} Applies when generator, EF, equality, and bounded-SR charges are subthreshold and FormulaEval row support is the remaining pressure term.

\paragraph{Next analysis.} Instantiate \(RealSATInfluenceRealizer_{N}\) for one concrete full semantic FormulaEval rule, or prove a rule-specific private-row countermodel showing that the full \(\mathrm{Psys}_{sym}\) extension must remain quarantined.

\paragraph{Falsifiers.}
\begin{itemize}
\item private-row falsifier: an uncharged influential row makes two visible transcripts agree while \(H_{N}\) parity differs
\item unread-enable-bit falsifier: a FormulaEval selector can hide whether a row is active unless the enable bit is visible and charged
\item threshold falsifier: \(N^{\rho+o(1)}+\Gamma_{trans}(pi)+E_{pi}\) N\textasciicircum{}\{1+o(1)\}<N\textasciicircum{}\{1+epsilon-o(1)\} refutes LB\_N(1+epsilon) in the same SIZE model
\item \(q<\epsilon\) density falsifier: required support density exceeds one when FormulaEval is the only remaining pressure term
\item semantic-advice falsifier: \(RealSATInfluenceRealizer_{N}\) is invalid if it smuggles nonuniform row labels into RowSet output
\end{itemize}

\subsection*{cycle-0115: Hardness-Magnification Bridge Realization}

\paragraph{Execution.} research; azure\_responses\_api; attempt primary; 970.044 seconds; 3 lemmas, 3 equations, 5 derivation steps, 6 falsifiers.

\paragraph{Learning.} Cycle 115 adds a minimality proof obligation to the promoted transparent-support theorem pair. Transparent FormulaEval is not just sufficient for a valid \(\Gamma_{trans}\) parser; hidden influence is exactly what the private-row/unread-enable-bit falsifier exploits. Thus the full-system boundary is a named nonrelativizing \(RealSATInfluenceRealizer_{N}\) lemma, while the paper-level contribution remains a scoped parser/falsifier result rather than evidence for \(\mathrm{P} \ne \mathrm{NP}\).

\paragraph{Pressure.} The transparent theorem and hidden-support countermodel are proof-shaped for occurrence-visible fragments, but the full semantic case is blocked by \(RealSATInfluenceRealizer_{N}\): one must show that real SAT semantics forces every influential row into a charged support set.

\paragraph{Derivation steps.}
\begin{enumerate}
\item Define the row universe \(R_{N}={(x,j):j\le M_{N}}\) and treat a FormulaEval line L as transparent only when its possible dependence set Infl\_pi(L) is listed by visible selector code or charged by full scan.
\item For transparent L, add enumerator cost U\_N(L) and row-expansion cost X\_N(L) to \(\Gamma_{trans}\); the existing bit-explicit line induction then compiles L into a circuit fragment without changing the SIZE model.
\item For a hidden row r with \(Delta_{r}(L)=1\) and r not in the charged support, build two occurrence-label worlds that agree on pi, visible selectors, EF support, SATeq links, and bounded \(SR_{h}\) neighborhoods but differ on SAT(r).
\item If the proof transcript is identical in the two worlds while the carrier parity changes, the parser-to-circuit inference is invalid; this is the private-row/unread-enable-bit falsifier for hidden FormulaEval support.
\item Contrapositively, any valid parser theorem for full semantic FormulaEval must prove that no such r exists, which is exactly the \(RealSATInfluenceRealizer_{N}\) obligation with charged output size.
\end{enumerate}

\paragraph{Equation 1: \(GammaTrans\).}
\[
Gamma_trans(pi)=Gamma_bEF(pi)+U_N(pi)+X_N(pi)+O(E_N N)
\]
\emph{Valid when:} Valid when every FormulaEval influence is surfaced by \(TransparentRowSupport_{N}\) or charged by full scan and residual row errors are explicitly counted.

\paragraph{Equation 2: ParsedCircuitUpperBound.}
\[
C_SIZE(H_N^SAT)\le N^{rho+o(1)}+Gamma_trans(pi)
\]
\emph{Valid when:} Applies after the generator, all parsed proof lines, and all transparent FormulaEval row expansions are compiled into the same nonuniform SIZE model.

\paragraph{Equation 3: \(\mathrm{LB}_{Falsifier,Window}\).}
\[
N^{rho+o(1)}+Gamma_trans(pi)+O(E_N N)< N^{1+epsilon-o(1)} => not LB_N(1+epsilon)
\]
\emph{Valid when:} Requires the source lower-bound claim to be in the same SIZE model and the parser transcript to compute \(H_{N}^\mathrm{SAT}\) on all but \(E_{N}\) patchable rows.

\paragraph{Next analysis.} Formalize the private-row independence lemma for one full symbolic FormulaEval rule with an unread enable bit: write the two transcripts, the charged constraint set, the flipped row r, and the proof that all visible lines agree while \(H_{N}^\mathrm{SAT}\) parity differs; if independence fails, record the exact same-theory constraint as a \(RealSATInfluenceRealizer_{N}\) candidate.

\paragraph{Falsifiers.}
\begin{itemize}
\item private-row falsifier: an uncharged influential FormulaEval row can be flipped while the visible proof transcript is fixed
\item unread-enable-bit falsifier: a selector contains a dormant semantic enable bit whose active row is not surfaced by \(TransparentRowSupport_{N}\)
\item \(q<\epsilon\) support-density falsifier: absolute support below \(N^{\epsilon-o(1)}\) gives a subthreshold parsed circuit after patching
\item full-scan collapse: unenumerated transparent FormulaEval pays \(F_{FE}\) \(N^{1+q+o(1)}\), which destroys the barely-superlinear window when q is too large
\item nonuniform leakage falsifier: hidden selector advice names rows outside the proof bit string and invalidates the parser theorem
\item positive-beta substitution cap: any black-box target-query exponent \(\beta>0\) leaves the transfer window q+beta c+o(1) and cannot exclude all SAT polynomial exponents
\end{itemize}

\subsection*{cycle-0116: Hardness-Magnification Bridge Realization: Transparent FormulaEval Minimality Audit}

\paragraph{Execution.} research; azure\_responses\_api; attempt primary; 148.454 seconds; 3 lemmas, 4 equations, 10 derivation steps, 5 falsifiers.

\paragraph{Learning.} Cycle 116 closes the minimality audit for the transparent FormulaEval theorem pair. The valid paper claim is now a restricted parser/falsifier boundary: every surfaced influence is charged in \(\Gamma_{trans}\), and every unsurfaced nonsemantic influence admits a private-row falsifier. The only remaining full-system route is the named nonrelativizing \(RealSATInfluenceRealizer_{N}\); this is not evidence for \(\mathrm{P} \ne \mathrm{NP}\).

\paragraph{Pressure.} Full semantic FormulaEval remains blocked at \(RealSATInfluenceRealizer_{N}\): private-row spoofing is a relativizing falsifier for hidden support, but it cannot refute genuinely real-SAT semantic influence without a same-theory realization or a concrete semantic no-go.

\paragraph{Derivation steps.}
\begin{enumerate}
\item Fix N, \(M_{N}=N^{q+o(1)}\), \(\beta=0\), generator \(G_{N}\) of size \(N^{\rho+o(1)}\), and a \(\mathrm{Psys}_{sym}\) proof pi whose non-FormulaEval lines are parsed by the \(\mathrm{Psys}_{bEF}^{-}\) theorem.
\item For each transparent FormulaEval line l, compute \(RowSet_{l}\) from visible selector code. If enumeration is output-sensitive charge \(U_{l}\); otherwise full scan costs \(M_{N}\) \(N^{1+o(1)}\) and is recorded in \(X_{FE}\).
\item Expand every j in \(RowSet_{l}\) into the same proof theory by inserting the occurrence neighborhood for SAT(G\_N(x,j)) over all source rows x; charge \(N^{1+o(1)}\) per row index and add this to \(\Gamma_{trans}\).
\item Run line induction: bit-explicit EF dependencies, SATeq links, \(bounded-SR_{h}\) leaves, and transparent FormulaEval expansions each yield circuits for the line predicates with total size bounded by \(\Gamma_{trans}\) plus generator overhead.
\item If the proof certifies LB\_N(1+epsilon) but \(N^{\rho+o(1)}+\Gamma_{trans}+E_{N}\) \(N^{1+o(1)}\) is below \(N^{1+\epsilon-o(1)}\), the derived circuit is a valid falsifier for the asserted source lower bound.
\item For minimality, suppose a FormulaEval line has a hidden influential row r not in the charged RowSet. Build two label worlds agreeing on all charged rows and proof-visible bits but differing on the SAT label of r; carrier parity differs after choosing x with r's coefficient one, so occurrence completion was not sound.
\item If the selector has an unread enable bit, choose two settings with identical parsed transcript and different inclusion of r. This gives the same two-world falsifier and shows extensional enumeration is necessary.
\item If influence closure fails, downstream FormulaEval can depend on r through an unlisted derived predicate. Flipping r preserves the transcript but changes the carrier-relevant line, so the parser misses a row charge.
\item If hidden advice selects r, the proof transcript is unchanged while the selected support changes. The advice must either be charged as \(U_{FE}\) or the private-row falsifier applies.
\item If same-theory expansion fails because real SAT semantics determine r, the two-world spoof may violate SAT. This is not a valid countermodel; it is exactly the \(RealSATInfluenceRealizer_{N}\) blocked lemma.
\end{enumerate}

\paragraph{Equation 1: carrier.}
\[
H_N^SAT(x)=XOR_{j< M_N} SAT(G_N(x,j)), M_N=N^{q+o(1)}, |G_N|=N^{rho+o(1)}
\]
\emph{Valid when:} Valid when the source carrier is the beta-zero batched SAT parity problem with polylogarithmic generated instances and explicit generator overhead.

\paragraph{Equation 2: transparent\_charge.}
\[
Gamma_trans(pi)=Gamma_bEF(pi)+U_FE(pi)+X_FE(pi), X_FE(pi)=N^{1+o(1)} sum_l |RowSet_l|
\]
\emph{Valid when:} Applies after every FormulaEval line has visible advice-free row support and same-theory occurrence expansion charged by row index.

\paragraph{Equation 3: parser\_upper\_bound.}
\[
C_SIZE(H_N^SAT)\le N^{rho+o(1)}+Gamma_trans(pi)+E_N N^{1+o(1)}
\]
\emph{Valid when:} Valid when the parsed proof transcript computes the carrier except for \(E_{N}\) residual source rows patchable by equality gates.

\paragraph{Equation 4: lb\_falsifier\_window.}
\[
N^{rho+o(1)}+Gamma_trans(pi)+E_N N^{1+o(1)}< N^{1+epsilon-o(1)} => not LB_N(1+epsilon)
\]
\emph{Valid when:} Requires the target-preserving SIZE simulation, beta-zero substitution, and exact lower-bound threshold to be in the same length regime.

\paragraph{Next analysis.} Instantiate \(RealSATInfluenceRealizer_{N}\) or refute it for one concrete strong rule: FormulaEval with Tseitin-extension equality plus bounded self-reduction; list exact bit fields, RowSet extraction, same-theory expansion certificates, and the first semantic inference that cannot be charged.

\paragraph{Falsifiers.}
\begin{itemize}
\item Private-row falsifier: an uncharged influential row r has two label worlds agreeing on the parsed transcript but differing on \(H_{N}^\mathrm{SAT}\).
\item Unread-enable-bit falsifier: selector inclusion of r is controlled by a bit not read by the parser, so RowSet is not extensionally valid.
\item Hidden-advice falsifier: nonuniform row-selection advice can encode the needed support while evading \(U_{FE}\) and \(X_{FE}\) charges.
\item Subthreshold \(\Gamma\) falsifier: \(N^{\rho}+\Gamma_{trans}+E_{N}\) \(N<N^{1+\epsilon}\) gives an explicit small SIZE circuit against LB\_N(1+epsilon).
\item Same-theory failure is not a valid oracle falsifier; it quarantines the route at \(RealSATInfluenceRealizer_{N}\) rather than at generic barrier folklore.
\end{itemize}

\subsection*{cycle-0117: Hardness-Magnification Bridge Realization: Transparent FormulaEval Boundary}

\paragraph{Execution.} research; azure\_responses\_api; attempt primary; 118.652 seconds; 3 lemmas, 4 equations, 8 derivation steps, 5 falsifiers.

\paragraph{Learning.} Cycle 117 upgrades the promoted theorem pair from sufficiency plus informal minimality to a local safety iff for extensional FormulaEval: surfaced influence is exactly what makes the \(\Gamma_{trans}\) parser valid, and unsurfaced nonsemantic influence is exactly what the private-row/unread-enable-bit falsifier exploits. The full semantic case is not dismissed by a barrier slogan; it is isolated as the nonrelativizing \(RealSATInfluenceRealizer_{N}\) obligation. This is a paper-level route-boundary result, not a proof of \(\mathrm{P} \ne \mathrm{NP}\).

\paragraph{Pressure.} The derivation is complete for extensional FormulaEval rules but still blocked for full semantic FormulaEval at \(RealSATInfluenceRealizer_{N}\), which must nonrelativizingly force real SAT influence into charged same-theory support.

\paragraph{Derivation steps.}
\begin{enumerate}
\item Fix a \(\mathrm{Psys}_{sym}\) proof pi of LB\_N(1+epsilon) for \(H_{N}^\mathrm{SAT}\) and parse all \(\mathrm{Psys}_{bEF}^{-}\) lines using the existing bit-explicit EF, SATeq, \(bounded-SR_{h}\), and sparse-row charges.
\item For each FormulaEval line, require \(TransparentRowSupport_{N}\): the parser receives visible selector code, an advice-free enumerator or full scan, row expansion into occurrence neighborhoods, and influence closure under enable bits and projections.
\item Add the enumerator cost \(U_{l}\), expansion cost \(X_{l}\), or full-scan cost \(N^{1+q+o(1)}\) for the line; summing over lines yields Gamma\_trans(pi).
\item Run the line induction: every accepted proof line has a circuit/transcript implementation whose size is bounded by the accumulated \(\Gamma_{trans}\) prefix, and \(bounded-SR_{h}\) contributes exactly its previous \(B_{SR}\) charge.
\item Compose the generator circuit of size \(N^{\rho+o(1)}\) with the parsed transcript; if the residual exceptional rows are \(E_{N}\), patch them by equality tests at cost O(E\_N N).
\item Contrapose: when the total size is below \(N^{1+\epsilon-o(1)}\), the alleged LB\_N(1+epsilon) certificate is falsified inside the same nonuniform SIZE model.
\item For the converse minimality check, suppose an extensional FormulaEval line has a hidden influencing row r. Construct two row-label worlds equal on all surfaced rows and proof text but differing at r; the parser output is identical while the FormulaEval value or parity carrier differs.
\item Therefore no occurrence parser that omits r can be sound for extensional rules; this proves local necessity of support visibility and identifies real SAT semantics as the only possible escape.
\end{enumerate}

\paragraph{Equation 1: carrier.}
\[
H_N^{SAT}(x)=\bigoplus_{j\le M_N} SAT(G_N(x,j)),\quad M_N=N^{q+o(1)},\quad |G_N|=N^{\rho+o(1)}
\]
\emph{Valid when:} Valid when the source carrier is the beta-zero generated parity family and the generator produces polylogarithmic SAT instances with uniform decoding.

\paragraph{Equation 2: transparent\_charge.}
\[
\Gamma_{trans}(\pi)=\Gamma_{bEF}(\pi)+U_N(\pi)+X_N(\pi)+F_{FE}(\pi)N^{1+q+o(1)}
\]
\emph{Valid when:} Applies after every FormulaEval line either supplies an advice-free row enumerator or is charged by full syntactic scan in the same proof theory.

\paragraph{Equation 3: falsifier\_window.}
\[
N^{\rho+o(1)}+\Gamma_{trans}(\pi)+E_NN< N^{1+\epsilon-o(1)}\;\Longrightarrow\; \neg LB_N(1+\epsilon)
\]
\emph{Valid when:} Requires the parser-to-circuit simulation to preserve the nonuniform SIZE model and the residual error set to have at most \(E_{N}\) source rows.

\paragraph{Equation 4: support\_density\_corollary.}
\[
|K_{idx}(\pi)|/M_N\ge N^{\epsilon-q-o(1)}
\]
\emph{Valid when:} Valid when LB\_N(1+epsilon) is claimed and the only surfaced FormulaEval support is row-index support among \(M_{N}=N^{q+o(1)}\) slots.

\paragraph{Next analysis.} Write the paper-ready theorem statement as three clauses: \(TransparentRowSupport_{N}\) positive \(\Gamma_{trans}\) induction; extensional hidden-support private-row/unread-enable-bit countermodel; \(RealSATInfluenceRealizer_{N}\) open semantic escape with overhead \(R_{N}\) and explicit non-claims.

\paragraph{Falsifiers.}
\begin{itemize}
\item Private-row falsifier: an uncharged active row outside K\_idx(pi) yields two visible-identical transcripts with different row labels, invalidating occurrence completion.
\item Unread-enable-bit falsifier: a hidden guard selects a row only in one world, so selector extensionality fails unless the guard is surfaced and charged.
\item Threshold falsifier: if N\textasciicircum{}\{rho+o(1)\}+Gamma\_trans(pi)+E\_NN<N\textasciicircum{}\{1+epsilon-o(1)\}, any claimed LB\_N(1+epsilon) certificate collapses to a small circuit.
\item Positive-beta falsifier: any repair introducing target queries of length \(N^\beta\) with \(\beta>0\) reactivates the q+beta c transfer cap and cannot exclude all SAT polynomial exponents.
\item Nonuniform-leakage falsifier: if row enumeration uses advice depending on SAT labels, the parser theorem is invalid even when its apparent \(U_{N}\) charge is small.
\end{itemize}

\subsection*{cycle-0118: Hardness-Magnification Bridge Realization}

\paragraph{Execution.} research; azure\_responses\_api; attempt primary; 111.25 seconds; 3 lemmas, 6 equations, 8 derivation steps, 5 falsifiers.

\paragraph{Learning.} Cycle 118 adds a compositional repair to the FormulaEval boundary: local surfaced influence is not enough unless hidden influences are closed under later proof use. The valid restricted theorem is \(\mathrm{Psys}_{sym}^{ext}\) plus charged \(RealSATInfluenceRealizer_{N}\) rows implies \(\Gamma_{comp}\) line induction and the \(\mathrm{LB}_{N}\) falsifier window. The necessity direction is now a last-unsafe-line private-row theorem. Full semantic FormulaEval remains open exactly at a nonrelativizing realizer lemma, not at generic relativization, natural-proofs, or P vs NP folklore.

\paragraph{Pressure.} The last-unsafe-line argument is valid in the extensional label model. For real SAT semantics, a hidden row may be forced by nonrelativizing formula identities, self-reduction constraints, or EqSAT links without being syntactically visible. That exact gap is \(RealSATInfluenceRealizer_{N}\).

\paragraph{Derivation steps.}
\begin{enumerate}
\item Define the influence cone \(I_{l}\) for each FormulaEval line l as the set of generated row labels whose flip can change the line value with all visible syntax fixed.
\item For transparent lines, enumerate \(I_{l}\) by the visible selector code or charge a full scan over \(M_{N}\) rows; add \(U_{l}+X_{l}\) to \(\Gamma_{comp}\).
\item For bEF, SATeq, and \(bounded-SR_{h}\) lines, reuse the bit-explicit dependency parser and \(SR_{h}\) field parser; add \(\Gamma_{bEF}\) and \(B_{SR,h}\) to \(\Gamma_{comp}\).
\item For hidden FormulaEval lines, either a same-theory \(RealSATInfluenceRealizer_{N}\) lists the forced influence rows and contributes \(R_{FE}\), or mark the line unsafe.
\item If no unsafe line remains, perform line induction: every proof line is simulated by a circuit using the generator, charged row labels, charged SR leaves, and charged FormulaEval expansions.
\item The final line gives a circuit for \(H_{N}\) except on \(E_{N}\) explicitly residual rows; equality patches yield \(C(H_{N})\le N^{\rho+o(1)}+\Gamma_{comp}(pi)+E_{NN}\).
\item If an unsafe line remains, choose the last unsafe line whose hidden row reaches the final output without later charged realization. Flip that row in an extensional two-world model while fixing all charged rows; the transcript is identical but the induced parity differs.
\item Conclude that surfaced influence plus realizer closure is sufficient for the parser, and failure of both is a private-row falsifier for extensional FormulaEval.
\end{enumerate}

\paragraph{Equation 1: carrier.}
\[
H_N^{SAT}(x)=\bigoplus_{j=1}^{M_N} SAT(G_N(x,j))
\]
\emph{Valid when:} Valid when the generated instance family \(G_{N}\) is fixed, advice-free, and has \(M_{N}=N^{q+o(1)}\) indexed rows of polylogarithmic target length.

\paragraph{Equation 2: composite\_transcript\_charge.}
\[
\Gamma_{comp}(\pi)=\Gamma_{bEF}(\pi)+B_{SR_h}(\pi)+U_{FE}(\pi)+X_{FE}(\pi)+R_{FE}(\pi)
\]
\emph{Valid when:} Applies after bit-explicit EF support, bounded \(SR_{h}\) fields, transparent FormulaEval enumeration, expansion charges, and semantic realizer rows are all separately charged.

\paragraph{Equation 3: sr\_charge.}
\[
B_{SR_h}(\pi)\le N^{q+\eta+\tau+o(1)}
\]
\emph{Valid when:} Valid when the bounded self-reduction rule has at most \(N^{\eta+o(1)}\) leaves per active row and leaf verifier overhead \(N^{\tau+o(1)}\).

\paragraph{Equation 4: parser\_upper\_bound.}
\[
C(H_N^{SAT})\le N^{\rho+o(1)}+\Gamma_{comp}(\pi)+E_N N
\]
\emph{Valid when:} Requires the parser to produce a same-model circuit for all but \(E_{N}\) exceptional source rows, with equality patches charged nonuniformly per row.

\paragraph{Equation 5: falsifier\_window.}
\[
N^{\rho+o(1)}+\Gamma_{comp}(\pi)+E_NN< N^{1+\epsilon-o(1)}\Rightarrow \neg LB_N(1+\epsilon)
\]
\emph{Valid when:} Valid when \(\mathrm{LB}_{N}\) is stated in the identical nonuniform SIZE model and the parsed circuit computes the carrier on every valid input length.

\paragraph{Equation 6: support\_density\_corollary.}
\[
|K_{idx}(\pi)|/M_N\ge N^{\epsilon-q-o(1)}
\]
\emph{Valid when:} Follows only when all other transcript terms are subthreshold and FormulaEval support is the remaining necessary charged source of hardness.

\paragraph{Next analysis.} Run a finite proof audit on \(\mathrm{Psys}_{sym}^{ext}\): instantiate one guarded FormulaEval line, compute \(I_{l}\), \(K_{idx}\), \(R_{FE}\), and \(\Gamma_{comp}\), then verify the last-unsafe-line two-world construction or identify the first real SAT-semantic rule that forces a nonrelativizing realizer charge.

\paragraph{Falsifiers.}
\begin{itemize}
\item Subthreshold \(\Gamma_{comp}\) falsifier: N\textasciicircum{}\{rho+o(1)\}+Gamma\_comp(pi)+E\_NN<N\textasciicircum{}\{1+epsilon-o(1)\} contradicts any claimed LB\_N(1+epsilon) in the same SIZE model.
\item Private-row falsifier: an unsurfaced nonsemantic FormulaEval influence outside \(K_{idx}\) union \(R_{FE}\) yields two label worlds with identical transcript and different carrier parity.
\item Unread-enable-bit falsifier: a selector whose enabling bit is never surfaced can switch a row into the influence cone without changing the charged transcript.
\item Threshold-collapse falsifier: if q+eta+tau, \(U_{FE}\), \(X_{FE}\), or \(R_{FE}\) reaches 1+epsilon, the parser theorem may remain valid but no barely-superlinear lower-bound contradiction follows.
\item Model-mismatch falsifier: a proof that treats semantic SAT implications as occurrence support without a same-theory realizer leaks nonuniform advice into the source lower-bound certificate.
\end{itemize}

\subsection*{cycle-0119: Hardness-Magnification Bridge Realization, cycle 119: compositional transparent support closure}

\paragraph{Execution.} research; azure\_responses\_api; attempt primary; 1005.976 seconds; 3 lemmas, 5 equations, 11 derivation steps, 5 falsifiers.

\paragraph{Learning.} Cycle 119 turns compositional closure from a warning into a theorem-shaped invariant: transparent FormulaEval support must be prefix-closed under later symbolic use, and subthreshold \(\Gamma_{comp}\) then gives a same-model \(\mathrm{LB}_{N}\) falsifier. The necessity direction is a last-unsafe-line private-row theorem for extensional support. Full semantic FormulaEval remains open exactly at \(RealSATInfluenceRealizer_{N}\), not at generic barrier folklore or the main P vs NP conjecture.

\paragraph{Pressure.} The derivation does not prove full \(\mathrm{Psys}_{sym}\) semantic FormulaEval parsing; it is blocked exactly at \(RealSATInfluenceRealizer_{N}\) for rules whose influence is forced by real SAT semantics but not surfaced as generated-row support.

\paragraph{Derivation steps.}
\begin{enumerate}
\item Live source\_statement: C\_SIZE(H\_N\textasciicircum{}SAT)>N\textasciicircum{}\{1+epsilon-o(1)\} for the beta-zero parity carrier is the lower-bound statement being audited, not assumed true.
\item Live target\_statement: a support-closed transparent proof pi with \(N^{\rho+o(1)}+\Gamma_{comp}(pi)+E_{pi}\) \(N<N^{1+\epsilon-o(1)}\) yields an explicit same-model circuit for \(H_{N}^\mathrm{SAT}\), hence is a valid falsifier for that \(\mathrm{LB}_{N}\) claim.
\item \(Model_{context}\): source carrier \(H_{N}^\mathrm{SAT}\) uses generated polylog SAT formulas; target class is nonuniform SIZE; proof model is restricted \(\mathrm{Psys}_{sym}^{ext}\); reduction type is parser-to-circuit simulation plus equality patching, not a proof of \(\mathrm{P} \ne \mathrm{NP}\).
\item \(Quantifier_{context}\): N ranges over valid carrier lengths, \(\epsilon>0\) is fixed, q and \(\rho\) are fixed below 1+epsilon, \(\beta=0\), parser uniformity is advice-free, and pi may be nonuniform only as the alleged proof/certificate object being audited.
\item Define \(Cl_{t}\), the prefix row-influence closure before line t. Initialize it from bit-explicit EF dependencies, SATeq links, \(bounded-SR_{h}\) leaves, and declared extension supports already covered by \(\Gamma_{bEF}\).
\item For a transparent FormulaEval line, add either the enumerated \(RowSet_{l}\) with cost \(U_{l}\) and occurrence expansion \(X_{l}\), or the full-scan support over all \(M_{N}\) rows with cost \(N^{1+q+o(1)}\); propagate this set to every later line that reads the symbolic value.
\item For an admitted semantic line, require a \(RealSATInfluenceRealizer_{N}\) object before the line is usable; add its generated row indices, occurrence neighborhoods, and verifier overhead to \(R_{N}\) and to \(Cl_{t}\).
\item Induct on proof lines: given \(Cl_{t}\) and the prefix circuit of size Gamma\_prefix(t), bit-explicit EF and \(SR_{h}\) rules preserve locality by the cycle-102 normal form; transparent FormulaEval preserves locality by explicit RowSet expansion; inherited closures prevent later hidden support reuse.
\item At the final lower-bound line, assemble a circuit for \(H_{N}^\mathrm{SAT}\) from the generator, the charged transcript circuit, and equality patches for \(E_{pi}\) exceptional source rows; its size is at most \(N^{\rho+o(1)}+\Gamma_{comp}(pi)+E_{pi}\) N.
\item Contrapositive: if pi is presented as a proof of LB\_N(1+epsilon) while the displayed cost is subthreshold, the constructed circuit contradicts the \(\mathrm{LB}_{N}\) statement in the same SIZE model; this is a route-level falsifier theorem, not evidence for \(\mathrm{P} \ne \mathrm{NP}\).
\item Reviewer findings: reduction-consistency passes for \(\mathrm{Psys}_{sym}^{ext}\) because source and target circuits are the same nonuniform model; quantifier-uniformity passes only with advice-free parser clauses; barrier-specificity says extensional hidden support is relativizing and refuted by private rows, while \(RealSATInfluenceRealizer_{N}\) is the named nonrelativizing open obligation.
\end{enumerate}

\paragraph{Equation 1: carrier\_batch\_size.}
\[
M_N=N^{q+o(1)}
\]
\emph{Valid when:} Valid when the beta-zero parity carrier has a fixed batched generator family with polynomially many source rows and q is measured on valid N lengths.

\paragraph{Equation 2: compositional\_cost.}
\[
Gamma_comp(pi)=Gamma_bEF(pi)+U_N(pi)+X_N(pi)+R_N(pi)
\]
\emph{Valid when:} Applies after every proof line is parsed with prefix-closed support inheritance and all FormulaEval or semantic row influences are surfaced and charged.

\paragraph{Equation 3: parser\_upper\_bound.}
\[
C_SIZE(H_N^SAT) \le N^{rho+o(1)}+Gamma_comp(pi)+E_pi N
\]
\emph{Valid when:} Requires the support-closed transcript, same-theory expansion, generator size \(N^{\rho+o(1)}\), and equality patching for at most \(E_{pi}\) exceptional source rows.

\paragraph{Equation 4: subthreshold\_falsifier\_window.}
\[
N^{rho+o(1)}+Gamma_comp(pi)+E_pi N < N^{1+epsilon-o(1)} => not LB_N(1+epsilon)
\]
\emph{Valid when:} Valid when \(\mathrm{LB}_{N}\) is a nonuniform SIZE lower-bound claim for the same \(H_{N}^\mathrm{SAT}\) and the constructed circuit family is advice-free except for pi.

\paragraph{Equation 5: support\_density\_pressure.}
\[
|K_idx(pi)|/M_N \ge N^{epsilon-q-o(1)}
\]
\emph{Valid when:} Applies to any valid \(\mathrm{LB}_{N}\) proof in this parser model after non-row costs are subthreshold and charged complete rows are the remaining support bottleneck.

\paragraph{Next analysis.} Write the finite prefix-closure verifier for \(TransparentRowSupport_{N}\): input pi,N; output \(Cl_{t}\) for each line, Gamma\_prefix(t), and either a valid support-closed derivation certificate or the first unsafe line with its private-row/unread-enable-bit countermodel.

\paragraph{Falsifiers.}
\begin{itemize}
\item Extensional hidden support falsifier: a private row absent from \(Cl_{t}\) is flipped while all visible transcript bits are fixed.
\item Unread-enable-bit falsifier: a symbolic FormulaEval gate contains an unenumerated guard bit whose later use changes the accepted transcript without paying \(U_{N}\) or \(X_{N}\).
\item Subthreshold transcript falsifier: \(N^{\rho+o(1)}+\Gamma_{comp}(pi)+E_{pi}\) \(N<N^{1+\epsilon-o(1)}\) constructs a same-model circuit contradicting LB\_N(1+epsilon).
\item \(q<\epsilon\) density falsifier: the required relative support \(|K_{idx}|/M_{N}\) exceeds one when all non-row costs are already subthreshold.
\item Semantic-oracle falsifier: a FormulaEval-SAT rule whose accepted behavior changes with a hidden generated row but supplies no \(RealSATInfluenceRealizer_{N}\) lies outside the transparent parser theorem.
\end{itemize}

\subsection*{cycle-0120: Hardness-Magnification Bridge Realization: least-closed transparent FormulaEval support}

\paragraph{Execution.} research; azure\_responses\_api; attempt primary; 167.458 seconds; 3 lemmas, 5 equations, 8 derivation steps, 5 falsifiers.

\paragraph{Learning.} Cycle 120 strengthens the promoted theorem pair by identifying the minimal valid charged object: the least prefix-closed reachable FormulaEval support cone K\_dag(pi), not the full selector zoo. The positive line induction uses Gamma\_dag=Gamma\_bEF+U(K\_dag)+X(K\_dag), and the necessity side is a first-hidden-reachable-row falsifier. Full semantic FormulaEval remains quarantined exactly at \(RealSATInfluenceRealizer_{N}\); this is a scoped parser/falsifier contribution, not progress toward proving \(\mathrm{P} \ne \mathrm{NP}\).

\paragraph{Pressure.} Full semantic FormulaEval over real SAT is still blocked at \(RealSATInfluenceRealizer_{N}\): no proof currently converts semantic influence of an uncharged formula row into an advice-free, prefix-closed, charged support certificate.

\paragraph{Derivation steps.}
\begin{enumerate}
\item Define the live inference source statement: \(\mathrm{Psys}_{sym}^{dag}\) proves LB\_N(1+epsilon) for \(H_{N}^\mathrm{SAT}\) with \(\beta=0\), \(q,\rho<1+\epsilon\), bit-explicit bEF support, \(bounded-SR_{h}\) normal form, and FormulaEval dependency DAGs.
\item Define the live inference target statement: if pi has N\textasciicircum{}\{rho+o(1)\}+Gamma\_dag(pi)+O(E\_pi N)<N\textasciicircum{}\{1+epsilon-o(1)\}, then the same nonuniform SIZE model has circuits for \(H_{N}^\mathrm{SAT}\) below the claimed lower-bound threshold; hence that pi is a falsifier for LB\_N(1+epsilon), not evidence for \(\mathrm{P} \ne \mathrm{NP}\).
\item Build K\_dag(pi) as a least fixed point: seed with FormulaEval tokens used by the final \(\mathrm{LB}_{N}\) line, close under explicit parent pointers, row enumerator certificates, \(SR_{h}\) leaf rows, and symbolic aliases that feed any retained line.
\item Run line induction only on the retained cone. Bit-explicit bEF lines use the existing \(\Gamma_{bEF}\) charge; \(bounded-SR_{h}\) lines add eta/tau-normalized row costs; transparent FormulaEval lines add U(K)+X(K). Unreachable FormulaEval rows are ignored by \(UnreachableSupportDeletion_{N}\).
\item Substitute a hypothetical SAT circuit of size \(L_{N}^{c}\) into each generated formula. Because \(L_{N}=polylog\) N, every fixed c contributes \(N^{o(1)}\) per surfaced row; generator plus transcript costs give C(H\_N)<=N\textasciicircum{}\{rho+o(1)\}+Gamma\_dag(pi)+O(E\_pi N).
\item Derive the density corollary: any valid proof certificate for LB\_N(1+epsilon) inside this parser model must force at least \(N^{\epsilon-o(1)}\) reachable complete parity rows, equivalently density \(N^{\epsilon-q-o(1)}\) among \(M_{N}=N^{q+o(1)}\) slots when q is fixed.
\item Prove necessity for the extensional case by choosing the earliest reachable uncharged influence. Flip its private row or unread enable bit while keeping charged transcript data fixed; the verifier sees the same prior cone but the final induced parity differs, contradicting sound surface occurrence completion.
\item Reviewer packet status: reduction consistency is valid for \(\beta=0\) same-model substitution; quantifier uniformity is fixed over q,rho,epsilon,h,d; barrier specificity says the countermodel relativizes only for extensional/oracle FormulaEval, while the real SAT escape is the named nonrelativizing \(RealSATInfluenceRealizer_{N}\).
\end{enumerate}

\paragraph{Equation 1: carrier.}
\[
H_N^{SAT}(x)=\bigoplus_{j\le M_N} SAT(G_N(x,j)),\quad M_N=N^{q+o(1)}
\]
\emph{Valid when:} Valid when the source object is the beta-zero generated parity carrier with uniformly computable generator slots.

\paragraph{Equation 2: least\_closed\_cost.}
\[
\Gamma_{dag}(\pi)=\Gamma_{bEF}(\pi)+U(K_{dag}(\pi))+X(K_{dag}(\pi))+O(|\pi|\,\mathrm{polylog}\,N)
\]
\emph{Valid when:} Valid when \(\mathrm{Psys}_{sym}^{dag}\) exposes parent pointers, dependency edges, and advice-free row enumeration certificates.

\paragraph{Equation 3: falsifier\_upper\_bound.}
\[
C(H_N^{SAT})\le N^{\rho+o(1)}+\Gamma_{dag}(\pi)+O(E_\pi N)
\]
\emph{Valid when:} Applies after the line induction preserves same-model circuits and sparse row errors are equality patched.

\paragraph{Equation 4: subthreshold\_negation.}
\[
N^{\rho+o(1)}+\Gamma_{dag}(\pi)+O(E_\pi N)< N^{1+\epsilon-o(1)}\Rightarrow \neg LB_N(1+\epsilon)
\]
\emph{Valid when:} Requires the target-preserving same-model transfer and the claimed \(\mathrm{LB}_{N}\) threshold in nonuniform SIZE.

\paragraph{Equation 5: support\_density\_corollary.}
\[
|K_{dag}(\pi)|\ge N^{\epsilon-o(1)},\quad |K_{dag}(\pi)|/M_N\ge N^{\epsilon-q-o(1)}
\]
\emph{Valid when:} Valid when each complete surfaced parity row costs \(N^{1+o(1)}\) to realize and \(\rho\) is subthreshold.

\paragraph{Next analysis.} Write the proof-audit version of L120.1: formalize the fixed-point computation of K\_dag(pi), give pseudocode for the verifier-visible dependency closure, and check one projected-dDNNF FormulaEval example plus one unread-enable-bit counterexample against the \(\Gamma_{dag}\) inequality.

\paragraph{Falsifiers.}
\begin{itemize}
\item Reachable hidden support falsifier: an uncharged row in the final dependency cone permits two extensional transcripts agreeing on all charged data while differing on the induced parity carrier.
\item Unread-enable-bit falsifier: a selector bit that can activate a row after its enumerator certificate is checked invalidates \(\Gamma_{dag}\) unless the activation edge is prefix-closed and charged.
\item Subthreshold transcript falsifier: N\textasciicircum{}\{rho+o(1)\}+Gamma\_dag(pi)+O(E\_pi N)<N\textasciicircum{}\{1+epsilon-o(1)\} gives an explicit small circuit for \(H_{N}^\mathrm{SAT}\) in the same model, refuting the source \(\mathrm{LB}_{N}\) certificate.
\item \(q<\epsilon\) support-density falsifier: if fewer than \(N^{\epsilon-o(1)}\) reachable complete rows are forced among \(M_{N}=N^{q+o(1)}\) slots, equality patching and the transcript circuit stay below the claimed threshold.
\item Semantic-realizer falsifier probe: if a claimed full FormulaEval rule influences a proof without producing \(RealSATInfluenceRealizer_{N}\) rows, the branch treats it as outside the valid parser theorem rather than as transparent support.
\end{itemize}

\subsection*{cycle-0121: Hardness-Magnification Bridge Realization: L121 least-closed support audit}

\paragraph{Execution.} research; azure\_responses\_api; attempt primary; 157.908 seconds; 3 lemmas, 5 equations, 10 derivation steps, 6 falsifiers.

\paragraph{Learning.} Cycle 121 turns K\_dag(pi) into a finite auditable fixed-point object. The valid contribution is sharper than the previous transparency slogan: least prefix-closed reachable support is sufficient for the \(\Gamma_{dag}\) parser and necessary for extensional safety via a first-hidden-row falsifier. The remaining full-system gap is exactly \(RealSATInfluenceRealizer_{N}\) for real SAT-semantic FormulaEval, not generic relativization, natural proofs, or P vs NP folklore.

\paragraph{Pressure.} The restricted proof is complete only for transparent extensional FormulaEval. Full semantic SAT FormulaEval remains blocked at \(RealSATInfluenceRealizer_{N}\): proof syntax must nonrelativizingly determine and charge every real SAT influence before use.

\paragraph{Derivation steps.}
\begin{enumerate}
\item Normalize pi into a finite DAG of lines with IDs, parent lists, and rule tags; reject lines whose parents are implicit semantic advice.
\item Initialize the active line set A with the final \(\mathrm{LB}_{N}\) line and initialize \(K_{idx}\), \(K_{occ}\), U, X to empty ledgers.
\item Process a worklist: when a line is popped, add all proof parents to A; for a bEF line add its bit-explicit dependency charge to \(\Gamma_{bEF}\).
\item For each reachable FormulaEval line, run its visible support enumerator or certified full scan; add each row identifier, occurrence-neighborhood, and parent support certificate to \(K_{dag}\).
\item Close under later symbolic use: if a reachable line reads a FormulaEval value, add the row certificate and all support parents before the line is counted valid.
\item Terminate because each step adds a proof line, support certificate, or row occurrence from the finite transcript; the fixed point is least by monotonicity of the closure operator.
\item Prove the line-induction invariant: every reachable line has a same-theory circuit/transcript of size at most the accumulated prefix of \(\Gamma_{bEF}+U+X\), with exceptional rows recorded in \(E_{pi}\).
\item At the final \(\mathrm{LB}_{N}\) line, compose the generator for \(G_{N}\) with the parsed transcript and equality patches, obtaining C(H\_N\textasciicircum{}SAT)<=N\textasciicircum{}\{rho+o(1)\}+Gamma\_dag+O(E\_pi N).
\item Derive the falsifier inequality: a proof of LB\_N(1+epsilon) with subthreshold \(\Gamma_{dag}\) contradicts the same-model circuit upper bound, so any valid proof must pay the support-density threshold.
\item For minimality, remove a reachable influencing row and choose two row-label worlds agreeing on \(K_{dag}\) but differing there; the first line whose value changes gives the private-row/unread-enable-bit countermodel.
\end{enumerate}

\paragraph{Equation 1: fixed\_point\_closure.}
\[
K_{0}={\ell_{final}},\quad K_{t+1}=K_t\cup Par(K_t)\cup Supp_{FE}(K_t)\cup Infl(K_t),\quad K_{dag}=\bigcup_t K_t.
\]
\emph{Valid when:} Valid when \(\mathrm{Psys}_{sym}^{dag}\) exposes parent pointers and FormulaEval support certificates syntactically.

\paragraph{Equation 2: gamma\_dag\_ledger.}
\[
\Gamma_{dag}(\pi)=\Gamma_{bEF}(\pi)+U_N(K_{dag}(\pi))+X_N(K_{dag}(\pi)).
\]
\emph{Valid when:} Applies after the closure algorithm has surfaced every reachable FormulaEval influence.

\paragraph{Equation 3: parser\_upper\_bound.}
\[
C(H_N^{SAT})\le N^{\rho+o(1)}+\Gamma_{dag}(\pi)+O(E_\pi N).
\]
\emph{Valid when:} Requires \(\beta\) zero, same-model row expansion, and equality patching for exceptional rows.

\paragraph{Equation 4: lb\_falsifier\_window.}
\[
N^{\rho+o(1)}+\Gamma_{dag}(\pi)+O(E_\pi N)< N^{1+\epsilon-o(1)}\Rightarrow \neg LB_N(1+\epsilon).
\]
\emph{Valid when:} Valid when \(\mathrm{LB}_{N}\) is the lower-bound claim for the same carrier and circuit model.

\paragraph{Equation 5: support\_density\_corollary.}
\[
|K_{idx}(\pi)|\ge N^{\epsilon-o(1)},\quad |K_{idx}(\pi)|/M_N\ge N^{\epsilon-q-o(1)}.
\]
\emph{Valid when:} Applies when each charged FormulaEval slot has expansion cost \(\widetilde{O}(N)\).

\paragraph{Next analysis.} Write the formal L121.1 proof skeleton as theorem text: fixed-point pseudocode for \(K_{dag}\), induction invariant per rule type, charge table for \(bEF/SR_{h}/\mathrm{FormulaEval}/projected-dDNNF\), and the unread-enable-bit counterexample with \(\theta<\epsilon\).

\paragraph{Falsifiers.}
\begin{itemize}
\item first-hidden-reachable-row falsifier: a later proof line depends on a row absent from \(K_{dag}\)
\item unread-enable-bit falsifier: selector output is disabled in syntax but semantically gates an uncharged row
\item \(q<\epsilon\) enumeration falsifier: subthreshold reachable support density yields an immediate small circuit for the claimed \(\mathrm{LB}_{N}\)
\item semantic-advice falsifier: RowSet(a) is chosen using SAT truth or nonuniform advice rather than proof syntax
\item same-theory failure falsifier: row expansion is asserted externally and cannot be represented inside \(\mathrm{Psys}_{sym}^{dag}\)
\item positive-beta transfer falsifier remains quarantined: q+beta c caps all-exponent target transfer unless \(\beta=0\) or subpolynomial
\end{itemize}

\subsection*{cycle-0122: Hardness-Magnification Bridge Realization}

\paragraph{Execution.} research; azure\_responses\_api; attempt primary; 120.494 seconds; 3 lemmas, 6 equations, 10 derivation steps, 6 falsifiers.

\paragraph{Learning.} Cycle 122 upgrades \(K_{dag}\) from a named minimal support object to a referee-auditable fixed-point proof packet. The valid theorem is restricted but paper-relevant: least prefix-closed reachable FormulaEval support is sufficient for \(\Gamma_{dag}\) line induction and necessary for extensional safety by a first-hidden-row falsifier. The full semantic case is still blocked exactly at \(RealSATInfluenceRealizer_{N}\), not by generic barrier folklore or by the P vs NP conjecture itself.

\paragraph{Pressure.} The restricted theorem is valid for transparent extensional FormulaEval, but it does not settle full semantic FormulaEval. The named gap is \(RealSATInfluenceRealizer_{N}\): real SAT-semantic reasoning may determine row influence without a verifier-visible row set, and this requires a nonrelativizing charged realizer or a rule-specific no-go.

\paragraph{Derivation steps.}
\begin{enumerate}
\item Parse pi line by line into ordinary bEF lines and FormulaEval cells (a,r). For each FormulaEval cell record opcode, row index r, parent-cell pointers, selector code, expansion code, and an enumerator or full-scan certificate.
\item Build the backward dependency graph \(D_{pi}\) whose edges go from a used cell to its verifier-visible parent cells and to row occurrences that must be expanded in the same theory.
\item Initialize \(K_{0}\) by the final \(\mathrm{LB}_{N}\) assertion and the carrier rows whose parity is read by the final falsifier construction. This seed is syntactic; it contains no SAT-truth advice.
\item Iterate \(K_{t+1}=K_{t}\) union BackParents(K\_t) union RowSupp(K\_t). Termination is bounded by the finite set of declared cells and rows, so \(K_{dag}\) is auditor-computable.
\item Define U\_N(K\_dag) as the sum of enumerator or full-scan costs needed to list reachable rows, and X\_N(K\_dag) as the same-theory occurrence-expansion cost for exactly those rows.
\item Run the existing \(\Gamma_{bEF}\) line induction for ordinary EF, SATeq, and \(bounded-SR_{h}\) lines. For FormulaEval lines, use closure of \(K_{dag}\) to guarantee that every later symbolic use already has its parent rows and expansions charged.
\item Construct circuits for all reachable proof-line predicates with total size Gamma\_bEF(pi)+U\_N(K\_dag)+X\_N(K\_dag), then compose with the generator circuit of size \(N^{\rho+o(1)}\).
\item If the induced circuit is wrong on \(E_{pi}\) exceptional source inputs, patch them by equality circuits at cost O(E\_pi N), giving the displayed upper bound for \(H_{N}^\mathrm{SAT}\).
\item Derive the q-trichotomy: subthreshold \(\Gamma_{dag}\) forces \(|K_{idx}|\ge N^{\epsilon-o(1)}\); hence qepsilon leaves weak but absolute support pressure.
\item For minimality, choose the first downstream use of an uncharged reachable row and construct two row-label worlds that agree on all visible fields but differ on that row, producing the private-row or unread-enable-bit contradiction.
\end{enumerate}

\paragraph{Equation 1: fixed\_point\_definition.}
\[
K_0=\operatorname{Seed}(\operatorname{FinalLB}(\pi)),\quad K_{t+1}=K_t\cup\operatorname{BackParents}(K_t)\cup\operatorname{RowSupp}(K_t),\quad K_{\dagger}=\operatorname{lfp}(K_t)
\]
\emph{Valid when:} Valid when every FormulaEval use exposes verifier-readable parent maps and row support certificates.

\paragraph{Equation 2: finite\_convergence\_bound.}
\[
T_{\dagger}\le |\operatorname{Cells}(\pi)|+|\operatorname{RowsDeclared}(\pi)|
\]
\emph{Valid when:} Applies when the dependency graph is finite and rows are explicitly indexed.

\paragraph{Equation 3: charged\_gamma.}
\[
\Gamma_{\dagger}(\pi)=\Gamma_{bEF}(\pi)+U_N(K_{\dagger})+X_N(K_{\dagger})
\]
\emph{Valid when:} Valid when bEF charges are already parsed and FormulaEval support is transparent.

\paragraph{Equation 4: parser\_to\_circuit\_falsifier.}
\[
C(H_N^{SAT})\le N^{\rho+o(1)}+\Gamma_{\dagger}(\pi)+O(E_{\pi}N)
\]
\emph{Valid when:} Requires the generator, parsed proof transcript, and sparse residual patches to be same-model circuits.

\paragraph{Equation 5: subthreshold\_contradiction\_window.}
\[
N^{\rho+o(1)}+\Gamma_{\dagger}(\pi)+O(E_{\pi}N)< N^{1+\epsilon-o(1)}\Rightarrow \neg LB_N(1+\epsilon)
\]
\emph{Valid when:} Valid when \(\mathrm{LB}_{N}\) is a lower bound against the same nonuniform SIZE model.

\paragraph{Equation 6: support\_density\_corollary.}
\[
|K_{idx}(\pi)|\ge N^{\epsilon-o(1)},\qquad |K_{idx}(\pi)|/M_N\ge N^{\epsilon-q-o(1)}
\]
\emph{Valid when:} Applies after each charged row costs at most \(N^{1+o(1)}\) to patch.

\paragraph{Next analysis.} Write the paper-ready L122 module: formal theorem statement, fixed-point pseudocode, line-induction proof, projected-dDNNF charge example, unread-enable-bit counterexample, q-trichotomy corollary, and a separate \(RealSATInfluenceRealizer_{N}\) open-boundary paragraph.

\paragraph{Falsifiers.}
\begin{itemize}
\item first-hidden-reachable-row falsifier: a row affects the final \(\mathrm{LB}_{N}\) line but is absent from \(K_{dag}\) and all charges, so two visible-identical transcripts disagree on the carrier.
\item unread-enable-bit falsifier: a selector enable controls row reachability while not appearing in parent pointers, invalidating local transparency despite an apparently valid row list.
\item \(q<\epsilon\) enumeration falsifier: \(K_{idx}\) must have absolute size \(N^{\epsilon-o(1)}\) but the universe has only \(M_{N}=N^{q+o(1)}\) rows.
\item semantic-advice leakage: \(Enum_{a}\) or RowSupp(a,r) depends on SAT truth or nonuniform row advice rather than proof syntax and N.
\item positive-beta transfer falsifier: any black-box target-circuit substitution with query length \(n^\beta\) gives the cap q+beta c+o(1), so it cannot yield all-exponent target exclusion unless \(\beta=0\) or subpolynomial.
\item nonprefix local support falsifier: a row is surfaced where defined but hidden after a later symbolic use, breaking the line-induction invariant.
\end{itemize}

\subsection*{cycle-0123: Hardness-Magnification Bridge Realization}

\paragraph{Execution.} research; azure\_responses\_api; attempt primary; 435.479 seconds; 4 lemmas, 7 equations, 9 derivation steps, 7 falsifiers.

\paragraph{Learning.} Cycle 123 closes a proof-audit gap rather than repeating the magnification slogan: \(K_{dag}\) is an implementable least fixed point, projected-dDNNF gives a concrete transparent positive case, and the unread-enable-bit/private-row construction gives a precise necessity falsifier. The valid result is a scoped parser/falsifier theorem for beta-zero parity magnification; full semantic FormulaEval is still blocked exactly at \(RealSATInfluenceRealizer_{N}\).

\paragraph{Pressure.} \(RealSATInfluenceRealizer_{N}\) remains unproved for full semantic FormulaEval: a real SAT-semantic rule might force r* nonrelativizingly from charged rows even when no extensional support certificate lists it.

\paragraph{Derivation steps.}
\begin{enumerate}
\item Instantiate the source statement as LB\_N(1+epsilon) for \(H_{N}^\mathrm{SAT}\) in nonuniform SIZE, with \(M_{N}=N^{q+o(1)}\), L\_N=(log N)\textasciicircum{}d, \(\beta=0\), and generator overhead \(N^{\rho+o(1)}\).
\item Parse all \(\mathrm{Psys}_{bEF}^{-}\) lines first; EF dependencies, SATeq links, and \(bounded-SR_{h}\) leaves contribute Gamma\_bEF(pi) exactly as in the restricted bit-explicit theorem packet.
\item Define row atoms for FormulaEval lines and run the verifier-visible fixed point: start from rows reachable from the final \(\mathrm{LB}_{N}\) line, repeatedly add parent rows, selector rows, same-theory expansion rows, and earlier rows used by later reachable lines.
\item Charge each closure addition either by U\_N(K\_dag), an advice-free enumerator/full-scan cost, or by X\_N(K\_dag), the N-per-row occurrence expansion; reject any row whose selector is only semantic advice.
\item For the projected-dDNNF example, gate pointers and projection maps give the enumerator. The reachable cone follows output-gate dependencies, so irrelevant dDNNF components are not charged.
\item Extend the line induction: every proof line whose reachable rows are in \(K_{dag}\) has a simulator circuit of cost bounded by the prior simulator plus the charged local rule cost; the final line yields C(H\_N\textasciicircum{}SAT)<=N\textasciicircum{}\{rho+o(1)\}+Gamma\_dag(pi)+O(E\_pi N).
\item Contrapose the lower-bound claim: if this upper bound is below \(N^{1+\epsilon-o(1)}\), the proof transcript is a same-model falsifier for LB\_N(1+epsilon).
\item Prove minimality for extensional FormulaEval by choosing the first later-used row r* omitted from \(K_{dag}\); two transcripts agree on charged rows and all syntax but differ on r*, so the parser cannot be valid for both.
\item Derive the q-trichotomy from \(X_{N}(K)=\)\(\widetilde{O}(N|K|)\): a valid \(\mathrm{LB}_{N}\) certificate must force \(|K_{dag}|\ge N^{\epsilon-o(1)}\), hence density at least \(N^{\epsilon-q-o(1)}\) inside \(M_{N}\) slots.
\end{enumerate}

\paragraph{Equation 1: carrier.}
\[
H_N^{SAT}(x)=\bigoplus_{j< M_N} SAT(G_N(x,j)),\quad M_N=N^{q+o(1)},\quad L_N=(\log N)^d
\]
\emph{Valid when:} Valid when the source carrier is the beta-zero parity carrier with uniformly generated polylogarithmic SAT instances.

\paragraph{Equation 2: least\_closure.}
\[
K_0=Reach(final),\quad K_{t+1}=K_t\cup Parents(K_t)\cup SupportRows(K_t),\quad K_{dag}(\pi)=\bigcup_t K_t
\]
\emph{Valid when:} Valid when every dependency edge and support certificate is visible in the proof syntax and uses no semantic row advice.

\paragraph{Equation 3: charged\_gamma.}
\[
\Gamma_{dag}(\pi)=\Gamma_{bEF}(\pi)+U_N(K_{dag})+X_N(K_{dag})
\]
\emph{Valid when:} Applies after the fixed-point closure terminates on the finite proof DAG and all transparent rows are charged.

\paragraph{Equation 4: row\_expansion.}
\[
X_N(K_{dag})\le \widetilde{O}(N\cdot |K_{dag}|)
\]
\emph{Valid when:} Requires same-theory occurrence expansion where each surfaced row is materialized at N-bit source cost.

\paragraph{Equation 5: same\_model\_upper\_bound.}
\[
C_{SIZE}(H_N^{SAT})\le N^{\rho+o(1)}+\Gamma_{dag}(\pi)+O(E_\pi N)
\]
\emph{Valid when:} Applies after the reduction preserves the same nonuniform SIZE model, beta-zero query length, and charged support parameters.

\paragraph{Equation 6: lb\_falsifier\_window.}
\[
N^{\rho+o(1)}+\Gamma_{dag}(\pi)+O(E_\pi N)< N^{1+\epsilon-o(1)}\Rightarrow \neg LB_N(1+\epsilon)
\]
\emph{Valid when:} Valid when LB\_N(1+epsilon) denotes an asserted same-model lower bound above the displayed threshold.

\paragraph{Equation 7: support\_density\_corollary.}
\[
LB_N(1+\epsilon)\Rightarrow |K_{dag}|\ge N^{\epsilon-o(1)},\quad |K_{dag}|/M_N\ge N^{\epsilon-q-o(1)}
\]
\emph{Valid when:} Requires the row-expansion term to dominate uncharged support cost and \(M_{N}=N^{q+o(1)}\) slots.

\paragraph{Next analysis.} Write the referee-checkable trace for the projected-dDNNF example: list gate fields, parent pointers, projection maps, RowSet certificates, \(K_{0},K_{1}\),... closure additions, \(U_{N}\) and \(X_{N}\) charges, then flip one unread enable bit to exhibit the first-hidden-row countermodel and verify the \(\Gamma_{dag}\) threshold numerically for qepsilon.

\paragraph{Falsifiers.}
\begin{itemize}
\item first-hidden-reachable-row falsifier for extensional FormulaEval without charged support
\item unread-enable-bit falsifier for projected or guarded FormulaEval whose selector is not surfaced before later use
\item private-row two-world falsifier for symbolic FormulaEval that treats row influence as semantic oracle advice
\item \(q<\epsilon\) total-slot falsifier: \(M_{N}=N^{q+o(1)}\) cannot supply \(N^{\epsilon-o(1)}\) charged rows
\item positive-beta black-box substitution falsifier: \(\mathrm{SAT}\in \mathrm{SIZE}(n^{c})\) gives source upper bound N\textasciicircum{}\{q+beta c+o(1)\} before all-exponent exclusion
\item nonuniform row-selection leakage falsifier when RowSet(a) is supplied by advice rather than proof syntax
\item semantic \(RealSAT\) escape falsifier probe: if no advice-free influence realizer exists, full FormulaEval stays outside the valid parser theorem
\end{itemize}

\subsection*{cycle-0124: Hardness-Magnification Bridge Realization}

\paragraph{Execution.} research; azure\_responses\_api; attempt primary; 134.362 seconds; 3 lemmas, 6 equations, 6 derivation steps, 6 falsifiers.

\paragraph{Learning.} Cycle 124 makes the semantic boundary sharper: bounded Tseitin self-reduction is foldable into the least-support theorem only when every restriction child is declared and charged. The exact threshold is \(\kappa+\theta<\epsilon\) for the new row expansion. Hidden branch-pruning is a concrete private-row falsifier, so full semantic FormulaEval remains blocked at \(RealSATInfluenceRealizer_{N}^{arb}\) rather than at generic barrier folklore.

\paragraph{Pressure.} The proof does not cover arbitrary full semantic FormulaEval. The named remaining gap is \(RealSATInfluenceRealizer_{N}^{arb}\): every semantic inference must uniformly reveal all rows whose labels can affect later proof use before the inference is consumed.

\paragraph{Derivation steps.}
\begin{enumerate}
\item Fix \(\mathrm{Psys}_{sym}^{dag+\mathrm{TSR}_{decl}}\). Each proof line exposes a type tag: bEF, transparent FormulaEval, or \(\mathrm{TSR}_{decl}\). A \(\mathrm{TSR}_{decl}\) line contains row id a, formula/Tseitin code \(T_{a}\), pivot list \(p_{1},...,p_{h}\), and a uniform constructor for each restricted child \(T_{a,b}\).
\item Define K\_TSR(pi) by fixed-point iteration. Initialize with rows reachable from the final \(\mathrm{LB}_{N}\) line. Add proof-DAG parents. For a transparent FormulaEval line, add its surfaced support as in L120. For a \(\mathrm{TSR}_{decl}\) line, add all children child(a,b) for b in \({0,1}^{h(a)}\) and the EF-checkable Tseitin equivalence rows.
\item Termination and auditability are finite because the proof DAG and every declared child list are finite; the parser does not ask whether a child is satisfiable. Enumeration cost is U(K\_TSR)<=|pi|L\_N\textasciicircum{}\{O(1)\}+|K\_base|2\textasciicircum{}\{h\_N\}L\_N\textasciicircum{}\{O(1)\} before row expansion.
\item Extend the L120 line-induction invariant. bEF and transparent lines use the previous proof. For a \(\mathrm{TSR}_{decl}\) line, the same-theory EF derivation justifies \(\mathrm{SAT}(T_{a})=OR_{b}\) SAT(T\_\{a,b\}); all child row truth values are already in \(K_\mathrm{TSR}\), so the semantic line has no hidden influence.
\item The resulting transcript circuit has size at most N\textasciicircum{}\{rho+o(1)\}+Gamma\_bEF+U(K\_TSR)+X(K\_TSR)+R\_TSR+O(E\_pi N). If this is below \(N^{1+\epsilon-o(1)}\), it is a valid same-model falsifier for LB\_N(1+epsilon).
\item For a branch-pruning variant where the rule names only one child chosen by SAT truth, repeat the private-row construction: two completions agree on all surfaced rows and proof syntax but disagree on the hidden child and on one carrier parity bit. Thus no advice-free parser can charge only visible support.
\end{enumerate}

\paragraph{Equation 1: carrier.}
\[
H_N^{SAT}(x)=\bigoplus_{j< M_N} SAT(G_N(x,j)),\quad M_N=N^{q+o(1)},\quad |G_N(x,j)|=L_N=(\log N)^{d+O(1)}
\]
\emph{Valid when:} Valid when the beta-zero parity carrier and generator length regime are fixed before parsing proofs.

\paragraph{Equation 2: tsr\_closure.}
\[
K_{TSR}(\pi)=\mu K\,[Reach(final)\cup Cl_{dag}(K)\cup {child(a,b):a\in K,\ b\in{0,1}^{h(a)}}]
\]
\emph{Valid when:} Applies when every TSR line exposes its parent row and all restriction-child row constructors syntactically.

\paragraph{Equation 3: gamma\_charge.}
\[
\Gamma_{TSR}(\pi)=\Gamma_{bEF}(\pi)+U(K_{TSR})+X(K_{TSR})+R_{TSR},\quad X(K_{TSR})\le N\,|K_{TSR}|\,L_N^{O(1)}
\]
\emph{Valid when:} Requires the charged support cone to be uniformly enumerable and expanded in the same SAT row theory.

\paragraph{Equation 4: height\_threshold.}
\[
|K_{base}|=N^{\kappa+o(1)},\ 2^{h_N}=N^{\theta+o(1)}\ \Longrightarrow\ \Gamma_{TSR}\le \Gamma_{bEF}+N^{1+\kappa+\theta+o(1)}
\]
\emph{Valid when:} Valid when TSR closure is the only new source of FormulaEval rows beyond the base reachable cone.

\paragraph{Equation 5: same\_model\_falsifier.}
\[
C_{SIZE}(H_N^{SAT})\le N^{\rho+o(1)}+\Gamma_{TSR}(\pi)+O(E_\pi N)
\]
\emph{Valid when:} Applies after the parser line induction converts every charged proof dependency into a same-model circuit component.

\paragraph{Equation 6: subthreshold\_contradiction.}
\[
N^{\rho+o(1)}+\Gamma_{TSR}(\pi)+O(E_\pi N)< N^{1+\epsilon-o(1)}\ \Longrightarrow\ \neg LB_N(1+\epsilon)
\]
\emph{Valid when:} Requires the source lower-bound statement to be exactly the same nonuniform SIZE lower bound for \(H_{N}\).

\paragraph{Next analysis.} Write the \(FE_{\mathrm{TSR},decl}\) rule schema as bit fields: row id, Tseitin code, pivot list, child constructor, parent pointers, and EF derivation handles; then check the fixed-point \(K_\mathrm{TSR}\) enumerator on a one-line bounded-SR example and a hidden satisfying-branch counterexample.

\paragraph{Falsifiers.}
\begin{itemize}
\item Hidden branch bit: a self-reduction line chooses b*(a) from SAT truth without exposing both children; the private-row falsifier keeps the transcript fixed while changing \(H_{N}\) on one source input.
\item Height threshold collapse: if \(|K_{base}|=N^{\kappa+o(1)}\) and 2\textasciicircum{}\{h\_N\}=N\textasciicircum{}\{theta+o(1)\} with \(\kappa+\theta\ge \epsilon\), the row expansion \(N|K_\mathrm{TSR}|\) is not subthreshold.
\item Same-theory failure: a restricted child formula cannot be represented as a Psys row with EF-checkable Tseitin equivalence, so the line induction is invalid.
\item Advice leakage: the child-row constructor or reachable support list contains nonuniform SAT-truth advice rather than visible proof syntax.
\item Positive \(\beta\) regression: if the carrier interface reintroduces target query length \(n^\beta\) with \(\beta>0\), the old q+beta c substitution falsifier caps all-exponent transfer.
\item Naturalness status falsifier: any dense label-functional certificate over \(K_\mathrm{TSR}\) must still pass the earlier functional naturalization audit; declared TSR only handles semantic influence, not natural proofs.
\end{itemize}

\subsection*{cycle-0125: Hardness-Magnification Bridge Realization}

\paragraph{Execution.} research; azure\_responses\_api; attempt primary; 108.242 seconds; 3 lemmas, 4 equations, 7 derivation steps, 5 falsifiers.

\paragraph{Learning.} Cycle 125 converts the TSR semantic boundary into a theorem-shaped split. Declared bounded TSR is validly absorbed by the least-support parser with charge \(N^{1+\kappa+\theta+o(1)}\), giving a falsifier exactly when \(\kappa+\theta<\epsilon\). Hidden semantic branch-pruning is not a free extension of FormulaEval: it has a first-hidden-row private-row no-go unless a nonrelativizing \(RealSATInfluenceRealizer_{N}^{arb}\) supplies charged influence rows. This is a scoped parser/falsifier contribution, not evidence for \(\mathrm{P} \ne \mathrm{NP}\).

\paragraph{Pressure.} \(RealSATInfluenceRealizer_{N}^{arb}\) for hidden semantic branch-pruning: the derivation does not extract unlisted rows whose relevance is certified only by actual SAT semantics.

\paragraph{Derivation steps.}
\begin{enumerate}
\item Initialize K to the final \(\mathrm{LB}_{N}\) line, its bit-explicit \(\mathrm{Psys}_{bEF}^{-}\) dependencies, and the least reachable FormulaEval support from the existing \(K_{dag}\) fixed-point algorithm.
\item For each transparent FormulaEval line, add the enumerated row set or charge the full scan; add all later symbolic uses by prefix closure.
\item For each reachable \(\mathrm{TSR}_{decl}\) line, read its parent row id, branch-word field, child restricted-formula bit fields, child row ids, and occurrence-neighborhood certificates; add every child row to K before the TSR equality is evaluated.
\item Iterate until no newly added row can be used by a later proof line. Termination follows from the finite proof DAG and finite declared child DAGs, not from semantic SAT truth.
\item Run line induction: bEF lines use \(\Gamma_{bEF}\); transparent FormulaEval lines use U(K\_dag)+X(K\_dag); \(\mathrm{TSR}_{decl}\) lines use the declared children and the Boolean Tseitin self-reduction equality over charged SAT labels.
\item Build the falsifying circuit for \(H_{N}\) by composing the generator, the charged transcript evaluator, and equality patches for \(E_{pi}\) residual source rows. This gives C(H\_N)<=N\textasciicircum{}\{rho+o(1)\}+Gamma\_TSR+O(E\_pi N).
\item Apply the exponent audit: if \(\kappa+\theta<\epsilon\) and the previous \(\Gamma_{dag}\) and residual terms are subthreshold, the alleged LB\_N(1+epsilon) proof produces a smaller same-model circuit, so the proof certificate is invalid for that \(\mathrm{LB}_{N}\) claim.
\end{enumerate}

\paragraph{Equation 1: declared\_tsr\_closure.}
\[
K_{TSR}(\pi)=\operatorname{lfp}(K_{dag}(\pi)\cup Child_{TSR}^{decl}(K)),\qquad |K_{TSR}\setminus K_{dag}|\le T_N R_N=N^{\kappa+\theta+o(1)}.
\]
\emph{Valid when:} Valid when the TSR rule syntax declares every restriction child and the fixed-point operator follows only reachable proof dependencies.

\paragraph{Equation 2: gamma\_tsr\_charge.}
\[
\Gamma_{TSR}(\pi)=\Gamma_{dag}(\pi)+U(K_{TSR})+X(K_{TSR}\setminus K_{dag}),\qquad X(\Delta K)\le N\,|\Delta K|L_N^{O(1)}\le N^{1+\kappa+\theta+o(1)}.
\]
\emph{Valid when:} Valid when row enumeration is advice-free and every declared child is expanded in the same occurrence-level proof theory.

\paragraph{Equation 3: parser\_to\_circuit\_bound.}
\[
C(H_N^{SAT})\le N^{\rho+o(1)}+\Gamma_{TSR}(\pi)+O(E_\pi N).
\]
\emph{Valid when:} Applies after the line induction supplies a complete charged transcript for all proof values used by the final carrier circuit.

\paragraph{Equation 4: subthreshold\_falsifier\_window.}
\[
\rho< 1+\epsilon,\quad \Gamma_{dag}(\pi)< N^{1+\epsilon-o(1)},\quad E_\pi< N^{\epsilon-o(1)},\quad \kappa+\theta< \epsilon\ \Longrightarrow\ C(H_N^{SAT})< N^{1+\epsilon-o(1)}.
\]
\emph{Valid when:} Valid when all four exponent inequalities hold with a common positive margin along the same valid length subsequence.

\paragraph{Next analysis.} Write the referee-auditable \(\mathrm{TSR}_{decl}\) proof block: child-row bit fields, fixed-point pseudocode, \(\Gamma_\mathrm{TSR}\) induction case, and a two-world hidden-pruning countermodel with a named first hidden child row r\_star.

\paragraph{Falsifiers.}
\begin{itemize}
\item Declared TSR threshold collapse: \(\kappa+\theta\ge \epsilon\) makes the extra row expansion too large to refute LB\_N(1+epsilon).
\item First-hidden-branch falsifier: a pruned restriction child influences a later line but is absent from \(K_\mathrm{TSR}\), so two private-row worlds share the transcript and differ on the conclusion.
\item Unread-enable-bit falsifier: a selector bit can activate a row after the parser fixed point without being included in the visible support cone.
\item Nonuniform leakage falsifier: the row enumerator uses advice depending on SAT labels or on the final hardwired carrier truth table.
\item Positive-beta substitution falsifier: any return to target queries of length \(N^\beta\) obeys q+beta c and cannot exclude all SAT circuit exponents.
\end{itemize}

\subsection*{cycle-0126: Hardness-Magnification Bridge Realization}

\paragraph{Execution.} research; azure\_responses\_api; attempt primary; 112.689 seconds; 3 lemmas, 7 equations, 7 derivation steps, 5 falsifiers.

\paragraph{Learning.} Cycle 126 converts declared bounded TSR from a semantic warning into a charged fixed-point rule case. The valid contribution is a restricted theorem: declared TSR rows are absorbed by \(K_{dag}\) with threshold \(\kappa+\theta<\epsilon\), while hidden branch-pruning is a precise private-row falsifier. Full semantic FormulaEval is still not proved; it is isolated at \(RealSATInfluenceRealizer_{N}^{arb}\).

\paragraph{Pressure.} The derivation does not cover arbitrary semantic branch-pruning where real SAT semantics decide which children matter without a syntactic declaration. The named missing lemma is \(RealSATInfluenceRealizer_{N}^{arb}\): from proof syntax and N alone, uniformly list all semantically influencing TSR/FormulaEval rows before use and charge them in \(\Gamma\).

\paragraph{Derivation steps.}
\begin{enumerate}
\item Define proof-line nodes for bit-explicit EF, transparent FormulaEval, bounded \(SR_{h}\), and declared TSR. Each node stores row-index fields, parent pointers, local verifier bits, and a list of declared child rows when a TSR rule is used.
\item Initialize \(K_{0}\) with the row occurrences syntactically read by the final \(\mathrm{LB}_{N}\) line. Iteratively add visible FormulaEval parents, bEF extension supports, \(SR_{h}\) leaves, and declared TSR child rows. The proof DAG is finite, so the monotone closure reaches a least fixed point K\_dag\textasciicircum{}TSR(pi).
\item For each line i, prove an invariant: the parsed circuit \(C_{i}\) depends only on generator bits, previous parsed circuits, and the expanded rows in \(K_{i}\). bEF lines use the existing bit-explicit parser; FormulaEval lines use enumerator/full-scan charges; declared TSR lines compute the Tseitin restriction combination over their listed children.
\item Charge every declared TSR child row by N \(L_{N}^{O(1)}\) gates for occurrence expansion and local verification. With \(R_\mathrm{TSR}=N^{\kappa+\theta+o(1)}\), this adds \(X_\mathrm{TSR}=N^{1+\kappa+\theta+o(1)}\).
\item At the final line, splice the generator circuit of size \(N^{\rho+o(1)}\) with the parsed transcript circuit. Sparse residual mistakes on \(E_{pi}\) source rows are corrected by equality gates costing O(E\_pi N).
\item Therefore a proof transcript with N\textasciicircum{}\{rho+o(1)\}+Gamma\_dag\textasciicircum{}TSR(pi)+O(E\_pi N)<N\textasciicircum{}\{1+epsilon-o(1)\} yields a same-model circuit contradicting LB\_N(1+epsilon). This is a valid falsifier theorem for the restricted proof system.
\item For necessity, suppose a semantic TSR child \(r^*\) influences a conclusion but is absent from K\_dag\textasciicircum{}TSR(pi). Keep all visible rows fixed, flip the private row \(r^*\) or an unread enable bit, and choose the generated parity carrier so the final value changes. The parser cannot distinguish the worlds, so uncharged occurrence completion fails.
\end{enumerate}

\paragraph{Equation 1: carrier.}
\[
H_N^{SAT}(x)=\bigoplus_{j\le M_N} SAT(G_N(x,j)),\quad M_N=N^{q+o(1)},\quad |G_N(x,j)|=(\log N)^d.
\]
\emph{Valid when:} Valid when the source carrier is the beta-zero generated parity carrier with uniformly generated polylogarithmic SAT instances.

\paragraph{Equation 2: fixed\_point.}
\[
K_{t+1}=K_t\cup Parents_{FE}(K_t)\cup Supp_{bEF}(K_t)\cup Leaves_{SR_h}(K_t)\cup Children_{TSR}^{decl}(K_t),\quad K_dag^{TSR}=\bigcup_t K_t.
\]
\emph{Valid when:} Valid when all proof lines expose parent pointers and declared TSR child row indices before semantic use.

\paragraph{Equation 3: tsr\_charge.}
\[
X_{TSR}(\pi)\le N\,R_{TSR}(\pi)L_N^{O(1)}=N^{1+\kappa+\theta+o(1)}\quad\text{if }R_{TSR}(\pi)=N^{\kappa+\theta+o(1)}.
\]
\emph{Valid when:} Valid when each declared child row expands to N source-row occurrences and polylogarithmic formula-local overhead is absorbed.

\paragraph{Equation 4: gamma\_packet.}
\[
\Gamma_{dag}^{TSR}(\pi)=\Gamma_{bEF}(\pi)+U(K_dag^{TSR})+X(K_dag^{TSR})+X_{TSR}(\pi).
\]
\emph{Valid when:} Valid when FormulaEval, bounded self-reduction, bit-explicit EF, and declared TSR supports are all charged in the same ledger.

\paragraph{Equation 5: upper\_bound.}
\[
C_{SIZE}(H_N^{SAT})\le N^{\rho+o(1)}+\Gamma_{dag}^{TSR}(\pi)+O(E_\pi N).
\]
\emph{Valid when:} Applies after the parser line induction converts the charged transcript into a same-model circuit for the carrier.

\paragraph{Equation 6: falsifier\_window.}
\[
\max{\rho,\gamma_{dag}^{TSR},1+e}< 1+\epsilon\Rightarrow LB_N(1+\epsilon)\text{ is false},\quad \Gamma_{dag}^{TSR}=N^{\gamma_{dag}^{TSR}+o(1)},\ E_\pi=N^{e+o(1)}.
\]
\emph{Valid when:} Valid when the asserted source lower bound is in nonuniform SIZE and all transcript charges are subthreshold.

\paragraph{Equation 7: declared\_tsr\_threshold.}
\[
\kappa+\theta< \epsilon\Rightarrow X_{TSR}(\pi)=N^{1+\kappa+\theta+o(1)}< N^{1+\epsilon-o(1)}.
\]
\emph{Valid when:} Valid when \(R_\mathrm{TSR}\) counts only declared reachable child rows and no hidden semantic branch pruning is used.

\paragraph{Next analysis.} Formalize the \(\mathrm{Psys}_{sym}^{dag+declTSR(h)}\) grammar in one page: exact row-index fields, TSR parent/child encodings, local verifier predicate, \(K_{dag}^\mathrm{TSR}\) fixed-point pseudocode, and the declared TSR line-induction case; then instantiate with \(\kappa=\epsilon/3\), \(\theta=\epsilon/3\), q in ((1+epsilon)/2,1+epsilon), \(\rho<1+\epsilon\) to verify all charges remain subthreshold except the intended q-density pressure.

\paragraph{Falsifiers.}
\begin{itemize}
\item Subthreshold declared TSR charge: if \(\kappa+\theta<\epsilon\) and all other \(\Gamma\) terms are subthreshold, the proof transcript gives a same-model SIZE upper-bound falsifier for LB\_N(1+epsilon).
\item First hidden reachable row: an undeclared TSR child or unread enable bit can change a rule conclusion while all charged rows remain fixed, refuting free occurrence completion.
\item Nonuniform selector leakage: if the row list is chosen by advice or SAT-truth labels rather than proof-visible code, \(K_{dag}^\mathrm{TSR}\) is not a valid controlled interface.
\item \(q=N^{\epsilon-q-o(1)}\) exceeds unit density, so a transparent certificate in that regime must fail or pay another superthreshold charge.
\item Semantic pruning escape: arbitrary real-SAT simplification may be nonrelativizing and outside the private-row spoof, but then it must be isolated as \(RealSATInfluenceRealizer_{N}^{arb}\) rather than folded into the parser for free.
\end{itemize}

\subsection*{cycle-0127: Hardness-Magnification Bridge Realization}

\paragraph{Execution.} research; azure\_responses\_api; attempt primary; 120.457 seconds; 3 lemmas, 6 equations, 8 derivation steps, 6 falsifiers.

\paragraph{Learning.} Cycle 127 adds a concrete rule-level proof packet: declared bounded TSR is absorbed by the least closed support cone with exact charge \(N^{1+\kappa+\theta+o(1)}\), while hidden branch pruning has a first-unlisted-child falsifier. This improves the paper theorem boundary without claiming \(\mathrm{P} \ne \mathrm{NP}\) or a general \(H_{N}\) lower bound.

\paragraph{Pressure.} The positive derivation is complete only for declared TSR syntax. Full real-SAT semantic branch pruning remains blocked at \(RealSATInfluenceRealizer_{N}^{arb}\): a uniform, advice-free, nonrelativizing procedure must list all live child rows and charge them before the semantic rule is used.

\paragraph{Derivation steps.}
\begin{enumerate}
\item Initialize \(K_{0}\) with the final \(\mathrm{LB}_{N}\) line, the final carrier row symbols, and any explicitly named residual-error rows.
\item Iterate the monotone closure operator: add bit-EF parents, SATeq links, transparent FormulaEval RowSet entries, and for every declared TSR node add both restriction children, the restriction path bits, and the Tseitin extension clauses.
\item Charge each insertion by its visible certificate: U(K) for enumerators, X(K) for occurrence or row expansion, \(\Gamma_{bEF}\) for bit-explicit proof dependencies, and \(X_\mathrm{TSR}=N\) \(R_\mathrm{TSR}\) for declared TSR child rows.
\item Prove minimality by induction on closure time: any safe extensional parser that validates the final line must include every predecessor inserted by the operator, otherwise the first omitted predecessor admits the private-row disagreement test.
\item Extend the existing \(\Gamma_{bEF}\) line induction. For ordinary bEF lines use the old invariant; for FormulaEval lines call \(TransparentRowSupport_{N}\); for declared TSR lines compute the parent SAT value from declared child values and charged Tseitin consistency circuits.
\item At the final line, splice the generator circuit of size \(N^{\rho+o(1)}\) with the transcript circuits of total size Gamma\_\{dag+TSR\}+O(E\_pi N) to obtain a same-model circuit for \(H_{N}^{\mathrm{SAT}}\).
\item Contrapose: if the final size is below \(N^{1+\epsilon-o(1)}\), the transcript is not a valid certificate of LB\_N(1+epsilon).
\item For hidden TSR pruning, choose the earliest proof use whose semantically live child \(r^*\) is absent from \(K_{dag}\); construct two label worlds agreeing on \(K_{dag}\) and flipping only \(r^*\), then place \(r^*\) at one parity slot \(x^*\) to flip \(H_{N}\) while preserving the visible transcript.
\end{enumerate}

\paragraph{Equation 1: carrier.}
\[
H_N^{\mathrm{SAT}}(x)=\bigoplus_{j< M_N} \mathrm{SAT}(G_N(x,j)),\quad M_N=N^{q+o(1)},\quad |G_N(x,j)|=(\log N)^d
\]
\emph{Valid when:} Valid when the source carrier is the beta-zero parity generator with uniformly indexed polylogarithmic SAT rows and no hidden row advice.

\paragraph{Equation 2: closure\_operator.}
\[
K_{t+1}=K_t\cup \mathrm{Par}_\pi(K_t)\cup \mathrm{RowSet}_{FE}(K_t)\cup \mathrm{ChildRows}_{TSR}(K_t),\quad K_{dag}=\mathrm{lfp}(K_t)
\]
\emph{Valid when:} Valid when every proof line exposes machine-checkable parent pointers and every TSR or FormulaEval influence has a syntactic row certificate.

\paragraph{Equation 3: \(declared_{\mathrm{TSR},charge}\).}
\[
R_{TSR}(\pi)=N^{\kappa+\theta+o(1)},\qquad X_{TSR}(\pi)=N\cdot R_{TSR}(\pi)=N^{1+\kappa+\theta+o(1)}
\]
\emph{Valid when:} Applies after declared TSR roots and children are uniformly enumerable and each declared row expansion is charged by an N-size equality patch.

\paragraph{Equation 4: gamma\_total.}
\[
\Gamma_{dag+TSR}(\pi)=\Gamma_{bEF}(\pi)+U(K_{dag})+X(K_{dag})+X_{TSR}(\pi)
\]
\emph{Valid when:} Valid when bit-explicit EF costs, transparent FormulaEval costs, and declared TSR expansion costs are all in the same circuit-size ledger.

\paragraph{Equation 5: parser\_upper\_bound.}
\[
C_{SIZE}(H_N^{\mathrm{SAT}})\le N^{\rho+o(1)}+\Gamma_{dag+TSR}(\pi)+O(E_\pi N)
\]
\emph{Valid when:} Requires the target-preserving parser induction to build the final carrier circuit in the same nonuniform SIZE model from the charged transcript.

\paragraph{Equation 6: lb\_falsifier\_window.}
\[
N^{\rho+o(1)}+\Gamma_{dag+TSR}(\pi)+O(E_\pi N)< N^{1+\epsilon-o(1)}\Rightarrow \neg LB_N(1+\epsilon)
\]
\emph{Valid when:} Valid when \(\mathrm{LB}_{N}\) is the claimed same-model lower bound for \(H_{N}\) and the constructed upper bound has all advice and row costs charged.

\paragraph{Next analysis.} Write the verifier bit fields for one declared TSR line: root row id, restriction-node path, pivot variable, child row ids, Tseitin map hash, occurrence-neighborhood list, parent pointers, and expansion certificate; then prove the single-line induction inequality adding exactly N times the declared child count to \(\Gamma_{dag+\mathrm{TSR}}\).

\paragraph{Falsifiers.}
\begin{itemize}
\item First-unlisted-child private-row falsifier: hidden TSR pruning lets two transcripts agree on all charged rows while the omitted restriction child flips the carrier parity.
\item Unread-enable-bit falsifier: a symbolic FormulaEval line whose enabling bit is never surfaced can select an uncharged row after the parser commits to \(K_{dag}\).
\item Threshold falsifier: if \(\kappa+\theta<\epsilon\) and all other \(\Gamma\) terms are subthreshold, declared TSR support produces an explicit SIZE upper bound below \(N^{1+\epsilon}\).
\item \(q<\epsilon\) enumeration falsifier: a support cone smaller than \(N^{\epsilon-o(1)}\) cannot certify LB\_N(1+epsilon) because N|K| is patchable.
\item Model-mismatch falsifier: row labels selected by SAT-truth advice are outside the same nonuniform SIZE ledger and cannot be counted as transparent support.
\item Positive-beta carryover falsifier: any return to black-box target-circuit substitution with \(\beta>0\) is capped by q+beta c+o(1) and cannot exclude all SAT circuit exponents.
\end{itemize}

\subsection*{cycle-0128: Hardness-Magnification Bridge Realization}

\paragraph{Execution.} research; azure\_responses\_api; attempt primary; 155.835 seconds; 3 lemmas, 4 equations, 9 derivation steps, 5 falsifiers.

\paragraph{Learning.} Cycle 128 closes a sharper declared-TSR support sublemma: complete child declarations compose with the least support fixed point and cost exactly \(N^{1+\kappa+\theta+o(1)}\), while hidden branch pruning has a first-unlisted-child private-row falsifier. The contribution remains a scoped parser/falsifier theorem, not evidence for \(\mathrm{P} \ne \mathrm{NP}\).

\paragraph{Pressure.} The declared complete-child schema has no current derivation gap beyond writing the final syntax checker. The full semantic rule remains blocked exactly at \(RealSATInfluenceRealizer_{N}^{arb}\), which must list all influencing rows without SAT-truth advice.

\paragraph{Derivation steps.}
\begin{enumerate}
\item Initialize \(K_{0}\) with row ids and proof lines reachable from the final LB\_N(1+epsilon) line.
\item Define closure operator Cl(S): add bit-explicit EF parent dependencies, SATeq quotient links, transparent FormulaEval support rows, and for every TSR parent row in S add every declared child row from its syntactic child list.
\item Iterate K\_\{t+1\}=Cl(K\_t) until a fixed point. Termination is bounded by the finite proof-size row universe; canonical ids ensure no duplicate charge.
\item For non-TSR lines reuse the established \(\Gamma_{bEF}\) and \(\Gamma_{dag}\) line induction: every reachable line has a local circuit using only previously charged reachable rows.
\item For a declared TSR line, verify each child descriptor locally: parent id, path bits, restriction mask, Tseitin map id, and child slot id. Build the parent semantic value by a bounded OR/AND circuit over the declared child row values plus poly(L\_N) consistency checks.
\item Charge N*poly(L\_N) per declared child row expansion. With \(A_\mathrm{TSR}=N^{\kappa+o(1)}\) active parents and \(B_\mathrm{TSR}=N^{\theta+o(1)}\) children per parent, obtain \(R_\mathrm{TSR}\le N^{1+\kappa+\theta+o(1)}\).
\item Patch \(E_{pi}\) residual rows by equality tests on the N-bit source input, costing O(E\_pi N). Add the generator cost \(N^{\rho+o(1)}\) to obtain a same-model circuit for \(H_{N}^\mathrm{SAT}\).
\item If the resulting size is below \(N^{1+\epsilon-o(1)}\), the alleged source lower-bound certificate has a valid upper-bound falsifier.
\item For hidden pruning, choose the first child r* used semantically but absent from \(K_\mathrm{TSR}\). Flip its SAT label in a two-world oracle while preserving all visible rows; the proof transcript is unchanged but the induced reachable parity can flip.
\end{enumerate}

\paragraph{Equation 1: \(\Gamma_{\mathrm{TSR},ledger}\).}
\[
Gamma_TSR(pi)=Gamma_bEF(pi)+U(K_dag)+X(K_dag)+R_TSR(pi)+O(E_pi N)
\]
\emph{Valid when:} Valid when pi belongs to the declared support schema and each dependency edge is surfaced before the parser line induction uses it.

\paragraph{Equation 2: \(Declared_{\mathrm{TSR},charge}\).}
\[
R_TSR(pi) \le N*A_TSR*B_TSR*L_N^{O(1)} = N^{1+kappa+theta+o(1)}
\]
\emph{Valid when:} Applies when active TSR parent rows are \(N^{\kappa+o(1)}\) and every parent declares \(N^{\theta+o(1)}\) bounded children of polylog description size.

\paragraph{Equation 3: \(Parser_{upper,bound}\).}
\[
C(H_N^SAT) \le N^{rho+o(1)}+Gamma_TSR(pi)+O(E_pi N)
\]
\emph{Valid when:} Requires the target-preserving fixed-point parser, same-model row expansion, and equality patching for only \(E_{pi}\) residual rows.

\paragraph{Equation 4: \(\mathrm{TSR}_{subthreshold,window}\).}
\[
rho< 1+epsilon, kappa+theta< epsilon, and E_pi< N^{epsilon-o(1)} imply a subthreshold LB_N falsifier if other Gamma_dag terms are subthreshold
\]
\emph{Valid when:} Valid when the preexisting transparent-support charges are already below \(N^{1+\epsilon-o(1)}\) and TSR is the only new rule charge.

\paragraph{Next analysis.} Write the paper-ready TSR rule schema table and verify on one toy projected-dDNNF plus TSR line: compute \(K_\mathrm{TSR}\), \(A_\mathrm{TSR}\), \(B_\mathrm{TSR}\), \(R_\mathrm{TSR}\), and then flip the first unlisted child in the hidden-pruning variant.

\paragraph{Falsifiers.}
\begin{itemize}
\item falsifier: first-unlisted TSR child whose SAT label affects a reachable line but is absent from the charged support cone.
\item falsifier: \(\kappa+\theta\ge \epsilon\), causing declared TSR expansion to lose the subthreshold contradiction window.
\item falsifier: semantic branch pruning that depends on real satisfiability before any advice-free row realizer is emitted.
\item falsifier: positive \(\beta\) fallback, where the transfer upper bound q+beta*c+o(1) caps all-exponent P/poly exclusion.
\item falsifier: nonuniform child-list advice hidden in extension variables rather than verifier-visible row encodings.
\end{itemize}

\subsection*{cycle-0129: Hardness-Magnification Bridge Realization}

\paragraph{Execution.} research; azure\_responses\_api; attempt primary; 118.925 seconds; 3 lemmas, 5 equations, 7 derivation steps, 5 falsifiers.

\paragraph{Learning.} Cycle 129 turns declared bounded TSR from a semantic-warning slogan into a line-induction rule case for the least support theorem. Complete child declaration adds exactly \(N^{1+\kappa+\theta+o(1)}\) to \(\Gamma_\mathrm{DTSR}\) and is valid below \(\kappa+\theta<\epsilon\); hidden branch pruning is refuted by a first-unlisted-child falsifier unless a nonrelativizing real-SAT influence realizer is supplied.

\paragraph{Pressure.} No blocking gap for the restricted complete-child declared TSR rule case. The full semantic FormulaEval/TSR system remains blocked exactly at \(RealSATInfluenceRealizer_{N}^{arb}\), which must uniformly list and charge semantically influencing rows before use.

\paragraph{Derivation steps.}
\begin{enumerate}
\item Fix the source carrier \(H_{N}^\mathrm{SAT}\) with \(M_{N}=N^{q+o(1)}\), L\_N=(log N)\textasciicircum{}d, \(d>1\), and \(\beta=0\). Fix \(\mathrm{Psys}_{sym}^{dag+declTSR,h}\) over bit-explicit EF lines, transparent FormulaEval lines, and declared TSR lines.
\item For each proof line record verifier bit fields: kind, line index, parent pointers, source row a, formula pointer p, restriction code z, child row code r=(a,p,z), enable bit e\_r, and an explicit occurrence-neighborhood list Occ(r).
\item Initialize \(K_{0}\) from the final \(\mathrm{LB}_{N}\) assertion and compute \(K_\mathrm{DTSR}\) by the finite fixed point \(K<-K\) union parents union FormulaEval supports union declared TSR children. This is implementable because every addition is syntactically listed or full-scan charged.
\item Extend the prior \(\Gamma_{bEF}\) line induction. EF and SATeq lines use the existing dependency and quotient-link charges; bounded \(SR_{h}\) lines use B\_SR(pi); transparent FormulaEval lines use U(K)+X\_FE(K); declared TSR lines add exactly \(N^{1+o(1)}\) per declared child occurrence expansion.
\item For a declared TSR inference, the local verifier recomputes the Tseitin/restriction map on Occ(r), checks child-parent consistency, and charges the child row before the rule can feed any later line. Thus the induction invariant is prefix-closed under later symbolic use.
\item At the final \(\mathrm{LB}_{N}\) line the parser yields a nonuniform circuit for \(H_{N}\) of size at most N\textasciicircum{}\{rho+o(1)\}+Gamma\_DTSR(pi)+O(E\_pi N). If this is below \(N^{1+\epsilon-o(1)}\), the claimed source lower bound is falsified in the same SIZE model.
\item For hidden branch pruning, choose the first reachable child row r* not in \(K_{vis}\). Two row-label worlds agreeing on \(K_{vis}\) and differing at r* pass the same visible transcript while changing a reachable parity bit, so surface occurrence parsing is invalid without \(RealSATInfluenceRealizer_{N}^{arb}\).
\end{enumerate}

\paragraph{Equation 1: carrier.}
\[
H_N^{\mathrm{SAT}}(x)=\bigoplus_{j\le M_N}\mathrm{SAT}(G_N(x,j)),\quad M_N=N^{q+o(1)}
\]
\emph{Valid when:} Valid when the branch is restricted to beta-zero parity carriers with generated polylogarithmic SAT queries and fixed source length N.

\paragraph{Equation 2: least\_support\_fixed\_point.}
\[
K_{t+1}=K_t\cup \mathrm{Par}(K_t)\cup \mathrm{Supp}_{FE}(K_t)\cup \mathrm{Child}_{TSR}(K_t),\quad K_{DTSR}=\mathrm{lfp}(K_0)
\]
\emph{Valid when:} Applies after every proof line exposes parent pointers and declared child-row encodings before the line can influence the final \(\mathrm{LB}_{N}\) assertion.

\paragraph{Equation 3: declared\_tsr\_charge.}
\[
X_{TSR}(\pi)\le N^{1+o(1)}D_{TSR}(\pi)C_{TSR}(N)=N^{1+\kappa+\theta+o(1)}
\]
\emph{Valid when:} Requires declared TSR lines to list all restriction children and their occurrence-neighborhood expansions with no semantic branch pruning.

\paragraph{Equation 4: charged\_transcript.}
\[
\Gamma_{DTSR}(\pi)=\Gamma_{bEF}(\pi)+U(K_{DTSR})+X_{FE}(K_{DTSR})+X_{TSR}(\pi)
\]
\emph{Valid when:} Valid when EF dependencies, FormulaEval support, and declared TSR child rows are closed under the least reachable support fixed point.

\paragraph{Equation 5: source\_falsifier\_window.}
\[
N^{\rho+o(1)}+\Gamma_{DTSR}(\pi)+O(E_\pi N)< N^{1+\epsilon-o(1)}\Rightarrow \neg LB_N(1+\epsilon)
\]
\emph{Valid when:} Applies after the line induction constructs an explicit same-model circuit for \(H_{N}\) and sparse row errors are equality patched.

\paragraph{Next analysis.} Write the paper-ready L129 sublemma proof: give pseudocode for \(K_\mathrm{DTSR}\) fixed-point enumeration, a table of declared TSR bit fields, the local verifier for child occurrence neighborhoods, and the two-world first-unlisted-child countermodel; then run the live inference-quality packet validation.

\paragraph{Falsifiers.}
\begin{itemize}
\item First-unlisted-child private-row or unread-enable-bit falsifier for hidden TSR branch pruning.
\item Threshold-collapse falsifier \(\kappa+\theta\ge \epsilon\): declared TSR expansion is too large to refute \(\mathrm{LB}_{N}\) by a subthreshold parser circuit.
\item \(q<\epsilon\) sparse-support falsifier: too few reachable rows allow explicit enumeration and equality patching.
\item Nonuniform leakage falsifier: a row selector or semantic pruning decision supplied as advice instead of verifier-visible syntax invalidates \(\Gamma_\mathrm{DTSR}\).
\item Positive-beta substitution falsifier remains quarantined: any black-box target-query exponent \(\beta>0\) gives the transfer cap q+beta c+o(1).
\end{itemize}

\subsection*{cycle-0130: Hardness-Magnification Bridge Realization}

\paragraph{Execution.} research; azure\_responses\_api; attempt primary; 140.086 seconds; 3 lemmas, 7 equations, 9 derivation steps, 5 falsifiers.

\paragraph{Learning.} Cycle 130 closes a real proof-delivery subcase: declared bounded TSR is not merely compatible with transparent support but is an explicit fixed-point rule with charge \(N^{1+\kappa+\theta+o(1)}\). The exact boundary is now complete declaration versus first-unlisted-child falsifier; full semantic pruning remains open only through the nonrelativizing \(RealSATInfluenceRealizer_{N}^{arb}\).

\paragraph{Pressure.} Full semantic branch pruning remains blocked at \(RealSATInfluenceRealizer_{N}^{arb}\): the derivation does not show that real SAT semantics can uniformly list all influential pruned children without advice and before use.

\paragraph{Derivation steps.}
\begin{enumerate}
\item Define \(\mathrm{DTSR}-TransparentSupport_{N}\): every rule line has a bit-level tag, parent line numbers, row-index fields, local formula identifiers, enumerator code, and for each active TSR row a complete list of restriction-child descriptors before the line is used.
\item Compute K\_DTSR(pi) by a monotone fixed point from the final \(\mathrm{LB}_{N}\) line: add proof parents, EF extension support, SATeq links, \(bounded-SR_{h}\) leaves, FormulaEval dependencies, and declared TSR children until no new row descriptor appears.
\item Maintain induction invariant I(l): after processing line l, every row label that can affect l is in \(K_\mathrm{DTSR}\) restricted to the prefix, and the compiled same-model circuit cost is at most the corresponding prefix ledger.
\item Base cases are inherited from the bit-explicit \(\mathrm{Psys}_{bEF}^{-}\) parser: literal rows, EF definitions, SATeq quotient links, and \(bounded-SR_{h}\) leaves contribute exactly their \(\Gamma_{bEF}\) charges.
\item FormulaEval lines use either a visible enumerator charge U(K) or a full-scan charge, then pay same-theory row-expansion X\_FE(K); hidden FormulaEval is not admitted.
\item For a declared TSR line, enumerate \(A_{N}\) active descriptors and \(F_{N}\) children per descriptor, write each restricted child as an ordinary SAT row, and charge N \(A_{N}\) \(F_{N}\) L\_N\textasciicircum{}\{O(1)\}=N\textasciicircum{}\{1+theta+kappa+o(1)\}.
\item The induction composes because later symbolic uses can reach only rows already in the prefix-closed fixed point; hence no last-unsafe-line remains under the declared interface.
\item At the final \(\mathrm{LB}_{N}\) line, the compiled circuit plus residual equality patches has size N\textasciicircum{}\{rho+o(1)\}+Gamma\_DTSR(pi)+O(E\_pi N).
\item If \(\rho<1+\epsilon\) and the ledger is below \(N^{1+\epsilon-o(1)}\), the constructed circuit contradicts the claimed source lower bound LB\_N(1+epsilon), giving a route-level falsifier.
\end{enumerate}

\paragraph{Equation 1: carrier\_parameters.}
\[
H_N^{SAT}(x)=\bigoplus_{j\le M_N} SAT(G_N(x,j)),\quad M_N=N^{q+o(1)},\quad |G_N|=N^{\rho+o(1)},\quad L_N=(\log N)^d
\]
\emph{Valid when:} Valid when the source carrier is the beta-zero generated parity carrier with uniformly computable generator and fixed polylogarithmic SAT row length.

\paragraph{Equation 2: least\_fixed\_point.}
\[
K_{t+1}=K_t\cup Par(K_t)\cup FE(K_t)\cup SR_h(K_t)\cup DTSR_{decl}(K_t),\quad K_{DTSR}(\pi)=\bigcup_{t\le |\pi|}K_t
\]
\emph{Valid when:} Applies when every proof line has verifier-visible parent pointers and every semantic dependency class is either empty or listed as a charged row set.

\paragraph{Equation 3: ledger.}
\[
\Gamma_{DTSR}(\pi)=\Gamma_{bEF}(\pi)+U(K_{DTSR})+X_{FE}(K_{DTSR})+X_{DTSR}(K_{DTSR})+O(E_\pi N)
\]
\emph{Valid when:} Valid when the parser induction charges bit-explicit proof dependencies, FormulaEval enumeration, declared TSR expansions, and residual sparse-row patching in the same SIZE model.

\paragraph{Equation 4: declared\_tsr\_charge.}
\[
X_{DTSR}(K_{DTSR})\le N\cdot A_N\cdot F_N\cdot L_N^{O(1)}=N^{1+\theta+\kappa+o(1)}
\]
\emph{Valid when:} Requires the declared TSR schema where \(A_{N}\) active row descriptors and \(F_{N}\) restriction children are uniformly listed before semantic use.

\paragraph{Equation 5: parser\_upper\_bound.}
\[
C_{SIZE}(H_N^{SAT})\le N^{\rho+o(1)}+\Gamma_{DTSR}(\pi)+O(E_\pi N)
\]
\emph{Valid when:} Applies after the line induction compiles every charged support row into a same-theory circuit and no hidden semantic row remains reachable.

\paragraph{Equation 6: subthreshold\_falsifier.}
\[
\rho< 1+\epsilon\ \wedge\ \Gamma_{DTSR}(\pi)+O(E_\pi N)< N^{1+\epsilon-o(1)}\Rightarrow \neg LB_N(1+\epsilon)
\]
\emph{Valid when:} Valid when \(\mathrm{LB}_{N}\) is a same-model \(N^{1+\epsilon-o(1)}\) circuit lower-bound claim for the carrier \(H_{N}^{\mathrm{SAT}}\).

\paragraph{Equation 7: declared\_tsr\_window.}
\[
\theta+\kappa< \epsilon \Rightarrow X_{DTSR}(K_{DTSR})=o(N^{1+\epsilon})
\]
\emph{Valid when:} Applies in the fixed-exponent regime where declared TSR fanout and active-row exponents are independent of N and all polylog factors are absorbed.

\paragraph{Next analysis.} Run the bit-field audit for one concrete DTSR verifier line: list parent row id, active descriptor id, restriction mask, child row id, Tseitin extension map, local identity tag, and downstream use pointer; verify that \(K_\mathrm{DTSR}\) reaches every used child and compute the exact theta+kappa ledger for the example.

\paragraph{Falsifiers.}
\begin{itemize}
\item first-unlisted-child private-row falsifier for hidden TSR pruning
\item unread-enable-bit falsifier for FormulaEval influence not in K\_DTSR(pi)
\item threshold-collapse falsifier when \(\rho\) or \(\Gamma_\mathrm{DTSR}\) reaches \(N^{1+\epsilon-o(1)}\)
\item nonuniform row-selection advice falsifier for claimed transparent enumeration
\item positive-beta black-box substitution falsifier q+beta*c against all-exponent target transfer
\end{itemize}

\subsection*{cycle-0131: Hardness-Magnification Bridge Realization}

\paragraph{Execution.} research; azure\_responses\_api; attempt primary; 128.799 seconds; 3 lemmas, 6 equations, 9 derivation steps, 6 falsifiers.

\paragraph{Learning.} Cycle 131 turns declared bounded TSR into an auditable fixed-point rule case: complete child declarations are sufficient for the \(\Gamma_\mathrm{DTSR}\) parser and necessary up to the first-unlisted-child falsifier. The valid contribution is a route-level theorem/countermodel boundary for beta-zero parity magnification certificates, not progress on the main P vs NP conjecture.

\paragraph{Pressure.} The restricted declared-TSR theorem has no current internal gap. The full semantic pruning case remains blocked exactly at \(RealSATInfluenceRealizer_{N}^{arb}\): an advice-free nonrelativizing procedure must list all real-SAT influencing rows and charge them before use.

\paragraph{Derivation steps.}
\begin{enumerate}
\item Encode every proof-line dependency by visible fields: line pointer, rule tag, source key a, slot j, child path p, parent pointers, polarity, formula hash, and touched occurrence-neighborhood list.
\item Compute K\_DTSR(pi) by a worklist from the final \(\mathrm{LB}_{N}\) line: add EF parents, SATeq quotient links, \(TransparentRowSupport_{N}\) influences, \(bounded-SR_{h}\) leaves, and every declared TSR child; terminate because pi is finite and all selectors are syntactic or full-scan charged.
\item For a bEF or SATeq line, reuse the existing \(\Gamma_{bEF}\) induction: evaluating the parent row values and link certificates suffices to evaluate the child line on all rows in the cone.
\item For a \(bounded-SR_{h}\) line, charge \(B_{SR,h}\le N^{q+\eta+\tau+o(1)}\) because the verifier lists the bounded depth path, row-index encoding, and occurrence-neighborhoods touched by each restriction.
\item For a declared TSR line, the child list \(A_{r}\) is part of the syntax; expanding each child over N source keys costs \(N|A_{r}|N^{o(1)}\), so all TSR instances cost at most \(N^{1+\kappa+\theta+o(1)}\).
\item For a transparent FormulaEval line, use its advice-free enumerator or full scan to add U(K\_DTSR), then charge X(K\_DTSR) for same-theory occurrence expansion of the listed rows.
\item Induct over the proof DAG to build a circuit of size N\textasciicircum{}\{rho+o(1)\}+Gamma\_DTSR(pi) that agrees with \(H_{N}\) except possibly on \(E_{pi}\) residual source rows; add O(E\_pi N) equality patches.
\item Contrapose: if the resulting size is below \(N^{1+\epsilon-o(1)}\), the claimed LB\_N(1+epsilon) is false in the same source model, so any sound certificate must exceed the threshold.
\item Necessity check: if a later line depends on an unlisted TSR child, construct two worlds agreeing on \(K_\mathrm{DTSR}\) but differing on that child; the visible parser outputs the same value while the enabled carrier parity differs.
\end{enumerate}

\paragraph{Equation 1: carrier.}
\[
H_N^{SAT}(x)=\bigoplus_{j\le M_N}\mathrm{SAT}(G_N(x,j)),\quad M_N=N^{q+o(1)},\quad L_N=(\log N)^d.
\]
\emph{Valid when:} Valid when the source carrier is the beta-zero generated parity family with uniformly generated polylogarithmic SAT instances.

\paragraph{Equation 2: \(bounded_{SR,charge}\).}
\[
B_{SR_h}(\pi)\le N^{q+\eta+\tau+o(1)}.
\]
\emph{Valid when:} Valid when the bounded self-reduction rule exposes every touched occurrence neighborhood and the eta,tau overheads are syntactic parser costs.

\paragraph{Equation 3: \(declared_{\mathrm{TSR},charge}\).}
\[
B_{DTSR}(\pi)=\sum_{r\in DTSR(\pi)} N\cdot |A_r|\cdot N^{o(1)}\le N^{1+\kappa+\theta+o(1)}.
\]
\emph{Valid when:} Applies when every TSR child in every active rule family is declared before use and row expansion is charged in the source length N.

\paragraph{Equation 4: \(gamma_\mathrm{DTSR}\).}
\[
\Gamma_{DTSR}(\pi)=\Gamma_{bEF}(\pi)+U(K_{DTSR})+X(K_{DTSR})+B_{DTSR}(\pi).
\]
\emph{Valid when:} Valid when \(K_\mathrm{DTSR}\) is the least prefix-closed support fixed point and all FormulaEval, \(SR_{h}\), and TSR influences are surfaced.

\paragraph{Equation 5: upper\_bound\_falsifier.}
\[
C_{SIZE}(H_N^{SAT})\le N^{\rho+o(1)}+\Gamma_{DTSR}(\pi)+O(E_\pi N).
\]
\emph{Valid when:} Requires the line-induction parser to evaluate every proof line on surfaced rows and to patch only the \(E_{pi}\) residual source rows.

\paragraph{Equation 6: subthreshold\_contradiction.}
\[
N^{\rho+o(1)}+\Gamma_{DTSR}(\pi)+O(E_\pi N)< N^{1+\epsilon-o(1)}\Rightarrow \neg LB_N(1+\epsilon).
\]
\emph{Valid when:} Valid when \(\mathrm{LB}_{N}\) is interpreted in the same nonuniform SIZE model and no semantic rule uses hidden advice or uncharged SAT influence.

\paragraph{Next analysis.} Run the reviewer audit on L131.1: instantiate one toy projected-dDNNF FormulaEval line plus one declared TSR parent with two restriction children, compute \(K_\mathrm{DTSR}\) by the worklist, verify \(\Gamma_\mathrm{DTSR}\) numerically as \(N^{1+\kappa+\theta+o(1)}\), then flip the first unlisted child in the hidden-pruning variant to validate the private-row falsifier.

\paragraph{Falsifiers.}
\begin{itemize}
\item first-unlisted-child private-row falsifier for hidden TSR branch pruning
\item unread-enable-bit falsifier for non-closed FormulaEval influence
\item \(q<\epsilon\) support-density falsifier when reachable rows are too sparse
\item positive-beta substitution cap N\textasciicircum{}\{q+beta c+o(1)\} for black-box target-circuit interfaces
\item nonuniform leakage through SAT-truth advice in row selection
\item threshold collapse when \(\rho\), \(\Gamma_\mathrm{DTSR}\), or \(E_{pi}\) N reaches \(N^{1+\epsilon-o(1)}\)
\end{itemize}

\subsection*{cycle-0132: Hardness-Magnification Bridge Realization}

\paragraph{Execution.} research; azure\_responses\_api; attempt primary; 109.826 seconds; 3 lemmas, 5 equations, 8 derivation steps, 5 falsifiers.

\paragraph{Learning.} Cycle 132 upgrades declared bounded TSR from a rule-case slogan to a fixed-point parser theorem with explicit bit fields, closure algorithm, ledger, and necessity falsifier. The valid contribution is the restricted \(TransparentRowSupport_{N}^{\mathrm{DTSR}}\) theorem/countermodel pair; full semantic pruning remains quarantined exactly at \(RealSATInfluenceRealizer_{N}^{arb}\).

\paragraph{Pressure.} \(RealSATInfluenceRealizer_{N}^{arb}\): arbitrary semantic branch-pruning may be valid over real SAT yet not expose an advice-free row list in the proof syntax. The present theorem covers declared DTSR and transparent support only.

\paragraph{Derivation steps.}
\begin{enumerate}
\item Define \(K_\mathrm{DTSR}^{0}\) as the final \(\mathrm{LB}_{N}\) line and its explicit \(H_{N}\) row references.
\item Iterate a deterministic closure operator: for each bEF line add parent dependencies and extension supports; for each transparent FormulaEval line add enumerated or full-scan RowSet entries; for SATeq add quotient-linked rows; for \(SR_{h}\) add bounded leaf rows; for DTSR add every declared child tuple and its occurrence neighborhood.
\item Termination is immediate because the closure ranges over the finite proof DAG and finite declared child lists; the least fixed point is independent of traversal order by monotonicity.
\item Verifier bit fields for a DTSR child tuple are line\_id, parent\_row, restriction\_id, child\_row, child\_formula\_id, local\_clause\_map, polarity, enable\_bit, and next\_pointer; each field is read by the local checker before charging \(B_\mathrm{DTSR}\).
\item Line induction invariant: after processing line t, the partial transcript over \(K_\mathrm{DTSR}\) restricted to lines \(\le t\) suffices to build the corresponding circuit fragment with size bounded by the accumulated \(\Gamma\) ledger plus generator and residual costs.
\item The DTSR line case composes the already-built parent-row circuit with the declared child circuits and local Tseitin maps; because every enabled child is in \(K_\mathrm{DTSR}\), no semantic branch is consulted for free.
\item At the final line, the transcript yields C(H\_N)<=N\textasciicircum{}\{rho+o(1)\}+Gamma\_DTSR(pi)+O(E\_pi N); if this is below \(N^{1+\epsilon-o(1)}\), the alleged \(\mathrm{LB}_{N}\) certificate is refuted in the same model.
\item For hidden pruning, choose the first reachable unlisted child r*, hold all charged rows fixed, flip the row label or unread enable bit at r*, and propagate the parity difference; this proves the countermodel unless a real-SAT influence realizer charges r*.
\end{enumerate}

\paragraph{Equation 1: carrier.}
\[
H_N^{SAT}(x)=\bigoplus_{j< M_N} SAT(G_N(x,j)),\qquad M_N=N^{q+o(1)}
\]
\emph{Valid when:} Valid when the source carrier is the beta-zero generated parity family and \(G_{N}\) uniformly emits polylog-length SAT instances indexed by j.

\paragraph{Equation 2: dtsr\_charge.}
\[
B_{DTSR}(\pi)=N\cdot P_N\cdot C_N=N^{1+\kappa+\theta+o(1)}
\]
\emph{Valid when:} Valid when \(\mathrm{P}_{N}=N^{\kappa+o(1)}\) reachable DTSR parent rows each declare \(C_{N}=N^{\theta+o(1)}\) child rows with bit-visible indices.

\paragraph{Equation 3: ledger.}
\[
\Gamma_{DTSR}(\pi)=\Gamma_{bEF}(\pi)+U(K_{DTSR})+X(K_{DTSR})+B_{SR_h}(\pi)+B_{DTSR}(\pi)
\]
\emph{Valid when:} Applies after the fixed-point parser has closed every reachable transparent influence and every declared self-reduction child before semantic use.

\paragraph{Equation 4: parser\_upper\_bound.}
\[
C(H_N^{SAT})\le N^{\rho+o(1)}+\Gamma_{DTSR}(\pi)+O(E_\pi N)
\]
\emph{Valid when:} Requires the target-preserving same-model simulation, advice-free support lists, and equality patching only on \(E_{pi}\) explicitly named residual rows.

\paragraph{Equation 5: subthreshold\_falsifier.}
\[
\max{\rho,\gamma,1+e}< 1+\epsilon,\ \Gamma_{DTSR}(\pi)=N^{\gamma+o(1)},\ E_\pi=N^{e+o(1)}\Rightarrow \neg LB_N(1+\epsilon)
\]
\emph{Valid when:} Valid when the alleged proof certificate asserts a SIZE lower bound for exactly the same \(H_{N}\) and all parser charges are included.

\paragraph{Next analysis.} Write the paper-ready L132 proof block: pseudocode for K\_DTSR(pi), a table of DTSR bit fields and verifier checks, the five line-induction cases, and the first-unlisted-child countermodel; then test one projected-dDNNF+DTSR positive example and one unread-enable-bit hidden-pruning counterexample.

\paragraph{Falsifiers.}
\begin{itemize}
\item positive \(\beta\) falsifier: if SAT query length is \(N^\beta\) with \(\beta>0\), \(\mathrm{SAT}\in \mathrm{SIZE}(n^{c})\) gives source size N\textasciicircum{}\{q+beta c+o(1)\} and caps all-exponent transfer
\item first-unlisted-child falsifier: a reachable DTSR child not listed in \(K_\mathrm{DTSR}\) permits two charged-identical transcripts with different final parity
\item unread-enable-bit falsifier: a hidden branch-selection bit controls influence while remaining absent from the row ledger
\item subthreshold ledger falsifier: Gamma\_DTSR(pi)+N\textasciicircum{}\{rho+o(1)\}+O(E\_pi N)<N\textasciicircum{}\{1+epsilon-o(1)\} constructs a same-model circuit for \(H_{N}\)
\item nonuniform leakage falsifier: row sets chosen by semantic advice rather than verifier-visible syntax invalidate \(TransparentRowSupport_{N}\)
\end{itemize}

\subsection*{cycle-0133: Hardness-Magnification Bridge Realization}

\paragraph{Execution.} research; azure\_responses\_api; attempt primary; 114.836 seconds; 3 lemmas, 6 equations, 10 derivation steps, 5 falsifiers.

\paragraph{Learning.} Cycle 133 upgrades declared bounded TSR from a compatible rule case to a live inference-quality theorem packet: complete child declarations are sufficient for the \(\Gamma_\mathrm{DTSR}\) parser and necessary up to a first-unlisted-child falsifier. The valid contribution is a scoped \(TransparentRowSupport_{N}^\mathrm{DTSR}\) theorem/countermodel boundary; full semantic pruning remains exactly blocked at \(RealSATInfluenceRealizer_{N}^{arb}\).

\paragraph{Pressure.} Full semantic branch pruning remains blocked at \(RealSATInfluenceRealizer_{N}^{arb}\): no proof here uniformly lists all real-SAT-influencing rows when the branch choice is determined by semantic satisfiability rather than declared syntax.

\paragraph{Derivation steps.}
\begin{enumerate}
\item Source statement: LB\_N(1+epsilon) claims C\_SIZE(H\_N\textasciicircum{}SAT)>N\textasciicircum{}\{1+epsilon-o(1)\} for the named beta-zero parity carrier, in the same nonuniform circuit model used by the parser simulation.
\item Target statement: under \(TransparentRowSupport_{N}^\mathrm{DTSR}\), a proof pi with subthreshold \(\Gamma_\mathrm{DTSR}\) gives C\_SIZE(H\_N\textasciicircum{}SAT)<N\textasciicircum{}\{1+epsilon-o(1)\}, hence pi is a route-level falsifier of that source certificate.
\item Model context: source carrier is \(H_{N}^\mathrm{SAT}\), target class for the bridge audit is nonuniform SIZE, proof model is \(\mathrm{Psys}_{bEF}^{-}\) plus transparent FormulaEval and declared bounded TSR, and reduction type is parser-to-circuit simulation plus sparse equality patching.
\item Quantifier context: constants q,rho,kappa,theta,epsilon are fixed; N tends through valid lengths; \(\beta=0\) via polylogarithmic SAT queries; enumeration and row expansion are advice-free except for explicit proof bits.
\item Bit fields fixed for DTSR lines: line kind, parent ids, source row x, slot j, formula-node address a, restriction mask r, child id, occurrence-neighborhood interval list, selector program id, and closure pointers into \(K_{dag}\).
\item Closure algorithm: start from final \(\mathrm{LB}_{N}\) line; repeatedly add parent dependencies, visible FormulaEval supports, same-theory row expansions, bounded \(SR_{h}\) leaves, and every declared TSR child of every active parent row until no new row appears.
\item Line induction: bEF and extension lines charge \(\Gamma_{bEF}\); transparent FormulaEval charges U\_N(K\_dag)+X\_N(K\_dag); \(SR_{h}\) charges \(B_{SR,h}\); declared TSR charges N times the number of active parent-child certificates; sparse failures charge O(E\_pi N).
\item DTSR rule case: for each active parent row, the verifier recomputes the syntactic Tseitin restriction child and its touched occurrence neighborhood; because all children are declared, later semantic use cannot introduce an uncharged row.
\item Countermodel step: if the first later-used child is absent from \(K_{dag}\), construct two SAT-label environments agreeing on charged rows and differing on that child; the proof transcript is unchanged but the dependent final predicate changes, so the parser theorem is invalid for hidden pruning.
\item Barrier status: the positive restricted theorem is relativizing over visible labels and therefore cannot prove a real SAT lower bound; the escape for full semantic pruning must be nonrelativizing through \(RealSATInfluenceRealizer_{N}^{arb}\), not a generic natural-proofs or algebrization slogan.
\end{enumerate}

\paragraph{Equation 1: carrier.}
\[
H_N^{SAT}(x)=\bigoplus_{j< M_N} SAT(G_N(x,j)),\quad M_N=N^{q+o(1)}
\]
\emph{Valid when:} Valid when N is in the carrier length set and \(G_{N}\) uniformly outputs polylogarithmic SAT instances with \(\beta\) equal to zero.

\paragraph{Equation 2: ledger.}
\[
\Gamma_{DTSR}(\pi)=\Gamma_{bEF}(\pi)+U_N(K_{dag})+X_N(K_{dag})+B_{SR_h}(\pi)+B_{DTSR}(\pi)
\]
\emph{Valid when:} Applies when every proof line has visible dependencies and \(K_{dag}\) is the least prefix-closed reachable support fixed point.

\paragraph{Equation 3: \(\mathrm{DTSR}_{charge}\).}
\[
B_{DTSR}(\pi)=N\cdot |A_N|\cdot N^{\theta+o(1)}=N^{1+\kappa+\theta+o(1)}
\]
\emph{Valid when:} Requires complete child declaration for \(|A_{N}|=N^{\kappa+o(1)}\) active parent rows and syntactic same-theory expansion for every child.

\paragraph{Equation 4: parser\_upper\_bound.}
\[
C_{SIZE}(H_N^{SAT})\le N^{\rho+o(1)}+\Gamma_{DTSR}(\pi)+O(E_\pi N)
\]
\emph{Valid when:} Valid after line induction converts the charged transcript into a same-model circuit and only sparse residual rows remain.

\paragraph{Equation 5: source\_falsifier\_window.}
\[
N^{\rho+o(1)}+\Gamma_{DTSR}(\pi)+O(E_\pi N)< N^{1+\epsilon-o(1)}\Rightarrow \neg LB_N(1+\epsilon)
\]
\emph{Valid when:} Applies when \(\mathrm{LB}_{N}\) asserts a same-model \(N^{1+\epsilon-o(1)}\) lower bound for the identical carrier \(H_{N}^{\mathrm{SAT}}\).

\paragraph{Equation 6: beta\_zero\_target\_substitution.}
\[
SAT\in SIZE(n^c)\Rightarrow C(H_N^{SAT})\le N^{\max{\rho,q}+o(1)}
\]
\emph{Valid when:} Valid when query length is polylogarithmic, \(\beta\) is zero, and the SAT circuit is used only on generated instances.

\paragraph{Next analysis.} Attack \(RealSATInfluenceRealizer_{N}^{arb}\) for one concrete semantic pruning rule: bounded Tseitin self-reduction with a SAT-truth branch-elimination premise. Either produce a uniform charged list of all influencing children or formalize the first semantic line where the private-row countermodel survives.

\paragraph{Falsifiers.}
\begin{itemize}
\item first-unlisted-child DTSR private-row falsifier for hidden semantic pruning
\item unread-enable-bit FormulaEval falsifier when reachable influence is not prefix-closed
\item \(q<\epsilon\) support-density contradiction for slot-bounded transparent certificates
\item positive-beta substitution cap q+beta*c+o(1) for black-box hardness magnification claims
\item nonuniform SAT-truth selector hiding eta,tau,kappa,theta outside \(\Gamma_\mathrm{DTSR}\)
\end{itemize}

\subsection*{cycle-0134: Hardness-Magnification Bridge Realization}

\paragraph{Execution.} research; azure\_responses\_api; attempt primary; 105.592 seconds; 3 lemmas, 5 equations, 7 derivation steps, 5 falsifiers.

\paragraph{Learning.} Cycle 134 sharpens the DTSR boundary: complete declaration must be counterfactual, not merely actual-branch. This yields a valid restricted \(\Gamma_\mathrm{CTSR}\) parser theorem and a sharper private-row falsifier for hidden adaptive pruning. The remaining full semantic escape is exactly a nonrelativizing \(RealSATInfluenceRealizer_{N}^{arb}\), not generic barrier folklore or evidence for \(\mathrm{P} \ne \mathrm{NP}\).

\paragraph{Pressure.} The restricted counterfactual DTSR theorem is proof-shaped, but full semantic adaptive pruning remains blocked at \(RealSATInfluenceRealizer_{N}^{arb}\): one must prove that real SAT semantics uniformly determines and charges every hidden influential row, or accept the private-row no-go.

\paragraph{Derivation steps.}
\begin{enumerate}
\item Define the DTSR row bit fields as (parent row r, formula index j, restriction word w, branch guard b, Tseitin seed t, child formula id c, occurrence-neighborhood list \(O_{c}\), and proof-line parent pointers).
\item Define the closure operator Phi(S): add ordinary bEF dependencies, transparent FormulaEval row supports, and for every reachable DTSR parent in S add every syntactic child Child(r,w,b,t) allowed by the bounded rule, not only the child on the semantic actual branch.
\item Iterate Phi from the final \(\mathrm{LB}_{N}\) line to a least fixed point K\_CTSR(pi). Since each line has polynomially bounded visible fanout and bounded DTSR depth, the fixed point is finite and auditable from the transcript.
\item For each line, maintain a circuit invariant: the parsed circuit computes the line predicate for every completion of the SAT labels on K\_CTSR(pi). bEF and FormulaEval cases use the old \(\Gamma_{dag}\) proof; the DTSR case is valid because all counterfactual child labels used by the self-reduction identity are already present in K\_CTSR(pi).
\item Charge the DTSR additions by the number of reachable parent rows times the bounded child fanout and occurrence expansion, giving \(N^{1+\kappa+\theta+o(1)}\) in the row-expanded ledger.
\item If the semantic rule declares only the actual visited branch, choose the first omitted influential sibling. Build two row-label completions agreeing on the transcript and declared rows but differing on that private child; the parent conclusion and hence one carrier parity bit differ. This is the first-unlisted-child falsifier.
\item Combine the generator circuit, parsed row circuits, and equality patches for \(E_{pi}\) residual rows to obtain C(H\_N\textasciicircum{}SAT)<=N\textasciicircum{}\{rho+o(1)\}+Gamma\_CTSR(pi)+O(E\_pi N). A subthreshold ledger contradicts the claimed LB\_N(1+epsilon) source certificate.
\end{enumerate}

\paragraph{Equation 1: carrier\_parameters.}
\[
M_N=N^{q+o(1)}, L_N=(log N)^d
\]
\emph{Valid when:} Valid when the carrier length convention fixes \(M_{N}\) and \(L_{N}\) uniformly before parsing.

\paragraph{Equation 2: counterfactual\_child\_count.}
\[
|K_CTSR(pi) \ K_dag(pi)| \le N^{kappa+theta+o(1)}
\]
\emph{Valid when:} Applies after each declared TSR parent row exposes all branch words and Tseitin variants.

\paragraph{Equation 3: ctsr\_ledger.}
\[
Gamma_CTSR(pi)=Gamma_bEF(pi)+U(K_dag)+X(K_dag)+N^{1+kappa+theta+o(1)}
\]
\emph{Valid when:} Requires the support closure to contain every counterfactual child before semantic rule use.

\paragraph{Equation 4: parser\_upper\_bound.}
\[
C(H_N^SAT) \le N^{rho+o(1)}+Gamma_CTSR(pi)+O(E_pi N)
\]
\emph{Valid when:} Applies after the generator circuit and transcript parser are in the same nonuniform SIZE model.

\paragraph{Equation 5: subthreshold\_falsifier\_window.}
\[
N^{rho+o(1)}+Gamma_CTSR(pi)+O(E_pi N) < N^{1+epsilon-o(1)}
\]
\emph{Valid when:} Requires \(\rho<1+\epsilon\) and every charged ledger exponent below 1+epsilon.

\paragraph{Next analysis.} Instantiate the counterfactual DTSR closure on a two-level nested TSR proof with one Tseitin seed family; compute the exact child count, \(K_\mathrm{CTSR}\) fixed point, \(\Gamma_\mathrm{CTSR}\) ledger, and first-unlisted-child countermodel transcript.

\paragraph{Falsifiers.}
\begin{itemize}
\item Actual-branch-only DTSR transcript: the first unlisted influential sibling gives a private-row falsifier while preserving every visible proof bit.
\item Unread-enable-bit FormulaEval: a hidden enable row can switch support reachability after \(K_{dag}\) is computed, invalidating the charged parser unless prefix closure includes it.
\item Threshold collapse: if \(\kappa+\theta\ge \epsilon\), complete counterfactual declarations may be valid but no longer yield a subthreshold \(\mathrm{LB}_{N}\) falsifier.
\item Nonuniform leakage: if child declarations depend on SAT truth labels or advice, the support cone is not advice-free and the same-model transfer is invalid.
\item Positive \(\beta\) fallback: any reinterpretation with target-query exponent \(\beta>0\) reintroduces the q+beta c transfer cap and cannot claim all-exponent exclusion.
\end{itemize}

\subsection*{cycle-0135: Hardness-Magnification Bridge Realization}

\paragraph{Execution.} research; azure\_responses\_api; attempt primary; 136.06 seconds; 3 lemmas, 5 equations, 6 derivation steps, 5 falsifiers.

\paragraph{Learning.} Cycle 135 upgrades declared TSR to counterfactual declared TSR: complete counterfactual child support is sufficient for the \(\Gamma_\mathrm{CTSR}\) parser and necessary for extensional safety. The exact new falsifier is first-hidden-child adaptive pruning; the only full semantic escape is a charged nonrelativizing \(RealSATInfluenceRealizer_{N}^{arb}\). This is a route-level theorem/countermodel boundary, not evidence for \(\mathrm{P} \ne \mathrm{NP}\).

\paragraph{Pressure.} The two-world private-row argument is a no-go for extensional or relativizing hidden support, not a proof against real SAT-specific semantic pruning. Full semantic pruning remains blocked exactly at \(RealSATInfluenceRealizer_{N}^{arb}\).

\paragraph{Derivation steps.}
\begin{enumerate}
\item Define \(K_\mathrm{CTSR}^{0}\) from the final \(\mathrm{LB}_{N}\) line and iterate a finite monotone closure: add EF parent rows, SATeq quotient links, FormulaEval support rows certified by visible enumerators or full scan, and for each CTSR node add every child row in the counterfactual child list.
\item For a CTSR line with parent (x,j,a,sigma,t), the local verifier recomputes each child restriction formula, checks the child row encoding, and checks the bounded Tseitin OR/equality certificate in the same proof theory; the per-child bit overhead is polylogarithmic, hence \(N^{o(1)}\) after row expansion.
\item Line induction invariant: after processing any prefix of pi, every formula-label dependency reachable from that prefix lies in \(K_\mathrm{CTSR}\) and contributes to \(\Gamma_{bEF}\), U, X, or \(T_\mathrm{CTSR}\); no later symbolic use may introduce an uncharged row because prefix closure re-runs the fixed point.
\item Circuit extraction: combine the generator circuit of size \(N^{\rho+o(1)}\), the parsed proof fragments charged by \(\Gamma_\mathrm{CTSR}\), and equality patches for \(E_{pi}\) residual rows, giving C(H\_N)<=N\textasciicircum{}\{rho+o(1)\}+Gamma\_CTSR+O(E\_pi N).
\item Threshold comparison: if \(\kappa+\theta<\epsilon\), \(\rho<1+\epsilon\), and \(E_{pi}<N^{\epsilon-o(1)}\), the extracted circuit is below \(N^{1+\epsilon-o(1)}\), so the transcript is a falsifier for a claimed same-model LB\_N(1+epsilon) certificate.
\item Necessity attempt: if a first reachable child row c* is used semantically but not listed, build two label worlds agreeing on every visible row and proof bit while flipping an unread enable bit for c*; the parent CTSR value and therefore one carrier parity row can differ, so surface occurrence completion is not sound.
\end{enumerate}

\paragraph{Equation 1: carrier.}
\[
H_N^{\mathrm{SAT}}(x)=\bigoplus_{j\le M_N} \mathrm{SAT}(G_N(x,j)),\quad M_N=N^{q+o(1)},\quad |G_N(x,j)|=(\log N)^{O(1)}
\]
\emph{Valid when:} Valid when the beta-zero parity carrier and generator family are fixed uniformly for all source lengths N.

\paragraph{Equation 2: ctsr-ledger.}
\[
\Gamma_{\mathrm{CTSR}}(\pi)=\Gamma_{\mathrm{bEF}}(\pi)+U(K_{\mathrm{CTSR}})+X(K_{\mathrm{CTSR}})+T_{\mathrm{CTSR}}(\pi)
\]
\emph{Valid when:} Applies after \(K_\mathrm{CTSR}\) is computed as the least influence-closed fixed point visible from the proof syntax.

\paragraph{Equation 3: counterfactual-tsr-charge.}
\[
T_{\mathrm{CTSR}}(\pi)=N\cdot |\mathrm{ChildDecl}_{\mathrm{cf}}(\pi)|\cdot N^{o(1)}=N^{1+\kappa+\theta+o(1)}
\]
\emph{Valid when:} Requires complete counterfactual child declaration with \(N^{\kappa+\theta+o(1)}\) declared row children and polylogarithmic row encodings.

\paragraph{Equation 4: parser-upper-bound.}
\[
C_{\mathrm{SIZE}}(H_N)\le N^{\rho+o(1)}+\Gamma_{\mathrm{CTSR}}(\pi)+O(E_\pi N)
\]
\emph{Valid when:} Valid when the line induction has converted every visible proof dependency into a same-model circuit fragment and \(E_{pi}\) residual rows are equality patched.

\paragraph{Equation 5: subthreshold-falsifier.}
\[
\max{\rho,1+\kappa+\theta,1+\log_N E_\pi}< 1+\epsilon\Rightarrow C(H_N)< N^{1+\epsilon-o(1)}
\]
\emph{Valid when:} Applies in the nonuniform SIZE regime where all \(\Gamma\) components are charged and compared against the claimed \(\mathrm{LB}_{N}\) threshold.

\paragraph{Next analysis.} Resolve \(RealSATInfluenceRealizer_{N}^{arb}\) for one concrete semantic rule: bounded Tseitin self-reduction with adaptive unsat-branch pruning. Either produce a syntax-and-N computable influence row list \(R_{A}\) with charge in \(\Gamma_{real}\), or formalize the first unchargeable semantic inference as a private-row no-realizer theorem.

\paragraph{Falsifiers.}
\begin{itemize}
\item first-hidden-counterfactual-child private-row falsifier for actual-branch-only DTSR declarations
\item unread-enable-bit falsifier for FormulaEval lines whose later symbolic use reaches an uncharged row
\item threshold falsifier when \(\kappa+\theta\ge \epsilon\) or \(\rho\ge 1+\epsilon\) because \(\Gamma_\mathrm{CTSR}\) no longer contradicts LB\_N(1+epsilon)
\item nonuniform leakage falsifier when a row set is selected by SAT truth advice rather than proof syntax and N
\item model mismatch falsifier if the extracted upper bound changes circuit class, uses oracle labels, or leaves the same SIZE model
\end{itemize}

\subsection*{cycle-0136: Hardness-Magnification Bridge Realization}

\paragraph{Execution.} research; azure\_responses\_api; attempt primary; 112.595 seconds; 3 lemmas, 6 equations, 7 derivation steps, 6 falsifiers.

\paragraph{Learning.} Cycle 136 turns counterfactual declared TSR from a boundary slogan into a bitfield-level closed-support theorem: K\_CTSR(pi) is a finite fixed point, \(\Gamma_\mathrm{CTSR}\) is derived by line induction, and first-hidden-child adaptive pruning is the exact falsifier for extensional safety. The valid contribution remains a scoped parser/falsifier theorem for beta-zero parity magnification certificates; it does not prove \(H_{N}\) lower bounds or \(\mathrm{P} \ne \mathrm{NP}\).

\paragraph{Pressure.} The derivation does not handle arbitrary semantic branch pruning. The named open bridge is \(RealSATInfluenceRealizer_{N}^{arb}\): from real SAT semantics and proof syntax alone, uniformly list every influencing child row with charged overhead before the semantic rule is used.

\paragraph{Derivation steps.}
\begin{enumerate}
\item Define a CTSR proof line as a finite record with explicit parent\_row, formula\_id, restriction\_word, depth\_counter, branch-coordinate list, child-row identifiers, Tseitin-map id, enable-bit source, and occurrence-neighborhood pointers.
\item Compute K\_CTSR(pi) by a monotone fixed point from the final \(\mathrm{LB}_{N}\) line: add bEF dependencies, transparent FormulaEval support, same-theory row expansions, and every counterfactual CTSR child of each reached parent.
\item For each row in \(K_\mathrm{CTSR}\), expand the referenced SAT formula occurrences and attach the local verifier showing \(\mathrm{SAT}(parent)=OR_{b}\) SAT(child\_b) or the declared bounded self-reduction relation used by the rule.
\item Induct on proof lines: bEF and SATeq lines use \(\Gamma_{bEF}\), transparent FormulaEval lines use U(K)+X(K), and CTSR lines use only already charged counterfactual child rows.
\item Build the transcript circuit \(A_{pi}\) from the generator plus the parsed proof transcript; its size is bounded by N\textasciicircum{}\{rho+o(1)\}+Gamma\_CTSR(pi).
\item If \(A_{pi}\) has residual error on \(E_{pi}\) source rows, patch by equality gates of total size O(E\_pi N), obtaining the displayed upper bound for \(H_{N}^{\mathrm{SAT}}\).
\item Contrapose: any claimed LB\_N(1+epsilon) proof in this restricted model with subthreshold \(\Gamma_\mathrm{CTSR}\) is refuted in the same SIZE model.
\end{enumerate}

\paragraph{Equation 1: carrier.}
\[
H_N^{SAT}(x)=\bigoplus_{j< M_N}\mathrm{SAT}(G_N(x,j)),\quad M_N=N^{q+o(1)},\quad |G_N(x,j)|=(\log N)^d
\]
\emph{Valid when:} Valid when N ranges over the carrier lengths and \(G_{N}\) is the fixed advice-free generator.

\paragraph{Equation 2: least\_closed\_support.}
\[
K_{t+1}=K_t\cup \mathrm{Dep}_{bEF}(K_t)\cup \mathrm{FE}(K_t)\cup \mathrm{CTSR}_{cf}(K_t),\quad K_{CTSR}=\bigcup_t K_t
\]
\emph{Valid when:} Valid when proof syntax exposes parent pointers, FormulaEval support, and counterfactual CTSR child tables.

\paragraph{Equation 3: ledger.}
\[
\Gamma_{CTSR}(\pi)=\Gamma_{bEF}(\pi)+U(K_{CTSR})+X(K_{CTSR})+B_{CTSR}(\pi)
\]
\emph{Valid when:} Valid when each line is parsed in the same nonuniform SIZE model and all reachable influences are charged.

\paragraph{Equation 4: ctsr\_charge.}
\[
B_{CTSR}(\pi)\le N^{1+\kappa+\theta+o(1)}
\]
\emph{Valid when:} Valid when each active row has \(N^{\kappa+\theta+o(1)}\) declared children and row expansion costs Theta(N).

\paragraph{Equation 5: transcript\_upper\_bound.}
\[
C(H_N^{SAT})\le N^{\rho+o(1)}+\Gamma_{CTSR}(\pi)+O(E_\pi N)
\]
\emph{Valid when:} Valid when the generated proof transcript has \(E_{pi}\) residual bad rows and equality patching is allowed in SIZE.

\paragraph{Equation 6: source\_lb\_falsifier.}
\[
N^{\rho+o(1)}+\Gamma_{CTSR}(\pi)+O(E_\pi N)< N^{1+\epsilon-o(1)}\Rightarrow \neg LB_N(1+\epsilon)
\]
\emph{Valid when:} Valid when the inequality is evaluated in the same source circuit model as \(\mathrm{LB}_{N}\).

\paragraph{Next analysis.} Select one concrete full semantic pruning rule beyond CTSR declarations, preferably adaptive Tseitin branch pruning with real SAT contradiction tests, and attempt either a charged \(RealSATInfluenceRealizer_{N}^{arb}\) or a minimal first-unchargeable-row countermodel.

\paragraph{Falsifiers.}
\begin{itemize}
\item first-hidden-child adaptive pruning: actual-branch-only declaration lets two transcripts agree while the carrier parity differs
\item unread-enable-bit FormulaEval: the selected row is semantically live but absent from \(K_\mathrm{CTSR}\)
\item SAT-truth advice in child selection: destroys advice-free support enumeration and invalidates the parser
\item positive-beta substitution leakage: if target query length is \(N^\beta\) with \(\beta>0\), the upper bound q+beta c caps all-exponent transfer
\item threshold collapse: \(\kappa+\theta\ge \epsilon\) makes CTSR expansion alone too expensive for the subthreshold falsifier
\item nonuniform row selector leakage: a short selector hardwires the charged support without surfacing verifier-visible dependencies
\end{itemize}

\subsection*{cycle-0137: Hardness-Magnification Bridge Realization}

\paragraph{Execution.} research; azure\_responses\_api; attempt primary; 138.76 seconds; 3 lemmas, 4 equations, 6 derivation steps, 5 falsifiers.

\paragraph{Learning.} Cycle 137 improves the promoted CTSR packet by proving minimality pressure, not just sufficiency: complete counterfactual child support is the least valid transparent object for adaptive TSR. Removing one reachable child gives a concrete first-hidden-child falsifier; adding all children gives the exact \(\Gamma_\mathrm{CTSR}\) threshold \(\kappa+\theta<\epsilon\). The learning is a scoped parser/countermodel contribution, not a lower bound for \(H_{N}\) and not \(\mathrm{P} \ne \mathrm{NP}\) progress.

\paragraph{Pressure.} The derivation does not handle full semantic branch pruning where real SAT itself proves the omitted child determined by charged rows. That is exactly \(RealSATInfluenceRealizer_{N}^{arb}\): an advice-free, nonrelativizing realizer must list and charge all influencing rows before the semantic rule is used.

\paragraph{Derivation steps.}
\begin{enumerate}
\item Define \(K_{0}\) as the final \(\mathrm{LB}_{N}\) line; iteratively add visible proof-DAG parents, bit-explicit EF supports, SATeq-linked rows, \(bounded-SR_{h}\) leaves, FormulaEval enumerator outputs, and every counterfactual child row declared by a CTSR rule touching the current set.
\item Termination is immediate because pi has finitely many line records and each CTSR record declares at most \(N^{\kappa+\theta+o(1)}\) child rows with finite bit encodings; the least fixed point is the union of all iterations.
\item Line induction: for bEF lines use the cycle-98/102 \(\Gamma_{bEF}\) parser; for transparent FormulaEval lines replace the symbolic line by its \(K_\mathrm{CTSR}\) row expansion charged in U+X; for CTSR lines use the complete counterfactual child table to compute the same restricted self-reduction value without semantic pruning advice.
\item The constructed circuit first computes \(G_{N}\) with cost \(N^{\rho+o(1)}\), then evaluates the charged transcript fragments and combines them according to the proof's final \(parity/\mathrm{LB}_{N}\) line; sparse residual rows, if any, are equality-patched at O(E\_pi N).
\item If a reachable row r is omitted, choose the first proof line l whose extensional value depends on r. Build two row-oracle worlds agreeing on K\_CTSR(pi) minus r and all bit-visible records, but differing on r so that l and ultimately the asserted occurrence completion differ. The verifier cannot distinguish the worlds, proving the private-row falsifier.
\item The quantitative contradiction is only against the source LB\_N(1+epsilon) claim: if N\textasciicircum{}\{rho+o(1)\}+Gamma\_CTSR(pi)<N\textasciicircum{}\{1+epsilon-o(1)\}, the parsed circuit is a valid same-model upper bound for \(H_{N}\).
\end{enumerate}

\paragraph{Equation 1: least\_fixed\_point.}
\[
K_{CTSR}(\pi)=\mu K\,[Final(\pi)\cup Parents(K)\cup FE(K)\cup EF(K)\cup SATeq(K)\cup SR_h(K)\cup CChild_{TSR}(K)]
\]
\emph{Valid when:} Valid when pi is a finite \(\mathrm{Psys}_{sym}^\mathrm{CTSR}\) proof whose rule records expose dependency edges and counterfactual child descriptors advice-free.

\paragraph{Equation 2: ledger.}
\[
\Gamma_{CTSR}(\pi)=\Gamma_{bEF}(\pi)+U(K_{CTSR})+X(K_{CTSR})+N^{1+\kappa+\theta+o(1)}+O(E_\pi N)
\]
\emph{Valid when:} Applies after every emitted row is expanded in the same row theory and the declared TSR family has \(N^{\kappa+o(1)}\) bases and \(N^{\theta+o(1)}\) counterfactual children per base.

\paragraph{Equation 3: parser\_upper\_bound.}
\[
C(H_N^{SAT})\le N^{\rho+o(1)}+\Gamma_{CTSR}(\pi)
\]
\emph{Valid when:} Requires the beta-zero target-preserving substitution regime and a line-induction parser that turns each charged proof line into a circuit fragment for \(H_{N}\).

\paragraph{Equation 4: subthreshold\_falsifier.}
\[
N^{\rho+o(1)}+\Gamma_{CTSR}(\pi)< N^{1+\epsilon-o(1)}\Rightarrow \neg LB_N(1+\epsilon)
\]
\emph{Valid when:} Valid when LB\_N(1+epsilon) is the same nonuniform SIZE lower-bound assertion for \(H_{N}\) and sparse residual errors are already included in \(\Gamma_\mathrm{CTSR}\).

\paragraph{Next analysis.} Write the standalone L137 first-hidden-child theorem: specify the CTSR bit fields, define the two row-oracle worlds, prove transcript equality on \(K_\mathrm{CTSR}\) minus r, prove final parity disagreement, and compute the \(\Gamma_\mathrm{CTSR}\) threshold for \(\kappa+\theta<\epsilon\) with q-trichotomy support density.

\paragraph{Falsifiers.}
\begin{itemize}
\item first-hidden-child adaptive pruning: actual-branch declarations undercharge a counterfactual child row that can change the semantic TSR value
\item private-row/unread-enable-bit falsifier: two transcripts agree on all surfaced support but differ on an uncharged row influencing \(H_{N}\)
\item threshold collapse falsifier: \(\rho\) or kappa+theta reaches \(\epsilon\), making the parser upper bound no longer subthreshold
\item nonuniform leakage falsifier: a selector or row list is encoded as advice rather than uniformly recovered from pi and N
\item positive-beta black-box substitution falsifier remains active for any carrier extension with query length \(N^\beta\) and \(\beta>0\)
\end{itemize}

\subsection*{cycle-0138: Hardness-Magnification Bridge Realization}

\paragraph{Execution.} research; azure\_responses\_api; attempt primary; 124.711 seconds; 3 lemmas, 5 equations, 7 derivation steps, 6 falsifiers.

\paragraph{Learning.} Cycle 138 identifies the exact minimal CTSR object as reachable counterfactual frontier closure. This improves the promoted packet: actual-branch declarations are not valid transparent support, while full reachable child declarations give the \(\Gamma_\mathrm{CTSR}\) parser with threshold \(\kappa+\theta<\epsilon\). The result is a scoped parser/countermodel theorem and live inference-quality packet, not evidence for \(\mathrm{P} \ne \mathrm{NP}\).

\paragraph{Pressure.} The two-world necessity proof is valid for extensional row-label or occurrence-local semantics, but it does not refute a real SAT semantic derivation proving that the hidden child is impossible unless \(RealSATInfluenceRealizer_{N}^{arb}\) fails for that derivation.

\paragraph{Derivation steps.}
\begin{enumerate}
\item Initialize \(K_{0}\) with the final \(\mathrm{LB}_{N}\) line and all verifier-visible parent pointers used to parse its asserted carrier lower bound.
\item For each line in reverse topological order, add ordinary bEF dependencies using the existing \(\Gamma_{bEF}\) parser; add FormulaEval rows only through visible enumerators or full-scan charge.
\item For a CTSR line, read bit fields (\(i,j,a,b,parent,child_{list},restriction_{mask},Tseitin_{map},formula_{hash},occurrence_{neighborhoods}\)). The local verifier checks that each listed child is the correct restricted/Tseitin-transformed formula and that every branch b in the rule alphabet appears, not merely the actual semantic branch.
\item Charge \(U_{N}\) for enumerating child indices and \(X_{N}\) for same-theory occurrence expansion. With N source rows, \(N^{\kappa+o(1)}\) CTSR sites, and \(N^{\theta+o(1)}\) counterfactual children per site, this contributes \(B_\mathrm{CTSR}\le N^{1+\kappa+\theta+o(1)}\).
\item Line induction invariant: after processing a prefix, the hardwired transcript plus generator circuit computes every proof-visible row value using only \(K_\mathrm{CTSR}\) rows and charged expansions. CTSR parent values are reconstructed from all declared counterfactual children, so adaptive pruning does not introduce unseen support.
\item If N\textasciicircum{}\{rho+o(1)\}+Gamma\_CTSR(pi)+O(E\_pi N)<N\textasciicircum{}\{1+epsilon-o(1)\}, the constructed circuit is a same-model upper-bound falsifier for LB\_N(1+epsilon).
\item Necessity attempt: take the first reachable CTSR child absent from \(K_\mathrm{CTSR}\). Build worlds W0,W1 agreeing on all charged rows and all bit fields, toggle only the hidden child by private-row/unread-enable-bit assignment, and propagate the toggled value to the parent and final parity. The verifier sees identical transcripts, so occurrence completion was unsound.
\end{enumerate}

\paragraph{Equation 1: ctsr\_fixed\_point.}
\[
K_{0}=Dep(final),\quad K_{t+1}=K_t\cup Pred(K_t)\cup Child_{CTSR}(K_t),\quad K_{CTSR}=\bigcup_t K_t
\]
\emph{Valid when:} Valid when proof syntax exposes all parent pointers, row indices, and CTSR child declarations.

\paragraph{Equation 2: ctsr\_charge.}
\[
B_{CTSR}(\pi)\le N^{1+\kappa+\theta+o(1)}
\]
\emph{Valid when:} Applies when there are N source rows, \(N^{\kappa+o(1)}\) CTSR sites, and \(N^{\theta+o(1)}\) declared counterfactual children per site.

\paragraph{Equation 3: gamma\_ctsr.}
\[
\Gamma_{CTSR}(\pi)=\Gamma_{bEF}(\pi)+U_N(K_{CTSR})+X_N(K_{CTSR})+B_{CTSR}(\pi)
\]
\emph{Valid when:} Valid after transparent support enumeration and same-theory occurrence expansion are both charged in the proof ledger.

\paragraph{Equation 4: parser\_upper\_bound.}
\[
C_{SIZE}(H_N^{SAT})\le N^{\rho+o(1)}+\Gamma_{CTSR}(\pi)+O(E_\pi N)
\]
\emph{Valid when:} Requires the target-preserving parser simulation in the same nonuniform SIZE model with beta-zero SAT queries.

\paragraph{Equation 5: support\_density\_corollary.}
\[
N\cdot |K_{idx}(\pi)|\ge N^{1+\epsilon-o(1)}\Rightarrow |K_{idx}(\pi)|/M_N\ge N^{\epsilon-q-o(1)}
\]
\emph{Valid when:} Valid when \(M_{N}=N^{q+o(1)}\) slots and \(K_\mathrm{CTSR}\) has no hidden semantic row advice.

\paragraph{Next analysis.} Write the referee-facing CTSRFrontierCompleteness proof: give closure pseudocode for \(K_\mathrm{CTSR}\), prove termination and minimality, verify one adaptive TSR line bit-by-bit, and instantiate the first-hidden-child falsifier on a one-node unread-enable-bit example.

\paragraph{Falsifiers.}
\begin{itemize}
\item first-hidden-child falsifier: a reachable counterfactual CTSR child is omitted while later proof lines still depend on its parent value
\item actual-branch-only falsifier: the proof declares only the semantically followed branch and hides the unused branch behind an unread enable bit
\item nonuniform row-selector falsifier: \(K_\mathrm{CTSR}\) is chosen with advice not recoverable from proof syntax and N
\item subthreshold residual falsifier: \(E_{pi}<N^{\epsilon-o(1)}\) lets equality patching erase sparse row errors
\item positive-beta transfer falsifier: any attempted return to black-box SAT substitution with \(\beta>0\) is capped by q+beta c+o(1) before all exponents
\item semantic-pruning escape test: a valid \(RealSATInfluenceRealizer_{N}^{arb}\) listing and charging all rows neutralizes the private-row countermodel for that rule
\end{itemize}

\subsection*{cycle-0139: Hardness-Magnification Bridge Realization}

\paragraph{Execution.} research; azure\_responses\_api; attempt primary; 135.774 seconds; 3 lemmas, 5 equations, 7 derivation steps, 6 falsifiers.

\paragraph{Learning.} Cycle 139 improves the CTSR frontier packet by adding a safe EqSAT quotient rule: duplicate counterfactual child rows reduce \(\Gamma\) only when a same-theory certificate is surfaced and charged. Otherwise quotienting is exactly where the private-row falsifier attacks. This is a valid scoped parser/countermodel contribution and still not evidence for \(\mathrm{P} \ne \mathrm{NP}\).

\paragraph{Pressure.} \(RealSATInfluenceRealizer_{N}^{arb}\) remains open for fully semantic adaptive pruning that never exposes counterfactual children or EqSAT certificate edges yet may be sound for real SAT instances.

\paragraph{Derivation steps.}
\begin{enumerate}
\item Bit fields fixed for each relevant line: rule tag, line id, parent ids, slot j, formula handle h, restriction path sigma, branch bit b, enable bit e, frontier flag, EqSAT certificate pointer, and row-expansion handle.
\item Define K\_CTSR\textasciicircum{}=(pi) by a least fixed point from the final \(\mathrm{LB}_{N}\) line. Closure adds syntactic parents, FormulaEval row sets, EqSAT dependency rows, and both counterfactual children of every reachable TSR node. Quotienting is delayed until a charged EqSAT certificate edge is encountered.
\item Verifier audit: enumerate stabilized keys by BFS over the exposed proof DAG; charge \(U_{enum}\) for enumeration, \(B_\mathrm{EqSAT}\) for certificate verification, and \(N^{1+o(1)}\) for expanding each distinct quotient row into the occurrence-level parser.
\item Line induction: bEF and SATeq lines use the existing \(\Gamma_{bEF}\) invariant; FormulaEval lines read only keys already in \(K^=\); CTSR lines are sound because all reachable children, including unchosen branches, are present; EqSAT lines justify only the certified quotient edges.
\item Circuit extraction: replace every parsed row occurrence by the generated SAT subcircuit when available or by the transcript subcircuit charged in \(\Gamma_\mathrm{CTSR}^=\); patch \(E_{pi}\) exceptional source rows with equality gates, giving C(H\_N)<=N\textasciicircum{}\{rho+o(1)\}+Gamma\_CTSR\textasciicircum{}=+O(E\_pi N).
\item Contrapositive: if this upper bound is below \(N^{1+\epsilon-o(1)}\), the alleged LB\_N(1+epsilon) certificate is invalid in the same nonuniform SIZE model; this is a route-level falsifier, not a proof that \(H_{N}\) is easy unconditionally.
\item Necessity test: remove one reachable child or merge two uncertified keys. Choose the first unsafe use, keep all visible fields fixed, and flip the private row label behind the omitted child or quotient collision. The transcript agrees but the induced premise or final parity differs.
\end{enumerate}

\paragraph{Equation 1: carrier.}
\[
H_N^{SAT}(x)=\bigoplus_{j\le M_N} SAT(G_N(x,j)),\quad M_N=N^{q+o(1)}
\]
\emph{Valid when:} Valid when N ranges over the carrier lengths, \(M_{N}=N^{q+o(1)}\), and \(G_{N}\) is the fixed advice-free generator used by the beta-zero parity route.

\paragraph{Equation 2: safe\_frontier\_fixed\_point.}
\[
K_{0}=Rows(final),\quad K_{t+1}=K_t\cup Parents(K_t)\cup FE(K_t)\cup CTSR_{all}(K_t)\cup EqSAT(K_t)\;/\equiv_{cert}
\]
\emph{Valid when:} Valid when every rule line exposes parent pointers and every quotient edge has a bit-explicit EqSAT certificate charged before use.

\paragraph{Equation 3: ctsr\_gamma.}
\[
\Gamma_{CTSR}^{=}(\pi)=\Gamma_{bEF}(\pi)+U_{enum}(K^{=})+N\,|K^{=}_{CTSR}(\pi)|+B_{SR_h}(\pi)+B_{EqSAT}(\pi)
\]
\emph{Valid when:} Valid when K\textasciicircum{}\{=\}\_\{CTSR\}(pi) is the stabilized safe quotient frontier and row expansion costs \(N^{1+o(1)}\) per distinct charged row key.

\paragraph{Equation 4: parser\_upper\_bound.}
\[
C_{SIZE}(H_N^{SAT})\le N^{\rho+o(1)}+\Gamma_{CTSR}^{=}(\pi)+O(E_\pi N)
\]
\emph{Valid when:} Applies after the line induction parses every proof line into a same-theory circuit transcript and \(E_{pi}\) exceptional source rows are patched by equality gates.

\paragraph{Equation 5: subthreshold\_falsifier.}
\[
\rho< 1+\epsilon,\quad |K^{=}_{CTSR}|< N^{\epsilon-o(1)},\quad E_\pi< N^{\epsilon-o(1)}\Rightarrow C(H_N^{SAT})< N^{1+\epsilon-o(1)}
\]
\emph{Valid when:} Requires beta-zero target queries, charged EqSAT quotienting, and no hidden FormulaEval or TSR influence outside the stabilized support frontier.

\paragraph{Next analysis.} Instantiate \(\Gamma_\mathrm{CTSR}^=\) on a projected-dDNNF FormulaEval proof with two Tseitin-equivalent child formulas: compute K before and after certified EqSAT quotienting, then remove one certificate and exhibit the first uncertified quotient-collision private-row falsifier.

\paragraph{Falsifiers.}
\begin{itemize}
\item First-hidden-child adaptive pruning falsifier: an unlisted counterfactual CTSR child changes the induced row label while the visible transcript is unchanged.
\item Uncertified quotient-collision falsifier: two syntactically merged child keys can be independently flipped unless a same-theory EqSAT certificate is charged.
\item Unread-enable-bit falsifier: a FormulaEval row selected by a hidden enable bit influences a later line without entering K\_CTSR\textasciicircum{}=(pi).
\item \(q<\epsilon\) full-scan enumeration falsifier: if the whole visible slot universe has size \(N^{q+o(1)}\) with \(q<\epsilon\), expanding all rows is already subthreshold.
\item Positive-beta transfer-cap falsifier retained from the atlas: any return to black-box target queries of length \(n^\beta\) faces the upper-bound window q+beta c+o(1).
\item Nonuniform row-selector advice falsifier: a short proof that names a support set only through semantic advice is outside \(TransparentRowSupport_{N}\) and cannot be promoted.
\end{itemize}

\subsection*{cycle-0140: Hardness-Magnification Bridge Realization: QCTSR Certified-Quotient Parser Boundary}

\paragraph{Execution.} research; azure\_responses\_api; attempt primary; 106.619 seconds; 3 lemmas, 5 equations, 7 derivation steps, 6 falsifiers.

\paragraph{Learning.} Cycle 140 turns safe EqSAT quotienting from a slogan into a ledger transformation: duplicate counterfactual CTSR rows may be removed only through a surfaced same-theory certificate forest whose dependencies are in \(K_{Q}\). The exact new falsifier is the first uncertified quotient merge. This strengthens the scoped parser/countermodel contribution and remains separate from proving any lower bound or \(\mathrm{P} \ne \mathrm{NP}\).

\paragraph{Pressure.} The derivation does not handle arbitrary real-SAT semantic pruning. It requires either surfaced same-theory EqSAT certificates or a future \(RealSATInfluenceRealizer_{N}^{arb}\) that uniformly lists and charges all semantic influences before use.

\paragraph{Derivation steps.}
\begin{enumerate}
\item Define the visible row universe \(R_{N}\) and initialize K with the final \(\mathrm{LB}_{N}\) line, final carrier rows, and explicit bEF dependencies.
\item Run the least fixed point: add parent rows of every reachable FormulaEval line, add every counterfactual TSR child of each reachable representative, and add all dependencies of each EqSAT quotient certificate before allowing its merge.
\item Use union-find only after certificate insertion: a suppressed row r may be replaced by representative s in the parser iff the certificate forest contains a path proving same-theory SAT equivalence between r and s.
\item Line induction cases: bEF lines charge \(\Gamma_{bEF}\); transparent FormulaEval lines charge \(U(K_{Q})+N|RepIdx|\); CTSR lines charge all reachable counterfactual children; quotient lines charge \(C_{Eq}\) and then reuse the representative row.
\item Construct the \(H_{N}\) circuit by wiring \(G_{N}\), the charged row evaluators/transcripts, representative occurrence expansions, EqSAT certificate checkers if needed, and O(E\_pi N) equality patches for sparse residual rows.
\item Derive C(H\_N)<=N\textasciicircum{}\{rho+o(1)\}+Gamma\_QCTSR(pi). If this is below \(N^{1+\epsilon-o(1)}\), the asserted LB\_N(1+epsilon) is falsified in the same nonuniform model.
\item For necessity, choose the first reachable quotient merge or child-pruning step absent from \(K_{Q}\). Two private-row worlds agree on all charged data and the proof transcript but flip the hidden row's parity contribution, contradicting extensional parser soundness.
\end{enumerate}

\paragraph{Equation 1: carrier.}
\[
H_N(x)=\bigoplus_{j< M_N} SAT(G_N(x,j)),\qquad M_N=N^{q+o(1)},\quad |G_N(x,j)|=(\log N)^d
\]
\emph{Valid when:} Valid when \(H_{N}\) is the beta-zero generated parity carrier with polylogarithmic SAT instances and explicit generator \(G_{N}\).

\paragraph{Equation 2: safe\_quotient\_closure.}
\[
K_Q(\pi)=\mu K\,[Final(\pi)\cup Par(K)\cup CTSR_{cf}(Rep(K))\cup EqCertDep(K)]
\]
\emph{Valid when:} Applies when proof syntax exposes parent pointers, counterfactual TSR children, representatives, and dependencies of every quotient certificate.

\paragraph{Equation 3: ledger.}
\[
\Gamma_{QCTSR}(\pi)=\Gamma_{bEF}(\pi)+U(K_Q)+N\cdot |RepIdx(K_Q)|+C_{Eq}(K_Q)+O(E_\pi N)
\]
\emph{Valid when:} Requires the transparent parser ledger to charge enumeration, row expansion, EqSAT proof forests, and sparse residual patching.

\paragraph{Equation 4: parser\_upper\_bound.}
\[
C(H_N)\le N^{\rho+o(1)}+\Gamma_{QCTSR}(\pi)
\]
\emph{Valid when:} Applies after the line induction replaces every proof use by charged rows in the same nonuniform SIZE model.

\paragraph{Equation 5: subthreshold\_falsifier.}
\[
\rho< 1+\epsilon,\ \Gamma_{QCTSR}(\pi)< N^{1+\epsilon-o(1)}\ \Longrightarrow\ \neg LB_N(1+\epsilon)
\]
\emph{Valid when:} Valid when \(\mathrm{LB}_{N}\) is a same-model circuit lower-bound claim for \(H_{N}\) on the same valid length sequence.

\paragraph{Next analysis.} Instantiate \(K_{Q}\) on a toy projected-dDNNF FormulaEval proof with two CTSR levels and three EqSAT quotient classes; compute |RepIdx|, \(C_{Eq}\), \(\Gamma_\mathrm{QCTSR}\), then delete one quotient certificate to exhibit the first-hidden-merge falsifier.

\paragraph{Falsifiers.}
\begin{itemize}
\item Uncertified quotient merge: a suppressed row influences the final parity but has no surfaced same-theory certificate to a representative.
\item First-hidden-child adaptive pruning: a reachable counterfactual TSR child is omitted because only the actual branch was declared.
\item Unread-enable-bit FormulaEval: selector output hides a row whose later symbolic use reaches the final line.
\item Positive \(\beta\) target queries: black-box substitution gives q+beta*c+o(1) and cannot exclude all SAT circuit exponents.
\item Nonuniform leakage: representative choices or row sets are supplied as semantic advice instead of proof-visible enumeration.
\item Certificate-charge collapse: C\_Eq(K\_Q) itself is \(N^{1+\epsilon-o(1)}\) or larger, so the parser theorem remains valid but gives no subthreshold falsifier.
\end{itemize}

\subsection*{cycle-0141: Hardness-Magnification Bridge Realization}

\paragraph{Execution.} research; azure\_responses\_api; attempt primary; 120.354 seconds; 2 lemmas, 6 equations, 8 derivation steps, 6 falsifiers.

\paragraph{Learning.} The new contribution is that quotient compression in counterfactual TSR is not a free \(\Gamma\) reduction. It is valid only as a dependency-closed, acyclic same-theory certificate transformation; otherwise the first uncertified merge gives a private-row falsifier. This closes a proof-audit loophole in the restricted transparent-support theorem while preserving the separation between route-level certificate collapse and the main P vs NP conjecture.

\paragraph{Pressure.} The derivation is complete only for syntactic same-theory quotient certificates. It blocks for full semantic FormulaEval or adaptive pruning where the set of influencing rows is determined by real SAT truth rather than proof-visible dependencies; this is exactly \(RealSATInfluenceRealizer_{N}^{arb}\).

\paragraph{Derivation steps.}
\begin{enumerate}
\item Define proof-line records with fields line\_id, rule\_tag, parent\_ids, row\_slot, formula\_id, restriction\_word, quotient\_id, certificate\_root, and local verifier bits.
\item Initialize \(K_{Q}\) with the final \(\mathrm{LB}_{N}\) line and close backward under bit-explicit EF parents, SATeq links, bounded-SR leaves, and FormulaEval surfaced rows.
\item For each CTSR node in \(K_{Q}\), add every syntactically reachable counterfactual restriction child, not merely the actual branch; this preserves extensional safety under private-row changes.
\item For each quotient merge in \(K_{Q}\), require a same-theory certificate forest whose leaves are earlier members of \(K_{Q}\) and whose conclusion is EqSAT(child\_a,child\_b); reject cyclic certificates by topological order on line\_id and certificate\_root.
\item Charge N R\_Q(pi) for row expansion, \(N^{q+\eta+\tau+o(1)}\) for bounded-SR leaves, C\_Q(pi) for certificate checking, and inherit Gamma\_bEF(pi) for bit-explicit EF support.
\item Run line induction: each proof line in \(K_{Q}\) has a local circuit simulator using only charged parents; quotient lines reuse a representative only after the certificate verifier establishes same-theory equality inside the transcript.
\item Assemble the carrier circuit from the generator circuit, charged row simulators, quotient representatives, and equality patches for \(E_{N}\) residual rows, giving C(H\_N)<=N\textasciicircum{}\{rho+o(1)\}+Gamma\_Q+O(E\_NN).
\item Derive the exponent window: if \(\rho<1+\epsilon\), \(q+\eta+\tau<1+\epsilon\), \(\kappa+\theta<\epsilon\), \(chi<1+\epsilon\), and \(\delta<\epsilon\), then the claimed LB\_N(1+epsilon) certificate is invalid in the same SIZE model.
\end{enumerate}

\paragraph{Equation 1: carrier.}
\[
H_N^{SAT}(x)=\bigoplus_{j< M_N} SAT(G_N(x,j))
\]
\emph{Valid when:} Valid when N ranges over admissible carrier lengths and \(G_{N}\) is the fixed advice-free generator.

\paragraph{Equation 2: closed\_support.}
\[
K_Q(\pi)=\operatorname{lfp}(S\mapsto S\cup Parents(S)\cup CTSRChild(S)\cup EqCertDeps(S))
\]
\emph{Valid when:} Valid when every FormulaEval, CTSR, and quotient dependency is syntactically addressable before the line using it.

\paragraph{Equation 3: ledger.}
\[
\Gamma_Q(\pi)=\Gamma_{bEF}(\pi)+O(N^{q+\eta+\tau+o(1)}+N\,R_Q(\pi)+C_Q(\pi))
\]
\emph{Valid when:} Valid when SR leaves, row expansions, and certificate-forest checks are charged in the same SIZE model.

\paragraph{Equation 4: transcript\_upper\_bound.}
\[
C_{SIZE}(H_N^{SAT})\le N^{\rho+o(1)}+\Gamma_Q(\pi)+O(E_NN)
\]
\emph{Valid when:} Valid after the line induction has replaced every proof line by its charged transcript simulator.

\paragraph{Equation 5: threshold\_window.}
\[
R_Q\le N^{\kappa+\theta+o(1)},\ C_Q\le N^{\chi+o(1)},\ E_N\le N^{\delta+o(1)}\Rightarrow C(H_N)\le N^{\max{\rho,q+\eta+\tau,1+\kappa+\theta,\chi,1+\delta}+o(1)}
\]
\emph{Valid when:} Valid when \(R_{Q}\), \(C_{Q}\), and \(E_{N}\) have the stated polynomial-exponent bounds uniformly in N.

\paragraph{Equation 6: \(\mathrm{LB}_{falsifier,window}\).}
\[
\max{\rho,q+\eta+\tau,1+\kappa+\theta,\chi,1+\delta}< 1+\epsilon\Rightarrow \neg LB_N(1+\epsilon)
\]
\emph{Valid when:} Valid when \(\mathrm{LB}_{N}\) is a claimed same-model SIZE lower bound for the carrier \(H_{N}\).

\paragraph{Next analysis.} Instantiate one two-merge CTSR proof toy model: one certified EqSAT quotient with charged forest and one uncertified merge. Write the local verifier fields, compute \(R_{Q}\) and \(C_{Q}\), and demonstrate transcript agreement plus parity disagreement for the uncertified merge.

\paragraph{Falsifiers.}
\begin{itemize}
\item first uncertified quotient merge: two transcripts agree on \(K_{Q}\) but differ on the hidden child row
\item cyclic same-theory certificate forest: quotient equality depends on its own compressed representative
\item actual-branch-only TSR declaration: private-row flip at an unvisited counterfactual child changes extensional safety
\item hidden SAT-semantic branch pruning without \(RealSATInfluenceRealizer_{N}^{arb}\): influence rows are not uniformly listable before use
\item threshold collapse: max\{rho,q+eta+tau,1+kappa+theta,chi,1+delta\} reaches or exceeds 1+epsilon
\item nonuniform leakage: quotient representative chosen using SAT-truth advice rather than surfaced proof syntax
\end{itemize}

\subsection*{cycle-0142: Hardness-Magnification Bridge Realization}

\paragraph{Execution.} research; azure\_responses\_api; attempt compact\_retry; 132.632 seconds; 1 lemmas, 3 equations, 5 derivation steps, 4 falsifiers.

\paragraph{Learning.} Cycle 142 makes quotient compression proof-carrying: row deletion is valid only after a ranked same-theory certificate is charged in \(\Gamma_\mathrm{QCTSR}\). Acyclicity is the new finite audit condition; cyclic, adaptive, or advice-chosen merges are concrete countermodels, not generic barrier folklore.

\paragraph{Pressure.} The argument does not realize hidden semantic FormulaEval or adaptive SAT branch pruning; \(RealSATInfluenceRealizer_{N}^{arb}\) remains the named nonrelativizing open lemma.

\paragraph{Derivation steps.}
\begin{enumerate}
\item Definition/assumption context: this derivation is interpreted inside Hardness-Magnification Bridge Realization; source object: \(H_{N}^\mathrm{SAT}(x)=xor_{j\le M_{N}}\) SAT(G\_N(x,j)), \(M_{N}=N^{q+o(1)}\), polylog row length \(L_{N}\), and claimed nonuniform SIZE lower bound LB\_N(1+epsilon).; target object: a same-model upper-bound falsifier for \(\mathrm{LB}_{N}\) via a parsed proof transcript; this is a scoped certificate-collapse theorem, not \(\mathrm{P} \ne \mathrm{NP}\) evidence.; Circuit model is nonuniform Boolean SIZE; \(G_{N}\) has size \(N^{\rho+o(1)}\) and all parser circuits stay in the same model..
\item Initialize \(K_{Q}\) from the final \(\mathrm{LB}_{N}\) assertion; close under bit-explicit parents, FormulaEval row supports, all counterfactual TSR children, quotient certificate leaves, and certificate parents until a finite fixed point is reached.
\item Orient each quotient certificate by a visible rank. The induction hypothesis covers lower-rank dependencies before the merge deletes a duplicate row, so the kept representative is same-theory, not SAT-advice chosen.
\item For each line, add its local verifier and row-expansion cost to \(\Gamma_\mathrm{QCTSR}\). The bEF, transparent FormulaEval, CTSR, and certified quotient cases compose, giving the parser\_bound circuit with O(E\_pi N) equality patching.
\item If the first bad merge is uncertified or cyclic, vary a private duplicated row while preserving all surfaced syntax; the transcript remains accepted but the carrier parity changes, contradicting extensional safety.
\end{enumerate}

\paragraph{Equation 1: ledger.}
\[
\Gamma_{QCTSR}(\pi)=\Gamma_{bEF}(\pi)+U(K_Q)+X(K_Q)+C(F_Q)+N^{1+\kappa+\theta+o(1)}+O(E_\pi N)
\]
\emph{Valid when:} Valid when \(K_{Q}\) is the least closed support cone and \(F_{Q}\) is an acyclic same-theory certificate forest with all leaves surfaced before quotient deletion.

\paragraph{Equation 2: parser\_bound.}
\[
C(H_N^{SAT})\le N^{\rho+o(1)}+\Gamma_{QCTSR}(\pi)
\]
\emph{Valid when:} Applies after the generator, row verifiers, quotient certificates, CTSR children, and sparse-row patches are implemented in the same nonuniform SIZE model.

\paragraph{Equation 3: subthreshold\_falsifier.}
\[
N^{\rho+o(1)}+\Gamma_{QCTSR}(\pi)< N^{1+\epsilon-o(1)}\Rightarrow \neg LB_N(1+\epsilon)
\]
\emph{Valid when:} Requires the claimed \(\mathrm{LB}_{N}\) to be a lower bound for the same \(H_{N}\) carrier and the parser-bound circuit to preserve all length and advice conventions.

\paragraph{Next analysis.} Write the finite verifier for \(K_{Q}\) and ranked \(F_{Q}\) on a toy projected-dDNNF plus CTSR example with one duplicate child pair; compute \(\Gamma\) savings, then remove one certificate edge to exhibit the private-row falsifier.

\paragraph{Falsifiers.}
\begin{itemize}
\item falsifier: a first quotient merge whose representative is chosen by SAT truth or nonuniform advice rather than by a surfaced same-theory certificate.
\item falsifier: a directed cycle in \(F_{Q}\), since the line induction can no longer justify deleting a row before its equivalence proof is parsed.
\item falsifier: an omitted counterfactual TSR child; two transcripts agree on actual branches while a hidden child flips the quotient row.
\item falsifier: any claimed positive-beta extension with transfer cap C(H\_N)<=N\textasciicircum{}\{q+beta c+o(1)\} that cannot exclude all SAT exponents.
\end{itemize}

\subsection*{cycle-0143: Hardness-Magnification Bridge Realization}

\paragraph{Execution.} research; azure\_responses\_api; attempt compact\_retry; 178.729 seconds; 1 lemmas, 3 equations, 4 derivation steps, 4 falsifiers.

\paragraph{Learning.} Cycle 143 turns quotient compression into a proof-carrying ledger rule: \(\Gamma\) may shrink only through ranked acyclic same-theory certificates already in \(K_{Q}\). Unranked quotienting is not a harmless optimization; it is exactly a private-row falsifier.

\paragraph{Pressure.} This does not solve \(RealSATInfluenceRealizer_{N}^{arb}\) for full semantic FormulaEval; it only handles proof-carrying same-theory quotient compression.

\paragraph{Derivation steps.}
\begin{enumerate}
\item Initialize \(K_{Q}\) at the final \(\mathrm{LB}_{N}\) rows; repeatedly add bEF parents, transparent FormulaEval rows, all reachable counterfactual TSR frontier children, and all dependencies of quotient certificates.
\item For a quotient merge a\textasciitilde{}b, require a same-theory certificate DAG whose edges have strictly smaller rank; add its occurrence rows before allowing deletion of either duplicate row from the ledger.
\item Induct on proof lines ordered by proof index and quotient rank. Each consumed row is either explicit in \(K_{Q}\) or replaced by a certified same-theory row; hence the simulator pays exactly \(\Gamma_\mathrm{QCTSR}\).
\item Contraposition: at the first merge lacking a ranked certificate, build two transcripts agreeing on \(K_{Q}\) but differing on the private deleted row or unread enable bit; the proof transcript agrees while \(H_{N}\) parity can differ.
\end{enumerate}

\paragraph{Equation 1: QClosure.}
\[
K_Q(\pi)=\operatorname{lfp}(R\mapsto R\cup Dep_{bEF}(R)\cup FE(R)\cup CTSR(R)\cup Cert_Q(R))
\]
\emph{Valid when:} Valid when proof syntax exposes every FormulaEval, CTSR, and quotient dependency before the consuming line.

\paragraph{Equation 2: QLedger.}
\[
\Gamma_{QCTSR}(\pi)=\Gamma_{bEF}(\pi)+U(K_Q)+X(K_Q)+B_{CTSR}(\pi)+B_Q(\pi)
\]
\emph{Valid when:} Applies after ranked same-theory certificates are charged and duplicate row deletions are certified.

\paragraph{Equation 3: ParserUpperBound.}
\[
C(H_N^{SAT})\le N^{\rho+o(1)}+\Gamma_{QCTSR}(\pi)+O(E_\pi N)
\]
\emph{Valid when:} Requires beta-zero target calls and same nonuniform SIZE model for the reconstructed circuit.

\paragraph{Next analysis.} Instantiate the six-row QCTSR transcript, list row IDs, quotient certificate dependencies, rank labels, and show the first uncertified merge produces two agreeing visible transcripts with opposite \(H_{N}\) parity.

\paragraph{Falsifiers.}
\begin{itemize}
\item First uncertified quotient merge: a deleted duplicate row has no surfaced same-theory certificate in \(K_{Q}\), so private-row spoofing invalidates the parser.
\item Cyclic certificate SCC: ranked descent fails, so the line induction has no well-founded order and the quotient deletion is not valid.
\item Adaptive advice-chosen merge: the row pair is selected by semantic advice rather than N and pi, causing nonuniform leakage outside \(\Gamma_\mathrm{QCTSR}\).
\item Threshold collapse falsifier: if \(1+\kappa+\theta\ge 1+\epsilon\) or \(\Gamma_\mathrm{QCTSR}\ge N^{1+\epsilon-o(1)}\), the parser gives no contradiction to \(\mathrm{LB}_{N}\).
\end{itemize}

\subsection*{cycle-0144: Hardness-Magnification Bridge Realization: ranked QCTSR ledger validation}

\paragraph{Execution.} research; azure\_responses\_api; attempt primary; 121.463 seconds; 3 lemmas, 5 equations, 6 derivation steps, 6 falsifiers.

\paragraph{Learning.} Cycle 144 closes a proof-audit loophole: quotient compression in counterfactual TSR is valid exactly as proof-carrying ranked same-theory compression, not as semantic row deletion. The contribution is a scoped parser/falsifier theorem: ranked QCTSR preserves the \(\Gamma\) line induction, while the first unranked quotient merge is a concrete private-row falsifier. This improves the paper theorem boundary without giving any \(H_{N}\) lower bound or \(\mathrm{P} \ne \mathrm{NP}\) result.

\paragraph{Pressure.} The positive proof is complete only for restricted \(\mathrm{Psys}_\mathrm{QCTSR}\). Full semantic FormulaEval or hidden adaptive pruning still requires \(RealSATInfluenceRealizer_{N}^{arb}\), an advice-free nonrelativizing theorem that lists and charges all real SAT influence rows before use.

\paragraph{Derivation steps.}
\begin{enumerate}
\item Fix pi and compute \(K_{Q}\) by least fixed point: seed final \(\mathrm{LB}_{N}\) line; add bEF parent dependencies; add transparent FormulaEval support rows and row-expansion witnesses; add all reachable counterfactual TSR children; add quotient-certificate dependency rows in increasing rank until closure.
\item Define a quotient edge r\_del \(\to \) r\_keep to be admissible only if a certificate cert(e) proves same-theory SAT(r\_del)=SAT(r\_keep) using rows in \(K_{<rank(e)}\). The rank order makes the certificate forest acyclic and independently checkable.
\item Run line induction. bEF lines use the existing \(\Gamma_{bEF}\) parser. FormulaEval lines add U(K\_Q)+X(K\_Q). CTSR lines add all counterfactual children, so actual-branch pruning cannot hide influence. Quotient lines replace a deleted row by a kept row only after the equality of SAT labels is certified from already charged dependencies.
\item Construct \(C_{pi}\): compute \(G_{N}\) in size \(N^{\rho+o(1)}\), evaluate the parsed transcript and same-theory certificates using \(\Gamma_\mathrm{QCTSR}\) gates, and patch \(E_{pi}\) exceptional source rows by equality tests of cost O(E\_pi N).
\item If the displayed size is below \(N^{1+\epsilon-o(1)}\), then the claimed source lower bound LB\_N(1+epsilon) is falsified in the same nonuniform model. This is a valid route-level implication and not a lower bound for \(H_{N}\).
\item For necessity, choose the first quotient edge without a ranked certificate. Because all earlier dependencies are fixed, flip the private deleted row label while preserving the visible prefix. If its residual parity incidence is odd, the proof transcript is unchanged but \(H_{N}\) changes, contradicting extensional safety of the compressed ledger.
\end{enumerate}

\paragraph{Equation 1: carrier.}
\[
H_N^{SAT}(x)=\bigoplus_{j=1}^{M_N} SAT(G_N(x,j)),\quad M_N=N^{q+o(1)}
\]
\emph{Valid when:} Valid when the source carrier is the beta-zero parity construction with uniformly generated polylog-length SAT rows and no target query length exponent hidden in \(G_{N}\).

\paragraph{Equation 2: ledger.}
\[
\Gamma_{QCTSR}(\pi)=\Gamma_{bEF}(\pi)+U(K_Q)+X(K_Q)+B_{CTSR}(K_Q)+B_Q(K_Q)
\]
\emph{Valid when:} Applies after \(K_{Q}\) is computed as the least dependency-closed support cone containing all proof-reachable FormulaEval, counterfactual TSR, and ranked quotient-certificate rows.

\paragraph{Equation 3: ctsr\_charge.}
\[
B_{CTSR}(K_Q)\le N^{1+\kappa+\theta+o(1)}
\]
\emph{Valid when:} Valid when the reachable TSR frontier has exponent \(\kappa\), every frontier node exposes \(N^{\theta+o(1)}\) counterfactual children, and child encodings are charged before semantic pruning.

\paragraph{Equation 4: quotient\_charge.}
\[
B_Q(K_Q)=\sum_{e\in F_Q}\left(|cert(e)|+N\cdot |Rows(cert(e))\setminus K_{< rank(e)}|\right)
\]
\emph{Valid when:} Requires each quotient edge e to carry a ranked same-theory certificate whose previously unseen dependency rows are added and charged before the deleted row is merged.

\paragraph{Equation 5: parser\_falsifier.}
\[
N^{\rho+o(1)}+\Gamma_{QCTSR}(\pi)+O(E_\pi N)< N^{1+\epsilon-o(1)}\Rightarrow \neg LB_N(1+\epsilon)
\]
\emph{Valid when:} Applies in the same nonuniform SIZE model after the generator, parsed transcript, quotient certificates, TSR children, and sparse residual corrections are all counted in the circuit size.

\paragraph{Next analysis.} Write the referee-checkable \(K_{Q}\) closure pseudocode with rank ordering, then test it on a two-child Tseitin-equivalent TSR merge and on the same example with the certificate omitted to exhibit the first-unranked-merge falsifier.

\paragraph{Falsifiers.}
\begin{itemize}
\item First unranked quotient merge: a deleted row is identified with a kept row without a surfaced rank-decreasing same-theory certificate.
\item Cyclic quotient dependency: certificates for two row deletions depend on each other and therefore cannot be audited by a finite lower-rank induction.
\item Advice-chosen quotient forest: the merge list depends on SAT labels or hidden semantic advice rather than on N and pi.
\item Actual-branch-only TSR declaration: an unlisted counterfactual child can be made the private row, so occurrence completion is invalid.
\item \(q<\epsilon\) with full row scan: FormulaEval expansion costs \(N^{1+q+o(1)}\) below the claimed \(N^{1+\epsilon}\) lower bound, giving a direct same-model falsifier.
\item Positive \(\beta\) black-box substitution: \(\mathrm{SAT}\in \mathrm{SIZE}(n^{c})\) would give source circuits of size N\textasciicircum{}\{q+beta c+o(1)\}, so all-exponent target exclusion is invalid unless \(\beta=0\) or genuinely subpolynomial.
\end{itemize}

\subsection*{cycle-0145: Hardness-Magnification Bridge Realization}

\paragraph{Execution.} research; azure\_responses\_api; attempt primary; 121.597 seconds; 3 lemmas, 6 equations, 10 derivation steps, 6 falsifiers.

\paragraph{Learning.} Cycle 145 upgrades ranked QCTSR from a line-induction slogan to a finite verifier/minimality packet. Quotient compression is valid only when proof-carrying same-theory certificates are ranked, dependency-closed, and charged before use. The exact obstruction is the first rejected quotient merge, which yields a private-row falsifier. This strengthens the paper parser/countermodel boundary without proving \(H_{N}\) lower bounds or \(\mathrm{P} \ne \mathrm{NP}\).

\paragraph{Pressure.} The derivation is complete for surface QCTSR syntax with a specified EqSAT/CTSR same-theory checker. It does not handle arbitrary semantic FormulaEval or hidden adaptive pruning unless \(RealSATInfluenceRealizer_{N}^{arb}\) uniformly lists and charges the same influenced rows before use.

\paragraph{Derivation steps.}
\begin{enumerate}
\item Parse pi line by line and create bit fields: line\_id, rule\_tag, row\_id, formula\_hash, restriction\_string, quotient\_class\_id, certificate\_id, parent\_pointer\_list, rank, and local charge.
\item Initialize K with every row mentioned by the final \(\mathrm{LB}_{N}\) line; then iterate the closure operators Dep, \(\mathrm{CTSR}_{cf}\), QuotRep, and \(CertRows_{Q}\) until a fixed point is reached.
\item For a transparent FormulaEval line, either enumerate the surfaced row set and charge U(K) plus occurrence expansion X(K), or full-scan all \(M_{N}\) rows and charge the corresponding \(N^{1+q+o(1)}\) term.
\item For a counterfactual TSR line, add every syntactically reachable child row, not just the actual branch; charge \(B_\mathrm{CTSR}=N^{1+\kappa+\theta+o(1)}\) in the exponent regime where the child list is complete.
\item For a quotient merge, accept row deletion only if a same-theory certificate forest is present, all certificate dependencies are already in K or are added by closure, and ranks strictly decrease along every dependency edge.
\item The local verifier checks certificate acyclicity by rank comparison, checks formula hashes and parent pointers, and refuses any edge whose target row was not already surfaced or newly charged.
\item Line induction: assuming all parent lines have circuit simulations using only \(K_{Q}\) rows, each allowed rule produces the next line with additional size bounded by its ledger charge; quotient representatives are interchangeable only through accepted same-theory certificates.
\item At the final line, compose the generator circuit of size \(N^{\rho+o(1)}\), the parsed transcript circuit of size \(\Gamma_\mathrm{QCTSR}\), and equality patches for \(E_{pi}\) residual rows, obtaining C(H\_N)<=N\textasciicircum{}\{rho+o(1)\}+Gamma\_QCTSR+O(E\_pi N).
\item If this size is below \(N^{1+\epsilon-o(1)}\), the claimed LB\_N(1+epsilon) is falsified in the same nonuniform SIZE model.
\item For necessity, choose the first unaccepted quotient merge e. Toggle a private row behind e while keeping all charged prefix rows fixed. The transcript remains accepted but the represented parity or branch value changes, proving that the unranked merge was not transparent.
\end{enumerate}

\paragraph{Equation 1: carrier.}
\[
H_N^{SAT}(x)=\bigoplus_{j=1}^{M_N} SAT(G_N(x,j)),\quad M_N=N^{q+o(1)},\quad L_N=(\log N)^d
\]
\emph{Valid when:} Valid when N is a valid carrier length and beta-zero polylogarithmic target-query length is part of the branch interface.

\paragraph{Equation 2: least\_closed\_support.}
\[
K_Q(\pi)=\mu K\,[Rows(final)\cup Dep(K)\cup CTSR_{cf}(K)\cup QuotRep(K)\cup CertRows_Q(K)]
\]
\emph{Valid when:} Valid when every dependency operator is syntactically computable from the proof text and certificate forests are ranked acyclic objects.

\paragraph{Equation 3: charged\_gamma.}
\[
\Gamma_{QCTSR}(\pi)=\Gamma_{bEF}(\pi)+U(K_Q)+X(K_Q)+B_{CTSR}(K_Q)+B_Q(K_Q)
\]
\emph{Valid when:} Applies after all transparent rows, counterfactual TSR children, and quotient certificates have been surfaced and charged in the same ledger.

\paragraph{Equation 4: simulator\_bound.}
\[
C(H_N^{SAT})\le N^{\rho+o(1)}+\Gamma_{QCTSR}(\pi)+O(E_\pi N)
\]
\emph{Valid when:} Requires the parser line induction, same-model circuit simulation, and residual error rows \(E_{pi}\) patchable by equality tests on N-bit inputs.

\paragraph{Equation 5: falsifier\_window.}
\[
N^{\rho+o(1)}+\Gamma_{QCTSR}(\pi)+O(E_\pi N)< N^{1+\epsilon-o(1)}\Rightarrow \neg LB_N(1+\epsilon)
\]
\emph{Valid when:} Valid when \(\mathrm{LB}_{N}\) is asserted against nonuniform SIZE circuits and the constructed simulator uses no stronger advice than the proof permits.

\paragraph{Equation 6: support\_density\_corollary.}
\[
|K_{idx}(\pi)|\ge N^{\epsilon-o(1)},\quad |K_{idx}(\pi)|/M_N\ge N^{\epsilon-q-o(1)}
\]
\emph{Valid when:} Applies when each surfaced row index costs an N-sized occurrence expansion and all other \(\Gamma\) terms remain subthreshold.

\paragraph{Next analysis.} Instantiate the finite verifier on one toy QCTSR proof with two EqSAT-equivalent counterfactual TSR children: compute \(K_{Q}\), ranks, \(B_{Q}\), and \(\Gamma_\mathrm{QCTSR}\); then delete one certificate edge and explicitly write the two-world private-row falsifier.

\paragraph{Falsifiers.}
\begin{itemize}
\item First unranked quotient merge falsifier: a row deletion not backed by a ranked same-theory certificate admits a private-row two-world collision.
\item Cyclic certificate falsifier: a quotient proof with no decreasing rank can hide semantic advice in a dependency loop and is rejected by the finite verifier.
\item Advice-chosen representative falsifier: selecting quotient representatives using SAT labels or hidden row sets violates the target-preserving condition and invalidates \(\Gamma\) shrinkage.
\item Actual-branch-only TSR falsifier: charging only the followed restriction child leaves a first-hidden-child private row for the counterfactual branch.
\item Threshold-collapse falsifier: if \(\kappa+\theta\ge \epsilon\), \(chi\ge \epsilon\), or \(E_{pi}\ge N^{\epsilon-o(1)}\), the constructed simulator is not subthreshold.
\item Positive-beta falsifier for broader magnification claims: \(\mathrm{SAT}\in \mathrm{SIZE}(n^{c})\) still gives carrier size N\textasciicircum{}\{q+beta c+o(1)\}, so \(\beta>0\) cannot exclude all c through this interface.
\end{itemize}

\subsection*{cycle-0146: Hardness-Magnification Bridge Realization}

\paragraph{Execution.} research; azure\_responses\_api; attempt primary; 85.08 seconds; 3 lemmas, 5 equations, 9 derivation steps, 6 falsifiers.

\paragraph{Learning.} Cycle 146 upgrades ranked QCTSR from a finite local verifier to a composable fixed-point line-induction theorem: quotient compression is valid only when proof-carrying same-theory certificates are ranked, acyclic, dependency-closed, and charged before use. The first rejected merge is a concrete private-row falsifier. This strengthens the transparent-support parser/countermodel paper module while preserving all nonclaims: no \(H_{N}\) lower bound, no \(\mathrm{P} \ne \mathrm{NP}\) result, and no free hidden semantic FormulaEval.

\paragraph{Pressure.} The derivation does not fold arbitrary hidden semantic FormulaEval or adaptive SAT branch pruning into \(K_{Q}\). That remains the named \(RealSATInfluenceRealizer_{N}^{arb}\) problem.

\paragraph{Derivation steps.}
\begin{enumerate}
\item Define K\_Q\textasciicircum{}0(pi) as the row, proof-line, and certificate seeds syntactically referenced by the final \(\mathrm{LB}_{N}\) assertion.
\item Iterate closure: add EF parent rows and extension-support rows; for FormulaEval lines add the enumerated RowSet or charge the full scan; for CTSR lines add every counterfactual child row in the bounded rule schema.
\item For each quotient deletion r \(\to \) r', require a displayed same-theory certificate forest \(F_{r,r'}\) with a rank label. Add every leaf and internal dependency of \(F_{r,r'}\) to \(K_{Q}\) before accepting the deletion.
\item Check acyclicity by requiring rank(child)<rank(parent) along certificate-dependency edges. This makes the verifier finite and prevents cyclic quotient justification.
\item Define \(\Gamma_\mathrm{QCTSR}\) as the sum of existing \(\Gamma_{bEF}\) charges, support-enumeration cost \(U_{Q}\), row-expansion cost \(X_{Q}=N|K_{Q}|\) up to encoding overhead, complete CTSR child cost, certificate-checking cost \(C_{Q}\), and residual patch cost O(E\_pi N).
\item Prove the line invariant: after line t, the partial circuit using only rows in \(K_{Q}^{\le t}\) agrees with the proof transcript on all surfaced influences; quotient representatives are safe because same-theory certificates are already available and ranked.
\item At the final line, compose the generator circuit for \(G_{N}\) with the parsed transcript circuit and add equality patches for \(E_{pi}\) residual source rows. The total size is N\textasciicircum{}\{rho+o(1)\}+Gamma\_QCTSR(pi)+O(E\_pi N).
\item If this total is below \(N^{1+\epsilon-o(1)}\), the asserted LB\_N(1+epsilon) is false in the same nonuniform SIZE model; this is the route-level falsifier conclusion.
\item For necessity, take the first quotient merge that fails the finite verifier. Create two private-row worlds agreeing on \(K_{Q}\) and all accepted certificates but differing on the uncharged deleted row. The proof sees the same transcript while carrier parity can differ, so unranked compression is not a valid parser rule.
\end{enumerate}

\paragraph{Equation 1: carrier.}
\[
H_N^{SAT}(x)=\bigoplus_{j< M_N} SAT(G_N(x,j)), \quad M_N=N^{q+o(1)}.
\]
\emph{Valid when:} Valid when the source carrier is the beta-zero generated parity construction and the generator outputs polylogarithmic SAT instances.

\paragraph{Equation 2: ledger.}
\[
\Gamma_{QCTSR}(\pi)=\Gamma_{bEF}(\pi)+U_Q(\pi)+X_Q(\pi)+B_{CTSR}(\pi)+C_Q(\pi)+O(E_\pi N).
\]
\emph{Valid when:} Valid when every proof rule is parsed through visible support, complete counterfactual children, and ranked same-theory quotient certificates.

\paragraph{Equation 3: ctsr\_charge.}
\[
B_{CTSR}(\pi)\le N^{1+\kappa+\theta+o(1)} \text{ for } N^{\kappa+o(1)} \text{ frontier rows and } N^{\theta+o(1)} \text{ children per row}.
\]
\emph{Valid when:} Applies when the bounded TSR schema declares all counterfactual children uniformly before adaptive pruning or quotient compression occurs.

\paragraph{Equation 4: falsifier\_window.}
\[
N^{\rho+o(1)}+\Gamma_{QCTSR}(\pi)< N^{1+\epsilon-o(1)} \Rightarrow \neg LB_N(1+\epsilon).
\]
\emph{Valid when:} Requires the parser-to-circuit construction, same nonuniform SIZE model, and residual row patching below the \(N^\epsilon\) threshold.

\paragraph{Equation 5: density\_corollary.}
\[
|K_Q(\pi)|\ge N^{\epsilon-o(1)} \quad \text{or} \quad |K_Q(\pi)|/M_N\ge N^{\epsilon-q-o(1)}.
\]
\emph{Valid when:} Valid when each surfaced row has N-cost expansion and the claimed lower-bound proof is not already falsified by \(\Gamma_\mathrm{QCTSR}\).

\paragraph{Next analysis.} Run the finite verifier on one projected-dDNNF/Tseitin-renaming example: list \(K_{Q}\) seeds, accepted quotient certificates, rank labels, computed \(\Gamma_\mathrm{QCTSR}\), and then delete one certificate to exhibit the first-rejected-merge private-row falsifier.

\paragraph{Falsifiers.}
\begin{itemize}
\item first unranked quotient merge: a private row can be flipped while all charged rows and verifier bits stay fixed
\item cyclic certificate forest: rank check fails, allowing circular deletion with no same-theory support
\item advice-chosen quotient representative: nonuniform leakage selects the surviving row without paying \(K_{Q}\) or \(C_{Q}\)
\item actual-branch-only TSR declaration: a hidden counterfactual child yields the first-hidden-child falsifier
\item semantic FormulaEval without \(RealSATInfluenceRealizer_{N}^{arb}\): unread-enable-bit row affects the proof but is absent from \(\Gamma_\mathrm{QCTSR}\)
\item positive \(\beta\) black-box interface: q+beta*c+o(1) caps all-exponent transfer and invalidates P/poly conclusions
\end{itemize}

\subsection*{cycle-0147: Hardness-Magnification Bridge Realization}

\paragraph{Execution.} research; azure\_responses\_api; attempt primary; 81.535 seconds; 3 lemmas, 5 equations, 8 derivation steps, 7 falsifiers.

\paragraph{Learning.} Cycle 147 closes a real proof-audit upgrade: ranked QCTSR quotient compression has a least, order-independent, proof-carrying support closure, so \(\Gamma\) reduction is valid only after lower-rank same-theory dependencies are charged. The sharp falsifier is the first nondecreasing, cyclic, advice-chosen, or hidden semantic merge. This strengthens the paper parser/countermodel module without proving any \(H_{N}\) lower bound or \(\mathrm{P} \ne \mathrm{NP}\).

\paragraph{Pressure.} No blocking gap remains for restricted ranked QCTSR syntax. The broader full-semantic FormulaEval system remains blocked exactly at \(RealSATInfluenceRealizer_{N}^{arb}\), which would have to surface all real SAT-semantic influences with charges before use.

\paragraph{Derivation steps.}
\begin{enumerate}
\item Define the finite universe U(pi) of row nodes, FormulaEval support nodes, counterfactual TSR child nodes, quotient representatives, and quotient-certificate nodes visible in pi.
\item Define the monotone operator F(K) by adding proof parents, reachable transparent support, all counterfactual children of active TSR nodes, deterministic quotient representatives, and all lower-rank certificate dependencies for quotient merges touching K.
\item Because F is monotone on a finite powerset lattice, iterating from the final \(\mathrm{LB}_{N}\) line reaches a least fixed point \(K_{Q}\) independent of enumeration order.
\item Strict rank decrease prevents a quotient certificate from validating deletion of a row through itself; every certificate node used in a merge has its dependency cone inserted before the merge can reduce the ledger.
\item Minimality follows by induction on closure time: any support set sufficient for the line induction must contain the final line, then must contain each parent, child, same-theory dependency, and quotient representative demanded by the verifier, hence contains \(K_{Q}\).
\item Line induction extends the prior \(\Gamma_{bEF}\) proof: bEF steps charge \(\Gamma_{bEF}\), transparent FormulaEval charges U+X, declared CTSR charges \(B_\mathrm{CTSR}\), and ranked quotient merges charge \(B_{Q}\) instead of deleting rows for free.
\item The circuit simulator hardwires \(G_{N}\), evaluates only rows in \(K_{Q}\) through the parsed support circuit, applies quotient representatives with their charged certificates, and patches \(E_{pi}\) residual rows by equality tests of cost O(E\_pi N).
\item If the resulting upper bound is below \(N^{1+\epsilon-o(1)}\), it is a valid same-model falsifier of the asserted LB\_N(1+epsilon).
\end{enumerate}

\paragraph{Equation 1: closure\_operator.}
\[
K_{t+1}=K_t\cup Parents(K_t)\cup FE(K_t)\cup CTSR(K_t)\cup QRep(K_t)\cup QCertDeps_{< rank}(K_t)
\]
\emph{Valid when:} Valid when pi is a finite ranked QCTSR transcript whose dependency predicates are syntactically decidable from pi and N.

\paragraph{Equation 2: termination\_bound.}
\[
K_Q(\pi)=K_T \text{ for some } T\le |Rows(\pi)|+|CertNodes(\pi)|
\]
\emph{Valid when:} Valid when every closure insertion is drawn from the explicit finite row and certificate-node universe of the transcript.

\paragraph{Equation 3: ledger\_definition.}
\[
\Gamma_{QCTSR}(\pi)=\Gamma_{bEF}(\pi)+U(K_Q)+X(K_Q)+B_{SR}(\pi)+B_{CTSR}(\pi)+B_Q(\pi)
\]
\emph{Valid when:} Valid when all transparent FormulaEval, bounded self-reduction, counterfactual child, and quotient-certificate charges are surfaced before line induction uses them.

\paragraph{Equation 4: parser\_upper\_bound.}
\[
C(H_N^{SAT})\le N^{\rho+o(1)}+\Gamma_{QCTSR}(\pi)+O(E_\pi N)
\]
\emph{Valid when:} Applies after the generator, parsed transcript ledger, and sparse residual-row patching are all circuits in the same nonuniform SIZE model.

\paragraph{Equation 5: lb\_falsifier\_window.}
\[
N^{\rho+o(1)}+\Gamma_{QCTSR}(\pi)+O(E_\pi N)< N^{1+\epsilon-o(1)}\Rightarrow \neg LB_N(1+\epsilon)
\]
\emph{Valid when:} Requires that LB\_N(1+epsilon) asserts a same-model lower bound for the exact carrier \(H_{N}^{\mathrm{SAT}}\) on the same valid lengths.

\paragraph{Next analysis.} Instantiate the closure verifier on a three-row QCTSR toy transcript: two counterfactual child rows, one EqSAT certificate forest, and one attempted cyclic merge; compute \(K_{Q}\), \(\Gamma_\mathrm{QCTSR}\), and the first-rejected-merge private-row falsifier explicitly.

\paragraph{Falsifiers.}
\begin{itemize}
\item first unranked quotient merge with no lower-rank same-theory certificate
\item cyclic quotient-certificate dependency that deletes the row needed to certify itself
\item advice-chosen quotient representative not computable from pi and N
\item hidden semantic FormulaEval row or unread enable bit outside K\_Q(pi)
\item threshold-collapse falsifier when \(N^{\rho}+\Gamma_\mathrm{QCTSR}+E_{pi}\) N reaches \(N^{1+\epsilon-o(1)}\)
\item model-mismatch falsifier when \(\mathrm{LB}_{N}\) and the constructed circuit are not in the same nonuniform SIZE model
\item positive-beta black-box substitution falsifier q+beta c bounding the target exponent window
\end{itemize}

\subsection*{cycle-0148: Hardness-Magnification Bridge Realization}

\paragraph{Execution.} research; azure\_responses\_api; attempt primary; 93.761 seconds; 3 lemmas, 6 equations, 9 derivation steps, 5 falsifiers.

\paragraph{Learning.} Cycle 148 turns ranked QCTSR from an order-independent closure statement into a representative-stable \(TransparentRowSupport_{N}\) rule case. The valid contribution is a restricted theorem/countermodel pair: proof-carrying quotient compression preserves \(\Gamma\) line induction, while the first unranked, cyclic, advice-chosen, or hidden semantic merge yields a private-row falsifier. This closes a proof-audit loophole without proving any \(H_{N}\) lower bound or \(\mathrm{P} \ne \mathrm{NP}\); full semantic pruning remains exactly at \(RealSATInfluenceRealizer_{N}^{arb}\).

\paragraph{Pressure.} Full semantic FormulaEval with hidden adaptive pruning remains blocked at \(RealSATInfluenceRealizer_{N}^{arb}\): no advice-free nonrelativizing procedure is proved that lists all real SAT influences before the semantic rule is used.

\paragraph{Derivation steps.}
\begin{enumerate}
\item Source statement: a \(\mathrm{Psys}_{sym}^\mathrm{QCTSR}\) proof pi certifies LB\_N(1+epsilon) for the beta-zero parity carrier \(H_{N}^\mathrm{SAT}\) in the nonuniform SIZE model.
\item Target statement: if Gamma\_QCTSR(pi) is subthreshold, then construct a same-model SIZE circuit for \(H_{N}^\mathrm{SAT}\) of size below \(N^{1+\epsilon-o(1)}\), contradicting the certified lower-bound claim but not proving \(\mathrm{P} \ne \mathrm{NP}\).
\item Model context: source carrier is \(H_{N}^\mathrm{SAT}\), target class is nonuniform \(\mathrm{SIZE}_{N}\), proof model is bit-explicit bEF plus transparent FormulaEval plus counterfactual declared TSR plus ranked same-theory quotient certificates, and the simulation is a parser-to-circuit upper-bound reduction.
\item Quantifier context: N ranges over valid source lengths; q,rho,kappa,theta,epsilon,d are fixed; L\_N=(log N)\textasciicircum{}d; \(\beta=0\); all row enumerators and certificates are advice-free once pi and N are fixed.
\item Closure construction: compute K\_Q(pi) by iterating Seed, parent dependencies, visible FormulaEval influence, all counterfactual TSR children, and certificate dependencies of smaller rank until a fixed point is reached.
\item Line induction: for bEF and transparent FormulaEval lines use the prior \(\Gamma_{bEF}\) and \(U_{Q}+N|K_{Q}|\) charges; for CTSR lines use complete child declarations; for quotient merges replace rows only after their ranked same-theory certificate forest is in \(K_{Q}\).
\item Upper-bound assembly: combine generator overhead \(N^{\rho+o(1)}\), parsed row circuits, certificate verifier circuits, and sparse error correction O(E\_pi N) to obtain C(H\_N\textasciicircum{}SAT)<=N\textasciicircum{}\{rho+o(1)\}+Gamma\_QCTSR(pi)+O(E\_pi N).
\item Necessity step: if a merge is first unranked, cyclic, advice-chosen, or hidden, keep all visible proof bits fixed and flip the private hidden row label; the transcript remains accepted while the parity carrier changes, so the parser is not extensionally valid.
\item Reviewer findings: reduction consistency is valid for \(\beta=0\) only; quantifier uniformity is valid after certificate rank and row enumeration are syntax-visible; barrier specificity says the obstruction is a private-row/nonuniform-leakage countermodel, not a generic natural-proofs or relativization citation.
\end{enumerate}

\paragraph{Equation 1: carrier.}
\[
H_N^{SAT}(x)=\bigoplus_{j< M_N} SAT(F_{N,x,j}),\quad M_N=N^{q+o(1)},\quad |F_{N,x,j}|=(\log N)^d
\]
\emph{Valid when:} Valid when N is a valid source length and generator rows are syntactically enumerated.

\paragraph{Equation 2: least\_support\_fixed\_point.}
\[
K_Q(\pi)=\mu K\,[Seed(\pi)\cup Par(K)\cup FE(K)\cup CTSR(K)\cup Cert_Q^{< rank}(K)]
\]
\emph{Valid when:} Applies after all proof dependencies, TSR children, and quotient certificates are exposed before use.

\paragraph{Equation 3: ledger.}
\[
\Gamma_{QCTSR}(\pi)=\Gamma_{bEF}^{-}(\pi)+U_Q(\pi)+N|K_Q(\pi)|+B_{CTSR}(\pi)+C_Q(\pi)
\]
\emph{Valid when:} Requires visible row enumeration, complete counterfactual child declarations, and charged same-theory certificate forests.

\paragraph{Equation 4: parser\_upper\_bound.}
\[
C(H_N^{SAT})\le N^{\rho+o(1)}+\Gamma_{QCTSR}(\pi)+O(E_\pi N)
\]
\emph{Valid when:} Requires the target-preserving beta-zero substitution with advice-free row enumeration and same SIZE model.

\paragraph{Equation 5: subthreshold\_falsifier.}
\[
N^{\rho+o(1)}+\Gamma_{QCTSR}(\pi)+O(E_\pi N)< N^{1+\epsilon-o(1)}\Rightarrow \neg LB_N(1+\epsilon)
\]
\emph{Valid when:} Valid when the assumed \(\mathrm{LB}_{N}\) threshold is in the identical nonuniform circuit model.

\paragraph{Equation 6: support\_density\_pressure.}
\[
|K_Q(\pi)|\ge N^{\epsilon-o(1)},\quad |K_Q(\pi)|/M_N\ge N^{\epsilon-q-o(1)}
\]
\emph{Valid when:} Applies when all other ledger terms are subthreshold and every omitted row costs N to patch.

\paragraph{Next analysis.} Write the finite \(K_\mathrm{QCTSR}\) verifier on a projected-dDNNF example with one certified duplicate row and one unread-enable hidden merge; compute \(U_{Q},N|K_{Q}|,B_\mathrm{CTSR},C_{Q},E_{pi}\) and show exactly where the hidden merge fails the valid support ledger.

\paragraph{Falsifiers.}
\begin{itemize}
\item first unranked quotient merge falsifier: two completions agree on charged rows but differ on the hidden representative row
\item cyclic or nondecreasing certificate-rank falsifier: the alleged compression is self-justifying and cannot be parsed before use
\item advice-chosen representative falsifier: the quotient map carries nonuniform row-selection leakage not charged in \(\Gamma_\mathrm{QCTSR}\)
\item hidden adaptive TSR pruning falsifier: an unlisted counterfactual child can be made decisive by a private-row assignment
\item positive-beta black-box substitution falsifier remains active outside this beta-zero parser route: q+beta*c caps all-exponent transfer
\end{itemize}

\subsection*{cycle-0149: Hardness-Magnification Bridge Realization}

\paragraph{Execution.} research; azure\_responses\_api; attempt primary; 100.413 seconds; 3 lemmas, 6 equations, 8 derivation steps, 6 falsifiers.

\paragraph{Learning.} Cycle 149 folds representative-stable ranked QCTSR into \(TransparentRowSupport_{N}\) as a valid rule case. The new contribution is not another magnification slogan: quotient compression is proof-carrying, rank-stratified, dependency-closed, and charged, while the first rejected merge gives a concrete private-row falsifier. Full semantic SAT pruning remains exactly at \(RealSATInfluenceRealizer_{N}^{arb}\).

\paragraph{Pressure.} The derivation does not cover arbitrary real-SAT semantic pruning or FormulaEval comprehension unless \(RealSATInfluenceRealizer_{N}^{arb}\) uniformly lists the influencing rows and charges them before use.

\paragraph{Derivation steps.}
\begin{enumerate}
\item Encode each QCTSR proof line with fields (rule tag, parents, row id, child restriction, representative id, quotient-certificate id, rank, dependency list). These fields are verifier-visible and do not query SAT truth advice.
\item Define \(F_{Qrep}\) on sets of support objects: add bEF dependencies, FormulaEval support rows, every counterfactual TSR child row, all nodes of each same-theory quotient certificate forest, and the canonical representative chosen by the charged certificate. Iterate from the final \(\mathrm{LB}_{N}\) line to obtain K\_Qrep(pi).
\item Monotonicity and finiteness give a least fixed point. Rank decrease in quotient certificates gives acyclicity, so the same K\_Qrep(pi) is reached independently of line order or certificate-processing order.
\item Prove representative stability by induction on rank: if row a is replaced by representative b, then all lower-rank certificate dependencies proving EqSAT(a,b) are already in K\_Qrep(pi), so the parser may identify a and b without reading an uncharged label.
\item Extend the existing \(\Gamma_{bEF}\) line induction. bEF lines use \(\Gamma_{bEF}\); transparent FormulaEval uses \(U_{N}+X_{N}\); complete counterfactual TSR child expansion contributes \(C_\mathrm{CTSR}=N^{1+\kappa+\theta+o(1)}\) in the bounded schema; ranked quotient compression contributes \(C_{Qcert}+C_{rep}\) and may only reduce row multiplicity after the certificate is charged.
\item Build the source circuit for \(H_{N}\) from the generator circuit, the parsed transcript evaluator, parity aggregation over charged representatives, and equality patches for \(E_{pi}\) residual rows. This yields C\_SIZE(H\_N)<=N\textasciicircum{}\{rho+o(1)\}+Gamma\_Qrep(pi)+O(E\_pi N).
\item Contrapose: if LB\_N(1+epsilon) is asserted in the same SIZE model, then any valid \(\mathrm{Psys}_{sym}^\mathrm{QCTSR}\) certificate must have Gamma\_Qrep(pi)+O(E\_pi N)>=N\textasciicircum{}\{1+epsilon-o(1)\}.
\item For necessity, take the first quotient merge or pruning child not accepted by the ranked closure verifier. Toggle a private row label r* outside K\_Qrep(pi), keep the visible transcript fixed, and force the affected line or carrier parity to differ. Thus hidden quotienting is a falsifier for transparent parsing.
\end{enumerate}

\paragraph{Equation 1: carrier.}
\[
H_N^SAT(x)=xor_{j< M_N} SAT(G_N(x,j)), M_N=N^{q+o(1)}, L_N=(log N)^d
\]
\emph{Valid when:} Valid when the carrier is the beta-zero parity construction with polylogarithmic SAT queries and a uniform generator.

\paragraph{Equation 2: least\_closure.}
\[
K_Qrep(pi)=lfp(F_parent union F_FE union F_CTSR union F_Qcert union F_rep)
\]
\emph{Valid when:} Valid when QCTSR syntax exposes all parent pointers, child rows, quotient certificates, and FormulaEval influences.

\paragraph{Equation 3: ledger.}
\[
Gamma_Qrep(pi)=Gamma_bEF(pi)+U_N+X_N+C_CTSR+C_Qcert+C_rep
\]
\emph{Valid when:} Valid when each support row expansion costs Theta(N) in the same nonuniform SIZE model and all certificate bits are charged.

\paragraph{Equation 4: parser\_upper\_bound.}
\[
C_SIZE(H_N^SAT) \le N^{rho+o(1)} + Gamma_Qrep(pi) + O(E_pi N)
\]
\emph{Valid when:} Valid when the parser line induction has no hidden semantic use and patchable row errors are explicitly bounded.

\paragraph{Equation 5: contrapositive\_window.}
\[
LB_N(1+epsilon) implies Gamma_Qrep(pi)+O(E_pi N) \ge N^{1+epsilon-o(1)}
\]
\emph{Valid when:} Valid when \(\mathrm{LB}_{N}\) means a same-model lower-bound assertion for \(H_{N}\) rather than any NP lower bound.

\paragraph{Equation 6: support\_density\_corollary.}
\[
|K_idx(pi)|/M_N \ge N^{epsilon-q-o(1)}
\]
\emph{Valid when:} Valid when row expansion dominates the ledger and each visible source row contributes Theta(N) charged gates.

\paragraph{Next analysis.} Write the standalone paper theorem for \(TransparentRowSupport_{N}^{Qrep}\): include the \(K_{Qrep}\) fixed-point verifier pseudocode, the \(\Gamma_{Qrep}\) ledger proof, the first-rejected-merge countermodel, and a one-page nonclaim box separating this parser boundary from \(H_{N}\) lower bounds and \(\mathrm{P} \ne \mathrm{NP}\).

\paragraph{Falsifiers.}
\begin{itemize}
\item first unranked quotient merge falsifier: a row is deleted without a charged same-theory certificate
\item cyclic or nondecreasing-rank certificate falsifier: the representative map is not well-founded
\item advice-chosen representative falsifier: nonuniform leakage selects a merge invisible to K\_Qrep(pi)
\item hidden semantic pruning falsifier: an unlisted counterfactual child changes the proof while the transcript stays fixed
\item positive \(\beta\) falsifier for broader magnification claims: \(\mathrm{SAT}\in \mathrm{SIZE}(n^{c})\) transfers only to N\textasciicircum{}\{q+beta c+o(1)\} source size
\item \(q<\epsilon\) enumeration falsifier: too little absolute support density makes LB\_N(1+epsilon) incompatible with the parser upper bound
\end{itemize}

\subsection*{cycle-0150: Hardness-Magnification Bridge Realization}

\paragraph{Execution.} research; azure\_responses\_api; attempt primary; 107.027 seconds; 3 lemmas, 4 equations, 8 derivation steps, 6 falsifiers.

\paragraph{Learning.} Cycle 150 closes a proof-delivery upgrade rather than rediscovering magnification: Qrep quotient compression is valid exactly as a batched, ranked, proof-carrying same-theory dictionary charged in \(\Gamma_{Qrep}\). The sharp falsifier is the first unranked, cyclic, advice-chosen, or hidden semantic merge. The result is a scoped parser/falsifier theorem and Formula Footprints supporting lemma; it does not prove \(H_{N}\) lower bounds or \(\mathrm{P} \ne \mathrm{NP}\).

\paragraph{Pressure.} The derivation does not cover a semantic pruning step that obtains a quotient or child deletion from real SAT truth without a surfaced same-theory certificate; that remains exactly \(RealSATInfluenceRealizer_{N}^{arb}\).

\paragraph{Derivation steps.}
\begin{enumerate}
\item Initialize \(K_{0}\) with the final \(\mathrm{LB}_{N}\) line and the carrier definition \(H_{N}^\mathrm{SAT}\); add bit-explicit EF dependencies exactly as in \(\Gamma_{bEF}\).
\item For each reachable FormulaEval line, insert the visible row selector support and charge either U(K) for uniform enumeration or the full scan charge; add occurrence expansion to X(K).
\item For each counterfactual TSR line, insert every reachable child row, not merely the actual branch; this preserves extensional safety under private-row variations.
\item For each quotient merge, require a lower-rank same-theory certificate forest proving equality of SAT truth for the merged children; add all certificate dependencies before replacing children by the representative.
\item Because ranks strictly decrease along certificate dependencies, the fixed point is finite and order-independent; representative stability follows by choosing the least certified representative in each connected component.
\item Line induction: after each proof line, every formula value used by the parsed circuit is either bit-explicit, a charged row in \(K_{Qrep}\), or a value transported along a charged same-theory certificate, so the line has circuit cost bounded by the accumulated \(\Gamma_{Qrep}\) ledger.
\item Hardwire pi, the generator \(G_{N}\), \(K_{Qrep}\), and \(D_{Q}\) nonuniformly; compute \(H_{N}\) by the generator plus the parsed transcript and sparse residual patches, giving C\_SIZE(H\_N)<=N\textasciicircum{}\{rho+o(1)\}+Gamma\_Qrep(pi)+O(E\_pi N).
\item Contrapose: if LB\_N(1+epsilon) asserts \(C_\mathrm{SIZE}(H_{N})>N^{1+\epsilon-o(1)}\), any transcript with N\textasciicircum{}\{rho+o(1)\}+Gamma\_Qrep(pi)+O(E\_pi N)<N\textasciicircum{}\{1+epsilon-o(1)\} is a valid same-model falsifier for that claimed source lower-bound certificate.
\end{enumerate}

\paragraph{Equation 1: parity\_carrier.}
\[
H_N^{SAT}(x)=\bigoplus_{j< M_N} SAT(G_N(x,j)),\quad M_N=N^{q+o(1)},\quad |G_N(x,j)|=(\log N)^d.
\]
\emph{Valid when:} Valid when N is a carrier length and \(G_{N}\) is the fixed beta-zero polylog-query generator used by the proof transcript.

\paragraph{Equation 2: least\_qrep\_closure.}
\[
K_{t+1}=K_t\cup Dep_{bEF}(K_t)\cup FE(K_t)\cup CTSR_{cf}(K_t)\cup CertDeps_Q(K_t),\quad K_{Qrep}=\bigcup_t K_t.
\]
\emph{Valid when:} Valid when the proof syntax exposes parent pointers, counterfactual child lists, and ranked quotient-certificate dependencies before each compressed use.

\paragraph{Equation 3: qrep\_ledger.}
\[
\Gamma_{Qrep}(\pi)=\Gamma_{bEF}(\pi)+U(K_{Qrep})+X(K_{Qrep})+C_Q(D_Q)+O(E_\pi N).
\]
\emph{Valid when:} Valid when every row enumeration, occurrence expansion, quotient certificate, and sparse residual patch is charged in the same nonuniform SIZE ledger.

\paragraph{Equation 4: upper\_bound\_falsifier.}
\[
C_{SIZE}(H_N^{SAT})\le N^{\rho+o(1)}+\Gamma_{Qrep}(\pi)+O(E_\pi N).
\]
\emph{Valid when:} Valid when the parsed transcript is hardwired nonuniformly and each SAT influence row is represented by \(K_{Qrep}\) or by a charged same-theory certificate.

\paragraph{Next analysis.} Run the Tseitin-renaming Qrep thought experiment: instantiate two batches with one certified duplicate and one unlisted duplicate, compute \(K_{Qrep}\), C\_Q(D\_Q), and \(\Gamma_{Qrep}\) explicitly, then verify that the unlisted duplicate yields the first-hidden-row private-row falsifier while the certified duplicate preserves the line induction.

\paragraph{Falsifiers.}
\begin{itemize}
\item first unranked quotient merge falsifier: an omitted child row varies privately while all charged support stays fixed
\item cyclic or nondecreasing rank certificate falsifier: the alleged dictionary is not a finite dependency-closed proof object
\item advice-chosen representative falsifier: the representative map encodes hidden SAT-label information outside \(\Gamma_{Qrep}\)
\item hidden semantic pruning falsifier: actual-branch-only TSR support misses a counterfactual child reachable from the final line
\item positive \(\beta\) black-box substitution falsifier: q+beta c reaches 1+epsilon before all target exponents are excluded
\item \(q<\epsilon\) support-density falsifier: subthreshold reachable support can be equality-patched below the claimed \(\mathrm{LB}_{N}\) threshold
\end{itemize}

\subsection*{cycle-0151: Hardness-Magnification Bridge Realization}

\paragraph{Execution.} research; azure\_responses\_api; attempt compact\_retry; 226.702 seconds; 1 lemmas, 3 equations, 4 derivation steps, 4 falsifiers.

\paragraph{Learning.} Cycle 151 upgrades the batched Qrep patch with temporal prefix stability: quotient compression is valid only as append-only, ranked, proof-carrying dictionary use. The first pre-use, rebinding, cyclic, advice-chosen, or hidden semantic merge is an explicit falsifier, while full semantic pruning stays quarantined.

\paragraph{Pressure.} The derivation does not realize arbitrary semantic SAT pruning; \(RealSATInfluenceRealizer_{N}^{arb}\) remains the named nonrelativizing bridge for full \(\mathrm{Psys}_{sym}\).

\paragraph{Derivation steps.}
\begin{enumerate}
\item Initialize \(K_{pref}\) from the final \(\mathrm{LB}_{N}\) line; iteratively add ordinary parents, FormulaEval influences, and every earlier Qrep dictionary entry used by a reachable line.
\item For a quotient lookup r \(\to \) rep, add the lower-rank same-theory certificate forest and charge its row expansions in \(D_{Qrep}\); append-only order makes the representative stable for all later lines.
\item Run the existing \(\Gamma_{bEF}\) line induction with dictionary lookups replaced by charged representative circuits and certificate checks; no line reads an uncharged row.
\item If \(\Gamma_{pref}\) is subthreshold, combine the generator circuit with the transcript circuit and sparse-row patching to obtain a same-model circuit for \(H_{N}\), contradicting the asserted \(\mathrm{LB}_{N}\).
\end{enumerate}

\paragraph{Equation 1: carrier.}
\[
H_N^{SAT}(x)=\bigoplus_{j\le M_N} SAT(G_N(x,j)),\quad M_N=N^{q+o(1)}.
\]
\emph{Valid when:} Valid when \(G_{N}\) is uniform and bit-explicit, \(M_{N}=N^{q+o(1)}\), and queried SAT instances have polylog length so \(\beta\) equals zero.

\paragraph{Equation 2: prefix\_ledger.}
\[
\Gamma_{pref}(\pi)=\Gamma_{bEF}(\pi)+U(K_{pref})+X(K_{pref})+D_{Qrep}(\pi)+O(E_\pi N).
\]
\emph{Valid when:} Applies when every dictionary entry is appended before first use, certificate dependencies are lower rank, and row supports are influence closed.

\paragraph{Equation 3: subthreshold\_falsifier.}
\[
N^{\rho+o(1)}+\Gamma_{pref}(\pi)< N^{1+\epsilon-o(1)}\Rightarrow C_{SIZE}(H_N^{SAT})< N^{1+\epsilon-o(1)}.
\]
\emph{Valid when:} Requires the target-preserving same-model SIZE simulation, q and \(\rho\) below 1+epsilon, and no hidden semantic support or representative advice.

\paragraph{Next analysis.} Write the finite prefix-stable Qrep verifier: fields are batch\_id, rank, source\_row, representative\_row, certificate\_root, dependency list, first\_use\_line, and no\_rebind flag; prove the first rejected field yields the private-row falsifier.

\paragraph{Falsifiers.}
\begin{itemize}
\item Prefix-use falsifier: a quotient lookup before its batch certificate is introduced leaves a private row uncharged; flip it while preserving all charged transcript rows.
\item Rebinding falsifier: if a representative map changes after use, two worlds agree on the old charged cone but differ on the later representative parity.
\item Rank-cycle falsifier: a nondecreasing or cyclic same-theory certificate lets a merge justify itself, so \(\Gamma\) reduction is advice-like rather than valid proof-carrying compression.
\item Hidden-pruning falsifier: semantic branch selection that omits a counterfactual child is outside Qrep unless \(RealSATInfluenceRealizer_{N}^{arb}\) lists and charges the missing row.
\end{itemize}

\subsection*{cycle-0152: Hardness-Magnification Bridge Realization: prefix-stable Qrep support closure}

\paragraph{Execution.} research; azure\_responses\_api; attempt primary; 62.825 seconds; 2 lemmas, 5 equations, 7 derivation steps, 5 falsifiers.

\paragraph{Learning.} Cycle 152 upgrades Qrep from an append-only convention to a prefix-verifiable rule case: quotient compression preserves the transparent parser only when certificate charge strictly precedes representative use. The exact no-go is temporal, not just acyclicity; delayed same-theory evidence is already a private-row falsifier. This closes a proof-audit loophole while preserving the nonclaim boundary: no \(H_{N}\) lower bound and no \(\mathrm{P} \ne \mathrm{NP}\) result.

\paragraph{Pressure.} The restricted PrefixStableQrepTransparentSupport proof is theorem-shaped; the unresolved extension remains \(RealSATInfluenceRealizer_{N}^{arb}\) for hidden full semantic SAT pruning and branch selection.

\paragraph{Derivation steps.}
\begin{enumerate}
\item Define \(K_{PQ}^{0}\) as the final \(\mathrm{LB}_{N}\) line plus its parsed bit-explicit parents. Iterate: add EF dependencies, SATeq certificate rows, \(bounded-SR_{h}\) leaves, reachable FormulaEval support rows, all counterfactual TSR children, and for each Qrep representative use add the lower-rank same-theory certificate forest and its row dependencies.
\item Prove termination by the finite proof DAG size plus a strict lexicographic measure: proof-line index, dictionary append time, and rank. Prefix stability forbids cycles and rebinding, so the fixed point is order-independent.
\item For each proof line l, maintain invariant I(l): all rows semantically influencing l through transparent rules are in \(K_{PQ}\) restricted to prefix l, and the circuit transcript for l has size at most the accumulated \(\Gamma_{PQ}\) charge plus generator overhead already charged.
\item EF and SATeq cases reuse the prior \(\Gamma_{bEF}\) parser; \(bounded-SR_{h}\) contributes \(B_{SR}\le N^{q+\eta+\tau+o(1)}\); FormulaEval contributes U(K)+X(K); counterfactual TSR contributes \(N^{1+\kappa+\theta+o(1)}\) for complete child declarations.
\item Qrep case: replacing r\_old by representative r\_new is sound because the certificate \(C_{e}\) is same-theory, lower-rank, and already parsed into the prefix closure. The row deletion saves duplicate expansion only after D\_Q(pi) has charged the dictionary entry and certificate forest.
\item At the final \(\mathrm{LB}_{N}\) line, compose the generator circuit, parsed transcript circuit, and equality patches for \(E_{pi}\) exceptional rows to obtain C(H\_N)<=N\textasciicircum{}\{rho+o(1)\}+Gamma\_PQ(pi)+O(E\_pi N).
\item If this upper bound is below \(N^{1+\epsilon-o(1)}\), it is a same-model falsifier for LB\_N(1+epsilon). This proves only failure of that proof certificate or lower-bound assertion in the restricted parser model.
\end{enumerate}

\paragraph{Equation 1: carrier.}
\[
H_N^{SAT}(x)=xor_{j=1}^{M_N} SAT(G_N(x,j)), M_N=N^{q+o(1)}
\]
\emph{Valid when:} Valid when the source carrier is the beta-zero parity aggregator with a uniformly parsed generator and visible row indices.

\paragraph{Equation 2: ledger.}
\[
Gamma_PQ(pi)=Gamma_bEF(pi)+U(K_PQ)+X(K_PQ)+D_Q(pi)+B_CTSR(pi)
\]
\emph{Valid when:} Valid when \(K_{PQ}\) is the least prefix-stable closure and every dictionary, quotient, and counterfactual TSR dependency is charged before use.

\paragraph{Equation 3: upper\_bound.}
\[
C(H_N^{SAT}) \le N^{rho+o(1)}+Gamma_PQ(pi)+O(E_pi N)
\]
\emph{Valid when:} Applies after the parser line induction produces a circuit transcript in the same nonuniform SIZE model with all residual bad rows equality-patched.

\paragraph{Equation 4: falsifier\_threshold.}
\[
N^{rho+o(1)}+Gamma_PQ(pi)+O(E_pi N)< N^{1+epsilon-o(1)} => not LB_N(1+epsilon)
\]
\emph{Valid when:} Requires the asserted \(\mathrm{LB}_{N}\) source lower bound to be in the same circuit model and quantified over the same valid lengths.

\paragraph{Equation 5: q\_density\_corollary.}
\[
|K_idx(pi)|/M_N \ge N^{epsilon-q-o(1)} when per-row expansion costs N
\]
\emph{Valid when:} Valid when FormulaEval rows are expanded to occurrence rows with unit N charge and \(M_{N}=N^{q+o(1)}\) slots.

\paragraph{Next analysis.} Write the finite verifier table for PrefixStableQrep entries: fields entry\_id, append\_time, old\_row, rep\_row, rank, certificate\_root, dependency\_list, first\_use\_time, \(charge_{to,Gamma}\), and rejection reason; then run it on one EqSAT duplicate-child example and one delayed-certificate private-row counterexample.

\paragraph{Falsifiers.}
\begin{itemize}
\item falsifier: a Qrep representative used before its proof-carrying dictionary certificate is appended invalidates the charged parser and yields a private-row two-world countermodel.
\item falsifier: representative rebinding lets two prefixes share the same visible transcript while later quotient meaning changes, so \(\Gamma_{PQ}\) is not valid.
\item falsifier: a rank cycle in same-theory certificates destroys the least fixed point and blocks line induction.
\item falsifier: advice-chosen dictionary compression hides semantic row selection and is outside \(TransparentRowSupport_{N}\).
\item falsifier: full semantic branch pruning without \(RealSATInfluenceRealizer_{N}^{arb}\) remains quarantined, even if the restricted Qrep ledger is valid.
\end{itemize}

\subsection*{cycle-0153: Hardness-Magnification Bridge Realization: prefix-strict Qrep support audit}

\paragraph{Execution.} research; azure\_responses\_api; attempt primary; 79.095 seconds; 3 lemmas, 5 equations, 6 derivation steps, 7 falsifiers.

\paragraph{Learning.} Cycle 153 upgrades Qrep from prefix-verifiable syntax to a theorem-shaped \(TransparentRowSupport_{N}\) rule case: quotient compression is valid exactly as strict-prefix, ranked, proof-carrying, dependency-closed dictionary use. The new falsifier is temporal; delayed same-theory evidence is already hidden support. This is a scoped parser/falsifier contribution, not a lower bound for \(H_{N}\) and not \(\mathrm{P} \ne \mathrm{NP}\) progress.

\paragraph{Pressure.} The derivation does not cover arbitrary semantic FormulaEval or adaptive SAT-pruning rules that may realize influence through real SAT semantics without exposing row identifiers; that remains \(RealSATInfluenceRealizer_{N}^{arb}\).

\paragraph{Derivation steps.}
\begin{enumerate}
\item Define Dict\_Q(pi) as append-only entries e=(row a, representative rep(a), rank r, certificate forest \(F_{e}\), cert\_line, first\_use\_line). Reject entries with \(cert_{line}\ge first_{use,line}\), cyclic dependencies, nondecreasing dependency rank, or rebinding of a previously used representative.
\item Compute K\_Qrep(pi) by the least fixed point from the final \(\mathrm{LB}_{N}\) line: add ordinary bit-explicit dependencies, transparent FormulaEval row dependencies, CTSR counterfactual children, and for every Qrep use add the representative row plus all lower-rank certificate dependencies already present in the prefix.
\item Induct on proof lines. Ordinary \(\mathrm{Psys}_{bEF}^{-}\) lines use the existing \(\Gamma_{bEF}\) parser. Transparent FormulaEval lines contribute \(U_{N}+X_{N}\). Qrep lines are simulated by substituting the charged representative expansion and the prior same-theory certificate forest; no new uncharged occurrence neighborhood is introduced.
\item The circuit extracted from the parsed transcript evaluates \(H_{N}\) using the generator circuit, charged row support, SAT-substitution for exposed polylog queries, and equality patches for \(E_{pi}\) source-row failures, giving C\_N(H\_N)<=N\textasciicircum{}\{rho+o(1)\}+Gamma\_Qrep(pi)+O(E\_pi N).
\item Contrapose the same-model lower-bound claim LB\_N(1+epsilon). If the right side is subthreshold, the proof transcript itself is a valid upper-bound falsifier for the claimed source lower bound.
\item For necessity, take the least line u violating strict prefix visibility. Hold all visible prefix bits fixed, flip the hidden private row or unread enable bit, and observe that the transcript before u is unchanged while the representative use can change the selected support or parity. This refutes hidden Qrep as transparent parsing.
\end{enumerate}

\paragraph{Equation 1: Qrep ledger.}
\[
Gamma_Qrep(pi)=Gamma_bEF(pi)+U_N(pi)+X_N(pi)+D_Q(pi)
\]
\emph{Valid when:} Valid when \(\mathrm{Psys}_{sym}^{dag+Qrep}\) exposes row support, FormulaEval expansion, and prefix-strict Qrep certificates before every representative use.

\paragraph{Equation 2: Dictionary charge.}
\[
D_Q(pi)=sum_{e in Dict_Q(pi)} (|desc(e)|+|cert(e)|+sum_{a in dep(e)} row_cost(a))
\]
\emph{Valid when:} Valid when dictionary entries are ranked, acyclic, dependency-closed, and no dependency is introduced after the representative is used.

\paragraph{Equation 3: Parser-to-circuit upper bound.}
\[
C_N(H_N^SAT) \le N^{rho+o(1)}+Gamma_Qrep(pi)+O(E_pi N)
\]
\emph{Valid when:} Applies after the line induction parses every proof line into charged source-row support in the same nonuniform SIZE model.

\paragraph{Equation 4: Subthreshold falsifier inequality.}
\[
N^{rho+o(1)}+Gamma_Qrep(pi)+O(E_pi N)< N^{1+epsilon-o(1)} => not LB_N(1+epsilon)
\]
\emph{Valid when:} Requires LB\_N(1+epsilon) to mean a same-model \(N^{1+\epsilon-o(1)}\) lower bound for the explicit source carrier \(H_{N}^\mathrm{SAT}\).

\paragraph{Equation 5: Strict-prefix condition.}
\[
forall e, cert_line(e)< first_use_line(e) and rank(dep(e))< rank(e)
\]
\emph{Valid when:} Valid when Qrep compression is proof-carrying rather than a semantic row-deletion convention or advice-chosen quotient map.

\paragraph{Next analysis.} Instantiate the strict-prefix Qrep verifier on a two-duplicate projected-dDNNF FormulaEval example and a delayed-certificate counterexample; compute \(\Gamma_{bEF}\), \(U_{N}\), \(X_{N}\), \(D_{Q}\), \(E_{pi}\), and the q/epsilon support-density consequence line by line.

\paragraph{Falsifiers.}
\begin{itemize}
\item first representative use before certificate charge is a temporal private-row falsifier
\item representative rebinding after prior use destroys prefix stability and invalidates \(\Gamma_{Qrep}\)
\item cyclic or nondecreasing-rank dictionary dependencies prevent a finite valid verifier
\item advice-chosen quotient representatives create nonuniform leakage outside \(TransparentRowSupport_{N}\)
\item hidden semantic SAT-pruning remains quarantined unless \(RealSATInfluenceRealizer_{N}^{arb}\) lists and charges the influence rows
\item positive target-query exponent \(\beta\) would reintroduce the black-box transfer cap q+beta c+o(1) and block all-exponent magnification
\item subthreshold Gamma\_Qrep(pi)+O(E\_pi N) is an upper-bound falsifier for any claimed LB\_N(1+epsilon) proof transcript
\end{itemize}

\subsection*{cycle-0154: Hardness-Magnification Bridge Realization}

\paragraph{Execution.} research; azure\_responses\_api; attempt primary; 75.659 seconds; 3 lemmas, 5 equations, 7 derivation steps, 5 falsifiers.

\paragraph{Learning.} Cycle 154 converts strict-prefix Qrep from a syntax convention into a replay theorem with a matching temporal falsifier. The valid contribution is a restricted \(TransparentRowSupport_{N}\) rule case and Formula Footprints supporting lemma: charged, ranked, dependency-closed dictionary compression preserves \(\Gamma\) line induction; delayed or hidden compression is exactly hidden support. This is route-level certification only, not a lower bound for \(H_{N}\) and not P vs NP progress.

\paragraph{Pressure.} The derivation does not cover arbitrary semantic SAT branch pruning or FormulaEval rules whose influencing rows are not determined by proof syntax before use; this is exactly \(RealSATInfluenceRealizer_{N}^{arb}\).

\paragraph{Derivation steps.}
\begin{enumerate}
\item Define \(K_{Qrep}^{t}\) by prefix recursion: start with bEF, \(SR_{h}\), FormulaEval, and CTSR visible supports; when a Qrep dictionary entry e appears, add its row endpoints and certificate footprint only if every dependency is already in \(K_{Qrep}^{<t}\) and every rank decreases.
\item For non-Qrep lines, reuse the existing \(\Gamma_{bEF}/\Gamma_{dag}/\mathrm{CTSR}\) line-induction invariant: each parsed proof dependency either contributes a charged row or is already in the support cone.
\item For a Qrep use at line t, replace the quotient representative by the row rep\_e plus the strict-prefix same-theory certificate forest sigma\_e. Since sigma\_e is already charged, the replay adds no post-hoc support and preserves the carrier row semantics used by the parser.
\item Compose all replays in line order. Rank decrease and prefix ordering prevent cycles and representative rebinding, so the expanded proof has the same final \(\mathrm{LB}_{N}\) line and total support bounded by \(\Gamma_{Qrep}\).
\item Build the nonuniform circuit for \(H_{N}\) from the generator circuit, parsed row computations, charged FormulaEval/TSR/Qrep support, and equality patches for \(E_{pi}\) residual rows, obtaining C(H\_N)<=N\textasciicircum{}\{rho+o(1)\}+Gamma\_Qrep+O(E\_pi N).
\item Contrapose the claimed source lower bound LB\_N(1+epsilon): if the displayed upper bound is below \(N^{1+\epsilon-o(1)}\), the claimed proof-mediated lower-bound certificate is falsified in the same model.
\item For necessity, choose the first unsafe quotient use. Keep the visible prefix fixed, toggle an uncharged private row or unread enable bit that controls same-theory status, and obtain two worlds with identical accepted prefix but different uncovered parity or certificate truth.
\end{enumerate}

\paragraph{Equation 1: carrier.}
\[
H_N^{SAT}(x)=\bigoplus_{j\le M_N} SAT(G_N(x,j)),\quad M_N=N^{q+o(1)},\quad |G_N(x,j)|=(\log N)^d.
\]
\emph{Valid when:} Valid when the beta-zero parity carrier has explicit generator \(G_{N}\), polylog target formulas, and all slots are indexed without advice.

\paragraph{Equation 2: qrep\_ledger.}
\[
\Gamma_{Qrep}(\pi)=\Gamma_{bEF}(\pi)+U(K_{Qrep})+X(K_{Qrep})+\sum_{e\in Dict(\pi)}(|\sigma_e|+|row(e)|)\,N^{o(1)}.
\]
\emph{Valid when:} Valid when K\_Qrep(pi) is the least strict-prefix support closure and every dictionary certificate is charged before representative use.

\paragraph{Equation 3: parser\_upper\_bound.}
\[
C(H_N^{SAT})\le N^{\rho+o(1)}+\Gamma_{Qrep}(\pi)+O(E_\pi N).
\]
\emph{Valid when:} Valid when the line induction is performed in the same nonuniform SIZE model and equality patches charge every residual bad row.

\paragraph{Equation 4: lb\_falsifier\_window.}
\[
N^{\rho+o(1)}+\Gamma_{Qrep}(\pi)+O(E_\pi N)< N^{1+\epsilon-o(1)}\Rightarrow \neg LB_N(1+\epsilon).
\]
\emph{Valid when:} Applies when LB\_N(1+epsilon) asserts a same-model \(N^{1+\epsilon-o(1)}\) lower bound for the named carrier.

\paragraph{Equation 5: support\_density\_corollary.}
\[
LB_N(1+\epsilon)\wedge \rho< 1+\epsilon\Rightarrow |K_{idx}(\pi)|\ge N^{\epsilon-o(1)},\quad |K_{idx}|/M_N\ge N^{\epsilon-q-o(1)}.
\]
\emph{Valid when:} Requires the claimed LB\_N(1+epsilon) to survive the parser upper bound and \(M_{N}=N^{q+o(1)}\) slots.

\paragraph{Next analysis.} Run the prefix-scan verifier on one minimal proof with a valid two-row strict-prefix Qrep merge and on one identical proof with the certificate delayed by one line; compute \(\Gamma_{Qrep}\) in both cases and export the rejected line as the private-row falsifier pivot.

\paragraph{Falsifiers.}
\begin{itemize}
\item First late Qrep merge: representative used before its same-theory certificate is charged, producing a private-row falsifier.
\item Rank cycle or representative rebinding: replay is not well-founded, so \(\Gamma\) reduction is invalid.
\item Advice-chosen dictionary: \(K_{Qrep}\) is not determined by N and pi, causing nonuniform leakage.
\item Hidden semantic pruning: the proof uses real SAT structure without listing influencing rows, blocked at \(RealSATInfluenceRealizer_{N}^{arb}\).
\item Positive \(\beta\) black-box substitution remains a separate threshold-collapse falsifier: q+beta c can exceed 1+epsilon before all SAT exponents are excluded.
\end{itemize}


\section{Current Ranked Next-Analysis Queue}

\label{sec:report-ranked-analyses}

Ordered by target payoff, discriminating power, execution readiness, target integration, and anti-drift pressure. The score is a control signal for the next run, not a theorem-progress claim.

\subsection*{\#1 pvsnp.paper.bit-explicit-parser-formulaeval-boundary: Prove the \(TransparentRowSupport_{N}\) theorem pair and stop selector-case accumulation}

\paragraph{Analysis.} Proof-delivery task sharpened from cycles 103-107: stop adding isolated selector classes and consolidate interval-union, scan, ROBDD, guarded-dDNNF, projected-dDNNF, and influence-complete FormulaEval cases under one \(TransparentRowSupport_{N}\) abstraction. Prove the positive \(\Gamma-extension\) theorem for transparent row support, prove the private-row/unread-enable-bit countermodel theorem for hidden semantic support, and state \(RealSATInfluenceRealizer_{N}\) as the exact remaining escape rather than a vague semantic gap.

\paragraph{Why this is a good target.} \(TransparentRowSupport_{N}\) theorem pair: if a \(\mathrm{Psys}_{sym}\) FormulaEval mechanism has visible selector code, advice-free enumeration or full-scan charge, influence closure, same-theory occurrence expansion, and no hidden support advice, then \(\Gamma_{trans}(pi)=\Gamma_{bEF}(pi)+U_{N}+X_{N}\) gives the parser-to-circuit falsifier N\textasciicircum{}\{rho+o(1)\}+Gamma\_trans(pi)<N\textasciicircum{}\{1+epsilon-o(1)\} for LB\_N(1+epsilon). If any visibility, extensionality, or influence-closure condition fails, construct the private-row/unread-enable-bit countermodel to surface occurrence completion.

\paragraph{Risk.} Main risk: the item becomes new exploration instead of emitting a proof packet or blocked-lemma disposition.

\paragraph{Branch control.} This proof-delivery item must be discharged before broad exploration because it targets the current theorem-shaped blocker.

\paragraph{Source object.} Cycles 103-107 FormulaEval selector packets over \(H_{N}^\mathrm{SAT}(x)=xor_{j\le M_{N}}\) SAT(G\_N(x,j)), with \(M_{N}=N^{q+o(1)}\), generator overhead \(N^{\rho+o(1)}\), polylog formula length, bit-explicit \(\mathrm{Psys}_{bEF}\) parser costs, and selector cases including interval unions, full scans, ROBDD, guarded dDNNF, projected dDNNF, and influence-complete extension rules.

\paragraph{Target object.} A paper-core theorem/countermodel pair replacing the selector zoo: one transparent-support theorem, one hidden-support countermodel theorem, examples as corollaries, and \(RealSATInfluenceRealizer_{N}\) as an open problem.

\paragraph{Transfer condition.} The output must preserve the P vs NP target relation while remaining a scoped route-specific theorem, proof audit, or blocked extraction obligation, not a claim of \(\mathrm{P} \ne \mathrm{NP}\).

\paragraph{Validation steps.}
\begin{itemize}
\item Proof step: define \(TransparentRowSupport_{N}\) as a reusable finite proof-system interface over \(\mathrm{Psys}_{sym}\) FormulaEval lines.
\item Proof step: prove that every transparent FormulaEval line contributes either a certified enumerator charge \(U_{N}\) or a full-scan charge \(F_{FE}\) \(N^{1+q+o(1)}\), plus the row-expansion charge \(X_{N}\).
\item Proof step: extend the existing bit-explicit \(\mathrm{Psys}_{bEF}\) line induction to \(\Gamma_{trans}(pi)=\Gamma_{bEF}(pi)+U_{N}+X_{N}\).
\item Proof step: derive the upper-bound falsifier N\textasciicircum{}\{rho+o(1)\}+Gamma\_trans(pi) not LB\_N(1+epsilon) in the same nonuniform SIZE model.
\item Proof step: prove the hidden-support countermodel: if a FormulaEval row can influence the proof without being visible and charged, two transcripts agree while the selected support or carrier parity differs.
\item Proof step: demote additional selector classes to examples unless they prove or refute a \(TransparentRowSupport_{N}\) clause.
\item Proof step: run proof-quality, barrier-specificity, quantifier-uniformity, proof-complexity-transfer, and inference-quality reviewers before promotion.
\end{itemize}

\paragraph{Payoff criteria.}
\begin{itemize}
\item \(TransparentRowSupport_{N}\) is formal enough that an independent reader can decide whether a FormulaEval rule satisfies it.
\item Gamma\_trans(pi) is derived rather than asserted, and the falsifier inequality states \(\rho\), q, \(\epsilon\), \(U_{N}\), \(X_{N}\), \(\eta\), \(\tau\), \(E_{N}\), and all o(1) regimes.
\item The private-row/unread-enable-bit countermodel is stated as a theorem with hypotheses, not as a metaphor.
\item Selector cases are reorganized into examples/corollaries; no new selector family is promoted as a separate main contribution.
\item The output advances a paper-level subgoal: a clean theorem pair and open-problem boundary, not another restatement of the obstruction ladder.
\end{itemize}

\paragraph{Stop criteria.}
\begin{itemize}
\item Demote if the work adds another selector class without changing the \(TransparentRowSupport_{N}\) definition, theorem, or countermodel.
\item Demote if Gamma\_trans(pi), K\_idx(pi), \(U_{N}\), \(X_{N}\), or the transcript upper bound is named without a derivation.
\item Demote if semantic FormulaEval is described as generic power rather than handled by \(RealSATInfluenceRealizer_{N}\) or the private-row countermodel.
\item Demote if the branch broadens into MCSP/MKtP, natural-proofs folklore, or generic barriers before the theorem pair is written.
\end{itemize}

\paragraph{Do not stop if.}
\begin{itemize}
\item A named lemma can still be proved, repaired, or refuted locally.
\item The current run can emit a live pvsnp\_inference\_quality packet.
\item A countermodel would become a route-specific no-go theorem rather than mere failure.
\end{itemize}

\paragraph{Required skills.}
\begin{itemize}
\item lemma-chain-proof-builder
\item proof-delivery-workflow-orchestrator
\item proof-gap-countermodel-splitter
\item proof-obligation-formalizer
\item proof-quality-gatekeeper
\item pvsnp-barrier-specificity-reviewer
\item pvsnp-conceptual-pressure-localizer
\item pvsnp-forcedrows-proof-workflow
\item pvsnp-impossibility-theorem-forger
\item pvsnp-inference-quality-controller
\item pvsnp-negative-tactic-selector
\item pvsnp-pressure-bridge-forger
\item pvsnp-proof-complexity-transfer-reviewer
\item pvsnp-proof-derivation-builder
\item pvsnp-qualitative-complexity-reasoner
\item pvsnp-quantifier-uniformity-reviewer
\item pvsnp-reduction-consistency-reviewer
\item pvsnp-symbolic-complexity-auditor
\end{itemize}

\subsection*{\#2 pvsnp.paper.l1201-least-closed-support-proof-audit: Write the L120.1 least-closed support proof-audit packet}

\paragraph{Analysis.} Turn cycle 120's least prefix-closed reachable FormulaEval support cone into a referee-checkable proof packet. The task is to define the fixed-point computation of K\_dag(pi), give verifier-visible dependency-closure pseudocode, prove the \(\Gamma_{dag}\) line-induction invariant, and test the theorem on one projected-dDNNF FormulaEval example plus one unread-enable-bit counterexample.

\paragraph{Why this is a good target.} \(LeastClosedTransparentSupport_{N}\) / L120.1: for a \(\mathrm{Psys}_{sym}^{dag}\) proof pi whose FormulaEval lines expose parent pointers and support certificates, the least prefix-closed reachable cone K\_dag(pi) from the final \(\mathrm{LB}_{N}\) line is the minimal charged object needed for the parser-to-circuit falsifier N\textasciicircum{}\{rho+o(1)\}+Gamma\_dag(pi)+O(E\_pi N)<N\textasciicircum{}\{1+epsilon-o(1)\}.

\paragraph{Risk.} Main risk: the item becomes new exploration instead of emitting a proof packet or blocked-lemma disposition.

\paragraph{Branch control.} This proof-delivery item must be discharged before broad exploration because it targets the current theorem-shaped blocker.

\paragraph{Source object.} Cycle 120 theorem claims over \(H_{N}^\mathrm{SAT}(x)=xor_{j\le M_{N}}\) SAT(G\_N(x,j)), \(\mathrm{Psys}_{sym}^{dag}\) FormulaEval lines, dependency edges from later symbolic uses to earlier row supports, and the cost ledger Gamma\_dag(pi)=Gamma\_bEF(pi)+U(K\_dag(pi))+X(K\_dag(pi)).

\paragraph{Target object.} A self-contained proof-audit packet that an independent reader can run mentally or mechanically: input pi,N; compute K\_dag(pi); either certify prefix-closed support and \(\Gamma_{dag}\), or return the first hidden reachable row with its private-row/unread-enable-bit countermodel.

\paragraph{Transfer condition.} The output must preserve the P vs NP target relation while remaining a scoped route-specific theorem, proof audit, or blocked extraction obligation, not a claim of \(\mathrm{P} \ne \mathrm{NP}\).

\paragraph{Validation steps.}
\begin{itemize}
\item Proof step: define the line-dependency DAG, support-token universe, and ancestor relation from the final \(\mathrm{LB}_{N}\) line.
\item Proof step: prove the least fixed point K\_dag(pi) exists, is unique, and is computable from proof syntax and N without advice.
\item Proof step: prove unreachable FormulaEval rows can be deleted without changing the retained line induction.
\item Proof step: prove each retained FormulaEval support either pays U(K)+X(K) or triggers the first-hidden-row falsifier.
\item Proof step: instantiate the positive and negative examples and attach an inference-quality packet before promotion.
\end{itemize}

\paragraph{Payoff criteria.}
\begin{itemize}
\item The fixed-point algorithm for K\_dag(pi) is explicit enough to implement or audit line by line.
\item \(\Gamma_{dag}\) is derived from the closure invariant, not asserted by analogy with \(\Gamma_{trans}\).
\item The projected-dDNNF example and unread-enable-bit counterexample both exercise the theorem boundary.
\item The result states that it is a scoped parser/falsifier theorem and not evidence for \(\mathrm{P} \ne \mathrm{NP}\).
\end{itemize}

\paragraph{Stop criteria.}
\begin{itemize}
\item Demote if K\_dag(pi) is described informally without a computable closure algorithm.
\item Demote if the proof does not distinguish reachable support from the full selector zoo.
\item Demote if the examples do not compute charges or do not expose a first hidden reachable row.
\item Demote if the theorem silently assumes full semantic SAT influence instead of charging it.
\end{itemize}

\paragraph{Do not stop if.}
\begin{itemize}
\item A named lemma can still be proved, repaired, or refuted locally.
\item The current run can emit a live pvsnp\_inference\_quality packet.
\item A countermodel would become a route-specific no-go theorem rather than mere failure.
\end{itemize}

\paragraph{Required skills.}
\begin{itemize}
\item lemma-chain-proof-builder
\item proof-delivery-workflow-orchestrator
\item proof-gap-countermodel-splitter
\item proof-obligation-formalizer
\item proof-quality-gatekeeper
\item pvsnp-barrier-specificity-reviewer
\item pvsnp-conceptual-pressure-localizer
\item pvsnp-forcedrows-proof-workflow
\item pvsnp-impossibility-theorem-forger
\item pvsnp-inference-quality-controller
\item pvsnp-negative-tactic-selector
\item pvsnp-pressure-bridge-forger
\item pvsnp-proof-complexity-transfer-reviewer
\item pvsnp-proof-derivation-builder
\item pvsnp-qualitative-complexity-reasoner
\item pvsnp-quantifier-uniformity-reviewer
\item pvsnp-reduction-consistency-reviewer
\item pvsnp-symbolic-complexity-auditor
\end{itemize}

\subsection*{\#3 pvsnp.paper.realsat-influence-realizer-boundary: Resolve or quarantine \(RealSATInfluenceRealizer_{N}\) for one concrete semantic FormulaEval rule}

\paragraph{Analysis.} Attack the exact remaining full-semantic FormulaEval gap instead of naming it abstractly. Choose one concrete semantic FormulaEval rule, such as Tseitin-extension equality plus bounded self-reduction, and either produce an advice-free nonrelativizing realizer listing all influencing rows with charges, or prove a rule-specific private-row/no-realizer countermodel.

\paragraph{Why this is a good target.} \(RealSATInfluenceRealizer_{N}\) boundary theorem: full semantic FormulaEval can be folded into the transparent support theorem only when proof syntax and N uniformly determine an influence-closed row set \(R_{A}\), same-theory row expansions, and a charged overhead \(R_{N}\) before the semantic rule is used.

\paragraph{Risk.} Main risk: the item becomes new exploration instead of emitting a proof packet or blocked-lemma disposition.

\paragraph{Branch control.} This proof-delivery item must be discharged before broad exploration because it targets the current theorem-shaped blocker.

\paragraph{Source object.} Semantic FormulaEval-SAT rules over \(H_{N}^\mathrm{SAT}\) generated rows where real SAT semantics, formula equality, Tseitin extensions, or self-reduction may make a proof line depend on rows not visible in a syntactic \(TransparentRowSupport_{N}\) certificate.

\paragraph{Target object.} A yes/no boundary for one named semantic rule: either a valid \(RealSATInfluenceRealizer_{N}\) certificate with \(\Gamma_{real}=\Gamma_{bEF}+U+X+R_{N}\), or a finite countermodel showing why that semantic rule remains quarantined.

\paragraph{Transfer condition.} The output must preserve the P vs NP target relation while remaining a scoped route-specific theorem, proof audit, or blocked extraction obligation, not a claim of \(\mathrm{P} \ne \mathrm{NP}\).

\paragraph{Validation steps.}
\begin{itemize}
\item Proof step: pick exactly one semantic FormulaEval rule and freeze its proof-system syntax.
\item Proof step: define which generated SAT rows can influence the rule value and what counts as a charged witness.
\item Proof step: attempt to construct \(R_{A}\) and same-theory expansions from proof syntax and N alone.
\item Proof step: if the construction fails, build the private-row/unread-enable-bit countermodel for that rule.
\item Proof step: record whether the rule joins the transparent theorem or remains quarantined as semantic power.
\end{itemize}

\paragraph{Payoff criteria.}
\begin{itemize}
\item The selected semantic rule is concrete enough that the realizer or countermodel is independently checkable.
\item All influence rows are uniformly listed and charged, or the exact unchargeable semantic step is named.
\item The output eliminates vague language about semantic FormulaEval and replaces it with a rule-level theorem or no-go.
\end{itemize}

\paragraph{Stop criteria.}
\begin{itemize}
\item Demote if \(RealSATInfluenceRealizer_{N}\) is only restated without choosing a semantic rule.
\item Demote if the row set uses SAT-truth advice or hidden nonuniform row selection.
\item Demote if the claimed realizer omits \(R_{N}\) from the \(\Gamma\) ledger.
\item Demote if the negative result does not identify the first semantic inference that cannot be charged.
\end{itemize}

\paragraph{Do not stop if.}
\begin{itemize}
\item A named lemma can still be proved, repaired, or refuted locally.
\item The current run can emit a live pvsnp\_inference\_quality packet.
\item A countermodel would become a route-specific no-go theorem rather than mere failure.
\end{itemize}

\paragraph{Required skills.}
\begin{itemize}
\item lemma-chain-proof-builder
\item proof-delivery-workflow-orchestrator
\item proof-gap-countermodel-splitter
\item proof-obligation-formalizer
\item proof-quality-gatekeeper
\item pvsnp-barrier-specificity-reviewer
\item pvsnp-conceptual-pressure-localizer
\item pvsnp-forcedrows-proof-workflow
\item pvsnp-impossibility-theorem-forger
\item pvsnp-inference-quality-controller
\item pvsnp-negative-tactic-selector
\item pvsnp-pressure-bridge-forger
\item pvsnp-proof-complexity-transfer-reviewer
\item pvsnp-proof-derivation-builder
\item pvsnp-qualitative-complexity-reasoner
\item pvsnp-quantifier-uniformity-reviewer
\item pvsnp-reduction-consistency-reviewer
\item pvsnp-symbolic-complexity-auditor
\end{itemize}

\subsection*{\#4 pvsnp.paper.l120-theorem-countermodel-package: Package the cycle-120 theorem/countermodel pair as a paper-ready module}

\paragraph{Analysis.} Assemble the least-closed transparent-support theorem, hidden-support countermodel, \(RealSATInfluenceRealizer_{N}\) open boundary, q-trichotomy support-density corollary, and selector-example table into one coherent paper module. The package must be readable without marathon context.

\paragraph{Why this is a good target.} Paper module: a scoped parser/falsifier contribution for FormulaEval in beta-zero parity magnification, showing exactly when transparent FormulaEval yields a same-model upper-bound falsifier and exactly where hidden semantic influence escapes the controlled interface.

\paragraph{Risk.} Main risk: the item becomes new exploration instead of emitting a proof packet or blocked-lemma disposition.

\paragraph{Branch control.} This proof-delivery item must be discharged before broad exploration because it targets the current theorem-shaped blocker.

\paragraph{Source object.} Cycles 116-120, especially L120.1 \(LeastClosedTransparentSupport_{N}\), L120.2 \(UnreachableSupportDeletion_{N}\), L120.3 \(FirstHiddenReachableUse_{N}\), and the branch outputs that state \(\Gamma_{dag}\), hidden-row falsifiers, and qepsilon support-density cases.

\paragraph{Target object.} A paper-ready theorem/countermodel section with definitions, theorem statements, proof sketches, examples, non-claims, and a handoff to the broader parity magnification/full-row barrier draft.

\paragraph{Transfer condition.} The output must preserve the P vs NP target relation while remaining a scoped route-specific theorem, proof audit, or blocked extraction obligation, not a claim of \(\mathrm{P} \ne \mathrm{NP}\).

\paragraph{Validation steps.}
\begin{itemize}
\item Proof step: consolidate definitions and notation from cycles 116-120 into one local ledger.
\item Proof step: write theorem/countermodel statements before examples, so selector cases cannot become separate main routes.
\item Proof step: attach proof sketches that cite the L120.1 audit packet and \(RealSAT\) boundary item where needed.
\item Proof step: state non-claims and model boundaries explicitly.
\item Proof step: validate against proof-delivery and inference-quality gates.
\end{itemize}

\paragraph{Payoff criteria.}
\begin{itemize}
\item The module is understandable as a standalone paper section.
\item The positive theorem, negative countermodel, and semantic open boundary are separated cleanly.
\item Examples are corollaries/examples only; no new selector zoo is created.
\item The paper handoff preserves the scoped claim and does not oversell P vs NP progress.
\end{itemize}

\paragraph{Stop criteria.}
\begin{itemize}
\item Demote if the package is a report summary rather than theorem/countermodel text.
\item Demote if it depends on marathon-specific workflow language.
\item Demote if it blurs transparent extensional FormulaEval with full semantic SAT FormulaEval.
\item Demote if it omits the q-trichotomy or the non-claims.
\end{itemize}

\paragraph{Do not stop if.}
\begin{itemize}
\item A named lemma can still be proved, repaired, or refuted locally.
\item The current run can emit a live pvsnp\_inference\_quality packet.
\item A countermodel would become a route-specific no-go theorem rather than mere failure.
\end{itemize}

\paragraph{Required skills.}
\begin{itemize}
\item lemma-chain-proof-builder
\item proof-delivery-workflow-orchestrator
\item proof-gap-countermodel-splitter
\item proof-obligation-formalizer
\item proof-quality-gatekeeper
\item pvsnp-barrier-specificity-reviewer
\item pvsnp-proof-complexity-transfer-reviewer
\item pvsnp-quantifier-uniformity-reviewer
\item pvsnp-reduction-consistency-reviewer
\item pvsnp-symbolic-complexity-auditor
\end{itemize}

\subsection*{\#5 pvsnp.paper.l120-inference-quality-gate-closure: Close the L120 inference-quality and marker gaps before promotion}

\paragraph{Analysis.} Prepare the promotion gate for the L120 theorem work. Emit a fresh pvsnp\_inference\_quality packet for the least-closed support theorem, run the reviewer stack, and repair any source/target/model/validity marker gaps before treating the theorem package as a paper result.

\paragraph{Why this is a good target.} Validation theorem for promotion readiness: every L120 claim has explicit source object, target object, controlled interface, target-preserving condition, reduction direction, quantifiers, advice discipline, barrier scope, and proof-complexity transfer status.

\paragraph{Risk.} Main risk: the item becomes new exploration instead of emitting a proof packet or blocked-lemma disposition.

\paragraph{Branch control.} This proof-delivery item must be discharged before broad exploration because it targets the current theorem-shaped blocker.

\paragraph{Source object.} The live L120 claims, the 93 active inference-quality packets, and the discovery report warning that not all source/target/same-theory/validity markers are globally present.

\paragraph{Target object.} A repaired inference-quality packet and gate report that says exactly which L120 claims are promotable, which are conditional, and which remain blocked at \(RealSATInfluenceRealizer_{N}\).

\paragraph{Transfer condition.} The output must preserve the P vs NP target relation while remaining a scoped route-specific theorem, proof audit, or blocked extraction obligation, not a claim of \(\mathrm{P} \ne \mathrm{NP}\).

\paragraph{Validation steps.}
\begin{itemize}
\item Proof step: write the IQ packet from the actual theorem text, not from a marathon summary.
\item Proof step: run each reviewer lens and convert objections into proof obligations or demotion notes.
\item Proof step: repair missing markers in the queue/report artifacts if validation finds them.
\item Proof step: record a final promotion verdict before paper integration.
\end{itemize}

\paragraph{Payoff criteria.}
\begin{itemize}
\item The L120 theorem package has a live IQ packet that passes validation.
\item No source/target/model/validity marker is implicit.
\item Reviewer objections are either discharged, converted to named obligations, or marked as demotion blockers.
\item The promotion verdict is conservative and explicitly says what is not proved.
\end{itemize}

\paragraph{Stop criteria.}
\begin{itemize}
\item Demote if the IQ packet claims P vs NP progress rather than a scoped parser/falsifier result.
\item Demote if source and target models are not the same nonuniform SIZE model.
\item Demote if \(RealSATInfluenceRealizer_{N}\) is treated as solved without a rule-level proof.
\item Demote if validation fails and the failure is not converted into a concrete repair item.
\end{itemize}

\paragraph{Do not stop if.}
\begin{itemize}
\item A named lemma can still be proved, repaired, or refuted locally.
\item The current run can emit a live pvsnp\_inference\_quality packet.
\item A countermodel would become a route-specific no-go theorem rather than mere failure.
\end{itemize}

\paragraph{Required skills.}
\begin{itemize}
\item proof-obligation-formalizer
\item proof-quality-gatekeeper
\item pvsnp-barrier-specificity-reviewer
\item pvsnp-inference-quality-controller
\item pvsnp-meta-complexity-carrier-reviewer
\item pvsnp-proof-complexity-transfer-reviewer
\item pvsnp-quantifier-uniformity-reviewer
\item pvsnp-reduction-consistency-reviewer
\end{itemize}

\subsection*{\#6 pvsnp.paper.full-row-barrier-proof-audit: Audit and package the draft's proved beta-zero/full-row barrier theorem chain}

\paragraph{Analysis.} Proof-audit task extracted from the draft theorem stack: rebuild the exponent-window, batched substitution, beta-zero transfer, equality patching, naturalization, row-forcing, parity full-row, parser-to-query, incidence-sketch, and one-sided forcing proofs as a dependency-checked proof packet.

\paragraph{Why this is a good target.} Verified theorem-chain packet: beta-zero transfers remove exponent debt; functional certificates naturalize; occurrence-relativizing parity certificates pay Omega(M\_n \(s_{n}/n\)).

\paragraph{Risk.} Main risk: the item becomes new exploration instead of emitting a proof packet or blocked-lemma disposition.

\paragraph{Branch control.} This proof-delivery item must be discharged before broad exploration because it targets the current theorem-shaped blocker.

\paragraph{Source object.} The draft's proposition/lemma/theorem environments and their dependency graph.

\paragraph{Target object.} A paper-ready proof packet with explicit hypotheses, quantifiers, non-claims, and theorem dependencies.

\paragraph{Transfer condition.} The output must preserve the P vs NP target relation while remaining a scoped route-specific theorem, proof audit, or blocked extraction obligation, not a claim of \(\mathrm{P} \ne \mathrm{NP}\).

\paragraph{Validation steps.}
\begin{itemize}
\item Proof step: extract every theorem/proposition/corollary proof from the draft and list its exact assumptions.
\item Proof step: check exponent arithmetic and the beta-zero disappearance of the target circuit exponent c.
\item Proof step: verify that equality patching and residual threshold use only row exceptions and O(n) equality predicates.
\item Proof step: verify naturalization hypotheses are dense, constructive, and functional over arbitrary label tables.
\item Proof step: verify row-forcing uses occurrence-completion soundness and not real SAT semantics.
\end{itemize}

\paragraph{Payoff criteria.}
\begin{itemize}
\item Every proved draft statement has assumptions, conclusion, proof dependencies, and barrier scope recorded.
\item No step silently upgrades occurrence-label statements to formula-label or full proof-system statements.
\item The theorem chain is ready for paper insertion or referee-style proof checking.
\end{itemize}

\paragraph{Stop criteria.}
\begin{itemize}
\item Demote if any proof uses an unstated uniformity, density, constructivity, or formula-realizability assumption.
\item Demote if theorem dependencies are circular or if the draft's non-claims are not preserved.
\end{itemize}

\paragraph{Do not stop if.}
\begin{itemize}
\item A named lemma can still be proved, repaired, or refuted locally.
\item The current run can emit a live pvsnp\_inference\_quality packet.
\item A countermodel would become a route-specific no-go theorem rather than mere failure.
\end{itemize}

\paragraph{Required skills.}
\begin{itemize}
\item lemma-chain-proof-builder
\item proof-delivery-workflow-orchestrator
\item proof-obligation-formalizer
\item proof-quality-gatekeeper
\item pvsnp-barrier-specificity-reviewer
\item pvsnp-proof-complexity-transfer-reviewer
\item pvsnp-quantifier-uniformity-reviewer
\item pvsnp-reduction-consistency-reviewer
\item pvsnp-symbolic-complexity-auditor
\end{itemize}

\subsection*{\#7 pvsnp.paper.explicit-hard-generator-open-problem: Attack the explicit hard-generator/source-lower-bound open problem}

\paragraph{Analysis.} Open-problem task extracted from the draft discussion: construct or rule out a label-independent batched generator with polylogarithmic SAT instances, \(M_{n}=n^{q+o(1)}\), generator size \(n^{\rho+o(1)}\), \(q,rhon^{1+\epsilon-o(1)}\).

\paragraph{Why this is a good target.} Source-lower-bound theorem or impossibility diagnosis for a concrete beta-zero SAT parity carrier.

\paragraph{Risk.} Main risk: the item becomes new exploration instead of emitting a proof packet or blocked-lemma disposition.

\paragraph{Branch control.} This proof-delivery item must be discharged before broad exploration because it targets the current theorem-shaped blocker.

\paragraph{Source object.} Candidate batched generator \(G_{n}\) and parity carrier \(H_{n}^\mathrm{SAT}\) from polylog-length SAT formulas.

\paragraph{Target object.} A source lower bound strong enough to activate the beta-zero magnification corollary, or a precise obstruction explaining why the candidate collapses.

\paragraph{Transfer condition.} The output must preserve the P vs NP target relation while remaining a scoped route-specific theorem, proof audit, or blocked extraction obligation, not a claim of \(\mathrm{P} \ne \mathrm{NP}\).

\paragraph{Validation steps.}
\begin{itemize}
\item Proof step: define the generator and compute q, \(\rho\), d, \(|U_{n}|\), collision structure, and formula-realizability constraints.
\item Proof step: run incidence-sketch falsification before claiming lower-bound promise.
\item Proof step: attempt a lower-bound route or derive an impossibility/upper-bound certificate.
\item Proof step: verify that any lower bound is for general circuits on the scalar carrier, not only incidence rank or restricted models.
\end{itemize}

\paragraph{Payoff criteria.}
\begin{itemize}
\item A concrete generator is named with verified beta-zero transfer parameters.
\item The branch produces either a real lower-bound lemma, a reproducible upper-bound falsifier, or a sharp reason it cannot be the source problem.
\item Any claimed progress distinguishes scalar circuit lower bounds from incidence-linear diagnostics.
\end{itemize}

\paragraph{Stop criteria.}
\begin{itemize}
\item Demote if no explicit generator is specified.
\item Demote if the evidence is only high rank, randomness intuition, or incidence complexity without a general-circuit lower-bound bridge.
\item Demote if q or \(\rho\) fails the stated threshold needed by the transfer theorem.
\end{itemize}

\paragraph{Do not stop if.}
\begin{itemize}
\item A named lemma can still be proved, repaired, or refuted locally.
\item The current run can emit a live pvsnp\_inference\_quality packet.
\item A countermodel would become a route-specific no-go theorem rather than mere failure.
\end{itemize}

\paragraph{Required skills.}
\begin{itemize}
\item lemma-chain-proof-builder
\item proof-delivery-workflow-orchestrator
\item proof-gap-countermodel-splitter
\item proof-quality-gatekeeper
\item pvsnp-barrier-specificity-reviewer
\item pvsnp-conceptual-pressure-localizer
\item pvsnp-forcedrows-proof-workflow
\item pvsnp-impossibility-theorem-forger
\item pvsnp-qualitative-complexity-reasoner
\item pvsnp-symbolic-complexity-auditor
\end{itemize}

\subsection*{\#8 pvsnp.paper.strong-system-parser-instantiation: Instantiate the parser-to-query consequence for a concrete strong proof system}

\paragraph{Analysis.} Open-problem/proof task extracted from the draft's strong-system paragraph and Proposition parser-to-query: choose a concrete proof formalism for lower-bound statements, define exactly which formula equality and SAT semantic facts it can access, and prove or refute the parser property that maps accepted proofs pi to occurrence transcripts \(Q_{pi}\).

\paragraph{Why this is a good target.} Concrete parser-to-query instantiation theorem or counterexample for a named proof system.

\paragraph{Risk.} Main risk: the item becomes new exploration instead of emitting a proof packet or blocked-lemma disposition.

\paragraph{Branch control.} This proof-delivery item must be discharged before broad exploration because it targets the current theorem-shaped blocker.

\paragraph{Source object.} Named proof formalism, proof pi, verifier/parser, and allowed SAT/equality axioms.

\paragraph{Target object.} A proof-system-specific theorem linking proof length/access to Omega(M\_n \(s_{n}/n\)), or an escape principle that blocks the link.

\paragraph{Transfer condition.} The output must preserve the P vs NP target relation while remaining a scoped route-specific theorem, proof audit, or blocked extraction obligation, not a claim of \(\mathrm{P} \ne \mathrm{NP}\).

\paragraph{Validation steps.}
\begin{itemize}
\item Proof step: select the proof formalism and define the lower-bound sentence it proves.
\item Proof step: map each proof rule to functional, occurrence-local, or formula-semantic behavior.
\item Proof step: prove parser invariance under occurrence completions or construct a rule-level countermodel.
\item Proof step: if parser holds, derive the Omega(M\_n \(s_{n}/n\)) access/proof-size consequence.
\end{itemize}

\paragraph{Payoff criteria.}
\begin{itemize}
\item The selected proof system and its SAT access are explicit enough for independent checking.
\item The parser property is proved with a quantitative p(T,n), or a precise semantic escape is isolated.
\item The output is not a generic 'strong systems are hard' statement; it names the exact proof rule that matters.
\end{itemize}

\paragraph{Stop criteria.}
\begin{itemize}
\item Demote if the proof system remains unnamed or if its SAT facts are unspecified.
\item Demote if the parser theorem assumes occurrence locality instead of deriving it.
\item Demote if a negative answer does not identify the formula-semantic principle defeating occurrence-wise completion.
\end{itemize}

\paragraph{Do not stop if.}
\begin{itemize}
\item A named lemma can still be proved, repaired, or refuted locally.
\item The current run can emit a live pvsnp\_inference\_quality packet.
\item A countermodel would become a route-specific no-go theorem rather than mere failure.
\end{itemize}

\paragraph{Required skills.}
\begin{itemize}
\item lemma-chain-proof-builder
\item proof-delivery-workflow-orchestrator
\item proof-gap-countermodel-splitter
\item proof-obligation-formalizer
\item proof-quality-gatekeeper
\item pvsnp-forcedrows-proof-workflow
\item pvsnp-inference-quality-controller
\item pvsnp-negative-tactic-selector
\item pvsnp-proof-complexity-transfer-reviewer
\item pvsnp-quantifier-uniformity-reviewer
\item pvsnp-reduction-consistency-reviewer
\end{itemize}

\subsection*{\#9 pvsnp.paper.incidence-falsifier-normal-form: Use incidence sketches as falsifiers and define the missing normal-form bridge}

\paragraph{Analysis.} Proof/diagnostic task extracted from the appendix and discussion: for candidate generators, search for low-rank or sparse affine-incidence sketches \(p_{x}=sum_{i}\) \(a_{i}(x)b_{i}\) off a small exceptional set, and state the exact nonrelativizing normal-form theorem that would be needed to turn absence of such sketches into lower-bound evidence.

\paragraph{Why this is a good target.} Incidence-falsifier theorem plus normal-form open problem for translating diagnostics into lower bounds.

\paragraph{Risk.} Main risk: the item becomes new exploration instead of emitting a proof packet or blocked-lemma disposition.

\paragraph{Branch control.} This proof-delivery item must be discharged before broad exploration because it targets the current theorem-shaped blocker.

\paragraph{Source object.} Generated formula universe \(U_{n}\), incidence vectors p\_x in \(F^{U_{n}}\), coefficient functions a\_i, sketch vectors b\_i, and exception set \(B_{n}\).

\paragraph{Target object.} A reproducible upper-bound falsifier or a clearly scoped normal-form theorem needed before incidence diagnostics can support source lower bounds.

\paragraph{Transfer condition.} The output must preserve the P vs NP target relation while remaining a scoped route-specific theorem, proof audit, or blocked extraction obligation, not a claim of \(\mathrm{P} \ne \mathrm{NP}\).

\paragraph{Validation steps.}
\begin{itemize}
\item Proof step: instantiate \(U_{n}\) and p\_x for the candidate carrier.
\item Proof step: attempt affine/sparse incidence sketches and compute \(A_{n}\), K, and \(E_{n}\).
\item Proof step: if a sketch is found, derive the explicit upper-bound contradiction to the proposed source lower bound.
\item Proof step: if no sketch is found, record this only as nonproof evidence and state the missing scalar-circuit normal form.
\end{itemize}

\paragraph{Payoff criteria.}
\begin{itemize}
\item Every candidate generator is checked for the appendix's explicit upper-bound mechanism before promotion.
\item The output never treats absence of a sketch as a lower bound.
\item The normal-form bridge is stated as a theorem obligation with assumptions strong enough to imply scalar circuit hardness.
\end{itemize}

\paragraph{Stop criteria.}
\begin{itemize}
\item Demote if the branch uses incidence rank as a lower bound without a scalar-circuit bridge.
\item Demote if K, \(A_{n}\), \(E_{n}\), or exceptional-row patch cost are not computed or bounded.
\item Demote if the diagnostic does not feed either a falsifier or a specific normal-form theorem.
\end{itemize}

\paragraph{Do not stop if.}
\begin{itemize}
\item A named lemma can still be proved, repaired, or refuted locally.
\item The current run can emit a live pvsnp\_inference\_quality packet.
\item A countermodel would become a route-specific no-go theorem rather than mere failure.
\end{itemize}

\paragraph{Required skills.}
\begin{itemize}
\item lemma-chain-proof-builder
\item proof-gap-countermodel-splitter
\item proof-obligation-formalizer
\item proof-quality-gatekeeper
\item pvsnp-barrier-specificity-reviewer
\item pvsnp-meta-complexity-carrier-reviewer
\item pvsnp-reduction-consistency-reviewer
\item pvsnp-symbolic-complexity-auditor
\end{itemize}

\subsection*{\#10 pvsnp.paper.aggregator-design-open-problem: Search aggregator designs beyond parity with explicit transfer and certificate tradeoffs}

\paragraph{Analysis.} Open-direction task extracted from the draft: compare parity with alternative aggregators by computing transfer cost, C\_0(A\_M), C\_1(A\_M), one-sided certificate behavior, algebraic structure, and whether the aggregator preserves beta-zero all-exponent SAT transfer without reopening easy selector/separator certificates.

\paragraph{Why this is a good target.} Aggregator tradeoff theorem or no-go table identifying which aggregators preserve the desired certificate debt.

\paragraph{Risk.} Main risk: the item becomes new exploration instead of emitting a proof packet or blocked-lemma disposition.

\paragraph{Branch control.} This proof-delivery item must be discharged before broad exploration because it targets the current theorem-shaped blocker.

\paragraph{Source object.} Candidate row aggregator family \(A_{M}\) and batched SAT labels y(G\_\{n,j\}(x)).

\paragraph{Target object.} A ranked set of aggregator candidates with proof-relevant certificate-complexity and transfer-cost guarantees.

\paragraph{Transfer condition.} The output must preserve the P vs NP target relation while remaining a scoped route-specific theorem, proof audit, or blocked extraction obligation, not a claim of \(\mathrm{P} \ne \mathrm{NP}\).

\paragraph{Validation steps.}
\begin{itemize}
\item Proof step: compute certificate complexities for each candidate aggregator.
\item Proof step: verify size O(M) or record the exact transfer cost penalty.
\item Proof step: test whether one-sided SAT evidence can cheaply force rows.
\item Proof step: promote only candidates that improve some tradeoff without losing the all-polynomial-exponent transfer.
\end{itemize}

\paragraph{Payoff criteria.}
\begin{itemize}
\item Parity remains justified or a better aggregator is identified by explicit C\_b(A\_M) and transfer calculations.
\item Any proposed aggregator has proof-locality consequences stated as theorem obligations, not analogies.
\item The output explains whether it strengthens, weakens, or complements the full-row barrier.
\end{itemize}

\paragraph{Stop criteria.}
\begin{itemize}
\item Demote if aggregator comparison does not compute \(C_{0}\) and \(C_{1}\).
\item Demote if the aggregator creates a cheap one-sided certificate that defeats the intended full-row burden.
\item Demote if transfer cost reintroduces positive \(\beta\) exponent debt.
\end{itemize}

\paragraph{Do not stop if.}
\begin{itemize}
\item A named lemma can still be proved, repaired, or refuted locally.
\item The current run can emit a live pvsnp\_inference\_quality packet.
\item A countermodel would become a route-specific no-go theorem rather than mere failure.
\end{itemize}

\paragraph{Required skills.}
\begin{itemize}
\item proof-obligation-formalizer
\item proof-quality-gatekeeper
\item pvsnp-barrier-specificity-reviewer
\item pvsnp-conceptual-pressure-localizer
\item pvsnp-qualitative-complexity-reasoner
\item pvsnp-symbolic-complexity-auditor
\end{itemize}

\subsection*{\#11 pvsnp.paper.occurrence-row-forcing-proposition-delivery: Deliver the occurrence-local row-forcing theorem as a standalone scoped contribution}

\paragraph{Analysis.} Proof-delivery task extracted from the proved row-forcing section: package the occurrence-completion theorem, parity full-row corollary, adaptive verifier corollary, and one-sided variant as a standalone barrier theorem with exact scope.

\paragraph{Why this is a good target.} Standalone occurrence-local certificate-debt theorem for beta-zero parity carriers.

\paragraph{Risk.} Main risk: the item becomes new exploration instead of emitting a proof packet or blocked-lemma disposition.

\paragraph{Branch control.} This proof-delivery item must be discharged before broad exploration because it targets the current theorem-shaped blocker.

\paragraph{Source object.} Completion-sound partial transcript (Q,alpha), forced rows R\_b(Q,alpha), aggregator certificate complexity C\_b(A\_M), and patch circuits.

\paragraph{Target object.} A scoped publishable theorem: occurrence-relativizing lower-bound certificates must pay full-row cost for parity.

\paragraph{Transfer condition.} The output must preserve the P vs NP target relation while remaining a scoped route-specific theorem, proof audit, or blocked extraction obligation, not a claim of \(\mathrm{P} \ne \mathrm{NP}\).

\paragraph{Validation steps.}
\begin{itemize}
\item Proof step: restate completion soundness, forced rows, and certificate complexity without SAT-semantics leakage.
\item Proof step: reconstruct baseline completion z\textasciicircum{}(b) and verify it extends the transcript.
\item Proof step: patch forced rows and derive |R\_b|=Omega(s\_n/n).
\item Proof step: specialize to parity and adaptive verifiers.
\end{itemize}

\paragraph{Payoff criteria.}
\begin{itemize}
\item The theorem is independently understandable without the rest of the P vs NP discovery machinery.
\item The proof states exactly why formula-semantic SAT proofs are outside scope.
\item The theorem can be cited as the paper's local-certificate barrier contribution.
\end{itemize}

\paragraph{Stop criteria.}
\begin{itemize}
\item Demote if the proof assumes formula-realizability or SAT semantics.
\item Demote if the sparse patch construction does not produce a circuit below the claimed threshold.
\item Demote if the theorem cannot state its non-claims cleanly.
\end{itemize}

\paragraph{Do not stop if.}
\begin{itemize}
\item A named lemma can still be proved, repaired, or refuted locally.
\item The current run can emit a live pvsnp\_inference\_quality packet.
\item A countermodel would become a route-specific no-go theorem rather than mere failure.
\end{itemize}

\paragraph{Required skills.}
\begin{itemize}
\item lemma-chain-proof-builder
\item proof-obligation-formalizer
\item proof-quality-gatekeeper
\item pvsnp-barrier-specificity-reviewer
\item pvsnp-forcedrows-proof-workflow
\item pvsnp-quantifier-uniformity-reviewer
\item pvsnp-reduction-consistency-reviewer
\end{itemize}

\subsection*{\#12 cycle-0098: Hardness-Magnification Bridge Realization}

\paragraph{Analysis.} Write the formal \(FE_{Extract,N}\) proposition in two columns: left column row-expanded FormulaEval axiom schema with \(N^{1+q+o(1)}\) charge, right column semantic-oracle version with the smallest private-row countermodel; verify which column matches each \(\mathrm{Psys}_{sym}\) rule.

\paragraph{Why this is a good target.} It directly tests the current black-box substitution cap by computing \(\beta\), q, and the induced target-exponent window. It audits a meta-complexity carrier where the decisive questions are verifier overhead, density/naturalness, and advice leakage. It checks whether padding/projection transfers really magnify hardness or merely move the proof debt into sparse lower bounds. It keeps the barrier and model-boundary audit attached to the mathematical transfer rather than as after-the-fact commentary.

\paragraph{Risk.} Low drift risk: the item is framed as an explicit P vs NP transfer audit.

\paragraph{Branch control.} The item has high payoff and a concrete decision test; stopping now would be premature.

\paragraph{Premature-stop guard.} Do not stop before deriving the transfer inequality from \(\mathrm{SAT}\in \mathrm{SIZE}(n^{c})\) to the carrier upper bound, identifying whether \(\beta=0\), \(\beta>0\), or the argument is genuinely nonblackbox, and recording the barrier status of the proposed escape.

\paragraph{Payoff criteria.}
\begin{itemize}
\item For this route, state the exact barrier premise it uses or escapes: relativizing, natural, algebrizing, nonuniform, or advice-leaking; promotion requires the escape to be attached to the specific transfer theorem.
\item Compute \(\beta\) and q exactly; promotion requires \(\beta=0\), subpolynomial target queries, or a genuinely nonblackbox use of the SAT circuit.
\item For MCSP/MKtP, state N, n, s, verifier overhead, lower-bound model, and whether the property is natural or non-natural.
\item State whether the transfer excludes every polynomial SAT exponent or only a bounded interval \(c<c0\).
\item Record relativization, natural-proofs, algebrization, nonuniformity, and advice-leakage status for the specific transfer theorem.
\item Control gate: name the source carrier or language, target class, and circuit or resource model before claiming progress.
\end{itemize}

\paragraph{Stop criteria.}
\begin{itemize}
\item Every named carrier in the branch has \(\beta>0\) and yields only a bounded target-exponent exclusion.
\item The branch cannot state a source lower-bound model, target class, and transfer inequality in the same notation.
\item The proposed escape is only a black-box substitution restatement with no new carrier or nonblackbox mechanism.
\item The barrier audit shows the needed source lower bound is just as hard as the original P vs NP/P/poly target.
\end{itemize}

\paragraph{Do not stop if.}
\begin{itemize}
\item The branch has not yet computed \(\beta\), q, and the induced source upper-bound exponent under \(\mathrm{SAT}\in \mathrm{SIZE}(n^{c})\).
\item A subpolynomial-query or nonblackbox carrier remains available and has not been audited.
\item The latest failure sharpened the threshold-collapse lemma into a formal no-go statement.
\item The route can still isolate advice leakage, nonuniformity, naturalness, or algebrization as the exact obstruction.
\end{itemize}

\paragraph{Required skills.}
\begin{itemize}
\item abducer-failure-signal-cartographer
\item contradiction-route-auditor
\item counterexample-hunter
\item entailment-gap-auditor
\item kb-abduce
\item kb-concept-bridge-planner
\item kb-deduce
\item kb-end-to-end-proof-critic
\item kb-hypothesis-refutation-loop
\item kb-induce
\item long-horizon-route-portfolio-manager
\item main-problem-alignment-auditor
\item next-experiment-queue-ranker
\item paper-contribution-readiness-judge
\item pvsnp-barrier-specificity-reviewer
\item pvsnp-conceptual-pressure-localizer
\item pvsnp-impossibility-theorem-forger
\item pvsnp-inference-quality-controller
\item pvsnp-meta-complexity-carrier-reviewer
\item pvsnp-proof-derivation-builder
\item pvsnp-qualitative-complexity-reasoner
\item pvsnp-quantifier-uniformity-reviewer
\item pvsnp-reduction-consistency-reviewer
\item pvsnp-symbolic-complexity-auditor
\end{itemize}

\subsection*{\#13 cycle-0094: Hardness-Magnification Bridge Realization}

\paragraph{Analysis.} Formalize \(\mathrm{Psys}_{sym}\) on finite length N: specify whether FormulaEval-SAT comprehension is a proof line, axiom schema, or extension rule; then prove either an expansion-charge parser with cost \(N^{1+q+o(1)}\) or a minimal countermodel showing surface occurrence completion is unsound.

\paragraph{Why this is a good target.} It directly tests the current black-box substitution cap by computing \(\beta\), q, and the induced target-exponent window. It checks whether padding/projection transfers really magnify hardness or merely move the proof debt into sparse lower bounds. It points toward the only unquarantined escape: a subpolynomial-query, nonprojection, or nonblackbox carrier. It keeps the barrier and model-boundary audit attached to the mathematical transfer rather than as after-the-fact commentary. It is executable as a bounded complexity calculation, CAS-backed exponent audit, or formal transfer check.

\paragraph{Risk.} Low drift risk: the item is framed as an explicit P vs NP transfer audit.

\paragraph{Branch control.} The item has high payoff and a concrete decision test; stopping now would be premature.

\paragraph{Premature-stop guard.} Do not stop before deriving the transfer inequality from \(\mathrm{SAT}\in \mathrm{SIZE}(n^{c})\) to the carrier upper bound, identifying whether \(\beta=0\), \(\beta>0\), or the argument is genuinely nonblackbox, and recording the barrier status of the proposed escape.

\paragraph{Payoff criteria.}
\begin{itemize}
\item For this route, state the exact barrier premise it uses or escapes: relativizing, natural, algebrizing, nonuniform, or advice-leaking; promotion requires the escape to be attached to the specific transfer theorem.
\item Compute \(\beta\) and q exactly; promotion requires \(\beta=0\), subpolynomial target queries, or a genuinely nonblackbox use of the SAT circuit.
\item State whether the transfer excludes every polynomial SAT exponent or only a bounded interval \(c<c0\).
\item Promotion requires a target-preserving nonblackbox/subpolynomial mechanism, not another black-box Cook-Levin substitution.
\item Record relativization, natural-proofs, algebrization, nonuniformity, and advice-leakage status for the specific transfer theorem.
\item Control gate: name the source carrier or language, target class, and circuit or resource model before claiming progress.
\end{itemize}

\paragraph{Stop criteria.}
\begin{itemize}
\item Every named carrier in the branch has \(\beta>0\) and yields only a bounded target-exponent exclusion.
\item The branch cannot state a source lower-bound model, target class, and transfer inequality in the same notation.
\item The proposed escape is only a black-box substitution restatement with no new carrier or nonblackbox mechanism.
\item The barrier audit shows the needed source lower bound is just as hard as the original P vs NP/P/poly target.
\end{itemize}

\paragraph{Do not stop if.}
\begin{itemize}
\item The branch has not yet computed \(\beta\), q, and the induced source upper-bound exponent under \(\mathrm{SAT}\in \mathrm{SIZE}(n^{c})\).
\item A subpolynomial-query or nonblackbox carrier remains available and has not been audited.
\item The latest failure sharpened the threshold-collapse lemma into a formal no-go statement.
\item The route can still isolate advice leakage, nonuniformity, naturalness, or algebrization as the exact obstruction.
\end{itemize}

\paragraph{Required skills.}
\begin{itemize}
\item abducer-failure-signal-cartographer
\item abducer-refinement-route-ranker
\item contradiction-route-auditor
\item entailment-gap-auditor
\item kb-deduce
\item kb-end-to-end-proof-critic
\item kb-gap-abstraction-conjecturer
\item kb-hypothesis-refutation-loop
\item kb-induce
\item kb-proof-connector
\item long-horizon-route-portfolio-manager
\item main-problem-alignment-auditor
\item next-experiment-queue-ranker
\item paper-contribution-readiness-judge
\item pvsnp-barrier-specificity-reviewer
\item pvsnp-conceptual-pressure-localizer
\item pvsnp-impossibility-theorem-forger
\item pvsnp-inference-quality-controller
\item pvsnp-proof-derivation-builder
\item pvsnp-qualitative-complexity-reasoner
\item pvsnp-quantifier-uniformity-reviewer
\item pvsnp-reduction-consistency-reviewer
\item pvsnp-symbolic-complexity-auditor
\end{itemize}

\subsection*{\#14 cycle-0086: Hardness-Magnification Bridge Realization}

\paragraph{Analysis.} Work the \(\mathrm{EF}+SATeq_{N}\) normal-form proof rule by rule: padding, renaming, duplicate clauses, extension-variable abbreviations, and SAT self-reduction; either translate each into a bounded EqSAT-link transcript or isolate the first noncompressible SAT-semantic principle as the next theorem-shaped counterexample.

\paragraph{Why this is a good target.} It directly tests the current black-box substitution cap by computing \(\beta\), q, and the induced target-exponent window. It audits a meta-complexity carrier where the decisive questions are verifier overhead, density/naturalness, and advice leakage. It checks whether padding/projection transfers really magnify hardness or merely move the proof debt into sparse lower bounds. It keeps the barrier and model-boundary audit attached to the mathematical transfer rather than as after-the-fact commentary. It is executable as a bounded complexity calculation, CAS-backed exponent audit, or formal transfer check.

\paragraph{Risk.} Low drift risk: the item is framed as an explicit P vs NP transfer audit.

\paragraph{Branch control.} The item has high payoff and a concrete decision test; stopping now would be premature.

\paragraph{Premature-stop guard.} Do not stop before deriving the transfer inequality from \(\mathrm{SAT}\in \mathrm{SIZE}(n^{c})\) to the carrier upper bound, identifying whether \(\beta=0\), \(\beta>0\), or the argument is genuinely nonblackbox, and recording the barrier status of the proposed escape.

\paragraph{Payoff criteria.}
\begin{itemize}
\item For the proof-complexity branch, separate fixed-family proof hardness from language hardness by naming the proof system, shortcut algorithm, and semantic lower-bound target.
\item Promotion requires a proof-system lower bound that survives the Gaussian-elimination or recognition shortcut; otherwise the result is a restricted-proof-system separation only.
\item For this route, state the exact barrier premise it uses or escapes: relativizing, natural, algebrizing, nonuniform, or advice-leaking; promotion requires the escape to be attached to the specific transfer theorem.
\item Compute \(\beta\) and q exactly; promotion requires \(\beta=0\), subpolynomial target queries, or a genuinely nonblackbox use of the SAT circuit.
\item For MCSP/MKtP, state N, n, s, verifier overhead, lower-bound model, and whether the property is natural or non-natural.
\item State whether the transfer excludes every polynomial SAT exponent or only a bounded interval \(c<c0\).
\end{itemize}

\paragraph{Stop criteria.}
\begin{itemize}
\item Every named carrier in the branch has \(\beta>0\) and yields only a bounded target-exponent exclusion.
\item The branch cannot state a source lower-bound model, target class, and transfer inequality in the same notation.
\item The proposed escape is only a black-box substitution restatement with no new carrier or nonblackbox mechanism.
\item The barrier audit shows the needed source lower bound is just as hard as the original P vs NP/P/poly target.
\end{itemize}

\paragraph{Do not stop if.}
\begin{itemize}
\item The branch has not yet computed \(\beta\), q, and the induced source upper-bound exponent under \(\mathrm{SAT}\in \mathrm{SIZE}(n^{c})\).
\item A subpolynomial-query or nonblackbox carrier remains available and has not been audited.
\item The latest failure sharpened the threshold-collapse lemma into a formal no-go statement.
\item The route can still isolate advice leakage, nonuniformity, naturalness, or algebrization as the exact obstruction.
\end{itemize}

\paragraph{Required skills.}
\begin{itemize}
\item abducer-failure-signal-cartographer
\item contradiction-route-auditor
\item counterexample-hunter
\item entailment-gap-auditor
\item kb-abduce
\item kb-concept-bridge-planner
\item kb-deduce
\item kb-end-to-end-proof-critic
\item kb-hypothesis-refutation-loop
\item kb-induce
\item long-horizon-route-portfolio-manager
\item main-problem-alignment-auditor
\item next-experiment-queue-ranker
\item paper-contribution-readiness-judge
\item pvsnp-barrier-specificity-reviewer
\item pvsnp-conceptual-pressure-localizer
\item pvsnp-impossibility-theorem-forger
\item pvsnp-inference-quality-controller
\item pvsnp-meta-complexity-carrier-reviewer
\item pvsnp-proof-derivation-builder
\item pvsnp-qualitative-complexity-reasoner
\item pvsnp-quantifier-uniformity-reviewer
\item pvsnp-reduction-consistency-reviewer
\item pvsnp-symbolic-complexity-auditor
\end{itemize}

\subsection*{\#15 cycle-0080: Hardness-Magnification Bridge Realization}

\paragraph{Analysis.} Attempt the full \(\mathrm{EF}+\mathrm{SAT}_{U}\) boundary lemma: inspect a standard encoding of \(\mathrm{SAT}_{U}\) extension axioms and decide whether each semantic rule is occurrence-local, charges enough advice into \(S_{Q}\), or yields a nonrelativizing complete-row extraction.

\paragraph{Why this is a good target.} It directly tests the current black-box substitution cap by computing \(\beta\), q, and the induced target-exponent window. It audits a meta-complexity carrier where the decisive questions are verifier overhead, density/naturalness, and advice leakage. It checks whether padding/projection transfers really magnify hardness or merely move the proof debt into sparse lower bounds. It keeps the barrier and model-boundary audit attached to the mathematical transfer rather than as after-the-fact commentary.

\paragraph{Risk.} Low drift risk: the item is framed as an explicit P vs NP transfer audit.

\paragraph{Branch control.} The item has high payoff and a concrete decision test; stopping now would be premature.

\paragraph{Premature-stop guard.} Do not stop before deriving the transfer inequality from \(\mathrm{SAT}\in \mathrm{SIZE}(n^{c})\) to the carrier upper bound, identifying whether \(\beta=0\), \(\beta>0\), or the argument is genuinely nonblackbox, and recording the barrier status of the proposed escape.

\paragraph{Payoff criteria.}
\begin{itemize}
\item For the proof-complexity branch, separate fixed-family proof hardness from language hardness by naming the proof system, shortcut algorithm, and semantic lower-bound target.
\item Promotion requires a proof-system lower bound that survives the Gaussian-elimination or recognition shortcut; otherwise the result is a restricted-proof-system separation only.
\item For this route, state the exact barrier premise it uses or escapes: relativizing, natural, algebrizing, nonuniform, or advice-leaking; promotion requires the escape to be attached to the specific transfer theorem.
\item Compute \(\beta\) and q exactly; promotion requires \(\beta=0\), subpolynomial target queries, or a genuinely nonblackbox use of the SAT circuit.
\item For MCSP/MKtP, state N, n, s, verifier overhead, lower-bound model, and whether the property is natural or non-natural.
\item State whether the transfer excludes every polynomial SAT exponent or only a bounded interval \(c<c0\).
\end{itemize}

\paragraph{Stop criteria.}
\begin{itemize}
\item Every named carrier in the branch has \(\beta>0\) and yields only a bounded target-exponent exclusion.
\item The branch cannot state a source lower-bound model, target class, and transfer inequality in the same notation.
\item The proposed escape is only a black-box substitution restatement with no new carrier or nonblackbox mechanism.
\item The barrier audit shows the needed source lower bound is just as hard as the original P vs NP/P/poly target.
\end{itemize}

\paragraph{Do not stop if.}
\begin{itemize}
\item The branch has not yet computed \(\beta\), q, and the induced source upper-bound exponent under \(\mathrm{SAT}\in \mathrm{SIZE}(n^{c})\).
\item A subpolynomial-query or nonblackbox carrier remains available and has not been audited.
\item The latest failure sharpened the threshold-collapse lemma into a formal no-go statement.
\item The route can still isolate advice leakage, nonuniformity, naturalness, or algebrization as the exact obstruction.
\end{itemize}

\paragraph{Required skills.}
\begin{itemize}
\item abducer-failure-signal-cartographer
\item contradiction-route-auditor
\item counterexample-hunter
\item entailment-gap-auditor
\item kb-abduce
\item kb-concept-bridge-planner
\item kb-deduce
\item kb-end-to-end-proof-critic
\item kb-hypothesis-refutation-loop
\item kb-induce
\item long-horizon-route-portfolio-manager
\item main-problem-alignment-auditor
\item next-experiment-queue-ranker
\item paper-contribution-readiness-judge
\item pvsnp-barrier-specificity-reviewer
\item pvsnp-conceptual-pressure-localizer
\item pvsnp-impossibility-theorem-forger
\item pvsnp-inference-quality-controller
\item pvsnp-meta-complexity-carrier-reviewer
\item pvsnp-proof-derivation-builder
\item pvsnp-qualitative-complexity-reasoner
\item pvsnp-quantifier-uniformity-reviewer
\item pvsnp-reduction-consistency-reviewer
\item pvsnp-symbolic-complexity-auditor
\end{itemize}

\subsection*{\#16 cycle-0071: Hardness-Magnification Bridge Realization}

\paragraph{Analysis.} Instantiate \(V_{N}\) with one concrete proof-system or MCSP-style certificate, compute \(kappa_{N}\) for accepted labels, and decide whether soundness is functional over arbitrary y or essentially SAT-semantic.

\paragraph{Why this is a good target.} It directly tests the current black-box substitution cap by computing \(\beta\), q, and the induced target-exponent window. It audits a meta-complexity carrier where the decisive questions are verifier overhead, density/naturalness, and advice leakage. It checks whether padding/projection transfers really magnify hardness or merely move the proof debt into sparse lower bounds. It keeps the barrier and model-boundary audit attached to the mathematical transfer rather than as after-the-fact commentary. It is executable as a bounded complexity calculation, CAS-backed exponent audit, or formal transfer check.

\paragraph{Risk.} Low drift risk: the item is framed as an explicit P vs NP transfer audit.

\paragraph{Branch control.} The item has high payoff and a concrete decision test; stopping now would be premature.

\paragraph{Premature-stop guard.} Do not stop before deriving the transfer inequality from \(\mathrm{SAT}\in \mathrm{SIZE}(n^{c})\) to the carrier upper bound, identifying whether \(\beta=0\), \(\beta>0\), or the argument is genuinely nonblackbox, and recording the barrier status of the proposed escape.

\paragraph{Payoff criteria.}
\begin{itemize}
\item For the MCSP/MKtP row branch, decide whether the certificate proves direct row hardness or only aggregate/rank-index hardness; record q\_extract=epsilon/(1+D) and q\_force=(D+epsilon)/(1+D) in the same parameter regime.
\item Promotion requires an advice-free row-hard certificate; backtrack if the ranked decoder, canonical unranking, or aggregate proof selects a hard row nonuniformly.
\item For the proof-complexity branch, separate fixed-family proof hardness from language hardness by naming the proof system, shortcut algorithm, and semantic lower-bound target.
\item Promotion requires a proof-system lower bound that survives the Gaussian-elimination or recognition shortcut; otherwise the result is a restricted-proof-system separation only.
\item For this route, state the exact barrier premise it uses or escapes: relativizing, natural, algebrizing, nonuniform, or advice-leaking; promotion requires the escape to be attached to the specific transfer theorem.
\item Compute \(\beta\) and q exactly; promotion requires \(\beta=0\), subpolynomial target queries, or a genuinely nonblackbox use of the SAT circuit.
\end{itemize}

\paragraph{Stop criteria.}
\begin{itemize}
\item Every named carrier in the branch has \(\beta>0\) and yields only a bounded target-exponent exclusion.
\item The branch cannot state a source lower-bound model, target class, and transfer inequality in the same notation.
\item The proposed escape is only a black-box substitution restatement with no new carrier or nonblackbox mechanism.
\item The barrier audit shows the needed source lower bound is just as hard as the original P vs NP/P/poly target.
\end{itemize}

\paragraph{Do not stop if.}
\begin{itemize}
\item The branch has not yet computed \(\beta\), q, and the induced source upper-bound exponent under \(\mathrm{SAT}\in \mathrm{SIZE}(n^{c})\).
\item A subpolynomial-query or nonblackbox carrier remains available and has not been audited.
\item The latest failure sharpened the threshold-collapse lemma into a formal no-go statement.
\item The route can still isolate advice leakage, nonuniformity, naturalness, or algebrization as the exact obstruction.
\end{itemize}

\paragraph{Required skills.}
\begin{itemize}
\item abducer-failure-signal-cartographer
\item contradiction-route-auditor
\item counterexample-hunter
\item entailment-gap-auditor
\item kb-abduce
\item kb-concept-bridge-planner
\item kb-deduce
\item kb-end-to-end-proof-critic
\item kb-hypothesis-refutation-loop
\item kb-induce
\item long-horizon-route-portfolio-manager
\item main-problem-alignment-auditor
\item next-experiment-queue-ranker
\item paper-contribution-readiness-judge
\item pvsnp-barrier-specificity-reviewer
\item pvsnp-conceptual-pressure-localizer
\item pvsnp-impossibility-theorem-forger
\item pvsnp-inference-quality-controller
\item pvsnp-meta-complexity-carrier-reviewer
\item pvsnp-proof-derivation-builder
\item pvsnp-qualitative-complexity-reasoner
\item pvsnp-quantifier-uniformity-reviewer
\item pvsnp-reduction-consistency-reviewer
\item pvsnp-symbolic-complexity-auditor
\end{itemize}

\subsection*{\#17 cycle-0034: Hardness-Magnification Bridge Realization: tuple-feature residual gate}

\paragraph{Analysis.} Instantiate the pair-pattern thought experiment for \(t=2\): choose q with \(q<1+\epsilon\) and test whether \(2q+\alpha<1+\epsilon\) permits a subthreshold pair separator \(S_{N}\); compute the residual complexity C(H\_N xor \(S_{N}\)) and check for hidden SAT lookup.

\paragraph{Why this is a good target.} It directly tests the current black-box substitution cap by computing \(\beta\), q, and the induced target-exponent window. It audits a meta-complexity carrier where the decisive questions are verifier overhead, density/naturalness, and advice leakage. It checks whether padding/projection transfers really magnify hardness or merely move the proof debt into sparse lower bounds. It keeps the barrier and model-boundary audit attached to the mathematical transfer rather than as after-the-fact commentary. It is executable as a bounded complexity calculation, CAS-backed exponent audit, or formal transfer check.

\paragraph{Risk.} Low drift risk: the item is framed as an explicit P vs NP transfer audit.

\paragraph{Branch control.} The item has high payoff and a concrete decision test; stopping now would be premature.

\paragraph{Premature-stop guard.} Do not stop before deriving the transfer inequality from \(\mathrm{SAT}\in \mathrm{SIZE}(n^{c})\) to the carrier upper bound, identifying whether \(\beta=0\), \(\beta>0\), or the argument is genuinely nonblackbox, and recording the barrier status of the proposed escape.

\paragraph{Payoff criteria.}
\begin{itemize}
\item For the high-q generated-OR carrier, compute the exact leakage exponent of the remaining witness-cover or unary-residual mechanism and compare it with 1+epsilon-q.
\item Promotion requires \(q+a>1+\epsilon\) after every projection, padding, and witness-cover leakage term is accounted for; otherwise the branch becomes a bounded-window or no-go result.
\item For the proof-complexity branch, separate fixed-family proof hardness from language hardness by naming the proof system, shortcut algorithm, and semantic lower-bound target.
\item Promotion requires a proof-system lower bound that survives the Gaussian-elimination or recognition shortcut; otherwise the result is a restricted-proof-system separation only.
\item For this route, state the exact barrier premise it uses or escapes: relativizing, natural, algebrizing, nonuniform, or advice-leaking; promotion requires the escape to be attached to the specific transfer theorem.
\item Compute \(\beta\) and q exactly; promotion requires \(\beta=0\), subpolynomial target queries, or a genuinely nonblackbox use of the SAT circuit.
\end{itemize}

\paragraph{Stop criteria.}
\begin{itemize}
\item Every named carrier in the branch has \(\beta>0\) and yields only a bounded target-exponent exclusion.
\item The branch cannot state a source lower-bound model, target class, and transfer inequality in the same notation.
\item The proposed escape is only a black-box substitution restatement with no new carrier or nonblackbox mechanism.
\item The barrier audit shows the needed source lower bound is just as hard as the original P vs NP/P/poly target.
\end{itemize}

\paragraph{Do not stop if.}
\begin{itemize}
\item The branch has not yet computed \(\beta\), q, and the induced source upper-bound exponent under \(\mathrm{SAT}\in \mathrm{SIZE}(n^{c})\).
\item A subpolynomial-query or nonblackbox carrier remains available and has not been audited.
\item The latest failure sharpened the threshold-collapse lemma into a formal no-go statement.
\item The route can still isolate advice leakage, nonuniformity, naturalness, or algebrization as the exact obstruction.
\end{itemize}

\paragraph{Required skills.}
\begin{itemize}
\item abducer-failure-signal-cartographer
\item contradiction-route-auditor
\item counterexample-hunter
\item entailment-gap-auditor
\item kb-abduce
\item kb-concept-bridge-planner
\item kb-deduce
\item kb-end-to-end-proof-critic
\item kb-hypothesis-refutation-loop
\item kb-induce
\item long-horizon-route-portfolio-manager
\item main-problem-alignment-auditor
\item next-experiment-queue-ranker
\item paper-contribution-readiness-judge
\item pvsnp-barrier-specificity-reviewer
\item pvsnp-conceptual-pressure-localizer
\item pvsnp-impossibility-theorem-forger
\item pvsnp-inference-quality-controller
\item pvsnp-meta-complexity-carrier-reviewer
\item pvsnp-proof-derivation-builder
\item pvsnp-qualitative-complexity-reasoner
\item pvsnp-quantifier-uniformity-reviewer
\item pvsnp-reduction-consistency-reviewer
\item pvsnp-symbolic-complexity-auditor
\end{itemize}

\subsection*{\#18 cycle-0013: Hardness-Magnification Bridge Realization: Polylog Nonblackbox Compact Consumer}

\paragraph{Analysis.} Instantiate \(PolyLogAudit_{N,d,V}\) with a universal-circuit verifier, compute the exact overhead exponent \(\rho\) and auxiliary advice length, and reject the instance immediately if \(H_{N}\) depends on a hidden SAT circuit or collapses to a black-box SAT query family.

\paragraph{Why this is a good target.} It directly tests the current black-box substitution cap by computing \(\beta\), q, and the induced target-exponent window. It audits a meta-complexity carrier where the decisive questions are verifier overhead, density/naturalness, and advice leakage. It checks whether padding/projection transfers really magnify hardness or merely move the proof debt into sparse lower bounds. It points toward the only unquarantined escape: a subpolynomial-query, nonprojection, or nonblackbox carrier. It keeps the barrier and model-boundary audit attached to the mathematical transfer rather than as after-the-fact commentary. It is executable as a bounded complexity calculation, CAS-backed exponent audit, or formal transfer check.

\paragraph{Risk.} Low drift risk: the item is framed as an explicit P vs NP transfer audit.

\paragraph{Branch control.} The item has high payoff and a concrete decision test; stopping now would be premature.

\paragraph{Premature-stop guard.} Do not stop before deriving the transfer inequality from \(\mathrm{SAT}\in \mathrm{SIZE}(n^{c})\) to the carrier upper bound, identifying whether \(\beta=0\), \(\beta>0\), or the argument is genuinely nonblackbox, and recording the barrier status of the proposed escape.

\paragraph{Payoff criteria.}
\begin{itemize}
\item For this route, state the exact barrier premise it uses or escapes: relativizing, natural, algebrizing, nonuniform, or advice-leaking; promotion requires the escape to be attached to the specific transfer theorem.
\item Compute \(\beta\) and q exactly; promotion requires \(\beta=0\), subpolynomial target queries, or a genuinely nonblackbox use of the SAT circuit.
\item For MCSP/MKtP, state N, n, s, verifier overhead, lower-bound model, and whether the property is natural or non-natural.
\item State whether the transfer excludes every polynomial SAT exponent or only a bounded interval \(c<c0\).
\item Promotion requires a target-preserving nonblackbox/subpolynomial mechanism, not another black-box Cook-Levin substitution.
\item Record relativization, natural-proofs, algebrization, nonuniformity, and advice-leakage status for the specific transfer theorem.
\end{itemize}

\paragraph{Stop criteria.}
\begin{itemize}
\item Every named carrier in the branch has \(\beta>0\) and yields only a bounded target-exponent exclusion.
\item The branch cannot state a source lower-bound model, target class, and transfer inequality in the same notation.
\item The proposed escape is only a black-box substitution restatement with no new carrier or nonblackbox mechanism.
\item The barrier audit shows the needed source lower bound is just as hard as the original P vs NP/P/poly target.
\end{itemize}

\paragraph{Do not stop if.}
\begin{itemize}
\item The branch has not yet computed \(\beta\), q, and the induced source upper-bound exponent under \(\mathrm{SAT}\in \mathrm{SIZE}(n^{c})\).
\item A subpolynomial-query or nonblackbox carrier remains available and has not been audited.
\item The latest failure sharpened the threshold-collapse lemma into a formal no-go statement.
\item The route can still isolate advice leakage, nonuniformity, naturalness, or algebrization as the exact obstruction.
\end{itemize}

\paragraph{Required skills.}
\begin{itemize}
\item abducer-failure-signal-cartographer
\item abducer-refinement-route-ranker
\item contradiction-route-auditor
\item counterexample-hunter
\item entailment-gap-auditor
\item kb-abduce
\item kb-concept-bridge-planner
\item kb-deduce
\item kb-end-to-end-proof-critic
\item kb-gap-abstraction-conjecturer
\item kb-hypothesis-refutation-loop
\item kb-induce
\item kb-proof-connector
\item long-horizon-route-portfolio-manager
\item main-problem-alignment-auditor
\item next-experiment-queue-ranker
\item paper-contribution-readiness-judge
\item pvsnp-barrier-specificity-reviewer
\item pvsnp-conceptual-pressure-localizer
\item pvsnp-impossibility-theorem-forger
\item pvsnp-inference-quality-controller
\item pvsnp-meta-complexity-carrier-reviewer
\item pvsnp-proof-derivation-builder
\item pvsnp-qualitative-complexity-reasoner
\item pvsnp-quantifier-uniformity-reviewer
\item pvsnp-reduction-consistency-reviewer
\item pvsnp-symbolic-complexity-auditor
\end{itemize}

\subsection*{\#19 cycle-0012: Hardness-Magnification Bridge Realization}

\paragraph{Analysis.} Instantiate \(H_{N}\) as a proof-checking or self-reduction transcript language whose verifier can call a compact \(\mathrm{SAT}_{L}\) circuit nonblackbox; compute its exact overhead a and test whether full truth-table substitution or input-encoded D leakage gives a subthreshold upper bound.

\paragraph{Why this is a good target.} It directly tests the current black-box substitution cap by computing \(\beta\), q, and the induced target-exponent window. It audits a meta-complexity carrier where the decisive questions are verifier overhead, density/naturalness, and advice leakage. It points toward the only unquarantined escape: a subpolynomial-query, nonprojection, or nonblackbox carrier. It keeps the barrier and model-boundary audit attached to the mathematical transfer rather than as after-the-fact commentary. It is executable as a bounded complexity calculation, CAS-backed exponent audit, or formal transfer check.

\paragraph{Risk.} Low drift risk: the item is framed as an explicit P vs NP transfer audit.

\paragraph{Branch control.} The item has high payoff and a concrete decision test; stopping now would be premature.

\paragraph{Premature-stop guard.} Do not stop before deriving the transfer inequality from \(\mathrm{SAT}\in \mathrm{SIZE}(n^{c})\) to the carrier upper bound, identifying whether \(\beta=0\), \(\beta>0\), or the argument is genuinely nonblackbox, and recording the barrier status of the proposed escape.

\paragraph{Payoff criteria.}
\begin{itemize}
\item For the restricted-model branch, compute the compiler envelope from unrestricted SIZE(m\textasciicircum{}k) or P-time algorithms into the named model; promotion requires a bound above that envelope on the same SAT slice.
\item Falsify direct transfer if a P function already has the same restricted-model lower bound, such as PARITY for \(\mathrm{AC0}\), HWB for OBDD, or XOR/Tseitin for proof-system degree.
\item For this route, state the exact barrier premise it uses or escapes: relativizing, natural, algebrizing, nonuniform, or advice-leaking; promotion requires the escape to be attached to the specific transfer theorem.
\item Compute \(\beta\) and q exactly; promotion requires \(\beta=0\), subpolynomial target queries, or a genuinely nonblackbox use of the SAT circuit.
\item For MCSP/MKtP, state N, n, s, verifier overhead, lower-bound model, and whether the property is natural or non-natural.
\item Promotion requires a target-preserving nonblackbox/subpolynomial mechanism, not another black-box Cook-Levin substitution.
\end{itemize}

\paragraph{Stop criteria.}
\begin{itemize}
\item Every named carrier in the branch has \(\beta>0\) and yields only a bounded target-exponent exclusion.
\item The branch cannot state a source lower-bound model, target class, and transfer inequality in the same notation.
\item The proposed escape is only a black-box substitution restatement with no new carrier or nonblackbox mechanism.
\item The barrier audit shows the needed source lower bound is just as hard as the original P vs NP/P/poly target.
\end{itemize}

\paragraph{Do not stop if.}
\begin{itemize}
\item The branch has not yet computed \(\beta\), q, and the induced source upper-bound exponent under \(\mathrm{SAT}\in \mathrm{SIZE}(n^{c})\).
\item A subpolynomial-query or nonblackbox carrier remains available and has not been audited.
\item The latest failure sharpened the threshold-collapse lemma into a formal no-go statement.
\item The route can still isolate advice leakage, nonuniformity, naturalness, or algebrization as the exact obstruction.
\end{itemize}

\paragraph{Required skills.}
\begin{itemize}
\item abducer-refinement-route-ranker
\item contradiction-route-auditor
\item counterexample-hunter
\item entailment-gap-auditor
\item kb-abduce
\item kb-concept-bridge-planner
\item kb-deduce
\item kb-end-to-end-proof-critic
\item kb-gap-abstraction-conjecturer
\item kb-hypothesis-refutation-loop
\item kb-proof-connector
\item long-horizon-route-portfolio-manager
\item main-problem-alignment-auditor
\item next-experiment-queue-ranker
\item paper-contribution-readiness-judge
\item pvsnp-barrier-specificity-reviewer
\item pvsnp-conceptual-pressure-localizer
\item pvsnp-impossibility-theorem-forger
\item pvsnp-inference-quality-controller
\item pvsnp-meta-complexity-carrier-reviewer
\item pvsnp-proof-derivation-builder
\item pvsnp-qualitative-complexity-reasoner
\item pvsnp-quantifier-uniformity-reviewer
\item pvsnp-reduction-consistency-reviewer
\item pvsnp-symbolic-complexity-auditor
\end{itemize}

\subsection*{\#20 cycle-0009: Hardness-Magnification Bridge Realization: subpolynomial-interface demagnification}

\paragraph{Analysis.} Pivot to a nonblackbox object: test a proof-complexity or derandomization carrier where SAT in P/poly yields the required source upper bound through a named nonrelativizing theorem rather than through black-box SAT-query substitution.

\paragraph{Why this is a good target.} It directly tests the current black-box substitution cap by computing \(\beta\), q, and the induced target-exponent window. It audits a meta-complexity carrier where the decisive questions are verifier overhead, density/naturalness, and advice leakage. It points toward the only unquarantined escape: a subpolynomial-query, nonprojection, or nonblackbox carrier. It keeps the barrier and model-boundary audit attached to the mathematical transfer rather than as after-the-fact commentary. It is executable as a bounded complexity calculation, CAS-backed exponent audit, or formal transfer check.

\paragraph{Risk.} Low drift risk: the item is framed as an explicit P vs NP transfer audit.

\paragraph{Branch control.} The item has high payoff and a concrete decision test; stopping now would be premature.

\paragraph{Premature-stop guard.} Do not stop before deriving the transfer inequality from \(\mathrm{SAT}\in \mathrm{SIZE}(n^{c})\) to the carrier upper bound, identifying whether \(\beta=0\), \(\beta>0\), or the argument is genuinely nonblackbox, and recording the barrier status of the proposed escape.

\paragraph{Payoff criteria.}
\begin{itemize}
\item For the restricted-model branch, compute the compiler envelope from unrestricted SIZE(m\textasciicircum{}k) or P-time algorithms into the named model; promotion requires a bound above that envelope on the same SAT slice.
\item Falsify direct transfer if a P function already has the same restricted-model lower bound, such as PARITY for \(\mathrm{AC0}\), HWB for OBDD, or XOR/Tseitin for proof-system degree.
\item For this route, state the exact barrier premise it uses or escapes: relativizing, natural, algebrizing, nonuniform, or advice-leaking; promotion requires the escape to be attached to the specific transfer theorem.
\item Compute \(\beta\) and q exactly; promotion requires \(\beta=0\), subpolynomial target queries, or a genuinely nonblackbox use of the SAT circuit.
\item For MCSP/MKtP, state N, n, s, verifier overhead, lower-bound model, and whether the property is natural or non-natural.
\item Promotion requires a target-preserving nonblackbox/subpolynomial mechanism, not another black-box Cook-Levin substitution.
\end{itemize}

\paragraph{Stop criteria.}
\begin{itemize}
\item Every named carrier in the branch has \(\beta>0\) and yields only a bounded target-exponent exclusion.
\item The branch cannot state a source lower-bound model, target class, and transfer inequality in the same notation.
\item The proposed escape is only a black-box substitution restatement with no new carrier or nonblackbox mechanism.
\item The barrier audit shows the needed source lower bound is just as hard as the original P vs NP/P/poly target.
\end{itemize}

\paragraph{Do not stop if.}
\begin{itemize}
\item The branch has not yet computed \(\beta\), q, and the induced source upper-bound exponent under \(\mathrm{SAT}\in \mathrm{SIZE}(n^{c})\).
\item A subpolynomial-query or nonblackbox carrier remains available and has not been audited.
\item The latest failure sharpened the threshold-collapse lemma into a formal no-go statement.
\item The route can still isolate advice leakage, nonuniformity, naturalness, or algebrization as the exact obstruction.
\end{itemize}

\paragraph{Required skills.}
\begin{itemize}
\item abducer-refinement-route-ranker
\item contradiction-route-auditor
\item counterexample-hunter
\item entailment-gap-auditor
\item kb-abduce
\item kb-concept-bridge-planner
\item kb-deduce
\item kb-end-to-end-proof-critic
\item kb-gap-abstraction-conjecturer
\item kb-hypothesis-refutation-loop
\item kb-proof-connector
\item long-horizon-route-portfolio-manager
\item main-problem-alignment-auditor
\item next-experiment-queue-ranker
\item paper-contribution-readiness-judge
\item pvsnp-barrier-specificity-reviewer
\item pvsnp-conceptual-pressure-localizer
\item pvsnp-impossibility-theorem-forger
\item pvsnp-inference-quality-controller
\item pvsnp-meta-complexity-carrier-reviewer
\item pvsnp-proof-derivation-builder
\item pvsnp-qualitative-complexity-reasoner
\item pvsnp-quantifier-uniformity-reviewer
\item pvsnp-reduction-consistency-reviewer
\item pvsnp-symbolic-complexity-auditor
\end{itemize}

\subsection*{\#21 cycle-0008: Hardness-Magnification Bridge Realization, cycle 8: subpolynomial-query threshold repair}

\paragraph{Analysis.} Search for non-truth-table carriers in \(\mathrm{P}^\mathrm{SAT}\) with near-linear ordinary work and subpolynomial SAT queries, e.g. succinct consistency or decompression-check problems, then compute A,Q,s and test against L8.1 before discussing P vs NP.

\paragraph{Why this is a good target.} It directly tests the current black-box substitution cap by computing \(\beta\), q, and the induced target-exponent window. It audits a meta-complexity carrier where the decisive questions are verifier overhead, density/naturalness, and advice leakage. It points toward the only unquarantined escape: a subpolynomial-query, nonprojection, or nonblackbox carrier. It keeps the barrier and model-boundary audit attached to the mathematical transfer rather than as after-the-fact commentary. It is executable as a bounded complexity calculation, CAS-backed exponent audit, or formal transfer check.

\paragraph{Risk.} Low drift risk: the item is framed as an explicit P vs NP transfer audit.

\paragraph{Branch control.} The item has high payoff and a concrete decision test; stopping now would be premature.

\paragraph{Premature-stop guard.} Do not stop before deriving the transfer inequality from \(\mathrm{SAT}\in \mathrm{SIZE}(n^{c})\) to the carrier upper bound, identifying whether \(\beta=0\), \(\beta>0\), or the argument is genuinely nonblackbox, and recording the barrier status of the proposed escape.

\paragraph{Payoff criteria.}
\begin{itemize}
\item For the MCSP/MKtP row branch, decide whether the certificate proves direct row hardness or only aggregate/rank-index hardness; record q\_extract=epsilon/(1+D) and q\_force=(D+epsilon)/(1+D) in the same parameter regime.
\item Promotion requires an advice-free row-hard certificate; backtrack if the ranked decoder, canonical unranking, or aggregate proof selects a hard row nonuniformly.
\item For this route, state the exact barrier premise it uses or escapes: relativizing, natural, algebrizing, nonuniform, or advice-leaking; promotion requires the escape to be attached to the specific transfer theorem.
\item Compute \(\beta\) and q exactly; promotion requires \(\beta=0\), subpolynomial target queries, or a genuinely nonblackbox use of the SAT circuit.
\item For MCSP/MKtP, state N, n, s, verifier overhead, lower-bound model, and whether the property is natural or non-natural.
\item Promotion requires a target-preserving nonblackbox/subpolynomial mechanism, not another black-box Cook-Levin substitution.
\end{itemize}

\paragraph{Stop criteria.}
\begin{itemize}
\item Every named carrier in the branch has \(\beta>0\) and yields only a bounded target-exponent exclusion.
\item The branch cannot state a source lower-bound model, target class, and transfer inequality in the same notation.
\item The proposed escape is only a black-box substitution restatement with no new carrier or nonblackbox mechanism.
\item The barrier audit shows the needed source lower bound is just as hard as the original P vs NP/P/poly target.
\end{itemize}

\paragraph{Do not stop if.}
\begin{itemize}
\item The branch has not yet computed \(\beta\), q, and the induced source upper-bound exponent under \(\mathrm{SAT}\in \mathrm{SIZE}(n^{c})\).
\item A subpolynomial-query or nonblackbox carrier remains available and has not been audited.
\item The latest failure sharpened the threshold-collapse lemma into a formal no-go statement.
\item The route can still isolate advice leakage, nonuniformity, naturalness, or algebrization as the exact obstruction.
\end{itemize}

\paragraph{Required skills.}
\begin{itemize}
\item abducer-refinement-route-ranker
\item contradiction-route-auditor
\item counterexample-hunter
\item entailment-gap-auditor
\item kb-abduce
\item kb-concept-bridge-planner
\item kb-deduce
\item kb-end-to-end-proof-critic
\item kb-gap-abstraction-conjecturer
\item kb-hypothesis-refutation-loop
\item kb-proof-connector
\item long-horizon-route-portfolio-manager
\item main-problem-alignment-auditor
\item next-experiment-queue-ranker
\item paper-contribution-readiness-judge
\item pvsnp-barrier-specificity-reviewer
\item pvsnp-conceptual-pressure-localizer
\item pvsnp-impossibility-theorem-forger
\item pvsnp-inference-quality-controller
\item pvsnp-meta-complexity-carrier-reviewer
\item pvsnp-proof-derivation-builder
\item pvsnp-qualitative-complexity-reasoner
\item pvsnp-quantifier-uniformity-reviewer
\item pvsnp-reduction-consistency-reviewer
\item pvsnp-symbolic-complexity-auditor
\end{itemize}

\subsection*{\#22 cycle-0007: Hardness-Magnification Bridge Realization: black-box substitution cap}

\paragraph{Analysis.} Test a succinct local-consistency or proof-complexity carrier with \(N^{o(1)}\) SAT queries: compute a,q,beta, then run the projection-demagnification falsifier before claiming any magnification.

\paragraph{Why this is a good target.} It directly tests the current black-box substitution cap by computing \(\beta\), q, and the induced target-exponent window. It audits a meta-complexity carrier where the decisive questions are verifier overhead, density/naturalness, and advice leakage. It checks whether padding/projection transfers really magnify hardness or merely move the proof debt into sparse lower bounds. It points toward the only unquarantined escape: a subpolynomial-query, nonprojection, or nonblackbox carrier. It keeps the barrier and model-boundary audit attached to the mathematical transfer rather than as after-the-fact commentary. It is executable as a bounded complexity calculation, CAS-backed exponent audit, or formal transfer check.

\paragraph{Risk.} Low drift risk: the item is framed as an explicit P vs NP transfer audit.

\paragraph{Branch control.} The item has high payoff and a concrete decision test; stopping now would be premature.

\paragraph{Premature-stop guard.} Do not stop before deriving the transfer inequality from \(\mathrm{SAT}\in \mathrm{SIZE}(n^{c})\) to the carrier upper bound, identifying whether \(\beta=0\), \(\beta>0\), or the argument is genuinely nonblackbox, and recording the barrier status of the proposed escape.

\paragraph{Payoff criteria.}
\begin{itemize}
\item For the MCSP/MKtP row branch, decide whether the certificate proves direct row hardness or only aggregate/rank-index hardness; record q\_extract=epsilon/(1+D) and q\_force=(D+epsilon)/(1+D) in the same parameter regime.
\item Promotion requires an advice-free row-hard certificate; backtrack if the ranked decoder, canonical unranking, or aggregate proof selects a hard row nonuniformly.
\item For this route, state the exact barrier premise it uses or escapes: relativizing, natural, algebrizing, nonuniform, or advice-leaking; promotion requires the escape to be attached to the specific transfer theorem.
\item Compute \(\beta\) and q exactly; promotion requires \(\beta=0\), subpolynomial target queries, or a genuinely nonblackbox use of the SAT circuit.
\item For MCSP/MKtP, state N, n, s, verifier overhead, lower-bound model, and whether the property is natural or non-natural.
\item State whether the transfer excludes every polynomial SAT exponent or only a bounded interval \(c<c0\).
\end{itemize}

\paragraph{Stop criteria.}
\begin{itemize}
\item Every named carrier in the branch has \(\beta>0\) and yields only a bounded target-exponent exclusion.
\item The branch cannot state a source lower-bound model, target class, and transfer inequality in the same notation.
\item The proposed escape is only a black-box substitution restatement with no new carrier or nonblackbox mechanism.
\item The barrier audit shows the needed source lower bound is just as hard as the original P vs NP/P/poly target.
\end{itemize}

\paragraph{Do not stop if.}
\begin{itemize}
\item The branch has not yet computed \(\beta\), q, and the induced source upper-bound exponent under \(\mathrm{SAT}\in \mathrm{SIZE}(n^{c})\).
\item A subpolynomial-query or nonblackbox carrier remains available and has not been audited.
\item The latest failure sharpened the threshold-collapse lemma into a formal no-go statement.
\item The route can still isolate advice leakage, nonuniformity, naturalness, or algebrization as the exact obstruction.
\end{itemize}

\paragraph{Required skills.}
\begin{itemize}
\item abducer-failure-signal-cartographer
\item abducer-refinement-route-ranker
\item contradiction-route-auditor
\item counterexample-hunter
\item entailment-gap-auditor
\item kb-abduce
\item kb-concept-bridge-planner
\item kb-deduce
\item kb-end-to-end-proof-critic
\item kb-gap-abstraction-conjecturer
\item kb-hypothesis-refutation-loop
\item kb-induce
\item kb-proof-connector
\item long-horizon-route-portfolio-manager
\item main-problem-alignment-auditor
\item next-experiment-queue-ranker
\item paper-contribution-readiness-judge
\item pvsnp-barrier-specificity-reviewer
\item pvsnp-conceptual-pressure-localizer
\item pvsnp-impossibility-theorem-forger
\item pvsnp-inference-quality-controller
\item pvsnp-meta-complexity-carrier-reviewer
\item pvsnp-proof-derivation-builder
\item pvsnp-qualitative-complexity-reasoner
\item pvsnp-quantifier-uniformity-reviewer
\item pvsnp-reduction-consistency-reviewer
\item pvsnp-symbolic-complexity-auditor
\end{itemize}

\subsection*{\#23 cycle-0006: Hardness-Magnification Bridge Realization, Cycle 6: MCSP verifier transfer audit}

\paragraph{Analysis.} Test a self-referential meta-complexity carrier, such as Gap-MCSP or MKtP with succinctly generated truth tables, where the target SAT circuit is used nonblack-box. Compute whether the induced carrier upper bound can avoid the multiplicative c factor; if not, quarantine it by the same beta-ceiling lemma.

\paragraph{Why this is a good target.} It directly tests the current black-box substitution cap by computing \(\beta\), q, and the induced target-exponent window. It audits a meta-complexity carrier where the decisive questions are verifier overhead, density/naturalness, and advice leakage. It checks whether padding/projection transfers really magnify hardness or merely move the proof debt into sparse lower bounds. It points toward the only unquarantined escape: a subpolynomial-query, nonprojection, or nonblackbox carrier. It keeps the barrier and model-boundary audit attached to the mathematical transfer rather than as after-the-fact commentary. It is executable as a bounded complexity calculation, CAS-backed exponent audit, or formal transfer check.

\paragraph{Risk.} Low drift risk: the item is framed as an explicit P vs NP transfer audit.

\paragraph{Branch control.} The item has high payoff and a concrete decision test; stopping now would be premature.

\paragraph{Premature-stop guard.} Do not stop before deriving the transfer inequality from \(\mathrm{SAT}\in \mathrm{SIZE}(n^{c})\) to the carrier upper bound, identifying whether \(\beta=0\), \(\beta>0\), or the argument is genuinely nonblackbox, and recording the barrier status of the proposed escape.

\paragraph{Payoff criteria.}
\begin{itemize}
\item For the MCSP/MKtP row branch, decide whether the certificate proves direct row hardness or only aggregate/rank-index hardness; record q\_extract=epsilon/(1+D) and q\_force=(D+epsilon)/(1+D) in the same parameter regime.
\item Promotion requires an advice-free row-hard certificate; backtrack if the ranked decoder, canonical unranking, or aggregate proof selects a hard row nonuniformly.
\item For this route, state the exact barrier premise it uses or escapes: relativizing, natural, algebrizing, nonuniform, or advice-leaking; promotion requires the escape to be attached to the specific transfer theorem.
\item Compute \(\beta\) and q exactly; promotion requires \(\beta=0\), subpolynomial target queries, or a genuinely nonblackbox use of the SAT circuit.
\item For MCSP/MKtP, state N, n, s, verifier overhead, lower-bound model, and whether the property is natural or non-natural.
\item State whether the transfer excludes every polynomial SAT exponent or only a bounded interval \(c<c0\).
\end{itemize}

\paragraph{Stop criteria.}
\begin{itemize}
\item Every named carrier in the branch has \(\beta>0\) and yields only a bounded target-exponent exclusion.
\item The branch cannot state a source lower-bound model, target class, and transfer inequality in the same notation.
\item The proposed escape is only a black-box substitution restatement with no new carrier or nonblackbox mechanism.
\item The barrier audit shows the needed source lower bound is just as hard as the original P vs NP/P/poly target.
\end{itemize}

\paragraph{Do not stop if.}
\begin{itemize}
\item The branch has not yet computed \(\beta\), q, and the induced source upper-bound exponent under \(\mathrm{SAT}\in \mathrm{SIZE}(n^{c})\).
\item A subpolynomial-query or nonblackbox carrier remains available and has not been audited.
\item The latest failure sharpened the threshold-collapse lemma into a formal no-go statement.
\item The route can still isolate advice leakage, nonuniformity, naturalness, or algebrization as the exact obstruction.
\end{itemize}

\paragraph{Required skills.}
\begin{itemize}
\item abducer-failure-signal-cartographer
\item abducer-refinement-route-ranker
\item contradiction-route-auditor
\item counterexample-hunter
\item entailment-gap-auditor
\item kb-abduce
\item kb-concept-bridge-planner
\item kb-deduce
\item kb-end-to-end-proof-critic
\item kb-gap-abstraction-conjecturer
\item kb-hypothesis-refutation-loop
\item kb-induce
\item kb-proof-connector
\item long-horizon-route-portfolio-manager
\item main-problem-alignment-auditor
\item next-experiment-queue-ranker
\item paper-contribution-readiness-judge
\item pvsnp-barrier-specificity-reviewer
\item pvsnp-conceptual-pressure-localizer
\item pvsnp-impossibility-theorem-forger
\item pvsnp-inference-quality-controller
\item pvsnp-meta-complexity-carrier-reviewer
\item pvsnp-proof-derivation-builder
\item pvsnp-qualitative-complexity-reasoner
\item pvsnp-quantifier-uniformity-reviewer
\item pvsnp-reduction-consistency-reviewer
\item pvsnp-symbolic-complexity-auditor
\end{itemize}

\subsection*{\#24 cycle-0002: Hardness-Magnification Bridge Realization}

\paragraph{Analysis.} Test MCSP or proof-complexity carriers for producing the unbounded-k source family: formalize whether a weak lower bound for MCSP[2\textasciicircum{}\{N\textasciicircum{}\{1/k\}\}] transfers to \(C_{H_{k}}(N)>N^{1+\epsilon}\) without becoming natural or losing the strict exponent margin.

\paragraph{Why this is a good target.} It audits a meta-complexity carrier where the decisive questions are verifier overhead, density/naturalness, and advice leakage. It checks whether padding/projection transfers really magnify hardness or merely move the proof debt into sparse lower bounds. It points toward the only unquarantined escape: a subpolynomial-query, nonprojection, or nonblackbox carrier. It keeps the barrier and model-boundary audit attached to the mathematical transfer rather than as after-the-fact commentary. It is executable as a bounded complexity calculation, CAS-backed exponent audit, or formal transfer check.

\paragraph{Risk.} Low drift risk: the item is framed as an explicit P vs NP transfer audit.

\paragraph{Branch control.} The item has high payoff and a concrete decision test; stopping now would be premature.

\paragraph{Premature-stop guard.} Do not stop before naming the carrier variant, deriving its verifier/reduction overhead, and checking whether a \(\mathrm{SAT}\in \mathrm{SIZE}(n^{c})\) assumption gives carrier circuits below the claimed source lower-bound threshold.

\paragraph{Payoff criteria.}
\begin{itemize}
\item For the MCSP/MKtP row branch, decide whether the certificate proves direct row hardness or only aggregate/rank-index hardness; record q\_extract=epsilon/(1+D) and q\_force=(D+epsilon)/(1+D) in the same parameter regime.
\item Promotion requires an advice-free row-hard certificate; backtrack if the ranked decoder, canonical unranking, or aggregate proof selects a hard row nonuniformly.
\item For this route, state the exact barrier premise it uses or escapes: relativizing, natural, algebrizing, nonuniform, or advice-leaking; promotion requires the escape to be attached to the specific transfer theorem.
\item For MCSP/MKtP, state N, n, s, verifier overhead, lower-bound model, and whether the property is natural or non-natural.
\item State whether the transfer excludes every polynomial SAT exponent or only a bounded interval \(c<c0\).
\item Promotion requires a target-preserving nonblackbox/subpolynomial mechanism, not another black-box Cook-Levin substitution.
\end{itemize}

\paragraph{Stop criteria.}
\begin{itemize}
\item The induced carrier upper bound under \(\mathrm{SAT}\in \mathrm{SIZE}(n^{c})\) remains N\textasciicircum{}\{q+beta c+o(1)\} with \(\beta>0\) for all tested parameter choices.
\item The source lower bound required for MCSP/MKtP is as strong as the target separation or becomes a natural property without an escape.
\item The route cannot prevent advice leakage or model mismatch between the carrier lower bound and P/poly target.
\end{itemize}

\paragraph{Do not stop if.}
\begin{itemize}
\item The MCSP/MKtP parameter regime has not yet been fixed.
\item The truth-table length N and underlying input length n have not been related in the transfer inequality.
\item The natural-proof largeness/constructivity issue has not yet been audited.
\item A nonblackbox use of the target circuit could still avoid the multiplicative c factor.
\end{itemize}

\paragraph{Required skills.}
\begin{itemize}
\item abducer-failure-signal-cartographer
\item abducer-refinement-route-ranker
\item counterexample-hunter
\item entailment-gap-auditor
\item kb-abduce
\item kb-concept-bridge-planner
\item kb-end-to-end-proof-critic
\item kb-gap-abstraction-conjecturer
\item kb-induce
\item kb-proof-connector
\item long-horizon-route-portfolio-manager
\item main-problem-alignment-auditor
\item next-experiment-queue-ranker
\item paper-contribution-readiness-judge
\item pvsnp-barrier-specificity-reviewer
\item pvsnp-conceptual-pressure-localizer
\item pvsnp-impossibility-theorem-forger
\item pvsnp-inference-quality-controller
\item pvsnp-meta-complexity-carrier-reviewer
\item pvsnp-proof-derivation-builder
\item pvsnp-qualitative-complexity-reasoner
\item pvsnp-quantifier-uniformity-reviewer
\item pvsnp-reduction-consistency-reviewer
\item pvsnp-symbolic-complexity-auditor
\end{itemize}

\subsection*{\#25 cycle-0125: Hardness-Magnification Bridge Realization}

\paragraph{Analysis.} Write the referee-auditable \(\mathrm{TSR}_{decl}\) proof block: child-row bit fields, fixed-point pseudocode, \(\Gamma_\mathrm{TSR}\) induction case, and a two-world hidden-pruning countermodel with a named first hidden child row r\_star.

\paragraph{Why this is a good target.} It directly tests the current black-box substitution cap by computing \(\beta\), q, and the induced target-exponent window. It audits a meta-complexity carrier where the decisive questions are verifier overhead, density/naturalness, and advice leakage. It keeps the barrier and model-boundary audit attached to the mathematical transfer rather than as after-the-fact commentary. It is executable as a bounded complexity calculation, CAS-backed exponent audit, or formal transfer check.

\paragraph{Risk.} Low drift risk: the item is framed as an explicit P vs NP transfer audit.

\paragraph{Branch control.} The item has high payoff and a concrete decision test; stopping now would be premature.

\paragraph{Premature-stop guard.} Do not stop before deriving the transfer inequality from \(\mathrm{SAT}\in \mathrm{SIZE}(n^{c})\) to the carrier upper bound, identifying whether \(\beta=0\), \(\beta>0\), or the argument is genuinely nonblackbox, and recording the barrier status of the proposed escape.

\paragraph{Payoff criteria.}
\begin{itemize}
\item For this route, state the exact barrier premise it uses or escapes: relativizing, natural, algebrizing, nonuniform, or advice-leaking; promotion requires the escape to be attached to the specific transfer theorem.
\item Compute \(\beta\) and q exactly; promotion requires \(\beta=0\), subpolynomial target queries, or a genuinely nonblackbox use of the SAT circuit.
\item For MCSP/MKtP, state N, n, s, verifier overhead, lower-bound model, and whether the property is natural or non-natural.
\item Record relativization, natural-proofs, algebrization, nonuniformity, and advice-leakage status for the specific transfer theorem.
\item Control gate: name the source carrier or language, target class, and circuit or resource model before claiming progress.
\item Control gate: derive the exact transfer inequality from \(\mathrm{SAT}\in \mathrm{SIZE}(n^{c})\), NP subset P/poly, or the stated target assumption to the source-carrier upper bound.
\end{itemize}

\paragraph{Stop criteria.}
\begin{itemize}
\item Every named carrier in the branch has \(\beta>0\) and yields only a bounded target-exponent exclusion.
\item The branch cannot state a source lower-bound model, target class, and transfer inequality in the same notation.
\item The proposed escape is only a black-box substitution restatement with no new carrier or nonblackbox mechanism.
\item The barrier audit shows the needed source lower bound is just as hard as the original P vs NP/P/poly target.
\end{itemize}

\paragraph{Do not stop if.}
\begin{itemize}
\item The branch has not yet computed \(\beta\), q, and the induced source upper-bound exponent under \(\mathrm{SAT}\in \mathrm{SIZE}(n^{c})\).
\item A subpolynomial-query or nonblackbox carrier remains available and has not been audited.
\item The latest failure sharpened the threshold-collapse lemma into a formal no-go statement.
\item The route can still isolate advice leakage, nonuniformity, naturalness, or algebrization as the exact obstruction.
\end{itemize}

\paragraph{Required skills.}
\begin{itemize}
\item contradiction-route-auditor
\item counterexample-hunter
\item entailment-gap-auditor
\item kb-abduce
\item kb-concept-bridge-planner
\item kb-deduce
\item kb-end-to-end-proof-critic
\item kb-hypothesis-refutation-loop
\item long-horizon-route-portfolio-manager
\item main-problem-alignment-auditor
\item next-experiment-queue-ranker
\item paper-contribution-readiness-judge
\item pvsnp-barrier-specificity-reviewer
\item pvsnp-conceptual-pressure-localizer
\item pvsnp-impossibility-theorem-forger
\item pvsnp-inference-quality-controller
\item pvsnp-meta-complexity-carrier-reviewer
\item pvsnp-proof-derivation-builder
\item pvsnp-qualitative-complexity-reasoner
\item pvsnp-quantifier-uniformity-reviewer
\item pvsnp-reduction-consistency-reviewer
\item pvsnp-symbolic-complexity-auditor
\end{itemize}

\subsection*{\#26 cycle-0068: Hardness-Magnification Bridge Realization}

\paragraph{Analysis.} Formalize the extension problem: given \(R_{N}\) and a sound searchable label certificate, prove either that no poly(N)-time ambient extension of density \(2^{-\mathrm{poly}(N)}\) exists for this certificate class, or that such an extension exists and therefore naturalizes the route.

\paragraph{Why this is a good target.} It directly tests the current black-box substitution cap by computing \(\beta\), q, and the induced target-exponent window. It audits a meta-complexity carrier where the decisive questions are verifier overhead, density/naturalness, and advice leakage. It keeps the barrier and model-boundary audit attached to the mathematical transfer rather than as after-the-fact commentary. It is executable as a bounded complexity calculation, CAS-backed exponent audit, or formal transfer check.

\paragraph{Risk.} Low drift risk: the item is framed as an explicit P vs NP transfer audit.

\paragraph{Branch control.} The item has high payoff and a concrete decision test; stopping now would be premature.

\paragraph{Premature-stop guard.} Do not stop before deriving the transfer inequality from \(\mathrm{SAT}\in \mathrm{SIZE}(n^{c})\) to the carrier upper bound, identifying whether \(\beta=0\), \(\beta>0\), or the argument is genuinely nonblackbox, and recording the barrier status of the proposed escape.

\paragraph{Payoff criteria.}
\begin{itemize}
\item For the proof-complexity branch, separate fixed-family proof hardness from language hardness by naming the proof system, shortcut algorithm, and semantic lower-bound target.
\item Promotion requires a proof-system lower bound that survives the Gaussian-elimination or recognition shortcut; otherwise the result is a restricted-proof-system separation only.
\item For this route, state the exact barrier premise it uses or escapes: relativizing, natural, algebrizing, nonuniform, or advice-leaking; promotion requires the escape to be attached to the specific transfer theorem.
\item Compute \(\beta\) and q exactly; promotion requires \(\beta=0\), subpolynomial target queries, or a genuinely nonblackbox use of the SAT circuit.
\item For MCSP/MKtP, state N, n, s, verifier overhead, lower-bound model, and whether the property is natural or non-natural.
\item Record relativization, natural-proofs, algebrization, nonuniformity, and advice-leakage status for the specific transfer theorem.
\end{itemize}

\paragraph{Stop criteria.}
\begin{itemize}
\item Every named carrier in the branch has \(\beta>0\) and yields only a bounded target-exponent exclusion.
\item The branch cannot state a source lower-bound model, target class, and transfer inequality in the same notation.
\item The proposed escape is only a black-box substitution restatement with no new carrier or nonblackbox mechanism.
\item The barrier audit shows the needed source lower bound is just as hard as the original P vs NP/P/poly target.
\end{itemize}

\paragraph{Do not stop if.}
\begin{itemize}
\item The branch has not yet computed \(\beta\), q, and the induced source upper-bound exponent under \(\mathrm{SAT}\in \mathrm{SIZE}(n^{c})\).
\item A subpolynomial-query or nonblackbox carrier remains available and has not been audited.
\item The latest failure sharpened the threshold-collapse lemma into a formal no-go statement.
\item The route can still isolate advice leakage, nonuniformity, naturalness, or algebrization as the exact obstruction.
\end{itemize}

\paragraph{Required skills.}
\begin{itemize}
\item contradiction-route-auditor
\item counterexample-hunter
\item entailment-gap-auditor
\item kb-abduce
\item kb-concept-bridge-planner
\item kb-deduce
\item kb-end-to-end-proof-critic
\item kb-hypothesis-refutation-loop
\item long-horizon-route-portfolio-manager
\item main-problem-alignment-auditor
\item next-experiment-queue-ranker
\item paper-contribution-readiness-judge
\item pvsnp-barrier-specificity-reviewer
\item pvsnp-conceptual-pressure-localizer
\item pvsnp-impossibility-theorem-forger
\item pvsnp-inference-quality-controller
\item pvsnp-meta-complexity-carrier-reviewer
\item pvsnp-proof-derivation-builder
\item pvsnp-qualitative-complexity-reasoner
\item pvsnp-quantifier-uniformity-reviewer
\item pvsnp-reduction-consistency-reviewer
\item pvsnp-symbolic-complexity-auditor
\end{itemize}

\subsection*{\#27 cycle-0108: Hardness-Magnification Bridge Realization}

\paragraph{Analysis.} Write the referee-style \(TransparentRowSupport_{N}\) definition and audit six examples as corollaries: interval-union, full scan, ROBDD, guarded-dDNNF, projected-dDNNF, and influence-complete extension rules. Do not add new selector families unless they prove a clause, refute a clause, or improve \(AC0PrefixNonemptiness_{N}\).

\paragraph{Why this is a good target.} It directly tests the current black-box substitution cap by computing \(\beta\), q, and the induced target-exponent window. It points toward the only unquarantined escape: a subpolynomial-query, nonprojection, or nonblackbox carrier. It keeps the barrier and model-boundary audit attached to the mathematical transfer rather than as after-the-fact commentary. It is executable as a bounded complexity calculation, CAS-backed exponent audit, or formal transfer check.

\paragraph{Risk.} Low drift risk: the item is framed as an explicit P vs NP transfer audit.

\paragraph{Branch control.} The item has high payoff and a concrete decision test; stopping now would be premature.

\paragraph{Premature-stop guard.} Do not stop before deriving the transfer inequality from \(\mathrm{SAT}\in \mathrm{SIZE}(n^{c})\) to the carrier upper bound, identifying whether \(\beta=0\), \(\beta>0\), or the argument is genuinely nonblackbox, and recording the barrier status of the proposed escape.

\paragraph{Payoff criteria.}
\begin{itemize}
\item For the restricted-model branch, compute the compiler envelope from unrestricted SIZE(m\textasciicircum{}k) or P-time algorithms into the named model; promotion requires a bound above that envelope on the same SAT slice.
\item Falsify direct transfer if a P function already has the same restricted-model lower bound, such as PARITY for \(\mathrm{AC0}\), HWB for OBDD, or XOR/Tseitin for proof-system degree.
\item For this route, state the exact barrier premise it uses or escapes: relativizing, natural, algebrizing, nonuniform, or advice-leaking; promotion requires the escape to be attached to the specific transfer theorem.
\item Compute \(\beta\) and q exactly; promotion requires \(\beta=0\), subpolynomial target queries, or a genuinely nonblackbox use of the SAT circuit.
\item Promotion requires a target-preserving nonblackbox/subpolynomial mechanism, not another black-box Cook-Levin substitution.
\item Record relativization, natural-proofs, algebrization, nonuniformity, and advice-leakage status for the specific transfer theorem.
\end{itemize}

\paragraph{Stop criteria.}
\begin{itemize}
\item Every named carrier in the branch has \(\beta>0\) and yields only a bounded target-exponent exclusion.
\item The branch cannot state a source lower-bound model, target class, and transfer inequality in the same notation.
\item The proposed escape is only a black-box substitution restatement with no new carrier or nonblackbox mechanism.
\item The barrier audit shows the needed source lower bound is just as hard as the original P vs NP/P/poly target.
\end{itemize}

\paragraph{Do not stop if.}
\begin{itemize}
\item The branch has not yet computed \(\beta\), q, and the induced source upper-bound exponent under \(\mathrm{SAT}\in \mathrm{SIZE}(n^{c})\).
\item A subpolynomial-query or nonblackbox carrier remains available and has not been audited.
\item The latest failure sharpened the threshold-collapse lemma into a formal no-go statement.
\item The route can still isolate advice leakage, nonuniformity, naturalness, or algebrization as the exact obstruction.
\end{itemize}

\paragraph{Required skills.}
\begin{itemize}
\item abducer-refinement-route-ranker
\item contradiction-route-auditor
\item entailment-gap-auditor
\item kb-deduce
\item kb-end-to-end-proof-critic
\item kb-gap-abstraction-conjecturer
\item kb-hypothesis-refutation-loop
\item kb-proof-connector
\item long-horizon-route-portfolio-manager
\item main-problem-alignment-auditor
\item next-experiment-queue-ranker
\item paper-contribution-readiness-judge
\item pvsnp-barrier-specificity-reviewer
\item pvsnp-conceptual-pressure-localizer
\item pvsnp-impossibility-theorem-forger
\item pvsnp-inference-quality-controller
\item pvsnp-proof-derivation-builder
\item pvsnp-qualitative-complexity-reasoner
\item pvsnp-quantifier-uniformity-reviewer
\item pvsnp-reduction-consistency-reviewer
\item pvsnp-symbolic-complexity-auditor
\end{itemize}

\subsection*{\#28 cycle-0011: Hardness-Magnification Bridge Realization: lookup-lift collapse for restricted target repair}

\paragraph{Analysis.} Test a nonblackbox proof-complexity carrier: encode a SAT circuit as a verifier object whose correctness is checked globally rather than by short oracle leaves, and compute whether the verifier overhead stays below \(N^{1+\epsilon}\) without allowing \(\mathrm{SAT}_{\le L}\) truth-table substitution.

\paragraph{Why this is a good target.} It directly tests the current black-box substitution cap by computing \(\beta\), q, and the induced target-exponent window. It points toward the only unquarantined escape: a subpolynomial-query, nonprojection, or nonblackbox carrier. It keeps the barrier and model-boundary audit attached to the mathematical transfer rather than as after-the-fact commentary. It is executable as a bounded complexity calculation, CAS-backed exponent audit, or formal transfer check.

\paragraph{Risk.} Low drift risk: the item is framed as an explicit P vs NP transfer audit.

\paragraph{Branch control.} The item has high payoff and a concrete decision test; stopping now would be premature.

\paragraph{Premature-stop guard.} Do not stop before deriving the transfer inequality from \(\mathrm{SAT}\in \mathrm{SIZE}(n^{c})\) to the carrier upper bound, identifying whether \(\beta=0\), \(\beta>0\), or the argument is genuinely nonblackbox, and recording the barrier status of the proposed escape.

\paragraph{Payoff criteria.}
\begin{itemize}
\item For this route, state the exact barrier premise it uses or escapes: relativizing, natural, algebrizing, nonuniform, or advice-leaking; promotion requires the escape to be attached to the specific transfer theorem.
\item Compute \(\beta\) and q exactly; promotion requires \(\beta=0\), subpolynomial target queries, or a genuinely nonblackbox use of the SAT circuit.
\item Promotion requires a target-preserving nonblackbox/subpolynomial mechanism, not another black-box Cook-Levin substitution.
\item Record relativization, natural-proofs, algebrization, nonuniformity, and advice-leakage status for the specific transfer theorem.
\item Control gate: name the source carrier or language, target class, and circuit or resource model before claiming progress.
\item Control gate: derive the exact transfer inequality from \(\mathrm{SAT}\in \mathrm{SIZE}(n^{c})\), NP subset P/poly, or the stated target assumption to the source-carrier upper bound.
\end{itemize}

\paragraph{Stop criteria.}
\begin{itemize}
\item Every named carrier in the branch has \(\beta>0\) and yields only a bounded target-exponent exclusion.
\item The branch cannot state a source lower-bound model, target class, and transfer inequality in the same notation.
\item The proposed escape is only a black-box substitution restatement with no new carrier or nonblackbox mechanism.
\item The barrier audit shows the needed source lower bound is just as hard as the original P vs NP/P/poly target.
\end{itemize}

\paragraph{Do not stop if.}
\begin{itemize}
\item The branch has not yet computed \(\beta\), q, and the induced source upper-bound exponent under \(\mathrm{SAT}\in \mathrm{SIZE}(n^{c})\).
\item A subpolynomial-query or nonblackbox carrier remains available and has not been audited.
\item The latest failure sharpened the threshold-collapse lemma into a formal no-go statement.
\item The route can still isolate advice leakage, nonuniformity, naturalness, or algebrization as the exact obstruction.
\end{itemize}

\paragraph{Required skills.}
\begin{itemize}
\item abducer-refinement-route-ranker
\item contradiction-route-auditor
\item entailment-gap-auditor
\item kb-deduce
\item kb-end-to-end-proof-critic
\item kb-gap-abstraction-conjecturer
\item kb-hypothesis-refutation-loop
\item kb-proof-connector
\item long-horizon-route-portfolio-manager
\item main-problem-alignment-auditor
\item next-experiment-queue-ranker
\item paper-contribution-readiness-judge
\item pvsnp-barrier-specificity-reviewer
\item pvsnp-conceptual-pressure-localizer
\item pvsnp-impossibility-theorem-forger
\item pvsnp-inference-quality-controller
\item pvsnp-proof-derivation-builder
\item pvsnp-qualitative-complexity-reasoner
\item pvsnp-quantifier-uniformity-reviewer
\item pvsnp-reduction-consistency-reviewer
\item pvsnp-symbolic-complexity-auditor
\end{itemize}

\subsection*{\#29 cycle-0010: Hardness-Magnification Bridge Realization}

\paragraph{Analysis.} Test a non-natural nonblackbox carrier whose query length is \(N^{o(1)}\) and whose use of a SAT circuit is semantic rather than oracle substitution; compute whether its lift tax \(\delta\) is below the source slack \(\epsilon\).

\paragraph{Why this is a good target.} It directly tests the current black-box substitution cap by computing \(\beta\), q, and the induced target-exponent window. It points toward the only unquarantined escape: a subpolynomial-query, nonprojection, or nonblackbox carrier. It keeps the barrier and model-boundary audit attached to the mathematical transfer rather than as after-the-fact commentary. It is executable as a bounded complexity calculation, CAS-backed exponent audit, or formal transfer check.

\paragraph{Risk.} Low drift risk: the item is framed as an explicit P vs NP transfer audit.

\paragraph{Branch control.} The item has high payoff and a concrete decision test; stopping now would be premature.

\paragraph{Premature-stop guard.} Do not stop before deriving the transfer inequality from \(\mathrm{SAT}\in \mathrm{SIZE}(n^{c})\) to the carrier upper bound, identifying whether \(\beta=0\), \(\beta>0\), or the argument is genuinely nonblackbox, and recording the barrier status of the proposed escape.

\paragraph{Payoff criteria.}
\begin{itemize}
\item For this route, state the exact barrier premise it uses or escapes: relativizing, natural, algebrizing, nonuniform, or advice-leaking; promotion requires the escape to be attached to the specific transfer theorem.
\item Compute \(\beta\) and q exactly; promotion requires \(\beta=0\), subpolynomial target queries, or a genuinely nonblackbox use of the SAT circuit.
\item Promotion requires a target-preserving nonblackbox/subpolynomial mechanism, not another black-box Cook-Levin substitution.
\item Record relativization, natural-proofs, algebrization, nonuniformity, and advice-leakage status for the specific transfer theorem.
\item Control gate: name the source carrier or language, target class, and circuit or resource model before claiming progress.
\item Control gate: derive the exact transfer inequality from \(\mathrm{SAT}\in \mathrm{SIZE}(n^{c})\), NP subset P/poly, or the stated target assumption to the source-carrier upper bound.
\end{itemize}

\paragraph{Stop criteria.}
\begin{itemize}
\item Every named carrier in the branch has \(\beta>0\) and yields only a bounded target-exponent exclusion.
\item The branch cannot state a source lower-bound model, target class, and transfer inequality in the same notation.
\item The proposed escape is only a black-box substitution restatement with no new carrier or nonblackbox mechanism.
\item The barrier audit shows the needed source lower bound is just as hard as the original P vs NP/P/poly target.
\end{itemize}

\paragraph{Do not stop if.}
\begin{itemize}
\item The branch has not yet computed \(\beta\), q, and the induced source upper-bound exponent under \(\mathrm{SAT}\in \mathrm{SIZE}(n^{c})\).
\item A subpolynomial-query or nonblackbox carrier remains available and has not been audited.
\item The latest failure sharpened the threshold-collapse lemma into a formal no-go statement.
\item The route can still isolate advice leakage, nonuniformity, naturalness, or algebrization as the exact obstruction.
\end{itemize}

\paragraph{Required skills.}
\begin{itemize}
\item abducer-refinement-route-ranker
\item contradiction-route-auditor
\item entailment-gap-auditor
\item kb-deduce
\item kb-end-to-end-proof-critic
\item kb-gap-abstraction-conjecturer
\item kb-hypothesis-refutation-loop
\item kb-proof-connector
\item long-horizon-route-portfolio-manager
\item main-problem-alignment-auditor
\item next-experiment-queue-ranker
\item paper-contribution-readiness-judge
\item pvsnp-barrier-specificity-reviewer
\item pvsnp-conceptual-pressure-localizer
\item pvsnp-impossibility-theorem-forger
\item pvsnp-inference-quality-controller
\item pvsnp-proof-derivation-builder
\item pvsnp-qualitative-complexity-reasoner
\item pvsnp-quantifier-uniformity-reviewer
\item pvsnp-reduction-consistency-reviewer
\item pvsnp-symbolic-complexity-auditor
\end{itemize}

\subsection*{\#30 cycle-0003: Hardness-Magnification Bridge Realization}

\paragraph{Analysis.} Instantiate g(r)=ceil(log log(r+4)); prove the verifier and length test are poly(N), then analyze whether valid-slice lower bounds for \(H_{g}\) admit a proof-complexity or MCSP carrier formulation.

\paragraph{Why this is a good target.} It audits a meta-complexity carrier where the decisive questions are verifier overhead, density/naturalness, and advice leakage. It checks whether padding/projection transfers really magnify hardness or merely move the proof debt into sparse lower bounds. It points toward the only unquarantined escape: a subpolynomial-query, nonprojection, or nonblackbox carrier. It is executable as a bounded complexity calculation, CAS-backed exponent audit, or formal transfer check.

\paragraph{Risk.} Low drift risk: the item is framed as an explicit P vs NP transfer audit.

\paragraph{Branch control.} The item has high payoff and a concrete decision test; stopping now would be premature.

\paragraph{Premature-stop guard.} Do not stop before naming the carrier variant, deriving its verifier/reduction overhead, and checking whether a \(\mathrm{SAT}\in \mathrm{SIZE}(n^{c})\) assumption gives carrier circuits below the claimed source lower-bound threshold.

\paragraph{Payoff criteria.}
\begin{itemize}
\item For the MCSP/MKtP row branch, decide whether the certificate proves direct row hardness or only aggregate/rank-index hardness; record q\_extract=epsilon/(1+D) and q\_force=(D+epsilon)/(1+D) in the same parameter regime.
\item Promotion requires an advice-free row-hard certificate; backtrack if the ranked decoder, canonical unranking, or aggregate proof selects a hard row nonuniformly.
\item For MCSP/MKtP, state N, n, s, verifier overhead, lower-bound model, and whether the property is natural or non-natural.
\item State whether the transfer excludes every polynomial SAT exponent or only a bounded interval \(c<c0\).
\item Promotion requires a target-preserving nonblackbox/subpolynomial mechanism, not another black-box Cook-Levin substitution.
\item Control gate: name the source carrier or language, target class, and circuit or resource model before claiming progress.
\end{itemize}

\paragraph{Stop criteria.}
\begin{itemize}
\item The induced carrier upper bound under \(\mathrm{SAT}\in \mathrm{SIZE}(n^{c})\) remains N\textasciicircum{}\{q+beta c+o(1)\} with \(\beta>0\) for all tested parameter choices.
\item The source lower bound required for MCSP/MKtP is as strong as the target separation or becomes a natural property without an escape.
\item The route cannot prevent advice leakage or model mismatch between the carrier lower bound and P/poly target.
\end{itemize}

\paragraph{Do not stop if.}
\begin{itemize}
\item The MCSP/MKtP parameter regime has not yet been fixed.
\item The truth-table length N and underlying input length n have not been related in the transfer inequality.
\item The natural-proof largeness/constructivity issue has not yet been audited.
\item A nonblackbox use of the target circuit could still avoid the multiplicative c factor.
\end{itemize}

\paragraph{Required skills.}
\begin{itemize}
\item abducer-failure-signal-cartographer
\item abducer-refinement-route-ranker
\item counterexample-hunter
\item entailment-gap-auditor
\item kb-abduce
\item kb-concept-bridge-planner
\item kb-gap-abstraction-conjecturer
\item kb-induce
\item kb-proof-connector
\item long-horizon-route-portfolio-manager
\item main-problem-alignment-auditor
\item next-experiment-queue-ranker
\item pvsnp-conceptual-pressure-localizer
\item pvsnp-impossibility-theorem-forger
\item pvsnp-inference-quality-controller
\item pvsnp-meta-complexity-carrier-reviewer
\item pvsnp-proof-derivation-builder
\item pvsnp-qualitative-complexity-reasoner
\item pvsnp-quantifier-uniformity-reviewer
\item pvsnp-reduction-consistency-reviewer
\item pvsnp-symbolic-complexity-auditor
\end{itemize}

\subsection*{\#31 cycle-0138: Hardness-Magnification Bridge Realization}

\paragraph{Analysis.} Write the referee-facing CTSRFrontierCompleteness proof: give closure pseudocode for \(K_\mathrm{CTSR}\), prove termination and minimality, verify one adaptive TSR line bit-by-bit, and instantiate the first-hidden-child falsifier on a one-node unread-enable-bit example.

\paragraph{Why this is a good target.} It directly tests the current black-box substitution cap by computing \(\beta\), q, and the induced target-exponent window. It checks whether padding/projection transfers really magnify hardness or merely move the proof debt into sparse lower bounds. It keeps the barrier and model-boundary audit attached to the mathematical transfer rather than as after-the-fact commentary. It is executable as a bounded complexity calculation, CAS-backed exponent audit, or formal transfer check.

\paragraph{Risk.} Low drift risk: the item is framed as an explicit P vs NP transfer audit.

\paragraph{Branch control.} The item has high payoff and a concrete decision test; stopping now would be premature.

\paragraph{Premature-stop guard.} Do not stop before deriving the transfer inequality from \(\mathrm{SAT}\in \mathrm{SIZE}(n^{c})\) to the carrier upper bound, identifying whether \(\beta=0\), \(\beta>0\), or the argument is genuinely nonblackbox, and recording the barrier status of the proposed escape.

\paragraph{Payoff criteria.}
\begin{itemize}
\item For this route, state the exact barrier premise it uses or escapes: relativizing, natural, algebrizing, nonuniform, or advice-leaking; promotion requires the escape to be attached to the specific transfer theorem.
\item Compute \(\beta\) and q exactly; promotion requires \(\beta=0\), subpolynomial target queries, or a genuinely nonblackbox use of the SAT circuit.
\item State whether the transfer excludes every polynomial SAT exponent or only a bounded interval \(c<c0\).
\item Record relativization, natural-proofs, algebrization, nonuniformity, and advice-leakage status for the specific transfer theorem.
\item Control gate: name the source carrier or language, target class, and circuit or resource model before claiming progress.
\item Control gate: derive the exact transfer inequality from \(\mathrm{SAT}\in \mathrm{SIZE}(n^{c})\), NP subset P/poly, or the stated target assumption to the source-carrier upper bound.
\end{itemize}

\paragraph{Stop criteria.}
\begin{itemize}
\item Every named carrier in the branch has \(\beta>0\) and yields only a bounded target-exponent exclusion.
\item The branch cannot state a source lower-bound model, target class, and transfer inequality in the same notation.
\item The proposed escape is only a black-box substitution restatement with no new carrier or nonblackbox mechanism.
\item The barrier audit shows the needed source lower bound is just as hard as the original P vs NP/P/poly target.
\end{itemize}

\paragraph{Do not stop if.}
\begin{itemize}
\item The branch has not yet computed \(\beta\), q, and the induced source upper-bound exponent under \(\mathrm{SAT}\in \mathrm{SIZE}(n^{c})\).
\item A subpolynomial-query or nonblackbox carrier remains available and has not been audited.
\item The latest failure sharpened the threshold-collapse lemma into a formal no-go statement.
\item The route can still isolate advice leakage, nonuniformity, naturalness, or algebrization as the exact obstruction.
\end{itemize}

\paragraph{Required skills.}
\begin{itemize}
\item abducer-failure-signal-cartographer
\item contradiction-route-auditor
\item entailment-gap-auditor
\item kb-deduce
\item kb-end-to-end-proof-critic
\item kb-hypothesis-refutation-loop
\item kb-induce
\item long-horizon-route-portfolio-manager
\item main-problem-alignment-auditor
\item next-experiment-queue-ranker
\item paper-contribution-readiness-judge
\item pvsnp-barrier-specificity-reviewer
\item pvsnp-conceptual-pressure-localizer
\item pvsnp-impossibility-theorem-forger
\item pvsnp-inference-quality-controller
\item pvsnp-proof-derivation-builder
\item pvsnp-qualitative-complexity-reasoner
\item pvsnp-quantifier-uniformity-reviewer
\item pvsnp-reduction-consistency-reviewer
\item pvsnp-symbolic-complexity-auditor
\end{itemize}

\subsection*{\#32 cycle-0131: Hardness-Magnification Bridge Realization}

\paragraph{Analysis.} Run the reviewer audit on L131.1: instantiate one toy projected-dDNNF FormulaEval line plus one declared TSR parent with two restriction children, compute \(K_\mathrm{DTSR}\) by the worklist, verify \(\Gamma_\mathrm{DTSR}\) numerically as \(N^{1+\kappa+\theta+o(1)}\), then flip the first unlisted child in the hidden-pruning variant to validate the private-row falsifier.

\paragraph{Why this is a good target.} It directly tests the current black-box substitution cap by computing \(\beta\), q, and the induced target-exponent window. It checks whether padding/projection transfers really magnify hardness or merely move the proof debt into sparse lower bounds. It keeps the barrier and model-boundary audit attached to the mathematical transfer rather than as after-the-fact commentary. It is executable as a bounded complexity calculation, CAS-backed exponent audit, or formal transfer check.

\paragraph{Risk.} Low drift risk: the item is framed as an explicit P vs NP transfer audit.

\paragraph{Branch control.} The item has high payoff and a concrete decision test; stopping now would be premature.

\paragraph{Premature-stop guard.} Do not stop before deriving the transfer inequality from \(\mathrm{SAT}\in \mathrm{SIZE}(n^{c})\) to the carrier upper bound, identifying whether \(\beta=0\), \(\beta>0\), or the argument is genuinely nonblackbox, and recording the barrier status of the proposed escape.

\paragraph{Payoff criteria.}
\begin{itemize}
\item For this route, state the exact barrier premise it uses or escapes: relativizing, natural, algebrizing, nonuniform, or advice-leaking; promotion requires the escape to be attached to the specific transfer theorem.
\item Compute \(\beta\) and q exactly; promotion requires \(\beta=0\), subpolynomial target queries, or a genuinely nonblackbox use of the SAT circuit.
\item State whether the transfer excludes every polynomial SAT exponent or only a bounded interval \(c<c0\).
\item Record relativization, natural-proofs, algebrization, nonuniformity, and advice-leakage status for the specific transfer theorem.
\item Control gate: name the source carrier or language, target class, and circuit or resource model before claiming progress.
\item Control gate: derive the exact transfer inequality from \(\mathrm{SAT}\in \mathrm{SIZE}(n^{c})\), NP subset P/poly, or the stated target assumption to the source-carrier upper bound.
\end{itemize}

\paragraph{Stop criteria.}
\begin{itemize}
\item Every named carrier in the branch has \(\beta>0\) and yields only a bounded target-exponent exclusion.
\item The branch cannot state a source lower-bound model, target class, and transfer inequality in the same notation.
\item The proposed escape is only a black-box substitution restatement with no new carrier or nonblackbox mechanism.
\item The barrier audit shows the needed source lower bound is just as hard as the original P vs NP/P/poly target.
\end{itemize}

\paragraph{Do not stop if.}
\begin{itemize}
\item The branch has not yet computed \(\beta\), q, and the induced source upper-bound exponent under \(\mathrm{SAT}\in \mathrm{SIZE}(n^{c})\).
\item A subpolynomial-query or nonblackbox carrier remains available and has not been audited.
\item The latest failure sharpened the threshold-collapse lemma into a formal no-go statement.
\item The route can still isolate advice leakage, nonuniformity, naturalness, or algebrization as the exact obstruction.
\end{itemize}

\paragraph{Required skills.}
\begin{itemize}
\item abducer-failure-signal-cartographer
\item contradiction-route-auditor
\item entailment-gap-auditor
\item kb-deduce
\item kb-end-to-end-proof-critic
\item kb-hypothesis-refutation-loop
\item kb-induce
\item long-horizon-route-portfolio-manager
\item main-problem-alignment-auditor
\item next-experiment-queue-ranker
\item paper-contribution-readiness-judge
\item pvsnp-barrier-specificity-reviewer
\item pvsnp-conceptual-pressure-localizer
\item pvsnp-impossibility-theorem-forger
\item pvsnp-inference-quality-controller
\item pvsnp-proof-derivation-builder
\item pvsnp-qualitative-complexity-reasoner
\item pvsnp-quantifier-uniformity-reviewer
\item pvsnp-reduction-consistency-reviewer
\item pvsnp-symbolic-complexity-auditor
\end{itemize}

\subsection*{\#33 cycle-0123: Hardness-Magnification Bridge Realization}

\paragraph{Analysis.} Write the referee-checkable trace for the projected-dDNNF example: list gate fields, parent pointers, projection maps, RowSet certificates, \(K_{0},K_{1}\),... closure additions, \(U_{N}\) and \(X_{N}\) charges, then flip one unread enable bit to exhibit the first-hidden-row countermodel and verify the \(\Gamma_{dag}\) threshold numerically for qepsilon.

\paragraph{Why this is a good target.} It directly tests the current black-box substitution cap by computing \(\beta\), q, and the induced target-exponent window. It checks whether padding/projection transfers really magnify hardness or merely move the proof debt into sparse lower bounds. It keeps the barrier and model-boundary audit attached to the mathematical transfer rather than as after-the-fact commentary. It is executable as a bounded complexity calculation, CAS-backed exponent audit, or formal transfer check.

\paragraph{Risk.} Low drift risk: the item is framed as an explicit P vs NP transfer audit.

\paragraph{Branch control.} The item has high payoff and a concrete decision test; stopping now would be premature.

\paragraph{Premature-stop guard.} Do not stop before deriving the transfer inequality from \(\mathrm{SAT}\in \mathrm{SIZE}(n^{c})\) to the carrier upper bound, identifying whether \(\beta=0\), \(\beta>0\), or the argument is genuinely nonblackbox, and recording the barrier status of the proposed escape.

\paragraph{Payoff criteria.}
\begin{itemize}
\item For this route, state the exact barrier premise it uses or escapes: relativizing, natural, algebrizing, nonuniform, or advice-leaking; promotion requires the escape to be attached to the specific transfer theorem.
\item Compute \(\beta\) and q exactly; promotion requires \(\beta=0\), subpolynomial target queries, or a genuinely nonblackbox use of the SAT circuit.
\item State whether the transfer excludes every polynomial SAT exponent or only a bounded interval \(c<c0\).
\item Record relativization, natural-proofs, algebrization, nonuniformity, and advice-leakage status for the specific transfer theorem.
\item Control gate: name the source carrier or language, target class, and circuit or resource model before claiming progress.
\item Control gate: derive the exact transfer inequality from \(\mathrm{SAT}\in \mathrm{SIZE}(n^{c})\), NP subset P/poly, or the stated target assumption to the source-carrier upper bound.
\end{itemize}

\paragraph{Stop criteria.}
\begin{itemize}
\item Every named carrier in the branch has \(\beta>0\) and yields only a bounded target-exponent exclusion.
\item The branch cannot state a source lower-bound model, target class, and transfer inequality in the same notation.
\item The proposed escape is only a black-box substitution restatement with no new carrier or nonblackbox mechanism.
\item The barrier audit shows the needed source lower bound is just as hard as the original P vs NP/P/poly target.
\end{itemize}

\paragraph{Do not stop if.}
\begin{itemize}
\item The branch has not yet computed \(\beta\), q, and the induced source upper-bound exponent under \(\mathrm{SAT}\in \mathrm{SIZE}(n^{c})\).
\item A subpolynomial-query or nonblackbox carrier remains available and has not been audited.
\item The latest failure sharpened the threshold-collapse lemma into a formal no-go statement.
\item The route can still isolate advice leakage, nonuniformity, naturalness, or algebrization as the exact obstruction.
\end{itemize}

\paragraph{Required skills.}
\begin{itemize}
\item abducer-failure-signal-cartographer
\item contradiction-route-auditor
\item entailment-gap-auditor
\item kb-deduce
\item kb-end-to-end-proof-critic
\item kb-hypothesis-refutation-loop
\item kb-induce
\item long-horizon-route-portfolio-manager
\item main-problem-alignment-auditor
\item next-experiment-queue-ranker
\item paper-contribution-readiness-judge
\item pvsnp-barrier-specificity-reviewer
\item pvsnp-conceptual-pressure-localizer
\item pvsnp-impossibility-theorem-forger
\item pvsnp-inference-quality-controller
\item pvsnp-proof-derivation-builder
\item pvsnp-qualitative-complexity-reasoner
\item pvsnp-quantifier-uniformity-reviewer
\item pvsnp-reduction-consistency-reviewer
\item pvsnp-symbolic-complexity-auditor
\end{itemize}

\subsection*{\#34 cycle-0107: Hardness-Magnification Bridge Realization}

\paragraph{Analysis.} Formalize \(AC0PrefixNonemptiness_{N}\) as a finite lemma: for \(\mathrm{AC0}\) selectors on q log N row bits, test whether descriptor structure yields \(N^{o(1)}\) prefix emptiness; if not, state the selector-scan theorem as final and add \(\mathrm{AC0}\) output-sensitive enumeration to open problems.

\paragraph{Why this is a good target.} It directly tests the current black-box substitution cap by computing \(\beta\), q, and the induced target-exponent window. It checks whether padding/projection transfers really magnify hardness or merely move the proof debt into sparse lower bounds. It keeps the barrier and model-boundary audit attached to the mathematical transfer rather than as after-the-fact commentary. It is executable as a bounded complexity calculation, CAS-backed exponent audit, or formal transfer check.

\paragraph{Risk.} Low drift risk: the item is framed as an explicit P vs NP transfer audit.

\paragraph{Branch control.} The item has high payoff and a concrete decision test; stopping now would be premature.

\paragraph{Premature-stop guard.} Do not stop before deriving the transfer inequality from \(\mathrm{SAT}\in \mathrm{SIZE}(n^{c})\) to the carrier upper bound, identifying whether \(\beta=0\), \(\beta>0\), or the argument is genuinely nonblackbox, and recording the barrier status of the proposed escape.

\paragraph{Payoff criteria.}
\begin{itemize}
\item For the restricted-model branch, compute the compiler envelope from unrestricted SIZE(m\textasciicircum{}k) or P-time algorithms into the named model; promotion requires a bound above that envelope on the same SAT slice.
\item Falsify direct transfer if a P function already has the same restricted-model lower bound, such as PARITY for \(\mathrm{AC0}\), HWB for OBDD, or XOR/Tseitin for proof-system degree.
\item For this route, state the exact barrier premise it uses or escapes: relativizing, natural, algebrizing, nonuniform, or advice-leaking; promotion requires the escape to be attached to the specific transfer theorem.
\item Compute \(\beta\) and q exactly; promotion requires \(\beta=0\), subpolynomial target queries, or a genuinely nonblackbox use of the SAT circuit.
\item State whether the transfer excludes every polynomial SAT exponent or only a bounded interval \(c<c0\).
\item Record relativization, natural-proofs, algebrization, nonuniformity, and advice-leakage status for the specific transfer theorem.
\end{itemize}

\paragraph{Stop criteria.}
\begin{itemize}
\item Every named carrier in the branch has \(\beta>0\) and yields only a bounded target-exponent exclusion.
\item The branch cannot state a source lower-bound model, target class, and transfer inequality in the same notation.
\item The proposed escape is only a black-box substitution restatement with no new carrier or nonblackbox mechanism.
\item The barrier audit shows the needed source lower bound is just as hard as the original P vs NP/P/poly target.
\end{itemize}

\paragraph{Do not stop if.}
\begin{itemize}
\item The branch has not yet computed \(\beta\), q, and the induced source upper-bound exponent under \(\mathrm{SAT}\in \mathrm{SIZE}(n^{c})\).
\item A subpolynomial-query or nonblackbox carrier remains available and has not been audited.
\item The latest failure sharpened the threshold-collapse lemma into a formal no-go statement.
\item The route can still isolate advice leakage, nonuniformity, naturalness, or algebrization as the exact obstruction.
\end{itemize}

\paragraph{Required skills.}
\begin{itemize}
\item abducer-failure-signal-cartographer
\item contradiction-route-auditor
\item entailment-gap-auditor
\item kb-deduce
\item kb-end-to-end-proof-critic
\item kb-hypothesis-refutation-loop
\item kb-induce
\item long-horizon-route-portfolio-manager
\item main-problem-alignment-auditor
\item next-experiment-queue-ranker
\item paper-contribution-readiness-judge
\item pvsnp-barrier-specificity-reviewer
\item pvsnp-conceptual-pressure-localizer
\item pvsnp-impossibility-theorem-forger
\item pvsnp-inference-quality-controller
\item pvsnp-proof-derivation-builder
\item pvsnp-qualitative-complexity-reasoner
\item pvsnp-quantifier-uniformity-reviewer
\item pvsnp-reduction-consistency-reviewer
\item pvsnp-symbolic-complexity-auditor
\end{itemize}

\subsection*{\#35 cycle-0090: Hardness-Magnification Bridge Realization}

\paragraph{Analysis.} Build the minimal \(\mathrm{EF}+\mathrm{EqSAT}+SR_{1}\) parser countermodel: allow one extension variable Z(F)=SAT(F\_0) OR SAT(F\_1) and test whether it is representable by the SR DAG without advice; if yes, increase to two nested extensions, if no, use it as the semantic escape separating full EF from \(\mathrm{Psys}_{h}\).

\paragraph{Why this is a good target.} It directly tests the current black-box substitution cap by computing \(\beta\), q, and the induced target-exponent window. It checks whether padding/projection transfers really magnify hardness or merely move the proof debt into sparse lower bounds. It keeps the barrier and model-boundary audit attached to the mathematical transfer rather than as after-the-fact commentary. It is executable as a bounded complexity calculation, CAS-backed exponent audit, or formal transfer check.

\paragraph{Risk.} Low drift risk: the item is framed as an explicit P vs NP transfer audit.

\paragraph{Branch control.} The item has high payoff and a concrete decision test; stopping now would be premature.

\paragraph{Premature-stop guard.} Do not stop before deriving the transfer inequality from \(\mathrm{SAT}\in \mathrm{SIZE}(n^{c})\) to the carrier upper bound, identifying whether \(\beta=0\), \(\beta>0\), or the argument is genuinely nonblackbox, and recording the barrier status of the proposed escape.

\paragraph{Payoff criteria.}
\begin{itemize}
\item For this route, state the exact barrier premise it uses or escapes: relativizing, natural, algebrizing, nonuniform, or advice-leaking; promotion requires the escape to be attached to the specific transfer theorem.
\item Compute \(\beta\) and q exactly; promotion requires \(\beta=0\), subpolynomial target queries, or a genuinely nonblackbox use of the SAT circuit.
\item State whether the transfer excludes every polynomial SAT exponent or only a bounded interval \(c<c0\).
\item Record relativization, natural-proofs, algebrization, nonuniformity, and advice-leakage status for the specific transfer theorem.
\item Control gate: name the source carrier or language, target class, and circuit or resource model before claiming progress.
\item Control gate: derive the exact transfer inequality from \(\mathrm{SAT}\in \mathrm{SIZE}(n^{c})\), NP subset P/poly, or the stated target assumption to the source-carrier upper bound.
\end{itemize}

\paragraph{Stop criteria.}
\begin{itemize}
\item Every named carrier in the branch has \(\beta>0\) and yields only a bounded target-exponent exclusion.
\item The branch cannot state a source lower-bound model, target class, and transfer inequality in the same notation.
\item The proposed escape is only a black-box substitution restatement with no new carrier or nonblackbox mechanism.
\item The barrier audit shows the needed source lower bound is just as hard as the original P vs NP/P/poly target.
\end{itemize}

\paragraph{Do not stop if.}
\begin{itemize}
\item The branch has not yet computed \(\beta\), q, and the induced source upper-bound exponent under \(\mathrm{SAT}\in \mathrm{SIZE}(n^{c})\).
\item A subpolynomial-query or nonblackbox carrier remains available and has not been audited.
\item The latest failure sharpened the threshold-collapse lemma into a formal no-go statement.
\item The route can still isolate advice leakage, nonuniformity, naturalness, or algebrization as the exact obstruction.
\end{itemize}

\paragraph{Required skills.}
\begin{itemize}
\item abducer-failure-signal-cartographer
\item contradiction-route-auditor
\item entailment-gap-auditor
\item kb-deduce
\item kb-end-to-end-proof-critic
\item kb-hypothesis-refutation-loop
\item kb-induce
\item long-horizon-route-portfolio-manager
\item main-problem-alignment-auditor
\item next-experiment-queue-ranker
\item paper-contribution-readiness-judge
\item pvsnp-barrier-specificity-reviewer
\item pvsnp-conceptual-pressure-localizer
\item pvsnp-impossibility-theorem-forger
\item pvsnp-inference-quality-controller
\item pvsnp-proof-derivation-builder
\item pvsnp-qualitative-complexity-reasoner
\item pvsnp-quantifier-uniformity-reviewer
\item pvsnp-reduction-consistency-reviewer
\item pvsnp-symbolic-complexity-auditor
\end{itemize}

\subsection*{\#36 cycle-0087: Hardness-Magnification Bridge Realization}

\paragraph{Analysis.} Instantiate the padded-copy EqSAT quotient family \(E_{k}\), compute the row-incidence matrix \(A_{E}\), and decide whether its constrained cosets have low-weight representatives below \(N^\epsilon\) or certify a valid semantic-link residual lower bound.

\paragraph{Why this is a good target.} It directly tests the current black-box substitution cap by computing \(\beta\), q, and the induced target-exponent window. It checks whether padding/projection transfers really magnify hardness or merely move the proof debt into sparse lower bounds. It keeps the barrier and model-boundary audit attached to the mathematical transfer rather than as after-the-fact commentary. It is executable as a bounded complexity calculation, CAS-backed exponent audit, or formal transfer check.

\paragraph{Risk.} Low drift risk: the item is framed as an explicit P vs NP transfer audit.

\paragraph{Branch control.} The item has high payoff and a concrete decision test; stopping now would be premature.

\paragraph{Premature-stop guard.} Do not stop before deriving the transfer inequality from \(\mathrm{SAT}\in \mathrm{SIZE}(n^{c})\) to the carrier upper bound, identifying whether \(\beta=0\), \(\beta>0\), or the argument is genuinely nonblackbox, and recording the barrier status of the proposed escape.

\paragraph{Payoff criteria.}
\begin{itemize}
\item For this route, state the exact barrier premise it uses or escapes: relativizing, natural, algebrizing, nonuniform, or advice-leaking; promotion requires the escape to be attached to the specific transfer theorem.
\item Compute \(\beta\) and q exactly; promotion requires \(\beta=0\), subpolynomial target queries, or a genuinely nonblackbox use of the SAT circuit.
\item State whether the transfer excludes every polynomial SAT exponent or only a bounded interval \(c<c0\).
\item Record relativization, natural-proofs, algebrization, nonuniformity, and advice-leakage status for the specific transfer theorem.
\item Control gate: name the source carrier or language, target class, and circuit or resource model before claiming progress.
\item Control gate: derive the exact transfer inequality from \(\mathrm{SAT}\in \mathrm{SIZE}(n^{c})\), NP subset P/poly, or the stated target assumption to the source-carrier upper bound.
\end{itemize}

\paragraph{Stop criteria.}
\begin{itemize}
\item Every named carrier in the branch has \(\beta>0\) and yields only a bounded target-exponent exclusion.
\item The branch cannot state a source lower-bound model, target class, and transfer inequality in the same notation.
\item The proposed escape is only a black-box substitution restatement with no new carrier or nonblackbox mechanism.
\item The barrier audit shows the needed source lower bound is just as hard as the original P vs NP/P/poly target.
\end{itemize}

\paragraph{Do not stop if.}
\begin{itemize}
\item The branch has not yet computed \(\beta\), q, and the induced source upper-bound exponent under \(\mathrm{SAT}\in \mathrm{SIZE}(n^{c})\).
\item A subpolynomial-query or nonblackbox carrier remains available and has not been audited.
\item The latest failure sharpened the threshold-collapse lemma into a formal no-go statement.
\item The route can still isolate advice leakage, nonuniformity, naturalness, or algebrization as the exact obstruction.
\end{itemize}

\paragraph{Required skills.}
\begin{itemize}
\item abducer-failure-signal-cartographer
\item contradiction-route-auditor
\item entailment-gap-auditor
\item kb-deduce
\item kb-end-to-end-proof-critic
\item kb-hypothesis-refutation-loop
\item kb-induce
\item long-horizon-route-portfolio-manager
\item main-problem-alignment-auditor
\item next-experiment-queue-ranker
\item paper-contribution-readiness-judge
\item pvsnp-barrier-specificity-reviewer
\item pvsnp-conceptual-pressure-localizer
\item pvsnp-impossibility-theorem-forger
\item pvsnp-inference-quality-controller
\item pvsnp-proof-derivation-builder
\item pvsnp-qualitative-complexity-reasoner
\item pvsnp-quantifier-uniformity-reviewer
\item pvsnp-reduction-consistency-reviewer
\item pvsnp-symbolic-complexity-auditor
\end{itemize}

\subsection*{\#37 cycle-0085: Hardness-Magnification Bridge Realization}

\paragraph{Analysis.} Attempt the finite \(\mathrm{EF}+SATeq_{N}\) parser proof: formalize proof lines as Boolean circuits over SAT-atoms, prove label-functionality for EF rules, charge every EqSAT endpoint to SemRows, and test whether SAT self-reduction axioms in full \(\mathrm{EF}+\mathrm{SAT}_{U}\) admit the same charge or force a new semantic-escape lemma.

\paragraph{Why this is a good target.} It directly tests the current black-box substitution cap by computing \(\beta\), q, and the induced target-exponent window. It checks whether padding/projection transfers really magnify hardness or merely move the proof debt into sparse lower bounds. It keeps the barrier and model-boundary audit attached to the mathematical transfer rather than as after-the-fact commentary. It is executable as a bounded complexity calculation, CAS-backed exponent audit, or formal transfer check.

\paragraph{Risk.} Low drift risk: the item is framed as an explicit P vs NP transfer audit.

\paragraph{Branch control.} The item has high payoff and a concrete decision test; stopping now would be premature.

\paragraph{Premature-stop guard.} Do not stop before deriving the transfer inequality from \(\mathrm{SAT}\in \mathrm{SIZE}(n^{c})\) to the carrier upper bound, identifying whether \(\beta=0\), \(\beta>0\), or the argument is genuinely nonblackbox, and recording the barrier status of the proposed escape.

\paragraph{Payoff criteria.}
\begin{itemize}
\item For this route, state the exact barrier premise it uses or escapes: relativizing, natural, algebrizing, nonuniform, or advice-leaking; promotion requires the escape to be attached to the specific transfer theorem.
\item Compute \(\beta\) and q exactly; promotion requires \(\beta=0\), subpolynomial target queries, or a genuinely nonblackbox use of the SAT circuit.
\item State whether the transfer excludes every polynomial SAT exponent or only a bounded interval \(c<c0\).
\item Record relativization, natural-proofs, algebrization, nonuniformity, and advice-leakage status for the specific transfer theorem.
\item Control gate: name the source carrier or language, target class, and circuit or resource model before claiming progress.
\item Control gate: derive the exact transfer inequality from \(\mathrm{SAT}\in \mathrm{SIZE}(n^{c})\), NP subset P/poly, or the stated target assumption to the source-carrier upper bound.
\end{itemize}

\paragraph{Stop criteria.}
\begin{itemize}
\item Every named carrier in the branch has \(\beta>0\) and yields only a bounded target-exponent exclusion.
\item The branch cannot state a source lower-bound model, target class, and transfer inequality in the same notation.
\item The proposed escape is only a black-box substitution restatement with no new carrier or nonblackbox mechanism.
\item The barrier audit shows the needed source lower bound is just as hard as the original P vs NP/P/poly target.
\end{itemize}

\paragraph{Do not stop if.}
\begin{itemize}
\item The branch has not yet computed \(\beta\), q, and the induced source upper-bound exponent under \(\mathrm{SAT}\in \mathrm{SIZE}(n^{c})\).
\item A subpolynomial-query or nonblackbox carrier remains available and has not been audited.
\item The latest failure sharpened the threshold-collapse lemma into a formal no-go statement.
\item The route can still isolate advice leakage, nonuniformity, naturalness, or algebrization as the exact obstruction.
\end{itemize}

\paragraph{Required skills.}
\begin{itemize}
\item abducer-failure-signal-cartographer
\item contradiction-route-auditor
\item entailment-gap-auditor
\item kb-deduce
\item kb-end-to-end-proof-critic
\item kb-hypothesis-refutation-loop
\item kb-induce
\item long-horizon-route-portfolio-manager
\item main-problem-alignment-auditor
\item next-experiment-queue-ranker
\item paper-contribution-readiness-judge
\item pvsnp-barrier-specificity-reviewer
\item pvsnp-conceptual-pressure-localizer
\item pvsnp-impossibility-theorem-forger
\item pvsnp-inference-quality-controller
\item pvsnp-proof-derivation-builder
\item pvsnp-qualitative-complexity-reasoner
\item pvsnp-quantifier-uniformity-reviewer
\item pvsnp-reduction-consistency-reviewer
\item pvsnp-symbolic-complexity-auditor
\end{itemize}

\subsection*{\#38 cycle-0084: Hardness-Magnification Bridge Realization}

\paragraph{Analysis.} Audit one standard full \(\mathrm{EF}+\mathrm{SAT}_{U}\) encoding rule: attempt to extract its occurrence-label dependency into \(Q_{pi}\) preserving \(A_{Q}-invariance\); if a witness or clause-soundness axiom breaks invariance, record the minimal rule-level countermodel and restrict the theorem statement to \(\mathrm{EF}+\mathrm{SAT}_{U}^{occ,rn}\).

\paragraph{Why this is a good target.} It directly tests the current black-box substitution cap by computing \(\beta\), q, and the induced target-exponent window. It checks whether padding/projection transfers really magnify hardness or merely move the proof debt into sparse lower bounds. It keeps the barrier and model-boundary audit attached to the mathematical transfer rather than as after-the-fact commentary. It is executable as a bounded complexity calculation, CAS-backed exponent audit, or formal transfer check.

\paragraph{Risk.} Low drift risk: the item is framed as an explicit P vs NP transfer audit.

\paragraph{Branch control.} The item has high payoff and a concrete decision test; stopping now would be premature.

\paragraph{Premature-stop guard.} Do not stop before deriving the transfer inequality from \(\mathrm{SAT}\in \mathrm{SIZE}(n^{c})\) to the carrier upper bound, identifying whether \(\beta=0\), \(\beta>0\), or the argument is genuinely nonblackbox, and recording the barrier status of the proposed escape.

\paragraph{Payoff criteria.}
\begin{itemize}
\item For the proof-complexity branch, separate fixed-family proof hardness from language hardness by naming the proof system, shortcut algorithm, and semantic lower-bound target.
\item Promotion requires a proof-system lower bound that survives the Gaussian-elimination or recognition shortcut; otherwise the result is a restricted-proof-system separation only.
\item For this route, state the exact barrier premise it uses or escapes: relativizing, natural, algebrizing, nonuniform, or advice-leaking; promotion requires the escape to be attached to the specific transfer theorem.
\item Compute \(\beta\) and q exactly; promotion requires \(\beta=0\), subpolynomial target queries, or a genuinely nonblackbox use of the SAT circuit.
\item State whether the transfer excludes every polynomial SAT exponent or only a bounded interval \(c<c0\).
\item Record relativization, natural-proofs, algebrization, nonuniformity, and advice-leakage status for the specific transfer theorem.
\end{itemize}

\paragraph{Stop criteria.}
\begin{itemize}
\item Every named carrier in the branch has \(\beta>0\) and yields only a bounded target-exponent exclusion.
\item The branch cannot state a source lower-bound model, target class, and transfer inequality in the same notation.
\item The proposed escape is only a black-box substitution restatement with no new carrier or nonblackbox mechanism.
\item The barrier audit shows the needed source lower bound is just as hard as the original P vs NP/P/poly target.
\end{itemize}

\paragraph{Do not stop if.}
\begin{itemize}
\item The branch has not yet computed \(\beta\), q, and the induced source upper-bound exponent under \(\mathrm{SAT}\in \mathrm{SIZE}(n^{c})\).
\item A subpolynomial-query or nonblackbox carrier remains available and has not been audited.
\item The latest failure sharpened the threshold-collapse lemma into a formal no-go statement.
\item The route can still isolate advice leakage, nonuniformity, naturalness, or algebrization as the exact obstruction.
\end{itemize}

\paragraph{Required skills.}
\begin{itemize}
\item abducer-failure-signal-cartographer
\item contradiction-route-auditor
\item entailment-gap-auditor
\item kb-deduce
\item kb-end-to-end-proof-critic
\item kb-hypothesis-refutation-loop
\item kb-induce
\item long-horizon-route-portfolio-manager
\item main-problem-alignment-auditor
\item next-experiment-queue-ranker
\item paper-contribution-readiness-judge
\item pvsnp-barrier-specificity-reviewer
\item pvsnp-conceptual-pressure-localizer
\item pvsnp-impossibility-theorem-forger
\item pvsnp-inference-quality-controller
\item pvsnp-proof-derivation-builder
\item pvsnp-qualitative-complexity-reasoner
\item pvsnp-quantifier-uniformity-reviewer
\item pvsnp-reduction-consistency-reviewer
\item pvsnp-symbolic-complexity-auditor
\end{itemize}

\subsection*{\#39 cycle-0083: Hardness-Magnification Bridge Realization}

\paragraph{Analysis.} Instantiate the \(\mathrm{EF}+\mathrm{SAT}_{U}^{occ,row}\) rule ledger: for each extension, weakening, substitution, \(\mathrm{SAT}_{U}\) occurrence axiom, and derived parity line, mark whether it is \(Q_{pi}-local\), row-separable, or a charged cross-row constraint; the first uncharged cross-row rule becomes the countermodel or repair target.

\paragraph{Why this is a good target.} It directly tests the current black-box substitution cap by computing \(\beta\), q, and the induced target-exponent window. It checks whether padding/projection transfers really magnify hardness or merely move the proof debt into sparse lower bounds. It keeps the barrier and model-boundary audit attached to the mathematical transfer rather than as after-the-fact commentary. It is executable as a bounded complexity calculation, CAS-backed exponent audit, or formal transfer check.

\paragraph{Risk.} Low drift risk: the item is framed as an explicit P vs NP transfer audit.

\paragraph{Branch control.} The item has high payoff and a concrete decision test; stopping now would be premature.

\paragraph{Premature-stop guard.} Do not stop before deriving the transfer inequality from \(\mathrm{SAT}\in \mathrm{SIZE}(n^{c})\) to the carrier upper bound, identifying whether \(\beta=0\), \(\beta>0\), or the argument is genuinely nonblackbox, and recording the barrier status of the proposed escape.

\paragraph{Payoff criteria.}
\begin{itemize}
\item For the proof-complexity branch, separate fixed-family proof hardness from language hardness by naming the proof system, shortcut algorithm, and semantic lower-bound target.
\item Promotion requires a proof-system lower bound that survives the Gaussian-elimination or recognition shortcut; otherwise the result is a restricted-proof-system separation only.
\item For this route, state the exact barrier premise it uses or escapes: relativizing, natural, algebrizing, nonuniform, or advice-leaking; promotion requires the escape to be attached to the specific transfer theorem.
\item Compute \(\beta\) and q exactly; promotion requires \(\beta=0\), subpolynomial target queries, or a genuinely nonblackbox use of the SAT circuit.
\item State whether the transfer excludes every polynomial SAT exponent or only a bounded interval \(c<c0\).
\item Record relativization, natural-proofs, algebrization, nonuniformity, and advice-leakage status for the specific transfer theorem.
\end{itemize}

\paragraph{Stop criteria.}
\begin{itemize}
\item Every named carrier in the branch has \(\beta>0\) and yields only a bounded target-exponent exclusion.
\item The branch cannot state a source lower-bound model, target class, and transfer inequality in the same notation.
\item The proposed escape is only a black-box substitution restatement with no new carrier or nonblackbox mechanism.
\item The barrier audit shows the needed source lower bound is just as hard as the original P vs NP/P/poly target.
\end{itemize}

\paragraph{Do not stop if.}
\begin{itemize}
\item The branch has not yet computed \(\beta\), q, and the induced source upper-bound exponent under \(\mathrm{SAT}\in \mathrm{SIZE}(n^{c})\).
\item A subpolynomial-query or nonblackbox carrier remains available and has not been audited.
\item The latest failure sharpened the threshold-collapse lemma into a formal no-go statement.
\item The route can still isolate advice leakage, nonuniformity, naturalness, or algebrization as the exact obstruction.
\end{itemize}

\paragraph{Required skills.}
\begin{itemize}
\item abducer-failure-signal-cartographer
\item contradiction-route-auditor
\item entailment-gap-auditor
\item kb-deduce
\item kb-end-to-end-proof-critic
\item kb-hypothesis-refutation-loop
\item kb-induce
\item long-horizon-route-portfolio-manager
\item main-problem-alignment-auditor
\item next-experiment-queue-ranker
\item paper-contribution-readiness-judge
\item pvsnp-barrier-specificity-reviewer
\item pvsnp-conceptual-pressure-localizer
\item pvsnp-impossibility-theorem-forger
\item pvsnp-inference-quality-controller
\item pvsnp-proof-derivation-builder
\item pvsnp-qualitative-complexity-reasoner
\item pvsnp-quantifier-uniformity-reviewer
\item pvsnp-reduction-consistency-reviewer
\item pvsnp-symbolic-complexity-auditor
\end{itemize}

\subsection*{\#40 cycle-0079: Hardness-Magnification Bridge Realization: beta-zero XOR row-forcing audit}

\paragraph{Analysis.} Instantiate \(\mathrm{Psys}=\mathrm{EF}+\mathrm{SAT}_{U}^{occ}\): define the proof parser that extracts \(Q_{pi}\), prove acceptance transcript invariance under \(A_{Q}\), and compute whether known EF encodings force full rows or merely individual generated formula occurrences.

\paragraph{Why this is a good target.} It directly tests the current black-box substitution cap by computing \(\beta\), q, and the induced target-exponent window. It checks whether padding/projection transfers really magnify hardness or merely move the proof debt into sparse lower bounds. It keeps the barrier and model-boundary audit attached to the mathematical transfer rather than as after-the-fact commentary. It is executable as a bounded complexity calculation, CAS-backed exponent audit, or formal transfer check.

\paragraph{Risk.} Low drift risk: the item is framed as an explicit P vs NP transfer audit.

\paragraph{Branch control.} The item has high payoff and a concrete decision test; stopping now would be premature.

\paragraph{Premature-stop guard.} Do not stop before deriving the transfer inequality from \(\mathrm{SAT}\in \mathrm{SIZE}(n^{c})\) to the carrier upper bound, identifying whether \(\beta=0\), \(\beta>0\), or the argument is genuinely nonblackbox, and recording the barrier status of the proposed escape.

\paragraph{Payoff criteria.}
\begin{itemize}
\item For the proof-complexity branch, separate fixed-family proof hardness from language hardness by naming the proof system, shortcut algorithm, and semantic lower-bound target.
\item Promotion requires a proof-system lower bound that survives the Gaussian-elimination or recognition shortcut; otherwise the result is a restricted-proof-system separation only.
\item For this route, state the exact barrier premise it uses or escapes: relativizing, natural, algebrizing, nonuniform, or advice-leaking; promotion requires the escape to be attached to the specific transfer theorem.
\item Compute \(\beta\) and q exactly; promotion requires \(\beta=0\), subpolynomial target queries, or a genuinely nonblackbox use of the SAT circuit.
\item State whether the transfer excludes every polynomial SAT exponent or only a bounded interval \(c<c0\).
\item Record relativization, natural-proofs, algebrization, nonuniformity, and advice-leakage status for the specific transfer theorem.
\end{itemize}

\paragraph{Stop criteria.}
\begin{itemize}
\item Every named carrier in the branch has \(\beta>0\) and yields only a bounded target-exponent exclusion.
\item The branch cannot state a source lower-bound model, target class, and transfer inequality in the same notation.
\item The proposed escape is only a black-box substitution restatement with no new carrier or nonblackbox mechanism.
\item The barrier audit shows the needed source lower bound is just as hard as the original P vs NP/P/poly target.
\end{itemize}

\paragraph{Do not stop if.}
\begin{itemize}
\item The branch has not yet computed \(\beta\), q, and the induced source upper-bound exponent under \(\mathrm{SAT}\in \mathrm{SIZE}(n^{c})\).
\item A subpolynomial-query or nonblackbox carrier remains available and has not been audited.
\item The latest failure sharpened the threshold-collapse lemma into a formal no-go statement.
\item The route can still isolate advice leakage, nonuniformity, naturalness, or algebrization as the exact obstruction.
\end{itemize}

\paragraph{Required skills.}
\begin{itemize}
\item abducer-failure-signal-cartographer
\item contradiction-route-auditor
\item entailment-gap-auditor
\item kb-deduce
\item kb-end-to-end-proof-critic
\item kb-hypothesis-refutation-loop
\item kb-induce
\item long-horizon-route-portfolio-manager
\item main-problem-alignment-auditor
\item next-experiment-queue-ranker
\item paper-contribution-readiness-judge
\item pvsnp-barrier-specificity-reviewer
\item pvsnp-conceptual-pressure-localizer
\item pvsnp-impossibility-theorem-forger
\item pvsnp-inference-quality-controller
\item pvsnp-proof-derivation-builder
\item pvsnp-qualitative-complexity-reasoner
\item pvsnp-quantifier-uniformity-reviewer
\item pvsnp-reduction-consistency-reviewer
\item pvsnp-symbolic-complexity-auditor
\end{itemize}

\subsection*{\#41 cycle-0076: Hardness-Magnification Bridge Realization}

\paragraph{Analysis.} Formalize Robust\_N(s,e) for \(\mathrm{EF}+\mathrm{SAT}_{U}\) and test whether a short proof can remain sound after replacing \(H_{N}\) by \(D_{N}\) off a sparse row set.

\paragraph{Why this is a good target.} It directly tests the current black-box substitution cap by computing \(\beta\), q, and the induced target-exponent window. It checks whether padding/projection transfers really magnify hardness or merely move the proof debt into sparse lower bounds. It keeps the barrier and model-boundary audit attached to the mathematical transfer rather than as after-the-fact commentary. It is executable as a bounded complexity calculation, CAS-backed exponent audit, or formal transfer check.

\paragraph{Risk.} Low drift risk: the item is framed as an explicit P vs NP transfer audit.

\paragraph{Branch control.} The item has high payoff and a concrete decision test; stopping now would be premature.

\paragraph{Premature-stop guard.} Do not stop before deriving the transfer inequality from \(\mathrm{SAT}\in \mathrm{SIZE}(n^{c})\) to the carrier upper bound, identifying whether \(\beta=0\), \(\beta>0\), or the argument is genuinely nonblackbox, and recording the barrier status of the proposed escape.

\paragraph{Payoff criteria.}
\begin{itemize}
\item For this route, state the exact barrier premise it uses or escapes: relativizing, natural, algebrizing, nonuniform, or advice-leaking; promotion requires the escape to be attached to the specific transfer theorem.
\item Compute \(\beta\) and q exactly; promotion requires \(\beta=0\), subpolynomial target queries, or a genuinely nonblackbox use of the SAT circuit.
\item State whether the transfer excludes every polynomial SAT exponent or only a bounded interval \(c<c0\).
\item Record relativization, natural-proofs, algebrization, nonuniformity, and advice-leakage status for the specific transfer theorem.
\item Control gate: name the source carrier or language, target class, and circuit or resource model before claiming progress.
\item Control gate: derive the exact transfer inequality from \(\mathrm{SAT}\in \mathrm{SIZE}(n^{c})\), NP subset P/poly, or the stated target assumption to the source-carrier upper bound.
\end{itemize}

\paragraph{Stop criteria.}
\begin{itemize}
\item Every named carrier in the branch has \(\beta>0\) and yields only a bounded target-exponent exclusion.
\item The branch cannot state a source lower-bound model, target class, and transfer inequality in the same notation.
\item The proposed escape is only a black-box substitution restatement with no new carrier or nonblackbox mechanism.
\item The barrier audit shows the needed source lower bound is just as hard as the original P vs NP/P/poly target.
\end{itemize}

\paragraph{Do not stop if.}
\begin{itemize}
\item The branch has not yet computed \(\beta\), q, and the induced source upper-bound exponent under \(\mathrm{SAT}\in \mathrm{SIZE}(n^{c})\).
\item A subpolynomial-query or nonblackbox carrier remains available and has not been audited.
\item The latest failure sharpened the threshold-collapse lemma into a formal no-go statement.
\item The route can still isolate advice leakage, nonuniformity, naturalness, or algebrization as the exact obstruction.
\end{itemize}

\paragraph{Required skills.}
\begin{itemize}
\item abducer-failure-signal-cartographer
\item contradiction-route-auditor
\item entailment-gap-auditor
\item kb-deduce
\item kb-end-to-end-proof-critic
\item kb-hypothesis-refutation-loop
\item kb-induce
\item long-horizon-route-portfolio-manager
\item main-problem-alignment-auditor
\item next-experiment-queue-ranker
\item paper-contribution-readiness-judge
\item pvsnp-barrier-specificity-reviewer
\item pvsnp-conceptual-pressure-localizer
\item pvsnp-impossibility-theorem-forger
\item pvsnp-inference-quality-controller
\item pvsnp-proof-derivation-builder
\item pvsnp-qualitative-complexity-reasoner
\item pvsnp-quantifier-uniformity-reviewer
\item pvsnp-reduction-consistency-reviewer
\item pvsnp-symbolic-complexity-auditor
\end{itemize}

\subsection*{\#42 cycle-0073: Hardness-Magnification Bridge Realization}

\paragraph{Analysis.} Instantiate a named proof system Psys for \(\beta=0\) XOR certificates and test whether any accepted SAT-semantic certificate can be transformed into either a functional label property or a residual circuit wrong on fewer than \(N^{\epsilon-o(1)}\) rows.

\paragraph{Why this is a good target.} It directly tests the current black-box substitution cap by computing \(\beta\), q, and the induced target-exponent window. It checks whether padding/projection transfers really magnify hardness or merely move the proof debt into sparse lower bounds. It keeps the barrier and model-boundary audit attached to the mathematical transfer rather than as after-the-fact commentary. It is executable as a bounded complexity calculation, CAS-backed exponent audit, or formal transfer check.

\paragraph{Risk.} Low drift risk: the item is framed as an explicit P vs NP transfer audit.

\paragraph{Branch control.} The item has high payoff and a concrete decision test; stopping now would be premature.

\paragraph{Premature-stop guard.} Do not stop before deriving the transfer inequality from \(\mathrm{SAT}\in \mathrm{SIZE}(n^{c})\) to the carrier upper bound, identifying whether \(\beta=0\), \(\beta>0\), or the argument is genuinely nonblackbox, and recording the barrier status of the proposed escape.

\paragraph{Payoff criteria.}
\begin{itemize}
\item For the proof-complexity branch, separate fixed-family proof hardness from language hardness by naming the proof system, shortcut algorithm, and semantic lower-bound target.
\item Promotion requires a proof-system lower bound that survives the Gaussian-elimination or recognition shortcut; otherwise the result is a restricted-proof-system separation only.
\item For this route, state the exact barrier premise it uses or escapes: relativizing, natural, algebrizing, nonuniform, or advice-leaking; promotion requires the escape to be attached to the specific transfer theorem.
\item Compute \(\beta\) and q exactly; promotion requires \(\beta=0\), subpolynomial target queries, or a genuinely nonblackbox use of the SAT circuit.
\item State whether the transfer excludes every polynomial SAT exponent or only a bounded interval \(c<c0\).
\item Record relativization, natural-proofs, algebrization, nonuniformity, and advice-leakage status for the specific transfer theorem.
\end{itemize}

\paragraph{Stop criteria.}
\begin{itemize}
\item Every named carrier in the branch has \(\beta>0\) and yields only a bounded target-exponent exclusion.
\item The branch cannot state a source lower-bound model, target class, and transfer inequality in the same notation.
\item The proposed escape is only a black-box substitution restatement with no new carrier or nonblackbox mechanism.
\item The barrier audit shows the needed source lower bound is just as hard as the original P vs NP/P/poly target.
\end{itemize}

\paragraph{Do not stop if.}
\begin{itemize}
\item The branch has not yet computed \(\beta\), q, and the induced source upper-bound exponent under \(\mathrm{SAT}\in \mathrm{SIZE}(n^{c})\).
\item A subpolynomial-query or nonblackbox carrier remains available and has not been audited.
\item The latest failure sharpened the threshold-collapse lemma into a formal no-go statement.
\item The route can still isolate advice leakage, nonuniformity, naturalness, or algebrization as the exact obstruction.
\end{itemize}

\paragraph{Required skills.}
\begin{itemize}
\item abducer-failure-signal-cartographer
\item contradiction-route-auditor
\item entailment-gap-auditor
\item kb-deduce
\item kb-end-to-end-proof-critic
\item kb-hypothesis-refutation-loop
\item kb-induce
\item long-horizon-route-portfolio-manager
\item main-problem-alignment-auditor
\item next-experiment-queue-ranker
\item paper-contribution-readiness-judge
\item pvsnp-barrier-specificity-reviewer
\item pvsnp-conceptual-pressure-localizer
\item pvsnp-impossibility-theorem-forger
\item pvsnp-inference-quality-controller
\item pvsnp-proof-derivation-builder
\item pvsnp-qualitative-complexity-reasoner
\item pvsnp-quantifier-uniformity-reviewer
\item pvsnp-reduction-consistency-reviewer
\item pvsnp-symbolic-complexity-auditor
\end{itemize}

\subsection*{\#43 cycle-0072: Hardness-Magnification Bridge Realization}

\paragraph{Analysis.} Instantiate the selected-probe audit for one named proof-system candidate: require \(q=1+\epsilon/2\), \(\rho=1+\epsilon/3\), L=(log N)\textasciicircum{}2, then compute whether its certificates force \(A_{N}+S_{N}+R_{N}+E_{NN}\) or \(A_{N}+S_{N}+K_{NL}^{c}+E_{NN}\) below \(N^{1+\epsilon}\).

\paragraph{Why this is a good target.} It directly tests the current black-box substitution cap by computing \(\beta\), q, and the induced target-exponent window. It checks whether padding/projection transfers really magnify hardness or merely move the proof debt into sparse lower bounds. It keeps the barrier and model-boundary audit attached to the mathematical transfer rather than as after-the-fact commentary. It is executable as a bounded complexity calculation, CAS-backed exponent audit, or formal transfer check.

\paragraph{Risk.} Low drift risk: the item is framed as an explicit P vs NP transfer audit.

\paragraph{Branch control.} The item has high payoff and a concrete decision test; stopping now would be premature.

\paragraph{Premature-stop guard.} Do not stop before deriving the transfer inequality from \(\mathrm{SAT}\in \mathrm{SIZE}(n^{c})\) to the carrier upper bound, identifying whether \(\beta=0\), \(\beta>0\), or the argument is genuinely nonblackbox, and recording the barrier status of the proposed escape.

\paragraph{Payoff criteria.}
\begin{itemize}
\item For the proof-complexity branch, separate fixed-family proof hardness from language hardness by naming the proof system, shortcut algorithm, and semantic lower-bound target.
\item Promotion requires a proof-system lower bound that survives the Gaussian-elimination or recognition shortcut; otherwise the result is a restricted-proof-system separation only.
\item For this route, state the exact barrier premise it uses or escapes: relativizing, natural, algebrizing, nonuniform, or advice-leaking; promotion requires the escape to be attached to the specific transfer theorem.
\item Compute \(\beta\) and q exactly; promotion requires \(\beta=0\), subpolynomial target queries, or a genuinely nonblackbox use of the SAT circuit.
\item State whether the transfer excludes every polynomial SAT exponent or only a bounded interval \(c<c0\).
\item Record relativization, natural-proofs, algebrization, nonuniformity, and advice-leakage status for the specific transfer theorem.
\end{itemize}

\paragraph{Stop criteria.}
\begin{itemize}
\item Every named carrier in the branch has \(\beta>0\) and yields only a bounded target-exponent exclusion.
\item The branch cannot state a source lower-bound model, target class, and transfer inequality in the same notation.
\item The proposed escape is only a black-box substitution restatement with no new carrier or nonblackbox mechanism.
\item The barrier audit shows the needed source lower bound is just as hard as the original P vs NP/P/poly target.
\end{itemize}

\paragraph{Do not stop if.}
\begin{itemize}
\item The branch has not yet computed \(\beta\), q, and the induced source upper-bound exponent under \(\mathrm{SAT}\in \mathrm{SIZE}(n^{c})\).
\item A subpolynomial-query or nonblackbox carrier remains available and has not been audited.
\item The latest failure sharpened the threshold-collapse lemma into a formal no-go statement.
\item The route can still isolate advice leakage, nonuniformity, naturalness, or algebrization as the exact obstruction.
\end{itemize}

\paragraph{Required skills.}
\begin{itemize}
\item abducer-failure-signal-cartographer
\item contradiction-route-auditor
\item entailment-gap-auditor
\item kb-deduce
\item kb-end-to-end-proof-critic
\item kb-hypothesis-refutation-loop
\item kb-induce
\item long-horizon-route-portfolio-manager
\item main-problem-alignment-auditor
\item next-experiment-queue-ranker
\item paper-contribution-readiness-judge
\item pvsnp-barrier-specificity-reviewer
\item pvsnp-conceptual-pressure-localizer
\item pvsnp-impossibility-theorem-forger
\item pvsnp-inference-quality-controller
\item pvsnp-proof-derivation-builder
\item pvsnp-qualitative-complexity-reasoner
\item pvsnp-quantifier-uniformity-reviewer
\item pvsnp-reduction-consistency-reviewer
\item pvsnp-symbolic-complexity-auditor
\end{itemize}

\subsection*{\#44 cycle-0070: Hardness-Magnification Bridge Realization}

\paragraph{Analysis.} Instantiate the density-tier audit on a toy proof-certificate family: certificates describing affine anti-alignment witnesses over \(U_{N}\). Compute \(delta_{N}\), verifier time, and whether soundness reaches arbitrary \(C_{N}(H_{y})>N^{1+\epsilon}\) or only sketch hardness.

\paragraph{Why this is a good target.} It directly tests the current black-box substitution cap by computing \(\beta\), q, and the induced target-exponent window. It checks whether padding/projection transfers really magnify hardness or merely move the proof debt into sparse lower bounds. It keeps the barrier and model-boundary audit attached to the mathematical transfer rather than as after-the-fact commentary. It is executable as a bounded complexity calculation, CAS-backed exponent audit, or formal transfer check.

\paragraph{Risk.} Low drift risk: the item is framed as an explicit P vs NP transfer audit.

\paragraph{Branch control.} The item has high payoff and a concrete decision test; stopping now would be premature.

\paragraph{Premature-stop guard.} Do not stop before deriving the transfer inequality from \(\mathrm{SAT}\in \mathrm{SIZE}(n^{c})\) to the carrier upper bound, identifying whether \(\beta=0\), \(\beta>0\), or the argument is genuinely nonblackbox, and recording the barrier status of the proposed escape.

\paragraph{Payoff criteria.}
\begin{itemize}
\item For the proof-complexity branch, separate fixed-family proof hardness from language hardness by naming the proof system, shortcut algorithm, and semantic lower-bound target.
\item Promotion requires a proof-system lower bound that survives the Gaussian-elimination or recognition shortcut; otherwise the result is a restricted-proof-system separation only.
\item For this route, state the exact barrier premise it uses or escapes: relativizing, natural, algebrizing, nonuniform, or advice-leaking; promotion requires the escape to be attached to the specific transfer theorem.
\item Compute \(\beta\) and q exactly; promotion requires \(\beta=0\), subpolynomial target queries, or a genuinely nonblackbox use of the SAT circuit.
\item State whether the transfer excludes every polynomial SAT exponent or only a bounded interval \(c<c0\).
\item Record relativization, natural-proofs, algebrization, nonuniformity, and advice-leakage status for the specific transfer theorem.
\end{itemize}

\paragraph{Stop criteria.}
\begin{itemize}
\item Every named carrier in the branch has \(\beta>0\) and yields only a bounded target-exponent exclusion.
\item The branch cannot state a source lower-bound model, target class, and transfer inequality in the same notation.
\item The proposed escape is only a black-box substitution restatement with no new carrier or nonblackbox mechanism.
\item The barrier audit shows the needed source lower bound is just as hard as the original P vs NP/P/poly target.
\end{itemize}

\paragraph{Do not stop if.}
\begin{itemize}
\item The branch has not yet computed \(\beta\), q, and the induced source upper-bound exponent under \(\mathrm{SAT}\in \mathrm{SIZE}(n^{c})\).
\item A subpolynomial-query or nonblackbox carrier remains available and has not been audited.
\item The latest failure sharpened the threshold-collapse lemma into a formal no-go statement.
\item The route can still isolate advice leakage, nonuniformity, naturalness, or algebrization as the exact obstruction.
\end{itemize}

\paragraph{Required skills.}
\begin{itemize}
\item abducer-failure-signal-cartographer
\item contradiction-route-auditor
\item entailment-gap-auditor
\item kb-deduce
\item kb-end-to-end-proof-critic
\item kb-hypothesis-refutation-loop
\item kb-induce
\item long-horizon-route-portfolio-manager
\item main-problem-alignment-auditor
\item next-experiment-queue-ranker
\item paper-contribution-readiness-judge
\item pvsnp-barrier-specificity-reviewer
\item pvsnp-conceptual-pressure-localizer
\item pvsnp-impossibility-theorem-forger
\item pvsnp-inference-quality-controller
\item pvsnp-proof-derivation-builder
\item pvsnp-qualitative-complexity-reasoner
\item pvsnp-quantifier-uniformity-reviewer
\item pvsnp-reduction-consistency-reviewer
\item pvsnp-symbolic-complexity-auditor
\end{itemize}

\subsection*{\#45 cycle-0067: Hardness-Magnification Bridge Realization}

\paragraph{Analysis.} Instantiate \(U_{N}\) for one explicit polylog-query XOR generator, compute whether \(U_{N}\) is enumerable in poly(2\textasciicircum{}L), and classify a candidate certificate as dense-label, image-relative, or nonlarge.

\paragraph{Why this is a good target.} It directly tests the current black-box substitution cap by computing \(\beta\), q, and the induced target-exponent window. It checks whether padding/projection transfers really magnify hardness or merely move the proof debt into sparse lower bounds. It keeps the barrier and model-boundary audit attached to the mathematical transfer rather than as after-the-fact commentary. It is executable as a bounded complexity calculation, CAS-backed exponent audit, or formal transfer check.

\paragraph{Risk.} Low drift risk: the item is framed as an explicit P vs NP transfer audit.

\paragraph{Branch control.} The item has high payoff and a concrete decision test; stopping now would be premature.

\paragraph{Premature-stop guard.} Do not stop before deriving the transfer inequality from \(\mathrm{SAT}\in \mathrm{SIZE}(n^{c})\) to the carrier upper bound, identifying whether \(\beta=0\), \(\beta>0\), or the argument is genuinely nonblackbox, and recording the barrier status of the proposed escape.

\paragraph{Payoff criteria.}
\begin{itemize}
\item For the proof-complexity branch, separate fixed-family proof hardness from language hardness by naming the proof system, shortcut algorithm, and semantic lower-bound target.
\item Promotion requires a proof-system lower bound that survives the Gaussian-elimination or recognition shortcut; otherwise the result is a restricted-proof-system separation only.
\item For this route, state the exact barrier premise it uses or escapes: relativizing, natural, algebrizing, nonuniform, or advice-leaking; promotion requires the escape to be attached to the specific transfer theorem.
\item Compute \(\beta\) and q exactly; promotion requires \(\beta=0\), subpolynomial target queries, or a genuinely nonblackbox use of the SAT circuit.
\item State whether the transfer excludes every polynomial SAT exponent or only a bounded interval \(c<c0\).
\item Record relativization, natural-proofs, algebrization, nonuniformity, and advice-leakage status for the specific transfer theorem.
\end{itemize}

\paragraph{Stop criteria.}
\begin{itemize}
\item Every named carrier in the branch has \(\beta>0\) and yields only a bounded target-exponent exclusion.
\item The branch cannot state a source lower-bound model, target class, and transfer inequality in the same notation.
\item The proposed escape is only a black-box substitution restatement with no new carrier or nonblackbox mechanism.
\item The barrier audit shows the needed source lower bound is just as hard as the original P vs NP/P/poly target.
\end{itemize}

\paragraph{Do not stop if.}
\begin{itemize}
\item The branch has not yet computed \(\beta\), q, and the induced source upper-bound exponent under \(\mathrm{SAT}\in \mathrm{SIZE}(n^{c})\).
\item A subpolynomial-query or nonblackbox carrier remains available and has not been audited.
\item The latest failure sharpened the threshold-collapse lemma into a formal no-go statement.
\item The route can still isolate advice leakage, nonuniformity, naturalness, or algebrization as the exact obstruction.
\end{itemize}

\paragraph{Required skills.}
\begin{itemize}
\item abducer-failure-signal-cartographer
\item contradiction-route-auditor
\item entailment-gap-auditor
\item kb-deduce
\item kb-end-to-end-proof-critic
\item kb-hypothesis-refutation-loop
\item kb-induce
\item long-horizon-route-portfolio-manager
\item main-problem-alignment-auditor
\item next-experiment-queue-ranker
\item paper-contribution-readiness-judge
\item pvsnp-barrier-specificity-reviewer
\item pvsnp-conceptual-pressure-localizer
\item pvsnp-impossibility-theorem-forger
\item pvsnp-inference-quality-controller
\item pvsnp-proof-derivation-builder
\item pvsnp-qualitative-complexity-reasoner
\item pvsnp-quantifier-uniformity-reviewer
\item pvsnp-reduction-consistency-reviewer
\item pvsnp-symbolic-complexity-auditor
\end{itemize}

\subsection*{\#46 cycle-0063: Hardness-Magnification Bridge Realization}

\paragraph{Analysis.} Perform the finite prior-art boundary check for the label-space naturalness claim, then instantiate the thought experiment with a concrete pairwise-independent generator and compute \(r_{N}\), quotient rank, and any row-coordinate decoder cost \(D_{N}\).

\paragraph{Why this is a good target.} It directly tests the current black-box substitution cap by computing \(\beta\), q, and the induced target-exponent window. It checks whether padding/projection transfers really magnify hardness or merely move the proof debt into sparse lower bounds. It keeps the barrier and model-boundary audit attached to the mathematical transfer rather than as after-the-fact commentary. It is executable as a bounded complexity calculation, CAS-backed exponent audit, or formal transfer check.

\paragraph{Risk.} Low drift risk: the item is framed as an explicit P vs NP transfer audit.

\paragraph{Branch control.} The item has high payoff and a concrete decision test; stopping now would be premature.

\paragraph{Premature-stop guard.} Do not stop before deriving the transfer inequality from \(\mathrm{SAT}\in \mathrm{SIZE}(n^{c})\) to the carrier upper bound, identifying whether \(\beta=0\), \(\beta>0\), or the argument is genuinely nonblackbox, and recording the barrier status of the proposed escape.

\paragraph{Payoff criteria.}
\begin{itemize}
\item For this route, state the exact barrier premise it uses or escapes: relativizing, natural, algebrizing, nonuniform, or advice-leaking; promotion requires the escape to be attached to the specific transfer theorem.
\item Compute \(\beta\) and q exactly; promotion requires \(\beta=0\), subpolynomial target queries, or a genuinely nonblackbox use of the SAT circuit.
\item State whether the transfer excludes every polynomial SAT exponent or only a bounded interval \(c<c0\).
\item Record relativization, natural-proofs, algebrization, nonuniformity, and advice-leakage status for the specific transfer theorem.
\item Control gate: name the source carrier or language, target class, and circuit or resource model before claiming progress.
\item Control gate: derive the exact transfer inequality from \(\mathrm{SAT}\in \mathrm{SIZE}(n^{c})\), NP subset P/poly, or the stated target assumption to the source-carrier upper bound.
\end{itemize}

\paragraph{Stop criteria.}
\begin{itemize}
\item Every named carrier in the branch has \(\beta>0\) and yields only a bounded target-exponent exclusion.
\item The branch cannot state a source lower-bound model, target class, and transfer inequality in the same notation.
\item The proposed escape is only a black-box substitution restatement with no new carrier or nonblackbox mechanism.
\item The barrier audit shows the needed source lower bound is just as hard as the original P vs NP/P/poly target.
\end{itemize}

\paragraph{Do not stop if.}
\begin{itemize}
\item The branch has not yet computed \(\beta\), q, and the induced source upper-bound exponent under \(\mathrm{SAT}\in \mathrm{SIZE}(n^{c})\).
\item A subpolynomial-query or nonblackbox carrier remains available and has not been audited.
\item The latest failure sharpened the threshold-collapse lemma into a formal no-go statement.
\item The route can still isolate advice leakage, nonuniformity, naturalness, or algebrization as the exact obstruction.
\end{itemize}

\paragraph{Required skills.}
\begin{itemize}
\item abducer-failure-signal-cartographer
\item contradiction-route-auditor
\item entailment-gap-auditor
\item kb-deduce
\item kb-end-to-end-proof-critic
\item kb-hypothesis-refutation-loop
\item kb-induce
\item long-horizon-route-portfolio-manager
\item main-problem-alignment-auditor
\item next-experiment-queue-ranker
\item paper-contribution-readiness-judge
\item pvsnp-barrier-specificity-reviewer
\item pvsnp-conceptual-pressure-localizer
\item pvsnp-impossibility-theorem-forger
\item pvsnp-inference-quality-controller
\item pvsnp-proof-derivation-builder
\item pvsnp-qualitative-complexity-reasoner
\item pvsnp-quantifier-uniformity-reviewer
\item pvsnp-reduction-consistency-reviewer
\item pvsnp-symbolic-complexity-auditor
\end{itemize}

\subsection*{\#47 cycle-0062: Hardness-Magnification Bridge Realization}

\paragraph{Analysis.} Instantiate the random-affine generator with L=(log N)\textasciicircum{}2 and q in ((1+epsilon)/2,1+epsilon); compute \(r_{N}\) for the induced sparse incidence matrix and test whether the rank gap forces naturalness or low-rank falsification.

\paragraph{Why this is a good target.} It directly tests the current black-box substitution cap by computing \(\beta\), q, and the induced target-exponent window. It checks whether padding/projection transfers really magnify hardness or merely move the proof debt into sparse lower bounds. It keeps the barrier and model-boundary audit attached to the mathematical transfer rather than as after-the-fact commentary. It is executable as a bounded complexity calculation, CAS-backed exponent audit, or formal transfer check.

\paragraph{Risk.} Low drift risk: the item is framed as an explicit P vs NP transfer audit.

\paragraph{Branch control.} The item has high payoff and a concrete decision test; stopping now would be premature.

\paragraph{Premature-stop guard.} Do not stop before deriving the transfer inequality from \(\mathrm{SAT}\in \mathrm{SIZE}(n^{c})\) to the carrier upper bound, identifying whether \(\beta=0\), \(\beta>0\), or the argument is genuinely nonblackbox, and recording the barrier status of the proposed escape.

\paragraph{Payoff criteria.}
\begin{itemize}
\item For the proof-complexity branch, separate fixed-family proof hardness from language hardness by naming the proof system, shortcut algorithm, and semantic lower-bound target.
\item Promotion requires a proof-system lower bound that survives the Gaussian-elimination or recognition shortcut; otherwise the result is a restricted-proof-system separation only.
\item For this route, state the exact barrier premise it uses or escapes: relativizing, natural, algebrizing, nonuniform, or advice-leaking; promotion requires the escape to be attached to the specific transfer theorem.
\item Compute \(\beta\) and q exactly; promotion requires \(\beta=0\), subpolynomial target queries, or a genuinely nonblackbox use of the SAT circuit.
\item State whether the transfer excludes every polynomial SAT exponent or only a bounded interval \(c<c0\).
\item Record relativization, natural-proofs, algebrization, nonuniformity, and advice-leakage status for the specific transfer theorem.
\end{itemize}

\paragraph{Stop criteria.}
\begin{itemize}
\item Every named carrier in the branch has \(\beta>0\) and yields only a bounded target-exponent exclusion.
\item The branch cannot state a source lower-bound model, target class, and transfer inequality in the same notation.
\item The proposed escape is only a black-box substitution restatement with no new carrier or nonblackbox mechanism.
\item The barrier audit shows the needed source lower bound is just as hard as the original P vs NP/P/poly target.
\end{itemize}

\paragraph{Do not stop if.}
\begin{itemize}
\item The branch has not yet computed \(\beta\), q, and the induced source upper-bound exponent under \(\mathrm{SAT}\in \mathrm{SIZE}(n^{c})\).
\item A subpolynomial-query or nonblackbox carrier remains available and has not been audited.
\item The latest failure sharpened the threshold-collapse lemma into a formal no-go statement.
\item The route can still isolate advice leakage, nonuniformity, naturalness, or algebrization as the exact obstruction.
\end{itemize}

\paragraph{Required skills.}
\begin{itemize}
\item abducer-failure-signal-cartographer
\item contradiction-route-auditor
\item entailment-gap-auditor
\item kb-deduce
\item kb-end-to-end-proof-critic
\item kb-hypothesis-refutation-loop
\item kb-induce
\item long-horizon-route-portfolio-manager
\item main-problem-alignment-auditor
\item next-experiment-queue-ranker
\item paper-contribution-readiness-judge
\item pvsnp-barrier-specificity-reviewer
\item pvsnp-conceptual-pressure-localizer
\item pvsnp-impossibility-theorem-forger
\item pvsnp-inference-quality-controller
\item pvsnp-proof-derivation-builder
\item pvsnp-qualitative-complexity-reasoner
\item pvsnp-quantifier-uniformity-reviewer
\item pvsnp-reduction-consistency-reviewer
\item pvsnp-symbolic-complexity-auditor
\end{itemize}

\subsection*{\#48 cycle-0061: Hardness-Magnification Bridge Realization}

\paragraph{Analysis.} Instantiate a candidate generator and compute \(U_{N},A_{N},r_{N},E_{N}\) for its SAT-preserving quotient codebook; if the cost is subthreshold, quarantine that generator, otherwise formulate the exact extraction theorem it would require.

\paragraph{Why this is a good target.} It directly tests the current black-box substitution cap by computing \(\beta\), q, and the induced target-exponent window. It checks whether padding/projection transfers really magnify hardness or merely move the proof debt into sparse lower bounds. It keeps the barrier and model-boundary audit attached to the mathematical transfer rather than as after-the-fact commentary. It is executable as a bounded complexity calculation, CAS-backed exponent audit, or formal transfer check.

\paragraph{Risk.} Low drift risk: the item is framed as an explicit P vs NP transfer audit.

\paragraph{Branch control.} The item has high payoff and a concrete decision test; stopping now would be premature.

\paragraph{Premature-stop guard.} Do not stop before deriving the transfer inequality from \(\mathrm{SAT}\in \mathrm{SIZE}(n^{c})\) to the carrier upper bound, identifying whether \(\beta=0\), \(\beta>0\), or the argument is genuinely nonblackbox, and recording the barrier status of the proposed escape.

\paragraph{Payoff criteria.}
\begin{itemize}
\item For the proof-complexity branch, separate fixed-family proof hardness from language hardness by naming the proof system, shortcut algorithm, and semantic lower-bound target.
\item Promotion requires a proof-system lower bound that survives the Gaussian-elimination or recognition shortcut; otherwise the result is a restricted-proof-system separation only.
\item For this route, state the exact barrier premise it uses or escapes: relativizing, natural, algebrizing, nonuniform, or advice-leaking; promotion requires the escape to be attached to the specific transfer theorem.
\item Compute \(\beta\) and q exactly; promotion requires \(\beta=0\), subpolynomial target queries, or a genuinely nonblackbox use of the SAT circuit.
\item State whether the transfer excludes every polynomial SAT exponent or only a bounded interval \(c<c0\).
\item Record relativization, natural-proofs, algebrization, nonuniformity, and advice-leakage status for the specific transfer theorem.
\end{itemize}

\paragraph{Stop criteria.}
\begin{itemize}
\item Every named carrier in the branch has \(\beta>0\) and yields only a bounded target-exponent exclusion.
\item The branch cannot state a source lower-bound model, target class, and transfer inequality in the same notation.
\item The proposed escape is only a black-box substitution restatement with no new carrier or nonblackbox mechanism.
\item The barrier audit shows the needed source lower bound is just as hard as the original P vs NP/P/poly target.
\end{itemize}

\paragraph{Do not stop if.}
\begin{itemize}
\item The branch has not yet computed \(\beta\), q, and the induced source upper-bound exponent under \(\mathrm{SAT}\in \mathrm{SIZE}(n^{c})\).
\item A subpolynomial-query or nonblackbox carrier remains available and has not been audited.
\item The latest failure sharpened the threshold-collapse lemma into a formal no-go statement.
\item The route can still isolate advice leakage, nonuniformity, naturalness, or algebrization as the exact obstruction.
\end{itemize}

\paragraph{Required skills.}
\begin{itemize}
\item abducer-failure-signal-cartographer
\item contradiction-route-auditor
\item entailment-gap-auditor
\item kb-deduce
\item kb-end-to-end-proof-critic
\item kb-hypothesis-refutation-loop
\item kb-induce
\item long-horizon-route-portfolio-manager
\item main-problem-alignment-auditor
\item next-experiment-queue-ranker
\item paper-contribution-readiness-judge
\item pvsnp-barrier-specificity-reviewer
\item pvsnp-conceptual-pressure-localizer
\item pvsnp-impossibility-theorem-forger
\item pvsnp-inference-quality-controller
\item pvsnp-proof-derivation-builder
\item pvsnp-qualitative-complexity-reasoner
\item pvsnp-quantifier-uniformity-reviewer
\item pvsnp-reduction-consistency-reviewer
\item pvsnp-symbolic-complexity-auditor
\end{itemize}

\subsection*{\#49 cycle-0060: Constructive Label-Free Sketch Gate for XOR Magnification}

\paragraph{Analysis.} Define the first concrete compiler class \(T^{inc-lin}\): \(T_{N}(x)=B_{N}\) a\_N(x) for explicit SAT-preserving quotient matrices \(B_{N}\) of row sparsity polylog N, then compute the minimum \(S_{T}+A_{N}\) \(r_{N}\) for representing \(H_{N}\) by sparse \(\mathrm{P}_{N}\) in the high-q, \(d=2\) regime.

\paragraph{Why this is a good target.} It directly tests the current black-box substitution cap by computing \(\beta\), q, and the induced target-exponent window. It checks whether padding/projection transfers really magnify hardness or merely move the proof debt into sparse lower bounds. It keeps the barrier and model-boundary audit attached to the mathematical transfer rather than as after-the-fact commentary. It is executable as a bounded complexity calculation, CAS-backed exponent audit, or formal transfer check.

\paragraph{Risk.} Low drift risk: the item is framed as an explicit P vs NP transfer audit.

\paragraph{Branch control.} The item has high payoff and a concrete decision test; stopping now would be premature.

\paragraph{Premature-stop guard.} Do not stop before deriving the transfer inequality from \(\mathrm{SAT}\in \mathrm{SIZE}(n^{c})\) to the carrier upper bound, identifying whether \(\beta=0\), \(\beta>0\), or the argument is genuinely nonblackbox, and recording the barrier status of the proposed escape.

\paragraph{Payoff criteria.}
\begin{itemize}
\item For the high-q generated-OR carrier, compute the exact leakage exponent of the remaining witness-cover or unary-residual mechanism and compare it with 1+epsilon-q.
\item Promotion requires \(q+a>1+\epsilon\) after every projection, padding, and witness-cover leakage term is accounted for; otherwise the branch becomes a bounded-window or no-go result.
\item For the proof-complexity branch, separate fixed-family proof hardness from language hardness by naming the proof system, shortcut algorithm, and semantic lower-bound target.
\item Promotion requires a proof-system lower bound that survives the Gaussian-elimination or recognition shortcut; otherwise the result is a restricted-proof-system separation only.
\item For this route, state the exact barrier premise it uses or escapes: relativizing, natural, algebrizing, nonuniform, or advice-leaking; promotion requires the escape to be attached to the specific transfer theorem.
\item Compute \(\beta\) and q exactly; promotion requires \(\beta=0\), subpolynomial target queries, or a genuinely nonblackbox use of the SAT circuit.
\end{itemize}

\paragraph{Stop criteria.}
\begin{itemize}
\item Every named carrier in the branch has \(\beta>0\) and yields only a bounded target-exponent exclusion.
\item The branch cannot state a source lower-bound model, target class, and transfer inequality in the same notation.
\item The proposed escape is only a black-box substitution restatement with no new carrier or nonblackbox mechanism.
\item The barrier audit shows the needed source lower bound is just as hard as the original P vs NP/P/poly target.
\end{itemize}

\paragraph{Do not stop if.}
\begin{itemize}
\item The branch has not yet computed \(\beta\), q, and the induced source upper-bound exponent under \(\mathrm{SAT}\in \mathrm{SIZE}(n^{c})\).
\item A subpolynomial-query or nonblackbox carrier remains available and has not been audited.
\item The latest failure sharpened the threshold-collapse lemma into a formal no-go statement.
\item The route can still isolate advice leakage, nonuniformity, naturalness, or algebrization as the exact obstruction.
\end{itemize}

\paragraph{Required skills.}
\begin{itemize}
\item abducer-failure-signal-cartographer
\item contradiction-route-auditor
\item entailment-gap-auditor
\item kb-deduce
\item kb-end-to-end-proof-critic
\item kb-hypothesis-refutation-loop
\item kb-induce
\item long-horizon-route-portfolio-manager
\item main-problem-alignment-auditor
\item next-experiment-queue-ranker
\item paper-contribution-readiness-judge
\item pvsnp-barrier-specificity-reviewer
\item pvsnp-conceptual-pressure-localizer
\item pvsnp-impossibility-theorem-forger
\item pvsnp-inference-quality-controller
\item pvsnp-proof-derivation-builder
\item pvsnp-qualitative-complexity-reasoner
\item pvsnp-quantifier-uniformity-reviewer
\item pvsnp-reduction-consistency-reviewer
\item pvsnp-symbolic-complexity-auditor
\end{itemize}

\subsection*{\#50 cycle-0058: Hardness-magnification bridge via XOR carriers and sparse quotient-polynomial falsifiers}

\paragraph{Analysis.} Formalize the weakest nontrivial extraction candidate: restrict arbitrary \(H_{N}\) circuits to circuits whose wires are invariant under SAT-preserving quotient symmetries, and test whether such invariant circuits collapse to sparse quotient-polynomial sketches with subthreshold loss.

\paragraph{Why this is a good target.} It directly tests the current black-box substitution cap by computing \(\beta\), q, and the induced target-exponent window. It checks whether padding/projection transfers really magnify hardness or merely move the proof debt into sparse lower bounds. It keeps the barrier and model-boundary audit attached to the mathematical transfer rather than as after-the-fact commentary. It is executable as a bounded complexity calculation, CAS-backed exponent audit, or formal transfer check.

\paragraph{Risk.} Low drift risk: the item is framed as an explicit P vs NP transfer audit.

\paragraph{Branch control.} The item has high payoff and a concrete decision test; stopping now would be premature.

\paragraph{Premature-stop guard.} Do not stop before deriving the exponent transfer c \(\to \) c/k or its growth-padding analogue, and before stating whether the result excludes all polynomial SAT exponents or only a bounded exponent interval.

\paragraph{Payoff criteria.}
\begin{itemize}
\item For the proof-complexity branch, separate fixed-family proof hardness from language hardness by naming the proof system, shortcut algorithm, and semantic lower-bound target.
\item Promotion requires a proof-system lower bound that survives the Gaussian-elimination or recognition shortcut; otherwise the result is a restricted-proof-system separation only.
\item For this route, state the exact barrier premise it uses or escapes: relativizing, natural, algebrizing, nonuniform, or advice-leaking; promotion requires the escape to be attached to the specific transfer theorem.
\item Compute \(\beta\) and q exactly; promotion requires \(\beta=0\), subpolynomial target queries, or a genuinely nonblackbox use of the SAT circuit.
\item State whether the transfer excludes every polynomial SAT exponent or only a bounded interval \(c<c0\).
\item Record relativization, natural-proofs, algebrization, nonuniformity, and advice-leakage status for the specific transfer theorem.
\end{itemize}

\paragraph{Stop criteria.}
\begin{itemize}
\item The branch only proves a finite target-exponent lower bound and cannot scale to \(\mathrm{NP}\not\subset \mathrm{P}/\mathrm{poly}\).
\item The route requires sparse valid-length lower bounds that are not weaker than the original target problem.
\item The projection carrier has no nonprojection mechanism and no subpolynomial embedding repair.
\end{itemize}

\paragraph{Do not stop if.}
\begin{itemize}
\item The quantifier pattern has not been repaired between fixed-k, unbounded-k, and single growth-padding carriers.
\item The projection-sandwich inequality has not yet been written as an explicit circuit-size transfer.
\item The sparse valid-length lower-bound obligation has not been isolated as the real source debt.
\end{itemize}

\paragraph{Required skills.}
\begin{itemize}
\item abducer-failure-signal-cartographer
\item contradiction-route-auditor
\item entailment-gap-auditor
\item kb-deduce
\item kb-end-to-end-proof-critic
\item kb-hypothesis-refutation-loop
\item kb-induce
\item long-horizon-route-portfolio-manager
\item main-problem-alignment-auditor
\item next-experiment-queue-ranker
\item paper-contribution-readiness-judge
\item pvsnp-barrier-specificity-reviewer
\item pvsnp-conceptual-pressure-localizer
\item pvsnp-impossibility-theorem-forger
\item pvsnp-inference-quality-controller
\item pvsnp-proof-derivation-builder
\item pvsnp-qualitative-complexity-reasoner
\item pvsnp-quantifier-uniformity-reviewer
\item pvsnp-reduction-consistency-reviewer
\item pvsnp-symbolic-complexity-auditor
\end{itemize}

\subsection*{\#51 cycle-0053: Hardness-Magnification Bridge Realization: XOR carrier affine-incidence audit}

\paragraph{Analysis.} Instantiate \(G_{N,\theta}\) with linear, pairwise-independent, and extractor-style query maps; compute K, coefficient-circuit cost \(A_{N}\), duplicate parity cancellations, and whether any case survives the \(N^{1+\epsilon}\) affine-sketch falsifier threshold.

\paragraph{Why this is a good target.} It directly tests the current black-box substitution cap by computing \(\beta\), q, and the induced target-exponent window. It checks whether padding/projection transfers really magnify hardness or merely move the proof debt into sparse lower bounds. It keeps the barrier and model-boundary audit attached to the mathematical transfer rather than as after-the-fact commentary. It is executable as a bounded complexity calculation, CAS-backed exponent audit, or formal transfer check.

\paragraph{Risk.} Low drift risk: the item is framed as an explicit P vs NP transfer audit.

\paragraph{Branch control.} The item has high payoff and a concrete decision test; stopping now would be premature.

\paragraph{Premature-stop guard.} Do not stop before deriving the exponent transfer c \(\to \) c/k or its growth-padding analogue, and before stating whether the result excludes all polynomial SAT exponents or only a bounded exponent interval.

\paragraph{Payoff criteria.}
\begin{itemize}
\item For the proof-complexity branch, separate fixed-family proof hardness from language hardness by naming the proof system, shortcut algorithm, and semantic lower-bound target.
\item Promotion requires a proof-system lower bound that survives the Gaussian-elimination or recognition shortcut; otherwise the result is a restricted-proof-system separation only.
\item For this route, state the exact barrier premise it uses or escapes: relativizing, natural, algebrizing, nonuniform, or advice-leaking; promotion requires the escape to be attached to the specific transfer theorem.
\item Compute \(\beta\) and q exactly; promotion requires \(\beta=0\), subpolynomial target queries, or a genuinely nonblackbox use of the SAT circuit.
\item State whether the transfer excludes every polynomial SAT exponent or only a bounded interval \(c<c0\).
\item Record relativization, natural-proofs, algebrization, nonuniformity, and advice-leakage status for the specific transfer theorem.
\end{itemize}

\paragraph{Stop criteria.}
\begin{itemize}
\item The branch only proves a finite target-exponent lower bound and cannot scale to \(\mathrm{NP}\not\subset \mathrm{P}/\mathrm{poly}\).
\item The route requires sparse valid-length lower bounds that are not weaker than the original target problem.
\item The projection carrier has no nonprojection mechanism and no subpolynomial embedding repair.
\end{itemize}

\paragraph{Do not stop if.}
\begin{itemize}
\item The quantifier pattern has not been repaired between fixed-k, unbounded-k, and single growth-padding carriers.
\item The projection-sandwich inequality has not yet been written as an explicit circuit-size transfer.
\item The sparse valid-length lower-bound obligation has not been isolated as the real source debt.
\end{itemize}

\paragraph{Required skills.}
\begin{itemize}
\item abducer-failure-signal-cartographer
\item contradiction-route-auditor
\item entailment-gap-auditor
\item kb-deduce
\item kb-end-to-end-proof-critic
\item kb-hypothesis-refutation-loop
\item kb-induce
\item long-horizon-route-portfolio-manager
\item main-problem-alignment-auditor
\item next-experiment-queue-ranker
\item paper-contribution-readiness-judge
\item pvsnp-barrier-specificity-reviewer
\item pvsnp-conceptual-pressure-localizer
\item pvsnp-impossibility-theorem-forger
\item pvsnp-inference-quality-controller
\item pvsnp-proof-derivation-builder
\item pvsnp-qualitative-complexity-reasoner
\item pvsnp-quantifier-uniformity-reviewer
\item pvsnp-reduction-consistency-reviewer
\item pvsnp-symbolic-complexity-auditor
\end{itemize}

\subsection*{\#52 cycle-0052: Hardness-Magnification Bridge Realization: XOR Carrier Repair}

\paragraph{Analysis.} In the ideal random-label model for \(G_{N}\), count subthreshold affine sketches of size \(N^\alpha\) and estimate the tau,alpha,error threshold at which a union bound could certify parity-sketch resistance without pretending to prove the actual explicit carrier lower bound.

\paragraph{Why this is a good target.} It directly tests the current black-box substitution cap by computing \(\beta\), q, and the induced target-exponent window. It checks whether padding/projection transfers really magnify hardness or merely move the proof debt into sparse lower bounds. It keeps the barrier and model-boundary audit attached to the mathematical transfer rather than as after-the-fact commentary. It is executable as a bounded complexity calculation, CAS-backed exponent audit, or formal transfer check.

\paragraph{Risk.} Low drift risk: the item is framed as an explicit P vs NP transfer audit.

\paragraph{Branch control.} The item has high payoff and a concrete decision test; stopping now would be premature.

\paragraph{Premature-stop guard.} Do not stop before deriving the exponent transfer c \(\to \) c/k or its growth-padding analogue, and before stating whether the result excludes all polynomial SAT exponents or only a bounded exponent interval.

\paragraph{Payoff criteria.}
\begin{itemize}
\item For the proof-complexity branch, separate fixed-family proof hardness from language hardness by naming the proof system, shortcut algorithm, and semantic lower-bound target.
\item Promotion requires a proof-system lower bound that survives the Gaussian-elimination or recognition shortcut; otherwise the result is a restricted-proof-system separation only.
\item For this route, state the exact barrier premise it uses or escapes: relativizing, natural, algebrizing, nonuniform, or advice-leaking; promotion requires the escape to be attached to the specific transfer theorem.
\item Compute \(\beta\) and q exactly; promotion requires \(\beta=0\), subpolynomial target queries, or a genuinely nonblackbox use of the SAT circuit.
\item State whether the transfer excludes every polynomial SAT exponent or only a bounded interval \(c<c0\).
\item Record relativization, natural-proofs, algebrization, nonuniformity, and advice-leakage status for the specific transfer theorem.
\end{itemize}

\paragraph{Stop criteria.}
\begin{itemize}
\item The branch only proves a finite target-exponent lower bound and cannot scale to \(\mathrm{NP}\not\subset \mathrm{P}/\mathrm{poly}\).
\item The route requires sparse valid-length lower bounds that are not weaker than the original target problem.
\item The projection carrier has no nonprojection mechanism and no subpolynomial embedding repair.
\end{itemize}

\paragraph{Do not stop if.}
\begin{itemize}
\item The quantifier pattern has not been repaired between fixed-k, unbounded-k, and single growth-padding carriers.
\item The projection-sandwich inequality has not yet been written as an explicit circuit-size transfer.
\item The sparse valid-length lower-bound obligation has not been isolated as the real source debt.
\end{itemize}

\paragraph{Required skills.}
\begin{itemize}
\item abducer-failure-signal-cartographer
\item contradiction-route-auditor
\item entailment-gap-auditor
\item kb-deduce
\item kb-end-to-end-proof-critic
\item kb-hypothesis-refutation-loop
\item kb-induce
\item long-horizon-route-portfolio-manager
\item main-problem-alignment-auditor
\item next-experiment-queue-ranker
\item paper-contribution-readiness-judge
\item pvsnp-barrier-specificity-reviewer
\item pvsnp-conceptual-pressure-localizer
\item pvsnp-impossibility-theorem-forger
\item pvsnp-inference-quality-controller
\item pvsnp-proof-derivation-builder
\item pvsnp-qualitative-complexity-reasoner
\item pvsnp-quantifier-uniformity-reviewer
\item pvsnp-reduction-consistency-reviewer
\item pvsnp-symbolic-complexity-auditor
\end{itemize}

\subsection*{\#53 cycle-0039: High-q generated-OR carrier: partial-certificate cover gate}

\paragraph{Analysis.} Formalize \(tau_{N}^{term}\) as a set-cover problem over generated CNF ranges and test whether random-looking high-q generators can force tau\_N\textasciicircum{}\{term\}>N\textasciicircum{}\{1+epsilon\} without making \(G_{N}\) nonexplicit or violating \(\rho<1+\epsilon\).

\paragraph{Why this is a good target.} It directly tests the current black-box substitution cap by computing \(\beta\), q, and the induced target-exponent window. It checks whether padding/projection transfers really magnify hardness or merely move the proof debt into sparse lower bounds. It keeps the barrier and model-boundary audit attached to the mathematical transfer rather than as after-the-fact commentary. It is executable as a bounded complexity calculation, CAS-backed exponent audit, or formal transfer check.

\paragraph{Risk.} Low drift risk: the item is framed as an explicit P vs NP transfer audit.

\paragraph{Branch control.} The item has high payoff and a concrete decision test; stopping now would be premature.

\paragraph{Premature-stop guard.} Do not stop before deriving the exponent transfer c \(\to \) c/k or its growth-padding analogue, and before stating whether the result excludes all polynomial SAT exponents or only a bounded exponent interval.

\paragraph{Payoff criteria.}
\begin{itemize}
\item For the certificate-cover branch, compute K,E,L and the induced circuit upper bound C(H\_N)<=C(G\_N)+O(K(s+L\textasciicircum{}\{O(1)\})+EN); promotion requires every sound local cover to stay above the \(N^{1+\epsilon-o(1)}\) source threshold.
\item Falsify the branch if Horn, 2-SAT, bounded-width CNF, or bounded-backdoor emitted formulas give subthreshold satisfying-term or sound-local-certificate covers.
\item For the high-q generated-OR carrier, compute the exact leakage exponent of the remaining witness-cover or unary-residual mechanism and compare it with 1+epsilon-q.
\item Promotion requires \(q+a>1+\epsilon\) after every projection, padding, and witness-cover leakage term is accounted for; otherwise the branch becomes a bounded-window or no-go result.
\item For this route, state the exact barrier premise it uses or escapes: relativizing, natural, algebrizing, nonuniform, or advice-leaking; promotion requires the escape to be attached to the specific transfer theorem.
\item Compute \(\beta\) and q exactly; promotion requires \(\beta=0\), subpolynomial target queries, or a genuinely nonblackbox use of the SAT circuit.
\end{itemize}

\paragraph{Stop criteria.}
\begin{itemize}
\item The branch only proves a finite target-exponent lower bound and cannot scale to \(\mathrm{NP}\not\subset \mathrm{P}/\mathrm{poly}\).
\item The route requires sparse valid-length lower bounds that are not weaker than the original target problem.
\item The projection carrier has no nonprojection mechanism and no subpolynomial embedding repair.
\end{itemize}

\paragraph{Do not stop if.}
\begin{itemize}
\item The quantifier pattern has not been repaired between fixed-k, unbounded-k, and single growth-padding carriers.
\item The projection-sandwich inequality has not yet been written as an explicit circuit-size transfer.
\item The sparse valid-length lower-bound obligation has not been isolated as the real source debt.
\end{itemize}

\paragraph{Required skills.}
\begin{itemize}
\item abducer-failure-signal-cartographer
\item contradiction-route-auditor
\item entailment-gap-auditor
\item kb-deduce
\item kb-end-to-end-proof-critic
\item kb-hypothesis-refutation-loop
\item kb-induce
\item long-horizon-route-portfolio-manager
\item main-problem-alignment-auditor
\item next-experiment-queue-ranker
\item paper-contribution-readiness-judge
\item pvsnp-barrier-specificity-reviewer
\item pvsnp-conceptual-pressure-localizer
\item pvsnp-impossibility-theorem-forger
\item pvsnp-inference-quality-controller
\item pvsnp-proof-derivation-builder
\item pvsnp-qualitative-complexity-reasoner
\item pvsnp-quantifier-uniformity-reviewer
\item pvsnp-reduction-consistency-reviewer
\item pvsnp-symbolic-complexity-auditor
\end{itemize}

\subsection*{\#54 cycle-0038: Hardness-Magnification Bridge Realization: realized-witness projection gate}

\paragraph{Analysis.} Instantiate the planted-pool thought experiment and compute \(W_{N}^*\) for affine, low-degree, and pseudorandom assignment-pool generators; check whether any preserves \(q+a>1+\epsilon\) without introducing natural-proof or advice-leakage collapse.

\paragraph{Why this is a good target.} It directly tests the current black-box substitution cap by computing \(\beta\), q, and the induced target-exponent window. It checks whether padding/projection transfers really magnify hardness or merely move the proof debt into sparse lower bounds. It keeps the barrier and model-boundary audit attached to the mathematical transfer rather than as after-the-fact commentary. It is executable as a bounded complexity calculation, CAS-backed exponent audit, or formal transfer check.

\paragraph{Risk.} Low drift risk: the item is framed as an explicit P vs NP transfer audit.

\paragraph{Branch control.} The item has high payoff and a concrete decision test; stopping now would be premature.

\paragraph{Premature-stop guard.} Do not stop before deriving the transfer inequality from \(\mathrm{SAT}\in \mathrm{SIZE}(n^{c})\) to the carrier upper bound, identifying whether \(\beta=0\), \(\beta>0\), or the argument is genuinely nonblackbox, and recording the barrier status of the proposed escape.

\paragraph{Payoff criteria.}
\begin{itemize}
\item For the high-q generated-OR carrier, compute the exact leakage exponent of the remaining witness-cover or unary-residual mechanism and compare it with 1+epsilon-q.
\item Promotion requires \(q+a>1+\epsilon\) after every projection, padding, and witness-cover leakage term is accounted for; otherwise the branch becomes a bounded-window or no-go result.
\item For this route, state the exact barrier premise it uses or escapes: relativizing, natural, algebrizing, nonuniform, or advice-leaking; promotion requires the escape to be attached to the specific transfer theorem.
\item Compute \(\beta\) and q exactly; promotion requires \(\beta=0\), subpolynomial target queries, or a genuinely nonblackbox use of the SAT circuit.
\item State whether the transfer excludes every polynomial SAT exponent or only a bounded interval \(c<c0\).
\item Record relativization, natural-proofs, algebrization, nonuniformity, and advice-leakage status for the specific transfer theorem.
\end{itemize}

\paragraph{Stop criteria.}
\begin{itemize}
\item Every named carrier in the branch has \(\beta>0\) and yields only a bounded target-exponent exclusion.
\item The branch cannot state a source lower-bound model, target class, and transfer inequality in the same notation.
\item The proposed escape is only a black-box substitution restatement with no new carrier or nonblackbox mechanism.
\item The barrier audit shows the needed source lower bound is just as hard as the original P vs NP/P/poly target.
\end{itemize}

\paragraph{Do not stop if.}
\begin{itemize}
\item The branch has not yet computed \(\beta\), q, and the induced source upper-bound exponent under \(\mathrm{SAT}\in \mathrm{SIZE}(n^{c})\).
\item A subpolynomial-query or nonblackbox carrier remains available and has not been audited.
\item The latest failure sharpened the threshold-collapse lemma into a formal no-go statement.
\item The route can still isolate advice leakage, nonuniformity, naturalness, or algebrization as the exact obstruction.
\end{itemize}

\paragraph{Required skills.}
\begin{itemize}
\item abducer-failure-signal-cartographer
\item contradiction-route-auditor
\item entailment-gap-auditor
\item kb-deduce
\item kb-end-to-end-proof-critic
\item kb-hypothesis-refutation-loop
\item kb-induce
\item long-horizon-route-portfolio-manager
\item main-problem-alignment-auditor
\item next-experiment-queue-ranker
\item paper-contribution-readiness-judge
\item pvsnp-barrier-specificity-reviewer
\item pvsnp-conceptual-pressure-localizer
\item pvsnp-impossibility-theorem-forger
\item pvsnp-inference-quality-controller
\item pvsnp-proof-derivation-builder
\item pvsnp-qualitative-complexity-reasoner
\item pvsnp-quantifier-uniformity-reviewer
\item pvsnp-reduction-consistency-reviewer
\item pvsnp-symbolic-complexity-auditor
\end{itemize}

\subsection*{\#55 cycle-0033: Hardness-Magnification Bridge Realization}

\paragraph{Analysis.} Formalize the restricted case where B is a depth-2 threshold or DNF over \(t=N^\tau\) local features, compute the exact q+alpha+beta leakage threshold, and search for an explicit low-residual construction on expander-walk generated SAT bags.

\paragraph{Why this is a good target.} It directly tests the current black-box substitution cap by computing \(\beta\), q, and the induced target-exponent window. It checks whether padding/projection transfers really magnify hardness or merely move the proof debt into sparse lower bounds. It keeps the barrier and model-boundary audit attached to the mathematical transfer rather than as after-the-fact commentary. It is executable as a bounded complexity calculation, CAS-backed exponent audit, or formal transfer check.

\paragraph{Risk.} Low drift risk: the item is framed as an explicit P vs NP transfer audit.

\paragraph{Branch control.} The item has high payoff and a concrete decision test; stopping now would be premature.

\paragraph{Premature-stop guard.} Do not stop before deriving the transfer inequality from \(\mathrm{SAT}\in \mathrm{SIZE}(n^{c})\) to the carrier upper bound, identifying whether \(\beta=0\), \(\beta>0\), or the argument is genuinely nonblackbox, and recording the barrier status of the proposed escape.

\paragraph{Payoff criteria.}
\begin{itemize}
\item For the high-q generated-OR carrier, compute the exact leakage exponent of the remaining witness-cover or unary-residual mechanism and compare it with 1+epsilon-q.
\item Promotion requires \(q+a>1+\epsilon\) after every projection, padding, and witness-cover leakage term is accounted for; otherwise the branch becomes a bounded-window or no-go result.
\item For the proof-complexity branch, separate fixed-family proof hardness from language hardness by naming the proof system, shortcut algorithm, and semantic lower-bound target.
\item Promotion requires a proof-system lower bound that survives the Gaussian-elimination or recognition shortcut; otherwise the result is a restricted-proof-system separation only.
\item For this route, state the exact barrier premise it uses or escapes: relativizing, natural, algebrizing, nonuniform, or advice-leaking; promotion requires the escape to be attached to the specific transfer theorem.
\item Compute \(\beta\) and q exactly; promotion requires \(\beta=0\), subpolynomial target queries, or a genuinely nonblackbox use of the SAT circuit.
\end{itemize}

\paragraph{Stop criteria.}
\begin{itemize}
\item Every named carrier in the branch has \(\beta>0\) and yields only a bounded target-exponent exclusion.
\item The branch cannot state a source lower-bound model, target class, and transfer inequality in the same notation.
\item The proposed escape is only a black-box substitution restatement with no new carrier or nonblackbox mechanism.
\item The barrier audit shows the needed source lower bound is just as hard as the original P vs NP/P/poly target.
\end{itemize}

\paragraph{Do not stop if.}
\begin{itemize}
\item The branch has not yet computed \(\beta\), q, and the induced source upper-bound exponent under \(\mathrm{SAT}\in \mathrm{SIZE}(n^{c})\).
\item A subpolynomial-query or nonblackbox carrier remains available and has not been audited.
\item The latest failure sharpened the threshold-collapse lemma into a formal no-go statement.
\item The route can still isolate advice leakage, nonuniformity, naturalness, or algebrization as the exact obstruction.
\end{itemize}

\paragraph{Required skills.}
\begin{itemize}
\item abducer-failure-signal-cartographer
\item contradiction-route-auditor
\item entailment-gap-auditor
\item kb-deduce
\item kb-end-to-end-proof-critic
\item kb-hypothesis-refutation-loop
\item kb-induce
\item long-horizon-route-portfolio-manager
\item main-problem-alignment-auditor
\item next-experiment-queue-ranker
\item paper-contribution-readiness-judge
\item pvsnp-barrier-specificity-reviewer
\item pvsnp-conceptual-pressure-localizer
\item pvsnp-impossibility-theorem-forger
\item pvsnp-inference-quality-controller
\item pvsnp-proof-derivation-builder
\item pvsnp-qualitative-complexity-reasoner
\item pvsnp-quantifier-uniformity-reviewer
\item pvsnp-reduction-consistency-reviewer
\item pvsnp-symbolic-complexity-auditor
\end{itemize}

\subsection*{\#56 cycle-0019: Hardness-Magnification Bridge Realization: semantic range leakage gate}

\paragraph{Analysis.} Instantiate \(G_{N}\) with a concrete expander or extractor-based 3-CNF generator, compute \(R_{N,j}\) structure, and test whether SAT on those positioned ranges collapses to a known easy subclass or requires a new semantic lower-bound certificate.

\paragraph{Why this is a good target.} It directly tests the current black-box substitution cap by computing \(\beta\), q, and the induced target-exponent window. It checks whether padding/projection transfers really magnify hardness or merely move the proof debt into sparse lower bounds. It keeps the barrier and model-boundary audit attached to the mathematical transfer rather than as after-the-fact commentary. It is executable as a bounded complexity calculation, CAS-backed exponent audit, or formal transfer check.

\paragraph{Risk.} Low drift risk: the item is framed as an explicit P vs NP transfer audit.

\paragraph{Branch control.} The item has high payoff and a concrete decision test; stopping now would be premature.

\paragraph{Premature-stop guard.} Do not stop before deriving the transfer inequality from \(\mathrm{SAT}\in \mathrm{SIZE}(n^{c})\) to the carrier upper bound, identifying whether \(\beta=0\), \(\beta>0\), or the argument is genuinely nonblackbox, and recording the barrier status of the proposed escape.

\paragraph{Payoff criteria.}
\begin{itemize}
\item For the high-q generated-OR carrier, compute the exact leakage exponent of the remaining witness-cover or unary-residual mechanism and compare it with 1+epsilon-q.
\item Promotion requires \(q+a>1+\epsilon\) after every projection, padding, and witness-cover leakage term is accounted for; otherwise the branch becomes a bounded-window or no-go result.
\item For this route, state the exact barrier premise it uses or escapes: relativizing, natural, algebrizing, nonuniform, or advice-leaking; promotion requires the escape to be attached to the specific transfer theorem.
\item Compute \(\beta\) and q exactly; promotion requires \(\beta=0\), subpolynomial target queries, or a genuinely nonblackbox use of the SAT circuit.
\item State whether the transfer excludes every polynomial SAT exponent or only a bounded interval \(c<c0\).
\item Record relativization, natural-proofs, algebrization, nonuniformity, and advice-leakage status for the specific transfer theorem.
\end{itemize}

\paragraph{Stop criteria.}
\begin{itemize}
\item Every named carrier in the branch has \(\beta>0\) and yields only a bounded target-exponent exclusion.
\item The branch cannot state a source lower-bound model, target class, and transfer inequality in the same notation.
\item The proposed escape is only a black-box substitution restatement with no new carrier or nonblackbox mechanism.
\item The barrier audit shows the needed source lower bound is just as hard as the original P vs NP/P/poly target.
\end{itemize}

\paragraph{Do not stop if.}
\begin{itemize}
\item The branch has not yet computed \(\beta\), q, and the induced source upper-bound exponent under \(\mathrm{SAT}\in \mathrm{SIZE}(n^{c})\).
\item A subpolynomial-query or nonblackbox carrier remains available and has not been audited.
\item The latest failure sharpened the threshold-collapse lemma into a formal no-go statement.
\item The route can still isolate advice leakage, nonuniformity, naturalness, or algebrization as the exact obstruction.
\end{itemize}

\paragraph{Required skills.}
\begin{itemize}
\item abducer-failure-signal-cartographer
\item contradiction-route-auditor
\item entailment-gap-auditor
\item kb-deduce
\item kb-end-to-end-proof-critic
\item kb-hypothesis-refutation-loop
\item kb-induce
\item long-horizon-route-portfolio-manager
\item main-problem-alignment-auditor
\item next-experiment-queue-ranker
\item paper-contribution-readiness-judge
\item pvsnp-barrier-specificity-reviewer
\item pvsnp-conceptual-pressure-localizer
\item pvsnp-impossibility-theorem-forger
\item pvsnp-inference-quality-controller
\item pvsnp-proof-derivation-builder
\item pvsnp-qualitative-complexity-reasoner
\item pvsnp-quantifier-uniformity-reviewer
\item pvsnp-reduction-consistency-reviewer
\item pvsnp-symbolic-complexity-auditor
\end{itemize}

\subsection*{\#57 cycle-0018: Hardness-Magnification Bridge Realization: positioned entropy cover gate}

\paragraph{Analysis.} Instantiate the affine-disperser generated-OR family and compute \(B_{N}\) exactly as a function of \(\alpha\), q, and clause density; then audit whether any known lower-bound technique for this restricted \(H_{N}\) avoids the natural-proofs density condition.

\paragraph{Why this is a good target.} It directly tests the current black-box substitution cap by computing \(\beta\), q, and the induced target-exponent window. It checks whether padding/projection transfers really magnify hardness or merely move the proof debt into sparse lower bounds. It keeps the barrier and model-boundary audit attached to the mathematical transfer rather than as after-the-fact commentary. It is executable as a bounded complexity calculation, CAS-backed exponent audit, or formal transfer check.

\paragraph{Risk.} Low drift risk: the item is framed as an explicit P vs NP transfer audit.

\paragraph{Branch control.} The item has high payoff and a concrete decision test; stopping now would be premature.

\paragraph{Premature-stop guard.} Do not stop before deriving the transfer inequality from \(\mathrm{SAT}\in \mathrm{SIZE}(n^{c})\) to the carrier upper bound, identifying whether \(\beta=0\), \(\beta>0\), or the argument is genuinely nonblackbox, and recording the barrier status of the proposed escape.

\paragraph{Payoff criteria.}
\begin{itemize}
\item For the high-q generated-OR carrier, compute the exact leakage exponent of the remaining witness-cover or unary-residual mechanism and compare it with 1+epsilon-q.
\item Promotion requires \(q+a>1+\epsilon\) after every projection, padding, and witness-cover leakage term is accounted for; otherwise the branch becomes a bounded-window or no-go result.
\item For this route, state the exact barrier premise it uses or escapes: relativizing, natural, algebrizing, nonuniform, or advice-leaking; promotion requires the escape to be attached to the specific transfer theorem.
\item Compute \(\beta\) and q exactly; promotion requires \(\beta=0\), subpolynomial target queries, or a genuinely nonblackbox use of the SAT circuit.
\item State whether the transfer excludes every polynomial SAT exponent or only a bounded interval \(c<c0\).
\item Record relativization, natural-proofs, algebrization, nonuniformity, and advice-leakage status for the specific transfer theorem.
\end{itemize}

\paragraph{Stop criteria.}
\begin{itemize}
\item Every named carrier in the branch has \(\beta>0\) and yields only a bounded target-exponent exclusion.
\item The branch cannot state a source lower-bound model, target class, and transfer inequality in the same notation.
\item The proposed escape is only a black-box substitution restatement with no new carrier or nonblackbox mechanism.
\item The barrier audit shows the needed source lower bound is just as hard as the original P vs NP/P/poly target.
\end{itemize}

\paragraph{Do not stop if.}
\begin{itemize}
\item The branch has not yet computed \(\beta\), q, and the induced source upper-bound exponent under \(\mathrm{SAT}\in \mathrm{SIZE}(n^{c})\).
\item A subpolynomial-query or nonblackbox carrier remains available and has not been audited.
\item The latest failure sharpened the threshold-collapse lemma into a formal no-go statement.
\item The route can still isolate advice leakage, nonuniformity, naturalness, or algebrization as the exact obstruction.
\end{itemize}

\paragraph{Required skills.}
\begin{itemize}
\item abducer-failure-signal-cartographer
\item contradiction-route-auditor
\item entailment-gap-auditor
\item kb-deduce
\item kb-end-to-end-proof-critic
\item kb-hypothesis-refutation-loop
\item kb-induce
\item long-horizon-route-portfolio-manager
\item main-problem-alignment-auditor
\item next-experiment-queue-ranker
\item paper-contribution-readiness-judge
\item pvsnp-barrier-specificity-reviewer
\item pvsnp-conceptual-pressure-localizer
\item pvsnp-impossibility-theorem-forger
\item pvsnp-inference-quality-controller
\item pvsnp-proof-derivation-builder
\item pvsnp-qualitative-complexity-reasoner
\item pvsnp-quantifier-uniformity-reviewer
\item pvsnp-reduction-consistency-reviewer
\item pvsnp-symbolic-complexity-auditor
\end{itemize}

\subsection*{\#58 cycle-0014: Hardness-Magnification Bridge Realization: polylog generated-OR carrier}

\paragraph{Analysis.} Instantiate \(G_{N}\) as a near-linear Cook-Levin micro-instance generator with \(M=N\) and L=(log N)\textasciicircum{}d; compute exact \(\rho\), verify \(H_{N}\) is in NP with poly(N) uniform verification, and test whether any structural simplification gives an \(N^{1+\epsilon}\) upper-bound falsifier.

\paragraph{Why this is a good target.} It directly tests the current black-box substitution cap by computing \(\beta\), q, and the induced target-exponent window. It checks whether padding/projection transfers really magnify hardness or merely move the proof debt into sparse lower bounds. It keeps the barrier and model-boundary audit attached to the mathematical transfer rather than as after-the-fact commentary. It is executable as a bounded complexity calculation, CAS-backed exponent audit, or formal transfer check.

\paragraph{Risk.} Low drift risk: the item is framed as an explicit P vs NP transfer audit.

\paragraph{Branch control.} The item has high payoff and a concrete decision test; stopping now would be premature.

\paragraph{Premature-stop guard.} Do not stop before deriving the transfer inequality from \(\mathrm{SAT}\in \mathrm{SIZE}(n^{c})\) to the carrier upper bound, identifying whether \(\beta=0\), \(\beta>0\), or the argument is genuinely nonblackbox, and recording the barrier status of the proposed escape.

\paragraph{Payoff criteria.}
\begin{itemize}
\item For the high-q generated-OR carrier, compute the exact leakage exponent of the remaining witness-cover or unary-residual mechanism and compare it with 1+epsilon-q.
\item Promotion requires \(q+a>1+\epsilon\) after every projection, padding, and witness-cover leakage term is accounted for; otherwise the branch becomes a bounded-window or no-go result.
\item For the restricted-model branch, compute the compiler envelope from unrestricted SIZE(m\textasciicircum{}k) or P-time algorithms into the named model; promotion requires a bound above that envelope on the same SAT slice.
\item Falsify direct transfer if a P function already has the same restricted-model lower bound, such as PARITY for \(\mathrm{AC0}\), HWB for OBDD, or XOR/Tseitin for proof-system degree.
\item For this route, state the exact barrier premise it uses or escapes: relativizing, natural, algebrizing, nonuniform, or advice-leaking; promotion requires the escape to be attached to the specific transfer theorem.
\item Compute \(\beta\) and q exactly; promotion requires \(\beta=0\), subpolynomial target queries, or a genuinely nonblackbox use of the SAT circuit.
\end{itemize}

\paragraph{Stop criteria.}
\begin{itemize}
\item Every named carrier in the branch has \(\beta>0\) and yields only a bounded target-exponent exclusion.
\item The branch cannot state a source lower-bound model, target class, and transfer inequality in the same notation.
\item The proposed escape is only a black-box substitution restatement with no new carrier or nonblackbox mechanism.
\item The barrier audit shows the needed source lower bound is just as hard as the original P vs NP/P/poly target.
\end{itemize}

\paragraph{Do not stop if.}
\begin{itemize}
\item The branch has not yet computed \(\beta\), q, and the induced source upper-bound exponent under \(\mathrm{SAT}\in \mathrm{SIZE}(n^{c})\).
\item A subpolynomial-query or nonblackbox carrier remains available and has not been audited.
\item The latest failure sharpened the threshold-collapse lemma into a formal no-go statement.
\item The route can still isolate advice leakage, nonuniformity, naturalness, or algebrization as the exact obstruction.
\end{itemize}

\paragraph{Required skills.}
\begin{itemize}
\item abducer-failure-signal-cartographer
\item contradiction-route-auditor
\item entailment-gap-auditor
\item kb-deduce
\item kb-end-to-end-proof-critic
\item kb-hypothesis-refutation-loop
\item kb-induce
\item long-horizon-route-portfolio-manager
\item main-problem-alignment-auditor
\item next-experiment-queue-ranker
\item paper-contribution-readiness-judge
\item pvsnp-barrier-specificity-reviewer
\item pvsnp-conceptual-pressure-localizer
\item pvsnp-impossibility-theorem-forger
\item pvsnp-inference-quality-controller
\item pvsnp-proof-derivation-builder
\item pvsnp-qualitative-complexity-reasoner
\item pvsnp-quantifier-uniformity-reviewer
\item pvsnp-reduction-consistency-reviewer
\item pvsnp-symbolic-complexity-auditor
\end{itemize}

\subsection*{\#59 cycle-0005: Hardness-Magnification Bridge Realization}

\paragraph{Analysis.} Audit an MCSP-based route: set source object \(\mathrm{MCSP}_{n,s}\) on \(N=2^{n}-bit\) truth tables, choose a concrete circuit model for the lower bound, state the exact threshold \(N^{1+\epsilon}\) or smaller, and test whether the known magnification transfer is non-natural, target-preserving, and strong enough to imply \(\mathrm{SAT}\notin \mathrm{P}/\mathrm{poly}\) rather than only EXP or subexponential derandomization consequences.

\paragraph{Why this is a good target.} It audits a meta-complexity carrier where the decisive questions are verifier overhead, density/naturalness, and advice leakage. It checks whether padding/projection transfers really magnify hardness or merely move the proof debt into sparse lower bounds. It keeps the barrier and model-boundary audit attached to the mathematical transfer rather than as after-the-fact commentary. It is executable as a bounded complexity calculation, CAS-backed exponent audit, or formal transfer check.

\paragraph{Risk.} Low drift risk: the item is framed as an explicit P vs NP transfer audit.

\paragraph{Branch control.} The item has high payoff and a concrete decision test; stopping now would be premature.

\paragraph{Premature-stop guard.} Do not stop before naming the carrier variant, deriving its verifier/reduction overhead, and checking whether a \(\mathrm{SAT}\in \mathrm{SIZE}(n^{c})\) assumption gives carrier circuits below the claimed source lower-bound threshold.

\paragraph{Payoff criteria.}
\begin{itemize}
\item For the MCSP/MKtP row branch, decide whether the certificate proves direct row hardness or only aggregate/rank-index hardness; record q\_extract=epsilon/(1+D) and q\_force=(D+epsilon)/(1+D) in the same parameter regime.
\item Promotion requires an advice-free row-hard certificate; backtrack if the ranked decoder, canonical unranking, or aggregate proof selects a hard row nonuniformly.
\item For this route, state the exact barrier premise it uses or escapes: relativizing, natural, algebrizing, nonuniform, or advice-leaking; promotion requires the escape to be attached to the specific transfer theorem.
\item For MCSP/MKtP, state N, n, s, verifier overhead, lower-bound model, and whether the property is natural or non-natural.
\item State whether the transfer excludes every polynomial SAT exponent or only a bounded interval \(c<c0\).
\item Record relativization, natural-proofs, algebrization, nonuniformity, and advice-leakage status for the specific transfer theorem.
\end{itemize}

\paragraph{Stop criteria.}
\begin{itemize}
\item The induced carrier upper bound under \(\mathrm{SAT}\in \mathrm{SIZE}(n^{c})\) remains N\textasciicircum{}\{q+beta c+o(1)\} with \(\beta>0\) for all tested parameter choices.
\item The source lower bound required for MCSP/MKtP is as strong as the target separation or becomes a natural property without an escape.
\item The route cannot prevent advice leakage or model mismatch between the carrier lower bound and P/poly target.
\end{itemize}

\paragraph{Do not stop if.}
\begin{itemize}
\item The MCSP/MKtP parameter regime has not yet been fixed.
\item The truth-table length N and underlying input length n have not been related in the transfer inequality.
\item The natural-proof largeness/constructivity issue has not yet been audited.
\item A nonblackbox use of the target circuit could still avoid the multiplicative c factor.
\end{itemize}

\paragraph{Required skills.}
\begin{itemize}
\item abducer-failure-signal-cartographer
\item counterexample-hunter
\item entailment-gap-auditor
\item kb-abduce
\item kb-concept-bridge-planner
\item kb-end-to-end-proof-critic
\item kb-induce
\item long-horizon-route-portfolio-manager
\item main-problem-alignment-auditor
\item next-experiment-queue-ranker
\item paper-contribution-readiness-judge
\item pvsnp-barrier-specificity-reviewer
\item pvsnp-conceptual-pressure-localizer
\item pvsnp-impossibility-theorem-forger
\item pvsnp-inference-quality-controller
\item pvsnp-meta-complexity-carrier-reviewer
\item pvsnp-qualitative-complexity-reasoner
\item pvsnp-quantifier-uniformity-reviewer
\item pvsnp-reduction-consistency-reviewer
\item pvsnp-symbolic-complexity-auditor
\end{itemize}

\subsection*{\#60 cycle-0004: Hardness-Magnification Bridge Realization}

\paragraph{Analysis.} Test an MCSP or KT-complexity carrier with source model, threshold \(N^{1+eps}\), and target \(\mathrm{NP}\not\subset \mathrm{P}/\mathrm{poly}\); compute the reduction overhead and state whether the needed proof must be non-natural or non-algebrizing.

\paragraph{Why this is a good target.} It audits a meta-complexity carrier where the decisive questions are verifier overhead, density/naturalness, and advice leakage. It checks whether padding/projection transfers really magnify hardness or merely move the proof debt into sparse lower bounds. It keeps the barrier and model-boundary audit attached to the mathematical transfer rather than as after-the-fact commentary. It is executable as a bounded complexity calculation, CAS-backed exponent audit, or formal transfer check.

\paragraph{Risk.} Low drift risk: the item is framed as an explicit P vs NP transfer audit.

\paragraph{Branch control.} The item has high payoff and a concrete decision test; stopping now would be premature.

\paragraph{Premature-stop guard.} Do not stop before naming the carrier variant, deriving its verifier/reduction overhead, and checking whether a \(\mathrm{SAT}\in \mathrm{SIZE}(n^{c})\) assumption gives carrier circuits below the claimed source lower-bound threshold.

\paragraph{Payoff criteria.}
\begin{itemize}
\item For the MCSP/MKtP row branch, decide whether the certificate proves direct row hardness or only aggregate/rank-index hardness; record q\_extract=epsilon/(1+D) and q\_force=(D+epsilon)/(1+D) in the same parameter regime.
\item Promotion requires an advice-free row-hard certificate; backtrack if the ranked decoder, canonical unranking, or aggregate proof selects a hard row nonuniformly.
\item For this route, state the exact barrier premise it uses or escapes: relativizing, natural, algebrizing, nonuniform, or advice-leaking; promotion requires the escape to be attached to the specific transfer theorem.
\item For MCSP/MKtP, state N, n, s, verifier overhead, lower-bound model, and whether the property is natural or non-natural.
\item State whether the transfer excludes every polynomial SAT exponent or only a bounded interval \(c<c0\).
\item Record relativization, natural-proofs, algebrization, nonuniformity, and advice-leakage status for the specific transfer theorem.
\end{itemize}

\paragraph{Stop criteria.}
\begin{itemize}
\item The induced carrier upper bound under \(\mathrm{SAT}\in \mathrm{SIZE}(n^{c})\) remains N\textasciicircum{}\{q+beta c+o(1)\} with \(\beta>0\) for all tested parameter choices.
\item The source lower bound required for MCSP/MKtP is as strong as the target separation or becomes a natural property without an escape.
\item The route cannot prevent advice leakage or model mismatch between the carrier lower bound and P/poly target.
\end{itemize}

\paragraph{Do not stop if.}
\begin{itemize}
\item The MCSP/MKtP parameter regime has not yet been fixed.
\item The truth-table length N and underlying input length n have not been related in the transfer inequality.
\item The natural-proof largeness/constructivity issue has not yet been audited.
\item A nonblackbox use of the target circuit could still avoid the multiplicative c factor.
\end{itemize}

\paragraph{Required skills.}
\begin{itemize}
\item abducer-failure-signal-cartographer
\item counterexample-hunter
\item entailment-gap-auditor
\item kb-abduce
\item kb-concept-bridge-planner
\item kb-end-to-end-proof-critic
\item kb-induce
\item long-horizon-route-portfolio-manager
\item main-problem-alignment-auditor
\item next-experiment-queue-ranker
\item paper-contribution-readiness-judge
\item pvsnp-barrier-specificity-reviewer
\item pvsnp-conceptual-pressure-localizer
\item pvsnp-impossibility-theorem-forger
\item pvsnp-inference-quality-controller
\item pvsnp-meta-complexity-carrier-reviewer
\item pvsnp-qualitative-complexity-reasoner
\item pvsnp-quantifier-uniformity-reviewer
\item pvsnp-reduction-consistency-reviewer
\item pvsnp-symbolic-complexity-auditor
\end{itemize}

\subsection*{\#61 cycle-0001: Hardness-Magnification Bridge Realization}

\paragraph{Analysis.} Instantiate the same exponent audit for a padded MKtP or GapMCSP carrier: name the verifier, compute Cook-Levin or direct-composition length blow-up, and test whether assuming \(\mathrm{SAT}\in \mathrm{SIZE}(n^{c})\) yields source circuits below \(N^{1+\epsilon}\) without model mismatch or nonuniform leakage.

\paragraph{Why this is a good target.} It audits a meta-complexity carrier where the decisive questions are verifier overhead, density/naturalness, and advice leakage. It checks whether padding/projection transfers really magnify hardness or merely move the proof debt into sparse lower bounds. It keeps the barrier and model-boundary audit attached to the mathematical transfer rather than as after-the-fact commentary. It is executable as a bounded complexity calculation, CAS-backed exponent audit, or formal transfer check.

\paragraph{Risk.} Low drift risk: the item is framed as an explicit P vs NP transfer audit.

\paragraph{Branch control.} The item has high payoff and a concrete decision test; stopping now would be premature.

\paragraph{Premature-stop guard.} Do not stop before naming the carrier variant, deriving its verifier/reduction overhead, and checking whether a \(\mathrm{SAT}\in \mathrm{SIZE}(n^{c})\) assumption gives carrier circuits below the claimed source lower-bound threshold.

\paragraph{Payoff criteria.}
\begin{itemize}
\item For the MCSP/MKtP row branch, decide whether the certificate proves direct row hardness or only aggregate/rank-index hardness; record q\_extract=epsilon/(1+D) and q\_force=(D+epsilon)/(1+D) in the same parameter regime.
\item Promotion requires an advice-free row-hard certificate; backtrack if the ranked decoder, canonical unranking, or aggregate proof selects a hard row nonuniformly.
\item For the restricted-model branch, compute the compiler envelope from unrestricted SIZE(m\textasciicircum{}k) or P-time algorithms into the named model; promotion requires a bound above that envelope on the same SAT slice.
\item Falsify direct transfer if a P function already has the same restricted-model lower bound, such as PARITY for \(\mathrm{AC0}\), HWB for OBDD, or XOR/Tseitin for proof-system degree.
\item For this route, state the exact barrier premise it uses or escapes: relativizing, natural, algebrizing, nonuniform, or advice-leaking; promotion requires the escape to be attached to the specific transfer theorem.
\item For MCSP/MKtP, state N, n, s, verifier overhead, lower-bound model, and whether the property is natural or non-natural.
\end{itemize}

\paragraph{Stop criteria.}
\begin{itemize}
\item The induced carrier upper bound under \(\mathrm{SAT}\in \mathrm{SIZE}(n^{c})\) remains N\textasciicircum{}\{q+beta c+o(1)\} with \(\beta>0\) for all tested parameter choices.
\item The source lower bound required for MCSP/MKtP is as strong as the target separation or becomes a natural property without an escape.
\item The route cannot prevent advice leakage or model mismatch between the carrier lower bound and P/poly target.
\end{itemize}

\paragraph{Do not stop if.}
\begin{itemize}
\item The MCSP/MKtP parameter regime has not yet been fixed.
\item The truth-table length N and underlying input length n have not been related in the transfer inequality.
\item The natural-proof largeness/constructivity issue has not yet been audited.
\item A nonblackbox use of the target circuit could still avoid the multiplicative c factor.
\end{itemize}

\paragraph{Required skills.}
\begin{itemize}
\item abducer-failure-signal-cartographer
\item counterexample-hunter
\item entailment-gap-auditor
\item kb-abduce
\item kb-concept-bridge-planner
\item kb-end-to-end-proof-critic
\item kb-induce
\item long-horizon-route-portfolio-manager
\item main-problem-alignment-auditor
\item next-experiment-queue-ranker
\item paper-contribution-readiness-judge
\item pvsnp-barrier-specificity-reviewer
\item pvsnp-conceptual-pressure-localizer
\item pvsnp-impossibility-theorem-forger
\item pvsnp-inference-quality-controller
\item pvsnp-meta-complexity-carrier-reviewer
\item pvsnp-qualitative-complexity-reasoner
\item pvsnp-quantifier-uniformity-reviewer
\item pvsnp-reduction-consistency-reviewer
\item pvsnp-symbolic-complexity-auditor
\end{itemize}

\subsection*{\#62 cycle-0110: Hardness-Magnification Bridge Realization}

\paragraph{Analysis.} Attempt \(RealSATInfluenceRealizer_{N}\) for one influence-complete extension rule: define \(RealInfl_{l}\), produce row-index witnesses and \(EnumCost+N|K_{l}|\) charge, or exhibit a semantic rule where a real SAT-sensitive row remains unchargeable, thereby demoting full \(\mathrm{Psys}_{sym}\) beyond the transparent theorem pair.

\paragraph{Why this is a good target.} It directly tests the current black-box substitution cap by computing \(\beta\), q, and the induced target-exponent window. It points toward the only unquarantined escape: a subpolynomial-query, nonprojection, or nonblackbox carrier. It keeps the barrier and model-boundary audit attached to the mathematical transfer rather than as after-the-fact commentary.

\paragraph{Risk.} Low drift risk: the item is framed as an explicit P vs NP transfer audit.

\paragraph{Branch control.} The item has high payoff and a concrete decision test; stopping now would be premature.

\paragraph{Premature-stop guard.} Do not stop before deriving the transfer inequality from \(\mathrm{SAT}\in \mathrm{SIZE}(n^{c})\) to the carrier upper bound, identifying whether \(\beta=0\), \(\beta>0\), or the argument is genuinely nonblackbox, and recording the barrier status of the proposed escape.

\paragraph{Payoff criteria.}
\begin{itemize}
\item For this route, state the exact barrier premise it uses or escapes: relativizing, natural, algebrizing, nonuniform, or advice-leaking; promotion requires the escape to be attached to the specific transfer theorem.
\item Compute \(\beta\) and q exactly; promotion requires \(\beta=0\), subpolynomial target queries, or a genuinely nonblackbox use of the SAT circuit.
\item Promotion requires a target-preserving nonblackbox/subpolynomial mechanism, not another black-box Cook-Levin substitution.
\item Record relativization, natural-proofs, algebrization, nonuniformity, and advice-leakage status for the specific transfer theorem.
\item Control gate: name the source carrier or language, target class, and circuit or resource model before claiming progress.
\item Control gate: derive the exact transfer inequality from \(\mathrm{SAT}\in \mathrm{SIZE}(n^{c})\), NP subset P/poly, or the stated target assumption to the source-carrier upper bound.
\end{itemize}

\paragraph{Stop criteria.}
\begin{itemize}
\item Every named carrier in the branch has \(\beta>0\) and yields only a bounded target-exponent exclusion.
\item The branch cannot state a source lower-bound model, target class, and transfer inequality in the same notation.
\item The proposed escape is only a black-box substitution restatement with no new carrier or nonblackbox mechanism.
\item The barrier audit shows the needed source lower bound is just as hard as the original P vs NP/P/poly target.
\end{itemize}

\paragraph{Do not stop if.}
\begin{itemize}
\item The branch has not yet computed \(\beta\), q, and the induced source upper-bound exponent under \(\mathrm{SAT}\in \mathrm{SIZE}(n^{c})\).
\item A subpolynomial-query or nonblackbox carrier remains available and has not been audited.
\item The latest failure sharpened the threshold-collapse lemma into a formal no-go statement.
\item The route can still isolate advice leakage, nonuniformity, naturalness, or algebrization as the exact obstruction.
\end{itemize}

\paragraph{Required skills.}
\begin{itemize}
\item abducer-refinement-route-ranker
\item contradiction-route-auditor
\item entailment-gap-auditor
\item kb-deduce
\item kb-end-to-end-proof-critic
\item kb-gap-abstraction-conjecturer
\item kb-hypothesis-refutation-loop
\item kb-proof-connector
\item long-horizon-route-portfolio-manager
\item main-problem-alignment-auditor
\item next-experiment-queue-ranker
\item paper-contribution-readiness-judge
\item pvsnp-barrier-specificity-reviewer
\item pvsnp-conceptual-pressure-localizer
\item pvsnp-impossibility-theorem-forger
\item pvsnp-inference-quality-controller
\item pvsnp-proof-derivation-builder
\item pvsnp-qualitative-complexity-reasoner
\item pvsnp-quantifier-uniformity-reviewer
\item pvsnp-reduction-consistency-reviewer
\item pvsnp-symbolic-complexity-auditor
\end{itemize}

\subsection*{\#63 cycle-0129: Hardness-Magnification Bridge Realization}

\paragraph{Analysis.} Write the paper-ready L129 sublemma proof: give pseudocode for \(K_\mathrm{DTSR}\) fixed-point enumeration, a table of declared TSR bit fields, the local verifier for child occurrence neighborhoods, and the two-world first-unlisted-child countermodel; then run the live inference-quality packet validation.

\paragraph{Why this is a good target.} It directly tests the current black-box substitution cap by computing \(\beta\), q, and the induced target-exponent window. It checks whether padding/projection transfers really magnify hardness or merely move the proof debt into sparse lower bounds. It keeps the barrier and model-boundary audit attached to the mathematical transfer rather than as after-the-fact commentary.

\paragraph{Risk.} Low drift risk: the item is framed as an explicit P vs NP transfer audit.

\paragraph{Branch control.} The item has high payoff and a concrete decision test; stopping now would be premature.

\paragraph{Premature-stop guard.} Do not stop before deriving the transfer inequality from \(\mathrm{SAT}\in \mathrm{SIZE}(n^{c})\) to the carrier upper bound, identifying whether \(\beta=0\), \(\beta>0\), or the argument is genuinely nonblackbox, and recording the barrier status of the proposed escape.

\paragraph{Payoff criteria.}
\begin{itemize}
\item For this route, state the exact barrier premise it uses or escapes: relativizing, natural, algebrizing, nonuniform, or advice-leaking; promotion requires the escape to be attached to the specific transfer theorem.
\item Compute \(\beta\) and q exactly; promotion requires \(\beta=0\), subpolynomial target queries, or a genuinely nonblackbox use of the SAT circuit.
\item State whether the transfer excludes every polynomial SAT exponent or only a bounded interval \(c<c0\).
\item Record relativization, natural-proofs, algebrization, nonuniformity, and advice-leakage status for the specific transfer theorem.
\item Control gate: name the source carrier or language, target class, and circuit or resource model before claiming progress.
\item Control gate: derive the exact transfer inequality from \(\mathrm{SAT}\in \mathrm{SIZE}(n^{c})\), NP subset P/poly, or the stated target assumption to the source-carrier upper bound.
\end{itemize}

\paragraph{Stop criteria.}
\begin{itemize}
\item Every named carrier in the branch has \(\beta>0\) and yields only a bounded target-exponent exclusion.
\item The branch cannot state a source lower-bound model, target class, and transfer inequality in the same notation.
\item The proposed escape is only a black-box substitution restatement with no new carrier or nonblackbox mechanism.
\item The barrier audit shows the needed source lower bound is just as hard as the original P vs NP/P/poly target.
\end{itemize}

\paragraph{Do not stop if.}
\begin{itemize}
\item The branch has not yet computed \(\beta\), q, and the induced source upper-bound exponent under \(\mathrm{SAT}\in \mathrm{SIZE}(n^{c})\).
\item A subpolynomial-query or nonblackbox carrier remains available and has not been audited.
\item The latest failure sharpened the threshold-collapse lemma into a formal no-go statement.
\item The route can still isolate advice leakage, nonuniformity, naturalness, or algebrization as the exact obstruction.
\end{itemize}

\paragraph{Required skills.}
\begin{itemize}
\item abducer-failure-signal-cartographer
\item contradiction-route-auditor
\item entailment-gap-auditor
\item kb-deduce
\item kb-end-to-end-proof-critic
\item kb-hypothesis-refutation-loop
\item kb-induce
\item long-horizon-route-portfolio-manager
\item main-problem-alignment-auditor
\item next-experiment-queue-ranker
\item paper-contribution-readiness-judge
\item pvsnp-barrier-specificity-reviewer
\item pvsnp-conceptual-pressure-localizer
\item pvsnp-impossibility-theorem-forger
\item pvsnp-inference-quality-controller
\item pvsnp-proof-derivation-builder
\item pvsnp-qualitative-complexity-reasoner
\item pvsnp-quantifier-uniformity-reviewer
\item pvsnp-reduction-consistency-reviewer
\item pvsnp-symbolic-complexity-auditor
\end{itemize}

\subsection*{\#64 cycle-0093: Hardness-Magnification Bridge Realization}

\paragraph{Analysis.} Attempt an explicit \(NSP_{N}\) countermodel for unrestricted EF extension variables: construct a small row-mask extension family \(e_{R}\), prove that occurrence-only parsing misses its support, and decide whether support transparency can be repaired by charging C(Supp(e\_R)) without making the parser target-equivalent.

\paragraph{Why this is a good target.} It directly tests the current black-box substitution cap by computing \(\beta\), q, and the induced target-exponent window. It checks whether padding/projection transfers really magnify hardness or merely move the proof debt into sparse lower bounds. It keeps the barrier and model-boundary audit attached to the mathematical transfer rather than as after-the-fact commentary.

\paragraph{Risk.} Low drift risk: the item is framed as an explicit P vs NP transfer audit.

\paragraph{Branch control.} The item has high payoff and a concrete decision test; stopping now would be premature.

\paragraph{Premature-stop guard.} Do not stop before deriving the transfer inequality from \(\mathrm{SAT}\in \mathrm{SIZE}(n^{c})\) to the carrier upper bound, identifying whether \(\beta=0\), \(\beta>0\), or the argument is genuinely nonblackbox, and recording the barrier status of the proposed escape.

\paragraph{Payoff criteria.}
\begin{itemize}
\item For this route, state the exact barrier premise it uses or escapes: relativizing, natural, algebrizing, nonuniform, or advice-leaking; promotion requires the escape to be attached to the specific transfer theorem.
\item Compute \(\beta\) and q exactly; promotion requires \(\beta=0\), subpolynomial target queries, or a genuinely nonblackbox use of the SAT circuit.
\item State whether the transfer excludes every polynomial SAT exponent or only a bounded interval \(c<c0\).
\item Record relativization, natural-proofs, algebrization, nonuniformity, and advice-leakage status for the specific transfer theorem.
\item Control gate: name the source carrier or language, target class, and circuit or resource model before claiming progress.
\item Control gate: derive the exact transfer inequality from \(\mathrm{SAT}\in \mathrm{SIZE}(n^{c})\), NP subset P/poly, or the stated target assumption to the source-carrier upper bound.
\end{itemize}

\paragraph{Stop criteria.}
\begin{itemize}
\item Every named carrier in the branch has \(\beta>0\) and yields only a bounded target-exponent exclusion.
\item The branch cannot state a source lower-bound model, target class, and transfer inequality in the same notation.
\item The proposed escape is only a black-box substitution restatement with no new carrier or nonblackbox mechanism.
\item The barrier audit shows the needed source lower bound is just as hard as the original P vs NP/P/poly target.
\end{itemize}

\paragraph{Do not stop if.}
\begin{itemize}
\item The branch has not yet computed \(\beta\), q, and the induced source upper-bound exponent under \(\mathrm{SAT}\in \mathrm{SIZE}(n^{c})\).
\item A subpolynomial-query or nonblackbox carrier remains available and has not been audited.
\item The latest failure sharpened the threshold-collapse lemma into a formal no-go statement.
\item The route can still isolate advice leakage, nonuniformity, naturalness, or algebrization as the exact obstruction.
\end{itemize}

\paragraph{Required skills.}
\begin{itemize}
\item abducer-failure-signal-cartographer
\item contradiction-route-auditor
\item entailment-gap-auditor
\item kb-deduce
\item kb-end-to-end-proof-critic
\item kb-hypothesis-refutation-loop
\item kb-induce
\item long-horizon-route-portfolio-manager
\item main-problem-alignment-auditor
\item next-experiment-queue-ranker
\item paper-contribution-readiness-judge
\item pvsnp-barrier-specificity-reviewer
\item pvsnp-conceptual-pressure-localizer
\item pvsnp-impossibility-theorem-forger
\item pvsnp-inference-quality-controller
\item pvsnp-proof-derivation-builder
\item pvsnp-qualitative-complexity-reasoner
\item pvsnp-quantifier-uniformity-reviewer
\item pvsnp-reduction-consistency-reviewer
\item pvsnp-symbolic-complexity-auditor
\end{itemize}

\subsection*{\#65 cycle-0092: Hardness-Magnification Bridge Realization}

\paragraph{Analysis.} Try to prove \(NSP_{N}\) for the fragment EF+SATeq+SR whose extension definitions are row-affine and support-monotone; if the proof fails, extract the minimal extension rule that creates uncharged SAT-semantic support and state the final theorem scope around \(\mathrm{Psys}_{h}^{supp}\).

\paragraph{Why this is a good target.} It directly tests the current black-box substitution cap by computing \(\beta\), q, and the induced target-exponent window. It checks whether padding/projection transfers really magnify hardness or merely move the proof debt into sparse lower bounds. It keeps the barrier and model-boundary audit attached to the mathematical transfer rather than as after-the-fact commentary.

\paragraph{Risk.} Low drift risk: the item is framed as an explicit P vs NP transfer audit.

\paragraph{Branch control.} The item has high payoff and a concrete decision test; stopping now would be premature.

\paragraph{Premature-stop guard.} Do not stop before deriving the transfer inequality from \(\mathrm{SAT}\in \mathrm{SIZE}(n^{c})\) to the carrier upper bound, identifying whether \(\beta=0\), \(\beta>0\), or the argument is genuinely nonblackbox, and recording the barrier status of the proposed escape.

\paragraph{Payoff criteria.}
\begin{itemize}
\item For this route, state the exact barrier premise it uses or escapes: relativizing, natural, algebrizing, nonuniform, or advice-leaking; promotion requires the escape to be attached to the specific transfer theorem.
\item Compute \(\beta\) and q exactly; promotion requires \(\beta=0\), subpolynomial target queries, or a genuinely nonblackbox use of the SAT circuit.
\item State whether the transfer excludes every polynomial SAT exponent or only a bounded interval \(c<c0\).
\item Record relativization, natural-proofs, algebrization, nonuniformity, and advice-leakage status for the specific transfer theorem.
\item Control gate: name the source carrier or language, target class, and circuit or resource model before claiming progress.
\item Control gate: derive the exact transfer inequality from \(\mathrm{SAT}\in \mathrm{SIZE}(n^{c})\), NP subset P/poly, or the stated target assumption to the source-carrier upper bound.
\end{itemize}

\paragraph{Stop criteria.}
\begin{itemize}
\item Every named carrier in the branch has \(\beta>0\) and yields only a bounded target-exponent exclusion.
\item The branch cannot state a source lower-bound model, target class, and transfer inequality in the same notation.
\item The proposed escape is only a black-box substitution restatement with no new carrier or nonblackbox mechanism.
\item The barrier audit shows the needed source lower bound is just as hard as the original P vs NP/P/poly target.
\end{itemize}

\paragraph{Do not stop if.}
\begin{itemize}
\item The branch has not yet computed \(\beta\), q, and the induced source upper-bound exponent under \(\mathrm{SAT}\in \mathrm{SIZE}(n^{c})\).
\item A subpolynomial-query or nonblackbox carrier remains available and has not been audited.
\item The latest failure sharpened the threshold-collapse lemma into a formal no-go statement.
\item The route can still isolate advice leakage, nonuniformity, naturalness, or algebrization as the exact obstruction.
\end{itemize}

\paragraph{Required skills.}
\begin{itemize}
\item abducer-failure-signal-cartographer
\item contradiction-route-auditor
\item entailment-gap-auditor
\item kb-deduce
\item kb-end-to-end-proof-critic
\item kb-hypothesis-refutation-loop
\item kb-induce
\item long-horizon-route-portfolio-manager
\item main-problem-alignment-auditor
\item next-experiment-queue-ranker
\item paper-contribution-readiness-judge
\item pvsnp-barrier-specificity-reviewer
\item pvsnp-conceptual-pressure-localizer
\item pvsnp-impossibility-theorem-forger
\item pvsnp-inference-quality-controller
\item pvsnp-proof-derivation-builder
\item pvsnp-qualitative-complexity-reasoner
\item pvsnp-quantifier-uniformity-reviewer
\item pvsnp-reduction-consistency-reviewer
\item pvsnp-symbolic-complexity-auditor
\end{itemize}

\subsection*{\#66 cycle-0154: Hardness-Magnification Bridge Realization}

\paragraph{Analysis.} Run the prefix-scan verifier on one minimal proof with a valid two-row strict-prefix Qrep merge and on one identical proof with the certificate delayed by one line; compute \(\Gamma_{Qrep}\) in both cases and export the rejected line as the private-row falsifier pivot.

\paragraph{Why this is a good target.} It directly tests the current black-box substitution cap by computing \(\beta\), q, and the induced target-exponent window. It keeps the barrier and model-boundary audit attached to the mathematical transfer rather than as after-the-fact commentary. It is executable as a bounded complexity calculation, CAS-backed exponent audit, or formal transfer check.

\paragraph{Risk.} Low drift risk: the item is framed as an explicit P vs NP transfer audit.

\paragraph{Branch control.} The item has high payoff and a concrete decision test; stopping now would be premature.

\paragraph{Premature-stop guard.} Do not stop before deriving the transfer inequality from \(\mathrm{SAT}\in \mathrm{SIZE}(n^{c})\) to the carrier upper bound, identifying whether \(\beta=0\), \(\beta>0\), or the argument is genuinely nonblackbox, and recording the barrier status of the proposed escape.

\paragraph{Payoff criteria.}
\begin{itemize}
\item For this route, state the exact barrier premise it uses or escapes: relativizing, natural, algebrizing, nonuniform, or advice-leaking; promotion requires the escape to be attached to the specific transfer theorem.
\item Compute \(\beta\) and q exactly; promotion requires \(\beta=0\), subpolynomial target queries, or a genuinely nonblackbox use of the SAT circuit.
\item Record relativization, natural-proofs, algebrization, nonuniformity, and advice-leakage status for the specific transfer theorem.
\item Control gate: name the source carrier or language, target class, and circuit or resource model before claiming progress.
\item Control gate: derive the exact transfer inequality from \(\mathrm{SAT}\in \mathrm{SIZE}(n^{c})\), NP subset P/poly, or the stated target assumption to the source-carrier upper bound.
\item Control gate: compute \(\beta\) and q exactly, distinguishing all-exponent consequences from bounded target-exponent windows.
\end{itemize}

\paragraph{Stop criteria.}
\begin{itemize}
\item Every named carrier in the branch has \(\beta>0\) and yields only a bounded target-exponent exclusion.
\item The branch cannot state a source lower-bound model, target class, and transfer inequality in the same notation.
\item The proposed escape is only a black-box substitution restatement with no new carrier or nonblackbox mechanism.
\item The barrier audit shows the needed source lower bound is just as hard as the original P vs NP/P/poly target.
\end{itemize}

\paragraph{Do not stop if.}
\begin{itemize}
\item The branch has not yet computed \(\beta\), q, and the induced source upper-bound exponent under \(\mathrm{SAT}\in \mathrm{SIZE}(n^{c})\).
\item A subpolynomial-query or nonblackbox carrier remains available and has not been audited.
\item The latest failure sharpened the threshold-collapse lemma into a formal no-go statement.
\item The route can still isolate advice leakage, nonuniformity, naturalness, or algebrization as the exact obstruction.
\end{itemize}

\paragraph{Required skills.}
\begin{itemize}
\item contradiction-route-auditor
\item entailment-gap-auditor
\item kb-deduce
\item kb-end-to-end-proof-critic
\item kb-hypothesis-refutation-loop
\item long-horizon-route-portfolio-manager
\item main-problem-alignment-auditor
\item next-experiment-queue-ranker
\item paper-contribution-readiness-judge
\item pvsnp-barrier-specificity-reviewer
\item pvsnp-conceptual-pressure-localizer
\item pvsnp-impossibility-theorem-forger
\item pvsnp-inference-quality-controller
\item pvsnp-proof-derivation-builder
\item pvsnp-qualitative-complexity-reasoner
\item pvsnp-quantifier-uniformity-reviewer
\item pvsnp-reduction-consistency-reviewer
\item pvsnp-symbolic-complexity-auditor
\end{itemize}

\subsection*{\#67 cycle-0153: Hardness-Magnification Bridge Realization: prefix-strict Qrep support audit}

\paragraph{Analysis.} Instantiate the strict-prefix Qrep verifier on a two-duplicate projected-dDNNF FormulaEval example and a delayed-certificate counterexample; compute \(\Gamma_{bEF}\), \(U_{N}\), \(X_{N}\), \(D_{Q}\), \(E_{pi}\), and the q/epsilon support-density consequence line by line.

\paragraph{Why this is a good target.} It directly tests the current black-box substitution cap by computing \(\beta\), q, and the induced target-exponent window. It keeps the barrier and model-boundary audit attached to the mathematical transfer rather than as after-the-fact commentary. It is executable as a bounded complexity calculation, CAS-backed exponent audit, or formal transfer check.

\paragraph{Risk.} Low drift risk: the item is framed as an explicit P vs NP transfer audit.

\paragraph{Branch control.} The item has high payoff and a concrete decision test; stopping now would be premature.

\paragraph{Premature-stop guard.} Do not stop before deriving the transfer inequality from \(\mathrm{SAT}\in \mathrm{SIZE}(n^{c})\) to the carrier upper bound, identifying whether \(\beta=0\), \(\beta>0\), or the argument is genuinely nonblackbox, and recording the barrier status of the proposed escape.

\paragraph{Payoff criteria.}
\begin{itemize}
\item For this route, state the exact barrier premise it uses or escapes: relativizing, natural, algebrizing, nonuniform, or advice-leaking; promotion requires the escape to be attached to the specific transfer theorem.
\item Compute \(\beta\) and q exactly; promotion requires \(\beta=0\), subpolynomial target queries, or a genuinely nonblackbox use of the SAT circuit.
\item Record relativization, natural-proofs, algebrization, nonuniformity, and advice-leakage status for the specific transfer theorem.
\item Control gate: name the source carrier or language, target class, and circuit or resource model before claiming progress.
\item Control gate: derive the exact transfer inequality from \(\mathrm{SAT}\in \mathrm{SIZE}(n^{c})\), NP subset P/poly, or the stated target assumption to the source-carrier upper bound.
\item Control gate: compute \(\beta\) and q exactly, distinguishing all-exponent consequences from bounded target-exponent windows.
\end{itemize}

\paragraph{Stop criteria.}
\begin{itemize}
\item Every named carrier in the branch has \(\beta>0\) and yields only a bounded target-exponent exclusion.
\item The branch cannot state a source lower-bound model, target class, and transfer inequality in the same notation.
\item The proposed escape is only a black-box substitution restatement with no new carrier or nonblackbox mechanism.
\item The barrier audit shows the needed source lower bound is just as hard as the original P vs NP/P/poly target.
\end{itemize}

\paragraph{Do not stop if.}
\begin{itemize}
\item The branch has not yet computed \(\beta\), q, and the induced source upper-bound exponent under \(\mathrm{SAT}\in \mathrm{SIZE}(n^{c})\).
\item A subpolynomial-query or nonblackbox carrier remains available and has not been audited.
\item The latest failure sharpened the threshold-collapse lemma into a formal no-go statement.
\item The route can still isolate advice leakage, nonuniformity, naturalness, or algebrization as the exact obstruction.
\end{itemize}

\paragraph{Required skills.}
\begin{itemize}
\item contradiction-route-auditor
\item entailment-gap-auditor
\item kb-deduce
\item kb-end-to-end-proof-critic
\item kb-hypothesis-refutation-loop
\item long-horizon-route-portfolio-manager
\item main-problem-alignment-auditor
\item next-experiment-queue-ranker
\item paper-contribution-readiness-judge
\item pvsnp-barrier-specificity-reviewer
\item pvsnp-conceptual-pressure-localizer
\item pvsnp-impossibility-theorem-forger
\item pvsnp-inference-quality-controller
\item pvsnp-proof-derivation-builder
\item pvsnp-qualitative-complexity-reasoner
\item pvsnp-quantifier-uniformity-reviewer
\item pvsnp-reduction-consistency-reviewer
\item pvsnp-symbolic-complexity-auditor
\end{itemize}

\subsection*{\#68 cycle-0150: Hardness-Magnification Bridge Realization}

\paragraph{Analysis.} Run the Tseitin-renaming Qrep thought experiment: instantiate two batches with one certified duplicate and one unlisted duplicate, compute \(K_{Qrep}\), C\_Q(D\_Q), and \(\Gamma_{Qrep}\) explicitly, then verify that the unlisted duplicate yields the first-hidden-row private-row falsifier while the certified duplicate preserves the line induction.

\paragraph{Why this is a good target.} It directly tests the current black-box substitution cap by computing \(\beta\), q, and the induced target-exponent window. It keeps the barrier and model-boundary audit attached to the mathematical transfer rather than as after-the-fact commentary. It is executable as a bounded complexity calculation, CAS-backed exponent audit, or formal transfer check.

\paragraph{Risk.} Low drift risk: the item is framed as an explicit P vs NP transfer audit.

\paragraph{Branch control.} The item has high payoff and a concrete decision test; stopping now would be premature.

\paragraph{Premature-stop guard.} Do not stop before deriving the transfer inequality from \(\mathrm{SAT}\in \mathrm{SIZE}(n^{c})\) to the carrier upper bound, identifying whether \(\beta=0\), \(\beta>0\), or the argument is genuinely nonblackbox, and recording the barrier status of the proposed escape.

\paragraph{Payoff criteria.}
\begin{itemize}
\item For the proof-complexity branch, separate fixed-family proof hardness from language hardness by naming the proof system, shortcut algorithm, and semantic lower-bound target.
\item Promotion requires a proof-system lower bound that survives the Gaussian-elimination or recognition shortcut; otherwise the result is a restricted-proof-system separation only.
\item For this route, state the exact barrier premise it uses or escapes: relativizing, natural, algebrizing, nonuniform, or advice-leaking; promotion requires the escape to be attached to the specific transfer theorem.
\item Compute \(\beta\) and q exactly; promotion requires \(\beta=0\), subpolynomial target queries, or a genuinely nonblackbox use of the SAT circuit.
\item Record relativization, natural-proofs, algebrization, nonuniformity, and advice-leakage status for the specific transfer theorem.
\item Control gate: name the source carrier or language, target class, and circuit or resource model before claiming progress.
\end{itemize}

\paragraph{Stop criteria.}
\begin{itemize}
\item Every named carrier in the branch has \(\beta>0\) and yields only a bounded target-exponent exclusion.
\item The branch cannot state a source lower-bound model, target class, and transfer inequality in the same notation.
\item The proposed escape is only a black-box substitution restatement with no new carrier or nonblackbox mechanism.
\item The barrier audit shows the needed source lower bound is just as hard as the original P vs NP/P/poly target.
\end{itemize}

\paragraph{Do not stop if.}
\begin{itemize}
\item The branch has not yet computed \(\beta\), q, and the induced source upper-bound exponent under \(\mathrm{SAT}\in \mathrm{SIZE}(n^{c})\).
\item A subpolynomial-query or nonblackbox carrier remains available and has not been audited.
\item The latest failure sharpened the threshold-collapse lemma into a formal no-go statement.
\item The route can still isolate advice leakage, nonuniformity, naturalness, or algebrization as the exact obstruction.
\end{itemize}

\paragraph{Required skills.}
\begin{itemize}
\item contradiction-route-auditor
\item entailment-gap-auditor
\item kb-deduce
\item kb-end-to-end-proof-critic
\item kb-hypothesis-refutation-loop
\item long-horizon-route-portfolio-manager
\item main-problem-alignment-auditor
\item next-experiment-queue-ranker
\item paper-contribution-readiness-judge
\item pvsnp-barrier-specificity-reviewer
\item pvsnp-conceptual-pressure-localizer
\item pvsnp-impossibility-theorem-forger
\item pvsnp-inference-quality-controller
\item pvsnp-proof-derivation-builder
\item pvsnp-qualitative-complexity-reasoner
\item pvsnp-quantifier-uniformity-reviewer
\item pvsnp-reduction-consistency-reviewer
\item pvsnp-symbolic-complexity-auditor
\end{itemize}

\subsection*{\#69 cycle-0148: Hardness-Magnification Bridge Realization}

\paragraph{Analysis.} Write the finite \(K_\mathrm{QCTSR}\) verifier on a projected-dDNNF example with one certified duplicate row and one unread-enable hidden merge; compute \(U_{Q},N|K_{Q}|,B_\mathrm{CTSR},C_{Q},E_{pi}\) and show exactly where the hidden merge fails the valid support ledger.

\paragraph{Why this is a good target.} It directly tests the current black-box substitution cap by computing \(\beta\), q, and the induced target-exponent window. It keeps the barrier and model-boundary audit attached to the mathematical transfer rather than as after-the-fact commentary. It is executable as a bounded complexity calculation, CAS-backed exponent audit, or formal transfer check.

\paragraph{Risk.} Low drift risk: the item is framed as an explicit P vs NP transfer audit.

\paragraph{Branch control.} The item has high payoff and a concrete decision test; stopping now would be premature.

\paragraph{Premature-stop guard.} Do not stop before deriving the transfer inequality from \(\mathrm{SAT}\in \mathrm{SIZE}(n^{c})\) to the carrier upper bound, identifying whether \(\beta=0\), \(\beta>0\), or the argument is genuinely nonblackbox, and recording the barrier status of the proposed escape.

\paragraph{Payoff criteria.}
\begin{itemize}
\item For this route, state the exact barrier premise it uses or escapes: relativizing, natural, algebrizing, nonuniform, or advice-leaking; promotion requires the escape to be attached to the specific transfer theorem.
\item Compute \(\beta\) and q exactly; promotion requires \(\beta=0\), subpolynomial target queries, or a genuinely nonblackbox use of the SAT circuit.
\item Record relativization, natural-proofs, algebrization, nonuniformity, and advice-leakage status for the specific transfer theorem.
\item Control gate: name the source carrier or language, target class, and circuit or resource model before claiming progress.
\item Control gate: derive the exact transfer inequality from \(\mathrm{SAT}\in \mathrm{SIZE}(n^{c})\), NP subset P/poly, or the stated target assumption to the source-carrier upper bound.
\item Control gate: compute \(\beta\) and q exactly, distinguishing all-exponent consequences from bounded target-exponent windows.
\end{itemize}

\paragraph{Stop criteria.}
\begin{itemize}
\item Every named carrier in the branch has \(\beta>0\) and yields only a bounded target-exponent exclusion.
\item The branch cannot state a source lower-bound model, target class, and transfer inequality in the same notation.
\item The proposed escape is only a black-box substitution restatement with no new carrier or nonblackbox mechanism.
\item The barrier audit shows the needed source lower bound is just as hard as the original P vs NP/P/poly target.
\end{itemize}

\paragraph{Do not stop if.}
\begin{itemize}
\item The branch has not yet computed \(\beta\), q, and the induced source upper-bound exponent under \(\mathrm{SAT}\in \mathrm{SIZE}(n^{c})\).
\item A subpolynomial-query or nonblackbox carrier remains available and has not been audited.
\item The latest failure sharpened the threshold-collapse lemma into a formal no-go statement.
\item The route can still isolate advice leakage, nonuniformity, naturalness, or algebrization as the exact obstruction.
\end{itemize}

\paragraph{Required skills.}
\begin{itemize}
\item contradiction-route-auditor
\item entailment-gap-auditor
\item kb-deduce
\item kb-end-to-end-proof-critic
\item kb-hypothesis-refutation-loop
\item long-horizon-route-portfolio-manager
\item main-problem-alignment-auditor
\item next-experiment-queue-ranker
\item paper-contribution-readiness-judge
\item pvsnp-barrier-specificity-reviewer
\item pvsnp-conceptual-pressure-localizer
\item pvsnp-impossibility-theorem-forger
\item pvsnp-inference-quality-controller
\item pvsnp-proof-derivation-builder
\item pvsnp-qualitative-complexity-reasoner
\item pvsnp-quantifier-uniformity-reviewer
\item pvsnp-reduction-consistency-reviewer
\item pvsnp-symbolic-complexity-auditor
\end{itemize}

\subsection*{\#70 cycle-0147: Hardness-Magnification Bridge Realization}

\paragraph{Analysis.} Instantiate the closure verifier on a three-row QCTSR toy transcript: two counterfactual child rows, one EqSAT certificate forest, and one attempted cyclic merge; compute \(K_{Q}\), \(\Gamma_\mathrm{QCTSR}\), and the first-rejected-merge private-row falsifier explicitly.

\paragraph{Why this is a good target.} It directly tests the current black-box substitution cap by computing \(\beta\), q, and the induced target-exponent window. It keeps the barrier and model-boundary audit attached to the mathematical transfer rather than as after-the-fact commentary. It is executable as a bounded complexity calculation, CAS-backed exponent audit, or formal transfer check.

\paragraph{Risk.} Low drift risk: the item is framed as an explicit P vs NP transfer audit.

\paragraph{Branch control.} The item has high payoff and a concrete decision test; stopping now would be premature.

\paragraph{Premature-stop guard.} Do not stop before deriving the transfer inequality from \(\mathrm{SAT}\in \mathrm{SIZE}(n^{c})\) to the carrier upper bound, identifying whether \(\beta=0\), \(\beta>0\), or the argument is genuinely nonblackbox, and recording the barrier status of the proposed escape.

\paragraph{Payoff criteria.}
\begin{itemize}
\item For this route, state the exact barrier premise it uses or escapes: relativizing, natural, algebrizing, nonuniform, or advice-leaking; promotion requires the escape to be attached to the specific transfer theorem.
\item Compute \(\beta\) and q exactly; promotion requires \(\beta=0\), subpolynomial target queries, or a genuinely nonblackbox use of the SAT circuit.
\item Record relativization, natural-proofs, algebrization, nonuniformity, and advice-leakage status for the specific transfer theorem.
\item Control gate: name the source carrier or language, target class, and circuit or resource model before claiming progress.
\item Control gate: derive the exact transfer inequality from \(\mathrm{SAT}\in \mathrm{SIZE}(n^{c})\), NP subset P/poly, or the stated target assumption to the source-carrier upper bound.
\item Control gate: compute \(\beta\) and q exactly, distinguishing all-exponent consequences from bounded target-exponent windows.
\end{itemize}

\paragraph{Stop criteria.}
\begin{itemize}
\item Every named carrier in the branch has \(\beta>0\) and yields only a bounded target-exponent exclusion.
\item The branch cannot state a source lower-bound model, target class, and transfer inequality in the same notation.
\item The proposed escape is only a black-box substitution restatement with no new carrier or nonblackbox mechanism.
\item The barrier audit shows the needed source lower bound is just as hard as the original P vs NP/P/poly target.
\end{itemize}

\paragraph{Do not stop if.}
\begin{itemize}
\item The branch has not yet computed \(\beta\), q, and the induced source upper-bound exponent under \(\mathrm{SAT}\in \mathrm{SIZE}(n^{c})\).
\item A subpolynomial-query or nonblackbox carrier remains available and has not been audited.
\item The latest failure sharpened the threshold-collapse lemma into a formal no-go statement.
\item The route can still isolate advice leakage, nonuniformity, naturalness, or algebrization as the exact obstruction.
\end{itemize}

\paragraph{Required skills.}
\begin{itemize}
\item contradiction-route-auditor
\item entailment-gap-auditor
\item kb-deduce
\item kb-end-to-end-proof-critic
\item kb-hypothesis-refutation-loop
\item long-horizon-route-portfolio-manager
\item main-problem-alignment-auditor
\item next-experiment-queue-ranker
\item paper-contribution-readiness-judge
\item pvsnp-barrier-specificity-reviewer
\item pvsnp-conceptual-pressure-localizer
\item pvsnp-impossibility-theorem-forger
\item pvsnp-inference-quality-controller
\item pvsnp-proof-derivation-builder
\item pvsnp-qualitative-complexity-reasoner
\item pvsnp-quantifier-uniformity-reviewer
\item pvsnp-reduction-consistency-reviewer
\item pvsnp-symbolic-complexity-auditor
\end{itemize}

\subsection*{\#71 cycle-0146: Hardness-Magnification Bridge Realization}

\paragraph{Analysis.} Run the finite verifier on one projected-dDNNF/Tseitin-renaming example: list \(K_{Q}\) seeds, accepted quotient certificates, rank labels, computed \(\Gamma_\mathrm{QCTSR}\), and then delete one certificate to exhibit the first-rejected-merge private-row falsifier.

\paragraph{Why this is a good target.} It directly tests the current black-box substitution cap by computing \(\beta\), q, and the induced target-exponent window. It keeps the barrier and model-boundary audit attached to the mathematical transfer rather than as after-the-fact commentary. It is executable as a bounded complexity calculation, CAS-backed exponent audit, or formal transfer check.

\paragraph{Risk.} Low drift risk: the item is framed as an explicit P vs NP transfer audit.

\paragraph{Branch control.} The item has high payoff and a concrete decision test; stopping now would be premature.

\paragraph{Premature-stop guard.} Do not stop before deriving the transfer inequality from \(\mathrm{SAT}\in \mathrm{SIZE}(n^{c})\) to the carrier upper bound, identifying whether \(\beta=0\), \(\beta>0\), or the argument is genuinely nonblackbox, and recording the barrier status of the proposed escape.

\paragraph{Payoff criteria.}
\begin{itemize}
\item For the proof-complexity branch, separate fixed-family proof hardness from language hardness by naming the proof system, shortcut algorithm, and semantic lower-bound target.
\item Promotion requires a proof-system lower bound that survives the Gaussian-elimination or recognition shortcut; otherwise the result is a restricted-proof-system separation only.
\item For this route, state the exact barrier premise it uses or escapes: relativizing, natural, algebrizing, nonuniform, or advice-leaking; promotion requires the escape to be attached to the specific transfer theorem.
\item Compute \(\beta\) and q exactly; promotion requires \(\beta=0\), subpolynomial target queries, or a genuinely nonblackbox use of the SAT circuit.
\item Record relativization, natural-proofs, algebrization, nonuniformity, and advice-leakage status for the specific transfer theorem.
\item Control gate: name the source carrier or language, target class, and circuit or resource model before claiming progress.
\end{itemize}

\paragraph{Stop criteria.}
\begin{itemize}
\item Every named carrier in the branch has \(\beta>0\) and yields only a bounded target-exponent exclusion.
\item The branch cannot state a source lower-bound model, target class, and transfer inequality in the same notation.
\item The proposed escape is only a black-box substitution restatement with no new carrier or nonblackbox mechanism.
\item The barrier audit shows the needed source lower bound is just as hard as the original P vs NP/P/poly target.
\end{itemize}

\paragraph{Do not stop if.}
\begin{itemize}
\item The branch has not yet computed \(\beta\), q, and the induced source upper-bound exponent under \(\mathrm{SAT}\in \mathrm{SIZE}(n^{c})\).
\item A subpolynomial-query or nonblackbox carrier remains available and has not been audited.
\item The latest failure sharpened the threshold-collapse lemma into a formal no-go statement.
\item The route can still isolate advice leakage, nonuniformity, naturalness, or algebrization as the exact obstruction.
\end{itemize}

\paragraph{Required skills.}
\begin{itemize}
\item contradiction-route-auditor
\item entailment-gap-auditor
\item kb-deduce
\item kb-end-to-end-proof-critic
\item kb-hypothesis-refutation-loop
\item long-horizon-route-portfolio-manager
\item main-problem-alignment-auditor
\item next-experiment-queue-ranker
\item paper-contribution-readiness-judge
\item pvsnp-barrier-specificity-reviewer
\item pvsnp-conceptual-pressure-localizer
\item pvsnp-impossibility-theorem-forger
\item pvsnp-inference-quality-controller
\item pvsnp-proof-derivation-builder
\item pvsnp-qualitative-complexity-reasoner
\item pvsnp-quantifier-uniformity-reviewer
\item pvsnp-reduction-consistency-reviewer
\item pvsnp-symbolic-complexity-auditor
\end{itemize}

\subsection*{\#72 cycle-0145: Hardness-Magnification Bridge Realization}

\paragraph{Analysis.} Instantiate the finite verifier on one toy QCTSR proof with two EqSAT-equivalent counterfactual TSR children: compute \(K_{Q}\), ranks, \(B_{Q}\), and \(\Gamma_\mathrm{QCTSR}\); then delete one certificate edge and explicitly write the two-world private-row falsifier.

\paragraph{Why this is a good target.} It directly tests the current black-box substitution cap by computing \(\beta\), q, and the induced target-exponent window. It keeps the barrier and model-boundary audit attached to the mathematical transfer rather than as after-the-fact commentary. It is executable as a bounded complexity calculation, CAS-backed exponent audit, or formal transfer check.

\paragraph{Risk.} Low drift risk: the item is framed as an explicit P vs NP transfer audit.

\paragraph{Branch control.} The item has high payoff and a concrete decision test; stopping now would be premature.

\paragraph{Premature-stop guard.} Do not stop before deriving the transfer inequality from \(\mathrm{SAT}\in \mathrm{SIZE}(n^{c})\) to the carrier upper bound, identifying whether \(\beta=0\), \(\beta>0\), or the argument is genuinely nonblackbox, and recording the barrier status of the proposed escape.

\paragraph{Payoff criteria.}
\begin{itemize}
\item For this route, state the exact barrier premise it uses or escapes: relativizing, natural, algebrizing, nonuniform, or advice-leaking; promotion requires the escape to be attached to the specific transfer theorem.
\item Compute \(\beta\) and q exactly; promotion requires \(\beta=0\), subpolynomial target queries, or a genuinely nonblackbox use of the SAT circuit.
\item Record relativization, natural-proofs, algebrization, nonuniformity, and advice-leakage status for the specific transfer theorem.
\item Control gate: name the source carrier or language, target class, and circuit or resource model before claiming progress.
\item Control gate: derive the exact transfer inequality from \(\mathrm{SAT}\in \mathrm{SIZE}(n^{c})\), NP subset P/poly, or the stated target assumption to the source-carrier upper bound.
\item Control gate: compute \(\beta\) and q exactly, distinguishing all-exponent consequences from bounded target-exponent windows.
\end{itemize}

\paragraph{Stop criteria.}
\begin{itemize}
\item Every named carrier in the branch has \(\beta>0\) and yields only a bounded target-exponent exclusion.
\item The branch cannot state a source lower-bound model, target class, and transfer inequality in the same notation.
\item The proposed escape is only a black-box substitution restatement with no new carrier or nonblackbox mechanism.
\item The barrier audit shows the needed source lower bound is just as hard as the original P vs NP/P/poly target.
\end{itemize}

\paragraph{Do not stop if.}
\begin{itemize}
\item The branch has not yet computed \(\beta\), q, and the induced source upper-bound exponent under \(\mathrm{SAT}\in \mathrm{SIZE}(n^{c})\).
\item A subpolynomial-query or nonblackbox carrier remains available and has not been audited.
\item The latest failure sharpened the threshold-collapse lemma into a formal no-go statement.
\item The route can still isolate advice leakage, nonuniformity, naturalness, or algebrization as the exact obstruction.
\end{itemize}

\paragraph{Required skills.}
\begin{itemize}
\item contradiction-route-auditor
\item entailment-gap-auditor
\item kb-deduce
\item kb-end-to-end-proof-critic
\item kb-hypothesis-refutation-loop
\item long-horizon-route-portfolio-manager
\item main-problem-alignment-auditor
\item next-experiment-queue-ranker
\item paper-contribution-readiness-judge
\item pvsnp-barrier-specificity-reviewer
\item pvsnp-conceptual-pressure-localizer
\item pvsnp-impossibility-theorem-forger
\item pvsnp-inference-quality-controller
\item pvsnp-proof-derivation-builder
\item pvsnp-qualitative-complexity-reasoner
\item pvsnp-quantifier-uniformity-reviewer
\item pvsnp-reduction-consistency-reviewer
\item pvsnp-symbolic-complexity-auditor
\end{itemize}

\subsection*{\#73 cycle-0144: Hardness-Magnification Bridge Realization: ranked QCTSR ledger validation}

\paragraph{Analysis.} Write the referee-checkable \(K_{Q}\) closure pseudocode with rank ordering, then test it on a two-child Tseitin-equivalent TSR merge and on the same example with the certificate omitted to exhibit the first-unranked-merge falsifier.

\paragraph{Why this is a good target.} It directly tests the current black-box substitution cap by computing \(\beta\), q, and the induced target-exponent window. It keeps the barrier and model-boundary audit attached to the mathematical transfer rather than as after-the-fact commentary. It is executable as a bounded complexity calculation, CAS-backed exponent audit, or formal transfer check.

\paragraph{Risk.} Low drift risk: the item is framed as an explicit P vs NP transfer audit.

\paragraph{Branch control.} The item has high payoff and a concrete decision test; stopping now would be premature.

\paragraph{Premature-stop guard.} Do not stop before deriving the transfer inequality from \(\mathrm{SAT}\in \mathrm{SIZE}(n^{c})\) to the carrier upper bound, identifying whether \(\beta=0\), \(\beta>0\), or the argument is genuinely nonblackbox, and recording the barrier status of the proposed escape.

\paragraph{Payoff criteria.}
\begin{itemize}
\item For the proof-complexity branch, separate fixed-family proof hardness from language hardness by naming the proof system, shortcut algorithm, and semantic lower-bound target.
\item Promotion requires a proof-system lower bound that survives the Gaussian-elimination or recognition shortcut; otherwise the result is a restricted-proof-system separation only.
\item For this route, state the exact barrier premise it uses or escapes: relativizing, natural, algebrizing, nonuniform, or advice-leaking; promotion requires the escape to be attached to the specific transfer theorem.
\item Compute \(\beta\) and q exactly; promotion requires \(\beta=0\), subpolynomial target queries, or a genuinely nonblackbox use of the SAT circuit.
\item Record relativization, natural-proofs, algebrization, nonuniformity, and advice-leakage status for the specific transfer theorem.
\item Control gate: name the source carrier or language, target class, and circuit or resource model before claiming progress.
\end{itemize}

\paragraph{Stop criteria.}
\begin{itemize}
\item Every named carrier in the branch has \(\beta>0\) and yields only a bounded target-exponent exclusion.
\item The branch cannot state a source lower-bound model, target class, and transfer inequality in the same notation.
\item The proposed escape is only a black-box substitution restatement with no new carrier or nonblackbox mechanism.
\item The barrier audit shows the needed source lower bound is just as hard as the original P vs NP/P/poly target.
\end{itemize}

\paragraph{Do not stop if.}
\begin{itemize}
\item The branch has not yet computed \(\beta\), q, and the induced source upper-bound exponent under \(\mathrm{SAT}\in \mathrm{SIZE}(n^{c})\).
\item A subpolynomial-query or nonblackbox carrier remains available and has not been audited.
\item The latest failure sharpened the threshold-collapse lemma into a formal no-go statement.
\item The route can still isolate advice leakage, nonuniformity, naturalness, or algebrization as the exact obstruction.
\end{itemize}

\paragraph{Required skills.}
\begin{itemize}
\item contradiction-route-auditor
\item entailment-gap-auditor
\item kb-deduce
\item kb-end-to-end-proof-critic
\item kb-hypothesis-refutation-loop
\item long-horizon-route-portfolio-manager
\item main-problem-alignment-auditor
\item next-experiment-queue-ranker
\item paper-contribution-readiness-judge
\item pvsnp-barrier-specificity-reviewer
\item pvsnp-conceptual-pressure-localizer
\item pvsnp-impossibility-theorem-forger
\item pvsnp-inference-quality-controller
\item pvsnp-proof-derivation-builder
\item pvsnp-qualitative-complexity-reasoner
\item pvsnp-quantifier-uniformity-reviewer
\item pvsnp-reduction-consistency-reviewer
\item pvsnp-symbolic-complexity-auditor
\end{itemize}

\subsection*{\#74 cycle-0142: Hardness-Magnification Bridge Realization}

\paragraph{Analysis.} Write the finite verifier for \(K_{Q}\) and ranked \(F_{Q}\) on a toy projected-dDNNF plus CTSR example with one duplicate child pair; compute \(\Gamma\) savings, then remove one certificate edge to exhibit the private-row falsifier.

\paragraph{Why this is a good target.} It directly tests the current black-box substitution cap by computing \(\beta\), q, and the induced target-exponent window. It keeps the barrier and model-boundary audit attached to the mathematical transfer rather than as after-the-fact commentary. It is executable as a bounded complexity calculation, CAS-backed exponent audit, or formal transfer check.

\paragraph{Risk.} Low drift risk: the item is framed as an explicit P vs NP transfer audit.

\paragraph{Branch control.} The item has high payoff and a concrete decision test; stopping now would be premature.

\paragraph{Premature-stop guard.} Do not stop before deriving the transfer inequality from \(\mathrm{SAT}\in \mathrm{SIZE}(n^{c})\) to the carrier upper bound, identifying whether \(\beta=0\), \(\beta>0\), or the argument is genuinely nonblackbox, and recording the barrier status of the proposed escape.

\paragraph{Payoff criteria.}
\begin{itemize}
\item For this route, state the exact barrier premise it uses or escapes: relativizing, natural, algebrizing, nonuniform, or advice-leaking; promotion requires the escape to be attached to the specific transfer theorem.
\item Compute \(\beta\) and q exactly; promotion requires \(\beta=0\), subpolynomial target queries, or a genuinely nonblackbox use of the SAT circuit.
\item Record relativization, natural-proofs, algebrization, nonuniformity, and advice-leakage status for the specific transfer theorem.
\item Control gate: name the source carrier or language, target class, and circuit or resource model before claiming progress.
\item Control gate: derive the exact transfer inequality from \(\mathrm{SAT}\in \mathrm{SIZE}(n^{c})\), NP subset P/poly, or the stated target assumption to the source-carrier upper bound.
\item Control gate: compute \(\beta\) and q exactly, distinguishing all-exponent consequences from bounded target-exponent windows.
\end{itemize}

\paragraph{Stop criteria.}
\begin{itemize}
\item Every named carrier in the branch has \(\beta>0\) and yields only a bounded target-exponent exclusion.
\item The branch cannot state a source lower-bound model, target class, and transfer inequality in the same notation.
\item The proposed escape is only a black-box substitution restatement with no new carrier or nonblackbox mechanism.
\item The barrier audit shows the needed source lower bound is just as hard as the original P vs NP/P/poly target.
\end{itemize}

\paragraph{Do not stop if.}
\begin{itemize}
\item The branch has not yet computed \(\beta\), q, and the induced source upper-bound exponent under \(\mathrm{SAT}\in \mathrm{SIZE}(n^{c})\).
\item A subpolynomial-query or nonblackbox carrier remains available and has not been audited.
\item The latest failure sharpened the threshold-collapse lemma into a formal no-go statement.
\item The route can still isolate advice leakage, nonuniformity, naturalness, or algebrization as the exact obstruction.
\end{itemize}

\paragraph{Required skills.}
\begin{itemize}
\item contradiction-route-auditor
\item entailment-gap-auditor
\item kb-deduce
\item kb-end-to-end-proof-critic
\item kb-hypothesis-refutation-loop
\item long-horizon-route-portfolio-manager
\item main-problem-alignment-auditor
\item next-experiment-queue-ranker
\item paper-contribution-readiness-judge
\item pvsnp-barrier-specificity-reviewer
\item pvsnp-conceptual-pressure-localizer
\item pvsnp-impossibility-theorem-forger
\item pvsnp-inference-quality-controller
\item pvsnp-proof-derivation-builder
\item pvsnp-qualitative-complexity-reasoner
\item pvsnp-quantifier-uniformity-reviewer
\item pvsnp-reduction-consistency-reviewer
\item pvsnp-symbolic-complexity-auditor
\end{itemize}

\subsection*{\#75 cycle-0140: Hardness-Magnification Bridge Realization: QCTSR Certified-Quotient Parser Boundary}

\paragraph{Analysis.} Instantiate \(K_{Q}\) on a toy projected-dDNNF FormulaEval proof with two CTSR levels and three EqSAT quotient classes; compute |RepIdx|, \(C_{Eq}\), \(\Gamma_\mathrm{QCTSR}\), then delete one quotient certificate to exhibit the first-hidden-merge falsifier.

\paragraph{Why this is a good target.} It directly tests the current black-box substitution cap by computing \(\beta\), q, and the induced target-exponent window. It keeps the barrier and model-boundary audit attached to the mathematical transfer rather than as after-the-fact commentary. It is executable as a bounded complexity calculation, CAS-backed exponent audit, or formal transfer check.

\paragraph{Risk.} Low drift risk: the item is framed as an explicit P vs NP transfer audit.

\paragraph{Branch control.} The item has high payoff and a concrete decision test; stopping now would be premature.

\paragraph{Premature-stop guard.} Do not stop before deriving the transfer inequality from \(\mathrm{SAT}\in \mathrm{SIZE}(n^{c})\) to the carrier upper bound, identifying whether \(\beta=0\), \(\beta>0\), or the argument is genuinely nonblackbox, and recording the barrier status of the proposed escape.

\paragraph{Payoff criteria.}
\begin{itemize}
\item For this route, state the exact barrier premise it uses or escapes: relativizing, natural, algebrizing, nonuniform, or advice-leaking; promotion requires the escape to be attached to the specific transfer theorem.
\item Compute \(\beta\) and q exactly; promotion requires \(\beta=0\), subpolynomial target queries, or a genuinely nonblackbox use of the SAT circuit.
\item Record relativization, natural-proofs, algebrization, nonuniformity, and advice-leakage status for the specific transfer theorem.
\item Control gate: name the source carrier or language, target class, and circuit or resource model before claiming progress.
\item Control gate: derive the exact transfer inequality from \(\mathrm{SAT}\in \mathrm{SIZE}(n^{c})\), NP subset P/poly, or the stated target assumption to the source-carrier upper bound.
\item Control gate: compute \(\beta\) and q exactly, distinguishing all-exponent consequences from bounded target-exponent windows.
\end{itemize}

\paragraph{Stop criteria.}
\begin{itemize}
\item Every named carrier in the branch has \(\beta>0\) and yields only a bounded target-exponent exclusion.
\item The branch cannot state a source lower-bound model, target class, and transfer inequality in the same notation.
\item The proposed escape is only a black-box substitution restatement with no new carrier or nonblackbox mechanism.
\item The barrier audit shows the needed source lower bound is just as hard as the original P vs NP/P/poly target.
\end{itemize}

\paragraph{Do not stop if.}
\begin{itemize}
\item The branch has not yet computed \(\beta\), q, and the induced source upper-bound exponent under \(\mathrm{SAT}\in \mathrm{SIZE}(n^{c})\).
\item A subpolynomial-query or nonblackbox carrier remains available and has not been audited.
\item The latest failure sharpened the threshold-collapse lemma into a formal no-go statement.
\item The route can still isolate advice leakage, nonuniformity, naturalness, or algebrization as the exact obstruction.
\end{itemize}

\paragraph{Required skills.}
\begin{itemize}
\item contradiction-route-auditor
\item entailment-gap-auditor
\item kb-deduce
\item kb-end-to-end-proof-critic
\item kb-hypothesis-refutation-loop
\item long-horizon-route-portfolio-manager
\item main-problem-alignment-auditor
\item next-experiment-queue-ranker
\item paper-contribution-readiness-judge
\item pvsnp-barrier-specificity-reviewer
\item pvsnp-conceptual-pressure-localizer
\item pvsnp-impossibility-theorem-forger
\item pvsnp-inference-quality-controller
\item pvsnp-proof-derivation-builder
\item pvsnp-qualitative-complexity-reasoner
\item pvsnp-quantifier-uniformity-reviewer
\item pvsnp-reduction-consistency-reviewer
\item pvsnp-symbolic-complexity-auditor
\end{itemize}

\subsection*{\#76 cycle-0139: Hardness-Magnification Bridge Realization}

\paragraph{Analysis.} Instantiate \(\Gamma_\mathrm{CTSR}^=\) on a projected-dDNNF FormulaEval proof with two Tseitin-equivalent child formulas: compute K before and after certified EqSAT quotienting, then remove one certificate and exhibit the first uncertified quotient-collision private-row falsifier.

\paragraph{Why this is a good target.} It directly tests the current black-box substitution cap by computing \(\beta\), q, and the induced target-exponent window. It keeps the barrier and model-boundary audit attached to the mathematical transfer rather than as after-the-fact commentary. It is executable as a bounded complexity calculation, CAS-backed exponent audit, or formal transfer check.

\paragraph{Risk.} Low drift risk: the item is framed as an explicit P vs NP transfer audit.

\paragraph{Branch control.} The item has high payoff and a concrete decision test; stopping now would be premature.

\paragraph{Premature-stop guard.} Do not stop before deriving the transfer inequality from \(\mathrm{SAT}\in \mathrm{SIZE}(n^{c})\) to the carrier upper bound, identifying whether \(\beta=0\), \(\beta>0\), or the argument is genuinely nonblackbox, and recording the barrier status of the proposed escape.

\paragraph{Payoff criteria.}
\begin{itemize}
\item For the proof-complexity branch, separate fixed-family proof hardness from language hardness by naming the proof system, shortcut algorithm, and semantic lower-bound target.
\item Promotion requires a proof-system lower bound that survives the Gaussian-elimination or recognition shortcut; otherwise the result is a restricted-proof-system separation only.
\item For this route, state the exact barrier premise it uses or escapes: relativizing, natural, algebrizing, nonuniform, or advice-leaking; promotion requires the escape to be attached to the specific transfer theorem.
\item Compute \(\beta\) and q exactly; promotion requires \(\beta=0\), subpolynomial target queries, or a genuinely nonblackbox use of the SAT circuit.
\item Record relativization, natural-proofs, algebrization, nonuniformity, and advice-leakage status for the specific transfer theorem.
\item Control gate: name the source carrier or language, target class, and circuit or resource model before claiming progress.
\end{itemize}

\paragraph{Stop criteria.}
\begin{itemize}
\item Every named carrier in the branch has \(\beta>0\) and yields only a bounded target-exponent exclusion.
\item The branch cannot state a source lower-bound model, target class, and transfer inequality in the same notation.
\item The proposed escape is only a black-box substitution restatement with no new carrier or nonblackbox mechanism.
\item The barrier audit shows the needed source lower bound is just as hard as the original P vs NP/P/poly target.
\end{itemize}

\paragraph{Do not stop if.}
\begin{itemize}
\item The branch has not yet computed \(\beta\), q, and the induced source upper-bound exponent under \(\mathrm{SAT}\in \mathrm{SIZE}(n^{c})\).
\item A subpolynomial-query or nonblackbox carrier remains available and has not been audited.
\item The latest failure sharpened the threshold-collapse lemma into a formal no-go statement.
\item The route can still isolate advice leakage, nonuniformity, naturalness, or algebrization as the exact obstruction.
\end{itemize}

\paragraph{Required skills.}
\begin{itemize}
\item contradiction-route-auditor
\item entailment-gap-auditor
\item kb-deduce
\item kb-end-to-end-proof-critic
\item kb-hypothesis-refutation-loop
\item long-horizon-route-portfolio-manager
\item main-problem-alignment-auditor
\item next-experiment-queue-ranker
\item paper-contribution-readiness-judge
\item pvsnp-barrier-specificity-reviewer
\item pvsnp-conceptual-pressure-localizer
\item pvsnp-impossibility-theorem-forger
\item pvsnp-inference-quality-controller
\item pvsnp-proof-derivation-builder
\item pvsnp-qualitative-complexity-reasoner
\item pvsnp-quantifier-uniformity-reviewer
\item pvsnp-reduction-consistency-reviewer
\item pvsnp-symbolic-complexity-auditor
\end{itemize}

\subsection*{\#77 cycle-0137: Hardness-Magnification Bridge Realization}

\paragraph{Analysis.} Write the standalone L137 first-hidden-child theorem: specify the CTSR bit fields, define the two row-oracle worlds, prove transcript equality on \(K_\mathrm{CTSR}\) minus r, prove final parity disagreement, and compute the \(\Gamma_\mathrm{CTSR}\) threshold for \(\kappa+\theta<\epsilon\) with q-trichotomy support density.

\paragraph{Why this is a good target.} It directly tests the current black-box substitution cap by computing \(\beta\), q, and the induced target-exponent window. It keeps the barrier and model-boundary audit attached to the mathematical transfer rather than as after-the-fact commentary. It is executable as a bounded complexity calculation, CAS-backed exponent audit, or formal transfer check.

\paragraph{Risk.} Low drift risk: the item is framed as an explicit P vs NP transfer audit.

\paragraph{Branch control.} The item has high payoff and a concrete decision test; stopping now would be premature.

\paragraph{Premature-stop guard.} Do not stop before deriving the transfer inequality from \(\mathrm{SAT}\in \mathrm{SIZE}(n^{c})\) to the carrier upper bound, identifying whether \(\beta=0\), \(\beta>0\), or the argument is genuinely nonblackbox, and recording the barrier status of the proposed escape.

\paragraph{Payoff criteria.}
\begin{itemize}
\item For this route, state the exact barrier premise it uses or escapes: relativizing, natural, algebrizing, nonuniform, or advice-leaking; promotion requires the escape to be attached to the specific transfer theorem.
\item Compute \(\beta\) and q exactly; promotion requires \(\beta=0\), subpolynomial target queries, or a genuinely nonblackbox use of the SAT circuit.
\item Record relativization, natural-proofs, algebrization, nonuniformity, and advice-leakage status for the specific transfer theorem.
\item Control gate: name the source carrier or language, target class, and circuit or resource model before claiming progress.
\item Control gate: derive the exact transfer inequality from \(\mathrm{SAT}\in \mathrm{SIZE}(n^{c})\), NP subset P/poly, or the stated target assumption to the source-carrier upper bound.
\item Control gate: compute \(\beta\) and q exactly, distinguishing all-exponent consequences from bounded target-exponent windows.
\end{itemize}

\paragraph{Stop criteria.}
\begin{itemize}
\item Every named carrier in the branch has \(\beta>0\) and yields only a bounded target-exponent exclusion.
\item The branch cannot state a source lower-bound model, target class, and transfer inequality in the same notation.
\item The proposed escape is only a black-box substitution restatement with no new carrier or nonblackbox mechanism.
\item The barrier audit shows the needed source lower bound is just as hard as the original P vs NP/P/poly target.
\end{itemize}

\paragraph{Do not stop if.}
\begin{itemize}
\item The branch has not yet computed \(\beta\), q, and the induced source upper-bound exponent under \(\mathrm{SAT}\in \mathrm{SIZE}(n^{c})\).
\item A subpolynomial-query or nonblackbox carrier remains available and has not been audited.
\item The latest failure sharpened the threshold-collapse lemma into a formal no-go statement.
\item The route can still isolate advice leakage, nonuniformity, naturalness, or algebrization as the exact obstruction.
\end{itemize}

\paragraph{Required skills.}
\begin{itemize}
\item contradiction-route-auditor
\item entailment-gap-auditor
\item kb-deduce
\item kb-end-to-end-proof-critic
\item kb-hypothesis-refutation-loop
\item long-horizon-route-portfolio-manager
\item main-problem-alignment-auditor
\item next-experiment-queue-ranker
\item paper-contribution-readiness-judge
\item pvsnp-barrier-specificity-reviewer
\item pvsnp-conceptual-pressure-localizer
\item pvsnp-impossibility-theorem-forger
\item pvsnp-inference-quality-controller
\item pvsnp-proof-derivation-builder
\item pvsnp-qualitative-complexity-reasoner
\item pvsnp-quantifier-uniformity-reviewer
\item pvsnp-reduction-consistency-reviewer
\item pvsnp-symbolic-complexity-auditor
\end{itemize}

\subsection*{\#78 cycle-0136: Hardness-Magnification Bridge Realization}

\paragraph{Analysis.} Select one concrete full semantic pruning rule beyond CTSR declarations, preferably adaptive Tseitin branch pruning with real SAT contradiction tests, and attempt either a charged \(RealSATInfluenceRealizer_{N}^{arb}\) or a minimal first-unchargeable-row countermodel.

\paragraph{Why this is a good target.} It directly tests the current black-box substitution cap by computing \(\beta\), q, and the induced target-exponent window. It keeps the barrier and model-boundary audit attached to the mathematical transfer rather than as after-the-fact commentary. It is executable as a bounded complexity calculation, CAS-backed exponent audit, or formal transfer check.

\paragraph{Risk.} Low drift risk: the item is framed as an explicit P vs NP transfer audit.

\paragraph{Branch control.} The item has high payoff and a concrete decision test; stopping now would be premature.

\paragraph{Premature-stop guard.} Do not stop before deriving the transfer inequality from \(\mathrm{SAT}\in \mathrm{SIZE}(n^{c})\) to the carrier upper bound, identifying whether \(\beta=0\), \(\beta>0\), or the argument is genuinely nonblackbox, and recording the barrier status of the proposed escape.

\paragraph{Payoff criteria.}
\begin{itemize}
\item For the proof-complexity branch, separate fixed-family proof hardness from language hardness by naming the proof system, shortcut algorithm, and semantic lower-bound target.
\item Promotion requires a proof-system lower bound that survives the Gaussian-elimination or recognition shortcut; otherwise the result is a restricted-proof-system separation only.
\item For this route, state the exact barrier premise it uses or escapes: relativizing, natural, algebrizing, nonuniform, or advice-leaking; promotion requires the escape to be attached to the specific transfer theorem.
\item Compute \(\beta\) and q exactly; promotion requires \(\beta=0\), subpolynomial target queries, or a genuinely nonblackbox use of the SAT circuit.
\item Record relativization, natural-proofs, algebrization, nonuniformity, and advice-leakage status for the specific transfer theorem.
\item Control gate: name the source carrier or language, target class, and circuit or resource model before claiming progress.
\end{itemize}

\paragraph{Stop criteria.}
\begin{itemize}
\item Every named carrier in the branch has \(\beta>0\) and yields only a bounded target-exponent exclusion.
\item The branch cannot state a source lower-bound model, target class, and transfer inequality in the same notation.
\item The proposed escape is only a black-box substitution restatement with no new carrier or nonblackbox mechanism.
\item The barrier audit shows the needed source lower bound is just as hard as the original P vs NP/P/poly target.
\end{itemize}

\paragraph{Do not stop if.}
\begin{itemize}
\item The branch has not yet computed \(\beta\), q, and the induced source upper-bound exponent under \(\mathrm{SAT}\in \mathrm{SIZE}(n^{c})\).
\item A subpolynomial-query or nonblackbox carrier remains available and has not been audited.
\item The latest failure sharpened the threshold-collapse lemma into a formal no-go statement.
\item The route can still isolate advice leakage, nonuniformity, naturalness, or algebrization as the exact obstruction.
\end{itemize}

\paragraph{Required skills.}
\begin{itemize}
\item contradiction-route-auditor
\item entailment-gap-auditor
\item kb-deduce
\item kb-end-to-end-proof-critic
\item kb-hypothesis-refutation-loop
\item long-horizon-route-portfolio-manager
\item main-problem-alignment-auditor
\item next-experiment-queue-ranker
\item paper-contribution-readiness-judge
\item pvsnp-barrier-specificity-reviewer
\item pvsnp-conceptual-pressure-localizer
\item pvsnp-impossibility-theorem-forger
\item pvsnp-inference-quality-controller
\item pvsnp-proof-derivation-builder
\item pvsnp-qualitative-complexity-reasoner
\item pvsnp-quantifier-uniformity-reviewer
\item pvsnp-reduction-consistency-reviewer
\item pvsnp-symbolic-complexity-auditor
\end{itemize}

\subsection*{\#79 cycle-0134: Hardness-Magnification Bridge Realization}

\paragraph{Analysis.} Instantiate the counterfactual DTSR closure on a two-level nested TSR proof with one Tseitin seed family; compute the exact child count, \(K_\mathrm{CTSR}\) fixed point, \(\Gamma_\mathrm{CTSR}\) ledger, and first-unlisted-child countermodel transcript.

\paragraph{Why this is a good target.} It directly tests the current black-box substitution cap by computing \(\beta\), q, and the induced target-exponent window. It keeps the barrier and model-boundary audit attached to the mathematical transfer rather than as after-the-fact commentary. It is executable as a bounded complexity calculation, CAS-backed exponent audit, or formal transfer check.

\paragraph{Risk.} Low drift risk: the item is framed as an explicit P vs NP transfer audit.

\paragraph{Branch control.} The item has high payoff and a concrete decision test; stopping now would be premature.

\paragraph{Premature-stop guard.} Do not stop before deriving the transfer inequality from \(\mathrm{SAT}\in \mathrm{SIZE}(n^{c})\) to the carrier upper bound, identifying whether \(\beta=0\), \(\beta>0\), or the argument is genuinely nonblackbox, and recording the barrier status of the proposed escape.

\paragraph{Payoff criteria.}
\begin{itemize}
\item For the proof-complexity branch, separate fixed-family proof hardness from language hardness by naming the proof system, shortcut algorithm, and semantic lower-bound target.
\item Promotion requires a proof-system lower bound that survives the Gaussian-elimination or recognition shortcut; otherwise the result is a restricted-proof-system separation only.
\item For this route, state the exact barrier premise it uses or escapes: relativizing, natural, algebrizing, nonuniform, or advice-leaking; promotion requires the escape to be attached to the specific transfer theorem.
\item Compute \(\beta\) and q exactly; promotion requires \(\beta=0\), subpolynomial target queries, or a genuinely nonblackbox use of the SAT circuit.
\item Record relativization, natural-proofs, algebrization, nonuniformity, and advice-leakage status for the specific transfer theorem.
\item Control gate: name the source carrier or language, target class, and circuit or resource model before claiming progress.
\end{itemize}

\paragraph{Stop criteria.}
\begin{itemize}
\item Every named carrier in the branch has \(\beta>0\) and yields only a bounded target-exponent exclusion.
\item The branch cannot state a source lower-bound model, target class, and transfer inequality in the same notation.
\item The proposed escape is only a black-box substitution restatement with no new carrier or nonblackbox mechanism.
\item The barrier audit shows the needed source lower bound is just as hard as the original P vs NP/P/poly target.
\end{itemize}

\paragraph{Do not stop if.}
\begin{itemize}
\item The branch has not yet computed \(\beta\), q, and the induced source upper-bound exponent under \(\mathrm{SAT}\in \mathrm{SIZE}(n^{c})\).
\item A subpolynomial-query or nonblackbox carrier remains available and has not been audited.
\item The latest failure sharpened the threshold-collapse lemma into a formal no-go statement.
\item The route can still isolate advice leakage, nonuniformity, naturalness, or algebrization as the exact obstruction.
\end{itemize}

\paragraph{Required skills.}
\begin{itemize}
\item contradiction-route-auditor
\item entailment-gap-auditor
\item kb-deduce
\item kb-end-to-end-proof-critic
\item kb-hypothesis-refutation-loop
\item long-horizon-route-portfolio-manager
\item main-problem-alignment-auditor
\item next-experiment-queue-ranker
\item paper-contribution-readiness-judge
\item pvsnp-barrier-specificity-reviewer
\item pvsnp-conceptual-pressure-localizer
\item pvsnp-impossibility-theorem-forger
\item pvsnp-inference-quality-controller
\item pvsnp-proof-derivation-builder
\item pvsnp-qualitative-complexity-reasoner
\item pvsnp-quantifier-uniformity-reviewer
\item pvsnp-reduction-consistency-reviewer
\item pvsnp-symbolic-complexity-auditor
\end{itemize}

\subsection*{\#80 cycle-0133: Hardness-Magnification Bridge Realization}

\paragraph{Analysis.} Attack \(RealSATInfluenceRealizer_{N}^{arb}\) for one concrete semantic pruning rule: bounded Tseitin self-reduction with a SAT-truth branch-elimination premise. Either produce a uniform charged list of all influencing children or formalize the first semantic line where the private-row countermodel survives.

\paragraph{Why this is a good target.} It directly tests the current black-box substitution cap by computing \(\beta\), q, and the induced target-exponent window. It keeps the barrier and model-boundary audit attached to the mathematical transfer rather than as after-the-fact commentary. It is executable as a bounded complexity calculation, CAS-backed exponent audit, or formal transfer check.

\paragraph{Risk.} Low drift risk: the item is framed as an explicit P vs NP transfer audit.

\paragraph{Branch control.} The item has high payoff and a concrete decision test; stopping now would be premature.

\paragraph{Premature-stop guard.} Do not stop before deriving the transfer inequality from \(\mathrm{SAT}\in \mathrm{SIZE}(n^{c})\) to the carrier upper bound, identifying whether \(\beta=0\), \(\beta>0\), or the argument is genuinely nonblackbox, and recording the barrier status of the proposed escape.

\paragraph{Payoff criteria.}
\begin{itemize}
\item For the proof-complexity branch, separate fixed-family proof hardness from language hardness by naming the proof system, shortcut algorithm, and semantic lower-bound target.
\item Promotion requires a proof-system lower bound that survives the Gaussian-elimination or recognition shortcut; otherwise the result is a restricted-proof-system separation only.
\item For this route, state the exact barrier premise it uses or escapes: relativizing, natural, algebrizing, nonuniform, or advice-leaking; promotion requires the escape to be attached to the specific transfer theorem.
\item Compute \(\beta\) and q exactly; promotion requires \(\beta=0\), subpolynomial target queries, or a genuinely nonblackbox use of the SAT circuit.
\item Record relativization, natural-proofs, algebrization, nonuniformity, and advice-leakage status for the specific transfer theorem.
\item Control gate: name the source carrier or language, target class, and circuit or resource model before claiming progress.
\end{itemize}

\paragraph{Stop criteria.}
\begin{itemize}
\item Every named carrier in the branch has \(\beta>0\) and yields only a bounded target-exponent exclusion.
\item The branch cannot state a source lower-bound model, target class, and transfer inequality in the same notation.
\item The proposed escape is only a black-box substitution restatement with no new carrier or nonblackbox mechanism.
\item The barrier audit shows the needed source lower bound is just as hard as the original P vs NP/P/poly target.
\end{itemize}

\paragraph{Do not stop if.}
\begin{itemize}
\item The branch has not yet computed \(\beta\), q, and the induced source upper-bound exponent under \(\mathrm{SAT}\in \mathrm{SIZE}(n^{c})\).
\item A subpolynomial-query or nonblackbox carrier remains available and has not been audited.
\item The latest failure sharpened the threshold-collapse lemma into a formal no-go statement.
\item The route can still isolate advice leakage, nonuniformity, naturalness, or algebrization as the exact obstruction.
\end{itemize}

\paragraph{Required skills.}
\begin{itemize}
\item contradiction-route-auditor
\item entailment-gap-auditor
\item kb-deduce
\item kb-end-to-end-proof-critic
\item kb-hypothesis-refutation-loop
\item long-horizon-route-portfolio-manager
\item main-problem-alignment-auditor
\item next-experiment-queue-ranker
\item paper-contribution-readiness-judge
\item pvsnp-barrier-specificity-reviewer
\item pvsnp-conceptual-pressure-localizer
\item pvsnp-impossibility-theorem-forger
\item pvsnp-inference-quality-controller
\item pvsnp-proof-derivation-builder
\item pvsnp-qualitative-complexity-reasoner
\item pvsnp-quantifier-uniformity-reviewer
\item pvsnp-reduction-consistency-reviewer
\item pvsnp-symbolic-complexity-auditor
\end{itemize}

\subsection*{\#81 cycle-0132: Hardness-Magnification Bridge Realization}

\paragraph{Analysis.} Write the paper-ready L132 proof block: pseudocode for K\_DTSR(pi), a table of DTSR bit fields and verifier checks, the five line-induction cases, and the first-unlisted-child countermodel; then test one projected-dDNNF+DTSR positive example and one unread-enable-bit hidden-pruning counterexample.

\paragraph{Why this is a good target.} It directly tests the current black-box substitution cap by computing \(\beta\), q, and the induced target-exponent window. It keeps the barrier and model-boundary audit attached to the mathematical transfer rather than as after-the-fact commentary. It is executable as a bounded complexity calculation, CAS-backed exponent audit, or formal transfer check.

\paragraph{Risk.} Low drift risk: the item is framed as an explicit P vs NP transfer audit.

\paragraph{Branch control.} The item has high payoff and a concrete decision test; stopping now would be premature.

\paragraph{Premature-stop guard.} Do not stop before deriving the transfer inequality from \(\mathrm{SAT}\in \mathrm{SIZE}(n^{c})\) to the carrier upper bound, identifying whether \(\beta=0\), \(\beta>0\), or the argument is genuinely nonblackbox, and recording the barrier status of the proposed escape.

\paragraph{Payoff criteria.}
\begin{itemize}
\item For this route, state the exact barrier premise it uses or escapes: relativizing, natural, algebrizing, nonuniform, or advice-leaking; promotion requires the escape to be attached to the specific transfer theorem.
\item Compute \(\beta\) and q exactly; promotion requires \(\beta=0\), subpolynomial target queries, or a genuinely nonblackbox use of the SAT circuit.
\item Record relativization, natural-proofs, algebrization, nonuniformity, and advice-leakage status for the specific transfer theorem.
\item Control gate: name the source carrier or language, target class, and circuit or resource model before claiming progress.
\item Control gate: derive the exact transfer inequality from \(\mathrm{SAT}\in \mathrm{SIZE}(n^{c})\), NP subset P/poly, or the stated target assumption to the source-carrier upper bound.
\item Control gate: compute \(\beta\) and q exactly, distinguishing all-exponent consequences from bounded target-exponent windows.
\end{itemize}

\paragraph{Stop criteria.}
\begin{itemize}
\item Every named carrier in the branch has \(\beta>0\) and yields only a bounded target-exponent exclusion.
\item The branch cannot state a source lower-bound model, target class, and transfer inequality in the same notation.
\item The proposed escape is only a black-box substitution restatement with no new carrier or nonblackbox mechanism.
\item The barrier audit shows the needed source lower bound is just as hard as the original P vs NP/P/poly target.
\end{itemize}

\paragraph{Do not stop if.}
\begin{itemize}
\item The branch has not yet computed \(\beta\), q, and the induced source upper-bound exponent under \(\mathrm{SAT}\in \mathrm{SIZE}(n^{c})\).
\item A subpolynomial-query or nonblackbox carrier remains available and has not been audited.
\item The latest failure sharpened the threshold-collapse lemma into a formal no-go statement.
\item The route can still isolate advice leakage, nonuniformity, naturalness, or algebrization as the exact obstruction.
\end{itemize}

\paragraph{Required skills.}
\begin{itemize}
\item contradiction-route-auditor
\item entailment-gap-auditor
\item kb-deduce
\item kb-end-to-end-proof-critic
\item kb-hypothesis-refutation-loop
\item long-horizon-route-portfolio-manager
\item main-problem-alignment-auditor
\item next-experiment-queue-ranker
\item paper-contribution-readiness-judge
\item pvsnp-barrier-specificity-reviewer
\item pvsnp-conceptual-pressure-localizer
\item pvsnp-impossibility-theorem-forger
\item pvsnp-inference-quality-controller
\item pvsnp-proof-derivation-builder
\item pvsnp-qualitative-complexity-reasoner
\item pvsnp-quantifier-uniformity-reviewer
\item pvsnp-reduction-consistency-reviewer
\item pvsnp-symbolic-complexity-auditor
\end{itemize}

\subsection*{\#82 cycle-0130: Hardness-Magnification Bridge Realization}

\paragraph{Analysis.} Run the bit-field audit for one concrete DTSR verifier line: list parent row id, active descriptor id, restriction mask, child row id, Tseitin extension map, local identity tag, and downstream use pointer; verify that \(K_\mathrm{DTSR}\) reaches every used child and compute the exact theta+kappa ledger for the example.

\paragraph{Why this is a good target.} It directly tests the current black-box substitution cap by computing \(\beta\), q, and the induced target-exponent window. It keeps the barrier and model-boundary audit attached to the mathematical transfer rather than as after-the-fact commentary. It is executable as a bounded complexity calculation, CAS-backed exponent audit, or formal transfer check.

\paragraph{Risk.} Low drift risk: the item is framed as an explicit P vs NP transfer audit.

\paragraph{Branch control.} The item has high payoff and a concrete decision test; stopping now would be premature.

\paragraph{Premature-stop guard.} Do not stop before deriving the transfer inequality from \(\mathrm{SAT}\in \mathrm{SIZE}(n^{c})\) to the carrier upper bound, identifying whether \(\beta=0\), \(\beta>0\), or the argument is genuinely nonblackbox, and recording the barrier status of the proposed escape.

\paragraph{Payoff criteria.}
\begin{itemize}
\item For the proof-complexity branch, separate fixed-family proof hardness from language hardness by naming the proof system, shortcut algorithm, and semantic lower-bound target.
\item Promotion requires a proof-system lower bound that survives the Gaussian-elimination or recognition shortcut; otherwise the result is a restricted-proof-system separation only.
\item For this route, state the exact barrier premise it uses or escapes: relativizing, natural, algebrizing, nonuniform, or advice-leaking; promotion requires the escape to be attached to the specific transfer theorem.
\item Compute \(\beta\) and q exactly; promotion requires \(\beta=0\), subpolynomial target queries, or a genuinely nonblackbox use of the SAT circuit.
\item Record relativization, natural-proofs, algebrization, nonuniformity, and advice-leakage status for the specific transfer theorem.
\item Control gate: name the source carrier or language, target class, and circuit or resource model before claiming progress.
\end{itemize}

\paragraph{Stop criteria.}
\begin{itemize}
\item Every named carrier in the branch has \(\beta>0\) and yields only a bounded target-exponent exclusion.
\item The branch cannot state a source lower-bound model, target class, and transfer inequality in the same notation.
\item The proposed escape is only a black-box substitution restatement with no new carrier or nonblackbox mechanism.
\item The barrier audit shows the needed source lower bound is just as hard as the original P vs NP/P/poly target.
\end{itemize}

\paragraph{Do not stop if.}
\begin{itemize}
\item The branch has not yet computed \(\beta\), q, and the induced source upper-bound exponent under \(\mathrm{SAT}\in \mathrm{SIZE}(n^{c})\).
\item A subpolynomial-query or nonblackbox carrier remains available and has not been audited.
\item The latest failure sharpened the threshold-collapse lemma into a formal no-go statement.
\item The route can still isolate advice leakage, nonuniformity, naturalness, or algebrization as the exact obstruction.
\end{itemize}

\paragraph{Required skills.}
\begin{itemize}
\item contradiction-route-auditor
\item entailment-gap-auditor
\item kb-deduce
\item kb-end-to-end-proof-critic
\item kb-hypothesis-refutation-loop
\item long-horizon-route-portfolio-manager
\item main-problem-alignment-auditor
\item next-experiment-queue-ranker
\item paper-contribution-readiness-judge
\item pvsnp-barrier-specificity-reviewer
\item pvsnp-conceptual-pressure-localizer
\item pvsnp-impossibility-theorem-forger
\item pvsnp-inference-quality-controller
\item pvsnp-proof-derivation-builder
\item pvsnp-qualitative-complexity-reasoner
\item pvsnp-quantifier-uniformity-reviewer
\item pvsnp-reduction-consistency-reviewer
\item pvsnp-symbolic-complexity-auditor
\end{itemize}

\subsection*{\#83 cycle-0128: Hardness-Magnification Bridge Realization}

\paragraph{Analysis.} Write the paper-ready TSR rule schema table and verify on one toy projected-dDNNF plus TSR line: compute \(K_\mathrm{TSR}\), \(A_\mathrm{TSR}\), \(B_\mathrm{TSR}\), \(R_\mathrm{TSR}\), and then flip the first unlisted child in the hidden-pruning variant.

\paragraph{Why this is a good target.} It directly tests the current black-box substitution cap by computing \(\beta\), q, and the induced target-exponent window. It keeps the barrier and model-boundary audit attached to the mathematical transfer rather than as after-the-fact commentary. It is executable as a bounded complexity calculation, CAS-backed exponent audit, or formal transfer check.

\paragraph{Risk.} Low drift risk: the item is framed as an explicit P vs NP transfer audit.

\paragraph{Branch control.} The item has high payoff and a concrete decision test; stopping now would be premature.

\paragraph{Premature-stop guard.} Do not stop before deriving the transfer inequality from \(\mathrm{SAT}\in \mathrm{SIZE}(n^{c})\) to the carrier upper bound, identifying whether \(\beta=0\), \(\beta>0\), or the argument is genuinely nonblackbox, and recording the barrier status of the proposed escape.

\paragraph{Payoff criteria.}
\begin{itemize}
\item For this route, state the exact barrier premise it uses or escapes: relativizing, natural, algebrizing, nonuniform, or advice-leaking; promotion requires the escape to be attached to the specific transfer theorem.
\item Compute \(\beta\) and q exactly; promotion requires \(\beta=0\), subpolynomial target queries, or a genuinely nonblackbox use of the SAT circuit.
\item Record relativization, natural-proofs, algebrization, nonuniformity, and advice-leakage status for the specific transfer theorem.
\item Control gate: name the source carrier or language, target class, and circuit or resource model before claiming progress.
\item Control gate: derive the exact transfer inequality from \(\mathrm{SAT}\in \mathrm{SIZE}(n^{c})\), NP subset P/poly, or the stated target assumption to the source-carrier upper bound.
\item Control gate: compute \(\beta\) and q exactly, distinguishing all-exponent consequences from bounded target-exponent windows.
\end{itemize}

\paragraph{Stop criteria.}
\begin{itemize}
\item Every named carrier in the branch has \(\beta>0\) and yields only a bounded target-exponent exclusion.
\item The branch cannot state a source lower-bound model, target class, and transfer inequality in the same notation.
\item The proposed escape is only a black-box substitution restatement with no new carrier or nonblackbox mechanism.
\item The barrier audit shows the needed source lower bound is just as hard as the original P vs NP/P/poly target.
\end{itemize}

\paragraph{Do not stop if.}
\begin{itemize}
\item The branch has not yet computed \(\beta\), q, and the induced source upper-bound exponent under \(\mathrm{SAT}\in \mathrm{SIZE}(n^{c})\).
\item A subpolynomial-query or nonblackbox carrier remains available and has not been audited.
\item The latest failure sharpened the threshold-collapse lemma into a formal no-go statement.
\item The route can still isolate advice leakage, nonuniformity, naturalness, or algebrization as the exact obstruction.
\end{itemize}

\paragraph{Required skills.}
\begin{itemize}
\item contradiction-route-auditor
\item entailment-gap-auditor
\item kb-deduce
\item kb-end-to-end-proof-critic
\item kb-hypothesis-refutation-loop
\item long-horizon-route-portfolio-manager
\item main-problem-alignment-auditor
\item next-experiment-queue-ranker
\item paper-contribution-readiness-judge
\item pvsnp-barrier-specificity-reviewer
\item pvsnp-conceptual-pressure-localizer
\item pvsnp-impossibility-theorem-forger
\item pvsnp-inference-quality-controller
\item pvsnp-proof-derivation-builder
\item pvsnp-qualitative-complexity-reasoner
\item pvsnp-quantifier-uniformity-reviewer
\item pvsnp-reduction-consistency-reviewer
\item pvsnp-symbolic-complexity-auditor
\end{itemize}

\subsection*{\#84 cycle-0124: Hardness-Magnification Bridge Realization}

\paragraph{Analysis.} Write the \(FE_{\mathrm{TSR},decl}\) rule schema as bit fields: row id, Tseitin code, pivot list, child constructor, parent pointers, and EF derivation handles; then check the fixed-point \(K_\mathrm{TSR}\) enumerator on a one-line bounded-SR example and a hidden satisfying-branch counterexample.

\paragraph{Why this is a good target.} It directly tests the current black-box substitution cap by computing \(\beta\), q, and the induced target-exponent window. It keeps the barrier and model-boundary audit attached to the mathematical transfer rather than as after-the-fact commentary. It is executable as a bounded complexity calculation, CAS-backed exponent audit, or formal transfer check.

\paragraph{Risk.} Low drift risk: the item is framed as an explicit P vs NP transfer audit.

\paragraph{Branch control.} The item has high payoff and a concrete decision test; stopping now would be premature.

\paragraph{Premature-stop guard.} Do not stop before deriving the transfer inequality from \(\mathrm{SAT}\in \mathrm{SIZE}(n^{c})\) to the carrier upper bound, identifying whether \(\beta=0\), \(\beta>0\), or the argument is genuinely nonblackbox, and recording the barrier status of the proposed escape.

\paragraph{Payoff criteria.}
\begin{itemize}
\item For the proof-complexity branch, separate fixed-family proof hardness from language hardness by naming the proof system, shortcut algorithm, and semantic lower-bound target.
\item Promotion requires a proof-system lower bound that survives the Gaussian-elimination or recognition shortcut; otherwise the result is a restricted-proof-system separation only.
\item For this route, state the exact barrier premise it uses or escapes: relativizing, natural, algebrizing, nonuniform, or advice-leaking; promotion requires the escape to be attached to the specific transfer theorem.
\item Compute \(\beta\) and q exactly; promotion requires \(\beta=0\), subpolynomial target queries, or a genuinely nonblackbox use of the SAT circuit.
\item Record relativization, natural-proofs, algebrization, nonuniformity, and advice-leakage status for the specific transfer theorem.
\item Control gate: name the source carrier or language, target class, and circuit or resource model before claiming progress.
\end{itemize}

\paragraph{Stop criteria.}
\begin{itemize}
\item Every named carrier in the branch has \(\beta>0\) and yields only a bounded target-exponent exclusion.
\item The branch cannot state a source lower-bound model, target class, and transfer inequality in the same notation.
\item The proposed escape is only a black-box substitution restatement with no new carrier or nonblackbox mechanism.
\item The barrier audit shows the needed source lower bound is just as hard as the original P vs NP/P/poly target.
\end{itemize}

\paragraph{Do not stop if.}
\begin{itemize}
\item The branch has not yet computed \(\beta\), q, and the induced source upper-bound exponent under \(\mathrm{SAT}\in \mathrm{SIZE}(n^{c})\).
\item A subpolynomial-query or nonblackbox carrier remains available and has not been audited.
\item The latest failure sharpened the threshold-collapse lemma into a formal no-go statement.
\item The route can still isolate advice leakage, nonuniformity, naturalness, or algebrization as the exact obstruction.
\end{itemize}

\paragraph{Required skills.}
\begin{itemize}
\item contradiction-route-auditor
\item entailment-gap-auditor
\item kb-deduce
\item kb-end-to-end-proof-critic
\item kb-hypothesis-refutation-loop
\item long-horizon-route-portfolio-manager
\item main-problem-alignment-auditor
\item next-experiment-queue-ranker
\item paper-contribution-readiness-judge
\item pvsnp-barrier-specificity-reviewer
\item pvsnp-conceptual-pressure-localizer
\item pvsnp-impossibility-theorem-forger
\item pvsnp-inference-quality-controller
\item pvsnp-proof-derivation-builder
\item pvsnp-qualitative-complexity-reasoner
\item pvsnp-quantifier-uniformity-reviewer
\item pvsnp-reduction-consistency-reviewer
\item pvsnp-symbolic-complexity-auditor
\end{itemize}

\subsection*{\#85 cycle-0122: Hardness-Magnification Bridge Realization}

\paragraph{Analysis.} Write the paper-ready L122 module: formal theorem statement, fixed-point pseudocode, line-induction proof, projected-dDNNF charge example, unread-enable-bit counterexample, q-trichotomy corollary, and a separate \(RealSATInfluenceRealizer_{N}\) open-boundary paragraph.

\paragraph{Why this is a good target.} It directly tests the current black-box substitution cap by computing \(\beta\), q, and the induced target-exponent window. It keeps the barrier and model-boundary audit attached to the mathematical transfer rather than as after-the-fact commentary. It is executable as a bounded complexity calculation, CAS-backed exponent audit, or formal transfer check.

\paragraph{Risk.} Low drift risk: the item is framed as an explicit P vs NP transfer audit.

\paragraph{Branch control.} The item has high payoff and a concrete decision test; stopping now would be premature.

\paragraph{Premature-stop guard.} Do not stop before deriving the transfer inequality from \(\mathrm{SAT}\in \mathrm{SIZE}(n^{c})\) to the carrier upper bound, identifying whether \(\beta=0\), \(\beta>0\), or the argument is genuinely nonblackbox, and recording the barrier status of the proposed escape.

\paragraph{Payoff criteria.}
\begin{itemize}
\item For this route, state the exact barrier premise it uses or escapes: relativizing, natural, algebrizing, nonuniform, or advice-leaking; promotion requires the escape to be attached to the specific transfer theorem.
\item Compute \(\beta\) and q exactly; promotion requires \(\beta=0\), subpolynomial target queries, or a genuinely nonblackbox use of the SAT circuit.
\item Record relativization, natural-proofs, algebrization, nonuniformity, and advice-leakage status for the specific transfer theorem.
\item Control gate: name the source carrier or language, target class, and circuit or resource model before claiming progress.
\item Control gate: derive the exact transfer inequality from \(\mathrm{SAT}\in \mathrm{SIZE}(n^{c})\), NP subset P/poly, or the stated target assumption to the source-carrier upper bound.
\item Control gate: compute \(\beta\) and q exactly, distinguishing all-exponent consequences from bounded target-exponent windows.
\end{itemize}

\paragraph{Stop criteria.}
\begin{itemize}
\item Every named carrier in the branch has \(\beta>0\) and yields only a bounded target-exponent exclusion.
\item The branch cannot state a source lower-bound model, target class, and transfer inequality in the same notation.
\item The proposed escape is only a black-box substitution restatement with no new carrier or nonblackbox mechanism.
\item The barrier audit shows the needed source lower bound is just as hard as the original P vs NP/P/poly target.
\end{itemize}

\paragraph{Do not stop if.}
\begin{itemize}
\item The branch has not yet computed \(\beta\), q, and the induced source upper-bound exponent under \(\mathrm{SAT}\in \mathrm{SIZE}(n^{c})\).
\item A subpolynomial-query or nonblackbox carrier remains available and has not been audited.
\item The latest failure sharpened the threshold-collapse lemma into a formal no-go statement.
\item The route can still isolate advice leakage, nonuniformity, naturalness, or algebrization as the exact obstruction.
\end{itemize}

\paragraph{Required skills.}
\begin{itemize}
\item contradiction-route-auditor
\item entailment-gap-auditor
\item kb-deduce
\item kb-end-to-end-proof-critic
\item kb-hypothesis-refutation-loop
\item long-horizon-route-portfolio-manager
\item main-problem-alignment-auditor
\item next-experiment-queue-ranker
\item paper-contribution-readiness-judge
\item pvsnp-barrier-specificity-reviewer
\item pvsnp-conceptual-pressure-localizer
\item pvsnp-impossibility-theorem-forger
\item pvsnp-inference-quality-controller
\item pvsnp-proof-derivation-builder
\item pvsnp-qualitative-complexity-reasoner
\item pvsnp-quantifier-uniformity-reviewer
\item pvsnp-reduction-consistency-reviewer
\item pvsnp-symbolic-complexity-auditor
\end{itemize}

\subsection*{\#86 cycle-0121: Hardness-Magnification Bridge Realization: L121 least-closed support audit}

\paragraph{Analysis.} Write the formal L121.1 proof skeleton as theorem text: fixed-point pseudocode for \(K_{dag}\), induction invariant per rule type, charge table for \(bEF/SR_{h}/\mathrm{FormulaEval}/projected-dDNNF\), and the unread-enable-bit counterexample with \(\theta<\epsilon\).

\paragraph{Why this is a good target.} It directly tests the current black-box substitution cap by computing \(\beta\), q, and the induced target-exponent window. It keeps the barrier and model-boundary audit attached to the mathematical transfer rather than as after-the-fact commentary. It is executable as a bounded complexity calculation, CAS-backed exponent audit, or formal transfer check.

\paragraph{Risk.} Low drift risk: the item is framed as an explicit P vs NP transfer audit.

\paragraph{Branch control.} The item has high payoff and a concrete decision test; stopping now would be premature.

\paragraph{Premature-stop guard.} Do not stop before deriving the transfer inequality from \(\mathrm{SAT}\in \mathrm{SIZE}(n^{c})\) to the carrier upper bound, identifying whether \(\beta=0\), \(\beta>0\), or the argument is genuinely nonblackbox, and recording the barrier status of the proposed escape.

\paragraph{Payoff criteria.}
\begin{itemize}
\item For this route, state the exact barrier premise it uses or escapes: relativizing, natural, algebrizing, nonuniform, or advice-leaking; promotion requires the escape to be attached to the specific transfer theorem.
\item Compute \(\beta\) and q exactly; promotion requires \(\beta=0\), subpolynomial target queries, or a genuinely nonblackbox use of the SAT circuit.
\item Record relativization, natural-proofs, algebrization, nonuniformity, and advice-leakage status for the specific transfer theorem.
\item Control gate: name the source carrier or language, target class, and circuit or resource model before claiming progress.
\item Control gate: derive the exact transfer inequality from \(\mathrm{SAT}\in \mathrm{SIZE}(n^{c})\), NP subset P/poly, or the stated target assumption to the source-carrier upper bound.
\item Control gate: compute \(\beta\) and q exactly, distinguishing all-exponent consequences from bounded target-exponent windows.
\end{itemize}

\paragraph{Stop criteria.}
\begin{itemize}
\item Every named carrier in the branch has \(\beta>0\) and yields only a bounded target-exponent exclusion.
\item The branch cannot state a source lower-bound model, target class, and transfer inequality in the same notation.
\item The proposed escape is only a black-box substitution restatement with no new carrier or nonblackbox mechanism.
\item The barrier audit shows the needed source lower bound is just as hard as the original P vs NP/P/poly target.
\end{itemize}

\paragraph{Do not stop if.}
\begin{itemize}
\item The branch has not yet computed \(\beta\), q, and the induced source upper-bound exponent under \(\mathrm{SAT}\in \mathrm{SIZE}(n^{c})\).
\item A subpolynomial-query or nonblackbox carrier remains available and has not been audited.
\item The latest failure sharpened the threshold-collapse lemma into a formal no-go statement.
\item The route can still isolate advice leakage, nonuniformity, naturalness, or algebrization as the exact obstruction.
\end{itemize}

\paragraph{Required skills.}
\begin{itemize}
\item contradiction-route-auditor
\item entailment-gap-auditor
\item kb-deduce
\item kb-end-to-end-proof-critic
\item kb-hypothesis-refutation-loop
\item long-horizon-route-portfolio-manager
\item main-problem-alignment-auditor
\item next-experiment-queue-ranker
\item paper-contribution-readiness-judge
\item pvsnp-barrier-specificity-reviewer
\item pvsnp-conceptual-pressure-localizer
\item pvsnp-impossibility-theorem-forger
\item pvsnp-inference-quality-controller
\item pvsnp-proof-derivation-builder
\item pvsnp-qualitative-complexity-reasoner
\item pvsnp-quantifier-uniformity-reviewer
\item pvsnp-reduction-consistency-reviewer
\item pvsnp-symbolic-complexity-auditor
\end{itemize}

\subsection*{\#87 cycle-0115: Hardness-Magnification Bridge Realization}

\paragraph{Analysis.} Formalize the private-row independence lemma for one full symbolic FormulaEval rule with an unread enable bit: write the two transcripts, the charged constraint set, the flipped row r, and the proof that all visible lines agree while \(H_{N}^\mathrm{SAT}\) parity differs; if independence fails, record the exact same-theory constraint as a \(RealSATInfluenceRealizer_{N}\) candidate.

\paragraph{Why this is a good target.} It directly tests the current black-box substitution cap by computing \(\beta\), q, and the induced target-exponent window. It keeps the barrier and model-boundary audit attached to the mathematical transfer rather than as after-the-fact commentary. It is executable as a bounded complexity calculation, CAS-backed exponent audit, or formal transfer check.

\paragraph{Risk.} Low drift risk: the item is framed as an explicit P vs NP transfer audit.

\paragraph{Branch control.} The item has high payoff and a concrete decision test; stopping now would be premature.

\paragraph{Premature-stop guard.} Do not stop before deriving the transfer inequality from \(\mathrm{SAT}\in \mathrm{SIZE}(n^{c})\) to the carrier upper bound, identifying whether \(\beta=0\), \(\beta>0\), or the argument is genuinely nonblackbox, and recording the barrier status of the proposed escape.

\paragraph{Payoff criteria.}
\begin{itemize}
\item For this route, state the exact barrier premise it uses or escapes: relativizing, natural, algebrizing, nonuniform, or advice-leaking; promotion requires the escape to be attached to the specific transfer theorem.
\item Compute \(\beta\) and q exactly; promotion requires \(\beta=0\), subpolynomial target queries, or a genuinely nonblackbox use of the SAT circuit.
\item Record relativization, natural-proofs, algebrization, nonuniformity, and advice-leakage status for the specific transfer theorem.
\item Control gate: name the source carrier or language, target class, and circuit or resource model before claiming progress.
\item Control gate: derive the exact transfer inequality from \(\mathrm{SAT}\in \mathrm{SIZE}(n^{c})\), NP subset P/poly, or the stated target assumption to the source-carrier upper bound.
\item Control gate: compute \(\beta\) and q exactly, distinguishing all-exponent consequences from bounded target-exponent windows.
\end{itemize}

\paragraph{Stop criteria.}
\begin{itemize}
\item Every named carrier in the branch has \(\beta>0\) and yields only a bounded target-exponent exclusion.
\item The branch cannot state a source lower-bound model, target class, and transfer inequality in the same notation.
\item The proposed escape is only a black-box substitution restatement with no new carrier or nonblackbox mechanism.
\item The barrier audit shows the needed source lower bound is just as hard as the original P vs NP/P/poly target.
\end{itemize}

\paragraph{Do not stop if.}
\begin{itemize}
\item The branch has not yet computed \(\beta\), q, and the induced source upper-bound exponent under \(\mathrm{SAT}\in \mathrm{SIZE}(n^{c})\).
\item A subpolynomial-query or nonblackbox carrier remains available and has not been audited.
\item The latest failure sharpened the threshold-collapse lemma into a formal no-go statement.
\item The route can still isolate advice leakage, nonuniformity, naturalness, or algebrization as the exact obstruction.
\end{itemize}

\paragraph{Required skills.}
\begin{itemize}
\item contradiction-route-auditor
\item entailment-gap-auditor
\item kb-deduce
\item kb-end-to-end-proof-critic
\item kb-hypothesis-refutation-loop
\item long-horizon-route-portfolio-manager
\item main-problem-alignment-auditor
\item next-experiment-queue-ranker
\item paper-contribution-readiness-judge
\item pvsnp-barrier-specificity-reviewer
\item pvsnp-conceptual-pressure-localizer
\item pvsnp-impossibility-theorem-forger
\item pvsnp-inference-quality-controller
\item pvsnp-proof-derivation-builder
\item pvsnp-qualitative-complexity-reasoner
\item pvsnp-quantifier-uniformity-reviewer
\item pvsnp-reduction-consistency-reviewer
\item pvsnp-symbolic-complexity-auditor
\end{itemize}

\subsection*{\#88 cycle-0109: Hardness-Magnification Bridge Realization: \(TransparentRowSupport\) theorem pair}

\paragraph{Analysis.} Referee-check one projected-dDNNF FormulaEval rule against the five \(TransparentRowSupport_{N}\) clauses, compute its \(U_{N}\) and \(X_{N}\) charges, and verify that any projected hidden variable is either influence-closed or triggers the unread-enable-bit countermodel.

\paragraph{Why this is a good target.} It directly tests the current black-box substitution cap by computing \(\beta\), q, and the induced target-exponent window. It keeps the barrier and model-boundary audit attached to the mathematical transfer rather than as after-the-fact commentary. It is executable as a bounded complexity calculation, CAS-backed exponent audit, or formal transfer check.

\paragraph{Risk.} Low drift risk: the item is framed as an explicit P vs NP transfer audit.

\paragraph{Branch control.} The item has high payoff and a concrete decision test; stopping now would be premature.

\paragraph{Premature-stop guard.} Do not stop before deriving the transfer inequality from \(\mathrm{SAT}\in \mathrm{SIZE}(n^{c})\) to the carrier upper bound, identifying whether \(\beta=0\), \(\beta>0\), or the argument is genuinely nonblackbox, and recording the barrier status of the proposed escape.

\paragraph{Payoff criteria.}
\begin{itemize}
\item For this route, state the exact barrier premise it uses or escapes: relativizing, natural, algebrizing, nonuniform, or advice-leaking; promotion requires the escape to be attached to the specific transfer theorem.
\item Compute \(\beta\) and q exactly; promotion requires \(\beta=0\), subpolynomial target queries, or a genuinely nonblackbox use of the SAT circuit.
\item Record relativization, natural-proofs, algebrization, nonuniformity, and advice-leakage status for the specific transfer theorem.
\item Control gate: name the source carrier or language, target class, and circuit or resource model before claiming progress.
\item Control gate: derive the exact transfer inequality from \(\mathrm{SAT}\in \mathrm{SIZE}(n^{c})\), NP subset P/poly, or the stated target assumption to the source-carrier upper bound.
\item Control gate: compute \(\beta\) and q exactly, distinguishing all-exponent consequences from bounded target-exponent windows.
\end{itemize}

\paragraph{Stop criteria.}
\begin{itemize}
\item Every named carrier in the branch has \(\beta>0\) and yields only a bounded target-exponent exclusion.
\item The branch cannot state a source lower-bound model, target class, and transfer inequality in the same notation.
\item The proposed escape is only a black-box substitution restatement with no new carrier or nonblackbox mechanism.
\item The barrier audit shows the needed source lower bound is just as hard as the original P vs NP/P/poly target.
\end{itemize}

\paragraph{Do not stop if.}
\begin{itemize}
\item The branch has not yet computed \(\beta\), q, and the induced source upper-bound exponent under \(\mathrm{SAT}\in \mathrm{SIZE}(n^{c})\).
\item A subpolynomial-query or nonblackbox carrier remains available and has not been audited.
\item The latest failure sharpened the threshold-collapse lemma into a formal no-go statement.
\item The route can still isolate advice leakage, nonuniformity, naturalness, or algebrization as the exact obstruction.
\end{itemize}

\paragraph{Required skills.}
\begin{itemize}
\item contradiction-route-auditor
\item entailment-gap-auditor
\item kb-deduce
\item kb-end-to-end-proof-critic
\item kb-hypothesis-refutation-loop
\item long-horizon-route-portfolio-manager
\item main-problem-alignment-auditor
\item next-experiment-queue-ranker
\item paper-contribution-readiness-judge
\item pvsnp-barrier-specificity-reviewer
\item pvsnp-conceptual-pressure-localizer
\item pvsnp-impossibility-theorem-forger
\item pvsnp-inference-quality-controller
\item pvsnp-proof-derivation-builder
\item pvsnp-qualitative-complexity-reasoner
\item pvsnp-quantifier-uniformity-reviewer
\item pvsnp-reduction-consistency-reviewer
\item pvsnp-symbolic-complexity-auditor
\end{itemize}

\subsection*{\#89 cycle-0106: Hardness-Magnification Bridge Realization}

\paragraph{Analysis.} Run the finite validation audit on one concrete \(\mathrm{Psys}_{sym}\) rule: encode FormulaEval as (a) interval-union axiom schema, (b) \(\mathrm{AC0}\) selector with enumerator, and (c) opaque semantic oracle; compute \(\alpha\) and either attach \(N^{1+\alpha+o(1)}\) charge or exhibit the private-row countermodel.

\paragraph{Why this is a good target.} It directly tests the current black-box substitution cap by computing \(\beta\), q, and the induced target-exponent window. It keeps the barrier and model-boundary audit attached to the mathematical transfer rather than as after-the-fact commentary. It is executable as a bounded complexity calculation, CAS-backed exponent audit, or formal transfer check.

\paragraph{Risk.} Low drift risk: the item is framed as an explicit P vs NP transfer audit.

\paragraph{Branch control.} The item has high payoff and a concrete decision test; stopping now would be premature.

\paragraph{Premature-stop guard.} Do not stop before deriving the transfer inequality from \(\mathrm{SAT}\in \mathrm{SIZE}(n^{c})\) to the carrier upper bound, identifying whether \(\beta=0\), \(\beta>0\), or the argument is genuinely nonblackbox, and recording the barrier status of the proposed escape.

\paragraph{Payoff criteria.}
\begin{itemize}
\item For the restricted-model branch, compute the compiler envelope from unrestricted SIZE(m\textasciicircum{}k) or P-time algorithms into the named model; promotion requires a bound above that envelope on the same SAT slice.
\item Falsify direct transfer if a P function already has the same restricted-model lower bound, such as PARITY for \(\mathrm{AC0}\), HWB for OBDD, or XOR/Tseitin for proof-system degree.
\item For this route, state the exact barrier premise it uses or escapes: relativizing, natural, algebrizing, nonuniform, or advice-leaking; promotion requires the escape to be attached to the specific transfer theorem.
\item Compute \(\beta\) and q exactly; promotion requires \(\beta=0\), subpolynomial target queries, or a genuinely nonblackbox use of the SAT circuit.
\item Record relativization, natural-proofs, algebrization, nonuniformity, and advice-leakage status for the specific transfer theorem.
\item Control gate: name the source carrier or language, target class, and circuit or resource model before claiming progress.
\end{itemize}

\paragraph{Stop criteria.}
\begin{itemize}
\item Every named carrier in the branch has \(\beta>0\) and yields only a bounded target-exponent exclusion.
\item The branch cannot state a source lower-bound model, target class, and transfer inequality in the same notation.
\item The proposed escape is only a black-box substitution restatement with no new carrier or nonblackbox mechanism.
\item The barrier audit shows the needed source lower bound is just as hard as the original P vs NP/P/poly target.
\end{itemize}

\paragraph{Do not stop if.}
\begin{itemize}
\item The branch has not yet computed \(\beta\), q, and the induced source upper-bound exponent under \(\mathrm{SAT}\in \mathrm{SIZE}(n^{c})\).
\item A subpolynomial-query or nonblackbox carrier remains available and has not been audited.
\item The latest failure sharpened the threshold-collapse lemma into a formal no-go statement.
\item The route can still isolate advice leakage, nonuniformity, naturalness, or algebrization as the exact obstruction.
\end{itemize}

\paragraph{Required skills.}
\begin{itemize}
\item contradiction-route-auditor
\item entailment-gap-auditor
\item kb-deduce
\item kb-end-to-end-proof-critic
\item kb-hypothesis-refutation-loop
\item long-horizon-route-portfolio-manager
\item main-problem-alignment-auditor
\item next-experiment-queue-ranker
\item paper-contribution-readiness-judge
\item pvsnp-barrier-specificity-reviewer
\item pvsnp-conceptual-pressure-localizer
\item pvsnp-impossibility-theorem-forger
\item pvsnp-inference-quality-controller
\item pvsnp-proof-derivation-builder
\item pvsnp-qualitative-complexity-reasoner
\item pvsnp-quantifier-uniformity-reviewer
\item pvsnp-reduction-consistency-reviewer
\item pvsnp-symbolic-complexity-auditor
\end{itemize}

\subsection*{\#90 cycle-0105: Hardness-Magnification Bridge Realization}

\paragraph{Analysis.} Run the delta-varied interval-union selector audit: instantiate \(|RowSet|=N^{q-\delta}\), verify \(RowSetEnumerator_{N}\) gate cost, compute \(\Gamma_{bEF}^{IU}\) at \(\delta=q-\epsilon\) \(\pm\) o(1), and test whether any proposed semantic selector supplies a valid \(FE_{Extract,N}\) certificate rather than triggering the private-row falsifier.

\paragraph{Why this is a good target.} It directly tests the current black-box substitution cap by computing \(\beta\), q, and the induced target-exponent window. It keeps the barrier and model-boundary audit attached to the mathematical transfer rather than as after-the-fact commentary. It is executable as a bounded complexity calculation, CAS-backed exponent audit, or formal transfer check.

\paragraph{Risk.} Low drift risk: the item is framed as an explicit P vs NP transfer audit.

\paragraph{Branch control.} The item has high payoff and a concrete decision test; stopping now would be premature.

\paragraph{Premature-stop guard.} Do not stop before deriving the transfer inequality from \(\mathrm{SAT}\in \mathrm{SIZE}(n^{c})\) to the carrier upper bound, identifying whether \(\beta=0\), \(\beta>0\), or the argument is genuinely nonblackbox, and recording the barrier status of the proposed escape.

\paragraph{Payoff criteria.}
\begin{itemize}
\item For this route, state the exact barrier premise it uses or escapes: relativizing, natural, algebrizing, nonuniform, or advice-leaking; promotion requires the escape to be attached to the specific transfer theorem.
\item Compute \(\beta\) and q exactly; promotion requires \(\beta=0\), subpolynomial target queries, or a genuinely nonblackbox use of the SAT circuit.
\item Record relativization, natural-proofs, algebrization, nonuniformity, and advice-leakage status for the specific transfer theorem.
\item Control gate: name the source carrier or language, target class, and circuit or resource model before claiming progress.
\item Control gate: derive the exact transfer inequality from \(\mathrm{SAT}\in \mathrm{SIZE}(n^{c})\), NP subset P/poly, or the stated target assumption to the source-carrier upper bound.
\item Control gate: compute \(\beta\) and q exactly, distinguishing all-exponent consequences from bounded target-exponent windows.
\end{itemize}

\paragraph{Stop criteria.}
\begin{itemize}
\item Every named carrier in the branch has \(\beta>0\) and yields only a bounded target-exponent exclusion.
\item The branch cannot state a source lower-bound model, target class, and transfer inequality in the same notation.
\item The proposed escape is only a black-box substitution restatement with no new carrier or nonblackbox mechanism.
\item The barrier audit shows the needed source lower bound is just as hard as the original P vs NP/P/poly target.
\end{itemize}

\paragraph{Do not stop if.}
\begin{itemize}
\item The branch has not yet computed \(\beta\), q, and the induced source upper-bound exponent under \(\mathrm{SAT}\in \mathrm{SIZE}(n^{c})\).
\item A subpolynomial-query or nonblackbox carrier remains available and has not been audited.
\item The latest failure sharpened the threshold-collapse lemma into a formal no-go statement.
\item The route can still isolate advice leakage, nonuniformity, naturalness, or algebrization as the exact obstruction.
\end{itemize}

\paragraph{Required skills.}
\begin{itemize}
\item contradiction-route-auditor
\item entailment-gap-auditor
\item kb-deduce
\item kb-end-to-end-proof-critic
\item kb-hypothesis-refutation-loop
\item long-horizon-route-portfolio-manager
\item main-problem-alignment-auditor
\item next-experiment-queue-ranker
\item paper-contribution-readiness-judge
\item pvsnp-barrier-specificity-reviewer
\item pvsnp-conceptual-pressure-localizer
\item pvsnp-impossibility-theorem-forger
\item pvsnp-inference-quality-controller
\item pvsnp-proof-derivation-builder
\item pvsnp-qualitative-complexity-reasoner
\item pvsnp-quantifier-uniformity-reviewer
\item pvsnp-reduction-consistency-reviewer
\item pvsnp-symbolic-complexity-auditor
\end{itemize}

\subsection*{\#91 cycle-0103: Hardness-Magnification Bridge Realization}

\paragraph{Analysis.} Extend \(RowSetEnumerator_{N}\) from interval-union selectors to a named \(\mathrm{AC0}\) selector fragment by full scanning over [M\_N]; compute whether scan overhead \(N^{q+o(1)}\) is dominated by N|RowSet| and identify the smallest \(\mathrm{AC0}-private-row\) countermodel if enumeration depends on semantic advice.

\paragraph{Why this is a good target.} It directly tests the current black-box substitution cap by computing \(\beta\), q, and the induced target-exponent window. It keeps the barrier and model-boundary audit attached to the mathematical transfer rather than as after-the-fact commentary. It is executable as a bounded complexity calculation, CAS-backed exponent audit, or formal transfer check.

\paragraph{Risk.} Low drift risk: the item is framed as an explicit P vs NP transfer audit.

\paragraph{Branch control.} The item has high payoff and a concrete decision test; stopping now would be premature.

\paragraph{Premature-stop guard.} Do not stop before deriving the transfer inequality from \(\mathrm{SAT}\in \mathrm{SIZE}(n^{c})\) to the carrier upper bound, identifying whether \(\beta=0\), \(\beta>0\), or the argument is genuinely nonblackbox, and recording the barrier status of the proposed escape.

\paragraph{Payoff criteria.}
\begin{itemize}
\item For the restricted-model branch, compute the compiler envelope from unrestricted SIZE(m\textasciicircum{}k) or P-time algorithms into the named model; promotion requires a bound above that envelope on the same SAT slice.
\item Falsify direct transfer if a P function already has the same restricted-model lower bound, such as PARITY for \(\mathrm{AC0}\), HWB for OBDD, or XOR/Tseitin for proof-system degree.
\item For this route, state the exact barrier premise it uses or escapes: relativizing, natural, algebrizing, nonuniform, or advice-leaking; promotion requires the escape to be attached to the specific transfer theorem.
\item Compute \(\beta\) and q exactly; promotion requires \(\beta=0\), subpolynomial target queries, or a genuinely nonblackbox use of the SAT circuit.
\item Record relativization, natural-proofs, algebrization, nonuniformity, and advice-leakage status for the specific transfer theorem.
\item Control gate: name the source carrier or language, target class, and circuit or resource model before claiming progress.
\end{itemize}

\paragraph{Stop criteria.}
\begin{itemize}
\item Every named carrier in the branch has \(\beta>0\) and yields only a bounded target-exponent exclusion.
\item The branch cannot state a source lower-bound model, target class, and transfer inequality in the same notation.
\item The proposed escape is only a black-box substitution restatement with no new carrier or nonblackbox mechanism.
\item The barrier audit shows the needed source lower bound is just as hard as the original P vs NP/P/poly target.
\end{itemize}

\paragraph{Do not stop if.}
\begin{itemize}
\item The branch has not yet computed \(\beta\), q, and the induced source upper-bound exponent under \(\mathrm{SAT}\in \mathrm{SIZE}(n^{c})\).
\item A subpolynomial-query or nonblackbox carrier remains available and has not been audited.
\item The latest failure sharpened the threshold-collapse lemma into a formal no-go statement.
\item The route can still isolate advice leakage, nonuniformity, naturalness, or algebrization as the exact obstruction.
\end{itemize}

\paragraph{Required skills.}
\begin{itemize}
\item contradiction-route-auditor
\item entailment-gap-auditor
\item kb-deduce
\item kb-end-to-end-proof-critic
\item kb-hypothesis-refutation-loop
\item long-horizon-route-portfolio-manager
\item main-problem-alignment-auditor
\item next-experiment-queue-ranker
\item paper-contribution-readiness-judge
\item pvsnp-barrier-specificity-reviewer
\item pvsnp-conceptual-pressure-localizer
\item pvsnp-impossibility-theorem-forger
\item pvsnp-inference-quality-controller
\item pvsnp-proof-derivation-builder
\item pvsnp-qualitative-complexity-reasoner
\item pvsnp-quantifier-uniformity-reviewer
\item pvsnp-reduction-consistency-reviewer
\item pvsnp-symbolic-complexity-auditor
\end{itemize}

\subsection*{\#92 cycle-0091: Hardness-Magnification Bridge Realization: beta-zero XOR formula-semantic parser audit}

\paragraph{Analysis.} Instantiate \(\mathrm{Psys}_{h,Ext}^{schema}\): EF+EqSAT+depth-h SR with extension definitions generated by row-oblivious schemas. Try to prove ExtSupp(NF(pi))<=poly(|pi|), or construct a minimal nonrelativizing SAT-semantic extension that defeats this normal form.

\paragraph{Why this is a good target.} It directly tests the current black-box substitution cap by computing \(\beta\), q, and the induced target-exponent window. It keeps the barrier and model-boundary audit attached to the mathematical transfer rather than as after-the-fact commentary. It is executable as a bounded complexity calculation, CAS-backed exponent audit, or formal transfer check.

\paragraph{Risk.} Low drift risk: the item is framed as an explicit P vs NP transfer audit.

\paragraph{Branch control.} The item has high payoff and a concrete decision test; stopping now would be premature.

\paragraph{Premature-stop guard.} Do not stop before deriving the transfer inequality from \(\mathrm{SAT}\in \mathrm{SIZE}(n^{c})\) to the carrier upper bound, identifying whether \(\beta=0\), \(\beta>0\), or the argument is genuinely nonblackbox, and recording the barrier status of the proposed escape.

\paragraph{Payoff criteria.}
\begin{itemize}
\item For the proof-complexity branch, separate fixed-family proof hardness from language hardness by naming the proof system, shortcut algorithm, and semantic lower-bound target.
\item Promotion requires a proof-system lower bound that survives the Gaussian-elimination or recognition shortcut; otherwise the result is a restricted-proof-system separation only.
\item For this route, state the exact barrier premise it uses or escapes: relativizing, natural, algebrizing, nonuniform, or advice-leaking; promotion requires the escape to be attached to the specific transfer theorem.
\item Compute \(\beta\) and q exactly; promotion requires \(\beta=0\), subpolynomial target queries, or a genuinely nonblackbox use of the SAT circuit.
\item Record relativization, natural-proofs, algebrization, nonuniformity, and advice-leakage status for the specific transfer theorem.
\item Control gate: name the source carrier or language, target class, and circuit or resource model before claiming progress.
\end{itemize}

\paragraph{Stop criteria.}
\begin{itemize}
\item Every named carrier in the branch has \(\beta>0\) and yields only a bounded target-exponent exclusion.
\item The branch cannot state a source lower-bound model, target class, and transfer inequality in the same notation.
\item The proposed escape is only a black-box substitution restatement with no new carrier or nonblackbox mechanism.
\item The barrier audit shows the needed source lower bound is just as hard as the original P vs NP/P/poly target.
\end{itemize}

\paragraph{Do not stop if.}
\begin{itemize}
\item The branch has not yet computed \(\beta\), q, and the induced source upper-bound exponent under \(\mathrm{SAT}\in \mathrm{SIZE}(n^{c})\).
\item A subpolynomial-query or nonblackbox carrier remains available and has not been audited.
\item The latest failure sharpened the threshold-collapse lemma into a formal no-go statement.
\item The route can still isolate advice leakage, nonuniformity, naturalness, or algebrization as the exact obstruction.
\end{itemize}

\paragraph{Required skills.}
\begin{itemize}
\item contradiction-route-auditor
\item entailment-gap-auditor
\item kb-deduce
\item kb-end-to-end-proof-critic
\item kb-hypothesis-refutation-loop
\item long-horizon-route-portfolio-manager
\item main-problem-alignment-auditor
\item next-experiment-queue-ranker
\item paper-contribution-readiness-judge
\item pvsnp-barrier-specificity-reviewer
\item pvsnp-conceptual-pressure-localizer
\item pvsnp-impossibility-theorem-forger
\item pvsnp-inference-quality-controller
\item pvsnp-proof-derivation-builder
\item pvsnp-qualitative-complexity-reasoner
\item pvsnp-quantifier-uniformity-reviewer
\item pvsnp-reduction-consistency-reviewer
\item pvsnp-symbolic-complexity-auditor
\end{itemize}

\subsection*{\#93 cycle-0089: Hardness-Magnification Bridge Realization}

\paragraph{Analysis.} Fix h with 2\textasciicircum{}h=N\textasciicircum{}\{eta\}; formalize \(Delta_{SR,h}\) for the generated formula closure and try to prove the parser for \(\mathrm{EF}+SATeq+SR_{h}\) restricted to extension variables naming only quotient labels and SR clauses.

\paragraph{Why this is a good target.} It directly tests the current black-box substitution cap by computing \(\beta\), q, and the induced target-exponent window. It keeps the barrier and model-boundary audit attached to the mathematical transfer rather than as after-the-fact commentary. It is executable as a bounded complexity calculation, CAS-backed exponent audit, or formal transfer check.

\paragraph{Risk.} Low drift risk: the item is framed as an explicit P vs NP transfer audit.

\paragraph{Branch control.} The item has high payoff and a concrete decision test; stopping now would be premature.

\paragraph{Premature-stop guard.} Do not stop before deriving the transfer inequality from \(\mathrm{SAT}\in \mathrm{SIZE}(n^{c})\) to the carrier upper bound, identifying whether \(\beta=0\), \(\beta>0\), or the argument is genuinely nonblackbox, and recording the barrier status of the proposed escape.

\paragraph{Payoff criteria.}
\begin{itemize}
\item For the MCSP/MKtP row branch, decide whether the certificate proves direct row hardness or only aggregate/rank-index hardness; record q\_extract=epsilon/(1+D) and q\_force=(D+epsilon)/(1+D) in the same parameter regime.
\item Promotion requires an advice-free row-hard certificate; backtrack if the ranked decoder, canonical unranking, or aggregate proof selects a hard row nonuniformly.
\item For this route, state the exact barrier premise it uses or escapes: relativizing, natural, algebrizing, nonuniform, or advice-leaking; promotion requires the escape to be attached to the specific transfer theorem.
\item Compute \(\beta\) and q exactly; promotion requires \(\beta=0\), subpolynomial target queries, or a genuinely nonblackbox use of the SAT circuit.
\item Record relativization, natural-proofs, algebrization, nonuniformity, and advice-leakage status for the specific transfer theorem.
\item Control gate: name the source carrier or language, target class, and circuit or resource model before claiming progress.
\end{itemize}

\paragraph{Stop criteria.}
\begin{itemize}
\item Every named carrier in the branch has \(\beta>0\) and yields only a bounded target-exponent exclusion.
\item The branch cannot state a source lower-bound model, target class, and transfer inequality in the same notation.
\item The proposed escape is only a black-box substitution restatement with no new carrier or nonblackbox mechanism.
\item The barrier audit shows the needed source lower bound is just as hard as the original P vs NP/P/poly target.
\end{itemize}

\paragraph{Do not stop if.}
\begin{itemize}
\item The branch has not yet computed \(\beta\), q, and the induced source upper-bound exponent under \(\mathrm{SAT}\in \mathrm{SIZE}(n^{c})\).
\item A subpolynomial-query or nonblackbox carrier remains available and has not been audited.
\item The latest failure sharpened the threshold-collapse lemma into a formal no-go statement.
\item The route can still isolate advice leakage, nonuniformity, naturalness, or algebrization as the exact obstruction.
\end{itemize}

\paragraph{Required skills.}
\begin{itemize}
\item contradiction-route-auditor
\item entailment-gap-auditor
\item kb-deduce
\item kb-end-to-end-proof-critic
\item kb-hypothesis-refutation-loop
\item long-horizon-route-portfolio-manager
\item main-problem-alignment-auditor
\item next-experiment-queue-ranker
\item paper-contribution-readiness-judge
\item pvsnp-barrier-specificity-reviewer
\item pvsnp-conceptual-pressure-localizer
\item pvsnp-impossibility-theorem-forger
\item pvsnp-inference-quality-controller
\item pvsnp-proof-derivation-builder
\item pvsnp-qualitative-complexity-reasoner
\item pvsnp-quantifier-uniformity-reviewer
\item pvsnp-reduction-consistency-reviewer
\item pvsnp-symbolic-complexity-auditor
\end{itemize}

\subsection*{\#94 cycle-0088: Hardness-Magnification Bridge Realization}

\paragraph{Analysis.} Instantiate \(K_{pi}\) for local one-step SR only: bound \(|K_{pi}|\) by proof size, define its row-parity constraint-distance \(Delta_{K}\), and test whether \(Delta_{K}\) can be reduced to a known SAT promise-hardness statement or whether it is target-equivalent to the desired \(H_{N}\) lower bound.

\paragraph{Why this is a good target.} It directly tests the current black-box substitution cap by computing \(\beta\), q, and the induced target-exponent window. It keeps the barrier and model-boundary audit attached to the mathematical transfer rather than as after-the-fact commentary. It is executable as a bounded complexity calculation, CAS-backed exponent audit, or formal transfer check.

\paragraph{Risk.} Low drift risk: the item is framed as an explicit P vs NP transfer audit.

\paragraph{Branch control.} The item has high payoff and a concrete decision test; stopping now would be premature.

\paragraph{Premature-stop guard.} Do not stop before deriving the transfer inequality from \(\mathrm{SAT}\in \mathrm{SIZE}(n^{c})\) to the carrier upper bound, identifying whether \(\beta=0\), \(\beta>0\), or the argument is genuinely nonblackbox, and recording the barrier status of the proposed escape.

\paragraph{Payoff criteria.}
\begin{itemize}
\item For this route, state the exact barrier premise it uses or escapes: relativizing, natural, algebrizing, nonuniform, or advice-leaking; promotion requires the escape to be attached to the specific transfer theorem.
\item Compute \(\beta\) and q exactly; promotion requires \(\beta=0\), subpolynomial target queries, or a genuinely nonblackbox use of the SAT circuit.
\item Record relativization, natural-proofs, algebrization, nonuniformity, and advice-leakage status for the specific transfer theorem.
\item Control gate: name the source carrier or language, target class, and circuit or resource model before claiming progress.
\item Control gate: derive the exact transfer inequality from \(\mathrm{SAT}\in \mathrm{SIZE}(n^{c})\), NP subset P/poly, or the stated target assumption to the source-carrier upper bound.
\item Control gate: compute \(\beta\) and q exactly, distinguishing all-exponent consequences from bounded target-exponent windows.
\end{itemize}

\paragraph{Stop criteria.}
\begin{itemize}
\item Every named carrier in the branch has \(\beta>0\) and yields only a bounded target-exponent exclusion.
\item The branch cannot state a source lower-bound model, target class, and transfer inequality in the same notation.
\item The proposed escape is only a black-box substitution restatement with no new carrier or nonblackbox mechanism.
\item The barrier audit shows the needed source lower bound is just as hard as the original P vs NP/P/poly target.
\end{itemize}

\paragraph{Do not stop if.}
\begin{itemize}
\item The branch has not yet computed \(\beta\), q, and the induced source upper-bound exponent under \(\mathrm{SAT}\in \mathrm{SIZE}(n^{c})\).
\item A subpolynomial-query or nonblackbox carrier remains available and has not been audited.
\item The latest failure sharpened the threshold-collapse lemma into a formal no-go statement.
\item The route can still isolate advice leakage, nonuniformity, naturalness, or algebrization as the exact obstruction.
\end{itemize}

\paragraph{Required skills.}
\begin{itemize}
\item contradiction-route-auditor
\item entailment-gap-auditor
\item kb-deduce
\item kb-end-to-end-proof-critic
\item kb-hypothesis-refutation-loop
\item long-horizon-route-portfolio-manager
\item main-problem-alignment-auditor
\item next-experiment-queue-ranker
\item paper-contribution-readiness-judge
\item pvsnp-barrier-specificity-reviewer
\item pvsnp-conceptual-pressure-localizer
\item pvsnp-impossibility-theorem-forger
\item pvsnp-inference-quality-controller
\item pvsnp-proof-derivation-builder
\item pvsnp-qualitative-complexity-reasoner
\item pvsnp-quantifier-uniformity-reviewer
\item pvsnp-reduction-consistency-reviewer
\item pvsnp-symbolic-complexity-auditor
\end{itemize}

\subsection*{\#95 cycle-0081: Hardness-Magnification Bridge Realization}

\paragraph{Analysis.} Formalize the semantic-coupling-width thought experiment as a lemma: for paired formulas F and F OR z, show exactly which \(\mathrm{SAT}_{U}\) implication lines break private-label invariance and whether their violations can be charged to complete forced rows below the \(N^{\epsilon}\) patch threshold.

\paragraph{Why this is a good target.} It directly tests the current black-box substitution cap by computing \(\beta\), q, and the induced target-exponent window. It keeps the barrier and model-boundary audit attached to the mathematical transfer rather than as after-the-fact commentary. It is executable as a bounded complexity calculation, CAS-backed exponent audit, or formal transfer check.

\paragraph{Risk.} Low drift risk: the item is framed as an explicit P vs NP transfer audit.

\paragraph{Branch control.} The item has high payoff and a concrete decision test; stopping now would be premature.

\paragraph{Premature-stop guard.} Do not stop before deriving the transfer inequality from \(\mathrm{SAT}\in \mathrm{SIZE}(n^{c})\) to the carrier upper bound, identifying whether \(\beta=0\), \(\beta>0\), or the argument is genuinely nonblackbox, and recording the barrier status of the proposed escape.

\paragraph{Payoff criteria.}
\begin{itemize}
\item For this route, state the exact barrier premise it uses or escapes: relativizing, natural, algebrizing, nonuniform, or advice-leaking; promotion requires the escape to be attached to the specific transfer theorem.
\item Compute \(\beta\) and q exactly; promotion requires \(\beta=0\), subpolynomial target queries, or a genuinely nonblackbox use of the SAT circuit.
\item Record relativization, natural-proofs, algebrization, nonuniformity, and advice-leakage status for the specific transfer theorem.
\item Control gate: name the source carrier or language, target class, and circuit or resource model before claiming progress.
\item Control gate: derive the exact transfer inequality from \(\mathrm{SAT}\in \mathrm{SIZE}(n^{c})\), NP subset P/poly, or the stated target assumption to the source-carrier upper bound.
\item Control gate: compute \(\beta\) and q exactly, distinguishing all-exponent consequences from bounded target-exponent windows.
\end{itemize}

\paragraph{Stop criteria.}
\begin{itemize}
\item Every named carrier in the branch has \(\beta>0\) and yields only a bounded target-exponent exclusion.
\item The branch cannot state a source lower-bound model, target class, and transfer inequality in the same notation.
\item The proposed escape is only a black-box substitution restatement with no new carrier or nonblackbox mechanism.
\item The barrier audit shows the needed source lower bound is just as hard as the original P vs NP/P/poly target.
\end{itemize}

\paragraph{Do not stop if.}
\begin{itemize}
\item The branch has not yet computed \(\beta\), q, and the induced source upper-bound exponent under \(\mathrm{SAT}\in \mathrm{SIZE}(n^{c})\).
\item A subpolynomial-query or nonblackbox carrier remains available and has not been audited.
\item The latest failure sharpened the threshold-collapse lemma into a formal no-go statement.
\item The route can still isolate advice leakage, nonuniformity, naturalness, or algebrization as the exact obstruction.
\end{itemize}

\paragraph{Required skills.}
\begin{itemize}
\item contradiction-route-auditor
\item entailment-gap-auditor
\item kb-deduce
\item kb-end-to-end-proof-critic
\item kb-hypothesis-refutation-loop
\item long-horizon-route-portfolio-manager
\item main-problem-alignment-auditor
\item next-experiment-queue-ranker
\item paper-contribution-readiness-judge
\item pvsnp-barrier-specificity-reviewer
\item pvsnp-conceptual-pressure-localizer
\item pvsnp-impossibility-theorem-forger
\item pvsnp-inference-quality-controller
\item pvsnp-proof-derivation-builder
\item pvsnp-qualitative-complexity-reasoner
\item pvsnp-quantifier-uniformity-reviewer
\item pvsnp-reduction-consistency-reviewer
\item pvsnp-symbolic-complexity-auditor
\end{itemize}

\subsection*{\#96 cycle-0078: Hardness-Magnification Bridge Realization}

\paragraph{Analysis.} Instantiate \(\mathrm{EF}+\mathrm{SAT}_{U}\) row occurrence access formally: define the touched-row set of a proof verifier, test whether standard EF reflection over \(\mathrm{SAT}_{U}\) extracts to row-value functionality, and if not construct the private-row oracle countermodel as a finite theorem.

\paragraph{Why this is a good target.} It directly tests the current black-box substitution cap by computing \(\beta\), q, and the induced target-exponent window. It keeps the barrier and model-boundary audit attached to the mathematical transfer rather than as after-the-fact commentary. It is executable as a bounded complexity calculation, CAS-backed exponent audit, or formal transfer check.

\paragraph{Risk.} Low drift risk: the item is framed as an explicit P vs NP transfer audit.

\paragraph{Branch control.} The item has high payoff and a concrete decision test; stopping now would be premature.

\paragraph{Premature-stop guard.} Do not stop before deriving the transfer inequality from \(\mathrm{SAT}\in \mathrm{SIZE}(n^{c})\) to the carrier upper bound, identifying whether \(\beta=0\), \(\beta>0\), or the argument is genuinely nonblackbox, and recording the barrier status of the proposed escape.

\paragraph{Payoff criteria.}
\begin{itemize}
\item For the proof-complexity branch, separate fixed-family proof hardness from language hardness by naming the proof system, shortcut algorithm, and semantic lower-bound target.
\item Promotion requires a proof-system lower bound that survives the Gaussian-elimination or recognition shortcut; otherwise the result is a restricted-proof-system separation only.
\item For this route, state the exact barrier premise it uses or escapes: relativizing, natural, algebrizing, nonuniform, or advice-leaking; promotion requires the escape to be attached to the specific transfer theorem.
\item Compute \(\beta\) and q exactly; promotion requires \(\beta=0\), subpolynomial target queries, or a genuinely nonblackbox use of the SAT circuit.
\item Record relativization, natural-proofs, algebrization, nonuniformity, and advice-leakage status for the specific transfer theorem.
\item Control gate: name the source carrier or language, target class, and circuit or resource model before claiming progress.
\end{itemize}

\paragraph{Stop criteria.}
\begin{itemize}
\item Every named carrier in the branch has \(\beta>0\) and yields only a bounded target-exponent exclusion.
\item The branch cannot state a source lower-bound model, target class, and transfer inequality in the same notation.
\item The proposed escape is only a black-box substitution restatement with no new carrier or nonblackbox mechanism.
\item The barrier audit shows the needed source lower bound is just as hard as the original P vs NP/P/poly target.
\end{itemize}

\paragraph{Do not stop if.}
\begin{itemize}
\item The branch has not yet computed \(\beta\), q, and the induced source upper-bound exponent under \(\mathrm{SAT}\in \mathrm{SIZE}(n^{c})\).
\item A subpolynomial-query or nonblackbox carrier remains available and has not been audited.
\item The latest failure sharpened the threshold-collapse lemma into a formal no-go statement.
\item The route can still isolate advice leakage, nonuniformity, naturalness, or algebrization as the exact obstruction.
\end{itemize}

\paragraph{Required skills.}
\begin{itemize}
\item contradiction-route-auditor
\item entailment-gap-auditor
\item kb-deduce
\item kb-end-to-end-proof-critic
\item kb-hypothesis-refutation-loop
\item long-horizon-route-portfolio-manager
\item main-problem-alignment-auditor
\item next-experiment-queue-ranker
\item paper-contribution-readiness-judge
\item pvsnp-barrier-specificity-reviewer
\item pvsnp-conceptual-pressure-localizer
\item pvsnp-impossibility-theorem-forger
\item pvsnp-inference-quality-controller
\item pvsnp-proof-derivation-builder
\item pvsnp-qualitative-complexity-reasoner
\item pvsnp-quantifier-uniformity-reviewer
\item pvsnp-reduction-consistency-reviewer
\item pvsnp-symbolic-complexity-auditor
\end{itemize}

\subsection*{\#97 cycle-0077: Hardness-Magnification Bridge Realization}

\paragraph{Analysis.} Instantiate Touch(Pi) for a concrete \(\mathrm{EF}+\mathrm{SAT}_{U}\) encoding of circuit lower-bound certificates: count whether proofs of size \(N^{k}\) can mention \(sub-N^\epsilon\) rows, and test if formula-equality/SAT-realization axioms break the occurrence-oracle extension.

\paragraph{Why this is a good target.} It directly tests the current black-box substitution cap by computing \(\beta\), q, and the induced target-exponent window. It keeps the barrier and model-boundary audit attached to the mathematical transfer rather than as after-the-fact commentary. It is executable as a bounded complexity calculation, CAS-backed exponent audit, or formal transfer check.

\paragraph{Risk.} Low drift risk: the item is framed as an explicit P vs NP transfer audit.

\paragraph{Branch control.} The item has high payoff and a concrete decision test; stopping now would be premature.

\paragraph{Premature-stop guard.} Do not stop before deriving the transfer inequality from \(\mathrm{SAT}\in \mathrm{SIZE}(n^{c})\) to the carrier upper bound, identifying whether \(\beta=0\), \(\beta>0\), or the argument is genuinely nonblackbox, and recording the barrier status of the proposed escape.

\paragraph{Payoff criteria.}
\begin{itemize}
\item For this route, state the exact barrier premise it uses or escapes: relativizing, natural, algebrizing, nonuniform, or advice-leaking; promotion requires the escape to be attached to the specific transfer theorem.
\item Compute \(\beta\) and q exactly; promotion requires \(\beta=0\), subpolynomial target queries, or a genuinely nonblackbox use of the SAT circuit.
\item Record relativization, natural-proofs, algebrization, nonuniformity, and advice-leakage status for the specific transfer theorem.
\item Control gate: name the source carrier or language, target class, and circuit or resource model before claiming progress.
\item Control gate: derive the exact transfer inequality from \(\mathrm{SAT}\in \mathrm{SIZE}(n^{c})\), NP subset P/poly, or the stated target assumption to the source-carrier upper bound.
\item Control gate: compute \(\beta\) and q exactly, distinguishing all-exponent consequences from bounded target-exponent windows.
\end{itemize}

\paragraph{Stop criteria.}
\begin{itemize}
\item Every named carrier in the branch has \(\beta>0\) and yields only a bounded target-exponent exclusion.
\item The branch cannot state a source lower-bound model, target class, and transfer inequality in the same notation.
\item The proposed escape is only a black-box substitution restatement with no new carrier or nonblackbox mechanism.
\item The barrier audit shows the needed source lower bound is just as hard as the original P vs NP/P/poly target.
\end{itemize}

\paragraph{Do not stop if.}
\begin{itemize}
\item The branch has not yet computed \(\beta\), q, and the induced source upper-bound exponent under \(\mathrm{SAT}\in \mathrm{SIZE}(n^{c})\).
\item A subpolynomial-query or nonblackbox carrier remains available and has not been audited.
\item The latest failure sharpened the threshold-collapse lemma into a formal no-go statement.
\item The route can still isolate advice leakage, nonuniformity, naturalness, or algebrization as the exact obstruction.
\end{itemize}

\paragraph{Required skills.}
\begin{itemize}
\item contradiction-route-auditor
\item entailment-gap-auditor
\item kb-deduce
\item kb-end-to-end-proof-critic
\item kb-hypothesis-refutation-loop
\item long-horizon-route-portfolio-manager
\item main-problem-alignment-auditor
\item next-experiment-queue-ranker
\item paper-contribution-readiness-judge
\item pvsnp-barrier-specificity-reviewer
\item pvsnp-conceptual-pressure-localizer
\item pvsnp-impossibility-theorem-forger
\item pvsnp-inference-quality-controller
\item pvsnp-proof-derivation-builder
\item pvsnp-qualitative-complexity-reasoner
\item pvsnp-quantifier-uniformity-reviewer
\item pvsnp-reduction-consistency-reviewer
\item pvsnp-symbolic-complexity-auditor
\end{itemize}

\subsection*{\#98 cycle-0075: Hardness-Magnification Bridge Realization}

\paragraph{Analysis.} Instantiate \(\mathrm{EF}+\mathrm{SAT}_{U}\) proofs as circuits over \(\mathrm{SAT}_{U}\) label atoms, define functionalization formally, and try to construct a \(\theta<\epsilon\) residual-leaky transcript or prove extraction to the robustness gate.

\paragraph{Why this is a good target.} It directly tests the current black-box substitution cap by computing \(\beta\), q, and the induced target-exponent window. It keeps the barrier and model-boundary audit attached to the mathematical transfer rather than as after-the-fact commentary. It is executable as a bounded complexity calculation, CAS-backed exponent audit, or formal transfer check.

\paragraph{Risk.} Low drift risk: the item is framed as an explicit P vs NP transfer audit.

\paragraph{Branch control.} The item has high payoff and a concrete decision test; stopping now would be premature.

\paragraph{Premature-stop guard.} Do not stop before deriving the transfer inequality from \(\mathrm{SAT}\in \mathrm{SIZE}(n^{c})\) to the carrier upper bound, identifying whether \(\beta=0\), \(\beta>0\), or the argument is genuinely nonblackbox, and recording the barrier status of the proposed escape.

\paragraph{Payoff criteria.}
\begin{itemize}
\item For the proof-complexity branch, separate fixed-family proof hardness from language hardness by naming the proof system, shortcut algorithm, and semantic lower-bound target.
\item Promotion requires a proof-system lower bound that survives the Gaussian-elimination or recognition shortcut; otherwise the result is a restricted-proof-system separation only.
\item For this route, state the exact barrier premise it uses or escapes: relativizing, natural, algebrizing, nonuniform, or advice-leaking; promotion requires the escape to be attached to the specific transfer theorem.
\item Compute \(\beta\) and q exactly; promotion requires \(\beta=0\), subpolynomial target queries, or a genuinely nonblackbox use of the SAT circuit.
\item Record relativization, natural-proofs, algebrization, nonuniformity, and advice-leakage status for the specific transfer theorem.
\item Control gate: name the source carrier or language, target class, and circuit or resource model before claiming progress.
\end{itemize}

\paragraph{Stop criteria.}
\begin{itemize}
\item Every named carrier in the branch has \(\beta>0\) and yields only a bounded target-exponent exclusion.
\item The branch cannot state a source lower-bound model, target class, and transfer inequality in the same notation.
\item The proposed escape is only a black-box substitution restatement with no new carrier or nonblackbox mechanism.
\item The barrier audit shows the needed source lower bound is just as hard as the original P vs NP/P/poly target.
\end{itemize}

\paragraph{Do not stop if.}
\begin{itemize}
\item The branch has not yet computed \(\beta\), q, and the induced source upper-bound exponent under \(\mathrm{SAT}\in \mathrm{SIZE}(n^{c})\).
\item A subpolynomial-query or nonblackbox carrier remains available and has not been audited.
\item The latest failure sharpened the threshold-collapse lemma into a formal no-go statement.
\item The route can still isolate advice leakage, nonuniformity, naturalness, or algebrization as the exact obstruction.
\end{itemize}

\paragraph{Required skills.}
\begin{itemize}
\item contradiction-route-auditor
\item entailment-gap-auditor
\item kb-deduce
\item kb-end-to-end-proof-critic
\item kb-hypothesis-refutation-loop
\item long-horizon-route-portfolio-manager
\item main-problem-alignment-auditor
\item next-experiment-queue-ranker
\item paper-contribution-readiness-judge
\item pvsnp-barrier-specificity-reviewer
\item pvsnp-conceptual-pressure-localizer
\item pvsnp-impossibility-theorem-forger
\item pvsnp-inference-quality-controller
\item pvsnp-proof-derivation-builder
\item pvsnp-qualitative-complexity-reasoner
\item pvsnp-quantifier-uniformity-reviewer
\item pvsnp-reduction-consistency-reviewer
\item pvsnp-symbolic-complexity-auditor
\end{itemize}

\subsection*{\#99 cycle-0074: Hardness-Magnification Bridge Realization: beta-zero XOR row-residual audit}

\paragraph{Analysis.} Formalize \(\mathrm{EF}+\mathrm{SAT}_{U}\) syntax with label table \(U_{N}\), prove the functional dense naturalization lemma cleanly, then attempt a finite extraction theorem: every nonfunctional accepted proof either identifies a row set \(R_{N}\) and approximator \(D_{N}\) or uses a nonrelativizing SAT principle outside the row-residual audit.

\paragraph{Why this is a good target.} It directly tests the current black-box substitution cap by computing \(\beta\), q, and the induced target-exponent window. It keeps the barrier and model-boundary audit attached to the mathematical transfer rather than as after-the-fact commentary. It is executable as a bounded complexity calculation, CAS-backed exponent audit, or formal transfer check.

\paragraph{Risk.} Low drift risk: the item is framed as an explicit P vs NP transfer audit.

\paragraph{Branch control.} The item has high payoff and a concrete decision test; stopping now would be premature.

\paragraph{Premature-stop guard.} Do not stop before deriving the transfer inequality from \(\mathrm{SAT}\in \mathrm{SIZE}(n^{c})\) to the carrier upper bound, identifying whether \(\beta=0\), \(\beta>0\), or the argument is genuinely nonblackbox, and recording the barrier status of the proposed escape.

\paragraph{Payoff criteria.}
\begin{itemize}
\item For the proof-complexity branch, separate fixed-family proof hardness from language hardness by naming the proof system, shortcut algorithm, and semantic lower-bound target.
\item Promotion requires a proof-system lower bound that survives the Gaussian-elimination or recognition shortcut; otherwise the result is a restricted-proof-system separation only.
\item For this route, state the exact barrier premise it uses or escapes: relativizing, natural, algebrizing, nonuniform, or advice-leaking; promotion requires the escape to be attached to the specific transfer theorem.
\item Compute \(\beta\) and q exactly; promotion requires \(\beta=0\), subpolynomial target queries, or a genuinely nonblackbox use of the SAT circuit.
\item Record relativization, natural-proofs, algebrization, nonuniformity, and advice-leakage status for the specific transfer theorem.
\item Control gate: name the source carrier or language, target class, and circuit or resource model before claiming progress.
\end{itemize}

\paragraph{Stop criteria.}
\begin{itemize}
\item Every named carrier in the branch has \(\beta>0\) and yields only a bounded target-exponent exclusion.
\item The branch cannot state a source lower-bound model, target class, and transfer inequality in the same notation.
\item The proposed escape is only a black-box substitution restatement with no new carrier or nonblackbox mechanism.
\item The barrier audit shows the needed source lower bound is just as hard as the original P vs NP/P/poly target.
\end{itemize}

\paragraph{Do not stop if.}
\begin{itemize}
\item The branch has not yet computed \(\beta\), q, and the induced source upper-bound exponent under \(\mathrm{SAT}\in \mathrm{SIZE}(n^{c})\).
\item A subpolynomial-query or nonblackbox carrier remains available and has not been audited.
\item The latest failure sharpened the threshold-collapse lemma into a formal no-go statement.
\item The route can still isolate advice leakage, nonuniformity, naturalness, or algebrization as the exact obstruction.
\end{itemize}

\paragraph{Required skills.}
\begin{itemize}
\item contradiction-route-auditor
\item entailment-gap-auditor
\item kb-deduce
\item kb-end-to-end-proof-critic
\item kb-hypothesis-refutation-loop
\item long-horizon-route-portfolio-manager
\item main-problem-alignment-auditor
\item next-experiment-queue-ranker
\item paper-contribution-readiness-judge
\item pvsnp-barrier-specificity-reviewer
\item pvsnp-conceptual-pressure-localizer
\item pvsnp-impossibility-theorem-forger
\item pvsnp-inference-quality-controller
\item pvsnp-proof-derivation-builder
\item pvsnp-qualitative-complexity-reasoner
\item pvsnp-quantifier-uniformity-reviewer
\item pvsnp-reduction-consistency-reviewer
\item pvsnp-symbolic-complexity-auditor
\end{itemize}

\subsection*{\#100 cycle-0069: Hardness-Magnification Bridge Realization}

\paragraph{Analysis.} Instantiate \(U_{N}\) with L=(log N)\textasciicircum{}2 and q in ((1+epsilon)/2,1+epsilon); compute whether candidate proof systems give certificate density above \(2^{-\mathrm{poly}(L)}\) or force a nonlarge escape, then run the prior-art boundary check.

\paragraph{Why this is a good target.} It directly tests the current black-box substitution cap by computing \(\beta\), q, and the induced target-exponent window. It keeps the barrier and model-boundary audit attached to the mathematical transfer rather than as after-the-fact commentary. It is executable as a bounded complexity calculation, CAS-backed exponent audit, or formal transfer check.

\paragraph{Risk.} Low drift risk: the item is framed as an explicit P vs NP transfer audit.

\paragraph{Branch control.} The item has high payoff and a concrete decision test; stopping now would be premature.

\paragraph{Premature-stop guard.} Do not stop before deriving the transfer inequality from \(\mathrm{SAT}\in \mathrm{SIZE}(n^{c})\) to the carrier upper bound, identifying whether \(\beta=0\), \(\beta>0\), or the argument is genuinely nonblackbox, and recording the barrier status of the proposed escape.

\paragraph{Payoff criteria.}
\begin{itemize}
\item For the proof-complexity branch, separate fixed-family proof hardness from language hardness by naming the proof system, shortcut algorithm, and semantic lower-bound target.
\item Promotion requires a proof-system lower bound that survives the Gaussian-elimination or recognition shortcut; otherwise the result is a restricted-proof-system separation only.
\item For this route, state the exact barrier premise it uses or escapes: relativizing, natural, algebrizing, nonuniform, or advice-leaking; promotion requires the escape to be attached to the specific transfer theorem.
\item Compute \(\beta\) and q exactly; promotion requires \(\beta=0\), subpolynomial target queries, or a genuinely nonblackbox use of the SAT circuit.
\item Record relativization, natural-proofs, algebrization, nonuniformity, and advice-leakage status for the specific transfer theorem.
\item Control gate: name the source carrier or language, target class, and circuit or resource model before claiming progress.
\end{itemize}

\paragraph{Stop criteria.}
\begin{itemize}
\item Every named carrier in the branch has \(\beta>0\) and yields only a bounded target-exponent exclusion.
\item The branch cannot state a source lower-bound model, target class, and transfer inequality in the same notation.
\item The proposed escape is only a black-box substitution restatement with no new carrier or nonblackbox mechanism.
\item The barrier audit shows the needed source lower bound is just as hard as the original P vs NP/P/poly target.
\end{itemize}

\paragraph{Do not stop if.}
\begin{itemize}
\item The branch has not yet computed \(\beta\), q, and the induced source upper-bound exponent under \(\mathrm{SAT}\in \mathrm{SIZE}(n^{c})\).
\item A subpolynomial-query or nonblackbox carrier remains available and has not been audited.
\item The latest failure sharpened the threshold-collapse lemma into a formal no-go statement.
\item The route can still isolate advice leakage, nonuniformity, naturalness, or algebrization as the exact obstruction.
\end{itemize}

\paragraph{Required skills.}
\begin{itemize}
\item contradiction-route-auditor
\item entailment-gap-auditor
\item kb-deduce
\item kb-end-to-end-proof-critic
\item kb-hypothesis-refutation-loop
\item long-horizon-route-portfolio-manager
\item main-problem-alignment-auditor
\item next-experiment-queue-ranker
\item paper-contribution-readiness-judge
\item pvsnp-barrier-specificity-reviewer
\item pvsnp-conceptual-pressure-localizer
\item pvsnp-impossibility-theorem-forger
\item pvsnp-inference-quality-controller
\item pvsnp-proof-derivation-builder
\item pvsnp-qualitative-complexity-reasoner
\item pvsnp-quantifier-uniformity-reviewer
\item pvsnp-reduction-consistency-reviewer
\item pvsnp-symbolic-complexity-auditor
\end{itemize}

\subsection*{\#101 cycle-0066: Hardness-Magnification Bridge Realization}

\paragraph{Analysis.} Instantiate \(R_{n}\) for one high-q polylog XOR carrier or prove a finite impossibility: no label-free density-preserving realization can coexist with \(q,\rho<1+\epsilon\) and generator universe \(U_{N}\) of SAT formulas length (log N)\textasciicircum{}2.

\paragraph{Why this is a good target.} It directly tests the current black-box substitution cap by computing \(\beta\), q, and the induced target-exponent window. It keeps the barrier and model-boundary audit attached to the mathematical transfer rather than as after-the-fact commentary. It is executable as a bounded complexity calculation, CAS-backed exponent audit, or formal transfer check.

\paragraph{Risk.} Low drift risk: the item is framed as an explicit P vs NP transfer audit.

\paragraph{Branch control.} The item has high payoff and a concrete decision test; stopping now would be premature.

\paragraph{Premature-stop guard.} Do not stop before deriving the transfer inequality from \(\mathrm{SAT}\in \mathrm{SIZE}(n^{c})\) to the carrier upper bound, identifying whether \(\beta=0\), \(\beta>0\), or the argument is genuinely nonblackbox, and recording the barrier status of the proposed escape.

\paragraph{Payoff criteria.}
\begin{itemize}
\item For the high-q generated-OR carrier, compute the exact leakage exponent of the remaining witness-cover or unary-residual mechanism and compare it with 1+epsilon-q.
\item Promotion requires \(q+a>1+\epsilon\) after every projection, padding, and witness-cover leakage term is accounted for; otherwise the branch becomes a bounded-window or no-go result.
\item For the proof-complexity branch, separate fixed-family proof hardness from language hardness by naming the proof system, shortcut algorithm, and semantic lower-bound target.
\item Promotion requires a proof-system lower bound that survives the Gaussian-elimination or recognition shortcut; otherwise the result is a restricted-proof-system separation only.
\item For this route, state the exact barrier premise it uses or escapes: relativizing, natural, algebrizing, nonuniform, or advice-leaking; promotion requires the escape to be attached to the specific transfer theorem.
\item Compute \(\beta\) and q exactly; promotion requires \(\beta=0\), subpolynomial target queries, or a genuinely nonblackbox use of the SAT circuit.
\end{itemize}

\paragraph{Stop criteria.}
\begin{itemize}
\item Every named carrier in the branch has \(\beta>0\) and yields only a bounded target-exponent exclusion.
\item The branch cannot state a source lower-bound model, target class, and transfer inequality in the same notation.
\item The proposed escape is only a black-box substitution restatement with no new carrier or nonblackbox mechanism.
\item The barrier audit shows the needed source lower bound is just as hard as the original P vs NP/P/poly target.
\end{itemize}

\paragraph{Do not stop if.}
\begin{itemize}
\item The branch has not yet computed \(\beta\), q, and the induced source upper-bound exponent under \(\mathrm{SAT}\in \mathrm{SIZE}(n^{c})\).
\item A subpolynomial-query or nonblackbox carrier remains available and has not been audited.
\item The latest failure sharpened the threshold-collapse lemma into a formal no-go statement.
\item The route can still isolate advice leakage, nonuniformity, naturalness, or algebrization as the exact obstruction.
\end{itemize}

\paragraph{Required skills.}
\begin{itemize}
\item contradiction-route-auditor
\item entailment-gap-auditor
\item kb-deduce
\item kb-end-to-end-proof-critic
\item kb-hypothesis-refutation-loop
\item long-horizon-route-portfolio-manager
\item main-problem-alignment-auditor
\item next-experiment-queue-ranker
\item paper-contribution-readiness-judge
\item pvsnp-barrier-specificity-reviewer
\item pvsnp-conceptual-pressure-localizer
\item pvsnp-impossibility-theorem-forger
\item pvsnp-inference-quality-controller
\item pvsnp-proof-derivation-builder
\item pvsnp-qualitative-complexity-reasoner
\item pvsnp-quantifier-uniformity-reviewer
\item pvsnp-reduction-consistency-reviewer
\item pvsnp-symbolic-complexity-auditor
\end{itemize}

\subsection*{\#102 cycle-0065: Hardness-Magnification Bridge Realization}

\paragraph{Analysis.} Instantiate \(R_{N}\) with a named proof system, e.g. Extended Frege or IPS, and test whether an LB certificate for \(C(H_{N,y})>N^{1+\epsilon}\) has \(O(L_{N})-bit\) searchable descriptions, inverse-polynomial label density, and truth-table soundness; otherwise record the precise escape parameter.

\paragraph{Why this is a good target.} It directly tests the current black-box substitution cap by computing \(\beta\), q, and the induced target-exponent window. It keeps the barrier and model-boundary audit attached to the mathematical transfer rather than as after-the-fact commentary. It is executable as a bounded complexity calculation, CAS-backed exponent audit, or formal transfer check.

\paragraph{Risk.} Low drift risk: the item is framed as an explicit P vs NP transfer audit.

\paragraph{Branch control.} The item has high payoff and a concrete decision test; stopping now would be premature.

\paragraph{Premature-stop guard.} Do not stop before deriving the transfer inequality from \(\mathrm{SAT}\in \mathrm{SIZE}(n^{c})\) to the carrier upper bound, identifying whether \(\beta=0\), \(\beta>0\), or the argument is genuinely nonblackbox, and recording the barrier status of the proposed escape.

\paragraph{Payoff criteria.}
\begin{itemize}
\item For the proof-complexity branch, separate fixed-family proof hardness from language hardness by naming the proof system, shortcut algorithm, and semantic lower-bound target.
\item Promotion requires a proof-system lower bound that survives the Gaussian-elimination or recognition shortcut; otherwise the result is a restricted-proof-system separation only.
\item For this route, state the exact barrier premise it uses or escapes: relativizing, natural, algebrizing, nonuniform, or advice-leaking; promotion requires the escape to be attached to the specific transfer theorem.
\item Compute \(\beta\) and q exactly; promotion requires \(\beta=0\), subpolynomial target queries, or a genuinely nonblackbox use of the SAT circuit.
\item Record relativization, natural-proofs, algebrization, nonuniformity, and advice-leakage status for the specific transfer theorem.
\item Control gate: name the source carrier or language, target class, and circuit or resource model before claiming progress.
\end{itemize}

\paragraph{Stop criteria.}
\begin{itemize}
\item Every named carrier in the branch has \(\beta>0\) and yields only a bounded target-exponent exclusion.
\item The branch cannot state a source lower-bound model, target class, and transfer inequality in the same notation.
\item The proposed escape is only a black-box substitution restatement with no new carrier or nonblackbox mechanism.
\item The barrier audit shows the needed source lower bound is just as hard as the original P vs NP/P/poly target.
\end{itemize}

\paragraph{Do not stop if.}
\begin{itemize}
\item The branch has not yet computed \(\beta\), q, and the induced source upper-bound exponent under \(\mathrm{SAT}\in \mathrm{SIZE}(n^{c})\).
\item A subpolynomial-query or nonblackbox carrier remains available and has not been audited.
\item The latest failure sharpened the threshold-collapse lemma into a formal no-go statement.
\item The route can still isolate advice leakage, nonuniformity, naturalness, or algebrization as the exact obstruction.
\end{itemize}

\paragraph{Required skills.}
\begin{itemize}
\item contradiction-route-auditor
\item entailment-gap-auditor
\item kb-deduce
\item kb-end-to-end-proof-critic
\item kb-hypothesis-refutation-loop
\item long-horizon-route-portfolio-manager
\item main-problem-alignment-auditor
\item next-experiment-queue-ranker
\item paper-contribution-readiness-judge
\item pvsnp-barrier-specificity-reviewer
\item pvsnp-conceptual-pressure-localizer
\item pvsnp-impossibility-theorem-forger
\item pvsnp-inference-quality-controller
\item pvsnp-proof-derivation-builder
\item pvsnp-qualitative-complexity-reasoner
\item pvsnp-quantifier-uniformity-reviewer
\item pvsnp-reduction-consistency-reviewer
\item pvsnp-symbolic-complexity-auditor
\end{itemize}

\subsection*{\#103 cycle-0064: Hardness-Magnification Bridge Realization}

\paragraph{Analysis.} Instantiate L\_N=(log N)\textasciicircum{}2 and q in ((1+epsilon)/2,1+epsilon); compute achievable \(M_{N}\) and \(r_{N}\) for an explicit generator family, then test for a label-free coordinate transcript or prove any proposed SAT anti-alignment certificate has density below \(2^{-\omega(L_{N})}\).

\paragraph{Why this is a good target.} It directly tests the current black-box substitution cap by computing \(\beta\), q, and the induced target-exponent window. It keeps the barrier and model-boundary audit attached to the mathematical transfer rather than as after-the-fact commentary. It is executable as a bounded complexity calculation, CAS-backed exponent audit, or formal transfer check.

\paragraph{Risk.} Low drift risk: the item is framed as an explicit P vs NP transfer audit.

\paragraph{Branch control.} The item has high payoff and a concrete decision test; stopping now would be premature.

\paragraph{Premature-stop guard.} Do not stop before deriving the transfer inequality from \(\mathrm{SAT}\in \mathrm{SIZE}(n^{c})\) to the carrier upper bound, identifying whether \(\beta=0\), \(\beta>0\), or the argument is genuinely nonblackbox, and recording the barrier status of the proposed escape.

\paragraph{Payoff criteria.}
\begin{itemize}
\item For the proof-complexity branch, separate fixed-family proof hardness from language hardness by naming the proof system, shortcut algorithm, and semantic lower-bound target.
\item Promotion requires a proof-system lower bound that survives the Gaussian-elimination or recognition shortcut; otherwise the result is a restricted-proof-system separation only.
\item For this route, state the exact barrier premise it uses or escapes: relativizing, natural, algebrizing, nonuniform, or advice-leaking; promotion requires the escape to be attached to the specific transfer theorem.
\item Compute \(\beta\) and q exactly; promotion requires \(\beta=0\), subpolynomial target queries, or a genuinely nonblackbox use of the SAT circuit.
\item Record relativization, natural-proofs, algebrization, nonuniformity, and advice-leakage status for the specific transfer theorem.
\item Control gate: name the source carrier or language, target class, and circuit or resource model before claiming progress.
\end{itemize}

\paragraph{Stop criteria.}
\begin{itemize}
\item Every named carrier in the branch has \(\beta>0\) and yields only a bounded target-exponent exclusion.
\item The branch cannot state a source lower-bound model, target class, and transfer inequality in the same notation.
\item The proposed escape is only a black-box substitution restatement with no new carrier or nonblackbox mechanism.
\item The barrier audit shows the needed source lower bound is just as hard as the original P vs NP/P/poly target.
\end{itemize}

\paragraph{Do not stop if.}
\begin{itemize}
\item The branch has not yet computed \(\beta\), q, and the induced source upper-bound exponent under \(\mathrm{SAT}\in \mathrm{SIZE}(n^{c})\).
\item A subpolynomial-query or nonblackbox carrier remains available and has not been audited.
\item The latest failure sharpened the threshold-collapse lemma into a formal no-go statement.
\item The route can still isolate advice leakage, nonuniformity, naturalness, or algebrization as the exact obstruction.
\end{itemize}

\paragraph{Required skills.}
\begin{itemize}
\item contradiction-route-auditor
\item entailment-gap-auditor
\item kb-deduce
\item kb-end-to-end-proof-critic
\item kb-hypothesis-refutation-loop
\item long-horizon-route-portfolio-manager
\item main-problem-alignment-auditor
\item next-experiment-queue-ranker
\item paper-contribution-readiness-judge
\item pvsnp-barrier-specificity-reviewer
\item pvsnp-conceptual-pressure-localizer
\item pvsnp-impossibility-theorem-forger
\item pvsnp-inference-quality-controller
\item pvsnp-proof-derivation-builder
\item pvsnp-qualitative-complexity-reasoner
\item pvsnp-quantifier-uniformity-reviewer
\item pvsnp-reduction-consistency-reviewer
\item pvsnp-symbolic-complexity-auditor
\end{itemize}

\subsection*{\#104 cycle-0048: Relativizing locality-extraction gate for generated-OR carriers}

\paragraph{Analysis.} Try to formulate a real-SAT locality extractor using witness self-reducibility: from a small circuit for \(H_{N}\), isolate a generated satisfiable query and build a one-sided predicate; audit algebrization and \(N^{o(1)}\) loss.

\paragraph{Why this is a good target.} It directly tests the current black-box substitution cap by computing \(\beta\), q, and the induced target-exponent window. It keeps the barrier and model-boundary audit attached to the mathematical transfer rather than as after-the-fact commentary. It is executable as a bounded complexity calculation, CAS-backed exponent audit, or formal transfer check.

\paragraph{Risk.} Low drift risk: the item is framed as an explicit P vs NP transfer audit.

\paragraph{Branch control.} The item has high payoff and a concrete decision test; stopping now would be premature.

\paragraph{Premature-stop guard.} Do not stop before the branch has produced a concrete source-to-target transfer theorem, an explicit obstruction lemma, or a barrier-aware reason why the route cannot reach the P vs NP/P/poly target.

\paragraph{Payoff criteria.}
\begin{itemize}
\item For the high-q generated-OR carrier, compute the exact leakage exponent of the remaining witness-cover or unary-residual mechanism and compare it with 1+epsilon-q.
\item Promotion requires \(q+a>1+\epsilon\) after every projection, padding, and witness-cover leakage term is accounted for; otherwise the branch becomes a bounded-window or no-go result.
\item For this route, state the exact barrier premise it uses or escapes: relativizing, natural, algebrizing, nonuniform, or advice-leaking; promotion requires the escape to be attached to the specific transfer theorem.
\item Compute \(\beta\) and q exactly; promotion requires \(\beta=0\), subpolynomial target queries, or a genuinely nonblackbox use of the SAT circuit.
\item Record relativization, natural-proofs, algebrization, nonuniformity, and advice-leakage status for the specific transfer theorem.
\item Control gate: name the source carrier or language, target class, and circuit or resource model before claiming progress.
\end{itemize}

\paragraph{Stop criteria.}
\begin{itemize}
\item The branch cannot identify a source language/carrier and target complexity class in the same model.
\item The route repeatedly proves only a proxy lower bound with no target-preserving transfer.
\item The remaining work is a slogan about hardness magnification rather than a computable transfer or obstruction.
\end{itemize}

\paragraph{Do not stop if.}
\begin{itemize}
\item A transfer exponent, quantifier repair, or barrier status remains uncomputed.
\item The branch can still name a target-preserving carrier and a falsifiable lower-bound threshold.
\item The failure can become a formal no-go lemma for a route family.
\end{itemize}

\paragraph{Required skills.}
\begin{itemize}
\item contradiction-route-auditor
\item entailment-gap-auditor
\item kb-deduce
\item kb-end-to-end-proof-critic
\item kb-hypothesis-refutation-loop
\item long-horizon-route-portfolio-manager
\item main-problem-alignment-auditor
\item next-experiment-queue-ranker
\item paper-contribution-readiness-judge
\item pvsnp-barrier-specificity-reviewer
\item pvsnp-conceptual-pressure-localizer
\item pvsnp-impossibility-theorem-forger
\item pvsnp-inference-quality-controller
\item pvsnp-proof-derivation-builder
\item pvsnp-qualitative-complexity-reasoner
\item pvsnp-quantifier-uniformity-reviewer
\item pvsnp-reduction-consistency-reviewer
\item pvsnp-symbolic-complexity-auditor
\end{itemize}

\subsection*{\#105 cycle-0047: Hardness-Magnification Bridge Realization: locality-extraction obstruction}

\paragraph{Analysis.} Try a real-SAT extraction using self-reducibility or witness isolation: from a small circuit for \(H_{N}\), define a local predicate that accepts formulas whose forced witness path changes the global OR value, compute the overhead, and test it against the planted-control oracle falsifier.

\paragraph{Why this is a good target.} It directly tests the current black-box substitution cap by computing \(\beta\), q, and the induced target-exponent window. It keeps the barrier and model-boundary audit attached to the mathematical transfer rather than as after-the-fact commentary. It is executable as a bounded complexity calculation, CAS-backed exponent audit, or formal transfer check.

\paragraph{Risk.} Low drift risk: the item is framed as an explicit P vs NP transfer audit.

\paragraph{Branch control.} The item has high payoff and a concrete decision test; stopping now would be premature.

\paragraph{Premature-stop guard.} Do not stop before deriving the transfer inequality from \(\mathrm{SAT}\in \mathrm{SIZE}(n^{c})\) to the carrier upper bound, identifying whether \(\beta=0\), \(\beta>0\), or the argument is genuinely nonblackbox, and recording the barrier status of the proposed escape.

\paragraph{Payoff criteria.}
\begin{itemize}
\item For the high-q generated-OR carrier, compute the exact leakage exponent of the remaining witness-cover or unary-residual mechanism and compare it with 1+epsilon-q.
\item Promotion requires \(q+a>1+\epsilon\) after every projection, padding, and witness-cover leakage term is accounted for; otherwise the branch becomes a bounded-window or no-go result.
\item For this route, state the exact barrier premise it uses or escapes: relativizing, natural, algebrizing, nonuniform, or advice-leaking; promotion requires the escape to be attached to the specific transfer theorem.
\item Compute \(\beta\) and q exactly; promotion requires \(\beta=0\), subpolynomial target queries, or a genuinely nonblackbox use of the SAT circuit.
\item Record relativization, natural-proofs, algebrization, nonuniformity, and advice-leakage status for the specific transfer theorem.
\item Control gate: name the source carrier or language, target class, and circuit or resource model before claiming progress.
\end{itemize}

\paragraph{Stop criteria.}
\begin{itemize}
\item Every named carrier in the branch has \(\beta>0\) and yields only a bounded target-exponent exclusion.
\item The branch cannot state a source lower-bound model, target class, and transfer inequality in the same notation.
\item The proposed escape is only a black-box substitution restatement with no new carrier or nonblackbox mechanism.
\item The barrier audit shows the needed source lower bound is just as hard as the original P vs NP/P/poly target.
\end{itemize}

\paragraph{Do not stop if.}
\begin{itemize}
\item The branch has not yet computed \(\beta\), q, and the induced source upper-bound exponent under \(\mathrm{SAT}\in \mathrm{SIZE}(n^{c})\).
\item A subpolynomial-query or nonblackbox carrier remains available and has not been audited.
\item The latest failure sharpened the threshold-collapse lemma into a formal no-go statement.
\item The route can still isolate advice leakage, nonuniformity, naturalness, or algebrization as the exact obstruction.
\end{itemize}

\paragraph{Required skills.}
\begin{itemize}
\item contradiction-route-auditor
\item entailment-gap-auditor
\item kb-deduce
\item kb-end-to-end-proof-critic
\item kb-hypothesis-refutation-loop
\item long-horizon-route-portfolio-manager
\item main-problem-alignment-auditor
\item next-experiment-queue-ranker
\item paper-contribution-readiness-judge
\item pvsnp-barrier-specificity-reviewer
\item pvsnp-conceptual-pressure-localizer
\item pvsnp-impossibility-theorem-forger
\item pvsnp-inference-quality-controller
\item pvsnp-proof-derivation-builder
\item pvsnp-qualitative-complexity-reasoner
\item pvsnp-quantifier-uniformity-reviewer
\item pvsnp-reduction-consistency-reviewer
\item pvsnp-symbolic-complexity-auditor
\end{itemize}

\subsection*{\#106 cycle-0044: Hardness-Magnification Bridge Realization: residual cap-denominator gate}

\paragraph{Analysis.} Instantiate the cap audit for partial-assignment predicates of width w on L=(log N)\textasciicircum{}2 formulas: compute theta(w,sigma) for random-looking generated ranges, then check whether sigma+theta can stay above 1+epsilon while q,rho remain in the target-preserving window.

\paragraph{Why this is a good target.} It directly tests the current black-box substitution cap by computing \(\beta\), q, and the induced target-exponent window. It keeps the barrier and model-boundary audit attached to the mathematical transfer rather than as after-the-fact commentary. It is executable as a bounded complexity calculation, CAS-backed exponent audit, or formal transfer check.

\paragraph{Risk.} Low drift risk: the item is framed as an explicit P vs NP transfer audit.

\paragraph{Branch control.} The item has high payoff and a concrete decision test; stopping now would be premature.

\paragraph{Premature-stop guard.} Do not stop before deriving the transfer inequality from \(\mathrm{SAT}\in \mathrm{SIZE}(n^{c})\) to the carrier upper bound, identifying whether \(\beta=0\), \(\beta>0\), or the argument is genuinely nonblackbox, and recording the barrier status of the proposed escape.

\paragraph{Payoff criteria.}
\begin{itemize}
\item For the one-sided semantic-separator branch, compute the cap profile A\_N(s) for the named generator and predicate class; promotion requires a bound strong enough to force Sep\_N\textasciicircum{}\{1s\}(E)>N\textasciicircum{}\{1+epsilon-o(1)\}.
\item Exhibit or refute a subthreshold semantic planted-slice predicate; this decides whether the generated-OR carrier needs range-closure repair, stronger pseudorandomness, or quarantine.
\item For the high-q generated-OR carrier, compute the exact leakage exponent of the remaining witness-cover or unary-residual mechanism and compare it with 1+epsilon-q.
\item Promotion requires \(q+a>1+\epsilon\) after every projection, padding, and witness-cover leakage term is accounted for; otherwise the branch becomes a bounded-window or no-go result.
\item For this route, state the exact barrier premise it uses or escapes: relativizing, natural, algebrizing, nonuniform, or advice-leaking; promotion requires the escape to be attached to the specific transfer theorem.
\item Compute \(\beta\) and q exactly; promotion requires \(\beta=0\), subpolynomial target queries, or a genuinely nonblackbox use of the SAT circuit.
\end{itemize}

\paragraph{Stop criteria.}
\begin{itemize}
\item Every named carrier in the branch has \(\beta>0\) and yields only a bounded target-exponent exclusion.
\item The branch cannot state a source lower-bound model, target class, and transfer inequality in the same notation.
\item The proposed escape is only a black-box substitution restatement with no new carrier or nonblackbox mechanism.
\item The barrier audit shows the needed source lower bound is just as hard as the original P vs NP/P/poly target.
\end{itemize}

\paragraph{Do not stop if.}
\begin{itemize}
\item The branch has not yet computed \(\beta\), q, and the induced source upper-bound exponent under \(\mathrm{SAT}\in \mathrm{SIZE}(n^{c})\).
\item A subpolynomial-query or nonblackbox carrier remains available and has not been audited.
\item The latest failure sharpened the threshold-collapse lemma into a formal no-go statement.
\item The route can still isolate advice leakage, nonuniformity, naturalness, or algebrization as the exact obstruction.
\end{itemize}

\paragraph{Required skills.}
\begin{itemize}
\item contradiction-route-auditor
\item entailment-gap-auditor
\item kb-deduce
\item kb-end-to-end-proof-critic
\item kb-hypothesis-refutation-loop
\item long-horizon-route-portfolio-manager
\item main-problem-alignment-auditor
\item next-experiment-queue-ranker
\item paper-contribution-readiness-judge
\item pvsnp-barrier-specificity-reviewer
\item pvsnp-conceptual-pressure-localizer
\item pvsnp-impossibility-theorem-forger
\item pvsnp-inference-quality-controller
\item pvsnp-proof-derivation-builder
\item pvsnp-qualitative-complexity-reasoner
\item pvsnp-quantifier-uniformity-reviewer
\item pvsnp-reduction-consistency-reviewer
\item pvsnp-symbolic-complexity-auditor
\end{itemize}

\subsection*{\#107 cycle-0043: Generated-OR residual cap gate for hardness magnification}

\paragraph{Analysis.} Try to construct residual caps for the planted/random 3CNF generator using partial-assignment certificates, and compute the resulting exponent b.

\paragraph{Why this is a good target.} It directly tests the current black-box substitution cap by computing \(\beta\), q, and the induced target-exponent window. It keeps the barrier and model-boundary audit attached to the mathematical transfer rather than as after-the-fact commentary. It is executable as a bounded complexity calculation, CAS-backed exponent audit, or formal transfer check.

\paragraph{Risk.} Low drift risk: the item is framed as an explicit P vs NP transfer audit.

\paragraph{Branch control.} The item has high payoff and a concrete decision test; stopping now would be premature.

\paragraph{Premature-stop guard.} Do not stop before deriving the transfer inequality from \(\mathrm{SAT}\in \mathrm{SIZE}(n^{c})\) to the carrier upper bound, identifying whether \(\beta=0\), \(\beta>0\), or the argument is genuinely nonblackbox, and recording the barrier status of the proposed escape.

\paragraph{Payoff criteria.}
\begin{itemize}
\item For the one-sided semantic-separator branch, compute the cap profile A\_N(s) for the named generator and predicate class; promotion requires a bound strong enough to force Sep\_N\textasciicircum{}\{1s\}(E)>N\textasciicircum{}\{1+epsilon-o(1)\}.
\item Exhibit or refute a subthreshold semantic planted-slice predicate; this decides whether the generated-OR carrier needs range-closure repair, stronger pseudorandomness, or quarantine.
\item For the high-q generated-OR carrier, compute the exact leakage exponent of the remaining witness-cover or unary-residual mechanism and compare it with 1+epsilon-q.
\item Promotion requires \(q+a>1+\epsilon\) after every projection, padding, and witness-cover leakage term is accounted for; otherwise the branch becomes a bounded-window or no-go result.
\item For this route, state the exact barrier premise it uses or escapes: relativizing, natural, algebrizing, nonuniform, or advice-leaking; promotion requires the escape to be attached to the specific transfer theorem.
\item Compute \(\beta\) and q exactly; promotion requires \(\beta=0\), subpolynomial target queries, or a genuinely nonblackbox use of the SAT circuit.
\end{itemize}

\paragraph{Stop criteria.}
\begin{itemize}
\item Every named carrier in the branch has \(\beta>0\) and yields only a bounded target-exponent exclusion.
\item The branch cannot state a source lower-bound model, target class, and transfer inequality in the same notation.
\item The proposed escape is only a black-box substitution restatement with no new carrier or nonblackbox mechanism.
\item The barrier audit shows the needed source lower bound is just as hard as the original P vs NP/P/poly target.
\end{itemize}

\paragraph{Do not stop if.}
\begin{itemize}
\item The branch has not yet computed \(\beta\), q, and the induced source upper-bound exponent under \(\mathrm{SAT}\in \mathrm{SIZE}(n^{c})\).
\item A subpolynomial-query or nonblackbox carrier remains available and has not been audited.
\item The latest failure sharpened the threshold-collapse lemma into a formal no-go statement.
\item The route can still isolate advice leakage, nonuniformity, naturalness, or algebrization as the exact obstruction.
\end{itemize}

\paragraph{Required skills.}
\begin{itemize}
\item contradiction-route-auditor
\item entailment-gap-auditor
\item kb-deduce
\item kb-end-to-end-proof-critic
\item kb-hypothesis-refutation-loop
\item long-horizon-route-portfolio-manager
\item main-problem-alignment-auditor
\item next-experiment-queue-ranker
\item paper-contribution-readiness-judge
\item pvsnp-barrier-specificity-reviewer
\item pvsnp-conceptual-pressure-localizer
\item pvsnp-impossibility-theorem-forger
\item pvsnp-inference-quality-controller
\item pvsnp-proof-derivation-builder
\item pvsnp-qualitative-complexity-reasoner
\item pvsnp-quantifier-uniformity-reviewer
\item pvsnp-reduction-consistency-reviewer
\item pvsnp-symbolic-complexity-auditor
\end{itemize}

\subsection*{\#108 cycle-0042: Cycle 42: distributional cap certificate for one-sided semantic separators}

\paragraph{Analysis.} Formalize the cap profile for the unmarked random k-CNF generated-OR thought experiment: compute assignment-test hit exponents as a function of \(\beta\), then identify the first semantic predicate class beyond terms that can violate or support \(A_{N}(s)\le s/N^{1+\epsilon-o(1)}\).

\paragraph{Why this is a good target.} It directly tests the current black-box substitution cap by computing \(\beta\), q, and the induced target-exponent window. It keeps the barrier and model-boundary audit attached to the mathematical transfer rather than as after-the-fact commentary. It is executable as a bounded complexity calculation, CAS-backed exponent audit, or formal transfer check.

\paragraph{Risk.} Low drift risk: the item is framed as an explicit P vs NP transfer audit.

\paragraph{Branch control.} The item has high payoff and a concrete decision test; stopping now would be premature.

\paragraph{Premature-stop guard.} Do not stop before deriving the transfer inequality from \(\mathrm{SAT}\in \mathrm{SIZE}(n^{c})\) to the carrier upper bound, identifying whether \(\beta=0\), \(\beta>0\), or the argument is genuinely nonblackbox, and recording the barrier status of the proposed escape.

\paragraph{Payoff criteria.}
\begin{itemize}
\item For the one-sided semantic-separator branch, compute the cap profile A\_N(s) for the named generator and predicate class; promotion requires a bound strong enough to force Sep\_N\textasciicircum{}\{1s\}(E)>N\textasciicircum{}\{1+epsilon-o(1)\}.
\item Exhibit or refute a subthreshold semantic planted-slice predicate; this decides whether the generated-OR carrier needs range-closure repair, stronger pseudorandomness, or quarantine.
\item For the high-q generated-OR carrier, compute the exact leakage exponent of the remaining witness-cover or unary-residual mechanism and compare it with 1+epsilon-q.
\item Promotion requires \(q+a>1+\epsilon\) after every projection, padding, and witness-cover leakage term is accounted for; otherwise the branch becomes a bounded-window or no-go result.
\item For this route, state the exact barrier premise it uses or escapes: relativizing, natural, algebrizing, nonuniform, or advice-leaking; promotion requires the escape to be attached to the specific transfer theorem.
\item Compute \(\beta\) and q exactly; promotion requires \(\beta=0\), subpolynomial target queries, or a genuinely nonblackbox use of the SAT circuit.
\end{itemize}

\paragraph{Stop criteria.}
\begin{itemize}
\item Every named carrier in the branch has \(\beta>0\) and yields only a bounded target-exponent exclusion.
\item The branch cannot state a source lower-bound model, target class, and transfer inequality in the same notation.
\item The proposed escape is only a black-box substitution restatement with no new carrier or nonblackbox mechanism.
\item The barrier audit shows the needed source lower bound is just as hard as the original P vs NP/P/poly target.
\end{itemize}

\paragraph{Do not stop if.}
\begin{itemize}
\item The branch has not yet computed \(\beta\), q, and the induced source upper-bound exponent under \(\mathrm{SAT}\in \mathrm{SIZE}(n^{c})\).
\item A subpolynomial-query or nonblackbox carrier remains available and has not been audited.
\item The latest failure sharpened the threshold-collapse lemma into a formal no-go statement.
\item The route can still isolate advice leakage, nonuniformity, naturalness, or algebrization as the exact obstruction.
\end{itemize}

\paragraph{Required skills.}
\begin{itemize}
\item contradiction-route-auditor
\item entailment-gap-auditor
\item kb-deduce
\item kb-end-to-end-proof-critic
\item kb-hypothesis-refutation-loop
\item long-horizon-route-portfolio-manager
\item main-problem-alignment-auditor
\item next-experiment-queue-ranker
\item paper-contribution-readiness-judge
\item pvsnp-barrier-specificity-reviewer
\item pvsnp-conceptual-pressure-localizer
\item pvsnp-impossibility-theorem-forger
\item pvsnp-inference-quality-controller
\item pvsnp-proof-derivation-builder
\item pvsnp-qualitative-complexity-reasoner
\item pvsnp-quantifier-uniformity-reviewer
\item pvsnp-reduction-consistency-reviewer
\item pvsnp-symbolic-complexity-auditor
\end{itemize}

\subsection*{\#109 cycle-0040: Hardness-Magnification Bridge Realization: sound local certificate cover gate}

\paragraph{Analysis.} Compute the SLC cover exponent for generated formulas with bounded backdoor width \(b=\alpha\) log N and for small semantic classes such as Horn, 2-SAT, and bounded-width CNF; identify whether the threshold \(\alpha\ge 1+\epsilon-q\) is achievable without destroying the SAT-substitution bridge.

\paragraph{Why this is a good target.} It directly tests the current black-box substitution cap by computing \(\beta\), q, and the induced target-exponent window. It keeps the barrier and model-boundary audit attached to the mathematical transfer rather than as after-the-fact commentary. It is executable as a bounded complexity calculation, CAS-backed exponent audit, or formal transfer check.

\paragraph{Risk.} Low drift risk: the item is framed as an explicit P vs NP transfer audit.

\paragraph{Branch control.} The item has high payoff and a concrete decision test; stopping now would be premature.

\paragraph{Premature-stop guard.} Do not stop before deriving the transfer inequality from \(\mathrm{SAT}\in \mathrm{SIZE}(n^{c})\) to the carrier upper bound, identifying whether \(\beta=0\), \(\beta>0\), or the argument is genuinely nonblackbox, and recording the barrier status of the proposed escape.

\paragraph{Payoff criteria.}
\begin{itemize}
\item For the certificate-cover branch, compute K,E,L and the induced circuit upper bound C(H\_N)<=C(G\_N)+O(K(s+L\textasciicircum{}\{O(1)\})+EN); promotion requires every sound local cover to stay above the \(N^{1+\epsilon-o(1)}\) source threshold.
\item Falsify the branch if Horn, 2-SAT, bounded-width CNF, or bounded-backdoor emitted formulas give subthreshold satisfying-term or sound-local-certificate covers.
\item For this route, state the exact barrier premise it uses or escapes: relativizing, natural, algebrizing, nonuniform, or advice-leaking; promotion requires the escape to be attached to the specific transfer theorem.
\item Compute \(\beta\) and q exactly; promotion requires \(\beta=0\), subpolynomial target queries, or a genuinely nonblackbox use of the SAT circuit.
\item Record relativization, natural-proofs, algebrization, nonuniformity, and advice-leakage status for the specific transfer theorem.
\item Control gate: name the source carrier or language, target class, and circuit or resource model before claiming progress.
\end{itemize}

\paragraph{Stop criteria.}
\begin{itemize}
\item Every named carrier in the branch has \(\beta>0\) and yields only a bounded target-exponent exclusion.
\item The branch cannot state a source lower-bound model, target class, and transfer inequality in the same notation.
\item The proposed escape is only a black-box substitution restatement with no new carrier or nonblackbox mechanism.
\item The barrier audit shows the needed source lower bound is just as hard as the original P vs NP/P/poly target.
\end{itemize}

\paragraph{Do not stop if.}
\begin{itemize}
\item The branch has not yet computed \(\beta\), q, and the induced source upper-bound exponent under \(\mathrm{SAT}\in \mathrm{SIZE}(n^{c})\).
\item A subpolynomial-query or nonblackbox carrier remains available and has not been audited.
\item The latest failure sharpened the threshold-collapse lemma into a formal no-go statement.
\item The route can still isolate advice leakage, nonuniformity, naturalness, or algebrization as the exact obstruction.
\end{itemize}

\paragraph{Required skills.}
\begin{itemize}
\item contradiction-route-auditor
\item entailment-gap-auditor
\item kb-deduce
\item kb-end-to-end-proof-critic
\item kb-hypothesis-refutation-loop
\item long-horizon-route-portfolio-manager
\item main-problem-alignment-auditor
\item next-experiment-queue-ranker
\item paper-contribution-readiness-judge
\item pvsnp-barrier-specificity-reviewer
\item pvsnp-conceptual-pressure-localizer
\item pvsnp-impossibility-theorem-forger
\item pvsnp-inference-quality-controller
\item pvsnp-proof-derivation-builder
\item pvsnp-qualitative-complexity-reasoner
\item pvsnp-quantifier-uniformity-reviewer
\item pvsnp-reduction-consistency-reviewer
\item pvsnp-symbolic-complexity-auditor
\end{itemize}

\subsection*{\#110 cycle-0026: Hardness-Magnification Bridge Realization}

\paragraph{Analysis.} Instantiate the thought experiment with a block generator \(E_{N}\) whose coordinates are grouped into planted easy and hard ranges; compute the smallest \(\delta\) forced by coverage of positive inputs and test whether selected-range promise SAT stays subthreshold.

\paragraph{Why this is a good target.} It directly tests the current black-box substitution cap by computing \(\beta\), q, and the induced target-exponent window. It keeps the barrier and model-boundary audit attached to the mathematical transfer rather than as after-the-fact commentary. It is executable as a bounded complexity calculation, CAS-backed exponent audit, or formal transfer check.

\paragraph{Risk.} Low drift risk: the item is framed as an explicit P vs NP transfer audit.

\paragraph{Branch control.} The item has high payoff and a concrete decision test; stopping now would be premature.

\paragraph{Premature-stop guard.} Do not stop before deriving the transfer inequality from \(\mathrm{SAT}\in \mathrm{SIZE}(n^{c})\) to the carrier upper bound, identifying whether \(\beta=0\), \(\beta>0\), or the argument is genuinely nonblackbox, and recording the barrier status of the proposed escape.

\paragraph{Payoff criteria.}
\begin{itemize}
\item For this route, state the exact barrier premise it uses or escapes: relativizing, natural, algebrizing, nonuniform, or advice-leaking; promotion requires the escape to be attached to the specific transfer theorem.
\item Compute \(\beta\) and q exactly; promotion requires \(\beta=0\), subpolynomial target queries, or a genuinely nonblackbox use of the SAT circuit.
\item Record relativization, natural-proofs, algebrization, nonuniformity, and advice-leakage status for the specific transfer theorem.
\item Control gate: name the source carrier or language, target class, and circuit or resource model before claiming progress.
\item Control gate: derive the exact transfer inequality from \(\mathrm{SAT}\in \mathrm{SIZE}(n^{c})\), NP subset P/poly, or the stated target assumption to the source-carrier upper bound.
\item Control gate: compute \(\beta\) and q exactly, distinguishing all-exponent consequences from bounded target-exponent windows.
\end{itemize}

\paragraph{Stop criteria.}
\begin{itemize}
\item Every named carrier in the branch has \(\beta>0\) and yields only a bounded target-exponent exclusion.
\item The branch cannot state a source lower-bound model, target class, and transfer inequality in the same notation.
\item The proposed escape is only a black-box substitution restatement with no new carrier or nonblackbox mechanism.
\item The barrier audit shows the needed source lower bound is just as hard as the original P vs NP/P/poly target.
\end{itemize}

\paragraph{Do not stop if.}
\begin{itemize}
\item The branch has not yet computed \(\beta\), q, and the induced source upper-bound exponent under \(\mathrm{SAT}\in \mathrm{SIZE}(n^{c})\).
\item A subpolynomial-query or nonblackbox carrier remains available and has not been audited.
\item The latest failure sharpened the threshold-collapse lemma into a formal no-go statement.
\item The route can still isolate advice leakage, nonuniformity, naturalness, or algebrization as the exact obstruction.
\end{itemize}

\paragraph{Required skills.}
\begin{itemize}
\item contradiction-route-auditor
\item entailment-gap-auditor
\item kb-deduce
\item kb-end-to-end-proof-critic
\item kb-hypothesis-refutation-loop
\item long-horizon-route-portfolio-manager
\item main-problem-alignment-auditor
\item next-experiment-queue-ranker
\item paper-contribution-readiness-judge
\item pvsnp-barrier-specificity-reviewer
\item pvsnp-conceptual-pressure-localizer
\item pvsnp-impossibility-theorem-forger
\item pvsnp-inference-quality-controller
\item pvsnp-proof-derivation-builder
\item pvsnp-qualitative-complexity-reasoner
\item pvsnp-quantifier-uniformity-reviewer
\item pvsnp-reduction-consistency-reviewer
\item pvsnp-symbolic-complexity-auditor
\end{itemize}

\subsection*{\#111 cycle-0025: Generated-OR magnification: selector-range leakage gate}

\paragraph{Analysis.} Instantiate \(Sel_{N}\) for simple generator templates, including lexicographic first-satisfying-query generators and hash-bucketed query generators, and compute whether selected-image lookup already drops below \(N^{1+\epsilon}\).

\paragraph{Why this is a good target.} It directly tests the current black-box substitution cap by computing \(\beta\), q, and the induced target-exponent window. It keeps the barrier and model-boundary audit attached to the mathematical transfer rather than as after-the-fact commentary. It is executable as a bounded complexity calculation, CAS-backed exponent audit, or formal transfer check.

\paragraph{Risk.} Low drift risk: the item is framed as an explicit P vs NP transfer audit.

\paragraph{Branch control.} The item has high payoff and a concrete decision test; stopping now would be premature.

\paragraph{Premature-stop guard.} Do not stop before the branch has produced a concrete source-to-target transfer theorem, an explicit obstruction lemma, or a barrier-aware reason why the route cannot reach the P vs NP/P/poly target.

\paragraph{Payoff criteria.}
\begin{itemize}
\item For the high-q generated-OR carrier, compute the exact leakage exponent of the remaining witness-cover or unary-residual mechanism and compare it with 1+epsilon-q.
\item Promotion requires \(q+a>1+\epsilon\) after every projection, padding, and witness-cover leakage term is accounted for; otherwise the branch becomes a bounded-window or no-go result.
\item For this route, state the exact barrier premise it uses or escapes: relativizing, natural, algebrizing, nonuniform, or advice-leaking; promotion requires the escape to be attached to the specific transfer theorem.
\item Compute \(\beta\) and q exactly; promotion requires \(\beta=0\), subpolynomial target queries, or a genuinely nonblackbox use of the SAT circuit.
\item Record relativization, natural-proofs, algebrization, nonuniformity, and advice-leakage status for the specific transfer theorem.
\item Control gate: name the source carrier or language, target class, and circuit or resource model before claiming progress.
\end{itemize}

\paragraph{Stop criteria.}
\begin{itemize}
\item The branch cannot identify a source language/carrier and target complexity class in the same model.
\item The route repeatedly proves only a proxy lower bound with no target-preserving transfer.
\item The remaining work is a slogan about hardness magnification rather than a computable transfer or obstruction.
\end{itemize}

\paragraph{Do not stop if.}
\begin{itemize}
\item A transfer exponent, quantifier repair, or barrier status remains uncomputed.
\item The branch can still name a target-preserving carrier and a falsifiable lower-bound threshold.
\item The failure can become a formal no-go lemma for a route family.
\end{itemize}

\paragraph{Required skills.}
\begin{itemize}
\item contradiction-route-auditor
\item entailment-gap-auditor
\item kb-deduce
\item kb-end-to-end-proof-critic
\item kb-hypothesis-refutation-loop
\item long-horizon-route-portfolio-manager
\item main-problem-alignment-auditor
\item next-experiment-queue-ranker
\item paper-contribution-readiness-judge
\item pvsnp-barrier-specificity-reviewer
\item pvsnp-conceptual-pressure-localizer
\item pvsnp-impossibility-theorem-forger
\item pvsnp-inference-quality-controller
\item pvsnp-proof-derivation-builder
\item pvsnp-qualitative-complexity-reasoner
\item pvsnp-quantifier-uniformity-reviewer
\item pvsnp-reduction-consistency-reviewer
\item pvsnp-symbolic-complexity-auditor
\end{itemize}

\subsection*{\#112 cycle-0023: Hardness-magnification bridge realization: generated-OR range-balance gate}

\paragraph{Analysis.} Instantiate a template generator for polylog-length 3CNFs and estimate \(\theta=min\) label-support exponent; then compare it with 1+epsilon-q and with the best promise-SAT solver size on the generated template range.

\paragraph{Why this is a good target.} It directly tests the current black-box substitution cap by computing \(\beta\), q, and the induced target-exponent window. It keeps the barrier and model-boundary audit attached to the mathematical transfer rather than as after-the-fact commentary. It is executable as a bounded complexity calculation, CAS-backed exponent audit, or formal transfer check.

\paragraph{Risk.} Low drift risk: the item is framed as an explicit P vs NP transfer audit.

\paragraph{Branch control.} The item has high payoff and a concrete decision test; stopping now would be premature.

\paragraph{Premature-stop guard.} Do not stop before the branch has produced a concrete source-to-target transfer theorem, an explicit obstruction lemma, or a barrier-aware reason why the route cannot reach the P vs NP/P/poly target.

\paragraph{Payoff criteria.}
\begin{itemize}
\item For the high-q generated-OR carrier, compute the exact leakage exponent of the remaining witness-cover or unary-residual mechanism and compare it with 1+epsilon-q.
\item Promotion requires \(q+a>1+\epsilon\) after every projection, padding, and witness-cover leakage term is accounted for; otherwise the branch becomes a bounded-window or no-go result.
\item For this route, state the exact barrier premise it uses or escapes: relativizing, natural, algebrizing, nonuniform, or advice-leaking; promotion requires the escape to be attached to the specific transfer theorem.
\item Compute \(\beta\) and q exactly; promotion requires \(\beta=0\), subpolynomial target queries, or a genuinely nonblackbox use of the SAT circuit.
\item Record relativization, natural-proofs, algebrization, nonuniformity, and advice-leakage status for the specific transfer theorem.
\item Control gate: name the source carrier or language, target class, and circuit or resource model before claiming progress.
\end{itemize}

\paragraph{Stop criteria.}
\begin{itemize}
\item The branch cannot identify a source language/carrier and target complexity class in the same model.
\item The route repeatedly proves only a proxy lower bound with no target-preserving transfer.
\item The remaining work is a slogan about hardness magnification rather than a computable transfer or obstruction.
\end{itemize}

\paragraph{Do not stop if.}
\begin{itemize}
\item A transfer exponent, quantifier repair, or barrier status remains uncomputed.
\item The branch can still name a target-preserving carrier and a falsifiable lower-bound threshold.
\item The failure can become a formal no-go lemma for a route family.
\end{itemize}

\paragraph{Required skills.}
\begin{itemize}
\item contradiction-route-auditor
\item entailment-gap-auditor
\item kb-deduce
\item kb-end-to-end-proof-critic
\item kb-hypothesis-refutation-loop
\item long-horizon-route-portfolio-manager
\item main-problem-alignment-auditor
\item next-experiment-queue-ranker
\item paper-contribution-readiness-judge
\item pvsnp-barrier-specificity-reviewer
\item pvsnp-conceptual-pressure-localizer
\item pvsnp-impossibility-theorem-forger
\item pvsnp-inference-quality-controller
\item pvsnp-proof-derivation-builder
\item pvsnp-qualitative-complexity-reasoner
\item pvsnp-quantifier-uniformity-reviewer
\item pvsnp-reduction-consistency-reviewer
\item pvsnp-symbolic-complexity-auditor
\end{itemize}

\subsection*{\#113 cycle-0022: Hardness-Magnification Bridge Realization}

\paragraph{Analysis.} Instantiate \(G_{N}\) with a high-entropy family of polylog-size 3-CNF formulas and compute the best promise solver or block-cover upper bound; vary \(\tau\) to locate the first threshold-respecting candidate.

\paragraph{Why this is a good target.} It directly tests the current black-box substitution cap by computing \(\beta\), q, and the induced target-exponent window. It keeps the barrier and model-boundary audit attached to the mathematical transfer rather than as after-the-fact commentary. It is executable as a bounded complexity calculation, CAS-backed exponent audit, or formal transfer check.

\paragraph{Risk.} Low drift risk: the item is framed as an explicit P vs NP transfer audit.

\paragraph{Branch control.} The item has high payoff and a concrete decision test; stopping now would be premature.

\paragraph{Premature-stop guard.} Do not stop before deriving the transfer inequality from \(\mathrm{SAT}\in \mathrm{SIZE}(n^{c})\) to the carrier upper bound, identifying whether \(\beta=0\), \(\beta>0\), or the argument is genuinely nonblackbox, and recording the barrier status of the proposed escape.

\paragraph{Payoff criteria.}
\begin{itemize}
\item For the high-q generated-OR carrier, compute the exact leakage exponent of the remaining witness-cover or unary-residual mechanism and compare it with 1+epsilon-q.
\item Promotion requires \(q+a>1+\epsilon\) after every projection, padding, and witness-cover leakage term is accounted for; otherwise the branch becomes a bounded-window or no-go result.
\item For this route, state the exact barrier premise it uses or escapes: relativizing, natural, algebrizing, nonuniform, or advice-leaking; promotion requires the escape to be attached to the specific transfer theorem.
\item Compute \(\beta\) and q exactly; promotion requires \(\beta=0\), subpolynomial target queries, or a genuinely nonblackbox use of the SAT circuit.
\item Record relativization, natural-proofs, algebrization, nonuniformity, and advice-leakage status for the specific transfer theorem.
\item Control gate: name the source carrier or language, target class, and circuit or resource model before claiming progress.
\end{itemize}

\paragraph{Stop criteria.}
\begin{itemize}
\item Every named carrier in the branch has \(\beta>0\) and yields only a bounded target-exponent exclusion.
\item The branch cannot state a source lower-bound model, target class, and transfer inequality in the same notation.
\item The proposed escape is only a black-box substitution restatement with no new carrier or nonblackbox mechanism.
\item The barrier audit shows the needed source lower bound is just as hard as the original P vs NP/P/poly target.
\end{itemize}

\paragraph{Do not stop if.}
\begin{itemize}
\item The branch has not yet computed \(\beta\), q, and the induced source upper-bound exponent under \(\mathrm{SAT}\in \mathrm{SIZE}(n^{c})\).
\item A subpolynomial-query or nonblackbox carrier remains available and has not been audited.
\item The latest failure sharpened the threshold-collapse lemma into a formal no-go statement.
\item The route can still isolate advice leakage, nonuniformity, naturalness, or algebrization as the exact obstruction.
\end{itemize}

\paragraph{Required skills.}
\begin{itemize}
\item contradiction-route-auditor
\item entailment-gap-auditor
\item kb-deduce
\item kb-end-to-end-proof-critic
\item kb-hypothesis-refutation-loop
\item long-horizon-route-portfolio-manager
\item main-problem-alignment-auditor
\item next-experiment-queue-ranker
\item paper-contribution-readiness-judge
\item pvsnp-barrier-specificity-reviewer
\item pvsnp-conceptual-pressure-localizer
\item pvsnp-impossibility-theorem-forger
\item pvsnp-inference-quality-controller
\item pvsnp-proof-derivation-builder
\item pvsnp-qualitative-complexity-reasoner
\item pvsnp-quantifier-uniformity-reviewer
\item pvsnp-reduction-consistency-reviewer
\item pvsnp-symbolic-complexity-auditor
\end{itemize}

\subsection*{\#114 cycle-0021: Hardness-Magnification Bridge Realization: Shared-Generator Semantic Block-Cover Gate}

\paragraph{Analysis.} Run the block-size thought experiment for \(b=N^\theta\): derive upper bounds on U(theta) from range cardinality, feature compression, and OR compression, then identify whether any \(\theta\) gives N\textasciicircum{}\{rho\}+N\textasciicircum{}\{q-theta\}U(theta)<N\textasciicircum{}\{1+epsilon\}.

\paragraph{Why this is a good target.} It directly tests the current black-box substitution cap by computing \(\beta\), q, and the induced target-exponent window. It keeps the barrier and model-boundary audit attached to the mathematical transfer rather than as after-the-fact commentary. It is executable as a bounded complexity calculation, CAS-backed exponent audit, or formal transfer check.

\paragraph{Risk.} Low drift risk: the item is framed as an explicit P vs NP transfer audit.

\paragraph{Branch control.} The item has high payoff and a concrete decision test; stopping now would be premature.

\paragraph{Premature-stop guard.} Do not stop before the branch has produced a concrete source-to-target transfer theorem, an explicit obstruction lemma, or a barrier-aware reason why the route cannot reach the P vs NP/P/poly target.

\paragraph{Payoff criteria.}
\begin{itemize}
\item For the high-q generated-OR carrier, compute the exact leakage exponent of the remaining witness-cover or unary-residual mechanism and compare it with 1+epsilon-q.
\item Promotion requires \(q+a>1+\epsilon\) after every projection, padding, and witness-cover leakage term is accounted for; otherwise the branch becomes a bounded-window or no-go result.
\item For this route, state the exact barrier premise it uses or escapes: relativizing, natural, algebrizing, nonuniform, or advice-leaking; promotion requires the escape to be attached to the specific transfer theorem.
\item Compute \(\beta\) and q exactly; promotion requires \(\beta=0\), subpolynomial target queries, or a genuinely nonblackbox use of the SAT circuit.
\item Record relativization, natural-proofs, algebrization, nonuniformity, and advice-leakage status for the specific transfer theorem.
\item Control gate: name the source carrier or language, target class, and circuit or resource model before claiming progress.
\end{itemize}

\paragraph{Stop criteria.}
\begin{itemize}
\item The branch cannot identify a source language/carrier and target complexity class in the same model.
\item The route repeatedly proves only a proxy lower bound with no target-preserving transfer.
\item The remaining work is a slogan about hardness magnification rather than a computable transfer or obstruction.
\end{itemize}

\paragraph{Do not stop if.}
\begin{itemize}
\item A transfer exponent, quantifier repair, or barrier status remains uncomputed.
\item The branch can still name a target-preserving carrier and a falsifiable lower-bound threshold.
\item The failure can become a formal no-go lemma for a route family.
\end{itemize}

\paragraph{Required skills.}
\begin{itemize}
\item contradiction-route-auditor
\item entailment-gap-auditor
\item kb-deduce
\item kb-end-to-end-proof-critic
\item kb-hypothesis-refutation-loop
\item long-horizon-route-portfolio-manager
\item main-problem-alignment-auditor
\item next-experiment-queue-ranker
\item paper-contribution-readiness-judge
\item pvsnp-barrier-specificity-reviewer
\item pvsnp-conceptual-pressure-localizer
\item pvsnp-impossibility-theorem-forger
\item pvsnp-inference-quality-controller
\item pvsnp-proof-derivation-builder
\item pvsnp-qualitative-complexity-reasoner
\item pvsnp-quantifier-uniformity-reviewer
\item pvsnp-reduction-consistency-reviewer
\item pvsnp-symbolic-complexity-auditor
\end{itemize}

\subsection*{\#115 cycle-0017: Hardness-Magnification Bridge Realization: positioned-range entropy audit}

\paragraph{Analysis.} Instantiate \(K_{N}\) from a specific extractor-like generator of polylog-length CNFs and compute \(|R_{g}|\), \(v_{N}\), q, and \(\rho\) exactly to see whether the positioned entropy threshold can be met without advice leakage.

\paragraph{Why this is a good target.} It directly tests the current black-box substitution cap by computing \(\beta\), q, and the induced target-exponent window. It keeps the barrier and model-boundary audit attached to the mathematical transfer rather than as after-the-fact commentary. It is executable as a bounded complexity calculation, CAS-backed exponent audit, or formal transfer check.

\paragraph{Risk.} Low drift risk: the item is framed as an explicit P vs NP transfer audit.

\paragraph{Branch control.} The item has high payoff and a concrete decision test; stopping now would be premature.

\paragraph{Premature-stop guard.} Do not stop before deriving the transfer inequality from \(\mathrm{SAT}\in \mathrm{SIZE}(n^{c})\) to the carrier upper bound, identifying whether \(\beta=0\), \(\beta>0\), or the argument is genuinely nonblackbox, and recording the barrier status of the proposed escape.

\paragraph{Payoff criteria.}
\begin{itemize}
\item For the high-q generated-OR carrier, compute the exact leakage exponent of the remaining witness-cover or unary-residual mechanism and compare it with 1+epsilon-q.
\item Promotion requires \(q+a>1+\epsilon\) after every projection, padding, and witness-cover leakage term is accounted for; otherwise the branch becomes a bounded-window or no-go result.
\item For this route, state the exact barrier premise it uses or escapes: relativizing, natural, algebrizing, nonuniform, or advice-leaking; promotion requires the escape to be attached to the specific transfer theorem.
\item Compute \(\beta\) and q exactly; promotion requires \(\beta=0\), subpolynomial target queries, or a genuinely nonblackbox use of the SAT circuit.
\item Record relativization, natural-proofs, algebrization, nonuniformity, and advice-leakage status for the specific transfer theorem.
\item Control gate: name the source carrier or language, target class, and circuit or resource model before claiming progress.
\end{itemize}

\paragraph{Stop criteria.}
\begin{itemize}
\item Every named carrier in the branch has \(\beta>0\) and yields only a bounded target-exponent exclusion.
\item The branch cannot state a source lower-bound model, target class, and transfer inequality in the same notation.
\item The proposed escape is only a black-box substitution restatement with no new carrier or nonblackbox mechanism.
\item The barrier audit shows the needed source lower bound is just as hard as the original P vs NP/P/poly target.
\end{itemize}

\paragraph{Do not stop if.}
\begin{itemize}
\item The branch has not yet computed \(\beta\), q, and the induced source upper-bound exponent under \(\mathrm{SAT}\in \mathrm{SIZE}(n^{c})\).
\item A subpolynomial-query or nonblackbox carrier remains available and has not been audited.
\item The latest failure sharpened the threshold-collapse lemma into a formal no-go statement.
\item The route can still isolate advice leakage, nonuniformity, naturalness, or algebrization as the exact obstruction.
\end{itemize}

\paragraph{Required skills.}
\begin{itemize}
\item contradiction-route-auditor
\item entailment-gap-auditor
\item kb-deduce
\item kb-end-to-end-proof-critic
\item kb-hypothesis-refutation-loop
\item long-horizon-route-portfolio-manager
\item main-problem-alignment-auditor
\item next-experiment-queue-ranker
\item paper-contribution-readiness-judge
\item pvsnp-barrier-specificity-reviewer
\item pvsnp-conceptual-pressure-localizer
\item pvsnp-impossibility-theorem-forger
\item pvsnp-inference-quality-controller
\item pvsnp-proof-derivation-builder
\item pvsnp-qualitative-complexity-reasoner
\item pvsnp-quantifier-uniformity-reviewer
\item pvsnp-reduction-consistency-reviewer
\item pvsnp-symbolic-complexity-auditor
\end{itemize}

\subsection*{\#116 cycle-0016: Hardness-Magnification Bridge Realization}

\paragraph{Analysis.} Instantiate \(G_{N}\) with high-essential-variable polylog 3CNF queries and test whether any structural simplification gives \(\mathrm{SIZE}(H_{N})\le N^{1+\epsilon-o(1)}\) despite v(N) superlogarithmic.

\paragraph{Why this is a good target.} It directly tests the current black-box substitution cap by computing \(\beta\), q, and the induced target-exponent window. It keeps the barrier and model-boundary audit attached to the mathematical transfer rather than as after-the-fact commentary. It is executable as a bounded complexity calculation, CAS-backed exponent audit, or formal transfer check.

\paragraph{Risk.} Low drift risk: the item is framed as an explicit P vs NP transfer audit.

\paragraph{Branch control.} The item has high payoff and a concrete decision test; stopping now would be premature.

\paragraph{Premature-stop guard.} Do not stop before the branch has produced a concrete source-to-target transfer theorem, an explicit obstruction lemma, or a barrier-aware reason why the route cannot reach the P vs NP/P/poly target.

\paragraph{Payoff criteria.}
\begin{itemize}
\item For the high-q generated-OR carrier, compute the exact leakage exponent of the remaining witness-cover or unary-residual mechanism and compare it with 1+epsilon-q.
\item Promotion requires \(q+a>1+\epsilon\) after every projection, padding, and witness-cover leakage term is accounted for; otherwise the branch becomes a bounded-window or no-go result.
\item For this route, state the exact barrier premise it uses or escapes: relativizing, natural, algebrizing, nonuniform, or advice-leaking; promotion requires the escape to be attached to the specific transfer theorem.
\item Compute \(\beta\) and q exactly; promotion requires \(\beta=0\), subpolynomial target queries, or a genuinely nonblackbox use of the SAT circuit.
\item Record relativization, natural-proofs, algebrization, nonuniformity, and advice-leakage status for the specific transfer theorem.
\item Control gate: name the source carrier or language, target class, and circuit or resource model before claiming progress.
\end{itemize}

\paragraph{Stop criteria.}
\begin{itemize}
\item The branch cannot identify a source language/carrier and target complexity class in the same model.
\item The route repeatedly proves only a proxy lower bound with no target-preserving transfer.
\item The remaining work is a slogan about hardness magnification rather than a computable transfer or obstruction.
\end{itemize}

\paragraph{Do not stop if.}
\begin{itemize}
\item A transfer exponent, quantifier repair, or barrier status remains uncomputed.
\item The branch can still name a target-preserving carrier and a falsifiable lower-bound threshold.
\item The failure can become a formal no-go lemma for a route family.
\end{itemize}

\paragraph{Required skills.}
\begin{itemize}
\item contradiction-route-auditor
\item entailment-gap-auditor
\item kb-deduce
\item kb-end-to-end-proof-critic
\item kb-hypothesis-refutation-loop
\item long-horizon-route-portfolio-manager
\item main-problem-alignment-auditor
\item next-experiment-queue-ranker
\item paper-contribution-readiness-judge
\item pvsnp-barrier-specificity-reviewer
\item pvsnp-conceptual-pressure-localizer
\item pvsnp-impossibility-theorem-forger
\item pvsnp-inference-quality-controller
\item pvsnp-proof-derivation-builder
\item pvsnp-qualitative-complexity-reasoner
\item pvsnp-quantifier-uniformity-reviewer
\item pvsnp-reduction-consistency-reviewer
\item pvsnp-symbolic-complexity-auditor
\end{itemize}

\subsection*{\#117 cycle-0015: Hardness-Magnification Bridge Realization}

\paragraph{Analysis.} Instantiate the boundary cases \(d=1\) with \(A=1+\epsilon-q\) and \(d=2\); search for a generator whose \(H_{N}\) avoids known small-circuit constructions without smuggling SAT advice.

\paragraph{Why this is a good target.} It directly tests the current black-box substitution cap by computing \(\beta\), q, and the induced target-exponent window. It keeps the barrier and model-boundary audit attached to the mathematical transfer rather than as after-the-fact commentary. It is executable as a bounded complexity calculation, CAS-backed exponent audit, or formal transfer check.

\paragraph{Risk.} Low drift risk: the item is framed as an explicit P vs NP transfer audit.

\paragraph{Branch control.} The item has high payoff and a concrete decision test; stopping now would be premature.

\paragraph{Premature-stop guard.} Do not stop before deriving the transfer inequality from \(\mathrm{SAT}\in \mathrm{SIZE}(n^{c})\) to the carrier upper bound, identifying whether \(\beta=0\), \(\beta>0\), or the argument is genuinely nonblackbox, and recording the barrier status of the proposed escape.

\paragraph{Payoff criteria.}
\begin{itemize}
\item For this route, state the exact barrier premise it uses or escapes: relativizing, natural, algebrizing, nonuniform, or advice-leaking; promotion requires the escape to be attached to the specific transfer theorem.
\item Compute \(\beta\) and q exactly; promotion requires \(\beta=0\), subpolynomial target queries, or a genuinely nonblackbox use of the SAT circuit.
\item Record relativization, natural-proofs, algebrization, nonuniformity, and advice-leakage status for the specific transfer theorem.
\item Control gate: name the source carrier or language, target class, and circuit or resource model before claiming progress.
\item Control gate: derive the exact transfer inequality from \(\mathrm{SAT}\in \mathrm{SIZE}(n^{c})\), NP subset P/poly, or the stated target assumption to the source-carrier upper bound.
\item Control gate: compute \(\beta\) and q exactly, distinguishing all-exponent consequences from bounded target-exponent windows.
\end{itemize}

\paragraph{Stop criteria.}
\begin{itemize}
\item Every named carrier in the branch has \(\beta>0\) and yields only a bounded target-exponent exclusion.
\item The branch cannot state a source lower-bound model, target class, and transfer inequality in the same notation.
\item The proposed escape is only a black-box substitution restatement with no new carrier or nonblackbox mechanism.
\item The barrier audit shows the needed source lower bound is just as hard as the original P vs NP/P/poly target.
\end{itemize}

\paragraph{Do not stop if.}
\begin{itemize}
\item The branch has not yet computed \(\beta\), q, and the induced source upper-bound exponent under \(\mathrm{SAT}\in \mathrm{SIZE}(n^{c})\).
\item A subpolynomial-query or nonblackbox carrier remains available and has not been audited.
\item The latest failure sharpened the threshold-collapse lemma into a formal no-go statement.
\item The route can still isolate advice leakage, nonuniformity, naturalness, or algebrization as the exact obstruction.
\end{itemize}

\paragraph{Required skills.}
\begin{itemize}
\item contradiction-route-auditor
\item entailment-gap-auditor
\item kb-deduce
\item kb-end-to-end-proof-critic
\item kb-hypothesis-refutation-loop
\item long-horizon-route-portfolio-manager
\item main-problem-alignment-auditor
\item next-experiment-queue-ranker
\item paper-contribution-readiness-judge
\item pvsnp-barrier-specificity-reviewer
\item pvsnp-conceptual-pressure-localizer
\item pvsnp-impossibility-theorem-forger
\item pvsnp-inference-quality-controller
\item pvsnp-proof-derivation-builder
\item pvsnp-qualitative-complexity-reasoner
\item pvsnp-quantifier-uniformity-reviewer
\item pvsnp-reduction-consistency-reviewer
\item pvsnp-symbolic-complexity-auditor
\end{itemize}

\subsection*{\#118 cycle-0057: Hardness-Magnification Bridge Realization: XOR quotient carrier, low-degree feature-sketch gate}

\paragraph{Analysis.} Formalize the degree-2 case: characterize whether Tseitin-style parity, formula-size, variable-incidence, or clause-type features give \(r=N^\alpha\) with \(2alpha<1+\epsilon\) for the generated length-(log N)\textasciicircum{}d SAT queries.

\paragraph{Why this is a good target.} It audits a meta-complexity carrier where the decisive questions are verifier overhead, density/naturalness, and advice leakage. It keeps the barrier and model-boundary audit attached to the mathematical transfer rather than as after-the-fact commentary. It is executable as a bounded complexity calculation, CAS-backed exponent audit, or formal transfer check.

\paragraph{Risk.} Moderate risk: run only if source/target/model/barrier fields are explicit.

\paragraph{Branch control.} The item remains target-linked but needs tighter payoff criteria before consuming a long run.

\paragraph{Premature-stop guard.} Do not stop before naming the carrier variant, deriving its verifier/reduction overhead, and checking whether a \(\mathrm{SAT}\in \mathrm{SIZE}(n^{c})\) assumption gives carrier circuits below the claimed source lower-bound threshold.

\paragraph{Payoff criteria.}
\begin{itemize}
\item For the restricted-model branch, compute the compiler envelope from unrestricted SIZE(m\textasciicircum{}k) or P-time algorithms into the named model; promotion requires a bound above that envelope on the same SAT slice.
\item Falsify direct transfer if a P function already has the same restricted-model lower bound, such as PARITY for \(\mathrm{AC0}\), HWB for OBDD, or XOR/Tseitin for proof-system degree.
\item For the proof-complexity branch, separate fixed-family proof hardness from language hardness by naming the proof system, shortcut algorithm, and semantic lower-bound target.
\item Promotion requires a proof-system lower bound that survives the Gaussian-elimination or recognition shortcut; otherwise the result is a restricted-proof-system separation only.
\item For this route, state the exact barrier premise it uses or escapes: relativizing, natural, algebrizing, nonuniform, or advice-leaking; promotion requires the escape to be attached to the specific transfer theorem.
\item For MCSP/MKtP, state N, n, s, verifier overhead, lower-bound model, and whether the property is natural or non-natural.
\end{itemize}

\paragraph{Stop criteria.}
\begin{itemize}
\item The induced carrier upper bound under \(\mathrm{SAT}\in \mathrm{SIZE}(n^{c})\) remains N\textasciicircum{}\{q+beta c+o(1)\} with \(\beta>0\) for all tested parameter choices.
\item The source lower bound required for MCSP/MKtP is as strong as the target separation or becomes a natural property without an escape.
\item The route cannot prevent advice leakage or model mismatch between the carrier lower bound and P/poly target.
\end{itemize}

\paragraph{Do not stop if.}
\begin{itemize}
\item The MCSP/MKtP parameter regime has not yet been fixed.
\item The truth-table length N and underlying input length n have not been related in the transfer inequality.
\item The natural-proof largeness/constructivity issue has not yet been audited.
\item A nonblackbox use of the target circuit could still avoid the multiplicative c factor.
\end{itemize}

\paragraph{Required skills.}
\begin{itemize}
\item counterexample-hunter
\item entailment-gap-auditor
\item kb-abduce
\item kb-concept-bridge-planner
\item kb-end-to-end-proof-critic
\item long-horizon-route-portfolio-manager
\item main-problem-alignment-auditor
\item next-experiment-queue-ranker
\item paper-contribution-readiness-judge
\item pvsnp-barrier-specificity-reviewer
\item pvsnp-conceptual-pressure-localizer
\item pvsnp-impossibility-theorem-forger
\item pvsnp-inference-quality-controller
\item pvsnp-meta-complexity-carrier-reviewer
\item pvsnp-qualitative-complexity-reasoner
\item pvsnp-quantifier-uniformity-reviewer
\item pvsnp-symbolic-complexity-auditor
\end{itemize}

\subsection*{\#119 cycle-0046: Hardness-Magnification Bridge Realization: normal-form extraction gate}

\paragraph{Analysis.} Formalize a restricted no-control generator class, such as block-product independent seeds with no shared label bit, and test whether an S-size circuit for \(OR_{j}\) SAT(G(x,j)) can be localized to a one-sided separator cover; compute the exact overhead in q,rho,sigma,E.

\paragraph{Why this is a good target.} It audits a meta-complexity carrier where the decisive questions are verifier overhead, density/naturalness, and advice leakage. It keeps the barrier and model-boundary audit attached to the mathematical transfer rather than as after-the-fact commentary. It is executable as a bounded complexity calculation, CAS-backed exponent audit, or formal transfer check.

\paragraph{Risk.} Moderate risk: run only if source/target/model/barrier fields are explicit.

\paragraph{Branch control.} The item remains target-linked but needs tighter payoff criteria before consuming a long run.

\paragraph{Premature-stop guard.} Do not stop before naming the carrier variant, deriving its verifier/reduction overhead, and checking whether a \(\mathrm{SAT}\in \mathrm{SIZE}(n^{c})\) assumption gives carrier circuits below the claimed source lower-bound threshold.

\paragraph{Payoff criteria.}
\begin{itemize}
\item For the one-sided semantic-separator branch, compute the cap profile A\_N(s) for the named generator and predicate class; promotion requires a bound strong enough to force Sep\_N\textasciicircum{}\{1s\}(E)>N\textasciicircum{}\{1+epsilon-o(1)\}.
\item Exhibit or refute a subthreshold semantic planted-slice predicate; this decides whether the generated-OR carrier needs range-closure repair, stronger pseudorandomness, or quarantine.
\item For the high-q generated-OR carrier, compute the exact leakage exponent of the remaining witness-cover or unary-residual mechanism and compare it with 1+epsilon-q.
\item Promotion requires \(q+a>1+\epsilon\) after every projection, padding, and witness-cover leakage term is accounted for; otherwise the branch becomes a bounded-window or no-go result.
\item For this route, state the exact barrier premise it uses or escapes: relativizing, natural, algebrizing, nonuniform, or advice-leaking; promotion requires the escape to be attached to the specific transfer theorem.
\item For MCSP/MKtP, state N, n, s, verifier overhead, lower-bound model, and whether the property is natural or non-natural.
\end{itemize}

\paragraph{Stop criteria.}
\begin{itemize}
\item The induced carrier upper bound under \(\mathrm{SAT}\in \mathrm{SIZE}(n^{c})\) remains N\textasciicircum{}\{q+beta c+o(1)\} with \(\beta>0\) for all tested parameter choices.
\item The source lower bound required for MCSP/MKtP is as strong as the target separation or becomes a natural property without an escape.
\item The route cannot prevent advice leakage or model mismatch between the carrier lower bound and P/poly target.
\end{itemize}

\paragraph{Do not stop if.}
\begin{itemize}
\item The MCSP/MKtP parameter regime has not yet been fixed.
\item The truth-table length N and underlying input length n have not been related in the transfer inequality.
\item The natural-proof largeness/constructivity issue has not yet been audited.
\item A nonblackbox use of the target circuit could still avoid the multiplicative c factor.
\end{itemize}

\paragraph{Required skills.}
\begin{itemize}
\item counterexample-hunter
\item entailment-gap-auditor
\item kb-abduce
\item kb-concept-bridge-planner
\item kb-end-to-end-proof-critic
\item long-horizon-route-portfolio-manager
\item main-problem-alignment-auditor
\item next-experiment-queue-ranker
\item paper-contribution-readiness-judge
\item pvsnp-barrier-specificity-reviewer
\item pvsnp-conceptual-pressure-localizer
\item pvsnp-impossibility-theorem-forger
\item pvsnp-inference-quality-controller
\item pvsnp-meta-complexity-carrier-reviewer
\item pvsnp-qualitative-complexity-reasoner
\item pvsnp-quantifier-uniformity-reviewer
\item pvsnp-symbolic-complexity-auditor
\end{itemize}

\subsection*{\#120 cycle-0045: Hardness-Magnification Bridge Realization: Greedy Cap Dual Gate}

\paragraph{Analysis.} Work in a restricted circuit model such as De Morgan formulas or \(\mathrm{AC0}[p]\) and test whether exact circuits for \(H_{N}\) admit an approximate one-sided separator normal form; compute the resulting \(D_{sigma,E}\) threshold for the same high-q polylog generated-OR carrier.

\paragraph{Why this is a good target.} It audits a meta-complexity carrier where the decisive questions are verifier overhead, density/naturalness, and advice leakage. It keeps the barrier and model-boundary audit attached to the mathematical transfer rather than as after-the-fact commentary. It is executable as a bounded complexity calculation, CAS-backed exponent audit, or formal transfer check.

\paragraph{Risk.} Moderate risk: run only if source/target/model/barrier fields are explicit.

\paragraph{Branch control.} The item remains target-linked but needs tighter payoff criteria before consuming a long run.

\paragraph{Premature-stop guard.} Do not stop before naming the carrier variant, deriving its verifier/reduction overhead, and checking whether a \(\mathrm{SAT}\in \mathrm{SIZE}(n^{c})\) assumption gives carrier circuits below the claimed source lower-bound threshold.

\paragraph{Payoff criteria.}
\begin{itemize}
\item For the one-sided semantic-separator branch, compute the cap profile A\_N(s) for the named generator and predicate class; promotion requires a bound strong enough to force Sep\_N\textasciicircum{}\{1s\}(E)>N\textasciicircum{}\{1+epsilon-o(1)\}.
\item Exhibit or refute a subthreshold semantic planted-slice predicate; this decides whether the generated-OR carrier needs range-closure repair, stronger pseudorandomness, or quarantine.
\item For the certificate-cover branch, compute K,E,L and the induced circuit upper bound C(H\_N)<=C(G\_N)+O(K(s+L\textasciicircum{}\{O(1)\})+EN); promotion requires every sound local cover to stay above the \(N^{1+\epsilon-o(1)}\) source threshold.
\item Falsify the branch if Horn, 2-SAT, bounded-width CNF, or bounded-backdoor emitted formulas give subthreshold satisfying-term or sound-local-certificate covers.
\item For the high-q generated-OR carrier, compute the exact leakage exponent of the remaining witness-cover or unary-residual mechanism and compare it with 1+epsilon-q.
\item Promotion requires \(q+a>1+\epsilon\) after every projection, padding, and witness-cover leakage term is accounted for; otherwise the branch becomes a bounded-window or no-go result.
\end{itemize}

\paragraph{Stop criteria.}
\begin{itemize}
\item The induced carrier upper bound under \(\mathrm{SAT}\in \mathrm{SIZE}(n^{c})\) remains N\textasciicircum{}\{q+beta c+o(1)\} with \(\beta>0\) for all tested parameter choices.
\item The source lower bound required for MCSP/MKtP is as strong as the target separation or becomes a natural property without an escape.
\item The route cannot prevent advice leakage or model mismatch between the carrier lower bound and P/poly target.
\end{itemize}

\paragraph{Do not stop if.}
\begin{itemize}
\item The MCSP/MKtP parameter regime has not yet been fixed.
\item The truth-table length N and underlying input length n have not been related in the transfer inequality.
\item The natural-proof largeness/constructivity issue has not yet been audited.
\item A nonblackbox use of the target circuit could still avoid the multiplicative c factor.
\end{itemize}

\paragraph{Required skills.}
\begin{itemize}
\item counterexample-hunter
\item entailment-gap-auditor
\item kb-abduce
\item kb-concept-bridge-planner
\item kb-end-to-end-proof-critic
\item long-horizon-route-portfolio-manager
\item main-problem-alignment-auditor
\item next-experiment-queue-ranker
\item paper-contribution-readiness-judge
\item pvsnp-barrier-specificity-reviewer
\item pvsnp-conceptual-pressure-localizer
\item pvsnp-impossibility-theorem-forger
\item pvsnp-inference-quality-controller
\item pvsnp-meta-complexity-carrier-reviewer
\item pvsnp-qualitative-complexity-reasoner
\item pvsnp-quantifier-uniformity-reviewer
\item pvsnp-symbolic-complexity-auditor
\end{itemize}

\subsection*{\#121 cycle-0149: Hardness-Magnification Bridge Realization}

\paragraph{Analysis.} Write the standalone paper theorem for \(TransparentRowSupport_{N}^{Qrep}\): include the \(K_{Qrep}\) fixed-point verifier pseudocode, the \(\Gamma_{Qrep}\) ledger proof, the first-rejected-merge countermodel, and a one-page nonclaim box separating this parser boundary from \(H_{N}\) lower bounds and \(\mathrm{P} \ne \mathrm{NP}\).

\paragraph{Why this is a good target.} It directly tests the current black-box substitution cap by computing \(\beta\), q, and the induced target-exponent window. It keeps the barrier and model-boundary audit attached to the mathematical transfer rather than as after-the-fact commentary.

\paragraph{Risk.} Moderate risk: run only if source/target/model/barrier fields are explicit.

\paragraph{Branch control.} The item remains target-linked but needs tighter payoff criteria before consuming a long run.

\paragraph{Premature-stop guard.} Do not stop before deriving the transfer inequality from \(\mathrm{SAT}\in \mathrm{SIZE}(n^{c})\) to the carrier upper bound, identifying whether \(\beta=0\), \(\beta>0\), or the argument is genuinely nonblackbox, and recording the barrier status of the proposed escape.

\paragraph{Payoff criteria.}
\begin{itemize}
\item For this route, state the exact barrier premise it uses or escapes: relativizing, natural, algebrizing, nonuniform, or advice-leaking; promotion requires the escape to be attached to the specific transfer theorem.
\item Compute \(\beta\) and q exactly; promotion requires \(\beta=0\), subpolynomial target queries, or a genuinely nonblackbox use of the SAT circuit.
\item Record relativization, natural-proofs, algebrization, nonuniformity, and advice-leakage status for the specific transfer theorem.
\item Control gate: name the source carrier or language, target class, and circuit or resource model before claiming progress.
\item Control gate: derive the exact transfer inequality from \(\mathrm{SAT}\in \mathrm{SIZE}(n^{c})\), NP subset P/poly, or the stated target assumption to the source-carrier upper bound.
\item Control gate: compute \(\beta\) and q exactly, distinguishing all-exponent consequences from bounded target-exponent windows.
\end{itemize}

\paragraph{Stop criteria.}
\begin{itemize}
\item Every named carrier in the branch has \(\beta>0\) and yields only a bounded target-exponent exclusion.
\item The branch cannot state a source lower-bound model, target class, and transfer inequality in the same notation.
\item The proposed escape is only a black-box substitution restatement with no new carrier or nonblackbox mechanism.
\item The barrier audit shows the needed source lower bound is just as hard as the original P vs NP/P/poly target.
\end{itemize}

\paragraph{Do not stop if.}
\begin{itemize}
\item The branch has not yet computed \(\beta\), q, and the induced source upper-bound exponent under \(\mathrm{SAT}\in \mathrm{SIZE}(n^{c})\).
\item A subpolynomial-query or nonblackbox carrier remains available and has not been audited.
\item The latest failure sharpened the threshold-collapse lemma into a formal no-go statement.
\item The route can still isolate advice leakage, nonuniformity, naturalness, or algebrization as the exact obstruction.
\end{itemize}

\paragraph{Required skills.}
\begin{itemize}
\item contradiction-route-auditor
\item entailment-gap-auditor
\item kb-deduce
\item kb-end-to-end-proof-critic
\item kb-hypothesis-refutation-loop
\item long-horizon-route-portfolio-manager
\item main-problem-alignment-auditor
\item next-experiment-queue-ranker
\item paper-contribution-readiness-judge
\item pvsnp-barrier-specificity-reviewer
\item pvsnp-conceptual-pressure-localizer
\item pvsnp-impossibility-theorem-forger
\item pvsnp-inference-quality-controller
\item pvsnp-proof-derivation-builder
\item pvsnp-qualitative-complexity-reasoner
\item pvsnp-quantifier-uniformity-reviewer
\item pvsnp-reduction-consistency-reviewer
\item pvsnp-symbolic-complexity-auditor
\end{itemize}

\subsection*{\#122 cycle-0127: Hardness-Magnification Bridge Realization}

\paragraph{Analysis.} Write the verifier bit fields for one declared TSR line: root row id, restriction-node path, pivot variable, child row ids, Tseitin map hash, occurrence-neighborhood list, parent pointers, and expansion certificate; then prove the single-line induction inequality adding exactly N times the declared child count to \(\Gamma_{dag+\mathrm{TSR}}\).

\paragraph{Why this is a good target.} It directly tests the current black-box substitution cap by computing \(\beta\), q, and the induced target-exponent window. It keeps the barrier and model-boundary audit attached to the mathematical transfer rather than as after-the-fact commentary.

\paragraph{Risk.} Moderate risk: run only if source/target/model/barrier fields are explicit.

\paragraph{Branch control.} The item remains target-linked but needs tighter payoff criteria before consuming a long run.

\paragraph{Premature-stop guard.} Do not stop before deriving the transfer inequality from \(\mathrm{SAT}\in \mathrm{SIZE}(n^{c})\) to the carrier upper bound, identifying whether \(\beta=0\), \(\beta>0\), or the argument is genuinely nonblackbox, and recording the barrier status of the proposed escape.

\paragraph{Payoff criteria.}
\begin{itemize}
\item For the proof-complexity branch, separate fixed-family proof hardness from language hardness by naming the proof system, shortcut algorithm, and semantic lower-bound target.
\item Promotion requires a proof-system lower bound that survives the Gaussian-elimination or recognition shortcut; otherwise the result is a restricted-proof-system separation only.
\item For this route, state the exact barrier premise it uses or escapes: relativizing, natural, algebrizing, nonuniform, or advice-leaking; promotion requires the escape to be attached to the specific transfer theorem.
\item Compute \(\beta\) and q exactly; promotion requires \(\beta=0\), subpolynomial target queries, or a genuinely nonblackbox use of the SAT circuit.
\item Record relativization, natural-proofs, algebrization, nonuniformity, and advice-leakage status for the specific transfer theorem.
\item Control gate: name the source carrier or language, target class, and circuit or resource model before claiming progress.
\end{itemize}

\paragraph{Stop criteria.}
\begin{itemize}
\item Every named carrier in the branch has \(\beta>0\) and yields only a bounded target-exponent exclusion.
\item The branch cannot state a source lower-bound model, target class, and transfer inequality in the same notation.
\item The proposed escape is only a black-box substitution restatement with no new carrier or nonblackbox mechanism.
\item The barrier audit shows the needed source lower bound is just as hard as the original P vs NP/P/poly target.
\end{itemize}

\paragraph{Do not stop if.}
\begin{itemize}
\item The branch has not yet computed \(\beta\), q, and the induced source upper-bound exponent under \(\mathrm{SAT}\in \mathrm{SIZE}(n^{c})\).
\item A subpolynomial-query or nonblackbox carrier remains available and has not been audited.
\item The latest failure sharpened the threshold-collapse lemma into a formal no-go statement.
\item The route can still isolate advice leakage, nonuniformity, naturalness, or algebrization as the exact obstruction.
\end{itemize}

\paragraph{Required skills.}
\begin{itemize}
\item contradiction-route-auditor
\item entailment-gap-auditor
\item kb-deduce
\item kb-end-to-end-proof-critic
\item kb-hypothesis-refutation-loop
\item long-horizon-route-portfolio-manager
\item main-problem-alignment-auditor
\item next-experiment-queue-ranker
\item paper-contribution-readiness-judge
\item pvsnp-barrier-specificity-reviewer
\item pvsnp-conceptual-pressure-localizer
\item pvsnp-impossibility-theorem-forger
\item pvsnp-inference-quality-controller
\item pvsnp-proof-derivation-builder
\item pvsnp-qualitative-complexity-reasoner
\item pvsnp-quantifier-uniformity-reviewer
\item pvsnp-reduction-consistency-reviewer
\item pvsnp-symbolic-complexity-auditor
\end{itemize}

\subsection*{\#123 cycle-0117: Hardness-Magnification Bridge Realization: Transparent FormulaEval Boundary}

\paragraph{Analysis.} Write the paper-ready theorem statement as three clauses: \(TransparentRowSupport_{N}\) positive \(\Gamma_{trans}\) induction; extensional hidden-support private-row/unread-enable-bit countermodel; \(RealSATInfluenceRealizer_{N}\) open semantic escape with overhead \(R_{N}\) and explicit non-claims.

\paragraph{Why this is a good target.} It directly tests the current black-box substitution cap by computing \(\beta\), q, and the induced target-exponent window. It keeps the barrier and model-boundary audit attached to the mathematical transfer rather than as after-the-fact commentary.

\paragraph{Risk.} Moderate risk: run only if source/target/model/barrier fields are explicit.

\paragraph{Branch control.} The item remains target-linked but needs tighter payoff criteria before consuming a long run.

\paragraph{Premature-stop guard.} Do not stop before deriving the transfer inequality from \(\mathrm{SAT}\in \mathrm{SIZE}(n^{c})\) to the carrier upper bound, identifying whether \(\beta=0\), \(\beta>0\), or the argument is genuinely nonblackbox, and recording the barrier status of the proposed escape.

\paragraph{Payoff criteria.}
\begin{itemize}
\item For this route, state the exact barrier premise it uses or escapes: relativizing, natural, algebrizing, nonuniform, or advice-leaking; promotion requires the escape to be attached to the specific transfer theorem.
\item Compute \(\beta\) and q exactly; promotion requires \(\beta=0\), subpolynomial target queries, or a genuinely nonblackbox use of the SAT circuit.
\item Record relativization, natural-proofs, algebrization, nonuniformity, and advice-leakage status for the specific transfer theorem.
\item Control gate: name the source carrier or language, target class, and circuit or resource model before claiming progress.
\item Control gate: derive the exact transfer inequality from \(\mathrm{SAT}\in \mathrm{SIZE}(n^{c})\), NP subset P/poly, or the stated target assumption to the source-carrier upper bound.
\item Control gate: compute \(\beta\) and q exactly, distinguishing all-exponent consequences from bounded target-exponent windows.
\end{itemize}

\paragraph{Stop criteria.}
\begin{itemize}
\item Every named carrier in the branch has \(\beta>0\) and yields only a bounded target-exponent exclusion.
\item The branch cannot state a source lower-bound model, target class, and transfer inequality in the same notation.
\item The proposed escape is only a black-box substitution restatement with no new carrier or nonblackbox mechanism.
\item The barrier audit shows the needed source lower bound is just as hard as the original P vs NP/P/poly target.
\end{itemize}

\paragraph{Do not stop if.}
\begin{itemize}
\item The branch has not yet computed \(\beta\), q, and the induced source upper-bound exponent under \(\mathrm{SAT}\in \mathrm{SIZE}(n^{c})\).
\item A subpolynomial-query or nonblackbox carrier remains available and has not been audited.
\item The latest failure sharpened the threshold-collapse lemma into a formal no-go statement.
\item The route can still isolate advice leakage, nonuniformity, naturalness, or algebrization as the exact obstruction.
\end{itemize}

\paragraph{Required skills.}
\begin{itemize}
\item contradiction-route-auditor
\item entailment-gap-auditor
\item kb-deduce
\item kb-end-to-end-proof-critic
\item kb-hypothesis-refutation-loop
\item long-horizon-route-portfolio-manager
\item main-problem-alignment-auditor
\item next-experiment-queue-ranker
\item paper-contribution-readiness-judge
\item pvsnp-barrier-specificity-reviewer
\item pvsnp-conceptual-pressure-localizer
\item pvsnp-impossibility-theorem-forger
\item pvsnp-inference-quality-controller
\item pvsnp-proof-derivation-builder
\item pvsnp-qualitative-complexity-reasoner
\item pvsnp-quantifier-uniformity-reviewer
\item pvsnp-reduction-consistency-reviewer
\item pvsnp-symbolic-complexity-auditor
\end{itemize}

\subsection*{\#124 cycle-0113: Hardness-Magnification Bridge Realization}

\paragraph{Analysis.} Run the finite proof-quality checklist on the paper statement: verify each \(TransparentRowSupport_{N}\) clause for interval-union, full scan, ROBDD, guarded-dDNNF, projected-dDNNF, and influence-complete examples; then attempt or refute \(RealSATInfluenceRealizer_{N}\) for one concrete semantic FormulaEval axiom.

\paragraph{Why this is a good target.} It directly tests the current black-box substitution cap by computing \(\beta\), q, and the induced target-exponent window. It keeps the barrier and model-boundary audit attached to the mathematical transfer rather than as after-the-fact commentary.

\paragraph{Risk.} Moderate risk: run only if source/target/model/barrier fields are explicit.

\paragraph{Branch control.} The item remains target-linked but needs tighter payoff criteria before consuming a long run.

\paragraph{Premature-stop guard.} Do not stop before deriving the transfer inequality from \(\mathrm{SAT}\in \mathrm{SIZE}(n^{c})\) to the carrier upper bound, identifying whether \(\beta=0\), \(\beta>0\), or the argument is genuinely nonblackbox, and recording the barrier status of the proposed escape.

\paragraph{Payoff criteria.}
\begin{itemize}
\item For the restricted-model branch, compute the compiler envelope from unrestricted SIZE(m\textasciicircum{}k) or P-time algorithms into the named model; promotion requires a bound above that envelope on the same SAT slice.
\item Falsify direct transfer if a P function already has the same restricted-model lower bound, such as PARITY for \(\mathrm{AC0}\), HWB for OBDD, or XOR/Tseitin for proof-system degree.
\item For this route, state the exact barrier premise it uses or escapes: relativizing, natural, algebrizing, nonuniform, or advice-leaking; promotion requires the escape to be attached to the specific transfer theorem.
\item Compute \(\beta\) and q exactly; promotion requires \(\beta=0\), subpolynomial target queries, or a genuinely nonblackbox use of the SAT circuit.
\item Record relativization, natural-proofs, algebrization, nonuniformity, and advice-leakage status for the specific transfer theorem.
\item Control gate: name the source carrier or language, target class, and circuit or resource model before claiming progress.
\end{itemize}

\paragraph{Stop criteria.}
\begin{itemize}
\item Every named carrier in the branch has \(\beta>0\) and yields only a bounded target-exponent exclusion.
\item The branch cannot state a source lower-bound model, target class, and transfer inequality in the same notation.
\item The proposed escape is only a black-box substitution restatement with no new carrier or nonblackbox mechanism.
\item The barrier audit shows the needed source lower bound is just as hard as the original P vs NP/P/poly target.
\end{itemize}

\paragraph{Do not stop if.}
\begin{itemize}
\item The branch has not yet computed \(\beta\), q, and the induced source upper-bound exponent under \(\mathrm{SAT}\in \mathrm{SIZE}(n^{c})\).
\item A subpolynomial-query or nonblackbox carrier remains available and has not been audited.
\item The latest failure sharpened the threshold-collapse lemma into a formal no-go statement.
\item The route can still isolate advice leakage, nonuniformity, naturalness, or algebrization as the exact obstruction.
\end{itemize}

\paragraph{Required skills.}
\begin{itemize}
\item contradiction-route-auditor
\item entailment-gap-auditor
\item kb-deduce
\item kb-end-to-end-proof-critic
\item kb-hypothesis-refutation-loop
\item long-horizon-route-portfolio-manager
\item main-problem-alignment-auditor
\item next-experiment-queue-ranker
\item paper-contribution-readiness-judge
\item pvsnp-barrier-specificity-reviewer
\item pvsnp-conceptual-pressure-localizer
\item pvsnp-impossibility-theorem-forger
\item pvsnp-inference-quality-controller
\item pvsnp-proof-derivation-builder
\item pvsnp-qualitative-complexity-reasoner
\item pvsnp-quantifier-uniformity-reviewer
\item pvsnp-reduction-consistency-reviewer
\item pvsnp-symbolic-complexity-auditor
\end{itemize}

\subsection*{\#125 cycle-0112: Hardness-Magnification Bridge Realization}

\paragraph{Analysis.} Write the corollary table as examples only: interval-union, full scan, ROBDD, guarded-dDNNF, projected-dDNNF, and influence-complete extension rules must each be checked against the five \(TransparentRowSupport_{N}\) clauses or demoted.

\paragraph{Why this is a good target.} It directly tests the current black-box substitution cap by computing \(\beta\), q, and the induced target-exponent window. It keeps the barrier and model-boundary audit attached to the mathematical transfer rather than as after-the-fact commentary.

\paragraph{Risk.} Moderate risk: run only if source/target/model/barrier fields are explicit.

\paragraph{Branch control.} The item remains target-linked but needs tighter payoff criteria before consuming a long run.

\paragraph{Premature-stop guard.} Do not stop before deriving the transfer inequality from \(\mathrm{SAT}\in \mathrm{SIZE}(n^{c})\) to the carrier upper bound, identifying whether \(\beta=0\), \(\beta>0\), or the argument is genuinely nonblackbox, and recording the barrier status of the proposed escape.

\paragraph{Payoff criteria.}
\begin{itemize}
\item For the restricted-model branch, compute the compiler envelope from unrestricted SIZE(m\textasciicircum{}k) or P-time algorithms into the named model; promotion requires a bound above that envelope on the same SAT slice.
\item Falsify direct transfer if a P function already has the same restricted-model lower bound, such as PARITY for \(\mathrm{AC0}\), HWB for OBDD, or XOR/Tseitin for proof-system degree.
\item For this route, state the exact barrier premise it uses or escapes: relativizing, natural, algebrizing, nonuniform, or advice-leaking; promotion requires the escape to be attached to the specific transfer theorem.
\item Compute \(\beta\) and q exactly; promotion requires \(\beta=0\), subpolynomial target queries, or a genuinely nonblackbox use of the SAT circuit.
\item Record relativization, natural-proofs, algebrization, nonuniformity, and advice-leakage status for the specific transfer theorem.
\item Control gate: name the source carrier or language, target class, and circuit or resource model before claiming progress.
\end{itemize}

\paragraph{Stop criteria.}
\begin{itemize}
\item Every named carrier in the branch has \(\beta>0\) and yields only a bounded target-exponent exclusion.
\item The branch cannot state a source lower-bound model, target class, and transfer inequality in the same notation.
\item The proposed escape is only a black-box substitution restatement with no new carrier or nonblackbox mechanism.
\item The barrier audit shows the needed source lower bound is just as hard as the original P vs NP/P/poly target.
\end{itemize}

\paragraph{Do not stop if.}
\begin{itemize}
\item The branch has not yet computed \(\beta\), q, and the induced source upper-bound exponent under \(\mathrm{SAT}\in \mathrm{SIZE}(n^{c})\).
\item A subpolynomial-query or nonblackbox carrier remains available and has not been audited.
\item The latest failure sharpened the threshold-collapse lemma into a formal no-go statement.
\item The route can still isolate advice leakage, nonuniformity, naturalness, or algebrization as the exact obstruction.
\end{itemize}

\paragraph{Required skills.}
\begin{itemize}
\item contradiction-route-auditor
\item entailment-gap-auditor
\item kb-deduce
\item kb-end-to-end-proof-critic
\item kb-hypothesis-refutation-loop
\item long-horizon-route-portfolio-manager
\item main-problem-alignment-auditor
\item next-experiment-queue-ranker
\item paper-contribution-readiness-judge
\item pvsnp-barrier-specificity-reviewer
\item pvsnp-conceptual-pressure-localizer
\item pvsnp-impossibility-theorem-forger
\item pvsnp-inference-quality-controller
\item pvsnp-proof-derivation-builder
\item pvsnp-qualitative-complexity-reasoner
\item pvsnp-quantifier-uniformity-reviewer
\item pvsnp-reduction-consistency-reviewer
\item pvsnp-symbolic-complexity-auditor
\end{itemize}

\subsection*{\#126 cycle-0082: Hardness-Magnification Bridge Realization}

\paragraph{Analysis.} Construct the smallest EF extension-rule countermodel that hides row parity without a row-normal node, or prove a normalization that exposes such parity with subthreshold overhead.

\paragraph{Why this is a good target.} It directly tests the current black-box substitution cap by computing \(\beta\), q, and the induced target-exponent window. It keeps the barrier and model-boundary audit attached to the mathematical transfer rather than as after-the-fact commentary.

\paragraph{Risk.} Moderate risk: run only if source/target/model/barrier fields are explicit.

\paragraph{Branch control.} The item remains target-linked but needs tighter payoff criteria before consuming a long run.

\paragraph{Premature-stop guard.} Do not stop before deriving the transfer inequality from \(\mathrm{SAT}\in \mathrm{SIZE}(n^{c})\) to the carrier upper bound, identifying whether \(\beta=0\), \(\beta>0\), or the argument is genuinely nonblackbox, and recording the barrier status of the proposed escape.

\paragraph{Payoff criteria.}
\begin{itemize}
\item For this route, state the exact barrier premise it uses or escapes: relativizing, natural, algebrizing, nonuniform, or advice-leaking; promotion requires the escape to be attached to the specific transfer theorem.
\item Compute \(\beta\) and q exactly; promotion requires \(\beta=0\), subpolynomial target queries, or a genuinely nonblackbox use of the SAT circuit.
\item Record relativization, natural-proofs, algebrization, nonuniformity, and advice-leakage status for the specific transfer theorem.
\item Control gate: name the source carrier or language, target class, and circuit or resource model before claiming progress.
\item Control gate: derive the exact transfer inequality from \(\mathrm{SAT}\in \mathrm{SIZE}(n^{c})\), NP subset P/poly, or the stated target assumption to the source-carrier upper bound.
\item Control gate: compute \(\beta\) and q exactly, distinguishing all-exponent consequences from bounded target-exponent windows.
\end{itemize}

\paragraph{Stop criteria.}
\begin{itemize}
\item Every named carrier in the branch has \(\beta>0\) and yields only a bounded target-exponent exclusion.
\item The branch cannot state a source lower-bound model, target class, and transfer inequality in the same notation.
\item The proposed escape is only a black-box substitution restatement with no new carrier or nonblackbox mechanism.
\item The barrier audit shows the needed source lower bound is just as hard as the original P vs NP/P/poly target.
\end{itemize}

\paragraph{Do not stop if.}
\begin{itemize}
\item The branch has not yet computed \(\beta\), q, and the induced source upper-bound exponent under \(\mathrm{SAT}\in \mathrm{SIZE}(n^{c})\).
\item A subpolynomial-query or nonblackbox carrier remains available and has not been audited.
\item The latest failure sharpened the threshold-collapse lemma into a formal no-go statement.
\item The route can still isolate advice leakage, nonuniformity, naturalness, or algebrization as the exact obstruction.
\end{itemize}

\paragraph{Required skills.}
\begin{itemize}
\item contradiction-route-auditor
\item entailment-gap-auditor
\item kb-deduce
\item kb-end-to-end-proof-critic
\item kb-hypothesis-refutation-loop
\item long-horizon-route-portfolio-manager
\item main-problem-alignment-auditor
\item next-experiment-queue-ranker
\item paper-contribution-readiness-judge
\item pvsnp-barrier-specificity-reviewer
\item pvsnp-conceptual-pressure-localizer
\item pvsnp-impossibility-theorem-forger
\item pvsnp-inference-quality-controller
\item pvsnp-proof-derivation-builder
\item pvsnp-qualitative-complexity-reasoner
\item pvsnp-quantifier-uniformity-reviewer
\item pvsnp-reduction-consistency-reviewer
\item pvsnp-symbolic-complexity-auditor
\end{itemize}

\subsection*{\#127 cycle-0059: Hardness-Magnification Bridge Realization: XOR carrier sparse-sketch audit}

\paragraph{Analysis.} Formalize the local-linear incidence transcript class for the XOR SAT carrier and compute the minimum sparse \(F_{2}\) support needed under a random-generator model, while auditing whether any extraction from arbitrary circuits to that class can avoid the planted-control relativization falsifier.

\paragraph{Why this is a good target.} It checks whether padding/projection transfers really magnify hardness or merely move the proof debt into sparse lower bounds. It keeps the barrier and model-boundary audit attached to the mathematical transfer rather than as after-the-fact commentary. It is executable as a bounded complexity calculation, CAS-backed exponent audit, or formal transfer check.

\paragraph{Risk.} Moderate risk: run only if source/target/model/barrier fields are explicit.

\paragraph{Branch control.} The item remains target-linked but needs tighter payoff criteria before consuming a long run.

\paragraph{Premature-stop guard.} Do not stop before deriving the exponent transfer c \(\to \) c/k or its growth-padding analogue, and before stating whether the result excludes all polynomial SAT exponents or only a bounded exponent interval.

\paragraph{Payoff criteria.}
\begin{itemize}
\item For the proof-complexity branch, separate fixed-family proof hardness from language hardness by naming the proof system, shortcut algorithm, and semantic lower-bound target.
\item Promotion requires a proof-system lower bound that survives the Gaussian-elimination or recognition shortcut; otherwise the result is a restricted-proof-system separation only.
\item For this route, state the exact barrier premise it uses or escapes: relativizing, natural, algebrizing, nonuniform, or advice-leaking; promotion requires the escape to be attached to the specific transfer theorem.
\item State whether the transfer excludes every polynomial SAT exponent or only a bounded interval \(c<c0\).
\item Record relativization, natural-proofs, algebrization, nonuniformity, and advice-leakage status for the specific transfer theorem.
\item Control gate: name the source carrier or language, target class, and circuit or resource model before claiming progress.
\end{itemize}

\paragraph{Stop criteria.}
\begin{itemize}
\item The branch only proves a finite target-exponent lower bound and cannot scale to \(\mathrm{NP}\not\subset \mathrm{P}/\mathrm{poly}\).
\item The route requires sparse valid-length lower bounds that are not weaker than the original target problem.
\item The projection carrier has no nonprojection mechanism and no subpolynomial embedding repair.
\end{itemize}

\paragraph{Do not stop if.}
\begin{itemize}
\item The quantifier pattern has not been repaired between fixed-k, unbounded-k, and single growth-padding carriers.
\item The projection-sandwich inequality has not yet been written as an explicit circuit-size transfer.
\item The sparse valid-length lower-bound obligation has not been isolated as the real source debt.
\end{itemize}

\paragraph{Required skills.}
\begin{itemize}
\item abducer-failure-signal-cartographer
\item entailment-gap-auditor
\item kb-end-to-end-proof-critic
\item kb-induce
\item long-horizon-route-portfolio-manager
\item main-problem-alignment-auditor
\item next-experiment-queue-ranker
\item paper-contribution-readiness-judge
\item pvsnp-barrier-specificity-reviewer
\item pvsnp-conceptual-pressure-localizer
\item pvsnp-impossibility-theorem-forger
\item pvsnp-inference-quality-controller
\item pvsnp-qualitative-complexity-reasoner
\item pvsnp-quantifier-uniformity-reviewer
\item pvsnp-reduction-consistency-reviewer
\item pvsnp-symbolic-complexity-auditor
\end{itemize}

\subsection*{\#128 cycle-0055: Hardness-Magnification Bridge Realization}

\paragraph{Analysis.} Compute QSk\_\{sim,N\}(E) for quotient relations generated by equality, variable renaming, clause permutation, tautological padding, and duplicate-pair involutions on the current high-q XOR carrier template; record the smallest residual exponent u and pairing-certificate cost.

\paragraph{Why this is a good target.} It checks whether padding/projection transfers really magnify hardness or merely move the proof debt into sparse lower bounds. It keeps the barrier and model-boundary audit attached to the mathematical transfer rather than as after-the-fact commentary. It is executable as a bounded complexity calculation, CAS-backed exponent audit, or formal transfer check.

\paragraph{Risk.} Moderate risk: run only if source/target/model/barrier fields are explicit.

\paragraph{Branch control.} The item remains target-linked but needs tighter payoff criteria before consuming a long run.

\paragraph{Premature-stop guard.} Do not stop before deriving the exponent transfer c \(\to \) c/k or its growth-padding analogue, and before stating whether the result excludes all polynomial SAT exponents or only a bounded exponent interval.

\paragraph{Payoff criteria.}
\begin{itemize}
\item For the high-q generated-OR carrier, compute the exact leakage exponent of the remaining witness-cover or unary-residual mechanism and compare it with 1+epsilon-q.
\item Promotion requires \(q+a>1+\epsilon\) after every projection, padding, and witness-cover leakage term is accounted for; otherwise the branch becomes a bounded-window or no-go result.
\item For the proof-complexity branch, separate fixed-family proof hardness from language hardness by naming the proof system, shortcut algorithm, and semantic lower-bound target.
\item Promotion requires a proof-system lower bound that survives the Gaussian-elimination or recognition shortcut; otherwise the result is a restricted-proof-system separation only.
\item For this route, state the exact barrier premise it uses or escapes: relativizing, natural, algebrizing, nonuniform, or advice-leaking; promotion requires the escape to be attached to the specific transfer theorem.
\item State whether the transfer excludes every polynomial SAT exponent or only a bounded interval \(c<c0\).
\end{itemize}

\paragraph{Stop criteria.}
\begin{itemize}
\item The branch only proves a finite target-exponent lower bound and cannot scale to \(\mathrm{NP}\not\subset \mathrm{P}/\mathrm{poly}\).
\item The route requires sparse valid-length lower bounds that are not weaker than the original target problem.
\item The projection carrier has no nonprojection mechanism and no subpolynomial embedding repair.
\end{itemize}

\paragraph{Do not stop if.}
\begin{itemize}
\item The quantifier pattern has not been repaired between fixed-k, unbounded-k, and single growth-padding carriers.
\item The projection-sandwich inequality has not yet been written as an explicit circuit-size transfer.
\item The sparse valid-length lower-bound obligation has not been isolated as the real source debt.
\end{itemize}

\paragraph{Required skills.}
\begin{itemize}
\item abducer-failure-signal-cartographer
\item entailment-gap-auditor
\item kb-end-to-end-proof-critic
\item kb-induce
\item long-horizon-route-portfolio-manager
\item main-problem-alignment-auditor
\item next-experiment-queue-ranker
\item paper-contribution-readiness-judge
\item pvsnp-barrier-specificity-reviewer
\item pvsnp-conceptual-pressure-localizer
\item pvsnp-impossibility-theorem-forger
\item pvsnp-inference-quality-controller
\item pvsnp-qualitative-complexity-reasoner
\item pvsnp-quantifier-uniformity-reviewer
\item pvsnp-reduction-consistency-reviewer
\item pvsnp-symbolic-complexity-auditor
\end{itemize}

\subsection*{\#129 cycle-0051: Hardness-Magnification Bridge Realization - cycle 51 direct-sum demagnification audit}

\paragraph{Analysis.} Try to formalize an aggregate lower-bound normal form for \(H_{N}\) that distinguishes global OR computation from \(M_{N}\) independent \(\mathrm{SAT}_{L}\) copies, and test it against the planted-control oracle before using separator hardness.

\paragraph{Why this is a good target.} It checks whether padding/projection transfers really magnify hardness or merely move the proof debt into sparse lower bounds. It keeps the barrier and model-boundary audit attached to the mathematical transfer rather than as after-the-fact commentary. It is executable as a bounded complexity calculation, CAS-backed exponent audit, or formal transfer check.

\paragraph{Risk.} Moderate risk: run only if source/target/model/barrier fields are explicit.

\paragraph{Branch control.} The item remains target-linked but needs tighter payoff criteria before consuming a long run.

\paragraph{Premature-stop guard.} Do not stop before deriving the exponent transfer c \(\to \) c/k or its growth-padding analogue, and before stating whether the result excludes all polynomial SAT exponents or only a bounded exponent interval.

\paragraph{Payoff criteria.}
\begin{itemize}
\item For this route, state the exact barrier premise it uses or escapes: relativizing, natural, algebrizing, nonuniform, or advice-leaking; promotion requires the escape to be attached to the specific transfer theorem.
\item State whether the transfer excludes every polynomial SAT exponent or only a bounded interval \(c<c0\).
\item Record relativization, natural-proofs, algebrization, nonuniformity, and advice-leakage status for the specific transfer theorem.
\item Control gate: name the source carrier or language, target class, and circuit or resource model before claiming progress.
\item Control gate: derive the exact transfer inequality from \(\mathrm{SAT}\in \mathrm{SIZE}(n^{c})\), NP subset P/poly, or the stated target assumption to the source-carrier upper bound.
\item Control gate: record relativization, natural-proofs, algebrization, nonuniformity, and advice-leakage status for the specific theorem.
\end{itemize}

\paragraph{Stop criteria.}
\begin{itemize}
\item The branch only proves a finite target-exponent lower bound and cannot scale to \(\mathrm{NP}\not\subset \mathrm{P}/\mathrm{poly}\).
\item The route requires sparse valid-length lower bounds that are not weaker than the original target problem.
\item The projection carrier has no nonprojection mechanism and no subpolynomial embedding repair.
\end{itemize}

\paragraph{Do not stop if.}
\begin{itemize}
\item The quantifier pattern has not been repaired between fixed-k, unbounded-k, and single growth-padding carriers.
\item The projection-sandwich inequality has not yet been written as an explicit circuit-size transfer.
\item The sparse valid-length lower-bound obligation has not been isolated as the real source debt.
\end{itemize}

\paragraph{Required skills.}
\begin{itemize}
\item abducer-failure-signal-cartographer
\item entailment-gap-auditor
\item kb-end-to-end-proof-critic
\item kb-induce
\item long-horizon-route-portfolio-manager
\item main-problem-alignment-auditor
\item next-experiment-queue-ranker
\item paper-contribution-readiness-judge
\item pvsnp-barrier-specificity-reviewer
\item pvsnp-conceptual-pressure-localizer
\item pvsnp-impossibility-theorem-forger
\item pvsnp-inference-quality-controller
\item pvsnp-qualitative-complexity-reasoner
\item pvsnp-quantifier-uniformity-reviewer
\item pvsnp-reduction-consistency-reviewer
\item pvsnp-symbolic-complexity-auditor
\end{itemize}

\subsection*{\#130 cycle-0032: Generated-OR magnification: correction-residual separator gate}

\paragraph{Analysis.} Instantiate \(Phi_{N}\) for a tagged generated-OR carrier and compute whether simple tag predicates produce subthreshold residual circuits; then test an untagged high-entropy generator for the same residual leakage.

\paragraph{Why this is a good target.} It checks whether padding/projection transfers really magnify hardness or merely move the proof debt into sparse lower bounds. It keeps the barrier and model-boundary audit attached to the mathematical transfer rather than as after-the-fact commentary. It is executable as a bounded complexity calculation, CAS-backed exponent audit, or formal transfer check.

\paragraph{Risk.} Moderate risk: run only if source/target/model/barrier fields are explicit.

\paragraph{Branch control.} The item remains target-linked but needs tighter payoff criteria before consuming a long run.

\paragraph{Premature-stop guard.} Do not stop before deriving the exponent transfer c \(\to \) c/k or its growth-padding analogue, and before stating whether the result excludes all polynomial SAT exponents or only a bounded exponent interval.

\paragraph{Payoff criteria.}
\begin{itemize}
\item For the high-q generated-OR carrier, compute the exact leakage exponent of the remaining witness-cover or unary-residual mechanism and compare it with 1+epsilon-q.
\item Promotion requires \(q+a>1+\epsilon\) after every projection, padding, and witness-cover leakage term is accounted for; otherwise the branch becomes a bounded-window or no-go result.
\item For the proof-complexity branch, separate fixed-family proof hardness from language hardness by naming the proof system, shortcut algorithm, and semantic lower-bound target.
\item Promotion requires a proof-system lower bound that survives the Gaussian-elimination or recognition shortcut; otherwise the result is a restricted-proof-system separation only.
\item For this route, state the exact barrier premise it uses or escapes: relativizing, natural, algebrizing, nonuniform, or advice-leaking; promotion requires the escape to be attached to the specific transfer theorem.
\item State whether the transfer excludes every polynomial SAT exponent or only a bounded interval \(c<c0\).
\end{itemize}

\paragraph{Stop criteria.}
\begin{itemize}
\item The branch only proves a finite target-exponent lower bound and cannot scale to \(\mathrm{NP}\not\subset \mathrm{P}/\mathrm{poly}\).
\item The route requires sparse valid-length lower bounds that are not weaker than the original target problem.
\item The projection carrier has no nonprojection mechanism and no subpolynomial embedding repair.
\end{itemize}

\paragraph{Do not stop if.}
\begin{itemize}
\item The quantifier pattern has not been repaired between fixed-k, unbounded-k, and single growth-padding carriers.
\item The projection-sandwich inequality has not yet been written as an explicit circuit-size transfer.
\item The sparse valid-length lower-bound obligation has not been isolated as the real source debt.
\end{itemize}

\paragraph{Required skills.}
\begin{itemize}
\item abducer-failure-signal-cartographer
\item entailment-gap-auditor
\item kb-end-to-end-proof-critic
\item kb-induce
\item long-horizon-route-portfolio-manager
\item main-problem-alignment-auditor
\item next-experiment-queue-ranker
\item paper-contribution-readiness-judge
\item pvsnp-barrier-specificity-reviewer
\item pvsnp-conceptual-pressure-localizer
\item pvsnp-impossibility-theorem-forger
\item pvsnp-inference-quality-controller
\item pvsnp-qualitative-complexity-reasoner
\item pvsnp-quantifier-uniformity-reviewer
\item pvsnp-reduction-consistency-reviewer
\item pvsnp-symbolic-complexity-auditor
\end{itemize}

\subsection*{\#131 cycle-0152: Hardness-Magnification Bridge Realization: prefix-stable Qrep support closure}

\paragraph{Analysis.} Write the finite verifier table for PrefixStableQrep entries: fields entry\_id, append\_time, old\_row, rep\_row, rank, certificate\_root, dependency\_list, first\_use\_time, \(charge_{to,Gamma}\), and rejection reason; then run it on one EqSAT duplicate-child example and one delayed-certificate private-row counterexample.

\paragraph{Why this is a good target.} It keeps the barrier and model-boundary audit attached to the mathematical transfer rather than as after-the-fact commentary. It is executable as a bounded complexity calculation, CAS-backed exponent audit, or formal transfer check.

\paragraph{Risk.} High drift risk unless rewritten into a concrete complexity-theoretic transfer test.

\paragraph{Branch control.} The branch still touches a relevant object, but the next experiment must be narrowed before continuation.

\paragraph{Premature-stop guard.} Do not stop before the branch has produced a concrete source-to-target transfer theorem, an explicit obstruction lemma, or a barrier-aware reason why the route cannot reach the P vs NP/P/poly target.

\paragraph{Payoff criteria.}
\begin{itemize}
\item For this route, state the exact barrier premise it uses or escapes: relativizing, natural, algebrizing, nonuniform, or advice-leaking; promotion requires the escape to be attached to the specific transfer theorem.
\item Record relativization, natural-proofs, algebrization, nonuniformity, and advice-leakage status for the specific transfer theorem.
\item Control gate: name the source carrier or language, target class, and circuit or resource model before claiming progress.
\item Control gate: derive the exact transfer inequality from \(\mathrm{SAT}\in \mathrm{SIZE}(n^{c})\), NP subset P/poly, or the stated target assumption to the source-carrier upper bound.
\item Control gate: record relativization, natural-proofs, algebrization, nonuniformity, and advice-leakage status for the specific theorem.
\end{itemize}

\paragraph{Stop criteria.}
\begin{itemize}
\item The branch cannot identify a source language/carrier and target complexity class in the same model.
\item The route repeatedly proves only a proxy lower bound with no target-preserving transfer.
\item The remaining work is a slogan about hardness magnification rather than a computable transfer or obstruction.
\end{itemize}

\paragraph{Do not stop if.}
\begin{itemize}
\item A transfer exponent, quantifier repair, or barrier status remains uncomputed.
\item The branch can still name a target-preserving carrier and a falsifiable lower-bound threshold.
\item The failure can become a formal no-go lemma for a route family.
\end{itemize}

\paragraph{Required skills.}
\begin{itemize}
\item entailment-gap-auditor
\item kb-end-to-end-proof-critic
\item long-horizon-route-portfolio-manager
\item main-problem-alignment-auditor
\item next-experiment-queue-ranker
\item paper-contribution-readiness-judge
\item pvsnp-barrier-specificity-reviewer
\item pvsnp-conceptual-pressure-localizer
\item pvsnp-impossibility-theorem-forger
\item pvsnp-inference-quality-controller
\item pvsnp-qualitative-complexity-reasoner
\item pvsnp-quantifier-uniformity-reviewer
\item pvsnp-symbolic-complexity-auditor
\end{itemize}

\subsection*{\#132 cycle-0143: Hardness-Magnification Bridge Realization}

\paragraph{Analysis.} Instantiate the six-row QCTSR transcript, list row IDs, quotient certificate dependencies, rank labels, and show the first uncertified merge produces two agreeing visible transcripts with opposite \(H_{N}\) parity.

\paragraph{Why this is a good target.} It keeps the barrier and model-boundary audit attached to the mathematical transfer rather than as after-the-fact commentary. It is executable as a bounded complexity calculation, CAS-backed exponent audit, or formal transfer check.

\paragraph{Risk.} High drift risk unless rewritten into a concrete complexity-theoretic transfer test.

\paragraph{Branch control.} The branch still touches a relevant object, but the next experiment must be narrowed before continuation.

\paragraph{Premature-stop guard.} Do not stop before the branch has produced a concrete source-to-target transfer theorem, an explicit obstruction lemma, or a barrier-aware reason why the route cannot reach the P vs NP/P/poly target.

\paragraph{Payoff criteria.}
\begin{itemize}
\item For this route, state the exact barrier premise it uses or escapes: relativizing, natural, algebrizing, nonuniform, or advice-leaking; promotion requires the escape to be attached to the specific transfer theorem.
\item Record relativization, natural-proofs, algebrization, nonuniformity, and advice-leakage status for the specific transfer theorem.
\item Control gate: name the source carrier or language, target class, and circuit or resource model before claiming progress.
\item Control gate: derive the exact transfer inequality from \(\mathrm{SAT}\in \mathrm{SIZE}(n^{c})\), NP subset P/poly, or the stated target assumption to the source-carrier upper bound.
\item Control gate: record relativization, natural-proofs, algebrization, nonuniformity, and advice-leakage status for the specific theorem.
\end{itemize}

\paragraph{Stop criteria.}
\begin{itemize}
\item The branch cannot identify a source language/carrier and target complexity class in the same model.
\item The route repeatedly proves only a proxy lower bound with no target-preserving transfer.
\item The remaining work is a slogan about hardness magnification rather than a computable transfer or obstruction.
\end{itemize}

\paragraph{Do not stop if.}
\begin{itemize}
\item A transfer exponent, quantifier repair, or barrier status remains uncomputed.
\item The branch can still name a target-preserving carrier and a falsifiable lower-bound threshold.
\item The failure can become a formal no-go lemma for a route family.
\end{itemize}

\paragraph{Required skills.}
\begin{itemize}
\item entailment-gap-auditor
\item kb-end-to-end-proof-critic
\item long-horizon-route-portfolio-manager
\item main-problem-alignment-auditor
\item next-experiment-queue-ranker
\item paper-contribution-readiness-judge
\item pvsnp-barrier-specificity-reviewer
\item pvsnp-conceptual-pressure-localizer
\item pvsnp-impossibility-theorem-forger
\item pvsnp-inference-quality-controller
\item pvsnp-qualitative-complexity-reasoner
\item pvsnp-quantifier-uniformity-reviewer
\item pvsnp-symbolic-complexity-auditor
\end{itemize}

\subsection*{\#133 cycle-0141: Hardness-Magnification Bridge Realization}

\paragraph{Analysis.} Instantiate one two-merge CTSR proof toy model: one certified EqSAT quotient with charged forest and one uncertified merge. Write the local verifier fields, compute \(R_{Q}\) and \(C_{Q}\), and demonstrate transcript agreement plus parity disagreement for the uncertified merge.

\paragraph{Why this is a good target.} It keeps the barrier and model-boundary audit attached to the mathematical transfer rather than as after-the-fact commentary. It is executable as a bounded complexity calculation, CAS-backed exponent audit, or formal transfer check.

\paragraph{Risk.} High drift risk unless rewritten into a concrete complexity-theoretic transfer test.

\paragraph{Branch control.} The branch still touches a relevant object, but the next experiment must be narrowed before continuation.

\paragraph{Premature-stop guard.} Do not stop before the branch has produced a concrete source-to-target transfer theorem, an explicit obstruction lemma, or a barrier-aware reason why the route cannot reach the P vs NP/P/poly target.

\paragraph{Payoff criteria.}
\begin{itemize}
\item For this route, state the exact barrier premise it uses or escapes: relativizing, natural, algebrizing, nonuniform, or advice-leaking; promotion requires the escape to be attached to the specific transfer theorem.
\item Record relativization, natural-proofs, algebrization, nonuniformity, and advice-leakage status for the specific transfer theorem.
\item Control gate: name the source carrier or language, target class, and circuit or resource model before claiming progress.
\item Control gate: derive the exact transfer inequality from \(\mathrm{SAT}\in \mathrm{SIZE}(n^{c})\), NP subset P/poly, or the stated target assumption to the source-carrier upper bound.
\item Control gate: record relativization, natural-proofs, algebrization, nonuniformity, and advice-leakage status for the specific theorem.
\end{itemize}

\paragraph{Stop criteria.}
\begin{itemize}
\item The branch cannot identify a source language/carrier and target complexity class in the same model.
\item The route repeatedly proves only a proxy lower bound with no target-preserving transfer.
\item The remaining work is a slogan about hardness magnification rather than a computable transfer or obstruction.
\end{itemize}

\paragraph{Do not stop if.}
\begin{itemize}
\item A transfer exponent, quantifier repair, or barrier status remains uncomputed.
\item The branch can still name a target-preserving carrier and a falsifiable lower-bound threshold.
\item The failure can become a formal no-go lemma for a route family.
\end{itemize}

\paragraph{Required skills.}
\begin{itemize}
\item entailment-gap-auditor
\item kb-end-to-end-proof-critic
\item long-horizon-route-portfolio-manager
\item main-problem-alignment-auditor
\item next-experiment-queue-ranker
\item paper-contribution-readiness-judge
\item pvsnp-barrier-specificity-reviewer
\item pvsnp-conceptual-pressure-localizer
\item pvsnp-impossibility-theorem-forger
\item pvsnp-inference-quality-controller
\item pvsnp-qualitative-complexity-reasoner
\item pvsnp-quantifier-uniformity-reviewer
\item pvsnp-symbolic-complexity-auditor
\end{itemize}

\subsection*{\#134 cycle-0135: Hardness-Magnification Bridge Realization}

\paragraph{Analysis.} Resolve \(RealSATInfluenceRealizer_{N}^{arb}\) for one concrete semantic rule: bounded Tseitin self-reduction with adaptive unsat-branch pruning. Either produce a syntax-and-N computable influence row list \(R_{A}\) with charge in \(\Gamma_{real}\), or formalize the first unchargeable semantic inference as a private-row no-realizer theorem.

\paragraph{Why this is a good target.} It keeps the barrier and model-boundary audit attached to the mathematical transfer rather than as after-the-fact commentary. It is executable as a bounded complexity calculation, CAS-backed exponent audit, or formal transfer check.

\paragraph{Risk.} High drift risk unless rewritten into a concrete complexity-theoretic transfer test.

\paragraph{Branch control.} The branch still touches a relevant object, but the next experiment must be narrowed before continuation.

\paragraph{Premature-stop guard.} Do not stop before the branch has produced a concrete source-to-target transfer theorem, an explicit obstruction lemma, or a barrier-aware reason why the route cannot reach the P vs NP/P/poly target.

\paragraph{Payoff criteria.}
\begin{itemize}
\item For the proof-complexity branch, separate fixed-family proof hardness from language hardness by naming the proof system, shortcut algorithm, and semantic lower-bound target.
\item Promotion requires a proof-system lower bound that survives the Gaussian-elimination or recognition shortcut; otherwise the result is a restricted-proof-system separation only.
\item For this route, state the exact barrier premise it uses or escapes: relativizing, natural, algebrizing, nonuniform, or advice-leaking; promotion requires the escape to be attached to the specific transfer theorem.
\item Record relativization, natural-proofs, algebrization, nonuniformity, and advice-leakage status for the specific transfer theorem.
\item Control gate: name the source carrier or language, target class, and circuit or resource model before claiming progress.
\item Control gate: derive the exact transfer inequality from \(\mathrm{SAT}\in \mathrm{SIZE}(n^{c})\), NP subset P/poly, or the stated target assumption to the source-carrier upper bound.
\end{itemize}

\paragraph{Stop criteria.}
\begin{itemize}
\item The branch cannot identify a source language/carrier and target complexity class in the same model.
\item The route repeatedly proves only a proxy lower bound with no target-preserving transfer.
\item The remaining work is a slogan about hardness magnification rather than a computable transfer or obstruction.
\end{itemize}

\paragraph{Do not stop if.}
\begin{itemize}
\item A transfer exponent, quantifier repair, or barrier status remains uncomputed.
\item The branch can still name a target-preserving carrier and a falsifiable lower-bound threshold.
\item The failure can become a formal no-go lemma for a route family.
\end{itemize}

\paragraph{Required skills.}
\begin{itemize}
\item entailment-gap-auditor
\item kb-end-to-end-proof-critic
\item long-horizon-route-portfolio-manager
\item main-problem-alignment-auditor
\item next-experiment-queue-ranker
\item paper-contribution-readiness-judge
\item pvsnp-barrier-specificity-reviewer
\item pvsnp-conceptual-pressure-localizer
\item pvsnp-impossibility-theorem-forger
\item pvsnp-inference-quality-controller
\item pvsnp-qualitative-complexity-reasoner
\item pvsnp-quantifier-uniformity-reviewer
\item pvsnp-symbolic-complexity-auditor
\end{itemize}

\subsection*{\#135 cycle-0126: Hardness-Magnification Bridge Realization}

\paragraph{Analysis.} Formalize the \(\mathrm{Psys}_{sym}^{dag+declTSR(h)}\) grammar in one page: exact row-index fields, TSR parent/child encodings, local verifier predicate, \(K_{dag}^\mathrm{TSR}\) fixed-point pseudocode, and the declared TSR line-induction case; then instantiate with \(\kappa=\epsilon/3\), \(\theta=\epsilon/3\), q in ((1+epsilon)/2,1+epsilon), \(\rho<1+\epsilon\) to verify all charges remain subthreshold except the intended q-density pressure.

\paragraph{Why this is a good target.} It keeps the barrier and model-boundary audit attached to the mathematical transfer rather than as after-the-fact commentary. It is executable as a bounded complexity calculation, CAS-backed exponent audit, or formal transfer check.

\paragraph{Risk.} High drift risk unless rewritten into a concrete complexity-theoretic transfer test.

\paragraph{Branch control.} The branch still touches a relevant object, but the next experiment must be narrowed before continuation.

\paragraph{Premature-stop guard.} Do not stop before the branch has produced a concrete source-to-target transfer theorem, an explicit obstruction lemma, or a barrier-aware reason why the route cannot reach the P vs NP/P/poly target.

\paragraph{Payoff criteria.}
\begin{itemize}
\item For this route, state the exact barrier premise it uses or escapes: relativizing, natural, algebrizing, nonuniform, or advice-leaking; promotion requires the escape to be attached to the specific transfer theorem.
\item Record relativization, natural-proofs, algebrization, nonuniformity, and advice-leakage status for the specific transfer theorem.
\item Control gate: name the source carrier or language, target class, and circuit or resource model before claiming progress.
\item Control gate: derive the exact transfer inequality from \(\mathrm{SAT}\in \mathrm{SIZE}(n^{c})\), NP subset P/poly, or the stated target assumption to the source-carrier upper bound.
\item Control gate: record relativization, natural-proofs, algebrization, nonuniformity, and advice-leakage status for the specific theorem.
\end{itemize}

\paragraph{Stop criteria.}
\begin{itemize}
\item The branch cannot identify a source language/carrier and target complexity class in the same model.
\item The route repeatedly proves only a proxy lower bound with no target-preserving transfer.
\item The remaining work is a slogan about hardness magnification rather than a computable transfer or obstruction.
\end{itemize}

\paragraph{Do not stop if.}
\begin{itemize}
\item A transfer exponent, quantifier repair, or barrier status remains uncomputed.
\item The branch can still name a target-preserving carrier and a falsifiable lower-bound threshold.
\item The failure can become a formal no-go lemma for a route family.
\end{itemize}

\paragraph{Required skills.}
\begin{itemize}
\item entailment-gap-auditor
\item kb-end-to-end-proof-critic
\item long-horizon-route-portfolio-manager
\item main-problem-alignment-auditor
\item next-experiment-queue-ranker
\item paper-contribution-readiness-judge
\item pvsnp-barrier-specificity-reviewer
\item pvsnp-conceptual-pressure-localizer
\item pvsnp-impossibility-theorem-forger
\item pvsnp-inference-quality-controller
\item pvsnp-qualitative-complexity-reasoner
\item pvsnp-quantifier-uniformity-reviewer
\item pvsnp-symbolic-complexity-auditor
\end{itemize}

\subsection*{\#136 cycle-0120: Hardness-Magnification Bridge Realization: least-closed transparent FormulaEval support}

\paragraph{Analysis.} Write the proof-audit version of L120.1: formalize the fixed-point computation of K\_dag(pi), give pseudocode for the verifier-visible dependency closure, and check one projected-dDNNF FormulaEval example plus one unread-enable-bit counterexample against the \(\Gamma_{dag}\) inequality.

\paragraph{Why this is a good target.} It keeps the barrier and model-boundary audit attached to the mathematical transfer rather than as after-the-fact commentary. It is executable as a bounded complexity calculation, CAS-backed exponent audit, or formal transfer check.

\paragraph{Risk.} High drift risk unless rewritten into a concrete complexity-theoretic transfer test.

\paragraph{Branch control.} The branch still touches a relevant object, but the next experiment must be narrowed before continuation.

\paragraph{Premature-stop guard.} Do not stop before the branch has produced a concrete source-to-target transfer theorem, an explicit obstruction lemma, or a barrier-aware reason why the route cannot reach the P vs NP/P/poly target.

\paragraph{Payoff criteria.}
\begin{itemize}
\item For this route, state the exact barrier premise it uses or escapes: relativizing, natural, algebrizing, nonuniform, or advice-leaking; promotion requires the escape to be attached to the specific transfer theorem.
\item Record relativization, natural-proofs, algebrization, nonuniformity, and advice-leakage status for the specific transfer theorem.
\item Control gate: name the source carrier or language, target class, and circuit or resource model before claiming progress.
\item Control gate: derive the exact transfer inequality from \(\mathrm{SAT}\in \mathrm{SIZE}(n^{c})\), NP subset P/poly, or the stated target assumption to the source-carrier upper bound.
\item Control gate: record relativization, natural-proofs, algebrization, nonuniformity, and advice-leakage status for the specific theorem.
\end{itemize}

\paragraph{Stop criteria.}
\begin{itemize}
\item The branch cannot identify a source language/carrier and target complexity class in the same model.
\item The route repeatedly proves only a proxy lower bound with no target-preserving transfer.
\item The remaining work is a slogan about hardness magnification rather than a computable transfer or obstruction.
\end{itemize}

\paragraph{Do not stop if.}
\begin{itemize}
\item A transfer exponent, quantifier repair, or barrier status remains uncomputed.
\item The branch can still name a target-preserving carrier and a falsifiable lower-bound threshold.
\item The failure can become a formal no-go lemma for a route family.
\end{itemize}

\paragraph{Required skills.}
\begin{itemize}
\item entailment-gap-auditor
\item kb-end-to-end-proof-critic
\item long-horizon-route-portfolio-manager
\item main-problem-alignment-auditor
\item next-experiment-queue-ranker
\item paper-contribution-readiness-judge
\item pvsnp-barrier-specificity-reviewer
\item pvsnp-conceptual-pressure-localizer
\item pvsnp-impossibility-theorem-forger
\item pvsnp-inference-quality-controller
\item pvsnp-qualitative-complexity-reasoner
\item pvsnp-quantifier-uniformity-reviewer
\item pvsnp-symbolic-complexity-auditor
\end{itemize}

\subsection*{\#137 cycle-0118: Hardness-Magnification Bridge Realization}

\paragraph{Analysis.} Run a finite proof audit on \(\mathrm{Psys}_{sym}^{ext}\): instantiate one guarded FormulaEval line, compute \(I_{l}\), \(K_{idx}\), \(R_{FE}\), and \(\Gamma_{comp}\), then verify the last-unsafe-line two-world construction or identify the first real SAT-semantic rule that forces a nonrelativizing realizer charge.

\paragraph{Why this is a good target.} It keeps the barrier and model-boundary audit attached to the mathematical transfer rather than as after-the-fact commentary. It is executable as a bounded complexity calculation, CAS-backed exponent audit, or formal transfer check.

\paragraph{Risk.} High drift risk unless rewritten into a concrete complexity-theoretic transfer test.

\paragraph{Branch control.} The branch still touches a relevant object, but the next experiment must be narrowed before continuation.

\paragraph{Premature-stop guard.} Do not stop before the branch has produced a concrete source-to-target transfer theorem, an explicit obstruction lemma, or a barrier-aware reason why the route cannot reach the P vs NP/P/poly target.

\paragraph{Payoff criteria.}
\begin{itemize}
\item For this route, state the exact barrier premise it uses or escapes: relativizing, natural, algebrizing, nonuniform, or advice-leaking; promotion requires the escape to be attached to the specific transfer theorem.
\item Record relativization, natural-proofs, algebrization, nonuniformity, and advice-leakage status for the specific transfer theorem.
\item Control gate: name the source carrier or language, target class, and circuit or resource model before claiming progress.
\item Control gate: derive the exact transfer inequality from \(\mathrm{SAT}\in \mathrm{SIZE}(n^{c})\), NP subset P/poly, or the stated target assumption to the source-carrier upper bound.
\item Control gate: record relativization, natural-proofs, algebrization, nonuniformity, and advice-leakage status for the specific theorem.
\end{itemize}

\paragraph{Stop criteria.}
\begin{itemize}
\item The branch cannot identify a source language/carrier and target complexity class in the same model.
\item The route repeatedly proves only a proxy lower bound with no target-preserving transfer.
\item The remaining work is a slogan about hardness magnification rather than a computable transfer or obstruction.
\end{itemize}

\paragraph{Do not stop if.}
\begin{itemize}
\item A transfer exponent, quantifier repair, or barrier status remains uncomputed.
\item The branch can still name a target-preserving carrier and a falsifiable lower-bound threshold.
\item The failure can become a formal no-go lemma for a route family.
\end{itemize}

\paragraph{Required skills.}
\begin{itemize}
\item entailment-gap-auditor
\item kb-end-to-end-proof-critic
\item long-horizon-route-portfolio-manager
\item main-problem-alignment-auditor
\item next-experiment-queue-ranker
\item paper-contribution-readiness-judge
\item pvsnp-barrier-specificity-reviewer
\item pvsnp-conceptual-pressure-localizer
\item pvsnp-impossibility-theorem-forger
\item pvsnp-inference-quality-controller
\item pvsnp-qualitative-complexity-reasoner
\item pvsnp-quantifier-uniformity-reviewer
\item pvsnp-symbolic-complexity-auditor
\end{itemize}

\subsection*{\#138 cycle-0116: Hardness-Magnification Bridge Realization: Transparent FormulaEval Minimality Audit}

\paragraph{Analysis.} Instantiate \(RealSATInfluenceRealizer_{N}\) or refute it for one concrete strong rule: FormulaEval with Tseitin-extension equality plus bounded self-reduction; list exact bit fields, RowSet extraction, same-theory expansion certificates, and the first semantic inference that cannot be charged.

\paragraph{Why this is a good target.} It keeps the barrier and model-boundary audit attached to the mathematical transfer rather than as after-the-fact commentary. It is executable as a bounded complexity calculation, CAS-backed exponent audit, or formal transfer check.

\paragraph{Risk.} High drift risk unless rewritten into a concrete complexity-theoretic transfer test.

\paragraph{Branch control.} The branch still touches a relevant object, but the next experiment must be narrowed before continuation.

\paragraph{Premature-stop guard.} Do not stop before the branch has produced a concrete source-to-target transfer theorem, an explicit obstruction lemma, or a barrier-aware reason why the route cannot reach the P vs NP/P/poly target.

\paragraph{Payoff criteria.}
\begin{itemize}
\item For the proof-complexity branch, separate fixed-family proof hardness from language hardness by naming the proof system, shortcut algorithm, and semantic lower-bound target.
\item Promotion requires a proof-system lower bound that survives the Gaussian-elimination or recognition shortcut; otherwise the result is a restricted-proof-system separation only.
\item For this route, state the exact barrier premise it uses or escapes: relativizing, natural, algebrizing, nonuniform, or advice-leaking; promotion requires the escape to be attached to the specific transfer theorem.
\item Record relativization, natural-proofs, algebrization, nonuniformity, and advice-leakage status for the specific transfer theorem.
\item Control gate: name the source carrier or language, target class, and circuit or resource model before claiming progress.
\item Control gate: derive the exact transfer inequality from \(\mathrm{SAT}\in \mathrm{SIZE}(n^{c})\), NP subset P/poly, or the stated target assumption to the source-carrier upper bound.
\end{itemize}

\paragraph{Stop criteria.}
\begin{itemize}
\item The branch cannot identify a source language/carrier and target complexity class in the same model.
\item The route repeatedly proves only a proxy lower bound with no target-preserving transfer.
\item The remaining work is a slogan about hardness magnification rather than a computable transfer or obstruction.
\end{itemize}

\paragraph{Do not stop if.}
\begin{itemize}
\item A transfer exponent, quantifier repair, or barrier status remains uncomputed.
\item The branch can still name a target-preserving carrier and a falsifiable lower-bound threshold.
\item The failure can become a formal no-go lemma for a route family.
\end{itemize}

\paragraph{Required skills.}
\begin{itemize}
\item entailment-gap-auditor
\item kb-end-to-end-proof-critic
\item long-horizon-route-portfolio-manager
\item main-problem-alignment-auditor
\item next-experiment-queue-ranker
\item paper-contribution-readiness-judge
\item pvsnp-barrier-specificity-reviewer
\item pvsnp-conceptual-pressure-localizer
\item pvsnp-impossibility-theorem-forger
\item pvsnp-inference-quality-controller
\item pvsnp-qualitative-complexity-reasoner
\item pvsnp-quantifier-uniformity-reviewer
\item pvsnp-symbolic-complexity-auditor
\end{itemize}

\subsection*{\#139 cycle-0114: Hardness-Magnification Bridge Realization}

\paragraph{Analysis.} Instantiate \(RealSATInfluenceRealizer_{N}\) for one concrete full semantic FormulaEval rule, or prove a rule-specific private-row countermodel showing that the full \(\mathrm{Psys}_{sym}\) extension must remain quarantined.

\paragraph{Why this is a good target.} It keeps the barrier and model-boundary audit attached to the mathematical transfer rather than as after-the-fact commentary. It is executable as a bounded complexity calculation, CAS-backed exponent audit, or formal transfer check.

\paragraph{Risk.} High drift risk unless rewritten into a concrete complexity-theoretic transfer test.

\paragraph{Branch control.} The branch still touches a relevant object, but the next experiment must be narrowed before continuation.

\paragraph{Premature-stop guard.} Do not stop before the branch has produced a concrete source-to-target transfer theorem, an explicit obstruction lemma, or a barrier-aware reason why the route cannot reach the P vs NP/P/poly target.

\paragraph{Payoff criteria.}
\begin{itemize}
\item For this route, state the exact barrier premise it uses or escapes: relativizing, natural, algebrizing, nonuniform, or advice-leaking; promotion requires the escape to be attached to the specific transfer theorem.
\item Record relativization, natural-proofs, algebrization, nonuniformity, and advice-leakage status for the specific transfer theorem.
\item Control gate: name the source carrier or language, target class, and circuit or resource model before claiming progress.
\item Control gate: derive the exact transfer inequality from \(\mathrm{SAT}\in \mathrm{SIZE}(n^{c})\), NP subset P/poly, or the stated target assumption to the source-carrier upper bound.
\item Control gate: record relativization, natural-proofs, algebrization, nonuniformity, and advice-leakage status for the specific theorem.
\end{itemize}

\paragraph{Stop criteria.}
\begin{itemize}
\item The branch cannot identify a source language/carrier and target complexity class in the same model.
\item The route repeatedly proves only a proxy lower bound with no target-preserving transfer.
\item The remaining work is a slogan about hardness magnification rather than a computable transfer or obstruction.
\end{itemize}

\paragraph{Do not stop if.}
\begin{itemize}
\item A transfer exponent, quantifier repair, or barrier status remains uncomputed.
\item The branch can still name a target-preserving carrier and a falsifiable lower-bound threshold.
\item The failure can become a formal no-go lemma for a route family.
\end{itemize}

\paragraph{Required skills.}
\begin{itemize}
\item entailment-gap-auditor
\item kb-end-to-end-proof-critic
\item long-horizon-route-portfolio-manager
\item main-problem-alignment-auditor
\item next-experiment-queue-ranker
\item paper-contribution-readiness-judge
\item pvsnp-barrier-specificity-reviewer
\item pvsnp-conceptual-pressure-localizer
\item pvsnp-impossibility-theorem-forger
\item pvsnp-inference-quality-controller
\item pvsnp-qualitative-complexity-reasoner
\item pvsnp-quantifier-uniformity-reviewer
\item pvsnp-symbolic-complexity-auditor
\end{itemize}

\subsection*{\#140 cycle-0111: Hardness-Magnification Bridge Realization}

\paragraph{Analysis.} Write the final paper theorem text for \(TransparentRowSupport_{N}\) with clauses T1-T6, prove the line induction lemma in referee style, then test one genuinely semantic FormulaEval rule against \(RealSATInfluenceRealizer_{N}\) or instantiate the private-row countermodel.

\paragraph{Why this is a good target.} It keeps the barrier and model-boundary audit attached to the mathematical transfer rather than as after-the-fact commentary. It is executable as a bounded complexity calculation, CAS-backed exponent audit, or formal transfer check.

\paragraph{Risk.} High drift risk unless rewritten into a concrete complexity-theoretic transfer test.

\paragraph{Branch control.} The branch still touches a relevant object, but the next experiment must be narrowed before continuation.

\paragraph{Premature-stop guard.} Do not stop before the branch has produced a concrete source-to-target transfer theorem, an explicit obstruction lemma, or a barrier-aware reason why the route cannot reach the P vs NP/P/poly target.

\paragraph{Payoff criteria.}
\begin{itemize}
\item For this route, state the exact barrier premise it uses or escapes: relativizing, natural, algebrizing, nonuniform, or advice-leaking; promotion requires the escape to be attached to the specific transfer theorem.
\item Record relativization, natural-proofs, algebrization, nonuniformity, and advice-leakage status for the specific transfer theorem.
\item Control gate: name the source carrier or language, target class, and circuit or resource model before claiming progress.
\item Control gate: derive the exact transfer inequality from \(\mathrm{SAT}\in \mathrm{SIZE}(n^{c})\), NP subset P/poly, or the stated target assumption to the source-carrier upper bound.
\item Control gate: record relativization, natural-proofs, algebrization, nonuniformity, and advice-leakage status for the specific theorem.
\end{itemize}

\paragraph{Stop criteria.}
\begin{itemize}
\item The branch cannot identify a source language/carrier and target complexity class in the same model.
\item The route repeatedly proves only a proxy lower bound with no target-preserving transfer.
\item The remaining work is a slogan about hardness magnification rather than a computable transfer or obstruction.
\end{itemize}

\paragraph{Do not stop if.}
\begin{itemize}
\item A transfer exponent, quantifier repair, or barrier status remains uncomputed.
\item The branch can still name a target-preserving carrier and a falsifiable lower-bound threshold.
\item The failure can become a formal no-go lemma for a route family.
\end{itemize}

\paragraph{Required skills.}
\begin{itemize}
\item entailment-gap-auditor
\item kb-end-to-end-proof-critic
\item long-horizon-route-portfolio-manager
\item main-problem-alignment-auditor
\item next-experiment-queue-ranker
\item paper-contribution-readiness-judge
\item pvsnp-barrier-specificity-reviewer
\item pvsnp-conceptual-pressure-localizer
\item pvsnp-impossibility-theorem-forger
\item pvsnp-inference-quality-controller
\item pvsnp-qualitative-complexity-reasoner
\item pvsnp-quantifier-uniformity-reviewer
\item pvsnp-symbolic-complexity-auditor
\end{itemize}

\subsection*{\#141 cycle-0104: Hardness-Magnification Bridge Realization}

\paragraph{Analysis.} Instantiate \(RowSetEnumerator_{N}\) for one richer transparent selector class beyond interval-unions, preferably sorted \(\mathrm{AC0}\) threshold-free selectors with explicit scan charge, and decide whether it expands at N|RowSet|, collapses to an \(N^{q}\) scan bound, or admits a private-row countermodel.

\paragraph{Why this is a good target.} It keeps the barrier and model-boundary audit attached to the mathematical transfer rather than as after-the-fact commentary. It is executable as a bounded complexity calculation, CAS-backed exponent audit, or formal transfer check.

\paragraph{Risk.} High drift risk unless rewritten into a concrete complexity-theoretic transfer test.

\paragraph{Branch control.} The branch still touches a relevant object, but the next experiment must be narrowed before continuation.

\paragraph{Premature-stop guard.} Do not stop before the branch has produced a concrete source-to-target transfer theorem, an explicit obstruction lemma, or a barrier-aware reason why the route cannot reach the P vs NP/P/poly target.

\paragraph{Payoff criteria.}
\begin{itemize}
\item For the restricted-model branch, compute the compiler envelope from unrestricted SIZE(m\textasciicircum{}k) or P-time algorithms into the named model; promotion requires a bound above that envelope on the same SAT slice.
\item Falsify direct transfer if a P function already has the same restricted-model lower bound, such as PARITY for \(\mathrm{AC0}\), HWB for OBDD, or XOR/Tseitin for proof-system degree.
\item For this route, state the exact barrier premise it uses or escapes: relativizing, natural, algebrizing, nonuniform, or advice-leaking; promotion requires the escape to be attached to the specific transfer theorem.
\item Record relativization, natural-proofs, algebrization, nonuniformity, and advice-leakage status for the specific transfer theorem.
\item Control gate: name the source carrier or language, target class, and circuit or resource model before claiming progress.
\item Control gate: derive the exact transfer inequality from \(\mathrm{SAT}\in \mathrm{SIZE}(n^{c})\), NP subset P/poly, or the stated target assumption to the source-carrier upper bound.
\end{itemize}

\paragraph{Stop criteria.}
\begin{itemize}
\item The branch cannot identify a source language/carrier and target complexity class in the same model.
\item The route repeatedly proves only a proxy lower bound with no target-preserving transfer.
\item The remaining work is a slogan about hardness magnification rather than a computable transfer or obstruction.
\end{itemize}

\paragraph{Do not stop if.}
\begin{itemize}
\item A transfer exponent, quantifier repair, or barrier status remains uncomputed.
\item The branch can still name a target-preserving carrier and a falsifiable lower-bound threshold.
\item The failure can become a formal no-go lemma for a route family.
\end{itemize}

\paragraph{Required skills.}
\begin{itemize}
\item entailment-gap-auditor
\item kb-end-to-end-proof-critic
\item long-horizon-route-portfolio-manager
\item main-problem-alignment-auditor
\item next-experiment-queue-ranker
\item paper-contribution-readiness-judge
\item pvsnp-barrier-specificity-reviewer
\item pvsnp-conceptual-pressure-localizer
\item pvsnp-impossibility-theorem-forger
\item pvsnp-inference-quality-controller
\item pvsnp-qualitative-complexity-reasoner
\item pvsnp-quantifier-uniformity-reviewer
\item pvsnp-symbolic-complexity-auditor
\end{itemize}

\subsection*{\#142 cycle-0102: Hardness-Magnification Bridge Realization - cycle 102 \(SR_{h}\) normal form audit}

\paragraph{Analysis.} Instantiate \(RowSetEnumerator_{N}\) for one transparent \(\mathrm{Psys}_{sym}^{sel}\) fragment, preferably interval-union or \(\mathrm{AC0}\) row selectors, and decide whether FormulaEval-SAT is an expanded axiom schema with \(N^{1+q+o(1)}\) charge or a semantic oracle exposed by the private-row falsifier.

\paragraph{Why this is a good target.} It keeps the barrier and model-boundary audit attached to the mathematical transfer rather than as after-the-fact commentary. It is executable as a bounded complexity calculation, CAS-backed exponent audit, or formal transfer check.

\paragraph{Risk.} High drift risk unless rewritten into a concrete complexity-theoretic transfer test.

\paragraph{Branch control.} The branch still touches a relevant object, but the next experiment must be narrowed before continuation.

\paragraph{Premature-stop guard.} Do not stop before the branch has produced a concrete source-to-target transfer theorem, an explicit obstruction lemma, or a barrier-aware reason why the route cannot reach the P vs NP/P/poly target.

\paragraph{Payoff criteria.}
\begin{itemize}
\item For the restricted-model branch, compute the compiler envelope from unrestricted SIZE(m\textasciicircum{}k) or P-time algorithms into the named model; promotion requires a bound above that envelope on the same SAT slice.
\item Falsify direct transfer if a P function already has the same restricted-model lower bound, such as PARITY for \(\mathrm{AC0}\), HWB for OBDD, or XOR/Tseitin for proof-system degree.
\item For this route, state the exact barrier premise it uses or escapes: relativizing, natural, algebrizing, nonuniform, or advice-leaking; promotion requires the escape to be attached to the specific transfer theorem.
\item Record relativization, natural-proofs, algebrization, nonuniformity, and advice-leakage status for the specific transfer theorem.
\item Control gate: name the source carrier or language, target class, and circuit or resource model before claiming progress.
\item Control gate: derive the exact transfer inequality from \(\mathrm{SAT}\in \mathrm{SIZE}(n^{c})\), NP subset P/poly, or the stated target assumption to the source-carrier upper bound.
\end{itemize}

\paragraph{Stop criteria.}
\begin{itemize}
\item The branch cannot identify a source language/carrier and target complexity class in the same model.
\item The route repeatedly proves only a proxy lower bound with no target-preserving transfer.
\item The remaining work is a slogan about hardness magnification rather than a computable transfer or obstruction.
\end{itemize}

\paragraph{Do not stop if.}
\begin{itemize}
\item A transfer exponent, quantifier repair, or barrier status remains uncomputed.
\item The branch can still name a target-preserving carrier and a falsifiable lower-bound threshold.
\item The failure can become a formal no-go lemma for a route family.
\end{itemize}

\paragraph{Required skills.}
\begin{itemize}
\item entailment-gap-auditor
\item kb-end-to-end-proof-critic
\item long-horizon-route-portfolio-manager
\item main-problem-alignment-auditor
\item next-experiment-queue-ranker
\item paper-contribution-readiness-judge
\item pvsnp-barrier-specificity-reviewer
\item pvsnp-conceptual-pressure-localizer
\item pvsnp-impossibility-theorem-forger
\item pvsnp-inference-quality-controller
\item pvsnp-qualitative-complexity-reasoner
\item pvsnp-quantifier-uniformity-reviewer
\item pvsnp-symbolic-complexity-auditor
\end{itemize}

\subsection*{\#143 cycle-0100: Hardness-Magnification Bridge Realization}

\paragraph{Analysis.} Prove \(RowSetEnumerator_{N}\) for an interval-union or \(\mathrm{AC0}\) FormulaEval selector and compute alpha-charge; if enumeration fails, emit the minimal private-row countermodel.

\paragraph{Why this is a good target.} It keeps the barrier and model-boundary audit attached to the mathematical transfer rather than as after-the-fact commentary. It is executable as a bounded complexity calculation, CAS-backed exponent audit, or formal transfer check.

\paragraph{Risk.} High drift risk unless rewritten into a concrete complexity-theoretic transfer test.

\paragraph{Branch control.} The branch still touches a relevant object, but the next experiment must be narrowed before continuation.

\paragraph{Premature-stop guard.} Do not stop before the branch has produced a concrete source-to-target transfer theorem, an explicit obstruction lemma, or a barrier-aware reason why the route cannot reach the P vs NP/P/poly target.

\paragraph{Payoff criteria.}
\begin{itemize}
\item For the restricted-model branch, compute the compiler envelope from unrestricted SIZE(m\textasciicircum{}k) or P-time algorithms into the named model; promotion requires a bound above that envelope on the same SAT slice.
\item Falsify direct transfer if a P function already has the same restricted-model lower bound, such as PARITY for \(\mathrm{AC0}\), HWB for OBDD, or XOR/Tseitin for proof-system degree.
\item For this route, state the exact barrier premise it uses or escapes: relativizing, natural, algebrizing, nonuniform, or advice-leaking; promotion requires the escape to be attached to the specific transfer theorem.
\item Record relativization, natural-proofs, algebrization, nonuniformity, and advice-leakage status for the specific transfer theorem.
\item Control gate: name the source carrier or language, target class, and circuit or resource model before claiming progress.
\item Control gate: derive the exact transfer inequality from \(\mathrm{SAT}\in \mathrm{SIZE}(n^{c})\), NP subset P/poly, or the stated target assumption to the source-carrier upper bound.
\end{itemize}

\paragraph{Stop criteria.}
\begin{itemize}
\item The branch cannot identify a source language/carrier and target complexity class in the same model.
\item The route repeatedly proves only a proxy lower bound with no target-preserving transfer.
\item The remaining work is a slogan about hardness magnification rather than a computable transfer or obstruction.
\end{itemize}

\paragraph{Do not stop if.}
\begin{itemize}
\item A transfer exponent, quantifier repair, or barrier status remains uncomputed.
\item The branch can still name a target-preserving carrier and a falsifiable lower-bound threshold.
\item The failure can become a formal no-go lemma for a route family.
\end{itemize}

\paragraph{Required skills.}
\begin{itemize}
\item entailment-gap-auditor
\item kb-end-to-end-proof-critic
\item long-horizon-route-portfolio-manager
\item main-problem-alignment-auditor
\item next-experiment-queue-ranker
\item paper-contribution-readiness-judge
\item pvsnp-barrier-specificity-reviewer
\item pvsnp-conceptual-pressure-localizer
\item pvsnp-impossibility-theorem-forger
\item pvsnp-inference-quality-controller
\item pvsnp-qualitative-complexity-reasoner
\item pvsnp-quantifier-uniformity-reviewer
\item pvsnp-symbolic-complexity-auditor
\end{itemize}

\subsection*{\#144 cycle-0097: Hardness-Magnification Bridge Realization}

\paragraph{Analysis.} Instantiate \(\mathrm{Psys}_{sym}\) in two concrete variants: FormulaEval as an EF extension rule with row-index bits and FormulaEval as a semantic oracle. For each, compute K\_idx(pi), X\_FE(K), and whether the private-row countermodel survives SATeq and bounded-SR constraints.

\paragraph{Why this is a good target.} It keeps the barrier and model-boundary audit attached to the mathematical transfer rather than as after-the-fact commentary. It is executable as a bounded complexity calculation, CAS-backed exponent audit, or formal transfer check.

\paragraph{Risk.} High drift risk unless rewritten into a concrete complexity-theoretic transfer test.

\paragraph{Branch control.} The branch still touches a relevant object, but the next experiment must be narrowed before continuation.

\paragraph{Premature-stop guard.} Do not stop before the branch has produced a concrete source-to-target transfer theorem, an explicit obstruction lemma, or a barrier-aware reason why the route cannot reach the P vs NP/P/poly target.

\paragraph{Payoff criteria.}
\begin{itemize}
\item For this route, state the exact barrier premise it uses or escapes: relativizing, natural, algebrizing, nonuniform, or advice-leaking; promotion requires the escape to be attached to the specific transfer theorem.
\item Record relativization, natural-proofs, algebrization, nonuniformity, and advice-leakage status for the specific transfer theorem.
\item Control gate: name the source carrier or language, target class, and circuit or resource model before claiming progress.
\item Control gate: derive the exact transfer inequality from \(\mathrm{SAT}\in \mathrm{SIZE}(n^{c})\), NP subset P/poly, or the stated target assumption to the source-carrier upper bound.
\item Control gate: record relativization, natural-proofs, algebrization, nonuniformity, and advice-leakage status for the specific theorem.
\end{itemize}

\paragraph{Stop criteria.}
\begin{itemize}
\item The branch cannot identify a source language/carrier and target complexity class in the same model.
\item The route repeatedly proves only a proxy lower bound with no target-preserving transfer.
\item The remaining work is a slogan about hardness magnification rather than a computable transfer or obstruction.
\end{itemize}

\paragraph{Do not stop if.}
\begin{itemize}
\item A transfer exponent, quantifier repair, or barrier status remains uncomputed.
\item The branch can still name a target-preserving carrier and a falsifiable lower-bound threshold.
\item The failure can become a formal no-go lemma for a route family.
\end{itemize}

\paragraph{Required skills.}
\begin{itemize}
\item entailment-gap-auditor
\item kb-end-to-end-proof-critic
\item long-horizon-route-portfolio-manager
\item main-problem-alignment-auditor
\item next-experiment-queue-ranker
\item paper-contribution-readiness-judge
\item pvsnp-barrier-specificity-reviewer
\item pvsnp-conceptual-pressure-localizer
\item pvsnp-impossibility-theorem-forger
\item pvsnp-inference-quality-controller
\item pvsnp-qualitative-complexity-reasoner
\item pvsnp-quantifier-uniformity-reviewer
\item pvsnp-symbolic-complexity-auditor
\end{itemize}

\subsection*{\#145 cycle-0096: Hardness-Magnification Bridge Realization}

\paragraph{Analysis.} Fix one concrete \(\mathrm{Psys}_{sym}\) syntax: FormulaEval-SAT as either an axiom schema over A(i,j,u), an EF extension rule producing formula bits, a proof line with explicit \(D_{u}\), or a semantic oracle. For that syntax, prove the expansion parser or instantiate the minimal private-row countermodel with the smallest pi, one unexpanded pair (i*,j*), and a verified \(\Gamma_{sym}\) inequality.

\paragraph{Why this is a good target.} It keeps the barrier and model-boundary audit attached to the mathematical transfer rather than as after-the-fact commentary. It is executable as a bounded complexity calculation, CAS-backed exponent audit, or formal transfer check.

\paragraph{Risk.} High drift risk unless rewritten into a concrete complexity-theoretic transfer test.

\paragraph{Branch control.} The branch still touches a relevant object, but the next experiment must be narrowed before continuation.

\paragraph{Premature-stop guard.} Do not stop before the branch has produced a concrete source-to-target transfer theorem, an explicit obstruction lemma, or a barrier-aware reason why the route cannot reach the P vs NP/P/poly target.

\paragraph{Payoff criteria.}
\begin{itemize}
\item For this route, state the exact barrier premise it uses or escapes: relativizing, natural, algebrizing, nonuniform, or advice-leaking; promotion requires the escape to be attached to the specific transfer theorem.
\item Record relativization, natural-proofs, algebrization, nonuniformity, and advice-leakage status for the specific transfer theorem.
\item Control gate: name the source carrier or language, target class, and circuit or resource model before claiming progress.
\item Control gate: derive the exact transfer inequality from \(\mathrm{SAT}\in \mathrm{SIZE}(n^{c})\), NP subset P/poly, or the stated target assumption to the source-carrier upper bound.
\item Control gate: record relativization, natural-proofs, algebrization, nonuniformity, and advice-leakage status for the specific theorem.
\end{itemize}

\paragraph{Stop criteria.}
\begin{itemize}
\item The branch cannot identify a source language/carrier and target complexity class in the same model.
\item The route repeatedly proves only a proxy lower bound with no target-preserving transfer.
\item The remaining work is a slogan about hardness magnification rather than a computable transfer or obstruction.
\end{itemize}

\paragraph{Do not stop if.}
\begin{itemize}
\item A transfer exponent, quantifier repair, or barrier status remains uncomputed.
\item The branch can still name a target-preserving carrier and a falsifiable lower-bound threshold.
\item The failure can become a formal no-go lemma for a route family.
\end{itemize}

\paragraph{Required skills.}
\begin{itemize}
\item entailment-gap-auditor
\item kb-end-to-end-proof-critic
\item long-horizon-route-portfolio-manager
\item main-problem-alignment-auditor
\item next-experiment-queue-ranker
\item paper-contribution-readiness-judge
\item pvsnp-barrier-specificity-reviewer
\item pvsnp-conceptual-pressure-localizer
\item pvsnp-impossibility-theorem-forger
\item pvsnp-inference-quality-controller
\item pvsnp-qualitative-complexity-reasoner
\item pvsnp-quantifier-uniformity-reviewer
\item pvsnp-symbolic-complexity-auditor
\end{itemize}

\subsection*{\#146 cycle-0095: Hardness-Magnification Bridge Realization: bit-explicit parser and FormulaEval boundary}

\paragraph{Analysis.} Instantiate \(FE_{expand,N}\) on the smallest symbolic FormulaEval-SAT rule: one generator circuit G\_N(x,j), one asserted vector y\_j=SAT(G\_N(x,j)), and \(\alpha\) row instantiations. Prove the private-row countermodel for \(\alpha<q\) and the \(N^{1+q-o(1)}\) expansion charge for \(\alpha=q\).

\paragraph{Why this is a good target.} It keeps the barrier and model-boundary audit attached to the mathematical transfer rather than as after-the-fact commentary. It is executable as a bounded complexity calculation, CAS-backed exponent audit, or formal transfer check.

\paragraph{Risk.} High drift risk unless rewritten into a concrete complexity-theoretic transfer test.

\paragraph{Branch control.} The branch still touches a relevant object, but the next experiment must be narrowed before continuation.

\paragraph{Premature-stop guard.} Do not stop before the branch has produced a concrete source-to-target transfer theorem, an explicit obstruction lemma, or a barrier-aware reason why the route cannot reach the P vs NP/P/poly target.

\paragraph{Payoff criteria.}
\begin{itemize}
\item For this route, state the exact barrier premise it uses or escapes: relativizing, natural, algebrizing, nonuniform, or advice-leaking; promotion requires the escape to be attached to the specific transfer theorem.
\item Record relativization, natural-proofs, algebrization, nonuniformity, and advice-leakage status for the specific transfer theorem.
\item Control gate: name the source carrier or language, target class, and circuit or resource model before claiming progress.
\item Control gate: derive the exact transfer inequality from \(\mathrm{SAT}\in \mathrm{SIZE}(n^{c})\), NP subset P/poly, or the stated target assumption to the source-carrier upper bound.
\item Control gate: record relativization, natural-proofs, algebrization, nonuniformity, and advice-leakage status for the specific theorem.
\end{itemize}

\paragraph{Stop criteria.}
\begin{itemize}
\item The branch cannot identify a source language/carrier and target complexity class in the same model.
\item The route repeatedly proves only a proxy lower bound with no target-preserving transfer.
\item The remaining work is a slogan about hardness magnification rather than a computable transfer or obstruction.
\end{itemize}

\paragraph{Do not stop if.}
\begin{itemize}
\item A transfer exponent, quantifier repair, or barrier status remains uncomputed.
\item The branch can still name a target-preserving carrier and a falsifiable lower-bound threshold.
\item The failure can become a formal no-go lemma for a route family.
\end{itemize}

\paragraph{Required skills.}
\begin{itemize}
\item entailment-gap-auditor
\item kb-end-to-end-proof-critic
\item long-horizon-route-portfolio-manager
\item main-problem-alignment-auditor
\item next-experiment-queue-ranker
\item paper-contribution-readiness-judge
\item pvsnp-barrier-specificity-reviewer
\item pvsnp-conceptual-pressure-localizer
\item pvsnp-impossibility-theorem-forger
\item pvsnp-inference-quality-controller
\item pvsnp-qualitative-complexity-reasoner
\item pvsnp-quantifier-uniformity-reviewer
\item pvsnp-symbolic-complexity-auditor
\end{itemize}

\subsection*{\#147 cycle-0056: Hardness-Magnification Bridge Realization: XOR quotient codebook gate}

\paragraph{Analysis.} Instantiate a candidate high-q generator and compute the exponent of its SAT-preserving odd-support transcript range versus affine span rank; reject it if min\{theta,rank-exponent\}+locator-cost falls below 1+epsilon.

\paragraph{Why this is a good target.} It keeps the barrier and model-boundary audit attached to the mathematical transfer rather than as after-the-fact commentary. It is executable as a bounded complexity calculation, CAS-backed exponent audit, or formal transfer check.

\paragraph{Risk.} High drift risk unless rewritten into a concrete complexity-theoretic transfer test.

\paragraph{Branch control.} The branch still touches a relevant object, but the next experiment must be narrowed before continuation.

\paragraph{Premature-stop guard.} Do not stop before the branch has produced a concrete source-to-target transfer theorem, an explicit obstruction lemma, or a barrier-aware reason why the route cannot reach the P vs NP/P/poly target.

\paragraph{Payoff criteria.}
\begin{itemize}
\item For the high-q generated-OR carrier, compute the exact leakage exponent of the remaining witness-cover or unary-residual mechanism and compare it with 1+epsilon-q.
\item Promotion requires \(q+a>1+\epsilon\) after every projection, padding, and witness-cover leakage term is accounted for; otherwise the branch becomes a bounded-window or no-go result.
\item For the proof-complexity branch, separate fixed-family proof hardness from language hardness by naming the proof system, shortcut algorithm, and semantic lower-bound target.
\item Promotion requires a proof-system lower bound that survives the Gaussian-elimination or recognition shortcut; otherwise the result is a restricted-proof-system separation only.
\item For this route, state the exact barrier premise it uses or escapes: relativizing, natural, algebrizing, nonuniform, or advice-leaking; promotion requires the escape to be attached to the specific transfer theorem.
\item Record relativization, natural-proofs, algebrization, nonuniformity, and advice-leakage status for the specific transfer theorem.
\end{itemize}

\paragraph{Stop criteria.}
\begin{itemize}
\item The branch cannot identify a source language/carrier and target complexity class in the same model.
\item The route repeatedly proves only a proxy lower bound with no target-preserving transfer.
\item The remaining work is a slogan about hardness magnification rather than a computable transfer or obstruction.
\end{itemize}

\paragraph{Do not stop if.}
\begin{itemize}
\item A transfer exponent, quantifier repair, or barrier status remains uncomputed.
\item The branch can still name a target-preserving carrier and a falsifiable lower-bound threshold.
\item The failure can become a formal no-go lemma for a route family.
\end{itemize}

\paragraph{Required skills.}
\begin{itemize}
\item entailment-gap-auditor
\item kb-end-to-end-proof-critic
\item long-horizon-route-portfolio-manager
\item main-problem-alignment-auditor
\item next-experiment-queue-ranker
\item paper-contribution-readiness-judge
\item pvsnp-barrier-specificity-reviewer
\item pvsnp-conceptual-pressure-localizer
\item pvsnp-impossibility-theorem-forger
\item pvsnp-inference-quality-controller
\item pvsnp-qualitative-complexity-reasoner
\item pvsnp-quantifier-uniformity-reviewer
\item pvsnp-symbolic-complexity-auditor
\end{itemize}

\subsection*{\#148 cycle-0054: Hardness-Magnification Bridge Realization: XOR Carrier Approximate Affine-Sketch Gate}

\paragraph{Analysis.} Instantiate the hash-based XOR generator and compute whether its incidence map has an explicit low-cost piecewise-affine cover; if none is found, formulate the exact extraction theorem needed to turn small \(H_{N}\) circuits into such covers.

\paragraph{Why this is a good target.} It keeps the barrier and model-boundary audit attached to the mathematical transfer rather than as after-the-fact commentary. It is executable as a bounded complexity calculation, CAS-backed exponent audit, or formal transfer check.

\paragraph{Risk.} High drift risk unless rewritten into a concrete complexity-theoretic transfer test.

\paragraph{Branch control.} The branch still touches a relevant object, but the next experiment must be narrowed before continuation.

\paragraph{Premature-stop guard.} Do not stop before the branch has produced a concrete source-to-target transfer theorem, an explicit obstruction lemma, or a barrier-aware reason why the route cannot reach the P vs NP/P/poly target.

\paragraph{Payoff criteria.}
\begin{itemize}
\item For the proof-complexity branch, separate fixed-family proof hardness from language hardness by naming the proof system, shortcut algorithm, and semantic lower-bound target.
\item Promotion requires a proof-system lower bound that survives the Gaussian-elimination or recognition shortcut; otherwise the result is a restricted-proof-system separation only.
\item For this route, state the exact barrier premise it uses or escapes: relativizing, natural, algebrizing, nonuniform, or advice-leaking; promotion requires the escape to be attached to the specific transfer theorem.
\item Record relativization, natural-proofs, algebrization, nonuniformity, and advice-leakage status for the specific transfer theorem.
\item Control gate: name the source carrier or language, target class, and circuit or resource model before claiming progress.
\item Control gate: derive the exact transfer inequality from \(\mathrm{SAT}\in \mathrm{SIZE}(n^{c})\), NP subset P/poly, or the stated target assumption to the source-carrier upper bound.
\end{itemize}

\paragraph{Stop criteria.}
\begin{itemize}
\item The branch cannot identify a source language/carrier and target complexity class in the same model.
\item The route repeatedly proves only a proxy lower bound with no target-preserving transfer.
\item The remaining work is a slogan about hardness magnification rather than a computable transfer or obstruction.
\end{itemize}

\paragraph{Do not stop if.}
\begin{itemize}
\item A transfer exponent, quantifier repair, or barrier status remains uncomputed.
\item The branch can still name a target-preserving carrier and a falsifiable lower-bound threshold.
\item The failure can become a formal no-go lemma for a route family.
\end{itemize}

\paragraph{Required skills.}
\begin{itemize}
\item entailment-gap-auditor
\item kb-end-to-end-proof-critic
\item long-horizon-route-portfolio-manager
\item main-problem-alignment-auditor
\item next-experiment-queue-ranker
\item paper-contribution-readiness-judge
\item pvsnp-barrier-specificity-reviewer
\item pvsnp-conceptual-pressure-localizer
\item pvsnp-impossibility-theorem-forger
\item pvsnp-inference-quality-controller
\item pvsnp-qualitative-complexity-reasoner
\item pvsnp-quantifier-uniformity-reviewer
\item pvsnp-symbolic-complexity-auditor
\end{itemize}

\subsection*{\#149 cycle-0050: Hardness-Magnification Bridge Realization: Relativizing Extraction Falsified}

\paragraph{Analysis.} Audit real-SAT self-reduction or proof-certificate structure as a possible nonrelativizing extractor, and compute whether its verifier/certificate overhead stays below \(N^{1+\epsilon}\) rather than collapsing to restricted-SAT hardness.

\paragraph{Why this is a good target.} It keeps the barrier and model-boundary audit attached to the mathematical transfer rather than as after-the-fact commentary. It is executable as a bounded complexity calculation, CAS-backed exponent audit, or formal transfer check.

\paragraph{Risk.} High drift risk unless rewritten into a concrete complexity-theoretic transfer test.

\paragraph{Branch control.} The branch still touches a relevant object, but the next experiment must be narrowed before continuation.

\paragraph{Premature-stop guard.} Do not stop before the branch has produced a concrete source-to-target transfer theorem, an explicit obstruction lemma, or a barrier-aware reason why the route cannot reach the P vs NP/P/poly target.

\paragraph{Payoff criteria.}
\begin{itemize}
\item For this route, state the exact barrier premise it uses or escapes: relativizing, natural, algebrizing, nonuniform, or advice-leaking; promotion requires the escape to be attached to the specific transfer theorem.
\item Record relativization, natural-proofs, algebrization, nonuniformity, and advice-leakage status for the specific transfer theorem.
\item Control gate: name the source carrier or language, target class, and circuit or resource model before claiming progress.
\item Control gate: derive the exact transfer inequality from \(\mathrm{SAT}\in \mathrm{SIZE}(n^{c})\), NP subset P/poly, or the stated target assumption to the source-carrier upper bound.
\item Control gate: record relativization, natural-proofs, algebrization, nonuniformity, and advice-leakage status for the specific theorem.
\end{itemize}

\paragraph{Stop criteria.}
\begin{itemize}
\item The branch cannot identify a source language/carrier and target complexity class in the same model.
\item The route repeatedly proves only a proxy lower bound with no target-preserving transfer.
\item The remaining work is a slogan about hardness magnification rather than a computable transfer or obstruction.
\end{itemize}

\paragraph{Do not stop if.}
\begin{itemize}
\item A transfer exponent, quantifier repair, or barrier status remains uncomputed.
\item The branch can still name a target-preserving carrier and a falsifiable lower-bound threshold.
\item The failure can become a formal no-go lemma for a route family.
\end{itemize}

\paragraph{Required skills.}
\begin{itemize}
\item entailment-gap-auditor
\item kb-end-to-end-proof-critic
\item long-horizon-route-portfolio-manager
\item main-problem-alignment-auditor
\item next-experiment-queue-ranker
\item paper-contribution-readiness-judge
\item pvsnp-barrier-specificity-reviewer
\item pvsnp-conceptual-pressure-localizer
\item pvsnp-impossibility-theorem-forger
\item pvsnp-inference-quality-controller
\item pvsnp-qualitative-complexity-reasoner
\item pvsnp-quantifier-uniformity-reviewer
\item pvsnp-symbolic-complexity-auditor
\end{itemize}

\subsection*{\#150 cycle-0049: Hardness-Magnification Bridge Realization: locality extraction threshold}

\paragraph{Analysis.} Formalize one candidate extraction using SAT self-reducibility plus generator local decodability, then compute whether the required decodability already implies neutralizing embeddings and hence demagnification.

\paragraph{Why this is a good target.} It keeps the barrier and model-boundary audit attached to the mathematical transfer rather than as after-the-fact commentary. It is executable as a bounded complexity calculation, CAS-backed exponent audit, or formal transfer check.

\paragraph{Risk.} High drift risk unless rewritten into a concrete complexity-theoretic transfer test.

\paragraph{Branch control.} The branch still touches a relevant object, but the next experiment must be narrowed before continuation.

\paragraph{Premature-stop guard.} Do not stop before the branch has produced a concrete source-to-target transfer theorem, an explicit obstruction lemma, or a barrier-aware reason why the route cannot reach the P vs NP/P/poly target.

\paragraph{Payoff criteria.}
\begin{itemize}
\item For the high-q generated-OR carrier, compute the exact leakage exponent of the remaining witness-cover or unary-residual mechanism and compare it with 1+epsilon-q.
\item Promotion requires \(q+a>1+\epsilon\) after every projection, padding, and witness-cover leakage term is accounted for; otherwise the branch becomes a bounded-window or no-go result.
\item For this route, state the exact barrier premise it uses or escapes: relativizing, natural, algebrizing, nonuniform, or advice-leaking; promotion requires the escape to be attached to the specific transfer theorem.
\item Record relativization, natural-proofs, algebrization, nonuniformity, and advice-leakage status for the specific transfer theorem.
\item Control gate: name the source carrier or language, target class, and circuit or resource model before claiming progress.
\item Control gate: derive the exact transfer inequality from \(\mathrm{SAT}\in \mathrm{SIZE}(n^{c})\), NP subset P/poly, or the stated target assumption to the source-carrier upper bound.
\end{itemize}

\paragraph{Stop criteria.}
\begin{itemize}
\item The branch cannot identify a source language/carrier and target complexity class in the same model.
\item The route repeatedly proves only a proxy lower bound with no target-preserving transfer.
\item The remaining work is a slogan about hardness magnification rather than a computable transfer or obstruction.
\end{itemize}

\paragraph{Do not stop if.}
\begin{itemize}
\item A transfer exponent, quantifier repair, or barrier status remains uncomputed.
\item The branch can still name a target-preserving carrier and a falsifiable lower-bound threshold.
\item The failure can become a formal no-go lemma for a route family.
\end{itemize}

\paragraph{Required skills.}
\begin{itemize}
\item entailment-gap-auditor
\item kb-end-to-end-proof-critic
\item long-horizon-route-portfolio-manager
\item main-problem-alignment-auditor
\item next-experiment-queue-ranker
\item paper-contribution-readiness-judge
\item pvsnp-barrier-specificity-reviewer
\item pvsnp-conceptual-pressure-localizer
\item pvsnp-impossibility-theorem-forger
\item pvsnp-inference-quality-controller
\item pvsnp-qualitative-complexity-reasoner
\item pvsnp-quantifier-uniformity-reviewer
\item pvsnp-symbolic-complexity-auditor
\end{itemize}

\subsection*{\#151 cycle-0041: One-sided promise-separator gate for high-q generated-OR magnification}

\paragraph{Analysis.} Test a range-closure repair: augment \(G_{N}\) so restrictions of emitted formulas also appear in controlled positions, compute the q/rho overhead, and see whether this lets a one-sided separator self-reduce to a term cover without breaking the \(N^{1+\epsilon}\) threshold.

\paragraph{Why this is a good target.} It keeps the barrier and model-boundary audit attached to the mathematical transfer rather than as after-the-fact commentary. It is executable as a bounded complexity calculation, CAS-backed exponent audit, or formal transfer check.

\paragraph{Risk.} High drift risk unless rewritten into a concrete complexity-theoretic transfer test.

\paragraph{Branch control.} The branch still touches a relevant object, but the next experiment must be narrowed before continuation.

\paragraph{Premature-stop guard.} Do not stop before the branch has produced a concrete source-to-target transfer theorem, an explicit obstruction lemma, or a barrier-aware reason why the route cannot reach the P vs NP/P/poly target.

\paragraph{Payoff criteria.}
\begin{itemize}
\item For the one-sided semantic-separator branch, compute the cap profile A\_N(s) for the named generator and predicate class; promotion requires a bound strong enough to force Sep\_N\textasciicircum{}\{1s\}(E)>N\textasciicircum{}\{1+epsilon-o(1)\}.
\item Exhibit or refute a subthreshold semantic planted-slice predicate; this decides whether the generated-OR carrier needs range-closure repair, stronger pseudorandomness, or quarantine.
\item For the certificate-cover branch, compute K,E,L and the induced circuit upper bound C(H\_N)<=C(G\_N)+O(K(s+L\textasciicircum{}\{O(1)\})+EN); promotion requires every sound local cover to stay above the \(N^{1+\epsilon-o(1)}\) source threshold.
\item Falsify the branch if Horn, 2-SAT, bounded-width CNF, or bounded-backdoor emitted formulas give subthreshold satisfying-term or sound-local-certificate covers.
\item For the high-q generated-OR carrier, compute the exact leakage exponent of the remaining witness-cover or unary-residual mechanism and compare it with 1+epsilon-q.
\item Promotion requires \(q+a>1+\epsilon\) after every projection, padding, and witness-cover leakage term is accounted for; otherwise the branch becomes a bounded-window or no-go result.
\end{itemize}

\paragraph{Stop criteria.}
\begin{itemize}
\item The branch cannot identify a source language/carrier and target complexity class in the same model.
\item The route repeatedly proves only a proxy lower bound with no target-preserving transfer.
\item The remaining work is a slogan about hardness magnification rather than a computable transfer or obstruction.
\end{itemize}

\paragraph{Do not stop if.}
\begin{itemize}
\item A transfer exponent, quantifier repair, or barrier status remains uncomputed.
\item The branch can still name a target-preserving carrier and a falsifiable lower-bound threshold.
\item The failure can become a formal no-go lemma for a route family.
\end{itemize}

\paragraph{Required skills.}
\begin{itemize}
\item entailment-gap-auditor
\item kb-end-to-end-proof-critic
\item long-horizon-route-portfolio-manager
\item main-problem-alignment-auditor
\item next-experiment-queue-ranker
\item paper-contribution-readiness-judge
\item pvsnp-barrier-specificity-reviewer
\item pvsnp-conceptual-pressure-localizer
\item pvsnp-impossibility-theorem-forger
\item pvsnp-inference-quality-controller
\item pvsnp-qualitative-complexity-reasoner
\item pvsnp-quantifier-uniformity-reviewer
\item pvsnp-symbolic-complexity-auditor
\end{itemize}

\subsection*{\#152 cycle-0037: Hardness-Magnification Bridge Realization: approximate witness-cover gate}

\paragraph{Analysis.} Instantiate the thought experiment with a concrete planted/anti-planted CNF generator and compute the minimum position-indexed assignment hitting exponent h for its satisfiable range; compare h to 1+epsilon-q.

\paragraph{Why this is a good target.} It keeps the barrier and model-boundary audit attached to the mathematical transfer rather than as after-the-fact commentary. It is executable as a bounded complexity calculation, CAS-backed exponent audit, or formal transfer check.

\paragraph{Risk.} High drift risk unless rewritten into a concrete complexity-theoretic transfer test.

\paragraph{Branch control.} The branch still touches a relevant object, but the next experiment must be narrowed before continuation.

\paragraph{Premature-stop guard.} Do not stop before the branch has produced a concrete source-to-target transfer theorem, an explicit obstruction lemma, or a barrier-aware reason why the route cannot reach the P vs NP/P/poly target.

\paragraph{Payoff criteria.}
\begin{itemize}
\item For the high-q generated-OR carrier, compute the exact leakage exponent of the remaining witness-cover or unary-residual mechanism and compare it with 1+epsilon-q.
\item Promotion requires \(q+a>1+\epsilon\) after every projection, padding, and witness-cover leakage term is accounted for; otherwise the branch becomes a bounded-window or no-go result.
\item For this route, state the exact barrier premise it uses or escapes: relativizing, natural, algebrizing, nonuniform, or advice-leaking; promotion requires the escape to be attached to the specific transfer theorem.
\item Record relativization, natural-proofs, algebrization, nonuniformity, and advice-leakage status for the specific transfer theorem.
\item Control gate: name the source carrier or language, target class, and circuit or resource model before claiming progress.
\item Control gate: derive the exact transfer inequality from \(\mathrm{SAT}\in \mathrm{SIZE}(n^{c})\), NP subset P/poly, or the stated target assumption to the source-carrier upper bound.
\end{itemize}

\paragraph{Stop criteria.}
\begin{itemize}
\item The branch cannot identify a source language/carrier and target complexity class in the same model.
\item The route repeatedly proves only a proxy lower bound with no target-preserving transfer.
\item The remaining work is a slogan about hardness magnification rather than a computable transfer or obstruction.
\end{itemize}

\paragraph{Do not stop if.}
\begin{itemize}
\item A transfer exponent, quantifier repair, or barrier status remains uncomputed.
\item The branch can still name a target-preserving carrier and a falsifiable lower-bound threshold.
\item The failure can become a formal no-go lemma for a route family.
\end{itemize}

\paragraph{Required skills.}
\begin{itemize}
\item entailment-gap-auditor
\item kb-end-to-end-proof-critic
\item long-horizon-route-portfolio-manager
\item main-problem-alignment-auditor
\item next-experiment-queue-ranker
\item paper-contribution-readiness-judge
\item pvsnp-barrier-specificity-reviewer
\item pvsnp-conceptual-pressure-localizer
\item pvsnp-impossibility-theorem-forger
\item pvsnp-inference-quality-controller
\item pvsnp-qualitative-complexity-reasoner
\item pvsnp-quantifier-uniformity-reviewer
\item pvsnp-symbolic-complexity-auditor
\end{itemize}

\subsection*{\#153 cycle-0036: Hardness-Magnification Bridge Realization: position-indexed unary residual gate}

\paragraph{Analysis.} Formalize \(Iota_{N}^{pos}\) as a minimax transcript-residual measure and test whether expander-style slot mixing can rule out all low-budget position-indexed unary decompositions without becoming a natural-proof argument.

\paragraph{Why this is a good target.} It keeps the barrier and model-boundary audit attached to the mathematical transfer rather than as after-the-fact commentary. It is executable as a bounded complexity calculation, CAS-backed exponent audit, or formal transfer check.

\paragraph{Risk.} High drift risk unless rewritten into a concrete complexity-theoretic transfer test.

\paragraph{Branch control.} The branch still touches a relevant object, but the next experiment must be narrowed before continuation.

\paragraph{Premature-stop guard.} Do not stop before the branch has produced a concrete source-to-target transfer theorem, an explicit obstruction lemma, or a barrier-aware reason why the route cannot reach the P vs NP/P/poly target.

\paragraph{Payoff criteria.}
\begin{itemize}
\item For the high-q generated-OR carrier, compute the exact leakage exponent of the remaining witness-cover or unary-residual mechanism and compare it with 1+epsilon-q.
\item Promotion requires \(q+a>1+\epsilon\) after every projection, padding, and witness-cover leakage term is accounted for; otherwise the branch becomes a bounded-window or no-go result.
\item For this route, state the exact barrier premise it uses or escapes: relativizing, natural, algebrizing, nonuniform, or advice-leaking; promotion requires the escape to be attached to the specific transfer theorem.
\item Record relativization, natural-proofs, algebrization, nonuniformity, and advice-leakage status for the specific transfer theorem.
\item Control gate: name the source carrier or language, target class, and circuit or resource model before claiming progress.
\item Control gate: derive the exact transfer inequality from \(\mathrm{SAT}\in \mathrm{SIZE}(n^{c})\), NP subset P/poly, or the stated target assumption to the source-carrier upper bound.
\end{itemize}

\paragraph{Stop criteria.}
\begin{itemize}
\item The branch cannot identify a source language/carrier and target complexity class in the same model.
\item The route repeatedly proves only a proxy lower bound with no target-preserving transfer.
\item The remaining work is a slogan about hardness magnification rather than a computable transfer or obstruction.
\end{itemize}

\paragraph{Do not stop if.}
\begin{itemize}
\item A transfer exponent, quantifier repair, or barrier status remains uncomputed.
\item The branch can still name a target-preserving carrier and a falsifiable lower-bound threshold.
\item The failure can become a formal no-go lemma for a route family.
\end{itemize}

\paragraph{Required skills.}
\begin{itemize}
\item entailment-gap-auditor
\item kb-end-to-end-proof-critic
\item long-horizon-route-portfolio-manager
\item main-problem-alignment-auditor
\item next-experiment-queue-ranker
\item paper-contribution-readiness-judge
\item pvsnp-barrier-specificity-reviewer
\item pvsnp-conceptual-pressure-localizer
\item pvsnp-impossibility-theorem-forger
\item pvsnp-inference-quality-controller
\item pvsnp-qualitative-complexity-reasoner
\item pvsnp-quantifier-uniformity-reviewer
\item pvsnp-symbolic-complexity-auditor
\end{itemize}

\subsection*{\#154 cycle-0035: Hardness-Magnification Bridge Realization: high-q arity cutoff for generated-OR carriers}

\paragraph{Analysis.} Formalize unary transcripts as monotone OR-of-range predicates, then test whether random-looking high-q generators make every subthreshold unary transcript leave residual circuit complexity above \(N^{1+\epsilon-o(1)}\) or whether a selected-image separator falsifier appears.

\paragraph{Why this is a good target.} It keeps the barrier and model-boundary audit attached to the mathematical transfer rather than as after-the-fact commentary. It is executable as a bounded complexity calculation, CAS-backed exponent audit, or formal transfer check.

\paragraph{Risk.} High drift risk unless rewritten into a concrete complexity-theoretic transfer test.

\paragraph{Branch control.} The branch still touches a relevant object, but the next experiment must be narrowed before continuation.

\paragraph{Premature-stop guard.} Do not stop before the branch has produced a concrete source-to-target transfer theorem, an explicit obstruction lemma, or a barrier-aware reason why the route cannot reach the P vs NP/P/poly target.

\paragraph{Payoff criteria.}
\begin{itemize}
\item For the high-q generated-OR carrier, compute the exact leakage exponent of the remaining witness-cover or unary-residual mechanism and compare it with 1+epsilon-q.
\item Promotion requires \(q+a>1+\epsilon\) after every projection, padding, and witness-cover leakage term is accounted for; otherwise the branch becomes a bounded-window or no-go result.
\item For this route, state the exact barrier premise it uses or escapes: relativizing, natural, algebrizing, nonuniform, or advice-leaking; promotion requires the escape to be attached to the specific transfer theorem.
\item Record relativization, natural-proofs, algebrization, nonuniformity, and advice-leakage status for the specific transfer theorem.
\item Control gate: name the source carrier or language, target class, and circuit or resource model before claiming progress.
\item Control gate: derive the exact transfer inequality from \(\mathrm{SAT}\in \mathrm{SIZE}(n^{c})\), NP subset P/poly, or the stated target assumption to the source-carrier upper bound.
\end{itemize}

\paragraph{Stop criteria.}
\begin{itemize}
\item The branch cannot identify a source language/carrier and target complexity class in the same model.
\item The route repeatedly proves only a proxy lower bound with no target-preserving transfer.
\item The remaining work is a slogan about hardness magnification rather than a computable transfer or obstruction.
\end{itemize}

\paragraph{Do not stop if.}
\begin{itemize}
\item A transfer exponent, quantifier repair, or barrier status remains uncomputed.
\item The branch can still name a target-preserving carrier and a falsifiable lower-bound threshold.
\item The failure can become a formal no-go lemma for a route family.
\end{itemize}

\paragraph{Required skills.}
\begin{itemize}
\item entailment-gap-auditor
\item kb-end-to-end-proof-critic
\item long-horizon-route-portfolio-manager
\item main-problem-alignment-auditor
\item next-experiment-queue-ranker
\item paper-contribution-readiness-judge
\item pvsnp-barrier-specificity-reviewer
\item pvsnp-conceptual-pressure-localizer
\item pvsnp-impossibility-theorem-forger
\item pvsnp-inference-quality-controller
\item pvsnp-qualitative-complexity-reasoner
\item pvsnp-quantifier-uniformity-reviewer
\item pvsnp-symbolic-complexity-auditor
\end{itemize}

\subsection*{\#155 cycle-0030: Hardness-Magnification Bridge Realization}

\paragraph{Analysis.} Formalize \(SepHit_{N}\) as a bipartite hitting-avoidance problem between positive bags and negative formula occurrences, then test whether random-looking generators satisfy the required lower bound by a counting argument and where explicitness breaks the argument.

\paragraph{Why this is a good target.} It keeps the barrier and model-boundary audit attached to the mathematical transfer rather than as after-the-fact commentary. It is executable as a bounded complexity calculation, CAS-backed exponent audit, or formal transfer check.

\paragraph{Risk.} High drift risk unless rewritten into a concrete complexity-theoretic transfer test.

\paragraph{Branch control.} The branch still touches a relevant object, but the next experiment must be narrowed before continuation.

\paragraph{Premature-stop guard.} Do not stop before the branch has produced a concrete source-to-target transfer theorem, an explicit obstruction lemma, or a barrier-aware reason why the route cannot reach the P vs NP/P/poly target.

\paragraph{Payoff criteria.}
\begin{itemize}
\item For this route, state the exact barrier premise it uses or escapes: relativizing, natural, algebrizing, nonuniform, or advice-leaking; promotion requires the escape to be attached to the specific transfer theorem.
\item Record relativization, natural-proofs, algebrization, nonuniformity, and advice-leakage status for the specific transfer theorem.
\item Control gate: name the source carrier or language, target class, and circuit or resource model before claiming progress.
\item Control gate: derive the exact transfer inequality from \(\mathrm{SAT}\in \mathrm{SIZE}(n^{c})\), NP subset P/poly, or the stated target assumption to the source-carrier upper bound.
\item Control gate: record relativization, natural-proofs, algebrization, nonuniformity, and advice-leakage status for the specific theorem.
\end{itemize}

\paragraph{Stop criteria.}
\begin{itemize}
\item The branch cannot identify a source language/carrier and target complexity class in the same model.
\item The route repeatedly proves only a proxy lower bound with no target-preserving transfer.
\item The remaining work is a slogan about hardness magnification rather than a computable transfer or obstruction.
\end{itemize}

\paragraph{Do not stop if.}
\begin{itemize}
\item A transfer exponent, quantifier repair, or barrier status remains uncomputed.
\item The branch can still name a target-preserving carrier and a falsifiable lower-bound threshold.
\item The failure can become a formal no-go lemma for a route family.
\end{itemize}

\paragraph{Required skills.}
\begin{itemize}
\item entailment-gap-auditor
\item kb-end-to-end-proof-critic
\item long-horizon-route-portfolio-manager
\item main-problem-alignment-auditor
\item next-experiment-queue-ranker
\item paper-contribution-readiness-judge
\item pvsnp-barrier-specificity-reviewer
\item pvsnp-conceptual-pressure-localizer
\item pvsnp-impossibility-theorem-forger
\item pvsnp-inference-quality-controller
\item pvsnp-qualitative-complexity-reasoner
\item pvsnp-quantifier-uniformity-reviewer
\item pvsnp-symbolic-complexity-auditor
\end{itemize}

\subsection*{\#156 cycle-0029: Hardness-Magnification Bridge Realization: one-sided separating-ensemble leakage gate}

\paragraph{Analysis.} Try to lower-bound or construct \(Psi_{N}^{sep}\) for a random-like generator: estimate whether selected positive and negative images admit \(N^{b}-size\) separators when \(R=N^{r}\) and \(s=N^{ell}\) are below the magnification threshold.

\paragraph{Why this is a good target.} It keeps the barrier and model-boundary audit attached to the mathematical transfer rather than as after-the-fact commentary. It is executable as a bounded complexity calculation, CAS-backed exponent audit, or formal transfer check.

\paragraph{Risk.} High drift risk unless rewritten into a concrete complexity-theoretic transfer test.

\paragraph{Branch control.} The branch still touches a relevant object, but the next experiment must be narrowed before continuation.

\paragraph{Premature-stop guard.} Do not stop before the branch has produced a concrete source-to-target transfer theorem, an explicit obstruction lemma, or a barrier-aware reason why the route cannot reach the P vs NP/P/poly target.

\paragraph{Payoff criteria.}
\begin{itemize}
\item For the high-q generated-OR carrier, compute the exact leakage exponent of the remaining witness-cover or unary-residual mechanism and compare it with 1+epsilon-q.
\item Promotion requires \(q+a>1+\epsilon\) after every projection, padding, and witness-cover leakage term is accounted for; otherwise the branch becomes a bounded-window or no-go result.
\item For this route, state the exact barrier premise it uses or escapes: relativizing, natural, algebrizing, nonuniform, or advice-leaking; promotion requires the escape to be attached to the specific transfer theorem.
\item Record relativization, natural-proofs, algebrization, nonuniformity, and advice-leakage status for the specific transfer theorem.
\item Control gate: name the source carrier or language, target class, and circuit or resource model before claiming progress.
\item Control gate: derive the exact transfer inequality from \(\mathrm{SAT}\in \mathrm{SIZE}(n^{c})\), NP subset P/poly, or the stated target assumption to the source-carrier upper bound.
\end{itemize}

\paragraph{Stop criteria.}
\begin{itemize}
\item The branch cannot identify a source language/carrier and target complexity class in the same model.
\item The route repeatedly proves only a proxy lower bound with no target-preserving transfer.
\item The remaining work is a slogan about hardness magnification rather than a computable transfer or obstruction.
\end{itemize}

\paragraph{Do not stop if.}
\begin{itemize}
\item A transfer exponent, quantifier repair, or barrier status remains uncomputed.
\item The branch can still name a target-preserving carrier and a falsifiable lower-bound threshold.
\item The failure can become a formal no-go lemma for a route family.
\end{itemize}

\paragraph{Required skills.}
\begin{itemize}
\item entailment-gap-auditor
\item kb-end-to-end-proof-critic
\item long-horizon-route-portfolio-manager
\item main-problem-alignment-auditor
\item next-experiment-queue-ranker
\item paper-contribution-readiness-judge
\item pvsnp-barrier-specificity-reviewer
\item pvsnp-conceptual-pressure-localizer
\item pvsnp-impossibility-theorem-forger
\item pvsnp-inference-quality-controller
\item pvsnp-qualitative-complexity-reasoner
\item pvsnp-quantifier-uniformity-reviewer
\item pvsnp-symbolic-complexity-auditor
\end{itemize}

\subsection*{\#157 cycle-0028: Hardness-Magnification Bridge Realization}

\paragraph{Analysis.} Estimate \(Xi_{N}^{ens}\) for a random high-entropy generator by a union-bound model over ensembles, then isolate which part of the argument can survive explicitness and avoid natural-proof leakage.

\paragraph{Why this is a good target.} It keeps the barrier and model-boundary audit attached to the mathematical transfer rather than as after-the-fact commentary. It is executable as a bounded complexity calculation, CAS-backed exponent audit, or formal transfer check.

\paragraph{Risk.} High drift risk unless rewritten into a concrete complexity-theoretic transfer test.

\paragraph{Branch control.} The branch still touches a relevant object, but the next experiment must be narrowed before continuation.

\paragraph{Premature-stop guard.} Do not stop before the branch has produced a concrete source-to-target transfer theorem, an explicit obstruction lemma, or a barrier-aware reason why the route cannot reach the P vs NP/P/poly target.

\paragraph{Payoff criteria.}
\begin{itemize}
\item For the high-q generated-OR carrier, compute the exact leakage exponent of the remaining witness-cover or unary-residual mechanism and compare it with 1+epsilon-q.
\item Promotion requires \(q+a>1+\epsilon\) after every projection, padding, and witness-cover leakage term is accounted for; otherwise the branch becomes a bounded-window or no-go result.
\item For this route, state the exact barrier premise it uses or escapes: relativizing, natural, algebrizing, nonuniform, or advice-leaking; promotion requires the escape to be attached to the specific transfer theorem.
\item Record relativization, natural-proofs, algebrization, nonuniformity, and advice-leakage status for the specific transfer theorem.
\item Control gate: name the source carrier or language, target class, and circuit or resource model before claiming progress.
\item Control gate: derive the exact transfer inequality from \(\mathrm{SAT}\in \mathrm{SIZE}(n^{c})\), NP subset P/poly, or the stated target assumption to the source-carrier upper bound.
\end{itemize}

\paragraph{Stop criteria.}
\begin{itemize}
\item The branch cannot identify a source language/carrier and target complexity class in the same model.
\item The route repeatedly proves only a proxy lower bound with no target-preserving transfer.
\item The remaining work is a slogan about hardness magnification rather than a computable transfer or obstruction.
\end{itemize}

\paragraph{Do not stop if.}
\begin{itemize}
\item A transfer exponent, quantifier repair, or barrier status remains uncomputed.
\item The branch can still name a target-preserving carrier and a falsifiable lower-bound threshold.
\item The failure can become a formal no-go lemma for a route family.
\end{itemize}

\paragraph{Required skills.}
\begin{itemize}
\item entailment-gap-auditor
\item kb-end-to-end-proof-critic
\item long-horizon-route-portfolio-manager
\item main-problem-alignment-auditor
\item next-experiment-queue-ranker
\item paper-contribution-readiness-judge
\item pvsnp-barrier-specificity-reviewer
\item pvsnp-conceptual-pressure-localizer
\item pvsnp-impossibility-theorem-forger
\item pvsnp-inference-quality-controller
\item pvsnp-qualitative-complexity-reasoner
\item pvsnp-quantifier-uniformity-reviewer
\item pvsnp-symbolic-complexity-auditor
\end{itemize}

\subsection*{\#158 cycle-0027: Hardness-Magnification Bridge Realization: approximate list-selector leakage gate}

\paragraph{Analysis.} Instantiate \(G_{N}\) with a high-entropy expander/hash-generated SAT family and compute whether any tau=O(epsilon/2) list selector has a selected image admitting a subthreshold promise solver.

\paragraph{Why this is a good target.} It keeps the barrier and model-boundary audit attached to the mathematical transfer rather than as after-the-fact commentary. It is executable as a bounded complexity calculation, CAS-backed exponent audit, or formal transfer check.

\paragraph{Risk.} High drift risk unless rewritten into a concrete complexity-theoretic transfer test.

\paragraph{Branch control.} The branch still touches a relevant object, but the next experiment must be narrowed before continuation.

\paragraph{Premature-stop guard.} Do not stop before the branch has produced a concrete source-to-target transfer theorem, an explicit obstruction lemma, or a barrier-aware reason why the route cannot reach the P vs NP/P/poly target.

\paragraph{Payoff criteria.}
\begin{itemize}
\item For this route, state the exact barrier premise it uses or escapes: relativizing, natural, algebrizing, nonuniform, or advice-leaking; promotion requires the escape to be attached to the specific transfer theorem.
\item Record relativization, natural-proofs, algebrization, nonuniformity, and advice-leakage status for the specific transfer theorem.
\item Control gate: name the source carrier or language, target class, and circuit or resource model before claiming progress.
\item Control gate: derive the exact transfer inequality from \(\mathrm{SAT}\in \mathrm{SIZE}(n^{c})\), NP subset P/poly, or the stated target assumption to the source-carrier upper bound.
\item Control gate: record relativization, natural-proofs, algebrization, nonuniformity, and advice-leakage status for the specific theorem.
\end{itemize}

\paragraph{Stop criteria.}
\begin{itemize}
\item The branch cannot identify a source language/carrier and target complexity class in the same model.
\item The route repeatedly proves only a proxy lower bound with no target-preserving transfer.
\item The remaining work is a slogan about hardness magnification rather than a computable transfer or obstruction.
\end{itemize}

\paragraph{Do not stop if.}
\begin{itemize}
\item A transfer exponent, quantifier repair, or barrier status remains uncomputed.
\item The branch can still name a target-preserving carrier and a falsifiable lower-bound threshold.
\item The failure can become a formal no-go lemma for a route family.
\end{itemize}

\paragraph{Required skills.}
\begin{itemize}
\item entailment-gap-auditor
\item kb-end-to-end-proof-critic
\item long-horizon-route-portfolio-manager
\item main-problem-alignment-auditor
\item next-experiment-queue-ranker
\item paper-contribution-readiness-judge
\item pvsnp-barrier-specificity-reviewer
\item pvsnp-conceptual-pressure-localizer
\item pvsnp-impossibility-theorem-forger
\item pvsnp-inference-quality-controller
\item pvsnp-qualitative-complexity-reasoner
\item pvsnp-quantifier-uniformity-reviewer
\item pvsnp-symbolic-complexity-auditor
\end{itemize}

\subsection*{\#159 cycle-0024: Hardness-Magnification Bridge Realization: hitting-cover gate for generated-OR carriers}

\paragraph{Analysis.} Formalize the minimum cover parameter as a bipartite set-cover problem between input fibers and generated SAT labels; test whether known disperser/extractor-style generators can force \(kappa_{N}>N^{1+\epsilon-o(1)}\) without reintroducing nonuniform leakage.

\paragraph{Why this is a good target.} It keeps the barrier and model-boundary audit attached to the mathematical transfer rather than as after-the-fact commentary. It is executable as a bounded complexity calculation, CAS-backed exponent audit, or formal transfer check.

\paragraph{Risk.} High drift risk unless rewritten into a concrete complexity-theoretic transfer test.

\paragraph{Branch control.} The branch still touches a relevant object, but the next experiment must be narrowed before continuation.

\paragraph{Premature-stop guard.} Do not stop before the branch has produced a concrete source-to-target transfer theorem, an explicit obstruction lemma, or a barrier-aware reason why the route cannot reach the P vs NP/P/poly target.

\paragraph{Payoff criteria.}
\begin{itemize}
\item For the high-q generated-OR carrier, compute the exact leakage exponent of the remaining witness-cover or unary-residual mechanism and compare it with 1+epsilon-q.
\item Promotion requires \(q+a>1+\epsilon\) after every projection, padding, and witness-cover leakage term is accounted for; otherwise the branch becomes a bounded-window or no-go result.
\item For this route, state the exact barrier premise it uses or escapes: relativizing, natural, algebrizing, nonuniform, or advice-leaking; promotion requires the escape to be attached to the specific transfer theorem.
\item Record relativization, natural-proofs, algebrization, nonuniformity, and advice-leakage status for the specific transfer theorem.
\item Control gate: name the source carrier or language, target class, and circuit or resource model before claiming progress.
\item Control gate: derive the exact transfer inequality from \(\mathrm{SAT}\in \mathrm{SIZE}(n^{c})\), NP subset P/poly, or the stated target assumption to the source-carrier upper bound.
\end{itemize}

\paragraph{Stop criteria.}
\begin{itemize}
\item The branch cannot identify a source language/carrier and target complexity class in the same model.
\item The route repeatedly proves only a proxy lower bound with no target-preserving transfer.
\item The remaining work is a slogan about hardness magnification rather than a computable transfer or obstruction.
\end{itemize}

\paragraph{Do not stop if.}
\begin{itemize}
\item A transfer exponent, quantifier repair, or barrier status remains uncomputed.
\item The branch can still name a target-preserving carrier and a falsifiable lower-bound threshold.
\item The failure can become a formal no-go lemma for a route family.
\end{itemize}

\paragraph{Required skills.}
\begin{itemize}
\item entailment-gap-auditor
\item kb-end-to-end-proof-critic
\item long-horizon-route-portfolio-manager
\item main-problem-alignment-auditor
\item next-experiment-queue-ranker
\item paper-contribution-readiness-judge
\item pvsnp-barrier-specificity-reviewer
\item pvsnp-conceptual-pressure-localizer
\item pvsnp-impossibility-theorem-forger
\item pvsnp-inference-quality-controller
\item pvsnp-qualitative-complexity-reasoner
\item pvsnp-quantifier-uniformity-reviewer
\item pvsnp-symbolic-complexity-auditor
\end{itemize}

\subsection*{\#160 cycle-0020: Hardness-Magnification Bridge Realization: Semantic Block-Cover Gate}

\paragraph{Analysis.} Instantiate \(M_{N}\) for a nontrivial generator without dominated coordinates, e.g. a universal verifier family whose queries are Cook-Levin encodings of polylog-size computations, and compute whether block predicates admit compression below \(N^{1+\epsilon}\).

\paragraph{Why this is a good target.} It keeps the barrier and model-boundary audit attached to the mathematical transfer rather than as after-the-fact commentary. It is executable as a bounded complexity calculation, CAS-backed exponent audit, or formal transfer check.

\paragraph{Risk.} High drift risk unless rewritten into a concrete complexity-theoretic transfer test.

\paragraph{Branch control.} The branch still touches a relevant object, but the next experiment must be narrowed before continuation.

\paragraph{Premature-stop guard.} Do not stop before the branch has produced a concrete source-to-target transfer theorem, an explicit obstruction lemma, or a barrier-aware reason why the route cannot reach the P vs NP/P/poly target.

\paragraph{Payoff criteria.}
\begin{itemize}
\item For the high-q generated-OR carrier, compute the exact leakage exponent of the remaining witness-cover or unary-residual mechanism and compare it with 1+epsilon-q.
\item Promotion requires \(q+a>1+\epsilon\) after every projection, padding, and witness-cover leakage term is accounted for; otherwise the branch becomes a bounded-window or no-go result.
\item For this route, state the exact barrier premise it uses or escapes: relativizing, natural, algebrizing, nonuniform, or advice-leaking; promotion requires the escape to be attached to the specific transfer theorem.
\item Record relativization, natural-proofs, algebrization, nonuniformity, and advice-leakage status for the specific transfer theorem.
\item Control gate: name the source carrier or language, target class, and circuit or resource model before claiming progress.
\item Control gate: derive the exact transfer inequality from \(\mathrm{SAT}\in \mathrm{SIZE}(n^{c})\), NP subset P/poly, or the stated target assumption to the source-carrier upper bound.
\end{itemize}

\paragraph{Stop criteria.}
\begin{itemize}
\item The branch cannot identify a source language/carrier and target complexity class in the same model.
\item The route repeatedly proves only a proxy lower bound with no target-preserving transfer.
\item The remaining work is a slogan about hardness magnification rather than a computable transfer or obstruction.
\end{itemize}

\paragraph{Do not stop if.}
\begin{itemize}
\item A transfer exponent, quantifier repair, or barrier status remains uncomputed.
\item The branch can still name a target-preserving carrier and a falsifiable lower-bound threshold.
\item The failure can become a formal no-go lemma for a route family.
\end{itemize}

\paragraph{Required skills.}
\begin{itemize}
\item entailment-gap-auditor
\item kb-end-to-end-proof-critic
\item long-horizon-route-portfolio-manager
\item main-problem-alignment-auditor
\item next-experiment-queue-ranker
\item paper-contribution-readiness-judge
\item pvsnp-barrier-specificity-reviewer
\item pvsnp-conceptual-pressure-localizer
\item pvsnp-impossibility-theorem-forger
\item pvsnp-inference-quality-controller
\item pvsnp-qualitative-complexity-reasoner
\item pvsnp-quantifier-uniformity-reviewer
\item pvsnp-symbolic-complexity-auditor
\end{itemize}

\subsection*{\#161 cycle-0151: Hardness-Magnification Bridge Realization}

\paragraph{Analysis.} Write the finite prefix-stable Qrep verifier: fields are batch\_id, rank, source\_row, representative\_row, certificate\_root, dependency list, first\_use\_line, and no\_rebind flag; prove the first rejected field yields the private-row falsifier.

\paragraph{Why this is a good target.} It keeps the barrier and model-boundary audit attached to the mathematical transfer rather than as after-the-fact commentary.

\paragraph{Risk.} High drift risk unless rewritten into a concrete complexity-theoretic transfer test.

\paragraph{Branch control.} The branch still touches a relevant object, but the next experiment must be narrowed before continuation.

\paragraph{Premature-stop guard.} Do not stop before the branch has produced a concrete source-to-target transfer theorem, an explicit obstruction lemma, or a barrier-aware reason why the route cannot reach the P vs NP/P/poly target.

\paragraph{Payoff criteria.}
\begin{itemize}
\item For this route, state the exact barrier premise it uses or escapes: relativizing, natural, algebrizing, nonuniform, or advice-leaking; promotion requires the escape to be attached to the specific transfer theorem.
\item Record relativization, natural-proofs, algebrization, nonuniformity, and advice-leakage status for the specific transfer theorem.
\item Control gate: name the source carrier or language, target class, and circuit or resource model before claiming progress.
\item Control gate: derive the exact transfer inequality from \(\mathrm{SAT}\in \mathrm{SIZE}(n^{c})\), NP subset P/poly, or the stated target assumption to the source-carrier upper bound.
\item Control gate: record relativization, natural-proofs, algebrization, nonuniformity, and advice-leakage status for the specific theorem.
\end{itemize}

\paragraph{Stop criteria.}
\begin{itemize}
\item The branch cannot identify a source language/carrier and target complexity class in the same model.
\item The route repeatedly proves only a proxy lower bound with no target-preserving transfer.
\item The remaining work is a slogan about hardness magnification rather than a computable transfer or obstruction.
\end{itemize}

\paragraph{Do not stop if.}
\begin{itemize}
\item A transfer exponent, quantifier repair, or barrier status remains uncomputed.
\item The branch can still name a target-preserving carrier and a falsifiable lower-bound threshold.
\item The failure can become a formal no-go lemma for a route family.
\end{itemize}

\paragraph{Required skills.}
\begin{itemize}
\item entailment-gap-auditor
\item kb-end-to-end-proof-critic
\item long-horizon-route-portfolio-manager
\item main-problem-alignment-auditor
\item next-experiment-queue-ranker
\item paper-contribution-readiness-judge
\item pvsnp-barrier-specificity-reviewer
\item pvsnp-conceptual-pressure-localizer
\item pvsnp-impossibility-theorem-forger
\item pvsnp-inference-quality-controller
\item pvsnp-qualitative-complexity-reasoner
\item pvsnp-quantifier-uniformity-reviewer
\item pvsnp-symbolic-complexity-auditor
\end{itemize}

\subsection*{\#162 cycle-0119: Hardness-Magnification Bridge Realization, cycle 119: compositional transparent support closure}

\paragraph{Analysis.} Write the finite prefix-closure verifier for \(TransparentRowSupport_{N}\): input pi,N; output \(Cl_{t}\) for each line, Gamma\_prefix(t), and either a valid support-closed derivation certificate or the first unsafe line with its private-row/unread-enable-bit countermodel.

\paragraph{Why this is a good target.} It keeps the barrier and model-boundary audit attached to the mathematical transfer rather than as after-the-fact commentary.

\paragraph{Risk.} High drift risk unless rewritten into a concrete complexity-theoretic transfer test.

\paragraph{Branch control.} The branch still touches a relevant object, but the next experiment must be narrowed before continuation.

\paragraph{Premature-stop guard.} Do not stop before the branch has produced a concrete source-to-target transfer theorem, an explicit obstruction lemma, or a barrier-aware reason why the route cannot reach the P vs NP/P/poly target.

\paragraph{Payoff criteria.}
\begin{itemize}
\item For this route, state the exact barrier premise it uses or escapes: relativizing, natural, algebrizing, nonuniform, or advice-leaking; promotion requires the escape to be attached to the specific transfer theorem.
\item Record relativization, natural-proofs, algebrization, nonuniformity, and advice-leakage status for the specific transfer theorem.
\item Control gate: name the source carrier or language, target class, and circuit or resource model before claiming progress.
\item Control gate: derive the exact transfer inequality from \(\mathrm{SAT}\in \mathrm{SIZE}(n^{c})\), NP subset P/poly, or the stated target assumption to the source-carrier upper bound.
\item Control gate: record relativization, natural-proofs, algebrization, nonuniformity, and advice-leakage status for the specific theorem.
\end{itemize}

\paragraph{Stop criteria.}
\begin{itemize}
\item The branch cannot identify a source language/carrier and target complexity class in the same model.
\item The route repeatedly proves only a proxy lower bound with no target-preserving transfer.
\item The remaining work is a slogan about hardness magnification rather than a computable transfer or obstruction.
\end{itemize}

\paragraph{Do not stop if.}
\begin{itemize}
\item A transfer exponent, quantifier repair, or barrier status remains uncomputed.
\item The branch can still name a target-preserving carrier and a falsifiable lower-bound threshold.
\item The failure can become a formal no-go lemma for a route family.
\end{itemize}

\paragraph{Required skills.}
\begin{itemize}
\item entailment-gap-auditor
\item kb-end-to-end-proof-critic
\item long-horizon-route-portfolio-manager
\item main-problem-alignment-auditor
\item next-experiment-queue-ranker
\item paper-contribution-readiness-judge
\item pvsnp-barrier-specificity-reviewer
\item pvsnp-conceptual-pressure-localizer
\item pvsnp-impossibility-theorem-forger
\item pvsnp-inference-quality-controller
\item pvsnp-qualitative-complexity-reasoner
\item pvsnp-quantifier-uniformity-reviewer
\item pvsnp-symbolic-complexity-auditor
\end{itemize}

\subsection*{\#163 cycle-0101: Hardness-Magnification Bridge Realization}

\paragraph{Analysis.} Write the finite verifier table for one \(SR_{h}\) rule: tag layout, row j encoding, node/depth fields, parent pointers, occurrence-neighborhood list, and the exact charge contributing to B\_SR(pi).

\paragraph{Why this is a good target.} It keeps the barrier and model-boundary audit attached to the mathematical transfer rather than as after-the-fact commentary.

\paragraph{Risk.} High drift risk unless rewritten into a concrete complexity-theoretic transfer test.

\paragraph{Branch control.} The branch still touches a relevant object, but the next experiment must be narrowed before continuation.

\paragraph{Premature-stop guard.} Do not stop before the branch has produced a concrete source-to-target transfer theorem, an explicit obstruction lemma, or a barrier-aware reason why the route cannot reach the P vs NP/P/poly target.

\paragraph{Payoff criteria.}
\begin{itemize}
\item For this route, state the exact barrier premise it uses or escapes: relativizing, natural, algebrizing, nonuniform, or advice-leaking; promotion requires the escape to be attached to the specific transfer theorem.
\item Record relativization, natural-proofs, algebrization, nonuniformity, and advice-leakage status for the specific transfer theorem.
\item Control gate: name the source carrier or language, target class, and circuit or resource model before claiming progress.
\item Control gate: derive the exact transfer inequality from \(\mathrm{SAT}\in \mathrm{SIZE}(n^{c})\), NP subset P/poly, or the stated target assumption to the source-carrier upper bound.
\item Control gate: record relativization, natural-proofs, algebrization, nonuniformity, and advice-leakage status for the specific theorem.
\end{itemize}

\paragraph{Stop criteria.}
\begin{itemize}
\item The branch cannot identify a source language/carrier and target complexity class in the same model.
\item The route repeatedly proves only a proxy lower bound with no target-preserving transfer.
\item The remaining work is a slogan about hardness magnification rather than a computable transfer or obstruction.
\end{itemize}

\paragraph{Do not stop if.}
\begin{itemize}
\item A transfer exponent, quantifier repair, or barrier status remains uncomputed.
\item The branch can still name a target-preserving carrier and a falsifiable lower-bound threshold.
\item The failure can become a formal no-go lemma for a route family.
\end{itemize}

\paragraph{Required skills.}
\begin{itemize}
\item entailment-gap-auditor
\item kb-end-to-end-proof-critic
\item long-horizon-route-portfolio-manager
\item main-problem-alignment-auditor
\item next-experiment-queue-ranker
\item paper-contribution-readiness-judge
\item pvsnp-barrier-specificity-reviewer
\item pvsnp-conceptual-pressure-localizer
\item pvsnp-impossibility-theorem-forger
\item pvsnp-inference-quality-controller
\item pvsnp-qualitative-complexity-reasoner
\item pvsnp-quantifier-uniformity-reviewer
\item pvsnp-symbolic-complexity-auditor
\end{itemize}

\subsection*{\#164 cycle-0099: Hardness-Magnification Bridge Realization}

\paragraph{Analysis.} Write the formal \(FE_{Extract,N}\) blocked-lemma proof: define the quotient-label model after \(SATeq/SR_{h}\) closure, prove the unmentioned row-slot existence bound, construct the two assignments Y,Y', and verify that the carrier row parity changes while the parsed transcript remains identical.

\paragraph{Why this is a good target.} It keeps the barrier and model-boundary audit attached to the mathematical transfer rather than as after-the-fact commentary.

\paragraph{Risk.} High drift risk unless rewritten into a concrete complexity-theoretic transfer test.

\paragraph{Branch control.} The branch still touches a relevant object, but the next experiment must be narrowed before continuation.

\paragraph{Premature-stop guard.} Do not stop before the branch has produced a concrete source-to-target transfer theorem, an explicit obstruction lemma, or a barrier-aware reason why the route cannot reach the P vs NP/P/poly target.

\paragraph{Payoff criteria.}
\begin{itemize}
\item For this route, state the exact barrier premise it uses or escapes: relativizing, natural, algebrizing, nonuniform, or advice-leaking; promotion requires the escape to be attached to the specific transfer theorem.
\item Record relativization, natural-proofs, algebrization, nonuniformity, and advice-leakage status for the specific transfer theorem.
\item Control gate: name the source carrier or language, target class, and circuit or resource model before claiming progress.
\item Control gate: derive the exact transfer inequality from \(\mathrm{SAT}\in \mathrm{SIZE}(n^{c})\), NP subset P/poly, or the stated target assumption to the source-carrier upper bound.
\item Control gate: record relativization, natural-proofs, algebrization, nonuniformity, and advice-leakage status for the specific theorem.
\end{itemize}

\paragraph{Stop criteria.}
\begin{itemize}
\item The branch cannot identify a source language/carrier and target complexity class in the same model.
\item The route repeatedly proves only a proxy lower bound with no target-preserving transfer.
\item The remaining work is a slogan about hardness magnification rather than a computable transfer or obstruction.
\end{itemize}

\paragraph{Do not stop if.}
\begin{itemize}
\item A transfer exponent, quantifier repair, or barrier status remains uncomputed.
\item The branch can still name a target-preserving carrier and a falsifiable lower-bound threshold.
\item The failure can become a formal no-go lemma for a route family.
\end{itemize}

\paragraph{Required skills.}
\begin{itemize}
\item entailment-gap-auditor
\item kb-end-to-end-proof-critic
\item long-horizon-route-portfolio-manager
\item main-problem-alignment-auditor
\item next-experiment-queue-ranker
\item paper-contribution-readiness-judge
\item pvsnp-barrier-specificity-reviewer
\item pvsnp-conceptual-pressure-localizer
\item pvsnp-impossibility-theorem-forger
\item pvsnp-inference-quality-controller
\item pvsnp-qualitative-complexity-reasoner
\item pvsnp-quantifier-uniformity-reviewer
\item pvsnp-symbolic-complexity-auditor
\end{itemize}

\subsection*{\#165 cycle-0031: Hardness-Magnification Bridge Realization: approximate relaxed bag separators}

\paragraph{Analysis.} Try a collision-balanced generator and estimate whether majority/hash-based local predicates yield \(s=N^{o(1)}\) with \(E<N^\epsilon\); if yes, quarantine this carrier family.

\paragraph{Why this is a good target.} It keeps the barrier and model-boundary audit attached to the mathematical transfer rather than as after-the-fact commentary.

\paragraph{Risk.} High drift risk unless rewritten into a concrete complexity-theoretic transfer test.

\paragraph{Branch control.} The branch still touches a relevant object, but the next experiment must be narrowed before continuation.

\paragraph{Premature-stop guard.} Do not stop before the branch has produced a concrete source-to-target transfer theorem, an explicit obstruction lemma, or a barrier-aware reason why the route cannot reach the P vs NP/P/poly target.

\paragraph{Payoff criteria.}
\begin{itemize}
\item For this route, state the exact barrier premise it uses or escapes: relativizing, natural, algebrizing, nonuniform, or advice-leaking; promotion requires the escape to be attached to the specific transfer theorem.
\item Record relativization, natural-proofs, algebrization, nonuniformity, and advice-leakage status for the specific transfer theorem.
\item Control gate: name the source carrier or language, target class, and circuit or resource model before claiming progress.
\item Control gate: derive the exact transfer inequality from \(\mathrm{SAT}\in \mathrm{SIZE}(n^{c})\), NP subset P/poly, or the stated target assumption to the source-carrier upper bound.
\item Control gate: record relativization, natural-proofs, algebrization, nonuniformity, and advice-leakage status for the specific theorem.
\end{itemize}

\paragraph{Stop criteria.}
\begin{itemize}
\item The branch cannot identify a source language/carrier and target complexity class in the same model.
\item The route repeatedly proves only a proxy lower bound with no target-preserving transfer.
\item The remaining work is a slogan about hardness magnification rather than a computable transfer or obstruction.
\end{itemize}

\paragraph{Do not stop if.}
\begin{itemize}
\item A transfer exponent, quantifier repair, or barrier status remains uncomputed.
\item The branch can still name a target-preserving carrier and a falsifiable lower-bound threshold.
\item The failure can become a formal no-go lemma for a route family.
\end{itemize}

\paragraph{Required skills.}
\begin{itemize}
\item entailment-gap-auditor
\item kb-end-to-end-proof-critic
\item long-horizon-route-portfolio-manager
\item main-problem-alignment-auditor
\item next-experiment-queue-ranker
\item paper-contribution-readiness-judge
\item pvsnp-barrier-specificity-reviewer
\item pvsnp-conceptual-pressure-localizer
\item pvsnp-impossibility-theorem-forger
\item pvsnp-inference-quality-controller
\item pvsnp-qualitative-complexity-reasoner
\item pvsnp-quantifier-uniformity-reviewer
\item pvsnp-symbolic-complexity-auditor
\end{itemize}


\section{Qualitative Mechanistic Graph}

\label{sec:report-qualitative-graph}

The qualitative graph is a mechanism map, not a JSON registry. It states how variables, regimes, exact finite tests, and transfer obligations steer branch control.

\subsection*{Mechanism in plain terms}

\begin{enumerate}
\item A candidate bridge starts with a named source carrier, target class, circuit model, and length map.
\item The transfer envelope computes how SAT circuits of size \(n^{c}\) induce source-carrier circuits, usually through an exponent term q+beta c.
\item Positive \(\beta\) or hidden advice can collapse all-exponent consequences into a bounded-window theorem.
\item Barrier checks attach to the actual theorem shape: relativization, naturalness, algebrization, nonuniformity, and advice leakage are not generic slogans.
\item A branch is promoted only when it yields an exact transfer theorem, a scoped no-go, or a finite certificate/soundness obligation that changes route priority.
\end{enumerate}

\subsection*{Relations that matter now}

\subsection*{Quantities and regimes}


\section{Gedanken Experiments}

\label{sec:report-gedanken}

Concrete imagined situations are used only when they isolate a mechanism, expose a hidden assumption, or create a follow-up task. They are conceptual probes, not proof progress.


\section{Parsimonious Probe Control}

\label{sec:report-pmp-probes}

This section keeps symbolic-first computation separate from analytic obligations and rare bounded numeric probes.

No target-compatible parsimonious probe packet is active for this P vs NP report. A future packet should use symbolic complexity auditing: exact exponent envelopes, length maps, advice terms, certificate-search exponents, and barrier predicates.


\section{Scientific Contribution Monitor}

\label{sec:report-contributions}

This section keeps the marathon aimed at independently valuable subgoals and decides when a branch is good enough to hand off instead of continuing indefinitely.

\subsection*{Standing objective}

Every long-horizon branch must aim at a standalone scientific contribution candidate, a contribution-search move, or an explicit demotion test.

\subsection*{Good enough definition}

A candidate is good enough only when its claim, evidence, novelty boundary, validation/disqualifier checks, and overclaim boundary are explicit enough for publishable handoff.

\subsection*{Good enough paper relevance}

A paper subgoal is good enough only when conceptual depth, conceptual novelty, unblocking power, applicability, and understanding gain are at least adequate, the decisive dimensions are strong, and hard gates prevent overclaiming.

\subsection*{Paper relevance marker criteria}

\begin{itemize}
\item applicability: The result is reusable across a nontrivial family of carriers, reductions, restricted models, proof systems, or certificate constructions under explicit assumptions and transfer inequalities.
\item conceptual\_depth: The result works at the level of a real P vs NP proof obstruction or bridge: source carrier, target class, circuit model, reduction/transfer map, quantifier regime, and barrier status are all explicit.
\item conceptual\_novelty: The result adds a new exact obstruction, bridge, no-go, or certificate interface relative to known hardness magnification, MCSP/MKtP, proof-complexity, natural-proofs, relativization, or algebrization families.
\item unblocking\_power: The result changes route priority by proving a transfer, exposing a quantitative loss, ruling out a method class, or identifying the precise bounded calculation that would decide the branch.
\item understanding\_gain: The result clarifies why the route is hard: where proof debt moves, which exponent or advice term destroys the transfer, or which barrier premise actually binds.
\end{itemize}

\subsection*{Stepping-stone maturity levels}

\begin{itemize}
\item contribution\_shaped: The branch has an independently meaningful object or obstruction, but it is not yet a publishable stepping stone.
\item near\_stepping\_stone: One or two bounded prior-art, disqualifier, proof-shape, or validation tasks could make the candidate paper-note ready.
\item publishable\_handoff\_ready: Broad exploration should pause; the candidate should move to validation, prior-art, and manuscript packaging.
\end{itemize}

\subsection*{Anti-infinite progression rule}

A branch may continue only while the next cycle can improve a contribution metric or resolve a disqualifier; otherwise adapt, backtrack, quarantine, stop, or hand off.

\subsection*{Next contribution objective}
Run the best near-stepping-stone candidate through prior-art, disqualifier, or paper-shape certification before broad exploration.

\subsection*{cycle-0124: Hardness-Magnification Bridge Realization}

\paragraph{Contribution type.} quantitative\_transfer\_obstruction

\paragraph{Stepping-stone maturity.} near\_stepping\_stone

\paragraph{Smallest honest claim.} A route-specific quantitative obstruction can be isolated for Hardness-Magnification Bridge Realization: the transfer formula names the source exponent, overhead, target-query exponent, and exact threshold where magnification collapses.

\paragraph{Why it could stand alone.} A quantitative transfer obstruction is independently useful because it replaces barrier folklore with an exact parameter condition showing when a proof route loses the target exponent.

\paragraph{Why this budget now.} Run the next bounded certification or prior-art gate before another broad marathon branch.

\paragraph{Branch control.} The candidate is strong but still needs bounded validation or disqualifier checks before handoff.

\paragraph{Overclaim boundary.} This candidate is a scoped contribution candidate; it does not prove the main target unless all target commitments are discharged in the same theory.

\paragraph{Good-enough paper relevance.} The candidate is scientifically relevant enough to shape a paper, but certification/prior-art gates still control handoff.

\paragraph{Paper relevance status.} strong\_relevance\_certification\_pending with average marker score 4.8

\paragraph{Good enough when.}
\begin{itemize}
\item Smallest honest claim can be stated without proving the main target.
\item Mathematical object, assumptions, scope, and evidence basis are explicit.
\item At least one disqualifier or falsifier has been checked or bounded.
\item Remaining validation is finite enough for a paper note, obstruction note, benchmark, method note, or lemma packet.
\end{itemize}

\paragraph{Continue only if.}
\begin{itemize}
\item The next cycle can improve novelty, explanation, implications, validation readiness, target relevance, or disqualifier resolution.
\item The branch still touches a core target commitment or a reusable scientific object.
\item The next action is a bounded calculation, derivation, source check, or falsifier test rather than generic elaboration.
\end{itemize}

\paragraph{Stop when.}
\begin{itemize}
\item Two explicit adaptation or backtracking attempts fail to improve any contribution metric.
\item The branch cannot name a contribution candidate, contribution-disqualifying test, or target-preserving validation handoff.
\item The candidate is already good enough for handoff and further exploration would only inflate proof debt.
\end{itemize}

\paragraph{Validation burden.}
\begin{itemize}
\item Run prior-art comparison before claiming novelty.
\end{itemize}

\paragraph{Evidence basis.}
\begin{itemize}
\item cycle:cycle-0124
\item branch:Hardness-Magnification Bridge Realization
\item equations:6
\item \(derivation_{steps}:6\)
\item falsifiers:6
\end{itemize}

\paragraph{Missing to publishable.}
\begin{itemize}
\item Run the finite prior-art boundary check and record closest known results.
\item Resolve disqualifier tests or convert unknowns into finite validation tasks.
\end{itemize}

\paragraph{Missing to good-enough paper.}
\begin{itemize}
\item Close the publishable handoff hard gates before calling the candidate paper-ready.
\end{itemize}

\paragraph{Paper relevance markers.}
\begin{itemize}
\item conceptual\_depth: \(score=5\), \(status=strong\) - The result works at the level of a real P vs NP proof obstruction or bridge: source carrier, target class, circuit model, reduction/transfer map, quantifier regime, and barrier status are all explicit.
\item conceptual\_novelty: \(score=4\), \(status=strong\) - The result adds a new exact obstruction, bridge, no-go, or certificate interface relative to known hardness magnification, MCSP/MKtP, proof-complexity, natural-proofs, relativization, or algebrization families.
\item unblocking\_power: \(score=5\), \(status=strong\) - The result changes route priority by proving a transfer, exposing a quantitative loss, ruling out a method class, or identifying the precise bounded calculation that would decide the branch.
\item applicability: \(score=5\), \(status=strong\) - The result is reusable across a nontrivial family of carriers, reductions, restricted models, proof systems, or certificate constructions under explicit assumptions and transfer inequalities.
\item understanding\_gain: \(score=5\), \(status=strong\) - The result clarifies why the route is hard: where proof debt moves, which exponent or advice term destroys the transfer, or which barrier premise actually binds.
\end{itemize}

\paragraph{Hard gates.}
\begin{itemize}
\item claim\_statement\_ready: \(ready=True\), \(finite_{task}=False\)
\item evidence\_ready: \(ready=True\), \(finite_{task}=True\)
\item novelty\_boundary\_ready: \(ready=False\), \(finite_{task}=True\)
\item overclaim\_boundary\_ready: \(ready=True\), \(finite_{task}=False\)
\item validation\_ready: \(ready=False\), \(finite_{task}=True\)
\end{itemize}

\paragraph{Disqualifiers.}
\begin{itemize}
\item already\_known: finite\_check\_queued - Prior-art boundary is queued before publishable handoff.
\item evidence\_gap: pass - Certification readiness is derived from equations, derivation steps, and falsifiers.
\item overclaim\_risk: pass - Overclaim boundary explicitly states that the candidate does not prove the main target.
\item validation\_unbounded: pass - The next action is checked for bounded proof, calculation, audit, or falsifier language.
\item narrow\_bookkeeping: pass - The candidate is treated as scientific only if it has external mathematical content.
\item target\_equivalent\_assumption: finite\_check\_queued - P vs NP candidates must show that the source obligation is not already the desired lower bound.
\item advice\_leakage: finite\_check\_queued - Advice and nonuniform hard-row selection are explicit P vs NP disqualifiers.
\item barrier\_status: finite\_check\_queued - Barrier status must be attached to the specific transfer theorem.
\item restricted\_model\_proxy\_drift: finite\_check\_queued - Restricted-model results need a compiler envelope or explicit proxy boundary.
\end{itemize}

\paragraph{Prior-art boundary.} finite\_check\_queued

\begin{itemize}
\item Run prior-art pass for hardness magnification threshold theorems.
\item Run prior-art pass for natural proofs and algebrization barriers.
\item Run prior-art pass for padding/projection lower-bound transfer arguments.
\end{itemize}

\subsection*{cycle-0121: Hardness-Magnification Bridge Realization: L121 least-closed support audit}

\paragraph{Contribution type.} quantitative\_transfer\_obstruction

\paragraph{Stepping-stone maturity.} near\_stepping\_stone

\paragraph{Smallest honest claim.} A route-specific quantitative obstruction can be isolated for Hardness-Magnification Bridge Realization: L121 least-closed support audit: the transfer formula names the source exponent, overhead, target-query exponent, and exact threshold where magnification collapses.

\paragraph{Why it could stand alone.} A quantitative transfer obstruction is independently useful because it replaces barrier folklore with an exact parameter condition showing when a proof route loses the target exponent.

\paragraph{Why this budget now.} Run the next bounded certification or prior-art gate before another broad marathon branch.

\paragraph{Branch control.} The candidate is strong but still needs bounded validation or disqualifier checks before handoff.

\paragraph{Overclaim boundary.} This candidate is a scoped contribution candidate; it does not prove the main target unless all target commitments are discharged in the same theory.

\paragraph{Good-enough paper relevance.} The candidate is scientifically relevant enough to shape a paper, but certification/prior-art gates still control handoff.

\paragraph{Paper relevance status.} strong\_relevance\_certification\_pending with average marker score 4.8

\paragraph{Good enough when.}
\begin{itemize}
\item Smallest honest claim can be stated without proving the main target.
\item Mathematical object, assumptions, scope, and evidence basis are explicit.
\item At least one disqualifier or falsifier has been checked or bounded.
\item Remaining validation is finite enough for a paper note, obstruction note, benchmark, method note, or lemma packet.
\end{itemize}

\paragraph{Continue only if.}
\begin{itemize}
\item The next cycle can improve novelty, explanation, implications, validation readiness, target relevance, or disqualifier resolution.
\item The branch still touches a core target commitment or a reusable scientific object.
\item The next action is a bounded calculation, derivation, source check, or falsifier test rather than generic elaboration.
\end{itemize}

\paragraph{Stop when.}
\begin{itemize}
\item Two explicit adaptation or backtracking attempts fail to improve any contribution metric.
\item The branch cannot name a contribution candidate, contribution-disqualifying test, or target-preserving validation handoff.
\item The candidate is already good enough for handoff and further exploration would only inflate proof debt.
\end{itemize}

\paragraph{Validation burden.}
\begin{itemize}
\item Run prior-art comparison before claiming novelty.
\end{itemize}

\paragraph{Evidence basis.}
\begin{itemize}
\item cycle:cycle-0121
\item branch:Hardness-Magnification Bridge Realization: L121 least-closed support audit
\item equations:5
\item \(derivation_{steps}:10\)
\item falsifiers:6
\end{itemize}

\paragraph{Missing to publishable.}
\begin{itemize}
\item Run the finite prior-art boundary check and record closest known results.
\item Resolve disqualifier tests or convert unknowns into finite validation tasks.
\end{itemize}

\paragraph{Missing to good-enough paper.}
\begin{itemize}
\item Close the publishable handoff hard gates before calling the candidate paper-ready.
\end{itemize}

\paragraph{Paper relevance markers.}
\begin{itemize}
\item conceptual\_depth: \(score=5\), \(status=strong\) - The result works at the level of a real P vs NP proof obstruction or bridge: source carrier, target class, circuit model, reduction/transfer map, quantifier regime, and barrier status are all explicit.
\item conceptual\_novelty: \(score=4\), \(status=strong\) - The result adds a new exact obstruction, bridge, no-go, or certificate interface relative to known hardness magnification, MCSP/MKtP, proof-complexity, natural-proofs, relativization, or algebrization families.
\item unblocking\_power: \(score=5\), \(status=strong\) - The result changes route priority by proving a transfer, exposing a quantitative loss, ruling out a method class, or identifying the precise bounded calculation that would decide the branch.
\item applicability: \(score=5\), \(status=strong\) - The result is reusable across a nontrivial family of carriers, reductions, restricted models, proof systems, or certificate constructions under explicit assumptions and transfer inequalities.
\item understanding\_gain: \(score=5\), \(status=strong\) - The result clarifies why the route is hard: where proof debt moves, which exponent or advice term destroys the transfer, or which barrier premise actually binds.
\end{itemize}

\paragraph{Hard gates.}
\begin{itemize}
\item claim\_statement\_ready: \(ready=True\), \(finite_{task}=False\)
\item evidence\_ready: \(ready=True\), \(finite_{task}=True\)
\item novelty\_boundary\_ready: \(ready=False\), \(finite_{task}=True\)
\item overclaim\_boundary\_ready: \(ready=True\), \(finite_{task}=False\)
\item validation\_ready: \(ready=False\), \(finite_{task}=True\)
\end{itemize}

\paragraph{Disqualifiers.}
\begin{itemize}
\item already\_known: finite\_check\_queued - Prior-art boundary is queued before publishable handoff.
\item evidence\_gap: pass - Certification readiness is derived from equations, derivation steps, and falsifiers.
\item overclaim\_risk: pass - Overclaim boundary explicitly states that the candidate does not prove the main target.
\item validation\_unbounded: pass - The next action is checked for bounded proof, calculation, audit, or falsifier language.
\item narrow\_bookkeeping: pass - The candidate is treated as scientific only if it has external mathematical content.
\item target\_equivalent\_assumption: finite\_check\_queued - P vs NP candidates must show that the source obligation is not already the desired lower bound.
\item advice\_leakage: finite\_check\_queued - Advice and nonuniform hard-row selection are explicit P vs NP disqualifiers.
\item barrier\_status: finite\_check\_queued - Barrier status must be attached to the specific transfer theorem.
\item restricted\_model\_proxy\_drift: finite\_check\_queued - Restricted-model results need a compiler envelope or explicit proxy boundary.
\end{itemize}

\paragraph{Prior-art boundary.} finite\_check\_queued

\begin{itemize}
\item Run prior-art pass for hardness magnification threshold theorems.
\item Run prior-art pass for natural proofs and algebrization barriers.
\item Run prior-art pass for padding/projection lower-bound transfer arguments.
\end{itemize}

\subsection*{cycle-0073: Hardness-Magnification Bridge Realization}

\paragraph{Contribution type.} quantitative\_transfer\_obstruction

\paragraph{Stepping-stone maturity.} near\_stepping\_stone

\paragraph{Smallest honest claim.} A route-specific quantitative obstruction can be isolated for Hardness-Magnification Bridge Realization: the transfer formula names the source exponent, overhead, target-query exponent, and exact threshold where magnification collapses.

\paragraph{Why it could stand alone.} A quantitative transfer obstruction is independently useful because it replaces barrier folklore with an exact parameter condition showing when a proof route loses the target exponent.

\paragraph{Why this budget now.} Run the next bounded certification or prior-art gate before another broad marathon branch.

\paragraph{Branch control.} The candidate is strong but still needs bounded validation or disqualifier checks before handoff.

\paragraph{Overclaim boundary.} This candidate is a scoped contribution candidate; it does not prove the main target unless all target commitments are discharged in the same theory.

\paragraph{Good-enough paper relevance.} The candidate is scientifically relevant enough to shape a paper, but certification/prior-art gates still control handoff.

\paragraph{Paper relevance status.} strong\_relevance\_certification\_pending with average marker score 4.8

\paragraph{Good enough when.}
\begin{itemize}
\item Smallest honest claim can be stated without proving the main target.
\item Mathematical object, assumptions, scope, and evidence basis are explicit.
\item At least one disqualifier or falsifier has been checked or bounded.
\item Remaining validation is finite enough for a paper note, obstruction note, benchmark, method note, or lemma packet.
\end{itemize}

\paragraph{Continue only if.}
\begin{itemize}
\item The next cycle can improve novelty, explanation, implications, validation readiness, target relevance, or disqualifier resolution.
\item The branch still touches a core target commitment or a reusable scientific object.
\item The next action is a bounded calculation, derivation, source check, or falsifier test rather than generic elaboration.
\end{itemize}

\paragraph{Stop when.}
\begin{itemize}
\item Two explicit adaptation or backtracking attempts fail to improve any contribution metric.
\item The branch cannot name a contribution candidate, contribution-disqualifying test, or target-preserving validation handoff.
\item The candidate is already good enough for handoff and further exploration would only inflate proof debt.
\end{itemize}

\paragraph{Validation burden.}
\begin{itemize}
\item Run prior-art comparison before claiming novelty.
\end{itemize}

\paragraph{Evidence basis.}
\begin{itemize}
\item cycle:cycle-0073
\item branch:Hardness-Magnification Bridge Realization
\item equations:4
\item \(derivation_{steps}:7\)
\item falsifiers:5
\end{itemize}

\paragraph{Missing to publishable.}
\begin{itemize}
\item Run the finite prior-art boundary check and record closest known results.
\item Resolve disqualifier tests or convert unknowns into finite validation tasks.
\end{itemize}

\paragraph{Missing to good-enough paper.}
\begin{itemize}
\item Close the publishable handoff hard gates before calling the candidate paper-ready.
\end{itemize}

\paragraph{Paper relevance markers.}
\begin{itemize}
\item conceptual\_depth: \(score=5\), \(status=strong\) - The result works at the level of a real P vs NP proof obstruction or bridge: source carrier, target class, circuit model, reduction/transfer map, quantifier regime, and barrier status are all explicit.
\item conceptual\_novelty: \(score=4\), \(status=strong\) - The result adds a new exact obstruction, bridge, no-go, or certificate interface relative to known hardness magnification, MCSP/MKtP, proof-complexity, natural-proofs, relativization, or algebrization families.
\item unblocking\_power: \(score=5\), \(status=strong\) - The result changes route priority by proving a transfer, exposing a quantitative loss, ruling out a method class, or identifying the precise bounded calculation that would decide the branch.
\item applicability: \(score=5\), \(status=strong\) - The result is reusable across a nontrivial family of carriers, reductions, restricted models, proof systems, or certificate constructions under explicit assumptions and transfer inequalities.
\item understanding\_gain: \(score=5\), \(status=strong\) - The result clarifies why the route is hard: where proof debt moves, which exponent or advice term destroys the transfer, or which barrier premise actually binds.
\end{itemize}

\paragraph{Hard gates.}
\begin{itemize}
\item claim\_statement\_ready: \(ready=True\), \(finite_{task}=False\)
\item evidence\_ready: \(ready=True\), \(finite_{task}=True\)
\item novelty\_boundary\_ready: \(ready=False\), \(finite_{task}=True\)
\item overclaim\_boundary\_ready: \(ready=True\), \(finite_{task}=False\)
\item validation\_ready: \(ready=False\), \(finite_{task}=True\)
\end{itemize}

\paragraph{Disqualifiers.}
\begin{itemize}
\item already\_known: finite\_check\_queued - Prior-art boundary is queued before publishable handoff.
\item evidence\_gap: pass - Certification readiness is derived from equations, derivation steps, and falsifiers.
\item overclaim\_risk: pass - Overclaim boundary explicitly states that the candidate does not prove the main target.
\item validation\_unbounded: pass - The next action is checked for bounded proof, calculation, audit, or falsifier language.
\item narrow\_bookkeeping: pass - The candidate is treated as scientific only if it has external mathematical content.
\item target\_equivalent\_assumption: finite\_check\_queued - P vs NP candidates must show that the source obligation is not already the desired lower bound.
\item advice\_leakage: finite\_check\_queued - Advice and nonuniform hard-row selection are explicit P vs NP disqualifiers.
\item barrier\_status: finite\_check\_queued - Barrier status must be attached to the specific transfer theorem.
\item restricted\_model\_proxy\_drift: pass - Restricted-model results need a compiler envelope or explicit proxy boundary.
\end{itemize}

\paragraph{Prior-art boundary.} finite\_check\_queued

\begin{itemize}
\item Run prior-art pass for hardness magnification threshold theorems.
\item Run prior-art pass for natural proofs and algebrization barriers.
\item Run prior-art pass for padding/projection lower-bound transfer arguments.
\end{itemize}

\subsection*{cycle-0071: Hardness-Magnification Bridge Realization}

\paragraph{Contribution type.} quantitative\_transfer\_obstruction

\paragraph{Stepping-stone maturity.} near\_stepping\_stone

\paragraph{Smallest honest claim.} A route-specific quantitative obstruction can be isolated for Hardness-Magnification Bridge Realization: the transfer formula names the source exponent, overhead, target-query exponent, and exact threshold where magnification collapses.

\paragraph{Why it could stand alone.} A quantitative transfer obstruction is independently useful because it replaces barrier folklore with an exact parameter condition showing when a proof route loses the target exponent.

\paragraph{Why this budget now.} Run the next bounded certification or prior-art gate before another broad marathon branch.

\paragraph{Branch control.} The candidate is strong but still needs bounded validation or disqualifier checks before handoff.

\paragraph{Overclaim boundary.} This candidate is a scoped contribution candidate; it does not prove the main target unless all target commitments are discharged in the same theory.

\paragraph{Good-enough paper relevance.} The candidate is scientifically relevant enough to shape a paper, but certification/prior-art gates still control handoff.

\paragraph{Paper relevance status.} strong\_relevance\_certification\_pending with average marker score 4.8

\paragraph{Good enough when.}
\begin{itemize}
\item Smallest honest claim can be stated without proving the main target.
\item Mathematical object, assumptions, scope, and evidence basis are explicit.
\item At least one disqualifier or falsifier has been checked or bounded.
\item Remaining validation is finite enough for a paper note, obstruction note, benchmark, method note, or lemma packet.
\end{itemize}

\paragraph{Continue only if.}
\begin{itemize}
\item The next cycle can improve novelty, explanation, implications, validation readiness, target relevance, or disqualifier resolution.
\item The branch still touches a core target commitment or a reusable scientific object.
\item The next action is a bounded calculation, derivation, source check, or falsifier test rather than generic elaboration.
\end{itemize}

\paragraph{Stop when.}
\begin{itemize}
\item Two explicit adaptation or backtracking attempts fail to improve any contribution metric.
\item The branch cannot name a contribution candidate, contribution-disqualifying test, or target-preserving validation handoff.
\item The candidate is already good enough for handoff and further exploration would only inflate proof debt.
\end{itemize}

\paragraph{Validation burden.}
\begin{itemize}
\item Run prior-art comparison before claiming novelty.
\end{itemize}

\paragraph{Evidence basis.}
\begin{itemize}
\item cycle:cycle-0071
\item branch:Hardness-Magnification Bridge Realization
\item equations:5
\item \(derivation_{steps}:7\)
\item falsifiers:5
\end{itemize}

\paragraph{Missing to publishable.}
\begin{itemize}
\item Run the finite prior-art boundary check and record closest known results.
\item Resolve disqualifier tests or convert unknowns into finite validation tasks.
\end{itemize}

\paragraph{Missing to good-enough paper.}
\begin{itemize}
\item Close the publishable handoff hard gates before calling the candidate paper-ready.
\end{itemize}

\paragraph{Paper relevance markers.}
\begin{itemize}
\item conceptual\_depth: \(score=5\), \(status=strong\) - The result works at the level of a real P vs NP proof obstruction or bridge: source carrier, target class, circuit model, reduction/transfer map, quantifier regime, and barrier status are all explicit.
\item conceptual\_novelty: \(score=4\), \(status=strong\) - The result adds a new exact obstruction, bridge, no-go, or certificate interface relative to known hardness magnification, MCSP/MKtP, proof-complexity, natural-proofs, relativization, or algebrization families.
\item unblocking\_power: \(score=5\), \(status=strong\) - The result changes route priority by proving a transfer, exposing a quantitative loss, ruling out a method class, or identifying the precise bounded calculation that would decide the branch.
\item applicability: \(score=5\), \(status=strong\) - The result is reusable across a nontrivial family of carriers, reductions, restricted models, proof systems, or certificate constructions under explicit assumptions and transfer inequalities.
\item understanding\_gain: \(score=5\), \(status=strong\) - The result clarifies why the route is hard: where proof debt moves, which exponent or advice term destroys the transfer, or which barrier premise actually binds.
\end{itemize}

\paragraph{Hard gates.}
\begin{itemize}
\item claim\_statement\_ready: \(ready=True\), \(finite_{task}=False\)
\item evidence\_ready: \(ready=True\), \(finite_{task}=True\)
\item novelty\_boundary\_ready: \(ready=False\), \(finite_{task}=True\)
\item overclaim\_boundary\_ready: \(ready=True\), \(finite_{task}=False\)
\item validation\_ready: \(ready=False\), \(finite_{task}=True\)
\end{itemize}

\paragraph{Disqualifiers.}
\begin{itemize}
\item already\_known: finite\_check\_queued - Prior-art boundary is queued before publishable handoff.
\item evidence\_gap: pass - Certification readiness is derived from equations, derivation steps, and falsifiers.
\item overclaim\_risk: pass - Overclaim boundary explicitly states that the candidate does not prove the main target.
\item validation\_unbounded: pass - The next action is checked for bounded proof, calculation, audit, or falsifier language.
\item narrow\_bookkeeping: pass - The candidate is treated as scientific only if it has external mathematical content.
\item target\_equivalent\_assumption: finite\_check\_queued - P vs NP candidates must show that the source obligation is not already the desired lower bound.
\item advice\_leakage: finite\_check\_queued - Advice and nonuniform hard-row selection are explicit P vs NP disqualifiers.
\item barrier\_status: finite\_check\_queued - Barrier status must be attached to the specific transfer theorem.
\item restricted\_model\_proxy\_drift: pass - Restricted-model results need a compiler envelope or explicit proxy boundary.
\end{itemize}

\paragraph{Prior-art boundary.} finite\_check\_queued

\begin{itemize}
\item Run prior-art pass for hardness magnification threshold theorems.
\item Run prior-art pass for natural proofs and algebrization barriers.
\item Run prior-art pass for padding/projection lower-bound transfer arguments.
\end{itemize}

\subsection*{cycle-0065: Hardness-Magnification Bridge Realization}

\paragraph{Contribution type.} quantitative\_transfer\_obstruction

\paragraph{Stepping-stone maturity.} near\_stepping\_stone

\paragraph{Smallest honest claim.} A route-specific quantitative obstruction can be isolated for Hardness-Magnification Bridge Realization: the transfer formula names the source exponent, overhead, target-query exponent, and exact threshold where magnification collapses.

\paragraph{Why it could stand alone.} A quantitative transfer obstruction is independently useful because it replaces barrier folklore with an exact parameter condition showing when a proof route loses the target exponent.

\paragraph{Why this budget now.} Run the next bounded certification or prior-art gate before another broad marathon branch.

\paragraph{Branch control.} The candidate is strong but still needs bounded validation or disqualifier checks before handoff.

\paragraph{Overclaim boundary.} This candidate is a scoped contribution candidate; it does not prove the main target unless all target commitments are discharged in the same theory.

\paragraph{Good-enough paper relevance.} The candidate is scientifically relevant enough to shape a paper, but certification/prior-art gates still control handoff.

\paragraph{Paper relevance status.} strong\_relevance\_certification\_pending with average marker score 4.8

\paragraph{Good enough when.}
\begin{itemize}
\item Smallest honest claim can be stated without proving the main target.
\item Mathematical object, assumptions, scope, and evidence basis are explicit.
\item At least one disqualifier or falsifier has been checked or bounded.
\item Remaining validation is finite enough for a paper note, obstruction note, benchmark, method note, or lemma packet.
\end{itemize}

\paragraph{Continue only if.}
\begin{itemize}
\item The next cycle can improve novelty, explanation, implications, validation readiness, target relevance, or disqualifier resolution.
\item The branch still touches a core target commitment or a reusable scientific object.
\item The next action is a bounded calculation, derivation, source check, or falsifier test rather than generic elaboration.
\end{itemize}

\paragraph{Stop when.}
\begin{itemize}
\item Two explicit adaptation or backtracking attempts fail to improve any contribution metric.
\item The branch cannot name a contribution candidate, contribution-disqualifying test, or target-preserving validation handoff.
\item The candidate is already good enough for handoff and further exploration would only inflate proof debt.
\end{itemize}

\paragraph{Validation burden.}
\begin{itemize}
\item Run prior-art comparison before claiming novelty.
\end{itemize}

\paragraph{Evidence basis.}
\begin{itemize}
\item cycle:cycle-0065
\item branch:Hardness-Magnification Bridge Realization
\item equations:4
\item \(derivation_{steps}:7\)
\item falsifiers:4
\end{itemize}

\paragraph{Missing to publishable.}
\begin{itemize}
\item Run the finite prior-art boundary check and record closest known results.
\item Resolve disqualifier tests or convert unknowns into finite validation tasks.
\end{itemize}

\paragraph{Missing to good-enough paper.}
\begin{itemize}
\item Close the publishable handoff hard gates before calling the candidate paper-ready.
\end{itemize}

\paragraph{Paper relevance markers.}
\begin{itemize}
\item conceptual\_depth: \(score=5\), \(status=strong\) - The result works at the level of a real P vs NP proof obstruction or bridge: source carrier, target class, circuit model, reduction/transfer map, quantifier regime, and barrier status are all explicit.
\item conceptual\_novelty: \(score=4\), \(status=strong\) - The result adds a new exact obstruction, bridge, no-go, or certificate interface relative to known hardness magnification, MCSP/MKtP, proof-complexity, natural-proofs, relativization, or algebrization families.
\item unblocking\_power: \(score=5\), \(status=strong\) - The result changes route priority by proving a transfer, exposing a quantitative loss, ruling out a method class, or identifying the precise bounded calculation that would decide the branch.
\item applicability: \(score=5\), \(status=strong\) - The result is reusable across a nontrivial family of carriers, reductions, restricted models, proof systems, or certificate constructions under explicit assumptions and transfer inequalities.
\item understanding\_gain: \(score=5\), \(status=strong\) - The result clarifies why the route is hard: where proof debt moves, which exponent or advice term destroys the transfer, or which barrier premise actually binds.
\end{itemize}

\paragraph{Hard gates.}
\begin{itemize}
\item claim\_statement\_ready: \(ready=True\), \(finite_{task}=False\)
\item evidence\_ready: \(ready=True\), \(finite_{task}=True\)
\item novelty\_boundary\_ready: \(ready=False\), \(finite_{task}=True\)
\item overclaim\_boundary\_ready: \(ready=True\), \(finite_{task}=False\)
\item validation\_ready: \(ready=False\), \(finite_{task}=True\)
\end{itemize}

\paragraph{Disqualifiers.}
\begin{itemize}
\item already\_known: finite\_check\_queued - Prior-art boundary is queued before publishable handoff.
\item evidence\_gap: pass - Certification readiness is derived from equations, derivation steps, and falsifiers.
\item overclaim\_risk: pass - Overclaim boundary explicitly states that the candidate does not prove the main target.
\item validation\_unbounded: pass - The next action is checked for bounded proof, calculation, audit, or falsifier language.
\item narrow\_bookkeeping: pass - The candidate is treated as scientific only if it has external mathematical content.
\item target\_equivalent\_assumption: finite\_check\_queued - P vs NP candidates must show that the source obligation is not already the desired lower bound.
\item advice\_leakage: finite\_check\_queued - Advice and nonuniform hard-row selection are explicit P vs NP disqualifiers.
\item barrier\_status: finite\_check\_queued - Barrier status must be attached to the specific transfer theorem.
\item restricted\_model\_proxy\_drift: pass - Restricted-model results need a compiler envelope or explicit proxy boundary.
\end{itemize}

\paragraph{Prior-art boundary.} finite\_check\_queued

\begin{itemize}
\item Run prior-art pass for hardness magnification threshold theorems.
\item Run prior-art pass for natural proofs and algebrization barriers.
\item Run prior-art pass for padding/projection lower-bound transfer arguments.
\end{itemize}

\subsection*{cycle-0008: Hardness-Magnification Bridge Realization, cycle 8: subpolynomial-query threshold repair}

\paragraph{Contribution type.} quantitative\_transfer\_obstruction

\paragraph{Stepping-stone maturity.} near\_stepping\_stone

\paragraph{Smallest honest claim.} A route-specific quantitative obstruction can be isolated for Hardness-Magnification Bridge Realization, cycle 8: subpolynomial-query threshold repair: the transfer formula names the source exponent, overhead, target-query exponent, and exact threshold where magnification collapses.

\paragraph{Why it could stand alone.} A quantitative transfer obstruction is independently useful because it replaces barrier folklore with an exact parameter condition showing when a proof route loses the target exponent.

\paragraph{Why this budget now.} Run the next bounded certification or prior-art gate before another broad marathon branch.

\paragraph{Branch control.} The candidate is strong but still needs bounded validation or disqualifier checks before handoff.

\paragraph{Overclaim boundary.} This candidate is a scoped contribution candidate; it does not prove the main target unless all target commitments are discharged in the same theory.

\paragraph{Good-enough paper relevance.} The candidate is scientifically relevant enough to shape a paper, but certification/prior-art gates still control handoff.

\paragraph{Paper relevance status.} strong\_relevance\_certification\_pending with average marker score 4.8

\paragraph{Good enough when.}
\begin{itemize}
\item Smallest honest claim can be stated without proving the main target.
\item Mathematical object, assumptions, scope, and evidence basis are explicit.
\item At least one disqualifier or falsifier has been checked or bounded.
\item Remaining validation is finite enough for a paper note, obstruction note, benchmark, method note, or lemma packet.
\end{itemize}

\paragraph{Continue only if.}
\begin{itemize}
\item The next cycle can improve novelty, explanation, implications, validation readiness, target relevance, or disqualifier resolution.
\item The branch still touches a core target commitment or a reusable scientific object.
\item The next action is a bounded calculation, derivation, source check, or falsifier test rather than generic elaboration.
\end{itemize}

\paragraph{Stop when.}
\begin{itemize}
\item Two explicit adaptation or backtracking attempts fail to improve any contribution metric.
\item The branch cannot name a contribution candidate, contribution-disqualifying test, or target-preserving validation handoff.
\item The candidate is already good enough for handoff and further exploration would only inflate proof debt.
\end{itemize}

\paragraph{Validation burden.}
\begin{itemize}
\item Run prior-art comparison before claiming novelty.
\end{itemize}

\paragraph{Evidence basis.}
\begin{itemize}
\item cycle:cycle-0008
\item branch:Hardness-Magnification Bridge Realization, cycle 8: subpolynomial-query threshold repair
\item equations:4
\item \(derivation_{steps}:7\)
\item falsifiers:6
\end{itemize}

\paragraph{Missing to publishable.}
\begin{itemize}
\item Run the finite prior-art boundary check and record closest known results.
\item Resolve disqualifier tests or convert unknowns into finite validation tasks.
\end{itemize}

\paragraph{Missing to good-enough paper.}
\begin{itemize}
\item Close the publishable handoff hard gates before calling the candidate paper-ready.
\end{itemize}

\paragraph{Paper relevance markers.}
\begin{itemize}
\item conceptual\_depth: \(score=5\), \(status=strong\) - The result works at the level of a real P vs NP proof obstruction or bridge: source carrier, target class, circuit model, reduction/transfer map, quantifier regime, and barrier status are all explicit.
\item conceptual\_novelty: \(score=4\), \(status=strong\) - The result adds a new exact obstruction, bridge, no-go, or certificate interface relative to known hardness magnification, MCSP/MKtP, proof-complexity, natural-proofs, relativization, or algebrization families.
\item unblocking\_power: \(score=5\), \(status=strong\) - The result changes route priority by proving a transfer, exposing a quantitative loss, ruling out a method class, or identifying the precise bounded calculation that would decide the branch.
\item applicability: \(score=5\), \(status=strong\) - The result is reusable across a nontrivial family of carriers, reductions, restricted models, proof systems, or certificate constructions under explicit assumptions and transfer inequalities.
\item understanding\_gain: \(score=5\), \(status=strong\) - The result clarifies why the route is hard: where proof debt moves, which exponent or advice term destroys the transfer, or which barrier premise actually binds.
\end{itemize}

\paragraph{Hard gates.}
\begin{itemize}
\item claim\_statement\_ready: \(ready=True\), \(finite_{task}=False\)
\item evidence\_ready: \(ready=True\), \(finite_{task}=True\)
\item novelty\_boundary\_ready: \(ready=False\), \(finite_{task}=True\)
\item overclaim\_boundary\_ready: \(ready=True\), \(finite_{task}=False\)
\item validation\_ready: \(ready=False\), \(finite_{task}=True\)
\end{itemize}

\paragraph{Disqualifiers.}
\begin{itemize}
\item already\_known: finite\_check\_queued - Prior-art boundary is queued before publishable handoff.
\item evidence\_gap: pass - Certification readiness is derived from equations, derivation steps, and falsifiers.
\item overclaim\_risk: pass - Overclaim boundary explicitly states that the candidate does not prove the main target.
\item validation\_unbounded: pass - The next action is checked for bounded proof, calculation, audit, or falsifier language.
\item narrow\_bookkeeping: pass - The candidate is treated as scientific only if it has external mathematical content.
\item target\_equivalent\_assumption: finite\_check\_queued - P vs NP candidates must show that the source obligation is not already the desired lower bound.
\item advice\_leakage: finite\_check\_queued - Advice and nonuniform hard-row selection are explicit P vs NP disqualifiers.
\item barrier\_status: finite\_check\_queued - Barrier status must be attached to the specific transfer theorem.
\item restricted\_model\_proxy\_drift: finite\_check\_queued - Restricted-model results need a compiler envelope or explicit proxy boundary.
\end{itemize}

\paragraph{Prior-art boundary.} finite\_check\_queued

\begin{itemize}
\item Run prior-art pass for hardness magnification threshold theorems.
\item Run prior-art pass for natural proofs and algebrization barriers.
\item Run prior-art pass for padding/projection lower-bound transfer arguments.
\end{itemize}

\subsection*{cycle-0007: Hardness-Magnification Bridge Realization: black-box substitution cap}

\paragraph{Contribution type.} quantitative\_transfer\_obstruction

\paragraph{Stepping-stone maturity.} near\_stepping\_stone

\paragraph{Smallest honest claim.} A route-specific quantitative obstruction can be isolated for Hardness-Magnification Bridge Realization: black-box substitution cap: the transfer formula names the source exponent, overhead, target-query exponent, and exact threshold where magnification collapses.

\paragraph{Why it could stand alone.} A quantitative transfer obstruction is independently useful because it replaces barrier folklore with an exact parameter condition showing when a proof route loses the target exponent.

\paragraph{Why this budget now.} Run the next bounded certification or prior-art gate before another broad marathon branch.

\paragraph{Branch control.} The candidate is strong but still needs bounded validation or disqualifier checks before handoff.

\paragraph{Overclaim boundary.} This candidate is a scoped contribution candidate; it does not prove the main target unless all target commitments are discharged in the same theory.

\paragraph{Good-enough paper relevance.} The candidate is scientifically relevant enough to shape a paper, but certification/prior-art gates still control handoff.

\paragraph{Paper relevance status.} strong\_relevance\_certification\_pending with average marker score 4.8

\paragraph{Good enough when.}
\begin{itemize}
\item Smallest honest claim can be stated without proving the main target.
\item Mathematical object, assumptions, scope, and evidence basis are explicit.
\item At least one disqualifier or falsifier has been checked or bounded.
\item Remaining validation is finite enough for a paper note, obstruction note, benchmark, method note, or lemma packet.
\end{itemize}

\paragraph{Continue only if.}
\begin{itemize}
\item The next cycle can improve novelty, explanation, implications, validation readiness, target relevance, or disqualifier resolution.
\item The branch still touches a core target commitment or a reusable scientific object.
\item The next action is a bounded calculation, derivation, source check, or falsifier test rather than generic elaboration.
\end{itemize}

\paragraph{Stop when.}
\begin{itemize}
\item Two explicit adaptation or backtracking attempts fail to improve any contribution metric.
\item The branch cannot name a contribution candidate, contribution-disqualifying test, or target-preserving validation handoff.
\item The candidate is already good enough for handoff and further exploration would only inflate proof debt.
\end{itemize}

\paragraph{Validation burden.}
\begin{itemize}
\item Run prior-art comparison before claiming novelty.
\end{itemize}

\paragraph{Evidence basis.}
\begin{itemize}
\item cycle:cycle-0007
\item branch:Hardness-Magnification Bridge Realization: black-box substitution cap
\item equations:3
\item \(derivation_{steps}:6\)
\item falsifiers:4
\end{itemize}

\paragraph{Missing to publishable.}
\begin{itemize}
\item Run the finite prior-art boundary check and record closest known results.
\item Resolve disqualifier tests or convert unknowns into finite validation tasks.
\end{itemize}

\paragraph{Missing to good-enough paper.}
\begin{itemize}
\item Close the publishable handoff hard gates before calling the candidate paper-ready.
\end{itemize}

\paragraph{Paper relevance markers.}
\begin{itemize}
\item conceptual\_depth: \(score=5\), \(status=strong\) - The result works at the level of a real P vs NP proof obstruction or bridge: source carrier, target class, circuit model, reduction/transfer map, quantifier regime, and barrier status are all explicit.
\item conceptual\_novelty: \(score=4\), \(status=strong\) - The result adds a new exact obstruction, bridge, no-go, or certificate interface relative to known hardness magnification, MCSP/MKtP, proof-complexity, natural-proofs, relativization, or algebrization families.
\item unblocking\_power: \(score=5\), \(status=strong\) - The result changes route priority by proving a transfer, exposing a quantitative loss, ruling out a method class, or identifying the precise bounded calculation that would decide the branch.
\item applicability: \(score=5\), \(status=strong\) - The result is reusable across a nontrivial family of carriers, reductions, restricted models, proof systems, or certificate constructions under explicit assumptions and transfer inequalities.
\item understanding\_gain: \(score=5\), \(status=strong\) - The result clarifies why the route is hard: where proof debt moves, which exponent or advice term destroys the transfer, or which barrier premise actually binds.
\end{itemize}

\paragraph{Hard gates.}
\begin{itemize}
\item claim\_statement\_ready: \(ready=True\), \(finite_{task}=False\)
\item evidence\_ready: \(ready=True\), \(finite_{task}=True\)
\item novelty\_boundary\_ready: \(ready=False\), \(finite_{task}=True\)
\item overclaim\_boundary\_ready: \(ready=True\), \(finite_{task}=False\)
\item validation\_ready: \(ready=False\), \(finite_{task}=True\)
\end{itemize}

\paragraph{Disqualifiers.}
\begin{itemize}
\item already\_known: finite\_check\_queued - Prior-art boundary is queued before publishable handoff.
\item evidence\_gap: pass - Certification readiness is derived from equations, derivation steps, and falsifiers.
\item overclaim\_risk: pass - Overclaim boundary explicitly states that the candidate does not prove the main target.
\item validation\_unbounded: pass - The next action is checked for bounded proof, calculation, audit, or falsifier language.
\item narrow\_bookkeeping: pass - The candidate is treated as scientific only if it has external mathematical content.
\item target\_equivalent\_assumption: finite\_check\_queued - P vs NP candidates must show that the source obligation is not already the desired lower bound.
\item advice\_leakage: finite\_check\_queued - Advice and nonuniform hard-row selection are explicit P vs NP disqualifiers.
\item barrier\_status: pass - Barrier status must be attached to the specific transfer theorem.
\item restricted\_model\_proxy\_drift: pass - Restricted-model results need a compiler envelope or explicit proxy boundary.
\end{itemize}

\paragraph{Prior-art boundary.} finite\_check\_queued

\begin{itemize}
\item Run prior-art pass for hardness magnification threshold theorems.
\item Run prior-art pass for natural proofs and algebrization barriers.
\item Run prior-art pass for padding/projection lower-bound transfer arguments.
\end{itemize}

\subsection*{cycle-0006: Hardness-Magnification Bridge Realization, Cycle 6: MCSP verifier transfer audit}

\paragraph{Contribution type.} quantitative\_transfer\_obstruction

\paragraph{Stepping-stone maturity.} near\_stepping\_stone

\paragraph{Smallest honest claim.} A route-specific quantitative obstruction can be isolated for Hardness-Magnification Bridge Realization, Cycle 6: MCSP verifier transfer audit: the transfer formula names the source exponent, overhead, target-query exponent, and exact threshold where magnification collapses.

\paragraph{Why it could stand alone.} A quantitative transfer obstruction is independently useful because it replaces barrier folklore with an exact parameter condition showing when a proof route loses the target exponent.

\paragraph{Why this budget now.} Run the next bounded certification or prior-art gate before another broad marathon branch.

\paragraph{Branch control.} The candidate is strong but still needs bounded validation or disqualifier checks before handoff.

\paragraph{Overclaim boundary.} This candidate is a scoped contribution candidate; it does not prove the main target unless all target commitments are discharged in the same theory.

\paragraph{Good-enough paper relevance.} The candidate is scientifically relevant enough to shape a paper, but certification/prior-art gates still control handoff.

\paragraph{Paper relevance status.} strong\_relevance\_certification\_pending with average marker score 4.8

\paragraph{Good enough when.}
\begin{itemize}
\item Smallest honest claim can be stated without proving the main target.
\item Mathematical object, assumptions, scope, and evidence basis are explicit.
\item At least one disqualifier or falsifier has been checked or bounded.
\item Remaining validation is finite enough for a paper note, obstruction note, benchmark, method note, or lemma packet.
\end{itemize}

\paragraph{Continue only if.}
\begin{itemize}
\item The next cycle can improve novelty, explanation, implications, validation readiness, target relevance, or disqualifier resolution.
\item The branch still touches a core target commitment or a reusable scientific object.
\item The next action is a bounded calculation, derivation, source check, or falsifier test rather than generic elaboration.
\end{itemize}

\paragraph{Stop when.}
\begin{itemize}
\item Two explicit adaptation or backtracking attempts fail to improve any contribution metric.
\item The branch cannot name a contribution candidate, contribution-disqualifying test, or target-preserving validation handoff.
\item The candidate is already good enough for handoff and further exploration would only inflate proof debt.
\end{itemize}

\paragraph{Validation burden.}
\begin{itemize}
\item Run prior-art comparison before claiming novelty.
\end{itemize}

\paragraph{Evidence basis.}
\begin{itemize}
\item cycle:cycle-0006
\item branch:Hardness-Magnification Bridge Realization, Cycle 6: MCSP verifier transfer audit
\item equations:4
\item \(derivation_{steps}:7\)
\item falsifiers:4
\end{itemize}

\paragraph{Missing to publishable.}
\begin{itemize}
\item Run the finite prior-art boundary check and record closest known results.
\item Resolve disqualifier tests or convert unknowns into finite validation tasks.
\end{itemize}

\paragraph{Missing to good-enough paper.}
\begin{itemize}
\item Close the publishable handoff hard gates before calling the candidate paper-ready.
\end{itemize}

\paragraph{Paper relevance markers.}
\begin{itemize}
\item conceptual\_depth: \(score=5\), \(status=strong\) - The result works at the level of a real P vs NP proof obstruction or bridge: source carrier, target class, circuit model, reduction/transfer map, quantifier regime, and barrier status are all explicit.
\item conceptual\_novelty: \(score=4\), \(status=strong\) - The result adds a new exact obstruction, bridge, no-go, or certificate interface relative to known hardness magnification, MCSP/MKtP, proof-complexity, natural-proofs, relativization, or algebrization families.
\item unblocking\_power: \(score=5\), \(status=strong\) - The result changes route priority by proving a transfer, exposing a quantitative loss, ruling out a method class, or identifying the precise bounded calculation that would decide the branch.
\item applicability: \(score=5\), \(status=strong\) - The result is reusable across a nontrivial family of carriers, reductions, restricted models, proof systems, or certificate constructions under explicit assumptions and transfer inequalities.
\item understanding\_gain: \(score=5\), \(status=strong\) - The result clarifies why the route is hard: where proof debt moves, which exponent or advice term destroys the transfer, or which barrier premise actually binds.
\end{itemize}

\paragraph{Hard gates.}
\begin{itemize}
\item claim\_statement\_ready: \(ready=True\), \(finite_{task}=False\)
\item evidence\_ready: \(ready=True\), \(finite_{task}=True\)
\item novelty\_boundary\_ready: \(ready=False\), \(finite_{task}=True\)
\item overclaim\_boundary\_ready: \(ready=True\), \(finite_{task}=False\)
\item validation\_ready: \(ready=False\), \(finite_{task}=True\)
\end{itemize}

\paragraph{Disqualifiers.}
\begin{itemize}
\item already\_known: finite\_check\_queued - Prior-art boundary is queued before publishable handoff.
\item evidence\_gap: pass - Certification readiness is derived from equations, derivation steps, and falsifiers.
\item overclaim\_risk: pass - Overclaim boundary explicitly states that the candidate does not prove the main target.
\item validation\_unbounded: pass - The next action is checked for bounded proof, calculation, audit, or falsifier language.
\item narrow\_bookkeeping: pass - The candidate is treated as scientific only if it has external mathematical content.
\item target\_equivalent\_assumption: finite\_check\_queued - P vs NP candidates must show that the source obligation is not already the desired lower bound.
\item advice\_leakage: finite\_check\_queued - Advice and nonuniform hard-row selection are explicit P vs NP disqualifiers.
\item barrier\_status: finite\_check\_queued - Barrier status must be attached to the specific transfer theorem.
\item restricted\_model\_proxy\_drift: pass - Restricted-model results need a compiler envelope or explicit proxy boundary.
\end{itemize}

\paragraph{Prior-art boundary.} finite\_check\_queued

\begin{itemize}
\item Run prior-art pass for hardness magnification threshold theorems.
\item Run prior-art pass for natural proofs and algebrization barriers.
\item Run prior-art pass for padding/projection lower-bound transfer arguments.
\end{itemize}

\subsection*{cycle-0002: Hardness-Magnification Bridge Realization}

\paragraph{Contribution type.} quantitative\_transfer\_obstruction

\paragraph{Stepping-stone maturity.} near\_stepping\_stone

\paragraph{Smallest honest claim.} A route-specific quantitative obstruction can be isolated for Hardness-Magnification Bridge Realization: the transfer formula names the source exponent, overhead, target-query exponent, and exact threshold where magnification collapses.

\paragraph{Why it could stand alone.} A quantitative transfer obstruction is independently useful because it replaces barrier folklore with an exact parameter condition showing when a proof route loses the target exponent.

\paragraph{Why this budget now.} Run the next bounded certification or prior-art gate before another broad marathon branch.

\paragraph{Branch control.} The candidate is strong but still needs bounded validation or disqualifier checks before handoff.

\paragraph{Overclaim boundary.} This candidate is a scoped contribution candidate; it does not prove the main target unless all target commitments are discharged in the same theory.

\paragraph{Good-enough paper relevance.} The candidate is scientifically relevant enough to shape a paper, but certification/prior-art gates still control handoff.

\paragraph{Paper relevance status.} strong\_relevance\_certification\_pending with average marker score 4.8

\paragraph{Good enough when.}
\begin{itemize}
\item Smallest honest claim can be stated without proving the main target.
\item Mathematical object, assumptions, scope, and evidence basis are explicit.
\item At least one disqualifier or falsifier has been checked or bounded.
\item Remaining validation is finite enough for a paper note, obstruction note, benchmark, method note, or lemma packet.
\end{itemize}

\paragraph{Continue only if.}
\begin{itemize}
\item The next cycle can improve novelty, explanation, implications, validation readiness, target relevance, or disqualifier resolution.
\item The branch still touches a core target commitment or a reusable scientific object.
\item The next action is a bounded calculation, derivation, source check, or falsifier test rather than generic elaboration.
\end{itemize}

\paragraph{Stop when.}
\begin{itemize}
\item Two explicit adaptation or backtracking attempts fail to improve any contribution metric.
\item The branch cannot name a contribution candidate, contribution-disqualifying test, or target-preserving validation handoff.
\item The candidate is already good enough for handoff and further exploration would only inflate proof debt.
\end{itemize}

\paragraph{Validation burden.}
\begin{itemize}
\item Run prior-art comparison before claiming novelty.
\end{itemize}

\paragraph{Evidence basis.}
\begin{itemize}
\item cycle:cycle-0002
\item branch:Hardness-Magnification Bridge Realization
\item equations:5
\item \(derivation_{steps}:7\)
\item falsifiers:5
\end{itemize}

\paragraph{Missing to publishable.}
\begin{itemize}
\item Run the finite prior-art boundary check and record closest known results.
\item Resolve disqualifier tests or convert unknowns into finite validation tasks.
\end{itemize}

\paragraph{Missing to good-enough paper.}
\begin{itemize}
\item Close the publishable handoff hard gates before calling the candidate paper-ready.
\end{itemize}

\paragraph{Paper relevance markers.}
\begin{itemize}
\item conceptual\_depth: \(score=5\), \(status=strong\) - The result works at the level of a real P vs NP proof obstruction or bridge: source carrier, target class, circuit model, reduction/transfer map, quantifier regime, and barrier status are all explicit.
\item conceptual\_novelty: \(score=4\), \(status=strong\) - The result adds a new exact obstruction, bridge, no-go, or certificate interface relative to known hardness magnification, MCSP/MKtP, proof-complexity, natural-proofs, relativization, or algebrization families.
\item unblocking\_power: \(score=5\), \(status=strong\) - The result changes route priority by proving a transfer, exposing a quantitative loss, ruling out a method class, or identifying the precise bounded calculation that would decide the branch.
\item applicability: \(score=5\), \(status=strong\) - The result is reusable across a nontrivial family of carriers, reductions, restricted models, proof systems, or certificate constructions under explicit assumptions and transfer inequalities.
\item understanding\_gain: \(score=5\), \(status=strong\) - The result clarifies why the route is hard: where proof debt moves, which exponent or advice term destroys the transfer, or which barrier premise actually binds.
\end{itemize}

\paragraph{Hard gates.}
\begin{itemize}
\item claim\_statement\_ready: \(ready=True\), \(finite_{task}=False\)
\item evidence\_ready: \(ready=True\), \(finite_{task}=True\)
\item novelty\_boundary\_ready: \(ready=False\), \(finite_{task}=True\)
\item overclaim\_boundary\_ready: \(ready=True\), \(finite_{task}=False\)
\item validation\_ready: \(ready=False\), \(finite_{task}=True\)
\end{itemize}

\paragraph{Disqualifiers.}
\begin{itemize}
\item already\_known: finite\_check\_queued - Prior-art boundary is queued before publishable handoff.
\item evidence\_gap: pass - Certification readiness is derived from equations, derivation steps, and falsifiers.
\item overclaim\_risk: pass - Overclaim boundary explicitly states that the candidate does not prove the main target.
\item validation\_unbounded: pass - The next action is checked for bounded proof, calculation, audit, or falsifier language.
\item narrow\_bookkeeping: pass - The candidate is treated as scientific only if it has external mathematical content.
\item target\_equivalent\_assumption: finite\_check\_queued - P vs NP candidates must show that the source obligation is not already the desired lower bound.
\item advice\_leakage: finite\_check\_queued - Advice and nonuniform hard-row selection are explicit P vs NP disqualifiers.
\item barrier\_status: finite\_check\_queued - Barrier status must be attached to the specific transfer theorem.
\item restricted\_model\_proxy\_drift: pass - Restricted-model results need a compiler envelope or explicit proxy boundary.
\end{itemize}

\paragraph{Prior-art boundary.} finite\_check\_queued

\begin{itemize}
\item Run prior-art pass for hardness magnification threshold theorems.
\item Run prior-art pass for natural proofs and algebrization barriers.
\item Run prior-art pass for padding/projection lower-bound transfer arguments.
\end{itemize}

\subsection*{cycle-0001: Hardness-Magnification Bridge Realization}

\paragraph{Contribution type.} quantitative\_transfer\_obstruction

\paragraph{Stepping-stone maturity.} near\_stepping\_stone

\paragraph{Smallest honest claim.} A route-specific quantitative obstruction can be isolated for Hardness-Magnification Bridge Realization: the transfer formula names the source exponent, overhead, target-query exponent, and exact threshold where magnification collapses.

\paragraph{Why it could stand alone.} A quantitative transfer obstruction is independently useful because it replaces barrier folklore with an exact parameter condition showing when a proof route loses the target exponent.

\paragraph{Why this budget now.} Run the next bounded certification or prior-art gate before another broad marathon branch.

\paragraph{Branch control.} The candidate is strong but still needs bounded validation or disqualifier checks before handoff.

\paragraph{Overclaim boundary.} This candidate is a scoped contribution candidate; it does not prove the main target unless all target commitments are discharged in the same theory.

\paragraph{Good-enough paper relevance.} The candidate is scientifically relevant enough to shape a paper, but certification/prior-art gates still control handoff.

\paragraph{Paper relevance status.} strong\_relevance\_certification\_pending with average marker score 4.8

\paragraph{Good enough when.}
\begin{itemize}
\item Smallest honest claim can be stated without proving the main target.
\item Mathematical object, assumptions, scope, and evidence basis are explicit.
\item At least one disqualifier or falsifier has been checked or bounded.
\item Remaining validation is finite enough for a paper note, obstruction note, benchmark, method note, or lemma packet.
\end{itemize}

\paragraph{Continue only if.}
\begin{itemize}
\item The next cycle can improve novelty, explanation, implications, validation readiness, target relevance, or disqualifier resolution.
\item The branch still touches a core target commitment or a reusable scientific object.
\item The next action is a bounded calculation, derivation, source check, or falsifier test rather than generic elaboration.
\end{itemize}

\paragraph{Stop when.}
\begin{itemize}
\item Two explicit adaptation or backtracking attempts fail to improve any contribution metric.
\item The branch cannot name a contribution candidate, contribution-disqualifying test, or target-preserving validation handoff.
\item The candidate is already good enough for handoff and further exploration would only inflate proof debt.
\end{itemize}

\paragraph{Validation burden.}
\begin{itemize}
\item Run prior-art comparison before claiming novelty.
\end{itemize}

\paragraph{Evidence basis.}
\begin{itemize}
\item cycle:cycle-0001
\item branch:Hardness-Magnification Bridge Realization
\item equations:3
\item \(derivation_{steps}:6\)
\item falsifiers:5
\end{itemize}

\paragraph{Missing to publishable.}
\begin{itemize}
\item Run the finite prior-art boundary check and record closest known results.
\item Resolve disqualifier tests or convert unknowns into finite validation tasks.
\end{itemize}

\paragraph{Missing to good-enough paper.}
\begin{itemize}
\item Close the publishable handoff hard gates before calling the candidate paper-ready.
\end{itemize}

\paragraph{Paper relevance markers.}
\begin{itemize}
\item conceptual\_depth: \(score=5\), \(status=strong\) - The result works at the level of a real P vs NP proof obstruction or bridge: source carrier, target class, circuit model, reduction/transfer map, quantifier regime, and barrier status are all explicit.
\item conceptual\_novelty: \(score=4\), \(status=strong\) - The result adds a new exact obstruction, bridge, no-go, or certificate interface relative to known hardness magnification, MCSP/MKtP, proof-complexity, natural-proofs, relativization, or algebrization families.
\item unblocking\_power: \(score=5\), \(status=strong\) - The result changes route priority by proving a transfer, exposing a quantitative loss, ruling out a method class, or identifying the precise bounded calculation that would decide the branch.
\item applicability: \(score=5\), \(status=strong\) - The result is reusable across a nontrivial family of carriers, reductions, restricted models, proof systems, or certificate constructions under explicit assumptions and transfer inequalities.
\item understanding\_gain: \(score=5\), \(status=strong\) - The result clarifies why the route is hard: where proof debt moves, which exponent or advice term destroys the transfer, or which barrier premise actually binds.
\end{itemize}

\paragraph{Hard gates.}
\begin{itemize}
\item claim\_statement\_ready: \(ready=True\), \(finite_{task}=False\)
\item evidence\_ready: \(ready=True\), \(finite_{task}=True\)
\item novelty\_boundary\_ready: \(ready=False\), \(finite_{task}=True\)
\item overclaim\_boundary\_ready: \(ready=True\), \(finite_{task}=False\)
\item validation\_ready: \(ready=False\), \(finite_{task}=True\)
\end{itemize}

\paragraph{Disqualifiers.}
\begin{itemize}
\item already\_known: finite\_check\_queued - Prior-art boundary is queued before publishable handoff.
\item evidence\_gap: pass - Certification readiness is derived from equations, derivation steps, and falsifiers.
\item overclaim\_risk: pass - Overclaim boundary explicitly states that the candidate does not prove the main target.
\item validation\_unbounded: pass - The next action is checked for bounded proof, calculation, audit, or falsifier language.
\item narrow\_bookkeeping: pass - The candidate is treated as scientific only if it has external mathematical content.
\item target\_equivalent\_assumption: finite\_check\_queued - P vs NP candidates must show that the source obligation is not already the desired lower bound.
\item advice\_leakage: finite\_check\_queued - Advice and nonuniform hard-row selection are explicit P vs NP disqualifiers.
\item barrier\_status: finite\_check\_queued - Barrier status must be attached to the specific transfer theorem.
\item restricted\_model\_proxy\_drift: finite\_check\_queued - Restricted-model results need a compiler envelope or explicit proxy boundary.
\end{itemize}

\paragraph{Prior-art boundary.} finite\_check\_queued

\begin{itemize}
\item Run prior-art pass for hardness magnification threshold theorems.
\item Run prior-art pass for natural proofs and algebrization barriers.
\item Run prior-art pass for padding/projection lower-bound transfer arguments.
\end{itemize}

\subsection*{cycle-0136: Hardness-Magnification Bridge Realization}

\paragraph{Contribution type.} quantitative\_transfer\_obstruction

\paragraph{Stepping-stone maturity.} near\_stepping\_stone

\paragraph{Smallest honest claim.} A route-specific quantitative obstruction can be isolated for Hardness-Magnification Bridge Realization: the transfer formula names the source exponent, overhead, target-query exponent, and exact threshold where magnification collapses.

\paragraph{Why it could stand alone.} A quantitative transfer obstruction is independently useful because it replaces barrier folklore with an exact parameter condition showing when a proof route loses the target exponent.

\paragraph{Why this budget now.} Run the next bounded certification or prior-art gate before another broad marathon branch.

\paragraph{Branch control.} The candidate is strong but still needs bounded validation or disqualifier checks before handoff.

\paragraph{Overclaim boundary.} This candidate is a scoped contribution candidate; it does not prove the main target unless all target commitments are discharged in the same theory.

\paragraph{Good-enough paper relevance.} The candidate is scientifically relevant enough to shape a paper, but certification/prior-art gates still control handoff.

\paragraph{Paper relevance status.} strong\_relevance\_certification\_pending with average marker score 4.8

\paragraph{Good enough when.}
\begin{itemize}
\item Smallest honest claim can be stated without proving the main target.
\item Mathematical object, assumptions, scope, and evidence basis are explicit.
\item At least one disqualifier or falsifier has been checked or bounded.
\item Remaining validation is finite enough for a paper note, obstruction note, benchmark, method note, or lemma packet.
\end{itemize}

\paragraph{Continue only if.}
\begin{itemize}
\item The next cycle can improve novelty, explanation, implications, validation readiness, target relevance, or disqualifier resolution.
\item The branch still touches a core target commitment or a reusable scientific object.
\item The next action is a bounded calculation, derivation, source check, or falsifier test rather than generic elaboration.
\end{itemize}

\paragraph{Stop when.}
\begin{itemize}
\item Two explicit adaptation or backtracking attempts fail to improve any contribution metric.
\item The branch cannot name a contribution candidate, contribution-disqualifying test, or target-preserving validation handoff.
\item The candidate is already good enough for handoff and further exploration would only inflate proof debt.
\end{itemize}

\paragraph{Validation burden.}
\begin{itemize}
\item Run prior-art comparison before claiming novelty.
\end{itemize}

\paragraph{Evidence basis.}
\begin{itemize}
\item cycle:cycle-0136
\item branch:Hardness-Magnification Bridge Realization
\item equations:6
\item \(derivation_{steps}:7\)
\item falsifiers:6
\end{itemize}

\paragraph{Missing to publishable.}
\begin{itemize}
\item Run the finite prior-art boundary check and record closest known results.
\item Resolve disqualifier tests or convert unknowns into finite validation tasks.
\end{itemize}

\paragraph{Missing to good-enough paper.}
\begin{itemize}
\item Close the publishable handoff hard gates before calling the candidate paper-ready.
\end{itemize}

\paragraph{Paper relevance markers.}
\begin{itemize}
\item conceptual\_depth: \(score=5\), \(status=strong\) - The result works at the level of a real P vs NP proof obstruction or bridge: source carrier, target class, circuit model, reduction/transfer map, quantifier regime, and barrier status are all explicit.
\item conceptual\_novelty: \(score=4\), \(status=strong\) - The result adds a new exact obstruction, bridge, no-go, or certificate interface relative to known hardness magnification, MCSP/MKtP, proof-complexity, natural-proofs, relativization, or algebrization families.
\item unblocking\_power: \(score=5\), \(status=strong\) - The result changes route priority by proving a transfer, exposing a quantitative loss, ruling out a method class, or identifying the precise bounded calculation that would decide the branch.
\item applicability: \(score=5\), \(status=strong\) - The result is reusable across a nontrivial family of carriers, reductions, restricted models, proof systems, or certificate constructions under explicit assumptions and transfer inequalities.
\item understanding\_gain: \(score=5\), \(status=strong\) - The result clarifies why the route is hard: where proof debt moves, which exponent or advice term destroys the transfer, or which barrier premise actually binds.
\end{itemize}

\paragraph{Hard gates.}
\begin{itemize}
\item claim\_statement\_ready: \(ready=True\), \(finite_{task}=False\)
\item evidence\_ready: \(ready=True\), \(finite_{task}=True\)
\item novelty\_boundary\_ready: \(ready=False\), \(finite_{task}=True\)
\item overclaim\_boundary\_ready: \(ready=True\), \(finite_{task}=False\)
\item validation\_ready: \(ready=False\), \(finite_{task}=True\)
\end{itemize}

\paragraph{Disqualifiers.}
\begin{itemize}
\item already\_known: finite\_check\_queued - Prior-art boundary is queued before publishable handoff.
\item evidence\_gap: pass - Certification readiness is derived from equations, derivation steps, and falsifiers.
\item overclaim\_risk: pass - Overclaim boundary explicitly states that the candidate does not prove the main target.
\item validation\_unbounded: pass - The next action is checked for bounded proof, calculation, audit, or falsifier language.
\item narrow\_bookkeeping: pass - The candidate is treated as scientific only if it has external mathematical content.
\item target\_equivalent\_assumption: finite\_check\_queued - P vs NP candidates must show that the source obligation is not already the desired lower bound.
\item advice\_leakage: finite\_check\_queued - Advice and nonuniform hard-row selection are explicit P vs NP disqualifiers.
\item barrier\_status: pass - Barrier status must be attached to the specific transfer theorem.
\item restricted\_model\_proxy\_drift: pass - Restricted-model results need a compiler envelope or explicit proxy boundary.
\end{itemize}

\paragraph{Prior-art boundary.} finite\_check\_queued

\begin{itemize}
\item Run prior-art pass for hardness magnification threshold theorems.
\item Run prior-art pass for natural proofs and algebrization barriers.
\item Run prior-art pass for padding/projection lower-bound transfer arguments.
\end{itemize}

\subsection*{cycle-0134: Hardness-Magnification Bridge Realization}

\paragraph{Contribution type.} quantitative\_transfer\_obstruction

\paragraph{Stepping-stone maturity.} near\_stepping\_stone

\paragraph{Smallest honest claim.} A route-specific quantitative obstruction can be isolated for Hardness-Magnification Bridge Realization: the transfer formula names the source exponent, overhead, target-query exponent, and exact threshold where magnification collapses.

\paragraph{Why it could stand alone.} A quantitative transfer obstruction is independently useful because it replaces barrier folklore with an exact parameter condition showing when a proof route loses the target exponent.

\paragraph{Why this budget now.} Run the next bounded certification or prior-art gate before another broad marathon branch.

\paragraph{Branch control.} The candidate is strong but still needs bounded validation or disqualifier checks before handoff.

\paragraph{Overclaim boundary.} This candidate is a scoped contribution candidate; it does not prove the main target unless all target commitments are discharged in the same theory.

\paragraph{Good-enough paper relevance.} The candidate is scientifically relevant enough to shape a paper, but certification/prior-art gates still control handoff.

\paragraph{Paper relevance status.} strong\_relevance\_certification\_pending with average marker score 4.8

\paragraph{Good enough when.}
\begin{itemize}
\item Smallest honest claim can be stated without proving the main target.
\item Mathematical object, assumptions, scope, and evidence basis are explicit.
\item At least one disqualifier or falsifier has been checked or bounded.
\item Remaining validation is finite enough for a paper note, obstruction note, benchmark, method note, or lemma packet.
\end{itemize}

\paragraph{Continue only if.}
\begin{itemize}
\item The next cycle can improve novelty, explanation, implications, validation readiness, target relevance, or disqualifier resolution.
\item The branch still touches a core target commitment or a reusable scientific object.
\item The next action is a bounded calculation, derivation, source check, or falsifier test rather than generic elaboration.
\end{itemize}

\paragraph{Stop when.}
\begin{itemize}
\item Two explicit adaptation or backtracking attempts fail to improve any contribution metric.
\item The branch cannot name a contribution candidate, contribution-disqualifying test, or target-preserving validation handoff.
\item The candidate is already good enough for handoff and further exploration would only inflate proof debt.
\end{itemize}

\paragraph{Validation burden.}
\begin{itemize}
\item Run prior-art comparison before claiming novelty.
\end{itemize}

\paragraph{Evidence basis.}
\begin{itemize}
\item cycle:cycle-0134
\item branch:Hardness-Magnification Bridge Realization
\item equations:5
\item \(derivation_{steps}:7\)
\item falsifiers:5
\end{itemize}

\paragraph{Missing to publishable.}
\begin{itemize}
\item Run the finite prior-art boundary check and record closest known results.
\item Resolve disqualifier tests or convert unknowns into finite validation tasks.
\end{itemize}

\paragraph{Missing to good-enough paper.}
\begin{itemize}
\item Close the publishable handoff hard gates before calling the candidate paper-ready.
\end{itemize}

\paragraph{Paper relevance markers.}
\begin{itemize}
\item conceptual\_depth: \(score=5\), \(status=strong\) - The result works at the level of a real P vs NP proof obstruction or bridge: source carrier, target class, circuit model, reduction/transfer map, quantifier regime, and barrier status are all explicit.
\item conceptual\_novelty: \(score=4\), \(status=strong\) - The result adds a new exact obstruction, bridge, no-go, or certificate interface relative to known hardness magnification, MCSP/MKtP, proof-complexity, natural-proofs, relativization, or algebrization families.
\item unblocking\_power: \(score=5\), \(status=strong\) - The result changes route priority by proving a transfer, exposing a quantitative loss, ruling out a method class, or identifying the precise bounded calculation that would decide the branch.
\item applicability: \(score=5\), \(status=strong\) - The result is reusable across a nontrivial family of carriers, reductions, restricted models, proof systems, or certificate constructions under explicit assumptions and transfer inequalities.
\item understanding\_gain: \(score=5\), \(status=strong\) - The result clarifies why the route is hard: where proof debt moves, which exponent or advice term destroys the transfer, or which barrier premise actually binds.
\end{itemize}

\paragraph{Hard gates.}
\begin{itemize}
\item claim\_statement\_ready: \(ready=True\), \(finite_{task}=False\)
\item evidence\_ready: \(ready=True\), \(finite_{task}=True\)
\item novelty\_boundary\_ready: \(ready=False\), \(finite_{task}=True\)
\item overclaim\_boundary\_ready: \(ready=True\), \(finite_{task}=False\)
\item validation\_ready: \(ready=False\), \(finite_{task}=True\)
\end{itemize}

\paragraph{Disqualifiers.}
\begin{itemize}
\item already\_known: finite\_check\_queued - Prior-art boundary is queued before publishable handoff.
\item evidence\_gap: pass - Certification readiness is derived from equations, derivation steps, and falsifiers.
\item overclaim\_risk: pass - Overclaim boundary explicitly states that the candidate does not prove the main target.
\item validation\_unbounded: pass - The next action is checked for bounded proof, calculation, audit, or falsifier language.
\item narrow\_bookkeeping: pass - The candidate is treated as scientific only if it has external mathematical content.
\item target\_equivalent\_assumption: finite\_check\_queued - P vs NP candidates must show that the source obligation is not already the desired lower bound.
\item advice\_leakage: finite\_check\_queued - Advice and nonuniform hard-row selection are explicit P vs NP disqualifiers.
\item barrier\_status: finite\_check\_queued - Barrier status must be attached to the specific transfer theorem.
\item restricted\_model\_proxy\_drift: finite\_check\_queued - Restricted-model results need a compiler envelope or explicit proxy boundary.
\end{itemize}

\paragraph{Prior-art boundary.} finite\_check\_queued

\begin{itemize}
\item Run prior-art pass for hardness magnification threshold theorems.
\item Run prior-art pass for natural proofs and algebrization barriers.
\item Run prior-art pass for padding/projection lower-bound transfer arguments.
\end{itemize}


\section{Target Problem And Literature Evolution}

\label{sec:report-problem-literature}

\subsection*{P vs NP}

Decide whether every efficiently verifiable decision problem is also efficiently solvable, equivalently whether SAT has a polynomial-time algorithm.

\[
\mathrm{P}\stackrel{?}{=}\mathrm{NP}\quad\text{or, equivalently,}\quad \mathrm{SAT}\in\mathrm{P}
\]

\subsection*{Why this problem matters}
\begin{itemize}
\item It is the central test of whether verification and search have the same polynomial-time power.
\item It governs the meaning of NP-completeness: if P equals NP, thousands of complete search problems become efficiently solvable in principle; if not, their common hardness has a rigorous source.
\item It is tied to cryptography, proof complexity, optimization, algorithms, derandomization, circuit lower bounds, and the philosophy of what counts as an efficient explanation.
\end{itemize}

\subsection*{What would count as a solution}
\begin{itemize}
\item To prove P equals NP, exhibit a polynomial-time algorithm for an NP-complete problem such as SAT and prove its runtime and correctness.
\item To prove P differs from NP, prove that some NP language, equivalently SAT, cannot be decided by any polynomial-time deterministic Turing machine.
\item For nonuniform routes, a claim strong enough for \(\mathrm{NP}\not\subset \mathrm{P}/\mathrm{poly}\) must rule out polynomial-size circuit families for an NP-complete language without uncharged advice.
\item Any restricted-model, algebraic, proof-complexity, meta-complexity, or hardness-magnification result must include a target-preserving bridge back to the Boolean main problem.
\end{itemize}

\subsection*{How the literature evolved}

\subsubsection*{1971-1972 -- Verification becomes a universal computational phenomenon}

\paragraph{Conceptual shift.} Cook and Karp show that many apparently different problems are interreducible through polynomial reductions and share SAT as a complete core.

\paragraph{What became possible.} The problem becomes mathematically precise: it is enough to decide whether any NP-complete language, especially SAT, has polynomial-time algorithms.

\paragraph{What remained open.} Completeness explains universality of difficulty but not the source of lower bounds against all polynomial-time algorithms or circuits.

\subsubsection*{1975 -- Proof techniques become objects of study}

\paragraph{Conceptual shift.} Relativization shows that any proof surviving arbitrary oracle substitution cannot settle the main question.

\paragraph{What became possible.} A serious route must track what part of the argument is nonrelativizing rather than merely proving a plausible separation in an oracle-insensitive way.

\paragraph{What remained open.} The barrier diagnoses broad method families but does not identify a successful nonrelativizing replacement.

\subsubsection*{Late 1970s onward -- Complexity becomes a landscape of representations}

\paragraph{Conceptual shift.} Algebraic, proof-complexity, descriptive, and structural viewpoints re-express computation through polynomials, proofs, reductions, and uniformity conditions.

\paragraph{What became possible.} The problem can be attacked by translating Boolean hardness into more rigid mathematical objects with their own invariants.

\paragraph{What remained open.} Every translation risks losing target strength: a theorem in a restricted or algebraic model must still return to Boolean P vs NP or P/poly with controlled loss.

\subsubsection*{1990s -- Lower-bound intuitions meet pseudorandomness}

\paragraph{Conceptual shift.} Natural proofs explain why many broad, efficiently checkable combinatorial properties are unlikely to prove strong unrestricted circuit lower bounds.

\paragraph{What became possible.} Discovery can now audit lower-bound proposals for largeness, constructivity, usefulness, and cryptographic tension.

\paragraph{What remained open.} The field still needs useful lower-bound properties that are either not large, not constructive, not black-box, or otherwise outside the natural-proofs template.

\subsubsection*{2000s onward -- Fine-grained and magnification thinking exposes exponent debt}

\paragraph{Conceptual shift.} Hardness transfers are judged by exact exponents, target-query counts, advice leakage, and whether a weak source lower bound really magnifies to the main target.

\paragraph{What became possible.} Routes can be evaluated with precise inequalities instead of broad optimism about reductions or carriers.

\paragraph{What remained open.} Most bridges lose too much exponent, import nonuniform advice, or prove a source lower bound that is already as hard as the desired conclusion.

\subsubsection*{2009 onward -- Modern progress is real but carefully scoped}

\paragraph{Conceptual shift.} Algebrization and restricted-circuit lower bounds clarify both the reach and limits of current techniques.

\paragraph{What became possible.} Algorithmic, algebraic, and restricted-circuit breakthroughs provide templates for bridge design and falsification.

\paragraph{What remained open.} The unrestricted P vs NP question still demands a target-preserving lower bound or impossibility theorem that escapes known barriers without losing the main quantifiers.

\subsection*{References preserved by the report}

\paragraph{S. A. Cook (1971).} The Complexity of Theorem-Proving Procedures. Supports: Introduces NP-completeness through satisfiability and turns efficient verification into a universal problem.. Limitation: It defines the central complete problem and reduction method; it does not separate deterministic from nondeterministic polynomial time.

\paragraph{R. M. Karp (1972).} Reducibility Among Combinatorial Problems. Supports: Shows that many natural search and optimization problems share the same NP-complete core.. Limitation: NP-completeness transfers difficulty conditionally; it does not prove that any NP-complete problem lacks polynomial algorithms.

\paragraph{T. Baker, J. Gill, and R. Solovay (1975).} Relativizations of the P \(=\) ? NP Question. Supports: Shows that relativizing methods alone cannot settle P vs NP, because oracle worlds exist on both sides.. Limitation: Relativization is a diagnostic for proof techniques, not a proof that separation is impossible.

\paragraph{A. A. Razborov and S. Rudich (1997).} Natural Proofs. Supports: Explains why many combinatorial lower-bound properties are too large and constructive to prove strong general circuit lower bounds under cryptographic hardness assumptions.. Limitation: The barrier constrains a proof format; it leaves room for non-natural, nonconstructive, algebraic, or otherwise barrier-evading arguments.

\paragraph{S. Aaronson and A. Wigderson (2009).} Algebrization: A New Barrier in Complexity Theory. Supports: Shows that even many non-relativizing arithmetization methods still fall inside an algebrizing barrier.. Limitation: Algebrization does not rule out all algebraic or geometric methods; it forces proof attempts to state exactly what escapes the barrier.

\paragraph{R. Impagliazzo, R. Paturi, and F. Zane (2001).} Which Problems Have Strongly Exponential Complexity?. Supports: Develops reduction notions that preserve subexponential complexity and clarify how quantitative hardness transfers.. Limitation: Strong exponential-time hypotheses and SERF reductions illuminate hardness structure but do not by themselves separate P from NP.

\paragraph{R. Williams (2014).} Nonuniform ACC Circuit Lower Bounds. Supports: Proves that NEXP does not have polynomial-size nonuniform ACC circuits and demonstrates an algorithms-to-lower-bounds bridge.. Limitation: The result is a major restricted-circuit breakthrough, not a lower bound against unrestricted polynomial-size circuits for NP.

\paragraph{L. G. Valiant (1979).} Completeness Classes in Algebra. Supports: Creates an algebraic analogue of completeness and lower-bound questions, showing how representation and computation can be studied through polynomial families.. Limitation: Algebraic complexity gives powerful analogues and transfer pressure, but VP vs VNP is not literally the Boolean P vs NP problem.


\section{Macro Strategies}

\label{sec:report-macro-strategies}

Current macro strategy: compute the transfer exponents or turn the route into a no-go theorem.

\subsection*{Black-box transfer and exponent-window audit}

\paragraph{Physics picture.} For this complexity-theory target, read the physical picture as a proof-dynamics picture: the target circuit is a resource flowing through an interface. Positive query-length exponent \(\beta\) amplifies the target exponent c and prevents a barely-superlinear source lower bound from ruling out all polynomial-size SAT circuits.

\paragraph{Active skills.}
\begin{itemize}
\item pvsnp-symbolic-complexity-auditor
\item pvsnp-qualitative-complexity-reasoner
\item pvsnp-conceptual-pressure-localizer
\item entailment-gap-auditor
\item progression-payoff-judge
\end{itemize}

\paragraph{Evidence from story.}
\begin{itemize}
\item branch\_title: Hardness-Magnification Bridge Realization; cycle\_id: cycle-0149; learning: Cycle 149 folds representative-stable ranked QCTSR into \(TransparentRowSupport_{N}\) as a valid rule case. The new contribution is not another magnification slogan: quotient compression is proof-carrying, rank-stratified, dependency-closed, and charged, while the first rejected merge gives a concrete private-row falsifier. Full semantic SAT pruning remains exactly at \(RealSATInfluenceRealizer_{N}^{arb}\).; next\_experiment: Write the standalone paper theorem for \(TransparentRowSupport_{N}^{Qrep}\): include the \(K_{Qrep}\) fixed-point verifier pseudocode, the \(\Gamma_{Qrep}\) ledger proof, the first-rejected-merge countermodel, and a one-page nonclaim box separating this parser boundary from \(H_{N}\) lower bounds and \(\mathrm{P} \ne \mathrm{NP}\).
\item branch\_title: Hardness-Magnification Bridge Realization; cycle\_id: cycle-0150; learning: Cycle 150 closes a proof-delivery upgrade rather than rediscovering magnification: Qrep quotient compression is valid exactly as a batched, ranked, proof-carrying same-theory dictionary charged in \(\Gamma_{Qrep}\). The sharp falsifier is the first unranked, cyclic, advice-chosen, or hidden semantic merge. The result is a scoped parser/falsifier theorem and Formula Footprints supporting lemma; it does not prove \(H_{N}\) lower bounds or \(\mathrm{P} \ne \mathrm{NP}\).; next\_experiment: Run the Tseitin-renaming Qrep thought experiment: instantiate two batches with one certified duplicate and one unlisted duplicate, compute \(K_{Qrep}\), C\_Q(D\_Q), and \(\Gamma_{Qrep}\) explicitly, then verify that the unlisted duplicate yields the first-hidden-row private-row falsifier while the certified duplicate preserves the line induction.
\item branch\_title: Hardness-Magnification Bridge Realization; cycle\_id: cycle-0151; learning: Cycle 151 upgrades the batched Qrep patch with temporal prefix stability: quotient compression is valid only as append-only, ranked, proof-carrying dictionary use. The first pre-use, rebinding, cyclic, advice-chosen, or hidden semantic merge is an explicit falsifier, while full semantic pruning stays quarantined.; next\_experiment: Write the finite prefix-stable Qrep verifier: fields are batch\_id, rank, source\_row, representative\_row, certificate\_root, dependency list, first\_use\_line, and no\_rebind flag; prove the first rejected field yields the private-row falsifier.
\item branch\_title: Hardness-Magnification Bridge Realization: prefix-stable Qrep support closure; cycle\_id: cycle-0152; learning: Cycle 152 upgrades Qrep from an append-only convention to a prefix-verifiable rule case: quotient compression preserves the transparent parser only when certificate charge strictly precedes representative use. The exact no-go is temporal, not just acyclicity; delayed same-theory evidence is already a private-row falsifier. This closes a proof-audit loophole while preserving the nonclaim boundary: no \(H_{N}\) lower bound and no \(\mathrm{P} \ne \mathrm{NP}\) result.; next\_experiment: Write the finite verifier table for PrefixStableQrep entries: fields entry\_id, append\_time, old\_row, rep\_row, rank, certificate\_root, dependency\_list, first\_use\_time, \(charge_{to,Gamma}\), and rejection reason; then run it on one EqSAT duplicate-child example and one delayed-certificate private-row counterexample.
\item branch\_title: Hardness-Magnification Bridge Realization: prefix-strict Qrep support audit; cycle\_id: cycle-0153; learning: Cycle 153 upgrades Qrep from prefix-verifiable syntax to a theorem-shaped \(TransparentRowSupport_{N}\) rule case: quotient compression is valid exactly as strict-prefix, ranked, proof-carrying, dependency-closed dictionary use. The new falsifier is temporal; delayed same-theory evidence is already hidden support. This is a scoped parser/falsifier contribution, not a lower bound for \(H_{N}\) and not \(\mathrm{P} \ne \mathrm{NP}\) progress.; next\_experiment: Instantiate the strict-prefix Qrep verifier on a two-duplicate projected-dDNNF FormulaEval example and a delayed-certificate counterexample; compute \(\Gamma_{bEF}\), \(U_{N}\), \(X_{N}\), \(D_{Q}\), \(E_{pi}\), and the q/epsilon support-density consequence line by line.
\item branch\_title: Hardness-Magnification Bridge Realization; cycle\_id: cycle-0154; learning: Cycle 154 converts strict-prefix Qrep from a syntax convention into a replay theorem with a matching temporal falsifier. The valid contribution is a restricted \(TransparentRowSupport_{N}\) rule case and Formula Footprints supporting lemma: charged, ranked, dependency-closed dictionary compression preserves \(\Gamma\) line induction; delayed or hidden compression is exactly hidden support. This is route-level certification only, not a lower bound for \(H_{N}\) and not P vs NP progress.; next\_experiment: Run the prefix-scan verifier on one minimal proof with a valid two-row strict-prefix Qrep merge and on one identical proof with the certificate delayed by one line; compute \(\Gamma_{Qrep}\) in both cases and export the rejected line as the private-row falsifier pivot.
\end{itemize}

\subsection*{Padding and projection demagnification as a no-go route}

\paragraph{Physics picture.} The conceptual mechanism is proof-debt conservation. Padding changes the coordinate system, but projection back to SAT reveals that the supposedly weak source lower bound already encodes the hard target lower bound.

\paragraph{Active skills.}
\begin{itemize}
\item pvsnp-symbolic-complexity-auditor
\item pvsnp-impossibility-theorem-forger
\item contradiction-route-auditor
\item paper-contribution-readiness-judge
\item scientific-contribution-detector
\end{itemize}

\paragraph{Evidence from story.}
\begin{itemize}
\item branch\_title: Hardness-Magnification Bridge Realization; cycle\_id: cycle-0098; learning: Cycle 98 upgrades the prior packet from theorem slogan to dependency-checked inference structure: \(\Gamma_{bEF}\) is justified by separate EF-dependency, SATeq-link, bounded-SR, and sparse-row charges; \(\mathrm{Psys}_{bEF}\) gives a valid restricted falsifier theorem; full \(\mathrm{Psys}_{sym}\) is blocked exactly at \(FE_{Extract,N}\), not at generic naturalness or relativization folklore.; next\_experiment: Write the formal \(FE_{Extract,N}\) proposition in two columns: left column row-expanded FormulaEval axiom schema with \(N^{1+q+o(1)}\) charge, right column semantic-oracle version with the smallest private-row countermodel; verify which column matches each \(\mathrm{Psys}_{sym}\) rule.
\item branch\_title: Hardness-Magnification Bridge Realization; cycle\_id: cycle-0123; learning: Cycle 123 closes a proof-audit gap rather than repeating the magnification slogan: \(K_{dag}\) is an implementable least fixed point, projected-dDNNF gives a concrete transparent positive case, and the unread-enable-bit/private-row construction gives a precise necessity falsifier. The valid result is a scoped parser/falsifier theorem for beta-zero parity magnification; full semantic FormulaEval is still blocked exactly at \(RealSATInfluenceRealizer_{N}\).; next\_experiment: Write the referee-checkable trace for the projected-dDNNF example: list gate fields, parent pointers, projection maps, RowSet certificates, \(K_{0},K_{1}\),... closure additions, \(U_{N}\) and \(X_{N}\) charges, then flip one unread enable bit to exhibit the first-hidden-row countermodel and verify the \(\Gamma_{dag}\) threshold numerically for qepsilon.
\item branch\_title: Hardness-Magnification Bridge Realization; cycle\_id: cycle-0136; learning: Cycle 136 turns counterfactual declared TSR from a boundary slogan into a bitfield-level closed-support theorem: K\_CTSR(pi) is a finite fixed point, \(\Gamma_\mathrm{CTSR}\) is derived by line induction, and first-hidden-child adaptive pruning is the exact falsifier for extensional safety. The valid contribution remains a scoped parser/falsifier theorem for beta-zero parity magnification certificates; it does not prove \(H_{N}\) lower bounds or \(\mathrm{P} \ne \mathrm{NP}\).; next\_experiment: Select one concrete full semantic pruning rule beyond CTSR declarations, preferably adaptive Tseitin branch pruning with real SAT contradiction tests, and attempt either a charged \(RealSATInfluenceRealizer_{N}^{arb}\) or a minimal first-unchargeable-row countermodel.
\item branch\_title: Hardness-Magnification Bridge Realization; cycle\_id: cycle-0145; learning: Cycle 145 upgrades ranked QCTSR from a line-induction slogan to a finite verifier/minimality packet. Quotient compression is valid only when proof-carrying same-theory certificates are ranked, dependency-closed, and charged before use. The exact obstruction is the first rejected quotient merge, which yields a private-row falsifier. This strengthens the paper parser/countermodel boundary without proving \(H_{N}\) lower bounds or \(\mathrm{P} \ne \mathrm{NP}\).; next\_experiment: Instantiate the finite verifier on one toy QCTSR proof with two EqSAT-equivalent counterfactual TSR children: compute \(K_{Q}\), ranks, \(B_{Q}\), and \(\Gamma_\mathrm{QCTSR}\); then delete one certificate edge and explicitly write the two-world private-row falsifier.
\item branch\_title: Hardness-Magnification Bridge Realization; cycle\_id: cycle-0149; learning: Cycle 149 folds representative-stable ranked QCTSR into \(TransparentRowSupport_{N}\) as a valid rule case. The new contribution is not another magnification slogan: quotient compression is proof-carrying, rank-stratified, dependency-closed, and charged, while the first rejected merge gives a concrete private-row falsifier. Full semantic SAT pruning remains exactly at \(RealSATInfluenceRealizer_{N}^{arb}\).; next\_experiment: Write the standalone paper theorem for \(TransparentRowSupport_{N}^{Qrep}\): include the \(K_{Qrep}\) fixed-point verifier pseudocode, the \(\Gamma_{Qrep}\) ledger proof, the first-rejected-merge countermodel, and a one-page nonclaim box separating this parser boundary from \(H_{N}\) lower bounds and \(\mathrm{P} \ne \mathrm{NP}\).
\item branch\_title: Hardness-Magnification Bridge Realization; cycle\_id: cycle-0150; learning: Cycle 150 closes a proof-delivery upgrade rather than rediscovering magnification: Qrep quotient compression is valid exactly as a batched, ranked, proof-carrying same-theory dictionary charged in \(\Gamma_{Qrep}\). The sharp falsifier is the first unranked, cyclic, advice-chosen, or hidden semantic merge. The result is a scoped parser/falsifier theorem and Formula Footprints supporting lemma; it does not prove \(H_{N}\) lower bounds or \(\mathrm{P} \ne \mathrm{NP}\).; next\_experiment: Run the Tseitin-renaming Qrep thought experiment: instantiate two batches with one certified duplicate and one unlisted duplicate, compute \(K_{Qrep}\), C\_Q(D\_Q), and \(\Gamma_{Qrep}\) explicitly, then verify that the unlisted duplicate yields the first-hidden-row private-row falsifier while the certified duplicate preserves the line induction.
\end{itemize}

\subsection*{Meta-complexity carrier pressure}

\paragraph{Physics picture.} The meta-complexity carrier is a diagnostic medium: it can reveal whether a weak source lower bound really transfers, or whether constructivity, largeness, and nonuniform advice absorb the intended separation.

\paragraph{Active skills.}
\begin{itemize}
\item pvsnp-conceptual-pressure-localizer
\item pvsnp-symbolic-complexity-auditor
\item pvsnp-qualitative-complexity-reasoner
\item counterexample-hunter
\item kb-concept-bridge-planner
\end{itemize}

\paragraph{Evidence from story.}
\begin{itemize}
\item branch\_title: Hardness-Magnification Bridge Realization; cycle\_id: cycle-0149; learning: Cycle 149 folds representative-stable ranked QCTSR into \(TransparentRowSupport_{N}\) as a valid rule case. The new contribution is not another magnification slogan: quotient compression is proof-carrying, rank-stratified, dependency-closed, and charged, while the first rejected merge gives a concrete private-row falsifier. Full semantic SAT pruning remains exactly at \(RealSATInfluenceRealizer_{N}^{arb}\).; next\_experiment: Write the standalone paper theorem for \(TransparentRowSupport_{N}^{Qrep}\): include the \(K_{Qrep}\) fixed-point verifier pseudocode, the \(\Gamma_{Qrep}\) ledger proof, the first-rejected-merge countermodel, and a one-page nonclaim box separating this parser boundary from \(H_{N}\) lower bounds and \(\mathrm{P} \ne \mathrm{NP}\).
\item branch\_title: Hardness-Magnification Bridge Realization; cycle\_id: cycle-0150; learning: Cycle 150 closes a proof-delivery upgrade rather than rediscovering magnification: Qrep quotient compression is valid exactly as a batched, ranked, proof-carrying same-theory dictionary charged in \(\Gamma_{Qrep}\). The sharp falsifier is the first unranked, cyclic, advice-chosen, or hidden semantic merge. The result is a scoped parser/falsifier theorem and Formula Footprints supporting lemma; it does not prove \(H_{N}\) lower bounds or \(\mathrm{P} \ne \mathrm{NP}\).; next\_experiment: Run the Tseitin-renaming Qrep thought experiment: instantiate two batches with one certified duplicate and one unlisted duplicate, compute \(K_{Qrep}\), C\_Q(D\_Q), and \(\Gamma_{Qrep}\) explicitly, then verify that the unlisted duplicate yields the first-hidden-row private-row falsifier while the certified duplicate preserves the line induction.
\item branch\_title: Hardness-Magnification Bridge Realization; cycle\_id: cycle-0151; learning: Cycle 151 upgrades the batched Qrep patch with temporal prefix stability: quotient compression is valid only as append-only, ranked, proof-carrying dictionary use. The first pre-use, rebinding, cyclic, advice-chosen, or hidden semantic merge is an explicit falsifier, while full semantic pruning stays quarantined.; next\_experiment: Write the finite prefix-stable Qrep verifier: fields are batch\_id, rank, source\_row, representative\_row, certificate\_root, dependency list, first\_use\_line, and no\_rebind flag; prove the first rejected field yields the private-row falsifier.
\item branch\_title: Hardness-Magnification Bridge Realization: prefix-stable Qrep support closure; cycle\_id: cycle-0152; learning: Cycle 152 upgrades Qrep from an append-only convention to a prefix-verifiable rule case: quotient compression preserves the transparent parser only when certificate charge strictly precedes representative use. The exact no-go is temporal, not just acyclicity; delayed same-theory evidence is already a private-row falsifier. This closes a proof-audit loophole while preserving the nonclaim boundary: no \(H_{N}\) lower bound and no \(\mathrm{P} \ne \mathrm{NP}\) result.; next\_experiment: Write the finite verifier table for PrefixStableQrep entries: fields entry\_id, append\_time, old\_row, rep\_row, rank, certificate\_root, dependency\_list, first\_use\_time, \(charge_{to,Gamma}\), and rejection reason; then run it on one EqSAT duplicate-child example and one delayed-certificate private-row counterexample.
\item branch\_title: Hardness-Magnification Bridge Realization: prefix-strict Qrep support audit; cycle\_id: cycle-0153; learning: Cycle 153 upgrades Qrep from prefix-verifiable syntax to a theorem-shaped \(TransparentRowSupport_{N}\) rule case: quotient compression is valid exactly as strict-prefix, ranked, proof-carrying, dependency-closed dictionary use. The new falsifier is temporal; delayed same-theory evidence is already hidden support. This is a scoped parser/falsifier contribution, not a lower bound for \(H_{N}\) and not \(\mathrm{P} \ne \mathrm{NP}\) progress.; next\_experiment: Instantiate the strict-prefix Qrep verifier on a two-duplicate projected-dDNNF FormulaEval example and a delayed-certificate counterexample; compute \(\Gamma_{bEF}\), \(U_{N}\), \(X_{N}\), \(D_{Q}\), \(E_{pi}\), and the q/epsilon support-density consequence line by line.
\item branch\_title: Hardness-Magnification Bridge Realization; cycle\_id: cycle-0154; learning: Cycle 154 converts strict-prefix Qrep from a syntax convention into a replay theorem with a matching temporal falsifier. The valid contribution is a restricted \(TransparentRowSupport_{N}\) rule case and Formula Footprints supporting lemma: charged, ranked, dependency-closed dictionary compression preserves \(\Gamma\) line induction; delayed or hidden compression is exactly hidden support. This is route-level certification only, not a lower bound for \(H_{N}\) and not P vs NP progress.; next\_experiment: Run the prefix-scan verifier on one minimal proof with a valid two-row strict-prefix Qrep merge and on one identical proof with the certificate delayed by one line; compute \(\Gamma_{Qrep}\) in both cases and export the rejected line as the private-row falsifier pivot.
\end{itemize}

\subsection*{Nonblackbox or subpolynomial-interface escape search}

\paragraph{Physics picture.} The route needs a new interaction pattern with the target circuit. If the circuit is only queried as a black box at polynomially growing target length, the beta-ceiling reappears; a real escape must change that interface.

\paragraph{Active skills.}
\begin{itemize}
\item pvsnp-symbolic-complexity-auditor
\item pvsnp-impossibility-theorem-forger
\item pvsnp-qualitative-complexity-reasoner
\item kb-gap-abstraction-conjecturer
\item abducer-refinement-route-ranker
\end{itemize}

\paragraph{Evidence from story.}
\begin{itemize}
\item branch\_title: Hardness-Magnification Bridge Realization: Polylog Nonblackbox Compact Consumer; cycle\_id: cycle-0013; learning: Cycle 13 turns the superlogarithmic polylog window into an exact bridge: L=(log N)\textasciicircum{}d with \(d>1\) makes SAT circuits of size \(L^{c}\) amortize to \(N^{o(1)}\), while truth-table lookup costs 2\textasciicircum{}L=N\textasciicircum{}\{omega(1)\} and no longer supplies the old falsifier. The valid threshold is \(\rho<1+\epsilon\) for a target-clean nonblackbox compact consumer. The new hard obstruction is not folklore but the absence of an explicit carrier with no advice leakage and a non-natural \(N^{1+\epsilon}\) general-circuit lower bound.; next\_experiment: Instantiate \(PolyLogAudit_{N,d,V}\) with a universal-circuit verifier, compute the exact overhead exponent \(\rho\) and auxiliary advice length, and reject the instance immediately if \(H_{N}\) depends on a hidden SAT circuit or collapses to a black-box SAT query family.
\item branch\_title: Hardness-Magnification Bridge Realization: approximate relaxed bag separators; cycle\_id: cycle-0031; learning: Cycle 31 converts relaxed bag separators into an approximate threshold lemma: E bad bags cost NE gates, so errors below \(N^\epsilon\) are nonuniformly patchable. The valid source obligation is now an approximate distribution-local separator lower bound, not merely promise SAT hardness.; next\_experiment: Try a collision-balanced generator and estimate whether majority/hash-based local predicates yield \(s=N^{o(1)}\) with \(E<N^\epsilon\); if yes, quarantine this carrier family.
\item branch\_title: Relativizing locality-extraction gate for generated-OR carriers; cycle\_id: cycle-0048; learning: Cycle 48 makes the obstruction formal: separator-denominator hardness is only an upper-bound falsifier, not a valid lower-bound certificate for \(H_{N}\). The surviving route is direct C(H\_N) lower bounds or a nonrelativizing locality-extraction theorem with explicit threshold loss.; next\_experiment: Try to formulate a real-SAT locality extractor using witness self-reducibility: from a small circuit for \(H_{N}\), isolate a generated satisfiable query and build a one-sided predicate; audit algebrization and \(N^{o(1)}\) loss.
\item branch\_title: Hardness-Magnification Bridge Realization; cycle\_id: cycle-0081; learning: Cycle 81 separates the scoped occurrence-local theorem from the full \(\mathrm{EF}+\mathrm{SAT}_{U}\) ambition. The valid contribution is now sharper: private row spoofing is a genuine falsifier for occurrence-local functional certificates, while full \(\mathrm{EF}+\mathrm{SAT}_{U}\) requires a nonrelativizing semantic-to-occurrence extraction theorem because real \(\mathrm{SAT}_{U}\) axioms are not invariant under independent label flips.; next\_experiment: Formalize the semantic-coupling-width thought experiment as a lemma: for paired formulas F and F OR z, show exactly which \(\mathrm{SAT}_{U}\) implication lines break private-label invariance and whether their violations can be charged to complete forced rows below the \(N^{\epsilon}\) patch threshold.
\item branch\_title: Hardness-Magnification Bridge Realization; cycle\_id: cycle-0094; learning: Cycle 94 closes a concrete support-transparency subcase: bit-explicit EF extension variables can be parsed without support declarations, giving a valid transcript falsifier theorem. The full formula-semantic extraction fork is now sharper: the obstruction is not arbitrary EF extension variables but succinct FormulaEval-SAT comprehension, which defeats size-only occurrence parsing unless its row-index expansion is charged.; next\_experiment: Formalize \(\mathrm{Psys}_{sym}\) on finite length N: specify whether FormulaEval-SAT comprehension is a proof line, axiom schema, or extension rule; then prove either an expansion-charge parser with cost \(N^{1+q+o(1)}\) or a minimal countermodel showing surface occurrence completion is unsound.
\item branch\_title: Hardness-Magnification Bridge Realization; cycle\_id: cycle-0107; learning: Cycle 107 adds a real validation upgrade: transparent FormulaEval selectors no longer require interval-union structure for a valid theorem, because full syntactic scanning gives a certified \(B_{FE}=F_{FE}\) \(N^{1+q+o(1)}\) charge. This closes the \(q<\epsilon\) falsifier for a broader selector fragment and isolates the only missing improvement as \(AC0PrefixNonemptiness_{N}\) for output-sensitive enumeration. Full semantic FormulaEval remains quarantined by the private-row falsifier, not by generic barrier folklore.; next\_experiment: Formalize \(AC0PrefixNonemptiness_{N}\) as a finite lemma: for \(\mathrm{AC0}\) selectors on q log N row bits, test whether descriptor structure yields \(N^{o(1)}\) prefix emptiness; if not, state the selector-scan theorem as final and add \(\mathrm{AC0}\) output-sensitive enumeration to open problems.
\end{itemize}


\section{Strategy Progression Judge}

\label{sec:report-strategy-judge}

The 10-cycle payoff judge reads the target problem, literature evolution, latest analysis outputs, macro strategy, and contribution monitor.

\subsection*{Current strategy}

\paragraph{Title.} quantitative transfer obstruction for hardness magnification

\paragraph{Route family.} Complexity lower-bound transfer

\paragraph{Active object.} source exponent a, overhead q, target-query exponent \(\beta\), and the target SAT exponent c

\paragraph{Macro thesis.} Treat every proposed magnification route as an explicit transfer inequality from \(\mathrm{SAT}\in \mathrm{SIZE}(n^{c})\) or NP subset P/poly to a source-carrier upper bound.

\subsection*{Judgement summary}

\paragraph{Decision.} continue\_with\_targeted\_adaptation

\paragraph{Payoff potential.} high

\paragraph{Contribution pivot invariant.} If the next work only restates hardness magnification, MCSP promise, or barrier folklore, the branch is looping. If it computes a concrete transfer parameter, proves a route-specific no-go, or certifies a searchable certificate carrier with no advice leakage, it is moving toward a real scientific contribution.

\paragraph{Next decisive test.} Run the prefix-scan verifier on one minimal proof with a valid two-row strict-prefix Qrep merge and on one identical proof with the certificate delayed by one line; compute \(\Gamma_{Qrep}\) in both cases and export the rejected line as the private-row falsifier pivot.

\paragraph{When to stop.} Stop, backtrack, or quarantine when two sharpened attempts still yield only a bounded exponent window, advised hard strings, a natural/relativizing/algebrizing proof object, or a certificate-existence condition equivalent to the target lower bound.

\subsection*{Paper-level contribution readiness}

\paragraph{Readiness.}

\paragraph{Closest contribution.}

\paragraph{Missing for paper.}

\paragraph{Recommended paper move.}

\subsection*{Main problem alignment}

\paragraph{Alignment class.} bridge\_aligned

\paragraph{Alignment gain.} The window is valuable when it replaces slogans about hardness magnification or barriers with exact source objects, target objects, thresholds, and falsifiers.

\paragraph{Alignment loss.} The main danger is proving a bounded-window, advised, natural, relativizing, or target-equivalent proxy that cannot affect P vs NP route priority.

\paragraph{Proxy or transfer risk.} The branch remains aligned only if its carrier, barrier escape, or transfer formula preserves uniformity, valid lengths, and the all-exponent target rather than proving a proxy statement.

\paragraph{Repair action.} Every next analysis must name the complexity class target, source model, transfer formula or certificate object, and the falsifier that would force pivot.

\subsection*{Strategy payoff assessments}

\subsubsection*{macro.pvsnp-black-box-transfer-audit}

\paragraph{Payoff.}

\paragraph{Risk.}

\paragraph{Decision.}

\subsubsection*{macro.pvsnp-padding-projection-demagnification}

\paragraph{Payoff.}

\paragraph{Risk.}

\paragraph{Decision.}

\subsubsection*{macro.pvsnp-meta-complexity-carrier-pressure}

\paragraph{Payoff.}

\paragraph{Risk.}

\paragraph{Decision.}

\subsubsection*{macro.pvsnp-nonblackbox-escape-search}

\paragraph{Payoff.}

\paragraph{Risk.}

\paragraph{Decision.}


\section{Accumulated Proof Strategy Story}

\label{sec:report-physics-story}

Content-specific physics story from cycle-0001 through cycle-0154.

\paragraph{Latest learning.} Cycle 154 converts strict-prefix Qrep from a syntax convention into a replay theorem with a matching temporal falsifier. The valid contribution is a restricted \(TransparentRowSupport_{N}\) rule case and Formula Footprints supporting lemma: charged, ranked, dependency-closed dictionary compression preserves \(\Gamma\) line induction; delayed or hidden compression is exactly hidden support. This is route-level certification only, not a lower bound for \(H_{N}\) and not P vs NP progress.

\paragraph{Current obstruction.} The derivation does not cover arbitrary semantic SAT branch pruning or FormulaEval rules whose influencing rows are not determined by proof syntax before use; this is exactly \(RealSATInfluenceRealizer_{N}^{arb}\).

\paragraph{Current next analysis.} Run the prefix-scan verifier on one minimal proof with a valid two-row strict-prefix Qrep merge and on one identical proof with the certificate delayed by one line; compute \(\Gamma_{Qrep}\) in both cases and export the rejected line as the private-row falsifier pivot.

\subsection*{cycle-0001: Hardness-Magnification Bridge Realization}

\paragraph{Learning.} The exact quantitative transfer is C(Pad\_k(SAT)\_N)<=O(N+N\textasciicircum{}\{c/k\}); hence the barely-superlinear lower bound magnifies only up to target exponent k(1+epsilon). The valid bridge is useful as an exponent-calibrated route lemma, while the fixed-k P/poly claim is blocked by a concrete threshold falsifier rather than by vague folklore.

\paragraph{Pressure.} The bridge does not reach NP notin P/poly unless one proves a cofinal family of padded lower bounds or finds a non-artificial carrier whose interface compresses every polynomial target exponent below \(N^{1+\epsilon}\).

\paragraph{Next.} Instantiate the same exponent audit for a padded MKtP or GapMCSP carrier: name the verifier, compute Cook-Levin or direct-composition length blow-up, and test whether assuming \(\mathrm{SAT}\in \mathrm{SIZE}(n^{c})\) yields source circuits below \(N^{1+\epsilon}\) without model mismatch or nonuniform leakage.

\paragraph{When to stop.} Stop or quarantine this meta-complexity carrier branch only after MCSP/MKtP parameters, truth-table length N, size threshold s, verifier overhead, target SAT-circuit use, and natural-proof or advice-leakage status have been computed; quarantine if the transfer still produces only a finite exponent window or relies on a dense natural property without an escape certificate.

\paragraph{Do not stop prematurely if.} Do not stop before naming the carrier variant, deriving its verifier/reduction overhead, and checking whether a \(\mathrm{SAT}\in \mathrm{SIZE}(n^{c})\) assumption gives carrier circuits below the claimed source lower-bound threshold.

\paragraph{How it adds to the global problem.} This move adds to the P vs NP problem because it tests whether a meta-complexity carrier gives a weaker source obligation than the original P vs NP/P-poly target; it separates true hardness magnification from projection or padding arguments that only move proof debt; it keeps barrier status and nonuniform/advice leakage attached to the actual transfer theorem. The current branch-control decision is continue\_run\_next. Payoff criterion: For the MCSP/MKtP row branch, decide whether the certificate proves direct row hardness or only aggregate/rank-index hardness; record q\_extract=epsilon/(1+D) and q\_force=(D+epsilon)/(1+D) in the same parameter regime.

\subsection*{cycle-0002: Hardness-Magnification Bridge Realization}

\paragraph{Learning.} Cycle 2 turns the previous threshold collapse into an exact quantifier repair. The valid route lemma is: fixed k gives only a finite SAT exponent lower bound, while a common-epsilon lower bound for unbounded padded carriers \(H_{k}\) would imply \(\mathrm{NP}\not\subset \mathrm{P}/\mathrm{poly}\). The decisive falsifier is not folklore but the inequality \(c/k\ge 1+\epsilon\), plus the O(N) checker floor.

\paragraph{Pressure.} No known source theorem proves \(C_{H_{k}}(N)>N^{1+\epsilon}\) for unbounded k with a common \(\epsilon\) against general nonuniform circuits.

\paragraph{Next.} Test MCSP or proof-complexity carriers for producing the unbounded-k source family: formalize whether a weak lower bound for MCSP[2\textasciicircum{}\{N\textasciicircum{}\{1/k\}\}] transfers to \(C_{H_{k}}(N)>N^{1+\epsilon}\) without becoming natural or losing the strict exponent margin.

\paragraph{When to stop.} Stop or quarantine this meta-complexity carrier branch only after MCSP/MKtP parameters, truth-table length N, size threshold s, verifier overhead, target SAT-circuit use, and natural-proof or advice-leakage status have been computed; quarantine if the transfer still produces only a finite exponent window or relies on a dense natural property without an escape certificate.

\paragraph{Do not stop prematurely if.} Do not stop before naming the carrier variant, deriving its verifier/reduction overhead, and checking whether a \(\mathrm{SAT}\in \mathrm{SIZE}(n^{c})\) assumption gives carrier circuits below the claimed source lower-bound threshold.

\paragraph{How it adds to the global problem.} This move adds to the P vs NP problem because it tests whether a meta-complexity carrier gives a weaker source obligation than the original P vs NP/P-poly target; it separates true hardness magnification from projection or padding arguments that only move proof debt; it keeps barrier status and nonuniform/advice leakage attached to the actual transfer theorem; it searches the remaining plausible escape route: nonblackbox or subpolynomial-interface carriers. The current branch-control decision is continue\_run\_next. Payoff criterion: For the MCSP/MKtP row branch, decide whether the certificate proves direct row hardness or only aggregate/rank-index hardness; record q\_extract=epsilon/(1+D) and q\_force=(D+epsilon)/(1+D) in the same parameter regime.

\subsection*{cycle-0003: Hardness-Magnification Bridge Realization}

\paragraph{Learning.} Cycle 3 finds a valid single-carrier repair: growth padding \(N_{m}=(m+1)^{g(m+1)}\) converts an i.o. \(N^{1+\epsilon}\) lower bound for \(H_{g}\) into exclusion of every SAT exponent c. This is route-level progress, not P vs NP; the hard obligation is now the sparse valid-length lower bound.

\paragraph{Pressure.} The required source lower bound is a barely-superlinear lower bound for a sparse explicit NP carrier on valid lengths; proving it already encodes superpolynomial SAT lower bounds in the unpadded parameter.

\paragraph{Next.} Instantiate g(r)=ceil(log log(r+4)); prove the verifier and length test are poly(N), then analyze whether valid-slice lower bounds for \(H_{g}\) admit a proof-complexity or MCSP carrier formulation.

\paragraph{When to stop.} Stop or quarantine this meta-complexity carrier branch only after MCSP/MKtP parameters, truth-table length N, size threshold s, verifier overhead, target SAT-circuit use, and natural-proof or advice-leakage status have been computed; quarantine if the transfer still produces only a finite exponent window or relies on a dense natural property without an escape certificate.

\paragraph{Do not stop prematurely if.} Do not stop before naming the carrier variant, deriving its verifier/reduction overhead, and checking whether a \(\mathrm{SAT}\in \mathrm{SIZE}(n^{c})\) assumption gives carrier circuits below the claimed source lower-bound threshold.

\paragraph{How it adds to the global problem.} This move adds to the P vs NP problem because it tests whether a meta-complexity carrier gives a weaker source obligation than the original P vs NP/P-poly target; it separates true hardness magnification from projection or padding arguments that only move proof debt; it searches the remaining plausible escape route: nonblackbox or subpolynomial-interface carriers. The current branch-control decision is continue\_run\_next. Payoff criterion: For the MCSP/MKtP row branch, decide whether the certificate proves direct row hardness or only aggregate/rank-index hardness; record q\_extract=epsilon/(1+D) and q\_force=(D+epsilon)/(1+D) in the same parameter regime.

\subsection*{cycle-0004: Hardness-Magnification Bridge Realization}

\paragraph{Learning.} Cycle 4 proves a valid projection-sandwich lemma and demagnifies the sparse-padding route: \(N^{1+eps}\) lower bounds for \(H_{g}\) equal \(m^{(1+eps)g}\) SAT lower bounds. The bridge is correct but not a weak-source path toward P vs NP.

\paragraph{Pressure.} The bridge is valid, but its source obligation is not barely nontrivial in the original SAT length; it already asks for \(m^{\omega(1)}\) SAT lower bounds.

\paragraph{Next.} Test an MCSP or KT-complexity carrier with source model, threshold \(N^{1+eps}\), and target \(\mathrm{NP}\not\subset \mathrm{P}/\mathrm{poly}\); compute the reduction overhead and state whether the needed proof must be non-natural or non-algebrizing.

\paragraph{When to stop.} Stop or quarantine this meta-complexity carrier branch only after MCSP/MKtP parameters, truth-table length N, size threshold s, verifier overhead, target SAT-circuit use, and natural-proof or advice-leakage status have been computed; quarantine if the transfer still produces only a finite exponent window or relies on a dense natural property without an escape certificate.

\paragraph{Do not stop prematurely if.} Do not stop before naming the carrier variant, deriving its verifier/reduction overhead, and checking whether a \(\mathrm{SAT}\in \mathrm{SIZE}(n^{c})\) assumption gives carrier circuits below the claimed source lower-bound threshold.

\paragraph{How it adds to the global problem.} This move adds to the P vs NP problem because it tests whether a meta-complexity carrier gives a weaker source obligation than the original P vs NP/P-poly target; it separates true hardness magnification from projection or padding arguments that only move proof debt; it keeps barrier status and nonuniform/advice leakage attached to the actual transfer theorem. The current branch-control decision is continue\_run\_next. Payoff criterion: For the MCSP/MKtP row branch, decide whether the certificate proves direct row hardness or only aggregate/rank-index hardness; record q\_extract=epsilon/(1+D) and q\_force=(D+epsilon)/(1+D) in the same parameter regime.

\subsection*{cycle-0005: Hardness-Magnification Bridge Realization}

\paragraph{Learning.} Cycle 5 converts the padding quarantine into a general projection-carrier dichotomy. The valid bridge says \(N^{1+\epsilon}\) lower bounds transfer through ell(N) only when \(ell(N)^{c}\) is below the threshold. Positive polynomial embedding length gives only a bounded exponent window; subpolynomial embedding length gives all exponents but turns the source lower bound into \(ell(N)^{\omega(1)}\). Thus pure projection magnification fails as a weak-source path, with an explicit falsifier rather than a slogan.

\paragraph{Pressure.} Projection carriers cannot simultaneously keep the lower bound weak in the semantic source size and transfer it to exclusion of all polynomial SAT exponents.

\paragraph{Next.} Audit an MCSP-based route: set source object \(\mathrm{MCSP}_{n,s}\) on \(N=2^{n}-bit\) truth tables, choose a concrete circuit model for the lower bound, state the exact threshold \(N^{1+\epsilon}\) or smaller, and test whether the known magnification transfer is non-natural, target-preserving, and strong enough to imply \(\mathrm{SAT}\notin \mathrm{P}/\mathrm{poly}\) rather than only EXP or subexponential derandomization consequences.

\paragraph{When to stop.} Stop or quarantine this meta-complexity carrier branch only after MCSP/MKtP parameters, truth-table length N, size threshold s, verifier overhead, target SAT-circuit use, and natural-proof or advice-leakage status have been computed; quarantine if the transfer still produces only a finite exponent window or relies on a dense natural property without an escape certificate.

\paragraph{Do not stop prematurely if.} Do not stop before naming the carrier variant, deriving its verifier/reduction overhead, and checking whether a \(\mathrm{SAT}\in \mathrm{SIZE}(n^{c})\) assumption gives carrier circuits below the claimed source lower-bound threshold.

\paragraph{How it adds to the global problem.} This move adds to the P vs NP problem because it tests whether a meta-complexity carrier gives a weaker source obligation than the original P vs NP/P-poly target; it separates true hardness magnification from projection or padding arguments that only move proof debt; it keeps barrier status and nonuniform/advice leakage attached to the actual transfer theorem; it searches the remaining plausible escape route: nonblackbox or subpolynomial-interface carriers. The current branch-control decision is continue\_run\_next. Payoff criterion: For the MCSP/MKtP row branch, decide whether the certificate proves direct row hardness or only aggregate/rank-index hardness; record q\_extract=epsilon/(1+D) and q\_force=(D+epsilon)/(1+D) in the same parameter regime.

\subsection*{cycle-0006: Hardness-Magnification Bridge Realization, Cycle 6: MCSP verifier transfer audit}

\paragraph{Learning.} Cycle 6 shows that moving from projection padding to a nonprojection MCSP carrier does not by itself create major magnification. The valid transfer is \(C(\mathrm{MCSP}_{N^\alpha},N)=0\) caps any fixed barely-superlinear bridge below all-exponent P/poly separation. The failure is a concrete threshold falsifier, not barrier folklore.

\paragraph{Pressure.} For \(\mathrm{NP}\not\subset \mathrm{P}/\mathrm{poly}\) one must contradict every constant c, but fixed \(\epsilon\) and \(\alpha\) certify only c<(1+epsilon)/(1+alpha). Even \(\alpha=0\) leaves the input-reading ceiling \(c<1+\epsilon\).

\paragraph{Next.} Test a self-referential meta-complexity carrier, such as Gap-MCSP or MKtP with succinctly generated truth tables, where the target SAT circuit is used nonblack-box. Compute whether the induced carrier upper bound can avoid the multiplicative c factor; if not, quarantine it by the same beta-ceiling lemma.

\paragraph{When to stop.} Stop or quarantine this black-box substitution branch only after the source carrier, SAT-query length exponent \(\beta\), interface size exponent q, source-threshold exponent, and target class have been explicitly computed; quarantine if every sharpened interface still has positive \(\beta\) and therefore only excludes a bounded range of SAT circuit exponents rather than NP subset P/poly.

\paragraph{Do not stop prematurely if.} Do not stop before deriving the transfer inequality from \(\mathrm{SAT}\in \mathrm{SIZE}(n^{c})\) to the carrier upper bound, identifying whether \(\beta=0\), \(\beta>0\), or the argument is genuinely nonblackbox, and recording the barrier status of the proposed escape.

\paragraph{How it adds to the global problem.} This move adds to the P vs NP problem because it identifies whether the transfer loses all-exponent strength through a positive target-query exponent \(\beta\); it tests whether a meta-complexity carrier gives a weaker source obligation than the original P vs NP/P-poly target; it separates true hardness magnification from projection or padding arguments that only move proof debt; it keeps barrier status and nonuniform/advice leakage attached to the actual transfer theorem. The current branch-control decision is continue\_run\_next. Payoff criterion: For the MCSP/MKtP row branch, decide whether the certificate proves direct row hardness or only aggregate/rank-index hardness; record q\_extract=epsilon/(1+D) and q\_force=(D+epsilon)/(1+D) in the same parameter regime.

\subsection*{cycle-0007: Hardness-Magnification Bridge Realization: black-box substitution cap}

\paragraph{Learning.} Cycle 7 generalizes the MCSP and padding failures to an oracle-interface theorem. A barely-superlinear carrier lower bound transfers through a black-box SAT interface as max\{a,q+beta c\}; hence positive target-query length exponent \(\beta\) is the decisive threshold-collapse falsifier. The remaining valid route must either use subpolynomial target queries with a genuinely nonprojection carrier or abandon black-box substitution.

\paragraph{Pressure.} The only unquarantined escape is to construct a nonprojection carrier with \(N^{o(1)}\) target query length, \(sub-N^{1+\epsilon}\) interface size, and a valid \(N^{1+\epsilon}\) nonuniform lower bound.

\paragraph{Next.} Test a succinct local-consistency or proof-complexity carrier with \(N^{o(1)}\) SAT queries: compute a,q,beta, then run the projection-demagnification falsifier before claiming any magnification.

\paragraph{When to stop.} Stop or quarantine this black-box substitution branch only after the source carrier, SAT-query length exponent \(\beta\), interface size exponent q, source-threshold exponent, and target class have been explicitly computed; quarantine if every sharpened interface still has positive \(\beta\) and therefore only excludes a bounded range of SAT circuit exponents rather than NP subset P/poly.

\paragraph{Do not stop prematurely if.} Do not stop before deriving the transfer inequality from \(\mathrm{SAT}\in \mathrm{SIZE}(n^{c})\) to the carrier upper bound, identifying whether \(\beta=0\), \(\beta>0\), or the argument is genuinely nonblackbox, and recording the barrier status of the proposed escape.

\paragraph{How it adds to the global problem.} This move adds to the P vs NP problem because it identifies whether the transfer loses all-exponent strength through a positive target-query exponent \(\beta\); it tests whether a meta-complexity carrier gives a weaker source obligation than the original P vs NP/P-poly target; it separates true hardness magnification from projection or padding arguments that only move proof debt; it keeps barrier status and nonuniform/advice leakage attached to the actual transfer theorem. The current branch-control decision is continue\_run\_next. Payoff criterion: For the MCSP/MKtP row branch, decide whether the certificate proves direct row hardness or only aggregate/rank-index hardness; record q\_extract=epsilon/(1+D) and q\_force=(D+epsilon)/(1+D) in the same parameter regime.

\subsection*{cycle-0008: Hardness-Magnification Bridge Realization, cycle 8: subpolynomial-query threshold repair}

\paragraph{Learning.} Cycle 8 repairs the remaining route exactly: subpolynomial SAT-query length is sufficient to remove the c-dependent threshold term, so \(N^{1+\epsilon}\) lower bounds magnify only when \(max(a,q)<1+\epsilon\). The valid contribution is a quantified bridge plus a model-preservation gate; the unresolved hard task is finding a concrete carrier satisfying the interface and source lower bound without natural-proof or nonuniform-leakage failure.

\paragraph{Pressure.} No concrete non-natural carrier K is known with both a near-linear subpolynomial-query SAT interface and an \(N^{1+\epsilon}\) nonuniform circuit lower bound. MCSP-like truth-table carriers fail because their SAT-verification queries have polynomial, usually linear, dependence on N.

\paragraph{Next.} Search for non-truth-table carriers in \(\mathrm{P}^\mathrm{SAT}\) with near-linear ordinary work and subpolynomial SAT queries, e.g. succinct consistency or decompression-check problems, then compute A,Q,s and test against L8.1 before discussing P vs NP.

\paragraph{When to stop.} Stop or quarantine this black-box substitution branch only after the source carrier, SAT-query length exponent \(\beta\), interface size exponent q, source-threshold exponent, and target class have been explicitly computed; quarantine if every sharpened interface still has positive \(\beta\) and therefore only excludes a bounded range of SAT circuit exponents rather than NP subset P/poly.

\paragraph{Do not stop prematurely if.} Do not stop before deriving the transfer inequality from \(\mathrm{SAT}\in \mathrm{SIZE}(n^{c})\) to the carrier upper bound, identifying whether \(\beta=0\), \(\beta>0\), or the argument is genuinely nonblackbox, and recording the barrier status of the proposed escape.

\paragraph{How it adds to the global problem.} This move adds to the P vs NP problem because it identifies whether the transfer loses all-exponent strength through a positive target-query exponent \(\beta\); it tests whether a meta-complexity carrier gives a weaker source obligation than the original P vs NP/P-poly target; it keeps barrier status and nonuniform/advice leakage attached to the actual transfer theorem; it searches the remaining plausible escape route: nonblackbox or subpolynomial-interface carriers. The current branch-control decision is continue\_run\_next. Payoff criterion: For the MCSP/MKtP row branch, decide whether the certificate proves direct row hardness or only aggregate/rank-index hardness; record q\_extract=epsilon/(1+D) and q\_force=(D+epsilon)/(1+D) in the same parameter regime.

\subsection*{cycle-0009: Hardness-Magnification Bridge Realization: subpolynomial-interface demagnification}

\paragraph{Learning.} Cycle 9 converts the cycle-8 repair into a route failure certificate. Subpolynomial SAT queries remove the c-dependent upper-bound term, but the same transfer read backward says that proving the source \(N^{1+\epsilon}\) lower bound in general SIZE already proves \(\mathrm{SAT}_{\le L}\) has \(L^{\omega(1)}\) circuits lower bounds. The remaining valid work is not carrier search inside this black-box same-model interface; it is either a restricted-model closure theorem or a nonblackbox carrier.

\paragraph{Pressure.} The branch has no concrete carrier that is simultaneously a weak same-model source lower bound and target-preserving. Same-model sources demagnify; restricted sources need an unproved substitution-closure theorem.

\paragraph{Next.} Pivot to a nonblackbox object: test a proof-complexity or derandomization carrier where SAT in P/poly yields the required source upper bound through a named nonrelativizing theorem rather than through black-box SAT-query substitution.

\paragraph{When to stop.} Stop or quarantine this black-box substitution branch only after the source carrier, SAT-query length exponent \(\beta\), interface size exponent q, source-threshold exponent, and target class have been explicitly computed; quarantine if every sharpened interface still has positive \(\beta\) and therefore only excludes a bounded range of SAT circuit exponents rather than NP subset P/poly.

\paragraph{Do not stop prematurely if.} Do not stop before deriving the transfer inequality from \(\mathrm{SAT}\in \mathrm{SIZE}(n^{c})\) to the carrier upper bound, identifying whether \(\beta=0\), \(\beta>0\), or the argument is genuinely nonblackbox, and recording the barrier status of the proposed escape.

\paragraph{How it adds to the global problem.} This move adds to the P vs NP problem because it identifies whether the transfer loses all-exponent strength through a positive target-query exponent \(\beta\); it tests whether a meta-complexity carrier gives a weaker source obligation than the original P vs NP/P-poly target; it keeps barrier status and nonuniform/advice leakage attached to the actual transfer theorem; it searches the remaining plausible escape route: nonblackbox or subpolynomial-interface carriers. The current branch-control decision is continue\_run\_next. Payoff criterion: For the restricted-model branch, compute the compiler envelope from unrestricted SIZE(m\textasciicircum{}k) or P-time algorithms into the named model; promotion requires a bound above that envelope on the same SAT slice.

\subsection*{cycle-0010: Hardness-Magnification Bridge Realization}

\paragraph{Learning.} Cycle 10 proves a valid restricted-model closure bridge and names the precise target-preservation failure. Subpolynomial queries plus \(max(a,q)<1+\epsilon\) magnify a \(\Gamma\) lower bound to SAT notin \(\Gamma/\mathrm{poly}\) for every fixed exponent. This is real route-level progress, but it is not NP notin P/poly: the missing object is a nonuniform P/poly-to-Gamma lift with subthreshold tax, and its absence is the decisive falsifier for advertising the bridge as a P vs NP route.

\paragraph{Pressure.} The target-preserving lift from arbitrary P/poly SAT circuits to \(\Gamma-circuits\) is absent; without it, restricted-source magnification leaks the target model.

\paragraph{Next.} Test a non-natural nonblackbox carrier whose query length is \(N^{o(1)}\) and whose use of a SAT circuit is semantic rather than oracle substitution; compute whether its lift tax \(\delta\) is below the source slack \(\epsilon\).

\paragraph{When to stop.} Stop or quarantine this black-box substitution branch only after the source carrier, SAT-query length exponent \(\beta\), interface size exponent q, source-threshold exponent, and target class have been explicitly computed; quarantine if every sharpened interface still has positive \(\beta\) and therefore only excludes a bounded range of SAT circuit exponents rather than NP subset P/poly.

\paragraph{Do not stop prematurely if.} Do not stop before deriving the transfer inequality from \(\mathrm{SAT}\in \mathrm{SIZE}(n^{c})\) to the carrier upper bound, identifying whether \(\beta=0\), \(\beta>0\), or the argument is genuinely nonblackbox, and recording the barrier status of the proposed escape.

\paragraph{How it adds to the global problem.} This move adds to the P vs NP problem because it identifies whether the transfer loses all-exponent strength through a positive target-query exponent \(\beta\); it keeps barrier status and nonuniform/advice leakage attached to the actual transfer theorem; it searches the remaining plausible escape route: nonblackbox or subpolynomial-interface carriers. The current branch-control decision is continue\_run\_next. Payoff criterion: For this route, state the exact barrier premise it uses or escapes: relativizing, natural, algebrizing, nonuniform, or advice-leaking; promotion requires the escape to be attached to the specific transfer theorem.

\subsection*{cycle-0011: Hardness-Magnification Bridge Realization: lookup-lift collapse for restricted target repair}

\paragraph{Learning.} Cycle 11 turns the cycle-10 target split into a sharper impossibility triangle. To lift \(\Gamma\) lower bounds to P/poly by unrolling arbitrary SAT circuits, queries must be so short that every SAT query can also be answered by unconditional lookup inside \(\Gamma\). That gives \(H_{N}\) a subthreshold \(\Gamma\) upper bound and falsifies the source lower bound. Intermediate log-query lengths either collapse the source or preserve only exponent \(c\le 1\). Thus black-box formula-style P/poly repair is quarantined; progress now requires a genuinely nonblackbox carrier or a proof object where target circuits are used without independent short-query lookup.

\paragraph{Pressure.} The missing bridge is not merely an efficient P/poly-to-Gamma simulation; it must be efficient for SAT-derived query computations while not enabling small lookup for all short query predicates. Generic unrolling cannot satisfy both requirements.

\paragraph{Next.} Test a nonblackbox proof-complexity carrier: encode a SAT circuit as a verifier object whose correctness is checked globally rather than by short oracle leaves, and compute whether the verifier overhead stays below \(N^{1+\epsilon}\) without allowing \(\mathrm{SAT}_{\le L}\) truth-table substitution.

\paragraph{When to stop.} Stop or quarantine this black-box substitution branch only after the source carrier, SAT-query length exponent \(\beta\), interface size exponent q, source-threshold exponent, and target class have been explicitly computed; quarantine if every sharpened interface still has positive \(\beta\) and therefore only excludes a bounded range of SAT circuit exponents rather than NP subset P/poly.

\paragraph{Do not stop prematurely if.} Do not stop before deriving the transfer inequality from \(\mathrm{SAT}\in \mathrm{SIZE}(n^{c})\) to the carrier upper bound, identifying whether \(\beta=0\), \(\beta>0\), or the argument is genuinely nonblackbox, and recording the barrier status of the proposed escape.

\paragraph{How it adds to the global problem.} This move adds to the P vs NP problem because it identifies whether the transfer loses all-exponent strength through a positive target-query exponent \(\beta\); it keeps barrier status and nonuniform/advice leakage attached to the actual transfer theorem; it searches the remaining plausible escape route: nonblackbox or subpolynomial-interface carriers. The current branch-control decision is continue\_run\_next. Payoff criterion: For this route, state the exact barrier premise it uses or escapes: relativizing, natural, algebrizing, nonuniform, or advice-leaking; promotion requires the escape to be attached to the specific transfer theorem.

\subsection*{cycle-0012: Hardness-Magnification Bridge Realization}

\paragraph{Learning.} Cycle 12 identifies a real surviving window: superlogarithmic subpolynomial target lengths, especially L=(log N)\textasciicircum{}d with \(d>1\). In this window SAT circuits of size \(L^{c}\) amortize to \(N^{o(1)}\), but SAT truth tables cost \(2^{L}\) and no longer give the old subthreshold lookup falsifier. The bridge is valid only for nonblackbox compact-circuit consumers; it is not a P vs NP proof without the explicit carrier and lower bound.

\paragraph{Pressure.} No explicit carrier \(H_{N}\) and evaluator \(E_{N}\) are yet exhibited whose computation genuinely exploits compact SAT circuits rather than ordinary oracle answers, and no \(N^{1+\epsilon}\) lower bound for such an \(H_{N}\) is known.

\paragraph{Next.} Instantiate \(H_{N}\) as a proof-checking or self-reduction transcript language whose verifier can call a compact \(\mathrm{SAT}_{L}\) circuit nonblackbox; compute its exact overhead a and test whether full truth-table substitution or input-encoded D leakage gives a subthreshold upper bound.

\paragraph{When to stop.} Stop or quarantine this black-box substitution branch only after the source carrier, SAT-query length exponent \(\beta\), interface size exponent q, source-threshold exponent, and target class have been explicitly computed; quarantine if every sharpened interface still has positive \(\beta\) and therefore only excludes a bounded range of SAT circuit exponents rather than NP subset P/poly.

\paragraph{Do not stop prematurely if.} Do not stop before deriving the transfer inequality from \(\mathrm{SAT}\in \mathrm{SIZE}(n^{c})\) to the carrier upper bound, identifying whether \(\beta=0\), \(\beta>0\), or the argument is genuinely nonblackbox, and recording the barrier status of the proposed escape.

\paragraph{How it adds to the global problem.} This move adds to the P vs NP problem because it identifies whether the transfer loses all-exponent strength through a positive target-query exponent \(\beta\); it tests whether a meta-complexity carrier gives a weaker source obligation than the original P vs NP/P-poly target; it keeps barrier status and nonuniform/advice leakage attached to the actual transfer theorem; it searches the remaining plausible escape route: nonblackbox or subpolynomial-interface carriers. The current branch-control decision is continue\_run\_next. Payoff criterion: For the restricted-model branch, compute the compiler envelope from unrestricted SIZE(m\textasciicircum{}k) or P-time algorithms into the named model; promotion requires a bound above that envelope on the same SAT slice.

\subsection*{cycle-0013: Hardness-Magnification Bridge Realization: Polylog Nonblackbox Compact Consumer}

\paragraph{Learning.} Cycle 13 turns the superlogarithmic polylog window into an exact bridge: L=(log N)\textasciicircum{}d with \(d>1\) makes SAT circuits of size \(L^{c}\) amortize to \(N^{o(1)}\), while truth-table lookup costs 2\textasciicircum{}L=N\textasciicircum{}\{omega(1)\} and no longer supplies the old falsifier. The valid threshold is \(\rho<1+\epsilon\) for a target-clean nonblackbox compact consumer. The new hard obstruction is not folklore but the absence of an explicit carrier with no advice leakage and a non-natural \(N^{1+\epsilon}\) general-circuit lower bound.

\paragraph{Pressure.} No explicit carrier \(H_{N}\) is yet supplied with both a target-clean polylog nonblackbox compact consumer and an i.o. \(N^{1+\epsilon}\) general-circuit lower bound. The lower-bound proof also must be non-natural or otherwise outside the natural-proofs obstruction.

\paragraph{Next.} Instantiate \(PolyLogAudit_{N,d,V}\) with a universal-circuit verifier, compute the exact overhead exponent \(\rho\) and auxiliary advice length, and reject the instance immediately if \(H_{N}\) depends on a hidden SAT circuit or collapses to a black-box SAT query family.

\paragraph{When to stop.} Stop or quarantine this black-box substitution branch only after the source carrier, SAT-query length exponent \(\beta\), interface size exponent q, source-threshold exponent, and target class have been explicitly computed; quarantine if every sharpened interface still has positive \(\beta\) and therefore only excludes a bounded range of SAT circuit exponents rather than NP subset P/poly.

\paragraph{Do not stop prematurely if.} Do not stop before deriving the transfer inequality from \(\mathrm{SAT}\in \mathrm{SIZE}(n^{c})\) to the carrier upper bound, identifying whether \(\beta=0\), \(\beta>0\), or the argument is genuinely nonblackbox, and recording the barrier status of the proposed escape.

\paragraph{How it adds to the global problem.} This move adds to the P vs NP problem because it identifies whether the transfer loses all-exponent strength through a positive target-query exponent \(\beta\); it tests whether a meta-complexity carrier gives a weaker source obligation than the original P vs NP/P-poly target; it separates true hardness magnification from projection or padding arguments that only move proof debt; it keeps barrier status and nonuniform/advice leakage attached to the actual transfer theorem. The current branch-control decision is continue\_run\_next. Payoff criterion: For this route, state the exact barrier premise it uses or escapes: relativizing, natural, algebrizing, nonuniform, or advice-leaking; promotion requires the escape to be attached to the specific transfer theorem.

\subsection*{cycle-0014: Hardness-Magnification Bridge Realization: polylog generated-OR carrier}

\paragraph{Learning.} Cycle 14 identifies a concrete carrier template, \(H_{N}=OR\) of \(N^{q}\) generated SAT instances of length (log N)\textasciicircum{}d, and computes the exact transfer: \(\mathrm{SAT}\in \mathrm{SIZE}(n^{c})\) gives \(H_{N}\) in SIZE(N\textasciicircum{}\{max\{rho,q\}+o(1)\}). The valid magnification window is max\{rho,q\}1, exhaustive lookup is \(N^{q+\omega(1)}\), so the previous short-query falsifier is avoided, but the source lower bound remains the hard non-natural obligation.

\paragraph{Pressure.} The bridge is valid, but no explicit target-clean \(G_{N},A_{N}\) is known with a provable unrestricted circuit lower bound \(\mathrm{SIZE}(H_{N})>N^{1+\epsilon}\). A proof of that source bound must also avoid natural-proof density, relativizing oracle, and advice-leakage obstructions.

\paragraph{Next.} Instantiate \(G_{N}\) as a near-linear Cook-Levin micro-instance generator with \(M=N\) and L=(log N)\textasciicircum{}d; compute exact \(\rho\), verify \(H_{N}\) is in NP with poly(N) uniform verification, and test whether any structural simplification gives an \(N^{1+\epsilon}\) upper-bound falsifier.

\paragraph{When to stop.} Stop or quarantine this black-box substitution branch only after the source carrier, SAT-query length exponent \(\beta\), interface size exponent q, source-threshold exponent, and target class have been explicitly computed; quarantine if every sharpened interface still has positive \(\beta\) and therefore only excludes a bounded range of SAT circuit exponents rather than NP subset P/poly.

\paragraph{Do not stop prematurely if.} Do not stop before deriving the transfer inequality from \(\mathrm{SAT}\in \mathrm{SIZE}(n^{c})\) to the carrier upper bound, identifying whether \(\beta=0\), \(\beta>0\), or the argument is genuinely nonblackbox, and recording the barrier status of the proposed escape.

\paragraph{How it adds to the global problem.} This move adds to the P vs NP problem because it identifies whether the transfer loses all-exponent strength through a positive target-query exponent \(\beta\); it separates true hardness magnification from projection or padding arguments that only move proof debt; it keeps barrier status and nonuniform/advice leakage attached to the actual transfer theorem. The current branch-control decision is continue\_run\_next. Payoff criterion: For the high-q generated-OR carrier, compute the exact leakage exponent of the remaining witness-cover or unary-residual mechanism and compare it with 1+epsilon-q.

\subsection*{cycle-0015: Hardness-Magnification Bridge Realization}

\paragraph{Learning.} Cycle 15 adds the exact lookup boundary. Polylog queries always preserve all SAT polynomial exponents in the substitution term, but source plausibility requires q+A(log N)\textasciicircum{}\{d-1\} not fall below 1+epsilon. Thus d1 is safely valid from lookup.

\paragraph{Pressure.} No explicit no-advice generator \(G_{N}\) is known for which \(H_{N}\) has an i.o. \(N^{1+\epsilon}\) general-circuit lower bound by a non-natural method.

\paragraph{Next.} Instantiate the boundary cases \(d=1\) with \(A=1+\epsilon-q\) and \(d=2\); search for a generator whose \(H_{N}\) avoids known small-circuit constructions without smuggling SAT advice.

\paragraph{When to stop.} Stop or quarantine this black-box substitution branch only after the source carrier, SAT-query length exponent \(\beta\), interface size exponent q, source-threshold exponent, and target class have been explicitly computed; quarantine if every sharpened interface still has positive \(\beta\) and therefore only excludes a bounded range of SAT circuit exponents rather than NP subset P/poly.

\paragraph{Do not stop prematurely if.} Do not stop before deriving the transfer inequality from \(\mathrm{SAT}\in \mathrm{SIZE}(n^{c})\) to the carrier upper bound, identifying whether \(\beta=0\), \(\beta>0\), or the argument is genuinely nonblackbox, and recording the barrier status of the proposed escape.

\paragraph{How it adds to the global problem.} This move adds to the P vs NP problem because it identifies whether the transfer loses all-exponent strength through a positive target-query exponent \(\beta\); it keeps barrier status and nonuniform/advice leakage attached to the actual transfer theorem. The current branch-control decision is continue\_run\_next. Payoff criterion: For this route, state the exact barrier premise it uses or escapes: relativizing, natural, algebrizing, nonuniform, or advice-leaking; promotion requires the escape to be attached to the specific transfer theorem.

\subsection*{cycle-0016: Hardness-Magnification Bridge Realization}

\paragraph{Learning.} Cycle 16 sharpens the generated-OR route: the valid transfer to \(\mathrm{SAT}\notin \mathrm{P}/\mathrm{poly}\) needs only \(q,\rho<1+\epsilon\), but source plausibility additionally requires variable entropy q+v/log N at least 1+epsilon. The old \(2^{L}\) lookup audit was too weak; the real falsifier is low variable count.

\paragraph{Pressure.} No explicit target-clean generator is known whose OR carrier has a non-natural unrestricted \(N^{1+\epsilon}\) circuit lower bound while also passing the variable-entropy gate.

\paragraph{Next.} Instantiate \(G_{N}\) with high-essential-variable polylog 3CNF queries and test whether any structural simplification gives \(\mathrm{SIZE}(H_{N})\le N^{1+\epsilon-o(1)}\) despite v(N) superlogarithmic.

\paragraph{When to stop.} Stop or quarantine this P vs NP branch only after it states the source object, target class, complexity model, transfer inequality, quantifier strength, and barrier status; quarantine if two sharpened attempts still cannot say whether the route implies \(\mathrm{SAT}\notin \mathrm{P}/\mathrm{poly}\), \(\mathrm{NP}\not\subset \mathrm{P}/\mathrm{poly}\), or only a weaker/proxy lower bound.

\paragraph{Do not stop prematurely if.} Do not stop before the branch has produced a concrete source-to-target transfer theorem, an explicit obstruction lemma, or a barrier-aware reason why the route cannot reach the P vs NP/P/poly target.

\paragraph{How it adds to the global problem.} This move adds to the P vs NP problem because it identifies whether the transfer loses all-exponent strength through a positive target-query exponent \(\beta\); it keeps barrier status and nonuniform/advice leakage attached to the actual transfer theorem. The current branch-control decision is continue\_run\_next. Payoff criterion: For the high-q generated-OR carrier, compute the exact leakage exponent of the remaining witness-cover or unary-residual mechanism and compare it with 1+epsilon-q.

\subsection*{cycle-0017: Hardness-Magnification Bridge Realization: positioned-range entropy audit}

\paragraph{Learning.} Cycle 17 adds a sharper necessary condition for the generated-OR route: not only \(q,\rho<1+\epsilon\), but the positioned effective range entropy sum\_g \(min{|R_{g}|,2^{v}}\) must reach the lower-bound threshold. This converts low-range query generation into an explicit falsifier and preserves the valid target transfer only for high-entropy, target-clean carriers.

\paragraph{Pressure.} No explicit target-clean carrier is known with q,rho=1+epsilon-o(1), and a non-natural \(N^{1+\epsilon}\) lower bound against general circuits.

\paragraph{Next.} Instantiate \(K_{N}\) from a specific extractor-like generator of polylog-length CNFs and compute \(|R_{g}|\), \(v_{N}\), q, and \(\rho\) exactly to see whether the positioned entropy threshold can be met without advice leakage.

\paragraph{When to stop.} Stop or quarantine this black-box substitution branch only after the source carrier, SAT-query length exponent \(\beta\), interface size exponent q, source-threshold exponent, and target class have been explicitly computed; quarantine if every sharpened interface still has positive \(\beta\) and therefore only excludes a bounded range of SAT circuit exponents rather than NP subset P/poly.

\paragraph{Do not stop prematurely if.} Do not stop before deriving the transfer inequality from \(\mathrm{SAT}\in \mathrm{SIZE}(n^{c})\) to the carrier upper bound, identifying whether \(\beta=0\), \(\beta>0\), or the argument is genuinely nonblackbox, and recording the barrier status of the proposed escape.

\paragraph{How it adds to the global problem.} This move adds to the P vs NP problem because it identifies whether the transfer loses all-exponent strength through a positive target-query exponent \(\beta\); it keeps barrier status and nonuniform/advice leakage attached to the actual transfer theorem. The current branch-control decision is continue\_run\_next. Payoff criterion: For the high-q generated-OR carrier, compute the exact leakage exponent of the remaining witness-cover or unary-residual mechanism and compare it with 1+epsilon-q.

\subsection*{cycle-0018: Hardness-Magnification Bridge Realization: positioned entropy cover gate}

\paragraph{Learning.} Cycle 18 proves the exact positioned entropy cover upper bound. The route now has two independent gates: target transfer needs \(max{\rho,q}=N^{1+\epsilon-o(1)}\) after charging the generator. Low variable count or low positioned range is a concrete falsifier, not a folklore barrier.

\paragraph{Pressure.} The derivation proves only an upper-bound/falsifier lemma. It does not prove \(C(H_{N})>N^{1+\epsilon}\) for any explicit high-entropy carrier.

\paragraph{Next.} Instantiate the affine-disperser generated-OR family and compute \(B_{N}\) exactly as a function of \(\alpha\), q, and clause density; then audit whether any known lower-bound technique for this restricted \(H_{N}\) avoids the natural-proofs density condition.

\paragraph{When to stop.} Stop or quarantine this black-box substitution branch only after the source carrier, SAT-query length exponent \(\beta\), interface size exponent q, source-threshold exponent, and target class have been explicitly computed; quarantine if every sharpened interface still has positive \(\beta\) and therefore only excludes a bounded range of SAT circuit exponents rather than NP subset P/poly.

\paragraph{Do not stop prematurely if.} Do not stop before deriving the transfer inequality from \(\mathrm{SAT}\in \mathrm{SIZE}(n^{c})\) to the carrier upper bound, identifying whether \(\beta=0\), \(\beta>0\), or the argument is genuinely nonblackbox, and recording the barrier status of the proposed escape.

\paragraph{How it adds to the global problem.} This move adds to the P vs NP problem because it identifies whether the transfer loses all-exponent strength through a positive target-query exponent \(\beta\); it separates true hardness magnification from projection or padding arguments that only move proof debt; it keeps barrier status and nonuniform/advice leakage attached to the actual transfer theorem. The current branch-control decision is continue\_run\_next. Payoff criterion: For the high-q generated-OR carrier, compute the exact leakage exponent of the remaining witness-cover or unary-residual mechanism and compare it with 1+epsilon-q.

\subsection*{cycle-0019: Hardness-Magnification Bridge Realization: semantic range leakage gate}

\paragraph{Learning.} Cycle 19 sharpens the generated-OR route: positioned entropy is not the decisive source gate. The valid bridge uses the semantic cover \(S_{N}=sum_{j}\) C\_\{R\_\{N,j\}\}(SAT). Subthreshold \(S_{N}\) is a concrete falsifier for \(H_{N}\) lower bounds, while \(q,\rho<1+\epsilon\) still gives the target-preserving implication from a valid \(H_{N}\) lower bound to \(\mathrm{SAT}\notin \mathrm{P}/\mathrm{poly}\).

\paragraph{Pressure.} The derivation does not supply an explicit \(G_{N}\) for which both \(H_{N}\) has a general-circuit lower bound \(N^{1+\epsilon}\) and \(S_{N}=sum_{j}\) C\_\{R\_\{N,j\}\}(SAT) is provably above \(N^{1+\epsilon-o(1)}\).

\paragraph{Next.} Instantiate \(G_{N}\) with a concrete expander or extractor-based 3-CNF generator, compute \(R_{N,j}\) structure, and test whether SAT on those positioned ranges collapses to a known easy subclass or requires a new semantic lower-bound certificate.

\paragraph{When to stop.} Stop or quarantine this black-box substitution branch only after the source carrier, SAT-query length exponent \(\beta\), interface size exponent q, source-threshold exponent, and target class have been explicitly computed; quarantine if every sharpened interface still has positive \(\beta\) and therefore only excludes a bounded range of SAT circuit exponents rather than NP subset P/poly.

\paragraph{Do not stop prematurely if.} Do not stop before deriving the transfer inequality from \(\mathrm{SAT}\in \mathrm{SIZE}(n^{c})\) to the carrier upper bound, identifying whether \(\beta=0\), \(\beta>0\), or the argument is genuinely nonblackbox, and recording the barrier status of the proposed escape.

\paragraph{How it adds to the global problem.} This move adds to the P vs NP problem because it identifies whether the transfer loses all-exponent strength through a positive target-query exponent \(\beta\); it separates true hardness magnification from projection or padding arguments that only move proof debt; it keeps barrier status and nonuniform/advice leakage attached to the actual transfer theorem. The current branch-control decision is continue\_run\_next. Payoff criterion: For the high-q generated-OR carrier, compute the exact leakage exponent of the remaining witness-cover or unary-residual mechanism and compare it with 1+epsilon-q.

\subsection*{cycle-0020: Hardness-Magnification Bridge Realization: Semantic Block-Cover Gate}

\paragraph{Learning.} Cycle 20 sharpens the generated-OR route: the valid target bridge is unchanged, but the source plausibility gate must minimize over semantic block covers. Singleton range complexity is too weak because OR compression and nonuniform block advice can falsify an \(N^{1+\epsilon}\) source lower bound. This is route-level progress, not a proof of \(\mathrm{P} \ne \mathrm{NP}\).

\paragraph{Pressure.} No direct-sum theorem converts \(M_{N}\ge N^{1+\epsilon-o(1)}\) into \(C(H_{N})>N^{1+\epsilon}\). A circuit for \(H_{N}\) may compute the union of positive preimages without decomposing along any chosen partition.

\paragraph{Next.} Instantiate \(M_{N}\) for a nontrivial generator without dominated coordinates, e.g. a universal verifier family whose queries are Cook-Levin encodings of polylog-size computations, and compute whether block predicates admit compression below \(N^{1+\epsilon}\).

\paragraph{When to stop.} Stop or quarantine this P vs NP branch only after it states the source object, target class, complexity model, transfer inequality, quantifier strength, and barrier status; quarantine if two sharpened attempts still cannot say whether the route implies \(\mathrm{SAT}\notin \mathrm{P}/\mathrm{poly}\), \(\mathrm{NP}\not\subset \mathrm{P}/\mathrm{poly}\), or only a weaker/proxy lower bound.

\paragraph{Do not stop prematurely if.} Do not stop before the branch has produced a concrete source-to-target transfer theorem, an explicit obstruction lemma, or a barrier-aware reason why the route cannot reach the P vs NP/P/poly target.

\paragraph{How it adds to the global problem.} This move adds to the P vs NP problem because it keeps barrier status and nonuniform/advice leakage attached to the actual transfer theorem. The current branch-control decision is backtrack\_to\_sharpen. Payoff criterion: For the high-q generated-OR carrier, compute the exact leakage exponent of the remaining witness-cover or unary-residual mechanism and compare it with 1+epsilon-q.

\subsection*{cycle-0021: Hardness-Magnification Bridge Realization: Shared-Generator Semantic Block-Cover Gate}

\paragraph{Learning.} Cycle 21 replaces the previous block-cover warning by an exact shared-generator semantic-cover lemma. The valid target bridge is unchanged, but the source lower bound must beat \(C(G_{N})+min_{Pi}\) \(sum_{B}\) C\_\{R\_B\}(OR-SAT\_B). This creates a concrete falsifier: any subthreshold semantic block partition kills the generated-OR carrier before it can contribute to \(\mathrm{SAT}\notin \mathrm{P}/\mathrm{poly}\).

\paragraph{Pressure.} No candidate generator is known for which \(min_{Pi}\) \(sum_{B}\) C\_\{R\_B\}(OR-SAT\_B) can be proved at \(N^{1+\epsilon-o(1)}\) in general circuits.

\paragraph{Next.} Run the block-size thought experiment for \(b=N^\theta\): derive upper bounds on U(theta) from range cardinality, feature compression, and OR compression, then identify whether any \(\theta\) gives N\textasciicircum{}\{rho\}+N\textasciicircum{}\{q-theta\}U(theta)<N\textasciicircum{}\{1+epsilon\}.

\paragraph{When to stop.} Stop or quarantine this P vs NP branch only after it states the source object, target class, complexity model, transfer inequality, quantifier strength, and barrier status; quarantine if two sharpened attempts still cannot say whether the route implies \(\mathrm{SAT}\notin \mathrm{P}/\mathrm{poly}\), \(\mathrm{NP}\not\subset \mathrm{P}/\mathrm{poly}\), or only a weaker/proxy lower bound.

\paragraph{Do not stop prematurely if.} Do not stop before the branch has produced a concrete source-to-target transfer theorem, an explicit obstruction lemma, or a barrier-aware reason why the route cannot reach the P vs NP/P/poly target.

\paragraph{How it adds to the global problem.} This move adds to the P vs NP problem because it identifies whether the transfer loses all-exponent strength through a positive target-query exponent \(\beta\); it keeps barrier status and nonuniform/advice leakage attached to the actual transfer theorem. The current branch-control decision is continue\_run\_next. Payoff criterion: For the high-q generated-OR carrier, compute the exact leakage exponent of the remaining witness-cover or unary-residual mechanism and compare it with 1+epsilon-q.

\subsection*{cycle-0022: Hardness-Magnification Bridge Realization}

\paragraph{Learning.} Cycle 22 isolates the next exact gate: the generated-OR bridge is valid only if the source lower bound defeats nonuniform promise solvers on the actual SAT range. The new quantitative obstruction is \(T_{N}>N^{1+\epsilon-q-o(1)}\) for SAT restricted to \(U_{N}\), plus the earlier semantic block-cover threshold. This is route-level progress, not a proof of \(\mathrm{P} \ne \mathrm{NP}\).

\paragraph{Pressure.} No current argument proves the required promise-SAT or block-cover lower bound for an explicit high-entropy generated range \(U_{N}\); entropy and query length do not imply semantic circuit hardness.

\paragraph{Next.} Instantiate \(G_{N}\) with a high-entropy family of polylog-size 3-CNF formulas and compute the best promise solver or block-cover upper bound; vary \(\tau\) to locate the first threshold-respecting candidate.

\paragraph{When to stop.} Stop or quarantine this black-box substitution branch only after the source carrier, SAT-query length exponent \(\beta\), interface size exponent q, source-threshold exponent, and target class have been explicitly computed; quarantine if every sharpened interface still has positive \(\beta\) and therefore only excludes a bounded range of SAT circuit exponents rather than NP subset P/poly.

\paragraph{Do not stop prematurely if.} Do not stop before deriving the transfer inequality from \(\mathrm{SAT}\in \mathrm{SIZE}(n^{c})\) to the carrier upper bound, identifying whether \(\beta=0\), \(\beta>0\), or the argument is genuinely nonblackbox, and recording the barrier status of the proposed escape.

\paragraph{How it adds to the global problem.} This move adds to the P vs NP problem because it identifies whether the transfer loses all-exponent strength through a positive target-query exponent \(\beta\); it keeps barrier status and nonuniform/advice leakage attached to the actual transfer theorem. The current branch-control decision is continue\_run\_next. Payoff criterion: For the high-q generated-OR carrier, compute the exact leakage exponent of the remaining witness-cover or unary-residual mechanism and compare it with 1+epsilon-q.

\subsection*{cycle-0023: Hardness-magnification bridge realization: generated-OR range-balance gate}

\paragraph{Learning.} Cycle 23 adds an exact range-label balance gate. Even with q,rho in the valid target window, a generated-OR carrier is source-plausible only if both satisfiable and unsatisfiable parts of its actual generated range exceed \(N^{1+\epsilon-q-o(1)}\). This is a concrete falsifier derived from equality-DNF membership circuits, independent of vague barrier folklore.

\paragraph{Pressure.} No explicit generator is known whose \(H_{N}\) beats unrestricted \(N^{1+\epsilon}\) circuits while also defeating range-label imbalance, promise solvers on \(U_{N}\), and semantic block-cover compression.

\paragraph{Next.} Instantiate a template generator for polylog-length 3CNFs and estimate \(\theta=min\) label-support exponent; then compare it with 1+epsilon-q and with the best promise-SAT solver size on the generated template range.

\paragraph{When to stop.} Stop or quarantine this P vs NP branch only after it states the source object, target class, complexity model, transfer inequality, quantifier strength, and barrier status; quarantine if two sharpened attempts still cannot say whether the route implies \(\mathrm{SAT}\notin \mathrm{P}/\mathrm{poly}\), \(\mathrm{NP}\not\subset \mathrm{P}/\mathrm{poly}\), or only a weaker/proxy lower bound.

\paragraph{Do not stop prematurely if.} Do not stop before the branch has produced a concrete source-to-target transfer theorem, an explicit obstruction lemma, or a barrier-aware reason why the route cannot reach the P vs NP/P/poly target.

\paragraph{How it adds to the global problem.} This move adds to the P vs NP problem because it identifies whether the transfer loses all-exponent strength through a positive target-query exponent \(\beta\); it keeps barrier status and nonuniform/advice leakage attached to the actual transfer theorem. The current branch-control decision is continue\_run\_next. Payoff criterion: For the high-q generated-OR carrier, compute the exact leakage exponent of the remaining witness-cover or unary-residual mechanism and compare it with 1+epsilon-q.

\subsection*{cycle-0024: Hardness-Magnification Bridge Realization: hitting-cover gate for generated-OR carriers}

\paragraph{Learning.} Cycle 24 sharpens the generated-OR route from label balance to exact hitting-cover hardness. The valid target bridge is unchanged, but the source lower bound must now defeat nonuniform covers of positive fibers and negative product fibers of total size \(N^{1+\epsilon-o(1)}\). This gives a new concrete falsifier and shows that high semantic range entropy alone cannot justify the carrier.

\paragraph{Pressure.} No current argument lower-bounds the exact hitting-cover parameter \(\kappa\) for an explicit target-clean generator; raw semantic range size, variable entropy, and label balance all allow small-cover falsifiers.

\paragraph{Next.} Formalize the minimum cover parameter as a bipartite set-cover problem between input fibers and generated SAT labels; test whether known disperser/extractor-style generators can force \(kappa_{N}>N^{1+\epsilon-o(1)}\) without reintroducing nonuniform leakage.

\paragraph{When to stop.} Stop or quarantine this P vs NP branch only after it states the source object, target class, complexity model, transfer inequality, quantifier strength, and barrier status; quarantine if two sharpened attempts still cannot say whether the route implies \(\mathrm{SAT}\notin \mathrm{P}/\mathrm{poly}\), \(\mathrm{NP}\not\subset \mathrm{P}/\mathrm{poly}\), or only a weaker/proxy lower bound.

\paragraph{Do not stop prematurely if.} Do not stop before the branch has produced a concrete source-to-target transfer theorem, an explicit obstruction lemma, or a barrier-aware reason why the route cannot reach the P vs NP/P/poly target.

\paragraph{How it adds to the global problem.} This move adds to the P vs NP problem because it keeps barrier status and nonuniform/advice leakage attached to the actual transfer theorem. The current branch-control decision is backtrack\_to\_sharpen. Payoff criterion: For the high-q generated-OR carrier, compute the exact leakage exponent of the remaining witness-cover or unary-residual mechanism and compare it with 1+epsilon-q.

\subsection*{cycle-0025: Generated-OR magnification: selector-range leakage gate}

\paragraph{Learning.} Cycle 25 adds a shared-selector leakage gate. A generated-OR carrier can be easy via one compact selector and a promise SAT solver on only the selected image, even when additive semantic or hitting covers look expensive. Thus the valid source obligation is not merely high range entropy or label balance; it is lower bounding every selector-plus-selected-range-solver pair above \(N^{1+\epsilon-o(1)}\) while preserving \(q,\rho<1+\epsilon\) for the target bridge.

\paragraph{Pressure.} Need a direct-sum or shared-computation lower bound converting any small selector/promise-solver circuit into a subthreshold cycle-24 cover. No such theorem is available, and it is plausibly false for structured selectors.

\paragraph{Next.} Instantiate \(Sel_{N}\) for simple generator templates, including lexicographic first-satisfying-query generators and hash-bucketed query generators, and compute whether selected-image lookup already drops below \(N^{1+\epsilon}\).

\paragraph{When to stop.} Stop or quarantine this P vs NP branch only after it states the source object, target class, complexity model, transfer inequality, quantifier strength, and barrier status; quarantine if two sharpened attempts still cannot say whether the route implies \(\mathrm{SAT}\notin \mathrm{P}/\mathrm{poly}\), \(\mathrm{NP}\not\subset \mathrm{P}/\mathrm{poly}\), or only a weaker/proxy lower bound.

\paragraph{Do not stop prematurely if.} Do not stop before the branch has produced a concrete source-to-target transfer theorem, an explicit obstruction lemma, or a barrier-aware reason why the route cannot reach the P vs NP/P/poly target.

\paragraph{How it adds to the global problem.} This move adds to the P vs NP problem because it identifies whether the transfer loses all-exponent strength through a positive target-query exponent \(\beta\); it keeps barrier status and nonuniform/advice leakage attached to the actual transfer theorem. The current branch-control decision is continue\_run\_next. Payoff criterion: For the high-q generated-OR carrier, compute the exact leakage exponent of the remaining witness-cover or unary-residual mechanism and compare it with 1+epsilon-q.

\subsection*{cycle-0026: Hardness-Magnification Bridge Realization}

\paragraph{Learning.} Cycle 26 adds an approximate-selector gate. Exact witness selection was too weak as a source audit: even a selector wrong on \(N^{\delta}\) inputs collapses \(H_{N}\) below \(N^{1+\epsilon}\) when \(deltaN^{1+\epsilon-o(1)}\), while the clean target bridge \(q,\rho<1+\epsilon\) remains intact.

\paragraph{Pressure.} No concrete \(E_{N}\) is known for which \(Lambda_{N}>N^{1+\epsilon-o(1)}\) can be proved in unrestricted circuits; proving it appears to require non-natural, selector-resistant lower-bound machinery rather than entropy counting.

\paragraph{Next.} Instantiate the thought experiment with a block generator \(E_{N}\) whose coordinates are grouped into planted easy and hard ranges; compute the smallest \(\delta\) forced by coverage of positive inputs and test whether selected-range promise SAT stays subthreshold.

\paragraph{When to stop.} Stop or quarantine this black-box substitution branch only after the source carrier, SAT-query length exponent \(\beta\), interface size exponent q, source-threshold exponent, and target class have been explicitly computed; quarantine if every sharpened interface still has positive \(\beta\) and therefore only excludes a bounded range of SAT circuit exponents rather than NP subset P/poly.

\paragraph{Do not stop prematurely if.} Do not stop before deriving the transfer inequality from \(\mathrm{SAT}\in \mathrm{SIZE}(n^{c})\) to the carrier upper bound, identifying whether \(\beta=0\), \(\beta>0\), or the argument is genuinely nonblackbox, and recording the barrier status of the proposed escape.

\paragraph{How it adds to the global problem.} This move adds to the P vs NP problem because it identifies whether the transfer loses all-exponent strength through a positive target-query exponent \(\beta\); it keeps barrier status and nonuniform/advice leakage attached to the actual transfer theorem. The current branch-control decision is continue\_run\_next. Payoff criterion: For this route, state the exact barrier premise it uses or escapes: relativizing, natural, algebrizing, nonuniform, or advice-leaking; promotion requires the escape to be attached to the specific transfer theorem.

\subsection*{cycle-0027: Hardness-Magnification Bridge Realization: approximate list-selector leakage gate}

\paragraph{Learning.} Cycle 27 upgrades the source audit from approximate single selectors to approximate list selectors. The valid target-preserving bridge is unchanged, but a short list plus a promise solver can decide \(H_{N}\) without computing a canonical witness. Therefore any contribution-making lower-bound attempt must lower-bound \(Lambda_{N}^{list}\); otherwise the route has a concrete nonuniform-leakage falsifier.

\paragraph{Pressure.} Witness canonicalization is semantic; it is not available from the selector circuit alone and can be bypassed by list-wise OR evaluation.

\paragraph{Next.} Instantiate \(G_{N}\) with a high-entropy expander/hash-generated SAT family and compute whether any tau=O(epsilon/2) list selector has a selected image admitting a subthreshold promise solver.

\paragraph{When to stop.} Stop or quarantine this P vs NP branch only after it states the source object, target class, complexity model, transfer inequality, quantifier strength, and barrier status; quarantine if two sharpened attempts still cannot say whether the route implies \(\mathrm{SAT}\notin \mathrm{P}/\mathrm{poly}\), \(\mathrm{NP}\not\subset \mathrm{P}/\mathrm{poly}\), or only a weaker/proxy lower bound.

\paragraph{Do not stop prematurely if.} Do not stop before the branch has produced a concrete source-to-target transfer theorem, an explicit obstruction lemma, or a barrier-aware reason why the route cannot reach the P vs NP/P/poly target.

\paragraph{How it adds to the global problem.} This move adds to the P vs NP problem because it keeps barrier status and nonuniform/advice leakage attached to the actual transfer theorem. The current branch-control decision is backtrack\_to\_sharpen. Payoff criterion: For this route, state the exact barrier premise it uses or escapes: relativizing, natural, algebrizing, nonuniform, or advice-leaking; promotion requires the escape to be attached to the specific transfer theorem.

\subsection*{cycle-0028: Hardness-Magnification Bridge Realization}

\paragraph{Learning.} Cycle 28 identifies an additional leakage mode: many imperfect list selectors can collectively cover the positive fiber. The valid source audit is the ensemble cover complexity \(Xi_{N}^{ens}\); single-selector minimax lower bounds are insufficient for this generated-OR magnification route.

\paragraph{Pressure.} No lower-bound method is known for \(Xi_{N}^{ens}\) against fully nonuniform selected-image promise solvers; ensemble covers can exploit semantic clusters invisible to entropy counts.

\paragraph{Next.} Estimate \(Xi_{N}^{ens}\) for a random high-entropy generator by a union-bound model over ensembles, then isolate which part of the argument can survive explicitness and avoid natural-proof leakage.

\paragraph{When to stop.} Stop or quarantine this P vs NP branch only after it states the source object, target class, complexity model, transfer inequality, quantifier strength, and barrier status; quarantine if two sharpened attempts still cannot say whether the route implies \(\mathrm{SAT}\notin \mathrm{P}/\mathrm{poly}\), \(\mathrm{NP}\not\subset \mathrm{P}/\mathrm{poly}\), or only a weaker/proxy lower bound.

\paragraph{Do not stop prematurely if.} Do not stop before the branch has produced a concrete source-to-target transfer theorem, an explicit obstruction lemma, or a barrier-aware reason why the route cannot reach the P vs NP/P/poly target.

\paragraph{How it adds to the global problem.} This move adds to the P vs NP problem because it keeps barrier status and nonuniform/advice leakage attached to the actual transfer theorem. The current branch-control decision is backtrack\_to\_sharpen. Payoff criterion: For the high-q generated-OR carrier, compute the exact leakage exponent of the remaining witness-cover or unary-residual mechanism and compare it with 1+epsilon-q.

\subsection*{cycle-0029: Hardness-Magnification Bridge Realization: one-sided separating-ensemble leakage gate}

\paragraph{Learning.} Cycle 29 sharpens the generated-OR source audit from ensemble list selectors to one-sided selected-image separators. The valid target bridge still needs only \(q,\rho<1+\epsilon\), but the source lower bound must now beat \(Psi_{N}^{sep}\); small local labelers are a concrete falsifier even when they do not solve SAT globally.

\paragraph{Pressure.} No available lower-bound method rules out arbitrary selected-image separators \(B_{t}\). These labelers need not be SAT algorithms, so the target SAT hardness hypothesis does not exclude them.

\paragraph{Next.} Try to lower-bound or construct \(Psi_{N}^{sep}\) for a random-like generator: estimate whether selected positive and negative images admit \(N^{b}-size\) separators when \(R=N^{r}\) and \(s=N^{ell}\) are below the magnification threshold.

\paragraph{When to stop.} Stop or quarantine this P vs NP branch only after it states the source object, target class, complexity model, transfer inequality, quantifier strength, and barrier status; quarantine if two sharpened attempts still cannot say whether the route implies \(\mathrm{SAT}\notin \mathrm{P}/\mathrm{poly}\), \(\mathrm{NP}\not\subset \mathrm{P}/\mathrm{poly}\), or only a weaker/proxy lower bound.

\paragraph{Do not stop prematurely if.} Do not stop before the branch has produced a concrete source-to-target transfer theorem, an explicit obstruction lemma, or a barrier-aware reason why the route cannot reach the P vs NP/P/poly target.

\paragraph{How it adds to the global problem.} This move adds to the P vs NP problem because it keeps barrier status and nonuniform/advice leakage attached to the actual transfer theorem. The current branch-control decision is backtrack\_to\_sharpen. Payoff criterion: For the high-q generated-OR carrier, compute the exact leakage exponent of the remaining witness-cover or unary-residual mechanism and compare it with 1+epsilon-q.

\subsection*{cycle-0030: Hardness-Magnification Bridge Realization}

\paragraph{Learning.} Cycle 30 replaces the one-sided selected-image audit by the weaker and more dangerous relaxed bag-separator audit. The valid target bridge is unchanged, but the source lower bound must defeat predicates that merely hit one positive witness per positive bag and avoid all negative-bag formulas. This exposes a concrete nonuniform-leakage falsifier not captured by promise SAT hardness.

\paragraph{Pressure.} The converse is false or at least unproved: absence of small relaxed formula separators does not imply \(H_{N}\) is hard, because an arbitrary \(H_{N}\) circuit may exploit input-level structure that does not factor through a formula predicate.

\paragraph{Next.} Formalize \(SepHit_{N}\) as a bipartite hitting-avoidance problem between positive bags and negative formula occurrences, then test whether random-looking generators satisfy the required lower bound by a counting argument and where explicitness breaks the argument.

\paragraph{When to stop.} Stop or quarantine this P vs NP branch only after it states the source object, target class, complexity model, transfer inequality, quantifier strength, and barrier status; quarantine if two sharpened attempts still cannot say whether the route implies \(\mathrm{SAT}\notin \mathrm{P}/\mathrm{poly}\), \(\mathrm{NP}\not\subset \mathrm{P}/\mathrm{poly}\), or only a weaker/proxy lower bound.

\paragraph{Do not stop prematurely if.} Do not stop before the branch has produced a concrete source-to-target transfer theorem, an explicit obstruction lemma, or a barrier-aware reason why the route cannot reach the P vs NP/P/poly target.

\paragraph{How it adds to the global problem.} This move adds to the P vs NP problem because it keeps barrier status and nonuniform/advice leakage attached to the actual transfer theorem. The current branch-control decision is backtrack\_to\_sharpen. Payoff criterion: For this route, state the exact barrier premise it uses or escapes: relativizing, natural, algebrizing, nonuniform, or advice-leaking; promotion requires the escape to be attached to the specific transfer theorem.

\subsection*{cycle-0031: Hardness-Magnification Bridge Realization: approximate relaxed bag separators}

\paragraph{Learning.} Cycle 31 converts relaxed bag separators into an approximate threshold lemma: E bad bags cost NE gates, so errors below \(N^\epsilon\) are nonuniformly patchable. The valid source obligation is now an approximate distribution-local separator lower bound, not merely promise SAT hardness.

\paragraph{Pressure.} We still lack a lower bound against all distribution-local approximate bag separators; global promise-SAT hardness does not rule out local predicates tuned to the generated bags.

\paragraph{Next.} Try a collision-balanced generator and estimate whether majority/hash-based local predicates yield \(s=N^{o(1)}\) with \(E<N^\epsilon\); if yes, quarantine this carrier family.

\paragraph{When to stop.} Stop or quarantine this P vs NP branch only after it states the source object, target class, complexity model, transfer inequality, quantifier strength, and barrier status; quarantine if two sharpened attempts still cannot say whether the route implies \(\mathrm{SAT}\notin \mathrm{P}/\mathrm{poly}\), \(\mathrm{NP}\not\subset \mathrm{P}/\mathrm{poly}\), or only a weaker/proxy lower bound.

\paragraph{Do not stop prematurely if.} Do not stop before the branch has produced a concrete source-to-target transfer theorem, an explicit obstruction lemma, or a barrier-aware reason why the route cannot reach the P vs NP/P/poly target.

\paragraph{How it adds to the global problem.} This move adds to the P vs NP problem because it keeps barrier status and nonuniform/advice leakage attached to the actual transfer theorem. The current branch-control decision is backtrack\_to\_sharpen. Payoff criterion: For this route, state the exact barrier premise it uses or escapes: relativizing, natural, algebrizing, nonuniform, or advice-leaking; promotion requires the escape to be attached to the specific transfer theorem.

\subsection*{cycle-0032: Generated-OR magnification: correction-residual separator gate}

\paragraph{Learning.} Cycle 32 converts the approximate bag-separator gate into an exact correction-residual gate. The valid target bridge is unchanged, but the source lower bound must prove Phi\_N>1+epsilon-o(1): every small local formula predicate must leave a residual \(H_{N}\) xor \(S_{A}\) of high circuit complexity. Sparse error patching is only one falsifier; structured residual leakage is the stronger obstruction.

\paragraph{Pressure.} No known reduction converts a small residual correction \(R_{A}\) into a small SAT solver or separator on the generated formula range. This is the current formal gap.

\paragraph{Next.} Instantiate \(Phi_{N}\) for a tagged generated-OR carrier and compute whether simple tag predicates produce subthreshold residual circuits; then test an untagged high-entropy generator for the same residual leakage.

\paragraph{When to stop.} Stop or quarantine this padding/projection branch only after the length map N(n), projection overhead, source lower-bound quantifier and target exponent window are calculated; quarantine if the route merely moves the P vs NP difficulty into sparse valid-length lower bounds.

\paragraph{Do not stop prematurely if.} Do not stop before deriving the exponent transfer c \(\to \) c/k or its growth-padding analogue, and before stating whether the result excludes all polynomial SAT exponents or only a bounded exponent interval.

\paragraph{How it adds to the global problem.} This move adds to the P vs NP problem because it separates true hardness magnification from projection or padding arguments that only move proof debt; it keeps barrier status and nonuniform/advice leakage attached to the actual transfer theorem. The current branch-control decision is adapt\_then\_run. Payoff criterion: For the high-q generated-OR carrier, compute the exact leakage exponent of the remaining witness-cover or unary-residual mechanism and compare it with 1+epsilon-q.

\subsection*{cycle-0033: Hardness-Magnification Bridge Realization}

\paragraph{Learning.} Cycle 33 strengthens the generated-OR source obligation from single residual correction to feature-aggregator residual hardness. This is a contribution-making repair: it preserves the exact target bridge but exposes a broader nonuniform-leakage falsifier class that prior selector and separator audits did not cover.

\paragraph{Pressure.} The missing proof obligation is a lower bound on C(H\_N xor \(S_{A,B}\)) for all subthreshold feature families and aggregators. Promise SAT hardness on the generated range does not imply this residual lower bound, because B may exploit many weak local features jointly.

\paragraph{Next.} Formalize the restricted case where B is a depth-2 threshold or DNF over \(t=N^\tau\) local features, compute the exact q+alpha+beta leakage threshold, and search for an explicit low-residual construction on expander-walk generated SAT bags.

\paragraph{When to stop.} Stop or quarantine this black-box substitution branch only after the source carrier, SAT-query length exponent \(\beta\), interface size exponent q, source-threshold exponent, and target class have been explicitly computed; quarantine if every sharpened interface still has positive \(\beta\) and therefore only excludes a bounded range of SAT circuit exponents rather than NP subset P/poly.

\paragraph{Do not stop prematurely if.} Do not stop before deriving the transfer inequality from \(\mathrm{SAT}\in \mathrm{SIZE}(n^{c})\) to the carrier upper bound, identifying whether \(\beta=0\), \(\beta>0\), or the argument is genuinely nonblackbox, and recording the barrier status of the proposed escape.

\paragraph{How it adds to the global problem.} This move adds to the P vs NP problem because it identifies whether the transfer loses all-exponent strength through a positive target-query exponent \(\beta\); it separates true hardness magnification from projection or padding arguments that only move proof debt; it keeps barrier status and nonuniform/advice leakage attached to the actual transfer theorem. The current branch-control decision is continue\_run\_next. Payoff criterion: For the high-q generated-OR carrier, compute the exact leakage exponent of the remaining witness-cover or unary-residual mechanism and compare it with 1+epsilon-q.

\subsection*{cycle-0034: Hardness-Magnification Bridge Realization: tuple-feature residual gate}

\paragraph{Learning.} Cycle 34 adds a quantitative tuple-arity term to the generated-OR audit. The valid target bridge still needs only \(max{\rho,q}<1+\epsilon\), but the source lower bound must rule out every subthreshold tuple-feature transcript with exponent s+tq+alpha and residual exponent \(\lambda\). This yields a finite dense-arity falsifier search and a sharper nonuniform-leakage obstruction, without claiming \(\mathrm{P} \ne \mathrm{NP}\).

\paragraph{Pressure.} The derivation gives an exact upper-bound falsifier class, but it does not prove that arbitrary small circuits for \(H_{N}\) factor through low-arity tuple transcripts. The remaining source obligation is a genuine unrestricted-circuit lower bound or a non-natural theorem ruling out all residual systems and all non-factored circuits.

\paragraph{Next.} Instantiate the pair-pattern thought experiment for \(t=2\): choose q with \(q<1+\epsilon\) and test whether \(2q+\alpha<1+\epsilon\) permits a subthreshold pair separator \(S_{N}\); compute the residual complexity C(H\_N xor \(S_{N}\)) and check for hidden SAT lookup.

\paragraph{When to stop.} Stop or quarantine this black-box substitution branch only after the source carrier, SAT-query length exponent \(\beta\), interface size exponent q, source-threshold exponent, and target class have been explicitly computed; quarantine if every sharpened interface still has positive \(\beta\) and therefore only excludes a bounded range of SAT circuit exponents rather than NP subset P/poly.

\paragraph{Do not stop prematurely if.} Do not stop before deriving the transfer inequality from \(\mathrm{SAT}\in \mathrm{SIZE}(n^{c})\) to the carrier upper bound, identifying whether \(\beta=0\), \(\beta>0\), or the argument is genuinely nonblackbox, and recording the barrier status of the proposed escape.

\paragraph{How it adds to the global problem.} This move adds to the P vs NP problem because it identifies whether the transfer loses all-exponent strength through a positive target-query exponent \(\beta\); it tests whether a meta-complexity carrier gives a weaker source obligation than the original P vs NP/P-poly target; it separates true hardness magnification from projection or padding arguments that only move proof debt; it keeps barrier status and nonuniform/advice leakage attached to the actual transfer theorem. The current branch-control decision is continue\_run\_next. Payoff criterion: For the high-q generated-OR carrier, compute the exact leakage exponent of the remaining witness-cover or unary-residual mechanism and compare it with 1+epsilon-q.

\subsection*{cycle-0035: Hardness-Magnification Bridge Realization: high-q arity cutoff for generated-OR carriers}

\paragraph{Learning.} Cycle 35 finds a parameter repair: set q in ((1+epsilon)/2,1+epsilon) to preserve the generated-OR target bridge while forcing every dense tuple-feature transcript of arity at least two above the \(N^{1+\epsilon}\) threshold. The route-level lower-bound obligation is now the sharper unary residual minimax gate, not an infinite tuple-arity audit.

\paragraph{Pressure.} No proof is known that an explicit high-q generator defeats all unary range-local transcripts and subthreshold residual corrections. High q may even help unary predicates by presenting many chances to hit a satisfiable generated formula.

\paragraph{Next.} Formalize unary transcripts as monotone OR-of-range predicates, then test whether random-looking high-q generators make every subthreshold unary transcript leave residual circuit complexity above \(N^{1+\epsilon-o(1)}\) or whether a selected-image separator falsifier appears.

\paragraph{When to stop.} Stop or quarantine this P vs NP branch only after it states the source object, target class, complexity model, transfer inequality, quantifier strength, and barrier status; quarantine if two sharpened attempts still cannot say whether the route implies \(\mathrm{SAT}\notin \mathrm{P}/\mathrm{poly}\), \(\mathrm{NP}\not\subset \mathrm{P}/\mathrm{poly}\), or only a weaker/proxy lower bound.

\paragraph{Do not stop prematurely if.} Do not stop before the branch has produced a concrete source-to-target transfer theorem, an explicit obstruction lemma, or a barrier-aware reason why the route cannot reach the P vs NP/P/poly target.

\paragraph{How it adds to the global problem.} This move adds to the P vs NP problem because it keeps barrier status and nonuniform/advice leakage attached to the actual transfer theorem. The current branch-control decision is backtrack\_to\_sharpen. Payoff criterion: For the high-q generated-OR carrier, compute the exact leakage exponent of the remaining witness-cover or unary-residual mechanism and compare it with 1+epsilon-q.

\subsection*{cycle-0036: Hardness-Magnification Bridge Realization: position-indexed unary residual gate}

\paragraph{Learning.} Cycle 36 repairs the high-q unary reduction: binary and higher tuple features are threshold-blocked when \(q>(1+\epsilon)/2\), but shared-unary audits are not sound for general circuits. The valid remaining source obligation is a position-indexed unary transcript plus residual lower bound with total local budget \(U_{N}\). This is a sharper route-level gate and supplies an explicit falsifier class, without claiming \(\mathrm{P} \ne \mathrm{NP}\).

\paragraph{Pressure.} No lower-bound technique is supplied for the position-indexed unary residual minimax quantity. Proving it for an explicit \(H_{N}\) is essentially a general-circuit lower bound for a highly structured generated OR, not a consequence of the magnification bridge.

\paragraph{Next.} Formalize \(Iota_{N}^{pos}\) as a minimax transcript-residual measure and test whether expander-style slot mixing can rule out all low-budget position-indexed unary decompositions without becoming a natural-proof argument.

\paragraph{When to stop.} Stop or quarantine this P vs NP branch only after it states the source object, target class, complexity model, transfer inequality, quantifier strength, and barrier status; quarantine if two sharpened attempts still cannot say whether the route implies \(\mathrm{SAT}\notin \mathrm{P}/\mathrm{poly}\), \(\mathrm{NP}\not\subset \mathrm{P}/\mathrm{poly}\), or only a weaker/proxy lower bound.

\paragraph{Do not stop prematurely if.} Do not stop before the branch has produced a concrete source-to-target transfer theorem, an explicit obstruction lemma, or a barrier-aware reason why the route cannot reach the P vs NP/P/poly target.

\paragraph{How it adds to the global problem.} This move adds to the P vs NP problem because it keeps barrier status and nonuniform/advice leakage attached to the actual transfer theorem. The current branch-control decision is backtrack\_to\_sharpen. Payoff criterion: For the high-q generated-OR carrier, compute the exact leakage exponent of the remaining witness-cover or unary-residual mechanism and compare it with 1+epsilon-q.

\subsection*{cycle-0037: Hardness-Magnification Bridge Realization: approximate witness-cover gate}

\paragraph{Learning.} Cycle 37 sharpens the position-indexed unary residual obligation into an approximate witness-cover threshold. In the high-q generated-OR route, dense tuple leakage is threshold-blocked, but unary assignment covers of total size below \(N^{1+\epsilon}\) plus fewer than \(N^\epsilon\) missed positives still give exact circuits for \(H_{N}\) after equality patching. This is route-level progress: a valid candidate source lower bound must first prove high approximate witness-cover complexity, while the main \(\mathrm{P} \ne \mathrm{NP}\) conjecture remains untouched.

\paragraph{Pressure.} No explicit generated-OR family is known to have approximate position-indexed witness-cover complexity above this threshold, and defeating witness covers is still weaker than defeating arbitrary position-indexed unary predicates plus residual circuits.

\paragraph{Next.} Instantiate the thought experiment with a concrete planted/anti-planted CNF generator and compute the minimum position-indexed assignment hitting exponent h for its satisfiable range; compare h to 1+epsilon-q.

\paragraph{When to stop.} Stop or quarantine this P vs NP branch only after it states the source object, target class, complexity model, transfer inequality, quantifier strength, and barrier status; quarantine if two sharpened attempts still cannot say whether the route implies \(\mathrm{SAT}\notin \mathrm{P}/\mathrm{poly}\), \(\mathrm{NP}\not\subset \mathrm{P}/\mathrm{poly}\), or only a weaker/proxy lower bound.

\paragraph{Do not stop prematurely if.} Do not stop before the branch has produced a concrete source-to-target transfer theorem, an explicit obstruction lemma, or a barrier-aware reason why the route cannot reach the P vs NP/P/poly target.

\paragraph{How it adds to the global problem.} This move adds to the P vs NP problem because it keeps barrier status and nonuniform/advice leakage attached to the actual transfer theorem. The current branch-control decision is backtrack\_to\_sharpen. Payoff criterion: For the high-q generated-OR carrier, compute the exact leakage exponent of the remaining witness-cover or unary-residual mechanism and compare it with 1+epsilon-q.

\subsection*{cycle-0038: Hardness-Magnification Bridge Realization: realized-witness projection gate}

\paragraph{Learning.} Cycle 38 isolates a canonical unary leakage mechanism: all realized satisfying witnesses per generated position form a zero-error cover. Thus a valid high-q generated-OR route needs W\_N\textasciicircum{}*>N\textasciicircum{}\{1+epsilon-o(1)\} or a stronger approximate-cover lower bound. The target substitution bridge survives, but source plausibility is now governed by witness-projection entropy rather than formula-range entropy alone.

\paragraph{Pressure.} This proves only a necessary gate for the carrier lower bound. It does not supply the hard lower bound against all possible approximate witness covers.

\paragraph{Next.} Instantiate the planted-pool thought experiment and compute \(W_{N}^*\) for affine, low-degree, and pseudorandom assignment-pool generators; check whether any preserves \(q+a>1+\epsilon\) without introducing natural-proof or advice-leakage collapse.

\paragraph{When to stop.} Stop or quarantine this black-box substitution branch only after the source carrier, SAT-query length exponent \(\beta\), interface size exponent q, source-threshold exponent, and target class have been explicitly computed; quarantine if every sharpened interface still has positive \(\beta\) and therefore only excludes a bounded range of SAT circuit exponents rather than NP subset P/poly.

\paragraph{Do not stop prematurely if.} Do not stop before deriving the transfer inequality from \(\mathrm{SAT}\in \mathrm{SIZE}(n^{c})\) to the carrier upper bound, identifying whether \(\beta=0\), \(\beta>0\), or the argument is genuinely nonblackbox, and recording the barrier status of the proposed escape.

\paragraph{How it adds to the global problem.} This move adds to the P vs NP problem because it identifies whether the transfer loses all-exponent strength through a positive target-query exponent \(\beta\); it separates true hardness magnification from projection or padding arguments that only move proof debt; it keeps barrier status and nonuniform/advice leakage attached to the actual transfer theorem. The current branch-control decision is continue\_run\_next. Payoff criterion: For the high-q generated-OR carrier, compute the exact leakage exponent of the remaining witness-cover or unary-residual mechanism and compare it with 1+epsilon-q.

\subsection*{cycle-0039: High-q generated-OR carrier: partial-certificate cover gate}

\paragraph{Learning.} Cycle 39 sharpens the generated-OR source audit: full realized witness entropy is not the right gate because partial satisfying terms can compress many witnesses. The valid target bridge remains \(max{\rho,q}<1+\epsilon\), but the route now needs a superthreshold lower bound against position-indexed partial-term covers, and ultimately against sound local certificate covers.

\paragraph{Pressure.} No valid extraction theorem converts an arbitrary small circuit for \(H_{N}\) into a small full-witness set, and partial certificates give a concrete counter-mechanism.

\paragraph{Next.} Formalize \(tau_{N}^{term}\) as a set-cover problem over generated CNF ranges and test whether random-looking high-q generators can force tau\_N\textasciicircum{}\{term\}>N\textasciicircum{}\{1+epsilon\} without making \(G_{N}\) nonexplicit or violating \(\rho<1+\epsilon\).

\paragraph{When to stop.} Stop or quarantine this padding/projection branch only after the length map N(n), projection overhead, source lower-bound quantifier and target exponent window are calculated; quarantine if the route merely moves the P vs NP difficulty into sparse valid-length lower bounds.

\paragraph{Do not stop prematurely if.} Do not stop before deriving the exponent transfer c \(\to \) c/k or its growth-padding analogue, and before stating whether the result excludes all polynomial SAT exponents or only a bounded exponent interval.

\paragraph{How it adds to the global problem.} This move adds to the P vs NP problem because it identifies whether the transfer loses all-exponent strength through a positive target-query exponent \(\beta\); it separates true hardness magnification from projection or padding arguments that only move proof debt; it keeps barrier status and nonuniform/advice leakage attached to the actual transfer theorem. The current branch-control decision is continue\_run\_next. Payoff criterion: For the certificate-cover branch, compute K,E,L and the induced circuit upper bound C(H\_N)<=C(G\_N)+O(K(s+L\textasciicircum{}\{O(1)\})+EN); promotion requires every sound local cover to stay above the \(N^{1+\epsilon-o(1)}\) source threshold.

\subsection*{cycle-0040: Hardness-Magnification Bridge Realization: sound local certificate cover gate}

\paragraph{Learning.} Cycle 40 sharpens the source audit from partial certificate covers to approximate sound local certificate covers. The exact upper bound C(H\_N)<=C(G\_N)+O(K(s+L\textasciicircum{}\{O(1)\})+EN) makes any subthreshold SLC cover a concrete falsifier. The target bridge remains valid, but a contribution now requires a superthreshold lower bound against one-sided local predicates on the generated range, not just high witness or formula entropy.

\paragraph{Pressure.} No current argument proves that every one-sided local certificate predicate family of total size below \(N^{1+\epsilon-o(1)}\) misses at least \(N^{\epsilon-o(1)}\) positive bags on the actual generated range.

\paragraph{Next.} Compute the SLC cover exponent for generated formulas with bounded backdoor width \(b=\alpha\) log N and for small semantic classes such as Horn, 2-SAT, and bounded-width CNF; identify whether the threshold \(\alpha\ge 1+\epsilon-q\) is achievable without destroying the SAT-substitution bridge.

\paragraph{When to stop.} Stop or quarantine this black-box substitution branch only after the source carrier, SAT-query length exponent \(\beta\), interface size exponent q, source-threshold exponent, and target class have been explicitly computed; quarantine if every sharpened interface still has positive \(\beta\) and therefore only excludes a bounded range of SAT circuit exponents rather than NP subset P/poly.

\paragraph{Do not stop prematurely if.} Do not stop before deriving the transfer inequality from \(\mathrm{SAT}\in \mathrm{SIZE}(n^{c})\) to the carrier upper bound, identifying whether \(\beta=0\), \(\beta>0\), or the argument is genuinely nonblackbox, and recording the barrier status of the proposed escape.

\paragraph{How it adds to the global problem.} This move adds to the P vs NP problem because it identifies whether the transfer loses all-exponent strength through a positive target-query exponent \(\beta\); it keeps barrier status and nonuniform/advice leakage attached to the actual transfer theorem. The current branch-control decision is continue\_run\_next. Payoff criterion: For the certificate-cover branch, compute K,E,L and the induced circuit upper bound C(H\_N)<=C(G\_N)+O(K(s+L\textasciicircum{}\{O(1)\})+EN); promotion requires every sound local cover to stay above the \(N^{1+\epsilon-o(1)}\) source threshold.

\subsection*{cycle-0041: One-sided promise-separator gate for high-q generated-OR magnification}

\paragraph{Learning.} Cycle 41 isolates the exact semantic object behind the sound-local-certificate audit: Sep\_N\textasciicircum{}\{1s\}(E). Random term-cover calculations support the high-q polylog design, but they do not imply the needed separator lower bound. The main new obstruction is the term-cover versus one-sided semantic-separator gap.

\paragraph{Pressure.} No valid extractor converts arbitrary one-sided semantic separators into partial witnesses; range closure under restrictions is absent, so term-cover calibration cannot certify \(Sep_{N}^{1s}\) lower bounds.

\paragraph{Next.} Test a range-closure repair: augment \(G_{N}\) so restrictions of emitted formulas also appear in controlled positions, compute the q/rho overhead, and see whether this lets a one-sided separator self-reduce to a term cover without breaking the \(N^{1+\epsilon}\) threshold.

\paragraph{When to stop.} Stop or quarantine this P vs NP branch only after it states the source object, target class, complexity model, transfer inequality, quantifier strength, and barrier status; quarantine if two sharpened attempts still cannot say whether the route implies \(\mathrm{SAT}\notin \mathrm{P}/\mathrm{poly}\), \(\mathrm{NP}\not\subset \mathrm{P}/\mathrm{poly}\), or only a weaker/proxy lower bound.

\paragraph{Do not stop prematurely if.} Do not stop before the branch has produced a concrete source-to-target transfer theorem, an explicit obstruction lemma, or a barrier-aware reason why the route cannot reach the P vs NP/P/poly target.

\paragraph{How it adds to the global problem.} This move adds to the P vs NP problem because it keeps barrier status and nonuniform/advice leakage attached to the actual transfer theorem. The current branch-control decision is backtrack\_to\_sharpen. Payoff criterion: For the one-sided semantic-separator branch, compute the cap profile A\_N(s) for the named generator and predicate class; promotion requires a bound strong enough to force Sep\_N\textasciicircum{}\{1s\}(E)>N\textasciicircum{}\{1+epsilon-o(1)\}.

\subsection*{cycle-0042: Cycle 42: distributional cap certificate for one-sided semantic separators}

\paragraph{Learning.} Cycle 42 turns the one-sided semantic-separator obstruction into an exact distributional cap profile A\_N(s). The valid source route now needs a superthreshold cap-to-separator lower bound, not merely random witness or term-cover entropy. The target bridge is preserved, but any large one-sided cap is a concrete falsifier for the generated-OR lower bound.

\paragraph{Pressure.} Term-cover randomness does not imply a cap bound for arbitrary one-sided semantic predicates. The missing theorem is a one-sided conditional pseudorandomness statement for the distribution of generated formulas inside positive bags.

\paragraph{Next.} Formalize the cap profile for the unmarked random k-CNF generated-OR thought experiment: compute assignment-test hit exponents as a function of \(\beta\), then identify the first semantic predicate class beyond terms that can violate or support \(A_{N}(s)\le s/N^{1+\epsilon-o(1)}\).

\paragraph{When to stop.} Stop or quarantine this black-box substitution branch only after the source carrier, SAT-query length exponent \(\beta\), interface size exponent q, source-threshold exponent, and target class have been explicitly computed; quarantine if every sharpened interface still has positive \(\beta\) and therefore only excludes a bounded range of SAT circuit exponents rather than NP subset P/poly.

\paragraph{Do not stop prematurely if.} Do not stop before deriving the transfer inequality from \(\mathrm{SAT}\in \mathrm{SIZE}(n^{c})\) to the carrier upper bound, identifying whether \(\beta=0\), \(\beta>0\), or the argument is genuinely nonblackbox, and recording the barrier status of the proposed escape.

\paragraph{How it adds to the global problem.} This move adds to the P vs NP problem because it identifies whether the transfer loses all-exponent strength through a positive target-query exponent \(\beta\); it keeps barrier status and nonuniform/advice leakage attached to the actual transfer theorem. The current branch-control decision is continue\_run\_next. Payoff criterion: For the one-sided semantic-separator branch, compute the cap profile A\_N(s) for the named generator and predicate class; promotion requires a bound strong enough to force Sep\_N\textasciicircum{}\{1s\}(E)>N\textasciicircum{}\{1+epsilon-o(1)\}.

\subsection*{cycle-0043: Generated-OR residual cap gate for hardness magnification}

\paragraph{Learning.} Cycle 43 converts the cap profile into a residual fractional covering threshold. A valid source lower bound must rule out greedy one-sided cap covers with \(B_{N}(s,E)<N^{\epsilon-sigma-o(1)}\); absence of a large first cap is not enough.

\paragraph{Pressure.} No known method proves worst-case residual cap denominator \(B_{N}(s,E)>N^{\epsilon-sigma-o(1)}\) for semantic one-sided caps with nonuniform advice.

\paragraph{Next.} Try to construct residual caps for the planted/random 3CNF generator using partial-assignment certificates, and compute the resulting exponent b.

\paragraph{When to stop.} Stop or quarantine this black-box substitution branch only after the source carrier, SAT-query length exponent \(\beta\), interface size exponent q, source-threshold exponent, and target class have been explicitly computed; quarantine if every sharpened interface still has positive \(\beta\) and therefore only excludes a bounded range of SAT circuit exponents rather than NP subset P/poly.

\paragraph{Do not stop prematurely if.} Do not stop before deriving the transfer inequality from \(\mathrm{SAT}\in \mathrm{SIZE}(n^{c})\) to the carrier upper bound, identifying whether \(\beta=0\), \(\beta>0\), or the argument is genuinely nonblackbox, and recording the barrier status of the proposed escape.

\paragraph{How it adds to the global problem.} This move adds to the P vs NP problem because it identifies whether the transfer loses all-exponent strength through a positive target-query exponent \(\beta\); it keeps barrier status and nonuniform/advice leakage attached to the actual transfer theorem. The current branch-control decision is continue\_run\_next. Payoff criterion: For the one-sided semantic-separator branch, compute the cap profile A\_N(s) for the named generator and predicate class; promotion requires a bound strong enough to force Sep\_N\textasciicircum{}\{1s\}(E)>N\textasciicircum{}\{1+epsilon-o(1)\}.

\subsection*{cycle-0044: Hardness-Magnification Bridge Realization: residual cap-denominator gate}

\paragraph{Learning.} Cycle 44 converts the residual cap profile into a greedy cap-denominator theorem. A generated-OR source lower bound must now produce hard residual positive distributions against all one-sided local predicates; otherwise \(D_{N}\) \(N^{sigma}\) and E N give an explicit subthreshold upper-bound falsifier. This is route-level progress, not a main \(\mathrm{P} \ne \mathrm{NP}\) proof.

\paragraph{Pressure.} No lower bound is known showing that every residual positive distribution fools all one-sided generated-range predicates of size \(N^{sigma}\) at hit probability \(N^{-(1+\epsilon-sigma)}\).

\paragraph{Next.} Instantiate the cap audit for partial-assignment predicates of width w on L=(log N)\textasciicircum{}2 formulas: compute theta(w,sigma) for random-looking generated ranges, then check whether sigma+theta can stay above 1+epsilon while q,rho remain in the target-preserving window.

\paragraph{When to stop.} Stop or quarantine this black-box substitution branch only after the source carrier, SAT-query length exponent \(\beta\), interface size exponent q, source-threshold exponent, and target class have been explicitly computed; quarantine if every sharpened interface still has positive \(\beta\) and therefore only excludes a bounded range of SAT circuit exponents rather than NP subset P/poly.

\paragraph{Do not stop prematurely if.} Do not stop before deriving the transfer inequality from \(\mathrm{SAT}\in \mathrm{SIZE}(n^{c})\) to the carrier upper bound, identifying whether \(\beta=0\), \(\beta>0\), or the argument is genuinely nonblackbox, and recording the barrier status of the proposed escape.

\paragraph{How it adds to the global problem.} This move adds to the P vs NP problem because it identifies whether the transfer loses all-exponent strength through a positive target-query exponent \(\beta\); it keeps barrier status and nonuniform/advice leakage attached to the actual transfer theorem. The current branch-control decision is continue\_run\_next. Payoff criterion: For the one-sided semantic-separator branch, compute the cap profile A\_N(s) for the named generator and predicate class; promotion requires a bound strong enough to force Sep\_N\textasciicircum{}\{1s\}(E)>N\textasciicircum{}\{1+epsilon-o(1)\}.

\subsection*{cycle-0045: Hardness-Magnification Bridge Realization: Greedy Cap Dual Gate}

\paragraph{Learning.} Cycle 45 adds the missing logarithmic greedy-cover charge and converts the cap profile into an exact threshold: \(D_{sigma,E}<N^{\epsilon-sigma-o(1)}\) is a concrete upper-bound falsifier for the generated-OR source lower bound. A valid route must now prove hard residual positive distributions against all one-sided separators and also close the normal-form gap from arbitrary circuits to separator covers.

\paragraph{Pressure.} The derivation only rules out separator-cover upper bounds. It does not prove that every subthreshold exact circuit for \(H_{N}\) has a valid one-sided separator normal form.

\paragraph{Next.} Work in a restricted circuit model such as De Morgan formulas or \(\mathrm{AC0}[p]\) and test whether exact circuits for \(H_{N}\) admit an approximate one-sided separator normal form; compute the resulting \(D_{sigma,E}\) threshold for the same high-q polylog generated-OR carrier.

\paragraph{When to stop.} Stop or quarantine this meta-complexity carrier branch only after MCSP/MKtP parameters, truth-table length N, size threshold s, verifier overhead, target SAT-circuit use, and natural-proof or advice-leakage status have been computed; quarantine if the transfer still produces only a finite exponent window or relies on a dense natural property without an escape certificate.

\paragraph{Do not stop prematurely if.} Do not stop before naming the carrier variant, deriving its verifier/reduction overhead, and checking whether a \(\mathrm{SAT}\in \mathrm{SIZE}(n^{c})\) assumption gives carrier circuits below the claimed source lower-bound threshold.

\paragraph{How it adds to the global problem.} This move adds to the P vs NP problem because it tests whether a meta-complexity carrier gives a weaker source obligation than the original P vs NP/P-poly target; it keeps barrier status and nonuniform/advice leakage attached to the actual transfer theorem. The current branch-control decision is adapt\_then\_run. Payoff criterion: For the one-sided semantic-separator branch, compute the cap profile A\_N(s) for the named generator and predicate class; promotion requires a bound strong enough to force Sep\_N\textasciicircum{}\{1s\}(E)>N\textasciicircum{}\{1+epsilon-o(1)\}.

\subsection*{cycle-0046: Hardness-Magnification Bridge Realization: normal-form extraction gate}

\paragraph{Learning.} Cycle 46 converts the cycle-45 normal-form warning into a precise obstruction. Greedy cap and separator denominators give valid upper-bound falsifiers, but separator hardness is not by itself a source lower bound. A planted-control generated-OR carrier can have \(H_{N}\) easy while local formula separators remain hard, so the route now needs a nonrelativizing locality-extraction theorem or a direct arbitrary-circuit lower bound for \(H_{N}\).

\paragraph{Pressure.} No valid locality-extraction theorem maps arbitrary small circuits for \(H_{N}\) into one-sided formula-local separator covers with subthreshold overhead.

\paragraph{Next.} Formalize a restricted no-control generator class, such as block-product independent seeds with no shared label bit, and test whether an S-size circuit for \(OR_{j}\) SAT(G(x,j)) can be localized to a one-sided separator cover; compute the exact overhead in q,rho,sigma,E.

\paragraph{When to stop.} Stop or quarantine this meta-complexity carrier branch only after MCSP/MKtP parameters, truth-table length N, size threshold s, verifier overhead, target SAT-circuit use, and natural-proof or advice-leakage status have been computed; quarantine if the transfer still produces only a finite exponent window or relies on a dense natural property without an escape certificate.

\paragraph{Do not stop prematurely if.} Do not stop before naming the carrier variant, deriving its verifier/reduction overhead, and checking whether a \(\mathrm{SAT}\in \mathrm{SIZE}(n^{c})\) assumption gives carrier circuits below the claimed source lower-bound threshold.

\paragraph{How it adds to the global problem.} This move adds to the P vs NP problem because it tests whether a meta-complexity carrier gives a weaker source obligation than the original P vs NP/P-poly target; it keeps barrier status and nonuniform/advice leakage attached to the actual transfer theorem. The current branch-control decision is adapt\_then\_run. Payoff criterion: For the one-sided semantic-separator branch, compute the cap profile A\_N(s) for the named generator and predicate class; promotion requires a bound strong enough to force Sep\_N\textasciicircum{}\{1s\}(E)>N\textasciicircum{}\{1+epsilon-o(1)\}.

\subsection*{cycle-0047: Hardness-Magnification Bridge Realization: locality-extraction obstruction}

\paragraph{Learning.} Cycle 47 converts the locality-extraction warning into a formal relativizing obstruction. A planted-control oracle gives hard one-sided local separators but an O(1)-size generated-OR carrier, so separator hardness is not a valid source lower-bound certificate by itself. The target substitution bridge survives unchanged; the route must now prove a nonrelativizing extraction theorem for real SAT or attack C(H\_N) directly.

\paragraph{Pressure.} There is no valid black-box operation that extracts a formula-local one-sided separator from an arbitrary global circuit for \(H_{N}\). The missing step must exploit nonrelativizing real-SAT structure or be replaced by a direct arbitrary-circuit lower bound for \(H_{N}\).

\paragraph{Next.} Try a real-SAT extraction using self-reducibility or witness isolation: from a small circuit for \(H_{N}\), define a local predicate that accepts formulas whose forced witness path changes the global OR value, compute the overhead, and test it against the planted-control oracle falsifier.

\paragraph{When to stop.} Stop or quarantine this black-box substitution branch only after the source carrier, SAT-query length exponent \(\beta\), interface size exponent q, source-threshold exponent, and target class have been explicitly computed; quarantine if every sharpened interface still has positive \(\beta\) and therefore only excludes a bounded range of SAT circuit exponents rather than NP subset P/poly.

\paragraph{Do not stop prematurely if.} Do not stop before deriving the transfer inequality from \(\mathrm{SAT}\in \mathrm{SIZE}(n^{c})\) to the carrier upper bound, identifying whether \(\beta=0\), \(\beta>0\), or the argument is genuinely nonblackbox, and recording the barrier status of the proposed escape.

\paragraph{How it adds to the global problem.} This move adds to the P vs NP problem because it identifies whether the transfer loses all-exponent strength through a positive target-query exponent \(\beta\); it keeps barrier status and nonuniform/advice leakage attached to the actual transfer theorem. The current branch-control decision is continue\_run\_next. Payoff criterion: For the high-q generated-OR carrier, compute the exact leakage exponent of the remaining witness-cover or unary-residual mechanism and compare it with 1+epsilon-q.

\subsection*{cycle-0048: Relativizing locality-extraction gate for generated-OR carriers}

\paragraph{Learning.} Cycle 48 makes the obstruction formal: separator-denominator hardness is only an upper-bound falsifier, not a valid lower-bound certificate for \(H_{N}\). The surviving route is direct C(H\_N) lower bounds or a nonrelativizing locality-extraction theorem with explicit threshold loss.

\paragraph{Pressure.} No known nonrelativizing locality-extraction theorem converts an arbitrary small \(H_{N}\) circuit into local SAT-range separators with only \(N^{o(1)}\) loss.

\paragraph{Next.} Try to formulate a real-SAT locality extractor using witness self-reducibility: from a small circuit for \(H_{N}\), isolate a generated satisfiable query and build a one-sided predicate; audit algebrization and \(N^{o(1)}\) loss.

\paragraph{When to stop.} Stop or quarantine this P vs NP branch only after it states the source object, target class, complexity model, transfer inequality, quantifier strength, and barrier status; quarantine if two sharpened attempts still cannot say whether the route implies \(\mathrm{SAT}\notin \mathrm{P}/\mathrm{poly}\), \(\mathrm{NP}\not\subset \mathrm{P}/\mathrm{poly}\), or only a weaker/proxy lower bound.

\paragraph{Do not stop prematurely if.} Do not stop before the branch has produced a concrete source-to-target transfer theorem, an explicit obstruction lemma, or a barrier-aware reason why the route cannot reach the P vs NP/P/poly target.

\paragraph{How it adds to the global problem.} This move adds to the P vs NP problem because it identifies whether the transfer loses all-exponent strength through a positive target-query exponent \(\beta\); it keeps barrier status and nonuniform/advice leakage attached to the actual transfer theorem. The current branch-control decision is continue\_run\_next. Payoff criterion: For the high-q generated-OR carrier, compute the exact leakage exponent of the remaining witness-cover or unary-residual mechanism and compare it with 1+epsilon-q.

\subsection*{cycle-0049: Hardness-Magnification Bridge Realization: locality extraction threshold}

\paragraph{Learning.} Cycle 49 converts the surviving generated-OR route into an exact extraction-threshold problem. Separator-denominator hardness is a valid source certificate only if every small H circuit yields covers with denominator S/N\textasciicircum{}\{1+sigma\}+N\textasciicircum{}kappa and \(\kappa<\epsilon-sigma\). Neutralization gives a concrete extraction but is a demagnifying falsifier, because it asks for restricted-SAT hardness rather than a barely-superlinear H lower bound.

\paragraph{Pressure.} There is no valid operation converting an arbitrary global H circuit into local SAT-range separators without either neutralizing the generator or using a new real-SAT, nonrelativizing locality principle.

\paragraph{Next.} Formalize one candidate extraction using SAT self-reducibility plus generator local decodability, then compute whether the required decodability already implies neutralizing embeddings and hence demagnification.

\paragraph{When to stop.} Stop or quarantine this P vs NP branch only after it states the source object, target class, complexity model, transfer inequality, quantifier strength, and barrier status; quarantine if two sharpened attempts still cannot say whether the route implies \(\mathrm{SAT}\notin \mathrm{P}/\mathrm{poly}\), \(\mathrm{NP}\not\subset \mathrm{P}/\mathrm{poly}\), or only a weaker/proxy lower bound.

\paragraph{Do not stop prematurely if.} Do not stop before the branch has produced a concrete source-to-target transfer theorem, an explicit obstruction lemma, or a barrier-aware reason why the route cannot reach the P vs NP/P/poly target.

\paragraph{How it adds to the global problem.} This move adds to the P vs NP problem because it keeps barrier status and nonuniform/advice leakage attached to the actual transfer theorem. The current branch-control decision is backtrack\_to\_sharpen. Payoff criterion: For the high-q generated-OR carrier, compute the exact leakage exponent of the remaining witness-cover or unary-residual mechanism and compare it with 1+epsilon-q.

\subsection*{cycle-0050: Hardness-Magnification Bridge Realization: Relativizing Extraction Falsified}

\paragraph{Learning.} Cycle 50 gives a precise obstruction: small \(H_{N}\) circuits can coexist with superthreshold one-sided separator complexity in a structured singleton oracle. Thus the extraction bridge must be explicitly nonrelativizing and real-SAT-specific; separator hardness alone is not a valid source lower-bound certificate.

\paragraph{Pressure.} The argument is an oracle obstruction, not a real-SAT lower bound. Real CNF satisfiability may have nonrelativizing syntactic structure unavailable to singleton oracle labels.

\paragraph{Next.} Audit real-SAT self-reduction or proof-certificate structure as a possible nonrelativizing extractor, and compute whether its verifier/certificate overhead stays below \(N^{1+\epsilon}\) rather than collapsing to restricted-SAT hardness.

\paragraph{When to stop.} Stop or quarantine this P vs NP branch only after it states the source object, target class, complexity model, transfer inequality, quantifier strength, and barrier status; quarantine if two sharpened attempts still cannot say whether the route implies \(\mathrm{SAT}\notin \mathrm{P}/\mathrm{poly}\), \(\mathrm{NP}\not\subset \mathrm{P}/\mathrm{poly}\), or only a weaker/proxy lower bound.

\paragraph{Do not stop prematurely if.} Do not stop before the branch has produced a concrete source-to-target transfer theorem, an explicit obstruction lemma, or a barrier-aware reason why the route cannot reach the P vs NP/P/poly target.

\paragraph{How it adds to the global problem.} This move adds to the P vs NP problem because it keeps barrier status and nonuniform/advice leakage attached to the actual transfer theorem. The current branch-control decision is backtrack\_to\_sharpen. Payoff criterion: For this route, state the exact barrier premise it uses or escapes: relativizing, natural, algebrizing, nonuniform, or advice-leaking; promotion requires the escape to be attached to the specific transfer theorem.

\subsection*{cycle-0051: Hardness-Magnification Bridge Realization - cycle 51 direct-sum demagnification audit}

\paragraph{Learning.} Cycle 51 identifies a new quantitative conservation law: the same gap \(g_{N}\) that protects the target transfer against all SAT polynomial exponents becomes the per-slot lower-bound burden for any direct-sum source proof. Hence slot-wise SAT hardness cannot be the weak source; a valid route must prove aggregate \(H_{N}\) hardness or supply a nonrelativizing extraction theorem.

\paragraph{Pressure.} This proves only a quarantine for factorized source certificates. It does not rule out a genuinely aggregate arbitrary-circuit lower bound for \(H_{N}\) or a nonrelativizing extraction theorem using real SAT structure.

\paragraph{Next.} Try to formalize an aggregate lower-bound normal form for \(H_{N}\) that distinguishes global OR computation from \(M_{N}\) independent \(\mathrm{SAT}_{L}\) copies, and test it against the planted-control oracle before using separator hardness.

\paragraph{When to stop.} Stop or quarantine this padding/projection branch only after the length map N(n), projection overhead, source lower-bound quantifier and target exponent window are calculated; quarantine if the route merely moves the P vs NP difficulty into sparse valid-length lower bounds.

\paragraph{Do not stop prematurely if.} Do not stop before deriving the exponent transfer c \(\to \) c/k or its growth-padding analogue, and before stating whether the result excludes all polynomial SAT exponents or only a bounded exponent interval.

\paragraph{How it adds to the global problem.} This move adds to the P vs NP problem because it separates true hardness magnification from projection or padding arguments that only move proof debt; it keeps barrier status and nonuniform/advice leakage attached to the actual transfer theorem. The current branch-control decision is adapt\_then\_run. Payoff criterion: For this route, state the exact barrier premise it uses or escapes: relativizing, natural, algebrizing, nonuniform, or advice-leaking; promotion requires the escape to be attached to the specific transfer theorem.

\subsection*{cycle-0052: Hardness-Magnification Bridge Realization: XOR Carrier Repair}

\paragraph{Learning.} Cycle 52 contributes a concrete repair: XOR aggregation preserves the valid polylog-query magnification threshold while invalidating the OR-specific one-sided separator obstruction. The new source burden is sharper: prove aggregate parity-carrier hardness or rule out subthreshold parity sketches and sparse corrections. This is route-level progress only, with explicit falsifier tests.

\paragraph{Pressure.} There is no nonrelativizing theorem extracting a parity sketch, low-rank transcript, or other analyzable normal form from an arbitrary small circuit for \(H_{N}^{oplus}\); direct C(H\_N\textasciicircum{}oplus)>N\textasciicircum{}\{1+epsilon\} remains unproved.

\paragraph{Next.} In the ideal random-label model for \(G_{N}\), count subthreshold affine sketches of size \(N^\alpha\) and estimate the tau,alpha,error threshold at which a union bound could certify parity-sketch resistance without pretending to prove the actual explicit carrier lower bound.

\paragraph{When to stop.} Stop or quarantine this padding/projection branch only after the length map N(n), projection overhead, source lower-bound quantifier and target exponent window are calculated; quarantine if the route merely moves the P vs NP difficulty into sparse valid-length lower bounds.

\paragraph{Do not stop prematurely if.} Do not stop before deriving the exponent transfer c \(\to \) c/k or its growth-padding analogue, and before stating whether the result excludes all polynomial SAT exponents or only a bounded exponent interval.

\paragraph{How it adds to the global problem.} This move adds to the P vs NP problem because it identifies whether the transfer loses all-exponent strength through a positive target-query exponent \(\beta\); it separates true hardness magnification from projection or padding arguments that only move proof debt; it keeps barrier status and nonuniform/advice leakage attached to the actual transfer theorem. The current branch-control decision is continue\_run\_next. Payoff criterion: For the proof-complexity branch, separate fixed-family proof hardness from language hardness by naming the proof system, shortcut algorithm, and semantic lower-bound target.

\subsection*{cycle-0053: Hardness-Magnification Bridge Realization: XOR carrier affine-incidence audit}

\paragraph{Learning.} Cycle 53 converts the XOR repair into an exact affine-incidence conservation law. The carrier equals the \(F_{2}\) inner product of the SAT label vector with a generated parity-incidence vector; low-dimensional computable incidence structure makes the carrier easy by hardwiring basis parities. This is a new concrete falsifier for the source lower bound, while the polylog-query target bridge remains valid.

\paragraph{Pressure.} A small circuit for \(H_{N}^{oplus}\) computes only the scalar inner product ; it need not reveal p\_x, a low-rank coordinate map, or any local SAT separator. The extraction is false in planted-label oracle worlds and would need real-SAT, nonrelativizing content.

\paragraph{Next.} Instantiate \(G_{N,\theta}\) with linear, pairwise-independent, and extractor-style query maps; compute K, coefficient-circuit cost \(A_{N}\), duplicate parity cancellations, and whether any case survives the \(N^{1+\epsilon}\) affine-sketch falsifier threshold.

\paragraph{When to stop.} Stop or quarantine this padding/projection branch only after the length map N(n), projection overhead, source lower-bound quantifier and target exponent window are calculated; quarantine if the route merely moves the P vs NP difficulty into sparse valid-length lower bounds.

\paragraph{Do not stop prematurely if.} Do not stop before deriving the exponent transfer c \(\to \) c/k or its growth-padding analogue, and before stating whether the result excludes all polynomial SAT exponents or only a bounded exponent interval.

\paragraph{How it adds to the global problem.} This move adds to the P vs NP problem because it identifies whether the transfer loses all-exponent strength through a positive target-query exponent \(\beta\); it separates true hardness magnification from projection or padding arguments that only move proof debt; it keeps barrier status and nonuniform/advice leakage attached to the actual transfer theorem. The current branch-control decision is continue\_run\_next. Payoff criterion: For the proof-complexity branch, separate fixed-family proof hardness from language hardness by naming the proof system, shortcut algorithm, and semantic lower-bound target.

\subsection*{cycle-0054: Hardness-Magnification Bridge Realization: XOR Carrier Approximate Affine-Sketch Gate}

\paragraph{Learning.} Cycle 54 strengthens the XOR route with an approximate piecewise-affine incidence-sketch conservation law. Bad seeds cost EN gates to patch, and a sketch of total cost s+Kr+EN below \(N^{1+\epsilon}\) is a concrete falsifier for the source lower bound. High incidence rank alone is not valid lower-bound evidence because the SAT label functional may be simple on the generated image.

\paragraph{Pressure.} Need a direct lower bound for the SAT-label pullback x \(\to \) , or a non-natural real-SAT theorem showing that every subthreshold \(H_{N}\) circuit induces a subthreshold approximate affine-incidence sketch. Current rank and sketch audits are only falsifiers.

\paragraph{Next.} Instantiate the hash-based XOR generator and compute whether its incidence map has an explicit low-cost piecewise-affine cover; if none is found, formulate the exact extraction theorem needed to turn small \(H_{N}\) circuits into such covers.

\paragraph{When to stop.} Stop or quarantine this P vs NP branch only after it states the source object, target class, complexity model, transfer inequality, quantifier strength, and barrier status; quarantine if two sharpened attempts still cannot say whether the route implies \(\mathrm{SAT}\notin \mathrm{P}/\mathrm{poly}\), \(\mathrm{NP}\not\subset \mathrm{P}/\mathrm{poly}\), or only a weaker/proxy lower bound.

\paragraph{Do not stop prematurely if.} Do not stop before the branch has produced a concrete source-to-target transfer theorem, an explicit obstruction lemma, or a barrier-aware reason why the route cannot reach the P vs NP/P/poly target.

\paragraph{How it adds to the global problem.} This move adds to the P vs NP problem because it keeps barrier status and nonuniform/advice leakage attached to the actual transfer theorem. The current branch-control decision is backtrack\_to\_sharpen. Payoff criterion: For the proof-complexity branch, separate fixed-family proof hardness from language hardness by naming the proof system, shortcut algorithm, and semantic lower-bound target.

\subsection*{cycle-0055: Hardness-Magnification Bridge Realization}

\paragraph{Learning.} Cycle 55 adds the SAT-preserving quotient normal form for XOR carriers. The valid target bridge is still \(max{\rho,q}<1+\epsilon\), but the source lower bound must beat quotient piecewise-affine incidence sketches, not raw range or rank. Symmetry-pair cancellation is a concrete falsifier that can make a high-entropy generator compute a constant carrier; this is route-level progress only, not a proof of \(\mathrm{P} \ne \mathrm{NP}\).

\paragraph{Pressure.} Label-sensitive extraction gap: a small arbitrary circuit for the single function x \(\to \)  need not reveal a small quotient-incidence sketch, and high quotient rank alone does not preclude the SAT-labelled linear functional from being simple.

\paragraph{Next.} Compute QSk\_\{sim,N\}(E) for quotient relations generated by equality, variable renaming, clause permutation, tautological padding, and duplicate-pair involutions on the current high-q XOR carrier template; record the smallest residual exponent u and pairing-certificate cost.

\paragraph{When to stop.} Stop or quarantine this padding/projection branch only after the length map N(n), projection overhead, source lower-bound quantifier and target exponent window are calculated; quarantine if the route merely moves the P vs NP difficulty into sparse valid-length lower bounds.

\paragraph{Do not stop prematurely if.} Do not stop before deriving the exponent transfer c \(\to \) c/k or its growth-padding analogue, and before stating whether the result excludes all polynomial SAT exponents or only a bounded exponent interval.

\paragraph{How it adds to the global problem.} This move adds to the P vs NP problem because it separates true hardness magnification from projection or padding arguments that only move proof debt; it keeps barrier status and nonuniform/advice leakage attached to the actual transfer theorem. The current branch-control decision is adapt\_then\_run. Payoff criterion: For the high-q generated-OR carrier, compute the exact leakage exponent of the remaining witness-cover or unary-residual mechanism and compare it with 1+epsilon-q.

\subsection*{cycle-0056: Hardness-Magnification Bridge Realization: XOR quotient codebook gate}

\paragraph{Learning.} Cycle 56 adds a sharper conservation law for XOR carriers: after SAT-preserving parity quotienting, \(H_{N}\) is easy whenever the incidence map has either low affine rank or a low-cardinality computable transcript, with patchable errors charged by EN. Thus high quotient entropy or high rank alone is not valid source evidence; the source obligation is a combined codebook-or-rank minimax lower bound, while the target bridge remains unchanged.

\paragraph{Pressure.} The derivation does not show that every small arbitrary circuit for \(H_{N}\) yields such a quotient sketch. It is an upper-bound falsifier, not a lower-bound certificate.

\paragraph{Next.} Instantiate a candidate high-q generator and compute the exponent of its SAT-preserving odd-support transcript range versus affine span rank; reject it if min\{theta,rank-exponent\}+locator-cost falls below 1+epsilon.

\paragraph{When to stop.} Stop or quarantine this P vs NP branch only after it states the source object, target class, complexity model, transfer inequality, quantifier strength, and barrier status; quarantine if two sharpened attempts still cannot say whether the route implies \(\mathrm{SAT}\notin \mathrm{P}/\mathrm{poly}\), \(\mathrm{NP}\not\subset \mathrm{P}/\mathrm{poly}\), or only a weaker/proxy lower bound.

\paragraph{Do not stop prematurely if.} Do not stop before the branch has produced a concrete source-to-target transfer theorem, an explicit obstruction lemma, or a barrier-aware reason why the route cannot reach the P vs NP/P/poly target.

\paragraph{How it adds to the global problem.} This move adds to the P vs NP problem because it keeps barrier status and nonuniform/advice leakage attached to the actual transfer theorem. The current branch-control decision is backtrack\_to\_sharpen. Payoff criterion: For the high-q generated-OR carrier, compute the exact leakage exponent of the remaining witness-cover or unary-residual mechanism and compare it with 1+epsilon-q.

\subsection*{cycle-0057: Hardness-Magnification Bridge Realization: XOR quotient carrier, low-degree feature-sketch gate}

\paragraph{Learning.} Cycle 57 adds a nonlinear conservation law for the XOR carrier. The source lower bound must beat low-degree quotient-feature sketches with threshold \(S_{T}+binom(r,\le t)+EN\), so high rank or high codebook entropy is not valid evidence. The target-preserving bridge is unchanged, but the source obligation is narrower and more nonrelativizing.

\paragraph{Pressure.} No method is known to prove that the real SAT-labelled quotient functional avoids all subthreshold low-degree feature sketches. Random-incidence or random-label counting does not transfer to SAT.

\paragraph{Next.} Formalize the degree-2 case: characterize whether Tseitin-style parity, formula-size, variable-incidence, or clause-type features give \(r=N^\alpha\) with \(2alpha<1+\epsilon\) for the generated length-(log N)\textasciicircum{}d SAT queries.

\paragraph{When to stop.} Stop or quarantine this meta-complexity carrier branch only after MCSP/MKtP parameters, truth-table length N, size threshold s, verifier overhead, target SAT-circuit use, and natural-proof or advice-leakage status have been computed; quarantine if the transfer still produces only a finite exponent window or relies on a dense natural property without an escape certificate.

\paragraph{Do not stop prematurely if.} Do not stop before naming the carrier variant, deriving its verifier/reduction overhead, and checking whether a \(\mathrm{SAT}\in \mathrm{SIZE}(n^{c})\) assumption gives carrier circuits below the claimed source lower-bound threshold.

\paragraph{How it adds to the global problem.} This move adds to the P vs NP problem because it tests whether a meta-complexity carrier gives a weaker source obligation than the original P vs NP/P-poly target; it keeps barrier status and nonuniform/advice leakage attached to the actual transfer theorem. The current branch-control decision is adapt\_then\_run. Payoff criterion: For the restricted-model branch, compute the compiler envelope from unrestricted SIZE(m\textasciicircum{}k) or P-time algorithms into the named model; promotion requires a bound above that envelope on the same SAT slice.

\subsection*{cycle-0058: Hardness-magnification bridge via XOR carriers and sparse quotient-polynomial falsifiers}

\paragraph{Learning.} Cycle 58 repairs the XOR-carrier audit: low degree was the wrong nonlinear measure. The valid upper-bound falsifier is sparse quotient-polynomial support with threshold \(S_{T}+A_{N}\) \(r_{N}+E_{N}\) \(N<N^{1+\epsilon}\). This preserves the polylog-query target bridge but narrows the source obligation and exposes a new extraction gap: sparse-sketch hardness is not yet general-circuit hardness.

\paragraph{Pressure.} No valid derivation converts an arbitrary small \(H_{N}\) circuit into a charged sparse quotient-polynomial sketch. Without this extraction, sketch lower bounds are falsifier tests, not source lower bounds.

\paragraph{Next.} Formalize the weakest nontrivial extraction candidate: restrict arbitrary \(H_{N}\) circuits to circuits whose wires are invariant under SAT-preserving quotient symmetries, and test whether such invariant circuits collapse to sparse quotient-polynomial sketches with subthreshold loss.

\paragraph{When to stop.} Stop or quarantine this padding/projection branch only after the length map N(n), projection overhead, source lower-bound quantifier and target exponent window are calculated; quarantine if the route merely moves the P vs NP difficulty into sparse valid-length lower bounds.

\paragraph{Do not stop prematurely if.} Do not stop before deriving the exponent transfer c \(\to \) c/k or its growth-padding analogue, and before stating whether the result excludes all polynomial SAT exponents or only a bounded exponent interval.

\paragraph{How it adds to the global problem.} This move adds to the P vs NP problem because it identifies whether the transfer loses all-exponent strength through a positive target-query exponent \(\beta\); it separates true hardness magnification from projection or padding arguments that only move proof debt; it keeps barrier status and nonuniform/advice leakage attached to the actual transfer theorem. The current branch-control decision is continue\_run\_next. Payoff criterion: For the proof-complexity branch, separate fixed-family proof hardness from language hardness by naming the proof system, shortcut algorithm, and semantic lower-bound target.

\subsection*{cycle-0059: Hardness-Magnification Bridge Realization: XOR carrier sparse-sketch audit}

\paragraph{Learning.} Cycle 59 identifies a conservation law for the XOR-carrier audit: unrestricted sparse quotient-polynomial hardness is not a simplification of the source lower bound, since \(T=H_{N}\) is an admissible one-bit quotient transcript. The valid contribution is a sharper quarantine: sparse sketches remain concrete upper-bound falsifiers, while any proof-relevant sketch lower bound now needs a named label-free transcript class plus a nonrelativizing extraction theorem. The main \(\mathrm{P} \ne \mathrm{NP}\) conjecture is untouched.

\paragraph{Pressure.} No nonrelativizing locality-extraction theorem is known that converts an arbitrary small circuit for the real SAT-XOR carrier into a subthreshold label-free incidence transcript; relativizing planted-control examples block the naive argument.

\paragraph{Next.} Formalize the local-linear incidence transcript class for the XOR SAT carrier and compute the minimum sparse \(F_{2}\) support needed under a random-generator model, while auditing whether any extraction from arbitrary circuits to that class can avoid the planted-control relativization falsifier.

\paragraph{When to stop.} Stop or quarantine this padding/projection branch only after the length map N(n), projection overhead, source lower-bound quantifier and target exponent window are calculated; quarantine if the route merely moves the P vs NP difficulty into sparse valid-length lower bounds.

\paragraph{Do not stop prematurely if.} Do not stop before deriving the exponent transfer c \(\to \) c/k or its growth-padding analogue, and before stating whether the result excludes all polynomial SAT exponents or only a bounded exponent interval.

\paragraph{How it adds to the global problem.} This move adds to the P vs NP problem because it separates true hardness magnification from projection or padding arguments that only move proof debt; it keeps barrier status and nonuniform/advice leakage attached to the actual transfer theorem. The current branch-control decision is adapt\_then\_run. Payoff criterion: For the proof-complexity branch, separate fixed-family proof hardness from language hardness by naming the proof system, shortcut algorithm, and semantic lower-bound target.

\subsection*{cycle-0060: Constructive Label-Free Sketch Gate for XOR Magnification}

\paragraph{Learning.} Cycle 60 separates two notions previously conflated. Sparse quotient-polynomial sketches are valid upper-bound falsifiers only when the transcript is constructively label-free; otherwise \(T_{N}=H_{N}\) collapses the audit. The target bridge \(max{\rho,q}<1+\epsilon\) survives, but sketch-hardness is not a source lower bound without a nonrelativizing extraction theorem or a direct C(H\_N) lower bound.

\paragraph{Pressure.} No theorem converts an arbitrary small circuit for the parity of SAT labels on generated queries into a constructively label-free low-support incidence sketch with subthreshold residual cost.

\paragraph{Next.} Define the first concrete compiler class \(T^{inc-lin}\): \(T_{N}(x)=B_{N}\) a\_N(x) for explicit SAT-preserving quotient matrices \(B_{N}\) of row sparsity polylog N, then compute the minimum \(S_{T}+A_{N}\) \(r_{N}\) for representing \(H_{N}\) by sparse \(\mathrm{P}_{N}\) in the high-q, \(d=2\) regime.

\paragraph{When to stop.} Stop or quarantine this black-box substitution branch only after the source carrier, SAT-query length exponent \(\beta\), interface size exponent q, source-threshold exponent, and target class have been explicitly computed; quarantine if every sharpened interface still has positive \(\beta\) and therefore only excludes a bounded range of SAT circuit exponents rather than NP subset P/poly.

\paragraph{Do not stop prematurely if.} Do not stop before deriving the transfer inequality from \(\mathrm{SAT}\in \mathrm{SIZE}(n^{c})\) to the carrier upper bound, identifying whether \(\beta=0\), \(\beta>0\), or the argument is genuinely nonblackbox, and recording the barrier status of the proposed escape.

\paragraph{How it adds to the global problem.} This move adds to the P vs NP problem because it identifies whether the transfer loses all-exponent strength through a positive target-query exponent \(\beta\); it separates true hardness magnification from projection or padding arguments that only move proof debt; it keeps barrier status and nonuniform/advice leakage attached to the actual transfer theorem. The current branch-control decision is continue\_run\_next. Payoff criterion: For the high-q generated-OR carrier, compute the exact leakage exponent of the remaining witness-cover or unary-residual mechanism and compare it with 1+epsilon-q.

\subsection*{cycle-0061: Hardness-Magnification Bridge Realization}

\paragraph{Learning.} Cycle 61 upgrades the XOR-carrier audit from vague label-free transcripts to a costed global-codebook normal form. The valid new contribution is an explicit upper-bound falsifier \(C(H_{N})\le S_{T}+S_{D}+\)\(\widetilde{O}(U_{N}+A_{Nr,N})\)\(+E_{NN}\). This closes a leakage loophole in sparse quotient sketches, while exposing the exact remaining gap: absence of label-free sketches is not a general-circuit lower bound without a nonrelativizing extraction theorem.

\paragraph{Pressure.} No extraction theorem converts arbitrary \(H_{N}\) circuits into constructively label-free codebook sketches. The trivial transcript \(T_{N}=H_{N}\) is small whenever \(H_{N}\) is small but is not label-free and may encode SAT-label information semantically.

\paragraph{Next.} Instantiate a candidate generator and compute \(U_{N},A_{N},r_{N},E_{N}\) for its SAT-preserving quotient codebook; if the cost is subthreshold, quarantine that generator, otherwise formulate the exact extraction theorem it would require.

\paragraph{When to stop.} Stop or quarantine this black-box substitution branch only after the source carrier, SAT-query length exponent \(\beta\), interface size exponent q, source-threshold exponent, and target class have been explicitly computed; quarantine if every sharpened interface still has positive \(\beta\) and therefore only excludes a bounded range of SAT circuit exponents rather than NP subset P/poly.

\paragraph{Do not stop prematurely if.} Do not stop before deriving the transfer inequality from \(\mathrm{SAT}\in \mathrm{SIZE}(n^{c})\) to the carrier upper bound, identifying whether \(\beta=0\), \(\beta>0\), or the argument is genuinely nonblackbox, and recording the barrier status of the proposed escape.

\paragraph{How it adds to the global problem.} This move adds to the P vs NP problem because it identifies whether the transfer loses all-exponent strength through a positive target-query exponent \(\beta\); it separates true hardness magnification from projection or padding arguments that only move proof debt; it keeps barrier status and nonuniform/advice leakage attached to the actual transfer theorem. The current branch-control decision is continue\_run\_next. Payoff criterion: For the proof-complexity branch, separate fixed-family proof hardness from language hardness by naming the proof system, shortcut algorithm, and semantic lower-bound target.

\subsection*{cycle-0062: Hardness-Magnification Bridge Realization}

\paragraph{Learning.} Cycle 62 adds a new fork: high incidence rank can make the XOR carrier lower-bound property large, so an efficient semantic certificate would be natural, while low rank is an upper-bound falsifier. The polylog \(\beta=0\) bridge remains valid, but the surviving source proof must be non-natural, non-large, proof-mediated, or direct.

\paragraph{Pressure.} The derivation does not produce the actual source lower bound. It only says that high-rank, efficiently checkable, large certificates are barrier-exposed, while absence of label-free sketches is still not equivalent to arbitrary-circuit hardness without a nonrelativizing extraction theorem.

\paragraph{Next.} Instantiate the random-affine generator with L=(log N)\textasciicircum{}2 and q in ((1+epsilon)/2,1+epsilon); compute \(r_{N}\) for the induced sparse incidence matrix and test whether the rank gap forces naturalness or low-rank falsification.

\paragraph{When to stop.} Stop or quarantine this black-box substitution branch only after the source carrier, SAT-query length exponent \(\beta\), interface size exponent q, source-threshold exponent, and target class have been explicitly computed; quarantine if every sharpened interface still has positive \(\beta\) and therefore only excludes a bounded range of SAT circuit exponents rather than NP subset P/poly.

\paragraph{Do not stop prematurely if.} Do not stop before deriving the transfer inequality from \(\mathrm{SAT}\in \mathrm{SIZE}(n^{c})\) to the carrier upper bound, identifying whether \(\beta=0\), \(\beta>0\), or the argument is genuinely nonblackbox, and recording the barrier status of the proposed escape.

\paragraph{How it adds to the global problem.} This move adds to the P vs NP problem because it identifies whether the transfer loses all-exponent strength through a positive target-query exponent \(\beta\); it separates true hardness magnification from projection or padding arguments that only move proof debt; it keeps barrier status and nonuniform/advice leakage attached to the actual transfer theorem. The current branch-control decision is continue\_run\_next. Payoff criterion: For the proof-complexity branch, separate fixed-family proof hardness from language hardness by naming the proof system, shortcut algorithm, and semantic lower-bound target.

\subsection*{cycle-0063: Hardness-Magnification Bridge Realization}

\paragraph{Learning.} Cycle 63 upgrades the rank-density fork into a quantitative gate: low constructive rank is a direct upper-bound falsifier, while high rank makes any efficient useful certificate large in label space. The \(\beta=0\) magnification bridge is still valid, but rank-only evidence is quarantined unless the source proof is direct or explicitly non-natural and nonrelativizing.

\paragraph{Pressure.} The branch still lacks a direct non-natural lower bound for \(H_{N}^{\mathrm{SAT}}\) or a nonrelativizing theorem extracting a label-free sketch from every small \(H_{N}\) circuit.

\paragraph{Next.} Perform the finite prior-art boundary check for the label-space naturalness claim, then instantiate the thought experiment with a concrete pairwise-independent generator and compute \(r_{N}\), quotient rank, and any row-coordinate decoder cost \(D_{N}\).

\paragraph{When to stop.} Stop or quarantine this black-box substitution branch only after the source carrier, SAT-query length exponent \(\beta\), interface size exponent q, source-threshold exponent, and target class have been explicitly computed; quarantine if every sharpened interface still has positive \(\beta\) and therefore only excludes a bounded range of SAT circuit exponents rather than NP subset P/poly.

\paragraph{Do not stop prematurely if.} Do not stop before deriving the transfer inequality from \(\mathrm{SAT}\in \mathrm{SIZE}(n^{c})\) to the carrier upper bound, identifying whether \(\beta=0\), \(\beta>0\), or the argument is genuinely nonblackbox, and recording the barrier status of the proposed escape.

\paragraph{How it adds to the global problem.} This move adds to the P vs NP problem because it identifies whether the transfer loses all-exponent strength through a positive target-query exponent \(\beta\); it separates true hardness magnification from projection or padding arguments that only move proof debt; it keeps barrier status and nonuniform/advice leakage attached to the actual transfer theorem; it searches the remaining plausible escape route: nonblackbox or subpolynomial-interface carriers. The current branch-control decision is continue\_run\_next. Payoff criterion: For this route, state the exact barrier premise it uses or escapes: relativizing, natural, algebrizing, nonuniform, or advice-leaking; promotion requires the escape to be attached to the specific transfer theorem.

\subsection*{cycle-0064: Hardness-Magnification Bridge Realization}

\paragraph{Learning.} Cycle 64 upgrades the XOR route from rank-density intuition to a fork theorem: low constructive incidence rank is an explicit upper-bound falsifier, while any large truth-table constructive anti-alignment proof for the SAT label pulls back through the valid \(\beta=0\) bridge to a natural property. The remaining source route must be direct, nonlarge, nonconstructive, proof-mediated, or explicitly nonrelativizing; this is route-level progress, not \(\mathrm{P} \ne \mathrm{NP}\).

\paragraph{Pressure.} The route still lacks a nonlarge, nonconstructive, proof-mediated, or direct argument that the actual SAT label avoids the easy induced-XOR label set.

\paragraph{Next.} Instantiate L\_N=(log N)\textasciicircum{}2 and q in ((1+epsilon)/2,1+epsilon); compute achievable \(M_{N}\) and \(r_{N}\) for an explicit generator family, then test for a label-free coordinate transcript or prove any proposed SAT anti-alignment certificate has density below \(2^{-\omega(L_{N})}\).

\paragraph{When to stop.} Stop or quarantine this black-box substitution branch only after the source carrier, SAT-query length exponent \(\beta\), interface size exponent q, source-threshold exponent, and target class have been explicitly computed; quarantine if every sharpened interface still has positive \(\beta\) and therefore only excludes a bounded range of SAT circuit exponents rather than NP subset P/poly.

\paragraph{Do not stop prematurely if.} Do not stop before deriving the transfer inequality from \(\mathrm{SAT}\in \mathrm{SIZE}(n^{c})\) to the carrier upper bound, identifying whether \(\beta=0\), \(\beta>0\), or the argument is genuinely nonblackbox, and recording the barrier status of the proposed escape.

\paragraph{How it adds to the global problem.} This move adds to the P vs NP problem because it identifies whether the transfer loses all-exponent strength through a positive target-query exponent \(\beta\); it keeps barrier status and nonuniform/advice leakage attached to the actual transfer theorem. The current branch-control decision is continue\_run\_next. Payoff criterion: For the proof-complexity branch, separate fixed-family proof hardness from language hardness by naming the proof system, shortcut algorithm, and semantic lower-bound target.

\subsection*{cycle-0065: Hardness-Magnification Bridge Realization}

\paragraph{Learning.} Cycle 65 closes a finite validation gap: short searchable proof certificates do not automatically escape natural proofs in the \(\beta=0\) XOR carrier. If they are large and sound, they form a valid natural property; if rank is low, there is a separate upper-bound falsifier. The surviving route is narrower: direct \(H_{N}\) lower bounds or nonlarge, nonconstructive, explicitly nonrelativizing proof objects.

\paragraph{Pressure.} The argument is conditional as a barrier diagnosis: it uses the standard PRF-based natural-proofs obstruction and does not rule out nonlarge SAT-specific proofs, \(\omega(L_{N})-length\) proofs, nonconstructive verifiers, or nonrelativizing proof-system semantics.

\paragraph{Next.} Instantiate \(R_{N}\) with a named proof system, e.g. Extended Frege or IPS, and test whether an LB certificate for \(C(H_{N,y})>N^{1+\epsilon}\) has \(O(L_{N})-bit\) searchable descriptions, inverse-polynomial label density, and truth-table soundness; otherwise record the precise escape parameter.

\paragraph{When to stop.} Stop or quarantine this black-box substitution branch only after the source carrier, SAT-query length exponent \(\beta\), interface size exponent q, source-threshold exponent, and target class have been explicitly computed; quarantine if every sharpened interface still has positive \(\beta\) and therefore only excludes a bounded range of SAT circuit exponents rather than NP subset P/poly.

\paragraph{Do not stop prematurely if.} Do not stop before deriving the transfer inequality from \(\mathrm{SAT}\in \mathrm{SIZE}(n^{c})\) to the carrier upper bound, identifying whether \(\beta=0\), \(\beta>0\), or the argument is genuinely nonblackbox, and recording the barrier status of the proposed escape.

\paragraph{How it adds to the global problem.} This move adds to the P vs NP problem because it identifies whether the transfer loses all-exponent strength through a positive target-query exponent \(\beta\); it keeps barrier status and nonuniform/advice leakage attached to the actual transfer theorem. The current branch-control decision is continue\_run\_next. Payoff criterion: For the proof-complexity branch, separate fixed-family proof hardness from language hardness by naming the proof system, shortcut algorithm, and semantic lower-bound target.

\subsection*{cycle-0066: Hardness-Magnification Bridge Realization}

\paragraph{Learning.} Cycle 66 sharpens the proof-mediated XOR route: searchable certificates do not automatically trigger the standard natural-proofs barrier. The valid obstruction needs parametric soundness, constructivity, large pulled-back density, and a density-preserving label-realization map. Without these, the barrier claim is invalid, but the source proof must be direct, nonlarge, nonconstructive, or explicitly nonrelativizing.

\paragraph{Pressure.} No explicit density-preserving label-realization map is known for the generated SAT universe \(U_{N}\), and soundness of plausible proof predicates is usually shown only for the actual SAT label rather than for every artificial label y.

\paragraph{Next.} Instantiate \(R_{n}\) for one high-q polylog XOR carrier or prove a finite impossibility: no label-free density-preserving realization can coexist with \(q,\rho<1+\epsilon\) and generator universe \(U_{N}\) of SAT formulas length (log N)\textasciicircum{}2.

\paragraph{When to stop.} Stop or quarantine this black-box substitution branch only after the source carrier, SAT-query length exponent \(\beta\), interface size exponent q, source-threshold exponent, and target class have been explicitly computed; quarantine if every sharpened interface still has positive \(\beta\) and therefore only excludes a bounded range of SAT circuit exponents rather than NP subset P/poly.

\paragraph{Do not stop prematurely if.} Do not stop before deriving the transfer inequality from \(\mathrm{SAT}\in \mathrm{SIZE}(n^{c})\) to the carrier upper bound, identifying whether \(\beta=0\), \(\beta>0\), or the argument is genuinely nonblackbox, and recording the barrier status of the proposed escape.

\paragraph{How it adds to the global problem.} This move adds to the P vs NP problem because it identifies whether the transfer loses all-exponent strength through a positive target-query exponent \(\beta\); it keeps barrier status and nonuniform/advice leakage attached to the actual transfer theorem. The current branch-control decision is continue\_run\_next. Payoff criterion: For the high-q generated-OR carrier, compute the exact leakage exponent of the remaining witness-cover or unary-residual mechanism and compare it with 1+epsilon-q.

\subsection*{cycle-0067: Hardness-Magnification Bridge Realization}

\paragraph{Learning.} Cycle 67 upgrades the previous label-realization warning into an exact restriction-extension theorem. For \(\beta=0\) XOR carriers, dense constructive label-mediated source lower bounds naturalize automatically when \(U_{N}\) is uniformly enumerable and soundness is in the same circuit model. The valid bridge survives, but any proof-mediated route must now be direct, nonlarge, nonconstructive, or explicitly nonrelativizing; ambient source-property largeness is not enough because density must be image-relative.

\paragraph{Pressure.} No lower bound for C(H\_N) is proved. The derivation only certifies that one dense constructive label-mediated proof shape would collide with natural-proofs methodology, while image-relative and nonlarge certificates remain open.

\paragraph{Next.} Instantiate \(U_{N}\) for one explicit polylog-query XOR generator, compute whether \(U_{N}\) is enumerable in poly(2\textasciicircum{}L), and classify a candidate certificate as dense-label, image-relative, or nonlarge.

\paragraph{When to stop.} Stop or quarantine this black-box substitution branch only after the source carrier, SAT-query length exponent \(\beta\), interface size exponent q, source-threshold exponent, and target class have been explicitly computed; quarantine if every sharpened interface still has positive \(\beta\) and therefore only excludes a bounded range of SAT circuit exponents rather than NP subset P/poly.

\paragraph{Do not stop prematurely if.} Do not stop before deriving the transfer inequality from \(\mathrm{SAT}\in \mathrm{SIZE}(n^{c})\) to the carrier upper bound, identifying whether \(\beta=0\), \(\beta>0\), or the argument is genuinely nonblackbox, and recording the barrier status of the proposed escape.

\paragraph{How it adds to the global problem.} This move adds to the P vs NP problem because it identifies whether the transfer loses all-exponent strength through a positive target-query exponent \(\beta\); it separates true hardness magnification from projection or padding arguments that only move proof debt; it keeps barrier status and nonuniform/advice leakage attached to the actual transfer theorem. The current branch-control decision is continue\_run\_next. Payoff criterion: For the proof-complexity branch, separate fixed-family proof hardness from language hardness by naming the proof system, shortcut algorithm, and semantic lower-bound target.

\subsection*{cycle-0068: Hardness-Magnification Bridge Realization}

\paragraph{Learning.} Cycle 68 separates two notions previously conflated: label-mediated certificates for the \(\beta=0\) XOR carrier are large inside the realization image, but the image has ambient density at most 2\textasciicircum{}\{-2\textasciicircum{}N+2\textasciicircum{}\{o(N)\}\}. Thus the valid bridge to \(\mathrm{SAT}\notin \mathrm{P}/\mathrm{poly}\) survives, while the standard natural-proofs barrier requires a separate constructive ambient extension; this is a sharper route-status theorem, not progress on the main conjecture itself.

\paragraph{Pressure.} No direct lower bound for the real SAT-label carrier \(H_{N}\) is obtained, and no polynomial-time density-preserving extension from Im(R\_N) to ambient N-bit truth tables is known.

\paragraph{Next.} Formalize the extension problem: given \(R_{N}\) and a sound searchable label certificate, prove either that no poly(N)-time ambient extension of density \(2^{-\mathrm{poly}(N)}\) exists for this certificate class, or that such an extension exists and therefore naturalizes the route.

\paragraph{When to stop.} Stop or quarantine this black-box substitution branch only after the source carrier, SAT-query length exponent \(\beta\), interface size exponent q, source-threshold exponent, and target class have been explicitly computed; quarantine if every sharpened interface still has positive \(\beta\) and therefore only excludes a bounded range of SAT circuit exponents rather than NP subset P/poly.

\paragraph{Do not stop prematurely if.} Do not stop before deriving the transfer inequality from \(\mathrm{SAT}\in \mathrm{SIZE}(n^{c})\) to the carrier upper bound, identifying whether \(\beta=0\), \(\beta>0\), or the argument is genuinely nonblackbox, and recording the barrier status of the proposed escape.

\paragraph{How it adds to the global problem.} This move adds to the P vs NP problem because it identifies whether the transfer loses all-exponent strength through a positive target-query exponent \(\beta\); it tests whether a meta-complexity carrier gives a weaker source obligation than the original P vs NP/P-poly target; it keeps barrier status and nonuniform/advice leakage attached to the actual transfer theorem. The current branch-control decision is continue\_run\_next. Payoff criterion: For the proof-complexity branch, separate fixed-family proof hardness from language hardness by naming the proof system, shortcut algorithm, and semantic lower-bound target.

\subsection*{cycle-0069: Hardness-Magnification Bridge Realization}

\paragraph{Learning.} Cycle 69 turns the image-density ambiguity into a precise threshold theorem. A label-mediated \(\beta=0\) XOR certificate naturalizes through restriction to \(U_{N}\) exactly when its label density is at least \(2^{-\mathrm{poly}(L)}\) and its verifier is \(2^{O(L)}-constructive\). The earlier ambient-tiny source image is not by itself a valid escape; the real escapes are nonlarge density, nonconstructivity, advice leakage, or explicitly nonrelativizing direct lower bounds.

\paragraph{Pressure.} No construction is given of a \(\mathrm{P}_{N}\) that accepts \(\mathrm{SAT}|U_{N}\) while remaining sound; the derivation only classifies when such a proof-mediated route would naturalize.

\paragraph{Next.} Instantiate \(U_{N}\) with L=(log N)\textasciicircum{}2 and q in ((1+epsilon)/2,1+epsilon); compute whether candidate proof systems give certificate density above \(2^{-\mathrm{poly}(L)}\) or force a nonlarge escape, then run the prior-art boundary check.

\paragraph{When to stop.} Stop or quarantine this black-box substitution branch only after the source carrier, SAT-query length exponent \(\beta\), interface size exponent q, source-threshold exponent, and target class have been explicitly computed; quarantine if every sharpened interface still has positive \(\beta\) and therefore only excludes a bounded range of SAT circuit exponents rather than NP subset P/poly.

\paragraph{Do not stop prematurely if.} Do not stop before deriving the transfer inequality from \(\mathrm{SAT}\in \mathrm{SIZE}(n^{c})\) to the carrier upper bound, identifying whether \(\beta=0\), \(\beta>0\), or the argument is genuinely nonblackbox, and recording the barrier status of the proposed escape.

\paragraph{How it adds to the global problem.} This move adds to the P vs NP problem because it identifies whether the transfer loses all-exponent strength through a positive target-query exponent \(\beta\); it keeps barrier status and nonuniform/advice leakage attached to the actual transfer theorem. The current branch-control decision is continue\_run\_next. Payoff criterion: For the proof-complexity branch, separate fixed-family proof hardness from language hardness by naming the proof system, shortcut algorithm, and semantic lower-bound target.

\subsection*{cycle-0070: Hardness-Magnification Bridge Realization}

\paragraph{Learning.} Cycle 70 corrects the naturalness gate for the \(\beta=0\) XOR route. A dense label certificate cylinder-extends to a valid classical natural property exactly at density \(2^{-O(L)}\) with \(2^{O(L)}\) constructivity; source-image ambient sparsity is irrelevant after restriction. The real escapes are nonlarge density below the classical cutoff, nonconstructivity, advice leakage, or a direct nonrelativizing \(H_{N}\) lower bound.

\paragraph{Pressure.} This does not prove that such \(Q_{N}\) exists, and it does not turn sparse or nonconstructive certificates into barriers. It only gives a valid no-go for dense, constructive, same-model label certificates.

\paragraph{Next.} Instantiate the density-tier audit on a toy proof-certificate family: certificates describing affine anti-alignment witnesses over \(U_{N}\). Compute \(delta_{N}\), verifier time, and whether soundness reaches arbitrary \(C_{N}(H_{y})>N^{1+\epsilon}\) or only sketch hardness.

\paragraph{When to stop.} Stop or quarantine this black-box substitution branch only after the source carrier, SAT-query length exponent \(\beta\), interface size exponent q, source-threshold exponent, and target class have been explicitly computed; quarantine if every sharpened interface still has positive \(\beta\) and therefore only excludes a bounded range of SAT circuit exponents rather than NP subset P/poly.

\paragraph{Do not stop prematurely if.} Do not stop before deriving the transfer inequality from \(\mathrm{SAT}\in \mathrm{SIZE}(n^{c})\) to the carrier upper bound, identifying whether \(\beta=0\), \(\beta>0\), or the argument is genuinely nonblackbox, and recording the barrier status of the proposed escape.

\paragraph{How it adds to the global problem.} This move adds to the P vs NP problem because it identifies whether the transfer loses all-exponent strength through a positive target-query exponent \(\beta\); it separates true hardness magnification from projection or padding arguments that only move proof debt; it keeps barrier status and nonuniform/advice leakage attached to the actual transfer theorem; it searches the remaining plausible escape route: nonblackbox or subpolynomial-interface carriers. The current branch-control decision is continue\_run\_next. Payoff criterion: For the proof-complexity branch, separate fixed-family proof hardness from language hardness by naming the proof system, shortcut algorithm, and semantic lower-bound target.

\subsection*{cycle-0071: Hardness-Magnification Bridge Realization}

\paragraph{Learning.} Cycle 71 upgrades the \(\beta=0\) XOR naturalness audit: density alone is not the decisive condition. Dense searchable certificates naturalize only when their soundness is functional over arbitrary label vectors. Therefore the surviving proof-mediated route must be sparse, nonconstructive, SAT-semantic/nonrelativizing, advice-free in a stronger sense, or a direct \(H_{N}\) lower bound.

\paragraph{Pressure.} This derivation needs functional soundness for all label vectors and a real accepted-label density estimate. SAT-only certificates and nonlarge accepted sets are outside the no-go.

\paragraph{Next.} Instantiate \(V_{N}\) with one concrete proof-system or MCSP-style certificate, compute \(kappa_{N}\) for accepted labels, and decide whether soundness is functional over arbitrary y or essentially SAT-semantic.

\paragraph{When to stop.} Stop or quarantine this black-box substitution branch only after the source carrier, SAT-query length exponent \(\beta\), interface size exponent q, source-threshold exponent, and target class have been explicitly computed; quarantine if every sharpened interface still has positive \(\beta\) and therefore only excludes a bounded range of SAT circuit exponents rather than NP subset P/poly.

\paragraph{Do not stop prematurely if.} Do not stop before deriving the transfer inequality from \(\mathrm{SAT}\in \mathrm{SIZE}(n^{c})\) to the carrier upper bound, identifying whether \(\beta=0\), \(\beta>0\), or the argument is genuinely nonblackbox, and recording the barrier status of the proposed escape.

\paragraph{How it adds to the global problem.} This move adds to the P vs NP problem because it identifies whether the transfer loses all-exponent strength through a positive target-query exponent \(\beta\); it tests whether a meta-complexity carrier gives a weaker source obligation than the original P vs NP/P-poly target; it separates true hardness magnification from projection or padding arguments that only move proof debt; it keeps barrier status and nonuniform/advice leakage attached to the actual transfer theorem. The current branch-control decision is continue\_run\_next. Payoff criterion: For the MCSP/MKtP row branch, decide whether the certificate proves direct row hardness or only aggregate/rank-index hardness; record q\_extract=epsilon/(1+D) and q\_force=(D+epsilon)/(1+D) in the same parameter regime.

\subsection*{cycle-0072: Hardness-Magnification Bridge Realization}

\paragraph{Learning.} Cycle 72 upgrades the \(\beta=0\) XOR route from a vague naturalness warning to a dichotomy: dense functional searchable certificates become valid natural properties, while nonfunctional proof objects must pass an explicit row-selection leakage audit. The target bridge survives, but the source proof space is narrowed to sparse, nonconstructive, SAT-semantic nonrelativizing, or direct \(H_{N}\) lower bounds.

\paragraph{Pressure.} No theorem yet converts SAT-semantic, nonrelativizing certificates into functional certificates or into charged row-selection transcripts. That is the exact remaining gap, not a proof of \(\mathrm{P} \ne \mathrm{NP}\).

\paragraph{Next.} Instantiate the selected-probe audit for one named proof-system candidate: require \(q=1+\epsilon/2\), \(\rho=1+\epsilon/3\), L=(log N)\textasciicircum{}2, then compute whether its certificates force \(A_{N}+S_{N}+R_{N}+E_{NN}\) or \(A_{N}+S_{N}+K_{NL}^{c}+E_{NN}\) below \(N^{1+\epsilon}\).

\paragraph{When to stop.} Stop or quarantine this black-box substitution branch only after the source carrier, SAT-query length exponent \(\beta\), interface size exponent q, source-threshold exponent, and target class have been explicitly computed; quarantine if every sharpened interface still has positive \(\beta\) and therefore only excludes a bounded range of SAT circuit exponents rather than NP subset P/poly.

\paragraph{Do not stop prematurely if.} Do not stop before deriving the transfer inequality from \(\mathrm{SAT}\in \mathrm{SIZE}(n^{c})\) to the carrier upper bound, identifying whether \(\beta=0\), \(\beta>0\), or the argument is genuinely nonblackbox, and recording the barrier status of the proposed escape.

\paragraph{How it adds to the global problem.} This move adds to the P vs NP problem because it identifies whether the transfer loses all-exponent strength through a positive target-query exponent \(\beta\); it separates true hardness magnification from projection or padding arguments that only move proof debt; it keeps barrier status and nonuniform/advice leakage attached to the actual transfer theorem. The current branch-control decision is continue\_run\_next. Payoff criterion: For the proof-complexity branch, separate fixed-family proof hardness from language hardness by naming the proof system, shortcut algorithm, and semantic lower-bound target.

\subsection*{cycle-0073: Hardness-Magnification Bridge Realization}

\paragraph{Learning.} Cycle 73 adds a concrete row-deletion residual threshold to the surviving \(\beta=0\) XOR magnification route. A claimed \(N^{1+\epsilon}\) source lower bound is invalid if the carrier has a subthreshold circuit wrong on fewer than \(N^{\epsilon-o(1)}\) source rows, because equality patching costs O(E\_N N). Thus a proof-mediated nonfunctional escape from natural proofs must be robust against sparse-row leakage; this is route-level progress, not a P vs NP proof.

\paragraph{Pressure.} No theorem yet converts an arbitrary SAT-semantic nonfunctional certificate into either a functional natural property or a sparse-row residual circuit. The row gate is a necessary falsifier audit, not a complete proof-system lower-bound framework.

\paragraph{Next.} Instantiate a named proof system Psys for \(\beta=0\) XOR certificates and test whether any accepted SAT-semantic certificate can be transformed into either a functional label property or a residual circuit wrong on fewer than \(N^{\epsilon-o(1)}\) rows.

\paragraph{When to stop.} Stop or quarantine this black-box substitution branch only after the source carrier, SAT-query length exponent \(\beta\), interface size exponent q, source-threshold exponent, and target class have been explicitly computed; quarantine if every sharpened interface still has positive \(\beta\) and therefore only excludes a bounded range of SAT circuit exponents rather than NP subset P/poly.

\paragraph{Do not stop prematurely if.} Do not stop before deriving the transfer inequality from \(\mathrm{SAT}\in \mathrm{SIZE}(n^{c})\) to the carrier upper bound, identifying whether \(\beta=0\), \(\beta>0\), or the argument is genuinely nonblackbox, and recording the barrier status of the proposed escape.

\paragraph{How it adds to the global problem.} This move adds to the P vs NP problem because it identifies whether the transfer loses all-exponent strength through a positive target-query exponent \(\beta\); it separates true hardness magnification from projection or padding arguments that only move proof debt; it keeps barrier status and nonuniform/advice leakage attached to the actual transfer theorem. The current branch-control decision is continue\_run\_next. Payoff criterion: For the proof-complexity branch, separate fixed-family proof hardness from language hardness by naming the proof system, shortcut algorithm, and semantic lower-bound target.

\subsection*{cycle-0074: Hardness-Magnification Bridge Realization: beta-zero XOR row-residual audit}

\paragraph{Learning.} Cycle 74 upgrades the focus packet: the beta-zero XOR bridge and row-residual patching bound are theorem-shaped and valid under explicit batched-generator assumptions. The live obstruction is now narrower: prove an \(\mathrm{EF}+\mathrm{SAT}_{U}\) semantic extraction/no-extraction theorem, or quarantine proof-mediated certificates that cannot certify residual robustness.

\paragraph{Pressure.} No proof-analysis theorem currently forces a semantic \(\mathrm{EF}+\mathrm{SAT}_{U}\) certificate either to functionalize over all label tables or to localize its nonfunctionality to fewer than \(N^{\epsilon-o(1)}\) source rows. EF may use global SAT structure shared by many rows through the same formula labels.

\paragraph{Next.} Formalize \(\mathrm{EF}+\mathrm{SAT}_{U}\) syntax with label table \(U_{N}\), prove the functional dense naturalization lemma cleanly, then attempt a finite extraction theorem: every nonfunctional accepted proof either identifies a row set \(R_{N}\) and approximator \(D_{N}\) or uses a nonrelativizing SAT principle outside the row-residual audit.

\paragraph{When to stop.} Stop or quarantine this black-box substitution branch only after the source carrier, SAT-query length exponent \(\beta\), interface size exponent q, source-threshold exponent, and target class have been explicitly computed; quarantine if every sharpened interface still has positive \(\beta\) and therefore only excludes a bounded range of SAT circuit exponents rather than NP subset P/poly.

\paragraph{Do not stop prematurely if.} Do not stop before deriving the transfer inequality from \(\mathrm{SAT}\in \mathrm{SIZE}(n^{c})\) to the carrier upper bound, identifying whether \(\beta=0\), \(\beta>0\), or the argument is genuinely nonblackbox, and recording the barrier status of the proposed escape.

\paragraph{How it adds to the global problem.} This move adds to the P vs NP problem because it identifies whether the transfer loses all-exponent strength through a positive target-query exponent \(\beta\); it keeps barrier status and nonuniform/advice leakage attached to the actual transfer theorem. The current branch-control decision is continue\_run\_next. Payoff criterion: For the proof-complexity branch, separate fixed-family proof hardness from language hardness by naming the proof system, shortcut algorithm, and semantic lower-bound target.

\subsection*{cycle-0075: Hardness-Magnification Bridge Realization}

\paragraph{Learning.} Cycle 75 turns the Psys gap into a precise \(\mathrm{EF}+\mathrm{SAT}_{U}\) audit. Ordinary semantic soundness is insufficient: a valid proof-mediated beta-zero XOR route must either be nonlarge/nonconstructive/direct or certify row-residual robustness against subthreshold patching.

\paragraph{Pressure.} No theorem yet extracts row-residual robustness from short \(\mathrm{EF}+\mathrm{SAT}_{U}\) semantic proofs, and no construction yet proves such extraction impossible.

\paragraph{Next.} Instantiate \(\mathrm{EF}+\mathrm{SAT}_{U}\) proofs as circuits over \(\mathrm{SAT}_{U}\) label atoms, define functionalization formally, and try to construct a \(\theta<\epsilon\) residual-leaky transcript or prove extraction to the robustness gate.

\paragraph{When to stop.} Stop or quarantine this black-box substitution branch only after the source carrier, SAT-query length exponent \(\beta\), interface size exponent q, source-threshold exponent, and target class have been explicitly computed; quarantine if every sharpened interface still has positive \(\beta\) and therefore only excludes a bounded range of SAT circuit exponents rather than NP subset P/poly.

\paragraph{Do not stop prematurely if.} Do not stop before deriving the transfer inequality from \(\mathrm{SAT}\in \mathrm{SIZE}(n^{c})\) to the carrier upper bound, identifying whether \(\beta=0\), \(\beta>0\), or the argument is genuinely nonblackbox, and recording the barrier status of the proposed escape.

\paragraph{How it adds to the global problem.} This move adds to the P vs NP problem because it identifies whether the transfer loses all-exponent strength through a positive target-query exponent \(\beta\); it keeps barrier status and nonuniform/advice leakage attached to the actual transfer theorem. The current branch-control decision is continue\_run\_next. Payoff criterion: For the proof-complexity branch, separate fixed-family proof hardness from language hardness by naming the proof system, shortcut algorithm, and semantic lower-bound target.

\subsection*{cycle-0076: Hardness-Magnification Bridge Realization}

\paragraph{Learning.} Cycle 76 closes the Psys instantiation to \(\mathrm{EF}+\mathrm{SAT}_{U}\) and isolates the exact missing certificate strength: ordinary SAT-semantic soundness is not valid evidence unless it proves row-residual robustness above the patching threshold.

\paragraph{Pressure.} No extraction/no-extraction theorem yet maps short \(\mathrm{EF}+\mathrm{SAT}_{U}\) proofs to functional label properties or to row-residual robustness against all subthreshold \(D_{N}\).

\paragraph{Next.} Formalize Robust\_N(s,e) for \(\mathrm{EF}+\mathrm{SAT}_{U}\) and test whether a short proof can remain sound after replacing \(H_{N}\) by \(D_{N}\) off a sparse row set.

\paragraph{When to stop.} Stop or quarantine this black-box substitution branch only after the source carrier, SAT-query length exponent \(\beta\), interface size exponent q, source-threshold exponent, and target class have been explicitly computed; quarantine if every sharpened interface still has positive \(\beta\) and therefore only excludes a bounded range of SAT circuit exponents rather than NP subset P/poly.

\paragraph{Do not stop prematurely if.} Do not stop before deriving the transfer inequality from \(\mathrm{SAT}\in \mathrm{SIZE}(n^{c})\) to the carrier upper bound, identifying whether \(\beta=0\), \(\beta>0\), or the argument is genuinely nonblackbox, and recording the barrier status of the proposed escape.

\paragraph{How it adds to the global problem.} This move adds to the P vs NP problem because it identifies whether the transfer loses all-exponent strength through a positive target-query exponent \(\beta\); it separates true hardness magnification from projection or padding arguments that only move proof debt; it keeps barrier status and nonuniform/advice leakage attached to the actual transfer theorem. The current branch-control decision is continue\_run\_next. Payoff criterion: For this route, state the exact barrier premise it uses or escapes: relativizing, natural, algebrizing, nonuniform, or advice-leaking; promotion requires the escape to be attached to the specific transfer theorem.

\subsection*{cycle-0077: Hardness-Magnification Bridge Realization}

\paragraph{Learning.} Cycle 77 sharpens the \(\mathrm{EF}+\mathrm{SAT}_{U}\) packet: direct lower-bound certificates automatically pass the explicit row-residual audit, so the real gap is not patching but semantic extraction. A new occurrence-oracle row-touch no-go gives a concrete falsifier for relativizing local certificate methods and leaves only SAT-semantic nonrelativizing, nonlarge/nonconstructive, or direct routes alive.

\paragraph{Pressure.} No theorem yet shows that actual \(\mathrm{EF}+\mathrm{SAT}_{U}\) certificates for the intended \(\mathrm{SAT}_{U}\) carrier either reduce to the row-touch occurrence-local model or yield a nonrelativizing extraction from certificate acceptance to LB\_N(T\_N).

\paragraph{Next.} Instantiate Touch(Pi) for a concrete \(\mathrm{EF}+\mathrm{SAT}_{U}\) encoding of circuit lower-bound certificates: count whether proofs of size \(N^{k}\) can mention \(sub-N^\epsilon\) rows, and test if formula-equality/SAT-realization axioms break the occurrence-oracle extension.

\paragraph{When to stop.} Stop or quarantine this black-box substitution branch only after the source carrier, SAT-query length exponent \(\beta\), interface size exponent q, source-threshold exponent, and target class have been explicitly computed; quarantine if every sharpened interface still has positive \(\beta\) and therefore only excludes a bounded range of SAT circuit exponents rather than NP subset P/poly.

\paragraph{Do not stop prematurely if.} Do not stop before deriving the transfer inequality from \(\mathrm{SAT}\in \mathrm{SIZE}(n^{c})\) to the carrier upper bound, identifying whether \(\beta=0\), \(\beta>0\), or the argument is genuinely nonblackbox, and recording the barrier status of the proposed escape.

\paragraph{How it adds to the global problem.} This move adds to the P vs NP problem because it identifies whether the transfer loses all-exponent strength through a positive target-query exponent \(\beta\); it keeps barrier status and nonuniform/advice leakage attached to the actual transfer theorem. The current branch-control decision is continue\_run\_next. Payoff criterion: For this route, state the exact barrier premise it uses or escapes: relativizing, natural, algebrizing, nonuniform, or advice-leaking; promotion requires the escape to be attached to the specific transfer theorem.

\subsection*{cycle-0078: Hardness-Magnification Bridge Realization}

\paragraph{Learning.} Cycle 78 converts the row-residual audit into a theorem-shaped no-go: functional row-local certificates for the beta-zero XOR carrier require \(N^{\epsilon-o(1)}\) touched rows at source threshold \(N^{1+\epsilon-o(1)}\). This sharpens the \(\mathrm{EF}+\mathrm{SAT}_{U}\) obstruction from ordinary soundness to a concrete extraction fork: either the proof is direct or nonlarge/nonconstructive, or it must use real SAT semantics in a nonrelativizing way that defeats the private-row falsifier.

\paragraph{Pressure.} The derivation is complete only for functional row-value certificates. \(\mathrm{EF}+\mathrm{SAT}_{U}\) proofs are SAT-semantic over formula labels, so a nonrelativizing extraction theorem or a private-row no-go formalization is still needed.

\paragraph{Next.} Instantiate \(\mathrm{EF}+\mathrm{SAT}_{U}\) row occurrence access formally: define the touched-row set of a proof verifier, test whether standard EF reflection over \(\mathrm{SAT}_{U}\) extracts to row-value functionality, and if not construct the private-row oracle countermodel as a finite theorem.

\paragraph{When to stop.} Stop or quarantine this black-box substitution branch only after the source carrier, SAT-query length exponent \(\beta\), interface size exponent q, source-threshold exponent, and target class have been explicitly computed; quarantine if every sharpened interface still has positive \(\beta\) and therefore only excludes a bounded range of SAT circuit exponents rather than NP subset P/poly.

\paragraph{Do not stop prematurely if.} Do not stop before deriving the transfer inequality from \(\mathrm{SAT}\in \mathrm{SIZE}(n^{c})\) to the carrier upper bound, identifying whether \(\beta=0\), \(\beta>0\), or the argument is genuinely nonblackbox, and recording the barrier status of the proposed escape.

\paragraph{How it adds to the global problem.} This move adds to the P vs NP problem because it identifies whether the transfer loses all-exponent strength through a positive target-query exponent \(\beta\); it keeps barrier status and nonuniform/advice leakage attached to the actual transfer theorem. The current branch-control decision is continue\_run\_next. Payoff criterion: For the proof-complexity branch, separate fixed-family proof hardness from language hardness by naming the proof system, shortcut algorithm, and semantic lower-bound target.

\subsection*{cycle-0079: Hardness-Magnification Bridge Realization: beta-zero XOR row-forcing audit}

\paragraph{Learning.} Cycle 79 replaces vague SAT-semantic escape language by the full-row-forcing parameter ForcedRows(Q). A proof-mediated beta-zero XOR certificate that treats generated SAT instances as occurrence labels must force at least \(N^{\epsilon-o(1)}\) complete parity rows; otherwise a private occurrence-label spoof gives a subthreshold sparse-row circuit. The valid contribution is a scoped no-go and an exact extraction obligation for full \(\mathrm{EF}+\mathrm{SAT}_{U}\), not progress on the main conjecture itself.

\paragraph{Pressure.} The derivation is a proof only for \(\mathrm{EF}+\mathrm{SAT}_{U}^{occ}\) or functional label-local certificates. Full \(\mathrm{EF}+\mathrm{SAT}_{U}\) may use real formula syntax to constrain unmentioned generated instances, so an extraction theorem for \(Q_{pi}\) or a nonrelativizing semantic certificate is still missing.

\paragraph{Next.} Instantiate \(\mathrm{Psys}=\mathrm{EF}+\mathrm{SAT}_{U}^{occ}\): define the proof parser that extracts \(Q_{pi}\), prove acceptance transcript invariance under \(A_{Q}\), and compute whether known EF encodings force full rows or merely individual generated formula occurrences.

\paragraph{When to stop.} Stop or quarantine this black-box substitution branch only after the source carrier, SAT-query length exponent \(\beta\), interface size exponent q, source-threshold exponent, and target class have been explicitly computed; quarantine if every sharpened interface still has positive \(\beta\) and therefore only excludes a bounded range of SAT circuit exponents rather than NP subset P/poly.

\paragraph{Do not stop prematurely if.} Do not stop before deriving the transfer inequality from \(\mathrm{SAT}\in \mathrm{SIZE}(n^{c})\) to the carrier upper bound, identifying whether \(\beta=0\), \(\beta>0\), or the argument is genuinely nonblackbox, and recording the barrier status of the proposed escape.

\paragraph{How it adds to the global problem.} This move adds to the P vs NP problem because it identifies whether the transfer loses all-exponent strength through a positive target-query exponent \(\beta\); it separates true hardness magnification from projection or padding arguments that only move proof debt; it keeps barrier status and nonuniform/advice leakage attached to the actual transfer theorem. The current branch-control decision is continue\_run\_next. Payoff criterion: For the proof-complexity branch, separate fixed-family proof hardness from language hardness by naming the proof system, shortcut algorithm, and semantic lower-bound target.

\subsection*{cycle-0080: Hardness-Magnification Bridge Realization}

\paragraph{Learning.} Cycle 80 discharges the scoped \(\mathrm{EF}+\mathrm{SAT}_{U}^{occ}\) proof packet: occurrence-local functional certificates for the beta-zero XOR carrier must force at least \(N^{\epsilon-o(1)}\) complete parity rows, or a private occurrence-label spoof plus sparse-row equality patch gives a valid subthreshold falsifier. The live gap is no longer ordinary soundness; it is the full \(\mathrm{EF}+\mathrm{SAT}_{U}\) extraction of real SAT-semantic reasoning into syntactic \(ForcedRows\), or a deliberate restriction of the paper theorem to \(\mathrm{EF}+\mathrm{SAT}_{U}^{occ}\).

\paragraph{Pressure.} The derivation is valid for \(\mathrm{EF}+\mathrm{SAT}_{U}^{occ}\). It is blocked for full \(\mathrm{EF}+\mathrm{SAT}_{U}\) by the missing extraction lemma from real SAT semantic derivations to syntactic row forcing; full EF may use formula structure in a nonrelativizing way that private occurrence-label spoofing does not preserve.

\paragraph{Next.} Attempt the full \(\mathrm{EF}+\mathrm{SAT}_{U}\) boundary lemma: inspect a standard encoding of \(\mathrm{SAT}_{U}\) extension axioms and decide whether each semantic rule is occurrence-local, charges enough advice into \(S_{Q}\), or yields a nonrelativizing complete-row extraction.

\paragraph{When to stop.} Stop or quarantine this black-box substitution branch only after the source carrier, SAT-query length exponent \(\beta\), interface size exponent q, source-threshold exponent, and target class have been explicitly computed; quarantine if every sharpened interface still has positive \(\beta\) and therefore only excludes a bounded range of SAT circuit exponents rather than NP subset P/poly.

\paragraph{Do not stop prematurely if.} Do not stop before deriving the transfer inequality from \(\mathrm{SAT}\in \mathrm{SIZE}(n^{c})\) to the carrier upper bound, identifying whether \(\beta=0\), \(\beta>0\), or the argument is genuinely nonblackbox, and recording the barrier status of the proposed escape.

\paragraph{How it adds to the global problem.} This move adds to the P vs NP problem because it identifies whether the transfer loses all-exponent strength through a positive target-query exponent \(\beta\); it tests whether a meta-complexity carrier gives a weaker source obligation than the original P vs NP/P-poly target; it separates true hardness magnification from projection or padding arguments that only move proof debt; it keeps barrier status and nonuniform/advice leakage attached to the actual transfer theorem. The current branch-control decision is continue\_run\_next. Payoff criterion: For the proof-complexity branch, separate fixed-family proof hardness from language hardness by naming the proof system, shortcut algorithm, and semantic lower-bound target.

\subsection*{cycle-0081: Hardness-Magnification Bridge Realization}

\paragraph{Learning.} Cycle 81 separates the scoped occurrence-local theorem from the full \(\mathrm{EF}+\mathrm{SAT}_{U}\) ambition. The valid contribution is now sharper: private row spoofing is a genuine falsifier for occurrence-local functional certificates, while full \(\mathrm{EF}+\mathrm{SAT}_{U}\) requires a nonrelativizing semantic-to-occurrence extraction theorem because real \(\mathrm{SAT}_{U}\) axioms are not invariant under independent label flips.

\paragraph{Pressure.} Semantic-to-occurrence shadow extraction: no proof currently shows that full \(\mathrm{EF}+\mathrm{SAT}_{U}\) reasoning about actual SAT formulas can be converted into a functional occurrence-local \(Q_{pi}\) with subthreshold row residual and \(N^{o(1)}\) overhead.

\paragraph{Next.} Formalize the semantic-coupling-width thought experiment as a lemma: for paired formulas F and F OR z, show exactly which \(\mathrm{SAT}_{U}\) implication lines break private-label invariance and whether their violations can be charged to complete forced rows below the \(N^{\epsilon}\) patch threshold.

\paragraph{When to stop.} Stop or quarantine this black-box substitution branch only after the source carrier, SAT-query length exponent \(\beta\), interface size exponent q, source-threshold exponent, and target class have been explicitly computed; quarantine if every sharpened interface still has positive \(\beta\) and therefore only excludes a bounded range of SAT circuit exponents rather than NP subset P/poly.

\paragraph{Do not stop prematurely if.} Do not stop before deriving the transfer inequality from \(\mathrm{SAT}\in \mathrm{SIZE}(n^{c})\) to the carrier upper bound, identifying whether \(\beta=0\), \(\beta>0\), or the argument is genuinely nonblackbox, and recording the barrier status of the proposed escape.

\paragraph{How it adds to the global problem.} This move adds to the P vs NP problem because it identifies whether the transfer loses all-exponent strength through a positive target-query exponent \(\beta\); it keeps barrier status and nonuniform/advice leakage attached to the actual transfer theorem. The current branch-control decision is continue\_run\_next. Payoff criterion: For this route, state the exact barrier premise it uses or escapes: relativizing, natural, algebrizing, nonuniform, or advice-leaking; promotion requires the escape to be attached to the specific transfer theorem.

\subsection*{cycle-0082: Hardness-Magnification Bridge Realization}

\paragraph{Learning.} Cycle 82 upgrades the proof-delivery packet by isolating a row-normal parser theorem: occurrence-local functional certificates are valid only if they force \(N^{\epsilon-o(1)}\) complete rows. The remaining gap is exactly full \(\mathrm{EF}+\mathrm{SAT}_{U}\) semantic-to-occurrence extraction.

\paragraph{Pressure.} The derivation is valid for row-normal occurrence-local transcripts; full \(\mathrm{EF}+\mathrm{SAT}_{U}\) may use real SAT syntax to create nonrelativizing dependencies not represented by \(Q_{pi}\).

\paragraph{Next.} Construct the smallest EF extension-rule countermodel that hides row parity without a row-normal node, or prove a normalization that exposes such parity with subthreshold overhead.

\paragraph{When to stop.} Stop or quarantine this black-box substitution branch only after the source carrier, SAT-query length exponent \(\beta\), interface size exponent q, source-threshold exponent, and target class have been explicitly computed; quarantine if every sharpened interface still has positive \(\beta\) and therefore only excludes a bounded range of SAT circuit exponents rather than NP subset P/poly.

\paragraph{Do not stop prematurely if.} Do not stop before deriving the transfer inequality from \(\mathrm{SAT}\in \mathrm{SIZE}(n^{c})\) to the carrier upper bound, identifying whether \(\beta=0\), \(\beta>0\), or the argument is genuinely nonblackbox, and recording the barrier status of the proposed escape.

\paragraph{How it adds to the global problem.} This move adds to the P vs NP problem because it identifies whether the transfer loses all-exponent strength through a positive target-query exponent \(\beta\); it keeps barrier status and nonuniform/advice leakage attached to the actual transfer theorem. The current branch-control decision is adapt\_then\_run. Payoff criterion: For this route, state the exact barrier premise it uses or escapes: relativizing, natural, algebrizing, nonuniform, or advice-leaking; promotion requires the escape to be attached to the specific transfer theorem.

\subsection*{cycle-0083: Hardness-Magnification Bridge Realization}

\paragraph{Learning.} Cycle 83 upgrades the proof-delivery packet from a parser slogan to a conditional theorem: row-separable row-normal functional occurrence certificates must force \(N^{\epsilon-o(1)}\) complete XOR rows, or private occurrence-label spoofing yields a valid subthreshold circuit for the induced carrier. The remaining open obligation is no longer ordinary soundness but the exact row-normalization/extraction step for broader \(\mathrm{EF}+\mathrm{SAT}_{U}\).

\paragraph{Pressure.} The proof is conditional on a row-separable row-normal parser. Arbitrary full \(\mathrm{EF}+\mathrm{SAT}_{U}\) can use real formula semantics, and arbitrary occurrence-local EF may contain short cross-row parity constraints unless the calculus forbids or charges them.

\paragraph{Next.} Instantiate the \(\mathrm{EF}+\mathrm{SAT}_{U}^{occ,row}\) rule ledger: for each extension, weakening, substitution, \(\mathrm{SAT}_{U}\) occurrence axiom, and derived parity line, mark whether it is \(Q_{pi}-local\), row-separable, or a charged cross-row constraint; the first uncharged cross-row rule becomes the countermodel or repair target.

\paragraph{When to stop.} Stop or quarantine this black-box substitution branch only after the source carrier, SAT-query length exponent \(\beta\), interface size exponent q, source-threshold exponent, and target class have been explicitly computed; quarantine if every sharpened interface still has positive \(\beta\) and therefore only excludes a bounded range of SAT circuit exponents rather than NP subset P/poly.

\paragraph{Do not stop prematurely if.} Do not stop before deriving the transfer inequality from \(\mathrm{SAT}\in \mathrm{SIZE}(n^{c})\) to the carrier upper bound, identifying whether \(\beta=0\), \(\beta>0\), or the argument is genuinely nonblackbox, and recording the barrier status of the proposed escape.

\paragraph{How it adds to the global problem.} This move adds to the P vs NP problem because it identifies whether the transfer loses all-exponent strength through a positive target-query exponent \(\beta\); it separates true hardness magnification from projection or padding arguments that only move proof debt; it keeps barrier status and nonuniform/advice leakage attached to the actual transfer theorem. The current branch-control decision is continue\_run\_next. Payoff criterion: For the proof-complexity branch, separate fixed-family proof hardness from language hardness by naming the proof system, shortcut algorithm, and semantic lower-bound target.

\subsection*{cycle-0084: Hardness-Magnification Bridge Realization}

\paragraph{Learning.} Cycle 84 turns the \(ForcedRows\) packet into a proof-shaped conditional theorem: row-normal occurrence-local functional certificates must force \(N^{\epsilon-o(1)}\) complete XOR rows, because otherwise private label spoofing plus sparse equality patching is a valid falsifier. The remaining gap is exactly semantic-to-occurrence extraction for full \(\mathrm{EF}+\mathrm{SAT}_{U}\), not ordinary soundness or generic natural-proofs folklore.

\paragraph{Pressure.} The derivation is complete only for row-normal occurrence-local functional certificates. Full \(\mathrm{EF}+\mathrm{SAT}_{U}\) remains blocked at semantic-to-occurrence extraction because \(\mathrm{SAT}_{U}\) axioms can couple labels through real formula syntax and witnesses, so private occurrence-label spoofs may not preserve acceptance.

\paragraph{Next.} Audit one standard full \(\mathrm{EF}+\mathrm{SAT}_{U}\) encoding rule: attempt to extract its occurrence-label dependency into \(Q_{pi}\) preserving \(A_{Q}-invariance\); if a witness or clause-soundness axiom breaks invariance, record the minimal rule-level countermodel and restrict the theorem statement to \(\mathrm{EF}+\mathrm{SAT}_{U}^{occ,rn}\).

\paragraph{When to stop.} Stop or quarantine this black-box substitution branch only after the source carrier, SAT-query length exponent \(\beta\), interface size exponent q, source-threshold exponent, and target class have been explicitly computed; quarantine if every sharpened interface still has positive \(\beta\) and therefore only excludes a bounded range of SAT circuit exponents rather than NP subset P/poly.

\paragraph{Do not stop prematurely if.} Do not stop before deriving the transfer inequality from \(\mathrm{SAT}\in \mathrm{SIZE}(n^{c})\) to the carrier upper bound, identifying whether \(\beta=0\), \(\beta>0\), or the argument is genuinely nonblackbox, and recording the barrier status of the proposed escape.

\paragraph{How it adds to the global problem.} This move adds to the P vs NP problem because it identifies whether the transfer loses all-exponent strength through a positive target-query exponent \(\beta\); it separates true hardness magnification from projection or padding arguments that only move proof debt; it keeps barrier status and nonuniform/advice leakage attached to the actual transfer theorem. The current branch-control decision is continue\_run\_next. Payoff criterion: For the proof-complexity branch, separate fixed-family proof hardness from language hardness by naming the proof system, shortcut algorithm, and semantic lower-bound target.

\subsection*{cycle-0085: Hardness-Magnification Bridge Realization}

\paragraph{Learning.} Cycle 85 sharpens the formula-semantic extraction fork: full occurrence parsing is false already for \(\mathrm{EF}+SATeq_{N}\) because EqSAT padding identities are real SAT-semantic and nonrelativizing. The valid repair is a semantic-link trichotomy: nonfunctional certificates must pay either occurrence full rows or EqSAT-linked rows at the \(N^{\epsilon-o(1)}\) threshold, otherwise sparse-row residual patching falsifies the source lower-bound certificate.

\paragraph{Pressure.} The derivation is complete only for \(\mathrm{EF}+SATeq_{N}\) where all SAT semantics factors through EqSAT links. It does not yet cover full \(\mathrm{EF}+\mathrm{SAT}_{U}\) with witness, self-reduction, or minimization rules unless those rules are decomposed into label-functional inferences plus row-chargeable semantic links.

\paragraph{Next.} Attempt the finite \(\mathrm{EF}+SATeq_{N}\) parser proof: formalize proof lines as Boolean circuits over SAT-atoms, prove label-functionality for EF rules, charge every EqSAT endpoint to SemRows, and test whether SAT self-reduction axioms in full \(\mathrm{EF}+\mathrm{SAT}_{U}\) admit the same charge or force a new semantic-escape lemma.

\paragraph{When to stop.} Stop or quarantine this black-box substitution branch only after the source carrier, SAT-query length exponent \(\beta\), interface size exponent q, source-threshold exponent, and target class have been explicitly computed; quarantine if every sharpened interface still has positive \(\beta\) and therefore only excludes a bounded range of SAT circuit exponents rather than NP subset P/poly.

\paragraph{Do not stop prematurely if.} Do not stop before deriving the transfer inequality from \(\mathrm{SAT}\in \mathrm{SIZE}(n^{c})\) to the carrier upper bound, identifying whether \(\beta=0\), \(\beta>0\), or the argument is genuinely nonblackbox, and recording the barrier status of the proposed escape.

\paragraph{How it adds to the global problem.} This move adds to the P vs NP problem because it identifies whether the transfer loses all-exponent strength through a positive target-query exponent \(\beta\); it separates true hardness magnification from projection or padding arguments that only move proof debt; it keeps barrier status and nonuniform/advice leakage attached to the actual transfer theorem. The current branch-control decision is continue\_run\_next. Payoff criterion: For this route, state the exact barrier premise it uses or escapes: relativizing, natural, algebrizing, nonuniform, or advice-leaking; promotion requires the escape to be attached to the specific transfer theorem.

\subsection*{cycle-0086: Hardness-Magnification Bridge Realization}

\paragraph{Learning.} Cycle 86 resolves one part of the formula-semantic fork: occurrence-only extraction is false for \(\mathrm{EF}+SATeq_{N}\) because real EqSAT identities are nonrelativizing semantic links. The valid replacement object is a semantic-link residual transcript; subthreshold link coverage is a falsifier for source lower-bound certificates, while the open proof task is a normal form for arbitrary \(\mathrm{EF}+SATeq_{N}\) reasoning.

\paragraph{Pressure.} The normal form is not proved for arbitrary EF reasoning that invokes SAT self-reducibility, extension variables encoding witnesses, or nonlocal formula transformations not decomposed into explicit EqSAT links.

\paragraph{Next.} Work the \(\mathrm{EF}+SATeq_{N}\) normal-form proof rule by rule: padding, renaming, duplicate clauses, extension-variable abbreviations, and SAT self-reduction; either translate each into a bounded EqSAT-link transcript or isolate the first noncompressible SAT-semantic principle as the next theorem-shaped counterexample.

\paragraph{When to stop.} Stop or quarantine this black-box substitution branch only after the source carrier, SAT-query length exponent \(\beta\), interface size exponent q, source-threshold exponent, and target class have been explicitly computed; quarantine if every sharpened interface still has positive \(\beta\) and therefore only excludes a bounded range of SAT circuit exponents rather than NP subset P/poly.

\paragraph{Do not stop prematurely if.} Do not stop before deriving the transfer inequality from \(\mathrm{SAT}\in \mathrm{SIZE}(n^{c})\) to the carrier upper bound, identifying whether \(\beta=0\), \(\beta>0\), or the argument is genuinely nonblackbox, and recording the barrier status of the proposed escape.

\paragraph{How it adds to the global problem.} This move adds to the P vs NP problem because it identifies whether the transfer loses all-exponent strength through a positive target-query exponent \(\beta\); it tests whether a meta-complexity carrier gives a weaker source obligation than the original P vs NP/P-poly target; it separates true hardness magnification from projection or padding arguments that only move proof debt; it keeps barrier status and nonuniform/advice leakage attached to the actual transfer theorem. The current branch-control decision is continue\_run\_next. Payoff criterion: For the proof-complexity branch, separate fixed-family proof hardness from language hardness by naming the proof system, shortcut algorithm, and semantic lower-bound target.

\subsection*{cycle-0087: Hardness-Magnification Bridge Realization}

\paragraph{Learning.} Cycle 87 replaces the vague semantic-link normal-form obligation by an exact residual-code invariant: after EqSAT quotienting, a proof-mediated certificate that is functional over quotient labels is valid only if its touched-label coset has row-parity distance at least \(N^{\epsilon-o(1)}\). Low quotient-code distance is a concrete falsifier; arbitrary \(\mathrm{EF}+SATeq_{N}\) remains open because it may use genuinely SAT-semantic nonrelativizing principles outside the quotient parser.

\paragraph{Pressure.} The derivation does not produce a normal form for arbitrary \(\mathrm{EF}+SATeq_{N}\) proofs. Full EF can use nonrelativizing SAT-semantic principles, formula self-reductions, and derived equality chains whose acceptance is not a predicate over a bounded W,E quotient transcript.

\paragraph{Next.} Instantiate the padded-copy EqSAT quotient family \(E_{k}\), compute the row-incidence matrix \(A_{E}\), and decide whether its constrained cosets have low-weight representatives below \(N^\epsilon\) or certify a valid semantic-link residual lower bound.

\paragraph{When to stop.} Stop or quarantine this black-box substitution branch only after the source carrier, SAT-query length exponent \(\beta\), interface size exponent q, source-threshold exponent, and target class have been explicitly computed; quarantine if every sharpened interface still has positive \(\beta\) and therefore only excludes a bounded range of SAT circuit exponents rather than NP subset P/poly.

\paragraph{Do not stop prematurely if.} Do not stop before deriving the transfer inequality from \(\mathrm{SAT}\in \mathrm{SIZE}(n^{c})\) to the carrier upper bound, identifying whether \(\beta=0\), \(\beta>0\), or the argument is genuinely nonblackbox, and recording the barrier status of the proposed escape.

\paragraph{How it adds to the global problem.} This move adds to the P vs NP problem because it identifies whether the transfer loses all-exponent strength through a positive target-query exponent \(\beta\); it separates true hardness magnification from projection or padding arguments that only move proof debt; it keeps barrier status and nonuniform/advice leakage attached to the actual transfer theorem. The current branch-control decision is continue\_run\_next. Payoff criterion: For this route, state the exact barrier premise it uses or escapes: relativizing, natural, algebrizing, nonuniform, or advice-leaking; promotion requires the escape to be attached to the specific transfer theorem.

\subsection*{cycle-0088: Hardness-Magnification Bridge Realization}

\paragraph{Learning.} Cycle 88 identifies the first concrete SAT-semantic escape from the EqSAT quotient parser: self-reduction OR constraints are not invariant under quotient-label spoofing. The valid paper move is not another magnification slogan but a trichotomy repair: equality-only proofs use quotient residual codes, occurrence expansions pay row-forcing cost, and SR proofs require a constrained residual-code audit. Low constrained row distance is a falsifier for source certificates; no main \(\mathrm{P} \ne \mathrm{NP}\) claim is made.

\paragraph{Pressure.} The derivation does not normalize arbitrary \(\mathrm{EF}+\mathrm{SAT}_{U}\) or arbitrary SAT-semantic reasoning to EqSAT+SR constraints. A stronger proof may use search-to-decision, Cook-Levin witness encodings, or formula minimization principles not represented by the finite SR hypergraph \(K_{pi}\).

\paragraph{Next.} Instantiate \(K_{pi}\) for local one-step SR only: bound \(|K_{pi}|\) by proof size, define its row-parity constraint-distance \(Delta_{K}\), and test whether \(Delta_{K}\) can be reduced to a known SAT promise-hardness statement or whether it is target-equivalent to the desired \(H_{N}\) lower bound.

\paragraph{When to stop.} Stop or quarantine this black-box substitution branch only after the source carrier, SAT-query length exponent \(\beta\), interface size exponent q, source-threshold exponent, and target class have been explicitly computed; quarantine if every sharpened interface still has positive \(\beta\) and therefore only excludes a bounded range of SAT circuit exponents rather than NP subset P/poly.

\paragraph{Do not stop prematurely if.} Do not stop before deriving the transfer inequality from \(\mathrm{SAT}\in \mathrm{SIZE}(n^{c})\) to the carrier upper bound, identifying whether \(\beta=0\), \(\beta>0\), or the argument is genuinely nonblackbox, and recording the barrier status of the proposed escape.

\paragraph{How it adds to the global problem.} This move adds to the P vs NP problem because it identifies whether the transfer loses all-exponent strength through a positive target-query exponent \(\beta\); it keeps barrier status and nonuniform/advice leakage attached to the actual transfer theorem. The current branch-control decision is continue\_run\_next. Payoff criterion: For this route, state the exact barrier premise it uses or escapes: relativizing, natural, algebrizing, nonuniform, or advice-leaking; promotion requires the escape to be attached to the specific transfer theorem.

\subsection*{cycle-0089: Hardness-Magnification Bridge Realization}

\paragraph{Learning.} The cycle upgrades the SR semantic escape into a quantitative constrained-distance invariant. Occurrence spoofing is invalid after SR constraints, but any bounded-SR parser still faces the same row-patching threshold. Progress is route-level certification, not a P vs NP proof.

\paragraph{Pressure.} The missing proof is semantic-to-residual parsing for full \(\mathrm{EF}+SATeq+SR_{h}\); SR OR clauses couple labels across restrictions, so independent occurrence-label spoofing is not valid.

\paragraph{Next.} Fix h with 2\textasciicircum{}h=N\textasciicircum{}\{eta\}; formalize \(Delta_{SR,h}\) for the generated formula closure and try to prove the parser for \(\mathrm{EF}+SATeq+SR_{h}\) restricted to extension variables naming only quotient labels and SR clauses.

\paragraph{When to stop.} Stop or quarantine this black-box substitution branch only after the source carrier, SAT-query length exponent \(\beta\), interface size exponent q, source-threshold exponent, and target class have been explicitly computed; quarantine if every sharpened interface still has positive \(\beta\) and therefore only excludes a bounded range of SAT circuit exponents rather than NP subset P/poly.

\paragraph{Do not stop prematurely if.} Do not stop before deriving the transfer inequality from \(\mathrm{SAT}\in \mathrm{SIZE}(n^{c})\) to the carrier upper bound, identifying whether \(\beta=0\), \(\beta>0\), or the argument is genuinely nonblackbox, and recording the barrier status of the proposed escape.

\paragraph{How it adds to the global problem.} This move adds to the P vs NP problem because it identifies whether the transfer loses all-exponent strength through a positive target-query exponent \(\beta\); it keeps barrier status and nonuniform/advice leakage attached to the actual transfer theorem. The current branch-control decision is continue\_run\_next. Payoff criterion: For the MCSP/MKtP row branch, decide whether the certificate proves direct row hardness or only aggregate/rank-index hardness; record q\_extract=epsilon/(1+D) and q\_force=(D+epsilon)/(1+D) in the same parameter regime.

\subsection*{cycle-0090: Hardness-Magnification Bridge Realization}

\paragraph{Learning.} Cycle 90 converts the SR semantic escape into a costed leaf-solver falsifier: bounded self-reduction depth is not free, but if \(q+\eta+\tau<1+\epsilon\) and the induced row error is below \(N^\epsilon\) then the source lower bound collapses by equality patching. The live inference-quality packet is now theorem-shaped for \(\mathrm{Psys}_{h}\); the only remaining full-system gap is the advice-free parser normal form for arbitrary EF extension variables.

\paragraph{Pressure.} ParserNormalForm(EF+EqSAT+SR\_h): prove that every short proof of LB(H\_N\textasciicircum{}\{SAT\}) either gives a dense functional label property, an occurrence/full-row transcript, or a bounded-SR residual code with advice-free size bound poly(|pi|,L\_N).

\paragraph{Next.} Build the minimal \(\mathrm{EF}+\mathrm{EqSAT}+SR_{1}\) parser countermodel: allow one extension variable Z(F)=SAT(F\_0) OR SAT(F\_1) and test whether it is representable by the SR DAG without advice; if yes, increase to two nested extensions, if no, use it as the semantic escape separating full EF from \(\mathrm{Psys}_{h}\).

\paragraph{When to stop.} Stop or quarantine this black-box substitution branch only after the source carrier, SAT-query length exponent \(\beta\), interface size exponent q, source-threshold exponent, and target class have been explicitly computed; quarantine if every sharpened interface still has positive \(\beta\) and therefore only excludes a bounded range of SAT circuit exponents rather than NP subset P/poly.

\paragraph{Do not stop prematurely if.} Do not stop before deriving the transfer inequality from \(\mathrm{SAT}\in \mathrm{SIZE}(n^{c})\) to the carrier upper bound, identifying whether \(\beta=0\), \(\beta>0\), or the argument is genuinely nonblackbox, and recording the barrier status of the proposed escape.

\paragraph{How it adds to the global problem.} This move adds to the P vs NP problem because it identifies whether the transfer loses all-exponent strength through a positive target-query exponent \(\beta\); it separates true hardness magnification from projection or padding arguments that only move proof debt; it keeps barrier status and nonuniform/advice leakage attached to the actual transfer theorem. The current branch-control decision is continue\_run\_next. Payoff criterion: For this route, state the exact barrier premise it uses or escapes: relativizing, natural, algebrizing, nonuniform, or advice-leaking; promotion requires the escape to be attached to the specific transfer theorem.

\subsection*{cycle-0091: Hardness-Magnification Bridge Realization: beta-zero XOR formula-semantic parser audit}

\paragraph{Learning.} Cycle 91 upgrades the bounded-SR leaf-solver audit by adding an explicit EF extension-support charge and proving a no-go for extension-blind parsers. The valid theorem is restricted but contribution-shaped: subthreshold charged transcripts are upper-bound falsifiers, while arbitrary EF+SAT remains blocked at a precise advice-free normal-form lemma.

\paragraph{Pressure.} Prove or refute the advice-free extension-support normal form for arbitrary EF extension variables in EF+EqSAT+bounded-SR; the known derivation only handles proofs whose extension dependencies are syntactically charged.

\paragraph{Next.} Instantiate \(\mathrm{Psys}_{h,Ext}^{schema}\): EF+EqSAT+depth-h SR with extension definitions generated by row-oblivious schemas. Try to prove ExtSupp(NF(pi))<=poly(|pi|), or construct a minimal nonrelativizing SAT-semantic extension that defeats this normal form.

\paragraph{When to stop.} Stop or quarantine this black-box substitution branch only after the source carrier, SAT-query length exponent \(\beta\), interface size exponent q, source-threshold exponent, and target class have been explicitly computed; quarantine if every sharpened interface still has positive \(\beta\) and therefore only excludes a bounded range of SAT circuit exponents rather than NP subset P/poly.

\paragraph{Do not stop prematurely if.} Do not stop before deriving the transfer inequality from \(\mathrm{SAT}\in \mathrm{SIZE}(n^{c})\) to the carrier upper bound, identifying whether \(\beta=0\), \(\beta>0\), or the argument is genuinely nonblackbox, and recording the barrier status of the proposed escape.

\paragraph{How it adds to the global problem.} This move adds to the P vs NP problem because it identifies whether the transfer loses all-exponent strength through a positive target-query exponent \(\beta\); it keeps barrier status and nonuniform/advice leakage attached to the actual transfer theorem. The current branch-control decision is continue\_run\_next. Payoff criterion: For the proof-complexity branch, separate fixed-family proof hardness from language hardness by naming the proof system, shortcut algorithm, and semantic lower-bound target.

\subsection*{cycle-0092: Hardness-Magnification Bridge Realization}

\paragraph{Learning.} Cycle 92 turns the formula-semantic fork into a sharper proof packet. For support-declared EF+SATeq+bounded-SR, subthreshold parsed support plus bounded leaf solving and sparse residual rows gives a valid small-circuit falsifier for \(\mathrm{LB}_{N}\). The full EF+SAT claim is not proved; the precise remaining obstruction is \(NSP_{N}\), the advice-free extraction of arbitrary extension-variable reasoning into charged row supports and SR leaves. This is route-level certification, not progress on the main \(\mathrm{P} \ne \mathrm{NP}\) conjecture itself.

\paragraph{Pressure.} \(NSP_{N}\): every short arbitrary EF+SATeq+SR proof of \(\mathrm{LB}_{N}\) must admit a support-complete, advice-free transcript with row distance and leaf-solver costs small enough to preserve the \(N^{1+\epsilon}\) threshold. This is unproved and may fail for extension-compressed SAT-semantic reasoning.

\paragraph{Next.} Try to prove \(NSP_{N}\) for the fragment EF+SATeq+SR whose extension definitions are row-affine and support-monotone; if the proof fails, extract the minimal extension rule that creates uncharged SAT-semantic support and state the final theorem scope around \(\mathrm{Psys}_{h}^{supp}\).

\paragraph{When to stop.} Stop or quarantine this black-box substitution branch only after the source carrier, SAT-query length exponent \(\beta\), interface size exponent q, source-threshold exponent, and target class have been explicitly computed; quarantine if every sharpened interface still has positive \(\beta\) and therefore only excludes a bounded range of SAT circuit exponents rather than NP subset P/poly.

\paragraph{Do not stop prematurely if.} Do not stop before deriving the transfer inequality from \(\mathrm{SAT}\in \mathrm{SIZE}(n^{c})\) to the carrier upper bound, identifying whether \(\beta=0\), \(\beta>0\), or the argument is genuinely nonblackbox, and recording the barrier status of the proposed escape.

\paragraph{How it adds to the global problem.} This move adds to the P vs NP problem because it identifies whether the transfer loses all-exponent strength through a positive target-query exponent \(\beta\); it separates true hardness magnification from projection or padding arguments that only move proof debt; it keeps barrier status and nonuniform/advice leakage attached to the actual transfer theorem. The current branch-control decision is continue\_run\_next. Payoff criterion: For this route, state the exact barrier premise it uses or escapes: relativizing, natural, algebrizing, nonuniform, or advice-leaking; promotion requires the escape to be attached to the specific transfer theorem.

\subsection*{cycle-0093: Hardness-Magnification Bridge Realization}

\paragraph{Learning.} Cycle 93 upgrades the formula-semantic extraction fork by naming \(NSP_{N}\) as a support-transparency lemma for unrestricted EF extension variables and proving the conditional transcript falsifier around it. The valid restricted contribution is support-declared \(\mathrm{EF}+SATeq+bounded-SR_{h}\): any subthreshold parsed transcript collapses \(\mathrm{LB}_{N}\) by a small circuit plus sparse-row patching. The full EF+SAT claim remains blocked exactly at advice-free extraction of extension-variable support, not at generic soundness or barrier folklore.

\paragraph{Pressure.} \(NSP_{N}\): an advice-free normal-support parser for unrestricted \(\mathrm{EF}+SATeq+SR_{h}\) extension variables. The missing lemma must extract charged row supports and bounded-SR leaves from arbitrary extension definitions without SAT-label advice and without solving a target-equivalent semantic support problem.

\paragraph{Next.} Attempt an explicit \(NSP_{N}\) countermodel for unrestricted EF extension variables: construct a small row-mask extension family \(e_{R}\), prove that occurrence-only parsing misses its support, and decide whether support transparency can be repaired by charging C(Supp(e\_R)) without making the parser target-equivalent.

\paragraph{When to stop.} Stop or quarantine this black-box substitution branch only after the source carrier, SAT-query length exponent \(\beta\), interface size exponent q, source-threshold exponent, and target class have been explicitly computed; quarantine if every sharpened interface still has positive \(\beta\) and therefore only excludes a bounded range of SAT circuit exponents rather than NP subset P/poly.

\paragraph{Do not stop prematurely if.} Do not stop before deriving the transfer inequality from \(\mathrm{SAT}\in \mathrm{SIZE}(n^{c})\) to the carrier upper bound, identifying whether \(\beta=0\), \(\beta>0\), or the argument is genuinely nonblackbox, and recording the barrier status of the proposed escape.

\paragraph{How it adds to the global problem.} This move adds to the P vs NP problem because it identifies whether the transfer loses all-exponent strength through a positive target-query exponent \(\beta\); it separates true hardness magnification from projection or padding arguments that only move proof debt; it keeps barrier status and nonuniform/advice leakage attached to the actual transfer theorem. The current branch-control decision is continue\_run\_next. Payoff criterion: For this route, state the exact barrier premise it uses or escapes: relativizing, natural, algebrizing, nonuniform, or advice-leaking; promotion requires the escape to be attached to the specific transfer theorem.

\subsection*{cycle-0094: Hardness-Magnification Bridge Realization}

\paragraph{Learning.} Cycle 94 closes a concrete support-transparency subcase: bit-explicit EF extension variables can be parsed without support declarations, giving a valid transcript falsifier theorem. The full formula-semantic extraction fork is now sharper: the obstruction is not arbitrary EF extension variables but succinct FormulaEval-SAT comprehension, which defeats size-only occurrence parsing unless its row-index expansion is charged.

\paragraph{Pressure.} \(NSP_{sym}\): an advice-free parser for succinct formula-indexed SAT atoms must either expand \(N^{1+q+o(1)}\) concrete row-index labels or prove a new semantic-support theorem that remains completion-sound without occurrence expansion.

\paragraph{Next.} Formalize \(\mathrm{Psys}_{sym}\) on finite length N: specify whether FormulaEval-SAT comprehension is a proof line, axiom schema, or extension rule; then prove either an expansion-charge parser with cost \(N^{1+q+o(1)}\) or a minimal countermodel showing surface occurrence completion is unsound.

\paragraph{When to stop.} Stop or quarantine this black-box substitution branch only after the source carrier, SAT-query length exponent \(\beta\), interface size exponent q, source-threshold exponent, and target class have been explicitly computed; quarantine if every sharpened interface still has positive \(\beta\) and therefore only excludes a bounded range of SAT circuit exponents rather than NP subset P/poly.

\paragraph{Do not stop prematurely if.} Do not stop before deriving the transfer inequality from \(\mathrm{SAT}\in \mathrm{SIZE}(n^{c})\) to the carrier upper bound, identifying whether \(\beta=0\), \(\beta>0\), or the argument is genuinely nonblackbox, and recording the barrier status of the proposed escape.

\paragraph{How it adds to the global problem.} This move adds to the P vs NP problem because it identifies whether the transfer loses all-exponent strength through a positive target-query exponent \(\beta\); it separates true hardness magnification from projection or padding arguments that only move proof debt; it keeps barrier status and nonuniform/advice leakage attached to the actual transfer theorem; it searches the remaining plausible escape route: nonblackbox or subpolynomial-interface carriers. The current branch-control decision is continue\_run\_next. Payoff criterion: For this route, state the exact barrier premise it uses or escapes: relativizing, natural, algebrizing, nonuniform, or advice-leaking; promotion requires the escape to be attached to the specific transfer theorem.

\subsection*{cycle-0095: Hardness-Magnification Bridge Realization: bit-explicit parser and FormulaEval boundary}

\paragraph{Learning.} Cycle 95 upgrades cycle 94 from a parser slogan to a restricted theorem packet: \(\mathrm{Psys}_{bEF}\) proofs with bit-explicit EF support have a valid charged transcript and a concrete \(\mathrm{LB}_{N}\) falsifier inequality. The surviving strong-system gap is exactly symbolic FormulaEval-SAT: row expansion costs \(N^{1+q-o(1)}\), while free symbolic comprehension has a private-row countermodel to surface occurrence completion.

\paragraph{Pressure.} The remaining named gap is \(FE_{expand,N}\): prove that any sound use of symbolic FormulaEval-SAT either expands to Omega(N\textasciicircum{}\{1+q-o(1)\}) row charge or is explicitly outside the parser model as a semantic oracle. \(\mathrm{Psys}_{bEF}\) itself is not blocked under the stated bit-explicit assumptions.

\paragraph{Next.} Instantiate \(FE_{expand,N}\) on the smallest symbolic FormulaEval-SAT rule: one generator circuit G\_N(x,j), one asserted vector y\_j=SAT(G\_N(x,j)), and \(\alpha\) row instantiations. Prove the private-row countermodel for \(\alpha<q\) and the \(N^{1+q-o(1)}\) expansion charge for \(\alpha=q\).

\paragraph{When to stop.} Stop or quarantine this P vs NP branch only after it states the source object, target class, complexity model, transfer inequality, quantifier strength, and barrier status; quarantine if two sharpened attempts still cannot say whether the route implies \(\mathrm{SAT}\notin \mathrm{P}/\mathrm{poly}\), \(\mathrm{NP}\not\subset \mathrm{P}/\mathrm{poly}\), or only a weaker/proxy lower bound.

\paragraph{Do not stop prematurely if.} Do not stop before the branch has produced a concrete source-to-target transfer theorem, an explicit obstruction lemma, or a barrier-aware reason why the route cannot reach the P vs NP/P/poly target.

\paragraph{How it adds to the global problem.} This move adds to the P vs NP problem because it keeps barrier status and nonuniform/advice leakage attached to the actual transfer theorem. The current branch-control decision is backtrack\_to\_sharpen. Payoff criterion: For this route, state the exact barrier premise it uses or escapes: relativizing, natural, algebrizing, nonuniform, or advice-leaking; promotion requires the escape to be attached to the specific transfer theorem.

\subsection*{cycle-0096: Hardness-Magnification Bridge Realization}

\paragraph{Learning.} Cycle 96 closes a sharper proof-delivery subgoal: \(\mathrm{Psys}_{bEF}\) parsing survives bit-level audit, and \(\mathrm{Psys}_{sym}\) FormulaEval-SAT has a precise expansion-or-countermodel boundary. The new contribution is not a \(\mathrm{P} \ne \mathrm{NP}\) advance; it is a valid route-level theorem/countermodel packet showing that symbolic formula evaluation either pays \(N^{1+q+o(1)}\) row expansion or must be treated as a semantic-oracle escape outside surface occurrence parsing.

\paragraph{Pressure.} \(SemanticForcing_{FE}\) remains open: a real-SAT FormulaEval semantic oracle might force private row labels by nonrelativizing reasoning not recoverable from occurrence syntax. The branch now treats that as a named semantic-oracle escape, not as a valid \(\mathrm{Psys}_{bEF}\) parser step.

\paragraph{Next.} Fix one concrete \(\mathrm{Psys}_{sym}\) syntax: FormulaEval-SAT as either an axiom schema over A(i,j,u), an EF extension rule producing formula bits, a proof line with explicit \(D_{u}\), or a semantic oracle. For that syntax, prove the expansion parser or instantiate the minimal private-row countermodel with the smallest pi, one unexpanded pair (i*,j*), and a verified \(\Gamma_{sym}\) inequality.

\paragraph{When to stop.} Stop or quarantine this P vs NP branch only after it states the source object, target class, complexity model, transfer inequality, quantifier strength, and barrier status; quarantine if two sharpened attempts still cannot say whether the route implies \(\mathrm{SAT}\notin \mathrm{P}/\mathrm{poly}\), \(\mathrm{NP}\not\subset \mathrm{P}/\mathrm{poly}\), or only a weaker/proxy lower bound.

\paragraph{Do not stop prematurely if.} Do not stop before the branch has produced a concrete source-to-target transfer theorem, an explicit obstruction lemma, or a barrier-aware reason why the route cannot reach the P vs NP/P/poly target.

\paragraph{How it adds to the global problem.} This move adds to the P vs NP problem because it keeps barrier status and nonuniform/advice leakage attached to the actual transfer theorem. The current branch-control decision is backtrack\_to\_sharpen. Payoff criterion: For this route, state the exact barrier premise it uses or escapes: relativizing, natural, algebrizing, nonuniform, or advice-leaking; promotion requires the escape to be attached to the specific transfer theorem.

\subsection*{cycle-0097: Hardness-Magnification Bridge Realization}

\paragraph{Learning.} Cycle 97 closes the support-density subgoal for the \(\mathrm{Psys}_{bEF}\) packet: FormulaEval support must have absolute size \(N^{\epsilon-o(1)}\), hence density \(N^{\epsilon-q-o(1)}\) inside \(M_{N}=N^{q+o(1)}\) slots. This yields a sharp q-trichotomy and upgrades the FormulaEval boundary to an expansion-or-private-row-countermodel test. The result is a valid route-level falsifier theorem, not evidence for \(\mathrm{P} \ne \mathrm{NP}\).

\paragraph{Pressure.} The derivation does not prove a parser theorem for \(\mathrm{Psys}_{sym}\) when FormulaEval-SAT is a genuine semantic oracle; the named missing lemma is FormulaEval-forcing: real SAT semantics must force every unexpanded row/slot value from the charged transcript with subthreshold overhead.

\paragraph{Next.} Instantiate \(\mathrm{Psys}_{sym}\) in two concrete variants: FormulaEval as an EF extension rule with row-index bits and FormulaEval as a semantic oracle. For each, compute K\_idx(pi), X\_FE(K), and whether the private-row countermodel survives SATeq and bounded-SR constraints.

\paragraph{When to stop.} Stop or quarantine this P vs NP branch only after it states the source object, target class, complexity model, transfer inequality, quantifier strength, and barrier status; quarantine if two sharpened attempts still cannot say whether the route implies \(\mathrm{SAT}\notin \mathrm{P}/\mathrm{poly}\), \(\mathrm{NP}\not\subset \mathrm{P}/\mathrm{poly}\), or only a weaker/proxy lower bound.

\paragraph{Do not stop prematurely if.} Do not stop before the branch has produced a concrete source-to-target transfer theorem, an explicit obstruction lemma, or a barrier-aware reason why the route cannot reach the P vs NP/P/poly target.

\paragraph{How it adds to the global problem.} This move adds to the P vs NP problem because it keeps barrier status and nonuniform/advice leakage attached to the actual transfer theorem. The current branch-control decision is backtrack\_to\_sharpen. Payoff criterion: For this route, state the exact barrier premise it uses or escapes: relativizing, natural, algebrizing, nonuniform, or advice-leaking; promotion requires the escape to be attached to the specific transfer theorem.

\subsection*{cycle-0098: Hardness-Magnification Bridge Realization}

\paragraph{Learning.} Cycle 98 upgrades the prior packet from theorem slogan to dependency-checked inference structure: \(\Gamma_{bEF}\) is justified by separate EF-dependency, SATeq-link, bounded-SR, and sparse-row charges; \(\mathrm{Psys}_{bEF}\) gives a valid restricted falsifier theorem; full \(\mathrm{Psys}_{sym}\) is blocked exactly at \(FE_{Extract,N}\), not at generic naturalness or relativization folklore.

\paragraph{Pressure.} The restricted \(\mathrm{Psys}_{bEF}\) theorem is derivation-complete at proof-packet level, but the full \(\mathrm{Psys}_{sym}\) extension remains blocked at \(FE_{Extract,N}\): an advice-free theorem converting semantic FormulaEval-SAT comprehension into charged occurrence rows without paying \(N^{1+q+o(1)}\) or admitting the private-row countermodel.

\paragraph{Next.} Write the formal \(FE_{Extract,N}\) proposition in two columns: left column row-expanded FormulaEval axiom schema with \(N^{1+q+o(1)}\) charge, right column semantic-oracle version with the smallest private-row countermodel; verify which column matches each \(\mathrm{Psys}_{sym}\) rule.

\paragraph{When to stop.} Stop or quarantine this black-box substitution branch only after the source carrier, SAT-query length exponent \(\beta\), interface size exponent q, source-threshold exponent, and target class have been explicitly computed; quarantine if every sharpened interface still has positive \(\beta\) and therefore only excludes a bounded range of SAT circuit exponents rather than NP subset P/poly.

\paragraph{Do not stop prematurely if.} Do not stop before deriving the transfer inequality from \(\mathrm{SAT}\in \mathrm{SIZE}(n^{c})\) to the carrier upper bound, identifying whether \(\beta=0\), \(\beta>0\), or the argument is genuinely nonblackbox, and recording the barrier status of the proposed escape.

\paragraph{How it adds to the global problem.} This move adds to the P vs NP problem because it identifies whether the transfer loses all-exponent strength through a positive target-query exponent \(\beta\); it tests whether a meta-complexity carrier gives a weaker source obligation than the original P vs NP/P-poly target; it separates true hardness magnification from projection or padding arguments that only move proof debt; it keeps barrier status and nonuniform/advice leakage attached to the actual transfer theorem. The current branch-control decision is continue\_run\_next. Payoff criterion: For this route, state the exact barrier premise it uses or escapes: relativizing, natural, algebrizing, nonuniform, or advice-leaking; promotion requires the escape to be attached to the specific transfer theorem.

\subsection*{cycle-0099: Hardness-Magnification Bridge Realization}

\paragraph{Learning.} Cycle 99 converts the remaining \(\mathrm{Psys}_{sym}\) ambiguity into a sharper theorem/countermodel fork. FormulaEval-SAT cannot be treated as a free occurrence-completion rule: an expanded schema pays \(N^{1+q+o(1)}\) row-slot charge, while an unexpanded symbolic line has a private-row falsifier unless a nonrelativizing \(FE_{Extract,N}\) theorem is supplied. The valid contribution is a scoped parser-boundary packet, not evidence for \(\mathrm{P} \ne \mathrm{NP}\).

\paragraph{Pressure.} The derivation does not rule out a real-SAT-specific, nonrelativizing \(FE_{Extract,N}\) theorem showing that true formula evaluation semantics force the missing labels without row-slot expansion. It only proves the surface-occurrence parser no-go and the expansion charge for explicit schemas.

\paragraph{Next.} Write the formal \(FE_{Extract,N}\) blocked-lemma proof: define the quotient-label model after \(SATeq/SR_{h}\) closure, prove the unmentioned row-slot existence bound, construct the two assignments Y,Y', and verify that the carrier row parity changes while the parsed transcript remains identical.

\paragraph{When to stop.} Stop or quarantine this P vs NP branch only after it states the source object, target class, complexity model, transfer inequality, quantifier strength, and barrier status; quarantine if two sharpened attempts still cannot say whether the route implies \(\mathrm{SAT}\notin \mathrm{P}/\mathrm{poly}\), \(\mathrm{NP}\not\subset \mathrm{P}/\mathrm{poly}\), or only a weaker/proxy lower bound.

\paragraph{Do not stop prematurely if.} Do not stop before the branch has produced a concrete source-to-target transfer theorem, an explicit obstruction lemma, or a barrier-aware reason why the route cannot reach the P vs NP/P/poly target.

\paragraph{How it adds to the global problem.} This move adds to the P vs NP problem because it keeps barrier status and nonuniform/advice leakage attached to the actual transfer theorem. The current branch-control decision is backtrack\_to\_sharpen. Payoff criterion: For this route, state the exact barrier premise it uses or escapes: relativizing, natural, algebrizing, nonuniform, or advice-leaking; promotion requires the escape to be attached to the specific transfer theorem.

\subsection*{cycle-0100: Hardness-Magnification Bridge Realization}

\paragraph{Learning.} Cycle 100 isolates the \(bounded-SR_{h}\) normal form at bit level: SR contributes exactly B\_SR(pi) to \(\Gamma_{bEF}\) for \(\mathrm{Psys}_{bEF}^{-}\). The remaining full-system obstruction is not generic soundness but \(FE_{Extract,N}\) for symbolic FormulaEval rows.

\paragraph{Pressure.} The derivation is valid only for \(\mathrm{Psys}_{bEF}^{-}\). Full \(\mathrm{Psys}_{sym}\) remains blocked at \(FE_{Extract,N}\): a symbolic FormulaEval-SAT line may select rows without exposing K\_idx(pi) or Nbr(j,u).

\paragraph{Next.} Prove \(RowSetEnumerator_{N}\) for an interval-union or \(\mathrm{AC0}\) FormulaEval selector and compute alpha-charge; if enumeration fails, emit the minimal private-row countermodel.

\paragraph{When to stop.} Stop or quarantine this P vs NP branch only after it states the source object, target class, complexity model, transfer inequality, quantifier strength, and barrier status; quarantine if two sharpened attempts still cannot say whether the route implies \(\mathrm{SAT}\notin \mathrm{P}/\mathrm{poly}\), \(\mathrm{NP}\not\subset \mathrm{P}/\mathrm{poly}\), or only a weaker/proxy lower bound.

\paragraph{Do not stop prematurely if.} Do not stop before the branch has produced a concrete source-to-target transfer theorem, an explicit obstruction lemma, or a barrier-aware reason why the route cannot reach the P vs NP/P/poly target.

\paragraph{How it adds to the global problem.} This move adds to the P vs NP problem because it keeps barrier status and nonuniform/advice leakage attached to the actual transfer theorem. The current branch-control decision is backtrack\_to\_sharpen. Payoff criterion: For the restricted-model branch, compute the compiler envelope from unrestricted SIZE(m\textasciicircum{}k) or P-time algorithms into the named model; promotion requires a bound above that envelope on the same SAT slice.

\subsection*{cycle-0101: Hardness-Magnification Bridge Realization}

\paragraph{Learning.} Cycle 101 upgrades \(bounded-SR_{h}\) from a named cost to a bit-normal restricted parser theorem: SR contributes \(B_{SR}(pi)\le N^{q+\eta+\tau+o(1)}\) to \(\Gamma_{bEF}\). The remaining live obstruction is exactly \(FE_{Extract,N}\) for symbolic FormulaEval rows.

\paragraph{Pressure.} The induction excludes symbolic FormulaEval-SAT; \(FE_{Extract,N}\) is still needed to enumerate RowSet(a) or force the \(N^{1+q+o(1)}\) expansion charge.

\paragraph{Next.} Write the finite verifier table for one \(SR_{h}\) rule: tag layout, row j encoding, node/depth fields, parent pointers, occurrence-neighborhood list, and the exact charge contributing to B\_SR(pi).

\paragraph{When to stop.} Stop or quarantine this P vs NP branch only after it states the source object, target class, complexity model, transfer inequality, quantifier strength, and barrier status; quarantine if two sharpened attempts still cannot say whether the route implies \(\mathrm{SAT}\notin \mathrm{P}/\mathrm{poly}\), \(\mathrm{NP}\not\subset \mathrm{P}/\mathrm{poly}\), or only a weaker/proxy lower bound.

\paragraph{Do not stop prematurely if.} Do not stop before the branch has produced a concrete source-to-target transfer theorem, an explicit obstruction lemma, or a barrier-aware reason why the route cannot reach the P vs NP/P/poly target.

\paragraph{How it adds to the global problem.} This move adds to the P vs NP problem because it keeps barrier status and nonuniform/advice leakage attached to the actual transfer theorem. The current branch-control decision is backtrack\_to\_sharpen. Payoff criterion: For this route, state the exact barrier premise it uses or escapes: relativizing, natural, algebrizing, nonuniform, or advice-leaking; promotion requires the escape to be attached to the specific transfer theorem.

\subsection*{cycle-0102: Hardness-Magnification Bridge Realization - cycle 102 \(SR_{h}\) normal form audit}

\paragraph{Learning.} Cycle 102 upgrades the \(bounded-SR_{h}\) subgoal from an exponent slogan to a parser theorem for the explicit rule schema: \(SR_{h}\) contributes exactly B\_SR(pi) to \(\Gamma_{bEF}\), and the line induction gives a valid restricted \(\mathrm{LB}_{N}\) falsifier. The remaining full-system issue is sharply isolated as \(FE_{Extract,N}\) for symbolic FormulaEval-SAT, not generic naturalness, relativization, or P vs NP itself.

\paragraph{Pressure.} No blocking gap remains for the restricted \(SR_{h}^{bit}\) normal form and line-induction packet. The full \(\mathrm{Psys}_{sym}\) theorem is blocked at \(FE_{Extract,N}\) or \(RowSetEnumerator_{N}\) for symbolic FormulaEval-SAT.

\paragraph{Next.} Instantiate \(RowSetEnumerator_{N}\) for one transparent \(\mathrm{Psys}_{sym}^{sel}\) fragment, preferably interval-union or \(\mathrm{AC0}\) row selectors, and decide whether FormulaEval-SAT is an expanded axiom schema with \(N^{1+q+o(1)}\) charge or a semantic oracle exposed by the private-row falsifier.

\paragraph{When to stop.} Stop or quarantine this P vs NP branch only after it states the source object, target class, complexity model, transfer inequality, quantifier strength, and barrier status; quarantine if two sharpened attempts still cannot say whether the route implies \(\mathrm{SAT}\notin \mathrm{P}/\mathrm{poly}\), \(\mathrm{NP}\not\subset \mathrm{P}/\mathrm{poly}\), or only a weaker/proxy lower bound.

\paragraph{Do not stop prematurely if.} Do not stop before the branch has produced a concrete source-to-target transfer theorem, an explicit obstruction lemma, or a barrier-aware reason why the route cannot reach the P vs NP/P/poly target.

\paragraph{How it adds to the global problem.} This move adds to the P vs NP problem because it keeps barrier status and nonuniform/advice leakage attached to the actual transfer theorem. The current branch-control decision is backtrack\_to\_sharpen. Payoff criterion: For the restricted-model branch, compute the compiler envelope from unrestricted SIZE(m\textasciicircum{}k) or P-time algorithms into the named model; promotion requires a bound above that envelope on the same SAT slice.

\subsection*{cycle-0103: Hardness-Magnification Bridge Realization}

\paragraph{Learning.} Cycle 103 discharges a concrete FormulaEval subcase: transparent interval-union selectors admit a valid \(FE_{Extract,N}\) parser with N|RowSet| expansion charge, so global FormulaEval costs \(N^{1+q+o(1)}\) and yields the q/epsilon support trichotomy. The full \(\mathrm{Psys}_{sym}\) claim remains quarantined exactly at semantic FormulaEval extraction; unexpanded symbolic FormulaEval is refuted by the private-row falsifier, not by generic barrier folklore.

\paragraph{Pressure.} The derivation does not prove \(FE_{Extract,N}\) for arbitrary \(\mathrm{Psys}_{sym}\) selectors or for real FormulaEval-SAT comprehension whose row set is defined by semantic SAT facts rather than displayed interval endpoints.

\paragraph{Next.} Extend \(RowSetEnumerator_{N}\) from interval-union selectors to a named \(\mathrm{AC0}\) selector fragment by full scanning over [M\_N]; compute whether scan overhead \(N^{q+o(1)}\) is dominated by N|RowSet| and identify the smallest \(\mathrm{AC0}-private-row\) countermodel if enumeration depends on semantic advice.

\paragraph{When to stop.} Stop or quarantine this black-box substitution branch only after the source carrier, SAT-query length exponent \(\beta\), interface size exponent q, source-threshold exponent, and target class have been explicitly computed; quarantine if every sharpened interface still has positive \(\beta\) and therefore only excludes a bounded range of SAT circuit exponents rather than NP subset P/poly.

\paragraph{Do not stop prematurely if.} Do not stop before deriving the transfer inequality from \(\mathrm{SAT}\in \mathrm{SIZE}(n^{c})\) to the carrier upper bound, identifying whether \(\beta=0\), \(\beta>0\), or the argument is genuinely nonblackbox, and recording the barrier status of the proposed escape.

\paragraph{How it adds to the global problem.} This move adds to the P vs NP problem because it identifies whether the transfer loses all-exponent strength through a positive target-query exponent \(\beta\); it keeps barrier status and nonuniform/advice leakage attached to the actual transfer theorem. The current branch-control decision is continue\_run\_next. Payoff criterion: For the restricted-model branch, compute the compiler envelope from unrestricted SIZE(m\textasciicircum{}k) or P-time algorithms into the named model; promotion requires a bound above that envelope on the same SAT slice.

\subsection*{cycle-0104: Hardness-Magnification Bridge Realization}

\paragraph{Learning.} Cycle 104 pins the \(bounded-SR_{h}\) normal form to concrete verifier fields rather than an exponent slogan. The restricted \(\mathrm{Psys}_{bEF}^{-}\) plus interval-union FormulaEval theorem packet is now valid as a paper-level parser/falsifier result, while full semantic FormulaEval remains quarantined at \(FE_{Extract,N}\) by a private-row countermodel.

\paragraph{Pressure.} Full \(\mathrm{Psys}_{sym}\) remains blocked exactly at \(FE_{Extract,N}\) for semantic FormulaEval selectors and at arbitrary extension variables whose row supports are not bit-explicit. The restricted \(\mathrm{Psys}_{bEF}^{-}\) plus interval-union FormulaEval packet is valid; the full semantic proof-system parser is not proved.

\paragraph{Next.} Instantiate \(RowSetEnumerator_{N}\) for one richer transparent selector class beyond interval-unions, preferably sorted \(\mathrm{AC0}\) threshold-free selectors with explicit scan charge, and decide whether it expands at N|RowSet|, collapses to an \(N^{q}\) scan bound, or admits a private-row countermodel.

\paragraph{When to stop.} Stop or quarantine this P vs NP branch only after it states the source object, target class, complexity model, transfer inequality, quantifier strength, and barrier status; quarantine if two sharpened attempts still cannot say whether the route implies \(\mathrm{SAT}\notin \mathrm{P}/\mathrm{poly}\), \(\mathrm{NP}\not\subset \mathrm{P}/\mathrm{poly}\), or only a weaker/proxy lower bound.

\paragraph{Do not stop prematurely if.} Do not stop before the branch has produced a concrete source-to-target transfer theorem, an explicit obstruction lemma, or a barrier-aware reason why the route cannot reach the P vs NP/P/poly target.

\paragraph{How it adds to the global problem.} This move adds to the P vs NP problem because it keeps barrier status and nonuniform/advice leakage attached to the actual transfer theorem. The current branch-control decision is backtrack\_to\_sharpen. Payoff criterion: For the restricted-model branch, compute the compiler envelope from unrestricted SIZE(m\textasciicircum{}k) or P-time algorithms into the named model; promotion requires a bound above that envelope on the same SAT slice.

\subsection*{cycle-0105: Hardness-Magnification Bridge Realization}

\paragraph{Learning.} Cycle 105 converts the promoted restricted packet into an auditable inference-quality theorem: \(\mathrm{Psys}_{bEF}^{-}\) plus transparent interval-union FormulaEval has a valid parser-to-circuit falsifier with \(\Gamma_{bEF}^{IU}\), while full semantic FormulaEval is not merely hard or barrier-limited but blocked by the exact \(FE_{Extract,N}^{sem}\) private-row countermodel. The contribution is a paper-level parser/falsifier boundary, not evidence for \(\mathrm{P} \ne \mathrm{NP}\).

\paragraph{Pressure.} No remaining gap for \(\mathrm{Psys}_{bEF}^{-}\) with transparent interval-union FormulaEval. Full \(\mathrm{Psys}_{sym}\) remains blocked at \(FE_{Extract,N}^{sem}\): a semantic FormulaEval rule may compute SAT facts over hidden selected formulas without outputting RowSet indices or occurrence neighborhoods.

\paragraph{Next.} Run the delta-varied interval-union selector audit: instantiate \(|RowSet|=N^{q-\delta}\), verify \(RowSetEnumerator_{N}\) gate cost, compute \(\Gamma_{bEF}^{IU}\) at \(\delta=q-\epsilon\) \(\pm\) o(1), and test whether any proposed semantic selector supplies a valid \(FE_{Extract,N}\) certificate rather than triggering the private-row falsifier.

\paragraph{When to stop.} Stop or quarantine this black-box substitution branch only after the source carrier, SAT-query length exponent \(\beta\), interface size exponent q, source-threshold exponent, and target class have been explicitly computed; quarantine if every sharpened interface still has positive \(\beta\) and therefore only excludes a bounded range of SAT circuit exponents rather than NP subset P/poly.

\paragraph{Do not stop prematurely if.} Do not stop before deriving the transfer inequality from \(\mathrm{SAT}\in \mathrm{SIZE}(n^{c})\) to the carrier upper bound, identifying whether \(\beta=0\), \(\beta>0\), or the argument is genuinely nonblackbox, and recording the barrier status of the proposed escape.

\paragraph{How it adds to the global problem.} This move adds to the P vs NP problem because it identifies whether the transfer loses all-exponent strength through a positive target-query exponent \(\beta\); it keeps barrier status and nonuniform/advice leakage attached to the actual transfer theorem; it searches the remaining plausible escape route: nonblackbox or subpolynomial-interface carriers. The current branch-control decision is continue\_run\_next. Payoff criterion: For this route, state the exact barrier premise it uses or escapes: relativizing, natural, algebrizing, nonuniform, or advice-leaking; promotion requires the escape to be attached to the specific transfer theorem.

\subsection*{cycle-0106: Hardness-Magnification Bridge Realization}

\paragraph{Learning.} Cycle 106 upgrades the promoted restricted result into a live inference-quality theorem packet: \(\Gamma_{bEF}^{IU}\) is derived by line induction with explicit \(SR_{h}\) and interval-union FormulaEval charges, the q/epsilon support trichotomy is a contraposed corollary, and full semantic FormulaEval remains quarantined by the private-row falsifier rather than by generic barrier folklore.

\paragraph{Pressure.} No blocking gap remains for \(\mathrm{Psys}_{bEF}^{-}\) plus transparent interval-union FormulaEval. The named unresolved gap is \(FE_{Extract,N}^{sem}\): extracting complete row supports from unexpanded semantic FormulaEval-SAT without advice or nonrelativizing SAT-specific reasoning.

\paragraph{Next.} Run the finite validation audit on one concrete \(\mathrm{Psys}_{sym}\) rule: encode FormulaEval as (a) interval-union axiom schema, (b) \(\mathrm{AC0}\) selector with enumerator, and (c) opaque semantic oracle; compute \(\alpha\) and either attach \(N^{1+\alpha+o(1)}\) charge or exhibit the private-row countermodel.

\paragraph{When to stop.} Stop or quarantine this black-box substitution branch only after the source carrier, SAT-query length exponent \(\beta\), interface size exponent q, source-threshold exponent, and target class have been explicitly computed; quarantine if every sharpened interface still has positive \(\beta\) and therefore only excludes a bounded range of SAT circuit exponents rather than NP subset P/poly.

\paragraph{Do not stop prematurely if.} Do not stop before deriving the transfer inequality from \(\mathrm{SAT}\in \mathrm{SIZE}(n^{c})\) to the carrier upper bound, identifying whether \(\beta=0\), \(\beta>0\), or the argument is genuinely nonblackbox, and recording the barrier status of the proposed escape.

\paragraph{How it adds to the global problem.} This move adds to the P vs NP problem because it identifies whether the transfer loses all-exponent strength through a positive target-query exponent \(\beta\); it keeps barrier status and nonuniform/advice leakage attached to the actual transfer theorem. The current branch-control decision is continue\_run\_next. Payoff criterion: For the restricted-model branch, compute the compiler envelope from unrestricted SIZE(m\textasciicircum{}k) or P-time algorithms into the named model; promotion requires a bound above that envelope on the same SAT slice.

\subsection*{cycle-0107: Hardness-Magnification Bridge Realization}

\paragraph{Learning.} Cycle 107 adds a real validation upgrade: transparent FormulaEval selectors no longer require interval-union structure for a valid theorem, because full syntactic scanning gives a certified \(B_{FE}=F_{FE}\) \(N^{1+q+o(1)}\) charge. This closes the \(q<\epsilon\) falsifier for a broader selector fragment and isolates the only missing improvement as \(AC0PrefixNonemptiness_{N}\) for output-sensitive enumeration. Full semantic FormulaEval remains quarantined by the private-row falsifier, not by generic barrier folklore.

\paragraph{Pressure.} \(AC0PrefixNonemptiness_{N}\): no advice-free, output-sensitive enumeration theorem is proved for arbitrary \(\mathrm{AC0}\) row selectors over q log N row bits; brute force gives a valid but sometimes over-threshold scan charge.

\paragraph{Next.} Formalize \(AC0PrefixNonemptiness_{N}\) as a finite lemma: for \(\mathrm{AC0}\) selectors on q log N row bits, test whether descriptor structure yields \(N^{o(1)}\) prefix emptiness; if not, state the selector-scan theorem as final and add \(\mathrm{AC0}\) output-sensitive enumeration to open problems.

\paragraph{When to stop.} Stop or quarantine this black-box substitution branch only after the source carrier, SAT-query length exponent \(\beta\), interface size exponent q, source-threshold exponent, and target class have been explicitly computed; quarantine if every sharpened interface still has positive \(\beta\) and therefore only excludes a bounded range of SAT circuit exponents rather than NP subset P/poly.

\paragraph{Do not stop prematurely if.} Do not stop before deriving the transfer inequality from \(\mathrm{SAT}\in \mathrm{SIZE}(n^{c})\) to the carrier upper bound, identifying whether \(\beta=0\), \(\beta>0\), or the argument is genuinely nonblackbox, and recording the barrier status of the proposed escape.

\paragraph{How it adds to the global problem.} This move adds to the P vs NP problem because it identifies whether the transfer loses all-exponent strength through a positive target-query exponent \(\beta\); it separates true hardness magnification from projection or padding arguments that only move proof debt; it keeps barrier status and nonuniform/advice leakage attached to the actual transfer theorem. The current branch-control decision is continue\_run\_next. Payoff criterion: For the restricted-model branch, compute the compiler envelope from unrestricted SIZE(m\textasciicircum{}k) or P-time algorithms into the named model; promotion requires a bound above that envelope on the same SAT slice.

\subsection*{cycle-0108: Hardness-Magnification Bridge Realization}

\paragraph{Learning.} Cycle 108 replaces selector-case accumulation by a single \(TransparentRowSupport_{N}\) theorem pair. Transparent FormulaEval is valid only after \(U_{N}+X_{N}\) charging or full scan; hidden semantic FormulaEval is refuted by a private-row falsifier unless \(RealSATInfluenceRealizer_{N}\) supplies a nonrelativizing charged-support certificate. This is a paper-level parser/falsifier boundary, not evidence for \(\mathrm{P} \ne \mathrm{NP}\).

\paragraph{Pressure.} The positive theorem is valid for \(TransparentRowSupport_{N}\) and full-scan fallback. The full \(\mathrm{Psys}_{sym}\) semantic theorem remains blocked exactly at \(RealSATInfluenceRealizer_{N}\); \(AC0PrefixNonemptiness_{N}\) is only an optimization gap for output-sensitive enumeration.

\paragraph{Next.} Write the referee-style \(TransparentRowSupport_{N}\) definition and audit six examples as corollaries: interval-union, full scan, ROBDD, guarded-dDNNF, projected-dDNNF, and influence-complete extension rules. Do not add new selector families unless they prove a clause, refute a clause, or improve \(AC0PrefixNonemptiness_{N}\).

\paragraph{When to stop.} Stop or quarantine this black-box substitution branch only after the source carrier, SAT-query length exponent \(\beta\), interface size exponent q, source-threshold exponent, and target class have been explicitly computed; quarantine if every sharpened interface still has positive \(\beta\) and therefore only excludes a bounded range of SAT circuit exponents rather than NP subset P/poly.

\paragraph{Do not stop prematurely if.} Do not stop before deriving the transfer inequality from \(\mathrm{SAT}\in \mathrm{SIZE}(n^{c})\) to the carrier upper bound, identifying whether \(\beta=0\), \(\beta>0\), or the argument is genuinely nonblackbox, and recording the barrier status of the proposed escape.

\paragraph{How it adds to the global problem.} This move adds to the P vs NP problem because it identifies whether the transfer loses all-exponent strength through a positive target-query exponent \(\beta\); it keeps barrier status and nonuniform/advice leakage attached to the actual transfer theorem; it searches the remaining plausible escape route: nonblackbox or subpolynomial-interface carriers. The current branch-control decision is continue\_run\_next. Payoff criterion: For the restricted-model branch, compute the compiler envelope from unrestricted SIZE(m\textasciicircum{}k) or P-time algorithms into the named model; promotion requires a bound above that envelope on the same SAT slice.

\subsection*{cycle-0109: Hardness-Magnification Bridge Realization: \(TransparentRowSupport\) theorem pair}

\paragraph{Learning.} Cycle 109 upgrades the theorem pair from selector consolidation to a proof-audited interface: \(TransparentRowSupport_{N}\) is a finite checkable condition, \(\Gamma_{trans}\) is derived by line induction, hidden semantic support has a concrete private-row/unread-enable-bit falsifier, and the only full \(\mathrm{Psys}_{sym}\) escape is the nonrelativizing \(RealSATInfluenceRealizer_{N}\). This is paper-level route certification, not evidence for \(\mathrm{P} \ne \mathrm{NP}\).

\paragraph{Pressure.} \(RealSATInfluenceRealizer_{N}\) remains open for genuine real-SAT semantic FormulaEval rules; the private-row countermodel proves that a relativizing or surface-occurrence parser cannot fill this gap.

\paragraph{Next.} Referee-check one projected-dDNNF FormulaEval rule against the five \(TransparentRowSupport_{N}\) clauses, compute its \(U_{N}\) and \(X_{N}\) charges, and verify that any projected hidden variable is either influence-closed or triggers the unread-enable-bit countermodel.

\paragraph{When to stop.} Stop or quarantine this black-box substitution branch only after the source carrier, SAT-query length exponent \(\beta\), interface size exponent q, source-threshold exponent, and target class have been explicitly computed; quarantine if every sharpened interface still has positive \(\beta\) and therefore only excludes a bounded range of SAT circuit exponents rather than NP subset P/poly.

\paragraph{Do not stop prematurely if.} Do not stop before deriving the transfer inequality from \(\mathrm{SAT}\in \mathrm{SIZE}(n^{c})\) to the carrier upper bound, identifying whether \(\beta=0\), \(\beta>0\), or the argument is genuinely nonblackbox, and recording the barrier status of the proposed escape.

\paragraph{How it adds to the global problem.} This move adds to the P vs NP problem because it identifies whether the transfer loses all-exponent strength through a positive target-query exponent \(\beta\); it keeps barrier status and nonuniform/advice leakage attached to the actual transfer theorem. The current branch-control decision is continue\_run\_next. Payoff criterion: For this route, state the exact barrier premise it uses or escapes: relativizing, natural, algebrizing, nonuniform, or advice-leaking; promotion requires the escape to be attached to the specific transfer theorem.

\subsection*{cycle-0110: Hardness-Magnification Bridge Realization}

\paragraph{Learning.} Cycle 110 consolidates the selector zoo into a single theorem pair. Transparent FormulaEval is valid only when every influence is surfaced and charged in \(\Gamma_{trans}\); hidden semantic support is blocked by a private-row/unread-enable-bit falsifier. The exact full-system escape is \(RealSATInfluenceRealizer_{N}\), not a generic naturalness or relativization slogan.

\paragraph{Pressure.} No remaining gap for the restricted \(TransparentRowSupport_{N}\) theorem pair. The named open gap for full \(\mathrm{Psys}_{sym}\) is \(RealSATInfluenceRealizer_{N}\): convert real SAT-semantic FormulaEval influence into explicit charged row support without advice leakage or threshold collapse.

\paragraph{Next.} Attempt \(RealSATInfluenceRealizer_{N}\) for one influence-complete extension rule: define \(RealInfl_{l}\), produce row-index witnesses and \(EnumCost+N|K_{l}|\) charge, or exhibit a semantic rule where a real SAT-sensitive row remains unchargeable, thereby demoting full \(\mathrm{Psys}_{sym}\) beyond the transparent theorem pair.

\paragraph{When to stop.} Stop or quarantine this black-box substitution branch only after the source carrier, SAT-query length exponent \(\beta\), interface size exponent q, source-threshold exponent, and target class have been explicitly computed; quarantine if every sharpened interface still has positive \(\beta\) and therefore only excludes a bounded range of SAT circuit exponents rather than NP subset P/poly.

\paragraph{Do not stop prematurely if.} Do not stop before deriving the transfer inequality from \(\mathrm{SAT}\in \mathrm{SIZE}(n^{c})\) to the carrier upper bound, identifying whether \(\beta=0\), \(\beta>0\), or the argument is genuinely nonblackbox, and recording the barrier status of the proposed escape.

\paragraph{How it adds to the global problem.} This move adds to the P vs NP problem because it identifies whether the transfer loses all-exponent strength through a positive target-query exponent \(\beta\); it keeps barrier status and nonuniform/advice leakage attached to the actual transfer theorem; it searches the remaining plausible escape route: nonblackbox or subpolynomial-interface carriers. The current branch-control decision is continue\_run\_next. Payoff criterion: For this route, state the exact barrier premise it uses or escapes: relativizing, natural, algebrizing, nonuniform, or advice-leaking; promotion requires the escape to be attached to the specific transfer theorem.

\subsection*{cycle-0111: Hardness-Magnification Bridge Realization}

\paragraph{Learning.} Cycle 111 converts the transparent FormulaEval consolidation into a proof-delivery packet: \(\Gamma_{trans}\) is derived by line induction, hidden support has a precise private-row falsifier, and the only full semantic escape is the named nonrelativizing \(RealSATInfluenceRealizer_{N}\). This is a valid scoped parser/falsifier contribution, not progress on the main P vs NP conjecture itself.

\paragraph{Pressure.} No gap remains for the conditional \(TransparentRowSupport_{N}\) theorem pair. The full \(\mathrm{Psys}_{sym}\) extension is blocked exactly at \(RealSATInfluenceRealizer_{N}\) for hidden semantic FormulaEval lines.

\paragraph{Next.} Write the final paper theorem text for \(TransparentRowSupport_{N}\) with clauses T1-T6, prove the line induction lemma in referee style, then test one genuinely semantic FormulaEval rule against \(RealSATInfluenceRealizer_{N}\) or instantiate the private-row countermodel.

\paragraph{When to stop.} Stop or quarantine this P vs NP branch only after it states the source object, target class, complexity model, transfer inequality, quantifier strength, and barrier status; quarantine if two sharpened attempts still cannot say whether the route implies \(\mathrm{SAT}\notin \mathrm{P}/\mathrm{poly}\), \(\mathrm{NP}\not\subset \mathrm{P}/\mathrm{poly}\), or only a weaker/proxy lower bound.

\paragraph{Do not stop prematurely if.} Do not stop before the branch has produced a concrete source-to-target transfer theorem, an explicit obstruction lemma, or a barrier-aware reason why the route cannot reach the P vs NP/P/poly target.

\paragraph{How it adds to the global problem.} This move adds to the P vs NP problem because it keeps barrier status and nonuniform/advice leakage attached to the actual transfer theorem. The current branch-control decision is backtrack\_to\_sharpen. Payoff criterion: For this route, state the exact barrier premise it uses or escapes: relativizing, natural, algebrizing, nonuniform, or advice-leaking; promotion requires the escape to be attached to the specific transfer theorem.

\subsection*{cycle-0112: Hardness-Magnification Bridge Realization}

\paragraph{Learning.} Cycle 112 turns the transparent FormulaEval packet into a finite theorem pair: transparency is exactly a charged support-realization condition, and hidden semantic support is not merely difficult but has a two-world private-row falsifier. The contribution remains a route-level parser/falsifier boundary for beta-zero magnification; the main P vs NP conjecture is untouched.

\paragraph{Pressure.} No blocking gap remains for \(\mathrm{Psys}_{sym}^{trans}\) as defined. The extension to full semantic \(\mathrm{Psys}_{sym}\) is blocked exactly at \(RealSATInfluenceRealizer_{N}\), which must convert real-SAT semantic influence into charged same-theory row support.

\paragraph{Next.} Write the corollary table as examples only: interval-union, full scan, ROBDD, guarded-dDNNF, projected-dDNNF, and influence-complete extension rules must each be checked against the five \(TransparentRowSupport_{N}\) clauses or demoted.

\paragraph{When to stop.} Stop or quarantine this black-box substitution branch only after the source carrier, SAT-query length exponent \(\beta\), interface size exponent q, source-threshold exponent, and target class have been explicitly computed; quarantine if every sharpened interface still has positive \(\beta\) and therefore only excludes a bounded range of SAT circuit exponents rather than NP subset P/poly.

\paragraph{Do not stop prematurely if.} Do not stop before deriving the transfer inequality from \(\mathrm{SAT}\in \mathrm{SIZE}(n^{c})\) to the carrier upper bound, identifying whether \(\beta=0\), \(\beta>0\), or the argument is genuinely nonblackbox, and recording the barrier status of the proposed escape.

\paragraph{How it adds to the global problem.} This move adds to the P vs NP problem because it identifies whether the transfer loses all-exponent strength through a positive target-query exponent \(\beta\); it keeps barrier status and nonuniform/advice leakage attached to the actual transfer theorem. The current branch-control decision is adapt\_then\_run. Payoff criterion: For the restricted-model branch, compute the compiler envelope from unrestricted SIZE(m\textasciicircum{}k) or P-time algorithms into the named model; promotion requires a bound above that envelope on the same SAT slice.

\subsection*{cycle-0113: Hardness-Magnification Bridge Realization}

\paragraph{Learning.} Cycle 113 converts the promoted FormulaEval boundary into a paper-ready theorem pair: transparent support is exactly a charged support-realization interface yielding \(\Gamma_{trans}\) and the \(\mathrm{LB}_{N}\) falsifier window, while hidden semantic support is blocked by a concrete private-row/unread-enable-bit falsifier. The full-system escape is precisely \(RealSATInfluenceRealizer_{N}\), not generic barrier folklore.

\paragraph{Pressure.} No blocking gap remains for the restricted \(TransparentRowSupport_{N}\) theorem pair. The full semantic \(\mathrm{Psys}_{sym}\) extension is blocked exactly at \(RealSATInfluenceRealizer_{N}\), an advice-free nonrelativizing conversion from real SAT influence to charged occurrence support.

\paragraph{Next.} Run the finite proof-quality checklist on the paper statement: verify each \(TransparentRowSupport_{N}\) clause for interval-union, full scan, ROBDD, guarded-dDNNF, projected-dDNNF, and influence-complete examples; then attempt or refute \(RealSATInfluenceRealizer_{N}\) for one concrete semantic FormulaEval axiom.

\paragraph{When to stop.} Stop or quarantine this black-box substitution branch only after the source carrier, SAT-query length exponent \(\beta\), interface size exponent q, source-threshold exponent, and target class have been explicitly computed; quarantine if every sharpened interface still has positive \(\beta\) and therefore only excludes a bounded range of SAT circuit exponents rather than NP subset P/poly.

\paragraph{Do not stop prematurely if.} Do not stop before deriving the transfer inequality from \(\mathrm{SAT}\in \mathrm{SIZE}(n^{c})\) to the carrier upper bound, identifying whether \(\beta=0\), \(\beta>0\), or the argument is genuinely nonblackbox, and recording the barrier status of the proposed escape.

\paragraph{How it adds to the global problem.} This move adds to the P vs NP problem because it identifies whether the transfer loses all-exponent strength through a positive target-query exponent \(\beta\); it keeps barrier status and nonuniform/advice leakage attached to the actual transfer theorem. The current branch-control decision is adapt\_then\_run. Payoff criterion: For the restricted-model branch, compute the compiler envelope from unrestricted SIZE(m\textasciicircum{}k) or P-time algorithms into the named model; promotion requires a bound above that envelope on the same SAT slice.

\subsection*{cycle-0114: Hardness-Magnification Bridge Realization}

\paragraph{Learning.} Cycle 114 upgrades the promoted theorem pair into an auditable live inference packet: transparent FormulaEval is exactly charged support realization, hidden semantic FormulaEval has a two-world falsifier, and the full-system question is the named nonrelativizing \(RealSATInfluenceRealizer_{N}\) rather than generic barrier folklore.

\paragraph{Pressure.} The positive theorem is complete for \(TransparentRowSupport_{N}\), but full semantic FormulaEval remains blocked at \(RealSATInfluenceRealizer_{N}\): no advice-free nonrelativizing procedure is known that surfaces all real-SAT influences with subthreshold cost.

\paragraph{Next.} Instantiate \(RealSATInfluenceRealizer_{N}\) for one concrete full semantic FormulaEval rule, or prove a rule-specific private-row countermodel showing that the full \(\mathrm{Psys}_{sym}\) extension must remain quarantined.

\paragraph{When to stop.} Stop or quarantine this P vs NP branch only after it states the source object, target class, complexity model, transfer inequality, quantifier strength, and barrier status; quarantine if two sharpened attempts still cannot say whether the route implies \(\mathrm{SAT}\notin \mathrm{P}/\mathrm{poly}\), \(\mathrm{NP}\not\subset \mathrm{P}/\mathrm{poly}\), or only a weaker/proxy lower bound.

\paragraph{Do not stop prematurely if.} Do not stop before the branch has produced a concrete source-to-target transfer theorem, an explicit obstruction lemma, or a barrier-aware reason why the route cannot reach the P vs NP/P/poly target.

\paragraph{How it adds to the global problem.} This move adds to the P vs NP problem because it keeps barrier status and nonuniform/advice leakage attached to the actual transfer theorem. The current branch-control decision is backtrack\_to\_sharpen. Payoff criterion: For this route, state the exact barrier premise it uses or escapes: relativizing, natural, algebrizing, nonuniform, or advice-leaking; promotion requires the escape to be attached to the specific transfer theorem.

\subsection*{cycle-0115: Hardness-Magnification Bridge Realization}

\paragraph{Learning.} Cycle 115 adds a minimality proof obligation to the promoted transparent-support theorem pair. Transparent FormulaEval is not just sufficient for a valid \(\Gamma_{trans}\) parser; hidden influence is exactly what the private-row/unread-enable-bit falsifier exploits. Thus the full-system boundary is a named nonrelativizing \(RealSATInfluenceRealizer_{N}\) lemma, while the paper-level contribution remains a scoped parser/falsifier result rather than evidence for \(\mathrm{P} \ne \mathrm{NP}\).

\paragraph{Pressure.} The transparent theorem and hidden-support countermodel are proof-shaped for occurrence-visible fragments, but the full semantic case is blocked by \(RealSATInfluenceRealizer_{N}\): one must show that real SAT semantics forces every influential row into a charged support set.

\paragraph{Next.} Formalize the private-row independence lemma for one full symbolic FormulaEval rule with an unread enable bit: write the two transcripts, the charged constraint set, the flipped row r, and the proof that all visible lines agree while \(H_{N}^\mathrm{SAT}\) parity differs; if independence fails, record the exact same-theory constraint as a \(RealSATInfluenceRealizer_{N}\) candidate.

\paragraph{When to stop.} Stop or quarantine this black-box substitution branch only after the source carrier, SAT-query length exponent \(\beta\), interface size exponent q, source-threshold exponent, and target class have been explicitly computed; quarantine if every sharpened interface still has positive \(\beta\) and therefore only excludes a bounded range of SAT circuit exponents rather than NP subset P/poly.

\paragraph{Do not stop prematurely if.} Do not stop before deriving the transfer inequality from \(\mathrm{SAT}\in \mathrm{SIZE}(n^{c})\) to the carrier upper bound, identifying whether \(\beta=0\), \(\beta>0\), or the argument is genuinely nonblackbox, and recording the barrier status of the proposed escape.

\paragraph{How it adds to the global problem.} This move adds to the P vs NP problem because it identifies whether the transfer loses all-exponent strength through a positive target-query exponent \(\beta\); it keeps barrier status and nonuniform/advice leakage attached to the actual transfer theorem. The current branch-control decision is continue\_run\_next. Payoff criterion: For this route, state the exact barrier premise it uses or escapes: relativizing, natural, algebrizing, nonuniform, or advice-leaking; promotion requires the escape to be attached to the specific transfer theorem.

\subsection*{cycle-0116: Hardness-Magnification Bridge Realization: Transparent FormulaEval Minimality Audit}

\paragraph{Learning.} Cycle 116 closes the minimality audit for the transparent FormulaEval theorem pair. The valid paper claim is now a restricted parser/falsifier boundary: every surfaced influence is charged in \(\Gamma_{trans}\), and every unsurfaced nonsemantic influence admits a private-row falsifier. The only remaining full-system route is the named nonrelativizing \(RealSATInfluenceRealizer_{N}\); this is not evidence for \(\mathrm{P} \ne \mathrm{NP}\).

\paragraph{Pressure.} Full semantic FormulaEval remains blocked at \(RealSATInfluenceRealizer_{N}\): private-row spoofing is a relativizing falsifier for hidden support, but it cannot refute genuinely real-SAT semantic influence without a same-theory realization or a concrete semantic no-go.

\paragraph{Next.} Instantiate \(RealSATInfluenceRealizer_{N}\) or refute it for one concrete strong rule: FormulaEval with Tseitin-extension equality plus bounded self-reduction; list exact bit fields, RowSet extraction, same-theory expansion certificates, and the first semantic inference that cannot be charged.

\paragraph{When to stop.} Stop or quarantine this P vs NP branch only after it states the source object, target class, complexity model, transfer inequality, quantifier strength, and barrier status; quarantine if two sharpened attempts still cannot say whether the route implies \(\mathrm{SAT}\notin \mathrm{P}/\mathrm{poly}\), \(\mathrm{NP}\not\subset \mathrm{P}/\mathrm{poly}\), or only a weaker/proxy lower bound.

\paragraph{Do not stop prematurely if.} Do not stop before the branch has produced a concrete source-to-target transfer theorem, an explicit obstruction lemma, or a barrier-aware reason why the route cannot reach the P vs NP/P/poly target.

\paragraph{How it adds to the global problem.} This move adds to the P vs NP problem because it keeps barrier status and nonuniform/advice leakage attached to the actual transfer theorem. The current branch-control decision is backtrack\_to\_sharpen. Payoff criterion: For the proof-complexity branch, separate fixed-family proof hardness from language hardness by naming the proof system, shortcut algorithm, and semantic lower-bound target.

\subsection*{cycle-0117: Hardness-Magnification Bridge Realization: Transparent FormulaEval Boundary}

\paragraph{Learning.} Cycle 117 upgrades the promoted theorem pair from sufficiency plus informal minimality to a local safety iff for extensional FormulaEval: surfaced influence is exactly what makes the \(\Gamma_{trans}\) parser valid, and unsurfaced nonsemantic influence is exactly what the private-row/unread-enable-bit falsifier exploits. The full semantic case is not dismissed by a barrier slogan; it is isolated as the nonrelativizing \(RealSATInfluenceRealizer_{N}\) obligation. This is a paper-level route-boundary result, not a proof of \(\mathrm{P} \ne \mathrm{NP}\).

\paragraph{Pressure.} The derivation is complete for extensional FormulaEval rules but still blocked for full semantic FormulaEval at \(RealSATInfluenceRealizer_{N}\), which must nonrelativizingly force real SAT influence into charged same-theory support.

\paragraph{Next.} Write the paper-ready theorem statement as three clauses: \(TransparentRowSupport_{N}\) positive \(\Gamma_{trans}\) induction; extensional hidden-support private-row/unread-enable-bit countermodel; \(RealSATInfluenceRealizer_{N}\) open semantic escape with overhead \(R_{N}\) and explicit non-claims.

\paragraph{When to stop.} Stop or quarantine this black-box substitution branch only after the source carrier, SAT-query length exponent \(\beta\), interface size exponent q, source-threshold exponent, and target class have been explicitly computed; quarantine if every sharpened interface still has positive \(\beta\) and therefore only excludes a bounded range of SAT circuit exponents rather than NP subset P/poly.

\paragraph{Do not stop prematurely if.} Do not stop before deriving the transfer inequality from \(\mathrm{SAT}\in \mathrm{SIZE}(n^{c})\) to the carrier upper bound, identifying whether \(\beta=0\), \(\beta>0\), or the argument is genuinely nonblackbox, and recording the barrier status of the proposed escape.

\paragraph{How it adds to the global problem.} This move adds to the P vs NP problem because it identifies whether the transfer loses all-exponent strength through a positive target-query exponent \(\beta\); it keeps barrier status and nonuniform/advice leakage attached to the actual transfer theorem. The current branch-control decision is adapt\_then\_run. Payoff criterion: For this route, state the exact barrier premise it uses or escapes: relativizing, natural, algebrizing, nonuniform, or advice-leaking; promotion requires the escape to be attached to the specific transfer theorem.

\subsection*{cycle-0118: Hardness-Magnification Bridge Realization}

\paragraph{Learning.} Cycle 118 adds a compositional repair to the FormulaEval boundary: local surfaced influence is not enough unless hidden influences are closed under later proof use. The valid restricted theorem is \(\mathrm{Psys}_{sym}^{ext}\) plus charged \(RealSATInfluenceRealizer_{N}\) rows implies \(\Gamma_{comp}\) line induction and the \(\mathrm{LB}_{N}\) falsifier window. The necessity direction is now a last-unsafe-line private-row theorem. Full semantic FormulaEval remains open exactly at a nonrelativizing realizer lemma, not at generic relativization, natural-proofs, or P vs NP folklore.

\paragraph{Pressure.} The last-unsafe-line argument is valid in the extensional label model. For real SAT semantics, a hidden row may be forced by nonrelativizing formula identities, self-reduction constraints, or EqSAT links without being syntactically visible. That exact gap is \(RealSATInfluenceRealizer_{N}\).

\paragraph{Next.} Run a finite proof audit on \(\mathrm{Psys}_{sym}^{ext}\): instantiate one guarded FormulaEval line, compute \(I_{l}\), \(K_{idx}\), \(R_{FE}\), and \(\Gamma_{comp}\), then verify the last-unsafe-line two-world construction or identify the first real SAT-semantic rule that forces a nonrelativizing realizer charge.

\paragraph{When to stop.} Stop or quarantine this P vs NP branch only after it states the source object, target class, complexity model, transfer inequality, quantifier strength, and barrier status; quarantine if two sharpened attempts still cannot say whether the route implies \(\mathrm{SAT}\notin \mathrm{P}/\mathrm{poly}\), \(\mathrm{NP}\not\subset \mathrm{P}/\mathrm{poly}\), or only a weaker/proxy lower bound.

\paragraph{Do not stop prematurely if.} Do not stop before the branch has produced a concrete source-to-target transfer theorem, an explicit obstruction lemma, or a barrier-aware reason why the route cannot reach the P vs NP/P/poly target.

\paragraph{How it adds to the global problem.} This move adds to the P vs NP problem because it keeps barrier status and nonuniform/advice leakage attached to the actual transfer theorem. The current branch-control decision is backtrack\_to\_sharpen. Payoff criterion: For this route, state the exact barrier premise it uses or escapes: relativizing, natural, algebrizing, nonuniform, or advice-leaking; promotion requires the escape to be attached to the specific transfer theorem.

\subsection*{cycle-0119: Hardness-Magnification Bridge Realization, cycle 119: compositional transparent support closure}

\paragraph{Learning.} Cycle 119 turns compositional closure from a warning into a theorem-shaped invariant: transparent FormulaEval support must be prefix-closed under later symbolic use, and subthreshold \(\Gamma_{comp}\) then gives a same-model \(\mathrm{LB}_{N}\) falsifier. The necessity direction is a last-unsafe-line private-row theorem for extensional support. Full semantic FormulaEval remains open exactly at \(RealSATInfluenceRealizer_{N}\), not at generic barrier folklore or the main P vs NP conjecture.

\paragraph{Pressure.} The derivation does not prove full \(\mathrm{Psys}_{sym}\) semantic FormulaEval parsing; it is blocked exactly at \(RealSATInfluenceRealizer_{N}\) for rules whose influence is forced by real SAT semantics but not surfaced as generated-row support.

\paragraph{Next.} Write the finite prefix-closure verifier for \(TransparentRowSupport_{N}\): input pi,N; output \(Cl_{t}\) for each line, Gamma\_prefix(t), and either a valid support-closed derivation certificate or the first unsafe line with its private-row/unread-enable-bit countermodel.

\paragraph{When to stop.} Stop or quarantine this P vs NP branch only after it states the source object, target class, complexity model, transfer inequality, quantifier strength, and barrier status; quarantine if two sharpened attempts still cannot say whether the route implies \(\mathrm{SAT}\notin \mathrm{P}/\mathrm{poly}\), \(\mathrm{NP}\not\subset \mathrm{P}/\mathrm{poly}\), or only a weaker/proxy lower bound.

\paragraph{Do not stop prematurely if.} Do not stop before the branch has produced a concrete source-to-target transfer theorem, an explicit obstruction lemma, or a barrier-aware reason why the route cannot reach the P vs NP/P/poly target.

\paragraph{How it adds to the global problem.} This move adds to the P vs NP problem because it keeps barrier status and nonuniform/advice leakage attached to the actual transfer theorem. The current branch-control decision is backtrack\_to\_sharpen. Payoff criterion: For this route, state the exact barrier premise it uses or escapes: relativizing, natural, algebrizing, nonuniform, or advice-leaking; promotion requires the escape to be attached to the specific transfer theorem.

\subsection*{cycle-0120: Hardness-Magnification Bridge Realization: least-closed transparent FormulaEval support}

\paragraph{Learning.} Cycle 120 strengthens the promoted theorem pair by identifying the minimal valid charged object: the least prefix-closed reachable FormulaEval support cone K\_dag(pi), not the full selector zoo. The positive line induction uses Gamma\_dag=Gamma\_bEF+U(K\_dag)+X(K\_dag), and the necessity side is a first-hidden-reachable-row falsifier. Full semantic FormulaEval remains quarantined exactly at \(RealSATInfluenceRealizer_{N}\); this is a scoped parser/falsifier contribution, not progress toward proving \(\mathrm{P} \ne \mathrm{NP}\).

\paragraph{Pressure.} Full semantic FormulaEval over real SAT is still blocked at \(RealSATInfluenceRealizer_{N}\): no proof currently converts semantic influence of an uncharged formula row into an advice-free, prefix-closed, charged support certificate.

\paragraph{Next.} Write the proof-audit version of L120.1: formalize the fixed-point computation of K\_dag(pi), give pseudocode for the verifier-visible dependency closure, and check one projected-dDNNF FormulaEval example plus one unread-enable-bit counterexample against the \(\Gamma_{dag}\) inequality.

\paragraph{When to stop.} Stop or quarantine this P vs NP branch only after it states the source object, target class, complexity model, transfer inequality, quantifier strength, and barrier status; quarantine if two sharpened attempts still cannot say whether the route implies \(\mathrm{SAT}\notin \mathrm{P}/\mathrm{poly}\), \(\mathrm{NP}\not\subset \mathrm{P}/\mathrm{poly}\), or only a weaker/proxy lower bound.

\paragraph{Do not stop prematurely if.} Do not stop before the branch has produced a concrete source-to-target transfer theorem, an explicit obstruction lemma, or a barrier-aware reason why the route cannot reach the P vs NP/P/poly target.

\paragraph{How it adds to the global problem.} This move adds to the P vs NP problem because it keeps barrier status and nonuniform/advice leakage attached to the actual transfer theorem. The current branch-control decision is backtrack\_to\_sharpen. Payoff criterion: For this route, state the exact barrier premise it uses or escapes: relativizing, natural, algebrizing, nonuniform, or advice-leaking; promotion requires the escape to be attached to the specific transfer theorem.

\subsection*{cycle-0121: Hardness-Magnification Bridge Realization: L121 least-closed support audit}

\paragraph{Learning.} Cycle 121 turns K\_dag(pi) into a finite auditable fixed-point object. The valid contribution is sharper than the previous transparency slogan: least prefix-closed reachable support is sufficient for the \(\Gamma_{dag}\) parser and necessary for extensional safety via a first-hidden-row falsifier. The remaining full-system gap is exactly \(RealSATInfluenceRealizer_{N}\) for real SAT-semantic FormulaEval, not generic relativization, natural proofs, or P vs NP folklore.

\paragraph{Pressure.} The restricted proof is complete only for transparent extensional FormulaEval. Full semantic SAT FormulaEval remains blocked at \(RealSATInfluenceRealizer_{N}\): proof syntax must nonrelativizingly determine and charge every real SAT influence before use.

\paragraph{Next.} Write the formal L121.1 proof skeleton as theorem text: fixed-point pseudocode for \(K_{dag}\), induction invariant per rule type, charge table for \(bEF/SR_{h}/\mathrm{FormulaEval}/projected-dDNNF\), and the unread-enable-bit counterexample with \(\theta<\epsilon\).

\paragraph{When to stop.} Stop or quarantine this black-box substitution branch only after the source carrier, SAT-query length exponent \(\beta\), interface size exponent q, source-threshold exponent, and target class have been explicitly computed; quarantine if every sharpened interface still has positive \(\beta\) and therefore only excludes a bounded range of SAT circuit exponents rather than NP subset P/poly.

\paragraph{Do not stop prematurely if.} Do not stop before deriving the transfer inequality from \(\mathrm{SAT}\in \mathrm{SIZE}(n^{c})\) to the carrier upper bound, identifying whether \(\beta=0\), \(\beta>0\), or the argument is genuinely nonblackbox, and recording the barrier status of the proposed escape.

\paragraph{How it adds to the global problem.} This move adds to the P vs NP problem because it identifies whether the transfer loses all-exponent strength through a positive target-query exponent \(\beta\); it keeps barrier status and nonuniform/advice leakage attached to the actual transfer theorem; it searches the remaining plausible escape route: nonblackbox or subpolynomial-interface carriers. The current branch-control decision is continue\_run\_next. Payoff criterion: For this route, state the exact barrier premise it uses or escapes: relativizing, natural, algebrizing, nonuniform, or advice-leaking; promotion requires the escape to be attached to the specific transfer theorem.

\subsection*{cycle-0122: Hardness-Magnification Bridge Realization}

\paragraph{Learning.} Cycle 122 upgrades \(K_{dag}\) from a named minimal support object to a referee-auditable fixed-point proof packet. The valid theorem is restricted but paper-relevant: least prefix-closed reachable FormulaEval support is sufficient for \(\Gamma_{dag}\) line induction and necessary for extensional safety by a first-hidden-row falsifier. The full semantic case is still blocked exactly at \(RealSATInfluenceRealizer_{N}\), not by generic barrier folklore or by the P vs NP conjecture itself.

\paragraph{Pressure.} The restricted theorem is valid for transparent extensional FormulaEval, but it does not settle full semantic FormulaEval. The named gap is \(RealSATInfluenceRealizer_{N}\): real SAT-semantic reasoning may determine row influence without a verifier-visible row set, and this requires a nonrelativizing charged realizer or a rule-specific no-go.

\paragraph{Next.} Write the paper-ready L122 module: formal theorem statement, fixed-point pseudocode, line-induction proof, projected-dDNNF charge example, unread-enable-bit counterexample, q-trichotomy corollary, and a separate \(RealSATInfluenceRealizer_{N}\) open-boundary paragraph.

\paragraph{When to stop.} Stop or quarantine this black-box substitution branch only after the source carrier, SAT-query length exponent \(\beta\), interface size exponent q, source-threshold exponent, and target class have been explicitly computed; quarantine if every sharpened interface still has positive \(\beta\) and therefore only excludes a bounded range of SAT circuit exponents rather than NP subset P/poly.

\paragraph{Do not stop prematurely if.} Do not stop before deriving the transfer inequality from \(\mathrm{SAT}\in \mathrm{SIZE}(n^{c})\) to the carrier upper bound, identifying whether \(\beta=0\), \(\beta>0\), or the argument is genuinely nonblackbox, and recording the barrier status of the proposed escape.

\paragraph{How it adds to the global problem.} This move adds to the P vs NP problem because it identifies whether the transfer loses all-exponent strength through a positive target-query exponent \(\beta\); it keeps barrier status and nonuniform/advice leakage attached to the actual transfer theorem; it searches the remaining plausible escape route: nonblackbox or subpolynomial-interface carriers. The current branch-control decision is continue\_run\_next. Payoff criterion: For this route, state the exact barrier premise it uses or escapes: relativizing, natural, algebrizing, nonuniform, or advice-leaking; promotion requires the escape to be attached to the specific transfer theorem.

\subsection*{cycle-0123: Hardness-Magnification Bridge Realization}

\paragraph{Learning.} Cycle 123 closes a proof-audit gap rather than repeating the magnification slogan: \(K_{dag}\) is an implementable least fixed point, projected-dDNNF gives a concrete transparent positive case, and the unread-enable-bit/private-row construction gives a precise necessity falsifier. The valid result is a scoped parser/falsifier theorem for beta-zero parity magnification; full semantic FormulaEval is still blocked exactly at \(RealSATInfluenceRealizer_{N}\).

\paragraph{Pressure.} \(RealSATInfluenceRealizer_{N}\) remains unproved for full semantic FormulaEval: a real SAT-semantic rule might force r* nonrelativizingly from charged rows even when no extensional support certificate lists it.

\paragraph{Next.} Write the referee-checkable trace for the projected-dDNNF example: list gate fields, parent pointers, projection maps, RowSet certificates, \(K_{0},K_{1}\),... closure additions, \(U_{N}\) and \(X_{N}\) charges, then flip one unread enable bit to exhibit the first-hidden-row countermodel and verify the \(\Gamma_{dag}\) threshold numerically for qepsilon.

\paragraph{When to stop.} Stop or quarantine this black-box substitution branch only after the source carrier, SAT-query length exponent \(\beta\), interface size exponent q, source-threshold exponent, and target class have been explicitly computed; quarantine if every sharpened interface still has positive \(\beta\) and therefore only excludes a bounded range of SAT circuit exponents rather than NP subset P/poly.

\paragraph{Do not stop prematurely if.} Do not stop before deriving the transfer inequality from \(\mathrm{SAT}\in \mathrm{SIZE}(n^{c})\) to the carrier upper bound, identifying whether \(\beta=0\), \(\beta>0\), or the argument is genuinely nonblackbox, and recording the barrier status of the proposed escape.

\paragraph{How it adds to the global problem.} This move adds to the P vs NP problem because it identifies whether the transfer loses all-exponent strength through a positive target-query exponent \(\beta\); it separates true hardness magnification from projection or padding arguments that only move proof debt; it keeps barrier status and nonuniform/advice leakage attached to the actual transfer theorem. The current branch-control decision is continue\_run\_next. Payoff criterion: For this route, state the exact barrier premise it uses or escapes: relativizing, natural, algebrizing, nonuniform, or advice-leaking; promotion requires the escape to be attached to the specific transfer theorem.

\subsection*{cycle-0124: Hardness-Magnification Bridge Realization}

\paragraph{Learning.} Cycle 124 makes the semantic boundary sharper: bounded Tseitin self-reduction is foldable into the least-support theorem only when every restriction child is declared and charged. The exact threshold is \(\kappa+\theta<\epsilon\) for the new row expansion. Hidden branch-pruning is a concrete private-row falsifier, so full semantic FormulaEval remains blocked at \(RealSATInfluenceRealizer_{N}^{arb}\) rather than at generic barrier folklore.

\paragraph{Pressure.} The proof does not cover arbitrary full semantic FormulaEval. The named remaining gap is \(RealSATInfluenceRealizer_{N}^{arb}\): every semantic inference must uniformly reveal all rows whose labels can affect later proof use before the inference is consumed.

\paragraph{Next.} Write the \(FE_{\mathrm{TSR},decl}\) rule schema as bit fields: row id, Tseitin code, pivot list, child constructor, parent pointers, and EF derivation handles; then check the fixed-point \(K_\mathrm{TSR}\) enumerator on a one-line bounded-SR example and a hidden satisfying-branch counterexample.

\paragraph{When to stop.} Stop or quarantine this black-box substitution branch only after the source carrier, SAT-query length exponent \(\beta\), interface size exponent q, source-threshold exponent, and target class have been explicitly computed; quarantine if every sharpened interface still has positive \(\beta\) and therefore only excludes a bounded range of SAT circuit exponents rather than NP subset P/poly.

\paragraph{Do not stop prematurely if.} Do not stop before deriving the transfer inequality from \(\mathrm{SAT}\in \mathrm{SIZE}(n^{c})\) to the carrier upper bound, identifying whether \(\beta=0\), \(\beta>0\), or the argument is genuinely nonblackbox, and recording the barrier status of the proposed escape.

\paragraph{How it adds to the global problem.} This move adds to the P vs NP problem because it identifies whether the transfer loses all-exponent strength through a positive target-query exponent \(\beta\); it keeps barrier status and nonuniform/advice leakage attached to the actual transfer theorem. The current branch-control decision is continue\_run\_next. Payoff criterion: For the proof-complexity branch, separate fixed-family proof hardness from language hardness by naming the proof system, shortcut algorithm, and semantic lower-bound target.

\subsection*{cycle-0125: Hardness-Magnification Bridge Realization}

\paragraph{Learning.} Cycle 125 converts the TSR semantic boundary into a theorem-shaped split. Declared bounded TSR is validly absorbed by the least-support parser with charge \(N^{1+\kappa+\theta+o(1)}\), giving a falsifier exactly when \(\kappa+\theta<\epsilon\). Hidden semantic branch-pruning is not a free extension of FormulaEval: it has a first-hidden-row private-row no-go unless a nonrelativizing \(RealSATInfluenceRealizer_{N}^{arb}\) supplies charged influence rows. This is a scoped parser/falsifier contribution, not evidence for \(\mathrm{P} \ne \mathrm{NP}\).

\paragraph{Pressure.} \(RealSATInfluenceRealizer_{N}^{arb}\) for hidden semantic branch-pruning: the derivation does not extract unlisted rows whose relevance is certified only by actual SAT semantics.

\paragraph{Next.} Write the referee-auditable \(\mathrm{TSR}_{decl}\) proof block: child-row bit fields, fixed-point pseudocode, \(\Gamma_\mathrm{TSR}\) induction case, and a two-world hidden-pruning countermodel with a named first hidden child row r\_star.

\paragraph{When to stop.} Stop or quarantine this black-box substitution branch only after the source carrier, SAT-query length exponent \(\beta\), interface size exponent q, source-threshold exponent, and target class have been explicitly computed; quarantine if every sharpened interface still has positive \(\beta\) and therefore only excludes a bounded range of SAT circuit exponents rather than NP subset P/poly.

\paragraph{Do not stop prematurely if.} Do not stop before deriving the transfer inequality from \(\mathrm{SAT}\in \mathrm{SIZE}(n^{c})\) to the carrier upper bound, identifying whether \(\beta=0\), \(\beta>0\), or the argument is genuinely nonblackbox, and recording the barrier status of the proposed escape.

\paragraph{How it adds to the global problem.} This move adds to the P vs NP problem because it identifies whether the transfer loses all-exponent strength through a positive target-query exponent \(\beta\); it tests whether a meta-complexity carrier gives a weaker source obligation than the original P vs NP/P-poly target; it keeps barrier status and nonuniform/advice leakage attached to the actual transfer theorem. The current branch-control decision is continue\_run\_next. Payoff criterion: For this route, state the exact barrier premise it uses or escapes: relativizing, natural, algebrizing, nonuniform, or advice-leaking; promotion requires the escape to be attached to the specific transfer theorem.

\subsection*{cycle-0126: Hardness-Magnification Bridge Realization}

\paragraph{Learning.} Cycle 126 converts declared bounded TSR from a semantic warning into a charged fixed-point rule case. The valid contribution is a restricted theorem: declared TSR rows are absorbed by \(K_{dag}\) with threshold \(\kappa+\theta<\epsilon\), while hidden branch-pruning is a precise private-row falsifier. Full semantic FormulaEval is still not proved; it is isolated at \(RealSATInfluenceRealizer_{N}^{arb}\).

\paragraph{Pressure.} The derivation does not cover arbitrary semantic branch-pruning where real SAT semantics decide which children matter without a syntactic declaration. The named missing lemma is \(RealSATInfluenceRealizer_{N}^{arb}\): from proof syntax and N alone, uniformly list all semantically influencing TSR/FormulaEval rows before use and charge them in \(\Gamma\).

\paragraph{Next.} Formalize the \(\mathrm{Psys}_{sym}^{dag+declTSR(h)}\) grammar in one page: exact row-index fields, TSR parent/child encodings, local verifier predicate, \(K_{dag}^\mathrm{TSR}\) fixed-point pseudocode, and the declared TSR line-induction case; then instantiate with \(\kappa=\epsilon/3\), \(\theta=\epsilon/3\), q in ((1+epsilon)/2,1+epsilon), \(\rho<1+\epsilon\) to verify all charges remain subthreshold except the intended q-density pressure.

\paragraph{When to stop.} Stop or quarantine this P vs NP branch only after it states the source object, target class, complexity model, transfer inequality, quantifier strength, and barrier status; quarantine if two sharpened attempts still cannot say whether the route implies \(\mathrm{SAT}\notin \mathrm{P}/\mathrm{poly}\), \(\mathrm{NP}\not\subset \mathrm{P}/\mathrm{poly}\), or only a weaker/proxy lower bound.

\paragraph{Do not stop prematurely if.} Do not stop before the branch has produced a concrete source-to-target transfer theorem, an explicit obstruction lemma, or a barrier-aware reason why the route cannot reach the P vs NP/P/poly target.

\paragraph{How it adds to the global problem.} This move adds to the P vs NP problem because it keeps barrier status and nonuniform/advice leakage attached to the actual transfer theorem. The current branch-control decision is backtrack\_to\_sharpen. Payoff criterion: For this route, state the exact barrier premise it uses or escapes: relativizing, natural, algebrizing, nonuniform, or advice-leaking; promotion requires the escape to be attached to the specific transfer theorem.

\subsection*{cycle-0127: Hardness-Magnification Bridge Realization}

\paragraph{Learning.} Cycle 127 adds a concrete rule-level proof packet: declared bounded TSR is absorbed by the least closed support cone with exact charge \(N^{1+\kappa+\theta+o(1)}\), while hidden branch pruning has a first-unlisted-child falsifier. This improves the paper theorem boundary without claiming \(\mathrm{P} \ne \mathrm{NP}\) or a general \(H_{N}\) lower bound.

\paragraph{Pressure.} The positive derivation is complete only for declared TSR syntax. Full real-SAT semantic branch pruning remains blocked at \(RealSATInfluenceRealizer_{N}^{arb}\): a uniform, advice-free, nonrelativizing procedure must list all live child rows and charge them before the semantic rule is used.

\paragraph{Next.} Write the verifier bit fields for one declared TSR line: root row id, restriction-node path, pivot variable, child row ids, Tseitin map hash, occurrence-neighborhood list, parent pointers, and expansion certificate; then prove the single-line induction inequality adding exactly N times the declared child count to \(\Gamma_{dag+\mathrm{TSR}}\).

\paragraph{When to stop.} Stop or quarantine this black-box substitution branch only after the source carrier, SAT-query length exponent \(\beta\), interface size exponent q, source-threshold exponent, and target class have been explicitly computed; quarantine if every sharpened interface still has positive \(\beta\) and therefore only excludes a bounded range of SAT circuit exponents rather than NP subset P/poly.

\paragraph{Do not stop prematurely if.} Do not stop before deriving the transfer inequality from \(\mathrm{SAT}\in \mathrm{SIZE}(n^{c})\) to the carrier upper bound, identifying whether \(\beta=0\), \(\beta>0\), or the argument is genuinely nonblackbox, and recording the barrier status of the proposed escape.

\paragraph{How it adds to the global problem.} This move adds to the P vs NP problem because it identifies whether the transfer loses all-exponent strength through a positive target-query exponent \(\beta\); it keeps barrier status and nonuniform/advice leakage attached to the actual transfer theorem. The current branch-control decision is adapt\_then\_run. Payoff criterion: For the proof-complexity branch, separate fixed-family proof hardness from language hardness by naming the proof system, shortcut algorithm, and semantic lower-bound target.

\subsection*{cycle-0128: Hardness-Magnification Bridge Realization}

\paragraph{Learning.} Cycle 128 closes a sharper declared-TSR support sublemma: complete child declarations compose with the least support fixed point and cost exactly \(N^{1+\kappa+\theta+o(1)}\), while hidden branch pruning has a first-unlisted-child private-row falsifier. The contribution remains a scoped parser/falsifier theorem, not evidence for \(\mathrm{P} \ne \mathrm{NP}\).

\paragraph{Pressure.} The declared complete-child schema has no current derivation gap beyond writing the final syntax checker. The full semantic rule remains blocked exactly at \(RealSATInfluenceRealizer_{N}^{arb}\), which must list all influencing rows without SAT-truth advice.

\paragraph{Next.} Write the paper-ready TSR rule schema table and verify on one toy projected-dDNNF plus TSR line: compute \(K_\mathrm{TSR}\), \(A_\mathrm{TSR}\), \(B_\mathrm{TSR}\), \(R_\mathrm{TSR}\), and then flip the first unlisted child in the hidden-pruning variant.

\paragraph{When to stop.} Stop or quarantine this black-box substitution branch only after the source carrier, SAT-query length exponent \(\beta\), interface size exponent q, source-threshold exponent, and target class have been explicitly computed; quarantine if every sharpened interface still has positive \(\beta\) and therefore only excludes a bounded range of SAT circuit exponents rather than NP subset P/poly.

\paragraph{Do not stop prematurely if.} Do not stop before deriving the transfer inequality from \(\mathrm{SAT}\in \mathrm{SIZE}(n^{c})\) to the carrier upper bound, identifying whether \(\beta=0\), \(\beta>0\), or the argument is genuinely nonblackbox, and recording the barrier status of the proposed escape.

\paragraph{How it adds to the global problem.} This move adds to the P vs NP problem because it identifies whether the transfer loses all-exponent strength through a positive target-query exponent \(\beta\); it keeps barrier status and nonuniform/advice leakage attached to the actual transfer theorem. The current branch-control decision is continue\_run\_next. Payoff criterion: For this route, state the exact barrier premise it uses or escapes: relativizing, natural, algebrizing, nonuniform, or advice-leaking; promotion requires the escape to be attached to the specific transfer theorem.

\subsection*{cycle-0129: Hardness-Magnification Bridge Realization}

\paragraph{Learning.} Cycle 129 turns declared bounded TSR from a semantic-warning slogan into a line-induction rule case for the least support theorem. Complete child declaration adds exactly \(N^{1+\kappa+\theta+o(1)}\) to \(\Gamma_\mathrm{DTSR}\) and is valid below \(\kappa+\theta<\epsilon\); hidden branch pruning is refuted by a first-unlisted-child falsifier unless a nonrelativizing real-SAT influence realizer is supplied.

\paragraph{Pressure.} No blocking gap for the restricted complete-child declared TSR rule case. The full semantic FormulaEval/TSR system remains blocked exactly at \(RealSATInfluenceRealizer_{N}^{arb}\), which must uniformly list and charge semantically influencing rows before use.

\paragraph{Next.} Write the paper-ready L129 sublemma proof: give pseudocode for \(K_\mathrm{DTSR}\) fixed-point enumeration, a table of declared TSR bit fields, the local verifier for child occurrence neighborhoods, and the two-world first-unlisted-child countermodel; then run the live inference-quality packet validation.

\paragraph{When to stop.} Stop or quarantine this black-box substitution branch only after the source carrier, SAT-query length exponent \(\beta\), interface size exponent q, source-threshold exponent, and target class have been explicitly computed; quarantine if every sharpened interface still has positive \(\beta\) and therefore only excludes a bounded range of SAT circuit exponents rather than NP subset P/poly.

\paragraph{Do not stop prematurely if.} Do not stop before deriving the transfer inequality from \(\mathrm{SAT}\in \mathrm{SIZE}(n^{c})\) to the carrier upper bound, identifying whether \(\beta=0\), \(\beta>0\), or the argument is genuinely nonblackbox, and recording the barrier status of the proposed escape.

\paragraph{How it adds to the global problem.} This move adds to the P vs NP problem because it identifies whether the transfer loses all-exponent strength through a positive target-query exponent \(\beta\); it separates true hardness magnification from projection or padding arguments that only move proof debt; it keeps barrier status and nonuniform/advice leakage attached to the actual transfer theorem. The current branch-control decision is continue\_run\_next. Payoff criterion: For this route, state the exact barrier premise it uses or escapes: relativizing, natural, algebrizing, nonuniform, or advice-leaking; promotion requires the escape to be attached to the specific transfer theorem.

\subsection*{cycle-0130: Hardness-Magnification Bridge Realization}

\paragraph{Learning.} Cycle 130 closes a real proof-delivery subcase: declared bounded TSR is not merely compatible with transparent support but is an explicit fixed-point rule with charge \(N^{1+\kappa+\theta+o(1)}\). The exact boundary is now complete declaration versus first-unlisted-child falsifier; full semantic pruning remains open only through the nonrelativizing \(RealSATInfluenceRealizer_{N}^{arb}\).

\paragraph{Pressure.} Full semantic branch pruning remains blocked at \(RealSATInfluenceRealizer_{N}^{arb}\): the derivation does not show that real SAT semantics can uniformly list all influential pruned children without advice and before use.

\paragraph{Next.} Run the bit-field audit for one concrete DTSR verifier line: list parent row id, active descriptor id, restriction mask, child row id, Tseitin extension map, local identity tag, and downstream use pointer; verify that \(K_\mathrm{DTSR}\) reaches every used child and compute the exact theta+kappa ledger for the example.

\paragraph{When to stop.} Stop or quarantine this black-box substitution branch only after the source carrier, SAT-query length exponent \(\beta\), interface size exponent q, source-threshold exponent, and target class have been explicitly computed; quarantine if every sharpened interface still has positive \(\beta\) and therefore only excludes a bounded range of SAT circuit exponents rather than NP subset P/poly.

\paragraph{Do not stop prematurely if.} Do not stop before deriving the transfer inequality from \(\mathrm{SAT}\in \mathrm{SIZE}(n^{c})\) to the carrier upper bound, identifying whether \(\beta=0\), \(\beta>0\), or the argument is genuinely nonblackbox, and recording the barrier status of the proposed escape.

\paragraph{How it adds to the global problem.} This move adds to the P vs NP problem because it identifies whether the transfer loses all-exponent strength through a positive target-query exponent \(\beta\); it keeps barrier status and nonuniform/advice leakage attached to the actual transfer theorem. The current branch-control decision is continue\_run\_next. Payoff criterion: For the proof-complexity branch, separate fixed-family proof hardness from language hardness by naming the proof system, shortcut algorithm, and semantic lower-bound target.

\subsection*{cycle-0131: Hardness-Magnification Bridge Realization}

\paragraph{Learning.} Cycle 131 turns declared bounded TSR into an auditable fixed-point rule case: complete child declarations are sufficient for the \(\Gamma_\mathrm{DTSR}\) parser and necessary up to the first-unlisted-child falsifier. The valid contribution is a route-level theorem/countermodel boundary for beta-zero parity magnification certificates, not progress on the main P vs NP conjecture.

\paragraph{Pressure.} The restricted declared-TSR theorem has no current internal gap. The full semantic pruning case remains blocked exactly at \(RealSATInfluenceRealizer_{N}^{arb}\): an advice-free nonrelativizing procedure must list all real-SAT influencing rows and charge them before use.

\paragraph{Next.} Run the reviewer audit on L131.1: instantiate one toy projected-dDNNF FormulaEval line plus one declared TSR parent with two restriction children, compute \(K_\mathrm{DTSR}\) by the worklist, verify \(\Gamma_\mathrm{DTSR}\) numerically as \(N^{1+\kappa+\theta+o(1)}\), then flip the first unlisted child in the hidden-pruning variant to validate the private-row falsifier.

\paragraph{When to stop.} Stop or quarantine this black-box substitution branch only after the source carrier, SAT-query length exponent \(\beta\), interface size exponent q, source-threshold exponent, and target class have been explicitly computed; quarantine if every sharpened interface still has positive \(\beta\) and therefore only excludes a bounded range of SAT circuit exponents rather than NP subset P/poly.

\paragraph{Do not stop prematurely if.} Do not stop before deriving the transfer inequality from \(\mathrm{SAT}\in \mathrm{SIZE}(n^{c})\) to the carrier upper bound, identifying whether \(\beta=0\), \(\beta>0\), or the argument is genuinely nonblackbox, and recording the barrier status of the proposed escape.

\paragraph{How it adds to the global problem.} This move adds to the P vs NP problem because it identifies whether the transfer loses all-exponent strength through a positive target-query exponent \(\beta\); it separates true hardness magnification from projection or padding arguments that only move proof debt; it keeps barrier status and nonuniform/advice leakage attached to the actual transfer theorem. The current branch-control decision is continue\_run\_next. Payoff criterion: For this route, state the exact barrier premise it uses or escapes: relativizing, natural, algebrizing, nonuniform, or advice-leaking; promotion requires the escape to be attached to the specific transfer theorem.

\subsection*{cycle-0132: Hardness-Magnification Bridge Realization}

\paragraph{Learning.} Cycle 132 upgrades declared bounded TSR from a rule-case slogan to a fixed-point parser theorem with explicit bit fields, closure algorithm, ledger, and necessity falsifier. The valid contribution is the restricted \(TransparentRowSupport_{N}^{\mathrm{DTSR}}\) theorem/countermodel pair; full semantic pruning remains quarantined exactly at \(RealSATInfluenceRealizer_{N}^{arb}\).

\paragraph{Pressure.} \(RealSATInfluenceRealizer_{N}^{arb}\): arbitrary semantic branch-pruning may be valid over real SAT yet not expose an advice-free row list in the proof syntax. The present theorem covers declared DTSR and transparent support only.

\paragraph{Next.} Write the paper-ready L132 proof block: pseudocode for K\_DTSR(pi), a table of DTSR bit fields and verifier checks, the five line-induction cases, and the first-unlisted-child countermodel; then test one projected-dDNNF+DTSR positive example and one unread-enable-bit hidden-pruning counterexample.

\paragraph{When to stop.} Stop or quarantine this black-box substitution branch only after the source carrier, SAT-query length exponent \(\beta\), interface size exponent q, source-threshold exponent, and target class have been explicitly computed; quarantine if every sharpened interface still has positive \(\beta\) and therefore only excludes a bounded range of SAT circuit exponents rather than NP subset P/poly.

\paragraph{Do not stop prematurely if.} Do not stop before deriving the transfer inequality from \(\mathrm{SAT}\in \mathrm{SIZE}(n^{c})\) to the carrier upper bound, identifying whether \(\beta=0\), \(\beta>0\), or the argument is genuinely nonblackbox, and recording the barrier status of the proposed escape.

\paragraph{How it adds to the global problem.} This move adds to the P vs NP problem because it identifies whether the transfer loses all-exponent strength through a positive target-query exponent \(\beta\); it keeps barrier status and nonuniform/advice leakage attached to the actual transfer theorem. The current branch-control decision is continue\_run\_next. Payoff criterion: For this route, state the exact barrier premise it uses or escapes: relativizing, natural, algebrizing, nonuniform, or advice-leaking; promotion requires the escape to be attached to the specific transfer theorem.

\subsection*{cycle-0133: Hardness-Magnification Bridge Realization}

\paragraph{Learning.} Cycle 133 upgrades declared bounded TSR from a compatible rule case to a live inference-quality theorem packet: complete child declarations are sufficient for the \(\Gamma_\mathrm{DTSR}\) parser and necessary up to a first-unlisted-child falsifier. The valid contribution is a scoped \(TransparentRowSupport_{N}^\mathrm{DTSR}\) theorem/countermodel boundary; full semantic pruning remains exactly blocked at \(RealSATInfluenceRealizer_{N}^{arb}\).

\paragraph{Pressure.} Full semantic branch pruning remains blocked at \(RealSATInfluenceRealizer_{N}^{arb}\): no proof here uniformly lists all real-SAT-influencing rows when the branch choice is determined by semantic satisfiability rather than declared syntax.

\paragraph{Next.} Attack \(RealSATInfluenceRealizer_{N}^{arb}\) for one concrete semantic pruning rule: bounded Tseitin self-reduction with a SAT-truth branch-elimination premise. Either produce a uniform charged list of all influencing children or formalize the first semantic line where the private-row countermodel survives.

\paragraph{When to stop.} Stop or quarantine this black-box substitution branch only after the source carrier, SAT-query length exponent \(\beta\), interface size exponent q, source-threshold exponent, and target class have been explicitly computed; quarantine if every sharpened interface still has positive \(\beta\) and therefore only excludes a bounded range of SAT circuit exponents rather than NP subset P/poly.

\paragraph{Do not stop prematurely if.} Do not stop before deriving the transfer inequality from \(\mathrm{SAT}\in \mathrm{SIZE}(n^{c})\) to the carrier upper bound, identifying whether \(\beta=0\), \(\beta>0\), or the argument is genuinely nonblackbox, and recording the barrier status of the proposed escape.

\paragraph{How it adds to the global problem.} This move adds to the P vs NP problem because it identifies whether the transfer loses all-exponent strength through a positive target-query exponent \(\beta\); it keeps barrier status and nonuniform/advice leakage attached to the actual transfer theorem. The current branch-control decision is continue\_run\_next. Payoff criterion: For the proof-complexity branch, separate fixed-family proof hardness from language hardness by naming the proof system, shortcut algorithm, and semantic lower-bound target.

\subsection*{cycle-0134: Hardness-Magnification Bridge Realization}

\paragraph{Learning.} Cycle 134 sharpens the DTSR boundary: complete declaration must be counterfactual, not merely actual-branch. This yields a valid restricted \(\Gamma_\mathrm{CTSR}\) parser theorem and a sharper private-row falsifier for hidden adaptive pruning. The remaining full semantic escape is exactly a nonrelativizing \(RealSATInfluenceRealizer_{N}^{arb}\), not generic barrier folklore or evidence for \(\mathrm{P} \ne \mathrm{NP}\).

\paragraph{Pressure.} The restricted counterfactual DTSR theorem is proof-shaped, but full semantic adaptive pruning remains blocked at \(RealSATInfluenceRealizer_{N}^{arb}\): one must prove that real SAT semantics uniformly determines and charges every hidden influential row, or accept the private-row no-go.

\paragraph{Next.} Instantiate the counterfactual DTSR closure on a two-level nested TSR proof with one Tseitin seed family; compute the exact child count, \(K_\mathrm{CTSR}\) fixed point, \(\Gamma_\mathrm{CTSR}\) ledger, and first-unlisted-child countermodel transcript.

\paragraph{When to stop.} Stop or quarantine this black-box substitution branch only after the source carrier, SAT-query length exponent \(\beta\), interface size exponent q, source-threshold exponent, and target class have been explicitly computed; quarantine if every sharpened interface still has positive \(\beta\) and therefore only excludes a bounded range of SAT circuit exponents rather than NP subset P/poly.

\paragraph{Do not stop prematurely if.} Do not stop before deriving the transfer inequality from \(\mathrm{SAT}\in \mathrm{SIZE}(n^{c})\) to the carrier upper bound, identifying whether \(\beta=0\), \(\beta>0\), or the argument is genuinely nonblackbox, and recording the barrier status of the proposed escape.

\paragraph{How it adds to the global problem.} This move adds to the P vs NP problem because it identifies whether the transfer loses all-exponent strength through a positive target-query exponent \(\beta\); it keeps barrier status and nonuniform/advice leakage attached to the actual transfer theorem. The current branch-control decision is continue\_run\_next. Payoff criterion: For the proof-complexity branch, separate fixed-family proof hardness from language hardness by naming the proof system, shortcut algorithm, and semantic lower-bound target.

\subsection*{cycle-0135: Hardness-Magnification Bridge Realization}

\paragraph{Learning.} Cycle 135 upgrades declared TSR to counterfactual declared TSR: complete counterfactual child support is sufficient for the \(\Gamma_\mathrm{CTSR}\) parser and necessary for extensional safety. The exact new falsifier is first-hidden-child adaptive pruning; the only full semantic escape is a charged nonrelativizing \(RealSATInfluenceRealizer_{N}^{arb}\). This is a route-level theorem/countermodel boundary, not evidence for \(\mathrm{P} \ne \mathrm{NP}\).

\paragraph{Pressure.} The two-world private-row argument is a no-go for extensional or relativizing hidden support, not a proof against real SAT-specific semantic pruning. Full semantic pruning remains blocked exactly at \(RealSATInfluenceRealizer_{N}^{arb}\).

\paragraph{Next.} Resolve \(RealSATInfluenceRealizer_{N}^{arb}\) for one concrete semantic rule: bounded Tseitin self-reduction with adaptive unsat-branch pruning. Either produce a syntax-and-N computable influence row list \(R_{A}\) with charge in \(\Gamma_{real}\), or formalize the first unchargeable semantic inference as a private-row no-realizer theorem.

\paragraph{When to stop.} Stop or quarantine this P vs NP branch only after it states the source object, target class, complexity model, transfer inequality, quantifier strength, and barrier status; quarantine if two sharpened attempts still cannot say whether the route implies \(\mathrm{SAT}\notin \mathrm{P}/\mathrm{poly}\), \(\mathrm{NP}\not\subset \mathrm{P}/\mathrm{poly}\), or only a weaker/proxy lower bound.

\paragraph{Do not stop prematurely if.} Do not stop before the branch has produced a concrete source-to-target transfer theorem, an explicit obstruction lemma, or a barrier-aware reason why the route cannot reach the P vs NP/P/poly target.

\paragraph{How it adds to the global problem.} This move adds to the P vs NP problem because it keeps barrier status and nonuniform/advice leakage attached to the actual transfer theorem. The current branch-control decision is backtrack\_to\_sharpen. Payoff criterion: For the proof-complexity branch, separate fixed-family proof hardness from language hardness by naming the proof system, shortcut algorithm, and semantic lower-bound target.

\subsection*{cycle-0136: Hardness-Magnification Bridge Realization}

\paragraph{Learning.} Cycle 136 turns counterfactual declared TSR from a boundary slogan into a bitfield-level closed-support theorem: K\_CTSR(pi) is a finite fixed point, \(\Gamma_\mathrm{CTSR}\) is derived by line induction, and first-hidden-child adaptive pruning is the exact falsifier for extensional safety. The valid contribution remains a scoped parser/falsifier theorem for beta-zero parity magnification certificates; it does not prove \(H_{N}\) lower bounds or \(\mathrm{P} \ne \mathrm{NP}\).

\paragraph{Pressure.} The derivation does not handle arbitrary semantic branch pruning. The named open bridge is \(RealSATInfluenceRealizer_{N}^{arb}\): from real SAT semantics and proof syntax alone, uniformly list every influencing child row with charged overhead before the semantic rule is used.

\paragraph{Next.} Select one concrete full semantic pruning rule beyond CTSR declarations, preferably adaptive Tseitin branch pruning with real SAT contradiction tests, and attempt either a charged \(RealSATInfluenceRealizer_{N}^{arb}\) or a minimal first-unchargeable-row countermodel.

\paragraph{When to stop.} Stop or quarantine this black-box substitution branch only after the source carrier, SAT-query length exponent \(\beta\), interface size exponent q, source-threshold exponent, and target class have been explicitly computed; quarantine if every sharpened interface still has positive \(\beta\) and therefore only excludes a bounded range of SAT circuit exponents rather than NP subset P/poly.

\paragraph{Do not stop prematurely if.} Do not stop before deriving the transfer inequality from \(\mathrm{SAT}\in \mathrm{SIZE}(n^{c})\) to the carrier upper bound, identifying whether \(\beta=0\), \(\beta>0\), or the argument is genuinely nonblackbox, and recording the barrier status of the proposed escape.

\paragraph{How it adds to the global problem.} This move adds to the P vs NP problem because it identifies whether the transfer loses all-exponent strength through a positive target-query exponent \(\beta\); it keeps barrier status and nonuniform/advice leakage attached to the actual transfer theorem. The current branch-control decision is continue\_run\_next. Payoff criterion: For the proof-complexity branch, separate fixed-family proof hardness from language hardness by naming the proof system, shortcut algorithm, and semantic lower-bound target.

\subsection*{cycle-0137: Hardness-Magnification Bridge Realization}

\paragraph{Learning.} Cycle 137 improves the promoted CTSR packet by proving minimality pressure, not just sufficiency: complete counterfactual child support is the least valid transparent object for adaptive TSR. Removing one reachable child gives a concrete first-hidden-child falsifier; adding all children gives the exact \(\Gamma_\mathrm{CTSR}\) threshold \(\kappa+\theta<\epsilon\). The learning is a scoped parser/countermodel contribution, not a lower bound for \(H_{N}\) and not \(\mathrm{P} \ne \mathrm{NP}\) progress.

\paragraph{Pressure.} The derivation does not handle full semantic branch pruning where real SAT itself proves the omitted child determined by charged rows. That is exactly \(RealSATInfluenceRealizer_{N}^{arb}\): an advice-free, nonrelativizing realizer must list and charge all influencing rows before the semantic rule is used.

\paragraph{Next.} Write the standalone L137 first-hidden-child theorem: specify the CTSR bit fields, define the two row-oracle worlds, prove transcript equality on \(K_\mathrm{CTSR}\) minus r, prove final parity disagreement, and compute the \(\Gamma_\mathrm{CTSR}\) threshold for \(\kappa+\theta<\epsilon\) with q-trichotomy support density.

\paragraph{When to stop.} Stop or quarantine this black-box substitution branch only after the source carrier, SAT-query length exponent \(\beta\), interface size exponent q, source-threshold exponent, and target class have been explicitly computed; quarantine if every sharpened interface still has positive \(\beta\) and therefore only excludes a bounded range of SAT circuit exponents rather than NP subset P/poly.

\paragraph{Do not stop prematurely if.} Do not stop before deriving the transfer inequality from \(\mathrm{SAT}\in \mathrm{SIZE}(n^{c})\) to the carrier upper bound, identifying whether \(\beta=0\), \(\beta>0\), or the argument is genuinely nonblackbox, and recording the barrier status of the proposed escape.

\paragraph{How it adds to the global problem.} This move adds to the P vs NP problem because it identifies whether the transfer loses all-exponent strength through a positive target-query exponent \(\beta\); it keeps barrier status and nonuniform/advice leakage attached to the actual transfer theorem. The current branch-control decision is continue\_run\_next. Payoff criterion: For this route, state the exact barrier premise it uses or escapes: relativizing, natural, algebrizing, nonuniform, or advice-leaking; promotion requires the escape to be attached to the specific transfer theorem.

\subsection*{cycle-0138: Hardness-Magnification Bridge Realization}

\paragraph{Learning.} Cycle 138 identifies the exact minimal CTSR object as reachable counterfactual frontier closure. This improves the promoted packet: actual-branch declarations are not valid transparent support, while full reachable child declarations give the \(\Gamma_\mathrm{CTSR}\) parser with threshold \(\kappa+\theta<\epsilon\). The result is a scoped parser/countermodel theorem and live inference-quality packet, not evidence for \(\mathrm{P} \ne \mathrm{NP}\).

\paragraph{Pressure.} The two-world necessity proof is valid for extensional row-label or occurrence-local semantics, but it does not refute a real SAT semantic derivation proving that the hidden child is impossible unless \(RealSATInfluenceRealizer_{N}^{arb}\) fails for that derivation.

\paragraph{Next.} Write the referee-facing CTSRFrontierCompleteness proof: give closure pseudocode for \(K_\mathrm{CTSR}\), prove termination and minimality, verify one adaptive TSR line bit-by-bit, and instantiate the first-hidden-child falsifier on a one-node unread-enable-bit example.

\paragraph{When to stop.} Stop or quarantine this black-box substitution branch only after the source carrier, SAT-query length exponent \(\beta\), interface size exponent q, source-threshold exponent, and target class have been explicitly computed; quarantine if every sharpened interface still has positive \(\beta\) and therefore only excludes a bounded range of SAT circuit exponents rather than NP subset P/poly.

\paragraph{Do not stop prematurely if.} Do not stop before deriving the transfer inequality from \(\mathrm{SAT}\in \mathrm{SIZE}(n^{c})\) to the carrier upper bound, identifying whether \(\beta=0\), \(\beta>0\), or the argument is genuinely nonblackbox, and recording the barrier status of the proposed escape.

\paragraph{How it adds to the global problem.} This move adds to the P vs NP problem because it identifies whether the transfer loses all-exponent strength through a positive target-query exponent \(\beta\); it separates true hardness magnification from projection or padding arguments that only move proof debt; it keeps barrier status and nonuniform/advice leakage attached to the actual transfer theorem. The current branch-control decision is continue\_run\_next. Payoff criterion: For this route, state the exact barrier premise it uses or escapes: relativizing, natural, algebrizing, nonuniform, or advice-leaking; promotion requires the escape to be attached to the specific transfer theorem.

\subsection*{cycle-0139: Hardness-Magnification Bridge Realization}

\paragraph{Learning.} Cycle 139 improves the CTSR frontier packet by adding a safe EqSAT quotient rule: duplicate counterfactual child rows reduce \(\Gamma\) only when a same-theory certificate is surfaced and charged. Otherwise quotienting is exactly where the private-row falsifier attacks. This is a valid scoped parser/countermodel contribution and still not evidence for \(\mathrm{P} \ne \mathrm{NP}\).

\paragraph{Pressure.} \(RealSATInfluenceRealizer_{N}^{arb}\) remains open for fully semantic adaptive pruning that never exposes counterfactual children or EqSAT certificate edges yet may be sound for real SAT instances.

\paragraph{Next.} Instantiate \(\Gamma_\mathrm{CTSR}^=\) on a projected-dDNNF FormulaEval proof with two Tseitin-equivalent child formulas: compute K before and after certified EqSAT quotienting, then remove one certificate and exhibit the first uncertified quotient-collision private-row falsifier.

\paragraph{When to stop.} Stop or quarantine this black-box substitution branch only after the source carrier, SAT-query length exponent \(\beta\), interface size exponent q, source-threshold exponent, and target class have been explicitly computed; quarantine if every sharpened interface still has positive \(\beta\) and therefore only excludes a bounded range of SAT circuit exponents rather than NP subset P/poly.

\paragraph{Do not stop prematurely if.} Do not stop before deriving the transfer inequality from \(\mathrm{SAT}\in \mathrm{SIZE}(n^{c})\) to the carrier upper bound, identifying whether \(\beta=0\), \(\beta>0\), or the argument is genuinely nonblackbox, and recording the barrier status of the proposed escape.

\paragraph{How it adds to the global problem.} This move adds to the P vs NP problem because it identifies whether the transfer loses all-exponent strength through a positive target-query exponent \(\beta\); it keeps barrier status and nonuniform/advice leakage attached to the actual transfer theorem. The current branch-control decision is continue\_run\_next. Payoff criterion: For the proof-complexity branch, separate fixed-family proof hardness from language hardness by naming the proof system, shortcut algorithm, and semantic lower-bound target.

\subsection*{cycle-0140: Hardness-Magnification Bridge Realization: QCTSR Certified-Quotient Parser Boundary}

\paragraph{Learning.} Cycle 140 turns safe EqSAT quotienting from a slogan into a ledger transformation: duplicate counterfactual CTSR rows may be removed only through a surfaced same-theory certificate forest whose dependencies are in \(K_{Q}\). The exact new falsifier is the first uncertified quotient merge. This strengthens the scoped parser/countermodel contribution and remains separate from proving any lower bound or \(\mathrm{P} \ne \mathrm{NP}\).

\paragraph{Pressure.} The derivation does not handle arbitrary real-SAT semantic pruning. It requires either surfaced same-theory EqSAT certificates or a future \(RealSATInfluenceRealizer_{N}^{arb}\) that uniformly lists and charges all semantic influences before use.

\paragraph{Next.} Instantiate \(K_{Q}\) on a toy projected-dDNNF FormulaEval proof with two CTSR levels and three EqSAT quotient classes; compute |RepIdx|, \(C_{Eq}\), \(\Gamma_\mathrm{QCTSR}\), then delete one quotient certificate to exhibit the first-hidden-merge falsifier.

\paragraph{When to stop.} Stop or quarantine this black-box substitution branch only after the source carrier, SAT-query length exponent \(\beta\), interface size exponent q, source-threshold exponent, and target class have been explicitly computed; quarantine if every sharpened interface still has positive \(\beta\) and therefore only excludes a bounded range of SAT circuit exponents rather than NP subset P/poly.

\paragraph{Do not stop prematurely if.} Do not stop before deriving the transfer inequality from \(\mathrm{SAT}\in \mathrm{SIZE}(n^{c})\) to the carrier upper bound, identifying whether \(\beta=0\), \(\beta>0\), or the argument is genuinely nonblackbox, and recording the barrier status of the proposed escape.

\paragraph{How it adds to the global problem.} This move adds to the P vs NP problem because it identifies whether the transfer loses all-exponent strength through a positive target-query exponent \(\beta\); it keeps barrier status and nonuniform/advice leakage attached to the actual transfer theorem. The current branch-control decision is continue\_run\_next. Payoff criterion: For this route, state the exact barrier premise it uses or escapes: relativizing, natural, algebrizing, nonuniform, or advice-leaking; promotion requires the escape to be attached to the specific transfer theorem.

\subsection*{cycle-0141: Hardness-Magnification Bridge Realization}

\paragraph{Learning.} The new contribution is that quotient compression in counterfactual TSR is not a free \(\Gamma\) reduction. It is valid only as a dependency-closed, acyclic same-theory certificate transformation; otherwise the first uncertified merge gives a private-row falsifier. This closes a proof-audit loophole in the restricted transparent-support theorem while preserving the separation between route-level certificate collapse and the main P vs NP conjecture.

\paragraph{Pressure.} The derivation is complete only for syntactic same-theory quotient certificates. It blocks for full semantic FormulaEval or adaptive pruning where the set of influencing rows is determined by real SAT truth rather than proof-visible dependencies; this is exactly \(RealSATInfluenceRealizer_{N}^{arb}\).

\paragraph{Next.} Instantiate one two-merge CTSR proof toy model: one certified EqSAT quotient with charged forest and one uncertified merge. Write the local verifier fields, compute \(R_{Q}\) and \(C_{Q}\), and demonstrate transcript agreement plus parity disagreement for the uncertified merge.

\paragraph{When to stop.} Stop or quarantine this P vs NP branch only after it states the source object, target class, complexity model, transfer inequality, quantifier strength, and barrier status; quarantine if two sharpened attempts still cannot say whether the route implies \(\mathrm{SAT}\notin \mathrm{P}/\mathrm{poly}\), \(\mathrm{NP}\not\subset \mathrm{P}/\mathrm{poly}\), or only a weaker/proxy lower bound.

\paragraph{Do not stop prematurely if.} Do not stop before the branch has produced a concrete source-to-target transfer theorem, an explicit obstruction lemma, or a barrier-aware reason why the route cannot reach the P vs NP/P/poly target.

\paragraph{How it adds to the global problem.} This move adds to the P vs NP problem because it keeps barrier status and nonuniform/advice leakage attached to the actual transfer theorem. The current branch-control decision is backtrack\_to\_sharpen. Payoff criterion: For this route, state the exact barrier premise it uses or escapes: relativizing, natural, algebrizing, nonuniform, or advice-leaking; promotion requires the escape to be attached to the specific transfer theorem.

\subsection*{cycle-0142: Hardness-Magnification Bridge Realization}

\paragraph{Learning.} Cycle 142 makes quotient compression proof-carrying: row deletion is valid only after a ranked same-theory certificate is charged in \(\Gamma_\mathrm{QCTSR}\). Acyclicity is the new finite audit condition; cyclic, adaptive, or advice-chosen merges are concrete countermodels, not generic barrier folklore.

\paragraph{Pressure.} The argument does not realize hidden semantic FormulaEval or adaptive SAT branch pruning; \(RealSATInfluenceRealizer_{N}^{arb}\) remains the named nonrelativizing open lemma.

\paragraph{Next.} Write the finite verifier for \(K_{Q}\) and ranked \(F_{Q}\) on a toy projected-dDNNF plus CTSR example with one duplicate child pair; compute \(\Gamma\) savings, then remove one certificate edge to exhibit the private-row falsifier.

\paragraph{When to stop.} Stop or quarantine this black-box substitution branch only after the source carrier, SAT-query length exponent \(\beta\), interface size exponent q, source-threshold exponent, and target class have been explicitly computed; quarantine if every sharpened interface still has positive \(\beta\) and therefore only excludes a bounded range of SAT circuit exponents rather than NP subset P/poly.

\paragraph{Do not stop prematurely if.} Do not stop before deriving the transfer inequality from \(\mathrm{SAT}\in \mathrm{SIZE}(n^{c})\) to the carrier upper bound, identifying whether \(\beta=0\), \(\beta>0\), or the argument is genuinely nonblackbox, and recording the barrier status of the proposed escape.

\paragraph{How it adds to the global problem.} This move adds to the P vs NP problem because it identifies whether the transfer loses all-exponent strength through a positive target-query exponent \(\beta\); it keeps barrier status and nonuniform/advice leakage attached to the actual transfer theorem. The current branch-control decision is continue\_run\_next. Payoff criterion: For this route, state the exact barrier premise it uses or escapes: relativizing, natural, algebrizing, nonuniform, or advice-leaking; promotion requires the escape to be attached to the specific transfer theorem.

\subsection*{cycle-0143: Hardness-Magnification Bridge Realization}

\paragraph{Learning.} Cycle 143 turns quotient compression into a proof-carrying ledger rule: \(\Gamma\) may shrink only through ranked acyclic same-theory certificates already in \(K_{Q}\). Unranked quotienting is not a harmless optimization; it is exactly a private-row falsifier.

\paragraph{Pressure.} This does not solve \(RealSATInfluenceRealizer_{N}^{arb}\) for full semantic FormulaEval; it only handles proof-carrying same-theory quotient compression.

\paragraph{Next.} Instantiate the six-row QCTSR transcript, list row IDs, quotient certificate dependencies, rank labels, and show the first uncertified merge produces two agreeing visible transcripts with opposite \(H_{N}\) parity.

\paragraph{When to stop.} Stop or quarantine this P vs NP branch only after it states the source object, target class, complexity model, transfer inequality, quantifier strength, and barrier status; quarantine if two sharpened attempts still cannot say whether the route implies \(\mathrm{SAT}\notin \mathrm{P}/\mathrm{poly}\), \(\mathrm{NP}\not\subset \mathrm{P}/\mathrm{poly}\), or only a weaker/proxy lower bound.

\paragraph{Do not stop prematurely if.} Do not stop before the branch has produced a concrete source-to-target transfer theorem, an explicit obstruction lemma, or a barrier-aware reason why the route cannot reach the P vs NP/P/poly target.

\paragraph{How it adds to the global problem.} This move adds to the P vs NP problem because it keeps barrier status and nonuniform/advice leakage attached to the actual transfer theorem. The current branch-control decision is backtrack\_to\_sharpen. Payoff criterion: For this route, state the exact barrier premise it uses or escapes: relativizing, natural, algebrizing, nonuniform, or advice-leaking; promotion requires the escape to be attached to the specific transfer theorem.

\subsection*{cycle-0144: Hardness-Magnification Bridge Realization: ranked QCTSR ledger validation}

\paragraph{Learning.} Cycle 144 closes a proof-audit loophole: quotient compression in counterfactual TSR is valid exactly as proof-carrying ranked same-theory compression, not as semantic row deletion. The contribution is a scoped parser/falsifier theorem: ranked QCTSR preserves the \(\Gamma\) line induction, while the first unranked quotient merge is a concrete private-row falsifier. This improves the paper theorem boundary without giving any \(H_{N}\) lower bound or \(\mathrm{P} \ne \mathrm{NP}\) result.

\paragraph{Pressure.} The positive proof is complete only for restricted \(\mathrm{Psys}_\mathrm{QCTSR}\). Full semantic FormulaEval or hidden adaptive pruning still requires \(RealSATInfluenceRealizer_{N}^{arb}\), an advice-free nonrelativizing theorem that lists and charges all real SAT influence rows before use.

\paragraph{Next.} Write the referee-checkable \(K_{Q}\) closure pseudocode with rank ordering, then test it on a two-child Tseitin-equivalent TSR merge and on the same example with the certificate omitted to exhibit the first-unranked-merge falsifier.

\paragraph{When to stop.} Stop or quarantine this black-box substitution branch only after the source carrier, SAT-query length exponent \(\beta\), interface size exponent q, source-threshold exponent, and target class have been explicitly computed; quarantine if every sharpened interface still has positive \(\beta\) and therefore only excludes a bounded range of SAT circuit exponents rather than NP subset P/poly.

\paragraph{Do not stop prematurely if.} Do not stop before deriving the transfer inequality from \(\mathrm{SAT}\in \mathrm{SIZE}(n^{c})\) to the carrier upper bound, identifying whether \(\beta=0\), \(\beta>0\), or the argument is genuinely nonblackbox, and recording the barrier status of the proposed escape.

\paragraph{How it adds to the global problem.} This move adds to the P vs NP problem because it identifies whether the transfer loses all-exponent strength through a positive target-query exponent \(\beta\); it keeps barrier status and nonuniform/advice leakage attached to the actual transfer theorem; it searches the remaining plausible escape route: nonblackbox or subpolynomial-interface carriers. The current branch-control decision is continue\_run\_next. Payoff criterion: For the proof-complexity branch, separate fixed-family proof hardness from language hardness by naming the proof system, shortcut algorithm, and semantic lower-bound target.

\subsection*{cycle-0145: Hardness-Magnification Bridge Realization}

\paragraph{Learning.} Cycle 145 upgrades ranked QCTSR from a line-induction slogan to a finite verifier/minimality packet. Quotient compression is valid only when proof-carrying same-theory certificates are ranked, dependency-closed, and charged before use. The exact obstruction is the first rejected quotient merge, which yields a private-row falsifier. This strengthens the paper parser/countermodel boundary without proving \(H_{N}\) lower bounds or \(\mathrm{P} \ne \mathrm{NP}\).

\paragraph{Pressure.} The derivation is complete for surface QCTSR syntax with a specified EqSAT/CTSR same-theory checker. It does not handle arbitrary semantic FormulaEval or hidden adaptive pruning unless \(RealSATInfluenceRealizer_{N}^{arb}\) uniformly lists and charges the same influenced rows before use.

\paragraph{Next.} Instantiate the finite verifier on one toy QCTSR proof with two EqSAT-equivalent counterfactual TSR children: compute \(K_{Q}\), ranks, \(B_{Q}\), and \(\Gamma_\mathrm{QCTSR}\); then delete one certificate edge and explicitly write the two-world private-row falsifier.

\paragraph{When to stop.} Stop or quarantine this black-box substitution branch only after the source carrier, SAT-query length exponent \(\beta\), interface size exponent q, source-threshold exponent, and target class have been explicitly computed; quarantine if every sharpened interface still has positive \(\beta\) and therefore only excludes a bounded range of SAT circuit exponents rather than NP subset P/poly.

\paragraph{Do not stop prematurely if.} Do not stop before deriving the transfer inequality from \(\mathrm{SAT}\in \mathrm{SIZE}(n^{c})\) to the carrier upper bound, identifying whether \(\beta=0\), \(\beta>0\), or the argument is genuinely nonblackbox, and recording the barrier status of the proposed escape.

\paragraph{How it adds to the global problem.} This move adds to the P vs NP problem because it identifies whether the transfer loses all-exponent strength through a positive target-query exponent \(\beta\); it keeps barrier status and nonuniform/advice leakage attached to the actual transfer theorem. The current branch-control decision is continue\_run\_next. Payoff criterion: For this route, state the exact barrier premise it uses or escapes: relativizing, natural, algebrizing, nonuniform, or advice-leaking; promotion requires the escape to be attached to the specific transfer theorem.

\subsection*{cycle-0146: Hardness-Magnification Bridge Realization}

\paragraph{Learning.} Cycle 146 upgrades ranked QCTSR from a finite local verifier to a composable fixed-point line-induction theorem: quotient compression is valid only when proof-carrying same-theory certificates are ranked, acyclic, dependency-closed, and charged before use. The first rejected merge is a concrete private-row falsifier. This strengthens the transparent-support parser/countermodel paper module while preserving all nonclaims: no \(H_{N}\) lower bound, no \(\mathrm{P} \ne \mathrm{NP}\) result, and no free hidden semantic FormulaEval.

\paragraph{Pressure.} The derivation does not fold arbitrary hidden semantic FormulaEval or adaptive SAT branch pruning into \(K_{Q}\). That remains the named \(RealSATInfluenceRealizer_{N}^{arb}\) problem.

\paragraph{Next.} Run the finite verifier on one projected-dDNNF/Tseitin-renaming example: list \(K_{Q}\) seeds, accepted quotient certificates, rank labels, computed \(\Gamma_\mathrm{QCTSR}\), and then delete one certificate to exhibit the first-rejected-merge private-row falsifier.

\paragraph{When to stop.} Stop or quarantine this black-box substitution branch only after the source carrier, SAT-query length exponent \(\beta\), interface size exponent q, source-threshold exponent, and target class have been explicitly computed; quarantine if every sharpened interface still has positive \(\beta\) and therefore only excludes a bounded range of SAT circuit exponents rather than NP subset P/poly.

\paragraph{Do not stop prematurely if.} Do not stop before deriving the transfer inequality from \(\mathrm{SAT}\in \mathrm{SIZE}(n^{c})\) to the carrier upper bound, identifying whether \(\beta=0\), \(\beta>0\), or the argument is genuinely nonblackbox, and recording the barrier status of the proposed escape.

\paragraph{How it adds to the global problem.} This move adds to the P vs NP problem because it identifies whether the transfer loses all-exponent strength through a positive target-query exponent \(\beta\); it keeps barrier status and nonuniform/advice leakage attached to the actual transfer theorem. The current branch-control decision is continue\_run\_next. Payoff criterion: For the proof-complexity branch, separate fixed-family proof hardness from language hardness by naming the proof system, shortcut algorithm, and semantic lower-bound target.

\subsection*{cycle-0147: Hardness-Magnification Bridge Realization}

\paragraph{Learning.} Cycle 147 closes a real proof-audit upgrade: ranked QCTSR quotient compression has a least, order-independent, proof-carrying support closure, so \(\Gamma\) reduction is valid only after lower-rank same-theory dependencies are charged. The sharp falsifier is the first nondecreasing, cyclic, advice-chosen, or hidden semantic merge. This strengthens the paper parser/countermodel module without proving any \(H_{N}\) lower bound or \(\mathrm{P} \ne \mathrm{NP}\).

\paragraph{Pressure.} No blocking gap remains for restricted ranked QCTSR syntax. The broader full-semantic FormulaEval system remains blocked exactly at \(RealSATInfluenceRealizer_{N}^{arb}\), which would have to surface all real SAT-semantic influences with charges before use.

\paragraph{Next.} Instantiate the closure verifier on a three-row QCTSR toy transcript: two counterfactual child rows, one EqSAT certificate forest, and one attempted cyclic merge; compute \(K_{Q}\), \(\Gamma_\mathrm{QCTSR}\), and the first-rejected-merge private-row falsifier explicitly.

\paragraph{When to stop.} Stop or quarantine this black-box substitution branch only after the source carrier, SAT-query length exponent \(\beta\), interface size exponent q, source-threshold exponent, and target class have been explicitly computed; quarantine if every sharpened interface still has positive \(\beta\) and therefore only excludes a bounded range of SAT circuit exponents rather than NP subset P/poly.

\paragraph{Do not stop prematurely if.} Do not stop before deriving the transfer inequality from \(\mathrm{SAT}\in \mathrm{SIZE}(n^{c})\) to the carrier upper bound, identifying whether \(\beta=0\), \(\beta>0\), or the argument is genuinely nonblackbox, and recording the barrier status of the proposed escape.

\paragraph{How it adds to the global problem.} This move adds to the P vs NP problem because it identifies whether the transfer loses all-exponent strength through a positive target-query exponent \(\beta\); it keeps barrier status and nonuniform/advice leakage attached to the actual transfer theorem. The current branch-control decision is continue\_run\_next. Payoff criterion: For this route, state the exact barrier premise it uses or escapes: relativizing, natural, algebrizing, nonuniform, or advice-leaking; promotion requires the escape to be attached to the specific transfer theorem.

\subsection*{cycle-0148: Hardness-Magnification Bridge Realization}

\paragraph{Learning.} Cycle 148 turns ranked QCTSR from an order-independent closure statement into a representative-stable \(TransparentRowSupport_{N}\) rule case. The valid contribution is a restricted theorem/countermodel pair: proof-carrying quotient compression preserves \(\Gamma\) line induction, while the first unranked, cyclic, advice-chosen, or hidden semantic merge yields a private-row falsifier. This closes a proof-audit loophole without proving any \(H_{N}\) lower bound or \(\mathrm{P} \ne \mathrm{NP}\); full semantic pruning remains exactly at \(RealSATInfluenceRealizer_{N}^{arb}\).

\paragraph{Pressure.} Full semantic FormulaEval with hidden adaptive pruning remains blocked at \(RealSATInfluenceRealizer_{N}^{arb}\): no advice-free nonrelativizing procedure is proved that lists all real SAT influences before the semantic rule is used.

\paragraph{Next.} Write the finite \(K_\mathrm{QCTSR}\) verifier on a projected-dDNNF example with one certified duplicate row and one unread-enable hidden merge; compute \(U_{Q},N|K_{Q}|,B_\mathrm{CTSR},C_{Q},E_{pi}\) and show exactly where the hidden merge fails the valid support ledger.

\paragraph{When to stop.} Stop or quarantine this black-box substitution branch only after the source carrier, SAT-query length exponent \(\beta\), interface size exponent q, source-threshold exponent, and target class have been explicitly computed; quarantine if every sharpened interface still has positive \(\beta\) and therefore only excludes a bounded range of SAT circuit exponents rather than NP subset P/poly.

\paragraph{Do not stop prematurely if.} Do not stop before deriving the transfer inequality from \(\mathrm{SAT}\in \mathrm{SIZE}(n^{c})\) to the carrier upper bound, identifying whether \(\beta=0\), \(\beta>0\), or the argument is genuinely nonblackbox, and recording the barrier status of the proposed escape.

\paragraph{How it adds to the global problem.} This move adds to the P vs NP problem because it identifies whether the transfer loses all-exponent strength through a positive target-query exponent \(\beta\); it keeps barrier status and nonuniform/advice leakage attached to the actual transfer theorem. The current branch-control decision is continue\_run\_next. Payoff criterion: For this route, state the exact barrier premise it uses or escapes: relativizing, natural, algebrizing, nonuniform, or advice-leaking; promotion requires the escape to be attached to the specific transfer theorem.

\subsection*{cycle-0149: Hardness-Magnification Bridge Realization}

\paragraph{Learning.} Cycle 149 folds representative-stable ranked QCTSR into \(TransparentRowSupport_{N}\) as a valid rule case. The new contribution is not another magnification slogan: quotient compression is proof-carrying, rank-stratified, dependency-closed, and charged, while the first rejected merge gives a concrete private-row falsifier. Full semantic SAT pruning remains exactly at \(RealSATInfluenceRealizer_{N}^{arb}\).

\paragraph{Pressure.} The derivation does not cover arbitrary real-SAT semantic pruning or FormulaEval comprehension unless \(RealSATInfluenceRealizer_{N}^{arb}\) uniformly lists the influencing rows and charges them before use.

\paragraph{Next.} Write the standalone paper theorem for \(TransparentRowSupport_{N}^{Qrep}\): include the \(K_{Qrep}\) fixed-point verifier pseudocode, the \(\Gamma_{Qrep}\) ledger proof, the first-rejected-merge countermodel, and a one-page nonclaim box separating this parser boundary from \(H_{N}\) lower bounds and \(\mathrm{P} \ne \mathrm{NP}\).

\paragraph{When to stop.} Stop or quarantine this black-box substitution branch only after the source carrier, SAT-query length exponent \(\beta\), interface size exponent q, source-threshold exponent, and target class have been explicitly computed; quarantine if every sharpened interface still has positive \(\beta\) and therefore only excludes a bounded range of SAT circuit exponents rather than NP subset P/poly.

\paragraph{Do not stop prematurely if.} Do not stop before deriving the transfer inequality from \(\mathrm{SAT}\in \mathrm{SIZE}(n^{c})\) to the carrier upper bound, identifying whether \(\beta=0\), \(\beta>0\), or the argument is genuinely nonblackbox, and recording the barrier status of the proposed escape.

\paragraph{How it adds to the global problem.} This move adds to the P vs NP problem because it identifies whether the transfer loses all-exponent strength through a positive target-query exponent \(\beta\); it keeps barrier status and nonuniform/advice leakage attached to the actual transfer theorem. The current branch-control decision is adapt\_then\_run. Payoff criterion: For this route, state the exact barrier premise it uses or escapes: relativizing, natural, algebrizing, nonuniform, or advice-leaking; promotion requires the escape to be attached to the specific transfer theorem.

\subsection*{cycle-0150: Hardness-Magnification Bridge Realization}

\paragraph{Learning.} Cycle 150 closes a proof-delivery upgrade rather than rediscovering magnification: Qrep quotient compression is valid exactly as a batched, ranked, proof-carrying same-theory dictionary charged in \(\Gamma_{Qrep}\). The sharp falsifier is the first unranked, cyclic, advice-chosen, or hidden semantic merge. The result is a scoped parser/falsifier theorem and Formula Footprints supporting lemma; it does not prove \(H_{N}\) lower bounds or \(\mathrm{P} \ne \mathrm{NP}\).

\paragraph{Pressure.} The derivation does not cover a semantic pruning step that obtains a quotient or child deletion from real SAT truth without a surfaced same-theory certificate; that remains exactly \(RealSATInfluenceRealizer_{N}^{arb}\).

\paragraph{Next.} Run the Tseitin-renaming Qrep thought experiment: instantiate two batches with one certified duplicate and one unlisted duplicate, compute \(K_{Qrep}\), C\_Q(D\_Q), and \(\Gamma_{Qrep}\) explicitly, then verify that the unlisted duplicate yields the first-hidden-row private-row falsifier while the certified duplicate preserves the line induction.

\paragraph{When to stop.} Stop or quarantine this black-box substitution branch only after the source carrier, SAT-query length exponent \(\beta\), interface size exponent q, source-threshold exponent, and target class have been explicitly computed; quarantine if every sharpened interface still has positive \(\beta\) and therefore only excludes a bounded range of SAT circuit exponents rather than NP subset P/poly.

\paragraph{Do not stop prematurely if.} Do not stop before deriving the transfer inequality from \(\mathrm{SAT}\in \mathrm{SIZE}(n^{c})\) to the carrier upper bound, identifying whether \(\beta=0\), \(\beta>0\), or the argument is genuinely nonblackbox, and recording the barrier status of the proposed escape.

\paragraph{How it adds to the global problem.} This move adds to the P vs NP problem because it identifies whether the transfer loses all-exponent strength through a positive target-query exponent \(\beta\); it keeps barrier status and nonuniform/advice leakage attached to the actual transfer theorem. The current branch-control decision is continue\_run\_next. Payoff criterion: For the proof-complexity branch, separate fixed-family proof hardness from language hardness by naming the proof system, shortcut algorithm, and semantic lower-bound target.

\subsection*{cycle-0151: Hardness-Magnification Bridge Realization}

\paragraph{Learning.} Cycle 151 upgrades the batched Qrep patch with temporal prefix stability: quotient compression is valid only as append-only, ranked, proof-carrying dictionary use. The first pre-use, rebinding, cyclic, advice-chosen, or hidden semantic merge is an explicit falsifier, while full semantic pruning stays quarantined.

\paragraph{Pressure.} The derivation does not realize arbitrary semantic SAT pruning; \(RealSATInfluenceRealizer_{N}^{arb}\) remains the named nonrelativizing bridge for full \(\mathrm{Psys}_{sym}\).

\paragraph{Next.} Write the finite prefix-stable Qrep verifier: fields are batch\_id, rank, source\_row, representative\_row, certificate\_root, dependency list, first\_use\_line, and no\_rebind flag; prove the first rejected field yields the private-row falsifier.

\paragraph{When to stop.} Stop or quarantine this P vs NP branch only after it states the source object, target class, complexity model, transfer inequality, quantifier strength, and barrier status; quarantine if two sharpened attempts still cannot say whether the route implies \(\mathrm{SAT}\notin \mathrm{P}/\mathrm{poly}\), \(\mathrm{NP}\not\subset \mathrm{P}/\mathrm{poly}\), or only a weaker/proxy lower bound.

\paragraph{Do not stop prematurely if.} Do not stop before the branch has produced a concrete source-to-target transfer theorem, an explicit obstruction lemma, or a barrier-aware reason why the route cannot reach the P vs NP/P/poly target.

\paragraph{How it adds to the global problem.} This move adds to the P vs NP problem because it keeps barrier status and nonuniform/advice leakage attached to the actual transfer theorem. The current branch-control decision is backtrack\_to\_sharpen. Payoff criterion: For this route, state the exact barrier premise it uses or escapes: relativizing, natural, algebrizing, nonuniform, or advice-leaking; promotion requires the escape to be attached to the specific transfer theorem.

\subsection*{cycle-0152: Hardness-Magnification Bridge Realization: prefix-stable Qrep support closure}

\paragraph{Learning.} Cycle 152 upgrades Qrep from an append-only convention to a prefix-verifiable rule case: quotient compression preserves the transparent parser only when certificate charge strictly precedes representative use. The exact no-go is temporal, not just acyclicity; delayed same-theory evidence is already a private-row falsifier. This closes a proof-audit loophole while preserving the nonclaim boundary: no \(H_{N}\) lower bound and no \(\mathrm{P} \ne \mathrm{NP}\) result.

\paragraph{Pressure.} The restricted PrefixStableQrepTransparentSupport proof is theorem-shaped; the unresolved extension remains \(RealSATInfluenceRealizer_{N}^{arb}\) for hidden full semantic SAT pruning and branch selection.

\paragraph{Next.} Write the finite verifier table for PrefixStableQrep entries: fields entry\_id, append\_time, old\_row, rep\_row, rank, certificate\_root, dependency\_list, first\_use\_time, \(charge_{to,Gamma}\), and rejection reason; then run it on one EqSAT duplicate-child example and one delayed-certificate private-row counterexample.

\paragraph{When to stop.} Stop or quarantine this P vs NP branch only after it states the source object, target class, complexity model, transfer inequality, quantifier strength, and barrier status; quarantine if two sharpened attempts still cannot say whether the route implies \(\mathrm{SAT}\notin \mathrm{P}/\mathrm{poly}\), \(\mathrm{NP}\not\subset \mathrm{P}/\mathrm{poly}\), or only a weaker/proxy lower bound.

\paragraph{Do not stop prematurely if.} Do not stop before the branch has produced a concrete source-to-target transfer theorem, an explicit obstruction lemma, or a barrier-aware reason why the route cannot reach the P vs NP/P/poly target.

\paragraph{How it adds to the global problem.} This move adds to the P vs NP problem because it keeps barrier status and nonuniform/advice leakage attached to the actual transfer theorem. The current branch-control decision is backtrack\_to\_sharpen. Payoff criterion: For this route, state the exact barrier premise it uses or escapes: relativizing, natural, algebrizing, nonuniform, or advice-leaking; promotion requires the escape to be attached to the specific transfer theorem.

\subsection*{cycle-0153: Hardness-Magnification Bridge Realization: prefix-strict Qrep support audit}

\paragraph{Learning.} Cycle 153 upgrades Qrep from prefix-verifiable syntax to a theorem-shaped \(TransparentRowSupport_{N}\) rule case: quotient compression is valid exactly as strict-prefix, ranked, proof-carrying, dependency-closed dictionary use. The new falsifier is temporal; delayed same-theory evidence is already hidden support. This is a scoped parser/falsifier contribution, not a lower bound for \(H_{N}\) and not \(\mathrm{P} \ne \mathrm{NP}\) progress.

\paragraph{Pressure.} The derivation does not cover arbitrary semantic FormulaEval or adaptive SAT-pruning rules that may realize influence through real SAT semantics without exposing row identifiers; that remains \(RealSATInfluenceRealizer_{N}^{arb}\).

\paragraph{Next.} Instantiate the strict-prefix Qrep verifier on a two-duplicate projected-dDNNF FormulaEval example and a delayed-certificate counterexample; compute \(\Gamma_{bEF}\), \(U_{N}\), \(X_{N}\), \(D_{Q}\), \(E_{pi}\), and the q/epsilon support-density consequence line by line.

\paragraph{When to stop.} Stop or quarantine this black-box substitution branch only after the source carrier, SAT-query length exponent \(\beta\), interface size exponent q, source-threshold exponent, and target class have been explicitly computed; quarantine if every sharpened interface still has positive \(\beta\) and therefore only excludes a bounded range of SAT circuit exponents rather than NP subset P/poly.

\paragraph{Do not stop prematurely if.} Do not stop before deriving the transfer inequality from \(\mathrm{SAT}\in \mathrm{SIZE}(n^{c})\) to the carrier upper bound, identifying whether \(\beta=0\), \(\beta>0\), or the argument is genuinely nonblackbox, and recording the barrier status of the proposed escape.

\paragraph{How it adds to the global problem.} This move adds to the P vs NP problem because it identifies whether the transfer loses all-exponent strength through a positive target-query exponent \(\beta\); it keeps barrier status and nonuniform/advice leakage attached to the actual transfer theorem. The current branch-control decision is continue\_run\_next. Payoff criterion: For this route, state the exact barrier premise it uses or escapes: relativizing, natural, algebrizing, nonuniform, or advice-leaking; promotion requires the escape to be attached to the specific transfer theorem.

\subsection*{cycle-0154: Hardness-Magnification Bridge Realization}

\paragraph{Learning.} Cycle 154 converts strict-prefix Qrep from a syntax convention into a replay theorem with a matching temporal falsifier. The valid contribution is a restricted \(TransparentRowSupport_{N}\) rule case and Formula Footprints supporting lemma: charged, ranked, dependency-closed dictionary compression preserves \(\Gamma\) line induction; delayed or hidden compression is exactly hidden support. This is route-level certification only, not a lower bound for \(H_{N}\) and not P vs NP progress.

\paragraph{Pressure.} The derivation does not cover arbitrary semantic SAT branch pruning or FormulaEval rules whose influencing rows are not determined by proof syntax before use; this is exactly \(RealSATInfluenceRealizer_{N}^{arb}\).

\paragraph{Next.} Run the prefix-scan verifier on one minimal proof with a valid two-row strict-prefix Qrep merge and on one identical proof with the certificate delayed by one line; compute \(\Gamma_{Qrep}\) in both cases and export the rejected line as the private-row falsifier pivot.

\paragraph{When to stop.} Stop or quarantine this black-box substitution branch only after the source carrier, SAT-query length exponent \(\beta\), interface size exponent q, source-threshold exponent, and target class have been explicitly computed; quarantine if every sharpened interface still has positive \(\beta\) and therefore only excludes a bounded range of SAT circuit exponents rather than NP subset P/poly.

\paragraph{Do not stop prematurely if.} Do not stop before deriving the transfer inequality from \(\mathrm{SAT}\in \mathrm{SIZE}(n^{c})\) to the carrier upper bound, identifying whether \(\beta=0\), \(\beta>0\), or the argument is genuinely nonblackbox, and recording the barrier status of the proposed escape.

\paragraph{How it adds to the global problem.} This move adds to the P vs NP problem because it identifies whether the transfer loses all-exponent strength through a positive target-query exponent \(\beta\); it keeps barrier status and nonuniform/advice leakage attached to the actual transfer theorem. The current branch-control decision is continue\_run\_next. Payoff criterion: For this route, state the exact barrier premise it uses or escapes: relativizing, natural, algebrizing, nonuniform, or advice-leaking; promotion requires the escape to be attached to the specific transfer theorem.


\end{document}
